\documentclass[10pt,oneside]{book}
\pdftrailerid{<E4C70ECBED55D6112007E25F5BE74344><E4C70ECBED55D6112007E25F5BE74344>}

\usepackage[T1]{fontenc}
\usepackage[utf8]{inputenc}
\usepackage[spanish,es-nodecimaldot,es-noquoting]{babel}
\usepackage{amsmath,amssymb,mathrsfs}
\usepackage[all]{xy}
\usepackage{tikz-cd}
\usetikzlibrary{babel}
\usepackage[shortlabels]{enumitem}
\usepackage{xparse}
\usepackage{xstring}
\usepackage[a4paper,margin=22mm]{geometry}
\usepackage[
  colorlinks=true,
  linkcolor=blue,
  hypertexnames=false,
  unicode=true,
  pdfencoding=auto
]{hyperref}
\hypersetup{
  pdftitle={Elementos de geometría algebraica IV, cuarta parte -- edición española},
  pdfauthor={Alexandre Grothendieck y Jean Dieudonné},
  pdfsubject={Estudio local de los esquemas y de los morfismos de esquemas: secciones 16 a 21}
}

\renewcommand{\thesection}{\arabic{section}}
\renewcommand{\thesubsection}{\thesection.\arabic{subsection}}
\def\sh{\mathcal}                   % sheaf font
\def\bb{\mathbf}                    % bold font
\def\cat{\mathtt}                   % category font
\def\leq{\leqslant}                 % <=
\def\geq{\geqslant}                 % >=
\def\setmin{-}                      % set minus
\def\rad{\mathfrak{r}}              % radical
\def\nilrad{\mathfrak{N}}           % nilradical
\def\emp{\varnothing}               % empty set
\def\vphi{\phi}                     % for switching \phi and \varphi, change if needed
\def\HH{\mathrm{H}}                 % cohomology H
\def\CHH{\check{\HH}}               % Čech cohomology H
\def\RR{\mathrm{R}}                 % right derived R
\def\LL{\mathrm{L}}                 % left derived L
\def\dual#1{{#1}^\vee}              % dual
\def\kres{k}                        % residue field k
\def\C{\cat{C}}                     % category C
\def\op{^\cat{op}}                  % opposite category
\def\Set{\cat{Set}}                 % category of sets
\def\CHom{\cat{Hom}}                % functor category
\def\supertilde{{\,\widetilde{\,}\,}}   % use \supertilde instead of ^\sim
\def\GL{\bb{GL}}
\def\bbGamma{\boldsymbol\Gamma}
\def\red{\mathrm{red}}
\def\rg{\operatorname{rg}}
\def\gr{\operatorname{gr}}
\def\Gr{\operatorname{Gr}}
\def\Hom{\operatorname{Hom}}
\def\Proj{\operatorname{Proj}}
\def\Tor{\operatorname{Tor}}
\def\Ext{\operatorname{Ext}}
\def\Supp{\operatorname{Supp}}
\def\Ker{\operatorname{Ker}}
\def\Im{\operatorname{Im}}
\def\Coker{\operatorname{Coker}}
\def\Spec{\operatorname{Spec}}
\def\Spm{\operatorname{Spm}}
\def\Spf{\operatorname{Spf}}
\def\grad{\operatorname{grad}}
\def\dimc{\operatorname{dimc}}
\def\codim{\operatorname{codim}}
\def\id{\operatorname{id}}
\def\Der{\operatorname{Der}}
\def\Diff{\operatorname{Diff}}
\def\Hyp{\operatorname{Hyp}}
\def\Tr{\operatorname{Tr}}

\renewcommand{\to}{\mathchoice{\longrightarrow}{\rightarrow}{\rightarrow}{\rightarrow}}
\newcommand{\from}{\mathchoice{\longleftarrow}{\leftarrow}{\leftarrow}{\leftarrow}}
\let\mapstoo\mapsto
\renewcommand{\mapsto}{\mathchoice{\longmapsto}{\mapstoo}{\mapstoo}{\mapstoo}}
\def\isoto{\simeq} % isomorphism
\def\simto{\xrightarrow{\sim}}

\renewcommand{\hat}[1]{\widehat{#1}}
\renewcommand{\tilde}[1]{\widetilde{#1}}

\makeatletter
\newcommand{\EGAUnavailableHyperref}[2]{%
  \begingroup
  \let\@latex@warning\@gobble
  \let\@latex@warning@no@line\@gobble
  \let\G@refundefinedtrue\relax
  \hyperref[#1]{#2}%
  \endgroup
}
\NewDocumentCommand{\sref}{o m o}{%
  \@ifundefined{r@#2}%
    {\IfNoValueTF{#1}%
      {\IfNoValueTF{#3}%
        {\EGAUnavailableHyperref{#2}{\normalfont{(\StrBehind{#2}{.})}}}%
        {\EGAUnavailableHyperref{#2}{\normalfont{(\StrBehind{#2}{.},~{#3})}}}}%
      {\IfNoValueTF{#3}%
        {\EGAUnavailableHyperref{#2}{\normalfont{(\textbf{#1},~\StrBehind{#2}{.})}}}%
        {\EGAUnavailableHyperref{#2}{\normalfont{(\textbf{#1},~\StrBehind{#2}{.},~{#3})}}}}}%
    {\IfNoValueTF{#1}%
      {\IfNoValueTF{#3}%
        {\hyperref[#2]{\normalfont{(\StrBehind{#2}{.})}}}%
        {\hyperref[#2]{\normalfont{(\StrBehind{#2}{.},~{#3})}}}}%
      {\IfNoValueTF{#3}%
        {\hyperref[#2]{\normalfont{(\textbf{#1},~\StrBehind{#2}{.})}}}%
        {\hyperref[#2]{\normalfont{(\textbf{#1},~\StrBehind{#2}{.},~{#3})}}}}}%
}
\newcommand{\EGAIVFourExternalLabel}[1]{%
  \expandafter\gdef\csname r@#1\endcsname{{#1}{0}{}{#1}{}}}
\newcommand{\EGAIVFourExternalCitation}[2]{%
  \expandafter\gdef\csname b@#1\endcsname{#2}}
\makeatother
\newcommand{\oldpage}[2][]{%
  \setcounter{footnote}{0}%
  \marginpar{\scriptsize\textbf{#1}~|~#2}\ignorespaces}
\newenvironment{env}[1][]{\par\medskip\noindent\textbf{(#1)}\ }{\par}
\newenvironment{definition}[1][]{%
  \par\medskip\noindent\itshape Definición (#1). ---\ }{\par}
\newenvironment{examples}[1][]{%
  \par\medskip\noindent\itshape Ejemplos (#1). ---\ }{\par}
\newenvironment{proposition}[1][]{%
  \par\medskip\noindent\itshape Proposición (#1). ---\ }{\par}
\newenvironment{theorem}[1][]{%
  \par\medskip\noindent\itshape Teorema (#1). ---\ }{\par}
\newenvironment{corollary}[1][]{%
  \par\medskip\noindent\itshape Corolario (#1). ---\ }{\par}
\newenvironment{remark}[1][]{%
  \par\medskip\noindent\itshape Observación (#1). ---\ }{\par}
\newenvironment{remarks}[1][]{%
  \par\medskip\noindent\itshape Observaciones (#1). ---\ }{\par}
\newenvironment{lemma}[1][]{%
  \par\medskip\noindent\itshape Lema (#1). ---\ }{\par}
\newenvironment{proof}{%
  \par\medskip\noindent\textit{Demostración. ---\ }}{\par}

\makeatletter
\expandafter\def\csname EGAExternal@0.1.4.5\endcsname{https://github.com/KokunoYumeto/ega-fr/releases/download/ega-fr-2026-10-23-r16/EGA_French_Current_Cumulative_Reader.pdf\#page=12}
\expandafter\def\csname EGAExternal@0.1.5.1\endcsname{https://github.com/KokunoYumeto/ega-fr/releases/download/ega-fr-2026-10-23-r16/EGA_French_Current_Cumulative_Reader.pdf\#page=13}
\expandafter\def\csname EGAExternal@0.1.5.7\endcsname{https://github.com/KokunoYumeto/ega-fr/releases/download/ega-fr-2026-10-23-r16/EGA_French_Current_Cumulative_Reader.pdf\#page=13}
\expandafter\def\csname EGAExternal@0.12.2.1\endcsname{https://github.com/KokunoYumeto/ega-fr/releases/download/ega-fr-2026-10-23-r16/EGA_French_Current_Cumulative_Reader.pdf\#page=337}
\expandafter\def\csname EGAExternal@0.16.3.4\endcsname{https://github.com/KokunoYumeto/ega-fr/releases/download/ega-fr-2026-10-23-r16/EGA_French_Current_Cumulative_Reader.pdf\#page=509}
\expandafter\def\csname EGAExternal@0.16.4.6\endcsname{https://github.com/KokunoYumeto/ega-fr/releases/download/ega-fr-2026-10-23-r16/EGA_French_Current_Cumulative_Reader.pdf\#page=513}
\expandafter\def\csname EGAExternal@0.16.5.1\endcsname{https://github.com/KokunoYumeto/ega-fr/releases/download/ega-fr-2026-10-23-r16/EGA_French_Current_Cumulative_Reader.pdf\#page=515}
\expandafter\def\csname EGAExternal@0.17.1.1\endcsname{https://github.com/KokunoYumeto/ega-fr/releases/download/ega-fr-2026-10-23-r16/EGA_French_Current_Cumulative_Reader.pdf\#page=517}
\expandafter\def\csname EGAExternal@0.17.3.2\endcsname{https://github.com/KokunoYumeto/ega-fr/releases/download/ega-fr-2026-10-23-r16/EGA_French_Current_Cumulative_Reader.pdf\#page=522}
\expandafter\def\csname EGAExternal@0.23.2.6\endcsname{https://github.com/KokunoYumeto/ega-fr/releases/download/ega-fr-2026-10-23-r16/EGA_French_Current_Cumulative_Reader.pdf\#page=643}
\expandafter\def\csname EGAExternal@0.3.4.1\endcsname{https://github.com/KokunoYumeto/ega-fr/releases/download/ega-fr-2026-10-23-r16/EGA_French_Current_Cumulative_Reader.pdf\#page=22}
\expandafter\def\csname EGAExternal@0.3.6\endcsname{https://github.com/KokunoYumeto/ega-fr/releases/download/ega-fr-2026-10-23-r16/EGA_French_Current_Cumulative_Reader.pdf\#page=25}
\expandafter\def\csname EGAExternal@0.4.3.3\endcsname{https://github.com/KokunoYumeto/ega-fr/releases/download/ega-fr-2026-10-23-r16/EGA_French_Current_Cumulative_Reader.pdf\#page=31}
\expandafter\def\csname EGAExternal@0.5.4.1\endcsname{https://github.com/KokunoYumeto/ega-fr/releases/download/ega-fr-2026-10-23-r16/EGA_French_Current_Cumulative_Reader.pdf\#page=36}
\expandafter\def\csname EGAExternal@0.5.4.3\endcsname{https://github.com/KokunoYumeto/ega-fr/releases/download/ega-fr-2026-10-23-r16/EGA_French_Current_Cumulative_Reader.pdf\#page=36}
\expandafter\def\csname EGAExternal@0.5.4.5\endcsname{https://github.com/KokunoYumeto/ega-fr/releases/download/ega-fr-2026-10-23-r16/EGA_French_Current_Cumulative_Reader.pdf\#page=37}
\expandafter\def\csname EGAExternal@0.5.4.7\endcsname{https://github.com/KokunoYumeto/ega-fr/releases/download/ega-fr-2026-10-23-r16/EGA_French_Current_Cumulative_Reader.pdf\#page=37}
\expandafter\def\csname EGAExternal@0.6.1.2\endcsname{https://github.com/KokunoYumeto/ega-fr/releases/download/ega-fr-2026-10-23-r16/EGA_French_Current_Cumulative_Reader.pdf\#page=40}
\expandafter\def\csname EGAExternal@0.6.3.4\endcsname{https://github.com/KokunoYumeto/ega-fr/releases/download/ega-fr-2026-10-23-r16/EGA_French_Current_Cumulative_Reader.pdf\#page=41}
\expandafter\def\csname EGAExternal@0.6.5.1\endcsname{https://github.com/KokunoYumeto/ega-fr/releases/download/ega-fr-2026-10-23-r16/EGA_French_Current_Cumulative_Reader.pdf\#page=43}
\expandafter\def\csname EGAExternal@0.6.6.2\endcsname{https://github.com/KokunoYumeto/ega-fr/releases/download/ega-fr-2026-10-23-r16/EGA_French_Current_Cumulative_Reader.pdf\#page=43}
\expandafter\def\csname EGAExternal@0.6.7.1\endcsname{https://github.com/KokunoYumeto/ega-fr/releases/download/ega-fr-2026-10-23-r16/EGA_French_Current_Cumulative_Reader.pdf\#page=43}
\expandafter\def\csname EGAExternal@0.9.2.4\endcsname{https://github.com/KokunoYumeto/ega-fr/releases/download/ega-fr-2026-10-23-r16/EGA_French_Current_Cumulative_Reader.pdf\#page=307}
\expandafter\def\csname EGAExternal@I.1.4.1\endcsname{https://github.com/KokunoYumeto/ega-fr/releases/download/ega-fr-2026-10-23-r16/EGA_French_Current_Cumulative_Reader.pdf\#page=65}
\expandafter\def\csname EGAExternal@I.2.2\endcsname{https://github.com/KokunoYumeto/ega-fr/releases/download/ega-fr-2026-10-23-r16/EGA_French_Current_Cumulative_Reader.pdf\#page=70}
\expandafter\def\csname EGAExternal@I.2.3.1\endcsname{https://github.com/KokunoYumeto/ega-fr/releases/download/ega-fr-2026-10-23-r16/EGA_French_Current_Cumulative_Reader.pdf\#page=72}
\expandafter\def\csname EGAExternal@I.3.3.10.1\endcsname{https://github.com/KokunoYumeto/ega-fr/releases/download/ega-fr-2026-10-23-r16/EGA_French_Current_Cumulative_Reader.pdf\#page=78}
\expandafter\def\csname EGAExternal@I.3.3.12\endcsname{https://github.com/KokunoYumeto/ega-fr/releases/download/ega-fr-2026-10-23-r16/EGA_French_Current_Cumulative_Reader.pdf\#page=78}
\expandafter\def\csname EGAExternal@I.3.3.14\endcsname{https://github.com/KokunoYumeto/ega-fr/releases/download/ega-fr-2026-10-23-r16/EGA_French_Current_Cumulative_Reader.pdf\#page=79}
\expandafter\def\csname EGAExternal@I.3.3.15\endcsname{https://github.com/KokunoYumeto/ega-fr/releases/download/ega-fr-2026-10-23-r16/EGA_French_Current_Cumulative_Reader.pdf\#page=79}
\expandafter\def\csname EGAExternal@I.3.4.8\endcsname{https://github.com/KokunoYumeto/ega-fr/releases/download/ega-fr-2026-10-23-r16/EGA_French_Current_Cumulative_Reader.pdf\#page=80}
\expandafter\def\csname EGAExternal@I.3.6.4\endcsname{https://github.com/KokunoYumeto/ega-fr/releases/download/ega-fr-2026-10-23-r16/EGA_French_Current_Cumulative_Reader.pdf\#page=83}
\expandafter\def\csname EGAExternal@I.4.4.1\endcsname{https://github.com/KokunoYumeto/ega-fr/releases/download/ega-fr-2026-10-23-r16/EGA_French_Current_Cumulative_Reader.pdf\#page=88}
\expandafter\def\csname EGAExternal@I.5.2.1\endcsname{https://github.com/KokunoYumeto/ega-fr/releases/download/ega-fr-2026-10-23-r16/EGA_French_Current_Cumulative_Reader.pdf\#page=92}
\expandafter\def\csname EGAExternal@I.5.2.2\endcsname{https://github.com/KokunoYumeto/ega-fr/releases/download/ega-fr-2026-10-23-r16/EGA_French_Current_Cumulative_Reader.pdf\#page=92}
\expandafter\def\csname EGAExternal@I.5.3.11\endcsname{https://github.com/KokunoYumeto/ega-fr/releases/download/ega-fr-2026-10-23-r16/EGA_French_Current_Cumulative_Reader.pdf\#page=93}
\expandafter\def\csname EGAExternal@I.5.4.3\endcsname{https://github.com/KokunoYumeto/ega-fr/releases/download/ega-fr-2026-10-23-r16/EGA_French_Current_Cumulative_Reader.pdf\#page=95}
\expandafter\def\csname EGAExternal@I.5.5.1\endcsname{https://github.com/KokunoYumeto/ega-fr/releases/download/ega-fr-2026-10-23-r16/EGA_French_Current_Cumulative_Reader.pdf\#page=96}
\expandafter\def\csname EGAExternal@I.6.4.4\endcsname{https://github.com/KokunoYumeto/ega-fr/releases/download/ega-fr-2026-10-23-r16/EGA_French_Current_Cumulative_Reader.pdf\#page=103}
\expandafter\def\csname EGAExternal@I.7.1.2\endcsname{https://github.com/KokunoYumeto/ega-fr/releases/download/ega-fr-2026-10-23-r16/EGA_French_Current_Cumulative_Reader.pdf\#page=109}
\expandafter\def\csname EGAExternal@I.7.1.9\endcsname{https://github.com/KokunoYumeto/ega-fr/releases/download/ega-fr-2026-10-23-r16/EGA_French_Current_Cumulative_Reader.pdf\#page=109}
\expandafter\def\csname EGAExternal@I.7.1.9.1\endcsname{https://github.com/KokunoYumeto/ega-fr/releases/download/ega-fr-2026-10-23-r16/EGA_French_Current_Cumulative_Reader.pdf\#page=109}
\expandafter\def\csname EGAExternal@I.7.2.2\endcsname{https://github.com/KokunoYumeto/ega-fr/releases/download/ega-fr-2026-10-23-r16/EGA_French_Current_Cumulative_Reader.pdf\#page=111}
\expandafter\def\csname EGAExternal@I.7.3.2\endcsname{https://github.com/KokunoYumeto/ega-fr/releases/download/ega-fr-2026-10-23-r16/EGA_French_Current_Cumulative_Reader.pdf\#page=113}
\expandafter\def\csname EGAExternal@I.7.3.5\endcsname{https://github.com/KokunoYumeto/ega-fr/releases/download/ega-fr-2026-10-23-r16/EGA_French_Current_Cumulative_Reader.pdf\#page=113}
\expandafter\def\csname EGAExternal@I.7.4\endcsname{https://github.com/KokunoYumeto/ega-fr/releases/download/ega-fr-2026-10-23-r16/EGA_French_Current_Cumulative_Reader.pdf\#page=114}
\expandafter\def\csname EGAExternal@I.9.5.1\endcsname{https://github.com/KokunoYumeto/ega-fr/releases/download/ega-fr-2026-10-23-r16/EGA_French_Current_Cumulative_Reader.pdf\#page=123}
\expandafter\def\csname EGAExternal@I.9.5.10\endcsname{https://github.com/KokunoYumeto/ega-fr/releases/download/ega-fr-2026-10-23-r16/EGA_French_Current_Cumulative_Reader.pdf\#page=123}
\expandafter\def\csname EGAExternal@I.9.5.11\endcsname{https://github.com/KokunoYumeto/ega-fr/releases/download/ega-fr-2026-10-23-r16/EGA_French_Current_Cumulative_Reader.pdf\#page=123}
\expandafter\def\csname EGAExternal@I.9.5.3\endcsname{https://github.com/KokunoYumeto/ega-fr/releases/download/ega-fr-2026-10-23-r16/EGA_French_Current_Cumulative_Reader.pdf\#page=123}
\expandafter\def\csname EGAExternal@II.1.3\endcsname{https://github.com/KokunoYumeto/ega-fr/releases/download/ega-fr-2026-10-23-r16/EGA_French_Current_Cumulative_Reader.pdf\#page=152}
\expandafter\def\csname EGAExternal@II.1.7.9\endcsname{https://github.com/KokunoYumeto/ega-fr/releases/download/ega-fr-2026-10-23-r16/EGA_French_Current_Cumulative_Reader.pdf\#page=156}
\expandafter\def\csname EGAExternal@II.4.5.4\endcsname{https://github.com/KokunoYumeto/ega-fr/releases/download/ega-fr-2026-10-23-r16/EGA_French_Current_Cumulative_Reader.pdf\#page=204}
\expandafter\def\csname EGAExternal@II.5.1.4\endcsname{https://github.com/KokunoYumeto/ega-fr/releases/download/ega-fr-2026-10-23-r16/EGA_French_Current_Cumulative_Reader.pdf\#page=212}
\expandafter\def\csname EGAExternal@II.5.3.1\endcsname{https://github.com/KokunoYumeto/ega-fr/releases/download/ega-fr-2026-10-23-r16/EGA_French_Current_Cumulative_Reader.pdf\#page=215}
\expandafter\def\csname EGAExternal@II.5.5.3\endcsname{https://github.com/KokunoYumeto/ega-fr/releases/download/ega-fr-2026-10-23-r16/EGA_French_Current_Cumulative_Reader.pdf\#page=218}
\expandafter\def\csname EGAExternal@II.6.3.8\endcsname{https://github.com/KokunoYumeto/ega-fr/releases/download/ega-fr-2026-10-23-r16/EGA_French_Current_Cumulative_Reader.pdf\#page=228}
\expandafter\def\csname EGAExternal@II.7.1.6\endcsname{https://github.com/KokunoYumeto/ega-fr/releases/download/ega-fr-2026-10-23-r16/EGA_French_Current_Cumulative_Reader.pdf\#page=241}
\expandafter\def\csname EGAExternal@II.7.1.9\endcsname{https://github.com/KokunoYumeto/ega-fr/releases/download/ega-fr-2026-10-23-r16/EGA_French_Current_Cumulative_Reader.pdf\#page=241}
\expandafter\def\csname EGAExternal@II.7.2.3\endcsname{https://github.com/KokunoYumeto/ega-fr/releases/download/ega-fr-2026-10-23-r16/EGA_French_Current_Cumulative_Reader.pdf\#page=243}
\expandafter\def\csname EGAExternal@IV.1.10.3\endcsname{https://github.com/KokunoYumeto/ega-fr/releases/download/ega-fr-2026-10-23-r16/EGA_French_Current_Cumulative_Reader.pdf\#page=664}
\expandafter\def\csname EGAExternal@IV.1.6.2\endcsname{https://github.com/KokunoYumeto/ega-fr/releases/download/ega-fr-2026-10-23-r16/EGA_French_Current_Cumulative_Reader.pdf\#page=655}
\expandafter\def\csname EGAExternal@IV.1.7.4\endcsname{https://github.com/KokunoYumeto/ega-fr/releases/download/ega-fr-2026-10-23-r16/EGA_French_Current_Cumulative_Reader.pdf\#page=656}
\expandafter\def\csname EGAExternal@IV.1.8.7.1\endcsname{https://github.com/KokunoYumeto/ega-fr/releases/download/ega-fr-2026-10-23-r16/EGA_French_Current_Cumulative_Reader.pdf\#page=659}
\expandafter\def\csname EGAExternal@IV.1.9.6\endcsname{https://github.com/KokunoYumeto/ega-fr/releases/download/ega-fr-2026-10-23-r16/EGA_French_Current_Cumulative_Reader.pdf\#page=659}
\expandafter\def\csname EGAExternal@IV.10.2.1\endcsname{https://github.com/KokunoYumeto/ega-fr/releases/download/ega-fr-2026-10-23-r16/EGA_French_Current_Cumulative_Reader.pdf\#page=897}
\expandafter\def\csname EGAExternal@IV.10.4\endcsname{https://github.com/KokunoYumeto/ega-fr/releases/download/ega-fr-2026-10-23-r16/EGA_French_Current_Cumulative_Reader.pdf\#page=900}
\expandafter\def\csname EGAExternal@IV.10.4.11\endcsname{https://github.com/KokunoYumeto/ega-fr/releases/download/ega-fr-2026-10-23-r16/EGA_French_Current_Cumulative_Reader.pdf\#page=901}
\expandafter\def\csname EGAExternal@IV.11.10.1\endcsname{https://github.com/KokunoYumeto/ega-fr/releases/download/ega-fr-2026-10-23-r16/EGA_French_Current_Cumulative_Reader.pdf\#page=948}
\expandafter\def\csname EGAExternal@IV.11.10.10\endcsname{https://github.com/KokunoYumeto/ega-fr/releases/download/ega-fr-2026-10-23-r16/EGA_French_Current_Cumulative_Reader.pdf\#page=949}
\expandafter\def\csname EGAExternal@IV.11.10.2\endcsname{https://github.com/KokunoYumeto/ega-fr/releases/download/ega-fr-2026-10-23-r16/EGA_French_Current_Cumulative_Reader.pdf\#page=948}
\expandafter\def\csname EGAExternal@IV.11.10.3\endcsname{https://github.com/KokunoYumeto/ega-fr/releases/download/ega-fr-2026-10-23-r16/EGA_French_Current_Cumulative_Reader.pdf\#page=948}
\expandafter\def\csname EGAExternal@IV.11.10.4\endcsname{https://github.com/KokunoYumeto/ega-fr/releases/download/ega-fr-2026-10-23-r16/EGA_French_Current_Cumulative_Reader.pdf\#page=949}
\expandafter\def\csname EGAExternal@IV.11.10.5\endcsname{https://github.com/KokunoYumeto/ega-fr/releases/download/ega-fr-2026-10-23-r16/EGA_French_Current_Cumulative_Reader.pdf\#page=949}
\expandafter\def\csname EGAExternal@IV.11.10.8\endcsname{https://github.com/KokunoYumeto/ega-fr/releases/download/ega-fr-2026-10-23-r16/EGA_French_Current_Cumulative_Reader.pdf\#page=949}
\expandafter\def\csname EGAExternal@IV.11.10.9\endcsname{https://github.com/KokunoYumeto/ega-fr/releases/download/ega-fr-2026-10-23-r16/EGA_French_Current_Cumulative_Reader.pdf\#page=949}
\expandafter\def\csname EGAExternal@IV.11.2.9\endcsname{https://github.com/KokunoYumeto/ega-fr/releases/download/ega-fr-2026-10-23-r16/EGA_French_Current_Cumulative_Reader.pdf\#page=917}
\expandafter\def\csname EGAExternal@IV.11.3.1\endcsname{https://github.com/KokunoYumeto/ega-fr/releases/download/ega-fr-2026-10-23-r16/EGA_French_Current_Cumulative_Reader.pdf\#page=922}
\expandafter\def\csname EGAExternal@IV.11.3.10\endcsname{https://github.com/KokunoYumeto/ega-fr/releases/download/ega-fr-2026-10-23-r16/EGA_French_Current_Cumulative_Reader.pdf\#page=925}
\expandafter\def\csname EGAExternal@IV.11.3.4\endcsname{https://github.com/KokunoYumeto/ega-fr/releases/download/ega-fr-2026-10-23-r16/EGA_French_Current_Cumulative_Reader.pdf\#page=922}
\expandafter\def\csname EGAExternal@IV.11.3.7\endcsname{https://github.com/KokunoYumeto/ega-fr/releases/download/ega-fr-2026-10-23-r16/EGA_French_Current_Cumulative_Reader.pdf\#page=922}
\expandafter\def\csname EGAExternal@IV.11.3.8\endcsname{https://github.com/KokunoYumeto/ega-fr/releases/download/ega-fr-2026-10-23-r16/EGA_French_Current_Cumulative_Reader.pdf\#page=923}
\expandafter\def\csname EGAExternal@IV.11.7.5\endcsname{https://github.com/KokunoYumeto/ega-fr/releases/download/ega-fr-2026-10-23-r16/EGA_French_Current_Cumulative_Reader.pdf\#page=941}
\expandafter\def\csname EGAExternal@IV.12.1.1\endcsname{https://github.com/KokunoYumeto/ega-fr/releases/download/ega-fr-2026-10-23-r16/EGA_French_Current_Cumulative_Reader.pdf\#page=950}
\expandafter\def\csname EGAExternal@IV.12.1.2\endcsname{https://github.com/KokunoYumeto/ega-fr/releases/download/ega-fr-2026-10-23-r16/EGA_French_Current_Cumulative_Reader.pdf\#page=952}
\expandafter\def\csname EGAExternal@IV.13.1.4\endcsname{https://github.com/KokunoYumeto/ega-fr/releases/download/ega-fr-2026-10-23-r16/EGA_French_Current_Cumulative_Reader.pdf\#page=961}
\expandafter\def\csname EGAExternal@IV.13.3.9\endcsname{https://github.com/KokunoYumeto/ega-fr/releases/download/ega-fr-2026-10-23-r16/EGA_French_Current_Cumulative_Reader.pdf\#page=967}
\expandafter\def\csname EGAExternal@IV.14.1.5\endcsname{https://github.com/KokunoYumeto/ega-fr/releases/download/ega-fr-2026-10-23-r16/EGA_French_Current_Cumulative_Reader.pdf\#page=969}
\expandafter\def\csname EGAExternal@IV.14.2.1\endcsname{https://github.com/KokunoYumeto/ega-fr/releases/download/ega-fr-2026-10-23-r16/EGA_French_Current_Cumulative_Reader.pdf\#page=969}
\expandafter\def\csname EGAExternal@IV.14.3.8\endcsname{https://github.com/KokunoYumeto/ega-fr/releases/download/ega-fr-2026-10-23-r16/EGA_French_Current_Cumulative_Reader.pdf\#page=973}
\expandafter\def\csname EGAExternal@IV.14.4.1\endcsname{https://github.com/KokunoYumeto/ega-fr/releases/download/ega-fr-2026-10-23-r16/EGA_French_Current_Cumulative_Reader.pdf\#page=975}
\expandafter\def\csname EGAExternal@IV.14.4.4\endcsname{https://github.com/KokunoYumeto/ega-fr/releases/download/ega-fr-2026-10-23-r16/EGA_French_Current_Cumulative_Reader.pdf\#page=976}
\expandafter\def\csname EGAExternal@IV.14.5.11\endcsname{https://github.com/KokunoYumeto/ega-fr/releases/download/ega-fr-2026-10-23-r16/EGA_French_Current_Cumulative_Reader.pdf\#page=983}
\expandafter\def\csname EGAExternal@IV.14.5.7\endcsname{https://github.com/KokunoYumeto/ega-fr/releases/download/ega-fr-2026-10-23-r16/EGA_French_Current_Cumulative_Reader.pdf\#page=981}
\expandafter\def\csname EGAExternal@IV.15.4.2\endcsname{https://github.com/KokunoYumeto/ega-fr/releases/download/ega-fr-2026-10-23-r16/EGA_French_Current_Cumulative_Reader.pdf\#page=988}
\expandafter\def\csname EGAExternal@IV.15.5.2\endcsname{https://github.com/KokunoYumeto/ega-fr/releases/download/ega-fr-2026-10-23-r16/EGA_French_Current_Cumulative_Reader.pdf\#page=989}
\expandafter\def\csname EGAExternal@IV.15.6.10\endcsname{https://github.com/KokunoYumeto/ega-fr/releases/download/ega-fr-2026-10-23-r16/EGA_French_Current_Cumulative_Reader.pdf\#page=996}
\expandafter\def\csname EGAExternal@IV.15.6.2\endcsname{https://github.com/KokunoYumeto/ega-fr/releases/download/ega-fr-2026-10-23-r16/EGA_French_Current_Cumulative_Reader.pdf\#page=993}
\expandafter\def\csname EGAExternal@IV.15.6.6\endcsname{https://github.com/KokunoYumeto/ega-fr/releases/download/ega-fr-2026-10-23-r16/EGA_French_Current_Cumulative_Reader.pdf\#page=994}
\expandafter\def\csname EGAExternal@IV.17.1\endcsname{https://github.com/KokunoYumeto/ega-fr/releases/download/ega-fr-2026-10-23-r16/EGA_French_Current_Cumulative_Reader.pdf\#page=1044}
\expandafter\def\csname EGAExternal@IV.17.10\endcsname{https://github.com/KokunoYumeto/ega-fr/releases/download/ega-fr-2026-10-23-r16/EGA_French_Current_Cumulative_Reader.pdf\#page=1060}
\expandafter\def\csname EGAExternal@IV.17.11\endcsname{https://github.com/KokunoYumeto/ega-fr/releases/download/ega-fr-2026-10-23-r16/EGA_French_Current_Cumulative_Reader.pdf\#page=1061}
\expandafter\def\csname EGAExternal@IV.17.12\endcsname{https://github.com/KokunoYumeto/ega-fr/releases/download/ega-fr-2026-10-23-r16/EGA_French_Current_Cumulative_Reader.pdf\#page=1063}
\expandafter\def\csname EGAExternal@IV.17.14\endcsname{https://github.com/KokunoYumeto/ega-fr/releases/download/ega-fr-2026-10-23-r16/EGA_French_Current_Cumulative_Reader.pdf\#page=1071}
\expandafter\def\csname EGAExternal@IV.17.15\endcsname{https://github.com/KokunoYumeto/ega-fr/releases/download/ega-fr-2026-10-23-r16/EGA_French_Current_Cumulative_Reader.pdf\#page=1072}
\expandafter\def\csname EGAExternal@IV.17.16\endcsname{https://github.com/KokunoYumeto/ega-fr/releases/download/ega-fr-2026-10-23-r16/EGA_French_Current_Cumulative_Reader.pdf\#page=1076}
\expandafter\def\csname EGAExternal@IV.17.2\endcsname{https://github.com/KokunoYumeto/ega-fr/releases/download/ega-fr-2026-10-23-r16/EGA_French_Current_Cumulative_Reader.pdf\#page=1046}
\expandafter\def\csname EGAExternal@IV.17.3\endcsname{https://github.com/KokunoYumeto/ega-fr/releases/download/ega-fr-2026-10-23-r16/EGA_French_Current_Cumulative_Reader.pdf\#page=1047}
\expandafter\def\csname EGAExternal@IV.17.4\endcsname{https://github.com/KokunoYumeto/ega-fr/releases/download/ega-fr-2026-10-23-r16/EGA_French_Current_Cumulative_Reader.pdf\#page=1049}
\expandafter\def\csname EGAExternal@IV.17.5\endcsname{https://github.com/KokunoYumeto/ega-fr/releases/download/ega-fr-2026-10-23-r16/EGA_French_Current_Cumulative_Reader.pdf\#page=1051}
\expandafter\def\csname EGAExternal@IV.17.6\endcsname{https://github.com/KokunoYumeto/ega-fr/releases/download/ega-fr-2026-10-23-r16/EGA_French_Current_Cumulative_Reader.pdf\#page=1053}
\expandafter\def\csname EGAExternal@IV.18.1\endcsname{https://github.com/KokunoYumeto/ega-fr/releases/download/ega-fr-2026-10-23-r16/EGA_French_Current_Cumulative_Reader.pdf\#page=1079}
\expandafter\def\csname EGAExternal@IV.18.10\endcsname{https://github.com/KokunoYumeto/ega-fr/releases/download/ega-fr-2026-10-23-r16/EGA_French_Current_Cumulative_Reader.pdf\#page=1110}
\expandafter\def\csname EGAExternal@IV.18.11\endcsname{https://github.com/KokunoYumeto/ega-fr/releases/download/ega-fr-2026-10-23-r16/EGA_French_Current_Cumulative_Reader.pdf\#page=1118}
\expandafter\def\csname EGAExternal@IV.18.12\endcsname{https://github.com/KokunoYumeto/ega-fr/releases/download/ega-fr-2026-10-23-r16/EGA_French_Current_Cumulative_Reader.pdf\#page=1125}
\expandafter\def\csname EGAExternal@IV.18.2\endcsname{https://github.com/KokunoYumeto/ega-fr/releases/download/ega-fr-2026-10-23-r16/EGA_French_Current_Cumulative_Reader.pdf\#page=1080}
\expandafter\def\csname EGAExternal@IV.18.3\endcsname{https://github.com/KokunoYumeto/ega-fr/releases/download/ega-fr-2026-10-23-r16/EGA_French_Current_Cumulative_Reader.pdf\#page=1082}
\expandafter\def\csname EGAExternal@IV.18.4\endcsname{https://github.com/KokunoYumeto/ega-fr/releases/download/ega-fr-2026-10-23-r16/EGA_French_Current_Cumulative_Reader.pdf\#page=1085}
\expandafter\def\csname EGAExternal@IV.18.5\endcsname{https://github.com/KokunoYumeto/ega-fr/releases/download/ega-fr-2026-10-23-r16/EGA_French_Current_Cumulative_Reader.pdf\#page=1089}
\expandafter\def\csname EGAExternal@IV.18.6\endcsname{https://github.com/KokunoYumeto/ega-fr/releases/download/ega-fr-2026-10-23-r16/EGA_French_Current_Cumulative_Reader.pdf\#page=1096}
\expandafter\def\csname EGAExternal@IV.18.7\endcsname{https://github.com/KokunoYumeto/ega-fr/releases/download/ega-fr-2026-10-23-r16/EGA_French_Current_Cumulative_Reader.pdf\#page=1101}
\expandafter\def\csname EGAExternal@IV.18.8\endcsname{https://github.com/KokunoYumeto/ega-fr/releases/download/ega-fr-2026-10-23-r16/EGA_French_Current_Cumulative_Reader.pdf\#page=1102}
\expandafter\def\csname EGAExternal@IV.18.9\endcsname{https://github.com/KokunoYumeto/ega-fr/releases/download/ega-fr-2026-10-23-r16/EGA_French_Current_Cumulative_Reader.pdf\#page=1106}
\expandafter\def\csname EGAExternal@IV.19.1\endcsname{https://github.com/KokunoYumeto/ega-fr/releases/download/ega-fr-2026-10-23-r16/EGA_French_Current_Cumulative_Reader.pdf\#page=1130}
\expandafter\def\csname EGAExternal@IV.19.2\endcsname{https://github.com/KokunoYumeto/ega-fr/releases/download/ega-fr-2026-10-23-r16/EGA_French_Current_Cumulative_Reader.pdf\#page=1133}
\expandafter\def\csname EGAExternal@IV.19.3\endcsname{https://github.com/KokunoYumeto/ega-fr/releases/download/ega-fr-2026-10-23-r16/EGA_French_Current_Cumulative_Reader.pdf\#page=1136}
\expandafter\def\csname EGAExternal@IV.19.4\endcsname{https://github.com/KokunoYumeto/ega-fr/releases/download/ega-fr-2026-10-23-r16/EGA_French_Current_Cumulative_Reader.pdf\#page=1138}
\expandafter\def\csname EGAExternal@IV.19.5\endcsname{https://github.com/KokunoYumeto/ega-fr/releases/download/ega-fr-2026-10-23-r16/EGA_French_Current_Cumulative_Reader.pdf\#page=1143}
\expandafter\def\csname EGAExternal@IV.19.6\endcsname{https://github.com/KokunoYumeto/ega-fr/releases/download/ega-fr-2026-10-23-r16/EGA_French_Current_Cumulative_Reader.pdf\#page=1146}
\expandafter\def\csname EGAExternal@IV.2.1.1\endcsname{https://github.com/KokunoYumeto/ega-fr/releases/download/ega-fr-2026-10-23-r16/EGA_French_Current_Cumulative_Reader.pdf\#page=676}
\expandafter\def\csname EGAExternal@IV.2.1.6\endcsname{https://github.com/KokunoYumeto/ega-fr/releases/download/ega-fr-2026-10-23-r16/EGA_French_Current_Cumulative_Reader.pdf\#page=676}
\expandafter\def\csname EGAExternal@IV.2.3.1\endcsname{https://github.com/KokunoYumeto/ega-fr/releases/download/ega-fr-2026-10-23-r16/EGA_French_Current_Cumulative_Reader.pdf\#page=682}
\expandafter\def\csname EGAExternal@IV.2.3.10\endcsname{https://github.com/KokunoYumeto/ega-fr/releases/download/ega-fr-2026-10-23-r16/EGA_French_Current_Cumulative_Reader.pdf\#page=682}
\expandafter\def\csname EGAExternal@IV.2.3.4\endcsname{https://github.com/KokunoYumeto/ega-fr/releases/download/ega-fr-2026-10-23-r16/EGA_French_Current_Cumulative_Reader.pdf\#page=682}
\expandafter\def\csname EGAExternal@IV.2.3.7\endcsname{https://github.com/KokunoYumeto/ega-fr/releases/download/ega-fr-2026-10-23-r16/EGA_French_Current_Cumulative_Reader.pdf\#page=682}
\expandafter\def\csname EGAExternal@IV.2.4.11\endcsname{https://github.com/KokunoYumeto/ega-fr/releases/download/ega-fr-2026-10-23-r16/EGA_French_Current_Cumulative_Reader.pdf\#page=685}
\expandafter\def\csname EGAExternal@IV.2.4.3\endcsname{https://github.com/KokunoYumeto/ega-fr/releases/download/ega-fr-2026-10-23-r16/EGA_French_Current_Cumulative_Reader.pdf\#page=685}
\expandafter\def\csname EGAExternal@IV.2.4.5\endcsname{https://github.com/KokunoYumeto/ega-fr/releases/download/ega-fr-2026-10-23-r16/EGA_French_Current_Cumulative_Reader.pdf\#page=686}
\expandafter\def\csname EGAExternal@IV.2.5.2\endcsname{https://github.com/KokunoYumeto/ega-fr/releases/download/ega-fr-2026-10-23-r16/EGA_French_Current_Cumulative_Reader.pdf\#page=688}
\expandafter\def\csname EGAExternal@IV.2.6.1\endcsname{https://github.com/KokunoYumeto/ega-fr/releases/download/ega-fr-2026-10-23-r16/EGA_French_Current_Cumulative_Reader.pdf\#page=692}
\expandafter\def\csname EGAExternal@IV.21.10\endcsname{https://github.com/KokunoYumeto/ega-fr/releases/download/ega-fr-2026-10-23-r16/EGA_French_Current_Cumulative_Reader.pdf\#page=1200}
\expandafter\def\csname EGAExternal@IV.21.2\endcsname{https://github.com/KokunoYumeto/ega-fr/releases/download/ega-fr-2026-10-23-r16/EGA_French_Current_Cumulative_Reader.pdf\#page=1179}
\expandafter\def\csname EGAExternal@IV.21.3\endcsname{https://github.com/KokunoYumeto/ega-fr/releases/download/ega-fr-2026-10-23-r16/EGA_French_Current_Cumulative_Reader.pdf\#page=1182}
\expandafter\def\csname EGAExternal@IV.21.8\endcsname{https://github.com/KokunoYumeto/ega-fr/releases/download/ega-fr-2026-10-23-r16/EGA_French_Current_Cumulative_Reader.pdf\#page=1194}
\expandafter\def\csname EGAExternal@IV.3.1.5\endcsname{https://github.com/KokunoYumeto/ega-fr/releases/download/ega-fr-2026-10-23-r16/EGA_French_Current_Cumulative_Reader.pdf\#page=700}
\expandafter\def\csname EGAExternal@IV.3.1.6\endcsname{https://github.com/KokunoYumeto/ega-fr/releases/download/ega-fr-2026-10-23-r16/EGA_French_Current_Cumulative_Reader.pdf\#page=699}
\expandafter\def\csname EGAExternal@IV.3.1.7\endcsname{https://github.com/KokunoYumeto/ega-fr/releases/download/ega-fr-2026-10-23-r16/EGA_French_Current_Cumulative_Reader.pdf\#page=700}
\expandafter\def\csname EGAExternal@IV.3.1.9\endcsname{https://github.com/KokunoYumeto/ega-fr/releases/download/ega-fr-2026-10-23-r16/EGA_French_Current_Cumulative_Reader.pdf\#page=699}
\expandafter\def\csname EGAExternal@IV.3.2.6\endcsname{https://github.com/KokunoYumeto/ega-fr/releases/download/ega-fr-2026-10-23-r16/EGA_French_Current_Cumulative_Reader.pdf\#page=701}
\expandafter\def\csname EGAExternal@IV.3.3.1\endcsname{https://github.com/KokunoYumeto/ega-fr/releases/download/ega-fr-2026-10-23-r16/EGA_French_Current_Cumulative_Reader.pdf\#page=703}
\expandafter\def\csname EGAExternal@IV.4.1.4\endcsname{https://github.com/KokunoYumeto/ega-fr/releases/download/ega-fr-2026-10-23-r16/EGA_French_Current_Cumulative_Reader.pdf\#page=710}
\expandafter\def\csname EGAExternal@IV.4.2.1\endcsname{https://github.com/KokunoYumeto/ega-fr/releases/download/ega-fr-2026-10-23-r16/EGA_French_Current_Cumulative_Reader.pdf\#page=711}
\expandafter\def\csname EGAExternal@IV.4.4.2\endcsname{https://github.com/KokunoYumeto/ega-fr/releases/download/ega-fr-2026-10-23-r16/EGA_French_Current_Cumulative_Reader.pdf\#page=715}
\expandafter\def\csname EGAExternal@IV.4.5.11\endcsname{https://github.com/KokunoYumeto/ega-fr/releases/download/ega-fr-2026-10-23-r16/EGA_French_Current_Cumulative_Reader.pdf\#page=717}
\expandafter\def\csname EGAExternal@IV.4.5.19\endcsname{https://github.com/KokunoYumeto/ega-fr/releases/download/ega-fr-2026-10-23-r16/EGA_French_Current_Cumulative_Reader.pdf\#page=719}
\expandafter\def\csname EGAExternal@IV.4.7.1\endcsname{https://github.com/KokunoYumeto/ega-fr/releases/download/ega-fr-2026-10-23-r16/EGA_French_Current_Cumulative_Reader.pdf\#page=726}
\expandafter\def\csname EGAExternal@IV.4.7.3\endcsname{https://github.com/KokunoYumeto/ega-fr/releases/download/ega-fr-2026-10-23-r16/EGA_French_Current_Cumulative_Reader.pdf\#page=726}
\expandafter\def\csname EGAExternal@IV.5.1.2\endcsname{https://github.com/KokunoYumeto/ega-fr/releases/download/ega-fr-2026-10-23-r16/EGA_French_Current_Cumulative_Reader.pdf\#page=734}
\expandafter\def\csname EGAExternal@IV.5.1.3.1\endcsname{https://github.com/KokunoYumeto/ega-fr/releases/download/ega-fr-2026-10-23-r16/EGA_French_Current_Cumulative_Reader.pdf\#page=734}
\expandafter\def\csname EGAExternal@IV.5.10.1\endcsname{https://github.com/KokunoYumeto/ega-fr/releases/download/ega-fr-2026-10-23-r16/EGA_French_Current_Cumulative_Reader.pdf\#page=753}
\expandafter\def\csname EGAExternal@IV.5.10.2\endcsname{https://github.com/KokunoYumeto/ega-fr/releases/download/ega-fr-2026-10-23-r16/EGA_French_Current_Cumulative_Reader.pdf\#page=753}
\expandafter\def\csname EGAExternal@IV.5.10.5\endcsname{https://github.com/KokunoYumeto/ega-fr/releases/download/ega-fr-2026-10-23-r16/EGA_French_Current_Cumulative_Reader.pdf\#page=754}
\expandafter\def\csname EGAExternal@IV.5.11.2\endcsname{https://github.com/KokunoYumeto/ega-fr/releases/download/ega-fr-2026-10-23-r16/EGA_French_Current_Cumulative_Reader.pdf\#page=759}
\expandafter\def\csname EGAExternal@IV.5.11.4\endcsname{https://github.com/KokunoYumeto/ega-fr/releases/download/ega-fr-2026-10-23-r16/EGA_French_Current_Cumulative_Reader.pdf\#page=760}
\expandafter\def\csname EGAExternal@IV.5.2.3\endcsname{https://github.com/KokunoYumeto/ega-fr/releases/download/ega-fr-2026-10-23-r16/EGA_French_Current_Cumulative_Reader.pdf\#page=737}
\expandafter\def\csname EGAExternal@IV.5.5.4\endcsname{https://github.com/KokunoYumeto/ega-fr/releases/download/ega-fr-2026-10-23-r16/EGA_French_Current_Cumulative_Reader.pdf\#page=687}
\expandafter\def\csname EGAExternal@IV.5.6.1\endcsname{https://github.com/KokunoYumeto/ega-fr/releases/download/ega-fr-2026-10-23-r16/EGA_French_Current_Cumulative_Reader.pdf\#page=741}
\expandafter\def\csname EGAExternal@IV.5.6.11.1\endcsname{https://github.com/KokunoYumeto/ega-fr/releases/download/ega-fr-2026-10-23-r16/EGA_French_Current_Cumulative_Reader.pdf\#page=744}
\expandafter\def\csname EGAExternal@IV.5.6.5\endcsname{https://github.com/KokunoYumeto/ega-fr/releases/download/ega-fr-2026-10-23-r16/EGA_French_Current_Cumulative_Reader.pdf\#page=743}
\expandafter\def\csname EGAExternal@IV.5.7.2\endcsname{https://github.com/KokunoYumeto/ega-fr/releases/download/ega-fr-2026-10-23-r16/EGA_French_Current_Cumulative_Reader.pdf\#page=745}
\expandafter\def\csname EGAExternal@IV.5.7.5\endcsname{https://github.com/KokunoYumeto/ega-fr/releases/download/ega-fr-2026-10-23-r16/EGA_French_Current_Cumulative_Reader.pdf\#page=746}
\expandafter\def\csname EGAExternal@IV.5.8.2\endcsname{https://github.com/KokunoYumeto/ega-fr/releases/download/ega-fr-2026-10-23-r16/EGA_French_Current_Cumulative_Reader.pdf\#page=687}
\expandafter\def\csname EGAExternal@IV.5.8.5\endcsname{https://github.com/KokunoYumeto/ega-fr/releases/download/ega-fr-2026-10-23-r16/EGA_French_Current_Cumulative_Reader.pdf\#page=687}
\expandafter\def\csname EGAExternal@IV.5.8.6\endcsname{https://github.com/KokunoYumeto/ega-fr/releases/download/ega-fr-2026-10-23-r16/EGA_French_Current_Cumulative_Reader.pdf\#page=687}
\expandafter\def\csname EGAExternal@IV.5.9.8\endcsname{https://github.com/KokunoYumeto/ega-fr/releases/download/ega-fr-2026-10-23-r16/EGA_French_Current_Cumulative_Reader.pdf\#page=751}
\expandafter\def\csname EGAExternal@IV.6.1.1\endcsname{https://github.com/KokunoYumeto/ega-fr/releases/download/ega-fr-2026-10-23-r16/EGA_French_Current_Cumulative_Reader.pdf\#page=767}
\expandafter\def\csname EGAExternal@IV.6.1.3\endcsname{https://github.com/KokunoYumeto/ega-fr/releases/download/ega-fr-2026-10-23-r16/EGA_French_Current_Cumulative_Reader.pdf\#page=768}
\expandafter\def\csname EGAExternal@IV.6.1.5\endcsname{https://github.com/KokunoYumeto/ega-fr/releases/download/ega-fr-2026-10-23-r16/EGA_French_Current_Cumulative_Reader.pdf\#page=768}
\expandafter\def\csname EGAExternal@IV.6.10.1\endcsname{https://github.com/KokunoYumeto/ega-fr/releases/download/ega-fr-2026-10-23-r16/EGA_French_Current_Cumulative_Reader.pdf\#page=781}
\expandafter\def\csname EGAExternal@IV.6.14.1\endcsname{https://github.com/KokunoYumeto/ega-fr/releases/download/ega-fr-2026-10-23-r16/EGA_French_Current_Cumulative_Reader.pdf\#page=792}
\expandafter\def\csname EGAExternal@IV.6.14.4\endcsname{https://github.com/KokunoYumeto/ega-fr/releases/download/ega-fr-2026-10-23-r16/EGA_French_Current_Cumulative_Reader.pdf\#page=795}
\expandafter\def\csname EGAExternal@IV.6.15.2\endcsname{https://github.com/KokunoYumeto/ega-fr/releases/download/ega-fr-2026-10-23-r16/EGA_French_Current_Cumulative_Reader.pdf\#page=796}
\expandafter\def\csname EGAExternal@IV.6.15.4\endcsname{https://github.com/KokunoYumeto/ega-fr/releases/download/ega-fr-2026-10-23-r16/EGA_French_Current_Cumulative_Reader.pdf\#page=797}
\expandafter\def\csname EGAExternal@IV.6.15.7\endcsname{https://github.com/KokunoYumeto/ega-fr/releases/download/ega-fr-2026-10-23-r16/EGA_French_Current_Cumulative_Reader.pdf\#page=798}
\expandafter\def\csname EGAExternal@IV.6.2.3\endcsname{https://github.com/KokunoYumeto/ega-fr/releases/download/ega-fr-2026-10-23-r16/EGA_French_Current_Cumulative_Reader.pdf\#page=769}
\expandafter\def\csname EGAExternal@IV.6.3.1\endcsname{https://github.com/KokunoYumeto/ega-fr/releases/download/ega-fr-2026-10-23-r16/EGA_French_Current_Cumulative_Reader.pdf\#page=769}
\expandafter\def\csname EGAExternal@IV.6.5.1\endcsname{https://github.com/KokunoYumeto/ega-fr/releases/download/ega-fr-2026-10-23-r16/EGA_French_Current_Cumulative_Reader.pdf\#page=773}
\expandafter\def\csname EGAExternal@IV.6.5.5\endcsname{https://github.com/KokunoYumeto/ega-fr/releases/download/ega-fr-2026-10-23-r16/EGA_French_Current_Cumulative_Reader.pdf\#page=773}
\expandafter\def\csname EGAExternal@IV.6.7.2\endcsname{https://github.com/KokunoYumeto/ega-fr/releases/download/ega-fr-2026-10-23-r16/EGA_French_Current_Cumulative_Reader.pdf\#page=775}
\expandafter\def\csname EGAExternal@IV.7.2.9\endcsname{https://github.com/KokunoYumeto/ega-fr/releases/download/ega-fr-2026-10-23-r16/EGA_French_Current_Cumulative_Reader.pdf\#page=806}
\expandafter\def\csname EGAExternal@IV.7.3.19\endcsname{https://github.com/KokunoYumeto/ega-fr/releases/download/ega-fr-2026-10-23-r16/EGA_French_Current_Cumulative_Reader.pdf\#page=810}
\expandafter\def\csname EGAExternal@IV.7.5.4\endcsname{https://github.com/KokunoYumeto/ega-fr/releases/download/ega-fr-2026-10-23-r16/EGA_French_Current_Cumulative_Reader.pdf\#page=817}
\expandafter\def\csname EGAExternal@IV.7.6.4\endcsname{https://github.com/KokunoYumeto/ega-fr/releases/download/ega-fr-2026-10-23-r16/EGA_French_Current_Cumulative_Reader.pdf\#page=820}
\expandafter\def\csname EGAExternal@IV.7.7.2\endcsname{https://github.com/KokunoYumeto/ega-fr/releases/download/ega-fr-2026-10-23-r16/EGA_French_Current_Cumulative_Reader.pdf\#page=822}
\expandafter\def\csname EGAExternal@IV.7.8.2\endcsname{https://github.com/KokunoYumeto/ega-fr/releases/download/ega-fr-2026-10-23-r16/EGA_French_Current_Cumulative_Reader.pdf\#page=823}
\expandafter\def\csname EGAExternal@IV.7.8.3\endcsname{https://github.com/KokunoYumeto/ega-fr/releases/download/ega-fr-2026-10-23-r16/EGA_French_Current_Cumulative_Reader.pdf\#page=824}
\expandafter\def\csname EGAExternal@IV.7.8.4\endcsname{https://github.com/KokunoYumeto/ega-fr/releases/download/ega-fr-2026-10-23-r16/EGA_French_Current_Cumulative_Reader.pdf\#page=825}
\expandafter\def\csname EGAExternal@IV.7.8.6\endcsname{https://github.com/KokunoYumeto/ega-fr/releases/download/ega-fr-2026-10-23-r16/EGA_French_Current_Cumulative_Reader.pdf\#page=825}
\expandafter\def\csname EGAExternal@IV.8.1.2\endcsname{https://github.com/KokunoYumeto/ega-fr/releases/download/ega-fr-2026-10-23-r16/EGA_French_Current_Cumulative_Reader.pdf\#page=839}
\expandafter\def\csname EGAExternal@IV.8.10.1\endcsname{https://github.com/KokunoYumeto/ega-fr/releases/download/ega-fr-2026-10-23-r16/EGA_French_Current_Cumulative_Reader.pdf\#page=859}
\expandafter\def\csname EGAExternal@IV.8.10.3\endcsname{https://github.com/KokunoYumeto/ega-fr/releases/download/ega-fr-2026-10-23-r16/EGA_French_Current_Cumulative_Reader.pdf\#page=859}
\expandafter\def\csname EGAExternal@IV.8.10.5\endcsname{https://github.com/KokunoYumeto/ega-fr/releases/download/ega-fr-2026-10-23-r16/EGA_French_Current_Cumulative_Reader.pdf\#page=859}
\expandafter\def\csname EGAExternal@IV.8.11.1\endcsname{https://github.com/KokunoYumeto/ega-fr/releases/download/ega-fr-2026-10-23-r16/EGA_French_Current_Cumulative_Reader.pdf\#page=862}
\expandafter\def\csname EGAExternal@IV.8.12.10\endcsname{https://github.com/KokunoYumeto/ega-fr/releases/download/ega-fr-2026-10-23-r16/EGA_French_Current_Cumulative_Reader.pdf\#page=866}
\expandafter\def\csname EGAExternal@IV.8.12.11\endcsname{https://github.com/KokunoYumeto/ega-fr/releases/download/ega-fr-2026-10-23-r16/EGA_French_Current_Cumulative_Reader.pdf\#page=866}
\expandafter\def\csname EGAExternal@IV.8.12.3\endcsname{https://github.com/KokunoYumeto/ega-fr/releases/download/ega-fr-2026-10-23-r16/EGA_French_Current_Cumulative_Reader.pdf\#page=863}
\expandafter\def\csname EGAExternal@IV.8.2.11\endcsname{https://github.com/KokunoYumeto/ega-fr/releases/download/ega-fr-2026-10-23-r16/EGA_French_Current_Cumulative_Reader.pdf\#page=843}
\expandafter\def\csname EGAExternal@IV.8.2.12\endcsname{https://github.com/KokunoYumeto/ega-fr/releases/download/ega-fr-2026-10-23-r16/EGA_French_Current_Cumulative_Reader.pdf\#page=841}
\expandafter\def\csname EGAExternal@IV.8.2.2\endcsname{https://github.com/KokunoYumeto/ega-fr/releases/download/ega-fr-2026-10-23-r16/EGA_French_Current_Cumulative_Reader.pdf\#page=840}
\expandafter\def\csname EGAExternal@IV.8.2.9\endcsname{https://github.com/KokunoYumeto/ega-fr/releases/download/ega-fr-2026-10-23-r16/EGA_French_Current_Cumulative_Reader.pdf\#page=841}
\expandafter\def\csname EGAExternal@IV.8.3.2\endcsname{https://github.com/KokunoYumeto/ega-fr/releases/download/ega-fr-2026-10-23-r16/EGA_French_Current_Cumulative_Reader.pdf\#page=843}
\expandafter\def\csname EGAExternal@IV.8.5.2\endcsname{https://github.com/KokunoYumeto/ega-fr/releases/download/ega-fr-2026-10-23-r16/EGA_French_Current_Cumulative_Reader.pdf\#page=848}
\expandafter\def\csname EGAExternal@IV.8.5.2.5\endcsname{https://github.com/KokunoYumeto/ega-fr/releases/download/ega-fr-2026-10-23-r16/EGA_French_Current_Cumulative_Reader.pdf\#page=849}
\expandafter\def\csname EGAExternal@IV.8.5.5\endcsname{https://github.com/KokunoYumeto/ega-fr/releases/download/ega-fr-2026-10-23-r16/EGA_French_Current_Cumulative_Reader.pdf\#page=849}
\expandafter\def\csname EGAExternal@IV.8.5.8\endcsname{https://github.com/KokunoYumeto/ega-fr/releases/download/ega-fr-2026-10-23-r16/EGA_French_Current_Cumulative_Reader.pdf\#page=851}
\expandafter\def\csname EGAExternal@IV.8.6.3\endcsname{https://github.com/KokunoYumeto/ega-fr/releases/download/ega-fr-2026-10-23-r16/EGA_French_Current_Cumulative_Reader.pdf\#page=852}
\expandafter\def\csname EGAExternal@IV.8.7.1\endcsname{https://github.com/KokunoYumeto/ega-fr/releases/download/ega-fr-2026-10-23-r16/EGA_French_Current_Cumulative_Reader.pdf\#page=853}
\expandafter\def\csname EGAExternal@IV.8.8.1\endcsname{https://github.com/KokunoYumeto/ega-fr/releases/download/ega-fr-2026-10-23-r16/EGA_French_Current_Cumulative_Reader.pdf\#page=854}
\expandafter\def\csname EGAExternal@IV.8.8.2\endcsname{https://github.com/KokunoYumeto/ega-fr/releases/download/ega-fr-2026-10-23-r16/EGA_French_Current_Cumulative_Reader.pdf\#page=854}
\expandafter\def\csname EGAExternal@IV.8.9.1\endcsname{https://github.com/KokunoYumeto/ega-fr/releases/download/ega-fr-2026-10-23-r16/EGA_French_Current_Cumulative_Reader.pdf\#page=857}
\expandafter\def\csname EGAExternal@IV.9.5.1\endcsname{https://github.com/KokunoYumeto/ega-fr/releases/download/ega-fr-2026-10-23-r16/EGA_French_Current_Cumulative_Reader.pdf\#page=879}
\expandafter\def\csname EGAExternal@IV.9.6.4\endcsname{https://github.com/KokunoYumeto/ega-fr/releases/download/ega-fr-2026-10-23-r16/EGA_French_Current_Cumulative_Reader.pdf\#page=883}
\expandafter\def\csname EGAExternal@section:I.7\endcsname{https://github.com/KokunoYumeto/ega-fr/releases/download/ega-fr-2026-10-23-r16/EGA_French_Current_Cumulative_Reader.pdf\#page=109}
\providecommand{\sgaStableOrOriginalLink}[2]{}
\renewcommand{\sgaStableOrOriginalLink}[2]{%
  \ifcsname EGAExternal@#1\endcsname
    \href{\csname EGAExternal@#1\endcsname}{#2}%
  \else
    \@ifundefined{r@#1}{#2}{\hyperref[#1]{#2}}%
  \fi}
\newcommand{\EGAReadLink}[2]{\sgaStableOrOriginalLink{#1}{#2}}
\RenewDocumentCommand{\sref}{o m o}{%
  \IfNoValueTF{#1}%
    {\IfNoValueTF{#3}%
      {\EGAReadLink{#2}{\normalfont{(\StrBehind{#2}{.})}}}%
      {\EGAReadLink{#2}{\normalfont{(\StrBehind{#2}{.},\nobreakspace{}{#3})}}}}%
    {\IfNoValueTF{#3}%
      {\EGAReadLink{#2}{\normalfont{(\textbf{#1},\nobreakspace{}\StrBehind{#2}{.})}}}%
      {\EGAReadLink{#2}{\normalfont{(\textbf{#1},\nobreakspace{}\StrBehind{#2}{.},\nobreakspace{}{#3})}}}}}
\makeatother


\setlength{\emergencystretch}{6em}
\DeclareFontShape{T1}{cmr}{m}{scit}{<->ssub*cmr/m/scsl}{}
\begin{document}
% BEGIN RIGHTS AND AI TRANSLATION NOTICE
\begingroup
\small

% BEGIN INHERITED HUMAN SOURCE CREDIT
\noindent\textbf{Fuentes y contribuciones}

La obra original es de Alexandre Grothendieck, con la colaboración de Jean Dieudonné: Éléments de géométrie algébrique IV, quatrième partie, Publications Mathématiques de l’IHÉS 32 (1967).

Las fuentes inglesas utilizadas en la preparación de esta traducción española incorporan la traducción humana de \href{https://github.com/ryankeleti/ega}{ryankeleti/ega}, versión \href{https://github.com/ryankeleti/ega/tree/3d7afe28c99d1887ccdb85acf2ca1afd93fb1971}{3d7afe28}, en el § 16 y las partes iniciales de los §§ 17 y 20 del capítulo IV. El proyecto nombra a Mateo Carmona, Tim Holzschuh, Tim Hosgood, Ryan Keleti, Thiago Solovera \& Nery y Wayne Yang como colaboradores; sus \href{https://github.com/ryankeleti/ega/blob/3d7afe28c99d1887ccdb85acf2ca1afd93fb1971/README.md}{créditos por sección} y su \href{https://github.com/ryankeleti/ega/commits/3d7afe28c99d1887ccdb85acf2ca1afd93fb1971}{historial de contribuciones} precisan la autoría.

La traducción inglesa y la estructura TeX heredadas conservan esa procedencia. La traducción al español, las adiciones posteriores y las correcciones y composición producidas por IA se distinguen de esas contribuciones humanas. No se atribuye a los colaboradores de la traducción inglesa la traducción española ni el material posterior que no figura en su versión. La atribución no implica autorización ni aval.

\par\medskip
% END INHERITED HUMAN SOURCE CREDIT
\noindent\textbf{Derechos, disponibilidad y autoría}

Los derechos sobre las obras originales y sobre las contribuciones protegidas a estas traducciones y transcripciones derivadas corresponden a sus respectivos titulares. Todo aquello sobre lo que el editor tenga derechos en estas ediciones queda dedicado al dominio público bajo CC0 1.0 (https://creativecommons.org/publicdomain/zero/1.0/deed.es). Las menciones de derechos de terceros y las autorizaciones ya concedidas siguen siendo aplicables.

Estas ediciones se ofrecen gratuitamente, con fines no comerciales y exclusivamente educativos, bajo la presunción de que los respectivos titulares de derechos no se opondrían a que estas traducciones y transcripciones estuvieran disponibles gratuitamente junto a los originales franceses en NUMDAM (www.numdam.org/). Esta presunción no significa que se haya obtenido autorización; el acceso gratuito a los originales no autoriza por sí solo su traducción, adaptación o republicación.

Si usted es titular de derechos sobre algún material presentado aquí y se opone a su disponibilidad, puede abrir una incidencia en GitHub (github.com/KokunoYumeto/ega-fga-es/issues). El responsable de la publicación retirará el material afectado a petición del titular de los derechos.

La traducción al español está producida por IA. Las correcciones, la edición y la composición también están producidas por IA. El responsable de la publicación no tradujo el texto. Se conservan la autoría y las atribuciones de las obras originales y de las contribuciones heredadas; no se atribuyen a la IA. Los modelos documentados se identifican en cada edición.
\par
\endgroup
\clearpage
% END RIGHTS AND AI TRANSLATION NOTICE

\EGAIVFourExternalLabel{0.4.3.3}
\EGAIVFourExternalLabel{0.5.4.3}
\EGAIVFourExternalLabel{0.5.4.5}
\EGAIVFourExternalLabel{0.5.4.7}
\EGAIVFourExternalLabel{0.6.1.2}
\EGAIVFourExternalLabel{0.16.4.6}
\EGAIVFourExternalCitation{IV-37}{37}
\setcounter{section}{15}
% Mechanical subsection components; prose is translated in the component files.
\section{Invariantes diferenciales. Morfismos diferencialmente lisos}
\label{section:IV.16}

\oldpage[IV-4]{5}
En este parágrafo presentamos, en forma global, algunas nociones de cálculo
diferencial particularmente útiles en Geometría algebraica. Omitimos numerosos
desarrollos clásicos de la Geometría diferencial (conexiones, transformaciones
infinitesimales asociadas a un campo de vectores, jets, etc.), aunque estas nociones
se expresan de manera particularmente natural en el marco de los esquemas. Omitimos
asimismo aquí los fenómenos propios de la característica $p>0$ (algunos de los cuales
se estudian, en el marco afín, en (0, 21)). Para ciertos complementos sobre el
formalismo diferencial en los preesquemas, el lector podrá consultar las exposiciones
II y VII de [42], así como los capítulos posteriores de este Tratado.

\subsection{Invariantes normales de una inmersión}
\label{IV.16.1}

\begin{env}[16.1.1]
\label{IV.16.1.1}
Sean $(X,\mathcal{O}_X)$, $(Y,\mathcal{O}_Y)$ dos espacios anillados,
$f=(\psi,\theta):Y\to X$ un morfismo de espacios anillados (0\textsubscript{I}, 4.1.1)
tal que el homomorfismo
\[
  \theta^\#:\psi^*(\mathcal{O}_X)\longrightarrow\mathcal{O}_Y
\]
sea sobreyectivo, de modo que $\mathcal{O}_Y$ se identifica con un haz de anillos cociente
$\psi^*(\mathcal{O}_X)/\mathscr{I}_f$. Se puede entonces dotar
$\psi^*(\mathcal{O}_X)$ de la filtración $\mathscr{I}_f$-preádica.
\end{env}

\begin{definition}[16.1.2]
\label{IV.16.1.2}
El haz de anillos aumentado sobre $\mathcal{O}_Y$
$\psi^*(\mathcal{O}_X)/\mathscr{I}_f^{n+1}$ se llama el $n$-ésimo invariante normal
de $f$; el espacio anillado
$(Y,\psi^*(\mathcal{O}_X)/\mathscr{I}_f^{n+1})$ se llama la $n$-ésima vecindad
infinitesimal de $Y$ para el morfismo $f$, y se denota por $Y_f^{(n)}$ o simplemente
$Y^{(n)}$. El haz de anillos graduados asociado al haz de anillos filtrados
$\psi^*(\mathcal{O}_X)$
\[
  \mathscr{G}r_\bullet(f)=\bigoplus_{n\geq 0}
  (\mathscr{I}_f^n/\mathscr{I}_f^{n+1})
  \tag{16.1.2.1}
  \label{IV.16.1.2.1}
\]
se llama el haz de anillos graduados asociado a $f$. El haz
$\mathscr{G}r_1(f)=\mathscr{I}_f/\mathscr{I}_f^2$ se llama el haz conormal
de $f$ (que también se denota por $\mathscr{N}_{Y/X}$ si no hay riesgo de confusión).
\end{definition}

Está claro que los
$\mathcal{O}_{Y^{(n)}}=\psi^*(\mathcal{O}_X)/\mathscr{I}_f^{n+1}$
(que también se denotan por $\mathcal{O}_{Y_f^{(n)}}$) forman un
\oldpage[IV-4]{6}
sistema proyectivo de haces de anillos sobre $Y$; el homomorfismo de transición
$\varphi_{nm}:\mathcal{O}_{Y^{(m)}}\to\mathcal{O}_{Y^{(n)}}$ para $n\leq m$
identifica $\mathcal{O}_{Y^{(n)}}$ con el cociente de $\mathcal{O}_{Y^{(m)}}$ por la
potencia $(\mathscr{I}_f/\mathscr{I}_f^{n+1})^m$ del ideal de aumentación de
$\mathcal{O}_{Y^{(n)}}$, núcleo de
$\varphi_{0n}:\mathcal{O}_{Y^{(n)}}\to\mathcal{O}_Y$. Por consiguiente, los $Y^{(n)}$ forman
un sistema inductivo de espacios anillados, todos ellos con espacio subyacente
$Y$, y se tienen morfismos canónicos de espacios anillados
$h_n:Y^{(n)}\to X$ iguales a $(\psi,\theta_n)$, donde $\theta_n^\#$ es el morfismo
canónico
$\psi^*(\mathcal{O}_X)\to\psi^*(\mathcal{O}_X)/\mathscr{I}_f^{n+1}$.
Está claro que el haz $\mathscr{G}r_\bullet(f)$ es un haz de álgebras
graduadas sobre el haz de anillos
$\mathcal{O}_Y=\mathscr{G}r_0(f)$, y los $\mathscr{G}r_k(f)$ son
$\mathcal{O}_Y$-módulos.

Como para todo haz de anillos filtrado, se tiene un homomorfismo canónico
sobreyectivo de $\mathcal{O}_Y$-álgebras graduadas
\[
  \mathscr{S}_{\mathcal{O}_Y}^{\bullet}(\mathscr{G}r_1(f))
  \longrightarrow\mathscr{G}r_\bullet(f)
  \tag{16.1.2.2}
  \label{IV.16.1.2.2}
\]
que coincide en grados 0 y 1 con los homomorfismos identidad.

\begin{examples}[16.1.3]
\label{IV.16.1.3}
\begin{enumerate}
\item[(i)] Supongamos que $X$ es un espacio localmente anillado, que $Y$ está
reducido a un único punto $y$ (dotado de un anillo $\mathcal{O}_y$) y que, si
$x=\psi(y)$, $\theta^\#:\mathcal{O}_x\to\mathcal{O}_y$ es un homomorfismo
sobreyectivo de anillos cuyo núcleo es el ideal maximal $\mathfrak{m}_x$ de
$\mathcal{O}_x$. Entonces los $\mathcal{O}_{Y^{(n)}}$ se identifican con los anillos
$\mathcal{O}_x/\mathfrak{m}_x^{n+1}$, y $\mathscr{G}r_\bullet(f)$ con el anillo
graduado asociado al anillo local $\mathcal{O}_x$ dotado de su filtración
$\mathfrak{m}_x$-preádica.
\item[(ii)] Supongamos que $Y$ es un subconjunto cerrado de un subespacio abierto $U$
de $X$ y que $\mathcal{O}_Y$ está inducido sobre $Y$ por un haz cociente
$\mathcal{O}_U/\mathscr{I}$, donde $\mathscr{I}$ es un ideal de $\mathcal{O}_U$
tal que $\mathscr{I}_x=\mathcal{O}_x$ para todo $x\notin Y$; si $X$ es un
espacio localmente anillado, supondremos además que
$\mathscr{I}_x\neq\mathcal{O}_x$ para $x\in Y$, de modo que
$(Y,\mathcal{O}_Y)$ sigue siendo un espacio localmente anillado.

Sea $\psi_0:Y\to U$ la inyección canónica, y denotemos por
$\theta_0:\mathcal{O}_U\to(\psi_0)_*(\mathcal{O}_Y)$ el homomorfismo tal que
$\theta_0^\#$ sea el homomorfismo canónico
$\psi_0^*(\mathcal{O}_U)=\mathcal{O}_U|Y\to
(\mathcal{O}_U/\mathscr{I})|Y$, de modo que
$j_0=(\psi_0,\theta_0):Y\to U$ es un morfismo de espacios anillados (y de espacios
localmente anillados si $X$ es un espacio localmente anillado); si
$i:U\to X$ es la inyección canónica (morfismo de espacios anillados),
$j=i\circ j_0$ es el morfismo $(\psi,\theta)$ de $Y$ en $X$, donde
$\psi:Y\to X$ es la inyección canónica, y
$\theta:\mathcal{O}_X\to\psi_*(\mathcal{O}_Y)$ es el homomorfismo tal que
$\theta^\#=\theta_0^\#$. Como $\theta^\#$ es sobreyectivo, se pueden aplicar las
definiciones precedentes; $\mathcal{O}_{Y^{(n)}}$ es igual a
$\psi_0^*(\mathcal{O}_U/\mathscr{I}^{n+1})$, y se tiene
$(\psi_0)_*(\mathcal{O}_{Y^{(n)}})=\mathcal{O}_U/\mathscr{I}^{n+1}$, y
$\mathscr{G}r_n(j)=\mathscr{G}r_n(j_0)=
\psi_0^*(\mathscr{I}^n/\mathscr{I}^{n+1})=
j_0^*(\mathscr{I}^n/\mathscr{I}^{n+1})$.
\end{enumerate}
\end{examples}

\begin{env}[16.1.4]
\label{IV.16.1.4}
El ejemplo (16.1.3, (ii)) muestra que, en general, los
$\mathcal{O}_{Y^{(n)}}$ no están dotados canónicamente de una estructura de
$\mathcal{O}_Y$-módulo, ni, a fortiori, de una estructura de
$\mathcal{O}_Y$-álgebra. El dato de tal estructura equivale al de un
homomorfismo de haces de anillos
$\lambda_n:\mathcal{O}_Y\to\mathcal{O}_{Y^{(n)}}$, inverso por la derecha del
homomorfismo de aumentación $\varphi_{0n}$; también equivale al dato
de un morfismo de espacios anillados
$(1_Y,\lambda_n):Y^{(n)}\to Y$, inverso por la izquierda del morfismo canónico
$(1_Y,\varphi_{0n}):Y\to Y^{(n)}$.
\end{env}

\begin{proposition}[16.1.5]
\label{IV.16.1.5}
Sea $f=(\psi,\theta):Y\to X$ una inmersión de preesquemas. Entonces:
\begin{enumerate}
\item[(i)] $\mathscr{G}r_\bullet(f)$ es una $\mathcal{O}_Y$-álgebra graduada
cuasicoherente.
\oldpage[IV-4]{7}
\item[(ii)] Los $Y^{(n)}$ son preesquemas, canónicamente isomorfos a
subpreesquemas de $X$.
\item[(iii)] Todo homomorfismo de haces de anillos
$\lambda_n:\mathcal{O}_Y\to\mathcal{O}_{Y^{(n)}}$, inverso por la derecha del
homomorfismo de aumentación $\varphi_{0n}$, convierte
$\mathcal{O}_{Y^{(n)}}$ y los $\mathcal{O}_{Y^{(k)}}$ para $k\leq n$ en
$\mathcal{O}_Y$-álgebras cuasicoherentes; las estructuras de
$\mathcal{O}_Y$-módulo deducidas de las estructuras precedentes sobre los
$\mathscr{G}r_k(f)$ para $k\leq n$ coinciden con las definidas en
(16.1.2).
\end{enumerate}
\end{proposition}

\noindent(i) Como la cuestión es local sobre $X$ y sobre $Y$, basta considerar el caso
en que $Y$ es un subpreesquema cerrado de $X$ definido por un ideal cuasicoherente
$\mathscr{I}$ de $\mathcal{O}_X$; como $\mathcal{O}_Y$ es la restricción a
$Y$ de $\mathcal{O}_X/\mathscr{I}$, la afirmación (i) es evidente, y $Y^{(n)}$
es el subpreesquema cerrado de $X$ definido por el ideal cuasicoherente
$\mathscr{I}^{n+1}$ de $\mathcal{O}_X$. Finalmente, para probar (iii), observemos que el
dato de $\lambda_n$ dota al ideal
$\mathscr{I}/\mathscr{I}^{n+1}$ de la aumentación $\varphi_{0n}$ y a sus
cocientes $\mathscr{I}/\mathscr{I}^{k+1}$ ($1\leq k\leq n$) de estructuras de
$\mathcal{O}_Y$-módulo, y basta probar por inducción sobre $k$ que los
$\mathscr{I}/\mathscr{I}^{k+1}$ son $\mathcal{O}_Y$-módulos
cuasicoherentes y que la estructura de $\mathcal{O}_Y$-módulo cociente que de ello se
deduce sobre $\mathscr{I}^k/\mathscr{I}^{k+1}$ es la misma que la definida en
(16.1.2). La segunda afirmación es inmediata, pues
$\mathscr{I}^k/\mathscr{I}^{k+1}$ está anulado por
$\mathscr{I}/\mathscr{I}^{n+1}$; la primera resulta, por inducción sobre $k$,
del hecho de que es trivial para $k=1$ y de que
$\mathscr{I}/\mathscr{I}^{k+1}$ es una extensión de
$\mathscr{I}/\mathscr{I}^k$ por
$\mathscr{I}^k/\mathscr{I}^{k+1}$ (III, 1.4.17).

\begin{corollary}[16.1.6]
\label{IV.16.1.6}
Bajo las hipótesis generales de (16.1.5), si la inmersión $f$ es localmente de
presentación finita, los $\mathscr{G}r_n(f)$ son $\mathcal{O}_Y$-módulos
cuasicoherentes de tipo finito.
\end{corollary}

En efecto, con las notaciones de la demostración de (16.1.5), $\mathscr{I}$ es un
ideal de tipo finito de $\mathcal{O}_X$ (1.4.7), por tanto los
$\mathscr{I}^n/\mathscr{I}^{n+1}$ son $\mathcal{O}_Y$-módulos de tipo finito,
de donde se sigue la conclusión.

\begin{corollary}[16.1.7]
\label{IV.16.1.7}
Bajo las hipótesis generales de (16.1.5), sea $g:X\to Y$ un morfismo de
preesquemas que es inverso por la izquierda de $f$. Entonces, para todo $n$, el morfismo compuesto
$(1,\lambda_n):Y^{(n)}\xrightarrow{h_n}X\xrightarrow{g}Y$ define un
homomorfismo de haces de anillos
$\lambda_n:\mathcal{O}_Y\to\mathcal{O}_{Y^{(n)}}$ inverso por la derecha de
la aumentación $\varphi_{0n}$, que convierte a $\mathcal{O}_{Y^{(n)}}$ en una
$\mathcal{O}_Y$-álgebra cuasicoherente; para estos homomorfismos, los
homomorfismos de transición
$\varphi_{nm}:\mathcal{O}_{Y^{(m)}}\to\mathcal{O}_{Y^{(n)}}$ ($n\leq m$) son
homomorfismos de $\mathcal{O}_Y$-álgebras. Además, si $g$ es localmente
de tipo finito, los $\mathcal{O}_{Y^{(n)}}$ son $\mathcal{O}_Y$-módulos
cuasicoherentes de tipo finito.
\end{corollary}

La primera afirmación se sigue inmediatamente de las definiciones y de (16.1.5). Por otra
parte, si $g$ es localmente de tipo finito, entonces $f$ es localmente de
presentación finita (1.4.3, (v)); como los $\mathscr{G}r_n(f)$ son entonces
$\mathcal{O}_Y$-módulos cuasicoherentes de tipo finito por (16.1.6), lo mismo ocurre
con los $\mathcal{O}_Y$-módulos
$\mathscr{I}/\mathscr{I}^{n+1}$, que son extensiones de un número finito de
$\mathscr{G}r_k(f)$ (III, 1.4.17).

\begin{proposition}[16.1.8]
\label{IV.16.1.8}
Sean $X$ un preesquema localmente noetheriano, $j:Y\to X$ una inmersión. Entonces
los $Y^{(n)}$ son preesquemas localmente noetherianos, los
$\mathscr{G}r_n(j)$ son $\mathcal{O}_Y$-módulos coherentes, y
$\mathscr{G}r_\bullet(j)$ es un haz coherente de anillos sobre el espacio $Y$.
\end{proposition}

Como todo es local sobre $X$ y $Y$, se reduce al caso en que $X$ es afín y $j$
una inmersión cerrada; entonces todas las afirmaciones son evidentes salvo la
última, que resulta de que, si $A$ es un anillo noetheriano y
$\mathfrak{I}$ un ideal de $A$, $\operatorname{gr}_{\mathfrak{I}}^\bullet(A)$
es un anillo noetheriano, teniendo en cuenta la exactitud del funtor $\psi^*$ y
(0\textsubscript{I}, 5.3.7).

\begin{proposition}[16.1.9]
\label{IV.16.1.9}
\oldpage[IV-4]{8}
Sean $X$ un preesquema, $j:Y\to X$ una inmersión localmente de presentación
finita, $y$ un punto de $Y$. Las condiciones siguientes son equivalentes:
\begin{enumerate}
\item[(a)] Existe una vecindad abierta $U$ de $y$ en $Y$ tal que $j|U$ es
un homeomorfismo de $U$ sobre un abierto de $X$.
\item[(b)] Existe un entero $n>0$ tal que el homomorfismo canónico
\[
  (\varphi_{n,n-1})_y:
  \mathcal{O}_{Y^{(n)},y}\longrightarrow\mathcal{O}_{Y^{(n-1)},y}
\]
sea biyectivo.
\item[(c)] Existe un entero $n>0$ tal que
$(\mathscr{G}r_n(j))_y=0$.
\end{enumerate}
Además, si el entero $n$ satisface b) o c), existe una vecindad $V$ de $y$
en $Y$ tal que $\mathscr{G}r_m(j)|V=0$ para $m\geq n$ y que
$\varphi_{nm}|V:\mathcal{O}_{Y^{(m)}}|V\to\mathcal{O}_{Y^{(n)}}|V$ sea
biyectivo para $m\geq n$.
\end{proposition}

Como la cuestión es local sobre $Y$, basta considerar el caso en que $j$ es una
inmersión cerrada, estando $Y$ definido por un ideal cuasicoherente de tipo finito
$\mathscr{I}$ de $\mathcal{O}_X$. La equivalencia de b) y c), para un $n$ dado,
es entonces inmediata; además, como
$\mathscr{I}^n/\mathscr{I}^{n+1}$ es un $\mathcal{O}_X$-módulo de tipo finito,
existe una vecindad abierta $U$ de $y$ en $X$ tal que
$\mathscr{I}^n|U=\mathscr{I}^{n+1}|U$ (0\textsubscript{I}, 5.2.2), y por tanto también
$\mathscr{I}^n|U=\mathscr{I}^m|U$ para $m\geq n$, lo que prueba las últimas
afirmaciones. Para probar que a) implica b), basta considerar el caso en que
el espacio subyacente a $Y$ es igual al espacio subyacente a $X$, y en que
$\mathscr{I}$ está engendrado por un número finito de sus secciones sobre $X$:
como $\mathscr{I}$ está entonces contenido en el nilradical $\mathscr{N}$ de
$\mathcal{O}_X$ (I, 5.1.2), es nilpotente, lo que prueba b). Finalmente, para
probar que b) implica a), basta también considerar el caso en que
$\mathscr{I}^n=\mathscr{I}^{n+1}$; entonces, para todo $y\in Y$, como
$\mathscr{I}_y\subset\mathfrak{m}_y$, ideal maximal de
$\mathcal{O}_{X,y}$, se tiene necesariamente $\mathscr{I}_y^n=0$ en virtud del
lema de Nakayama, puesto que $\mathscr{I}_y$ es un ideal de tipo finito.
El conjunto de los $x\in X$ tales que $\mathscr{I}_x^n=0$ es por tanto un abierto $U$ de
$X$ que contiene a $Y$ (0\textsubscript{I}, 5.2.2); como, por otra parte,
$\mathscr{I}_x\neq 0$ para $x\notin Y$, se tiene necesariamente $U=Y$.

\begin{corollary}[16.1.10]
\label{IV.16.1.10}
Para que la restricción de la inmersión $j$ a una vecindad de $y$ en $Y$ sea
una inmersión abierta (es decir, para que $j$ sea un isomorfismo local en el
punto $y$), es necesario y suficiente que
$(\mathscr{G}r_1(j))_y=(\mathscr{N}_{Y/X})_y=0$.
\end{corollary}

La condición es evidentemente necesaria, y el razonamiento precedente, aplicado
para $n=1$, prueba que es suficiente.

\begin{remarks}[16.1.11]
\label{IV.16.1.11}
\begin{enumerate}
\item[(i)] Bajo las condiciones de la definición (16.1.1), el límite proyectivo
del sistema proyectivo
$(\mathcal{O}_{Y^{(n)}},\varphi_{nm})$ de haces de anillos sobre $Y$ se
llama el invariante normal de orden infinito de $f$, y se denota a veces por
$\mathcal{O}_{Y^{(\infty)}}$. Cuando $X$ es un preesquema localmente
noetheriano y $j:Y\to X$ es una inmersión cerrada, siendo por tanto $Y$ un subpreesquema
cerrado de $X$ definido por un ideal coherente $\mathscr{I}$,
$\mathcal{O}_{Y^{(\infty)}}$ no es sino la completación formal de
$\mathcal{O}_X$ a lo largo de $Y$ (I, 10.8.4), y
$Y^{(\infty)}=(Y,\mathcal{O}_{Y^{(\infty)}})$ es el preesquema formal completado de
$X$ a lo largo de $Y$ (I, 10.8.5). En todos los casos, se podrá decir que
$Y^{(\infty)}$ es la vecindad formal de $Y$ en $X$ (para el morfismo
$f$). En el caso particular que acabamos de considerar, es por tanto el preesquema
formal que es límite inductivo de las vecindades infinitesimales de orden $n$.
\item[(ii)] Obsérvese que, para un morfismo de preesquemas
$f=(\psi,\theta):Y\to X$, puede ocurrir que el homomorfismo
$\theta^\#:\psi^*(\mathcal{O}_X)\to\mathcal{O}_Y$ sea sobreyectivo sin que $f$
sea una inmersión
\oldpage[IV-4]{9}
local, y sin que $f$ sea inyectivo. Se obtiene un ejemplo tomando como $Y$
una suma de preesquemas $Y_\lambda$ todos isomorfos a
$\operatorname{Spec}(\mathcal{O}_x)$, donde $x\in X$, y como $f$ el morfismo
cuya restricción a cada uno de los $Y_\lambda$ es el morfismo canónico.
\end{enumerate}
\end{remarks}

\subsection{Propiedades funtoriales de los invariantes normales de una inmersión}
\label{IV.16.2}

\begin{env}[16.2.1]
\label{IV.16.2.1}
Sean $f=(\psi,\theta):Y\to X$,
$f'=(\psi',\theta'):Y'\to X'$ dos morfismos de espacios anillados tales que los
homomorfismos $\theta^\#$ y $\theta'^\#$ sean sobreyectivos; consideremos un
diagrama conmutativo de morfismos de espacios anillados
\[
  \xymatrix{
    Y \ar[r]^f & X \\
    Y' \ar[r]_{f'} \ar[u]^u & X' \ar[u]_v
  }
  \tag{16.2.1.1}
  \label{IV.16.2.1.1}
\]
Pongamos $u=(\rho,\lambda)$, $v=(\sigma,\mu)$. Se tiene
$\rho^*(\psi^*(\mathcal{O}_X))=
\psi'^*(\sigma^*(\mathcal{O}_X))$ y, por consiguiente, un diagrama conmutativo
de homomorfismos de haces de anillos sobre $Y'$
\[
  \xymatrix{
    \rho^*(\psi^*(\mathcal{O}_X))
      =\psi'^*(\sigma^*(\mathcal{O}_X))
      \ar[r]^-{\psi'^*(\mu^\#)}
      \ar[d]_{\rho^*(\theta^\#)}
      & \psi'^*(\mathcal{O}_{X'}) \ar[d]^{\theta'^\#} \\
    \rho^*(\mathcal{O}_Y) \ar[r]_-{\lambda^\#}
      & \mathcal{O}_{Y'}
  }
\]
de donde se concluye, si $\mathscr{I}$ y $\mathscr{I}'$ son los núcleos de
$\theta^\#$ y $\theta'^\#$, que se tiene
$\psi'^*(\mu^\#)(\rho^*(\mathscr{I}))\subset\mathscr{I}'$, teniendo en cuenta la
exactitud del funtor $\rho^*$. Se deduce inmediatamente que, para todo entero
$n$,
$\psi'^*(\mu^\#)(\rho^*(\mathscr{I}^n))\subset\mathscr{I}'^n$, lo que muestra
que $\psi'^*(\mu^\#)$ define, al pasar a los cocientes, un homomorfismo de
haces de anillos
\[
  \nu_n:
  \rho^*(\psi^*(\mathcal{O}_X)/\mathscr{I}^{n+1})
  \longrightarrow
  \psi'^*(\mathcal{O}_{X'})/\mathscr{I}'^{n+1}
  \tag{16.2.1.2}
  \label{IV.16.2.1.2}
\]
y, por consiguiente, un morfismo de espacios anillados
$w_n=(\rho,\nu_n):Y'^{(n)}\to Y^{(n)}$ (que, para $n=0$, no es sino
$u$). De esta definición se sigue inmediatamente que los diagramas
\[
  \xymatrix@R=1pc{
    Y^{(n)} \ar[r]^-{h_{mn}} & Y^{(m)} \ar[r]^-{h_m} & X \\
    & & & (n\leq m) \\
    Y'^{(n)} \ar[r]_-{h'_{mn}} \ar[uu]^-{w_n}
      & Y'^{(m)} \ar[r]_-{h'_m} \ar[uu]^-{w_m}
      & X' \ar[uu]_-v
  }
\]
(donde las flechas horizontales son los morfismos canónicos (16.1.2)) son
conmutativos.

Al pasar a los cocientes a partir de los homomorfismos (16.2.1.2), y teniendo en
\oldpage[IV-4]{10}
cuenta la exactitud del funtor $\rho^*$, se obtiene un dihomomorfismo
de álgebras graduadas (relativo al homomorfismo
$\lambda^\#:\rho^*(\mathcal{O}_Y)\to\mathcal{O}_{Y'}$)
\[
  \operatorname{gr}(u):
  \rho^*(\mathscr{G}r_\bullet(f))
  \longrightarrow\mathscr{G}r_\bullet(f')
  \tag{16.2.1.3}
  \label{IV.16.2.1.3}
\]
(o, si se prefiere, un $\rho$-morfismo (0\textsubscript{I}, 3.5.1)
$\mathscr{G}r_\bullet(f)\to\mathscr{G}r_\bullet(f')$), y en particular un
dihomomorfismo de los haces conormales
\[
  \operatorname{gr}_1(u):
  \rho^*(\mathscr{G}r_1(f))\longrightarrow\mathscr{G}r_1(f').
\]
Es además inmediato que estos homomorfismos dan lugar a un diagrama
conmutativo
\[
  \xymatrix{
    \rho^*(\mathscr{S}_{\mathcal{O}_Y}^{\bullet}
      (\mathscr{G}r_1(f)))
      \ar[r] \ar[d]_-{\mathscr{S}(\operatorname{gr}_1(u))}
      & \rho^*(\mathscr{G}r_\bullet(f))
        \ar[d]^-{\operatorname{gr}(u)} \\
    \mathscr{S}_{\mathcal{O}_{Y'}}^{\bullet}
      (\mathscr{G}r_1(f'))
      \ar[r]
      & \mathscr{G}r_\bullet(f')
  }
  \tag{16.2.1.4}
  \label{IV.16.2.1.4}
\]
donde las flechas horizontales son los homomorfismos canónicos (16.1.2.2).

Finalmente, si se tiene un diagrama conmutativo de morfismos de espacios anillados
\[
  \xymatrix{
    Y \ar[r]^f & X \\
    Y' \ar[r]_{f'} \ar[u]^u & X' \ar[u]_v \\
    Y'' \ar[r]_{f''} \ar[u]^{u'} & X'' \ar[u]_{v'}
  }
\]
donde $f''=(\psi'',\theta'')$ es tal que $\theta''^\#$ es sobreyectivo, y si
$w'_n$ y $w''_n$ se definen a partir de $u'$, $v'$, por una parte, y de
$u''=u\circ u'$, $v''=v\circ v'$, por otra, entonces se tiene
$w''_n=w_n\circ w'_n$, como se sigue inmediatamente de las definiciones y de
(0\textsubscript{I}, 3.5.5); se tiene asimismo
$\operatorname{gr}(u'')=
\operatorname{gr}(u')\circ\rho'^*(\operatorname{gr}(u))$ si
$u'=(\rho',\lambda')$. Se puede, por tanto, decir que los $Y^{(n)}$ y
$\mathscr{G}r_\bullet(f)$ dependen funtorialmente de $f$.
\end{env}

\begin{proposition}[16.2.2]
\label{IV.16.2.2}
Con las notaciones e hipótesis de (16.2.1), supongamos además que $f$, $f'$,
$u$ y $v$ son morfismos de preesquemas. Entonces:
\begin{enumerate}
\item[(i)] Los morfismos $w_n:Y'^{(n)}\to Y^{(n)}$ son morfismos de
preesquemas.
\item[(ii)] Si $Y'=Y\times_X X'$, siendo $u$ y $f'$ las proyecciones
canónicas, y si $f$ es una inmersión o si $v$ es plano, se tiene
$Y'^{(n)}=Y^{(n)}\times_X X'$.
\item[(iii)] Si $Y'=Y\times_X X'$ y si $v$ es plano (resp. si $f$ es una
inmersión), el homomorfismo
\[
  \operatorname{Gr}(u)=\operatorname{gr}(u)\otimes 1:
  \mathscr{G}r_\bullet(f)\otimes_{\mathcal{O}_Y}\mathcal{O}_{Y'}
  \longrightarrow\mathscr{G}r_\bullet(f')
\]
es biyectivo (resp. sobreyectivo).
\end{enumerate}
\end{proposition}

\noindent(i) Las hipótesis implican inmediatamente que, para todo $y'\in Y'$,
$\rho_{y'}^*(\theta_{\psi'(y')}^\#)$ es un homomorfismo local (I, 1.6.2),
por tanto $w_n$ es un morfismo de preesquemas (I, 2.2.1).

\oldpage[IV-4]{11}
\noindent(ii) y (iii) Si $f$ es una inmersión, basta considerar el caso en que
$f$ es una inmersión cerrada, estando entonces $Y$ definido por el ideal
cuasicoherente $\mathscr{I}$ de $\mathcal{O}_X$, y $Y^{(n)}$ por el ideal
$\mathscr{I}^{n+1}$; las afirmaciones resultan entonces de (I, 4.4.5).

Supongamos, en segundo lugar, que $v$ es plano; basta considerar el caso en que
$X=\operatorname{Spec}(A)$, $Y=\operatorname{Spec}(B)$,
$X'=\operatorname{Spec}(A')$ son afines, siendo $A'$ un $A$-módulo plano;
entonces $Y'=\operatorname{Spec}(B')$ con $B'=B\otimes_A A'$; además, si
$\mathfrak{I}$ es el núcleo del homomorfismo $A\to B$, el núcleo
$\mathfrak{I}'$ de $A'\to B'$ se identifica con $\mathfrak{I}\otimes_A A'$ por
planitud, y
\[
  \mathfrak{I}'^n/\mathfrak{I}'^{n+1}
  =(\mathfrak{I}^n/\mathfrak{I}^{n+1})\otimes_A A'.
\]
Se deduce inmediatamente, teniendo en cuenta (0\textsubscript{I}, 4.3.3), que el
$\mathcal{O}_{Y'}$-módulo
$\mathscr{I}'^n/\mathscr{I}'^{n+1}$ es igual a
\begin{align*}
&\psi'^*\!\left(\sigma^*
  \left((\mathfrak{I}^n/\mathfrak{I}^{n+1})^{\sim}\right)\right)
  \otimes_{\sigma^*(\mathcal{O}_X)}\mathcal{O}_{X'} \\
={}&
\psi'^*\!\left(\sigma^*
  \left((\mathfrak{I}^n/\mathfrak{I}^{n+1})^{\sim}\right)\right)
  \otimes_{\psi'^*(\sigma^*(\mathcal{O}_X))}
  \psi'^*(\mathcal{O}_{X'}) \\
={}&
\rho^*(\mathscr{I}^n/\mathscr{I}^{n+1})
  \otimes_{\rho^*(\psi^*(\mathcal{O}_X))}
  \psi'^*(\mathcal{O}_{X'}),
\end{align*}
y como, en particular, para $n=0$, se tiene
\[
  \mathcal{O}_{Y'}=
  \rho^*(\mathcal{O}_Y)
  \otimes_{\rho^*(\psi^*(\mathcal{O}_X))}
  \psi'^*(\mathcal{O}_{X'}),
\]
se obtiene un isomorfismo canónico de
$\mathscr{I}'^n/\mathscr{I}'^{n+1}$ con
\[
  \rho^*(\mathscr{I}^n/\mathscr{I}^{n+1})
  \otimes_{\rho^*(\mathcal{O}_Y)}\mathcal{O}_{Y'}
  =
  (\mathscr{I}^n/\mathscr{I}^{n+1})
  \otimes_{\mathcal{O}_Y}\mathcal{O}_{Y'},
\]
lo que prueba (iii). Pongamos ahora
$C_n=\Gamma(Y,\mathcal{O}_{Y^{(n)}})$,
$C'_n=\Gamma(Y',\mathcal{O}_{Y'^{(n)}})$. Puesto que $Y^{(n)}$ y
$Y'^{(n)}$ son esquemas afines (16.1.5), el núcleo
$\mathfrak{K}_n$ (resp. $\mathfrak{K}'_n$) del homomorfismo
$C_n\to C_{n-1}$ (resp. $C'_n\to C'_{n-1}$) es
$\Gamma(Y,\mathscr{I}^n/\mathscr{I}^{n+1})$ (resp.
$\Gamma(Y',\mathscr{I}'^n/\mathscr{I}'^{n+1})$); por tanto, de lo que
precede se deduce que $\mathfrak{K}'_n=\mathfrak{K}_n\otimes_A A'$. Ahora bien, se tiene un
diagrama conmutativo
\[
  \xymatrix{
    0 \ar[r] &
    \mathfrak{K}_n\otimes_A A' \ar[r]\ar[d]^r &
    C_n\otimes_A A' \ar[r]\ar[d]^{s_n} &
    C_{n-1}\otimes_A A' \ar[r]\ar[d]^{s_{n-1}} & 0 \\
    0 \ar[r] &
    \mathfrak{K}'_n \ar[r] &
    C'_n \ar[r] &
    C'_{n-1} \ar[r] & 0
  }
\]
donde la flecha vertical de la izquierda es biyectiva y las dos filas son exactas
(siendo $A'$ un $A$-módulo plano). Se deduce por inducción que $s_n$ es
biyectivo para todo $n$: esto es cierto por hipótesis para $n=0$, y el caso de un
$n$ cualquiera se deduce aplicando el lema de los cinco. Esto prueba la segunda
afirmación de (ii).

\begin{corollary}[16.2.3]
\label{IV.16.2.3}
Sean $g:X\to Y$, $u:Y'\to Y$ dos morfismos de preesquemas,
$X'=X\times_Y Y'$, $g':X'\to Y'$ y $v:X'\to X$ las proyecciones canónicas.
Sea $f:Y\to X$ una sección sobre $Y$ de $X$ (por tanto, una inmersión),
$f'=f_{(Y')}:Y'\to X'$ la sección sobre $Y'$ de $X'$ deducida de $f$ por el
cambio de base $u$. Entonces:
\begin{enumerate}
\item[(i)] El morfismo
$w_n:Y_{f'}'^{(n)}\to Y_f^{(n)}$ correspondiente a $f$, $f'$, $u$, $v$
(16.2.1) y el morfismo canónico
$h'_n:Y_{f'}'^{(n)}\to X'$ identifican $Y_{f'}'^{(n)}$ con el producto fibrado
$Y_f^{(n)}\times_X X'$.
\item[(ii)] Si se dota a $\mathcal{O}_{Y_f^{(n)}}$ (resp.
$\mathcal{O}_{Y_{f'}'^{(n)}}$) de la estructura de $\mathcal{O}_Y$-álgebra
definida por $g$ (resp. de la estructura de $\mathcal{O}_{Y'}$-álgebra definida
por $g'$) (16.1.6), el homomorfismo de $\mathcal{O}_{Y'}$-álgebras
\[
  \rho^*(\mathcal{O}_{Y_f^{(n)}})
  \otimes_{\mathcal{O}_Y}\mathcal{O}_{Y'}
  \longrightarrow
  \mathcal{O}_{Y_{f'}'^{(n)}}
  \tag{16.2.3.1}
  \label{IV.16.2.3.1}
\]
\oldpage[IV-4]{12}
deducido del homomorfismo $\nu_n$ (16.2.1.2) es biyectivo. Además,
el homomorfismo de $\mathcal{O}_{Y'}$-módulos
\[
  \operatorname{Gr}_1(u):
  \mathscr{G}r_1(f)\otimes_{\mathcal{O}_Y}\mathcal{O}_{Y'}
  \longrightarrow\mathscr{G}r_1(f')
  \tag{16.2.3.2}
  \label{IV.16.2.3.2}
\]
es biyectivo.
\end{enumerate}
\end{corollary}

\noindent(i) Observemos primero que los morfismos $f':Y'\to X'$ y $u:Y'\to Y$
identifican $Y'$ con el producto fibrado $Y\times_X X'$ (para los morfismos estructurales
$f:Y\to X$ y $v:X'\to X$) (14.5.12.1). La conclusión de (i) resulta entonces de
(16.2.2, (ii)), pues el morfismo $g$ es una inmersión.

\noindent(ii) El diagrama conmutativo
\[
  \xymatrix{
    Y_f^{(n)} \ar[d]^{h_n}
      & Y_{f'}'^{(n)} \ar[d]^-{h'_n} \ar[l]^{w_n} \\
    X \ar[d]^g
      & X' \ar[d]^{g'} \ar[l]^v \\
    Y
      & Y' \ar[l]^u
  }
\]
identifica $Y_{f'}'^{(n)}$ con el producto fibrado $Y_f^{(n)}\times_X X'$ y $X'$ con el producto fibrado
$X\times_Y Y'$; por tanto, (I, 3.3.9) identifica (para los morfismos
$g'\circ h'_n$ y $w_n$) $Y_{f'}'^{(n)}$ con el producto fibrado
$Y_f^{(n)}\times_Y Y'$. Como $Y_f^{(n)}$ (resp. $Y_{f'}'^{(n)}$) es el
preesquema afín sobre $Y$ (resp. $Y'$) asociado al
$\mathcal{O}_Y$-álgebra $\mathcal{O}_{Y_f^{(n)}}$ (resp. al
$\mathcal{O}_{Y'}$-álgebra $\mathcal{O}_{Y_{f'}'^{(n)}}$), la biyectividad del
homomorfismo canónico (16.2.3.1) resulta de (II, 1.5.2).
Finalmente, el homomorfismo canónico (16.2.3.1) es compatible con las
aumentaciones
$\mathcal{O}_{Y_f^{(n)}}\to\mathcal{O}_Y$ y
$\mathcal{O}_{Y_{f'}'^{(n)}}\to\mathcal{O}_{Y'}$; como
$\mathcal{O}_{Y_f^{(n)}}$ es suma directa (en tanto que
$\mathcal{O}_Y$-módulo) de $\mathcal{O}_Y$ y del ideal de aumentación
$\mathscr{I}/\mathscr{I}^{n+1}$, se ve por tanto que el homomorfismo canónico
(16.2.3.1), restringido a
$(\mathscr{I}/\mathscr{I}^{n+1})\otimes_{\mathcal{O}_Y}\mathcal{O}_{Y'}$,
es una biyección de este último sobre
$\mathscr{I}'/\mathscr{I}'^{n+1}$. Para $n=1$, esto muestra que
$\operatorname{Gr}_1(u)$ es biyectivo.

Obsérvese que, bajo las hipótesis de (16.2.3), los homomorfismos
$\operatorname{Gr}_n(u)$ son sobreyectivos en virtud de lo que precede, pero no
son biyectivos en general para $n\geq 2$. Sin embargo:

\begin{corollary}[16.2.4]
\label{IV.16.2.4}
Bajo las hipótesis de (16.2.3), supongamos que $u:Y'\to Y$ es un morfismo
plano (resp. que los $\mathscr{G}r_n(f)$ son $\mathcal{O}_Y$-módulos
planos para $n\leq m$). Entonces el homomorfismo
\[
  \operatorname{Gr}_n(u):
  \mathscr{G}r_n(f)\otimes_{\mathcal{O}_Y}\mathcal{O}_{Y'}
  \longrightarrow\mathscr{G}r_n(f')
\]
es biyectivo para todo $n$ (resp. para $n\leq m$).
\end{corollary}

Si $u$ es plano, lo mismo ocurre con $v:X'\to X$, que se obtiene de él por
cambio de base, y ya se sabe en este caso que $\operatorname{Gr}(u)$ es
biyectivo (16.2.2, (iii)). Si los $\mathscr{G}r_n(f)$ son planos para
$n\leq m$, se ve primero por inducción sobre $n$ que lo mismo ocurre con los
$\mathscr{I}/\mathscr{I}^{n+1}$ para $n\leq m$, en virtud de las sucesiones exactas
\[
  0\longrightarrow\mathscr{I}^n/\mathscr{I}^{n+1}
  \longrightarrow\mathscr{I}/\mathscr{I}^{n+1}
  \longrightarrow\mathscr{I}/\mathscr{I}^n
  \longrightarrow 0
\]
\oldpage[IV-4]{13}
(0\textsubscript{I}, 6.1.2); además, se tienen entonces diagramas conmutativos
\[
  \xymatrix{
    0 \ar[r] &
    (\mathscr{I}^n/\mathscr{I}^{n+1})
      \otimes_{\mathcal{O}_Y}\mathcal{O}_{Y'}
      \ar[r]\ar[d] &
    (\mathscr{I}/\mathscr{I}^{n+1})
      \otimes_{\mathcal{O}_Y}\mathcal{O}_{Y'}
      \ar[r]\ar[d] &
    (\mathscr{I}/\mathscr{I}^n)
      \otimes_{\mathcal{O}_Y}\mathcal{O}_{Y'}
      \ar[r]\ar[d] & 0 \\
    0 \ar[r] &
    \mathscr{I}'^n/\mathscr{I}'^{n+1}
      \ar[r] &
    \mathscr{I}'/\mathscr{I}'^{n+1}
      \ar[r] &
    \mathscr{I}'/\mathscr{I}'^n
      \ar[r] & 0
  }
\]
en los cuales las filas son exactas (la primera por planitud
(0\textsubscript{I}, 6.1.2)) y las dos últimas flechas verticales son
biyectivas en virtud de (16.2.3, (ii)); de donde se sigue la conclusión.

\begin{remarks}[16.2.5]
\label{IV.16.2.5}
\begin{enumerate}
\item[(i)] El razonamiento de (16.2.2, (i)) sigue siendo aplicable cuando en
(16.2.1.1) se trata de morfismos de espacios localmente anillados
($\mathbf{Err}_{\mathrm{II}}$, (1.8.2)).
\item[(ii)] En (16.2.2, (ii)), la conclusión ya no es necesariamente válida
cuando se supone solamente que $v$ y $f$ son morfismos de preesquemas
($f$ satisface la condición de (16.1.1)). Por ejemplo (con las notaciones de la
demostración de (16.2.2, (ii))), puede ocurrir que $\mathfrak{I}=0$ pero que
el núcleo $\mathfrak{I}'$ de $A'\to B'=B\otimes_A A'$ no sea nulo y que
$B'\neq 0$, en cuyo caso se tiene $Y^{(n)}=Y$ para todo $n$, pero
$Y'^{(n)}\neq Y'$. Se obtiene un ejemplo de este hecho tomando
$A=\mathbb{Z}$, $B=\mathbb{Q}$,
$A'=\prod_{h=1}^{\infty}(\mathbb{Z}/m^h\mathbb{Z})$ donde $m>1$.
\end{enumerate}
\end{remarks}

\begin{env}[16.2.6]
\label{IV.16.2.6}
Consideremos el caso particular del diagrama (16.2.1.1) en que $X'=X$, $v$ es
la identidad, $X$ es un preesquema, $Y$ un subpreesquema de $X$, $Y'$ un
subpreesquema de $Y$, y $f$, $u$ y $f'=f\circ u$ las inyecciones canónicas; el
dihomomorfismo (16.2.1.3) da, por tensorización con
$\mathcal{O}_{Y'}$ sobre $\rho^*(\mathcal{O}_Y)$, un homomorfismo de
$\mathcal{O}_{Y'}$-álgebras graduadas
\[
  u^*(\mathscr{G}r_\bullet(f))
  \longrightarrow\mathscr{G}r_\bullet(f')
  \tag{16.2.6.1}
  \label{IV.16.2.6.1}
\]
Por otra parte, se identifica $\mathcal{O}_Y$ con
$\psi^*(\mathcal{O}_X)/\mathscr{I}_f$, y $\mathcal{O}_{Y'}$ con
$\rho^*(\mathcal{O}_Y)/\mathscr{I}_u$; como $\rho^*$ es un funtor exacto,
se tiene
$\rho^*(\mathcal{O}_Y)=
\rho^*(\psi^*(\mathcal{O}_X))/\rho^*(\mathscr{I}_f)=
\psi'^*(\mathcal{O}_X)/\rho^*(\mathscr{I}_f)$, y como
$\mathcal{O}_{Y'}$ se identifica además con
$\psi'^*(\mathcal{O}_X)/\mathscr{I}_{f'}$, se ve que se tiene
$\mathscr{I}_u=\mathscr{I}_{f'}/\rho^*(\mathscr{I}_f)$. Se deduce para
todo entero $n$ un homomorfismo canónico
$\mathscr{I}_{f'}^n/\mathscr{I}_{f'}^{n+1}\to
\mathscr{I}_u^n/\mathscr{I}_u^{n+1}$, de donde se obtiene un homomorfismo canónico de
$\mathcal{O}_{Y'}$-álgebras graduadas
\[
  \mathscr{G}r_\bullet(f')
  \longrightarrow\mathscr{G}r_\bullet(u).
  \tag{16.2.6.2}
  \label{IV.16.2.6.2}
\]
\end{env}

\begin{proposition}[16.2.7]
\label{IV.16.2.7}
Sean $X$ un preesquema, $Y$ un subpreesquema de $X$, $Y'$ un
subpreesquema de $Y$, $j:Y'\to Y$ la inyección canónica. Se tiene entonces una sucesión
exacta de haces conormales ($\mathcal{O}_{Y'}$-módulos)
\[
  \xymatrix{
    j^*(\mathscr{N}_{Y/X}) \ar[r] &
    \mathscr{N}_{Y'/X} \ar[r] &
    \mathscr{N}_{Y'/Y} \ar[r] & 0
  }
  \tag{16.2.7.1}
  \label{IV.16.2.7.1}
\]
donde las flechas son las componentes de grado 1 de los homomorfismos canónicos
(16.2.6.1) y (16.2.6.2).
\end{proposition}

Como la cuestión es local, basta considerar el caso en que
$X=\operatorname{Spec}(A)$, $Y=\operatorname{Spec}(A/\mathfrak{I})$,
$Y'=\operatorname{Spec}(A/\mathfrak{K})$, $\mathfrak{I}$ y $\mathfrak{K}$
son dos ideales de $A$ tales que $\mathfrak{I}\subset\mathfrak{K}$; todo
se reduce entonces a ver
\oldpage[IV-4]{14}
que la sucesión de homomorfismos canónicos
$\mathfrak{I}/\mathfrak{K}\mathfrak{I}\to
\mathfrak{K}/\mathfrak{K}^2\to
(\mathfrak{K}/\mathfrak{I})/(\mathfrak{K}/\mathfrak{I})^2\to 0$
es exacta, lo que es inmediato, puesto que la imagen de
$\mathfrak{I}/\mathfrak{K}\mathfrak{I}$ en
$\mathfrak{K}/\mathfrak{K}^2$ es
$(\mathfrak{I}+\mathfrak{K}^2)/\mathfrak{K}^2$ y que
$(\mathfrak{K}/\mathfrak{I})/(\mathfrak{K}/\mathfrak{I})^2$ se identifica con
$\mathfrak{K}/(\mathfrak{I}+\mathfrak{K}^2)$.

Es fácil dar ejemplos en los que la sucesión (16.2.7.1), prolongada a la izquierda
por un 0, deja de ser exacta; con las notaciones precedentes, basta
tomar $A=k[T]$, $\mathfrak{I}=AT^2$, $\mathfrak{K}=AT$, pues se tiene entonces
$(\mathfrak{I}+\mathfrak{K}^2)/\mathfrak{K}^2=0$ y
$\mathfrak{I}/\mathfrak{K}\mathfrak{I}\neq 0$. Véanse, sin embargo, (16.9.13) y
(19.1.5) para casos útiles en los que la sucesión prolongada sigue siendo exacta.

\subsection{Invariantes diferenciales fundamentales de un morfismo de preesquemas}
\label{IV.16.3}

\begin{definition}[16.3.1]
\label{IV.16.3.1}
Sean $f:X\to S$ un morfismo de preesquemas, $\Delta_f:X\to X\times_S X$
el morfismo diagonal correspondiente, que es una inmersión (I, 5.3.9 y
$\mathbf{Err}_{\mathrm{III}}$, 10). Se denota por $\mathscr{P}_f^n$ o
$\mathscr{P}_{X/S}^n$, y se llama haz de partes principales de orden $n$
del $S$-preesquema $X$, al haz de anillos aumentado sobre $\mathcal{O}_X$, $n$-ésimo
invariante normal de $\Delta_f$ (16.1.2). Se pone
$\mathscr{P}_f^\infty=\mathscr{P}_{X/S}^\infty=
\varprojlim_n\mathscr{P}_{X/S}^n$,
$\mathscr{G}r_n(\mathscr{P}_f)=\mathscr{G}r_n(\mathscr{P}_{X/S})=
\mathscr{G}r_n(\Delta_f)$ (16.1.2); el $\mathcal{O}_X$-módulo
$\mathscr{G}r_1(\Delta_f)$, ideal de aumentación de $\mathscr{P}_{X/S}^1$, se
denota también por $\Omega_f^1$ o $\Omega_{X/S}^1$, y se llama el $\mathcal{O}_X$-módulo
de diferenciales de primer orden de $f$, o de $X$ respecto de $S$, o del $S$-preesquema $X$.
\end{definition}

De esta definición resulta que $\mathscr{P}_{X/S}^0$ se identifica
canónicamente con $\mathcal{O}_X$ (16.1.2).

Se tiene (16.1.2.2) un homomorfismo canónico sobreyectivo de
$\mathcal{O}_X$-álgebras graduadas
\[
  \mathscr{S}_{\mathcal{O}_X}^{\bullet}(\Omega_{X/S}^1)
  \longrightarrow\mathscr{G}r_\bullet(\mathscr{P}_{X/S}).
  \tag{16.3.1.1}
  \label{IV.16.3.1.1}
\]
Resulta asimismo de la definición (16.3.1) que, para todo abierto $U$ de $X$, se tiene
$\mathscr{P}_{f|U}^n=\mathscr{P}_f^n|U$,
$\mathscr{P}_{f|U}^\infty=\mathscr{P}_f^\infty|U$,
$\mathscr{G}r_n(\mathscr{P}_{f|U})=\mathscr{G}r_n(\mathscr{P}_f)|U$,
$\Omega_{f|U}^1=\Omega_f^1|U$ (es decir, las nociones introducidas son
locales sobre $X$).

\begin{env}[16.3.2]
\label{IV.16.3.2}
Denotemos por $p_1$, $p_2$ las dos proyecciones canónicas del producto fibrado $X\times_S X$;
como $\Delta_f$ es una sección sobre $X$ de $X\times_S X$ a la vez para los
morfismos $p_1$ y $p_2$, cada uno de estos morfismos define, para todo $n$, un
homomorfismo de haces de anillos $\mathcal{O}_X\to\mathscr{P}_{X/S}^n$,
inverso por la derecha de la aumentación $\mathscr{P}_{X/S}^n\to\mathcal{O}_X$
(16.1.7); también puede decirse que se definen así sobre
$\mathscr{P}_{X/S}^n$ dos estructuras cuasicoherentes de álgebra aumentada sobre
$\mathcal{O}_X$; las estructuras correspondientes de $\mathcal{O}_X$-módulo sobre
$\mathscr{G}r_n(\mathscr{P}_{X/S})$ son las mismas. Se tienen asimismo, al pasar al
límite, dos estructuras de $\mathcal{O}_X$-álgebra sobre
$\mathscr{P}_{X/S}^\infty$.
\end{env}

\begin{env}[16.3.3]
\label{IV.16.3.3}
El morfismo $s=(p_2,p_1)_S:X\times_S X\to X\times_S X$ es un automorfismo
involutivo de $X\times_S X$, llamado la simetría canónica, tal que
\[
  p_1\circ s=p_2,\qquad p_2\circ s=p_1,\qquad s\circ\Delta_f=\Delta_f.
  \tag{16.3.3.1}
  \label{IV.16.3.3.1}
\]
Si se pone $s=(\rho,\lambda)$, $p_i=(\pi_i,\mu_i)$ ($i=1,2$),
$\Delta_f=(\delta,\nu)$, $\lambda^\#$ es por tanto un isomorfismo de
$\rho^*(\pi_1^*(\mathcal{O}_X))$ en $\pi_2^*(\mathcal{O}_X)$, y
$\delta^*(\lambda^\#)$ deja invariantes $\delta^*(\mathcal{O}_{X\times_S X})$
y el núcleo $\mathscr{I}$ del homomorfismo
$\nu^\#:\delta^*(\mathcal{O}_{X\times_S X})\to\mathcal{O}_X$. Por consiguiente:
\end{env}

\begin{proposition}[16.3.4]
\label{IV.16.3.4}
El homomorfismo $\sigma=\delta^*(\lambda^\#)$ deducido de $s$ (y llamado también
simetría canónica) es un automorfismo involutivo del sistema proyectivo
$(\mathscr{P}_{X/S}^n)$ de haces
\oldpage[IV-4]{15}
de anillos aumentados sobre $\mathcal{O}_X$, y por consiguiente también de su límite proyectivo
$\mathscr{P}_{X/S}^\infty$. Este automorfismo intercambia las dos estructuras de
$\mathcal{O}_X$-álgebra sobre los $\mathscr{P}_{X/S}^n$ y sobre
$\mathscr{P}_{X/S}^\infty$.
\end{proposition}

\begin{env}[16.3.5]
\label{IV.16.3.5}
En lo que sigue, las dos estructuras de $\mathcal{O}_X$-álgebra definidas sobre los
$\mathscr{P}_{X/S}^n$ y sobre $\mathscr{P}_{X/S}^\infty$ desempeñarán papeles muy
distintos: convendremos en adelante, salvo mención expresa de lo contrario,
que cuando $\mathscr{P}_{X/S}^n$ o $\mathscr{P}_{X/S}^\infty$ se considere
como una $\mathcal{O}_X$-álgebra, se tratará de la estructura de álgebra definida por
$p_1$.
\end{env}

Para todo abierto $U$ de $X$ y toda sección $t\in\Gamma(U,\mathcal{O}_X)$, se
denotará simplemente por $t.1$ o incluso por $t$ la imagen de $t$ mediante el homomorfismo
estructural $\Gamma(U,\mathcal{O}_X)\to\Gamma(U,\mathscr{P}_{X/S}^n)$ (resp.
$\Gamma(U,\mathcal{O}_X)\to\Gamma(U,\mathscr{P}_{X/S}^\infty)$)
(es decir, el homomorfismo correspondiente a $p_1$).

\begin{definition}[16.3.6]
\label{IV.16.3.6}
Se denota por $d_f^n$, o $d_{X/S}^n$ (resp. $d_f^\infty$, o
$d_{X/S}^\infty$), o simplemente $d^n$ (resp. $d^\infty$), el homomorfismo de
haces de anillos $\mathcal{O}_X\to\mathscr{P}_f^n=\mathscr{P}_{X/S}^n$
(resp. $\mathcal{O}_X\to\mathscr{P}_f^\infty=\mathscr{P}_{X/S}^\infty$)
deducido de $p_2$. Para todo abierto $U$ de $X$, y todo
$t\in\Gamma(U,\mathcal{O}_X)$, $d^nt$ (resp. $d^\infty t$) se llama la parte
principal de orden $n$ (resp. parte principal de orden infinito) de $t$. Se pone
$dt=d^1t-t$, y se dice que $dt$ es el diferencial de $t$ (elemento de
$\Gamma(U,\Omega_{X/S}^1)$, denotado también por $d_{X/S}(t)$).
\end{definition}

De esta definición se sigue inmediatamente que
\[
  d(t_1t_2)=t_1dt_2+t_2dt_1
  \tag{16.3.6.1}
  \label{IV.16.3.6.1}
\]
para cualesquiera $t_1,t_2$ en $\Gamma(U,\mathcal{O}_X)$; es decir, $d$
es una derivación del anillo $\Gamma(U,\mathcal{O}_X)$ con valores en el
$\Gamma(U,\mathcal{O}_X)$-módulo $\Gamma(U,\Omega_{X/S}^1)$.

En todas las notaciones introducidas en (16.3.1) y (16.3.6), se reemplazará
a veces $S$ por $A$ cuando $S=\operatorname{Spec}(A)$.

\begin{env}[16.3.7]
\label{IV.16.3.7}
Supongamos, en particular, que $S=\operatorname{Spec}(A)$ y
$X=\operatorname{Spec}(B)$ son esquemas afines, siendo entonces $B$ una
$A$-álgebra. Entonces $\Delta_f$ corresponde al homomorfismo canónico sobreyectivo
$\pi:B\otimes_A B\to B$ tal que $\pi(b\otimes b')=bb'$, cuyo núcleo es
$\mathfrak{I}=\mathfrak{I}_{B/A}$ (0, 20.4.1); $\mathscr{P}_f^n$ es el haz
estructural del preesquema $\operatorname{Spec}(P_{B/A}^n)$, donde
\[
  P_{B/A}^n=(B\otimes_A B)/\mathfrak{I}^{n+1};
\]
$\mathscr{G}r_\bullet(\mathscr{P}_f)$ es el $\mathcal{O}_X$-módulo
cuasicoherente correspondiente al $B$-módulo graduado
\[
  \operatorname{gr}_{\mathfrak{I}}^\bullet(B\otimes_A B)
  =\bigoplus_{n\geq 0}(\mathfrak{I}^n/\mathfrak{I}^{n+1});
\]
en particular, $\Omega_f^1=\Omega_{X/S}^1$ es el $\mathcal{O}_X$-módulo
cuasicoherente correspondiente al $B$-módulo de diferenciales de primer orden de $B$ con
respecto a $A$, $\Omega_{B/A}^1$ (0, 20.4.3). Los morfismos de proyección
$p_1:X\times_S X\to X$, $p_2:X\times_S X\to X$ corresponden a los dos
homomorfismos de anillos $j_1:B\to B\otimes_A B$, $j_2:B\to B\otimes_A B$
tales que $j_1(b)=b\otimes 1$, $j_2(b)=1\otimes b$, de modo que (por la
convención de (16.3.5)), $P_{B/A}^n$ se considera siempre como una
$B$-álgebra mediante el homomorfismo compuesto
$B\xrightarrow{j_1}B\otimes_A B\to P_{B/A}^n$; el homomorfismo de anillos
$B\xrightarrow{j_2}B\otimes_A B\to P_{B/A}^n$ se denota por $d_{B/A}^n$ y corresponde
a $d_{X/S}^n$, que actúa sobre $\Gamma(X,\mathcal{O}_X)$; para todo $t\in B$, $dt$
es igual a $d_{B/A}t$, definido en (0, 20.4.6) (cf. $\mathbf{Err}_{\mathrm{IV}}$,
11).

Si $\pi_n:B\otimes_A B\to P_{B/A}^n$ es el homomorfismo canónico, se tiene, por tanto,
en virtud de las definiciones precedentes,
\[
  \pi_n(b\otimes b')=b.\pi_n(1\otimes b')=b.d_{B/A}^n(b')
  \quad\text{para }b\in B,\ b'\in B.
  \tag{16.3.7.1}
  \label{IV.16.3.7.1}
\]
\end{env}

\oldpage[IV-4]{16}
\begin{proposition}[16.3.8]
\label{IV.16.3.8}
La imagen del homomorfismo $d_{X/S}^n:\mathcal{O}_X\to\mathscr{P}_{X/S}^n$
genera el $\mathcal{O}_X$-módulo $\mathscr{P}_{X/S}^n$.
\end{proposition}

Se reduce inmediatamente al caso en que $X=\operatorname{Spec}(B)$ y
$S=\operatorname{Spec}(A)$ son afines, y la proposición resulta de (16.3.7.1),
puesto que $\pi_n$ es sobreyectivo. Obsérvese que, en general, $d_{X/S}^n$ no es
sobreyectivo (incluso para $n=1$).

\begin{proposition}[16.3.9]
\label{IV.16.3.9}
Supongamos que $f:X\to S$ es un morfismo localmente de tipo finito. Entonces los
$\mathscr{P}_f^n$ y los $\mathscr{G}r_n(\mathscr{P}_f)$ son
$\mathcal{O}_X$-módulos cuasicoherentes de tipo finito.
\end{proposition}

Esto resulta de (16.1.6) y del hecho de que $\Delta_f$ es localmente de presentación
finita (I, 4.3.1).

\subsection{Propiedades funtoriales de los invariantes diferenciales}
\label{IV.16.4}

\begin{env}[16.4.1]
\label{IV.16.4.1}
Consideremos un diagrama conmutativo de morfismos de preesquemas
\[
  \xymatrix{
    X \ar[d]_-{f} & X' \ar[l]_-u \ar[d]^-{f'} \\
    S & S' \ar[l]^-w
  }
  \tag{16.4.1.1}
  \label{IV.16.4.1.1}
\]
De él se deduce un diagrama conmutativo
\[
  \xymatrix{
    X \ar[d]_-{\Delta_f} & X' \ar[l]_-u \ar[d]^-{\Delta_{f'}} \\
    X\times_S X & X'\times_{S'}X' \ar[l]^-v
  }
\]
donde $v$ es el morfismo compuesto (I, 5.3.5 y 5.3.15)
\[
  X'\times_{S'}X'
  \xrightarrow{(p'_1,p'_2)_S}X'\times_S X'
  \xrightarrow{u\times_S u}X\times_S X.
  \tag{16.4.1.2}
  \label{IV.16.4.1.2}
\]
Se deducen, por tanto, de $u$ y $v$, como se explicó en (16.2.1),
homomorfismos de haces de anillos aumentados
\[
  \nu_n:\rho^*(\mathscr{P}_{X/S}^n)\longrightarrow
  \mathscr{P}_{X'/S'}^n
  \tag{16.4.1.3}
  \label{IV.16.4.1.3}
\]
(donde se ha puesto $u=(\rho,\lambda)$); estos homomorfismos forman un sistema
proyectivo y dan, por tanto, en el límite un homomorfismo de haces de anillos
aumentados
\[
  \nu_\infty:\rho^*(\mathscr{P}_{X/S}^\infty)\longrightarrow
  \mathscr{P}_{X'/S'}^\infty;
  \tag{16.4.1.4}
  \label{IV.16.4.1.4}
\]
por otra parte, al pasar a los cocientes, los homomorfismos $\nu_n$ dan un
dihomomorfismo de álgebras graduadas (relativo a $\lambda^\#$):
\[
  \operatorname{gr}(u):\rho^*(\mathscr{G}r_\bullet(\mathscr{P}_{X/S}))
  \longrightarrow\mathscr{G}r_\bullet(\mathscr{P}_{X'/S'}).
  \tag{16.4.1.5}
  \label{IV.16.4.1.5}
\]
\end{env}

\begin{env}[16.4.2]
\label{IV.16.4.2}
\oldpage[IV-4]{17}
Si se tiene un diagrama conmutativo
\[
  \xymatrix{
    X \ar[d]_-{f} & X' \ar[l]_-u \ar[d]^-{f'} &
      X'' \ar[l]_-{u'} \ar[d]^-{f''} \\
    S & S' \ar[l]^-w & S'' \ar[l]^-{w'}
  }
\]
se deduce un diagrama conmutativo
\[
  \xymatrix{
    X \ar[d]_-{\Delta_f} & X' \ar[l]_-u \ar[d]^-{\Delta_{f'}} &
      X'' \ar[l]_-{u'} \ar[d]^-{\Delta_{f''}} \\
    X\times_S X & X'\times_{S'}X' \ar[l]^-v &
      X''\times_{S''}X'' \ar[l]^-{v'}
  }
\]
donde $v'$ se define a partir de $u'$, $w'$, $f'$, $f''$ como $v$ a partir de
$u$, $w$, $f$, $f'$. Se verifica inmediatamente que, si $u''=u\circ u'$,
$w''=w\circ w'$, entonces el morfismo compuesto $v\circ v'$ es igual al morfismo
$v''$ definido a partir de $u''$, $w''$, $f$, $f''$ como $v$ a partir de $u$,
$w$, $f$, $f'$. Si se pone $u'=(\rho',\lambda')$,
$u''=(\rho'',\lambda'')$, resulta entonces de (16.2.1) que el homomorfismo
$\nu_n'':\rho''^*(\mathscr{P}_{X/S}^n)\to\mathscr{P}_{X''/S''}^n$ es igual al
compuesto
\[
  \rho'^*(\rho^*(\mathscr{P}_{X/S}^n))
  \xrightarrow{\rho'^*(\nu_n^\#)}
  \rho'^*(\mathscr{P}_{X'/S'}^n)
  \xrightarrow{\nu_n'}\mathscr{P}_{X''/S''}^n
\]
y se tienen propiedades de transitividad análogas para los homomorfismos
(16.4.1.4) y (16.4.1.5), lo que permite decir que los
$\mathscr{P}_{X/S}^n$, $\mathscr{P}_{X/S}^\infty$ y
$\mathscr{G}r_\bullet(\mathscr{P}_{X/S})$ dependen funtorialmente de $f$.
\end{env}

\begin{env}[16.4.3]
\label{IV.16.4.3}
Se verifica inmediatamente (por ejemplo, reduciéndose al caso afín con ayuda de
(16.3.7)) que, con las notaciones de (16.4.1), el diagrama
\[
  \xymatrix{
    \rho^*(\mathcal{O}_X) \ar[r]^{\lambda^\#} \ar[d] &
      \mathcal{O}_{X'} \ar[d] \\
    \rho^*(\mathscr{P}_{X/S}^n) \ar[r]_{\nu_n} &
      \mathscr{P}_{X'/S'}^n
  }
  \tag{16.4.3.1}
  \label{IV.16.4.3.1}
\]
donde las flechas verticales son las que definen las estructuras de álgebra
elegidas en (16.3.5) (es decir, las que provienen de las primeras proyecciones),
es conmutativo; lo mismo ocurre con el diagrama
\[
  \xymatrix{
    \rho^*(\mathcal{O}_X) \ar[r]^{\lambda^\#}
      \ar[d]_{\rho^*(d_{X/S}^n)} &
      \mathcal{O}_{X'} \ar[d]^{d_{X'/S'}^n} \\
    \rho^*(\mathscr{P}_{X/S}^n) \ar[r]_{\nu_n} &
      \mathscr{P}_{X'/S'}^n
  }
  \tag{16.4.3.2}
  \label{IV.16.4.3.2}
\]
\oldpage[IV-4]{18}
las flechas verticales definen aquí las estructuras de álgebra provenientes de las
segundas proyecciones; además, si $\sigma$ y $\sigma'$ son las simetrías
canónicas correspondientes a $f$ y $f'$ (16.3.4), se tiene
\[
  \nu_n\circ\rho^*(\sigma)=\sigma'\circ\nu_n
\]
que permite pasar de uno de los diagramas precedentes al otro. Se deduce, por tanto, de
(16.4.3.1) un homomorfismo canónico de $\mathcal{O}_{X'}$-álgebras aumentadas
\[
  P^n(u):u^*(\mathscr{P}_{X/S}^n)
  =\mathscr{P}_{X/S}^n\otimes_{\mathcal{O}_X}\mathcal{O}_{X'}
  \longrightarrow\mathscr{P}_{X'/S'}^n
  \tag{16.4.3.3}
  \label{IV.16.4.3.3}
\]
y resulta de (16.4.3.2) que el diagrama
\[
  \xymatrix{
    \mathcal{O}_{X'} \ar[r]^{\operatorname{id}}
      \ar[d]_{u^*(d_{X/S}^n)} &
      \mathcal{O}_{X'} \ar[d]^{d_{X'/S'}^n} \\
    u^*(\mathscr{P}_{X/S}^n) \ar[r]_{P^n(u)} &
      \mathscr{P}_{X'/S'}^n
  }
  \tag{16.4.3.4}
  \label{IV.16.4.3.4}
\]
es conmutativo. Se deduce un homomorfismo de $\mathcal{O}_{X'}$-álgebras
graduadas
\[
  \operatorname{Gr}_\bullet(u):u^*(\mathscr{G}r_\bullet(\mathscr{P}_{X/S}))
  \longrightarrow\mathscr{G}r_\bullet(\mathscr{P}_{X'/S'})
  \tag{16.4.3.5}
  \label{IV.16.4.3.5}
\]
y, en particular, un homomorfismo de $\mathcal{O}_{X'}$-módulos
\[
  \operatorname{Gr}_1(u):
  \Omega_{X/S}^1\otimes_{\mathcal{O}_X}\mathcal{O}_{X'}
  \longrightarrow\Omega_{X'/S'}^1
  \tag{16.4.3.6}
  \label{IV.16.4.3.6}
\]
que da lugar a un diagrama conmutativo
\[
  \xymatrix{
    \mathcal{O}_{X'} \ar[r]^{\operatorname{id}}
      \ar[d]_{d_{X/S}\otimes 1} &
      \mathcal{O}_{X'} \ar[d]^{d_{X'/S'}} \\
    \Omega_{X/S}^1\otimes_{\mathcal{O}_X}\mathcal{O}_{X'} \ar[r] &
      \Omega_{X'/S'}^1
  }
  \tag{16.4.3.7}
  \label{IV.16.4.3.7}
\]
\end{env}

\begin{env}[16.4.4]
\label{IV.16.4.4}
Cuando $S=\operatorname{Spec}(A)$, $S'=\operatorname{Spec}(A')$,
$X=\operatorname{Spec}(B)$, $X'=\operatorname{Spec}(B')$ son afines, de modo
que se tiene un diagrama conmutativo de homomorfismos de anillos
\[
  \xymatrix{
    B \ar[r] & B' \\
    A \ar[u] \ar[r] & A' \ar[u]
  }
\]
la imagen de $\mathfrak{I}_{B/A}$ en $B'\otimes_{A'}B'$ está contenida en
$\mathfrak{I}_{B'/A'}$, y el homomorfismo $\nu_n$ corresponde al
homomorfismo de anillos $P_{B/A}^n\to P_{B'/A'}^n$ deducido del homomorfismo
$B\otimes_A B\to B'\otimes_{A'}B'$ al pasar a los cocientes. El homomorfismo
(16.4.3.6) corresponde al homomorfismo definido en (0, 20.5.4.1), y el
diagrama conmutativo (16.4.3.7) al diagrama (0, 20.5.4.2).
\end{env}

\oldpage[IV-4]{19}
\begin{proposition}[16.4.5]
\label{IV.16.4.5}
Supongamos que $X'=X\times_S S'$, siendo $f'$ y $u$ las proyecciones canónicas.
Entonces los homomorfismos canónicos $P^n(u)$ (16.4.3.3) y
$\operatorname{Gr}_1(u)$ (16.4.3.6) son biyectivos.
\end{proposition}

En efecto, se tiene entonces $X'\times_{S'}X'=(X\times_S X)\times_S S'$, y se puede, por tanto,
aplicar (16.2.3, (ii)) reemplazando $g$ por la primera proyección
$p_1:X\times_S X\to X$ y $f$ por la diagonal $\Delta_f$.

Obsérvese también que, bajo las hipótesis de (16.4.5), el homomorfismo
$\operatorname{Gr}_\bullet(u)$ (16.4.3.5) es sobreyectivo, pero no biyectivo en
general. Sin embargo (16.2.4):

\begin{corollary}[16.4.6]
\label{IV.16.4.6}
Siendo las hipótesis las de (16.4.5), supongamos además que $w:S'\to S$ es
plano (resp. que los $\mathscr{G}r_n(\mathscr{P}_{X/S})$ son
$\mathcal{O}_X$-módulos planos para $n\leq m$); entonces el homomorfismo
\[
  \operatorname{Gr}_n(u):u^*(\mathscr{G}r_n(\mathscr{P}_{X/S}))
  \longrightarrow\mathscr{G}r_n(\mathscr{P}_{X'/S'})
\]
es biyectivo para todo $n$ (resp. para $n\leq m$).
\end{corollary}

En efecto, si $w$ es plano, lo mismo ocurre con
$v:X'\times_{S'}X'\to X\times_S X$, por lo que la conclusión resulta de (16.2.4).

\begin{env}[16.4.7]
\label{IV.16.4.7}
Sean $S$ un preesquema, $\mathscr{E}$ un $\mathcal{O}_S$-módulo
cuasicoherente, y pongamos $X=\mathbf{V}(\mathscr{E})$ (II, 1.7.8), fibrado
vectorial asociado a $\mathscr{E}$, igual a
$\operatorname{Spec}(\mathscr{S}_{\mathcal{O}_S}(\mathscr{E}))$. Sea
$f:X\to S$ el morfismo estructural. Para todo abierto $U$ de $S$, y toda
sección $t\in\Gamma(U,\mathscr{E})$, $t$ se identifica con una sección de
$\mathscr{S}_{\mathcal{O}_S}(\mathscr{E})$ sobre $U$; sea $t'$ su
imagen en
$\Gamma(f^{-1}(U),\mathcal{O}_X)=\Gamma(U,f_*(\mathcal{O}_X))=
\Gamma(U,\mathscr{S}_{\mathcal{O}_S}(\mathscr{E}))$, y pongamos
\[
  \delta(t)=d_{X/S}^n(t')-t'\in
  \Gamma(f^{-1}(U),\mathscr{P}_{X/S}^n);
  \tag{16.4.7.1}
  \label{IV.16.4.7.1}
\]
está claro que $\delta$ es un dihomomorfismo de módulos (correspondiente al
homomorfismo de anillos
$\Gamma(U,\mathcal{O}_S)\to\Gamma(f^{-1}(U),\mathcal{O}_X)$) de
$\Gamma(U,\mathscr{E})$ en $\Gamma(f^{-1}(U),\mathscr{P}_{X/S}^n)$, cuya
imagen pertenece además al ideal de aumentación de
$\Gamma(f^{-1}(U),\mathscr{P}_{X/S}^n)$. Se deduce (haciendo variar $U$)
un homomorfismo canónico de $\mathcal{O}_X$-álgebras
\[
  f^*(\mathscr{S}_{\mathcal{O}_S}(\mathscr{E}))
  \longrightarrow\mathscr{P}_{X/S}^n
  \tag{16.4.7.2}
  \label{IV.16.4.7.2}
\]
y, según la observación precedente, si $\mathscr{K}$ es el ideal que es núcleo de
la aumentación $\mathscr{S}_{\mathcal{O}_S}(\mathscr{E})\to\mathcal{O}_S$,
la imagen de $\mathscr{K}^{n+1}$ por (16.4.7.2) es nula, de modo que, al
factorizar por $\mathscr{K}^{n+1}$, se obtiene finalmente un homomorfismo canónico
\[
  \delta_n:f^*(\mathscr{S}_{\mathcal{O}_S}(\mathscr{E})/
  \mathscr{K}^{n+1})\longrightarrow\mathscr{P}_{X/S}^n.
  \tag{16.4.7.3}
  \label{IV.16.4.7.3}
\]
\end{env}

\begin{proposition}[16.4.8]
\label{IV.16.4.8}
Bajo las condiciones de (16.4.7), los homomorfismos $\delta_n$ son biyectivos y
forman un sistema proyectivo de isomorfismos; se deduce un isomorfismo de
$\mathcal{O}_X$-álgebras graduadas
\[
  f^*(\mathscr{S}_{\mathcal{O}_S}^{\bullet}(\mathscr{E}))
  \longrightarrow\mathscr{G}r_\bullet(\mathscr{P}_{X/S}).
  \tag{16.4.8.1}
  \label{IV.16.4.8.1}
\]
\end{proposition}

El hecho de que los homomorfismos (16.4.7.3) formen un sistema proyectivo resulta
inmediatamente de su definición. Para probar que son isomorfismos, bastará
\oldpage[IV-4]{20}
probar que (16.4.8.1) es un isomorfismo, pues las filtraciones de los dos miembros de
(16.4.7.3) son finitas (Bourbaki, \emph{Alg. comm.}, cap. III, \S~2,
n\textsuperscript{o} 8, cor. 3 del teo. 1). Para ello, consideremos la sucesión exacta
escindida de $\mathcal{O}_S$-módulos
\[
  \xymatrix{
    0 \ar[r] & \mathscr{E} \ar[r]^{u} &
    \mathscr{E}\oplus\mathscr{E} \ar[r]^{v} &
    \mathscr{E} \ar[r] & 0
  }
  \tag{16.4.8.2}
  \label{IV.16.4.8.2}
\]
donde, para todo par de secciones $s$, $t$ de $\mathscr{E}$ sobre un abierto
$U$ de $S$, se toma $u(s)=(-s,s)$ y $v(s,t)=s+t$. Se tiene
\[
  X\times_S X=\operatorname{Spec}
  (\mathscr{S}_{\mathcal{O}_S}(\mathscr{E})\otimes_{\mathcal{O}_S}
  \mathscr{S}_{\mathcal{O}_S}(\mathscr{E}))
  =\operatorname{Spec}(\mathscr{S}_{\mathcal{O}_S}
  (\mathscr{E}\oplus\mathscr{E}))
\]
(II, 1.4.6 y 1.7.11), y el morfismo diagonal $X\to X\times_S X$ corresponde
(II, 1.2.7) al homomorfismo de $\mathcal{O}_S$-álgebras
$\mathscr{S}(v):\mathscr{S}_{\mathcal{O}_S}(\mathscr{E}\oplus\mathscr{E})
\to\mathscr{S}_{\mathcal{O}_S}(\mathscr{E})$ (II, 1.7.4), de modo que, si
$\mathscr{I}$ es el núcleo de este homomorfismo, se tiene
\[
  \mathscr{P}_{X/S}^n=f^*(\mathscr{S}_{\mathcal{O}_S}
  (\mathscr{E}\oplus\mathscr{E})/\mathscr{I}^{n+1}).
\]

La proposición será consecuencia del lema siguiente:

\begin{lemma}[16.4.8.3]
\label{IV.16.4.8.3}
Sean $Y$ un espacio anillado,
$0\to\mathscr{F}'\xrightarrow{u}\mathscr{F}\xrightarrow{v}\mathscr{F}''\to0$
una sucesión exacta de $\mathcal{O}_Y$-módulos tal que todo punto $y\in Y$
tiene una vecindad abierta $V$ tal que la sucesión
$0\to\mathscr{F}'|V\to\mathscr{F}|V\to\mathscr{F}''|V\to0$ es escindida.
Sea $\mathscr{I}$ el ideal que es núcleo de
$\mathscr{S}(v):\mathscr{S}_{\mathcal{O}_Y}(\mathscr{F})\to
\mathscr{S}_{\mathcal{O}_Y}(\mathscr{F}'')$, y sea
$\operatorname{gr}_{\mathscr{I}}^\bullet
(\mathscr{S}_{\mathcal{O}_Y}(\mathscr{F}))$ la $\mathcal{O}_Y$-álgebra graduada
asociada a la $\mathcal{O}_Y$-álgebra
$\mathscr{S}_{\mathcal{O}_Y}(\mathscr{F})$ dotada de la filtración
$\mathscr{I}$-preádica. Entonces el homomorfismo de $\mathcal{O}_Y$-álgebras
graduadas
\[
  \mathscr{S}_{\mathcal{O}_Y}^{\bullet}(\mathscr{F}')
  \otimes_{\mathcal{O}_Y}
  \mathscr{S}_{\mathcal{O}_Y}^{\bullet}(\mathscr{F}'')
  \longrightarrow
  \operatorname{gr}_{\mathscr{I}}^\bullet
  (\mathscr{S}_{\mathcal{O}_Y}(\mathscr{F}))
  \tag{16.4.8.4}
  \label{IV.16.4.8.4}
\]
(donde el primer miembro es el producto tensorial graduado de las
$\mathcal{O}_Y$-álgebras simétricas dotadas de su graduación canónica
(II, 1.7.4 y 2.1.2)), proveniente de la inyección canónica
\[
  \mathscr{F}'\longrightarrow\mathscr{I}
  =\operatorname{gr}_{\mathscr{I}}^1
  (\mathscr{S}_{\mathcal{O}_Y}(\mathscr{F})),
\]
es biyectivo.
\end{lemma}

En efecto, la inyección $\mathscr{F}'\to\mathscr{I}$ da canónicamente un
homomorfismo de $\mathcal{O}_Y$-álgebras graduadas\linebreak
$\mathscr{S}_{\mathcal{O}_Y}^{\bullet}(\mathscr{F}')\to\allowbreak
\operatorname{gr}_{\mathscr{I}}^\bullet
(\mathscr{S}_{\mathcal{O}_Y}(\mathscr{F}))$, y como el segundo miembro es por
definición una $\mathscr{S}_{\mathcal{O}_Y}^{\bullet}(\mathscr{F}'')$-álgebra
graduada, se deduce el homomorfismo canónico (16.4.8.4) tensorizando el
precedente con $\mathscr{S}_{\mathcal{O}_Y}^{\bullet}(\mathscr{F}'')$. Para
probar el lema, como la cuestión es local, basta considerar el caso en que
$\mathscr{F}=\mathscr{F}'\oplus\mathscr{F}''$, siendo $u$ y $v$ los
homomorfismos canónicos. Entonces el álgebra graduada
$\mathscr{S}_{\mathcal{O}_Y}^{\bullet}(\mathscr{F})$ se identifica canónicamente con el
producto tensorial graduado
$\mathscr{S}_{\mathcal{O}_Y}^{\bullet}(\mathscr{F}')
\otimes_{\mathcal{O}_Y}
\mathscr{S}_{\mathcal{O}_Y}^{\bullet}(\mathscr{F}'')$ (II, 1.7.4), y es
inmediato que $\mathscr{I}$ es entonces el ideal
$\mathscr{I}'\otimes_{\mathcal{O}_Y}
\mathscr{S}_{\mathcal{O}_Y}^{\bullet}(\mathscr{F}'')$, donde $\mathscr{I}'$ es
el ideal de aumentación de
$\mathscr{S}_{\mathcal{O}_Y}^{\bullet}(\mathscr{F}')$, es decir, la suma
(directa) de los $\mathscr{S}_{\mathcal{O}_Y}^m(\mathscr{F}')$ para $m\geq1$. Se
concluye que $\mathscr{I}^n=\mathscr{I}'^n\otimes_{\mathcal{O}_Y}
\mathscr{S}_{\mathcal{O}_Y}^{\bullet}(\mathscr{F}'')$, donde esta vez
$\mathscr{I}'^n$ es la suma (directa) de los
$\mathscr{S}_{\mathcal{O}_Y}^m(\mathscr{F}')$ para $m\geq n$; se tiene por consiguiente
$\mathscr{I}^n/\mathscr{I}^{n+1}=
\mathscr{S}_{\mathcal{O}_Y}^n(\mathscr{F}')\otimes_{\mathcal{O}_Y}
\mathscr{S}_{\mathcal{O}_Y}^{\bullet}(\mathscr{F}'')$, lo que prueba que
(16.4.8.4) es biyectivo.

\oldpage[IV-4]{21}
Demostrado este lema, queda por ver que el homomorfismo (16.4.8.1) es
efectivamente la imagen por $f^*$ del homomorfismo (16.4.8.4) correspondiente a la sucesión
exacta (16.4.8.2); se comprueba fácilmente que esto resulta de la definición de
$u$ (16.4.8.2) y de la de $\delta$ (16.4.7.1), teniendo en cuenta la definición
de la estructura de $\mathcal{O}_X$-álgebra sobre $\mathscr{P}_{X/S}^n$ y de
la de $d_{X/S}^n$ (16.3.5 y 16.3.6).

En particular:

\begin{corollary}[16.4.9]
\label{IV.16.4.9}
Bajo las condiciones de (16.4.7), se tiene un isomorfismo canónico
\[
  \operatorname{gr}_1(\delta):f^*(\mathscr{E})
  \xrightarrow{\sim}\Omega_{X/S}^1.
  \tag{16.4.9.1}
  \label{IV.16.4.9.1}
\]
\end{corollary}

\begin{corollary}[16.4.10]
\label{IV.16.4.10}
Si $S=\operatorname{Spec}(A)$, $\mathscr{E}=\mathcal{O}_S^m$, de modo que
\[
  X=\operatorname{Spec}(A[T_1,\ldots,T_m]),
\]
$\mathscr{P}_{X/S}^n$ se identifica canónicamente con el
$\mathcal{O}_X$-álgebra correspondiente al $A[T_1,\ldots,T_m]$-álgebra cociente
\[
  A[T_1,\ldots,T_m,U_1,\ldots,U_m]/\mathfrak{K}^{n+1},
\]
donde las $U_i$ ($1\leq i\leq m$) son $m$ nuevas indeterminadas y
$\mathfrak{K}$ es el ideal engendrado por $U_1,\ldots,U_m$.
\end{corollary}

Se recupera así, en particular, la estructura de $\Omega_{X/S}^1$ en este caso
(0, 20.5.13).

Obsérvese además que $d_{X/S}^n$ hace corresponder entonces a un polinomio
$F(T_1,\ldots,T_m)$ la clase mód.\ $\mathfrak{K}^{n+1}$ de
$F(T_1+U_1,\ldots,T_m+U_m)$, como resulta de la definición (16.4.7.1).

\begin{proposition}[16.4.11]
\label{IV.16.4.11}
Sean $f:X\to S$ un morfismo, $g:S\to X$ una sección sobre $S$ de $X$,
$S_g^{(n)}$ la $n$-ésima vecindad infinitesimal de $S$ para la inmersión $g$
(16.1.2). Existe entonces un único isomorfismo de
$\mathcal{O}_S$-álgebras
\[
  \varpi_n:g^*(\mathscr{P}_{X/S}^n)\longrightarrow
  \mathcal{O}_{S_g^{(n)}}
  \tag{16.4.11.1}
  \label{IV.16.4.11.1}
\]
(para la estructura de $\mathcal{O}_S$-álgebra sobre
$\mathcal{O}_{S_g^{(n)}}$ definida por $f$ (16.1.7)), que hace conmutativo el
diagrama
\[
  \xymatrix{
    \mathcal{O}_S=g^*(\mathcal{O}_X)
      \ar[rr]^-{\lambda_n}\ar[dr]_-{g^*(d_{X/S}^n)}
      &&\mathcal{O}_{S_g^{(n)}}\\
    &g^*(\mathscr{P}_{X/S}^n)\ar[ur]_-{\varpi_n}
  }
  \tag{16.4.11.2}
  \label{IV.16.4.11.2}
\]
(donde $\lambda_n$ es el homomorfismo estructural).
\end{proposition}

En virtud de (I, 5.3.7), donde se reemplazan $X$, $Y$, $S$ por $S$, $X$, $S$,
respectivamente, y $f$ por $g$, los diagramas
\[
  \xymatrix@=3pc{
    S\ar[r]^-g\ar[d]_-g&X\ar[d]^-{\Delta_f}
      &S\ar[r]^-g\ar[d]_-g&X\ar[d]^-{\Delta_f}\\
    X\ar[r]_{(g\circ f,1_X)_S}&X\times_S X
      &X\ar[r]_{(1_X,g\circ f)_S}&X\times_S X
  }
  \tag{16.4.11.3}
  \label{IV.16.4.11.3}
\]
\oldpage[IV-4]{22}
identifican $S$ con el producto fibrado de los $(X\times_S X)$-preesquemas $X$ y $X$ para los
morfismos $\Delta_f$ y $(g\circ f,1_X)_S$ (resp.\ $(1_X,g\circ f)_S$).
Por otra parte, los diagramas
\[
  \xymatrix@=3pc{
    X\ar[r]^{(g\circ f,1_X)_S}\ar[d]_-f&X\times_S X\ar[d]^-{p_1}
      &X\ar[r]^{(1_X,g\circ f)_S}\ar[d]_-f&X\times_S X\ar[d]^-{p_2}\\
    S\ar[r]_{g}&X&S\ar[r]_{g}&X
  }
  \tag{16.4.11.4}
  \label{IV.16.4.11.4}
\]
identifican $X$ con el producto fibrado de los $X$-preesquemas $S$ y $X\times_S X$ para los
morfismos $g$ y $p_1$ (resp.\ $p_2$) (caso particular de la fórmula
de asociatividad (I, 3.3.9.1)). Puede decirse que $\Delta_f$, considerado como
sección sobre $X$ de $X\times_S X$ (relativamente a $p_1$ o $p_2$), desempeña el papel
de una \emph{sección universal} para las secciones sobre $S$ de $X$: cada una de estas
secciones $g$ se deduce, en efecto, de ella por el \emph{cambio de base}
$(g\circ f,1_X)_S:X\to X\times_S X$. La definición del homomorfismo
$\varpi_n$ y el hecho de que sea biyectivo resultan, por tanto, de estas observaciones y de
(16.2.3, (ii)) aplicado al primer diagrama (16.4.11.4). La conmutatividad del
diagrama (16.4.11.2) resulta asimismo de (16.2.3, (ii)) aplicado esta vez al
segundo diagrama (16.4.11.4). Para explicitar $\varpi_n$, basta considerar el
caso en que $g$ es una inmersión \emph{cerrada}: en efecto, para todo $s\in S$, hay
una vecindad abierta $W$ de $s$ en $S$ tal que $g(W)$ es cerrado en un
abierto $U$ de $X$, y está claro que $g|W$ es una sección sobre $W$ del morfismo
$U\cap f^{-1}(W)\to W$, restricción de $f$, y $g(W)$ es \emph{a fortiori}
cerrado en $U\cap f^{-1}(W)$. Se puede, por tanto, suponer que $S$ es un
subpreesquema cerrado de $X$ definido por un ideal cuasicoherente $\mathscr{K}$.
Entonces las definiciones precedentes muestran que, si $W$ es un abierto de $S$,
$t$ una sección de $\mathcal{O}_X$ sobre $f^{-1}(W)$,
$\varpi_n(d^n t|W)$ es igual a la imagen canónica de $t$ en
$\Gamma(W,(\mathcal{O}_X/\mathscr{K}^{n+1})|W)$. La unicidad de $\varpi_n$
resulta, por tanto, de que la imagen de $\mathcal{O}_X$ por $d_{X/S}^n$ genera el
$\mathcal{O}_X$-módulo $\mathscr{P}_{X/S}^n$ (16.3.8).

\begin{corollary}[16.4.12]
\label{IV.16.4.12}
Sean $k$ un cuerpo, $X$ un $k$-preesquema, $x$ un punto de $X$ racional sobre
$k$. Entonces $(\mathscr{P}_{X/k}^n)_x\otimes_{\mathcal{O}_x}k(x)$ es
canónicamente isomorfo (como $k(x)$-álgebra aumentada) a
$\mathcal{O}_x/\mathfrak{m}_x^{n+1}$.
\end{corollary}

Basta considerar la única sección sobre $k$, $g$, de $X$ tal que
$g(\operatorname{Spec}(k))=\{x\}$.

\begin{corollary}[16.4.13]
\label{IV.16.4.13}
Sean $f:X\to S$ un morfismo, $s$ un punto de $S$,
$X_s=X\times_S\operatorname{Spec}(k(s))$ la fibra de $f$ en $s$. Si
$x\in X_s$ es racional sobre $k(s)$,
$(\mathscr{P}_{X/S}^n)_x\otimes_{\mathcal{O}_s}k(s)$ es canónicamente
isomorfo a $\mathcal{O}_{X_s,x}/\mathfrak{m}'_x{}^{\,n+1}$, donde
$\mathfrak{m}'_x$ es el ideal maximal de $\mathcal{O}_{X_s,x}$; más
precisamente, este isomorfismo hace corresponder a $(d^n t)_x\otimes1$ (donde $t$ es
una sección de $\mathcal{O}_X$ sobre una vecindad abierta de $x$ en
$X$) la clase de $t_x\otimes1$ módulo $\mathfrak{m}'_x{}^{\,n+1}$.
\end{corollary}

Esto resulta de (16.4.5) y (16.4.12).

Los corolarios precedentes justifican la terminología de «haz de partes
principales de orden $n$».

\begin{proposition}[16.4.14]
\label{IV.16.4.14}
Sea $\rho:A\to B$ un homomorfismo de anillos, $S$ un subconjunto
multiplicativamente estable de $B$. Entonces los homomorfismos canónicos
\[
  S^{-1}P_{B/A}^n\longrightarrow P_{S^{-1}B/A}^n
  \tag{16.4.14.1}
  \label{IV.16.4.14.1}
\]
\oldpage[IV-4]{23}
deducidos de los homomorfismos canónicos
$P_{B/A}^n\to P_{S^{-1}B/A}^n$ (16.4.4), forman un sistema proyectivo y son
biyectivos.
\end{proposition}

Basta observar que
\[
  S^{-1}((B\otimes_A B)/\mathfrak{I}^{n+1})
  =S^{-1}(B\otimes_A B)/(S^{-1}\mathfrak{I})^{n+1}
\]
por planitud, y que
$S^{-1}(B\otimes_A B)=(S^{-1}B)\otimes_A(S^{-1}B)$ (I, 1.3.4).

\begin{corollary}[16.4.15]
\label{IV.16.4.15}
Con las notaciones de (16.4.14), sea $R$ un subconjunto multiplicativo de
$A$ tal que $\rho(R)\subset S$. Entonces se tienen isomorfismos canónicos
\[
  S^{-1}P_{B/A}^n\xrightarrow{\sim}P_{S^{-1}B/R^{-1}A}^n
  \tag{16.4.15.1}
  \label{IV.16.4.15.1}
\]
que forman un sistema proyectivo.
\end{corollary}

Basta, evidentemente, definir isomorfismos canónicos
\[
  P_{S^{-1}B/A}^n\xrightarrow{\sim}P_{S^{-1}B/R^{-1}A}^n,
  \tag{16.4.15.2}
  \label{IV.16.4.15.2}
\]
es decir, se reduce al caso en que $\rho(R)$ está formado por elementos
\emph{invertibles} de $B$. Pero entonces el isomorfismo (16.4.15.2) se
deduce simplemente del isomorfismo canónico
$B\otimes_A B\to B\otimes_{R^{-1}A}B$ al pasar a los cocientes
(0$_{\rm I}$, 1.5.3).

\begin{corollary}[16.4.16]
\label{IV.16.4.16}
Sean $f:X\to S$ un morfismo de preesquemas, $x$ un punto de $X$,
$s=f(x)$. Entonces se tienen isomorfismos canónicos
\[
  (\mathscr{P}_{X/S}^n)_x\xrightarrow{\sim}
  P_{\mathcal{O}_x/\mathcal{O}_s}^n
  \tag{16.4.16.1}
  \label{IV.16.4.16.1}
\]
que forman un sistema proyectivo.
\end{corollary}

Se deducen de aquí isomorfismos para los graduados asociados y, en particular,
un isomorfismo canónico
\[
  (\Omega_{X/S}^1)_x\xrightarrow{\sim}
  \Omega_{\mathcal{O}_x/\mathcal{O}_s}^1.
  \tag{16.4.16.2}
  \label{IV.16.4.16.2}
\]

\begin{corollary}[16.4.17]
\label{IV.16.4.17}
Sean $k$ un cuerpo, $K$ el cuerpo de fracciones racionales
$k(T_1,\ldots,T_r)$. Entonces, para todo entero $n$, el homomorfismo de
$K[U_1,\ldots,U_r]$ ($U_i$ indeterminadas) en $P_{K/k}^n$ que hace corresponder a cada $U_i$
$d^nT_i-T_i\cdot1$ es sobreyectivo y define un isomorfismo
del cociente $K[U_1,\ldots,U_r]/\mathfrak{m}^{n+1}$ (donde $\mathfrak{m}$ es
el ideal engendrado por las $U_i$) sobre $P_{K/k}^n$.
\end{corollary}

Esto resulta de (16.4.8), (16.4.10) y (16.4.14), donde se toma $A=k$,
$B=k[T_1,\ldots,T_r]$ y $S=B-\{0\}$.

Se recupera así el hecho de que los $dT_i$ forman una base del $K$-espacio
vectorial $\Omega_{K/k}^1$ (0, 20.5.10).

\begin{proposition}[16.4.18]
\label{IV.16.4.18}
Sean $f:X\to Y$, $g:Y\to Z$ dos morfismos de preesquemas; consideremos
los homomorfismos canónicos de $\mathcal{O}_X$-álgebras aumentadas
(16.4.3.3)
\[
  g_{X/Y/Z}:\mathscr{P}_{X/Z}^n\longrightarrow\mathscr{P}_{X/Y}^n
  \tag{16.4.18.1}
  \label{IV.16.4.18.1}
\]
\[
  f_{X/Y/Z}:f^*(\mathscr{P}_{Y/Z}^n)
  \longrightarrow\mathscr{P}_{X/Z}^n.
  \tag{16.4.18.2}
  \label{IV.16.4.18.2}
\]
Entonces $g_{X/Y/Z}$ es sobreyectivo, y su núcleo es el ideal engendrado por la imagen
por $f_{X/Y/Z}$ del ideal de aumentación de
$f^*(\mathscr{P}_{Y/Z}^n)$.
\end{proposition}

\oldpage[IV-4]{24}
Observemos primero que $g_{X/Y/Z}$ corresponde al caso en que, en (16.4.3.3), se toma
$X'=X$, $S'=Y$, $S=Z$ y $u=1_X$, y $f_{X/Y/Z}$ al caso en que se reemplazan
$X'$, $X$, $S$, $S'$ por $X$, $Y$, $Z$, $Z$, respectivamente, y $u$, $f$ por
$f$, $g$, respectivamente.

Se tiene un diagrama conmutativo (I, 5.3.5)
\[
  \xymatrix{
    X\ar[r]^-{\Delta_f}\ar[dr]_-f
      &X\times_Y X\ar[d]^-p\ar[r]^-j
      &X\times_Z X\ar[d]^-{f\times_Z f}\\
    &Y\ar[r]_-{\Delta_g}&Y\times_Z Y
  }
  \tag{16.4.18.3}
  \label{IV.16.4.18.3}
\]
donde $j=(1_X,1_X)_Z$ es una inmersión, $j\circ\Delta_f=\Delta_{g\circ f}$ y
$p$ es el morfismo estructural. Como basta considerar el caso en que $X$, $Y$,
$Z$ son afines, se puede suponer que las inmersiones $\Delta_f$, $\Delta_g$ y
$j$ son cerradas, de modo que $\mathcal{O}_X$ y
$\mathcal{O}_{X\times_Y X}$ se identifican, respectivamente, con
$\mathcal{O}_{X\times_Z X}/\mathscr{J}$ y
$\mathcal{O}_{X\times_Z X}/\mathscr{L}$, donde
$\mathscr{J}\supset\mathscr{L}$ son los dos ideales cuasicoherentes
correspondientes, respectivamente, a las inmersiones $\Delta_{g\circ f}$ y $j$. El
$\mathcal{O}_X$-álgebra $\mathscr{P}_{X/Z}^n$ se identifica, por tanto, con
$\mathcal{O}_{X\times_Z X}/\mathscr{J}^{n+1}$, y
$\mathscr{P}_{X/Y}^n$ se identifica con
$\mathcal{O}_{X\times_Y X}/(\mathscr{J}/\mathscr{L})^{n+1}$, es decir, con
$\mathcal{O}_{X\times_Z X}/(\mathscr{J}^{n+1}+\mathscr{L})$, y por consiguiente con el
cociente de $\mathscr{P}_{X/Z}^n$ por
$(\mathscr{J}^{n+1}+\mathscr{L})/\mathscr{J}^{n+1}$. Pero se sabe
(\emph{loc. cit.}) que $p$ y $j$ hacen de $X\times_Y X$ el producto fibrado de los
$(Y\times_Z Y)$-preesquemas $Y$ y $X\times_Z X$; por tanto, si $\mathcal{O}_Y$ se
identifica con $\mathcal{O}_{Y\times_Z Y}/\mathscr{K}$, donde $\mathscr{K}$ es
el ideal correspondiente a $\Delta_g$, $\mathscr{L}$ es igual a
$(f\times_Z f)^*(\mathscr{K})\cdot\mathcal{O}_{X\times_Z X}$ (I, 4.4.5).
Como $(\mathscr{J}^{n+1}+\mathscr{L})/\mathscr{J}^{n+1}$ es el ideal de
$\mathscr{P}_{X/Z}^n$ engendrado por la imagen de $\mathscr{L}$, se deduce la
proposición.

\begin{corollary}[16.4.19]
\label{IV.16.4.19}
Con las notaciones de (16.4.18), se tiene una sucesión exacta de
$\mathcal{O}_X$-módulos cuasicoherentes
\[
  \xymatrix{
    f^*(\Omega_{Y/Z}^1)\ar[r]^-{f_{X/Y/Z}}&
    \Omega_{X/Z}^1\ar[r]^-{g_{X/Y/Z}}&
    \Omega_{X/Y}^1\ar[r]&0.
  }
  \tag{16.4.19.1}
  \label{IV.16.4.19.1}
\]
\end{corollary}

Cuando $X$, $Y$, $Z$ son afines, se recupera así la sucesión exacta
(0, 20.5.7.1).

\begin{proposition}[16.4.20]
\label{IV.16.4.20}
Sean $f:Y\to Z$ un morfismo, $j:X\to Y$ una inmersión cerrada,
$\mathscr{K}$ el ideal cuasicoherente de $\mathcal{O}_Y$ correspondiente a $j$.
Entonces se tiene
$\mathscr{P}_{X/Y}^n=\mathcal{O}_X=\mathcal{O}_Y/\mathscr{K}$,
el homomorfismo canónico
$j_{X/Y/Z}:j^*(\mathscr{P}_{Y/Z}^n)\to\mathscr{P}_{X/Z}^n$ es sobreyectivo, y
su núcleo es el ideal de $j^*(\mathscr{P}_{Y/Z}^n)$ engendrado por
$j^*(\mathcal{O}_Y\cdot d_{Y/Z}^n(\mathscr{K}))$ (debe observarse que
$d_{Y/Z}^n(\mathscr{K})$ es un subhaz de grupos conmutativos de
$\mathscr{P}_{Y/Z}^n$, pero no un $\mathcal{O}_Y$-módulo en general).
\end{proposition}

Se sabe (I, 5.3.8) que la diagonal
$\Delta_j:X\to X\times_Y X$ es un isomorfismo, de donde se sigue la primera afirmación.
Si $\varpi_1$ y $\varpi_2$ son los dos homomorfismos de álgebras
$\mathcal{O}_Y\to\mathscr{P}_{Y/Z}^n$ correspondientes, respectivamente, a las dos
proyecciones canónicas $p_1$, $p_2$ de $Y\times_Z Y$ sobre $Y$, recordemos que,
por definición (16.3.5 y 16.3.6), $\varpi_1$ es el homomorfismo estructural del
$\mathcal{O}_Y$-álgebra $\mathscr{P}_{Y/Z}^n$ y
$\varpi_2=d_{Y/Z}^n$. El $\mathcal{O}_X$-álgebra
$j^*(\mathscr{P}_{Y/Z}^n)$ se identifica, por tanto, con
$\mathscr{P}_{Y/Z}^n/\varpi_1(\mathscr{K})\mathscr{P}_{Y/Z}^n$ y su cociente
por el ideal engendrado por $j^*(d_{Y/Z}^n(\mathscr{K}))$ con
\[
  \mathscr{P}_{Y/Z}^n/
  (\varpi_1(\mathscr{K})+\varpi_2(\mathscr{K}))\mathscr{P}_{Y/Z}^n.
\]
Observemos ahora que se tiene un diagrama conmutativo
\oldpage[IV-4]{25}
\[
  \xymatrix{
    Y\ar[d]_-{\Delta_f}&X\ar[l]_-j\ar[d]^-{\Delta_{f\circ j}}\\
    Y\times_Z Y&X\times_Z X\ar[l]^-{j\times_Z j}
  }
\]
que identifica $X$ con el producto fibrado de los $(Y\times_Z Y)$-preesquemas $Y$ y
$X\times_Z X$ (I, 5.3.7). Como $j\times_Z j$ es una inmersión, se deduce
de esta observación y de (16.2.2) que, si $\Delta_{Y/Z}^n$ y
$\Delta_{X/Z}^n$ denotan las vecindades infinitesimales de orden $n$ de $Y$
y $X$ para las inmersiones canónicas $\Delta_f$ y $\Delta_{f\circ j}$,
respectivamente, se tiene un diagrama
\[
  \xymatrix{
    \Delta_{Y/Z}^n\ar[d]&\Delta_{X/Z}^n\ar[l]\ar[d]\\
    Y\times_Z Y&X\times_Z X\ar[l]^-{j\times_Z j}
  }
\]
que hace de $\Delta_{X/Z}^n$ el producto fibrado de los $(Y\times_Z Y)$-preesquemas
$\Delta_{Y/Z}^n$ y $X\times_Z X$. También puede decirse que
$\mathscr{P}_{X/Z}^n$ se identifica con el haz de anillos
$\mathscr{P}_{Y/Z}^n\otimes_{\mathcal{O}_{Y\times_Z Y}}
\mathcal{O}_{X\times_Z X}$. Pero se ve inmediatamente (por ejemplo, reduciéndose
al caso afín) que
\[
  \mathcal{O}_{X\times_Z X}=
  \mathcal{O}_{Y\times_Z Y}/
  (p_1^*(\mathscr{K})+p_2^*(\mathscr{K}))
  \mathcal{O}_{Y\times_Z Y}.
\]
Por tanto, $\mathscr{P}_{X/Z}^n$ se identifica con el cociente de
$\mathscr{P}_{Y/Z}^n$ por el ideal engendrado por la imagen en
$\mathscr{P}_{Y/Z}^n$ de $p_1^*(\mathscr{K})+p_2^*(\mathscr{K})$. Pero, por
definición, este ideal también está engendrado por
$\varpi_1(\mathscr{K})+\varpi_2(\mathscr{K})$. Q.E.D.

\begin{corollary}[16.4.21]
\label{IV.16.4.21}
Sean $f:Y\to Z$ un morfismo, $j:X\to Y$ una inmersión. Se tiene una sucesión
exacta de $\mathcal{O}_X$-módulos cuasicoherentes
\[
  \xymatrix{
    \mathscr{N}_{X/Y}\ar[r]&j^*(\Omega_{Y/Z}^1)\ar[r]&
    \Omega_{X/Z}^1\ar[r]&0.
  }
  \tag{16.4.21.1}
  \label{IV.16.4.21.1}
\]
\end{corollary}

Cuando $X$, $Y$ y $Z$ son afines, se recupera así la sucesión exacta
(0, 20.5.12.1).

\begin{corollary}[16.4.22]
\label{IV.16.4.22}
Si $f:X\to S$ es un morfismo localmente de presentación finita,
$\mathscr{P}_{X/S}^n$ y $\Omega_{X/S}^1$ son
$\mathcal{O}_X$-módulos cuasicoherentes de presentación finita.
\end{corollary}

Se reduce inmediatamente al caso en que $S=\operatorname{Spec}(A)$ es afín,
$X=\operatorname{Spec}(B)$, donde
$B=A[T_1,\ldots,T_r]/\mathfrak{K}$, siendo $\mathfrak{K}$ un ideal de tipo
finito de $C=A[T_1,\ldots,T_r]$. Se aplica entonces (16.4.20) con $Z=S$,
$Y=\operatorname{Spec}(C)$ y $\mathscr{K}=\widetilde{\mathfrak{K}}$. Entonces
$j^*(\mathscr{P}_{Y/Z}^n)$ es un $\mathcal{O}_X$-módulo libre de rango finito
(16.4.10), y la hipótesis sobre $\mathfrak{K}$ implica que
$j^*(\mathcal{O}_Y\cdot d_{Y/Z}^n(\mathscr{K}))$ genera un
$\mathcal{O}_X$-módulo cuasicoherente de tipo finito; de donde se sigue la conclusión.

\begin{proposition}[16.4.23]
\label{IV.16.4.23}
Sean $X$, $Y$ dos $S$-preesquemas, $Z=X\times_S Y$ su producto fibrado,
$p:X\times_S Y\to X$ y $q:X\times_S Y\to Y$ las proyecciones canónicas.
Entonces el homomorfismo canónico
\[
  p_{Z/X/S}\oplus q_{Z/Y/S}:
  p^*(\Omega_{X/S}^1)\oplus q^*(\Omega_{Y/S}^1)
  \longrightarrow\Omega_{(X\times_S Y)/S}^1
  \tag{16.4.23.1}
  \label{IV.16.4.23.1}
\]
es biyectivo.
\end{proposition}

\oldpage[IV-4]{26}
El diagrama conmutativo
\[
  \xymatrix{
    Y\ar[d]_-g&X\times_S Y\ar[l]_-q\ar[d]_-h&
      X\times_S Y\ar[l]_-{\operatorname{id}}\ar[d]^-p\\
    S&S\ar[l]^-{\operatorname{id}}&X\ar[l]^-f
  }
\]
da una factorización del \emph{isomorfismo} canónico $P^n(p)$ (16.4.5)
\[
  p^*(\mathscr{P}_{X/S}^n)\longrightarrow
  \mathscr{P}_{Z/S}^n\longrightarrow\mathscr{P}_{Z/Y}^n
\]
y, del mismo modo, intercambiando los papeles de $X$ y $Y$, se tiene una factorización
del \emph{isomorfismo} $P^n(q)$
\[
  q^*(\mathscr{P}_{Y/S}^n)\longrightarrow
  \mathscr{P}_{Z/S}^n\longrightarrow\mathscr{P}_{Z/X}^n.
\]
Esto prueba que el homomorfismo canónico (16.4.18.1)
\[
  p_{Z/X/S}:p^*(\mathscr{P}_{X/S}^n)\longrightarrow\mathscr{P}_{Z/S}^n
  \quad\text{(resp.\
  $q_{Z/Y/S}:q^*(\mathscr{P}_{Y/S}^n)\to\mathscr{P}_{Z/S}^n$)}
\]
es \emph{inyectivo}, y que el núcleo del homomorfismo canónico sobreyectivo
(16.4.18.2)
\[
  \mathscr{P}_{Z/S}^n\longrightarrow\mathscr{P}_{Z/Y}^n
  \quad\text{(resp.\
  $\mathscr{P}_{Z/S}^n\to\mathscr{P}_{Z/X}^n$)}
\]
es suplementario de la imagen de $p_{Z/X/S}$ (resp.\ $q_{Z/Y/S}$). Pero,
por otra parte, este núcleo está, en virtud de (16.4.18), engendrado por la imagen por
$q_{Z/Y/S}$ (resp.\ $p_{Z/X/S}$) del ideal de aumentación de
$q^*(\mathscr{P}_{Y/S}^n)$ (resp.\ $p^*(\mathscr{P}_{X/S}^n)$). Se concluye
la proposición considerando el caso $n=1$.

Se generaliza inmediatamente (16.4.23) al caso de un producto fibrado de un número finito
arbitrario de $S$-preesquemas.

\begin{remarks}[16.4.24]
\label{IV.16.4.24}
\begin{enumerate}
\item[(i)] Veremos (17.2.3) que, cuando el morfismo $f:X\to Y$ de
(16.4.18) es \emph{liso}, el homomorfismo $f_{X/Y/Z}$ de (16.4.19.1) es
localmente \emph{invertible por la izquierda} y, en particular, inyectivo. Del mismo modo,
cuando el morfismo $f\circ j:X\to Z$ de (16.4.20) es \emph{liso},
el homomorfismo de la izquierda en (16.4.21.2) es localmente
\emph{invertible por la izquierda} y \emph{a fortiori} inyectivo (17.2.5). En el cap. V,
daremos también una variante, para módulos sobre
preesquemas, de los «módulos de imperfección» estudiados en (0, 20.6), y de las sucesiones
exactas en las que figuran.

\item[(ii)] Sean $X$ un espacio topológico, $\mathscr{A}$ un haz
de anillos sobre $X$ y $\mathscr{B}$ una $\mathscr{A}$-álgebra sobre $X$. Entonces está
claro que
\[
  U\longmapsto
  P_{\Gamma(U,\mathscr{B})/\Gamma(U,\mathscr{A})}^n
  \qquad (U\text{ abierto en }X)
\]
es un prehaz de $\Gamma(U,\mathscr{B})$-álgebras aumentadas; por tanto, el
haz asociado $\mathscr{P}_{\mathscr{B}/\mathscr{A}}^n$ es una
$\mathscr{B}$-álgebra aumentada. En el caso particular en que $X$ es un
preesquema, $f=(\psi,\theta):X\to S$ un morfismo de preesquemas, resulta
fácilmente de (16.4.16) y de la exactitud del funtor $\varinjlim$ que
$\mathscr{P}_{X/S}^n$ es canónicamente isomorfo a
$\mathscr{P}_{\mathcal{O}_X/\psi^*(\mathcal{O}_S)}^n$. De ello se sigue que el
formalismo desarrollado en el presente parágrafo podría considerarse como
un
\oldpage[IV-4]{27}
caso particular de un formalismo diferencial para espacios anillados dotados
de un haz de álgebras sobre el haz estructural. Sin embargo, no hemos
querido partir de este punto de vista, menos intuitivo y menos cómodo para las
aplicaciones. Parece, además, que, para las diversas especies de
«variedades», la construcción «global» de los $\mathscr{P}^n$, análoga a la que
utilizamos aquí, está también mejor adaptada a las aplicaciones.
\end{enumerate}
\end{remarks}

\subsection{Haces y fibrados tangentes relativos; derivaciones.}
\label{IV.16.5}

\begin{env}[16.5.1]
\label{IV.16.5.1}
Sea $f=(\psi,\theta):X\to S$ un morfismo de espacios anillados. Para todo
$\mathcal{O}_X$-módulo $\mathscr{F}$, se llama \emph{$S$-derivación}
(o \emph{$(X/S)$-derivación}, o \emph{$f$-derivación}) \emph{de
$\mathcal{O}_X$ en $\mathscr{F}$} a todo homomorfismo de \emph{haces de
grupos aditivos} $D:\mathcal{O}_X\to\mathscr{F}$ que satisface las condiciones
siguientes:
\begin{enumerate}
\item[(a)] para todo abierto $V$ de $X$ y todo par de secciones
$(t_1,t_2)$ de $\mathcal{O}_X$ sobre $V$, se tiene
\[
  D(t_1t_2)=t_1D(t_2)+D(t_1)t_2;
  \tag{16.5.1.1}
  \label{IV.16.5.1.1}
\]
\item[(b)] para todo abierto $V$ de $X$, toda sección $t$ de $\mathcal{O}_X$
sobre $V$ y toda sección $s$ de $\mathcal{O}_S$ sobre un abierto
$U$ de $S$ tal que $V\subset f^{-1}(U)$, se tiene
\[
  D((s|V)t)=(s|V)D(t).
  \tag{16.5.1.2}
  \label{IV.16.5.1.2}
\]
\end{enumerate}
Está claro que equivale a decir que, para todo $x\in X$,
el homomorfismo de grupos aditivos
$D_x:\mathcal{O}_x\to\mathscr{F}_x$ es una
$\mathcal{O}_{f(x)}$-derivación.

Otra interpretación consiste en considerar la $\mathcal{O}_X$-álgebra
$\mathscr{D}_{\mathcal{O}_X}(\mathscr{F})$ igual a
$\mathcal{O}_X\oplus\mathscr{F}$, cuya estructura de álgebra se define por la
condición de que, para todo abierto $V$ de $X$, el producto de dos secciones de
$\mathcal{O}_X$ (resp. de una sección de $\mathcal{O}_X$ y una sección de
$\mathscr{F}$) sobre $V$ viene dado por la estructura de anillo de
$\Gamma(V,\mathcal{O}_X)$ (resp. por la estructura de
$\Gamma(V,\mathcal{O}_X)$-módulo de $\Gamma(V,\mathscr{F})$), y el producto de
dos secciones de $\mathscr{F}$ sobre $V$ se toma igual a 0; entonces,
$\mathscr{F}$ es un ideal de
$\mathscr{D}_{\mathcal{O}_X}(\mathscr{F})$, núcleo del aumento canónico
$\mathscr{D}_{\mathcal{O}_X}(\mathscr{F})\to\mathcal{O}_X$, y decir que $D$
es una $S$-derivación de $\mathcal{O}_X$ en $\mathscr{F}$ significa que
$1_{\mathcal{O}_X}+D$ es un $\mathcal{O}_S$-homomorfismo de álgebras de
$\mathcal{O}_X$ en $\mathscr{D}_{\mathcal{O}_X}(\mathscr{F})$ que, compuesto
con el aumento, da $1_{\mathcal{O}_X}$.

Las $S$-derivaciones de $\mathcal{O}_X$ en $\mathscr{F}$ forman evidentemente
un $\Gamma(X,\mathcal{O}_X)$-módulo
$\Der_S(\mathcal{O}_X,\mathscr{F})$.

Cuando $\mathscr{F}=\mathcal{O}_X$, una $S$-derivación de $\mathcal{O}_X$ en
sí mismo se llama simplemente una \emph{$S$-derivación de $\mathcal{O}_X$}.
\end{env}

\begin{proposition}[16.5.2]
\label{IV.16.5.2}
Sean $A$ un anillo, $B$ una $A$-álgebra, $L$ un $B$-módulo; pongamos
$S=\operatorname{Spec}(A)$, $X=\operatorname{Spec}(B)$,
$\mathscr{F}=\widetilde{L}$. Entonces la aplicación $D\mapsto\Gamma(D)$ que a
toda $S$-derivación $D$ de $\mathcal{O}_X$ en $\mathscr{F}$ hace
corresponder la aplicación $\Gamma(D):t\mapsto D(t)$ de $B$ en $L$ es un
isomorfismo de $B$-módulos de
$\Der_S(\mathcal{O}_X,\mathscr{F})$ en
$\Der_A(B,L)$ (cf. (0, 20.1.2)).
\end{proposition}

Esto resulta inmediatamente de la interpretación dada anteriormente de las $S$-derivaciones
en términos
\oldpage[IV-4]{28}
de homomorfismos de álgebras, de la interpretación análoga dada en
(0, 20.1.6) y de la correspondencia canónica entre homomorfismos de
$\mathcal{O}_X$-álgebras y homomorfismos de $B$-álgebras (I, 1.3.13 y
1.3.8).

\begin{proposition}[16.5.3]
\label{IV.16.5.3}
Sea $f=(\psi,\theta):X\to S$ un morfismo de preesquemas.
\begin{enumerate}
\item[(i)] La diferencial
$d_{X/S}:\mathcal{O}_X\to\Omega_{X/S}^1$ (16.3.6) es una $S$-derivación.
\item[(ii)] Para todo $\mathcal{O}_X$-módulo $\mathscr{F}$, la aplicación
$u\mapsto u\circ d_{X/S}$ es un isomorfismo de
$\Gamma(X,\mathcal{O}_X)$-módulos
\[
  \operatorname{Hom}_{\mathcal{O}_X}(\Omega_{X/S}^1,\mathscr{F})
  \xrightarrow{\sim}
  \Der_S(\mathcal{O}_X,\mathscr{F}).
  \tag{16.5.3.1}
  \label{IV.16.5.3.1}
\]
\end{enumerate}
\end{proposition}

La afirmación (i) ya se señaló en (16.3.6). Por otra parte, es inmediato (en
virtud de (0, 20.4.8)) que $u\mapsto u\circ d_{X/S}$ es inyectiva, considerando
las restricciones de ambos miembros al tallo $\mathcal{O}_x$ y utilizando
(16.4.16.2). Para ver que el homomorfismo (16.5.3.1) es sobreyectivo,
consideremos una $S$-derivación $D:\mathcal{O}_X\to\mathscr{F}$; para todo
abierto afín $V=\operatorname{Spec}(B)$ de $X$ tal que $f(V)$ esté contenido
en un abierto afín $U=\operatorname{Spec}(A)$ de $S$,
$D_V:B\to\Gamma(V,\mathscr{F})$ es una $A$-derivación, por lo que existe un
\emph{único} $B$-homomorfismo
$u_V:\Omega_{B/A}^1\to\Gamma(V,\mathscr{F})$ tal que
$D_V=u_V\circ d_{B/A}$ (0, 20.4.8); además, la unicidad de $u_V$ muestra
inmediatamente que, para todo abierto afín $W\subset V$, se tiene $u_W=u_V|W$, por lo que los
$u_V$ definen un homomorfismo de $\mathcal{O}_X$-módulos
$u:\mathcal{O}_X\to\mathscr{F}$ que responde a la cuestión.

\begin{env}[16.5.4]
\label{IV.16.5.4}
Con las notaciones de (16.5.1), para todo abierto $U$ de $X$,
$\Der_S(\mathcal{O}_U,\mathscr{F}|U)$ es un
$\Gamma(U,\mathcal{O}_X)$-módulo y está claro que la aplicación
$U\mapsto\Der_S(\mathcal{O}_U,\mathscr{F}|U)$ es un
prehaz; de hecho, es incluso un \emph{haz} (y, por tanto, un
$\mathcal{O}_X$-módulo), en virtud de la caracterización puntual de las
$S$-derivaciones vista en (16.5.1). Este $\mathcal{O}_X$-módulo se denota por
$\mathscr{D}\!{\rm er}_S(\mathcal{O}_X,\mathscr{F})$ y se llama
\emph{haz de $S$-derivaciones de $\mathcal{O}_X$ en $\mathscr{F}$}; lo
que acabamos de ver se expresa también mediante el corolario siguiente:
\end{env}

\begin{corollary}[16.5.5]
\label{IV.16.5.5}
Para todo $\mathcal{O}_X$-módulo $\mathscr{F}$, el homomorfismo de
$\mathcal{O}_X$-módulos deducido de $u\mapsto u\circ d_{X/S}$
\[
  \mathscr{H}\!om_{\mathcal{O}_X}(\Omega_{X/S}^1,\mathscr{F})
  \longrightarrow
  \mathscr{D}\!{\rm er}_S(\mathcal{O}_X,\mathscr{F})
  \tag{16.5.5.1}
  \label{IV.16.5.5.1}
\]
es biyectivo.
\end{corollary}

\begin{corollary}[16.5.6]
\label{IV.16.5.6}
\begin{enumerate}
\item[(i)] Si el morfismo $f:X\to S$ es localmente de presentación finita y
si $\mathscr{F}$ es un $\mathcal{O}_X$-módulo cuasicoherente, entonces
$\mathscr{D}\!{\rm er}_S(\mathcal{O}_X,\mathscr{F})$ es un
$\mathcal{O}_X$-módulo cuasicoherente.
\item[(ii)] Si además $S$ es localmente noetheriano y si $\mathscr{F}$ es
coherente, entonces $\mathscr{D}\!{\rm er}_S(\mathcal{O}_X,\mathscr{F})$ es un
$\mathcal{O}_X$-módulo coherente.
\end{enumerate}
\end{corollary}

La afirmación (i) resulta del isomorfismo (16.5.5.1), de (16.4.22) y de
(I, 1.3.12); la afirmación (ii) resulta de (0$_{\rm I}$, 5.3.5).

\begin{env}[16.5.7]
\label{IV.16.5.7}
Pongamos
\[
  \mathfrak{G}_{X/S}=
  \mathscr{H}\!om_{\mathcal{O}_X}(\Omega_{X/S}^1,\mathcal{O}_X)
  =\mathscr{D}\!{\rm er}_S(\mathcal{O}_X,\mathcal{O}_X)
  \tag{16.5.7.1}
  \label{IV.16.5.7.1}
\]
y se dice que es el \emph{haz de $S$-derivaciones de
$\mathcal{O}_X$}, o también el \emph{haz tangente de $X$ relativo a
$S$}: es, por tanto, el \emph{dual} del $\mathcal{O}_X$-módulo
$\Omega_{X/S}^1$. Si $f$ es localmente de presentación
\oldpage[IV-4]{29}
finita, $\mathfrak{G}_{X/S}$ es un $\mathcal{O}_X$-módulo cuasicoherente; si
además $S$ es localmente noetheriano, $\mathfrak{G}_{X/S}$ es coherente
(16.5.6).
\end{env}

\begin{env}[16.5.8]
\label{IV.16.5.8}
Supongamos, más particularmente, que $\Omega_{X/S}^1$ es un
$\mathcal{O}_X$-módulo \emph{localmente libre} (de rango finito) (lo que ocurrirá
cuando $f$ sea \emph{liso} (17.2.3)); entonces $\mathfrak{G}_{X/S}$ es un
$\mathcal{O}_X$-módulo localmente libre y tiene el mismo rango que $\Omega_{X/S}^1$
en cada punto. Más precisamente, supongamos que $\Omega_{X/S}^1$ tiene
rango $n$ en un punto $x$; existen entonces $n$ secciones $s_i$ ($1\leq i\leq n$)
de $\mathcal{O}_X$ sobre una vecindad afín $U$ de $x$ tales que las
imágenes canónicas de los $ds_i$ en
$\Omega_{X/S}^1\otimes_{\mathcal{O}_X}k(x)$ forman una base de este
$k(x)$-espacio vectorial; en virtud del lema de Nakayama, los gérmenes
$(ds_i)_x=d(s_i)_x$ de los $ds_i$ en el punto $x$ forman una base del
$\mathcal{O}_x$-módulo $(\Omega_{X/S}^1)_x$, por lo que, restringiendo $U$, se
puede suponer que los $ds_i$ forman una \emph{base} del
$\Gamma(U,\mathcal{O}_X)$-módulo $\Gamma(U,\Omega_{X/S}^1)$. Entonces el
$\Gamma(U,\mathcal{O}_X)$-módulo $\Gamma(U,\mathfrak{G}_{X/S})$ es dual del
anterior; se denota por $(D_i)_{1\leq i\leq n}$ o por
$\left(\frac{\partial}{\partial s_i}\right)_{1\leq i\leq n}$ la base dual de
$(ds_i)_{1\leq i\leq n}$, de modo que, por (16.5.3), se tiene
\[
  D_is_j=\langle D_i,ds_j\rangle
  =\left\langle\frac{\partial}{\partial s_i},ds_j\right\rangle
  =\delta_{ij}\qquad\text{(índice de Kronecker).}
  \tag{16.5.8.1}
  \label{IV.16.5.8.1}
\]
Toda $\Gamma(S,\mathcal{O}_S)$-derivación de la
$\Gamma(S,\mathcal{O}_S)$-álgebra $\Gamma(U,\mathcal{O}_X)$ se escribe, por tanto,
de manera única
\[
  D=\sum_{i=1}^n a_iD_i=\sum_{i=1}^n a_i\frac{\partial}{\partial s_i}
\]
donde los $a_i$ ($1\leq i\leq n$) son secciones de $\mathcal{O}_X$ sobre
$U$. Para toda sección $g\in\Gamma(U,\mathcal{O}_X)$, si se pone
$dg=\sum_{i=1}^n c_i ds_i$, se tiene
$c_i=\langle D_i,dg\rangle=D_ig$ en virtud de (16.5.8.1), es decir,
\[
  dg=\sum_{i=1}^n(D_ig)ds_i
  =\sum_{i=1}^n\frac{\partial g}{\partial s_i}ds_i.
  \tag{16.5.8.2}
  \label{IV.16.5.8.2}
\]
\end{env}

\begin{env}[16.5.9]
\label{IV.16.5.9}
Sean $D_1$, $D_2$ dos $S$-derivaciones de $\mathcal{O}_X$. Para todo abierto
$U$ de $X$, si $D_1^U$, $D_2^U$ son las derivaciones correspondientes del
anillo $\Gamma(U,\mathcal{O}_X)$, el \emph{corchete}
\[
  [D_1^U,D_2^U]=D_1^U\circ D_2^U-D_2^U\circ D_1^U
\]
es también una derivación de este anillo, por lo que el
$\psi^*(\mathcal{O}_S)$-endomorfismo de $\mathcal{O}_X$
\[
  [D_1,D_2]=D_1\circ D_2-D_2\circ D_1
  \tag{16.5.9.1}
  \label{IV.16.5.9.1}
\]
es de nuevo una $S$-derivación; como se comprueba inmediatamente que este corchete
satisface la identidad de Jacobi, se ve que se ha definido así sobre
$\Der_S(\mathcal{O}_X,\mathcal{O}_X)$ una estructura de
$\Gamma(S,\mathcal{O}_S)$-álgebra de Lie. Como la definición de esta
estructura conmuta con la restricción a un abierto de $X$, se ve que
$\mathfrak{G}_{X/S}$ está canónicamente dotado de una estructura de
$\psi^*(\mathcal{O}_S)$-álgebra de Lie. Obsérvese que la aplicación
$(D_1,D_2)\mapsto[D_1,D_2]$ no es
$\Gamma(X,\mathcal{O}_X)$-bilineal.
\end{env}

\begin{env}[16.5.10]
\label{IV.16.5.10}
Para todo cambio de base $g:S'\to S$, si se pone $X'=X\times_S S'$, se
ha visto en (16.4.5) que se tiene un isomorfismo canónico
\[
  \Omega_{X/S}^1\otimes_S S'\xrightarrow{\sim}\Omega_{X'/S'}^1
  \tag{16.5.10.1}
  \label{IV.16.5.10.1}
\]
\oldpage[IV-4]{30}
del que se deduce, en virtud de (16.5.10.1), un homomorfismo canónico
(Bourbaki, \emph{Alg.}, cap. II, 3\textsuperscript{a} ed., \S~5,
n\textsuperscript{o} 3)
\[
  \mathfrak{G}_{X/S}\otimes_{\mathcal{O}_S}\mathcal{O}_{S'}
  \longrightarrow\mathfrak{G}_{X'/S'}
  \tag{16.5.10.2}
  \label{IV.16.5.10.2}
\]
que, en general, no es ni inyectivo ni sobreyectivo. Sin embargo:
\end{env}

\begin{proposition}[16.5.11]
\label{IV.16.5.11}
\begin{enumerate}
\item[(i)] Si $g:S'\to S$ es un morfismo plano y si $f$ es localmente de
tipo finito (resp.\ localmente de presentación finita), el homomorfismo
(16.5.10.2) es inyectivo (resp.\ biyectivo).
\item[(ii)] Si $\Omega_{X/S}^1$ es un $\mathcal{O}_X$-módulo localmente libre
de tipo finito, el homomorfismo (16.5.10.2) es biyectivo.
\end{enumerate}
\end{proposition}

En efecto, la afirmación (ii) resulta de Bourbaki, \emph{Alg.}, cap. II,
3\textsuperscript{a} ed., \S~5, n\textsuperscript{o} 3, prop. 7. La afirmación
(i) resulta asimismo de Bourbaki, \emph{Alg. comm.}, cap. I, \S~2,
n\textsuperscript{o} 10, prop. 11 y del hecho de que, si $f$ es localmente de tipo
finito (resp.\ localmente de presentación finita), $\Omega_{X/S}^1$ es un
$\mathcal{O}_X$-módulo de tipo finito (resp.\ de presentación finita) ((16.3.9)
y (16.4.22)).

\begin{env}[16.5.12]
\label{IV.16.5.12}
Puesto que $\Omega_{X/S}^1$ es un $\mathcal{O}_X$-módulo cuasicoherente, se puede
considerar el fibrado vectorial sobre $X$ definido por $\Omega_{X/S}^1$
(II, 1.7.8)
\[
  T_{X/S}=\mathbf{V}(\Omega_{X/S}^1)
  \tag{16.5.12.1}
  \label{IV.16.5.12.1}
\]
que se llama \emph{fibrado tangente de $X$ relativo a $S$}. Se tiene, por tanto, una
biyección canónica (II, 1.7.9)
\[
  \Gamma(T_{X/S}/S)\xrightarrow{\sim}
  \operatorname{Hom}_{\mathcal{O}_X}(\Omega_{X/S}^1,\mathcal{O}_X)
  =\Gamma(X,\mathfrak{G}_{X/S})
\]
por definición de $\mathfrak{G}_{X/S}$, y en este isomorfismo se puede
reemplazar $X$ por un abierto cualquiera $U$ de $X$; se puede decir, por tanto, que el
\emph{haz tangente} de $X$ relativo a $S$ es isomorfo al
\emph{haz de gérmenes de secciones sobre $S$} del fibrado tangente de $X$
relativo a $S$. Si $f:X\to Y$ es un $S$-morfismo, se ha visto en (16.4.19)
que se tiene un homomorfismo canónico
$f_{X/Y/S}:f^*(\Omega_{Y/S}^1)\to\Omega_{X/S}^1$, lo que da, teniendo en cuenta
que
\[
  \mathbf{V}(f^*(\Omega_{Y/S}^1))
  =\mathbf{V}(\Omega_{Y/S}^1)\times_Y X\qquad\text{(II, 1.7.11),}
\]
un $X$-morfismo
$T_{X/S}(f):T_{X/S}\to T_{Y/S}\times_Y X$. Si $g:Y\to Z$ es un segundo
$S$-morfismo, se tiene
$T_{X/S}(g\circ f)=(T_{Y/S}(g)\times1_X)\circ T_{X/S}(f)$ (0, 20.5.4.1).

Resulta de (16.5.10.1) y de (II, 1.7.11) que, para todo cambio de base
$g:S'\to S$, se tiene un isomorfismo canónico
\[
  T_{X'/S'}\xrightarrow{\sim}T_{X/S}\times_S S'
  =T_{X/S}\times_X X'.
  \tag{16.5.12.2}
  \label{IV.16.5.12.2}
\]
\end{env}

\begin{env}[16.5.13]
\label{IV.16.5.13}
Para todo punto $x\in X$, se llama \emph{espacio tangente a $X$ en el punto $x$}
(relativo a $S$) al conjunto de puntos de la fibra
$T_{X/S}\times_X\operatorname{Spec}(k(x))$ \emph{racionales sobre $k(x)$}, es decir,
al conjunto
\[
  T_{X/S}(x)=
  \operatorname{Hom}_{k(x)}
  (\Omega_{X/S}^1\otimes_{\mathcal{O}_x}k(x),k(x))
  \tag{16.5.13.1}
  \label{IV.16.5.13.1}
\]
que es el \emph{dual} del $k(x)$-espacio vectorial
$\Omega_{\mathcal{O}_x/\mathcal{O}_s}^1/
\mathfrak{m}_x\cdot\Omega_{\mathcal{O}_x/\mathcal{O}_s}^1$. Cuando
$\Omega_{X/S}^1$ es un $\mathcal{O}_X$-módulo de \emph{tipo finito},
$T_{X/S}(x)$ es, por tanto, un espacio vectorial de rango finito sobre $k(x)$ y, para
todo
\oldpage[IV-4]{31}
cambio de base $g:S'\to S$ y todo punto
$x'\in X'=X\times_S S'$ sobre $x$, se tiene un isomorfismo canónico
\[
  T_{X'/S'}(x')\xrightarrow{\sim}
  T_{X/S}(x)\otimes_{k(x)}k(x').
  \tag{16.5.13.2}
  \label{IV.16.5.13.2}
\]

Si $x$ es \emph{racional sobre $k(s)$}, donde $s=f(x)$ (de modo que
$k(s)\to k(x)$ es un isomorfismo), resulta de (16.4.13) que se tiene un
isomorfismo canónico
\[
  T_{X/S}(x)=T_{X_s/k(s)}(x)
  =\operatorname{Hom}_{k(s)}
  (\mathfrak{m}'_x/\mathfrak{m}'_x{}^2,k(s))
  \tag{16.5.13.3}
  \label{IV.16.5.13.3}
\]
donde $\mathfrak{m}'_x$ es el ideal maximal de
$\mathcal{O}_{X_s,x}=\mathcal{O}_{X,x}/
\mathfrak{m}_s\mathcal{O}_{X,x}$. En el caso en que $S$ es el espectro de un cuerpo
$k$, se recupera así la definición de Zariski del espacio tangente en un
punto $x\in X$ \emph{racional sobre $k$}, como dual de
$\mathfrak{m}_x/\mathfrak{m}_x^2$.

Sea $Y$ un segundo $S$-preesquema y sea $g:Y\to X$ un $S$-morfismo; se tiene
entonces un homomorfismo canónico de $\mathcal{O}_Y$-módulos (16.4.19)
\[
  g_{Y/X/S}:g^*(\Omega_{X/S}^1)\longrightarrow\Omega_{Y/S}^1.
  \tag{16.5.13.4}
  \label{IV.16.5.13.4}
\]
Observemos ahora que, si $y\in Y$ y $x=g(y)$, se tiene
\[
  g^*(\Omega_{X/S}^1)\otimes_{\mathcal{O}_Y}k(y)
  =(\Omega_{X/S}^1\otimes_{\mathcal{O}_X}k(x))
  \otimes_{k(x)}k(y)
\]
y, por consiguiente, si $\Omega_{X/S}^1$ es un $\mathcal{O}_X$-módulo
\emph{de tipo finito}, se puede identificar
\[
  \operatorname{Hom}_{k(y)}
  (g^*(\Omega_{X/S}^1)\otimes_{\mathcal{O}_Y}k(y),k(y))
\]
con $T_{X/S}(x)\otimes_{k(x)}k(y)$. Se deduce entonces del homomorfismo
(16.5.13.4) un homomorfismo de \emph{$k(y)$-espacios vectoriales}
\[
  T_y(g):T_{Y/S}(y)\longrightarrow
  T_{X/S}(x)\otimes_{k(x)}k(y)
  \tag{16.5.13.5}
  \label{IV.16.5.13.5}
\]
llamado \emph{aplicación lineal tangente a $g$ en el punto $y$}. Cuando $y$ es
\emph{racional sobre $k(s)$}, se pueden identificar $k(s)$, $k(y)$ y $k(x)$ y
$T_y(g)$ es entonces un homomorfismo de $k(s)$-espacios vectoriales
$T_{Y/S}(y)\to T_{X/S}(x)$; obsérvese además que, en este caso,
$g^*(\Omega_{X/S}^1)\otimes_{\mathcal{O}_Y}k(y)$ se identifica con
$\Omega_{X/S}^1\otimes_{\mathcal{O}_X}k(x)$, y el homomorfismo anterior queda
así definido sin condición alguna de finitud sobre $\Omega_{X/S}^1$ y no es
otro que el homomorfismo $T_{Y/S}(g)$ (16.5.12) restringido a la fibra en el punto
$y$ de $T_{Y/S}$.
\end{env}

\begin{env}[16.5.14]
\label{IV.16.5.14}
La interpretación de las derivaciones de una $A$-álgebra $B$ en un $B$-módulo
$L$, dada en (0, 20.1.1), se traduce al lenguaje de los preesquemas de la
manera siguiente.

Consideremos dos morfismos de preesquemas $f:X\to S$, $g:Y\to S$, y un
subpreesquema cerrado $Y_0$ de $Y$ definido por un ideal de cuadrado nulo
$\mathscr{J}$ de $\mathcal{O}_Y$ (de modo que $Y$ y $Y_0$ tienen el
\emph{mismo espacio topológico subyacente}). Supongamos dado un $S$-morfismo
$u_0:Y_0\to X$, de modo que se tiene un diagrama conmutativo
\[
  \xymatrix{
    X\ar[d]_-f & Y_0\ar[l]_-{u_0}\ar[d]^j\\
    S & Y\ar@{-->}[ul]_-u\ar[l]^-g
  }
  \tag{16.5.14.1}
  \label{IV.16.5.14.1}
\]
\oldpage[IV-4]{32}
y nos proponemos determinar si existen \emph{$S$-morfismos}
$u:Y\to X$ tales que $u_0=u\circ j$ (es decir, si es posible
completar el diagrama anterior con la flecha punteada $u$, de manera que siga
siendo \emph{conmutativo}).

Para ello, consideremos un abierto afín $U=\operatorname{Spec}(C)$ de $Y$;
su imagen inversa $j^{-1}(U)$ es el abierto afín
$U_0=\operatorname{Spec}(C/\mathfrak{L})$, donde
$\mathfrak{L}=\Gamma(U,\mathscr{J})$, ideal \emph{de cuadrado nulo} en $C$;
supondremos $U$ suficientemente pequeño para que $u_0(U_0)$ esté contenido en un
abierto afín $V=\operatorname{Spec}(B)$ de $X$, y
$g(U)=f(u_0(U_0))$ contenido en un abierto afín
$W=\operatorname{Spec}(A)$ de $S$, de modo que $B$ y $C$ son
$A$-álgebras y que $u_0|U_0$ corresponde a un $A$-homomorfismo $\psi$ de
$B$ en $C/\mathfrak{L}$; sea $P(U_0)$ el conjunto de las restricciones $u|U$
de los morfismos buscados, que corresponden canónicamente a los
$A$-homomorfismos de álgebras $\varphi:B\to C$ tales que el \emph{compuesto}
\[
  B\xrightarrow{\varphi}C\longrightarrow C/\mathfrak{L}
\]
\emph{sea igual a $\psi$}. Se sabe entonces por (0, 20.1.1) que el conjunto de estos
homomorfismos está vacío o es de la forma
$\varphi_1+\operatorname{D\acute{e}r}_A(B,\mathfrak{L})$; cuando $P(U_0)$ no es
vacío, el grupo aditivo $\operatorname{D\acute{e}r}_A(B,\mathfrak{L})$ actúa por
adición en $P(U_0)$, que es, por tanto, un \emph{espacio afín} para el grupo
aditivo $\operatorname{D\acute{e}r}_A(B,\mathfrak{L})$ (o también un
\emph{espacio homogéneo principal} (o \emph{torsor}) \emph{bajo}
$\operatorname{D\acute{e}r}_A(B,\mathfrak{L})$).

Observemos ahora que, puesto que $\mathfrak{L}$ está dotado de una estructura de
$B$-módulo por medio de $\psi$, se tiene un \emph{isomorfismo}
$v\mapsto v\circ d_{B/A}$ de
$\operatorname{Hom}_B(\Omega_{B/A}^1,\mathfrak{L})$ en
$\operatorname{D\acute{e}r}_A(B,\mathfrak{L})$ (0, 20.4.8). Además, como
$\mathfrak{L}$ es de cuadrado nulo y, por tanto, un $(C/\mathfrak{L})$-módulo, todo
$B$-homomorfismo $v:\Omega_{B/A}^1\to\mathfrak{L}$ puede considerarse
como un $(C/\mathfrak{L})$-homomorfismo
$\Omega_{B/A}^1\otimes_B(C/\mathfrak{L})\to\mathfrak{L}$. Como
$\mathscr{J}$ es de cuadrado nulo, puede considerarse como un
$\mathcal{O}_{Y_0}$-módulo cuasicoherente; introduzcamos el
$\mathcal{O}_{Y_0}$-módulo
\[
  \mathscr{G}=\mathscr{H}om(u_0^*(\Omega_{X/S}^1),\mathscr{J});
  \tag{16.5.14.2}
  \label{IV.16.5.14.2}
\]
resulta entonces del hecho de que
$\Omega_{B/A}^1=\Gamma(V,\Omega_{X/S}^1)$ (16.3.7) que se puede escribir
$\operatorname{D\acute{e}r}_A(B,\mathfrak{L})=\Gamma(U_0,\mathscr{G})$.

Como $P(U_0)$ se define como conjunto de $S$-morfismos $U\to X$, está
claro que $U_0\mapsto P(U_0)$ es un \emph{haz de conjuntos}
$\mathscr{P}$ sobre $Y_0$. Utilicemos este hecho para probar que la aplicación
\[
  h:\Gamma(U_0,\mathscr{G})\times P(U_0)\longrightarrow P(U_0)
\]
que define la estructura de torsor sobre $P(U_0)$ es también independiente de la
elección de $V$ y $W$, y además que, si $U'\subset U$ es un segundo abierto
afín de $Y$, y $U'_0$ su imagen inversa en $Y_0$, el diagrama
\[
  \xymatrix{
    \Gamma(U_0,\mathscr{G})\times P(U_0)\ar[d]\ar[r]^h
      & P(U_0)\ar[d]\\
    \Gamma(U'_0,\mathscr{G})\times P(U'_0)\ar[r]_{h'}
      & P(U'_0)
  }
  \tag{16.5.14.3}
  \label{IV.16.5.14.3}
\]
es conmutativo (siendo las flechas verticales los operadores de restricción).
En virtud de la observación anterior, se reduce a probar la conmutatividad
del diagrama precedente cuando $h$ se define como arriba a partir de los
abiertos afines $V$, $W$ y $h'$ a partir de abiertos
\oldpage[IV-4]{33}
afines $V'\subset V$ y $W'\subset W$. Pero, en virtud de la descripción
anterior de $h$, esto resulta de la conmutatividad del diagrama (0, 20.5.4.2).

Las aplicaciones
$\Gamma(U_0,\mathscr{G})\times P(U_0)\to P(U_0)$ definen, por tanto, un
homomorfismo de \emph{haces de conjuntos}
$m:\mathscr{G}\times\mathscr{P}\to\mathscr{P}$ tal que, para todo abierto
$U_0$ para el cual $\Gamma(U_0,\mathscr{P})\neq\varnothing$,
\[
  m_{U_0}:\Gamma(U_0,\mathscr{G})\times\Gamma(U_0,\mathscr{P})
  \longrightarrow\Gamma(U_0,\mathscr{P})
\]
sea una ley externa que define sobre $\Gamma(U_0,\mathscr{P})$ una estructura
de torsor bajo el grupo $\Gamma(U_0,\mathscr{G})$.
\end{env}

\begin{env}[16.5.15]
\label{IV.16.5.15}
De manera general, cuando se dan sobre un espacio topológico $Z$ un
haz de conjuntos $\mathscr{P}$, un haz de grupos $\mathscr{G}$ (no
necesariamente conmutativo) y un homomorfismo de haces de conjuntos
$m:\mathscr{G}\times\mathscr{P}\to\mathscr{P}$ tal que, para todo abierto
$U\subset Z$ para el cual $\Gamma(U,\mathscr{P})\neq\varnothing$,
\[
  m_U:\Gamma(U,\mathscr{G})\times\Gamma(U,\mathscr{P})
  \longrightarrow\Gamma(U,\mathscr{P})
\]
haga de $\Gamma(U,\mathscr{P})$ un \emph{torsor} bajo el grupo
$\Gamma(U,\mathscr{G})$, se dice que $\mathscr{P}$ es un
\emph{pseudotorsor} (o \emph{haz homogéneo formalmente principal})
bajo el haz de grupos $\mathscr{G}$. Se dice que $\mathscr{P}$ es un
\emph{torsor} (o \emph{haz homogéneo principal}) bajo $\mathscr{G}$ si,
además, $\Gamma(U,\mathscr{P})\neq\varnothing$ para todo abierto
$U\neq\varnothing$ de una base adecuada de la topología de $Z$.

Para la teoría general de los torsores, remitimos a [42]; aquí nos
limitaremos a recordar la correspondencia canónica entre las clases
de isomorfismo de estos torsores (para un $\mathscr{G}$ dado) y los elementos del
conjunto de cohomología $H^1(Z,\mathscr{G})$. Consideremos, en efecto, un
torsor $\mathscr{P}$ bajo $\mathscr{G}$ y un recubrimiento abierto
$(U_\lambda)$ de $Z$ tal que
$\Gamma(U_\lambda,\mathscr{P})\neq\varnothing$ para todo $\lambda$;
designemos por $p_\lambda$ un elemento de
$\Gamma(U_\lambda,\mathscr{P})$. Para todo par de índices $\lambda$, $\mu$
tal que $U_\lambda\cap U_\mu\neq\varnothing$, existe entonces un único elemento
$\gamma_{\lambda\mu}$ de
$\Gamma(U_\lambda\cap U_\mu,\mathscr{G})$ tal que
\[
  \gamma_{\lambda\mu}\cdot
  (p_\mu|U_\lambda\cap U_\mu)=p_\lambda|U_\lambda\cap U_\mu;
\]
además, si $\lambda$, $\mu$, $\nu$ son tres índices tales que
$U_\lambda\cap U_\mu\cap U_\nu\neq\varnothing$, las restricciones
$\gamma'_{\lambda\mu}$, $\gamma'_{\mu\nu}$, $\gamma'_{\lambda\nu}$ de
$\gamma_{\lambda\mu}$, $\gamma_{\mu\nu}$, $\gamma_{\lambda\nu}$ a
$U_\lambda\cap U_\mu\cap U_\nu$ satisfacen la condición
$\gamma'_{\lambda\nu}=\gamma'_{\lambda\mu}\gamma'_{\mu\nu}$; es decir,
$(\lambda,\mu)\mapsto\gamma_{\lambda\mu}$ es un \emph{1-cociclo} del
recubrimiento $(U_\lambda)$ con valores en $\mathscr{G}$. Si, para todo
$\lambda$, $p'_\lambda$ es un segundo elemento de
$\Gamma(U_\lambda,\mathscr{P})$, existe un único elemento
$\beta_\lambda\in\Gamma(U_\lambda,\mathscr{G})$ tal que
$p'_\lambda=\beta_\lambda\cdot p_\lambda$ y el 1-cociclo
$(\gamma'_{\lambda\mu})$ correspondiente a la familia $(p'_\lambda)$ viene dado
por
\[
  \gamma'_{\lambda\mu}=\beta_\lambda\gamma_{\lambda\mu}\beta_\mu^{-1},
\]
es decir, es \emph{cohomólogo} a $(\gamma_{\lambda\mu})$.
Recíprocamente, el dato de un 1-cociclo $(\gamma_{\lambda\mu})$ define, para
todo par $(\lambda,\mu)$, un automorfismo $\theta_{\lambda\mu}$ del
haz \emph{de conjuntos}
$\mathscr{G}|U_\lambda\cap U_\mu$, a saber, la traslación a la derecha por
$\gamma_{\lambda\mu}$, y el hecho de que se trate de un cociclo muestra que se
pueden \emph{pegar} los haces de conjuntos $\mathscr{G}|U_\lambda$ con
ayuda de los automorfismos $\theta_{\lambda\mu}$ (0, 3.3.1); se obtiene así
de manera evidente un torsor bajo $\mathscr{G}$, sea $\mathscr{P}$, y si
se toma como $p_\lambda$ la sección unidad sobre $U_\lambda$, el
1-cociclo correspondiente no es otro que el 1-cociclo dado
$(\gamma_{\lambda\mu})$; además, si se reemplaza
$(\gamma_{\lambda\mu})$ por un 1-cociclo
$\gamma'_{\lambda\mu}=\beta_\lambda\gamma_{\lambda\mu}\beta_\mu^{-1}$ que
le es cohomólogo, se comprueba inmediatamente que el torsor obtenido es isomorfo
a $\mathscr{P}$.

En particular, si $(\gamma_{\lambda\mu})$ es un \emph{1-coborde}, es decir,
de la forma
$\gamma_{\lambda\mu}=\beta_\lambda\beta_\mu^{-1}$, el torsor obtenido
$\mathscr{P}$ es \emph{isomorfo a $\mathscr{G}$} (considerado como torsor
bajo sí mismo para las traslaciones a la izquierda); en este caso se dice que
$\mathscr{P}$ es \emph{trivial}, y la recíproca es evidente.

Más particularmente, resulta de (III, 1.3.1) que se tiene:
\end{env}

\begin{proposition}[16.5.16]
\label{IV.16.5.16}
Sean $Z$ un esquema afín, $\mathscr{G}$ un $\mathcal{O}_Z$-módulo
cuasicoherente; entonces todo torsor bajo $\mathscr{G}$ es trivial.
\end{proposition}

\oldpage[IV-4]{34}
Volviendo al problema considerado en (16.5.13), se obtiene, por tanto:

\begin{proposition}[16.5.17]
\label{IV.16.5.17}
Sean $X$, $Y$ dos $S$-preesquemas, $Y_0$ un subpreesquema cerrado de $Y$
definido por un ideal cuasicoherente $\mathscr{J}$ de $\mathcal{O}_Y$ tal que
$\mathscr{J}^2=0$, y $j:Y_0\to Y$ la inyección canónica. Sea
$u_0:Y_0\to X$ un $S$-morfismo, y sea $\mathscr{P}$ el haz de conjuntos sobre
$Y$ tal que, para todo abierto $U$ de $Y$, $\Gamma(U,\mathscr{P})$ sea
el conjunto de los $S$-morfismos $u:U\to X$ tales que
$u_0|U_0=u\circ(j|U_0)$, donde $U_0=j^{-1}(U)$. Entonces existe sobre
$\mathscr{P}$ una estructura de pseudotorsor bajo el
$\mathcal{O}_{Y_0}$-módulo
\[
  \mathscr{G}=\mathscr{H}om_{\mathcal{O}_{Y_0}}
  (u_0^*(\Omega_{X/S}^1),\mathscr{J}).
\]
\end{proposition}

En particular:

\begin{corollary}[16.5.18]
\label{IV.16.5.18}
Con las notaciones de (16.5.16), supongamos $Y$ afín y $\Omega_{X/S}^1$ de
presentación finita; si existe un recubrimiento abierto $(U_\alpha)$ de $Y$
y, para cada índice $\alpha$, un $S$-morfismo $v_\alpha:U_\alpha\to X$
tal que, si $U_\alpha^0=j^{-1}(U_\alpha)$, se tenga
$v_\alpha\circ(j|U_\alpha^0)=u_0|U_\alpha^0$, entonces existe un
$S$-morfismo $u:Y\to X$ tal que $u\circ j=u_0$.
\end{corollary}

En efecto, $\mathscr{G}$ es entonces un $\mathcal{O}_{Y_0}$-módulo
cuasicoherente (I, 1.3.12); en virtud de (16.5.16) y del hecho de que $Y_0$ es
entonces afín, el haz $\mathscr{P}$, que por hipótesis es un torsor bajo
$\mathscr{G}$, y no solamente un pseudotorsor, es \emph{trivial}; pero si
$w$ es un isomorfismo de $\mathscr{G}$ en $\mathscr{P}$ (como
torsores bajo $\mathscr{G}$), la imagen por $w$ de la sección 0 de
$\mathscr{G}$ es el $S$-morfismo $u$ buscado.

\subsection{Haces de $p$-diferenciales y diferencial exterior.}
\label{IV.16.6}

\begin{env}[16.6.1]
\label{IV.16.6.1}
Sea $f:X\to S$ un morfismo de preesquemas. Se llama \emph{haz de
$p$-diferenciales de $X$ relativo a $S$} ($p$ entero) a la
\emph{$p$-ésima potencia exterior} (0, 4.1.5) del $\mathcal{O}_X$-módulo
$\Omega_{X/S}^1$, denotada por
\[
  \Omega_{X/S}^p=\bigwedge^p(\Omega_{X/S}^1).
  \tag{16.6.1.1}
  \label{IV.16.6.1.1}
\]
Se tiene, por tanto, $\Omega_{X/S}^0=\mathcal{O}_X$, y $\Omega_{X/S}^p=0$ para $p<0$;
los $\Omega_{X/S}^p$ son los componentes homogéneos del álgebra exterior
sobre $\Omega_{X/S}^1$
\[
  \Omega_{X/S}^{\bullet}=\bigwedge(\Omega_{X/S}^1)
  =\bigoplus_{p\in\mathbf{Z}}\bigwedge^p(\Omega_{X/S}^1),
  \tag{16.6.1.2}
  \label{IV.16.6.1.2}
\]
que es, por tanto, una $\mathcal{O}_X$-álgebra graduada cuasicoherente,
anticonmutativa y cuyos elementos de grado 1 son de cuadrado nulo. Para todo
abierto afín $U$ de $X$, se tiene
\[
  \Gamma(U,\Omega_{X/S}^{\bullet})
  =\bigwedge(\Gamma(U,\Omega_{X/S}^1)),
\]
donde $\Gamma(U,\Omega_{X/S}^1)$ se considera como
$\Gamma(U,\mathcal{O}_X)$-módulo.

Cuando $S=\operatorname{Spec}(A)$ y $X=\operatorname{Spec}(B)$ son
afines, siendo $B$ una $A$-álgebra, se tiene (0, 4.1.5)
$\Omega_{X/S}^p=(\Omega_{B/A}^p)^{\sim}$, poniendo
$\Omega_{B/A}^p=\bigwedge^p\Omega_{B/A}^1$.
\end{env}

\begin{theorem}[16.6.2]
\label{IV.16.6.2}
Existe un único endomorfismo $d$ del haz de grupos aditivos
$\Omega_{X/S}^{\bullet}$ que tiene las propiedades siguientes:
\begin{enumerate}
\item[(i)] $d\circ d=0$.
\item[(ii)] Para todo abierto $U$ de $X$ y toda sección
$f\in\Gamma(U,\mathcal{O}_X)$, se tiene $df=d_{X/S}f$.
\oldpage[IV-4]{35}
\item[(iii)] Para todo abierto $U$ de $X$, todo par de enteros $p$, $q$ y
todo par de secciones
$\omega'_p\in\Gamma(U,\Omega_{X/S}^p)$,
$\omega''_q\in\Gamma(U,\Omega_{X/S}^q)$, se tiene
\[
  d(\omega'_p\wedge\omega''_q)
  =(d\omega'_p)\wedge\omega''_q
  +(-1)^p\omega'_p\wedge d\omega''_q.
  \tag{16.6.2.1}
  \label{IV.16.6.2.1}
\]
\end{enumerate}
Además, $d$ es un endomorfismo de $\psi^*(\mathcal{O}_X)$-módulos
graduados, de grado $+1$.
\end{theorem}

Supongamos probada la existencia del endomorfismo $d$. Para todo abierto afín
$U$ de $X$, toda sección de $\Omega_{X/S}^p$ sobre $U$ es (en virtud
de (i)) combinación lineal de un número finito de elementos de la forma
$g(df_1\wedge df_2\wedge\cdots\wedge df_p)$, donde $g$ y los $f_i$ son
secciones de $\mathcal{O}_X$ sobre $U$ (0, 20.4.7). Las condiciones (i)
y (iii) muestran entonces, por inducción sobre $p$, que necesariamente se tiene
\[
  d(g(df_1\wedge df_2\wedge\cdots\wedge df_p))
  =dg\wedge df_1\wedge df_2\wedge\cdots\wedge df_p.
  \tag{16.6.2.2}
  \label{IV.16.6.2.2}
\]
Esto prueba, por tanto, la \emph{unicidad} de $d$ y la última afirmación del teorema.
En virtud de esta propiedad de unicidad, para demostrar la existencia de $d$, se
puede reducir al caso en que $S=\operatorname{Spec}(A)$ y
$X=\operatorname{Spec}(B)$ son afines. Ahora bien (Bourbaki, \emph{Alg.}, cap.
III, 3\textsuperscript{a} ed., \S~10), para definir una $A$-antiderivación
$D$ de grado $+1$ de un álgebra exterior $\bigwedge(M)$ (donde $M$ es un
$B$-módulo y $B$ una $A$-álgebra), con valores en una
$A$-álgebra \emph{graduada anticonmutativa}
$C=\bigoplus_{n=0}^{\infty}C_n$, cuyos elementos de grado 1 son de cuadrado
nulo, basta elegir arbitrariamente una $A$-derivación $D_0$ de $B$
en $C_1$ y un $A$-homomorfismo $D_1$ de $M$ en $C_2$; existe entonces
una única $A$-antiderivación $D$ de $\bigwedge(M)$ en $C$
que coincide con $D_0$ en $B$ y con $D_1$ en $M$.

En el caso actual, $D_0$ es necesariamente igual a $d_{B/A}$ en virtud de
(ii); todo se reduce a ver, teniendo en cuenta (16.6.2.2), que existe un
$A$-homomorfismo $u$ de $\Omega_{B/A}^1$ en $\Omega_{B/A}^2$ tal que
se tenga
\[
  u(g.df)=dg\wedge df
  \tag{16.6.2.3}
  \label{IV.16.6.2.3}
\]
para cualesquiera $f$, $g$ en $A$; para ello bastará mostrar que
existe un $A$-homomorfismo
$v:B\otimes_A\Omega_{B/A}^1\to\Omega_{B/A}^2$ tal que
\[
  v(g.\omega)=dg\wedge\omega
  \tag{16.6.2.4}
  \label{IV.16.6.2.4}
\]
para $g\in B$ y $\omega\in\Omega_{B/A}^1$. Finalmente, como
$\Omega_{B/A}^1=\mathfrak{I}/\mathfrak{I}^2$ (donde
$\mathfrak{I}=\mathfrak{I}_{B/A}$ es el núcleo del homomorfismo canónico
$B\otimes_A B\to B$) y $\Omega_{B/A}^1$ está engendrado por los elementos
de la forma $g.df$, basta definir un $A$-homomorfismo
$w:B\otimes_A(B\otimes_A B)\to\Omega_{B/A}^2$ tal que
\[
  w(g'\otimes g\otimes f)=dg'\wedge(g.df)
  \tag{16.6.2.5}
  \label{IV.16.6.2.5}
\]
y tal que $w$ sea nulo sobre la imagen de $B\otimes_A\mathfrak{I}^2$. Ahora bien,
como el segundo miembro de (16.6.2.5) es $A$-trilineal en $g'$, $g$, $f$,
la existencia de $w$ que satisface (16.6.2.5) es inmediata. Como, por otra parte,
$\mathfrak{I}$ está engendrado por los elementos $1\otimes x-x\otimes1$
($x\in B$), se reduce a comprobar que, cuando
$z=(1\otimes x-x\otimes1)(1\otimes y-y\otimes1)$, se tiene
$w(g'\otimes z)=0$. Ahora bien, como
$z=1\otimes(xy)+(xy)\otimes1-x\otimes y-y\otimes x$, la fórmula (16.6.2.4)
muestra que
\oldpage[IV-4]{36}
basta ver que se tiene $d(xy)-x.dy-y.dx=0$, lo que expresa que $d$ es una
derivación.

Queda por probar que $d$ satisface la condición (i). Ahora bien, el cuadrado de una
antiderivación es una derivación (Bourbaki, \emph{loc. cit.}), y como
$\Omega_{B/A}^{\bullet}$ está engendrado por $\Omega_{B/A}^1$ como
$B$-álgebra, basta comprobar que $d(dz)=0$ para $z\in B$ y
$z\in\Omega_{B/A}^1$; en el primer caso, esto resulta de la fórmula
(16.6.2.3) con $g=1$; en el segundo, se puede reducir al caso en que $z=g.df$
con $f$, $g$ en $B$, y entonces, en virtud de (16.6.2.1) y (16.6.2.3), se tiene
\[
  d(d(g.df))=d(dg\wedge df)
  =(d(dg))\wedge(df)-(dg)\wedge(d(df))=0.
\]
\hfill Q.E.D.

\begin{definition}[16.6.3]
\label{IV.16.6.3}
La antiderivación $d$ definida en (16.6.2) (denotada también por $d_{X/S}$) se
llama \emph{diferencial exterior sobre $X$} (relativa a $S$).
\end{definition}

\begin{proposition}[16.6.4]
\label{IV.16.6.4}
Para todo cambio de base $g:S'\to S$, si se pone
$X'=X\times_S S'$, el homomorfismo canónico
\[
  \Omega_{X/S}^{\bullet}\otimes_S S'
  \longrightarrow\Omega_{X'/S'}^{\bullet}
  \tag{16.6.4.1}
  \label{IV.16.6.4.1}
\]
deducido del isomorfismo (16.5.9.1) es biyectivo. Además, si $s$ es una
sección de $\Omega_{X/S}^{\bullet}$ sobre un abierto $U$ de $X$,
$s\otimes1$ su imagen inversa, sección de $\Omega_{X'/S'}^{\bullet}$
sobre la imagen inversa $U'$ de $U$ en $X'$, se tiene
$d_{X'/S'}(s\otimes1)=d_{X/S}(s)\otimes1$.
\end{proposition}

La primera afirmación es inmediata, pues la formación del álgebra exterior
de un módulo conmuta con toda extensión del anillo de escalares. Para
probar la segunda, se puede, en virtud de (16.6.2.2), reducir al caso en que
$s\in\Gamma(U,\mathcal{O}_X)$, y en este caso la afirmación ya fue probada
(16.4.3.7).

\begin{env}[16.6.5]
\label{IV.16.6.5}
Supongamos que $\Omega_{X/S}^1$ es un $\mathcal{O}_X$-módulo localmente libre
de rango $n$ en un punto $x$, de modo que hay $n$ secciones
$s_i\in\Gamma(U,\mathcal{O}_X)$ tales que los $ds_i$ forman una base del
$\Gamma(U,\mathcal{O}_X)$-módulo $\Gamma(U,\Omega_{X/S}^1)$ (16.5.8).
Entonces, para todo entero $p\geq1$, los $p$-diferenciales
$ds_{i_1}\wedge ds_{i_2}\wedge\cdots\wedge ds_{i_p}$ (para
$i_1<i_2<\cdots<i_p$, elementos de $[1,n]$) forman una base de
$\binom{n}{p}$ elementos de $\Gamma(U,\Omega_{X/S}^p)$ sobre
$\Gamma(U,\mathcal{O}_X)$. Además, la fórmula (16.6.2.2) muestra que, para
toda sección $g\in\Gamma(U,\mathcal{O}_X)$, se tiene
\[
  d(g.ds_{i_1}\wedge ds_{i_2}\wedge\cdots\wedge ds_{i_p})
  =\sum_k(-1)^r\frac{\partial g}{\partial s_k}
  ds_{i_1}\wedge\cdots\wedge ds_{i_r}\wedge ds_k\wedge
  ds_{i_{r+1}}\wedge\cdots\wedge ds_{i_p}
  \tag{16.6.5.1}
  \label{IV.16.6.5.1}
\]
donde, en el segundo miembro, $k$ recorre el conjunto de los $n-p$ índices distintos de los
$i_h$, siendo $i_r$ el mayor índice $<k$.

Obsérvese que la relación $d(dg)=0$ para toda sección
$g\in\Gamma(U,\mathcal{O}_X)$ se expresa en la forma
\[
  D_i(D_jg)=D_j(D_ig)\qquad\text{para $i\neq j$};
\]
es decir, las derivaciones $D_i$ definidas en (16.5.7) conmutan dos a dos.
\end{env}

\subsection{Los $\mathscr{P}_{X/S}^n(\mathscr{F})$.}
\label{IV.16.7}

\begin{env}[16.7.1]
\label{IV.16.7.1}
Sean $f:X\to S$ un morfismo de preesquemas, $\mathscr{F}$ un
$\mathcal{O}_X$-módulo. Denotemos por $X_{\Delta_f}^{(n)}$ la \emph{$n$-ésima
vecindad infinitesimal de $X$} para el morfismo diagonal
\oldpage[IV-4]{37}
$\Delta_f:X\to X\times_S X$, por
$h_n:X_{\Delta_f}^{(n)}\to X\times_S X$ el morfismo canónico (16.1.2), y
consideremos los dos morfismos compuestos
\[
  p_1^{(n)}:X_{\Delta_f}^{(n)}\xrightarrow{h_n}X\times_S X
  \xrightarrow{p_1}X,\qquad
  p_2^{(n)}:X_{\Delta_f}^{(n)}\xrightarrow{h_n}X\times_S X
  \xrightarrow{p_2}X
\]
de modo que, por definición, $p_1^{(n)}$ corresponde al homomorfismo de
haces de anillos $\mathcal{O}_X\to\mathscr{P}_{X/S}^n$ que hemos
elegido para definir la estructura de $\mathcal{O}_X$-álgebra sobre
$\mathscr{P}_{X/S}^n$ (16.3.5), y $p_2^{(n)}$ al homomorfismo de haces
de anillos $d_{X/S}^n:\mathcal{O}_X\to\mathscr{P}_{X/S}^n$ (16.3.6). Como
$X_{\Delta_f}^{(n)}$ y $X$ tienen el mismo espacio subyacente, se puede escribir
\[
  \mathscr{P}_{X/S}^n
  =(p_1^{(n)})_*((p_2^{(n)})^*(\mathcal{O}_X)).
  \tag{16.7.1.1}
  \label{IV.16.7.1.1}
\]
Más generalmente, pondremos
\[
  \mathscr{P}_{X/S}^n(\mathscr{F})
  =(p_1^{(n)})_*((p_2^{(n)})^*(\mathscr{F}))
  \tag{16.7.1.2}
  \label{IV.16.7.1.2}
\]
de modo que
$\mathscr{P}_{X/S}^n=\mathscr{P}_{X/S}^n(\mathcal{O}_X)$; por definición,
$\mathscr{P}_{X/S}^n(\mathscr{F})$ es, por tanto, un $\mathcal{O}_X$-módulo.
\end{env}

\begin{env}[16.7.2]
\label{IV.16.7.2}
Si se vuelve a las definiciones de las imágenes inversas de módulos sobre los
espacios anillados (0, 4.3.1) y se tiene en cuenta que
$X_{\Delta_f}^{(n)}$ y $X$ tienen el mismo espacio subyacente, se ve que se puede
escribir también la definición (16.7.1.2) en la forma
\[
  \mathscr{P}_{X/S}^n(\mathscr{F})
  =\mathscr{P}_{X/S}^n\otimes_{\mathcal{O}_X}\mathscr{F}
  \tag{16.7.2.1}
  \label{IV.16.7.2.1}
\]
pero debe tenerse cuidado de que, al interpretar el signo $\otimes$,
$\mathscr{P}_{X/S}^n$ está dotado de su estructura de $\mathcal{O}_X$-módulo
definida por el homomorfismo de haces de anillos
$d_{X/S}^n:\mathcal{O}_X\to\mathscr{P}_{X/S}^n$. Resulta inmediatamente de esta
fórmula (o directamente de (16.7.1.2)) que
$\mathscr{P}_{X/S}^n(\mathscr{F})$ está canónicamente dotado de una estructura de
$\mathscr{P}_{X/S}^n$-módulo.
\end{env}

\begin{proposition}[16.7.3]
\label{IV.16.7.3}
\begin{enumerate}
\item[(i)] El funtor
$\mathscr{F}\mapsto\mathscr{P}_{X/S}^n(\mathscr{F})$ de la categoría de los
$\mathcal{O}_X$-módulos en la de los $\mathscr{P}_{X/S}^n$-módulos es
exacto por la derecha y conmuta con los límites inductivos cualesquiera; es exacto
cuando $\mathscr{P}_{X/S}^n$ es un $\mathcal{O}_X$-módulo plano.
\item[(ii)] Si $\mathscr{F}$ es un $\mathcal{O}_X$-módulo cuasicoherente
(resp. de tipo finito, resp. de presentación finita),
$\mathscr{P}_{X/S}^n(\mathscr{F})$ es un $\mathscr{P}_{X/S}^n$-módulo
cuasicoherente (resp. de tipo finito, resp. de presentación finita).
\end{enumerate}
\end{proposition}

Las afirmaciones de (i) resultan inmediatamente de la fórmula (16.7.2.1) y de la
consideración de la simetría de $\mathscr{P}_{X/S}^n$ (16.3.4). Las afirmaciones
de (ii) se deducen de la exactitud por la derecha del funtor
$\mathscr{F}\mapsto\mathscr{P}_{X/S}^n(\mathscr{F})$.

\begin{env}[16.7.4]
\label{IV.16.7.4}
Las dos estructuras de $\mathcal{O}_X$-módulo de $\mathscr{P}_{X/S}^n$
definen sobre $\mathscr{P}_{X/S}^n(\mathscr{F})$ dos estructuras de
$\mathcal{O}_X$-módulo que, por lo demás, conmutan y, por tanto, una estructura de
$\mathcal{O}_X$-bimódulo. Conviene llamar \emph{izquierda} a la estructura
proveniente del homomorfismo estructural
$\mathcal{O}_X\to\mathscr{P}_{X/S}^n$ (elegido en (16.3.5)) y
\emph{derecha} a la proveniente del homomorfismo
$d_{X/S}^n:\mathcal{O}_X\to\mathscr{P}_{X/S}^n$. Es decir, para todo
abierto $U$ de $X$ y toda terna de elementos
$a\in\Gamma(U,\mathcal{O}_X)$, $b\in\Gamma(U,\mathscr{P}_{X/S}^n)$,
$t\in\Gamma(U,\mathscr{F})$, se tiene por definición
\[
  a(b\otimes t)=(ab)\otimes t,\qquad
  (b\otimes t)a=(b.d^na)\otimes t=b\otimes(at)
  =(d^na).(b\otimes t).
  \tag{16.7.4.1}
  \label{IV.16.7.4.1}
\]
La estructura de $\mathcal{O}_X$-módulo proveniente de la definición (16.7.1.2)
es, por tanto, con estas convenciones, la estructura de $\mathcal{O}_X$-módulo a la
izquierda.

\oldpage[IV-4]{38}
Si $\mathscr{F}$ es un $\mathcal{O}_X$-módulo cuasicoherente, lo mismo ocurre
con $\mathscr{P}_{X/S}^n(\mathscr{F})$ para cualquiera de sus
estructuras de $\mathcal{O}_X$-módulo. Si además $\mathscr{F}$ es de tipo
finito (resp. de presentación finita) y $f:X\to S$ es localmente de tipo finito
(resp. localmente de presentación finita),
$\mathscr{P}_{X/S}^n(\mathscr{F})$ es (para cualquiera de sus
estructuras de $\mathcal{O}_X$-módulo) de tipo finito (resp. de presentación
finita), como resulta de (16.3.9) y de (16.4.22).
\end{env}

\begin{env}[16.7.5]
\label{IV.16.7.5}
La definición (16.7.2.1) implica la existencia de un homomorfismo de haces
de grupos conmutativos
\[
  d_{X/S,\mathscr{F}}^n:\mathscr{F}
  \longrightarrow\mathscr{P}_{X/S}^n(\mathscr{F})
  \qquad\text{(también denotado por $d_{X/S}^n$)}
  \tag{16.7.5.1}
  \label{IV.16.7.5.1}
\]
tal que, con las notaciones de (16.7.4), se tenga
\[
  d_{X/S,\mathscr{F}}^n(t)=1\otimes t
  \tag{16.7.5.2}
  \label{IV.16.7.5.2}
\]
y, por consiguiente, en virtud de (16.7.4.1),
\[
  d_{X/S,\mathscr{F}}^n(at)
  =(1\otimes t)a=(d_{X/S,\mathscr{F}}^n(t)).a
  \tag{16.7.5.3}
  \label{IV.16.7.5.3}
\]
\[
  d_{X/S,\mathscr{F}}^n(at)
  =(d_{X/S}^n(a)).(1\otimes t)
  =(d_{X/S}^n(a))(d_{X/S,\mathscr{F}}^n(t)).
  \tag{16.7.5.4}
  \label{IV.16.7.5.4}
\]
Es, por tanto, $\mathcal{O}_X$-lineal para la estructura de
$\mathcal{O}_X$-módulo a la derecha sobre $\mathscr{P}_{X/S}^n(\mathscr{F})$, y
\emph{semilineal} (respecto del automorfismo $\sigma$ (16.3.4)) para
la estructura de $\mathcal{O}_X$-módulo a la izquierda.
\end{env}

\begin{proposition}[16.7.6]
\label{IV.16.7.6}
El $\mathcal{O}_X$-módulo a la izquierda $\mathscr{P}_{X/S}^n(\mathscr{F})$ está
engendrado por la imagen de $\mathscr{F}$ por el homomorfismo canónico
$d_{X/S,\mathscr{F}}^n$.
\end{proposition}

Esto resulta inmediatamente de (16.7.5.3) y del caso particular correspondiente a
$\mathscr{F}=\mathcal{O}_X$ (16.3.8).

\begin{env}[16.7.7]
\label{IV.16.7.7}
Los homomorfismos canónicos de haces de anillos
\[
  \varphi_{nm}:\mathscr{P}_{X/S}^m\longrightarrow\mathscr{P}_{X/S}^n
\]
para $n\leq m$ (16.1.2) definen, en virtud de (16.7.2.1),
homomorfismos canónicos
\[
  \mathscr{P}_{X/S}^m(\mathscr{F})
  \longrightarrow\mathscr{P}_{X/S}^n(\mathscr{F})\qquad(n\leq m)
\]
que son homomorfismos de $\mathcal{O}_X$-bimódulos en virtud de (16.1.6)
y (16.7.4.1); además, se tienen diagramas conmutativos
\[
  \xymatrix{
    \mathscr{P}_{X/S}^m(\mathscr{F})\ar[rr]
      &&\mathscr{P}_{X/S}^n(\mathscr{F})\\
    &\mathscr{F}\ar[ul]^-{d_{X/S,\mathscr{F}}^m}
      \ar[ur]_-{d_{X/S,\mathscr{F}}^n}
  }
\]
Se tiene así un sistema proyectivo de $\mathcal{O}_X$-bimódulos
$(\mathscr{P}_{X/S}^n(\mathscr{F}))$, y se pone
\[
  \mathscr{P}_{X/S}^{\infty}(\mathscr{F})
  =\varprojlim\mathscr{P}_{X/S}^n(\mathscr{F}).
  \tag{16.7.7.1}
  \label{IV.16.7.7.1}
\]
Además, lo que precede muestra que los homomorfismos (16.7.5.1) forman un
sistema proyectivo de homomorfismos y definen, por tanto, un homomorfismo
canónico
\[
  d_{X/S,\mathscr{F}}^{\infty}:\mathscr{F}
  \longrightarrow\mathscr{P}_{X/S}^{\infty}(\mathscr{F}).
  \tag{16.7.7.2}
  \label{IV.16.7.7.2}
\]
\end{env}

\oldpage[IV-4]{39}

\begin{env}[16.7.8]
\label{IV.16.7.8}
Sean $\mathscr{F}$, $\mathscr{G}$ dos $\mathcal{O}_X$-módulos; resulta
inmediatamente de la definición (16.7.2.1) que se tiene un isomorfismo canónico de
$\mathscr{P}_{X/S}^n$-módulos
\[
  \mathscr{P}_{X/S}^n(\mathscr{F}\otimes_{\mathcal{O}_X}\mathscr{G})
  \xrightarrow{\sim}
  \mathscr{P}_{X/S}^n(\mathscr{F})
  \otimes_{\mathscr{P}_{X/S}^n}\mathscr{P}_{X/S}^n(\mathscr{G})
  \tag{16.7.8.1}
  \label{IV.16.7.8.1}
\]
(Bourbaki, \emph{Alg.}, cap.~II, 3\textsuperscript{a} ed., §~5,
n\textsuperscript{o}~1, prop.~3).

Se concluye en particular (o se ve directamente en la definición
(16.7.2.1)) que, si $\mathscr{F}$ está dotado de una estructura de
$\mathcal{O}_X$-álgebra (no necesariamente asociativa),
$\mathscr{P}_{X/S}^n(\mathscr{F})$ está canónicamente dotado de una estructura de
$\mathscr{P}_{X/S}^n$-álgebra; esta última es asociativa (resp.
conmutativa, resp. unitaria, resp. un álgebra de Lie) cuando así lo es
$\mathscr{F}$. Además, los homomorfismos canónicos
$\mathscr{P}_{X/S}^m(\mathscr{F})\to\mathscr{P}_{X/S}^n(\mathscr{F})$ para
$n\leq m$ (16.7.7) son entonces dihomomorfismos de álgebras; asimismo,
(16.7.5.1) es entonces un homomorfismo de $\mathcal{O}_X$-álgebras cuando
$\mathscr{P}_{X/S}^n(\mathscr{F})$ está dotado de su estructura de
$\mathcal{O}_X$-álgebra proveniente de su estructura de $\mathcal{O}_X$-módulo
\emph{a la derecha}.

Con las mismas notaciones, se tiene también un homomorfismo canónico de
$\mathscr{P}_{X/S}^n$-módulos
\[
  \mathscr{P}_{X/S}^n
  (\mathscr{H}\!om_{\mathcal{O}_X}(\mathscr{F},\mathscr{G}))
  \longrightarrow
  \mathscr{H}\!om_{\mathscr{P}_{X/S}^n}
  (\mathscr{P}_{X/S}^n(\mathscr{F}),\mathscr{P}_{X/S}^n(\mathscr{G}))
  \tag{16.7.8.2}
  \label{IV.16.7.8.2}
\]
(Bourbaki, \emph{Alg.}, 3\textsuperscript{a} ed., §~5,
n\textsuperscript{o}~3), que es biyectivo cuando $\mathscr{P}_{X/S}^n$ es
un $\mathcal{O}_X$-módulo \emph{localmente libre de tipo finito}
(\emph{loc. cit.}, prop.~7).
\end{env}

\begin{env}[16.7.9]
\label{IV.16.7.9}
Supongamos que se está en la situación descrita en (16.4.1); entonces, del
homomorfismo canónico $P^n(u)$ (16.4.3.3), se deduce inmediatamente un
homomorfismo canónico de $\mathcal{O}_{X'}$-bimódulos
\[
  u^*(\mathscr{P}_{X/S}^n(\mathscr{F}))
  \longrightarrow
  \mathscr{P}_{X'/S'}^n(u^*(\mathscr{F})).
  \tag{16.7.9.1}
  \label{IV.16.7.9.1}
\]
Dejamos al lector la tarea de extender a este homomorfismo las propiedades
vistas en (16.4) para el caso $\mathscr{F}=\mathcal{O}_X$.
\end{env}

\begin{remark}[16.7.10]
\label{IV.16.7.10}
La definición de $\mathscr{P}_{X/S}^n(\mathscr{F})$ en la forma (16.7.1.2)
sigue teniendo sentido cuando $\mathscr{F}$ es un haz de conjuntos cualquiera
(estando definida en (0\textsubscript{I}, 3.7.1) la imagen inversa de un haz
de conjuntos por $p_2^{(n)}$); una variante de esta definición permite
definir el «esquema de jets» (relativo a $S$) de un $X$-preesquema
cualquiera.
\end{remark}

\subsection{Operadores diferenciales.}
\footnote{Para un formalismo más general, véase la Exposición VII de [42]
(debida a P. Gabriel).}
\label{IV.16.8}

\begin{definition}[16.8.1]
\label{IV.16.8.1}
Sean $f=(\psi,\theta):X\to S$ un morfismo de preesquemas,
$\mathscr{F}$, $\mathscr{G}$ dos $\mathcal{O}_X$-módulos, y $n$ un entero
$\geq0$. Se dice que un homomorfismo de haces de grupos aditivos
$D:\mathscr{F}\to\mathscr{G}$ es un \emph{operador diferencial de orden
$\leq n$} (relativamente a $S$) si existe un homomorfismo de
$\mathcal{O}_X$-módulos
$u:\mathscr{P}_{X/S}^n(\mathscr{F})\to\mathscr{G}$ (donde
$\mathscr{P}_{X/S}^n(\mathscr{F})$ está provisto de su estructura de
$\mathcal{O}_X$-módulo a la izquierda (16.7.4)) tal que
$D=u\circ d_{X/S,\mathscr{F}}^n$.
\end{definition}

Es claro, en virtud de la existencia de los homomorfismos canónicos
\[
  \mathscr{P}_{X/S}^m(\mathscr{F})
  \longrightarrow\mathscr{P}_{X/S}^n(\mathscr{F})
\]
\oldpage[IV-4]{40}
para $n\leq m$ (16.7.7), que un operador diferencial de orden $\leq n$ es
también un operador diferencial de orden $\leq m$ para todo $m\geq n$. Si
$D:\mathscr{F}\to\mathscr{G}$ es un operador diferencial de orden
$\leq n$, entonces, para todo abierto $U$ de $X$,
$D|U:\mathscr{F}|U\to\mathscr{G}|U$ es también un operador diferencial
de orden $\leq n$.

Se dice que un homomorfismo $D:\mathscr{F}\to\mathscr{G}$ de los haces de
grupos aditivos subyacentes a $\mathscr{F}$ y $\mathscr{G}$ es un
\emph{operador diferencial} (relativamente a $S$) si, para todo $x\in X$,
existen una vecindad abierta $U$ de $x$ y un entero $n\geq0$ tales que
$D|U:\mathscr{F}|U\to\mathscr{G}|U$ sea un operador diferencial de orden
$\leq n$. El \emph{orden} de un operador diferencial
$D:\mathscr{F}\to\mathscr{G}$ es el ínfimo de los enteros $n$ tales que
$D$ sea un operador diferencial de orden $\leq n$ (y, por tanto, $+\infty$
si no hay tales enteros); este orden es siempre finito si $X$ es
\emph{cuasicompacto}. Los operadores diferenciales de orden $0$ no son sino
los homomorfismos de $\mathcal{O}_X$-módulos
$\mathscr{F}\to\mathscr{G}$; se conviene en que todo operador diferencial
de orden $<0$ es \emph{nulo}. Para $n\geq0$, un operador diferencial no es,
en general, un homomorfismo de $\mathcal{O}_X$-módulos, pero es siempre
un homomorfismo de $\psi^*(\mathcal{O}_S)$-módulos.

Cuando $\mathscr{F}=\mathcal{O}_X$, un operador diferencial de orden
$\leq1$ de $\mathcal{O}_X$ en $\mathscr{G}$ se escribe de manera única
en la forma $v+D$, donde $v:\mathcal{O}_X\to\mathscr{G}$ es un
$\mathcal{O}_X$-homomorfismo, y $D$ una $S$-\emph{derivación} (16.5.1) de
$\mathcal{O}_X$ en $\mathscr{G}$: esto resulta de la estructura de
$P_{B/A}^1$ (0, 20.4.8).

\begin{env}[16.8.2]
\label{IV.16.8.2}
Para describir de manera más explícita un operador diferencial de orden
$\leq n$, $D:\mathscr{F}\to\mathscr{G}$, basta, para todo abierto afín
$U$ de $X$ cuya imagen en $S$ esté contenida en un abierto afín $V$,
caracterizar el homomorfismo
$D=D_U:\Gamma(U,\mathscr{F})\to\Gamma(U,\mathscr{G})$. Si se pone
$\Gamma(V,\mathcal{O}_S)=A$, $\Gamma(U,\mathcal{O}_X)=B$, de modo que $B$
es una $A$-álgebra, se tiene
\[
  \Gamma(U,\mathscr{P}_{X/S}^n)
  =(B\otimes_A B)/\mathfrak{I}^{n+1},
\]
donde se ha puesto, para abreviar, $\mathfrak{I}=\mathfrak{I}_{B/A}$.
Pongamos además $M=\Gamma(U,\mathscr{F})$, $N=\Gamma(U,\mathscr{G})$;
entonces la definición de $D$ significa que, para cada par $(U,V)$ que
satisface las condiciones precedentes, el $A$-homomorfismo $D:M\to N$ se
factoriza como
\[
  M\longrightarrow
  ((B\otimes_A B)/\mathfrak{I}^{n+1})\otimes_B M
  \xrightarrow{v}N
\]
donde la primera flecha es el homomorfismo canónico $t\mapsto1\otimes t$, y
$v$ es un $B$-homomorfismo; la estructura de $B$-módulo de
$((B\otimes_A B)/\mathfrak{I}^{n+1})\otimes_B M$ proviene del primer
factor $B$ (mientras que, recordemos, al formar el producto tensorial
sobre $B$, la estructura de $B$-módulo de
$(B\otimes_A B)/\mathfrak{I}^{n+1}$ proviene del segundo factor $B$).
Observemos ahora que el $B$-módulo
$((B\otimes_A B)/\mathfrak{I}^{n+1})\otimes_B M$ es isomorfo a
$(B\otimes_A M)/\mathfrak{I}^{n+1}(B\otimes_A M)$, donde $B\otimes_A M$ se
considera como $(B\otimes_A B)$-módulo y su estructura de $B$-módulo
proviene del homomorfismo $b\mapsto b\otimes1$ de $B$ en $B\otimes_A B$.
Sea entonces $D'$ el $B$-homomorfismo de $B\otimes_A M$ en $N$ tal que
$D'(b\otimes t)=bD(t)$; la condición de factorización de $D$ se expresa
también diciendo que $D'$ debe \emph{anularse sobre el $B$-submódulo}
$\mathfrak{I}^{n+1}(B\otimes_A M)$.
\end{env}

\begin{env}[16.8.3]
\label{IV.16.8.3}
Es claro que el conjunto de los operadores diferenciales de orden $\leq n$ de
$\mathscr{F}$ en $\mathscr{G}$ es un grupo aditivo, denotado por
$\operatorname{Diff}_{X/S}^n(\mathscr{F},\mathscr{G})$; cuando
$\mathscr{F}=\mathscr{G}=\mathcal{O}_X$, se escribe también
$\operatorname{Diff}_{X/S}^n$ en lugar de
$\operatorname{Diff}_{X/S}^n(\mathcal{O}_X,\mathcal{O}_X)$.

Se ha visto en (16.8.1) que, para dos abiertos $U\supset V$ de $X$, se tiene
un homomorfismo canónico de restricción
\[
  \operatorname{Diff}_{U/S}^n(\mathscr{F}|U,\mathscr{G}|U)
  \longrightarrow
  \operatorname{Diff}_{V/S}^n(\mathscr{F}|V,\mathscr{G}|V)
\]
\oldpage[IV-4]{41}
por lo que $U\mapsto\operatorname{Diff}_{U/S}^n(\mathscr{F}|U,\mathscr{G}|U)$
es un prehaz de grupos aditivos; de hecho, es incluso un \emph{haz}, pues,
cuando $U$ recorre los abiertos de $X$, los homomorfismos
$u\mapsto u\circ d_{U/S,\mathscr{F}|U}^n$ son \emph{isomorfismos} de grupos
aditivos
\[
  \mathscr{H}\!om_{\mathcal{O}_U}
  (\mathscr{P}_{U/S}^n(\mathscr{F}|U),\mathscr{G}|U)
  \xrightarrow{\sim}
  \operatorname{Diff}_{U/S}^n(\mathscr{F}|U,\mathscr{G}|U),
  \tag{16.8.3.1}
  \label{IV.16.8.3.1}
\]
en virtud de que la imagen de $\mathscr{F}$ por
$d_{X/S,\mathscr{F}}^n$ engendra $\mathscr{P}_{X/S}^n(\mathscr{F})$
(16.7.6). Se denota este haz por
$\mathscr{D}\!iff_{X/S}^n(\mathscr{F},\mathscr{G})$, y se tiene, por tanto:
\end{env}

\begin{proposition}[16.8.4]
\label{IV.16.8.4}
Los isomorfismos (16.8.3.1) definen un isomorfismo de haces de grupos
aditivos
\[
  \mathscr{H}\!om_{\mathcal{O}_X}
  (\mathscr{P}_{X/S}^n(\mathscr{F}),\mathscr{G})
  \xrightarrow{\sim}
  \mathscr{D}\!iff_{X/S}^n(\mathscr{F},\mathscr{G}).
  \tag{16.8.4.1}
  \label{IV.16.8.4.1}
\]
\end{proposition}

Cuando $\mathscr{F}=\mathscr{G}=\mathcal{O}_X$, se escribe también
$\mathscr{D}\!iff_{X/S}^n$ en lugar de
$\mathscr{D}\!iff_{X/S}^n(\mathcal{O}_X,\mathcal{O}_X)$; de (16.8.4) resulta
que $\mathscr{D}\!iff_{X/S}^n$ se identifica canónicamente con el
\emph{dual} del $\mathcal{O}_X$-módulo $\mathscr{P}_{X/S}^n$; por ello se
escribe $\langle t,D\rangle$ en lugar de $u(t)$ si $t$ es una sección de
$\mathscr{P}_{X/S}^n$ sobre un abierto y si $u$ es el homomorfismo de
$\mathscr{P}_{X/S}^n$ en $\mathcal{O}_X$ correspondiente a $D$.

\begin{env}[16.8.5]
\label{IV.16.8.5}
Puesto que $\mathscr{P}_{X/S}^n(\mathscr{F})$ está provisto de una estructura
de $\mathcal{O}_X$-bimódulo (16.7.4), se deduce canónicamente una estructura
de $\mathcal{O}_X$-bimódulo sobre
$\mathscr{H}\!om_{\mathcal{O}_X}
(\mathscr{P}_{X/S}^n(\mathscr{F}),\mathscr{G})$, y, por tanto, también sobre
$\mathscr{D}\!iff_{X/S}^n(\mathscr{F},\mathscr{G})$ en virtud de (16.8.4.1).
Más precisamente, a la estructura de $\mathcal{O}_X$-módulo a la izquierda
de $\mathscr{P}_{X/S}^n(\mathscr{F})$ corresponde, en virtud de la definición
(16.8.1), la estructura de $\mathcal{O}_X$-módulo a la izquierda sobre
$\mathscr{D}\!iff_{X/S}^n(\mathscr{F},\mathscr{G})$ que se explicita como
sigue: para todo abierto $U$ de $X$, toda sección
$a\in\Gamma(U,\mathcal{O}_X)$ y todo operador diferencial
$D:\mathscr{F}|U\to\mathscr{G}|U$, $aD$ es el operador diferencial que a
toda sección $t\in\Gamma(U,\mathscr{F})$ hace corresponder la sección
\[
  (aD)(t)=a(D(t))
  \tag{16.8.5.1}
  \label{IV.16.8.5.1}
\]
de $\Gamma(U,\mathscr{G})$. Análogamente, a la estructura de
$\mathcal{O}_X$-módulo a la derecha de $\mathscr{P}_{X/S}^n(\mathscr{F})$
corresponde la estructura de $\mathcal{O}_X$-módulo a la derecha sobre
$\mathscr{D}\!iff_{X/S}^n(\mathscr{F},\mathscr{G})$ que se explicita como
sigue: con las mismas notaciones que antes, $Da$ es el operador diferencial
que a $t\in\Gamma(U,\mathscr{F})$ hace corresponder la sección
\[
  (Da)(t)=D(at).
  \tag{16.8.5.2}
  \label{IV.16.8.5.2}
\]
\end{env}

\begin{proposition}[16.8.6]
\label{IV.16.8.6}
Si $f:X\to S$ es un morfismo localmente de presentación finita,
$\mathscr{F}$ un $\mathcal{O}_X$-módulo cuasicoherente de presentación finita
y $\mathscr{G}$ un $\mathcal{O}_X$-módulo cuasicoherente, entonces
$\mathscr{D}\!iff_{X/S}^n(\mathscr{F},\mathscr{G})$ es un
$\mathcal{O}_X$-módulo cuasicoherente para una u otra de las estructuras
definidas en (16.8.5).
\end{proposition}

La proposición resulta de que, bajo las hipótesis formuladas,
$\mathscr{P}_{X/S}^n(\mathscr{F})$ es un $\mathcal{O}_X$-módulo
cuasicoherente de presentación finita (16.7.4), y de (I, 1.3.12).

\begin{env}[16.8.7]
\label{IV.16.8.7}
El conjunto de los operadores diferenciales de $\mathscr{F}$ en
$\mathscr{G}$ (de orden no especificado (16.8.1)) se denota por
$\operatorname{Diff}_{X/S}(\mathscr{F},\mathscr{G})$; se ve también, como en
(16.8.3), que
$U\mapsto\operatorname{Diff}_{U/S}(\mathscr{F}|U,\mathscr{G}|U)$ es un haz de
grupos aditivos, que se designará por
$\mathscr{D}\!iff_{X/S}(\mathscr{F},\mathscr{G})$. Es inmediato que
$\mathscr{D}\!iff_{X/S}(\mathscr{F},\mathscr{G})$ es la unión de la familia
filtrante creciente de sus subhaces
$\mathscr{D}\!iff_{X/S}^n(\mathscr{F},\mathscr{G})$; si $X$ es
cuasicompacto, $\operatorname{Diff}_{X/S}(\mathscr{F},\mathscr{G})$
\oldpage[IV-4]{42}
es asimismo la unión de sus subgrupos
$\operatorname{Diff}_{X/S}^n(\mathscr{F},\mathscr{G})$ (16.8.1). Las
estructuras de $\mathcal{O}_X$-bimódulo sobre los
$\mathscr{D}\!iff_{X/S}^n(\mathscr{F},\mathscr{G})$ definen, por tanto, una
estructura de $\mathcal{O}_X$-bimódulo sobre
$\mathscr{D}\!iff_{X/S}(\mathscr{F},\mathscr{G})$, explicitada también por
(16.8.5.1) y (16.8.5.2).

Obsérvese que, para $n\leq m$, se tiene un diagrama conmutativo
\[
  \xymatrix{
    \mathscr{H}\!om_{\mathcal{O}_X}
      (\mathscr{P}_{X/S}^n(\mathscr{F}),\mathscr{G})
      \ar[r]^-{\sim}\ar[d]
      &\mathscr{D}\!iff_{X/S}^n(\mathscr{F},\mathscr{G})\ar[d]\\
    \mathscr{H}\!om_{\mathcal{O}_X}
      (\mathscr{P}_{X/S}^m(\mathscr{F}),\mathscr{G})
      \ar[r]^-{\sim}
      &\mathscr{D}\!iff_{X/S}^m(\mathscr{F},\mathscr{G})
  }
  \tag{16.8.7.1}
  \label{IV.16.8.7.1}
\]
donde las flechas horizontales son los isomorfismos
(16.8.4.1)\footnote{El texto impreso remite a (16.8.4.2), que no existe; la
remisión correcta es a (16.8.4.1).} y la flecha vertical de la izquierda
proviene del homomorfismo canónico
$\mathscr{P}_{X/S}^m(\mathscr{F})\to
\mathscr{P}_{X/S}^n(\mathscr{F})$ (16.7.7). Para todo abierto $U$ de $X$,
provistamos entonces
\[
  \Gamma(U,\mathscr{P}_{X/S}^{\infty}(\mathscr{F}))
  =\varprojlim\Gamma(U,\mathscr{P}_{X/S}^n(\mathscr{F}))
\]
de la topología de límite proyectivo de las topologías discretas sobre los
$\Gamma(U,\mathscr{P}_{X/S}^n(\mathscr{F}))$, lo que define
$\Gamma(U,\mathscr{P}_{X/S}^{\infty}(\mathscr{F}))$ como un
$\Gamma(U,\mathcal{O}_X)$-bimódulo topológico, de modo que
$\mathscr{P}_{X/S}^{\infty}(\mathscr{F})$ aparece como un haz con valores en
la categoría de los grupos conmutativos topológicos
(0\textsubscript{I}, 3.2.6). Entonces, por (G, II, 1.11), el límite del sistema
inductivo de haces de grupos conmutativos
\[
  (\mathscr{H}\!om_{\mathcal{O}_X}
  (\mathscr{P}_{X/S}^n(\mathscr{F}),\mathscr{G}))
\]
no es sino el haz de gérmenes de homomorfismos \emph{continuos} de
$\mathscr{P}_{X/S}^{\infty}(\mathscr{F})$ en $\mathscr{G}$ (este último
provisto de la topología discreta): en efecto, los homomorfismos continuos de
$\Gamma(U,\mathscr{P}_{X/S}^{\infty}(\mathscr{F}))$ en el grupo discreto
$\Gamma(U,\mathscr{G})$ corresponden biunívocamente a los sistemas inductivos
de homomorfismos de grupos
$\Gamma(U,\mathscr{P}_{X/S}^n(\mathscr{F}))\to\Gamma(U,\mathscr{G})$. Se
puede, por tanto, expresar también (16.8.4) diciendo que se tiene un
isomorfismo canónico
\[
  \mathscr{H}\!om.\mathrm{cont}_{\mathcal{O}_X}
  (\mathscr{P}_{X/S}^{\infty}(\mathscr{F}),\mathscr{G})
  \xrightarrow{\sim}
  \mathscr{D}\!iff_{X/S}(\mathscr{F},\mathscr{G})
\]
donde el primer miembro designa el haz de gérmenes de homomorfismos continuos
de $\mathscr{P}_{X/S}^{\infty}(\mathscr{F})$ en $\mathscr{G}$.
\end{env}

\begin{proposition}[16.8.8]
\label{IV.16.8.8}
Sean $\mathscr{F}$, $\mathscr{G}$ dos $\mathcal{O}_X$-módulos,
$D:\mathscr{F}\to\mathscr{G}$ un homomorfismo de
$\psi^*(\mathcal{O}_S)$-módulos, y $n$ un entero $\geq0$. Las condiciones
siguientes son equivalentes:
\begin{enumerate}
  \item[(a)] $D$ es un operador diferencial de orden $\leq n$.
  \item[(b)] Para toda sección $a$ de $\mathcal{O}_X$ sobre un abierto
  $U$, el homomorfismo $D_a:\mathscr{F}|U\to\mathscr{G}|U$ tal que, para
  toda sección $t$ de $\mathscr{F}$ sobre un abierto $V\subset U$, se tenga
  \[
    D_a(t)=D(at)-aD(t)
    \tag{16.8.8.1}
    \label{IV.16.8.8.1}
  \]
  es un operador diferencial de orden $\leq n-1$.
  \item[(c)] Para todo abierto $U$ de $X$, toda familia
  $(a_i)_{1\leq i\leq n+1}$ de $n+1$ secciones de $\mathcal{O}_X$ sobre
  $U$ y toda sección $t$ de $\mathscr{F}$ sobre $U$, se tiene la identidad
  \[
    \sum_{H\subset I_{n+1}}(-1)^{\operatorname{Card}(H)}
    \left(\prod_{i\in H}a_i\right)
    D\!\left(\left(\prod_{i\notin H}a_i\right)t\right)=0
    \tag{16.8.8.2}
    \label{IV.16.8.8.2}
  \]
  (donde $I_{n+1}$ es el intervalo $1\leq i\leq n+1$ de $\mathbb{N}$).
\end{enumerate}
\oldpage[IV-4]{43}
\end{proposition}

Demostremos primero la equivalencia entre a) y c). Por definición, para
demostrar que $D$ es un operador diferencial de orden $\leq n$, basta
demostrarlo para la restricción
$D|U:\mathscr{F}|U\to\mathscr{G}|U$ a todo abierto afín $U$ de $X$; por
otra parte, la propiedad c) es válida para todo abierto $U$ de $X$ si lo es
para todo abierto afín. Se puede, por tanto, reducir al caso en que
$S=\operatorname{Spec}(A)$ y $X=\operatorname{Spec}(B)$ son afines. En
virtud de (16.8.2) (cuyas notaciones se conservan), la condición a) significa
entonces que el $A$-homomorfismo $D':B\otimes_A M\to N$ tal que
$D'(b\otimes t)=bD(t)$ se anula sobre
$\mathfrak{I}^{n+1}(B\otimes_A M)$, lo que, en virtud de (0, 20.4.4),
equivale a decir que $D'$ se anula para todos los elementos de la forma
\[
  \left(\prod_{i=1}^{n+1}(a_i\otimes1-1\otimes a_i)\right).(1\otimes t)
\]
donde $a_i\in B$ y $t\in M$. Ahora bien, este elemento se escribe, salvo el
signo común $(-1)^{n+1}$, como\footnote{El texto impreso omite en esta
expansión los signos $(-1)^{\operatorname{Card}(H)}$; sin ellos, la imagen
por $D'$ no es el primer miembro de (16.8.8.2).}
\[
  \sum_{H\subset I_{n+1}}(-1)^{\operatorname{Card}(H)}
  \left(\prod_{i\in H}a_i\right)\otimes
  \left(\left(\prod_{i\notin H}a_i\right)t\right),
\]
y el valor de $D'$ sobre este elemento no es sino, salvo ese signo común, el
primer miembro de (16.8.8.2), lo que demuestra la equivalencia entre a) y c).

Demostremos ahora la equivalencia entre b) y c). Razonemos por inducción
sobre $n$, siendo trivial la afirmación para $n=0$. Escribiendo $a_{n+1}$ en
lugar de $a$ en la condición b), se ve, en virtud de la hipótesis de
inducción, que la condición b) significa que, para toda familia
$(a_i)_{1\leq i\leq n}$ de $n$ secciones de $\mathcal{O}_X$ sobre $U$ y
toda sección $t$ de $\mathscr{F}$ sobre $U$, se tiene
\[
  \sum_{H'\subset I_n}(-1)^{\operatorname{Card}(H')}
  \left(\prod_{i\in H'}a_i\right)
  D_{a_{n+1}}\!\left(\left(\prod_{i\notin H'}a_i\right)t\right)=0.
\]
Pero, si en esta relación se sustituye $D_{a_{n+1}}$ por su definición
(16.8.8.1), se comprueba inmediatamente que se obtiene, salvo el signo, el
primer miembro de (16.8.8.2); de donde se sigue la conclusión.

\begin{proposition}[16.8.9]
\label{IV.16.8.9}
Si $D:\mathscr{F}\to\mathscr{G}$ es un operador diferencial de orden
$\leq n$, y $D':\mathscr{G}\to\mathscr{H}$ un operador diferencial de orden
$\leq n'$, entonces $D'\circ D:\mathscr{F}\to\mathscr{H}$ es un operador
diferencial de orden $\leq n+n'$.
\end{proposition}

Por hipótesis, se puede escribir $D=u\circ d_{X/S,\mathscr{F}}^n$ y
$D'=v\circ d_{X/S,\mathscr{G}}^{n'}$, donde
$u:\mathscr{P}_{X/S}^n\otimes_{\mathcal{O}_X}\mathscr{F}\to\mathscr{G}$ y
$v:\mathscr{P}_{X/S}^{n'}\otimes_{\mathcal{O}_X}\mathscr{G}\to\mathscr{H}$
son $\mathcal{O}_X$-homomorfismos. Todo se reduce a demostrar que el
homomorfismo compuesto de haces de grupos aditivos
\[
  \xymatrix{
    \mathscr{F}\ar[r]^-{d_{X/S,\mathscr{F}}^n}
      &\mathscr{P}_{X/S}^n\otimes_{\mathcal{O}_X}\mathscr{F}
        \ar[r]^-u
      &\mathscr{G}\ar[r]^-{d_{X/S,\mathscr{G}}^{n'}}
      &\mathscr{P}_{X/S}^{n'}\otimes_{\mathcal{O}_X}\mathscr{G}
  }
\]
se factoriza como
\[
  \xymatrix{
    \mathscr{F}\ar[r]^-{d_{X/S,\mathscr{F}}^{n+n'}}
      &\mathscr{P}_{X/S}^{n+n'}\otimes_{\mathcal{O}_X}\mathscr{F}
        \ar[r]^-w
      &\mathscr{P}_{X/S}^{n'}\otimes_{\mathcal{O}_X}\mathscr{G}
  }
\]
donde $w$ es un $\mathcal{O}_X$-homomorfismo. Bastará demostrar el siguiente

\begin{lemma}[16.8.9.1]
\label{IV.16.8.9.1}
Existe un único $\mathcal{O}_X$-homomorfismo
\[
  \delta:\mathscr{P}_{X/S}^{n+n'}
  \longrightarrow
  \mathscr{P}_{X/S}^{n'}(\mathscr{P}_{X/S}^n)
  =\mathscr{P}_{X/S}^{n'}\otimes_{\mathcal{O}_X}\mathscr{P}_{X/S}^n
  \tag{16.8.9.2}
  \label{IV.16.8.9.2}
\]
\oldpage[IV-4]{44}
que hace conmutativo el diagrama
\[
  \xymatrix{
    \mathcal{O}_X\ar[r]^-{d_{X/S}^{n+n'}}\ar[d]_-{d_{X/S}^n}
      &\mathscr{P}_{X/S}^{n+n'}\ar[d]^-\delta\\
    \mathscr{P}_{X/S}^n
      \ar[r]_-{d_{X/S,\mathscr{P}_{X/S}^n}^{n'}}
      &\mathscr{P}_{X/S}^{n'}(\mathscr{P}_{X/S}^n)
  }
  \tag{16.8.9.3}
  \label{IV.16.8.9.3}
\]
\end{lemma}

Se tendrá entonces, en efecto, un diagrama conmutativo deducido de (16.8.9.3)
por tensorización con $\mathscr{F}$
\[
  \xymatrix{
    \mathscr{F}\ar[r]^-{d_{X/S,\mathscr{F}}^{n+n'}}
      \ar[d]_-{d_{X/S,\mathscr{F}}^n}
      &\mathscr{P}_{X/S}^{n+n'}(\mathscr{F})
        \ar[d]^-{\delta\otimes1}\\
    \mathscr{P}_{X/S}^n(\mathscr{F})
      \ar[r]_-{d_{X/S,\mathscr{P}_{X/S}^n(\mathscr{F})}^{n'}}
      &\mathscr{P}_{X/S}^{n'}(\mathscr{P}_{X/S}^n(\mathscr{F}))
  }
\]
y, por otra parte, se verifica inmediatamente a partir de la definición
(16.7.5) que el diagrama
\[
  \xymatrix{
    \mathscr{P}_{X/S}^n(\mathscr{F})\ar[r]^-u
      \ar[d]_-{d_{X/S,\mathscr{P}_{X/S}^n(\mathscr{F})}^{n'}}
      &\mathscr{G}\ar[d]^-{d_{X/S,\mathscr{G}}^{n'}}\\
    \mathscr{P}_{X/S}^{n'}(\mathscr{P}_{X/S}^n(\mathscr{F}))
      \ar[r]_-{1\otimes u}
      &\mathscr{P}_{X/S}^{n'}(\mathscr{G})
  }
\]
es conmutativo. Se resuelve, por tanto, la cuestión tomando como $w$ el
$\mathcal{O}_X$-homomorfismo compuesto
\[
  \xymatrix{
    \mathscr{P}_{X/S}^{n+n'}(\mathscr{F})
      \ar[r]^-{\delta\otimes1}
      &\mathscr{P}_{X/S}^{n'}(\mathscr{P}_{X/S}^n(\mathscr{F}))
        \ar[r]^-{1\otimes u}
      &\mathscr{P}_{X/S}^{n'}(\mathscr{G}).
  }
\]

Queda por demostrar el lema (16.8.9.1). Teniendo en cuenta (16.7.6), que
demuestra la unicidad de $\delta$, se reduce al caso en que
$S=\operatorname{Spec}(A)$ y $X=\operatorname{Spec}(B)$ son afines; poniendo
$\mathfrak{I}=\mathfrak{I}_{B/A}$, se trata de definir un homomorfismo
canónico de $B$-módulos
\[
  \varphi:(B\otimes_A B)/\mathfrak{I}^{n+n'+1}
  \longrightarrow
  ((B\otimes_A B)/\mathfrak{I}^{n'+1})
  \otimes_B((B\otimes_A B)/\mathfrak{I}^{n+1})
\]
donde las estructuras de $B$-módulo de ambos miembros provienen del
primer factor $B$; recordemos que, en el producto tensorial del segundo
miembro, $(B\otimes_A B)/\mathfrak{I}^{n'+1}$ debe considerarse
\oldpage[IV-4]{45}
como $B$-módulo a la derecha mediante su segundo factor $B$, y
$(B\otimes_A B)/\mathfrak{I}^{n+1}$ como $B$-módulo a la izquierda
mediante su primer factor $B$ (16.7.2). Esto equivale a definir un
homomorfismo de $B$-módulos
\[
  \varphi_0:B\otimes_A B\longrightarrow
  ((B\otimes_A B)/\mathfrak{I}^{n'+1})
  \otimes_B((B\otimes_A B)/\mathfrak{I}^{n+1})
\]
y demostrar que se anula sobre $\mathfrak{I}^{n+n'+1}$. Ahora bien, se
define inmediatamente tal homomorfismo mediante la condición
\[
  \varphi_0(b\otimes b')
  =\pi_{n'}(b\otimes1)\otimes\pi_n(1\otimes b')
  \qquad\text{para $b,b'$ en $B$}
\]
con las notaciones de (16.3.7). Además, es inmediato que $\varphi_0$ es un
homomorfismo de \emph{anillos}. Ahora bien, se puede escribir
\[
  \begin{aligned}
  \varphi_0(b\otimes1-1\otimes b)
  ={}&\pi_{n'}(b\otimes1-1\otimes b)\otimes\pi_n(1\otimes1)\\
  &+\pi_{n'}(1\otimes b)\otimes\pi_n(1\otimes1)
  -\pi_{n'}(1\otimes1)\otimes\pi_n(1\otimes b)
  \end{aligned}
\]
y se tiene
\[
  \begin{aligned}
  \pi_{n'}(1\otimes b)\otimes\pi_n(1\otimes1)
  &=\pi_{n'}(1\otimes1)b\otimes\pi_n(1\otimes1)\\
  &=\pi_{n'}(1\otimes1)\otimes b\pi_n(1\otimes1)\\
  &=\pi_{n'}(1\otimes1)\otimes\pi_n(b\otimes1)
  \end{aligned}
\]
de donde, finalmente,
\[
  \begin{aligned}
  \varphi_0(b\otimes1-1\otimes b)
  ={}&\pi_{n'}(b\otimes1-1\otimes b)\otimes\pi_n(1\otimes1)\\
  &+\pi_{n'}(1\otimes1)\otimes\pi_n(b\otimes1-1\otimes b).
  \end{aligned}
  \tag{16.8.9.4}
  \label{IV.16.8.9.4}
\]
Un producto de $n+n'+1$ términos de la forma (16.8.9.4) es, por tanto,
necesariamente nulo, pues así sucede con un producto de $n+1$ términos de la
forma $\pi_n(b\otimes1-1\otimes b)$ y con un producto de $n'+1$ términos de
la forma $\pi_{n'}(b\otimes1-1\otimes b)$. La conclusión resulta, por tanto,
de (0, 20.4.4).

\begin{corollary}[16.8.10]
\label{IV.16.8.10}
El haz $\mathscr{D}\!iff_{X/S}(\mathcal{O}_X,\mathcal{O}_X)$ (denotado también
por $\mathscr{D}\!iff_{X/S}$) está provisto canónicamente de una estructura de
haz de anillos, y los $\mathscr{D}\!iff_{X/S}^n$ forman una filtración
creciente compatible con esta estructura.
\end{corollary}

En particular, $\mathscr{D}\!iff_{X/S}^0$ es un haz de subanillos de
$\mathscr{D}\!iff_{X/S}$ que se identifica canónicamente con $\mathcal{O}_X$
(16.8.1). Las fórmulas (16.8.5.1) y (16.8.5.2) muestran que la estructura de
$\mathcal{O}_X$-bimódulo de $\mathscr{D}\!iff_{X/S}$ proviene de la
multiplicación por la izquierda y por la derecha mediante las secciones de
$\mathcal{O}_X$, considerado como haz de subanillos de
$\mathscr{D}\!iff_{X/S}$.

\begin{remarks}[16.8.11]
\label{IV.16.8.11}
\begin{enumerate}
  \item[(i)] Supongamos que
  $\mathscr{F}=\bigoplus_{\lambda\in L}\mathscr{F}_\lambda$; entonces es
  claro por (16.7.2.1) que
  $\mathscr{P}_{X/S}^n(\mathscr{F})=
  \bigoplus_{\lambda\in L}\mathscr{P}_{X/S}^n(\mathscr{F}_\lambda)$; como el
  funtor $\mathscr{F}\mapsto\Gamma(U,\mathscr{F})$ conmuta con la formación
  de sumas directas cualesquiera, $d_{X/S,\mathscr{F}}^n$ es el homomorfismo
  cuya restricción a cada $\mathscr{F}_\lambda$ es
  \[
    d_{X/S,\mathscr{F}_\lambda}^n:\mathscr{F}_\lambda
    \longrightarrow\mathscr{P}_{X/S}^n(\mathscr{F}_\lambda);
  \]
  se concluye inmediatamente que se tiene
  \[
    \operatorname{Diff}_{X/S}^n(\mathscr{F},\mathscr{G})
    =\prod_{\lambda\in L}
    \operatorname{Diff}_{X/S}^n(\mathscr{F}_\lambda,\mathscr{G}),
  \]
  y, por consiguiente, también (0\textsubscript{I}, 3.2.6)
  \[
    \mathscr{D}\!iff_{X/S}^n(\mathscr{F},\mathscr{G})
    =\prod_{\lambda\in L}
    \mathscr{D}\!iff_{X/S}^n(\mathscr{F}_\lambda,\mathscr{G}).
  \]
  Por otra parte, si $\mathscr{G}=\prod_{\mu\in M}\mathscr{G}_\mu$
  (0\textsubscript{I}, 3.2.6), se tiene
  \[
    \operatorname{Hom}_{\mathcal{O}_X}
    (\mathscr{P}_{X/S}^n(\mathscr{F}),\mathscr{G})
    =\prod_{\mu\in M}\operatorname{Hom}_{\mathcal{O}_X}
    (\mathscr{P}_{X/S}^n(\mathscr{F}),\mathscr{G}_\mu),
  \]
  \oldpage[IV-4]{46}
  pues todo homomorfismo $u$ de $\mathscr{P}_{X/S}^n(\mathscr{F})$ en
  $\mathscr{G}$ corresponde biunívocamente a la familia de sus compuestos
  \[
    u_\mu:\mathscr{P}_{X/S}^n(\mathscr{F})
    \longrightarrow\mathscr{G}\longrightarrow\mathscr{G}_\mu.
  \]
  Se tiene, por tanto,
  \[
    \operatorname{Diff}_{X/S}^n(\mathscr{F},\mathscr{G})
    =\prod_{\mu\in M}
    \operatorname{Diff}_{X/S}^n(\mathscr{F},\mathscr{G}_\mu),
  \]
  y, por consiguiente, también
  \[
    \mathscr{D}\!iff_{X/S}^n(\mathscr{F},\mathscr{G})
    =\prod_{\mu\in M}
    \mathscr{D}\!iff_{X/S}^n(\mathscr{F},\mathscr{G}_\mu).
  \]
  \item[(ii)] Hasta aquí apenas se han encontrado operadores diferenciales
  $\mathscr{F}\to\mathscr{G}$ salvo cuando $\mathscr{F}$ y $\mathscr{G}$ son
  $\mathcal{O}_X$-módulos localmente libres de rango finito, en cuyo caso, en
  virtud de (i), su estructura se reduce localmente a la del haz
  $\mathscr{D}\!iff_{X/S}$; esta última se estudiará más adelante (16.11) en
  un caso particular.
\end{enumerate}
\end{remarks}

\subsection{Inmersiones regulares y cuasirregulares.}
\label{IV.16.9}

\begin{definition}[16.9.1]
\label{IV.16.9.1}
Sea $X$ un espacio anillado. Se dice que un ideal $\mathscr{I}$ de
$\mathcal{O}_X$ es \emph{regular} (resp. \emph{cuasirregular}) si, para todo
punto $x\in\operatorname{Supp}(\mathcal{O}_X/\mathscr{I})$, existen una
vecindad abierta $U$ de $x$ en $X$ y una sucesión regular (0, 15.2.2) (resp.
cuasirregular (0, 15.2.2)) de elementos de $\Gamma(U,\mathcal{O}_X)$ que
engendra $\mathscr{I}|U$.
\end{definition}

Se dirá que una sucesión regular (resp. cuasirregular) de secciones de
$\mathcal{O}_X$ sobre $U$ que engendra $\mathscr{I}|U$ es un
\emph{sistema regular} (resp. \emph{cuasirregular}) \emph{de generadores}
de $\mathscr{I}|U$.

\begin{definition}[16.9.2]
\label{IV.16.9.2}
Sea $j:Y\to X$ una inmersión de preesquemas y sea $U$ un abierto de $X$ tal
que $j(Y)\subset U$ y que $j$ sea una inmersión cerrada de $Y$ en $U$. Se
dice que $j$ es \emph{regular} (resp. \emph{cuasirregular}) si el
subpreesquema cerrado $j(Y)$ de $U$ asociado a $j$ está definido por un ideal
regular (resp. cuasirregular) de $\mathcal{O}_U$ (condición independiente del
abierto $U$ elegido).
\end{definition}

Se dice que un subpreesquema $Y$ de un preesquema $X$ está \emph{inmerso
regularmente} (resp. \emph{inmerso cuasirregularmente}) si la inyección
canónica $j:Y\to X$ es una inmersión regular (resp. cuasirregular). Si $Y$ es
un subpreesquema cerrado de $X$ y $\mathscr{I}$ el ideal de $\mathcal{O}_X$
que define $Y$, esto equivale a decir que $\mathscr{I}$ es \emph{regular}
(resp. \emph{cuasirregular}).

Por ejemplo, si $A$ es un anillo \emph{íntegro} y $f$ un elemento $\neq0$ de
$A$, el subpreesquema cerrado $V(f)$ de $\operatorname{Spec}(A)$ (isomorfo a
$\operatorname{Spec}(A/fA)$) está \emph{inmerso regularmente} en
$\operatorname{Spec}(A)$.

Todo ideal regular es cuasirregular (0, 15.2.2); toda inmersión regular es
cuasirregular (cf. (16.9.11) para un enunciado recíproco).

\begin{proposition}[16.9.3]
\label{IV.16.9.3}
Sean $X$ un espacio anillado, $\mathscr{I}$ un ideal de $\mathcal{O}_X$ y
$(f_i)_{1\leq i\leq m}$ una sucesión finita de secciones de $\mathcal{O}_X$
sobre $X$ que engendra $\mathscr{I}$. Para que $(f_i)$ sea una sucesión
cuasirregular (0, 15.2.2), es necesario y suficiente que se satisfagan las
condiciones siguientes:
\begin{enumerate}
  \item[(i)] Las imágenes canónicas de los $f_i$ en
  $\mathscr{I}/\mathscr{I}^2$ forman una base de este
  $\mathcal{O}_X/\mathscr{I}$-módulo.
  \item[(ii)] El homomorfismo sobreyectivo canónico (16.1.2.2)
  \[
    \mathbf{S}_{\mathcal{O}_X/\mathscr{I}}^{\bullet}
    (\mathscr{I}/\mathscr{I}^2)
    \longrightarrow\mathscr{G}r_{\mathscr{I}}^{\bullet}(\mathcal{O}_X)
  \]
  es biyectivo.
\end{enumerate}
Además, si así ocurre, toda sucesión $(f'_i)_{1\leq i\leq m}$ de $m$
secciones de $\mathscr{I}$ sobre $X$ que engendra $\mathscr{I}$ es
cuasirregular.\footnote{El texto impreso usa $n$ en ambos lugares; debe ser
$m$, el número de elementos de la base dada por $(f_i)_{1\leq i\leq m}$. La
demostración exige que la nueva familia tenga el mismo número de elementos que
una base de $\mathscr{I}/\mathscr{I}^2$.}
\end{proposition}

\oldpage[IV-4]{47}

Las dos condiciones del enunciado no hacen sino expresar la definición dada
en (0, 15.2.2), teniendo en cuenta la definición de los homomorfismos
canónicos (0, 15.2.1.1). La última afirmación resulta de que, si un módulo $M$
sobre un anillo conmutativo $A$ admite una base de $n$ elementos, todo sistema
de generadores de $M$ que tenga $n$ elementos es una base de $M$ (Bourbaki,
\emph{Alg. comm.}, cap.~II, §~3, cor.~5 del teorema~1).

\begin{corollary}[16.9.4]
\label{IV.16.9.4}
Sean $X$ un espacio localmente anillado y $\mathscr{I}$ un ideal de
$\mathcal{O}_X$. Para que $\mathscr{I}$ sea cuasirregular, es necesario y
suficiente que satisfaga las condiciones siguientes:
\begin{enumerate}
  \item[(i)] $\mathscr{I}$ es de tipo finito.
  \item[(ii)] $\mathscr{I}/\mathscr{I}^2$ es un
  $(\mathcal{O}_X/\mathscr{I})$-módulo localmente libre.
  \item[(iii)] El homomorfismo canónico
  \[
    \mathbf{S}_{\mathcal{O}_X/\mathscr{I}}^{\bullet}
    (\mathscr{I}/\mathscr{I}^2)
    \longrightarrow\mathscr{G}r_{\mathscr{I}}^{\bullet}(\mathcal{O}_X)
    \tag{16.9.4.1}
    \label{IV.16.9.4.1}
  \]
  es biyectivo.
\end{enumerate}
\end{corollary}

La necesidad de las condiciones resulta inmediatamente de (16.9.3). Para ver
que son suficientes, hay que demostrar, en virtud de (16.9.3), que, si en un
punto $x\in\operatorname{Supp}(\mathcal{O}_X/\mathscr{I})$ existen una
vecindad abierta $U$ de $x$ en $X$ y $n$ secciones $f_i$
($1\leq i\leq n$) de $\mathscr{I}$ sobre $U$ cuyas imágenes canónicas en
$\mathscr{I}/\mathscr{I}^2$ forman una base de
$(\mathscr{I}/\mathscr{I}^2)|U$ sobre
$(\mathcal{O}_X/\mathscr{I})|U$, entonces existe una vecindad abierta
$V\subset U$ de $x$ tal que los $f_i|V$ engendran $\mathscr{I}|V$. Ahora
bien, por hipótesis, $\mathscr{I}_x\neq\mathcal{O}_x$, por lo que
$\mathscr{I}_x$ está contenido en el ideal maximal de $\mathcal{O}_x$; como
$\mathscr{I}_x$ es un $\mathcal{O}_x$-módulo de tipo finito y las clases de
los $(f_i)_x$ en $\mathscr{I}_x/\mathscr{I}_x^2$ engendran este
$(\mathcal{O}_x/\mathscr{I}_x)$-módulo, el lema de Nakayama muestra que los
$(f_i)_x$ engendran $\mathscr{I}_x$. Como $\mathscr{I}$ es de tipo finito,
se concluye por (0\textsubscript{I}, 5.2.2).

\begin{corollary}[16.9.5]
\label{IV.16.9.5}
Sean $X$ un espacio localmente anillado, $\mathscr{I}$ un ideal cuasirregular
de $\mathcal{O}_X$, $(f_i)_{1\leq i\leq n}$ una sucesión de secciones de
$\mathscr{I}$ sobre $X$, y
$x$ un punto de $\operatorname{Supp}(\mathcal{O}_X/\mathscr{I})$. Las
condiciones siguientes son equivalentes:
\begin{enumerate}
  \item[(a)] Existe una vecindad abierta $U$ de $x$ en $X$ tal que los
  $f_i|U$ forman una sucesión cuasirregular de elementos de
  $\Gamma(U,\mathcal{O}_X)$ que engendra $\mathscr{I}|U$.
  \item[(b)] Los $(f_i)_x$ forman un sistema de generadores de
  $\mathscr{I}_x$ cuyo número de elementos es el menor posible.
  \item[(b')] Los $(f_i)_x$ forman un sistema minimal de generadores de
  $\mathscr{I}_x$.
  \item[(c)] Si $\bar f_i$ es la imagen canónica de $f_i$ en
  $\Gamma(X,\mathscr{I}/\mathscr{I}^2)$, los $(\bar f_i)_x$ forman una base
  del $(\mathcal{O}_x/\mathscr{I}_x)$-módulo
  $\mathscr{I}_x/\mathscr{I}_x^2$.
\end{enumerate}
\end{corollary}

Por hipótesis, $\mathcal{O}_x$ es un anillo local y $\mathscr{I}_x$ un ideal
de tipo finito de $\mathcal{O}_x$ contenido en el ideal maximal de
$\mathcal{O}_x$; la equivalencia entre b), b') y c) resulta, por tanto, del
lema de Nakayama (Bourbaki, \emph{Alg. comm.}, cap.~II, §~3,
n\textsuperscript{o}~2, prop.~5). Es claro que a) implica c) en virtud de
(16.9.3); por otra parte, de (0\textsubscript{I}, 5.2.2) resulta que, si se
satisface la condición c) (y, por tanto, también b)), existe una vecindad $U$
de $x$ en $X$ tal que $(\mathscr{I}/\mathscr{I}^2)|U$ tiene rango constante
igual a $n$, y que los $f_i|U$ engendran $\mathscr{I}|U$; basta, por tanto,
aplicar en $U$ la última afirmación de (16.9.3).

\begin{remarks}[16.9.6]
\label{IV.16.9.6}
\begin{enumerate}
  \item[(i)] Bajo las hipótesis generales de (16.9.5), para que la sucesión
  $(f_i)$ engendre $\mathscr{I}$ no basta con que los $(\bar f_i)_y$ formen
  una base del $(\mathcal{O}_y/\mathscr{I}_y)$-módulo
  $\mathscr{I}_y/\mathscr{I}_y^2$ para todo $y\in X$. Se obtiene un ejemplo
  tomando
\oldpage[IV-4]{48}
  $X=\operatorname{Spec}(A)$, donde $A$ es un anillo de Dedekind, y
  $\mathscr{I}=\widetilde{\mathfrak{I}}$, donde $\mathfrak{I}$ es un ideal
  primo no principal de $A$; se tiene entonces, en efecto,
  $\mathscr{I}_y/\mathscr{I}_y^2=0$ en todo punto $y$ distinto del punto
  $x\in X$ correspondiente a $\mathfrak{I}$, y
  $\mathscr{I}_x/\mathscr{I}_x^2$ es de rango $1$ sobre el cuerpo
  $\mathcal{O}_x/\mathscr{I}_x$; además, $\mathscr{I}$ es evidentemente un
  ideal regular.
  \item[(ii)] En (16.9.5) no se puede sustituir «cuasirregular» por
  «regular», ni siquiera cuando $X$ es un preesquema (cf. (16.9.12)). En
  efecto, designemos por $B$ el anillo de los gérmenes de funciones
  indefinidamente diferenciables en el punto $0$ de $\mathbb{R}$; tiene un
  ideal maximal $\mathfrak{m}$ engendrado por el germen $t$ de la aplicación
  identidad de $\mathbb{R}$ en el punto $0$, y la intersección $\mathfrak{n}$
  de los $\mathfrak{m}^k$ para $k>0$ no se reduce a $0$. Sea ahora $A$ el
  anillo cociente $B[T]/\mathfrak{n}TB[T]$, y sean $f_1,f_2$ las imágenes
  canónicas en $A$ de los elementos $t$ y $T$ de $B[T]$. La sucesión
  $(f_1,f_2)$ es regular en $A$: en efecto, $f_1$ no es divisor de $0$ en
  $A$, pues la relación $tP[T]\in\mathfrak{n}TB[T]$, para un polinomio
  $P\in B[T]$, implica que los productos de $t$ por los coeficientes de $P$
  pertenecen al ideal $\mathfrak{n}$, de lo que resulta inmediatamente que
  estos coeficientes están a su vez en $\mathfrak{n}$, y, por tanto,
  $P[T]\in\mathfrak{n}TB[T]$. Como $B/tB$ es isomorfo a $\mathbb{R}$,
  $A/f_1A$ es isomorfo al anillo de polinomios $\mathbb{R}[T]$, y, por tanto,
  es íntegro; la imagen de $f_2$ en $A/f_1A$, al ser igual a $T$, no es, por
  consiguiente, divisor de $0$, lo que demuestra la afirmación. Sin embargo,
  $f_2$ es divisor de $0$ en $A$, pues para todo elemento no nulo
  $x\in\mathfrak{n}$, la imagen de $x$ en $A$ es $\ne0$, mientras que la
  imagen de $xT$ es nula. Se concluye que la sucesión $(f_2,f_1)$ no es
  regular en $A$; por otra parte, el ideal $\mathfrak{I}=f_1A+f_2A$ es
  distinto de $A$, por lo que las condiciones b), b') y c) de (16.9.5) no
  implican la condición a) cuando en ella se sustituye «cuasirregular» por
  «regular».
\end{enumerate}
\end{remarks}

\begin{env}[16.9.7]
\label{IV.16.9.7}
Si $X=\operatorname{Spec}(A)$ es un esquema afín, se dirá que un ideal
$\mathfrak{I}$ de $A$ es regular (resp. cuasirregular) si el ideal
$\mathscr{I}=\widetilde{\mathfrak{I}}$ de $\mathcal{O}_X$ es regular
(resp. cuasirregular); obsérvese que esta noción es local y no implica en
absoluto la existencia de un sistema de generadores de $\mathfrak{I}$ que
forme en $A$ una sucesión regular (resp. cuasirregular), como muestra el
ejemplo (16.9.6, (i)); sin embargo, así ocurre cuando $A$ es local (16.9.5).

La proposición (16.9.4) se expresa en términos de inmersiones cuasirregulares
de la manera siguiente:
\end{env}

\begin{proposition}[16.9.8]
\label{IV.16.9.8}
Sea $j:Y\to X$ un morfismo de preesquemas; para que $j$ sea una inmersión
cuasirregular, es necesario y suficiente que $j$ satisfaga las condiciones
siguientes:
\begin{enumerate}
  \item[(i)] $j$ es una inmersión localmente de presentación finita.
  \item[(ii)] El haz conormal
  $\operatorname{Gr}^1(j)=\mathcal{N}_{Y/X}$ (16.1.2) es un
  $\mathcal{O}_Y$-módulo localmente libre.
  \item[(iii)] El homomorfismo canónico
  \[
    \operatorname{S}_{\mathcal{O}_Y}^{\bullet}
      (\operatorname{Gr}^1(j))\longrightarrow\operatorname{Gr}^{\bullet}(j)
  \]
  (16.1.2.2) es biyectivo.
\end{enumerate}
\end{proposition}

Como la cuestión es local sobre $Y$, se puede reducir al caso en que $j$ es
la inyección canónica de un subpreesquema cerrado $Y$ de $X$, y entonces la
expresión de (16.9.4) en la forma (16.9.8) resulta de la explicitación de
$\operatorname{Gr}^1(j)$ y $\operatorname{Gr}^{\bullet}(j)$ en términos del
ideal $\mathscr{I}$ de $\mathcal{O}_X$ que define el subpreesquema $Y$
(16.1.3, (ii)).

\begin{corollary}[16.9.9]
\label{IV.16.9.9}
Sean $Y$ un preesquema, $X$ un $Y$-preesquema y $j:Y\to X$ una $Y$-sección
de $X$, de modo que el $n$-ésimo invariante normal $\mathscr{A}^{(n)}$ de $j$
(16.1.2) es una $\mathcal{O}_Y$-álgebra aumentada (16.1.7); pongamos
$\mathscr{A}^{(\infty)}=\varprojlim\mathscr{A}^{(n)}$. Para que $j$ sea una
inmersión cuasirregular, es necesario y suficiente que $j$ sea localmente de
presentación finita y que todo $y\in Y$ admita una vecindad abierta afín $U$
cuyo anillo es $C$ tal que $\mathscr{A}^{(\infty)}|U$ sea isomorfa, como
$\mathcal{O}_U$-álgebra topológica aumentada, a
$\mathcal{O}_U[[T_1,\ldots,T_n]]$.
\end{corollary}

Se puede reducir al caso en que $j$ es una inmersión cerrada restringiéndose
a una vecindad suficientemente pequeña de $y$ (véase el razonamiento de
(16.4.11)), y entonces $\mathcal{O}_Y$ se identifica con el álgebra cociente
$\mathcal{O}_X/\mathscr{I}$ y el homomorfismo sobreyectivo canónico
$\mathcal{O}_X\to\mathcal{O}_Y$ admite un
\oldpage[IV-4]{49}
inverso por la derecha (16.1.7). Se puede, por tanto, suponer que
$X=\operatorname{Spec}(B)$ e $Y=\operatorname{Spec}(A)$ son afines, que $B$
es una $A$-álgebra aumentada y que el ideal de aumento $\mathfrak{I}$ es de
tipo finito. Como $\mathscr{A}^{(n)}$ se identifica entonces con
$(B/\mathfrak{I}^{n+1})^{\sim}$, el corolario resulta de la equivalencia entre
b) y c) en (0, 19.5.4), puesto que $B/\mathfrak{I}=A$.

Obsérvese que, en el caso afín considerado, el hecho de que $j$ sea una
inmersión cuasirregular equivale también, en virtud de (0, 19.5.4), a decir
que $B$ es una $A$-álgebra formalmente lisa para la topología
$\mathfrak{I}$-preádica.

Obsérvese asimismo que la condición de que $j$ sea una inmersión localmente de
presentación finita se satisface siempre cuando el morfismo $X\to Y$ es
localmente de tipo finito (1.4.3, (v)).

\begin{proposition}[16.9.10]
\label{IV.16.9.10}
Sean $X$ un preesquema localmente noetheriano, $Y$ un subpreesquema de $X$,
$j:Y\to X$ la inyección canónica y $y$ un punto de $Y$.
\begin{enumerate}
  \item[(i)] Para que exista una vecindad abierta $U$ de $y$ en $X$ tal que
  la restricción $Y\cap U\to U$ de $j$ sea una inmersión regular, es necesario
  y suficiente que el núcleo $\mathscr{I}_y$ del homomorfismo sobreyectivo
  $\mathcal{O}_{X,y}\to\mathcal{O}_{Y,y}$ esté engendrado por una sucesión
  regular de elementos de $\mathcal{O}_{X,y}$.
  \item[(ii)] Para que la inmersión $j$ sea regular, es necesario y suficiente
  que sea cuasirregular.
\end{enumerate}
\end{proposition}

\begin{enumerate}
  \item[(i)] Se puede reducir al caso en que $Y$ es un subpreesquema cerrado
  de $X$ definido por un ideal \emph{coherente} $\mathscr{I}$ de
  $\mathcal{O}_X$. La condición es evidentemente necesaria. Recíprocamente,
  si $\mathscr{I}_y$ está engendrado por una sucesión regular $(s_i)_y$, donde
  los $s_i$ son secciones de $\mathscr{I}$ sobre una vecindad abierta $U$ de
  $y$ en $X$, se puede suponer que los $s_i$ engendran $\mathscr{I}|U$
  (0\textsubscript{I}, 5.2.2) y forman una sucesión regular (0, 15.2.4), de
  donde se sigue la afirmación.
  \item[(ii)] El hecho de que una inmersión cuasirregular sea regular resulta
  de (i) y de la identidad entre las sucesiones cuasirregulares y las
  sucesiones regulares de $\mathcal{O}_{X,y}$ formadas por elementos del ideal
  maximal (0, 15.1.11).
\end{enumerate}

Cuando (sin hipótesis noetheriana sobre $X$) el núcleo $\mathscr{I}_y$ de
$\mathcal{O}_{X,y}\to\mathcal{O}_{Y,y}$ está engendrado por una sucesión
regular de elementos de $\mathcal{O}_{X,y}$, se dice que la inmersión $j$ es
\emph{regular en el punto $y$}.

\begin{corollary}[16.9.11]
\label{IV.16.9.11}
Sea $X$ un preesquema localmente noetheriano; entonces todo ideal
cuasirregular de $\mathcal{O}_X$ es regular.
\end{corollary}

\begin{remarks}[16.9.12]
\label{IV.16.9.12}
\begin{enumerate}
  \item[(i)] Obsérvese que una inmersión regular no es, en general, un
  morfismo plano, ni \emph{a fortiori} un morfismo \emph{regular} en el sentido
  de (6.8.1).
  \item[(ii)] Sea $A$ un anillo local noetheriano; de (16.9.4) y (0, 17.1.1)
  resulta inmediatamente que, para que $A$ sea \emph{regular}, es necesario y
  suficiente que su ideal maximal $\mathfrak{m}$ sea \emph{cuasirregular}
  (o \emph{regular}, lo que equivale a lo mismo puesto que $A$ es
  noetheriano). Para que un esquema afín noetheriano $X$ sea \emph{regular},
  es necesario y suficiente que, para todo punto cerrado $x\in X$, la
  inyección canónica $\operatorname{Spec}(k(x))\to X$ sea una inmersión
  \emph{regular}.
\end{enumerate}
\end{remarks}

\begin{proposition}[16.9.13]
\label{IV.16.9.13}
Sean $X$ un preesquema localmente noetheriano, $Y$ un subpreesquema de $X$ e
$Y'$ un subpreesquema de $Y$, tal que la inyección canónica $j:Y'\to Y$ sea
regular. Entonces la sucesión de $\mathcal{O}_{Y'}$-módulos
\[
  \xymatrix{
    0\ar[r] & j^*(\mathcal{N}_{Y/X})\ar[r]
      & \mathcal{N}_{Y'/X}\ar[r] & \mathcal{N}_{Y'/Y}\ar[r] & 0
  }
  \tag{16.9.13.1}
  \label{IV.16.9.13.1}
\]
\oldpage[IV-4]{50}
es exacta; además, para todo $x\in X$, existe una vecindad abierta $U$ de $x$
tal que las restricciones a $U$ de los homomorfismos de (16.9.13.1) forman
una sucesión exacta y escindida.
\end{proposition}

Demostremos primero el lema siguiente:

\begin{lemma}[16.9.13.2]
\label{IV.16.9.13.2}
Sean $A$ un anillo, $\mathfrak{I}$ un ideal de $A$,
$A'=A/\mathfrak{I}$, $(f_i)_{1\leq i\leq r}$ una sucesión de elementos de
$A$ que es $A'$-regular, $\mathfrak{K}=\sum_i f_iA$,
$\mathfrak{L}=\mathfrak{I}+\mathfrak{K}$ y
$\mathfrak{K}'=\sum_i f_iA'$, de modo que $C=A/\mathfrak{L}$ es isomorfo a
$A'/\mathfrak{K}'$. Entonces, para todo entero $n>0$ y todo entero $N\geq n$,
se tiene la relación
\[
  \mathfrak{I}\cap\mathfrak{K}^n
  =\mathfrak{I}\mathfrak{K}^n+(\mathfrak{I}\cap\mathfrak{K}^N).
  \tag{16.9.13.3}
  \label{IV.16.9.13.3}
\]
\end{lemma}

Evidentemente, basta demostrar que todo elemento del primer miembro está
contenido en el segundo y, por inducción sobre $N-n$,\footnote{El texto
impreso atribuye la inducción a $n$; la iteración que reduce el enunciado al
caso $N=n+1$ es una inducción sobre la diferencia $N-n$.} se reduce al caso en
que $N=n+1$. Un elemento del primer miembro de (16.9.13.3), al pertenecer a
$\mathfrak{K}^n$, se escribe $P(f_1,\ldots,f_r)$, donde
$P\in A[T_1,\ldots,T_r]$ es homogéneo y de grado $n$. Si $f_i'$ es la imagen
canónica de $f_i$ en $A'$, la hipótesis
$P(f_1,\ldots,f_r)\in\mathfrak{I}$ significa que
$P(f_1',\ldots,f_r')=0$. Pero
$P(f_1',\ldots,f_r')\in\mathfrak{K}'^n$, por lo que la imagen canónica de
$P(f_1',\ldots,f_r')$ en
$\mathfrak{K}'^n/\mathfrak{K}'^{n+1}$ es nula. Ahora bien, la hipótesis de que
la sucesión $(f_i)$ es $A'$-regular implica que el homomorfismo canónico
\[
  \mathbf{S}_C^n(\mathfrak{K}'/\mathfrak{K}'^2)
  \longrightarrow\mathfrak{K}'^n/\mathfrak{K}'^{n+1}
\]
es biyectivo (0, 15.1.9); se concluye que los coeficientes de $P$ pertenecen
a $\mathfrak{L}=\mathfrak{I}+\mathfrak{K}$. De ello resulta inmediatamente
que $P(f_1,\ldots,f_r)\in\mathfrak{I}\mathfrak{K}^n+\mathfrak{K}^{n+1}$, y,
como $P(f_1,\ldots,f_r)\in\mathfrak{I}$, se tiene finalmente
\[
  P(f_1,\ldots,f_r)
  \in\mathfrak{I}\mathfrak{K}^n
    +(\mathfrak{I}\cap\mathfrak{K}^{n+1}),
\]
lo que demuestra el lema.

Tomando los cocientes de ambos miembros de (16.9.13.3) por
$\mathfrak{I}\mathfrak{K}^n$, se ve que las relaciones (16.9.13.3) para
$N\geq n$ implican
\[
  (\mathfrak{I}\cap\mathfrak{K}^n)/\mathfrak{I}\mathfrak{K}^n
  \subset\bigcap_{N\geq n}\mathfrak{K}^N
    \cdot(A/(\mathfrak{I}\mathfrak{K}^n)).
  \tag{16.9.13.4}
  \label{IV.16.9.13.4}
\]
De ello se deduce el

\begin{corollary}[16.9.13.5]
\label{IV.16.9.13.5}
Supongamos satisfechas las hipótesis de (16.9.13.2) y supongamos, además, que
el anillo $A$ es noetheriano y que $\mathfrak{K}$ está contenido en el radical
de $A$. Entonces, para todo entero $n>0$, se tiene
\[
  \mathfrak{I}\cap\mathfrak{K}^n=\mathfrak{I}\mathfrak{K}^n.
  \tag{16.9.13.6}
  \label{IV.16.9.13.6}
\]
\end{corollary}

En efecto, el segundo miembro de (16.9.13.4) es entonces nulo, puesto que
$A/\mathfrak{I}\mathfrak{K}^n$ es un $A$-módulo de tipo finito (Bourbaki,
\emph{Alg. comm.}, cap.~III, §~3, n\textsuperscript{o}~3, prop.~6).

Tomemos, en particular, $n=2$ en (16.9.13.6), y observemos que se tiene
$\mathfrak{L}^2=\mathfrak{I}^2+\mathfrak{I}\mathfrak{K}+\mathfrak{K}^2
=\mathfrak{I}\mathfrak{L}+\mathfrak{K}^2$; puesto que
$\mathfrak{I}\mathfrak{L}\subset\mathfrak{L}^2$, se deduce que
\[
  \mathfrak{I}\cap\mathfrak{L}^2
  =\mathfrak{I}\mathfrak{L}+(\mathfrak{I}\cap\mathfrak{K}^2)
  =\mathfrak{I}\mathfrak{L}+\mathfrak{I}\mathfrak{K}^2
  =\mathfrak{I}\mathfrak{L},
\]
es decir,
\[
  \mathfrak{I}\cap\mathfrak{L}^2=\mathfrak{I}\mathfrak{L},
  \tag{16.9.13.7}
  \label{IV.16.9.13.7}
\]
lo que también puede expresarse diciendo que el homomorfismo canónico
\[
  \mathfrak{I}/\mathfrak{I}\mathfrak{L}
  \longrightarrow(\mathfrak{I}+\mathfrak{L}^2)/\mathfrak{L}^2
\]
es biyectivo.

Demostrados estos lemas, demostremos la primera afirmación de (16.9.13): basta
\oldpage[IV-4]{51}
evidentemente demostrar que la sucesión de los tallos, en un punto $x\in Y'$,
de los haces que figuran en (16.9.13.1) es exacta. Ahora bien, si se pone
$A=\mathcal{O}_{X,x}$, se puede escribir
$\mathcal{O}_{Y,x}=A'=A/\mathfrak{I}$, donde $\mathfrak{I}$ es un ideal
contenido en el ideal maximal de $A$, y luego
$\mathcal{O}_{Y',x}=A'/\mathfrak{K}'$, donde $\mathfrak{K}'$ está engendrado
por una sucesión $A'$-regular de elementos de $A'$, que son a su vez imágenes
de los elementos de una sucesión $A'$-regular de elementos de $A$
pertenecientes al ideal maximal de $A$. Si $\mathfrak{K}$ es el ideal
engendrado por estos últimos y $\mathfrak{L}=\mathfrak{I}+\mathfrak{K}$, se
tiene $\mathcal{O}_{Y',x}=A/\mathfrak{L}$ y, como se está en la situación de
(16.9.13.5), el homomorfismo canónico
\[
  \mathfrak{I}/\mathfrak{I}\mathfrak{L}
  \longrightarrow(\mathfrak{I}+\mathfrak{L}^2)/\mathfrak{L}^2
\]
es biyectivo. Pero esto muestra que la sucesión
\[
  0\longrightarrow\mathfrak{I}/\mathfrak{I}\mathfrak{L}
  \longrightarrow\mathfrak{L}/\mathfrak{L}^2
  \longrightarrow(\mathfrak{L}/\mathfrak{I})/
    (\mathfrak{L}/\mathfrak{I})^2\longrightarrow0
\]
es exacta (véase la demostración de (16.2.7)), y los módulos que figuran en
esta sucesión son precisamente los tallos en $x$ de los haces de (16.9.13.1).
La segunda afirmación resulta de que $\mathcal{N}_{Y'/Y}$ es un
$\mathcal{O}_{Y'}$-módulo localmente \emph{libre} (16.9.8) y de Bourbaki,
\emph{Alg.}, cap.~II, 3\textsuperscript{a} ed., §~1, n\textsuperscript{o}~11,
prop.~21.

\subsection{Morfismos diferencialmente lisos}
\label{IV.16.10}

\begin{definition}[16.10.1]
\label{IV.16.10.1}
Se dice que un morfismo de preesquemas $f:X\to S$ es
\emph{diferencialmente liso} (o que $X$ es \emph{diferencialmente
liso sobre $S$}) si satisface las condiciones siguientes:
\begin{enumerate}
  \item[(i)] $\Omega_{X/S}^1$ es un $\mathcal{O}_X$-módulo localmente
  proyectivo, es decir, todo punto de $X$ admite una vecindad abierta
  afín $U$ tal que $\Gamma(U,\Omega_{X/S}^1)$ sea un
  $\Gamma(U,\mathcal{O}_X)$-módulo proyectivo (no necesariamente de tipo
  finito).
  \item[(ii)] El homomorfismo canónico (16.3.1.1)
  \[
    \mathbf{S}_{\mathcal{O}_X}^{\bullet}(\Omega_{X/S}^1)
    \longrightarrow\mathscr{G}r_{\bullet}(\mathscr{P}_{X/S})
  \]
  es biyectivo.
\end{enumerate}
En particular, si $\Omega_{X/S}^1$ es localmente libre de rango finito, los
$\mathscr{P}_{X/S}^n$ son $\mathcal{O}_X$-módulos localmente libres de
rango finito (por ser extensiones de tales módulos).
\end{definition}

Se dice que $f$ es \emph{diferencialmente liso en un punto $x\in X$} (o
que $X$ es diferencialmente liso sobre $S$ en el punto $x$) si existe una
vecindad abierta $U$ de $x$ en $X$ tal que $f|U$ sea diferencialmente
liso.

Veremos más adelante (17.12.4) que un morfismo liso es diferencialmente
liso, lo que justifica la terminología; pero el recíproco es falso; en
efecto, un \emph{monomorfismo} $f:X\to S$ es diferencialmente liso,
puesto que $\Omega_{X/S}^1=0$ en virtud de (I, 5.3.8), y por consiguiente
el homomorfismo sobreyectivo (16.3.1.1) es evidentemente biyectivo; ahora
bien, un monomorfismo ni siquiera es necesariamente plano, ni \emph{a fortiori}
liso. Limitémonos aquí a señalar la proposición siguiente:

\begin{proposition}[16.10.2]
\label{IV.16.10.2}
Sean $A$ un anillo, $B$ una $A$-álgebra formalmente lisa para las
topologías discretas (0, 19.3.1). Entonces $\operatorname{Spec}(B)$ es
diferencialmente liso sobre $\operatorname{Spec}(A)$.
\end{proposition}

En efecto, $B\otimes_A B$ es entonces (para las topologías discretas) una
$B$-álgebra formalmente lisa (para uno u otro de los homomorfismos
canónicos $b\mapsto b\otimes1$, $b\mapsto1\otimes b$ de $B$
\oldpage[IV-4]{52}
en $B\otimes_A B$) (0, 19.3.5, (iii)); por tanto, $B\otimes_A B$ es también
una $A$-álgebra formalmente lisa para las topologías discretas (0, 19.3.5,
(ii)). Poniendo $\mathfrak{I}=\mathfrak{I}_{B/A}$, se deduce que
$B\otimes_A B$ es también una $A$-álgebra formalmente lisa para la topología
$\mathfrak{I}$-preádica (0, 19.3.8); como, por hipótesis,
$B=(B\otimes_A B)/\mathfrak{I}$ es una $A$-álgebra formalmente lisa para
las topologías discretas, la proposición resulta de la equivalencia entre a)
y b) en (0, 19.5.4).

\begin{proposition}[16.10.3]
\label{IV.16.10.3}
Para que un morfismo $f:X\to S$ sea diferencialmente liso, es necesario y
suficiente que para todo $x\in X$ exista una vecindad abierta afín $U$ de
$x$, de anillo $A$, tal que $\Gamma(U,\mathscr{P}_{X/S}^{\infty})$ sea una
$A$-álgebra topológica aumentada isomorfa al álgebra completada $\widehat B$,
donde $B=\mathbf{S}_A(V)$, siendo $V$ un $A$-módulo proyectivo y estando
$B$ provista de la topología $B^+$-preádica (donde $B^+$ es el ideal de aumento). Si
$\Omega_{X/S}^1$ es localmente libre de rango finito, se puede sustituir
$\widehat B$ por el álgebra de series formales $A[[T_1,\ldots,T_n]]$.
\end{proposition}

La noción de morfismo diferencialmente liso es evidentemente local sobre
$X$, por lo que podemos reducirnos al caso en que $S=\operatorname{Spec}(B)$,
$X=\operatorname{Spec}(C)$. Consideremos $C\otimes_B C$ como una $C$-álgebra
(para el primer factor); pongamos $\mathfrak{I}=\mathfrak{I}_{C/B}$ y
dotemos $C\otimes_B C$ de la topología $\mathfrak{I}$-preádica; podemos
aplicar a la $C$-álgebra topológica $C\otimes_B C$ y al ideal
$\mathfrak{I}$ de $C\otimes_B C$ la equivalencia entre b) y c) en (0, 19.5.4),
puesto que $(C\otimes_B C)/\mathfrak{I}=C$ es evidentemente una $C$-álgebra
formalmente lisa para las topologías discretas. La topología sobre
$\Gamma(U,\mathscr{P}_{X/S}^{\infty})$ es evidentemente la topología de límite
proyectivo de este anillo (16.1.11).

Se observará que el entero $n$ del enunciado de (16.10.3) es el \emph{rango} de
$\Omega_{X/S}^1$ en el punto $x$. Veremos más adelante (17.13.5) que cuando
$f$ es diferencialmente liso y localmente de tipo finito, $n$ es también
igual a la dimensión de la fibra $f^{-1}(f(x))$ en el punto $x$.

\begin{proposition}[16.10.4]
\label{IV.16.10.4}
Sean $f:X\to S$, $g:S'\to S$ dos morfismos, y pongamos
$X'=X\times_S S'$, $f'=f_{(S')}:X'\to S'$.
\begin{enumerate}
  \item[(i)] Si $f$ es diferencialmente liso, lo mismo ocurre con $f'$.
  \item[(ii)] Recíprocamente, si $g$ es fielmente plano y cuasicompacto, y si
  $f'$ es diferencialmente liso y $\Omega_{X'/S'}^1$ es un
  $\mathcal{O}_{X'}$-módulo de tipo finito, $f$ es diferencialmente liso y
  $\Omega_{X/S}^1$ es un $\mathcal{O}_X$-módulo de tipo finito.
\end{enumerate}
\end{proposition}

En efecto, si $f$ es diferencialmente liso, los
$\mathscr{G}r_n(\mathscr{P}_{X/S})$ son $\mathcal{O}_X$-módulos
\emph{planos}; por consiguiente, en virtud de (16.4.6), el homomorfismo
\[
  \mathscr{G}r_n(\mathscr{P}_{X/S})\otimes_{\mathcal{O}_X}\mathcal{O}_{X'}
  \longrightarrow\mathscr{G}r_n(\mathscr{P}_{X'/S'})
\]
es biyectivo para todo $n$, y en virtud de la conmutatividad del diagrama
(16.2.1.3), se deduce de la definición (16.10.1) que $f'$ es
diferencialmente liso. Por otra parte, si $g$ es fielmente plano y
cuasicompacto, se deduce también de (16.4.6) que
\[
  \mathscr{G}r_n(\mathscr{P}_{X/S})\otimes_{\mathcal{O}_X}\mathcal{O}_{X'}
  \longrightarrow\mathscr{G}r_n(\mathscr{P}_{X'/S'})
\]
es biyectivo para todo $n$. Supongamos entonces que $f'$ es diferencialmente liso y
$\Omega_{X'/S'}^1$ tiene rango finito. Como la proyección canónica $X'\to X$ es
un morfismo fielmente plano y cuasicompacto, se deduce primero de (2.5.2)
que $\Omega_{X/S}^1$ es un $\mathcal{O}_X$-módulo localmente libre de rango
finito, y luego de (2.2.7) que el homomorfismo canónico (16.3.1.1) es biyectivo,
por lo que $f$ es diferencialmente liso.

\begin{proposition}[16.10.5]
\label{IV.16.10.5}
Para que un morfismo localmente de tipo finito $f:X\to S$ sea
diferencialmente liso, es necesario y suficiente que la inmersión diagonal
$\Delta_f:X\to X\times_S X$ sea cuasirregular.
\end{proposition}

\oldpage[IV-4]{53}

La cuestión es local, por lo que podemos reducirnos al caso en que $S$ y $X$
son afines, y por consiguiente el subpreesquema diagonal de $X\times_S X$ es
cerrado. La hipótesis de que $f$ es localmente de tipo finito implica que
$\Delta_f$ es localmente de presentación finita (I, 4.3.1), por lo que el
subpreesquema diagonal de $X\times_S X$ está definido por un ideal de tipo
finito $\mathscr{I}$, y $\Omega_{X/S}^1=\mathscr{I}/\mathscr{I}^2$ es un
$\mathcal{O}_X$-módulo de tipo finito. La proposición resulta entonces
inmediatamente de la comparación de las condiciones en (16.10.1) y (16.9.4).

\begin{remark}[16.10.6]
\label{IV.16.10.6}
Sea $f:X\to S$ un morfismo tal que el $\mathcal{O}_X$-módulo
$\Omega_{X/S}^1$ sea localmente libre de rango finito. Se deduce entonces de
(0, 20.4.7) que todo $x\in X$ posee una vecindad abierta $U$ tal que existe
una familia finita $(z_\lambda)_{\lambda\in L}$ de secciones de
$\mathcal{O}_X$ sobre $U$ para las cuales $(dz_\lambda)_{\lambda\in L}$
forma una \emph{base} del $\Gamma(U,\mathcal{O}_X)$-módulo
$\Gamma(U,\Omega_{X/S}^1)$.
\end{remark}

\subsection{Operadores diferenciales sobre un $S$-preesquema diferencialmente
liso}
\label{IV.16.11}

\begin{env}[16.11.1]
\label{IV.16.11.1}
Sean $f:X\to S$ un morfismo, $U$ un abierto de $X$,
$(z_\lambda)_{\lambda\in L}$ una familia de secciones de $\mathcal{O}_X$
sobre $U$ tal que los $dz_\lambda$ formen un sistema de generadores
de $\Omega_{X/S}^1|U=\Omega_{U/S}^1$. Sea $m$ un entero o el símbolo
$\infty$, y pongamos, para todo $\lambda$
\[
  \zeta_\lambda=\delta z_\lambda=d^m z_\lambda-z_\lambda
  \in\Gamma(U,\mathscr{P}_{X/S}^m).
  \tag{16.11.1.1}
  \label{IV.16.11.1.1}
\]

Usaremos también las notaciones usuales del análisis; para todo
$\mathbf{p}=(p_\lambda)\in\mathbb{N}^{(L)}$ (con $p_\lambda=0$ salvo para un
número finito de índices), pondremos
\[
  |\mathbf{p}|=\sum_\lambda p_\lambda,
  \qquad \mathbf{p}!=\prod_\lambda(p_\lambda!)
  \tag{16.11.1.2}
  \label{IV.16.11.1.2}
\]
\[
  \binom{\mathbf{p}}{\mathbf{q}}
  =\mathbf{p}!/(\mathbf{q}!(\mathbf{p}-\mathbf{q})!)
  \quad\text{para $\mathbf{p},\mathbf{q}$ en $\mathbb{N}^{(L)}$,
  $\mathbf{q}\leq\mathbf{p}$}
  \tag{16.11.1.3}
  \label{IV.16.11.1.3}
\]
y convenimos que $\binom{\mathbf{p}}{\mathbf{q}}=0$ si
$\mathbf{q}\nleq\mathbf{p}$,
\[
  \mathbf{z}^{\mathbf{p}}=\prod_\lambda(z_\lambda)^{p_\lambda},
  \qquad
  \boldsymbol{\zeta}^{\mathbf{p}}=\prod_\lambda(\zeta_\lambda)^{p_\lambda}.
  \tag{16.11.1.4}
  \label{IV.16.11.1.4}
\]

Por tanto, con estas notaciones
\[
  d^m(\mathbf{z}^{\mathbf{p}})
  =(d^m(\mathbf{z}))^{\mathbf{p}}
  =(\boldsymbol{\zeta}+\mathbf{z})^{\mathbf{p}}
  =\sum_{\mathbf{q}\leq\mathbf{p}}
    \binom{\mathbf{p}}{\mathbf{q}}
    \mathbf{z}^{\mathbf{p}-\mathbf{q}}\boldsymbol{\zeta}^{\mathbf{q}}
  \tag{16.11.1.5}
  \label{IV.16.11.1.5}
\]
\[
  \boldsymbol{\zeta}^{\mathbf{p}}
  =(d^m\mathbf{z}-\mathbf{z})^{\mathbf{p}}
  =\sum_{\mathbf{q}\leq\mathbf{p}}
    (-1)^{|\mathbf{p}-\mathbf{q}|}
    \binom{\mathbf{p}}{\mathbf{q}}
    \mathbf{z}^{\mathbf{p}-\mathbf{q}}d^m(\mathbf{z}^{\mathbf{q}}).
  \tag{16.11.1.6}
  \label{IV.16.11.1.6}
\]

Como los $dz_\lambda$ generan $\Omega_{X/S}^1$ y son las imágenes de
$\delta z_\lambda$, y puesto que el homomorfismo canónico (16.3.1.1) es
sobreyectivo, se concluye que para $m$ finito los $\delta z_\lambda$ generan
la $\mathcal{O}_U$-álgebra $\mathscr{P}_{U/S}^m$ (Bourbaki,
\emph{Alg. comm.}, cap.~III, §~2, n\textsuperscript{o}~8, cor.~2 del teorema~1).
Por tanto, los $\boldsymbol{\zeta}^{\mathbf{p}}$ (para $|\mathbf{p}|\leq m$)
generan el $\mathcal{O}_U$-módulo $\mathscr{P}_{U/S}^m$. Un operador
diferencial $D\in\operatorname{Diff}_{U/S}^m$ queda por consiguiente
completamente determinado por los valores de
$\langle\boldsymbol{\zeta}^{\mathbf{p}},D\rangle$ para
$|\mathbf{p}|\leq m$, o, lo que es equivalente en virtud de
(16.11.1.5) y (16.11.1.6), por los valores
\oldpage[IV-4]{54}
de $\langle d^m(\mathbf{z}^{\mathbf{p}}),D\rangle=D(\mathbf{z}^{\mathbf{p}})$
para $|\mathbf{p}|\leq m$; más precisamente, se deduce de (16.11.1.5) que
se tiene
\[
  D(\mathbf{z}^{\mathbf{p}})
  =\langle d^m(\mathbf{z}^{\mathbf{p}}),D\rangle
  =\sum_{\mathbf{q}\leq\mathbf{p}}
    \binom{\mathbf{p}}{\mathbf{q}}
    \langle\boldsymbol{\zeta}^{\mathbf{q}},D\rangle
    \mathbf{z}^{\mathbf{p}-\mathbf{q}}.
  \tag{16.11.1.7}
  \label{IV.16.11.1.7}
\]
\end{env}

\begin{theorem}[16.11.2]
\label{IV.16.11.2}
Sean $f:X\to S$ un morfismo, $U$ un abierto de $X$,
$(z_\lambda)_{\lambda\in L}$ una familia de secciones de $\mathcal{O}_X$
sobre $U$ tal que la familia $(dz_\lambda)_{\lambda\in L}$ genere
$\Omega_{X/S}^1|U=\Omega_{U/S}^1$. Las condiciones siguientes son
equivalentes:
\begin{enumerate}
  \item[a)] $f|U$ es diferencialmente liso y $(dz_\lambda)$ es una base
  del $\mathcal{O}_U$-módulo $\Omega_{U/S}^1$.
  \item[b)] Existe una familia
  $(D_{\mathbf{p}})_{\mathbf{p}\in\mathbb{N}^{(L)}}$ de operadores
  diferenciales de $\mathcal{O}_U$ en sí mismo que satisface las condiciones
  \[
    D_{\mathbf{p}}(\mathbf{z}^{\mathbf{q}})
    =\binom{\mathbf{q}}{\mathbf{p}}
      \mathbf{z}^{\mathbf{q}-\mathbf{p}}
    \quad(\mathbf{p},\mathbf{q}\text{ en }\mathbb{N}^{(L)}).
    \tag{16.11.2.1}
    \label{IV.16.11.2.1}
  \]
\end{enumerate}
Además, cuando estas condiciones se satisfacen, la familia
$(D_{\mathbf{p}})$ queda determinada de forma única por las condiciones
(16.11.2.1) y satisface las relaciones
\[
  D_{\mathbf{p}}\circ D_{\mathbf{q}}
  =D_{\mathbf{q}}\circ D_{\mathbf{p}}
  =\frac{(\mathbf{p}+\mathbf{q})!}{\mathbf{p}!\mathbf{q}!}
    D_{\mathbf{p}+\mathbf{q}}
  \quad(\mathbf{p},\mathbf{q}\text{ en }\mathbb{N}^{(L)}).
  \tag{16.11.2.2}
  \label{IV.16.11.2.2}
\]
Finalmente, si $L$ es finito, para todo entero $m$, los $D_{\mathbf{p}}$ tales que
$|\mathbf{p}|\leq m$ forman una base del $\mathcal{O}_U$-módulo
$\operatorname{Diff}_{U/S}^m$, es decir, todo operador diferencial
de orden $\leq m$ sobre $U$ se escribe de una sola manera bajo la forma
\[
  D=\sum_{|\mathbf{p}|\leq m}a_{\mathbf{p}}D_{\mathbf{p}}
\]
donde los $a_{\mathbf{p}}$ son secciones de $\mathcal{O}_X$ sobre
$U$.
\end{theorem}

Observemos primero que, en virtud de (16.11.1.6) y (16.11.1.5), se verifica
inmediatamente que las condiciones (16.11.2.1) son equivalentes a
\[
  \langle\boldsymbol{\zeta}^{\mathbf{p}},D_{\mathbf{q}}\rangle
  =\delta_{\mathbf{p}\mathbf{q}}
  \quad\text{(índice de Kronecker)}.
  \tag{16.11.2.3}
  \label{IV.16.11.2.3}
\]
La existencia de la familia $(D_{\mathbf{p}})$ que satisface estas relaciones
implica primero (tomando $|\mathbf{p}|=1$) que los $dz_\lambda$ son
linealmente independientes y, por tanto, forman una \emph{base} del
$\mathcal{O}_U$-módulo $\Omega_{U/S}^1$. Luego, para todo entero $m\geq1$, se
deduce del mismo modo de (16.11.2.3) que los $\boldsymbol{\zeta}^{\mathbf{p}}$
tales que $|\mathbf{p}|\leq m$ son linealmente independientes; por consiguiente
el homomorfismo canónico (16.3.1.1) es inyectivo, luego biyectivo, y esto
demuestra que b) implica a). La recíproca se deduce inmediatamente de la
definición (16.10.1), pues el hecho de que los $\boldsymbol{\zeta}^{\mathbf{p}}$
formen una base de $\mathscr{P}_{U/S}^m$ para $|\mathbf{p}|\leq m$ implica la
existencia y unicidad de una familia de homomorfismos
$u_{\mathbf{q},m}:\mathscr{P}_{U/S}^m\to\mathcal{O}_U$
($|\mathbf{q}|\leq m$) tales que
$\langle\boldsymbol{\zeta}^{\mathbf{p}},u_{\mathbf{q},m}\rangle
=\delta_{\mathbf{p}\mathbf{q}}$ para $|\mathbf{p}|\leq m$,
$|\mathbf{q}|\leq m$. Para un valor dado de $\mathbf{q}$, los operadores
diferenciales correspondientes a los $u_{\mathbf{q},m}$ para
$m\geq|\mathbf{q}|$ se identifican con un mismo operador $D_{\mathbf{q}}$.
Esto demuestra, pues, que a) implica b), y además que la familia
$(D_{\mathbf{p}})$ queda determinada de forma única y que, si $L$ es finito,
para $|\mathbf{p}|\leq m$, los $D_{\mathbf{p}}$ forman una base del dual
$\operatorname{Diff}_{U/S}^m$ de $\mathscr{P}_{U/S}^m$. Finalmente, las relaciones
(16.11.2.2) se deducen inmediatamente de la expresión de los valores de los tres
operadores considerados sobre los $\mathbf{z}^{\mathbf{r}}$, y del hecho de que
los $\boldsymbol{\zeta}^{\mathbf{r}}$ para $|\mathbf{r}|\leq m$ generan
$\mathscr{P}_{U/S}^m$.

\oldpage[IV-4]{55}

\begin{remarks}[16.11.3]
\label{IV.16.11.3}
\begin{enumerate}
  \item[(i)] El hecho de que los $D_{\mathbf{p}}$ sean permutables dos a dos en
  virtud de (16.11.2.2) no implica naturalmente que la
  $\mathcal{O}_U$-álgebra $\operatorname{Diff}_{U/S}$ sea conmutativa, pues los
  $D_{\mathbf{p}}$ no conmutan con la multiplicación por secciones de
  $\mathcal{O}_U$ salvo si $\mathbf{p}=0$.\footnote{El texto impreso dice
  $n=0$; debe ser $\mathbf{p}=0$, pues $D_{\mathbf{p}}$ es de orden
  $|\mathbf{p}|$.}
  \item[(ii)] Los índices $\mathbf{p}$ tales que $|\mathbf{p}|=1$ son los
  $\boldsymbol{\varepsilon}_\lambda
  =(\varepsilon_{\lambda\mu})_{\mu\in L}$ con
  $\varepsilon_{\lambda\mu}=0$ si $\mu\neq\lambda$ y
  $\varepsilon_{\lambda\lambda}=1$; cuando $L$ es finito, los operadores
  $D_{\boldsymbol{\varepsilon}_\lambda}$ no son otros que las
  $S$-derivaciones $D_i$ introducidas en (16.5.7). Obsérvese que, en general
  (y contrariamente a lo que ocurre en el análisis clásico), no es cierto que
  un operador diferencial de orden cualquiera pueda escribirse como combinación
  lineal de potencias de los $D_i$ (cf. (16.12)).
  \item[(iii)] Para todo entero $r\geq1$, se puede definir la noción de
  morfismo \emph{diferencialmente liso hasta el orden $r$} sustituyendo
  en (16.10.1) la condición (ii) por la condición de que los homomorfismos
  \[
    \mathbf{S}_{\mathcal{O}_X}^m(\Omega_{X/S}^1)
    \longrightarrow\mathscr{G}r_m(\mathscr{P}_{X/S})
  \]
  sean biyectivos \emph{para todo $m\leq r$}. El razonamiento de (16.11.2)
  demuestra entonces que si, en la condición a), se sustituye «diferencialmente
  liso» por «diferencialmente liso hasta el orden $r$», esta condición es
  equivalente a la condición b) en la que nos restringimos a los
  $\mathbf{p}\in\mathbb{N}^{(L)}$, $\mathbf{q}\in\mathbb{N}^{(L)}$ tales que
  $|\mathbf{p}|\leq r$, $|\mathbf{q}|\leq r$.
\end{enumerate}
\end{remarks}

\subsection{Caso de característica cero: criterio jacobiano para los
morfismos diferencialmente lisos}
\label{IV.16.12}

\begin{env}[16.12.1]
\label{IV.16.12.1}
Diremos que un preesquema $X$ es \emph{de característica $p$} ($p$ igual a $0$
o a un número primo) si, para todo abierto afín $U$ de $X$, el anillo
$\Gamma(U,\mathcal{O}_X)$ es de característica $p$ (0, 21.1.1). Se deduce
de (0, 21.1.3) que, para que $X$ sea de característica $0$, es necesario y
suficiente que, para todo punto \emph{cerrado} $x$ de $X$, el cuerpo residual
$k(x)$ sea de característica $0$, o, equivalentemente, que $X$ pueda estar
provisto de una estructura de preesquema sobre $\mathbb{Q}$ (necesariamente
única).
\end{env}

\begin{theorem}[16.12.2]
\label{IV.16.12.2}
Sean $X$ un preesquema de característica $0$ y $f:X\to S$ un morfismo. Si
$\Omega_{X/S}^1$ es un $\mathcal{O}_X$-módulo localmente libre (no
necesariamente de tipo finito), $f$ es diferencialmente liso.
\end{theorem}

La cuestión es local sobre $X$, así que podemos suponer que existe una familia
$(z_\lambda)$ de secciones de $\mathcal{O}_X$ sobre $X$ tal que $(dz_\lambda)$
sea una base del $\mathcal{O}_X$-módulo $\Omega_{X/S}^1$. Aplicando el criterio
(16.11.2), basta observar que los operadores
\[
  D_{\mathbf{p}}=(\mathbf{p}!)^{-1}
  \prod_\lambda D_\lambda^{p_\lambda}
\]
(donde los $D_\lambda$ son las formas coordenadas correspondientes a la base
$(dz_\lambda)$) verifican las relaciones (16.11.2.1), lo cual se deduce del
hecho de que los $D_\lambda$ son derivaciones.

\begin{env}[16.12.3]
\label{IV.16.12.3}
El teorema anterior deja de ser cierto cuando se abandona la hipótesis de que
$X$ es de característica $0$. Por ejemplo, si
$S=\operatorname{Spec}(k)$, donde $k$ es un cuerpo de característica $p>0$,
$X=\operatorname{Spec}(K)$, donde $K=k(\alpha)$ con
$\alpha\notin k$ y $\alpha^p\in k$, se comprueba inmediatamente
\oldpage[IV-4]{56}
que $\Omega_{X/S}^1$ es de rango $1$ y que el morfismo $X\to S$ es
diferencialmente liso hasta el orden $p-1$ (16.11.3, (iii)), pero no hasta
el orden $p$. Sin embargo, la demostración de (16.12.2) prueba que, si
$\Omega_{X/S}^1$ es localmente libre y $n!1_{\mathcal{O}_X}$ es invertible en
$\Gamma(X,\mathcal{O}_X)$, entonces $X$ es diferencialmente liso sobre $S$
hasta el orden $n$.
\end{env}


% Mechanical subsection components; prose is translated in the component files.
\section[Morfismos lisos, morfismos no ramificados (o netos), morfismos étale]{Morfismos lisos, morfismos no ramificados (o netos)\\morfismos étale}
\label{section:IV.17}

En este parágrafo retomamos las nociones estudiadas en (0, 19), expresadas mediante el lenguaje geométrico de los esquemas y desde el punto de vista global, para los preesquemas localmente de presentación finita sobre un preesquema base dado. La mayoría de los resultados (con excepción de los n\textsuperscript{os} 17.7, 17.8, 17.9, 17.13 y 17.16) se reducen, por lo demás, a variantes de propiedades ya encontradas en (0, 19). Para resultados más específicos sobre los morfismos étale, el lector consultará el § 18.

\subsection{Morfismos formalmente lisos, morfismos formalmente no ramificados, morfismos formalmente étale}
\label{subsection:IV.17.1}

\begin{definition}[17.1.1]
\label{IV.17.1.1}
Sea $f:X\to Y$ un morfismo de preesquemas. Se dice que $f$ es \emph{formalmente liso} (resp. \emph{formalmente no ramificado} o \emph{formalmente neto}, resp. \emph{formalmente étale}) si, para todo esquema afín $Y'$, todo subesquema cerrado $Y'_0$ de $Y'$ definido por un ideal nilpotente $\sh{J}$ de $\sh{O}_{Y'}$, y todo morfismo $Y'\to Y$, la aplicación
\[
\label{IV.17.1.1.1}
  \Hom_Y(Y',X)\longrightarrow\Hom_Y(Y'_0,X)
  \tag{17.1.1.1}
\]
deducida de la inyección canónica $Y'_0\to Y'$, es \emph{sobreyectiva} (resp. \emph{inyectiva}, resp. \emph{biyectiva}).
\end{definition}

Se dice también entonces que $X$ es \emph{formalmente liso} (resp. \emph{formalmente no ramificado} o \emph{formalmente neto}, resp. \emph{formalmente étale}) sobre $Y$.

Es claro que decir que $f$ es formalmente étale significa que es \emph{a la vez} formalmente liso y formalmente no ramificado.

\begin{remarks}[17.1.2]
\label{IV.17.1.2}
\begin{enumerate}
  \item[(i)] Supongamos que $Y=\Spec(A)$ y $X=\Spec(B)$ son afines, de modo que $f$ proviene de un homomorfismo de anillos $\varphi:A\to B$. En virtud de (0, 19.3.1 y 0, 19.10.1), decir que $f$ es formalmente liso (resp. formalmente no ramificado, resp. formalmente étale) significa que $\varphi$ hace de $B$ una $A$-álgebra \emph{formalmente lisa} (resp. \emph{formalmente no ramificada}, resp. \emph{formalmente étale}) para las topologías \emph{discretas} sobre $A$ y $B$.
  \item[(ii)] Para comprobar que $f$ es formalmente liso (resp. formalmente no ramificado, resp. formalmente étale), se puede, en la definición (17.1.1), restringirse al caso en que $\sh{J}^2=0$. En efecto, si $f$ satisface la condición correspondiente de la definición (17.1.1) en este caso particular, y si se tiene $\sh{J}^n=0$, se considera el subesquema cerrado $Y'_j$ de $Y'$ definido por el ideal $\sh{J}^{j+1}$ para $0\leq j\leq n-1$, de modo que $Y'_j$ es un subesquema cerrado de $Y'_{j+1}$ definido por un ideal de cuadrado cero; la hipótesis implica que cada una de las aplicaciones
  \[
    \Hom_Y(Y'_{j+1},X)\longrightarrow\Hom_Y(Y'_j,X)qquad(0\leq j\leq n-1)
  \]
  \oldpage[IV-4]{57}
  es sobreyectiva (resp. inyectiva, resp. biyectiva); por composición, se concluye que lo mismo ocurre con (17.1.1.1).
  \item[(iii)] Obsérvese que las propiedades del morfismo $f$ definidas en (17.1.1) son propiedades del \emph{funtor representable} $(0_{\mathrm{III}}, 8.1.8)$
  \[
    Y'\longmapsto\Hom_Y(Y',X)
  \]
  de la categoría de los $Y$-preesquemas en la categoría de los conjuntos; siguen teniendo sentido para \emph{cualquier} funtor contravariante con las mismas categorías de partida y de llegada, representable o no.
  \item[(iv)] Supongamos que el morfismo $f$ es formalmente no ramificado (resp. formalmente étale); consideremos un $Y$-preesquema \emph{cualquiera} $Z$ y un subpreesquema cerrado $Z_0$ de $Z$ definido por un ideal \emph{localmente nilpotente} $\sh{J}$ de $\sh{O}_Z$. Entonces la aplicación
  \[
  \label{IV.17.1.2.1}
    \Hom_Y(Z,X)\longrightarrow\Hom_Y(Z_0,X)
    \tag{17.1.2.1}
  \]
  deducida de la inyección canónica $Z_0\to Z$, sigue siendo inyectiva (resp. biyectiva). En efecto, sea $(U_\alpha)$ un recubrimiento abierto afín de $Z$ tal que los ideales $\sh{J}|U_\alpha$ sean nilpotentes, y para todo $\alpha$, sea $U^0_\alpha$ la imagen inversa de $U_\alpha$ en $Z_0$, que es el subesquema cerrado de $U_\alpha$ definido por $\sh{J}|U_\alpha$. Sea $f_0:Z_0\to X$ un $Y$-morfismo; por hipótesis, para todo $\alpha$, existe a lo sumo un (resp. uno y solo uno) $Y$-morfismo $f_\alpha:U_\alpha\to X$ cuya restricción a $U^0_\alpha$ sea igual a $f_0|U^0_\alpha$. Se concluye inmediatamente que, si $f_\alpha$ y $f_\beta$ están definidos, entonces, para todo abierto afín $V\subset U_\alpha\cap U_\beta$, se tiene $f_\alpha|V=f_\beta|V$, pues las restricciones de estos morfismos a la imagen inversa $V_0$ de $V$ en $Z_0$ coinciden. Por tanto, existe a lo sumo un (resp. uno y solo uno) $Y$-morfismo $f:Z\to X$ cuya restricción a $Z_0$ coincide con $f_0$.
\end{enumerate}
\end{remarks}

\begin{proposition}[17.1.3]
\label{IV.17.1.3}
\begin{enumerate}
  \item[(i)] Un monomorfismo de preesquemas es formalmente no ramificado; una inmersión abierta es formalmente étale.
  \item[(ii)] La composición de dos morfismos formalmente lisos (resp. formalmente no ramificados, resp. formalmente étale) es formalmente lisa (resp. formalmente no ramificada, resp. formalmente étale).
  \item[(iii)] Si $f:X\to Y$ es un $S$-morfismo formalmente liso (resp. formalmente no ramificado, resp. formalmente étale), lo mismo ocurre con $f_{(S')}:X_{(S')}\to Y_{(S')}$ para todo cambio $S'\to S$ del preesquema de base.
  \item[(iv)] Si $f:X\to X'$ y $g:Y\to Y'$ son dos $S$-morfismos formalmente lisos (resp. formalmente no ramificados, resp. formalmente étale), lo mismo ocurre con $f\times_S g:X\times_S Y\to X'\times_S Y'$.
  \item[(v)] Sean $f:X\to Y$, $g:Y\to Z$ dos morfismos; si $g\circ f$ es formalmente no ramificado, lo mismo ocurre con $f$.
  \item[(vi)] Si $f:X\to Y$ es un morfismo formalmente no ramificado, lo mismo ocurre con $f_{\mathrm{red}}:X_{\mathrm{red}}\to Y_{\mathrm{red}}$.
\end{enumerate}
\end{proposition}

En virtud de (I, 5.5.12), basta demostrar (i), (ii) y (iii). Las dos afirmaciones de (i) son triviales. Para demostrar (ii), consideremos dos morfismos $f:X\to Y$, $g:Y\to Z$, un esquema afín $Z'$, un subesquema cerrado $Z'_0$ de $Z'$ definido por un ideal nilpotente y un morfismo $Z'\to Z$. Supongamos que $f$ y $g$ son formalmente lisos, y consideremos un
\oldpage[IV-4]{58}
$Z$-morfismo $u_0:Z'_0\to X$; la hipótesis sobre $g$ implica que existe un $Z$-morfismo $v:Z'\to Y$ tal que $f\circ u_0=v\circ j$ (donde $j:Z'_0\to Z'$ es la inyección canónica); la hipótesis sobre $f$ implica entonces que existe un morfismo $u:Z'\to X$ tal que $f\circ u=v$ y $u\circ j=u_0$, por tanto $(g\circ f)\circ u$ es igual al morfismo dado $Z'\to Z$ y $u\circ j=u_0$, lo que demuestra que $g\circ f$ es formalmente liso; se razona del mismo modo cuando se supone que $f$ y $g$ son formalmente no ramificados.

Por último, para demostrar (iii), pongamos $X'=X_{(S')}$, $Y'=Y_{(S')}$, $f'=f_{(S')}$; consideremos un esquema afín $Y''$, un subesquema cerrado $Y''_0$ de $Y''$ definido por un ideal nilpotente y un morfismo $g:Y''\to Y'$ que hace de $Y''$ un $Y'$-preesquema; se sabe entonces por (I, 3.3.8) que $\Hom_{Y'}(Y'',X')$ se identifica canónicamente con $\Hom_Y(Y'',X)$ y $\Hom_{Y'}(Y''_0,X')$ con $\Hom_Y(Y''_0,X)$, y la conclusión resulta entonces inmediatamente de la definición (17.1.1).

Obsérvese que una \emph{inmersión cerrada} no es necesariamente un morfismo formalmente liso.

\begin{proposition}[17.1.4]
\label{IV.17.1.4}
Sean $f:X\to Y$, $g:Y\to Z$ dos morfismos, y supongamos que $g$ es formalmente no ramificado. Entonces, si $g\circ f$ es formalmente liso (resp. formalmente étale), lo mismo ocurre con $f$.
\end{proposition}

En efecto, sean $Y'$ un esquema afín, $Y'_0$ un subesquema cerrado de $Y'$ definido por un ideal nilpotente, $h:Y'\to Y$ un morfismo, $j:Y'_0\to Y'$ la inyección canónica, $u_0:Y'_0\to X$ un $Y$-morfismo, de modo que $f\circ u_0=h\circ j$. Supongamos que $g\circ f$ es formalmente liso; entonces existe un morfismo $u:Y'\to X$ tal que $u\circ j=u_0$ y $(g\circ f)\circ u=g\circ h$. Pero estas relaciones implican que $f\circ u$ y $h$ son dos $Z$-morfismos de $Y'$ en $Y$ tales que $(f\circ u)\circ j=h\circ j$; en virtud de la hipótesis de que $g$ es formalmente no ramificado, se deduce que $f\circ u=h$, es decir, $u$ es un $Y$-morfismo; por tanto, $f$ es formalmente liso. Teniendo en cuenta (17.1.3, (v)), esto demuestra la proposición.

\begin{corollary}[17.1.5]
\label{IV.17.1.5}
Supongamos que $g$ es formalmente étale; entonces, para que $g\circ f$ sea formalmente liso (resp. formalmente no ramificado, resp. formalmente étale), es necesario y suficiente que $f$ lo sea.
\end{corollary}

Esto resulta de (17.1.4) y de (17.1.3, (ii) y (v)).

\begin{proposition}[17.1.6]
\label{IV.17.1.6}
Sea $f:X\to Y$ un morfismo de preesquemas.
\begin{enumerate}
  \item[(i)] Sea $(U_\alpha)$ un recubrimiento abierto de $X$ y, para todo $\alpha$, sea $i_\alpha:U_\alpha\to X$ la inyección canónica. Para que $f$ sea formalmente liso (resp. formalmente no ramificado, resp. formalmente étale), es necesario y suficiente que cada uno de los morfismos $f\circ i_\alpha$ lo sea.
  \item[(ii)] Sea $(V_\lambda)$ un recubrimiento abierto de $Y$. Para que $f$ sea formalmente liso (resp. formalmente no ramificado, resp. formalmente étale), es necesario y suficiente que cada una de las restricciones $f^{-1}(V_\lambda)\to V_\lambda$ de $f$ lo sea.
\end{enumerate}
\end{proposition}

Observemos primero que (ii) es consecuencia de (i): en efecto, si $j_\lambda:V_\lambda\to Y$ e $i_\lambda:f^{-1}(V_\lambda)\to X$ son las inyecciones canónicas, la restricción $f_\lambda:f^{-1}(V_\lambda)\to V_\lambda$ de $f$ es tal que $j_\lambda\circ f_\lambda=f\circ i_\lambda$; si $f$ es formalmente liso (resp. formalmente no ramificado), lo mismo ocurre con $f\circ i_\lambda$, pues $i_\lambda$ es formalmente étale (17.1.3); pero, como $j_\lambda$ es formalmente étale, esto implica que $f_\lambda$ es formalmente liso (resp. formalmente no ramificado) en virtud de (17.1.5). Recíprocamente, si todos los $f_\lambda$ son formalmente lisos (resp. formalmente no ramificados), lo mismo ocurre con los $j_\lambda\circ f_\lambda$ (17.1.3), y por tanto también con $f$ en virtud de (i).

\oldpage[IV-4]{59}
Teniendo en cuenta que los $i_\alpha$ son formalmente étale, todo se reduce, pues, a demostrar que si los $f\circ i_\alpha$ son formalmente lisos (resp. formalmente no ramificados), lo mismo ocurre con $f$.

Sean, pues, $Y'$ un esquema afín, $Y'_0$ un subesquema cerrado de $Y'$ definido por un ideal nilpotente $\sh{J}$, que se puede suponer tal que $\sh{J}^2=0$ (17.1.2, (ii)), y finalmente sea $g:Y'\to Y$ un morfismo. Supongamos dado un $Y$-morfismo $u_0:Y'_0\to X$; designemos por $W_\alpha$ (resp. $W^0_\alpha$) el preesquema inducido por $Y'$ (resp. $Y'_0$) sobre el abierto $u_0^{-1}(U_\alpha)$ (recuérdese que $Y'$ e $Y'_0$ tienen \emph{el mismo espacio topológico subyacente}). Supongamos primero que los $f\circ i_\alpha$ son \emph{formalmente no ramificados}, y demostremos que, si $u'$ y $u''$ son dos $Y$-morfismos de $Y'$ en $X$ cuyas restricciones a $Y'_0$ coinciden, entonces se tiene $u'=u''$. En efecto, teniendo en cuenta (17.1.2, (iv)), la hipótesis de que los $f\circ i_\alpha$ son no ramificados implica que, para todo $\alpha$, se tiene $u'|W_\alpha=u''|W_\alpha$, pues las restricciones de estos dos $Y$-morfismos a $W^0_\alpha$ coinciden. De aquí se sigue la conclusión en este caso.

Supongamos ahora que todos los $f\circ i_\alpha$ son \emph{formalmente lisos} y demostremos que existe un $Y$-morfismo $u:Y'\to X$ del cual $u_0$ es la restricción a $Y'_0$. Ahora bien, puesto que $Y'$ es un \emph{esquema afín}, se puede aplicar (16.5.17), cuyas hipótesis se satisfacen y cuya conclusión demuestra precisamente la existencia de $u$.

Por tanto, se puede decir que las nociones introducidas en (17.1.1) son \emph{locales} sobre $X$ y sobre $Y$, lo que permite siempre, en virtud de (17.1.2, (i)), reducirse al estudio de las \emph{álgebras} formalmente lisas (resp. formalmente no ramificadas, resp. formalmente étale).

\subsection{Propiedades generales de los diferenciales}
\label{subsection:IV.17.2}

\begin{proposition}[17.2.1]
\label{IV.17.2.1}
Para que un morfismo $f:X\to Y$ sea formalmente no ramificado, es necesario y suficiente que $\Omega_f^1=0$ (lo que también se escribe $\Omega_{X/Y}^1=0$ (16.3.1)).
\end{proposition}

Teniendo en cuenta (17.1.6), se reduce al caso en que $Y=\Spec(A)$ y $X=\Spec(B)$ son afines, y la conclusión resulta entonces de (0, 20.7.4) y de la interpretación de $\Omega_{X/Y}^1$ en este caso (16.3.7).

\begin{corollary}[17.2.2]
\label{IV.17.2.2}
Sean $f:X\to Y$, $g:Y\to Z$ dos morfismos. Para que $f$ sea formalmente no ramificado, es necesario y suficiente que el homomorfismo canónico (16.4.19)
\[
  f^*(\Omega_{Y/Z}^1)\longrightarrow\Omega_{X/Z}^1
\]
sea sobreyectivo.
\end{corollary}

Es una consecuencia inmediata de (17.2.1) y de la sucesión exacta (16.4.19.1).

\begin{proposition}[17.2.3]
\label{IV.17.2.3}
Sea $f:X\to Y$ un morfismo formalmente liso.
\begin{enumerate}
  \item[(i)] El $\sh{O}_X$-módulo $\Omega_{X/Y}^1$ es localmente proyectivo (16.10.1). Si $f$ es localmente de tipo finito, $\Omega_{X/Y}^1$ es localmente libre de tipo finito.
  \item[(ii)] Para todo morfismo $g:Y\to Z$, la sucesión (16.4.19) de $\sh{O}_X$-módulos
  \[
  \label{IV.17.2.3.1}
    0\longrightarrow f^*(\Omega_{Y/Z}^1)\longrightarrow\Omega_{X/Z}^1\longrightarrow\Omega_{X/Y}^1\longrightarrow0
    \tag{17.2.3.1}
  \]
  es exacta; además, para todo $x\in X$, existe una vecindad abierta $U$ de $x$ tal que las restricciones a $U$ de los homomorfismos de (17.2.3.1) forman una sucesión exacta escindida.
\end{enumerate}
\end{proposition}

\oldpage[IV-4]{60}
\begin{enumerate}
  \item[(i)] Se sabe (16.3.9) que si $f$ es localmente de tipo finito, $\Omega_f^1$ es un $\sh{O}_X$-módulo de tipo finito. Para demostrar que, en todos los casos, es localmente proyectivo, se puede restringir, en virtud de (17.1.6), al caso en que $Y=\Spec(A)$ y $X=\Spec(B)$ son afines, y esto resulta de la hipótesis sobre $f$ y de (0, 20.4.9 y 0, 19.2.1).
  \item[(ii)] También aquí se puede restringir al caso en que $X$, $Y$ y $Z$ son afines (17.1.6), y la conclusión resulta en este caso de la interpretación de los módulos que figuran en la sucesión (17.2.3.1) y de (0, 20.5.7).
\end{enumerate}

\begin{corollary}[17.2.4]
\label{IV.17.2.4}
Si $f:X\to Y$ es un morfismo formalmente étale, entonces, para todo morfismo $g:Y\to Z$, el homomorfismo canónico de $\sh{O}_X$-módulos
\[
  f^*(\Omega_{Y/Z}^1)\longrightarrow\Omega_{X/Z}^1
\]
es biyectivo.
\end{corollary}

Esto resulta de la exactitud de la sucesión (17.2.3.1) y del hecho de que entonces se tiene $\Omega_{X/Y}^1=0$ (17.2.1).

\begin{proposition}[17.2.5]
\label{IV.17.2.5}
Sean $f:X\to Y$ un morfismo, $X'$ un subpreesquema de $X$ tal que el morfismo compuesto $X'\xrightarrow{j}X\to Y$ (donde $j$ es la inyección canónica) sea formalmente liso. Entonces la sucesión de $\sh{O}_{X'}$-módulos (16.4.21)
\[
\label{IV.17.2.5.1}
  0\longrightarrow\mathcal{N}_{X'/X}\longrightarrow\Omega_{X/Y}^1\otimes_{\sh{O}_X}\sh{O}_{X'}\longrightarrow\Omega_{X'/Y}^1\longrightarrow0
  \tag{17.2.5.1}
\]
es exacta; además, para todo $x\in X$, existe una vecindad abierta $U$ de $x$ tal que las restricciones a $U$ de los homomorfismos de (17.2.5.1) forman una sucesión exacta escindida.
\end{proposition}

De nuevo en virtud de (17.1.6), se puede restringir al caso en que $Y=\Spec(A)$ y $X=\Spec(B)$ son afines, y $X'=\Spec(B/\mathfrak{J})$, donde $\mathfrak{J}$ es un ideal de $B$. Entonces el haz conormal $\mathcal{N}_{X'/X}$ corresponde al $B$-módulo $\mathfrak{J}/\mathfrak{J}^2$ (16.1.3), y la conclusión se deduce de (0, 20.5.14).

\begin{proposition}[17.2.6]
\label{IV.17.2.6}
Sean $X$, $Y$ dos preesquemas, $f:X\to Y$ un morfismo localmente de tipo finito. Las condiciones siguientes son equivalentes:
\begin{enumerate}
  \item[a)] $f$ es un monomorfismo.
  \item[b)] $f$ es radicial y formalmente no ramificado.
  \item[c)] Para todo $y\in Y$, la fibra $f^{-1}(y)$ es vacía o $k(y)$-isomorfa a $\Spec(k(y))$ (es decir, está reducida a un solo punto $z$ tal que $k(y)\to\sh{O}_z/\mathfrak{m}_y\sh{O}_z$ sea un isomorfismo).
\end{enumerate}
\end{proposition}

El hecho de que a) implique c) resulta de (8.11.5.1). Es claro que c) implica que $f$ es radicial; demostremos que de c) también resulta que $\Omega_{X/Y}^1=0$, lo que probará que c) implica b) (17.2.1). Obsérvese que el $\sh{O}_X$-módulo $\Omega_{X/Y}^1$ es cuasicoherente de tipo finito (16.3.9). Resulta, pues, de (I, 9.1.13.1) que, para que $(\Omega_{X/Y}^1)_x=0$, es necesario y suficiente que, si se pone $Y_1=\Spec(k(y))$, $X_1=f^{-1}(y)=X\times_Y Y_1$, se tenga $(\Omega_{X_1/Y_1}^1)_x=0$; pero, como el morfismo $f_1:X_1\to Y_1$ deducido de $f$ es formalmente no ramificado en virtud de la hipótesis c) (17.1.3), la conclusión resulta de (17.2.1). Demostremos finalmente que b) implica a); para ello, consideremos el morfismo diagonal $g=\Delta_f:X\to X\times_Y X$; puesto que $f$ es radicial, $g$ es sobreyectivo (I, 8.7.1); por otra parte, $\Omega_{X/Y}^1$ es por definición el haz conormal $\mathop{Gr}\nolimits_1(g)$ de la inmersión $g$ (16.3.1), y decir que $f$ es formalmente no ramificado significa, por tanto, que

\oldpage[IV-4]{61}
$\mathop{Gr}\nolimits_1(g)=0$ (17.2.1). Además, $g$ es localmente de presentación finita (I, 4.3.1); por tanto, la hipótesis $\mathop{Gr}\nolimits_1(g)=0$ implica que $g$ es una inmersión abierta (16.1.10); al ser sobreyectiva, esta inmersión es un isomorfismo, por lo que $f$ es un monomorfismo (I, 5.3.8).

\subsection{Morfismos lisos, morfismos no ramificados, morfismos étale}
\label{subsection:IV.17.3}

\begin{definition}[17.3.1]
\label{IV.17.3.1}
Se dice que un morfismo $f:X\to Y$ es \emph{liso} (resp. \emph{no ramificado}, o \emph{neto}\footnote{Las palabras «neto» y «formalmente neto» parecen claramente preferibles a la terminología consagrada de «no ramificado» (resp. «formalmente no ramificado») y se emplearán casi exclusivamente a partir del cap. V. En este capítulo hemos conservado la terminología antigua para no entrar en conflicto con (0, 19.10).}, resp. \emph{étale}) si es localmente de presentación finita y formalmente liso (resp. formalmente no ramificado, resp. formalmente étale).
\end{definition}

Se dice también entonces que $X$ es \emph{liso} (resp. \emph{no ramificado} o \emph{neto}, resp. \emph{étale}) sobre $Y$.

Veremos más adelante (17.5.2) que esta definición de morfismo liso coincide con la ya dada en (6.8.1); hasta entonces, utilizaremos exclusivamente la definición de (17.3.1).

Es claro que decir que $f$ es étale significa que es \emph{a la vez} liso y no ramificado.

\begin{remarks}[17.3.2]
\label{IV.17.3.2}
\begin{enumerate}
  \item[(i)] Obsérvese que la definición (17.3.1) puede expresarse únicamente mediante el funtor
  \[
    Y'\longmapsto\Hom_Y(Y',X)
  \]
  considerado en (17.1.2, (iii)), pues decir que $f$ es localmente de presentación finita equivale a decir que el funtor anterior \emph{conmuta con los límites proyectivos de esquemas afines} (8.14.2).
  \item[(ii)] Sean $A$ un anillo, $B$ una $A$-álgebra. Se dice que $B$ es una $A$-álgebra \emph{lisa} (resp. \emph{no ramificada}, resp. \emph{étale}) si el morfismo correspondiente $\Spec(B)\to\Spec(A)$ es liso (resp. no ramificado, resp. étale). Equivale a decir que $B$ es una $A$-álgebra \emph{de presentación finita} (I, 4.4.6) y formalmente lisa (resp. formalmente no ramificada, resp. formalmente étale) para las topologías discretas.
  \item[(iii)] De (17.1.6) y de la definición de morfismo localmente de presentación finita (I, 4.4.2) resulta que las nociones de morfismo liso, no ramificado y étale son \emph{locales sobre $X$ y sobre $Y$}.
\end{enumerate}
\end{remarks}

\begin{proposition}[17.3.3]
\label{IV.17.3.3}
\begin{enumerate}
  \item[(i)] Una inmersión abierta es étale. Para que una inmersión sea no ramificada, es necesario y suficiente que sea localmente de presentación finita.
  \item[(ii)] La composición de dos morfismos lisos (resp. no ramificados, resp. étale) es lisa (resp. no ramificada, resp. étale).
  \item[(iii)] Si $f:X\to Y$ es un $S$-morfismo liso (resp. no ramificado, resp. étale), lo mismo ocurre con $f_{(S')}:X_{(S')}\to Y_{(S')}$ para toda extensión $S'\to S$ del preesquema base.
  \item[(iv)] Si $f:X\to X'$ y $g:Y\to Y'$ son dos $S$-morfismos lisos (resp. no ramificados, resp. étale), lo mismo ocurre con $f\times_S g:X\times_S Y\to X'\times_S Y'$.
\end{enumerate}

\oldpage[IV-4]{62}

\begin{enumerate}
  \item[(v)] Sean $f:X\to Y$, $g:Y\to Z$ dos morfismos; si $g$ es localmente de tipo finito y si $g\circ f$ es no ramificado, entonces $f$ es no ramificado.
\end{enumerate}
\end{proposition}

Esto resulta inmediatamente de (I, 4.3) y de (17.1.3).

\begin{proposition}[17.3.4]
\label{IV.17.3.4}
Sean $f:X\to Y$, $g:Y\to Z$ dos morfismos, y supongamos que $g$ es no ramificado. Entonces, si $g\circ f$ es liso (resp. no ramificado, resp. étale), lo mismo ocurre con $f$.
\end{proposition}

En efecto, como $g$ y $g\circ f$ son localmente de presentación finita, lo mismo ocurre con $f$ (I, 4.4.3, (v)); la conclusión resulta, pues, de (17.1.4) y (17.1.3, (v)).

\begin{corollary}[17.3.5]
\label{IV.17.3.5}
Supongamos que $g$ es étale; entonces, para que $f$ sea liso (resp. no ramificado, resp. étale), es necesario y suficiente que $g\circ f$ lo sea.
\end{corollary}

Esto resulta de (17.3.4) y de (17.3.3, (iii)).

\begin{proposition}[17.3.6]
\label{IV.17.3.6}
Sean $g:Y\to S$, $h:X\to S$ dos morfismos localmente de presentación finita. Para que un $S$-morfismo $f:X\to Y$ sea no ramificado, es necesario y suficiente que el homomorfismo canónico (16.4.19)
\[
  f^*(\Omega_{Y/S}^1)\longrightarrow\Omega_{X/S}^1
\]
sea sobreyectivo.
\end{proposition}

Como $f$ es entonces localmente de presentación finita (I, 4.4.3, (v)), la proposición resulta inmediatamente de (17.2.2).

\begin{definition}[17.3.7]
\label{IV.17.3.7}
Sea $f:X\to Y$ un morfismo. Se dice que $f$ es \emph{liso} (resp. \emph{no ramificado}, resp. \emph{étale}) \emph{en un punto} $x\in X$ si existe una vecindad abierta $U$ de $x$ en $X$ tal que la restricción $f|U$ sea un morfismo liso (resp. no ramificado, resp. étale) de $U$ en $Y$.
\end{definition}

Se dice también entonces que $X$ es liso (resp. no ramificado, resp. étale) sobre $Y$ en el punto $x$.

Teniendo en cuenta la observación (17.3.2, (iii)), equivale a decir que $f$ es un morfismo liso (resp. no ramificado, resp. étale) o que es liso (resp. no ramificado, resp. étale) en cada punto de $X$.

Es claro que el conjunto de los puntos de $X$ en los que un morfismo $f:X\to Y$ es liso (resp. no ramificado, resp. étale) es \emph{abierto} en $X$.

\begin{proposition}[17.3.8]
\label{IV.17.3.8}
Para todo preesquema $Y$ y todo $\sh{O}_Y$-módulo localmente libre $\sh{E}$ de tipo finito, el preesquema fibrado vectorial $\mathbf{V}(\sh{E})$ (II, 1.7.8) es un $Y$-preesquema liso.
\end{proposition}

En efecto (17.3.2, (iii)), se puede restringir al caso en que $Y=\Spec(A)$ es afín y $\mathbf{V}(\sh{E})=\Spec(A[T_1,\ldots,T_r])$; como $A[T_1,\ldots,T_r]$ es una $A$-álgebra formalmente lisa para las topologías discretas (0, 19.3.2) y de presentación finita, esto demuestra la proposición (17.3.2, (ii)).

\begin{corollary}[17.3.9]
\label{IV.17.3.9}
Bajo las hipótesis de (17.3.8), el preesquema fibrado proyectivo $\mathbf{P}(\sh{E})$ (II, 4.1.1) es un $Y$-preesquema liso.
\end{corollary}

También se puede restringir al caso en que $Y=\Spec(A)$ es afín y $\mathbf{P}(\sh{E})=\mathbf{P}^r_Y$. Se sabe entonces (II, 2.3.14) que se tiene un recubrimiento abierto finito de $\mathbf{P}^r_A$ tomando los $D_+(T_i)$ ($0\leq i\leq r$), iguales respectivamente a los espectros de los anillos $S_{(f)}$, donde se sustituye $S$ por $A[T_0,T_1,\ldots,T_r]$ y $f$ por $T_i$; pero de la definición de $S_{(f)}$ (II, 2.2.1) resulta inmediatamente que este anillo, en el caso considerado, es isomorfo a $A[T_0,\ldots,T_{i-1},T_{i+1},\ldots,T_r]$; por tanto, el corolario resulta de (17.3.8).

\oldpage[IV-4]{63}


\subsection{Caracterizaciones de los morfismos no ramificados}
\label{subsection:IV.17.4}

\begin{theorem}[17.4.1]
\label{IV.17.4.1}
Sean $f:X\to Y$ un morfismo localmente de presentación finita y $x$ un punto de $X$. Las propiedades siguientes son equivalentes:
\begin{enumerate}
  \item[a)] $f$ es no ramificado en el punto $x$.
  \item[b)] El morfismo diagonal $\Delta_f:X\to X\times_Y X$ es un isomorfismo local en el punto $x$.
  \item[b')] Si se pone $\Delta_f=(\psi,\theta)$, $Z=X\times_Y X$ y $z=\psi(x)$, el homomorfismo $\theta_z^\#:\sh{O}_{Z,z}\to\sh{O}_{X,x}$ es biyectivo.
  \item[b'')] Para todo morfismo $g:Y'\to Y$, y todo punto $y'\in Y'$ situado sobre $y=f(x)$, toda $Y'$-sección $s'$ de $X'=X\times_Y Y'$ tal que $x'=s'(y')$ esté situado sobre $x$ es un isomorfismo local en el punto $y'$.
  \item[c)] Se tiene $(\Omega^1_{X/Y})_x=0$.
  \item[d)] El $k(y)$-preesquema $f^{-1}(y)$ es no ramificado sobre $k(y)$ en el punto $x$.
  \item[d')] El punto $x$ está aislado en $X_y=f^{-1}(y)$ (es decir (Err$_{\mathrm{III}}$, 20), el morfismo $f$ es cuasifinito en el punto $x$) y el anillo $\sh{O}_{X_y,x}$ es un cuerpo, extensión separable de $k(y)$.
  \item[d'')] El anillo $\sh{O}_{X_y,x}=\sh{O}_{X,x}/\mathfrak{m}_y\sh{O}_{X,x}$ es un cuerpo, extensión finita separable de $k(y)$.
  \item[e)] El anillo $\sh{O}_{X,x}$ es una $\sh{O}_{Y,y}$-álgebra formalmente no ramificada para las topologías discretas.
\end{enumerate}
\end{theorem}

Como $f$ es localmente de tipo finito, el $\sh{O}_X$-módulo $\Omega^1_{X/Y}$ es de tipo finito (16.3.9), por lo que equivale a decir que $(\Omega^1_{X/Y})_x=0$ o que existe una vecindad abierta $U$ de $x$ tal que $\Omega^1_{X/Y}|U=0$. Teniendo en cuenta (17.2.1), esto demuestra la equivalencia de a) y c). Por otra parte, si se pone $A=\sh{O}_{Y,y}$, $B=\sh{O}_{X,x}$, se tiene $(\Omega^1_{X/Y})_x=\Omega^1_{B/A}$ (16.4.15), y la equivalencia de c) y e) resulta, pues, de (0, 20.7.4).

Como d') solo concierne a propiedades del morfismo $X_y\to\Spec(k(y))$, la equivalencia de a) y d') implicará \emph{ipso facto} la de d) y d'). Por otra parte, d') y d'') son equivalentes, pues equivale a decir que $\sh{O}_{X_y,x}$ es una $k(y)$-álgebra \emph{finita} o que $x$ es un punto aislado de $X_y$, ya que $X_y$ es un $k(y)$-preesquema localmente de tipo finito (I, 6.4.4).

Demostremos ahora la equivalencia de b) y b'). Se puede restringir al caso en que $Y=\Spec(R)$ y $X=\Spec(S)$ son afines y $f$ es de presentación finita; entonces se tiene $Z=\Spec(S\otimes_R S)$ y $\Delta_f$ corresponde al homomorfismo canónico sobreyectivo $S\otimes_R S\to S$, cuyo núcleo $\mathfrak{J}$ es, como se sabe, un ideal de tipo finito (0, 20.4.4). Si se pone $\sh{J}=\widetilde{\mathfrak{J}}$, el $\sh{O}_Z$-módulo $\psi_*(\sh{O}_X)=\sh{O}_Z/\sh{J}$ es, por tanto, de presentación finita, y la hipótesis de que el homomorfismo $\theta_z^\#: \sh{O}_{Z,z}\to(\psi_*(\sh{O}_X))_z$ es biyectivo implica que, reemplazando si es preciso $X$ por una vecindad abierta de $x$, el propio homomorfismo $\theta:\sh{O}_Z\to\psi_*(\sh{O}_X)$ es biyectivo ($0_{\mathrm I}$, 5.2.7). Esto demuestra que b') implica b); el recíproco es evidente.

Por otra parte, la equivalencia de b) y b'') resulta de (I, 5.3.7) sin hipótesis de finitud sobre $f$: dar una $Y'$-sección $s':Y'\to X'$ equivale, en efecto, a dar un

\oldpage[IV-4]{64}
$Y$-morfismo $h=g'\circ s':Y'\to X$ (donde $g':X'\to X$ es la proyección canónica), de modo que $s'=(1_{Y'},h)_X$, y entonces el diagrama
\[
\tag{17.4.1.1}
\label{IV.17.4.1.1}
\xymatrix{
Y' \ar[r]^{s'} \ar[d]_h & X'=Y'\times_Y X \ar[d]^{h\times 1_X}\\
X \ar[r]_{\Delta_f} & X\times_Y X
}
\]
identifica $Y'$ con el producto de los $(X\times_Y X)$-preesquemas $X$ y $X'$. Por consiguiente (I, 4.3.2), si $\Delta_f$ es un isomorfismo local en el punto $x$, $s'$ es un isomorfismo local en el punto $y'$ (pues $x=h(y')$), lo que demuestra que b) implica b''). El recíproco se obtiene aplicando b'') al caso en que se toman $Y'=X$, $y'=x$, $g=f$ y $s'=\Delta_f$.

Para terminar la demostración de (17.4.1), basta probar las implicaciones
\[
d'')\Rightarrow c)\Rightarrow b)\Rightarrow d'').
\]

$d'')\Rightarrow c)$: Como $\Omega^1_{X/Y}$ es un $\sh{O}_X$-módulo de tipo finito, del lema de Nakayama resulta que la condición c) equivale a $(\Omega^1_{X/Y})_x/\mathfrak{m}_y(\Omega^1_{X/Y})_x=0$, es decir (16.4.5), a $(\Omega^1_{X_y/\Spec(k(y))})_x=0$. Se reduce, pues, al caso en que $Y$ es el espectro de un cuerpo $k$ y $X$ un $k$-preesquema algebraico. La hipótesis de que $\sh{O}_{X,x}$ es un cuerpo $k'$, extensión finita de $k$, implica primero que $x$ es \emph{cerrado} en $X$ (I, 6.4.2), y después que $x$ es un punto \emph{maximal} del preesquema noetheriano $X$, por lo que es un punto aislado de $X$. Reemplazando $X$ por el abierto $\{x\}$ de $X$, se puede suponer que $X=\Spec(k')$; pero entonces la hipótesis de que $k'$ es una extensión \emph{finita separable} de $k$ implica $\Omega^1_{k'/k}=0$ (0, 20.6.20), lo que demuestra c).

$c)\Rightarrow b)$: Se ha visto más arriba que entonces se tiene $\Omega^1_{X/Y}|U=0$ para una vecindad abierta $U$ de $x$ en $X$; la afirmación b) resulta entonces de la definición de $\Omega^1_{X/Y}$ (16.3.1) y de (16.1.9).

$b)\Rightarrow d'')$: Reemplazando $X$ por una vecindad abierta de $x$, se puede suponer que $\Delta_f$ es una inmersión abierta; si se designa por $f_y:X_y\to\Spec(k(y))$ el morfismo deducido de $f$ por cambio de base, $\Delta_{f_y}$ también es entonces una inmersión abierta (I, 5.3.4), y como la condición d'') solo concierne al preesquema $X_y$, se ve que se puede restringir al caso en que $Y$ es el espectro de un cuerpo $k$, $X$ el espectro de una $k$-álgebra $A$ de tipo finito; la propiedad d'') quedará establecida si se demuestra que $A$ es una $k$-álgebra \emph{finita y separable}, pues tal álgebra es un producto directo de extensiones finitas separables de $k$. Si $K$ es una extensión algebraicamente cerrada de $k$, equivale a decir que $A\otimes_k K$ es una $K$-álgebra finita y separable (4.6.1), por lo que se ve que se puede restringir al caso en que $k$ es algebraicamente cerrado. Demostremos primero que $A$ es una $k$-álgebra \emph{finita}: bastará demostrar que todo punto \emph{cerrado} $x$ de $X$ está aislado, pues entonces el conjunto de estos puntos es abierto en $X$ y discreto, y por tanto finito, ya que es cuasicompacto ($X$ es noetheriano), lo que establecerá nuestra afirmación en virtud de (I, 6.4.4). Ahora bien, se tiene entonces $k(x)=k$ porque $k$ es algebraicamente cerrado (I, 6.4.3), por lo que existe una $Y$-sección $s$ de $X$ tal que $s(Y)=\{x\}$, y en virtud de (17.4.1.1), $\{x\}$ es la imagen inversa de la diagonal $\Delta_X(X)$ por un morfismo $X\to X\times_Y X$; por tanto, $\{x\}$ es abierto en $X$ en virtud de la hipótesis b). Se ha demostrado así que $A$ es una

\oldpage[IV-4]{65}
$k$-álgebra \emph{finita}, producto directo de $k$-álgebras finitas locales. Para expresar que $\Delta_f$ es una inmersión abierta, se puede, pues, restringir al caso en que $A$ es una $k$-álgebra local finita, de modo que $X=\Spec(A)$ está reducido a un solo punto; el cuerpo residual de $A$, al ser una extensión finita de $k$, es necesariamente idéntico a $k$, y por consiguiente (I, 3.4.9) $X\times_k X$ está reducido a un solo punto y $\Delta_f$ es necesariamente un \emph{isomorfismo}. Ahora bien, puesto que $A$ es una $k$-álgebra, el homomorfismo canónico $A\otimes_k A\to A$ solo puede ser biyectivo si $A=k$. C.Q.D.

\begin{remark}[17.4.1.2]
\label{IV.17.4.1.2}
Supongamos únicamente que $f$ sea \emph{localmente de tipo finito}. Entonces $\Omega^1_{X/Y}$ sigue siendo un $\sh{O}_X$-módulo de tipo finito (16.3.9), y $\Delta_f$ un morfismo localmente de presentación finita (I, 4.3.1). Toda la demostración de (17.4.1) es, por tanto, válida, con la condición de reemplazar a) por: \emph{la restricción de $f$ a una vecindad adecuada de $x$ es un morfismo formalmente no ramificado}. Se ve además que en este caso la restricción de $f$ a una vecindad adecuada de $x$ es un morfismo \emph{localmente cuasifinito}.
\end{remark}

\begin{corollary}[17.4.2]
\label{IV.17.4.2}
Sea $f:X\to Y$ un morfismo localmente de presentación finita. Las propiedades siguientes son equivalentes:
\begin{enumerate}
  \item[a)] $f$ es no ramificado.
  \item[b)] El morfismo diagonal $\Delta_f:X\to X\times_Y X$ es una inmersión abierta.
  \item[b')] Para todo morfismo $Y'\to Y$, toda $Y'$-sección de $X'=X\times_Y Y'$ es una inmersión abierta.
  \item[c)] Se tiene $\Omega^1_{X/Y}=0$.
  \item[d)] Para todo $y\in Y$, el $k(y)$-preesquema $f^{-1}(y)$ es no ramificado sobre $k(y)$.
  \item[d')] Para todo $y\in Y$, el $k(y)$-preesquema $f^{-1}(y)$ es isomorfo a un preesquema de la forma $\coprod_{\lambda\in L}\Spec(K_\lambda)$, donde, para todo $\lambda\in L$, $K_\lambda$ es una extensión finita y separable de $k(y)$.
  \item[e)] Para todo $x\in X$, el anillo $\sh{O}_{X,x}$ es una $\sh{O}_{Y,f(x)}$-álgebra formalmente no ramificada para las topologías discretas.
\end{enumerate}
\end{corollary}

\begin{corollary}[17.4.3]
\label{IV.17.4.3}
Si $f:X\to Y$ es no ramificado, entonces $f$ es localmente cuasifinito (Err$_{\mathrm{III}}$, 20).
\end{corollary}

\begin{proposition}[17.4.4]
\label{IV.17.4.4}
Sean $Y$ un preesquema localmente noetheriano, $f:X\to Y$ un morfismo localmente de tipo finito, $x$ un punto de $X$, $y=f(x)$. Pongamos $A=\sh{O}_{Y,y}$, $B=\sh{O}_{X,x}$, que son anillos locales noetherianos, y sea $k$ el cuerpo residual de $A$. Entonces las condiciones equivalentes a) a e) del teorema (17.4.1) son también equivalentes a cada una de las siguientes:
\begin{enumerate}
  \item[f)] $\widehat{B}\otimes_{\widehat{A}}k$ es un cuerpo, extensión finita separable de $k$ (lo que implica que $\widehat{B}$ es una $\widehat{A}$-álgebra finita).
  \item[f')] $\widehat{B}$ es una $\widehat{A}$-álgebra formalmente no ramificada para las topologías ádicas.
\end{enumerate}
Si además $k(x)=k(y)$, o si $k$ es separablemente cerrado, estas condiciones también equivalen a:
\begin{enumerate}
  \item[f'')] El homomorfismo $\widehat{A}\to\widehat{B}$ es sobreyectivo.
\end{enumerate}
\end{proposition}

Observemos primero, mediante el mismo razonamiento que en (0, 19.3.6), que equivale a decir que $B$ es una $A$-álgebra formalmente no ramificada para las topologías preádicas, o que $\widehat{B}$ es una $\widehat{A}$-álgebra formalmente no ramificada para las topologías ádicas. Por otra parte, la hipótesis de que $f$ es localmente de tipo finito implica que $\Omega^1_{B/A}$

\oldpage[IV-4]{66}
es un $B$-módulo de tipo finito (16.3.9), por tanto \emph{separado} para la topología $\mathfrak n$-preádica (donde $\mathfrak n$ es el ideal maximal de $B$) ($0_{\mathrm I}$, 7.3.5); por consiguiente, equivale a decir que $\Omega^1_{B/A}=0$ o que $\widehat{\Omega}^1_{B/A}=0$; así pues (0, 20.7.4), equivale a decir que $B$ es una $A$-álgebra formalmente no ramificada para las topologías \emph{discretas}, o para las topologías \emph{preádicas}. Esto demuestra la equivalencia de las condiciones e) y f'). Si $\mathfrak m$ es el ideal maximal de $A$, se tiene
\[
k=A/\mathfrak m=\widehat A/\mathfrak m\widehat A,
\qquad
\widehat B\otimes_{\widehat A}k=\widehat B/\mathfrak m\widehat B
=\widehat B\otimes_B(B/\mathfrak m B),
\]
y por tanto ($0_{\mathrm I}$, 7.3.5), $\widehat B/\mathfrak m\widehat B$ es la completación de $B/\mathfrak m B=B\otimes_A k$ para la topología $\mathfrak n$-preádica; esto demuestra la equivalencia de d'') y f). Finalmente, cuando $k(x)=k(y)$ o cuando $k$ es separablemente cerrado, la condición f) implica que el homomorfismo $\widehat A/\mathfrak m\widehat A\to\widehat B/\mathfrak m\widehat B$ es biyectivo; por otra parte, la condición f) implica que $\widehat B$ es una $\widehat A$-álgebra \emph{cuasifinita} ($0_{\mathrm I}$, 7.4.4), por tanto \emph{finita}, ya que $\widehat A$ es completo y $\widehat B$ está separado para la topología $\mathfrak m$-preádica, siendo $\mathfrak m\widehat B$ un ideal de definición de $\widehat B$ ($0_{\mathrm I}$, 7.4.1). El homomorfismo $\widehat A\to\widehat B$ es, pues, sobreyectivo en virtud del lema de Nakayama. Por tanto, f) implica f''), y el recíproco es evidente.

\begin{env}[17.4.5]
\label{IV.17.4.5}
Dados un $S$-preesquema $Y$ y dos $S$-morfismos $f:X\to Y$, $g:X\to Y$, se deduce canónicamente un $S$-morfismo $(f,g)_S:X\to Y\times_S Y$. Llamaremos \emph{preesquema de coincidencias} de $f$ y $g$ a la imagen inversa por $(f,g)_S$ de la diagonal $\Delta_{Y|S}$; es, por tanto, un subpreesquema de $X$, que es \emph{cerrado} cuando $Y$ es un $S$-esquema (I, 5.4.1).
\end{env}

\begin{proposition}[17.4.6]
\label{IV.17.4.6}
Sean $h:Y\to S$ un morfismo no ramificado, y sean $f:X\to Y$, $g:X\to Y$ dos $S$-morfismos. Entonces el preesquema $C$ de coincidencias de $f$ y $g$ es un subpreesquema inducido sobre un abierto de $X$; si además $Y$ es un $S$-esquema (I, 5.4.1), $C$ es un subpreesquema cerrado de $X$.
\end{proposition}

En efecto, puesto que $\Delta_h:Y\to Y\times_S Y$ es una inmersión abierta (17.4.2), la imagen inversa por $(f,g)_S$ de $\Delta_{Y|S}$ es un subpreesquema inducido sobre un abierto de $X$ (I, 4.4.1). La última afirmación resulta de (17.4.5).

\begin{corollary}[17.4.7]
\label{IV.17.4.7}
Bajo las hipótesis de (17.4.6), sea $x$ un punto de $X$ tal que los dos morfismos compuestos $\Spec(k(x))\to X\xrightarrow{f}Y$ y $\Spec(k(x))\to X\xrightarrow{g}Y$ sean iguales. Entonces existe una vecindad abierta $U$ de $x$ tal que $f|U=g|U$. Si además $Y$ es un $S$-esquema, existe una vecindad abierta y cerrada $X'$ de $x$ en $X$ tal que $f|X'=g|X'$. Si finalmente se supone además que $X$ es conexo, se tiene $f=g$.
\end{corollary}

Esto resulta de (17.4.6) y de (I, 5.3.17).

\begin{corollary}[17.4.8]
\label{IV.17.4.8}
Bajo las hipótesis de (17.4.6), supongamos que el morfismo estructural $\varphi=h\circ f=h\circ g$ de $X$ en $S$ sea cerrado. Sea $s$ un punto de $S$; sea $X_s$ el $k(s)$-preesquema $\varphi^{-1}(s)$ y supongamos que los dos morfismos compuestos $X_s\to X\xrightarrow{f}Y$ y $X_s\to X\xrightarrow{g}Y$ sean iguales. Entonces existe una vecindad abierta $V$ de $s$ en $S$ tal que $f|\varphi^{-1}(V)=g|\varphi^{-1}(V)$. Si además $Y$ es un $S$-esquema y si $\varphi$ es abierto, se puede tomar $V$ abierto y cerrado. Si finalmente se supone además que $S$ es conexo, se tiene $f=g$.
\end{corollary}

De (17.4.7) resulta que el preesquema $C$ de coincidencias de $f$ y $g$ está inducido sobre un abierto de $X$ y contiene $X_s$. Como $\varphi$ es cerrado, existe una vecindad abierta $V$ de $s$ tal que $\varphi^{-1}(V)\subset C$. Si además $Y$ es un $S$-esquema, $C$ es cerrado, por lo que $\varphi(X-C)$ es a la

\oldpage[IV-4]{67}
vez abierto y cerrado en $S$, y su complemento $V$ en $S$ es, por tanto, una vecindad abierta y cerrada de $s$ tal que $\varphi^{-1}(V)\subset C$.

\begin{proposition}[17.4.9]
\label{IV.17.4.9}
Sean $Y$ un preesquema conexo y $f:X\to Y$ un morfismo no ramificado y separado. Entonces toda $Y$-sección $g$ de $X$ es un isomorfismo de $Y$ sobre una componente conexa abierta de $X$, y la aplicación $g\mapsto g(Y)$ es una biyección de $\Gamma(X/Y)$ sobre el conjunto de las componentes conexas $Z$ de $X$ (necesariamente abiertas en $X$) tales que la restricción de $f$ a $Z$ sea un isomorfismo de $Z$ sobre $Y$. En particular, si $g'$ y $g''$ son dos $Y$-secciones de $X$ tales que $g'(y)=g''(y)$ para un $y\in Y$, se tiene $g'=g''$.
\end{proposition}

En efecto, de (17.4.1, b'')) resulta que una $Y$-sección $s$ de $X$ es una inmersión abierta, y como $X$ es un $Y$-esquema, $s$ es también una inmersión cerrada (I, 5.4.7); de ello resulta que $s$ es un isomorfismo de $Y$ sobre un subpreesquema de $X$ inducido sobre una parte abierta y cerrada de $X$, y como $s(Y)$ es conexo, se trata necesariamente de una componente conexa de $X$. El resto de la proposición es inmediato.

\begin{remark}[17.4.10]
\label{IV.17.4.10}
Teniendo en cuenta la observación (17.4.1.2), se ve que, en los enunciados (17.4.6) a (17.4.9), se pueden reemplazar en todas partes las palabras «no ramificado» por «formalmente no ramificado y localmente de tipo finito».
\end{remark}

\subsection{Caracterizaciones de los morfismos lisos}
\label{subsection:IV.17.5}

\begin{theorem}[17.5.1]
\label{IV.17.5.1}
Sean $f:X\to Y$ un morfismo localmente de presentación finita, $x$ un punto de $X$, $y=f(x)$. Las condiciones siguientes son equivalentes:
\begin{enumerate}
  \item[a)] $f$ es liso en el punto $x$.
  \item[b)] $f$ es plano en el punto $x$ y el $k(y)$-preesquema $f^{-1}(y)$ es liso sobre $k(y)$ en el punto $x$.
  \item[b')] $f$ es regular en el punto $x$ (6.8.1).
  \item[c)] El anillo $\sh{O}_{X,x}$ es una $\sh{O}_{Y,y}$-álgebra formalmente lisa para las topologías discretas.
\end{enumerate}
\end{theorem}

Se puede restringir al caso en que $Y=\Spec(A)$, $X=\Spec(C)$, donde $C=B/\mathfrak{J}$, siendo $B=A[T_1,\ldots,T_n]$ un álgebra de polinomios y $\mathfrak{J}$ un ideal de $B$ de tipo finito. La equivalencia de a) y c) resulta entonces de la equivalencia de a) y c) en (0, 22.6.4). Por otra parte, aplicando este resultado al morfismo localmente de tipo finito $f^{-1}(y)\to\Spec(k(y))$, se ve que la equivalencia de b) y b') resulta de la equivalencia de a) y b) en (6.8.6). Resta, pues, demostrar la equivalencia de a) y b).

Demostremos primero que a) implica b); denotemos por $\mathfrak p$ el ideal primo $\mathfrak i_x$ en $C$, y por $\mathfrak r$ el ideal primo $\mathfrak i_y$ en $A$; se tiene $\mathfrak p=\mathfrak q/\mathfrak J$, donde $\mathfrak q$ es un ideal primo de $B$ y $\mathfrak r$ es la imagen inversa de $\mathfrak q$ en $A$. La hipótesis a) implica primero que $f^{-1}(y)$ es liso sobre $k(y)$ en el punto $x$ por (17.3.3), y se trata de demostrar además que $C_{\mathfrak p}=\sh{O}_{X,x}$ es un $A_{\mathfrak r}$-módulo \emph{plano}. Como $C_{\mathfrak p}$ es una $A_{\mathfrak r}$-álgebra formalmente lisa y $B$ es una $A$-álgebra formalmente lisa (para las topologías discretas), el criterio jacobiano (0, 22.6.4), unido a (0, 19.1.12), implica que existe en $\mathfrak J$ un sistema de $r$ polinomios $u_i$ ($1\leq i\leq r$) y $r$ índices $j_i$ ($1\leq i\leq r$) tales que las imágenes de los $u_i$ en $\mathfrak J_{\mathfrak q}/\mathfrak J_{\mathfrak q}^2$ generan este $B_{\mathfrak q}$-módulo y que se tiene
\[
\label{IV.17.5.1.1}
\det(\partial u_i/\partial T_{j_k})\notin\mathfrak q.
\tag{17.5.1.1}
\]

\oldpage[IV-4]{68}
Observemos ahora que, si $g:\Spec(B)\to\Spec(A)$ es el morfismo estructural, la fibra $g^{-1}(y)$ es el espectro del anillo regular $k(y)[T_1,\ldots,T_n]$ (0, 17.3.7), por lo que el anillo local noetheriano $B_{\mathfrak q}/\mathfrak rB_{\mathfrak q}$ en un punto de esta fibra es regular. Ahora bien, la condición (17.5.1.1) implica que las imágenes canónicas $v_i$ de los $u_i$ en el ideal maximal $\mathfrak m$ de $B_{\mathfrak q}/\mathfrak rB_{\mathfrak q}$ son \emph{linealmente independientes} módulo $\mathfrak m^2$: en caso contrario, existirían, en efecto, polinomios $w_i\in B$ ($1\leq i\leq r$), no todos pertenecientes a $\mathfrak q$, tales que $\sum_{i=1}^r w_i u_i\in\mathfrak q^2$. Derivando respecto de los $T_{j_k}$, se concluiría que $\sum_{i=1}^r w_i(\partial u_i/\partial T_{j_k})\in\mathfrak q$ para $1\leq k\leq r$, lo que contradiría (17.5.1.1), puesto que $\mathfrak q$ es primo. Se concluye, pues, de (0, 17.1.7) que $(v_i)$ es una \emph{sucesión regular} en $B_{\mathfrak q}/\mathfrak rB_{\mathfrak q}$. Pero como el morfismo $g$ es localmente de presentación finita y $B$ es un $A$-módulo plano, de (11.3.8) resulta que las imágenes canónicas $u'_i$ de los $u_i$ en $B_{\mathfrak q}$ forman también una \emph{sucesión regular} y que $B_{\mathfrak q}/(\sum_i u'_iB_{\mathfrak q})$ es un $A_{\mathfrak r}$-módulo \emph{plano}. Como las imágenes de los $u'_i$ en $\mathfrak J_{\mathfrak q}/\mathfrak J_{\mathfrak q}^2$ generan este $B_{\mathfrak q}$-módulo, del lema de Nakayama resulta que se tiene $\sum_i u'_iB_{\mathfrak q}=\mathfrak J_{\mathfrak q}$, y $C_{\mathfrak p}=B_{\mathfrak q}/\mathfrak J_{\mathfrak q}$ es, por tanto, efectivamente un $A_{\mathfrak r}$-módulo plano.

Demostremos finalmente que b) implica a). Con las mismas notaciones, la hipótesis de que $B_{\mathfrak q}/\mathfrak J_{\mathfrak q}$ es un $A_{\mathfrak r}$-módulo plano implica que el homomorfismo canónico $\mathfrak J_{\mathfrak q}/\mathfrak r\mathfrak J_{\mathfrak q}\to B_{\mathfrak q}/\mathfrak rB_{\mathfrak q}$ es \emph{inyectivo} ($0_{\mathrm I}$, 6.1.2), de modo que $\mathfrak J_{\mathfrak q}/\mathfrak r\mathfrak J_{\mathfrak q}$ se identifica con un ideal de $B_{\mathfrak q}/\mathfrak rB_{\mathfrak q}$. Como $B_{\mathfrak q}/\mathfrak rB_{\mathfrak q}$ es una $A_{\mathfrak r}$-álgebra formalmente lisa para las topologías discretas, se puede aplicar a $C_{\mathfrak p}=(B_{\mathfrak q}/\mathfrak rB_{\mathfrak q})/(\mathfrak J_{\mathfrak q}/\mathfrak r\mathfrak J_{\mathfrak q})$ el criterio jacobiano (0, 22.6.4); unido a (0, 19.1.12), este último demuestra, en virtud de la hipótesis b), la existencia de $r$ polinomios $v_i\in k(y)[T_1,\ldots,T_n]$ tales que sus imágenes en $(\mathfrak J_{\mathfrak q}/\mathfrak r\mathfrak J_{\mathfrak q})/(\mathfrak J_{\mathfrak q}/\mathfrak r\mathfrak J_{\mathfrak q})^2$ generan este $(B_{\mathfrak q}/\mathfrak rB_{\mathfrak q})$-módulo y que se tiene
\[
\label{IV.17.5.1.2}
\det(\partial v_i/\partial T_{j_k})\notin\mathfrak qB_{\mathfrak q}/\mathfrak rB_{\mathfrak q}.
\tag{17.5.1.2}
\]
Si, para todo $i$, se designa entonces por $u_i$ un elemento de $\mathfrak J$ cuya imagen canónica es $v_i$, de (17.5.1.2) resulta que los $u_i$ satisfacen la condición (17.5.1.1); por otra parte, en virtud del lema de Nakayama, las imágenes $u'_i$ de los $u_i$ en $\mathfrak J_{\mathfrak q}$ generan este $B_{\mathfrak q}$-módulo. El criterio jacobiano (0, 22.6.4), unido a (0, 19.1.12), demuestra entonces que $C_{\mathfrak p}=B_{\mathfrak q}/\mathfrak J_{\mathfrak q}$ es una $A_{\mathfrak r}$-álgebra formalmente lisa para las topologías discretas. C.Q.D.

\begin{corollary}[17.5.2]
\label{IV.17.5.2}
Sea $f:X\to Y$ un morfismo localmente de presentación finita. Para que $f$ sea liso (en el sentido de (17.3.1)), es necesario y suficiente que $f$ sea regular (6.8.1), es decir, que $f$ sea plano y que, para todo $y\in Y$, $f^{-1}(y)$ sea un $k(y)$-preesquema geométricamente regular (6.7.6).
\end{corollary}

Se ha establecido así la equivalencia de las dos definiciones de «morfismo liso» dadas en (6.8.1) y (17.3.1).

\begin{proposition}[17.5.3]
\label{IV.17.5.3}
Sean $Y$ un preesquema localmente noetheriano, $f:X\to Y$ un morfismo localmente de tipo finito, $x$ un punto de $X$, $y=f(x)$. Pongamos $A=\sh{O}_{Y,y}$, $B=\sh{O}_{X,x}$, que son anillos locales noetherianos. Entonces las condiciones equivalentes a) a c) de (17.5.1) son también equivalentes a cada una de las siguientes:
\begin{enumerate}
  \item[d)] $B$ es una $A$-álgebra formalmente lisa para las topologías preádicas.
  \item[d')] $\widehat B$ es una $\widehat A$-álgebra formalmente lisa para las topologías ádicas.
\end{enumerate}
Si además $k(x)=k(y)$, estas condiciones también equivalen a:
\begin{enumerate}
  \item[d'')] $\widehat B$ es una $\widehat A$-álgebra isomorfa a un álgebra de series formales $\widehat A[[T_1,\ldots,T_n]]$.
\end{enumerate}
\end{proposition}

\oldpage[IV-4]{69}
La equivalencia de la condición c) de (17.5.1) y de d) resulta de la equivalencia de a) y d) en el criterio jacobiano (0, 22.6.4), y la equivalencia de d) y d') resulta de (0, 19.3.6). Por otra parte, d'') implica d') sin hipótesis sobre los cuerpos residuales (0, 19.3.4). Finalmente, si $\mathfrak m$ denota el ideal maximal de $\widehat A$, la hipótesis d') implica que $\widehat B/\mathfrak m\widehat B$ es una $k(y)$-álgebra local noetheriana completa, formalmente lisa para su topología ádica (0, 19.3.5); la hipótesis $k(y)=k(x)$ implica, pues, que $\widehat B/\mathfrak m\widehat B$ es $k(y)$-isomorfa a un álgebra de series formales $k(y)[[T_1,\ldots,T_n]]$ (0, 19.6.4). Como, por otra parte, $\widehat A[[T_1,\ldots,T_n]]$ es un $\widehat A$-módulo plano y una $\widehat A$-álgebra local noetheriana completa, se concluye de (0, 19.7.1.5) que esta álgebra es isomorfa a $\widehat B$. Por tanto, d') implica d'') bajo la hipótesis adicional $k(x)=k(y)$.

\begin{remark}[17.5.4]
\label{IV.17.5.4}
Supongamos que $Y$ es un preesquema localmente noetheriano, y $f:X\to Y$ un morfismo localmente de tipo finito. El criterio (17.5.3, d)), unido a (0, 22.1.4), muestra que para demostrar que $f$ es liso se puede aplicar la definición (17.1.1), restringiéndose al caso en que el esquema afín $Y'$ es el espectro de un anillo local artiniano.
\end{remark}

\begin{proposition}[17.5.5]
\label{IV.17.5.5}
Sean $Y$ un preesquema localmente noetheriano, $f:X\to Y$ un morfismo localmente de tipo finito, $x$ un punto de $X$, $y=f(x)$. Supongamos que $Y$ sea reducido en el punto $y$. Entonces, para que $f$ sea liso en el punto $x$, es necesario y suficiente que $f$ sea universalmente abierto en una vecindad de $x$ en $f^{-1}(y)$ y que $f^{-1}(y)$ sea un $k(y)$-preesquema geométricamente regular en el punto $x$.
\end{proposition}

Teniendo en cuenta (17.5.1), todo se reduce a ver que, si $f^{-1}(y)$ es un $k(y)$-preesquema geométricamente regular en el punto $x$, equivale a decir que $f$ es plano en el punto $x$ o universalmente abierto en una vecindad de $x$ en $f^{-1}(y)$. Ahora bien, si $f$ es plano en el punto $x$, lo es en una vecindad de $x$ en $X$ (11.1.1) y, por tanto, es universalmente abierto en esta vecindad (2.4.6). Recíprocamente, la hipótesis de que $f$ es universalmente abierto en una vecindad de $x$ en $f^{-1}(y)$ y de que $f^{-1}(y)$ es un $k(y)$-preesquema geométricamente regular en el punto $x$ implica que $f$ es plano en el punto $x$, puesto que $\sh{O}_{Y,y}$ es reducido (15.2.2).

\begin{corollary}[17.5.6]
\label{IV.17.5.6}
Sean $Y$ un preesquema localmente noetheriano, $f:X\to Y$ un morfismo localmente de tipo finito, $x$ un punto de $X$, $y=f(x)$. Se supone que $Y$ es reducido y geométricamente unirrama (6.15.1) en el punto $y$. Entonces, para que $f$ sea liso en el punto $x$, es necesario y suficiente que $f$ sea equidimensional en el punto $x$ y que $f^{-1}(y)$ sea un $k(y)$-preesquema geométricamente regular en el punto $x$.
\end{corollary}

Observando que el conjunto de los puntos donde $f$ es equidimensional es abierto (13.3.2), se ve que el corolario resulta de (17.5.5) y del criterio de Chevalley (14.4.4).

El hecho de que $f$ sea un morfismo liso en un punto implica, en particular, que $f$ satisface

\oldpage[IV-4]{70}
en ese punto todas las propiedades definidas en (6.8.1). Se tienen, pues, las propiedades siguientes, que recordamos para facilitar las referencias:

\begin{proposition}[17.5.7]
\label{IV.17.5.7}
Sea $f:X\to Y$ un morfismo localmente de presentación finita, liso en un punto $x\in X$; pongamos $y=f(x)$. Entonces, para que el anillo $\sh{O}_{X,x}$ sea reducido (resp. íntegramente cerrado, resp. geométricamente unirrama), es necesario y suficiente que $\sh{O}_{Y,y}$ lo sea.
\end{proposition}

Esto se ha demostrado, en efecto, en (11.3.13) y en (11.3.14), este último completado por Err$_{\mathrm{IV}}$, 30.

\begin{proposition}[17.5.8]
\label{IV.17.5.8}
Sean $Y$ un preesquema localmente noetheriano, $f:X\to Y$ un morfismo localmente de tipo finito, liso en un punto $x\in X$; pongamos $y=f(x)$. Entonces:
\begin{enumerate}
  \item[(i)] Se tiene $\dim(\sh{O}_{X,x})=\dim(\sh{O}_{Y,y})+\dim(\sh{O}_{X,x}\otimes_{\sh{O}_{Y,y}}k(y))$.
  \item[(ii)] Se tiene $\operatorname{coprof}(\sh{O}_{X,x})=\operatorname{coprof}(\sh{O}_{Y,y})$.
  \item[(iii)] Para que el anillo $\sh{O}_{X,x}$ posea la propiedad $(S_n)$ (5.7.2) (resp. $(R_n)$ (5.8.2)), es necesario y suficiente que el anillo $\sh{O}_{Y,y}$ la posea. En particular, para que $\sh{O}_{X,x}$ sea regular, es necesario y suficiente que $\sh{O}_{Y,y}$ lo sea.
\end{enumerate}
\end{proposition}

Se trata de casos particulares de (6.1.2), (6.3.2), (6.4.1) y (6.5.3).

\subsection{Caracterizaciones de los morfismos étale}
\label{subsection:IV.17.6}

\begin{theorem}[17.6.1]
\label{IV.17.6.1}
Sean $f:X\to Y$ un morfismo localmente de presentación finita, $x$ un punto de $X$, $y=f(x)$. Las condiciones siguientes son equivalentes:
\begin{enumerate}
  \item[a)] $f$ es étale en el punto $x$.
  \item[a')] $f$ es liso en el punto $x$ y no ramificado en el punto $x$.
  \item[b)] $f$ es liso en el punto $x$ y cuasifinito en el punto $x$ (Err$_{\mathrm{III}}$, 20).
  \item[c)] $f$ es plano en el punto $x$ y no ramificado en el punto $x$.
  \item[c')] $f$ es plano en el punto $x$ y el anillo $\sh{O}_{X,x}/\mathfrak m_y\sh{O}_{X,x}$ es un cuerpo, extensión finita separable de $k(y)$.
  \item[d)] El anillo $\sh{O}_{X,x}$ es una $\sh{O}_{Y,y}$-álgebra formalmente étale para las topologías discretas.
\end{enumerate}
\end{theorem}

La equivalencia de a) y a') resulta inmediatamente de las definiciones; la de a) y d) resulta de la equivalencia de a) y e) en (17.4.1) y de la equivalencia de a) y d) en (17.5.1). La equivalencia de c) y c') resulta de la equivalencia de a) y d') en (17.4.1). El hecho de que a') implique c') se deduce de (17.5.1); recíprocamente, si se verifica c'), $f$ es regular (por tanto liso por (17.5.1)) en el punto $x$, pues si $K$ es una extensión finita separable de un cuerpo $k$, entonces, para toda extensión $k'$ de $k$, $\Spec(K\otimes_k k')$ es regular, ya que es suma de un número finito de espectros de cuerpos. El hecho de que a') implique b) resulta de (17.4.1, d')) y de (17.5.1, b)). Resta, pues, ver que b) implica c), y como ya se sabe que $f$ es plano en el punto $x$ por (17.5.1), basta demostrar que $f^{-1}(y)$ es un $k(y)$-preesquema no ramificado sobre $k(y)$; es decir, se reduce a demostrar que b) implica c) cuando $Y=\Spec(k)$ es el espectro de un cuerpo $k$. Como la cuestión es local sobre $X$, se puede restringir al caso en que $X=\Spec(A)$, donde $A$ es una $k$-álgebra finita y local ($0_{\mathrm I}$, 7.4.1). En virtud de la hipótesis b), $A$ es una $k$-álgebra formalmente lisa para las topologías discretas, que aquí coinciden con las topologías preádicas; por tanto (0, 19.6.5), $A$ es un anillo local regular, y así

\oldpage[IV-4]{71}
un cuerpo, puesto que es artiniano, y resulta entonces de (0, 19.6.5.1) que $A$ debe ser una extensión finita y separable de $k$, lo que concluye la demostración (17.4.1).

\begin{corollary}[17.6.2]
\label{IV.17.6.2}
Sea $f:X\to Y$ un morfismo localmente de presentación finita. Las condiciones siguientes son equivalentes:
\begin{enumerate}
  \item[a)] $f$ es étale.
  \item[a')] $f$ es liso y no ramificado.
  \item[b)] $f$ es liso y localmente cuasifinito (Err$_{\mathrm{III}}$, 20).
  \item[c)] $f$ es plano y no ramificado.
  \item[c')] $f$ es plano, y toda fibra $f^{-1}(y)$ es un $k(y)$-esquema que es suma de espectros de cuerpos, extensiones finitas y separables de $k(y)$.
  \item[c'')] $f$ es plano, y para todo $y\in Y$ y toda extensión algebraicamente cerrada $k'$ de $k(y)$, la «fibra geométrica» $f^{-1}(y)\otimes_{k(y)}k'$ es suma de espectros de cuerpos isomorfos a $k'$.
\end{enumerate}
\end{corollary}

El único punto que queda por demostrar es la equivalencia de c') y c''). Es claro que c') implica c'') por cambio de base (17.3.3). Por otra parte, como el morfismo de proyección $f^{-1}(y)\otimes_{k(y)}k'\to f^{-1}(y)$ es abierto (2.4.10), la hipótesis c'') implica que el espacio $f^{-1}(y)$ es discreto; por tanto, para todo punto $x\in X_y=f^{-1}(y)$, el anillo local $\sh{O}_{X_y,x}=A$ es una $k(y)$-álgebra finita y, por tanto, un anillo local artiniano; además, de c'') resulta que $\Spec(A\otimes_{k(y)}k')$ es suma de espectros de cuerpos isomorfos a $k'$, lo cual solo es posible si $A$ es un cuerpo, extensión finita separable de $k(y)$ (4.6.1).

\begin{proposition}[17.6.3]
\label{IV.17.6.3}
Sean $Y$ un preesquema localmente noetheriano, $f:X\to Y$ un morfismo localmente de tipo finito, $x$ un punto de $X$, $y=f(x)$. Pongamos $A=\sh{O}_{Y,y}$, $B=\sh{O}_{X,x}$, que son anillos locales noetherianos, y sea $k$ el cuerpo residual de $A$. Entonces las condiciones equivalentes a) a d) de (17.6.1) son también equivalentes a cada una de las siguientes:
\begin{enumerate}
  \item[e)] $\widehat B$ es una $\widehat A$-álgebra formalmente étale para las topologías ádicas.
  \item[e')] $\widehat B$ es un $\widehat A$-módulo libre y $\widehat B\otimes_{\widehat A}k$ es un cuerpo, extensión finita y separable de $k$ (lo que implica que $\widehat B$ es una $\widehat A$-álgebra finita).
\end{enumerate}
Si además $k(x)=k(y)$, o si $k$ es separablemente cerrado, estas condiciones también equivalen a:
\begin{enumerate}
  \item[e'')] El homomorfismo canónico $\widehat A\to\widehat B$ es biyectivo.
\end{enumerate}
\end{proposition}

La equivalencia de e) y de cada una de las condiciones de (17.6.1) resulta inmediatamente de (17.4.4, f')) y (17.5.3, d')). El hecho de que e) implique e') resulta de (17.4.4, f)) y de (0, 19.7.1), teniendo en cuenta que $\widehat B$ es entonces una $\widehat A$-álgebra finita (17.4.4) y que, por tanto, equivale a decir que $\widehat B$ es un $\widehat A$-módulo plano o un $\widehat A$-módulo libre ($0_{\mathrm{III}}$, 10.1.3). Recíprocamente, el hecho de que e') implique e) resulta de (17.4.4) y de (0, 19.7.1). Finalmente, e') implica que el homomorfismo $\widehat A\to\widehat B$ es inyectivo, y si $k(x)=k(y)$ o si $k$ es separablemente cerrado, este homomorfismo es sobreyectivo por (17.4.4). El recíproco es inmediato.

\begin{proposition}[17.6.4]
\label{IV.17.6.4}
Bajo las hipótesis de (17.6.3), si $f$ es étale en el punto $x$, se tiene $\dim(\sh{O}_{X,x})=\dim(\sh{O}_{Y,y})$.
\end{proposition}

Es un caso particular de (17.5.8, (i)), puesto que $x$ está aislado en su fibra $f^{-1}(y)$.

\oldpage[IV-4]{72}

\subsection{Propiedades de descenso, de paso al límite y de constructibilidad}
\label{subsection:IV.17.7}
\label{IV.17.7}

\begin{proposition}[17.7.1]
\label{IV.17.7.1}
Sean $f:X\to Y$ un morfismo localmente de presentación finita, $g:Y'\to Y$ un morfismo, $X'=X\times_Y Y'$, $f'=f_{(Y')}:X'\to Y'$ y $g':X'\to X$ las proyecciones canónicas. Sea $x'$ un punto de $X'$ y pongamos $x=g'(x')$, $y'=f'(x')$.
\begin{enumerate}
  \item[(i)] Si $f'$ es no ramificado en el punto $x'$, entonces $f$ es no ramificado en el punto $x$.
  \item[(ii)] Supongamos además que $g$ sea plano en el punto $y'$. Entonces, si $f'$ es liso (resp. étale) en el punto $x'$, $f$ es liso (resp. étale) en el punto $x$.
\end{enumerate}
\end{proposition}

Pongamos $y=f(x)=g(y')$, de modo que se tiene $f'^{-1}(y')=f^{-1}(y)\otimes_{k(y)}k(y')$; observemos que $f'$ es localmente de presentación finita.

(i) Como la propiedad de un morfismo (localmente de presentación finita) de ser no ramificado en un punto solo depende de la fibra del morfismo en dicho punto (17.4.1, $d$), podemos limitarnos al caso en que $Y$ e $Y'$ son espectros de cuerpos. Pero entonces equivale decir que $f$ (resp. $f'$) es no ramificado en $x$ (resp. $x'$) o que es étale en ese punto (17.6.1, $c$), por lo que (i) es consecuencia de (ii).

(ii) Como $f'$ es plano en el punto $x'$ (17.5.1), la hipótesis de que $g$ es plano en el punto $y'$ implica que $f$ es plano en el punto $x$, pues la proyección $g':X'\to X$ es un morfismo plano en el punto $x'$, y $f\circ g'=g\circ f'$ es un morfismo plano en el punto $X'$, de donde se sigue la conclusión (2.2.11, (iv)). El hecho de que $f$ sea liso (resp. étale) en el punto $x$ solo depende ya de la fibra $f^{-1}(y)$ (17.5.1 y 17.6.1), y nos reducimos, por tanto, una vez más al caso en que $Y=\Spec(k)$ e $Y'=\Spec(k')$ son espectros de cuerpos. Decir que $f'$ es liso en el punto $x'$ significa entonces (17.5.1) que $X'$ es un $k'$-preesquema geométricamente regular en el punto $x'$, y esto implica (6.7.8) que $X$ es un $k$-preesquema geométricamente regular en el punto $x$, y por tanto que $f$ es liso en el punto $x$. Supongamos además que $f'$ sea étale en el punto $x'$, de modo que $x'$ es aislado en $f'^{-1}(y')$ (17.6.1); como la proyección $f'^{-1}(y')\to f^{-1}(y)$ es un morfismo abierto (2.4.10), $x$ es aislado en $f^{-1}(y)$; como ya sabemos que $f$ es liso en el punto $x$, es étale en dicho punto (17.6.1).

\begin{corollary}[17.7.2]
\label{IV.17.7.2}
\begin{enumerate}
  \item[(i)] Con las notaciones de (17.7.1), sea $U$ (resp. $U'$) el conjunto de puntos de $X$ (resp. $X'$) donde $f$ (resp. $f'$) es no ramificado; entonces se tiene $U'=g'^{-1}(U)$.
  \item[(ii)] Supongamos además que $g$ sea plano, y sea $V$ (resp. $V'$) el conjunto de puntos de $X$ donde $f$ (resp. $f'$) es liso; entonces $V'=g'^{-1}(V)$.
\end{enumerate}
\end{corollary}

\begin{corollary}[17.7.3]
\label{IV.17.7.3}
\begin{enumerate}
  \item[(i)] Supongamos que $g$ es sobreyectivo; entonces, para que $f$ sea no ramificado, es necesario y suficiente que $f'$ lo sea.
  \item[(ii)] Sean $f:X\to Y$ un $S$-morfismo y $g:S'\to S$ un morfismo fielmente plano. Supongamos que $f$ es localmente de presentación finita, o que $g$ es cuasicompacto. Entonces, para que $f$ sea liso (resp. no ramificado, resp. étale), es necesario y suficiente que $f'$ lo sea.
\end{enumerate}
\end{corollary}

En (ii), el caso en que $f$ es localmente de presentación finita resulta de (17.7.1). Si $g$ es cuasicompacto y $f'$ es liso (resp. no ramificado, resp. étale), $f'$ es localmente de presentación finita por (2.7.1, (iv)) y nos reducimos al primer caso.

\begin{proposition}[17.7.4]
\label{IV.17.7.4}
Sean $f:X\to Y$ un morfismo, $g:Y'\to Y$ un morfismo plano y localmente de presentación finita, $X'=X\times_Y Y'$, $f'=f_{(Y')}:X'\to Y'$ y $g':X'\to X$ las

\oldpage[IV-4]{73}

proyecciones canónicas. Sea $V$ (resp. $V'$) el conjunto de los $x\in X$ (resp. $x'\in X'$) donde $f$ posee una de las propiedades siguientes (resp. donde $f'$ posee la misma propiedad): ser:
\begin{enumerate}
  \item[(i)] localmente de tipo finito;
  \item[(ii)] localmente de presentación finita;
  \item[(iii)] plano;
  \item[(iv)] no ramificado;
  \item[(v)] liso;
  \item[(vi)] étale.
\end{enumerate}
Entonces se tiene $V'=g'^{-1}(V)$ (dicho de otro modo, para que $f'$ tenga la propiedad considerada en un punto $x'$, es necesario y suficiente que $f$ tenga dicha propiedad en el punto $x$).
\end{proposition}

La propiedad (iii) solo se menciona como recordatorio, y no requiere la hipótesis de que $g$ sea localmente de presentación finita ((2.2.11, (iv)), teniendo en cuenta que la proyección $g':X'\to X$ es un morfismo plano). En virtud de (17.7.2), las afirmaciones relativas a las propiedades (iv), (v) y (vi) son consecuencias de la afirmación relativa a (ii). Basta, por tanto, considerar los casos (i) y (ii). Es claro que $V'\supset g'^{-1}(V)$; queda, pues, por demostrar que $V'\subset g'^{-1}(V)$; observemos que los conjuntos $V$ y $V'$ son abiertos, y $W=g'(V')$ también es abierto en $X$ en virtud de (2.4.6). Se trata, por tanto, de probar que el morfismo $f|W:W\to Y$ es localmente de tipo finito (resp. localmente de presentación finita); como, por hipótesis, el compuesto
\[
V'\xrightarrow{g''}W\xrightarrow{f|W}Y
\]
(donde $g''$ es la restricción de $g'$), igual a $V'\xrightarrow{f'}Y'\xrightarrow{g}Y$, es localmente de tipo finito (resp. localmente de presentación finita) y $g''$ es sobreyectivo (y, por tanto, fielmente plano), nos reducimos a probar el lema siguiente, que mejora (11.3.16):

\begin{lemma}[17.7.5]
\label{IV.17.7.5}
Sean $f:X\to Y$ un morfismo fielmente plano y localmente de presentación finita, $g:Y\to Z$ un morfismo tal que $g\circ f:X\to Z$ tenga una de las propiedades siguientes: ser:
\begin{enumerate}
  \item[(i)] localmente de tipo finito;
  \item[(ii)] localmente de presentación finita;
  \item[(iii)] de tipo finito.
\end{enumerate}
Entonces $g$ tiene la misma propiedad. Si, además, $f$ es cuasicompacto o $g$ cuasiseparado, la misma conclusión es válida para la propiedad:
\begin{enumerate}
  \item[(iv)] ser de presentación finita.
\end{enumerate}
\end{lemma}

En los casos (i) y (ii), se trata de ver que, para todo $y\in Y$, existen una vecindad abierta afín $V$ de $y$ en $Y$ y una vecindad abierta afín $W$ de $z=g(y)$ en $Z$ que contiene $g(V)$, tales que el morfismo $V\to W$, restricción de $g$, sea de tipo finito (resp. de presentación finita). Ahora bien, existen por hipótesis un $x\in X$ tal que $f(x)=y$, una vecindad abierta afín $U$ de $x$ en $X$, una vecindad abierta afín $V'$ de $y$ en $Y$ que contiene $f(U)$ y una vecindad abierta afín $W$ de $z$ en $Z$ que contiene $g(V')$, tales que el morfismo $f_1:U\to V'$, restricción de $f$, sea plano y de presentación finita, y el morfismo $g_1:V'\to W$, restricción de $g$, sea tal que $g_1\circ f_1$ sea de tipo finito (resp. de presentación finita). Entonces $f(U)$ es abierto en $Y$ (2.4.6), y si $V\subset f(U)$ es una vecindad afín de $y$, el morfismo $f_2:f_1^{-1}(V)\to V$,

\oldpage[IV-4]{74}

restricción de $f_1$, sigue siendo de presentación finita y es, además, fielmente plano; asimismo, si $g_2=g_1|V$, $g_2\circ f_2$ es de tipo finito (resp. de presentación finita), pues el abierto $f_1^{-1}(V)$ es cuasicompacto (resp. cuasicompacto y cuasiseparado). Nos reducimos entonces a las hipótesis de (11.3.16), de donde se concluye que $g_2$ es un morfismo de tipo finito (resp. de presentación finita).

En el caso (iii), la cuestión es local sobre $Z$, por lo que podemos suponer que $Z$ es afín, y resulta entonces que $X$ es cuasicompacto, y por tanto también $Y=f(X)$. El caso (iii) es, pues, consecuencia de (i).

En el caso (iv) (con las hipótesis adicionales sobre $f$ o $g$), también podemos suponer que $Z$ es afín, y por tanto que $X$ e $Y$ son cuasicompactos; además, ya sabemos que $g$ es localmente de presentación finita y cuasicompacto (1.1.3), de modo que todo se reduce a ver que $g$ es cuasiseparado, y basta, por tanto, mostrar que esta propiedad es cierta cuando se supone que $f$ es cuasicompacto. Ahora bien, puesto que $g\circ f$ es cuasiseparado, lo mismo ocurre con $f$ (1.2.2, (v)); como $f$ es cuasicompacto y localmente de presentación finita, es de presentación finita (1.6.1); basta entonces repetir el razonamiento del primer párrafo de la demostración de (11.3.16).

\begin{remark}[17.7.6]
\label{IV.17.7.6}
En la afirmación relativa a la propiedad (iv), no puede suprimirse la hipótesis de que $f$ es cuasicompacto; de lo contrario, ello implicaría que todo morfismo $g:Y\to Z$ cuasicompacto y localmente de presentación finita sería de presentación finita, conclusión que se sabe errónea (1.6.4). En efecto, podemos limitarnos al caso en que $Z$ es afín, y por tanto $Y$ cuasicompacto; existe entonces un recubrimiento finito $(U_i)$ de $Y$ por abiertos afines tales que las restricciones $g|U_i$ sean de presentación finita; bastaría entonces tomar como $X$ el preesquema \emph{suma} de los $U_i$, y como $f:X\to Y$ el morfismo canónico, que es evidentemente fielmente plano y localmente de presentación finita (1.4.3); $g\circ f$ sería de presentación finita en virtud de la elección de los $U_i$ y de (1.6.5), de donde se sigue nuestra afirmación.
\end{remark}

\begin{proposition}[17.7.7]
\label{IV.17.7.7}
Sean $f:Y\to S$ y $h:X\to S$ dos morfismos localmente de presentación finita, $g:X\to Y$ un $S$-morfismo, $x$ un punto de $X$, $y=g(x)$. Supongamos que $g$ sea plano en el punto $x$. Entonces, si $h$ es liso (resp. no ramificado, resp. étale) en el punto $x$, $f$ es liso (resp. no ramificado, resp. étale) en el punto $y$.
\end{proposition}

Pongamos $s=h(x)=f(y)$. Decir que $h$ es no ramificado en el punto $x$ (resp. que $f$ es no ramificado en el punto $y$) equivale a decir que $h^{-1}(s)$ es étale sobre $k(s)$ en el punto $x$ (resp. que $f^{-1}(s)$ es étale sobre $k(s)$ en el punto $y$) (17.4.1 y 17.6.1); como el morfismo $g_s:h^{-1}(s)\to f^{-1}(s)$ deducido de $g$ es plano en el punto $x$, vemos que podemos limitarnos a probar la proposición cuando $h$ es liso o étale en el punto $x$. Además, como $h$ es entonces plano en el punto $x$ (17.5.1), $f$ es plano en el punto $y$, como resulta de (2.2.11, (iv)). Equivale, por tanto, (17.5.1) decir que $f$ es liso (resp. étale) en el punto $y$, o que $f^{-1}(s)$ es liso (resp. étale) sobre $k(s)$ en el punto $x$. Nos reducimos así al caso en que $S=\Spec(k)$ es el espectro de un cuerpo.

(i) \emph{Caso de los morfismos lisos}. Como $g$ es un morfismo localmente de presentación finita (1.4.3, (v)), existe una vecindad abierta $U$ de $x$ en $X$ en la cual $g$ es plano (11.3.1) y $h$ liso. Además, $g(U)$ es una vecindad abierta de $y$ en $Y$ (2.4.6); reemplazando $X$ por $U$ e $Y$ por $g(U)$, podemos suponer, por tanto, que $g$ es \emph{fielmente plano} y que $h$ es liso, y nos reducimos a probar que $f$ es entonces liso. Si $k'$ es una extensión algebraicamente

\oldpage[IV-4]{75}

cerrada de $k$, $X\otimes_k k'$ es entonces liso sobre $k'$, y en virtud de (17.7.3, (ii)), basta probar que $Y\otimes_k k'$ es liso sobre $k'$ (puesto que $\Spec(k')\to\Spec(k)$ es fielmente plano); podemos, por tanto, limitarnos al caso en que $k$ es \emph{algebraicamente cerrado}. Como el conjunto de los puntos de $Y$ racionales sobre $k$ es entonces muy denso en $Y$ (10.4.8) y como el conjunto de los puntos de $Y$ donde $Y$ es liso sobre $k$ es abierto, vemos que basta probar que $Y$ es liso sobre $k$ en todo punto racional sobre $k$. Pero, en tal punto $y$, decir que $Y$ es liso sobre $k$ en dicho punto equivale a decir que $Y$ es \emph{regular} en $y$ (17.5.1 y 6.7.8). Ahora bien, se tiene $y=g(x)$ para un $x\in X$ y, por hipótesis, (17.5.1) $X$ es regular en el punto $x$; como $X$ e $Y$ son entonces localmente noetherianos y $g$ es plano, $Y$ es, en efecto, regular en el punto $y$ (6.5.1, (i)).

(ii) \emph{Caso de los morfismos étale}. Por (i), ya sabemos que $f$ es liso en el punto $y$; en virtud de (17.6.1), basta, por tanto, mostrar que $f$ es cuasifinito en el punto $y$, o también que $\sh{O}_{Y,y}$ es una $k$-álgebra finita. Ahora bien, como $\sh{O}_{X,x}$ es un $\sh{O}_{Y,y}$-módulo fielmente plano ($0_{\mathrm I}$, 6.6.2), $\sh{O}_{Y,y}$ se identifica con un $k$-submódulo de $\sh{O}_{X,x}$ y, como $\sh{O}_{X,x}$ es por hipótesis una $k$-álgebra finita, lo mismo ocurre con $\sh{O}_{Y,y}$.

\begin{proposition}[17.7.8]
\label{IV.17.7.8}
Con las notaciones de (8.8.1), se supone que $X_\alpha$ e $Y_\alpha$ son localmente de presentación finita sobre $S_\alpha$. Sean $f_\alpha:X_\alpha\to Y_\alpha$ un $S_\alpha$-morfismo y $f:X\to Y$ el $S$-morfismo correspondiente.
\begin{enumerate}
  \item[(i)] Sean $x$ un punto de $X$ y $x_\lambda$ su proyección canónica en $X_\lambda$. Para que $f$ sea liso (resp. no ramificado, resp. étale) en el punto $x$, es necesario y suficiente que exista $\lambda\geq\alpha$ tal que $f_\lambda$ sea liso (resp. no ramificado, resp. étale) en el punto $x_\lambda$.
  \item[(ii)] Supongamos además que $X_\alpha$ es cuasicompacto. Para que $f$ sea liso (resp. no ramificado, resp. étale), es necesario y suficiente que exista $\lambda\geq\alpha$ tal que $f_\lambda$ sea liso (resp. no ramificado, resp. étale).
\end{enumerate}
\end{proposition}

(i) Si $y=f(x)$, $y_\lambda=f_\lambda(x_\lambda)$ es la proyección canónica de $y$ en $Y_\lambda$, y se tiene $f^{-1}(y)=f_\lambda^{-1}(y_\lambda)\otimes_{k(y_\lambda)}k(y)$; la parte del enunciado relativa a los morfismos no ramificados resulta, por tanto, de (17.7.1, (i)), y basta, pues, considerar el caso de los morfismos lisos. Como $f$ y $f_\lambda$ son localmente de presentación finita, equivale decir que $f^{-1}(y)$ es geométricamente regular en el punto $x$ o que $f_\lambda^{-1}(y_\lambda)$ es geométricamente regular en el punto $x_\lambda$ (6.7.8). La proposición resulta, por tanto, de (17.5.1) y de (11.2.6).

(ii) Para todo $\lambda$, sea $U_\lambda$ el conjunto abierto de los $x_\lambda\in X_\lambda$ tales que $f_\lambda$ sea liso (resp. no ramificado, resp. étale) en el punto $x_\lambda$; sea $V_\lambda$ su imagen inversa en $X$. Como, por hipótesis, para todo $x\in X$ existe, en virtud de (i), un $\lambda$ tal que $f_\lambda$ sea liso (resp. no ramificado, resp. étale) en el punto $x_\lambda$, $X$ es unión de los $V_\lambda$. Por otra parte (17.3.3), para $\lambda\leq\mu$ se tiene $V_\lambda\subset V_\mu$, de modo que, como $X$ es cuasicompacto, existe un índice $\mu$ tal que $X=V_\mu$. Como los $X_\lambda$ son cuasicompactos, resulta entonces de (8.3.4) que existe un índice $\nu\geq\mu$ tal que la imagen inversa de $U_\mu$ en $X_\nu$ sea todo $X_\nu$, lo que significa que $f_\nu$ es liso (resp. no ramificado, resp. étale) por (17.3.3).

\begin{corollary}[17.7.9]
\label{IV.17.7.9}
Sean $S=\Spec(A)$ un esquema afín y $f:X\to S$ un morfismo. Las condiciones siguientes son equivalentes:
\begin{enumerate}
  \item[a)] $f$ es un morfismo de presentación finita y liso (resp. no ramificado, resp. étale).

  \oldpage[IV-4]{76}

  \item[b)] Existen un esquema afín noetheriano $S_0=\Spec(A_0)$, un morfismo de tipo finito $f_0:X_0\to S_0$ y un morfismo $S\to S_0$ tales que el $S$-preesquema $X_0\otimes_{S_0}S$ sea $S$-isomorfo a $X$ y que $f_0$ sea liso (resp. no ramificado, resp. étale).
  \item[c)] Se verifican las condiciones de b), y además $A_0$ es una $\mathbb{Z}$-subálgebra de tipo finito de $A$, correspondiendo el morfismo $S\to S_0$ a la inyección canónica $A_0\to A$.
\end{enumerate}
\end{corollary}

La demostración a partir de (17.7.8) es la misma que la de (11.2.7) a partir de (11.2.6).

\begin{proposition}[17.7.10]
\label{IV.17.7.10}
Sean $f:Y\to S$, $h:X\to S$ dos morfismos localmente de presentación finita, $g:X\to Y$ un $S$-morfismo, $x$ un punto de $X$, $y=g(x)$. Supongamos que $h$ sea plano en el punto $x$ y que $f$ sea no ramificado en el punto $y$. Entonces $f$ es étale en el punto $y$ y $g$ es plano en el punto $x$.
\end{proposition}

(Más adelante se verá (18.4.9) que, de hecho, se puede prescindir de la hipótesis de que $h$ es localmente de presentación finita).

Como la cuestión es local sobre $S$, $X$ e $Y$, podemos suponer que $S$, $X$ e $Y$ son afines, que $f,g,h$ son morfismos de presentación finita, que $h$ es plano y que $f$ es no ramificado. Teniendo en cuenta (11.2.7) y (17.7.9), podemos suponer además que $S$, $X$ e $Y$ son noetherianos; finalmente, podemos limitarnos al caso en que $S=\Spec(\sh{O}_{S,s})$ (donde $s=f(y)=h(x)$). Como el anillo $A=\sh{O}_{S,s}$ es un anillo local noetheriano, existen un anillo local noetheriano $B$ completo, de cuerpo residual algebraicamente cerrado, y un homomorfismo local $A\to B$ que hace de $B$ un $A$-módulo fielmente plano ($0_{\mathrm{III}}$, 10.3.1). Reemplazando $X$ e $Y$ por $X\times_S\Spec(B)$ e $Y\times_S\Spec(B)$, se concluye por (2.5.1) y (17.7.1) que podemos limitarnos al caso en que $A$ es completo y tiene un cuerpo residual algebraicamente cerrado. La hipótesis de que $f$ es no ramificado en el punto $y$ implica entonces (17.4.4) que $\sh{O}_{Y,y}$ es una $\sh{O}_{S,s}$-álgebra finita y un anillo local noetheriano completo ($0_{\mathrm I}$, 7.4.2), y como el cuerpo residual de $\sh{O}_{S,s}$ es algebraicamente cerrado, el homomorfismo $\sh{O}_{S,s}\to\sh{O}_{Y,y}$ es sobreyectivo (17.4.4). Pero, por otra parte, la hipótesis de que $h$ es plano en el punto $x$ implica que el homomorfismo compuesto $\sh{O}_{S,s}\to\sh{O}_{Y,y}\to\sh{O}_{X,x}$ es inyectivo ($0_{\mathrm I}$, 6.5.1), y por tanto el homomorfismo $\sh{O}_{S,s}\to\sh{O}_{Y,y}$ es biyectivo, lo que muestra que $f$ es étale en el punto $y$ (17.6.3); además, $\sh{O}_{X,x}$ es un $\sh{O}_{Y,y}$-módulo plano, y por tanto $g$ es plano en el punto $x$.

\begin{proposition}[17.7.11]
\label{IV.17.7.11}
Sea $S$ un preesquema, sean $X$ e $Y$ dos $S$-preesquemas localmente de presentación finita sobre $S$, y sea $f:X\to Y$ un $S$-morfismo. Para todo $s\in S$, se denotan por $X_s$, $Y_s$, $f_s$ los preesquemas y el morfismo deducidos de $X$, $Y$, $f$ por el cambio de base $\Spec(k(s))\to S$. Entonces:
\begin{enumerate}
  \item[(i)] El conjunto de los $x\in X$ tales que, si $s$ es la imagen de $x$ en $S$, $f_s:X_s\to Y_s$ sea liso (resp. no ramificado, resp. étale, resp. diferencialmente liso) en el punto $x$, es localmente constructible.
  \item[(ii)] Supongamos que $f$ es de presentación finita. El conjunto de los $y\in Y$ tales que, si $s$ es la imagen de $y$ en $S$, $f_s$ sea liso (resp. no ramificado, resp. étale, resp. diferencialmente liso) en todos los puntos de $f_s^{-1}(y)$, es localmente constructible.
  \item[(iii)] Supongamos que $X$ e $Y$ son de presentación finita sobre $S$. El conjunto de los $s\in S$ tales que $f_s$ sea liso (resp. no ramificado, resp. étale, resp. diferencialmente liso) es localmente constructible.
\end{enumerate}
\end{proposition}

Sea $E$ el conjunto de los $x\in X$ que verifican la propiedad considerada en (i). Entonces el conjunto $F$ de los $y\in Y$ que verifican la propiedad correspondiente considerada en (ii) no es

\oldpage[IV-4]{77}

otro que $Y-f(X-E)$, por lo que (ii) resulta de (i) y del teorema de Chevalley cuando $f$ es de presentación finita (1.8.4). De igual modo, si $h:Y\to S$ es el morfismo estructural, el conjunto de los $s\in S$ que verifican la propiedad correspondiente considerada en (iii) es $S-h(Y-F)$, por lo que (iii) se deduce también de (ii) y del teorema de Chevalley cuando $f$ y $h$ son de presentación finita. Basta, por tanto, probar las afirmaciones de (i).

Probemos primero (i) cuando se trata de la propiedad de ser \emph{liso}. Como la cuestión es local sobre $X$, podemos limitarnos al caso en que $S=\Spec(A)$, $X=\Spec(B)$ e $Y=\Spec(C)$ son afines, siendo $B$ y $C$ $A$-álgebras de presentación finita. Razonando como al principio de (9.9.1) y utilizando (17.7.2, (ii)), nos reducimos al caso en que $A$ es noetheriano. En virtud de ($0_{\mathrm{III}}$, 9.2.3), nos reducimos a ver que, si $x\in E$ (resp. si $x\notin E$), existe una vecindad $V$ de $x$ en $\overline{\{x\}}$ contenida en $E$ (resp. en $X-E$). Denotando por $g:X\to S$ y $h:Y\to S$ los morfismos estructurales, podemos primero reemplazar $S$ por el subpreesquema reducido $S'$ de $S$ que tiene $\overline{\{g(x)\}}$ por espacio subyacente, $X$ por $g^{-1}(S')$, e $Y$ por $h^{-1}(S')$, pues las fibras de $X$ y $X'$ (resp. $Y$ e $Y'$) en los puntos de $S'$ son las mismas. Dicho de otro modo, podemos limitarnos al caso en que $S$ es \emph{íntegro} y en que $\eta=g(x)=h(y)$ (donde $y=f(x)$) es su punto genérico.

$1^\circ$ Supongamos primero que $x\in E$. Los anillos locales $\sh{O}_{X_\eta,x}$ y $\sh{O}_{Y_\eta,y}$ son respectivamente iguales a $\sh{O}_{X,x}$ y $\sh{O}_{Y,y}$; como la propiedad de lisitud de un morfismo de presentación finita en un punto solo depende del anillo local de dicho punto y del anillo local de su imagen (17.5.1), vemos que la hipótesis $x\in E$ equivale a decir que el morfismo $f$ es liso en el punto $x$; conserva entonces esta propiedad en los puntos de una vecindad abierta de $x$ en $X$, y basta aplicar (17.3.3, (iii)) para obtener la conclusión.

$2^\circ$ Supongamos, en segundo lugar, que $x\in X-E$, y que el morfismo $f_\eta$ no sea plano en el punto $x$. La conclusión resulta entonces del lema siguiente, que precisa (11.2.8):

\begin{lemma}[17.7.11.1]
\label{IV.17.7.11.1}
Sean $g:X\to S$, $h:Y\to S$ dos morfismos localmente de presentación finita, $f:X\to Y$ un $S$-morfismo, y $\sh{F}$ un $\sh{O}_X$-módulo cuasicoherente de presentación finita. Entonces el conjunto $E$ de los $x\in X$ tales que $\sh{F}_{g(x)}$ sea $f_{g(x)}$-plano en el punto $x$ es localmente constructible.
\end{lemma}

Razonando una vez más como al principio de (9.9.1) y utilizando (2.5.1), nos reducimos al caso en que $S$, $X$ e $Y$ son noetherianos; después nos reducimos, como antes, al caso en que $S$ es íntegro, en que $\eta=g(x)=h(f(x))$ es su punto genérico, y se trata de mostrar que si $x\in E$ (resp. $x\notin E$) existe una vecindad $V$ de $x$ en $\overline{\{x\}}$ contenida en $E$ (resp. en $X-E$). El caso en que $x\in E$ es consecuencia inmediata de (11.1.1). Para tratar el caso en que $x\notin E$, se razona como en (9.4.7.1), cuyas notaciones conservamos, de modo que podemos suponer, reemplazando eventualmente $Y$ y $X$ por vecindades de $f(x)$ y $x$, respectivamente, que existen dos $\sh{O}_Y$-módulos coherentes $\sh{G}$, $\sh{H}$ y un $\sh{O}_Y$-homomorfismo $u:\sh{G}\to\sh{H}$ tales que, para todo $s\in S$, $u_s:\sh{G}_s\to\sh{H}_s$ sea inyectivo, pero que el homomorfismo $1\otimes u_\eta:\sh{F}_\eta\otimes_{\sh{O}_{Y_\eta}}\sh{G}_\eta\to\sh{F}_\eta\otimes_{\sh{O}_{Y_\eta}}\sh{H}_\eta$ no sea inyectivo en el punto $x$, dicho de otro modo, $x\in\operatorname{Supp}(\operatorname{Ker}(1\otimes u_\eta))$; si se pone $T=\overline{\{x\}}$, se tiene, por tanto, $\operatorname{Supp}(\operatorname{Ker}(1\otimes u_\eta))\supset T_\eta$. Pero podemos suponer que, para todo $s\in S$, se tiene $\operatorname{Ker}(1\otimes u_s)=(\operatorname{Ker}(1\otimes u))_s$ (9.4.2), y por tanto

\oldpage[IV-4]{78}

$\operatorname{Supp}(\operatorname{Ker}(1\otimes u_s))=(\operatorname{Supp}(\operatorname{Ker}(1\otimes u)))_s$ (I, 9.1.13.1); resulta finalmente de (9.5.2) que, para $s$ en una vecindad de $\eta$, se tiene $(\operatorname{Supp}(\operatorname{Ker}(1\otimes u)))_s\supset T_s$, lo que demuestra el lema.

$3^\circ$ Supongamos ahora que $x\in X-E$, que el morfismo $f_\eta$ sea plano en el punto $x$, pero que $f_\eta$ no sea liso en el punto $x$. Observemos que decir que $f_\eta$ es plano en el punto $x$ equivale a decir que el propio $f$ es plano en el punto $x$ y, reemplazando $X$ por una vecindad de $x$, podemos suponer que $f$ es plano (11.1.1); de ello se concluye que lo mismo ocurre con $f_s$ para todo $s\in S$, y como para todo $y\in Y$, $f^{-1}(y)=f_{h(y)}^{-1}(y)$, equivale decir que $f_{g(x')}$ es liso en el punto $x'$ o decir que $f$ es liso en el punto $x'$. Pero el conjunto de los $x'\in X$ donde $f$ es liso es \emph{abierto} en $X$ (12.1.7), por lo que el conjunto de los $x'\in X$ donde $f$ no es liso es cerrado, y puesto que contiene a $x$ por hipótesis, también contiene a $\overline{\{x\}}$, lo que termina la demostración para la primera propiedad considerada en (i).

Probemos, en segundo lugar, (i) cuando se trata de la propiedad de ser \emph{étale}. Observemos para ello que esta propiedad para $f_s$ en el punto $x$ equivale a decir que $f_s$ es a la vez liso y cuasifinito en el punto $x$ (17.6.1). Ahora bien, equivale decir que $f_{g(x)}$ es cuasifinito en el punto $x$ o que el propio $f$ es cuasifinito en dicho punto; resulta, por tanto, de (13.1.4) que el conjunto de los puntos $x$ tales que $f_{g(x)}$ sea cuasifinito en el punto $x$ es abierto en $X$, y \emph{a fortiori} localmente constructible; la conclusión resulta entonces de que el conjunto de los $x$ tales que $f_{g(x)}$ sea liso en el punto $x$ es también localmente constructible.

Pasemos a la demostración de (i) cuando se trata de la propiedad de ser \emph{diferencialmente liso}. Sea $p:X\times_Y X\to X$ la segunda proyección canónica; la segunda proyección $X_s\times_{Y_s}X_s\to X_s$ no es otra que $p_s$ para todo $s\in S$, y resulta, por tanto, de (17.12.5.1)\footnote{El lector comprobará que el resultado de (17.7.11) no se utiliza en lo que sigue del § 17 y que, por tanto, no hay círculo vicioso.} que para que $f_{g(x)}$ sea diferencialmente liso en el punto $x$, es necesario y suficiente que $p_{g(x)}$ sea liso en el punto $\Delta_f(x)$. Como $p$ es localmente de presentación finita, el conjunto de los puntos $z\in X\times_Y X$ tales que $p_{g(p(z))}$ sea liso en el punto $z$ es localmente constructible, y la intersección de este conjunto con el conjunto localmente cerrado $\Delta_f(X)$ es, por tanto, también localmente constructible en $\Delta_f(X)$ (1.8.2); como la restricción de $p$ a $\Delta_f(X)$ es un isomorfismo sobre $X$, vemos que el conjunto de los $x\in X$ tales que $f_{g(x)}$ sea diferencialmente liso en el punto $x$ es localmente constructible.

Consideremos, finalmente, la propiedad de ser no ramificado; observemos que el morfismo diagonal $\Delta_f:X\to X\times_Y X$ es una inmersión localmente de presentación finita (1.4.3.1) y, para todo $s\in S$, el morfismo diagonal $\Delta_{f_s}:X_s\to X_s\times_{Y_s}X_s$ no es otro que $(\Delta_f)_s$; decir que $f_{g(x)}$ es no ramificado en el punto $x$ equivale a decir que $(\Delta_f)_{g(x)}$ es un isomorfismo local en el punto $x$ (17.4.1), y puesto que es una inmersión localmente de presentación finita, equivale (17.9.1)\footnote{El lector comprobará que el resultado de (17.7.11) no se utiliza en el resto del § 17 y que, por tanto, no hay círculo vicioso.} a decir que $(\Delta_f)_{g(x)}$ es étale en el punto $x$; basta, por tanto, aplicar lo visto más arriba para la propiedad de ser étale.

\oldpage[IV-4]{79}

\subsection{Criterios de lisitud y no ramificación por las fibras}
\label{subsection:IV.17.8}
\label{IV.17.8}

\begin{proposition}[17.8.1]
\label{IV.17.8.1}
Sean $g:Y\to S$, $h:X\to S$ dos morfismos localmente de presentación finita. Para que un $S$-morfismo $f:X\to Y$ sea no ramificado, es necesario y suficiente que, para todo $s\in S$, el morfismo $f_s:h^{-1}(s)\to g^{-1}(s)$ deducido de $f$ por el cambio de base $\Spec(k(s))\to S$ sea no ramificado.
\end{proposition}

Esto resulta inmediatamente del hecho de que, para un morfismo localmente de presentación finita, ser no ramificado es una propiedad de las \emph{fibras} de dicho morfismo (17.4.1, $d$)), y de que, para todo $y\in Y$, se tiene $f^{-1}(y)=f_s^{-1}(y)$ si $s=g(y)$.

\begin{proposition}[17.8.2]
\label{IV.17.8.2}
Sean $g:Y\to S$, $h:X\to S$ dos morfismos localmente de presentación finita; se supone además que $h$ es plano. Para que un $S$-morfismo $f:X\to Y$ sea liso (resp. étale), es necesario y suficiente que, para todo $s\in S$, el morfismo $f_s:h^{-1}(s)\to g^{-1}(s)$ deducido de $f$ por el cambio de base $\Spec(k(s))\to S$ sea liso (resp. étale). Cuando esto ocurre, el morfismo $g$ es plano en los puntos de $f(X)$.
\end{proposition}

En efecto, se sabe (11.3.10) que, para que $f$ sea plano, es necesario y suficiente que $f_s$ lo sea para todo $s\in S$, y que entonces $g$ es plano en los puntos de $f(X)$. Pero, para un morfismo plano localmente de presentación finita, ser liso es una propiedad de las fibras de dicho morfismo (17.5.1, $b$)), y para todo $y\in Y$, se tiene $f^{-1}(y)=f_s^{-1}(y)$ si $s=g(y)$.

\begin{remark}[17.8.3]
\label{IV.17.8.3}
Las demostraciones anteriores muestran (teniendo en cuenta (11.3.10)) que, bajo las mismas hipótesis anteriores sobre $g$ y $h$, para que $f$ sea no ramificado (resp. liso, resp. étale) \emph{en un punto} $x\in X$, basta que, si se pone $s=h(x)$, $f_s$ sea no ramificado (resp. liso, resp. étale) en el punto $x$.
\end{remark}

\subsection{Morfismos étale e inmersiones abiertas}
\label{subsection:IV.17.9}
\label{IV.17.9}

\begin{theorem}[17.9.1]
\label{IV.17.9.1}
Sea $f:X\to Y$ un morfismo. Las propiedades siguientes son equivalentes:
\begin{enumerate}
  \item[a)] $f$ es una inmersión abierta.
  \item[b)] $f$ es un monomorfismo plano localmente de presentación finita.
  \item[c)] $f$ es étale y radicial.
\end{enumerate}
\end{theorem}

Resulta de (1.4.3, (i)) que a) implica b). La condición b) implica que, para todo $y\in Y$, la fibra $f^{-1}(y)$ es vacía o isomorfa a $\Spec(k(y))$ (8.11.5.1), de modo que b) implica c), en virtud de (17.6.2, $c$)). Queda por ver que c) implica a).

Como la cuestión es local sobre $X$ y sobre $Y$ (puesto que $f$ es inyectivo), podemos limitarnos al caso en que $Y$ es afín y $f$ de presentación finita. Como $f$ es plano, es un morfismo \emph{abierto} (2.4.6), por lo que, reemplazando $Y$ por $f(X)$, podemos suponer que $f$ es \emph{sobreyectivo}. Para todo morfismo $Y'\to Y$, $f'=f_{(Y')}:X_{(Y')}\to Y'$ sigue siendo étale, radicial, sobreyectivo y de presentación finita, y por tanto abierto, y en consecuencia un homeomorfismo; dicho de otro modo, $f$ es un \emph{homeomorfismo universal}, y, al ser de tipo finito y separado (1.8.7.1), $f$ es \emph{propio}. La hipótesis de que $f$ es radicial y de tipo finito implica entonces que $f$ es cuasifinito; por tanto, (8.11.1) $f$ es un morfismo \emph{finito}. Para probar que $f$ es un isomorfismo, podemos limitarnos al caso

\oldpage[IV-4]{80}

en que $Y=\Spec(A)$, siendo $A$ un anillo local. Como $f$ es de presentación finita, se tiene $X=\Spec(B)$, donde $B$ es un $A$-módulo plano de presentación finita (1.4.7), y por tanto \emph{libre} (Bourbaki, \emph{Alg. comm.}, cap. II, § 5, n\textsuperscript{o} 2, cor. 2 del teorema 1). Además, si $\mathfrak m$ es el ideal maximal de $A$ y $k$ su cuerpo residual, $B/\mathfrak mB$ es por hipótesis un cuerpo, a la vez extensión radicial y extensión finita separable de $k$, puesto que $f$ es étale y radicial (17.6.1); por tanto, $B/\mathfrak mB$ es isomorfo a $k$, y como $B$ es un $A$-módulo libre, $B$ es isomorfo a $A$. C.Q.D.

\begin{corollary}[17.9.2]
\label{IV.17.9.2}
Sea $X$ un preesquema conexo; si $f:X\to Y$ es una inmersión cerrada étale, entonces $f$ es un isomorfismo de $X$ sobre una componente conexa abierta de $Y$.
\end{corollary}

En efecto, $f$ es étale y radicial, y por tanto un isomorfismo de $X$ sobre un preesquema inducido en un subconjunto abierto de $Y$; pero, por hipótesis, $f(X)$ es cerrado en $Y$, y por tanto a la vez abierto y cerrado, y como $f(X)$ es conexo, es una componente conexa de $Y$.

\begin{corollary}[17.9.3]
\label{IV.17.9.3}
Sea $f:X\to Y$ un morfismo étale (resp. étale y separado). Entonces toda $Y$-sección $g:Y\to X$ de $X$ es una inmersión abierta (resp. abierta y cerrada); además, la aplicación $g\mapsto g(Y)$ es una biyección del conjunto $\Gamma(X/Y)$ de las $Y$-secciones de $X$ sobre el conjunto de los subconjuntos abiertos $Z$ (resp. abiertos y cerrados) de $X$ tales que la restricción de $f$ a $Z$ sea un morfismo sobreyectivo y radicial de $Z$ sobre $Y$.
\end{corollary}

En efecto, el hecho de que $g$ sea una inmersión abierta resulta ya de que $f$ es no ramificado (17.4.1, $b''$)), y la restricción de $f$ al abierto $g(Y)$ de $X$ es un isomorfismo. Recíprocamente, si $Z$ es un abierto de $X$ tal que $f|Z$ sea un morfismo sobreyectivo y radicial de $Z$ sobre $Y$, $f|Z$ es un isomorfismo en virtud de (17.9.1), puesto que es étale. Si $f$ es además separado, se sabe que $g$ es una inmersión cerrada (1, 5.4.6), lo que termina de probar el corolario.

\begin{corollary}[17.9.4]
\label{IV.17.9.4}
Sean $Y$ un preesquema conexo y $f:X\to Y$ un morfismo étale y separado. Entonces toda $Y$-sección $g$ de $X$ es un isomorfismo de $Y$ sobre una componente conexa abierta de $X$, y la aplicación $g\mapsto g(Y)$ es una biyección de $\Gamma(X/Y)$ sobre el conjunto de las componentes conexas abiertas $Z$ de $X$ tales que la restricción de $f$ a $Z$ sea un morfismo sobreyectivo y radicial de $Z$ sobre $Y$.
\end{corollary}

\begin{corollary}[17.9.5]
\label{IV.17.9.5}
Sean $g:Y\to S$, $h:X\to S$ dos morfismos localmente de presentación finita; se supone además que $h$ es plano. Para que un $S$-morfismo $f:X\to Y$ sea una inmersión abierta (resp. un isomorfismo), es necesario y suficiente que, para todo $s\in S$, el morfismo $f_s:h^{-1}(s)\to g^{-1}(s)$ deducido de $f$ por el cambio de base $\Spec(k(s))\to S$ sea una inmersión abierta (resp. un isomorfismo).
\end{corollary}

En efecto, si $f_s$ es una inmersión abierta para todo $s\in S$, resulta de (17.8.3) que $f$ es un morfismo étale; como, para todo $y\in Y$, se tiene $f^{-1}(y)=f_s^{-1}(y)$ con $s=g(y)$, $f$ es radicial; por tanto, $f$ es una inmersión abierta en virtud de (17.9.1). Si, además, $f_s$ es sobreyectivo para todo $s\in S$, $f$ es sobreyectivo, y por tanto un isomorfismo.

La proposición siguiente precisa (10.4.11):

\begin{proposition}[17.9.6]
\label{IV.17.9.6}
Sean $S$ un preesquema y $X$ un $S$-preesquema de presentación finita. Todo $S$-endomorfismo de $X$ que sea un monomorfismo es un automorfismo de $X$.
\end{proposition}

Sean $f:X\to S$ el morfismo estructural y $g$ el $S$-endomorfismo considerado. La cuestión es local sobre $S$, por lo que podemos suponer que $S$ es afín. Utilizando (8.9.1) y (8.10.5, (i \emph{bis})), nos reducimos al caso en que $S=\Spec(A)$, donde $A$ es una $\mathbb{Z}$-álgebra de tipo finito, y por tanto $X$ es un $\mathbb{Z}$-preesquema de tipo finito. Resulta ya de (10.4.11) que $g$ es un morfismo biyectivo, puesto que un monomorfismo es radicial (8.11.5.1); bastará, por tanto, mostrar

\oldpage[IV-4]{81}

que $g$ es una inmersión abierta, y, puesto que $g$ es radicial, bastará, en virtud de (17.9.1), probar que $g$ es étale. Además, como el conjunto de los puntos donde $g$ es étale es abierto y $X$ es un preesquema de Jacobson (10.4.7), bastará mostrar que $g$ es étale en todo punto \emph{cerrado} $z$ de $X$ (10.3.1). Pongamos $z'=g(z)$; resulta de (10.4.11.1, (i)) que $z'$ también es un punto cerrado de $X$ y, puesto que $g$ es un monomorfismo, la aplicación $k(z')\to k(z)$ deducida de $g$ es un \emph{isomorfismo}. Para que $g$ sea étale en el punto $z$, es necesario y suficiente, por tanto, que el homomorfismo canónico $\widehat{\sh{O}}_{X,z'}\to\widehat{\sh{O}}_{X,z}$ sea biyectivo (17.6.3, $e''$)). Para ello, probaremos que, para todo entero $n$, el homomorfismo canónico $\sh{O}_{X,z'}/\mathfrak m_{z'}^{n+1}\to\sh{O}_{X,z}/\mathfrak m_z^{n+1}$ es biyectivo, de donde resultará inmediatamente la conclusión. Podemos suponer que $z$ pertenece al conjunto finito $T_{p,d}$ de los puntos cerrados $t$ de $X$ tales que $k(t)$ sea una extensión de $\mathbb{F}_p$ cuyo grado divide a $d$; en cuyo caso lo mismo ocurre con $z'$. Designemos por $\sh{J}$ el ideal coherente de $\sh{O}_X$ tal que $\sh{J}_x=\sh{O}_x$ para $x\notin T_{p,d}$, y $\sh{J}_x=\mathfrak m_x^{n+1}$ para $x\in T_{p,d}$. Sean $T_{p,d,n}$ el subpreesquema cerrado de $X$ definido por $\sh{J}$, y $j:T_{p,d,n}\to X$ la inyección canónica; el morfismo compuesto $g\circ j:T_{p,d,n}\to X$ lleva el conjunto $T_{p,d}$ a sí mismo, y para todo $x\in T_{p,d}$, el homomorfismo $\sh{O}_{X,g(x)}\to\sh{O}_{X,x}/\mathfrak m_x^{n+1}$ se factoriza como
\[
\sh{O}_{X,g(x)}\to\sh{O}_{X,g(x)}/\mathfrak m_{g(x)}^{n+1}\to\sh{O}_{X,x}/\mathfrak m_x^{n+1}.
\]
Por tanto (I, 4.9), existe un único endomorfismo $g'$ de $T_{p,d,n}$ tal que $j\circ g'=g\circ j$. Como $j$ es un monomorfismo, y lo mismo ocurre con $g$ por hipótesis, se deduce que $g'$ también es un monomorfismo, y la proposición quedará probada si se muestra que $g'$ es un \emph{automorfismo} de $T_{p,d,n}$. Ahora bien, para todo $x\in T_{p,d,n}$, hemos visto que $k(x)$ es un cuerpo finito, y como $\sh{O}_{X,x}$ es noetheriano, cada uno de los $\mathfrak m_x^h/\mathfrak m_x^{h+1}$ es un $k(x)$-espacio vectorial de dimensión finita; de donde se concluye inmediatamente que el anillo local $\sh{O}_{X,x}/\mathfrak m_x^{n+1}$ de $T_{p,d,n}$ en el punto $x$ tiene un número \emph{finito} de elementos y, \emph{a fortiori}, es un $\mathbb{Z}$-módulo de tipo finito. Como $T_{p,d,n}$ es suma de un número finito de preesquemas $\Spec(\sh{O}_{X,x}/\mathfrak m_x^{n+1})$, es un $\mathbb{Z}$-preesquema \emph{finito}; el endomorfismo $g'$ es, por tanto, también \emph{finito} (II, 6.1.5, (v)), y por consiguiente propio. Pero un monomorfismo propio es una \emph{inmersión cerrada} (8.11.5); si $T_{p,d,n}=\Spec(B)$, $g'$ corresponde entonces a un endomorfismo sobreyectivo $\varphi:B\to B$ del anillo $B$. Ahora bien, el conjunto $B$ es finito, por lo que $\varphi$ es necesariamente biyectivo. C.Q.D.

\subsection{Dimensión relativa de un preesquema liso sobre otro}
\label{subsection:IV.17.10}

\begin{definition}[17.10.1]
\label{IV.17.10.1}
Sea $f:X\to Y$ un morfismo localmente de tipo finito. Se llama \emph{dimensión relativa de $f$ en el punto $x\in X$} (o \emph{dimensión relativa de $X$ sobre $Y$ en el punto $x$}) y se denota por $\dim_x f$ el entero no negativo $\dim_x(f^{-1}(f(x)))$.
\end{definition}

Decir que $f$ es \emph{cuasifinito} en el punto $x$ (Err$_{\mathrm{III}}$, 20) equivale, por tanto, a decir que $\dim_x f=0$. Hemos visto (13.1.3) que la función $x\mapsto\dim_x f$ es semicontinua superiormente. Obsérvese que, incluso cuando el morfismo $f$ posee la propiedad $(S_1)$ (es decir, por (6.8.1), es plano y sus fibras no tienen ningún ciclo primo asociado inmerso), la función $x\mapsto\dim_x f$ no es necesariamente continua, como muestra el ejemplo en que $Y=\Spec(k)$, donde $k$ es un cuerpo, y $X=\Spec(k[U,V,W]/\mathfrak p\mathfrak q)$, donde $\mathfrak p=(W)$ y $\mathfrak q=(U)+(V-W)$ son ideales primos de $k[U,V,W]$ (siendo, por tanto, $X$ la unión, en el espacio de 3 dimensiones, de un plano y una recta no paralela a dicho plano).

Decir que un morfismo localmente de presentación finita $f:X\to Y$ es étale en el punto $x\in X$ significa asimismo que $f$ es liso en el punto $x$ y que $\dim_x f=0$ (17.6.1).

\begin{proposition}[17.10.2]
\label{IV.17.10.2}
Sea $f:X\to Y$ un morfismo liso. Para todo $x\in X$, el $\sh{O}_X$-módulo localmente libre $\Omega_f^1$ (17.2.3) tiene en $x$ un rango igual a $\dim_x f$ (lo que implica que $x\mapsto\dim_x f$ es una función continua en $X$).
\end{proposition}

En efecto, si $y=f(x)$, $X_y=f^{-1}(y)$ es liso sobre $k(y)$, y si $f_y:X_y\to\Spec(k(y))$ es el morfismo estructural, se tiene $\Omega_{f_y}^1=\Omega_f^1\otimes_{\sh{O}_X}\sh{O}_{X_y}$ (16.4.5); podemos, por tanto, limitarnos al caso en que $Y=\Spec(k)$ es el espectro de un cuerpo; además, en virtud de (16.4.5) y (17.7.1), podemos reemplazar $k$ por una extensión algebraicamente cerrada, dicho de otro modo,

\oldpage[IV-4]{82}

suponer que $k$ es algebraicamente cerrado. Como el conjunto de puntos de $X$ racionales sobre $k$ es entonces denso en $X$ (10.4.8), podemos limitarnos al caso en que $x$ es racional sobre $k$. Ahora bien, (16.4.12) $(\Omega^1_{X/k})_x\otimes_{\sh{O}_{X,x}}k(x)$ es entonces $k$-isomorfo a $\mathfrak m_x/\mathfrak m_x^2$, y puesto que $\sh{O}_x$ es un anillo local regular (17.5.1), $\operatorname{rg}_{k(x)}(\mathfrak m_x/\mathfrak m_x^2)=\dim(\sh{O}_x)$ (0, 17.1.1); por tanto, $\operatorname{rg}_{\sh{O}_x}(\Omega^1_{X/k})_x=\dim(\sh{O}_x)$. Pero, puesto que $k(x)=k$, se tiene $\dim_x f=\dim(\sh{O}_x)$ (5.2.3). C.Q.D.

Más adelante demostraremos un recíproco de este resultado (17.15.5).

\begin{corollary}[17.10.3]
\label{IV.17.10.3}
Sean $f:X\to Y$, $g:Y\to Z$ dos morfismos lisos. Entonces, para todo $x\in X$, se tiene
\begin{equation}
\label{IV.17.10.3.1}
\tag{17.10.3.1}
\dim_x(g\circ f)=\dim_x f+\dim_{f(x)}g.
\end{equation}
\end{corollary}

En efecto, $g\circ f$ es liso (17.3.3), por lo que los tres $\sh{O}_X$-módulos $\Omega^1_{X/Y}$, $\Omega^1_{X/Z}$ y $f^*(\Omega^1_{Y/Z})$ son localmente libres (17.2.3 y $0_{\mathrm I}$, 5.4.5); además, el rango en $x$ de $f^*(\Omega^1_{Y/Z})$ es igual al rango en $y=f(x)$ de $\Omega^1_{Y/Z}$. La igualdad (17.10.3.1) es, por tanto, consecuencia de (17.10.2) y de la exactitud de la sucesión (17.2.3.1).

\begin{corollary}[17.10.4]
\label{IV.17.10.4}
Sean $f:X\to Y$ un morfismo liso y $X'$ un subpreesquema de $X$ tal que el morfismo compuesto $X'\xrightarrow{j}X\xrightarrow{f}Y$ (donde $j$ es la inyección canónica) sea liso. Entonces el haz conormal $\sh{N}_{X'/X}$ es un $\sh{O}_{X'}$-módulo localmente libre, y para todo $x\in X'$, se tiene
\begin{equation}
\label{IV.17.10.4.1}
\tag{17.10.4.1}
\dim_x f=\dim_x(f\circ j)+\operatorname{rg}_{\sh{O}_{X',x}}(\sh{N}_{X'/X})_x.
\end{equation}
\end{corollary}

En efecto, $\Omega^1_{X/Y}\otimes_{\sh{O}_X}\sh{O}_{X'}$ y $\Omega^1_{X'/Y}$ son entonces localmente libres y la sucesión exacta (17.2.5.1) se escinde en una vecindad adecuada de cada punto de $X'$, por lo que $\sh{N}_{X'/X}$ es localmente libre (Bourbaki, \emph{Alg. comm.}, cap. II, § 5, n\textsuperscript{o} 2, teorema 1), y la relación (17.10.4.1) resulta inmediatamente de la exactitud de la sucesión (17.2.5.1).

\subsection{Morfismos lisos de preesquemas lisos}
\label{subsection:IV.17.11}

\begin{theorem}[17.11.1]
\label{IV.17.11.1}
Sean $f:Y\to S$ y $h:X\to S$ dos morfismos localmente de presentación finita, $g:X\to Y$ un $S$-morfismo y $x$ un punto de $X$; pongamos $y=g(x)$, $s=f(y)=h(x)$. Las condiciones siguientes son equivalentes:
\begin{enumerate}
  \item[a)] $f$ es liso en el punto $y$ y $g$ es liso en el punto $x$.
  \item[b)] $g$ y $h$ son lisos en el punto $x$.
  \item[c)] $h$ es liso en el punto $x$, y el homomorfismo canónico (16.4.18)
  \begin{equation}
  \label{IV.17.11.1.1}
  \tag{17.11.1.1}
  (g^*(\Omega^1_{Y/S}))_x\to(\Omega^1_{X/S})_x
  \end{equation}
  admite un inverso por la izquierda (dicho de otro modo, es un isomorfismo sobre un \emph{factor directo} de $(\Omega^1_{X/S})_x$).
  \item[c')] $h$ es liso en el punto $x$, y el homomorfismo canónico
  \begin{equation}
  \label{IV.17.11.1.2}
  \tag{17.11.1.2}
  ((\Omega^1_{Y/S})_y\otimes_{\sh{O}_{Y,y}}k(y))\otimes_{k(y)}k(x)\to(\Omega^1_{X/S})_x\otimes_{\sh{O}_{X,x}}k(x)
  \end{equation}
  es inyectivo.
\end{enumerate}
Supongamos, además, que el homomorfismo $k(y)\to k(x)$ sea biyectivo. Entonces las condiciones anteriores son también equivalentes a la siguiente:
\begin{enumerate}
  \item[d)] $h$ es liso en el punto $x$, y la aplicación canónica $\mathrm{T}_{X/S}(x)\to\mathrm{T}_{Y/S}(y)$ que va del espacio vectorial tangente en $x$ a $X$ al espacio vectorial tangente en $y$ a $Y$ (16.5.12) es sobreyectiva.
\end{enumerate}
\end{theorem}

El hecho de que a) implique b) resulta trivialmente de (17.3.3, (iii)); b) implica c) por aplicación de (17.2.3, (ii)). Para ver que c) equivale a c'), observemos que, en virtud de (17.2.3, (i)), $\Omega^1_{X/S}$ es un $\sh{O}_X$-módulo localmente libre de tipo finito; basta entonces aplicar (0, 19.1.12) al anillo local $\sh{O}_{X,x}$ y al homomorfismo (17.11.1.1) de $\sh{O}_{X,x}$-módulos de tipo finito. Cuando $k(x)=k(y)$, la aplicación lineal tangente $\mathrm{T}_{X/S}(x)\to\mathrm{T}_{Y/S}(y)$ es la transpuesta de (17.11.1.2), por (16.5.12), lo que establece en este caso la equivalencia de c') y d).

\oldpage[IV-4]{83}

Queda, por tanto, probar que c) implica a). Podemos limitarnos al caso en que $S=\Spec(A)$, $Y=\Spec(B)$, $X=\Spec(C)$ son afines. La hipótesis implica (17.2.3, (i)) que $(\Omega^1_{X/S})_x=\Omega^1_{C_{\mathfrak m}/A_{\mathfrak s}}$ es un $C_{\mathfrak m}$-módulo libre: en efecto, (16.10.6) asegura que existen elementos $t_i$ ($1\leq i\leq r$) de $B_{\mathfrak p}=\sh{O}_{Y,y}$ tales que las diferenciales $d_{B_{\mathfrak p}/A_{\mathfrak s}}(t_i)$ generan el $B_{\mathfrak p}$-módulo $\Omega^1_{B_{\mathfrak p}/A_{\mathfrak s}}$ y que sus imágenes en $\Omega^1_{C_{\mathfrak m}/A_{\mathfrak s}}$ forman parte de una base de este $C_{\mathfrak m}$-módulo libre. Como $\Omega^1_{Y/S}$ (resp. $\Omega^1_{X/S}$) es un $\sh{O}_Y$-módulo (resp. un $\sh{O}_X$-módulo) de presentación finita (16.4.22), podemos, reemplazando si es preciso $X$ e $Y$ por vecindades abiertas afines adecuadas de $x$ e $y$, respectivamente, suponer que $\Omega^1_{C/A}$ es un $C$-módulo libre y que los $t_i$ son las imágenes de elementos $s_i$ ($1\leq i\leq r$) de $B$ tales que los $d_{B/A}(s_i)$ generan el $B$-módulo $\Omega^1_{B/A}$ y que sus imágenes en $\Omega^1_{C/A}$ forman parte de una base de este $C$-módulo (Bourbaki, \emph{Alg. comm.}, cap. II, § 5, n\textsuperscript{o} 1, prop. 2). Sea $\varphi$ el $A$-homomorfismo de $B'=A[T_1,\ldots,T_r]$ en $B$ tal que $\varphi(T_i)=s_i$ para todo $i$; el dihomomorfismo correspondiente $\Omega^1_{B'/A}\to\Omega^1_{B/A}$ (0, 20.5.2) envía los $d_{B'/A}(T_i)$, que forman una base de $\Omega^1_{B'/A}$ (0, 20.4.13), a los $d_{B/A}(s_i)$ y es, por tanto, \emph{sobreyectivo}; si $Y'=\Spec(B')$, y si $u:Y\to Y'$ es el $S$-morfismo correspondiente a $\varphi$, se concluye de (17.2.2) que $u$ es \emph{no ramificado}. Si se prueba que el morfismo compuesto $u\circ g:X\to Y'$ es liso en el punto $x$, se seguirá de (17.7.10) que $u$ es étale en el punto $y$, y después de (17.3.5) que $g$ es liso en el punto $x$; finalmente, como el morfismo estructural $f':Y'\to S$ es liso (17.3.8), $f=f'\circ u$ será liso en el punto $y$. Como las imágenes canónicas de los $d_{B'/A}(T_i)$ en $\Omega^1_{C/A}$ son las de los $d_{B/A}(s_i)$, forman parte de una base de $\Omega^1_{C/A}$. Vemos así que, para probar que c) implica a), podemos limitarnos al caso en que $Y'=Y$, y por tanto suponer que $f$ es un morfismo \emph{liso}.

En virtud de (17.8.2), podemos entonces limitarnos al caso en que $S=\Spec(k)$ es el espectro de un cuerpo; además, gracias a (17.7.1, (ii)), podemos suponer que $k$ es algebraicamente cerrado; finalmente, reemplazando si es preciso $X$ e $Y$ por vecindades abiertas de $x$ e $y$, respectivamente, podemos suponer que $f$ y $h$ son ambos lisos, y que el homomorfismo canónico
\[
g^*(\Omega^1_{Y/S})\to\Omega^1_{X/S}
\]
admite un inverso por la izquierda (0, 19.1.12). Entonces el conjunto de los puntos de $X$ racionales sobre $k$ es muy denso en $X$ (10.4.8), de modo que, para probar que $g$ es liso, basta probar que $g$ es liso en todo punto $x\in X$ racional sobre $k$, o también que, en tal punto, $\sh{O}_{X,x}$ es un $\sh{O}_{Y,y}$-módulo plano y que $\sh{O}_{X,x}/\mathfrak m_y\sh{O}_{X,x}$ es un anillo regular (17.5.1). Ahora bien, $y=g(x)$ es

\oldpage[IV-4]{84}

\emph{a fortiori} racional sobre $k$, y como $\sh{O}_{Y,y}$ es entonces una $k$-álgebra formalmente lisa para las topologías discretas y $\sh{O}_{Y,y}/\mathfrak m_y=k$, resulta de (0, 20.5.14) que el homomorfismo canónico $\mathfrak m_y/\mathfrak m_y^2\to(\Omega^1_{Y/S})_y\otimes_{\sh{O}_{Y,y}}k(y)$ es biyectivo; de igual modo, el homomorfismo canónico $\mathfrak m_x/\mathfrak m_x^2\to(\Omega^1_{X/S})_x\otimes_{\sh{O}_{X,x}}k(x)$ es biyectivo. La hipótesis c'), equivalente a c), significa, por tanto, aquí que el homomorfismo canónico
\[
(\mathfrak m_y/\mathfrak m_y^2)\otimes_{k(y)}k(x)\to\mathfrak m_x/\mathfrak m_x^2
\]
es inyectivo. Como el anillo $\sh{O}_{X,x}$ es regular, la conclusión resulta de (0, 17.3.3). C.Q.D.

\begin{corollary}[17.11.2]
\label{IV.17.11.2}
Bajo las hipótesis generales de (17.11.1), las condiciones siguientes son equivalentes:
\begin{enumerate}
  \item[a)] $f$ es liso en el punto $y$ y $g$ es étale en el punto $x$.
  \item[b)] $h$ es liso en el punto $x$ y $g$ es étale en el punto $x$.
  \item[c)] $h$ es liso en el punto $x$, y el homomorfismo canónico
  \[
  (g^*(\Omega^1_{Y/S}))_x\to(\Omega^1_{X/S})_x
  \]
  es biyectivo.
  \item[c')] $h$ es liso en el punto $x$, y el homomorfismo canónico
  \[
  ((\Omega^1_{Y/S})_y\otimes_{\sh{O}_{Y,y}}k(y))\otimes_{k(y)}k(x)\to(\Omega^1_{X/S})_x\otimes_{\sh{O}_{X,x}}k(x)
  \]
  es biyectivo.
\end{enumerate}
Supongamos, además, que el homomorfismo $k(y)\to k(x)$ sea biyectivo. Entonces las condiciones anteriores son también equivalentes a la siguiente:
\begin{enumerate}
  \item[d)] $h$ es liso en el punto $x$, y la aplicación canónica $\mathrm{T}_{X/S}(x)\to\mathrm{T}_{Y/S}(y)$ (16.5.12) es biyectiva.
\end{enumerate}
\end{corollary}

Cada una de las condiciones a), b), c) de (17.11.2) equivale a la conjunción de la condición correspondiente de (17.11.1) y del hecho de que $g$ es no ramificado en el punto $x$, teniendo en cuenta (17.2.2) por lo que respecta a la condición c); de donde resulta la equivalencia de a), b) y c). La equivalencia de c) y c') resulta de la equivalencia de las condiciones correspondientes de (17.11.1) y del lema de Nakayama; la equivalencia de c') y d) cuando $k(x)=k(y)$ es inmediata por transposición (16.5.12).

\begin{corollary}[17.11.3]
\label{IV.17.11.3}
Sean $h:X\to S$ un morfismo liso, $s_i$ ($1\leq i\leq r$) secciones de $\sh{O}_X$ sobre $X$ (que son también secciones de $h_*(\sh{O}_X)$ sobre $S$), y $g:X\to S[T_1,\ldots,T_r]=\mathbf{V}_S^r$ el $S$-morfismo correspondiente al homomorfismo de $\sh{O}_S$-módulos $\sh{O}_S^r\to h_*(\sh{O}_X)$ definido por dichas secciones (II, 1.2.7). Para que $g$ sea liso (resp. étale) en un punto $x\in X$, es necesario y suficiente que los $(d_{X/S}(s_i))_x$ formen parte de una base (resp. constituyan una base) del $\sh{O}_{X,x}$-módulo $(\Omega^1_{X/S})_x$.
\end{corollary}

Basta aplicar (17.11.1) (resp. (17.11.2)) tomando como $f$ el morfismo estructural $S[T_1,\ldots,T_r]\to S$.

\begin{corollary}[17.11.4]
\label{IV.17.11.4}
Para que un morfismo $h:X\to S$ sea liso en un punto $x\in X$, es necesario y suficiente que existan una vecindad abierta $U$ de $x$, un entero $r$ y un $S$-morfismo étale $U\to S[T_1,\ldots,T_r]$.
\end{corollary}

\oldpage[IV-4]{85}

La condición es evidentemente suficiente, puesto que el morfismo estructural $S[T_1,\ldots,T_r]\to S$ es liso (17.3.8). Para mostrar que es necesaria, observemos que, puesto que $h$ es liso en el punto $x$, existe una vecindad abierta $U$ de $x$ tal que $\Omega^1_{X/S}|U$ sea localmente libre (17.2.3); basta entonces utilizar (16.10.6) para obtener (restringiendo $U$ si es preciso) secciones $s_i$ de $\sh{O}_X$ sobre $U$ tales que los $d_{X/S}(s_i)$ formen una base de $\Omega^1_{X/S}|U$; se concluye aplicando (17.11.3).

\begin{proposition}[17.11.5]
\label{IV.17.11.5}
Sean $f:Y\to S$, $h:X\to S$ dos morfismos lisos. Para que un $S$-morfismo $g:X\to Y$ sea una inmersión abierta, es necesario y suficiente que $g$ sea un monomorfismo de preesquemas y que, para todo $x\in X$, si se pone $y=g(x)$, se tenga $\dim_y(f)=\dim_x(h)$ (para una generalización de esta proposición, véase (18.10.5)).
\end{proposition}

La condición es evidentemente necesaria; probemos que es suficiente.

En virtud de (17.9.5), nos reducimos de inmediato al caso en que $S=\Spec(k)$ es el espectro de un cuerpo, y en virtud de (2.7.1, (x)), podemos suponer que $k$ es algebraicamente cerrado. Teniendo en cuenta (17.9.1), basta entonces probar que $g$ es étale, y como el conjunto de los puntos donde $g$ es étale es abierto, basta mostrar que $g$ es étale en los puntos \emph{cerrados} (o también, racionales sobre $k$) de $X$ (10.4.8). Sea, pues, $x$ un punto tal, y pongamos $y=g(x)$, que también es racional sobre $k$; por hipótesis, el anillo $A=\sh{O}_{Y,y}$ es regular y tiene cuerpo residual $k$; sean $d=\dim(A)$ y $(t_i)_{1\leq i\leq d}$ un sistema regular de parámetros para $A$. Pongamos $B=\sh{O}_{X,x}$, $C=B/\mathfrak m_yB$. Como $g$ es un monomorfismo, lo mismo ocurre con el morfismo $\Spec(C)\to\Spec(k(y))=\Spec(k)$ deducido de $g$ por cambio de base (I, 3.3.12); pero esto significa que el homomorfismo correspondiente $u:k\to C$ es \emph{sobreyectivo} (y por tanto biyectivo), pues $u$ admite un inverso por la izquierda $v:C\to k$, y $u\circ v$ y la identidad de $C$, compuestos con $u$, dan el mismo morfismo $u:k\to C$. Como, por hipótesis, $B$ es un anillo regular de dimensión $d$, las imágenes de los $t_i$ en $B$ forman un sistema regular de parámetros para $B$ (0, 17.1.7); la condición (17.6.3, $e''$) se verifica, por tanto, para $g$ en el punto $x$ (0, 17.1.1 y Bourbaki, \emph{Alg. comm.}, cap. III, § 2, n\textsuperscript{o} 8, cor. 3 del teorema 1), lo que termina la demostración.

\subsection{Subpreesquemas lisos de un preesquema liso. Morfismos lisos y morfismos diferencialmente lisos}
\label{subsection:IV.17.12}

\begin{theorem}[17.12.1]
\label{IV.17.12.1}
Sean $f:X\to S$, $h:Y\to S$ dos morfismos localmente de presentación finita, $j:Y\to X$ una inmersión, $y$ un punto de $Y$, $x=j(y)$. Las condiciones siguientes son equivalentes:
\begin{enumerate}
  \item[a)] $h$ es liso en el punto $y$ y $f$ es liso en el punto $x$.
  \item[b)] $f$ es liso en el punto $x$, y el homomorfismo canónico (16.4.21)
  \[
  (\sh{N}_{Y/X})_y\to(j^*(\Omega^1_{X/S}))_y
  \]
  admite un inverso por la izquierda.
  \item[c)] $h$ es liso en el punto $y$ y existe una vecindad abierta $U$ de $y$ en $Y$ tal que $j|U:U\to X$ sea una inmersión regular (16.9.2).

\oldpage[IV-4]{86}

  \item[c')] $h$ es liso en el punto $y$ y existe una vecindad abierta $U$ de $y$ en $Y$ tal que $j|U:U\to X$ sea una inmersión cuasi-regular (es decir (16.9.8), existe una vecindad $U$ de $y$ en $Y$ en la que $\sh{N}_{Y/X}$ es un $\sh{O}_Y$-módulo localmente libre y el homomorfismo canónico
  \[
  \mathbf{S}_{\sh{O}_Y}(\sh{N}_{Y/X})\to\operatorname{gr}^{\bullet}(j)
  \]
  es biyectivo).
\end{enumerate}
\end{theorem}

Para demostrar la equivalencia de a) y b), podemos reducirnos al caso en que $S=\Spec(A)$ y $X=\Spec(B)$ son afines, con $Y=\Spec(B/\mathfrak J)$, donde $\mathfrak J$ es un ideal de tipo finito de $B$, y $B$ es una $A$-álgebra lisa (17.3.2, (ii)). Basta entonces aplicar el criterio jacobiano (0, 22.6.1), así como (0, 19.1.14), teniendo en cuenta que $\Omega^1_{X/S}$ es un $\sh{O}_X$-módulo localmente libre en una vecindad de $x$ y $\sh{N}_{Y/X}$ un $\sh{O}_Y$-módulo de tipo finito (16.1.6).

En segundo lugar, probemos que a) implica c'). Podemos reducirnos nuevamente al caso en que $S$, $X$, $Y$ son afines, $B$ y $B/\mathfrak J$ son $A$-álgebras lisas, y resulta entonces de (17.7.9) y (8.10.5, (iv)) que existen un subanillo noetheriano $A'$ de $A$, una $A'$-álgebra de tipo finito $B'$ y un ideal $\mathfrak J'$ de $B'$ tales que $B=B'\otimes_{A'}A$, $\mathfrak J=\mathfrak J'B$ y que $B'$ y $B'/\mathfrak J'$ son $A'$-álgebras lisas. Observemos que, según (0, 19.3.8), $B'$ sigue siendo una $A'$-álgebra formalmente lisa cuando se dota a $A'$ de la topología discreta y a $B'$ de la topología $\mathfrak J'$-preádica. Por consiguiente (0, 19.5.4), $\mathfrak J'/\mathfrak J'^2$ es un $(B'/\mathfrak J')$-módulo proyectivo y el homomorfismo canónico
\[
\mathbf{S}_{B'/\mathfrak J'}(\mathfrak J'/\mathfrak J'^2)\to\operatorname{gr}_{\mathfrak J'}^{\bullet}(B')
\]
es biyectivo. En otros términos (16.9.8), la inmersión $\Spec(B'/\mathfrak J')\to\Spec(B')$ es cuasi-regular, por tanto regular puesto que $B'$ es noetheriano (16.9.10). Pero, como por hipótesis $B'/\mathfrak J'$ es un $A'$-módulo plano (17.5.1) y $B'$ una $A'$-álgebra de presentación finita, podemos aplicar (11.3.8) reemplazando $\sh{F}$ por $\widetilde{\mathfrak J'}$, y vemos así (por cambio de base) que $j$ es una inmersión regular, y \emph{a fortiori} cuasi-regular. El hecho de que $B$ sea una $A$-álgebra de presentación finita y un $A$-módulo plano muestra además, en virtud de (11.3.8), que las condiciones c') y c) son equivalentes, y son estables por cambio de base.

Demostremos finalmente que c) implica a). Del hecho de que, en (11.3.8), la condición b) implica c), se ve ya que, si c) se cumple, $f$ es \emph{plano} en el punto $x$; en virtud de (17.5.1), todo se reduce, pues, a ver que, bajo la hipótesis c), $f^{-1}(s)$ es liso sobre $k(s)$ en el punto $x$, poniendo $s=h(x)$; como se ha observado que la condición c) es estable por cambio de base, vemos que nos reducimos al caso en que $S=\Spec(k)$ es el espectro de un cuerpo, y teniendo en cuenta (17.7.1, (ii)) podemos suponer que $k$ es \emph{algebraicamente cerrado}. Basta entonces probar que $\sh{O}_{X,x}$ es un anillo \emph{regular} (17.5.1). Ahora bien, por hipótesis se tiene $\sh{O}_{Y,y}=\sh{O}_{X,x}/\mathfrak J_x$, donde $\mathfrak J_x$ es un ideal engendrado por una sucesión $\sh{O}_{X,x}$-regular $(t_i)$, y $\sh{O}_{Y,y}$ es un anillo regular; como los $t_i$ forman parte de un sistema de parámetros para $\sh{O}_{X,x}$ (0, 16.4.1), la conclusión resulta de (0, 17.1.7).

\begin{corollary}[17.12.2]
\label{IV.17.12.2}
Con las notaciones de (17.12.1), supongamos que $f$ es liso en el punto $x$, y sea $Y$ un subpreesquema cerrado de $X$ definido por un ideal $\sh{J}$ de $\sh{O}_X$, y liso sobre $S$ en el punto $x$. Sean $(g_i)_{1\leq i\leq r}$ secciones de $\sh{J}$ sobre $X$. Las condiciones siguientes son equivalentes:
\begin{enumerate}
  \item[a)] Los $(g_i)_x$ forman un sistema de generadores del $\sh{O}_{X,x}$-módulo $\sh{J}_x$ cuyo número de elementos es el menor posible.
  \item[b)] Si $g_i'$ es la imagen canónica de $g_i$ en $\Gamma(X,\sh{J}/\sh{J}^2)$, los $(g_i')_x$ forman una base del $(\sh{O}_{X,x}/\sh{J}_x)$-módulo $\sh{J}_x/\sh{J}_x^2$.

\oldpage[IV-4]{87}

  \item[c)] Las imágenes de los $d_{X/S}g_i$ en $(\Omega^1_{X/S})_x\otimes_{\sh{O}_{X,x}}k(x)$ son elementos linealmente independientes sobre $k(x)$, y los $(g_i)_x$ engendran $\sh{J}_x$.
  \item[d)] Existe una vecindad abierta $U$ de $x$ en $X$ y secciones $g_j$ ($r+1\leq j\leq n$) de $\sh{O}_X$ sobre $U$ tales que el $S$-morfismo $u:U\to S[T_1,\ldots,T_n]$ correspondiente al homomorfismo $\sh{O}_S^n\to f_*(\sh{O}_U)$ definido por las secciones $g_i|U$ ($1\leq i\leq r$) y $g_j$ ($r+1\leq j\leq n$) (II, 1.2.7) sea un morfismo étale, y que la imagen inversa del subpreesquema $Y'=S[T_{r+1},\ldots,T_n]$ por este morfismo sea el preesquema inducido por $j(Y)$ sobre el abierto $U\cap j(Y)$.
\end{enumerate}
\end{corollary}

Como $(g_i')_x$ es la imagen canónica de $(g_i)_x$, la equivalencia de a) y b) resulta del lema de Nakayama, pues $\sh{J}_x$ es de tipo finito y $\sh{J}_x/\sh{J}_x^2$ un $(\sh{O}_{X,x}/\sh{J}_x)$-módulo libre (17.10.4) (Bourbaki, \emph{Alg. comm.}, cap. II, § 3, n\textsuperscript{o} 2, prop. 5). En virtud de (17.12.1, b)), $\sh{J}_x/\sh{J}_x^2$ se identifica canónicamente con un factor directo del $(\sh{O}_{X,x}/\sh{J}_x)$-módulo libre de rango $n$, $(\Omega^1_{X/S})_x\otimes_{\sh{O}_{X,x}}(\sh{O}_{X,x}/\sh{J}_x)$, y la equivalencia de b) y c) resulta de Bourbaki, \emph{loc. cit.} Además, si a) se cumple, $(g_i')_x$ se identifica así con $(d_{X/S}g_i)_x\otimes1$ para $1\leq i\leq r$; como el $\sh{O}_{X,x}$-módulo $(\Omega^1_{X/S})_x$ es libre de rango $n$, resulta nuevamente de Bourbaki, \emph{loc. cit.}, que, reemplazando si es preciso $X$ por una vecindad abierta de $x$, se puede suponer que existen $n-r$ secciones $g_i$ ($r+1\leq i\leq n$) de $\sh{O}_X$ sobre $X$, tales que los $(d_{X/S}g_i)_x$ para $1\leq i\leq n$ forman una base de $(\Omega^1_{X/S})_x$; el hecho de que el morfismo correspondiente $u$ sea étale en el punto $x$ resulta entonces de (17.11.3); por la misma razón, el morfismo $u^{-1}(Y')\to Y'$, restricción de $u$, es étale en el punto $x$; reemplazando $X$ por una vecindad de $x$, se puede suponer que ambos morfismos son étale. Además, es inmediato que $Y$ (identificado, para simplificar, con el subpreesquema cerrado $j(Y)$ de $X$) es un subpreesquema cerrado de $u^{-1}(Y')$, y, en virtud de la elección de los $g_i$ para $i>r$ y de (17.2.5) y (17.11.2), se puede suponer también que la restricción a $Y$ de $u$ es étale. Se deduce que, para todo $y'\in Y'$, la inmersión $Y\cap u^{-1}(y')\to u^{-1}(y')$ es abierta, de donde se concluye, con ayuda de (17.9.6), que $Y\to u^{-1}(Y')$ es una inmersión abierta; esto prueba que a) implica d). El recíproco se deduce inmediatamente de (17.11.3).

\begin{corollary}[17.12.3]
\label{IV.17.12.3}
Sean $f:X\to S$ un morfismo localmente de presentación finita, $u:S\to X$ una $S$-sección de $X$. Para que $f$ sea liso en un punto $x\in X$, es necesario y suficiente que $u$ sea una inmersión cuasi-regular en una vecindad de $s=f(x)$.
\end{corollary}

La conclusión resulta de (17.12.1), puesto que $f\circ u=1_S$ es liso.

\begin{proposition}[17.12.4]
\label{IV.17.12.4}
Todo morfismo liso $f:X\to S$ es diferencialmente liso (es decir (16.10.5), $\Delta_f:X\to X\times_SX$ es una inmersión cuasi-regular). \emph{En particular, los $\sh{P}_f^n=\sh{P}_{X/S}^n$ y los $\operatorname{gr}_n(\sh{P}_{X/S})$ son $\sh{O}_X$-módulos localmente libres de tipo finito.}
\end{proposition}

Basta observar que el morfismo estructural $p_1:X\times_SX\to X$ es liso (17.3.3, (iii)) y que $\Delta_f$ es una $X$-sección de $p_1$; la conclusión resulta de (17.12.3).

\begin{proposition}[17.12.5]
\label{IV.17.12.5}
Sea $f:X\to Y$ un morfismo localmente de presentación finita. Las condiciones siguientes son equivalentes:
\begin{enumerate}
  \item[a)] $f$ es diferencialmente liso.
  \item[b)] Para todo morfismo $g:Y'\to Y$, si se pone $X'=X\times_Y Y'$ y $f'=f_{(Y')}:X'\to Y'$, entonces, para toda $Y'$-sección $s'$ de $X'$, $f'$ es liso en todos los puntos de $s'(Y')$.
  \item[c)] La segunda proyección $p_2:X\times_YX\to X$ es lisa en todos los puntos de la diagonal $\Delta_f(X)$.
\end{enumerate}
\end{proposition}

\oldpage[IV-4]{88}

La condición c) es un caso particular de b): basta, en efecto, tomar $Y'=X$ y $g=f$ en b), pues entonces $f'$ no es sino la segunda proyección $p_2$, y $\Delta_f$ es una $X$-sección de $X\times_YX$. En segundo lugar, c) implica a), pues $p_2$ es un morfismo localmente de presentación finita, y si $p_2$ es liso en los puntos de $\Delta_f(X)$, $\Delta_f$ es una inmersión cuasi-regular (17.12.3), por tanto $f$ es diferencialmente liso. Finalmente, demostremos que a) implica b); si $\Delta_f$ es una inmersión cuasi-regular, $p_2$ es liso en todos los puntos de $\Delta_f(X)$ (17.12.3).

Ahora bien, si $g':X'\to X$ es la proyección canónica, y $v=(g',g')_Y:X'\to X\times_YX$, se comprueba inmediatamente que el diagrama
\[
\xymatrix{
  X\times_Y X \ar[d]_{p_2} & X' \ar[l]_-v \ar[d]^{f'}\\
  X & Y' \ar[l]^{h=g'\circ s'}
}
\]
es conmutativo e identifica $X'$ con el producto $(X\times_YX)\times_XY'$ (I, 3.3.9); teniendo en cuenta que el diagrama (17.4.1.1) identifica $Y'$ con el producto de los $(X\times_YX)$-preesquemas $X$ y $X'$, se concluye, del hecho de que $p_2$ es liso en todo punto de $\Delta_f(X)$, que $f'$ es liso en todo punto de $s'(Y')$ (17.3.3, (iii)).

\begin{corollary}[17.12.6]
\label{IV.17.12.6}
Con las notaciones de (8.8.1), se supone $X_\alpha$ cuasicompacto, y $f_\alpha:X_\alpha\to S_\alpha$ localmente de presentación finita. Para que $f:X\to S$ sea diferencialmente liso, es necesario y suficiente que exista $\lambda\geq\alpha$ tal que $f_\lambda:X_\lambda\to S_\lambda$ sea diferencialmente liso (en cuyo caso $f_\mu:X_\mu\to S_\mu$ es diferencialmente liso para $\mu\geq\lambda$).
\end{corollary}

La suficiencia de la condición y la última afirmación resultan de (16.10.4). Para probar que la condición es necesaria, observemos que $\Delta_f:X\to X\times_SX$ y $p_2:X\times_SX\to X$ se deducen por cambio de base de $\Delta_{f_\lambda}:X_\lambda\to X_\lambda\times_{S_\lambda}X_\lambda$ y $p_{2,\lambda}:X_\lambda\times_{S_\lambda}X_\lambda\to X_\lambda$. En virtud de (17.12.5), para todo $z\in\Delta_f(X)$, existe una vecindad abierta afín $U(z)$ de $z$ en $X\times_SX$ tal que $p_2$ sea liso en $U(z)$; en virtud de (8.2.11) y (17.7.8), existen un índice $\lambda(z)$ y una vecindad $U_{\lambda(z)}$ de la proyección de $z$ en $X_{\lambda(z)}\times_{S_{\lambda(z)}}X_{\lambda(z)}$ tales que $U(z)$ sea la imagen inversa de $U_{\lambda(z)}$, y que $p_{2,\lambda(z)}$ sea liso en $U_{\lambda(z)}$. Como $X$ es cuasicompacto, se puede recubrir $\Delta_f(X)$ por un número finito de vecindades $U(z_i)$; en virtud de (8.3.4), existe un $\lambda$ superior a todos los $\lambda(z_i)$ tal que las imágenes inversas de los $U_{\lambda(z_i)}$ en $X_\lambda\times_{S_\lambda}X_\lambda$ forman un recubrimiento de $\Delta_{f_\lambda}(X_\lambda)$; resulta entonces de (17.12.5) que este índice $\lambda$ cumple lo requerido.

\begin{examples}[17.12.7]
\label{IV.17.12.7}
Sean $S$ un preesquema, $G$ un $S$-\emph{preesquema en grupos}, localmente de presentación finita sobre $S$. Entonces, para que $G$ sea diferencialmente liso, es necesario y suficiente que $G$ sea liso sobre $S$ en los puntos de la «sección unidad» $e$ de $G$. La condición es, en efecto, necesaria por (17.12.5). Recíprocamente, supongamos que se cumple; para todo morfismo $S'\to S$, $G'=G\times_SS'$ es entonces un $S'$-preesquema en grupos localmente de presentación finita sobre $S'$; podemos, por tanto, en virtud de (17.12.5), reducirnos a probar que, para toda $S$-sección $s$ de $G$, $f$ es liso en los puntos de $s(S)$. Ahora bien, si $m:G\times_SG\to G$ es el morfismo que define la estructura de preesquema en grupos de $G$, $m\circ(s\times1_G):S\times_SG\to G\times_SG\to G$ es un

\oldpage[IV-4]{89}

$S$-isomorfismo del preesquema $G$, la «traslación por la izquierda» $\tau_s$, que transforma los puntos de la sección unidad de $G$ en los puntos de $s(S)$; se concluye que $G$ es liso sobre $S$ en los puntos de $s(S)$.
\end{examples}

\begin{remark}[17.12.8]
\label{IV.17.12.8}
Un $S$-preesquema en grupos $G$, de tipo finito (e incluso finito) sobre un preesquema localmente noetheriano $S$, puede ser diferencialmente liso sobre $S$ sin ser liso (ni siquiera \emph{plano}) sobre $S$. Tomemos, por ejemplo, como $S$ el espectro del álgebra de números duales $D=k[T]/T^2k[T]$ sobre un cuerpo $k$, y como $G$ el espectro de la $D$-álgebra producto directo $E=D\oplus k$. Se define en $G$ una estructura de $S$-esquema en grupos definiendo una «aplicación diagonal» $\delta$, homomorfismo de $E$ en $E\otimes_DE$: si $e'$, $e''$ son los idempotentes, imágenes canónicas de la unidad $1$ de $D$ en los factores $D$ y $k$ de $E$, se comprueba fácilmente que se obtiene una aplicación diagonal de este tipo tomando
\[
\delta(e')=e'\otimes e'+e''\otimes e'',\qquad
\delta(e'')=e'\otimes e''+e''\otimes e'.
\]
La sección unidad del esquema en grupos $G$ corresponde al homomorfismo $E\to D$ que es la identidad en $D$ y $0$ en $k$; es un isomorfismo de $S$ sobre una componente conexa de $G$, y \emph{a fortiori} $G$ es diferencialmente liso sobre $S$ (17.12.6), pero está claro que $G$ no es $S$-plano.
\end{remark}

\subsection{Morfismos transversales.}
\label{subsection:IV.17.13}
\label{IV.17.13}

\begin{env}[17.13.1]
\label{IV.17.13.1}
Sean $S$ un preesquema, $X$, $Y$, $X'$ tres $S$-preesquemas, $i:Y\to X$ una $S$-inmersión, $f:X'\to X$ un $S$-morfismo. Pongamos $Y'=Y\times_XX'$, y sean $g:Y'\to Y$, $j:Y'\to X'$ las proyecciones canónicas, de modo que se tiene el diagrama conmutativo
\[
\label{IV.17.13.1.1}
\tag{17.13.1.1}
\xymatrix{
  Y \ar[d]_i & Y' \ar[l]_g \ar[d]^j\\
  X & X' \ar[l]^f
}
\]
y que $j$ es una $S$-inmersión. Se tiene entonces un diagrama conmutativo de $\sh{O}_{Y'}$-módulos cuasicoherentes
\[
\label{IV.17.13.1.2}
\tag{17.13.1.2}
\begin{tikzcd}
  g^*(\sh{N}_{Y/X}) \ar{r} \ar{d}{\operatorname{gr}_1(g)} & f^*(\Omega^1_{X/S})\otimes_{\sh{O}_{X'}}\sh{O}_{Y'} \ar{r} \ar{d} & g^*(\Omega^1_{Y/S}) \ar{r} \ar{d} & 0\\
  \sh{N}_{Y'/X'} \ar{r} & \Omega^1_{X'/S}\otimes_{\sh{O}_{X'}}\sh{O}_{Y'} \ar{r} & \Omega^1_{Y'/S} \ar{r} & 0
\end{tikzcd}
\]
donde la fila inferior es la sucesión exacta (16.4.21) aplicada a $j$, la fila superior proviene de la misma sucesión exacta para $i$, mediante la aplicación del funtor exacto por la derecha $g^*$ (que, por tanto, la deja exacta); $\operatorname{gr}_1(g)$ se define en (16.2.1), y la conmutatividad de los dos cuadrados resulta de (0, 20.5.7.3) y (0, 20.5.11.3).

\oldpage[IV-4]{90}

Además, se ha visto en (16.2.2, (iii)) que $\operatorname{gr}_1(g)$ es aquí sobreyectivo, por lo que se deduce de (17.13.1.2) la sucesión exacta
\[
\label{IV.17.13.1.3}
\tag{17.13.1.3}
g^*(\sh{N}_{Y/X})\xrightarrow{\alpha}\Omega^1_{X'/S}\otimes_{\sh{O}_{X'}}\sh{O}_{Y'}\to\Omega^1_{Y'/S}\to0.
\]
\end{env}

\begin{proposition}[17.13.2]
\label{IV.17.13.2}
Con las notaciones de (17.13.1), sea $x'$ un punto de $Y'$, $x=g(x')$ su imagen en $Y$; supongamos que $X$ e $Y$ son lisos sobre $S$ en el punto $x$, y que $X'$ es liso sobre $S$ en el punto $x'$. Sean $m$, $m-c$ las dimensiones relativas de $X$ e $Y$ sobre $S$ en el punto $x$ (17.10.1), y $n$ la de $X'$ sobre $S$ en el punto $x'$. Las condiciones siguientes son equivalentes:
\begin{enumerate}
  \item[a)] $Y'$ es liso sobre $S$ en el punto $x'$ y de dimensión relativa $n-c$.
  \item[b)] El homomorfismo $\alpha\otimes1:g^*(\sh{N}_{Y/X})\otimes_{\sh{O}_{Y'}}k(x')\to\Omega^1_{X'/S}\otimes_{\sh{O}_{X'}}k(x')$ deducido del homomorfismo $\alpha$ de (17.13.1.3) es inyectivo.
\end{enumerate}
Sea $s$ la imagen canónica de $x$ en $S$, y supongamos que $x'$ es racional sobre $k(s)$. Entonces las condiciones a) y b) son también equivalentes a la siguiente:
\begin{enumerate}
  \item[b')] El homomorfismo transpuesto de $\alpha\otimes1$
  \[
  \mathrm{T}_{X'/S}(x')\to(\sh{N}_{Y/X}\otimes_{\sh{O}_Y}k(x'))^\vee=\mathrm{T}_{X/S}(x)/\mathrm{T}_{Y/S}(x)
  \]
  (cf. (16.5.12)) es sobreyectivo.
\end{enumerate}
Además, cuando las condiciones a) y b) se cumplen en $x'$, se cumplen en una vecindad de $x'$ en $Y'$; sustituyendo, si es preciso, $X'$ por una vecindad de $x'$, el homomorfismo
\[
\operatorname{gr}_1(g):g^*(\sh{N}_{Y/X})\to\sh{N}_{Y'/X'}
\]
es biyectivo, y la sucesión
\[
\label{IV.17.13.2.1}
\tag{17.13.2.1}
0\to g^*(\sh{N}_{Y/X})\xrightarrow{\alpha}\Omega^1_{X'/S}\otimes_{\sh{O}_{X'}}\sh{O}_{Y'}\to\Omega^1_{Y'/S}\to0
\]
obtenida añadiendo un $0$ a (17.13.1.3), es exacta.
\end{proposition}

El hecho de que, si las condiciones equivalentes a) y b) se cumplen en el punto $x'$, se cumplen también en una vecindad de $x'$ en $Y'$, resulta de que el conjunto de los puntos donde un morfismo es liso es abierto (17.3.7) y de (17.10.2).

En virtud de (17.12.1), aplicado a $X'$ e $Y'$, decir que $Y'$ es liso sobre $S$ en el punto $x'$ equivale a decir que el homomorfismo
\[
\delta\otimes1:\sh{N}_{Y'/X'}\otimes_{\sh{O}_{Y'}}k(x')\to\Omega^1_{X'/S}\otimes_{\sh{O}_{X'}}k(x')
\]
es inyectivo, teniendo en cuenta (0, 19.1.12) y que $\Omega^1_{X'/S}$ es un $\sh{O}_{X'}$-módulo localmente libre en el punto $x'$ (17.2.3); como entonces $\Omega^1_{Y'/S}$ es también un $\sh{O}_{Y'}$-módulo localmente libre en una vecindad de $x'$ y la sucesión
\[
\label{IV.17.13.2.2}
\tag{17.13.2.2}
0\to\sh{N}_{Y'/X'}\to\Omega^1_{X'/S}\otimes_{\sh{O}_{X'}}\sh{O}_{Y'}\to\Omega^1_{Y'/S}\to0
\]
es exacta (17.2.5), decir que $Y'$ es de dimensión relativa $n-c$ sobre $S$ en el punto $x'$ significa, por (17.10.2), que $\sh{N}_{Y'/X'}$ (que es localmente libre en una vecindad de $x'$) es de rango $c$ en el punto $x'$. Pero, en virtud de las hipótesis sobre $X$ e $Y$ en el punto $x$, y por el mismo razonamiento, $\sh{N}_{Y/X}$ es localmente libre y de rango $c$ en el punto $x$, por lo que $g^*(\sh{N}_{Y/X})$ es también localmente libre y de rango $c$ en el punto $x'$. Como el homomorfismo $\operatorname{gr}_1(g):g^*(\sh{N}_{Y/X})\to\sh{N}_{Y'/X'}$ es sobreyectivo, las condiciones anteriores equivalen a decir que este homomorfismo es

\oldpage[IV-4]{91}

biyectivo en el punto $x'$ (Bourbaki, \emph{Alg. comm.}, cap. II, § 3, n\textsuperscript{o} 2, cor. de la prop. 6), y, por tanto, también en una vecindad de $x'$ (0$_{\mathrm I}$, 5.2.7); esto implica evidentemente b), así como la última afirmación del enunciado, en virtud de la exactitud de (17.13.2.2). Recíprocamente, como $\alpha\otimes1$ se factoriza en
\[
g^*(\sh{N}_{Y/X})\otimes k(x')\xrightarrow{\operatorname{gr}_1(g)\otimes1}\sh{N}_{Y'/X'}\otimes k(x')\xrightarrow{\delta\otimes1}\Omega^1_{X'/S}\otimes k(x')
\]
y $\operatorname{gr}_1(g)$ es sobreyectivo, decir que $\alpha\otimes1$ es inyectivo implica que $\delta\otimes1$ lo es y que $\operatorname{gr}_1(g)\otimes1$ es biyectivo. Se concluye de (17.12.1) que $Y'$ es liso sobre $S$ en el punto $x'$, y que $\operatorname{gr}_1(g)$ es biyectivo en una vecindad de $x'$ (Bourbaki, \emph{loc. cit.}). Además, como la sucesión (17.13.2.2) es entonces exacta, y $\sh{N}_{Y'/X'}$, isomorfo a $g^*(\sh{N}_{Y/X})$, es de rango $c$ en el punto $x'$, $\Omega^1_{Y'/S}$ es de rango $n-c$ en este punto, lo que termina de demostrar la equivalencia de a) y b), en virtud de (17.10.2).

Queda por demostrar la equivalencia de b) y b') cuando $x'$ es racional sobre $k(s)$ (lo que implica que también lo es $x$); entonces $g^*(\sh{N}_{Y/X})\otimes_{\sh{O}_{Y'}}k(x')$ se identifica con $\sh{N}_{Y/X}\otimes_{\sh{O}_Y}k(x')$, y puesto que la sucesión
\[
0\to\sh{N}_{Y/X}\to\Omega^1_{X/S}\otimes_{\sh{O}_X}\sh{O}_Y\to\Omega^1_{Y/S}\to0
\]
es exacta en el punto $x$ y está formada por $\sh{O}_Y$-módulos localmente libres, el dual de $\sh{N}_{Y/X}\otimes_{\sh{O}_Y}k(x)$ se identifica con el espacio cociente $\mathrm{T}_{X/S}(x)/\mathrm{T}_{Y/S}(x)$ (16.5.12); de ahí la equivalencia de b) y b').

\begin{definition}[17.13.3]
\label{IV.17.13.3}
Con las notaciones de (17.13.1), se dice que el morfismo $f$ es \emph{transversal} a $Y$ en el punto $x'$, relativamente a $S$, si $X$ e $Y$ son lisos sobre $S$ en el punto $x$, $X'$ es liso sobre $S$ en el punto $x'$, y si se satisfacen las condiciones equivalentes a), b) de (17.13.2).
\end{definition}

Cuando no hay riesgo de confusión, se omite la mención del preesquema $S$ y se dice simplemente que $f$ es transversal a $Y$ en el punto $x'$.

\begin{remarks}[17.13.4]
\label{IV.17.13.4}
(i) Supongamos que $X$, $Y$ y $X'$ son \emph{planos y localmente de presentación finita} sobre $S$. Para todo $s\in S$, denotemos por $X_s$, $Y_s$, $X'_s$, $Y'_s$ las fibras de $X$, $Y$, $X'$, $Y'$ en el punto $s$, y por $f_s:X'_s\to X_s$ el morfismo deducido de $f$ por cambio de base. Resulta entonces de (17.5.1)\footnote{El texto impreso cita (17.5.9), pero no existe tal apartado; el criterio de lisura por fibras utilizado aquí es (17.5.1).} que, para que $f$ sea un morfismo transversal a $Y$ en el punto $x'$, es necesario y suficiente que, si $s$ es la imagen de $x$ en $S$, $f_s$ sea transversal a $Y_s$ en el punto $x'$, relativamente a $k(s)$.

(ii) Se ha visto en (17.13.2) que el conjunto de los puntos $x'\in Y'$ donde $f$ es transversal a $Y$ es \emph{abierto} en $Y'$; si, además, $Y'$ es \emph{propio} sobre $S$, se deduce que el conjunto de los $s\in S$ tales que $f$ sea transversal a $Y$ en todos los puntos de $Y'_s$, relativamente a $S$, es \emph{abierto} en $S$.

Cuando $X$, $Y$ y $X'$ son planos y localmente de presentación finita sobre $S$, resulta de (i) que el conjunto de los $x'\in Y'$ tales que $f_s$ sea transversal a $Y_s$ en el punto $x'$ (siendo $s$ la imagen de $x'$ en $S$), relativamente a $k(s)$, es \emph{abierto} en $Y'$. Si, además, $Y'$ es propio sobre $S$, el conjunto de los $s\in S$ tales que $f_s$ sea transversal a $Y_s$ en todos los puntos de $Y'_s$ (relativamente a $k(s)$) es \emph{abierto} en $S$.

(iii) Puesto que la propiedad de ser liso sobre $S$, así como la noción de dimensión relativa a $S$ en un punto, son estables por todo cambio de base $S'\to S$ ((17.3.3) y (4.2.7)), lo mismo ocurre con la propiedad de que un morfismo $f:X'\to X$ sea transversal a un subpreesquema de $X$ en un punto, relativamente a $S$.

\oldpage[IV-4]{92}

(iv) La condición b) de (17.13.2) expresa también que el homomorfismo $\alpha:g^*(\sh{N}_{Y/X})\to\Omega^1_{X'/S}\otimes_{\sh{O}_{X'}}\sh{O}_{Y'}$ es \emph{universalmente inyectivo} relativamente a $Y'$ (11.9.18).\footnote{El texto fuente cita (II, 9.18); la caracterización de los homomorfismos universalmente inyectivos utilizada aquí es la observación (11.9.18).}
\end{remarks}

\begin{env}[17.13.5]
\label{IV.17.13.5}
Consideremos ahora un preesquema $S$, tres $S$-preesquemas $X$, $Y$, $Z$, dos $S$-morfismos $f:Y\to X$, $g:Z\to X$; pongamos $T=Y\times_XZ$; se sabe entonces (I, 5.3.5) que se tiene un diagrama conmutativo
\[
\label{IV.17.13.5.1}
\tag{17.13.5.1}
\xymatrix{
  X \ar[d]^{\Delta} & T=Y\times_X Z \ar[l] \ar[d]^{j}\\
  X\times_S X & Y\times_S Z \ar[l]^{u}
}
\]
que hace de $T$ el producto de los $(X\times_SX)$-preesquemas $X$ e $Y\times_SZ$, donde $u=f\times_Sg$. Como $\Delta$ es una $S$-inmersión (I, 5.3.9), se está en la situación del diagrama (17.13.1.1); lo que corresponde a $\Omega^1_{X'/S}$ en (17.13.1) es entonces $(\Omega^1_{Y/S}\otimes_{\sh{O}_Y}\sh{O}_{Y\times_SZ})\oplus(\Omega^1_{Z/S}\otimes_{\sh{O}_Z}\sh{O}_{Y\times_SZ})$, en virtud de (16.4.23). Por otra parte, lo que corresponde al $\sh{O}_Y$-módulo $\sh{N}_{Y/X}$ en (17.13.1) es aquí, por definición, $\Omega^1_{X/S}$ (16.3.1); corresponde, por tanto, a (17.13.1.3) una sucesión exacta
\[
\label{IV.17.13.5.2}
\tag{17.13.5.2}
\Omega^1_{X/S}\otimes_{\sh{O}_X}\sh{O}_T\xrightarrow{\rho}
(\Omega^1_{Y/S}\otimes_{\sh{O}_Y}\sh{O}_T)\oplus
(\Omega^1_{Z/S}\otimes_{\sh{O}_Z}\sh{O}_T)
\xrightarrow{\sigma}\Omega^1_{T/S}\to0
\]
donde quedan por precisar los homomorfismos $\rho$ y $\sigma$. Teniendo en cuenta, en primer lugar, (16.4.23) y (0, 20.5.2), se ve que, si $p:T\to Y$, $q:T\to Z$ son las proyecciones canónicas, se tiene, con las notaciones de (16.4.19),
\[
\label{IV.17.13.5.3}
\tag{17.13.5.3}
\sigma=p_{T/Y/S}+q_{T/Z/S}.
\]
Por otra parte, para calcular $\rho$, utilicemos la conmutatividad del cuadrado izquierdo de (17.13.1.2), lo que, en el caso presente, reduce primero a explicitar el homomorfismo canónico
\[
\rho':\Omega^1_{X/S}\to\Omega^1_{(X\times_SX)/S}\otimes_{\sh{O}_{X\times_SX}}\sh{O}_X
\]
definido en (16.4.21), aplicado a la inmersión $\Delta$. Se puede reducir al caso en que $S=\Spec(A)$, $X=\Spec(B)$ son afines, y entonces $\rho'$ corresponde al homomorfismo $\delta$ de (0, 20.5.11.2), donde hay que sustituir $B$ por $B\otimes_AB$ y $C$ por $B=(B\otimes_AB)/\mathfrak J_{B/A}$. Se ve entonces que $\delta$ hace corresponder a la clase de $x\otimes1-1\otimes x\pmod{\mathfrak J_{B/A}^2}$ (para un $x\in B$) la imagen de
\[
(x\otimes1-1\otimes x)\otimes(1\otimes1)-(1\otimes1)\otimes(x\otimes1-1\otimes x)
\]
en $\Omega^1_{(B\otimes_AB)/A}\otimes_{B\otimes_AB}B$; pero el elemento anterior puede escribirse
\[
\begin{split}
&((x\otimes1)\otimes(1\otimes1)-(1\otimes1)\otimes(x\otimes1))\\
&\quad-((1\otimes x)\otimes(1\otimes1)-(1\otimes1)\otimes(1\otimes x))
\end{split}
\]
y se ve, por tanto, que $\rho'$ es la diferencia de los dos homomorfismos $\pi^{(1)}$ y $\pi^{(2)}$ de $\Omega^1_{X/S}$ en $\Omega^1_{(X\times_SX)/S}\otimes_{\sh{O}_{X\times_SX}}\sh{O}_X$, correspondientes, respectivamente, a la primera y a la segunda proyección de $X\times_SX$ en $X$, por (16.4.3.3). Para obtener $\rho$, hay que considerar primero el homomorfismo
\[
\rho'':\Omega^1_{(X\times_SX)/S}\otimes_{\sh{O}_{X\times_SX}}\sh{O}_{Y\times_SZ}\to\Omega^1_{(Y\times_SZ)/S}
\]
correspondiente, por (16.4.3.3), al morfismo $u$ de (17.13.5.1), y luego, tras tensorizar por $\sh{O}_T$, formar el compuesto $(\rho''\otimes1)\circ(\rho'\otimes1)$; resulta de lo anterior que se tiene, con las notaciones de (16.4.18),
\[
\label{IV.17.13.5.4}
\tag{17.13.5.4}
\rho=(f_{Y/X/S}\otimes1_{\sh{O}_T})\oplus(-g_{Z/X/S}\otimes1_{\sh{O}_T}).
\]

\oldpage[IV-4]{93}

Dicho esto, la aplicación de (17.13.2) a la situación del diagrama (17.13.5.1) (teniendo en cuenta (17.3.3, (iv)), que implica que $X\times_SX$ es liso sobre $S$ en $(x,x)$ si $X$ lo es)\footnote{El texto impreso dice que $X\times_SX$ es liso «en $x$»; el punto correspondiente de $X\times_SX$ es $\Delta(x)=(x,x)$.} da el corolario siguiente.
\end{env}

\begin{corollary}[17.13.6]
\label{IV.17.13.6}
Sean $S$ un preesquema, $X$, $Y$, $Z$ tres $S$-preesquemas, $f:Y\to X$, $g:Z\to X$ dos $S$-morfismos; pongamos $T=Y\times_XZ$, y sean $p:T\to Y$, $q:T\to Z$ las proyecciones canónicas. Sea $t$ un punto de $T$, $y=p(t)$, $z=q(t)$, $x=f(y)=g(z)$. Supongamos que $X$ es liso sobre $S$ en el punto $x$, de dimensión relativa $m$, que $Y$ (resp. $Z$) es liso sobre $S$ en el punto $y$ (resp. $z$), de dimensión relativa $m+a$ (resp. $m+b$), siendo $a$ y $b$ positivos o negativos. Entonces las condiciones siguientes son equivalentes:
\begin{enumerate}
  \item[a)] $T$ es liso sobre $S$ en el punto $t$, de dimensión relativa $m+a+b$.
  \item[b)] El homomorfismo
  \[
  \rho\otimes1:\Omega^1_{X/S}\otimes_{\sh{O}_X}k(t)\to
  (\Omega^1_{Y/S}\otimes_{\sh{O}_Y}k(t))\oplus
  (\Omega^1_{Z/S}\otimes_{\sh{O}_Z}k(t))
  \]
  donde $\rho$ viene dado por (17.13.5.4), es inyectivo.
  \item[c)] El morfismo $u:Y\times_SZ\to X\times_SX$ es transversal a la diagonal de $X\times_SX$ en el punto $t$.\footnote{El texto impreso escribe $T=Y\times_SZ\to X\times_SX$; aquí $T$ ya denota $Y\times_XZ$, mientras que el morfismo transversal es $u:Y\times_SZ\to X\times_SX$ del diagrama (17.13.5.1).}
\end{enumerate}
Si $t$ es racional sobre el cuerpo residual de su imagen $s$ en $S$, estas condiciones son también equivalentes a la siguiente:
\begin{enumerate}
  \item[b')] El homomorfismo
  \[
  \label{IV.17.13.6.1}
  \tag{17.13.6.1}
  \mathrm{T}_y(f)-\mathrm{T}_z(g):\mathrm{T}_{Y/S}(y)\oplus\mathrm{T}_{Z/S}(z)\to\mathrm{T}_{X/S}(x)
  \]
  (cf. (16.5.13.5))\footnote{El texto impreso cita (16.5.12.5); el morfismo tangente empleado está definido en (16.5.13.5).} es sobreyectivo.
\end{enumerate}
Además, cuando las condiciones a) y b) se cumplen en $t$, se cumplen en una vecindad de $t$ en $T$, y, restringiendo $T$ a tal vecindad, la sucesión
\[
\label{IV.17.13.6.2}
\tag{17.13.6.2}
0\to\Omega^1_{X/S}\otimes_{\sh{O}_X}\sh{O}_T\xrightarrow{\rho}
(\Omega^1_{Y/S}\otimes_{\sh{O}_Y}\sh{O}_T)\oplus
(\Omega^1_{Z/S}\otimes_{\sh{O}_Z}\sh{O}_T)
\xrightarrow{\sigma}\Omega^1_{T/S}\to0
\]
(donde $\sigma$ viene dado por (17.13.5.3)) es exacta.

El único punto que queda por demostrar en (17.13.6) concierne a b'), y resulta de que, bajo la hipótesis de que $t$ es racional sobre $k(s)$, el homomorfismo $\rho\otimes1$ es \emph{transpuesto} de $\mathrm{T}_y(f)-\mathrm{T}_z(g)$, en virtud de (17.13.5.4).

Cuando se satisfacen las condiciones equivalentes de (17.13.6), se dice que $f$ y $g$ forman un \emph{par de morfismos transversales en el punto $t$, relativamente} a $S$; también aquí se omite a menudo la mención de $S$.
\end{corollary}

\begin{env}[17.13.7]
\label{IV.17.13.7}
Consideremos, en particular, el caso en que $Y$ y $Z$ son \emph{subpreesquemas} de $X$, siendo $f$ y $g$ las inyecciones canónicas; $T$ es entonces el subpreesquema «intersección»

\oldpage[IV-4]{94}

$\inf(Y,Z)$ de $X$ (I, 4.4.3), y se tiene $t=y=z=x$. En lugar de decir que el par $(f,g)$ es transversal en el punto $x$, se dice entonces que $Y$ y $Z$ \emph{se cortan transversalmente en el punto $x$} (relativamente a $S$). Denotando por $X_s$, $Y_s$, $Z_s$, $T_s$ las fibras de $X$, $Y$, $Z$, $T$ en el punto $s$, imagen de $x$ en $S$, y teniendo en cuenta (5.2.3), (5.1.9) y (0, 16.5.12), se ve que, para que esto ocurra (cuando $X$, $Y$, $Z$ son lisos sobre $S$ en el punto $x$), es necesario y suficiente que $T$ sea liso sobre $S$ en el punto $x$, y que se tenga la relación
\[
\label{IV.17.13.7.1}
\tag{17.13.7.1}
\operatorname{codim}_x(T_s,X_s)=\operatorname{codim}_x(Y_s,X_s)+\operatorname{codim}_x(Z_s,X_s).
\]
\end{env}

\begin{proposition}[17.13.8]
\label{IV.17.13.8}
Sean $S$ un preesquema, $X$ un $S$-preesquema localmente de presentación finita sobre $S$, $Y$, $Z$ dos subpreesquemas de $X$, $T=\inf(Y,Z)$ el subpreesquema «intersección» de $Y$ y $Z$, $x$ un punto de $T$. Las condiciones siguientes son equivalentes:
\begin{enumerate}
  \item[a)] La inyección canónica $i:Y\to X$ es un morfismo transversal a $Z$ en el punto $x$, relativamente a $S$ (17.13.3).
  \item[a')] La inyección canónica $j:Z\to X$ es un morfismo transversal a $Y$ en el punto $x$, relativamente a $S$.
  \item[a'')] $Y$ y $Z$ se cortan transversalmente en el punto $x$, relativamente a $S$.
  \item[b)] $X$, $Y$, $Z$ son lisos sobre $S$ en el punto $x$, y el homomorfismo $\rho$ (17.13.5.4) es tal que
  \[
  \rho\otimes1:\Omega^1_{X/S}\otimes_{\sh{O}_X}k(x)\to
  (\Omega^1_{Y/S}\otimes_{\sh{O}_Y}k(x))\oplus
  (\Omega^1_{Z/S}\otimes_{\sh{O}_Z}k(x))
  \]
  es inyectivo.
\end{enumerate}
Cuando $x$ es racional sobre $k(s)$ (siendo $s$ la imagen de $x$ en $S$), estas condiciones son también equivalentes a la siguiente:
\begin{enumerate}
  \item[b')] El homomorfismo
  \[
  \mathrm{T}_x(i)-\mathrm{T}_x(j):\mathrm{T}_{Y/S}(x)\oplus\mathrm{T}_{Z/S}(x)\to\mathrm{T}_{X/S}(x)
  \]
  es sobreyectivo.
\end{enumerate}
Además, cuando las condiciones equivalentes a) a b) se cumplen en el punto $x$, se cumplen en una vecindad de $x$ en $T$, y, restringiendo $X$ a una vecindad de $x$, la sucesión (17.13.6.2) es exacta, y se tiene un isomorfismo canónico
\[
\label{IV.17.13.8.1}
\tag{17.13.8.1}
\sh{N}_{T/X}\xrightarrow{\sim}
(\sh{N}_{Y/X}\otimes_{\sh{O}_Y}\sh{O}_T)\oplus
(\sh{N}_{Z/X}\otimes_{\sh{O}_Z}\sh{O}_T).
\]
\end{proposition}

Las condiciones a), a'), a'') implican todas que $X$, $Y$, $Z$ son lisos sobre $S$ en el punto $x$. Sean además $m$, $m-a$, $m-b$ las dimensiones relativas de $X$, $Y$, $Z$ sobre $S$ en el punto $x$. Resulta entonces de (17.13.2), aplicado sustituyendo $X'$ por $Z$ e $Y'$ por $T=Y\times_XZ$, que las condiciones a), a') equivalen ambas a decir que $T$ es liso sobre $S$ en el punto $x$ y de dimensión relativa $m-a-b$ en este punto; pero, en virtud de (17.13.7), esto significa precisamente que se cumple la condición a''), de donde resulta la equivalencia de a), a') y a''). La equivalencia de a'') y b) (o b') cuando $x$ es racional sobre $k(s)$) se ha demostrado en (17.13.6), así como el hecho de que, si estas condiciones se satisfacen en el punto $x$, se satisfacen en una vecindad de $x$, y la exactitud de la sucesión (17.13.6.2) en tal vecindad. Queda por definir el isomorfismo canónico (17.13.8.1).

\oldpage[IV-4]{95}

Denotemos por $\alpha$ y $-\beta$ los homomorfismos que figuran en el segundo miembro de (17.13.5.4), y por $\gamma$ y $\delta$ los que figuran en el segundo miembro de (17.13.5.2). Decir que la sucesión (17.13.6.2)
\[
0\to\Omega^1_{X/S}\otimes_{\sh{O}_X}\sh{O}_T
\xrightarrow{\alpha\oplus(-\beta)}
(\Omega^1_{Y/S}\otimes_{\sh{O}_Y}\sh{O}_T)\oplus
(\Omega^1_{Z/S}\otimes_{\sh{O}_Z}\sh{O}_T)
\xrightarrow{\gamma+\delta}\Omega^1_{T/S}\to0
\]
es exacta significa que, en la categoría de los $\sh{O}_T$-módulos, $\Omega^1_{X/S}\otimes_{\sh{O}_X}\sh{O}_T$ se identifica canónicamente con el \emph{producto fibrado} de $\Omega^1_{Y/S}\otimes_{\sh{O}_Y}\sh{O}_T$ y $\Omega^1_{Z/S}\otimes_{\sh{O}_Z}\sh{O}_T$ sobre $\Omega^1_{T/S}$, para los homomorfismos $\gamma$ y $\delta$. El mismo razonamiento que en (0, 18.1.2 y 18.1.3), sustituyendo respectivamente los anillos y los ideales biláteros por módulos y submódulos, proporciona un diagrama conmutativo
{\scriptsize
\[
\begin{tikzcd}[column sep=small,row sep=large]
0 \ar[dr] & & 0 \ar[d] & 0 \ar[d] & \\
& (\sh{N}_{Y/X}\otimes\sh{O}_T)\oplus(\sh{N}_{Z/X}\otimes\sh{O}_T)
  \ar[r] \ar[d] \ar[dr,"\iota" description]
& \sh{N}_{Z/X}\otimes\sh{O}_T \ar[r,"\sim"] \ar[d]
& \sh{N}_{T/Y} \ar[d] & \\
0 \ar[r] & \sh{N}_{Y/X}\otimes\sh{O}_T \ar[r] \ar[d]
& \Omega^1_{X/S}\otimes\sh{O}_T \ar[r,"\alpha"] \ar[d,"\beta"] \ar[dr,"\varepsilon" description]
& \Omega^1_{Y/S}\otimes\sh{O}_T \ar[r] \ar[d,"\gamma"] & 0\\
0 \ar[r] & \sh{N}_{T/Z} \ar[r]
& \Omega^1_{Z/S}\otimes\sh{O}_T \ar[r,"\delta"] \ar[d]
& \Omega^1_{T/S} \ar[r] \ar[d] \ar[dr] & 0\\
& & 0 & 0 & 0
\end{tikzcd}
\]
}
donde las 3\textsuperscript{a} y 4\textsuperscript{a} filas y columnas son exactas en virtud de las hipótesis de lisura. El hecho de que el homomorfismo compuesto $\varepsilon=\gamma\circ\alpha=\delta\circ\beta$ sea el homomorfismo correspondiente a la inyección canónica $T\to X$ resulta de (0, 20.5.2.7). Como la diagonal del diagrama anterior es exacta, $\operatorname{Ker}(\varepsilon)$ se identifica canónicamente con $\sh{N}_{T/X}$, teniendo en cuenta que $T$ es liso sobre $S$ (17.2.5); por tanto, se tiene un isomorfismo canónico (17.13.8.1), inverso de $\iota$.

\begin{env}[17.13.9]
\label{IV.17.13.9}
Se pueden generalizar los resultados de (17.13.5) y (17.13.6) a un número finito cualquiera de $S$-morfismos $f_i:Y_i\to X$ ($i\in I$, siendo $I$ un conjunto finito). Para ello, denotemos por $X^I$ el producto de la familia $(X_i)_{i\in I}$ de $S$-preesquemas todos idénticos a $X$ (I, 3.3.5), y sean $p_i$ ($i\in I$) las proyecciones canónicas. El \emph{morfismo diagonal} $\Delta:X\to X^I$ se define (como en (I, 5.3.1)) como el único $S$-morfismo tal que $p_i\circ\Delta=1_X$ para todo $i\in I$. Es una $X$-sección de $X^I$ para el morfismo $p_1$, y, por tanto, (I, 5.3.11) una \emph{inmersión}.

\oldpage[IV-4]{96}

Sea entonces $Y$ el producto de los $S$-preesquemas $Y_i$ (para los morfismos compuestos $Y_i\xrightarrow{f_i}X\to S$), y $T$ el producto de los $X$-preesquemas $Y_i$ (para los morfismos $f_i$). Se demuestra, como en (I, 5.3.5), que se tiene un diagrama conmutativo
\[
\label{IV.17.13.9.1}
\tag{17.13.9.1}
\xymatrix{
  X \ar[d]^{\Delta} & T \ar[l]_-v \ar[d]^{j}\\
  X^I & Y \ar[l]^{u}
}
\]
(donde $u$ es el producto (sobre $S$) de los $f_i$), que hace de $T$ el producto de los $X^I$-preesquemas $X$ e $Y$.

Por recurrencia sobre el número de elementos de $I$, resulta de (16.4.23) que $\Omega^1_{X^I/S}$ (resp. $\Omega^1_{Y/S}$) se identifica canónicamente con la suma directa de los $p_i^*(\Omega^1_{X/S})$ (resp. de los $q_i^*(\Omega^1_{Y_i/S})$, si $q_i:Y\to Y_i$ son las proyecciones canónicas).

Supongamos ahora que $X$ es \emph{liso} sobre $S$ en un punto $x$ (y, por tanto, en una vecindad de este punto). Lo mismo ocurre entonces con $X^I$ en este punto (17.3.3, (iv)), y, restringiendo $X$ a una vecindad de $x$, se tiene, por (17.2.5), una sucesión exacta de $\sh{O}_X$-módulos localmente libres
\[
0\to\operatorname{gr}_1(\Delta)\to(\Omega^1_{X/S})^I\to\Omega^1_{X/S}\to0.
\]
Pongamos $\sh{J}=\operatorname{gr}_1(\Delta)$; se deduce de la sucesión anterior una sucesión exacta de $\sh{O}_T$-módulos localmente libres
\[
\label{IV.17.13.9.1-seq}
\tag{17.13.9.1}
0\to\sh{J}\otimes_{\sh{O}_X}\sh{O}_T\to(\Omega^1_{X/S})^I\otimes_{\sh{O}_X}\sh{O}_T\xrightarrow{\sigma}\Omega^1_{X/S}\otimes_{\sh{O}_X}\sh{O}_T\to0
\]
que corresponde a la primera fila del diagrama (17.13.1.2), y donde $\sigma$ no es otro que el homomorfismo canónico que hace corresponder a cada familia de secciones $(t_i)_{i\in I}$ de $\Omega^1_{X/S}\otimes_{\sh{O}_X}\sh{O}_T$ sobre un abierto de $T$ su \emph{suma}.

Por otra parte, lo que corresponde aquí a la segunda flecha vertical del diagrama (17.13.1.2) es el homomorfismo
\[
\label{IV.17.13.9.2}
\tag{17.13.9.2}
\tau:(\Omega^1_{X/S})^I\otimes_{\sh{O}_X}\sh{O}_T\to\bigoplus_{i\in I}(\Omega^1_{Y_i/S}\otimes_{\sh{O}_{Y_i}}\sh{O}_T)
\]
que, a toda familia $(t'_i\otimes1)_{i\in I}$, donde los $t'_i$ son aquí secciones sobre un abierto de $X$ de $\Omega^1_{X/S}$, asocia la suma de los $((f_i)_{Y_i/X/S}(t'_i))\otimes1$, con la notación de (16.4.18). El homomorfismo $\rho$ que corresponde al homomorfismo $\alpha$ de (17.13.1.3) es, por tanto, la \emph{restricción} a $\sh{J}\otimes_{\sh{O}_X}\sh{O}_T$ del homomorfismo $\tau$ anterior.

Precisadas estas hipótesis y notaciones, se puede aplicar a la situación de (17.13.9.1) la prop. (17.13.2), lo que da el
\end{env}

\begin{corollary}[17.13.10]
\label{IV.17.13.10}
Bajo las hipótesis generales de (17.13.9), sea $t$ un punto de $T$, $y_i$ ($i\in I$) sus proyecciones en los $Y_i$, y $x$ el punto de $X$ igual a cada uno de los $f_i(y_i)$. Supongamos que $X$ es liso sobre $S$ en el punto $x$ y de dimensión relativa $d$, y que cada uno de los $Y_i$ es liso sobre $S$ en el punto $y_i$; sea $d+c_i$ (siendo $c_i$ un entero positivo o negativo) la dimensión relativa de $Y_i$ en el punto $y_i$. Entonces las condiciones siguientes son equivalentes:

\oldpage[IV-4]{97}

\begin{enumerate}
  \item[a)] $T$ es liso sobre $S$ en el punto $t$, de dimensión relativa $d+\sum_{i\in I}c_i$.
  \item[b)] El homomorfismo
  \[
  \rho\otimes1:\sh{J}\otimes_{\sh{O}_X}k(t)\to\bigoplus_{i\in I}(\Omega^1_{Y_i/S}\otimes_{\sh{O}_{Y_i}}k(t))
  \]
  es inyectivo.
\end{enumerate}
Cuando $t$ es racional sobre el cuerpo $k(s)$ (donde $s$ es la imagen de $t$ en $S$), estas condiciones son también equivalentes a la siguiente:
\begin{enumerate}
  \item[b')] El homomorfismo transpuesto de $\rho\otimes1$
  \[
  \bigoplus_{i\in I}\mathrm{T}_{Y_i/S}(y_i)\to(\mathrm{T}_{X/S}(x))^I/\delta(\mathrm{T}_{X/S}(x))
  \]
  (donde $\delta$ es la aplicación diagonal) es sobreyectivo.\footnote{El texto impreso pone $\mathrm{T}_{Y_i/X}(y_i)$ en el miembro izquierdo; el transpuesto de $\rho\otimes1$ tiene como dominio la suma directa de los espacios tangentes relativos $\mathrm{T}_{Y_i/S}(y_i)$.}
\end{enumerate}
Basta observar que $X^I$ es liso sobre $S$ y de dimensión relativa $d\cdot\operatorname{Card}(I)$, e $Y$ es liso sobre $S$ y de dimensión relativa $d\cdot\operatorname{Card}(I)+\sum_{i\in I}c_i$.

Cuando se cumplen las condiciones equivalentes de (17.13.10), se dice también que los $f_i$ forman una \emph{familia de morfismos transversales en el punto $t$, relativamente} a $S$. Se ve asimismo que el conjunto de los puntos $t\in T$ donde esto ocurre es \emph{abierto} en $T$. La observación (17.13.4, (ii))\footnote{El texto impreso remite a (17.13.4, (iii)); la afirmación de apertura bajo la hipótesis de ser propio es la observación (ii).} muestra entonces que, si $X$ y los $Y_i$ son planos y localmente de presentación finita sobre $S$, y, además, $T$ es \emph{propio} sobre $S$, el conjunto de los $s\in S$ tales que los $f_i$ formen una familia de morfismos transversales en \emph{todos} los puntos de $T_s$, relativamente a $S$, es \emph{abierto} en $S$.
\end{corollary}

\begin{env}[17.13.11]
\label{IV.17.13.11}
Consideremos, en particular, el caso, que generaliza (17.13.7), en que los $Y_i$ ($i\in I$) son \emph{subpreesquemas} de $X$, siendo los $f_i$ las inyecciones canónicas, de modo que $T=\inf_{i\in I}(Y_i)$ es también el subpreesquema «intersección» de los $Y_i$, y $t=y_i=x$ para todo $i$; en lugar de decir que los $f_i$ forman una familia de morfismos transversales en el punto $x$, se dice también que \emph{los $Y_i$ se cortan transversalmente en el punto $x$} (relativamente a $S$). La condición a) de (17.13.10) se expresa asimismo mediante la relación que generaliza (17.13.7.1)
\[
\label{IV.17.13.11.1}
\tag{17.13.11.1}
\operatorname{codim}_x(T_s,X_s)=\sum_{i\in I}\operatorname{codim}_x((Y_i)_s,X_s).
\]
Además, se tiene la propiedad siguiente, que extiende (17.13.8.1), y proporciona otra demostración de ella cuando $I$ tiene 2 elementos:
\end{env}

\begin{corollary}[17.13.12]
\label{IV.17.13.12}
Cuando los $Y_i$ se cortan transversalmente en el punto $x$, se tiene un isomorfismo canónico
\[
\label{IV.17.13.12.1}
\tag{17.13.12.1}
\sh{N}_{T/X}\xrightarrow{\sim}\bigoplus_{i\in I}(\sh{N}_{Y_i/X}\otimes_{\sh{O}_{Y_i}}\sh{O}_T).
\]
\end{corollary}

Se puede reducir al caso en que los $Y_i$ son subpreesquemas cerrados de $X$, definidos por ideales cuasicoherentes $\sh{J}_i$, de modo que $T$ está definido por el ideal $\sh{J}=\sum_i\sh{J}_i$. Por definición del haz conormal a una inmersión (16.1.3), el homomorfismo canónico $\bigoplus_i\sh{J}_i\to\sh{J}$ da, al pasar a los cocientes, un homomorfismo \emph{sobreyectivo}
\[
\label{IV.17.13.12.2}
\tag{17.13.12.2}
\bigoplus_{i\in I}(\sh{N}_{Y_i/X}\otimes_{\sh{O}_{Y_i}}\sh{O}_T)\to\sh{N}_{T/X}.
\]

\oldpage[IV-4]{98}

Pero aquí los $\sh{O}_T$-módulos de los dos miembros de (17.13.12.2) son localmente libres y \emph{del mismo rango} $\sum_i c_i$ (si $c_i$ es el rango de $\sh{N}_{Y_i/X}$), en virtud de (17.2.5) y de la condición a) de (17.13.10); se concluye, por tanto, de Bourbaki, \emph{Alg. comm.}, cap. II, § 3, n\textsuperscript{o} 2, cor. de la prop. 6, que (17.13.12.2) es biyectivo, y (17.13.12.1) es el isomorfismo inverso.

\subsection{Caracterizaciones locales e infinitesimales de los morfismos lisos, de los morfismos no ramificados y de los morfismos étale.}
\label{subsection:IV.17.14}

\begin{proposition}[17.14.1]
\label{IV.17.14.1}
Sea $f:X\to Y$ un morfismo localmente de presentación finita, $x$ un punto de $X$, $y=f(x)$. Para que $f$ sea liso (resp. no ramificado, resp. étale) en el punto $x$, es necesario y suficiente que se satisfaga la condición siguiente:

Para todo esquema local $Y'=\operatorname{Spec}(A')$ con punto cerrado $y'$, todo morfismo $h:Y'\to Y$ tal que $h(y')=y$, todo subesquema cerrado $Y'_0=\operatorname{Spec}(A'/\mathfrak{J}')$ de $Y'$, donde el ideal $\mathfrak{J}'$ es de cuadrado nulo, y todo $Y$-morfismo $g_0:Y'_0\to X$ tal que $g_0(y')=x$, existe al menos un (resp. a lo sumo uno, resp. uno y solo uno) $Y$-morfismo $g:Y'\to X$ tal que $g_0$ sea su restricción a $Y'_0$.
\end{proposition}

Teniendo en cuenta las definiciones (17.1.1 y 17.3.1), se trata de demostrar que la condición del enunciado es \emph{suficiente} para que $f$ sea liso (resp. no ramificado) en el punto $x$.

\emph{(i) Caso de los morfismos lisos.} Se puede reducir al caso en que $Y=\operatorname{Spec}(A)$ y $X=\operatorname{Spec}(C)$ son afines, con $C=B/\mathfrak{R}$, donde $B=A[T_1,\ldots,T_n]$ es una $A$-álgebra de polinomios y $\mathfrak{R}$ un ideal de tipo finito de $B$. Para demostrar que $f$ es liso en el punto $x$, basta establecer que $\sh{O}_{X,x}$ es una $\sh{O}_{Y,y}$-álgebra formalmente lisa para las topologías discretas (17.5.1). Ahora bien, se tiene $\sh{O}_{X,x}=B_{\mathfrak q}/\mathfrak{R}B_{\mathfrak q}$, donde $\mathfrak q$ es un ideal primo de $B$, y, si $\mathfrak p$ es la imagen inversa de $\mathfrak q$ en $A$, $B_{\mathfrak q}$ es una $A_{\mathfrak p}$-álgebra formalmente lisa para las topologías discretas (0, 19.3.2 y 19.3.5). Por aplicación del criterio jacobiano (0, 22.6.1 y 20.5.12), basta, pues, ver que $B_{\mathfrak q}/\mathfrak{R}^2B_{\mathfrak q}$ es una extensión $A_{\mathfrak p}$-trivial de $B_{\mathfrak q}/\mathfrak{R}B_{\mathfrak q}$ por $\mathfrak{R}B_{\mathfrak q}/\mathfrak{R}^2B_{\mathfrak q}$. Pero esto resulta precisamente de la hipótesis aplicada a $A'=B_{\mathfrak q}/\mathfrak{R}^2B_{\mathfrak q}$ y $\mathfrak{J}'=\mathfrak{R}B_{\mathfrak q}/\mathfrak{R}^2B_{\mathfrak q}$ y al morfismo $g_0$ correspondiente al automorfismo identidad de $A'/\mathfrak{J}'=\sh{O}_{X,x}$.

\emph{(ii) Caso de los morfismos no ramificados.} Pongamos aquí $A=\sh{O}_{Y,y}$, $B=\sh{O}_{X,x}$, y observemos que, en virtud de (17.4.1, c')), basta demostrar que $\Omega^1_{B/A}=(\Omega^1_{X/Y})_x=0$. Ahora bien, con las notaciones de (0, 20.4.1), el anillo $C=\sh{P}^1_{B/A}=(B\otimes_A B)/\mathfrak{J}^2$ es local, puesto que lo es $C/(\mathfrak{J}/\mathfrak{J}^2)$, isomorfo a $B$. La hipótesis, aplicada a $A'=C$ y $\mathfrak{J}'=\mathfrak{J}/\mathfrak{J}^2$, muestra, por definición, que la aplicación $d_{B/A}:B\to\Omega^1_{B/A}$ es nula (0, 20.4.6), y, por tanto, que $\Omega^1_{B/A}=0$ por (0, 20.4.7).

\begin{proposition}[17.14.2]
\label{IV.17.14.2}
Sean $Y$ un preesquema localmente noetheriano, $f:X\to Y$ un morfismo localmente de tipo finito, $x$ un punto de $X$, $y=f(x)$. Para que $f$ sea liso (resp. no ramificado, resp. étale) en el punto $x$, es necesario y suficiente que satisfaga la condición de (17.14.1), suponiendo además que el anillo local $A'$ es artiniano, que $\mathfrak m'\mathfrak{J}'=0$, donde $\mathfrak m'$ es el ideal maximal de $A'$, y que el cuerpo residual $A'/\mathfrak m'$ de $A'$ es igual a $k(x)$.
\end{proposition}

Se trata nuevamente de demostrar que estas condiciones son suficientes en el caso de los morfismos lisos y en el de los morfismos no ramificados.

\oldpage[IV-4]{99}

\emph{(i) Caso de los morfismos lisos.} Si se pone $A=\sh{O}_{Y,y}$ y $B=\sh{O}_{X,x}$, se trata aquí de demostrar, teniendo en cuenta (17.5.3), que $B$ es una $A$-álgebra formalmente lisa para las topologías preádicas. Ahora bien, en virtud de (0, 22.1.4), esto se deduce de la hipótesis cuando se sustituye en ella la condición $\mathfrak m'\mathfrak{J}'=0$ por $\mathfrak{J}'\subset\mathfrak m'$. Pero, como $\mathfrak m'$ es nilpotente, se ve que la condición del enunciado ya es suficiente considerando los anillos $A'/\mathfrak m'^j\mathfrak{J}'$\footnote{El texto fuente escribe $\mathfrak m'^{,j}$ aquí y $\mathfrak m'^{,j}/\mathfrak m'^{,j+1}$ en (17.14.3); la coma espuria se elimina, como exigen la filtración por potencias y el testigo inglés mantenido.} y razonando por recurrencia sobre $j$, como en (0, 19.4.3).

\emph{(ii) Caso de los morfismos no ramificados.} Retomemos las notaciones de la demostración de (17.14.1, (ii)); para demostrar que el ideal $\mathfrak{J}/\mathfrak{J}^2$ de $C$ es nulo, observemos que $C$ es entonces un anillo local noetheriano, pues es un anillo local del anillo de un abierto afín de un subpreesquema de $X\times_YX$, que es localmente de tipo finito sobre $X$, y, por tanto, también sobre $Y$. Si $\mathfrak r$ es el ideal maximal de $C$, basta, pues, verificar que se tiene $\mathfrak{J}/\mathfrak{J}^2\subset\mathfrak r^n$ para todo $n$ (0$_{\mathrm I}$, 7.3.5), o también que $\mathfrak r^n=\mathfrak r^n+(\mathfrak{J}/\mathfrak{J}^2)$. Con las notaciones de (0, 20.4.6), esto significa también que, para todo $n$, las dos aplicaciones compuestas $B\xrightarrow{p_1}C\to C/\mathfrak r^n$ y $B\xrightarrow{p_2}C\to C/\mathfrak r^n$ son idénticas; pero, como $C/\mathfrak r^n$ es artiniano, esta identidad resulta de la hipótesis aplicada, por recurrencia sobre $n$, a $A'=C/\mathfrak r^n$ y $\mathfrak{J}'=(\mathfrak r^n+(\mathfrak{J}/\mathfrak{J}^2))/\mathfrak r^n$.

\begin{remark}[17.14.3]
\label{IV.17.14.3}
Obsérvese que, bajo las hipótesis de (17.14.2), si, además, el cuerpo $k(x)$ es una extensión \emph{finita} de $k(y)$, entonces los $A'$-módulos $\mathfrak m'^j/\mathfrak m'^{j+1}$ son $k(y)$-espacios vectoriales de dimensión finita, y, por tanto, \emph{a fortiori}, $A'$ es una $\sh{O}_{Y,y}$-álgebra finita.
\end{remark}

\subsection{Caso de los preesquemas sobre un cuerpo de base.}
\label{subsection:IV.17.15}

Recordemos primero (6.7.7, 6.7.8 y 6.8.1) la

\begin{proposition}[17.15.1]
\label{IV.17.15.1}
Sean $k$ un cuerpo y $X$ un preesquema localmente de tipo finito sobre $k$. Para que $X$ sea liso en un punto $x$, es necesario y suficiente que, para toda extensión radicial $k'$ de $k$ (o solamente para toda extensión finita radicial de $k$), el anillo local $\sh{O}_{X,x}\otimes_k k'$ sea regular. Si $k(x)$ es una extensión separable de $k$ (en particular, si $k$ es perfecto), esto equivale a decir que $\sh{O}_{X,x}$ es regular.
\end{proposition}

\begin{corollary}[17.15.2]
\label{IV.17.15.2}
Bajo las condiciones de (17.15.1), para que $X$ sea liso sobre $k$, es necesario y suficiente que $X$ sea geométricamente regular sobre $k$ (6.7.6), lo que implica que $X$ es regular. Si $k$ es perfecto, entonces $X$ es liso sobre $k$ si y solo si $X$ es regular.
\end{corollary}

\begin{proposition}[17.15.3]
\label{IV.17.15.3}
Sean $k$ un cuerpo, $X$ un preesquema localmente de tipo finito sobre $k$, $x$ un punto de $X$, $n=\dim_x(X)$. Sea $(g_i)_{1\leq i\leq n}$ una familia de $n$ secciones de $\sh{O}_X$ sobre $X$, y sea $f:X\to Y=\operatorname{Spec}(k[T_1,\ldots,T_n])$ el morfismo correspondiente al $k$-homomorfismo $k[T_1,\ldots,T_n]\to\Gamma(X,\sh{O}_X)$ que transforma $T_i$ en $g_i$ (I, 2.2.4). Las condiciones siguientes son equivalentes (e implican que $X$ es liso sobre $k$ en el punto $x$ y, por consiguiente, regular en el punto $x$):
\begin{enumerate}
  \item[a)] $f$ es étale en el punto $x$.
  \item[b)] Las imágenes de los $d_{X/k}g_i$ forman una base del $\sh{O}_{X,x}$-módulo $(\Omega^1_{X/k})_x$.
  \item[c)] Las imágenes de los $d_{X/k}g_i$ generan el $\sh{O}_{X,x}$-módulo $(\Omega^1_{X/k})_x$.
\end{enumerate}
\end{proposition}

Si $f$ es étale en el punto $x$, $X$ es liso sobre $k$ en el punto $x$ puesto que $Y=\operatorname{Spec}(k[T_1,\ldots,T_n])$ es liso sobre $k$ (17.3.8); el hecho de que a) implique b) es, pues, un caso particular de (17.11.3). Como b) implica trivialmente c), queda por ver que c) implica a). Observemos primero

\oldpage[IV-4]{100}

que la hipótesis entraña que el homomorfismo $(f^*(\Omega^1_{Y/k}))_x\to(\Omega^1_{X/k})_x$ es sobreyectivo; por tanto, reemplazando $X$ por una vecindad abierta de $x$, podemos suponer que el homomorfismo $f^*(\Omega^1_{Y/k})\to\Omega^1_{X/k}$ es sobreyectivo (Bourbaki, \emph{Alg. comm.}, cap. II, § 5, n\textsuperscript{o} 1, prop. 2); por consiguiente $f$ es \emph{no ramificado} (17.2.2). Veamos primero que podemos reducirnos al caso en que $x$ es \emph{racional sobre $k$}. En efecto, si ponemos $k'=k(x)$, $X'=X\otimes_k k'$, $Y'=Y\otimes_k k'=\operatorname{Spec}(k'[T_1,\ldots,T_n])$, existe un punto $x'\in X'$ por encima de $x$ tal que $k(x')=k'$. Para probar que $f$ es étale en el punto $x$, basta demostrar que $f'=f_{(k')}:X'\to Y'$ es étale en el punto $x'$ (17.7.1, (ii)); además $f'$ es no ramificado (17.3.3, (iii)) y se tiene $\dim_{x'}(X')=n$ (4.2.7). De la misma manera, reemplazando $k'$ por una extensión algebraicamente cerrada de $k'$, podemos suponer que $k$ es \emph{algebraicamente cerrado}. Pongamos entonces $y=f(x)$, $A=\sh{O}_{Y,y}$, $B=\sh{O}_{X,x}$; como el cuerpo residual de $B$ es igual a $k$, lo mismo ocurre con el de $A$, de modo que $x$ (resp. $y$) es un punto cerrado de $X$ (resp. $Y$) (I, 6.4.2) y, por consiguiente, $\dim(A)=\dim(B)=n$. Puesto que $f$ es no ramificado, el homomorfismo $\widehat A\to\widehat B$ es sobreyectivo (17.4.4, f'')); pero como $\dim(\widehat A)=\dim(\widehat B)=n$, y $\widehat A$, al ser un anillo regular, es íntegro, el homomorfismo $\widehat A\to\widehat B$ es también inyectivo (0, 16.3.10). Por tanto, este homomorfismo es biyectivo, lo que termina de demostrar que $f$ es étale en el punto $x$ (17.6.3, e'')).

\begin{corollary}[17.15.4]
\label{IV.17.15.4}
Bajo las hipótesis generales de (17.15.3), supongamos además que $k(x)$ sea una extensión finita y separable de $k$ y que los gérmenes $(g_i)_x$ pertenezcan a $\mathfrak m_x$. Entonces las condiciones a), b), c) de (17.15.3) son también equivalentes a cada una de las siguientes:
\begin{enumerate}
  \item[d)] Los $n$ gérmenes $(g_i)_x$ generan el ideal maximal $\mathfrak m_x$ de $\sh{O}_{X,x}$.
  \item[d')] El anillo $\sh{O}_{X,x}$ es regular y los $(g_i)_x$ forman un sistema regular de parámetros para este anillo (0, 17.1.6).
\end{enumerate}
\end{corollary}

En efecto, d') implica trivialmente d). Por otra parte, como $k(x)$ es una extensión finita y separable de $k$, se tiene $\Omega^1_{k(x)/k}=0$ (0, 20.6.20) y $k(x)=\sh{O}_{X,x}/\mathfrak m_x$ es una $k$-álgebra formalmente lisa para las topologías discretas (0, 19.6.1), por lo que la sucesión exacta (0, 20.5.14.1) se aplica a $A=k$, $B=\sh{O}_{X,x}$, $\mathfrak{R}=\mathfrak m_x$, y proporciona un isomorfismo canónico
\[
\delta:\mathfrak m_x/\mathfrak m_x^2\xrightarrow{\sim}(\Omega^1_{X/k})_x\otimes_{\sh{O}_{X,x}}k(x).
\]
De aquí se deduce primero la equivalencia de las condiciones c) y d), teniendo en cuenta el lema de Nakayama. Por otra parte, si $f$ es étale en el punto $x$, el anillo $\sh{O}_{X,x}$ es regular y de dimensión $n$, puesto que $x$ es un punto cerrado de $X$ (I, 6.4.2), y $n$ elementos de $\mathfrak m_x$ que generan este ideal maximal forman entonces necesariamente un sistema regular de parámetros para $\sh{O}_{X,x}$ (0, 17.1.6), lo que prueba que a) implica d').

\begin{proposition}[17.15.5]
\label{IV.17.15.5}
Sean $k$ un cuerpo, $X$ un preesquema localmente de tipo finito sobre $k$, $x$ un punto de $X$, $n=\dim_x(X)$. Las condiciones siguientes son equivalentes:
\begin{enumerate}
  \item[a)] $X$ es liso sobre $k$ en el punto $x$.
  \item[b)] $X$ es diferencialmente liso sobre $k$ en el punto $x$.
  \item[c)] $(\Omega^1_{X/k})_x$ es un $\sh{O}_{X,x}$-módulo libre de rango $n$.
  \item[d)] $(\Omega^1_{X/k})_x$ es un $\sh{O}_{X,x}$-módulo que admite un sistema de $n$ generadores.

\oldpage[IV-4]{101}

  \item[e)] Existe una extensión perfecta $k'$ de $k$ y una vecindad abierta $U$ de $x$ en $X$ tales que el preesquema $U\otimes_k k'$ sea regular.
\end{enumerate}
\end{proposition}

El hecho de que $X$ sea liso sobre $k$ en el punto $x$ entraña la existencia de una vecindad abierta $U$ de $x$ que sea lisa sobre $k$, por lo que $U\otimes_k k'$ es regular para toda extensión $k'$ de $k$ (17.15.2), lo que prueba que a) implica e). Recíprocamente, e) implica a), pues entonces $U\otimes_k k'$ es liso sobre $k'$ (17.15.2), y por tanto $U$ es liso sobre $k$ (17.7.1, (ii)).

Ya hemos demostrado que a) implica c) (17.10.2); c) implica d) trivialmente y el hecho de que d) implique a) resulta de (0, 20.4.7) y (17.15.3, c)).

Finalmente, ya hemos visto que a) implica b) (17.12.4). Recíprocamente, para demostrar que b) implica a), podemos reducirnos al caso en que $x$ es racional sobre $k$, considerando, como en la demostración de (17.15.3), un punto $x'$ de $X'=X\otimes_k k'$ por encima de $x$ y tal que $k(x')=k'=k(x)$, utilizando de nuevo (17.7.1, (ii)) y el hecho de que, si $X$ es diferencialmente liso en el punto $x$, $X'$ es diferencialmente liso en el punto $x'$ (16.10.4). Suponiendo, pues, que $x$ es racional sobre $k$, el hecho de que $f$ sea liso en $x$ resulta entonces de la hipótesis b) y de (17.12.5), aplicado a la $k$-sección $u:\operatorname{Spec}(k)\to X$ tal que $u(\operatorname{Spec}(k))=\{x\}$.

\begin{corollary}[17.15.6]
\label{IV.17.15.6}
Sea $X$ un preesquema de tipo finito sobre un cuerpo $k$. Para que $X$ sea liso sobre $k$, es necesario y suficiente que el $\sh{O}_X$-Módulo $\Omega^1_{X/k}$ sea localmente libre y que los anillos locales en los puntos maximales de $X$ sean cuerpos, extensiones separables de $k$ (esta última condición se verifica automáticamente si $k$ es un cuerpo perfecto y $X$ un preesquema reducido).
\end{corollary}

Las condiciones son necesarias, pues si $X$ es liso sobre $k$, resulta de (17.10.2) que $\Omega^1_{X/k}$ es localmente libre; por otra parte, $X$ es regular, por tanto \emph{a fortiori} reducido, y en todo punto maximal $x$ de $X$, $\sh{O}_{X,x}$ es un cuerpo, que debe ser una $k$-álgebra formalmente lisa para las topologías discretas (17.5.1), y por consiguiente una extensión separable de $k$ (0, 19.6.1).

Las condiciones son suficientes. En efecto, podemos reducirnos al caso en que $X$ es conexo, por lo que $\Omega^1_{X/k}$ es localmente libre de rango constante $n$. Para todo punto maximal $x$ de $X$, se tiene entonces $\operatorname{rg}_{k(x)}\Omega^1_{k(x)/k}=n$, puesto que $\sh{O}_{X,x}=k(x)$; como por hipótesis $k(x)$ es una extensión separable de $k$, se tiene $\Upsilon_{k(x)/k}=0$ (0, 20.6.3), por lo que $\deg.\operatorname{tr}_k k(x)=n$ por la igualdad de Cartier (0, 21.7.1). Todas las componentes irreducibles de $X$ tienen, pues, la misma dimensión $n$ (5.2.1), y concluimos que $X$ es liso sobre $k$ en todo punto en virtud de que c) implica a) en (17.15.5).

\begin{corollary}[17.15.7]
\label{IV.17.15.7}
Si $k$ es de característica 0, entonces, para que $X$ sea liso sobre $k$ en el punto $x$, es necesario y suficiente que $(\Omega^1_{X/k})_x$ sea un $\sh{O}_{X,x}$-módulo libre.
\end{corollary}

Esto resulta de (16.12.2) y de la equivalencia de b) y a) en (17.15.5).

\begin{proposition}[17.15.8]
\label{IV.17.15.8}
Sean $k$ un cuerpo, $X$ un preesquema localmente de tipo finito sobre $k$, $x$ un punto de $X$; pongamos $n=\dim_x(X)$, $r=\dim(\sh{O}_{X,x})$, de modo que $n-r=\deg.\operatorname{tr}_k k(x)$ (5.2.3). Sea $(g_i)_{1\leq i\leq n}$ una familia de $n$ secciones de $\sh{O}_X$ sobre $X$, tales que $(g_i)_x\in\mathfrak m_x$ para $1\leq i\leq r$; sea $f:X\to Y=\operatorname{Spec}(k[T_1,\ldots,T_n])$ el morfismo correspondiente al $k$-homomorfismo $k[T_1,\ldots,T_n]\to\Gamma(X,\sh{O}_X)$ que transforma $T_i$ en $g_i$. Las condiciones siguientes son equivalentes (e implican que $X$ es liso sobre $k$ en el punto $x$ y que $k(x)$ es una extensión separable de $k$):
\begin{enumerate}
  \item[a)] $f$ es étale en el punto $x$.

\oldpage[IV-4]{102}

  \item[b)] Los gérmenes $(g_i)_x$ para $1\leq i\leq r$ generan $\mathfrak m_x$ (y por consiguiente forman un sistema regular de parámetros para $\sh{O}_{X,x}$ (0, 17.1.6)) y las imágenes en $\Omega^1_{k(x)/k}$ de los elementos $(d_{X/k}g_i)_x$ para $r+1\leq i\leq n$ generan $\Omega^1_{k(x)/k}$.
\end{enumerate}
\end{proposition}

Se tiene en efecto (0, 20.5.12.1) la sucesión exacta de $k(x)$-módulos
\[
\label{IV.17.15.8.1}
\tag{17.15.8.1}
\mathfrak m_x/\mathfrak m_x^2\to(\Omega^1_{X/k})_x\otimes_{\sh{O}_{X,x}}k(x)\to\Omega^1_{k(x)/k}\to0
\]
y la condición b) entraña, por tanto, que los $(d_{X/k}g_i)_x$ generan $(\Omega^1_{X/k})_x$ teniendo en cuenta el lema de Nakayama; el hecho de que b) implica a) resulta, pues, de (17.15.3). Recíprocamente, si a) se verifica, los $(d_{X/k}g_i)_x$ para $1\leq i\leq n$ forman una base de $(\Omega^1_{X/k})_x$ en virtud de (17.15.3). Si $t_i$ ($1\leq i\leq n$) es la imagen de $(g_i)_x$ en $k(x)$, se deduce de lo anterior que los $dt_i$ generan $\Omega^1_{k(x)/k}$ para $1\leq i\leq n$, y como por hipótesis $t_i=0$ para $1\leq i\leq r$, los $dt_i$ para $r+1\leq i\leq n$ generan ya $\Omega^1_{k(x)/k}$. Como $\deg.\operatorname{tr}_k k(x)=n-r$, se sigue de la igualdad de Cartier (0, 21.7.1) que $\Upsilon_{k(x)/k}=0$, por lo que $k(x)$ es una extensión separable de $k$ (0, 20.6.3), y la sucesión de $k(x)$-espacios vectoriales
\[
\label{IV.17.15.8.2}
\tag{17.15.8.2}
0\to\mathfrak m_x/\mathfrak m_x^2\to(\Omega^1_{X/k})_x\otimes_{\sh{O}_{X,x}}k(x)\to\Omega^1_{k(x)/k}\to0
\]
es exacta ((0, 20.5.14) y (0, 19.6.1)); además los $dt_i$ para $r+1\leq i\leq n$ forman una base de $\Omega^1_{k(x)/k}$, por lo que ninguno de los $t_i$ tales que $r+1\leq i\leq n$ puede ser nulo. Esto demuestra que las imágenes de los $(g_i)_x$ en $\mathfrak m_x/\mathfrak m_x^2$ para $1\leq i\leq r$ generan necesariamente $\mathfrak m_x/\mathfrak m_x^2$, y por consiguiente los $(g_i)_x$ para $1\leq i\leq r$ generan $\mathfrak m_x$ en virtud del lema de Nakayama, lo que termina de probar que a) implica b).

\begin{corollary}[17.15.9]
\label{IV.17.15.9}
Sean $X$ un preesquema localmente de tipo finito sobre un cuerpo $k$, $x$ un punto de $X$ y pongamos $n=\dim_x(X)$, $r=\dim(\sh{O}_{X,x})$. Las condiciones siguientes son equivalentes:
\begin{enumerate}
  \item[a)] El anillo $\sh{O}_{X,x}$ es regular y $k(x)$ es una extensión separable de $k$.
  \item[b)] $X$ es liso sobre $k$ en el punto $x$, y el homomorfismo canónico
  \[
  \mathfrak m_x/\mathfrak m_x^2\to(\Omega^1_{X/k})_x\otimes_{\sh{O}_{X,x}}k(x)
  \]
  es inyectivo.
  \item[c)] Existen $n$ secciones $g_i$ de $\sh{O}_X$ sobre una vecindad abierta $U$ de $x$ tales que $(g_i)_x\in\mathfrak m_x$ para $1\leq i\leq r$ y que el morfismo $U\to\operatorname{Spec}(k[T_1,\ldots,T_n])$ correspondiente a los $g_i$ (cf. (17.15.8)) sea étale en el punto $x$.
  \item[d)] Existen $n$ secciones $g_i$ ($1\leq i\leq n$) de $\sh{O}_X$ sobre una vecindad abierta $U$ de $x$, tales que los $(g_i)_x$ para $1\leq i\leq r$ generan $\mathfrak m_x$ y que las imágenes en $\Omega^1_{k(x)/k}$ de los $(d_{U/k}g_i)_x$ para $r+1\leq i\leq n$ generan $\Omega^1_{k(x)/k}$.
\end{enumerate}
\end{corollary}

La equivalencia de c) y d) resulta de (17.15.8). Además, se ha visto en la demostración de (17.15.8) que la condición c) implica que $X$ es liso sobre $k$ en el punto $x$ y que la sucesión (17.15.8.2) es exacta, por lo que c) implica b). La condición b) entraña que el anillo $\sh{O}_{X,x}$ es regular, por lo que $\mathfrak m_x/\mathfrak m_x^2$ es de rango $r$; además, $(\Omega^1_{X/k})_x$ es entonces un $\sh{O}_{X,x}$-módulo libre de rango $n$ (17.10.2); como la sucesión (17.15.7.2) es exacta por hipótesis, $\Omega^1_{k(x)/k}$ es de rango $n-r=\deg.\operatorname{tr}_k k(x)$ y la igualdad de Cartier muestra que $\Upsilon_{k(x)/k}=0$, por lo que $k(x)$ es separable sobre $k$ (0, 20.6.3); así b) implica a). Finalmente, si a)

\oldpage[IV-4]{103}

se verifica, se deduce de nuevo de la igualdad de Cartier que $\Omega^1_{k(x)/k}$ es de rango $n-r$. Como, por otra parte, el anillo $\sh{O}_{X,x}$ es regular, la existencia de los $g_i$ que verifican las condiciones de d) es inmediata.

\begin{remarks}[17.15.10]
\label{IV.17.15.10}
\begin{enumerate}
  \item[(i)] Para que un preesquema de tipo finito $X$ sobre $k$ sea liso sobre $k$, no basta que $\Omega^1_{X/k}$ sea localmente libre, como muestra el ejemplo en que $X=\operatorname{Spec}(K)$, siendo $K$ una extensión finita no separable de $k$.
  \item[(ii)] Cuando $k$ no es perfecto, puede ocurrir que $X$ sea liso sobre $k$ sin que $k(x)$ sea separable sobre $k$. Tenemos un ejemplo tomando $X=\operatorname{Spec}(k[T])$ y para $x$ el punto correspondiente al ideal primo principal $(T^p-\lambda)$ ($p>0$ es la característica de $k$, $\lambda\in k-k^p$).
  \item[(iii)] Sin embargo, si $X$ es un preesquema liso sobre $k$, el conjunto de los puntos cerrados de $X$ tales que $k(x)$ sea una extensión separable (y finita) de $k$ es denso en $X$. En efecto, sea $f:X\to\operatorname{Spec}(k)$ el morfismo estructural; para todo $x_0\in X$, existe una vecindad abierta $U$ de $x_0$ en $X$ y una factorización de $f|U:U\xrightarrow{g}\operatorname{Spec}(k[T_1,\ldots,T_n])\xrightarrow{h}\operatorname{Spec}(k)$ en la que $g$ es étale (17.11.4). Como podemos reducirnos al caso en que $k$ no es perfecto, y por tanto es infinito, el conjunto de los puntos del abierto $g(U)$ que son racionales sobre $k$ no es vacío; si $y$ es un punto de este tipo y $x\in U$ un punto por encima de $y$, $x$ es cerrado en $X$ y $k(x)$ es separable sobre $k(y)=k$ (17.6.2).
\end{enumerate}
\end{remarks}

\begin{proposition}[17.15.11]
\label{IV.17.15.11}
Sea $X$ un preesquema de tipo finito sobre un cuerpo $k$. Las condiciones siguientes son equivalentes:
\begin{enumerate}
  \item[a)] $X$ es étale sobre $k$.
  \item[b)] $X$ es no ramificado sobre $k$.
  \item[c)] $X$ es isomorfo a $\operatorname{Spec}(A)$, donde $A$ es una $k$-álgebra finita y separable.
\end{enumerate}
\end{proposition}

Esto resulta de (17.6.2) y (17.4.2), teniendo en cuenta que la condición b) implica aquí que $X$ es finito sobre $k$.

\begin{proposition}[17.15.12]
\label{IV.17.15.12}
Sean $k$ un cuerpo y $X$ un preesquema localmente de tipo finito sobre $k$. Para que exista un abierto en todas partes denso $U$ de $X$ que sea liso sobre $k$, es necesario y suficiente que para todo punto maximal $x$ de $X$, $X$ sea reducido en este punto y que $k(x)$ sea una extensión separable de $k$. Para que exista un abierto en todas partes denso $U$ de $X$ tal que $U_{\mathrm{red}}$ sea liso sobre $k$, es necesario y suficiente que para todo punto maximal $x$ de $X$, $k(x)$ sea una extensión separable de $k$.
\end{proposition}

La segunda afirmación resulta evidentemente de la primera. Decir que existe un abierto en todas partes denso $U$ liso sobre $k$ significa que $X$ es liso sobre $k$ en cada uno de sus puntos maximales $x$. Para ello es necesario que $\sh{O}_{X,x}$ sea regular (17.15.1) y, \emph{a fortiori}, reducido, por tanto un cuerpo, igual a $k(x)$; además (17.15.1), $k(x)\otimes_k k'$ debe ser regular para toda extensión radicial $k'$ de $k$, por lo que $k(x)$ debe ser una extensión separable de $k$ (4.3.5); el recíproco es inmediato, en virtud de (17.15.1).

\begin{corollary}[17.15.13]
\label{IV.17.15.13}
Sea $X$ un preesquema algebraico sobre un cuerpo $k$. Entonces existen un abierto en todas partes denso $U$ en $X$ y una extensión finita radicial $k'$ de $k$ tales que $(U_{(k')})_{\mathrm{red}}$ sea liso sobre $k'$.
\end{corollary}

En efecto, existe una extensión finita radicial $k'$ tal que $(X_{(k')})_{\mathrm{red}}$ sea geométricamente reducido sobre $k'$ (4.6.6), lo que equivale a decir (4.6.1) que para todo punto maximal $x'$ de $X_{(k')}$, $k(x')$ sea una extensión separable de $k'$. Se concluye, pues, de (17.15.12) que existe un abierto en todas partes denso $U'$ de $X_{(k')}$ tal que $U'_{\mathrm{red}}$ sea liso sobre $k'$. Pero el \oldpage[IV-4]{104}morfismo $X_{(k')}\to X$ es un homeomorfismo (2.4.5), por lo que $U'$ es de la forma $U_{(k')}$, donde $U$ es un abierto en todas partes denso de $X$.

\begin{proposition}[17.15.14]
\label{IV.17.15.14}
Sea $X$ un preesquema algebraico sobre $k$, de dimensión $\leq1$. Entonces existe una extensión finita radicial $k'$ de $k$ tal que el normalizado (II, 6.3.8) de $(X_{(k')})_{\mathrm{red}}$ sea liso sobre $k'$.
\end{proposition}

Demostremos primero los dos lemas siguientes:

\begin{lemma}[17.15.14.1]
\label{IV.17.15.14.1}
Sean $X$ un preesquema reducido cuyo conjunto de componentes irreducibles sea localmente finito, $f:Y\to X$ un morfismo. Para que $Y$ sea $X$-isomorfo al normalizado $X'$ de $X$ (II, 6.3.8), es necesario y suficiente que se verifiquen las condiciones siguientes:
\begin{enumerate}
  \item[(i)] $Y$ es normal;
  \item[(ii)] $f$ es un morfismo entero y birracional.
\end{enumerate}
Cuando existe en $X$ un abierto denso y normal (lo que siempre ocurre cuando $X$ es excelente (7.8.3), en particular cuando $X$ es localmente de tipo finito sobre un cuerpo), se puede reemplazar la condición (ii) por la siguiente:
\begin{enumerate}
  \item[(ii bis)] $f$ es entero y existe un abierto denso $U$ en $X$ tal que $f^{-1}(U)$ sea denso en $Y$ y que la restricción $f^{-1}(U)\to U$ de $f$ sea un isomorfismo.
\end{enumerate}
\end{lemma}

Como la cuestión es local sobre $X$, podemos suponer que $X$ tiene un número finito de componentes irreducibles; sean $X_i$ ($1\leq i\leq m$) los subpreesquemas reducidos de $X$ que tienen esas componentes como espacios subyacentes. Sabemos (II, 6.3.6) que $X'$ es la suma de los normalizados $X'_i$ de los $X_i$; cada uno de los morfismos estructurales $f_i:X'_i\to X_i$ es, pues, entero y birracional, por lo que $f$ es entero y birracional; es inmediato que (ii) implica (ii bis) cuando existe un abierto $V$ denso y normal en $X$, tomando $U=V$. Recíprocamente, si se cumplen las condiciones (i) y (ii), $Y$ tiene un número finito de componentes irreducibles $Y_i$ ($1\leq i\leq m$), con $Y_i$ dominante sobre $X_i$ para cada índice $i$ (6.15.4); $Y$ es la suma de los $Y_i$ puesto que es normal, y el morfismo $g_i:Y_i\to X_i$ en el que se factoriza la restricción $Y_i\to X$ de $f$ (I, 5.2.2) es birracional; como $X_i$ y $Y_i$ son íntegros, se deduce del hecho de que $g_i$ es entero y $Y_i$ normal que, para todo abierto afín $V$ de $X_i$, $\Gamma(g_i^{-1}(V),\sh{O}_{Y_i})$ es la clausura entera de $\Gamma(V,\sh{O}_{X_i})$, por lo que $Y$ se identifica canónicamente con $X'$ (II, 6.3.4). Finalmente, si suponemos verificada la condición (ii bis), podemos reducirnos al caso en que $U$ es unión de abiertos irreducibles disjuntos $U_i$ ($1\leq i\leq m$), por lo que $f^{-1}(U)$ es unión disjunta de los $V_i=f^{-1}(U_i)$. Como los $V_i$ son irreducibles y $f^{-1}(U)$ es denso en $Y$, las clausuras $\overline{V_i}$ son las componentes irreducibles de $Y$ y, por consiguiente, $f$ es birracional.

\begin{lemma}[17.15.14.2]
\label{IV.17.15.14.2}
Sean $k$ un cuerpo y $X$ un preesquema algebraico sobre $k$. Entonces existe una extensión finita radicial $k'$ de $k$ tal que $(X_{(k')})_{\mathrm{red}}$ sea geométricamente reducido sobre $k'$ y su normalizado sea geométricamente normal sobre $k'$.
\end{lemma}

Teniendo en cuenta (4.6.6) y (I, 5.1.8), podemos reducirnos al caso en que $X$ ya sea geométricamente reducido sobre $k$. Sean $p$ el exponente característico de $k$ y $k_1=k^{p^{-\infty}}$ la menor extensión perfecta de $k$; $k_1$ es, pues, límite inductivo de las extensiones finitas radiciales de $k$. Pongamos $X_1=X_{(k_1)}$, que es reducido por hipótesis, y sea $Y_1$ su normalizado; si $f_1:Y_1\to X_1$ es el morfismo estructural, $f_1$ es, por tanto, finito (7.8.3, (vi)) y sobreyectivo, y existe un abierto denso $U_1$ en $X_1$ tal que $f_1^{-1}(U_1)=V_1$ sea denso en $Y_1$ y que \oldpage[IV-4]{105}la restricción $V_1\to U_1$ de $f_1$ sea un isomorfismo (17.15.14.1). Aplicando (8.9.1) y (8.10.5, (vi) y (x)), se ve primero que existen una extensión finita radicial $k''$ de $k$ y un morfismo finito sobreyectivo $Y''\to X''=X_{(k'')}$ tales que $Y_1=Y''\times_{X''}X_1=Y''\otimes_{k''}k_1$; puesto que $Y_1$ es normal y $k_1$ perfecto, se deduce de (6.7.7) que $Y''$ es geométricamente normal sobre $k''$. Como las proyecciones $X_1\to X''$ y $Y_1\to Y''$ son morfismos enteros, sobreyectivos y radiciales, son homeomorfismos (2.4.5), y si $U''$ y $V''$ son las imágenes de $U_1$ y $V_1$ en $X''$ y $Y''$, respectivamente, son abiertos densos en $X''$ y $Y''$, respectivamente, tales que $V''=f''^{-1}(U'')$, donde $f'':Y''\to X''$ es el morfismo estructural. Como $U_1=U''\otimes_{k''}k_1$, resulta de (8.10.5, (i)) que existe una extensión finita radicial $k'$ de $k''$ tal que, si ponemos $X'=X\otimes_k k'=X''\otimes_{k''}k'$, $Y'=Y''\otimes_{k''}k'$, $f'=f''_{(k')}:Y'\to X'$, y si $U'$ y $V'$ son las imágenes de $U''$ y $V''$ en $X'$ y $Y'$ respectivamente, la restricción $V'\to U'$ de $f'$ sea un isomorfismo. Como $Y'$ es normal y $f'$ entero y birracional, se concluye de (17.15.14.1) que $Y'$ es isomorfo al normalizado de $X'$, lo que prueba el lema, puesto que $Y'$ es geométricamente normal.

Volvamos ahora a la demostración de (17.15.14), y apliquemos a $k$ y $X$ el lema (17.15.14.2); como $\dim(X)\leq1$, también se tiene $\dim(X_{(k')})\leq1$ (4.1.4), y el normalizado $Y'$ de $X_{(k')}$ tiene asimismo dimensión $\leq1$ (5.4.2 y II, 7.4.6). Decir que $Y'$ es geométricamente normal sobre $k'$ equivale entonces, en virtud de las definiciones y de (II, 7.4.5), a que $Y'$ sea geométricamente regular sobre $k'$, y por tanto liso sobre $k'$ (17.5.1), lo que prueba (17.15.14).

\begin{proposition}[17.15.15]
\label{IV.17.15.15}
Sea $f:X\to Y$ un morfismo localmente de presentación finita. Para que $f$ sea liso en un punto $x\in X$, es necesario y suficiente que $f$ sea plano en el punto $x$ y que $\Omega^1_f$ sea un $\sh{O}_X$-Módulo localmente libre en una vecindad de $x$, de rango en $x$ igual a $\dim_x f$.
\end{proposition}

La necesidad de las condiciones resulta de (17.5.1) y (17.10.2). Recíprocamente, si se verifican estas condiciones, y si $y=f(x)$, basta demostrar (17.5.1) que $f^{-1}(y)$ es liso sobre $k(y)$ en el punto $x$; pero esto resulta de la definición de $\dim_x f$ (17.10.1), de (16.4.5) y de (17.15.5).

\subsection{Cuasi-secciones de morfismos planos o lisos.}
\label{subsection:IV.17.16}

Los enunciados de este número completan los de (14.5), mediante hipótesis de planitud.

\begin{proposition}[17.16.1]
\label{IV.17.16.1}
Sea $f:X\to S$ un morfismo plano localmente de presentación finita. Sean $s\in S$, $x$ un punto cerrado de $X_s$ tal que $\sh{O}_{X_s,x}$ sea un anillo de Cohen--Macaulay; entonces existe una vecindad abierta $U$ de $s$ en $S$ y un subpreesquema $X'\subset f^{-1}(U)$ de $X$ tal que $x\in X'$ y que el morfismo $X'\to U$, restricción de $f$, sea plano, cuasi-finito y de presentación finita.
\end{proposition}

La cuestión es local sobre $X$ y $S$, por lo que podemos suponer que $f$ es de presentación finita.\footnote{La transcripción francesa imprime «sur $X$ et $Y$»; el morfismo y el testigo inglés muestran que la base correcta es $S$.} Sea $(t_i)_{1\leq i\leq m}$ un sistema de parámetros del anillo local $\sh{O}_{X_s,x}$; existe una vecindad abierta afín $V=\Spec(A)$ de $x$ en $X$ y $m$ secciones $g_i$ de $\sh{O}_X$ sobre $V$ tales que

\oldpage[IV-4]{106}

las imágenes de los $(g_i)_x\in\sh{O}_{X,x}$ en $\sh{O}_{X_s,x}=\sh{O}_{X,x}/\mathfrak m_s\sh{O}_{X,x}$ sean iguales a los $t_i$. Sea $X'$ el subpreesquema cerrado de $V$ definido por el ideal de $A$ generado por los $g_i$; como la sucesión $(t_i)$ es regular por hipótesis (0, 16.5.7), resulta de (11.3.8) que, reemplazando si es necesario $V$ por una vecindad abierta más pequeña, podemos suponer que el morfismo $X'\to S$, restricción de $f$, es plano y de presentación finita. Por otra parte, puesto que los $t_i$ forman un sistema de parámetros de $\sh{O}_{X_s,x}$, el anillo $\sh{O}_{X'_s,x}$ es artiniano por definición, y como $x$ es cerrado en $X'_s$, se concluye que está aislado en $X'_s$. Reemplazando $V$ si es necesario por una vecindad más pequeña de $x$, se concluye, gracias a (13.1.4), que el morfismo $X'\to S$ es cuasi-finito.

\begin{corollary}[17.16.2]
\label{IV.17.16.2}
Sea $f:X\to S$ un morfismo fielmente plano y localmente de presentación finita. Entonces existe un morfismo $g:S'\to S$, fielmente plano, localmente de presentación finita y localmente cuasi-finito, tal que exista un $S$-morfismo $S'\to X$ (en otras palabras, que exista una $S'$-sección de $X'=X\times_S S'$ (I, 3.3.14)). Si $S$ es cuasi-compacto (resp. cuasi-compacto y cuasi-separado), se puede suponer que $S'$ es afín (resp. $S'$ afín y el morfismo $g$ de presentación finita y cuasi-finito).\footnote{La frase francesa omite «de présentation finie»; se restituye según el testigo inglés y la condición de la proposición.}
\end{corollary}

Para todo $s\in S$, la fibra $X_s$ no es vacía por hipótesis y es un preesquema localmente de tipo finito sobre $k(s)$; el conjunto $U_s$ de los puntos $x$ de $X_s$ donde $\sh{O}_{X_s,x}$ es un anillo de Cohen--Macaulay es abierto en $X_s$ (6.11.3) y no es vacío, puesto que contiene los puntos maximales de $X_s$ (0, 16.5.1); contiene, por tanto, un punto $x_s$ cerrado en $X_s$ (10.4.7). Sea $X'(s)$ un subpreesquema afín de $X$ que contenga $x_s$ y verifique las condiciones de (17.16.1). Para obtener un preesquema $S'$ que verifique las condiciones del enunciado, basta tomar la suma de los $X'(s)$, recorriendo $s$ el conjunto $S$. Puesto que el morfismo $X'(s)\to S$ es plano y localmente de presentación finita, la imagen $U(s)$ de $X'(s)$ es abierta en $S$ (2.4.6); cuando $S$ es cuasi-compacto, podemos, pues, recubrir $S$ mediante un número finito de los $U(s_i)$ y el preesquema $S'$ suma de los $X'(s_i)$ sigue respondiendo a la cuestión y es afín. Si además $S$ es cuasi-separado, podemos suponer que las inmersiones abiertas $U(s_i)\to S$ son cuasi-compactas (I, 2.7), por tanto de presentación finita (I, 6.2), y entonces los morfismos $X'(s_i)\to S$ son de presentación finita (I, 6.2), y también lo es, por consiguiente, el morfismo $S'\to S$.

\begin{corollary}[17.16.3]
\label{IV.17.16.3}
\begin{enumerate}
  \item[(i)] Sea $f:X\to S$ un morfismo liso. Sean $s\in S$, $x$ un punto cerrado de $X_s$ tal que el cuerpo residual $k(x)$ sea separable sobre $k(s)$; entonces, en la conclusión de (17.16.1), se puede tomar $X'$ de modo que el morfismo $X'\to U$, restricción de $f$, sea étale.
  \item[(ii)] Sea $f:X\to S$ un morfismo liso sobreyectivo. Entonces, en las conclusiones de (17.16.2), se puede suponer además que $g:S'\to S$ sea étale.
\end{enumerate}
\end{corollary}

Es claro que (ii) se deduce de (i) del mismo modo que (17.16.2) de (17.16.1), teniendo en cuenta que, para todo $s\in S$, como $X_s$ no es vacío y es liso sobre $k(s)$, existe un punto cerrado $x\in X_s$ tal que $k(x)$ sea separable sobre $k(s)$ (17.15.10, (iii)). Basta, pues, demostrar (i). Se observará que el anillo $\sh{O}_{X_s,x}$ es entonces regular; si se repite la construcción hecha en (17.16.1) tomando para $(t_i)$ un sistema regular de parámetros de $\sh{O}_{X_s,x}$, el anillo $\sh{O}_{X'_s,x}$ es un cuerpo isomorfo a $k(x)$, por tanto separable sobre $k(s)$ por hipótesis. La conclusión resulta entonces de (17.6.1, $c'$)).

\begin{proposition}[17.16.4]
\label{IV.17.16.4}
Sean $S$ un preesquema cuasi-compacto y cuasi-separado, $f:X\to S$ un morfismo sobreyectivo localmente de presentación finita. Entonces existe una familia finita $(S_i)_{i\in I}$

\oldpage[IV-4]{107}

de subpreesquemas afines de $S$, de presentación finita sobre $S$, dos a dos disjuntos, cuya unión es $S$, y que tiene la propiedad siguiente: para cada $i\in I$, existen dos morfismos finitos, de presentación finita y sobreyectivos
\[
  S_i''\xrightarrow{g_i}S_i'\xrightarrow{h_i}S_i,
\]
donde $h_i$ es étale y $g_i$ es plano y radicial, y un $S$-morfismo $S_i''\to X$ (en otras palabras, una $S_i''$-sección de $X_i''=X\times_S S_i''$).
\end{proposition}

Hay un recubrimiento finito $(U_i)_{1\leq i\leq n}$ de $S$ por abiertos afines; para todo $i$, la intersección
\[
  W_i=U_i\cap\left(\bigcup_{j<i}U_j\right)
\]
es cuasi-compacta (I, 2.7), por lo que existe un subpreesquema cerrado $V_i$ de $U_i$ cuyo espacio subyacente es $U_i-W_i$ y que está definido por un ideal de tipo finito del anillo de $U_i$; por tanto, $V_i$ es un subpreesquema afín de $S$ que es de presentación finita sobre $S$ (I, 6.2). Es claro que podemos limitarnos a demostrar la proposición reemplazando $S$ por $V_i$ y $X$ por $X\times_S V_i$. En otras palabras, ya podemos suponer que $S$ es afín. Por otra parte, $X$ es unión de abiertos afines $W_\alpha$ y los $f(W_\alpha)$ son constructibles en $S$ (I, 8.4) y forman un recubrimiento de $S$; por consiguiente (I, 9.9), existe una subfamilia finita $(W_{\alpha_j})_{1\leq j\leq m}$ tal que los $f(W_{\alpha_j})$ ya forman un recubrimiento de $S$. Sea entonces $X'$ el preesquema suma de los subpreesquemas inducidos sobre los abiertos $W_{\alpha_j}$ de $X$; es inmediato que basta demostrar la proposición reemplazando $X$ por $X'$, puesto que un $S$-morfismo $S_i''\to X'$ da por composición un $S$-morfismo $S_i''\to X'\to X$.

Podemos, pues, suponer $X$ afín y $f$ de presentación finita. Se puede entonces (8.9.1) escribir $X$ bajo la forma $X_0\times_{S_0}S$, donde $S_0$ es noetheriano, $X_0\to S_0$ un morfismo de tipo finito, que podemos suponer sobreyectivo (8.10.5, (vi)). Si se prueba la proposición para este morfismo, resultará inmediatamente para $X$ por cambio de base $S\to S_0$. Podemos, pues, suponer además que $S$ es noetheriano y que $f$ es de tipo finito.

Sean $s_h$ $(1\leq h\leq r)$ los puntos maximales de $S$. Demostremos que basta probar el enunciado reemplazando $S$ por una vecindad abierta suficientemente pequeña $U_h$ de cada uno de los $s_h$, y $X$ por la imagen inversa de $U_h$ en $X$. En efecto, supongamos establecida la proposición en este caso, y razonemos por recurrencia noetheriana ($0_{\mathrm I}$, 2.2.2), suponiendo establecido el enunciado para todo subpreesquema cerrado de $S$ cuyo espacio subyacente sea $\neq S$. Podemos suponer que los $U_h$ son dos a dos disjuntos. Si $T$ es un subpreesquema cerrado cuyo espacio subyacente es $S-(\bigcup_h U_h)$, la hipótesis de recurrencia entraña que el enunciado es cierto para $T$; como también es cierto para cada uno de los $U_h$, es evidentemente cierto para $S$.

Como evidentemente podemos (reemplazando $S$ por $S_{\mathrm{red}}$ y $X$ por $X\times_S S_{\mathrm{red}}$) suponer que $S$ es reducido, cada uno de los $\sh{O}_{S,s_h}$ es un cuerpo $k(s_h)$. Supongamos que se ha demostrado la proposición cuando $S$ es el espectro de un cuerpo y $f$ es de tipo finito. Entonces la existencia de los $U_h$ resulta del método de (8.1.2, $a$)) y de (8.8.2, (i) y (ii)), (8.10.5, (vi), (vii) y (x)), (11.2.6, (ii)) y (17.7.8, (ii)).

Supongamos, pues, $S=\Spec(k)$, donde $k$ es un cuerpo, siendo $X$ de tipo finito sobre $k$ y $X\neq\varnothing$. Como $X$ es noetheriano, existe en $X$ un punto cerrado $x$ ($0_{\mathrm I}$, 2.1.3), por lo que $k(x)$ es una extensión finita de $k$ (I, 6.4.2). Existe, por tanto, una extensión finita separable $k'$ de $k$ tal que $k''=k(x)$ sea una extensión finita radicial de $k'$. Respondemos entonces a la cuestión tomando $I$ reducido a un elemento, $S_i'=\Spec(k')$, $S_i''=\Spec(k'')$.

\oldpage[IV-4]{108}

\begin{corollary}[17.16.5]
\label{IV.17.16.5}
Sea $f:X\to S$ un morfismo sobreyectivo localmente de presentación finita. Entonces existe un morfismo $g:S'\to S$ sobreyectivo, localmente de presentación finita y localmente cuasi-finito, tal que exista un $S$-morfismo $S'\to X$ (en otras palabras, que exista una $S'$-sección de $X'=X\times_S S'$). Si $S$ es cuasi-compacto (resp. cuasi-compacto y cuasi-separado), se puede suponer que $S'$ es afín (resp. $S'$ afín y el morfismo $g$ de presentación finita y cuasi-finito).
\end{corollary}

Es suficiente demostrar el corolario suponiendo que $S$ es afín: en efecto, $S$ es unión de una familia $(S_\alpha)$ de abiertos afines, y si para cada $\alpha$, $S'_\alpha$ es afín y el morfismo $g_\alpha:S'_\alpha\to S_\alpha$ responde a la cuestión y es de presentación finita y cuasi-finito, tomando para $S'$ el preesquema suma de los $S'_\alpha$, con $g:S'\to S$ coincidiendo con $g_\alpha$ en cada uno de los $S'_\alpha$, este morfismo responde a la cuestión; además, si $S$ es cuasi-compacto, podemos suponer finita la familia $(S_\alpha)$, por lo que $S'$ es afín; si además $S$ es cuasi-separado, las inmersiones $S_\alpha\to S$ son de presentación finita (I, 6.2), por lo que también lo es $g$.

Consideremos, pues, el caso en que $S$ es afín: formamos entonces los morfismos finitos y de presentación finita $S_i''\to S_i$ de (17.16.4); como las inmersiones $S_i\to S$ son de presentación finita (siendo los $S_i$ afines), los morfismos compuestos $g_i'':S_i''\to S_i\to S$ son de presentación finita y cuasi-finitos, y respondemos a la cuestión tomando $S'$ igual a la suma de los $S_i''$ y $g$ coincidiendo con $g_i''$ en cada uno de los $S_i''$.

\begin{proposition}[17.16.6]
\label{IV.17.16.6}
Sean $S$ un preesquema cuasi-compacto y cuasi-separado, $f:X\to S$ un morfismo de presentación finita tal que, para todo $s\in S$, $X_s=f^{-1}(s)$ sea un preesquema propio sobre $k(s)$. Entonces existe una familia finita $(S_i)_{i\in I}$ de subpreesquemas afines de $S$, de presentación finita sobre $S$, dos a dos disjuntos, cuya unión es $S$, y tal que, para todo $i$, el morfismo $X_i=X\times_S S_i\to S_i$ sea propio y plano. Si además, para todo $s\in S$, $X_s$ es finito sobre $k(s)$, se pueden tomar los $S_i$ de modo que cada uno de los morfismos $X_i\to S_i$ se factorice como
\[
  X_i\xrightarrow{u_i}X_i'\xrightarrow{h_i}S_i
\]
donde $u_i$ y $h_i$ son finitos y localmente libres (18.2.7), $h_i$ es étale, y $u_i$ es radicial y sobreyectivo.
\end{proposition}

Seguimos un procedimiento análogo al de (17.16.4). Nos reducimos primero al caso en que $S$ es afín y $f$ es plano, utilizando el teorema de planitud genérica (8.9.5) (obsérvese que los subpreesquemas $S_i'$ definidos en la demostración de (8.9.5) son afines). Puesto que $f$ es de presentación finita, podemos escribir a continuación $X=X_0\times_{S_0}S$, donde $S_0$ es noetheriano y $X_0\to S_0$ un morfismo de tipo finito y plano (11.2.7); además, cada fibra $(X_0)_{s_0}$ sigue siendo propia sobre $k(s_0)$, como resulta de (2.7.1, (vii)), y nos reducimos, por tanto, al caso en que $S$ es noetheriano. Por recurrencia noetheriana y aplicación del procedimiento de (8.1.2, $a$)) (utilizando también esta vez (8.10.5, (xii))), para demostrar la primera afirmación nos reducimos finalmente al caso en que $S$ es el espectro de un cuerpo, que resulta trivialmente de la hipótesis (con $S_i=S$). Se trata de la misma manera el caso en que se supone que $X_s$ es finito sobre $k(s)$ para todo $s\in S$ (esta vez hay que utilizar (2.7.1, (xv)), (8.10.5, (x), (iv), (vi) y (vii)), (17.7.8) y (2.1.12)), y nos reducimos a demostrar la última afirmación del enunciado cuando $S=\Spec(k)$ es el espectro de un cuerpo.

Pero entonces $X=\Spec(A)$, donde $A$ es una $k$-álgebra finita (I, 6.4.4); por tanto, $A$ es un producto directo de $k$-álgebras finitas $A_j$ que son anillos locales $(1\leq j\leq m)$;\footnote{La transcripción francesa dice «composée directe»; aquí se expresa el producto directo indicado por el testigo inglés y por la descomposición de una álgebra artiniana.} puesto que $A_j$ es un anillo artiniano y una

\oldpage[IV-4]{109}

$k$-álgebra, contiene un subcuerpo $K'_j$ isomorfo canónicamente al cuerpo residual de $A_j$ (0, 19.6.2). Tomamos entonces $X_i=X$, $X'_i$ suma de los $\Spec(K'_j)$ y es claro que respondemos así a la cuestión, pues para todo $j$ tenemos dos homomorfismos $K'_j\to A_j\to k(A_j)$ cuyo compuesto es un isomorfismo.


% Mechanical subsection components; prose is translated in the component files.
\section[Complementos sobre los morfismos étale. Anillos locales henselianos y anillos estrictamente locales]{Complementos sobre los morfismos étale\\Anillos locales henselianos y anillos estrictamente locales}
\label{section:IV.18}

En esta sección estudiamos diversas propiedades específicas de los morfismos étale. Además, la noción de morfismo étale permite desarrollar de manera muy natural la teoría de Nagata de los anillos henselianos, así como la de los anillos estrictamente locales. Estos anillos desempeñan un papel importante en numerosos desarrollos recientes, pues aparecen cada vez que se necesita un procedimiento de «localización» más fino que el proporcionado por la topología de Zariski (cf., por ejemplo, [43], mientras se espera la publicación del capítulo de este Tratado dedicado al estudio de la «topología étale»).

\subsection{Una equivalencia notable de categorías.}
\label{subsection:IV.18.1}

\begin{proposition}[18.1.1]
\label{IV.18.1.1}
Sean $S$ un preesquema, $S_0$ un subpreesquema cerrado de $S$, $X_0$ un $S_0$-preesquema liso (resp. étale) sobre $S_0$, y $x_0$ un punto de $X_0$. Existen entonces un entorno abierto $U_0$ de $x_0$ en $X_0$, un $S$-preesquema $U$ liso (resp. étale) sobre $S$ y un $S_0$-isomorfismo
\[
  U\times_S S_0\xrightarrow{\sim}U_0.
\]
\end{proposition}

Observemos que si $X_0$ es étale sobre $S_0$ en el punto $x_0$, es con mayor razón no ramificado sobre $S_0$ en dicho punto; si se ha construido un $S$-preesquema $U$ liso sobre $S$ tal que $U\times_S S_0$ sea isomorfo a $U_0$, como las fibras de los morfismos $U_0\to S_0$ y $U\to S$ que contienen a $x_0$ son entonces isomorfas, se seguirá que $U$ es no ramificado sobre $S$ en el punto $x_0$ (17.4.1, $d$)), y por tanto también en un entorno de $x_0$; sustituyendo $U$ por este entorno, se concluye que $U$ será étale sobre $S$. Basta, pues, demostrar la proposición cuando solo se supone que $X_0$ es liso sobre $S_0$.

Como la cuestión es local sobre $S$ y sobre $X_0$, se puede suponer que $S=\Spec(A)$ y $X_0=\Spec(C_0)$ son afines, de modo que $S_0=\Spec(A_0)$, donde $A_0$ es un anillo cociente de $A$, $C_0=B_0/\mathfrak J_0$, donde $B_0=A_0[T_1,\ldots,T_n]$ y $\mathfrak J_0$ es un ideal de tipo finito de $B_0$; por último, $C_0$ es una $A_0$-álgebra formalmente lisa para las topologías discretas. Sea $\mathfrak p_0$ el ideal $\mathfrak j_{x_0}$ en $C_0$; se tiene $\mathfrak p_0=\mathfrak q_0/\mathfrak J_0$, donde $\mathfrak q_0$ es un ideal primo de $B_0$. El criterio jacobiano (0, 22.6.4), junto con (0, 19.1.12), muestra que existen en $\mathfrak J_0$ una familia de $r$ polinomios $u_i$ $(1\leq i\leq r)$ y $r$ índices $j_h$ $(1\leq h\leq r)$ tales que las imágenes de los $u_i$ en $(\mathfrak J_0)_{\mathfrak q_0}/(\mathfrak J_0^2)_{\mathfrak q_0}$ generan este $(B_0)_{\mathfrak q_0}$-módulo y se tiene
\begin{equation}
  \det\left(\frac{\partial u_i}{\partial T_{j_h}}\right)\notin\mathfrak q_0.
  \tag{18.1.1.1}\label{IV.18.1.1.1}
\end{equation}

Como $(B_0)_{\mathfrak q_0}$ es un anillo local, se sigue del lema de Nakayama que se puede suponer que las imágenes de los $u_i$ en $(\mathfrak J_0)_{\mathfrak q_0}$ generan este $(B_0)_{\mathfrak q_0}$-módulo y, después, sustituyendo si es preciso $X_0$ por un entorno abierto afín de $x_0$, que los $u_i$ generan $\mathfrak J_0$ ($0_{\mathrm I}$, 5.2.2). Pongamos entonces $B=A[T_1,\ldots,T_n]$; $B_0$ es, por tanto, un anillo cociente

\oldpage[IV-4]{110}

de $B$, $\mathfrak q_0$ es la imagen de un ideal primo $\mathfrak q$ de $B$ y $\mathfrak q$ es la imagen inversa de $\mathfrak q_0$. Para todo $i$, sea $v_i\in B$ un elemento del que $u_i$ es la imagen en $B_0$, y sea $\mathfrak J$ el ideal de $B$ generado por los $v_i$, de modo que $\mathfrak J_0$ es la imagen de $\mathfrak J$ en $B_0$. La proposición quedará demostrada tomando como $U$ un entorno abierto del punto de $\Spec(B/\mathfrak J)$ correspondiente al ideal primo $\mathfrak p=\mathfrak q/\mathfrak J$, con tal de que se demuestre que $B_{\mathfrak q}/\mathfrak J_{\mathfrak q}$ es una $A$-álgebra formalmente lisa para las topologías discretas. Ahora bien, esto se sigue del criterio jacobiano, pues las imágenes de los $v_i$ en $\mathfrak J_{\mathfrak q}/\mathfrak J_{\mathfrak q}^2$ generan este $B_{\mathfrak q}$-módulo, y de (18.1.1.1) se sigue que se tiene $\det(\partial v_i/\partial T_{j_h})\notin\mathfrak q$.

\begin{theorem}[18.1.2]
\label{IV.18.1.2}
Sean $S$ un preesquema y $S_0$ un subpreesquema cerrado de $S$ cuyo espacio subyacente es idéntico al de $S$. Entonces el funtor
\[
  X\longmapsto X\times_S S_0
\]
de la categoría de los $S$-preesquemas étales sobre $S$ a la categoría de los $S_0$-preesquemas étales sobre $S_0$ es una equivalencia de categorías.
\end{theorem}

Demostremos primero que este funtor es plenamente fiel. Sean $X$, $Y$ dos $S$-preesquemas étales sobre $S$, y pongamos $X_0=X\times_S S_0$, $Y_0=Y\times_S S_0$. Si $Z=X\times_S Y$, el conjunto $\Hom_S(X,Y)$ está en correspondencia biyectiva canónica con el conjunto de las $X$-secciones $\Gamma(Z/X)$, y del mismo modo $\Hom_{S_0}(X_0,Y_0)$ está en correspondencia biyectiva canónica con $\Gamma(Z_0/X_0)$, donde $Z_0=Z\times_S S_0=X_0\times_{S_0}Y_0$. Ahora bien, $Z$ es étale sobre $X$, $Z_0$ es étale sobre $X_0$, y $X_0$ (resp. $Z_0$) es un subpreesquema cerrado de $X$ (resp. $Z$) con el mismo espacio subyacente. Los abiertos de $Z$ tales que la restricción del morfismo $Z\to X$ sea sobreyectiva y radicial son, por tanto, los mismos que los abiertos de $Z_0$ que tienen las propiedades correspondientes, y nuestra afirmación se sigue así de (17.9.3).

Para terminar la demostración, basta ver que para todo $S_0$-preesquema $X_0$ étale sobre $S_0$, existe un $S$-preesquema $X$ étale sobre $S$ y un $S_0$-isomorfismo $X_0\xrightarrow{\sim}X\times_S S_0$. En virtud de la prop. (18.1.1), existe un recubrimiento abierto $(U_\alpha)$ de $X_0$ y, para todo $\alpha$, un $S$-preesquema $V_\alpha$ que es étale sobre $S$, y finalmente un $S_0$-isomorfismo $\theta_\alpha:U_\alpha\xrightarrow{\sim}V_\alpha\times_S S_0$. Además, en virtud de la primera parte de la demostración, existe un único $S_0$-isomorfismo $\varphi_{\alpha\beta}$ de $\theta_\alpha(U_\alpha\cap U_\beta)$ sobre $\theta_\beta(U_\alpha\cap U_\beta)$, correspondiente al automorfismo identidad de $U_\alpha\cap U_\beta$, y es inmediato, por la misma razón, que estos isomorfismos satisfacen la condición de pegado ($0_{\mathrm I}$, 4.1.7). Existe, por tanto, un $S$-preesquema $X$ tal que los $V_\alpha$ se identifican canónicamente con subpreesquemas inducidos sobre abiertos de $X$, y los $\theta_\alpha$ se identifican con $S_0$-isomorfismos que coinciden en las intersecciones $U_\alpha\cap U_\beta$ y definen, por consiguiente, un $S_0$-isomorfismo $X_0\xrightarrow{\sim}X\times_S S_0$. Es claro que $X$ es étale sobre $S$ (17.3.2), lo que termina la demostración.

\begin{corollary}[18.1.3]
\label{IV.18.1.3}
Sean $S$ un preesquema, $X$ un $S$-preesquema étale sobre $S$, $S'$ un $S$-preesquema, y $S'_0$ un subpreesquema cerrado de $S'$ con el mismo espacio subyacente. Entonces la aplicación canónica $X(S')_S\to X(S'_0)_S$ (I, 3.4.3) es biyectiva.
\end{corollary}

En efecto, pongamos $X'=X\times_S S'$, $X'_0=X\times_S S'_0$, de modo que $X(S')_S=\Gamma(X'/S')$ y $X(S'_0)_S=\Gamma(X'_0/S'_0)$; el corolario se sigue del hecho de que el funtor definido en (18.1.2) (sustituyendo $S$ y $S_0$ por $S'$ y $S'_0$, respectivamente) es plenamente fiel (o directamente de (17.9.3)).

\subsection{Recubrimientos étales.}
\label{subsection:IV.18.2}

\oldpage[IV-4]{111}

\begin{env}[18.2.1]
\label{IV.18.2.1}
Dados un anillo $A$ y una $A$-álgebra conmutativa $B$ que es \emph{finita} y es un $A$-módulo \emph{libre}, recordemos (Bourbaki, \emph{Alg.}, cap. VIII, \textsection~12, n.\textsuperscript{o}~2) que se define sobre $B$ una \emph{forma $A$-lineal} $\operatorname{Tr}_{B/A}$, la «forma traza»; de ella se deduce la definición de una \emph{forma $A$-bilineal simétrica} (también llamada «forma traza»)
\begin{equation}
  (x,y)\longmapsto\operatorname{Tr}_{B/A}(xy)
  \tag{18.2.1.1}\label{IV.18.2.1.1}
\end{equation}
cuyo dato equivale al de la aplicación $A$-lineal asociada $\operatorname{astr}_{B/A}:B\to\check B$ del $A$-módulo $B$ a su \emph{dual} $\check B$, igual a su traspuesta. Cuando $A$ es un \emph{cuerpo}, es equivalente decir que esta forma bilineal es \emph{no degenerada} o que $B$ es una $A$-álgebra \emph{separable} (Bourbaki, \emph{Alg.}, cap. IX, \textsection~2, prop. 5).

Sea $f:A\to A'$ un homomorfismo de anillos; pongamos $B'=B\otimes_A A'$, y sea $g:B\to B'$ el homomorfismo canónico; la imagen por $g$ de una base del $A$-módulo $B$ es entonces una base del $A'$-módulo $B'$, y se sigue de las definiciones que se tiene, para todo $x\in B$,
\begin{equation}
  \operatorname{Tr}_{B'/A'}(g(x))=f(\operatorname{Tr}_{B/A}(x)).
  \tag{18.2.1.2}\label{IV.18.2.1.2}
\end{equation}
\end{env}

\begin{env}[18.2.2]
\label{IV.18.2.2}
Consideremos ahora un \emph{espacio anillado} $(X,\sh{O}_X)$ y sea $\sh{B}$ una $\sh{O}_X$-álgebra que, como $\sh{O}_X$-módulo, es \emph{localmente libre de rango finito}; entonces, para todo abierto $U\subset X$ tal que $\sh{B}|U$ sea (como $\sh{O}_U$-módulo) isomorfa a $\sh{O}_U^n$ (para un $n$ que depende de $U$), $\Gamma(U,\sh{B})$ es una $\Gamma(U,\sh{O}_X)$-álgebra que, como $\Gamma(U,\sh{O}_X)$-módulo, es libre de rango finito y define, por tanto, una forma $\Gamma(U,\sh{O}_X)$-lineal $\operatorname{Tr}_{\Gamma(U,\sh{B})/\Gamma(U,\sh{O}_X)}$, que denotaremos también por $\operatorname{Tr}_{\sh{B}|\sh{O}_X,U}$; de ella se deduce una aplicación lineal asociada
\[
  \operatorname{astr}_{\sh{B}|\sh{O}_X,U}:\Gamma(U,\sh{B})\longrightarrow\Gamma(U,\sh{B})^\vee=\Gamma(U,\check{\sh{B}}).
\]
Además, se sigue de (18.2.1.2) que estas aplicaciones lineales son compatibles con las operaciones de restricción de $U$ a un abierto menor y definen, por una parte, un \emph{homomorfismo de $\sh{O}_X$-módulos}, llamado también \emph{homomorfismo traza}:
\begin{equation}
  \operatorname{Tr}_{\sh{B}|\sh{O}_X}:\sh{B}\longrightarrow\sh{O}_X
  \tag{18.2.2.1}\label{IV.18.2.2.1}
\end{equation}
y, por otra parte, un \emph{homomorfismo de $\sh{O}_X$-módulos}
\begin{equation}
  \operatorname{astr}_{\sh{B}|\sh{O}_X}:\sh{B}\longrightarrow\check{\sh{B}}
  \tag{18.2.2.2}\label{IV.18.2.2.2}
\end{equation}
llamado \emph{asociado a la traza}, e igual a su traspuesto. Se sigue también de (18.2.1.2) que, para todo $x\in X$, se tiene
\begin{equation}
  (\operatorname{Tr}_{\sh{B}|\sh{O}_X})_x=\operatorname{Tr}_{\sh{B}_x|\sh{O}_{X,x}}
  \tag{18.2.2.3}\label{IV.18.2.2.3}
\end{equation}
\begin{equation}
  (\operatorname{astr}_{\sh{B}|\sh{O}_X})_x=\operatorname{astr}_{\sh{B}_x|\sh{O}_{X,x}}.
  \tag{18.2.2.4}\label{IV.18.2.2.4}
\end{equation}
Finalmente, bajo las condiciones de (18.2.1), si se pone $X=\Spec(A)$ y si $\sh{B}=\widetilde B$ es la $\sh{O}_X$-álgebra correspondiente a $B$, la forma $\operatorname{Tr}_{\sh{B}|\sh{O}_X}$ (resp. el homomorfismo $\operatorname{astr}_{\sh{B}|\sh{O}_X}$) corresponde a la forma $\operatorname{Tr}_{B/A}$ (resp. al homomorfismo de $A$-módulos $\operatorname{astr}_{B/A}$), como se sigue igualmente de (18.2.1.2).
\end{env}


\oldpage[IV-4]{112}

\begin{proposition}[18.2.3]
\label{IV.18.2.3}
Sea $f:X\to Y$ un morfismo finito de preesquemas y sea $\sh{B}=f_*(\sh{O}_X)$. Las condiciones siguientes son equivalentes:
\begin{enumerate}
  \item[a)] $f$ es étale.
  \item[a$'$)] $f$ es un morfismo plano de presentación finita, y para todo $x\in X$, si se pone $y=f(x)$, $\sh{O}_{X,x}/\mathfrak m_y\sh{O}_{X,x}$ es un cuerpo, extensión finita separable de $\kappa(y)$.
  \item[b)] $\sh{B}$ es un $\sh{O}_Y$-módulo localmente libre y, para todo $y\in Y$, $\sh{B}_y\otimes_{\sh{O}_{Y,y}}\kappa(y)$ es una $\kappa(y)$-álgebra finita separable (por tanto, el producto directo de un número finito de cuerpos, extensiones finitas separables de $\kappa(y)$).
  \item[c)] $\sh{B}$ es un $\sh{O}_Y$-módulo localmente libre, y el homomorfismo $\operatorname{astr}_{\sh{B}|\sh{O}_Y}:\sh{B}\to\check{\sh{B}}$ (18.2.2) es biyectivo.
\end{enumerate}
\end{proposition}

Habida cuenta de que $f$ es cuasicompacto, la equivalencia de $a)$ y $a')$ ya ha sido demostrada (17.6.2). Para demostrar el resto de la proposición, se puede reducir al caso en que $Y=\Spec(A)$ y $X=\Spec(B)$ son afines, siendo $B$ una $A$-álgebra finita y $\sh{B}=\widetilde B$. Decir que $f$ es un morfismo de presentación finita equivale entonces a decir que $B$ es un $A$-módulo de presentación finita (1.4.7). Si además $f$ es plano, y por tanto $B$ un $A$-módulo plano, se sabe (Bourbaki, \emph{Alg. comm.}, cap. II, \textsection~5, n.\textsuperscript{o}~2, cor. 2 del th. 1) que $B$ es un $A$-módulo proyectivo, por lo que $\sh{B}$ es un $\sh{O}_Y$-módulo localmente libre (\emph{loc. cit.}, n.\textsuperscript{o}~2, th. 1), y la recíproca es inmediata. Por otra parte, $f^{-1}(y)$ no es sino el espectro de la $\kappa(y)$-álgebra $\sh{B}(y)=\sh{B}_y\otimes_{\sh{O}_{Y,y}}\kappa(y)$, lo que termina de demostrar la equivalencia de $a')$ y $b)$. Para ver que $b)$ equivale a $c)$, observemos que la segunda afirmación de $b)$ equivale a que el homomorfismo $\operatorname{astr}_{\sh{B}(y)|\kappa(y)}:\sh{B}(y)\to\check{\sh{B}(y)}$ sea biyectivo; como $\sh{B}(y)=\sh{B}_y/\mathfrak m_y\sh{B}_y$ y $\check{\sh{B}(y)}=\check{\sh{B}}_y/\mathfrak m_y\check{\sh{B}}_y$ y $\sh{B}_y$ y $\check{\sh{B}}_y$ son $\sh{O}_{Y,y}$-módulos libres, se sigue de (18.2.2.4) y de Bourbaki, \emph{Alg. comm.}, cap. II, \textsection~3, n.\textsuperscript{o}~3, cor. de la prop. 6, que el homomorfismo $\operatorname{astr}_{\sh{B}_y|\sh{O}_{Y,y}}:\sh{B}_y\to\check{\sh{B}}_y$ es también biyectivo; siendo evidente la recíproca, esto termina la demostración.

Cuando una $\sh{O}_X$-álgebra $\sh{B}$ satisface las condiciones equivalentes $b)$ y $c)$ de (18.2.3), se dice que $\sh{B}$ es una $\sh{O}_X$-álgebra \emph{finita y étale}. Cuando $X=\Spec(A)$ es afín y se tiene, por tanto, $\sh{B}=\widetilde B$, donde $B$ es una $A$-álgebra, en virtud de (18.2.3) equivale decir que $\sh{B}$ es una $\sh{O}_X$-álgebra finita y étale o que $B$ es una $A$-álgebra finita y étale (en el sentido de (17.3.2)).

\begin{corollary}[18.2.4]
\label{IV.18.2.4}
Sea $f:X\to Y$ un morfismo finito y de presentación finita, y pongamos $\sh{B}=f_*(\sh{O}_X)$. Sea $y$ un punto de $Y$. Las condiciones siguientes son equivalentes:
\begin{enumerate}
  \item[a)] Existe un entorno abierto $U$ de $y$ en $Y$ tal que la restricción $f^{-1}(U)\to U$ de $f$ sea un morfismo étale.
  \item[b)] $\sh{B}_y$ es un $\sh{O}_{Y,y}$-módulo libre de tipo finito y $\sh{B}_y\otimes_{\sh{O}_{Y,y}}\kappa(y)$ es una $\kappa(y)$-álgebra separable.
\end{enumerate}
\end{corollary}

Es claro que $a)$ implica $b)$ en virtud de (18.2.3). Por otra parte, $\sh{B}$ es un $\sh{O}_Y$-módulo de presentación finita (1.6.3 y 1.4.7), de modo que, si $\sh{B}_y$ es un $\sh{O}_{Y,y}$-módulo libre, existe un entorno abierto $U$ de $y$ en $Y$ tal que $\sh{B}|U$ sea un $\sh{O}_U$-módulo localmente libre ($0_{\mathrm I}$, 5.2.7); además, por hipótesis, como el homomorfismo $\operatorname{astr}_{\sh{B}_y|\sh{O}_{Y,y}}:\sh{B}_y\to\check{\sh{B}}_y$ es biyectivo, se sigue también de ($0_{\mathrm I}$, 5.2.7) que se puede suponer elegido $U$ de modo que

\oldpage[IV-4]{113}

el homomorfismo $\operatorname{astr}_{\sh{B}|U/\sh{O}_U}$ sea biyectivo. El hecho de que $b)$ implique $a)$ se sigue entonces de (18.2.3).

\begin{corollary}[18.2.5]
\label{IV.18.2.5}
Sean $Y$ un preesquema cuasicompacto o localmente noetheriano y $f:X\to Y$ un morfismo finito y de presentación finita; pongamos $\sh{B}=f_*(\sh{O}_X)$. Supongamos que, para todo punto cerrado $y$ de $Y$, $\sh{B}_y$ sea un $\sh{O}_{Y,y}$-módulo libre y $\sh{B}_y\otimes_{\sh{O}_{Y,y}}\kappa(y)$ una $\kappa(y)$-álgebra separable. Entonces $f$ es étale.
\end{corollary}

En efecto, se sigue de (18.2.4) que todo punto cerrado $y$ de $Y$ posee un entorno abierto $U$ tal que la restricción $f^{-1}(U)\to U$ de $f$ sea étale; la conclusión se sigue de que, en los dos casos considerados, todo subconjunto cerrado no vacío de $Y$ contiene un punto cerrado (5.1.11 y $0_{\mathrm I}$, 2.1.3).

\begin{corollary}[18.2.6]
\label{IV.18.2.6}
Si $f:X\to Y$ es un morfismo finito y étale, y se pone $\sh{B}=f_*(\sh{O}_X)$, entonces el $\sh{O}_Y$-homomorfismo $\operatorname{Tr}_{\sh{B}|\sh{O}_Y}:\sh{B}\to\sh{O}_Y$ (18.2.2) (denotado también por $\operatorname{Tr}_f$) es sobreyectivo.
\end{corollary}

Como la cuestión es local, en virtud de (18.2.3) se puede suponer que $Y=\Spec(A)$, $X=\Spec(B)$ con $\sh{B}=\widetilde B$, siendo $B$ un $A$-módulo libre; como en virtud de (18.2.3) la forma bilineal (18.2.1.1) es no degenerada, esto implica, en particular, que la forma lineal $\operatorname{Tr}_{B/A}$ es sobreyectiva.

\begin{remarks}[18.2.7]
\label{IV.18.2.7}
\begin{enumerate}
  \item[(i)] Cuando $f:X\to Y$ es un morfismo finito tal que $f_*(\sh{O}_X)$ sea un $\sh{O}_Y$-módulo localmente libre (resp. localmente libre de rango $n$), se dice también que $f$ es un morfismo \emph{finito localmente libre} (resp. \emph{localmente libre de rango $n$}). Esta condición, en virtud de (18.2.3), se satisface si $f$ es un morfismo finito étale, pero no implica por sí sola que $f$ sea étale, como muestra el ejemplo en que $X=\Spec(K)$ e $Y=\Spec(k)$ son espectros de cuerpos, siendo $K$ una extensión finita no separable de $k$. Cuando $f:X\to Y$ es un morfismo finito étale, se dice también que $X$ es un \emph{recubrimiento étale} de $Y$.

Obsérvese que, en este caso, $f$ es universalmente abierto y universalmente cerrado y, en particular, $f(X)$ es un subconjunto a la vez abierto y cerrado de $Y$.

Se dice que un recubrimiento étale $X$ de $Y$ es \emph{trivial} si $X$ es \emph{suma} de un número finito de preesquemas \emph{isomorfos} a $Y$. Se dice que un recubrimiento étale $X$ de $Y$ es \emph{localmente trivial} si el morfismo $f:X\to Y$ es tal que todo punto $y\in Y$ posee un entorno abierto $U$ para el cual el recubrimiento $f^{-1}(U)$ de $U$ es trivial.

  \item[(ii)] Sea $f:X\to Y$ un morfismo finito, localmente libre de rango $n$; pongamos $\sh{B}=f_*(\sh{O}_X)$ y sea $u=\operatorname{astr}_{\sh{B}|\sh{O}_Y}:\sh{B}\to\check{\sh{B}}$; de él se deduce un homomorfismo sobre la $n$-ésima potencia exterior
\[
  \bigwedge^n u:\bigwedge^n\sh{B}\longrightarrow\bigwedge^n\check{\sh{B}}=(\bigwedge^n\sh{B})^\vee
\]
de $\sh{O}_Y$-módulos \emph{invertibles} y, por consiguiente, ($0_{\mathrm I}$, 5.4.2), un elemento
\begin{equation}
  d_{X/Y}\in\Gamma\left(Y,(\bigwedge^n\check{\sh{B}})\otimes_{\sh{O}_Y}(\bigwedge^n\check{\sh{B}})\right)
  \tag{18.2.7.1}\label{IV.18.2.7.1}
\end{equation}
que se llama el \emph{discriminante de $X$ sobre $Y$}. Además, como $(\bigwedge^n\check{\sh{B}})\otimes_{\sh{O}_Y}(\bigwedge^n\check{\sh{B}})$ es el dual de $(\bigwedge^n\sh{B})\otimes_{\sh{O}_Y}(\bigwedge^n\sh{B})$, $d_{X/Y}$ puede identificarse también con un homomorfismo
\begin{equation}
  (\bigwedge^n\sh{B})\otimes_{\sh{O}_Y}(\bigwedge^n\sh{B})\longrightarrow\sh{O}_Y
  \tag{18.2.7.2}\label{IV.18.2.7.2}
\end{equation}
y se denota por $\mathscr{D}_{X/Y}$ el ideal cuasicoherente de tipo finito de $\sh{O}_Y$, \emph{imagen} del homomorfismo (18.2.7.2), llamado también \emph{ideal discriminante de $X$ sobre $Y$}.
\end{enumerate}
\end{remarks}

\oldpage[IV-4]{114}

Dicho esto, para que el homomorfismo $u$ sea biyectivo, es necesario y suficiente que $\bigwedge^n u$ sea biyectivo, o también que la sección $d_{X/Y}$ tenga un germen \emph{invertible} en todo punto $y\in Y$, lo que se escribe asimismo $d_{X/Y}(y)\neq 0$ para todo $y$ ($0_{\mathrm I}$, 5.5.2). Equivale también decir que el ideal discriminante $\mathscr{D}_{X/Y}$ es \emph{igual a $\sh{O}_Y$}.

La terminología de «recubrimiento» introducida en (18.2.7, (i)) se justifica por la proposición siguiente:

\begin{proposition}[18.2.8]
\label{IV.18.2.8}
Sea $f:X\to Y$ un morfismo étale, separado y de tipo finito y, para todo $y\in Y$, sea $n(y)$ el número geométrico de puntos de $f^{-1}(y)$. Entonces la función $y\mapsto n(y)$ es semicontinua inferiormente en $Y$. Para que sea continua en un punto $y$ (y, por tanto, constante en un entorno de $y$), es necesario y suficiente que exista un entorno abierto $U$ de $y$ tal que la restricción $f^{-1}(U)\to U$ de $f$ sea un morfismo finito (étale).
\end{proposition}

Como $f$ es cuasifinito (17.6.1) y localmente de presentación finita, decir que $f$ es finito equivale a decir que $f$ es propio (8.11.1); además, cada fibra $f^{-1}(y)$ es geométricamente reducida sobre $\kappa(y)$. Las conclusiones se siguen, por tanto, de (15.5.9, (i) y (ii)) y de que $f$ es plano.

\begin{corollary}[18.2.9]
\label{IV.18.2.9}
Sean $Y$ un preesquema conexo y $f:X\to Y$ un morfismo étale, separado y de tipo finito. Para que $f$ sea finito (en otras palabras, para que $X$ sea un recubrimiento étale de $Y$), es necesario y suficiente que todas las fibras de $f$ tengan el mismo número geométrico de puntos.
\end{corollary}

\begin{remarks}[18.2.10]
\label{IV.18.2.10}
\begin{enumerate}
  \item[(i)] El ejemplo de la «recta afín con un punto desdoblado» (I, 5.5.11) muestra que un morfismo étale de tipo finito de preesquemas noetherianos puede ser \emph{no separado}; para este ejemplo, la primera afirmación de (18.2.8) ya no es exacta.

  \item[(ii)] Para que un morfismo separado, de tipo finito y étale $f:X\to Y$ haga de $X$ un recubrimiento \emph{localmente trivial}, es necesario y suficiente que, para todo $x\in X$, exista un entorno abierto $U$ de $y=f(x)$ y una $U$-sección $g$ de $f^{-1}(U)$ tal que $g(y)=x$. En efecto, la condición es evidentemente necesaria; que sea suficiente se sigue de que toda fibra $f^{-1}(y)$ es finita (17.6.1) y (I, 6.2.2), de la caracterización de las secciones de un $Y$-preesquema étale (17.9.3) y de la prop. (18.2.8).
\end{enumerate}
\end{remarks}

\subsection{Álgebras finitas y étales.}
\label{subsection:IV.18.3}

\begin{proposition}[18.3.1]
\label{IV.18.3.1}
Sean $A$ un anillo y $B$ una $A$-álgebra de presentación finita.
\begin{enumerate}
  \item[(i)] Para que $B$ sea una $A$-álgebra no ramificada, es necesario y suficiente que $B$ sea un $(B\otimes_A B)$-módulo proyectivo.
  \item[(ii)] Supongamos además que $B$ es una $A$-álgebra finita. Para que $B$ sea una $A$-álgebra étale, es necesario y suficiente que $B$ sea un $A$-módulo proyectivo y un $(B\otimes_A B)$-módulo proyectivo.
\end{enumerate}
\end{proposition}

Por supuesto, la estructura de $(B\otimes_A B)$-módulo sobre $B$ es la que proviene de la estructura de $(B\otimes_A B)$-álgebra sobre $B$ correspondiente al $A$-homomorfismo canónico de anillos $B\otimes_A B\to B$, que es sobreyectivo y tiene núcleo $\mathfrak J_{B/A}$ (0, 20.4.1).

(i) Decir que el morfismo $\Spec(B)\to\Spec(A)$ es localmente de presentación finita equivale a decir que $B$ es una $A$-álgebra de presentación finita (1.4.6). Decir

\oldpage[IV-4]{115}

que $B$ es una $A$-álgebra no ramificada significa entonces (17.4.2) que $\Spec((B\otimes_A B)/\mathfrak J_{B/A})$ es un subesquema inducido sobre un subconjunto \emph{abierto y cerrado} de $\Spec(B\otimes_A B)$, y se sabe que, para que esto ocurra, es necesario y suficiente que $\mathfrak J_{B/A}$ sea un ideal \emph{factor directo} de $B\otimes_A B$ (Bourbaki, \emph{Alg. comm.}, cap. II, \textsection~4, n.\textsuperscript{o}~3, prop. 15); pero equivale a decir que el $(B\otimes_A B)$-módulo cociente $(B\otimes_A B)/\mathfrak J_{B/A}$ es proyectivo (Bourbaki, \emph{Alg.}, cap. II, 3.\textsuperscript{a} ed., \textsection~2, n.\textsuperscript{o}~2, prop. 4).

(ii) Si se recuerda que un $A$-módulo plano de presentación finita es proyectivo y recíprocamente (Bourbaki, \emph{Alg. comm.}, cap. II, \textsection~5, n.\textsuperscript{o}~2, cor. 2 del th. 1), la afirmación de (ii) se sigue de la de (i) y de (17.6.2).

\begin{proposition}[18.3.2]
\label{IV.18.3.2}
Sean $A$ un anillo y $\mathfrak J$ un ideal de $A$ tal que, para la topología $\mathfrak J$-preádica, $A$ sea separado y completo; pongamos $A_0=A/\mathfrak J$. Entonces el funtor
\[
  B\longmapsto B\otimes_A A_0
\]
es una equivalencia de la categoría de las $A$-álgebras finitas y étales con la categoría de las $A_0$-álgebras finitas y étales.
\end{proposition}

Demostraremos primero el lema siguiente:

\begin{lemma}[18.3.2.1]
\label{IV.18.3.2.1}
Sean $A$ un anillo y $\mathfrak J$ un ideal de $A$ tal que, para la topología $\mathfrak J$-preádica, $A$ sea separado y completo.
\begin{enumerate}
  \item[(i)] Todo $A$-módulo proyectivo de tipo finito $M$ es separado y completo para la topología $\mathfrak J$-preádica y, por tanto, límite proyectivo de los $(A/\mathfrak J^{n+1})$-módulos proyectivos $M/\mathfrak J^{n+1}M$.
  \item[(ii)] Recíprocamente, pongamos $A_n=A/\mathfrak J^{n+1}$, y sea $(M_n)$ un sistema proyectivo de $A_n$-módulos tal que, para todo $n$, el homomorfismo $M_{n+1}\otimes_{A_{n+1}}A_n\to M_n$ deducido del homomorfismo de transición $M_{n+1}\to M_n$ sea biyectivo. Supongamos además que los $M_n$ son proyectivos y que $M_0$ es de tipo finito. Entonces $M=\varprojlim M_n$ es un $A$-módulo proyectivo de tipo finito tal que el homomorfismo canónico $M\otimes_A A_0\to M_0$ es biyectivo.
\end{enumerate}
\end{lemma}

(i) Existe un $A$-módulo libre de tipo finito $L$ tal que $M$ es isomorfo a un sumando directo de $L$; como $L$ es separado para la topología $\mathfrak J$-preádica, lo mismo ocurre con todo submódulo $N$ de $L$, puesto que $\mathfrak J^{n+1}N\subset\mathfrak J^{n+1}L$; en particular, $M$ es separado para esta topología y, como el homomorfismo sobreyectivo $f:L\to M$ es continuo para la topología $\mathfrak J$-preádica, su núcleo $N$ es cerrado para la topología inducida por la de $L$; puesto que $L$ es completo y $f$ un morfismo estricto, se concluye que $M$ es completo (Bourbaki, \emph{Top. gén.}, cap. IX, 2.\textsuperscript{a} ed., \textsection~3, n.\textsuperscript{o}~1, prop. 4).

(ii) Se sigue del lema de Nakayama que, si $M_0$ está generado por una familia finita $(x_{i0})$ de $r$ elementos y si, para todo $n\geq 1$, $x_{in}$ es un elemento de $M_n$ del que $x_{i,n-1}$ es la imagen en $M_{n-1}$, entonces $(x_{in})$ $(1\leq i\leq r)$ es un sistema de generadores de $M_n$ (Bourbaki, \emph{Alg. comm.}, cap. II, \textsection~3, cor. 2 de la prop. 4). Dicho esto, para todo $n$, pongamos $L_n=A_n^r$; si $(e_{in})_{1\leq i\leq r}$ es la base canónica de $L_n$, sea $u_n:L_n\to M_n$ la aplicación $A_n$-lineal tal que $u_n(e_{in})=x_{in}$ para todo $i$. Por hipótesis, se tiene una sucesión exacta escindida
\[
  0\longrightarrow N_n\xrightarrow{v_n}L_n\xrightarrow{u_n}M_n\longrightarrow 0
\]

\oldpage[IV-4]{116}

y, como $L_n=L_{n+1}/\mathfrak J^{n+1}L_{n+1}$ y $M_n=M_{n+1}/\mathfrak J^{n+1}M_{n+1}$, las flechas verticales del diagrama conmutativo
\[
\xymatrix{
0\ar[r] & N_{n+1}\ar[r]^{v_{n+1}}\ar[d] & L_{n+1}\ar[r]^{u_{n+1}}\ar[d] & M_{n+1}\ar[r]\ar[d] & 0\\
0\ar[r] & N_n\ar[r]_{v_n} & L_n\ar[r]_{u_n} & M_n\ar[r] & 0
}
\]
son las tres sobreyectivas. Ahora bien, se tiene $M=\varprojlim M_n$ y $L=A^r=\varprojlim L_n$; si se pone $N=\varprojlim N_n$, $u=\varprojlim u_n$, $v=\varprojlim v_n$, se tiene, en virtud de ($0_{\mathrm{III}}$, 13.2.2), la sucesión exacta
\begin{equation}
  0\longrightarrow N\xrightarrow{v}L\xrightarrow{u}M\longrightarrow 0.
  \tag{18.3.2.2}\label{IV.18.3.2.2}
\end{equation}
Además, como para todo $n$, $v_n$ es invertible por la izquierda y $M_n$ es un $A_n$-módulo proyectivo, se sigue de (0, 19.1.8) que la sucesión exacta (18.3.2.2) es escindida, lo que demuestra el lema.

Dicho esto, demostremos primero que el funtor del enunciado de (18.3.2) es \emph{plenamente fiel}. Pongamos, como en el lema, $A_n=A/\mathfrak J^{n+1}$; sean $B$, $C$ dos $A$-álgebras finitas y étales, y pongamos, para todo $n$, $B_n=B\otimes_A A_n$, $C_n=C\otimes_A A_n$; en virtud de (18.3.1) y (18.3.2.1), $B$ y $C$ son separadas y completas para la topología $\mathfrak J$-preádica, y se tiene $B=\varprojlim B_n$, $C=\varprojlim C_n$; además, todo homomorfismo de $A$-álgebras $u:B\to C$ es continuo para las topologías $\mathfrak J$-preádicas y proporciona, por tanto, un sistema proyectivo de homomorfismos de $A_n$-álgebras $u_n=u\otimes 1:B_n\to C_n$, del que es el límite proyectivo; siendo evidente la recíproca, se tiene así una biyección canónica
\[
  \Hom_{A\text{-alg.}}(B,C)\xrightarrow{\sim}\varprojlim\Hom_{A_n\text{-alg.}}(B_n,C_n).
\]
Pero como $B$ y $C$ son $A$-álgebras étales, se sigue también de inmediato de (18.1.2) que la aplicación canónica
\[
  \Hom_{A_{n+1}\text{-alg.}}(B_{n+1},C_{n+1})\longrightarrow\Hom_{A_n\text{-alg.}}(B_n,C_n)
\]
es biyectiva para $n\geq 0$, lo que termina de demostrar que la aplicación canónica
\[
  \Hom_{A\text{-alg.}}(B,C)\longrightarrow\Hom_{A_0\text{-alg.}}(B_0,C_0)
\]
es biyectiva.

Para terminar de demostrar (18.3.2), basta ver que para toda $A_0$-álgebra finita y étale $B_0$, existe una $A$-álgebra finita y étale $B$ y un $A_0$-isomorfismo $B_0\xrightarrow{\sim}B\otimes_A A_0$. Ahora bien, se sigue de (18.1.2) que existe un sistema proyectivo $(B_n)$ tal que $B_n$ es una $A_n$-álgebra étale y los homomorfismos $B_{n+1}\otimes_{A_{n+1}}A_n\to B_n$ son biyectivos. Se sigue de (18.3.1) y (18.3.2.1) que la $A$-álgebra $B=\varprojlim B_n$ es un $A$-módulo proyectivo de tipo finito y que $B_0$ es isomorfa a $B\otimes_A A_0$. Para demostrar que $B$ es una $A$-álgebra étale, basta, en virtud de (18.2.5), mostrar que para todo ideal maximal $\mathfrak m$ de $A$, $B_{\mathfrak m}/\mathfrak mB_{\mathfrak m}$ es una $(A/\mathfrak m)$-álgebra separable. Ahora bien, como $\mathfrak J$ está contenido en el radical de $A$ ($0_{\mathrm I}$, 7.1.10),

\oldpage[IV-4]{117}

se tiene $\mathfrak J\subset\mathfrak m$ y, si $\mathfrak m_0=\mathfrak m/\mathfrak J$, se tiene $A_0/\mathfrak m_0=A/\mathfrak m$ y $B_{\mathfrak m}/\mathfrak mB_{\mathfrak m}=(B_0)_{\mathfrak m_0}/\mathfrak m_0(B_0)_{\mathfrak m_0}$; la conclusión se sigue, por tanto, del hecho de que $B_0$ es una $A_0$-álgebra étale (18.2.5).

\begin{env}[18.3.3]
\label{IV.18.3.3}
\emph{Ejemplo.} --- La prop. (18.3.2) se aplica, en particular, cuando $A$ es un \emph{anillo local} separado y completo, siendo $\mathfrak J$ el ideal maximal de $A$, de modo que $A_0$ es un \emph{cuerpo} y la categoría de las $A_0$-álgebras finitas y étales es idéntica a la de las $A_0$-álgebras finitas y separables, que son, por tanto, isomorfas a productos directos de cuerpos, extensiones separables y finitas de $A_0$. En particular, si el cuerpo $A_0$ es \emph{separablemente cerrado}, estas extensiones son todas idénticas a $A_0$ y, por consiguiente, todo recubrimiento étale de $\Spec(A)$ es \emph{trivial} (18.2.7) en virtud de (18.3.2).
\end{env}

\begin{theorem}[18.3.4]
\label{IV.18.3.4}
Sean $A$ un anillo noetheriano y $\mathfrak J$ un ideal de $A$ tal que $A$ sea separado y completo para la topología $\mathfrak J$-preádica, $A_0=A/\mathfrak J$. Pongamos $S=\Spec(A)$, $S_0=\Spec(A_0)$. Sea $X$ un $S$-esquema propio sobre $S$, y pongamos $X_0=X\times_S S_0$. Entonces el funtor
\[
  Z\longmapsto Z\times_X X_0
\]
de la categoría de los $X$-esquemas finitos y étales sobre $X$ a la categoría de los $X_0$-esquemas finitos y étales sobre $X_0$ es una equivalencia de categorías.
\end{theorem}

Demostremos primero que este funtor es \emph{plenamente fiel}. Sean, pues, $Z'$ y $Z''$ dos $X$-esquemas finitos y étales sobre $X$. Pongamos $S_n=\Spec(A/\mathfrak J^{n+1})$, $X_n=X\times_S S_n$, $Z'_n=Z'\times_S S_n$, $Z''_n=Z''\times_S S_n$ para todo $n\geq 0$. Se sigue de (III, 5.4.1) que se tiene una biyección canónica $\Hom_X(Z',Z'')\xrightarrow{\sim}\varprojlim\Hom_{X_n}(Z'_n,Z''_n)$. Ahora bien, en virtud de (18.1.2), la aplicación canónica $\Hom_{X_{n+1}}(Z'_{n+1},Z''_{n+1})\to\Hom_{X_n}(Z'_n,Z''_n)$ es biyectiva, lo que termina de demostrar nuestra afirmación.

Queda por demostrar que, si $\sh{B}_0$ es una $\sh{O}_{X_0}$-álgebra finita y étale, existe una $\sh{O}_X$-álgebra finita y étale $\sh{B}$ y un isomorfismo $\sh{B}_0\xrightarrow{\sim}\sh{B}\otimes_{\sh{O}_X}\sh{O}_{X_0}$. Se sigue de (18.1.2) que existe un sistema proyectivo $(\sh{B}_n)$, donde $\sh{B}_n$ es una $\sh{O}_{X_n}$-álgebra finita y étale, y el segundo teorema de comparación (III, 5.1.4) demuestra la existencia de un $\sh{O}_X$-módulo coherente $\sh{B}$ y de un sistema proyectivo de isomorfismos $\sh{B}_n\xrightarrow{\sim}\sh{B}\otimes_{\sh{O}_X}\sh{O}_{X_n}$. El dato de una estructura de álgebra sobre un módulo $\sh{F}$ equivale al de un homomorfismo $\sh{O}_X\to\sh{F}$ que define la unidad y de un homomorfismo de multiplicación $\sh{F}\otimes_{\sh{O}_X}\sh{F}\to\sh{F}$ que hacen conmutar los diagramas de unidad, asociatividad y conmutatividad en los que solo intervienen potencias tensoriales de $\sh{F}$; se sigue de (III, 5.1.3), (I, 10.11.4) y (I, 10.11.7) que $\sh{B}$ está dotada de manera natural de una estructura de $\sh{O}_X$-álgebra para la cual el isomorfismo $\sh{B}_0\xrightarrow{\sim}\sh{B}\otimes_{\sh{O}_X}\sh{O}_{X_0}$ es un isomorfismo de álgebras. Además, $\sh{B}$ es un $\sh{O}_X$-módulo \emph{localmente libre}; esto se sigue también de (III, 5.1.3), (I, 10.11.4), (I, 10.11.7) y del hecho de que, en la categoría de los $\sh{O}_X$-módulos coherentes, los $\sh{O}_X$-módulos localmente libres $\sh{F}$ pueden definirse como aquellos para los que el funtor $\sh{G}\mapsto\mathscr{H}\!om_{\sh{O}_X}(\sh{F},\sh{G})$ es \emph{exacto}. Finalmente, para ver que $\sh{B}$ es una $\sh{O}_X$-álgebra étale, basta (18.2.5) mostrar que, para todo punto \emph{cerrado} $x\in X$, $\sh{B}_x\otimes_{\sh{O}_{X,x}}\kappa(x)$ es una $\kappa(x)$-álgebra separable. Pero como el morfismo estructural $f:X\to S$ es \emph{propio}, $f(x)$ es un punto cerrado de $S$ y, por tanto, pertenece a $S_0$, ya que $\mathfrak J$ está contenido en el radical de $A$ ($0_{\mathrm I}$, 7.1.10); la conclusión se sigue, por tanto, de que $X_0=f^{-1}(S_0)$ y de que $\sh{B}_0$ es una $\sh{O}_{X_0}$-álgebra finita y étale.

\subsection{Estructura local de los morfismos no ramificados y de los morfismos étales.}
\label{subsection:IV.18.4}

\oldpage[IV-4]{118}

\begin{lemma}[18.4.1]
\label{IV.18.4.1}
Sean $A$ un anillo, $B$ una $A$-álgebra finita monógena, $u$ un generador de la $A$-álgebra $B$, $F\in A[T]$ un polinomio tal que $F(u)=0$, y $F'$ el polinomio derivado; pongamos $u'=F'(u)$. Entonces el ideal de $B$ que anula $\Omega^1_{B/A}$ contiene $u'B$; es igual a $u'B$ si el ideal $\mathfrak J$ de $A[T]$ formado por los polinomios $G$ tales que $G(u)=0$ está generado por $F$, es decir, si el homomorfismo canónico sobreyectivo $\varphi:A[T]/F.A[T]\to B$ que envía la imagen de $T$ a $u$ es biyectivo.
\end{lemma}

Pongamos $C=A[T]$, de modo que $B=C/\mathfrak J$. Se tiene la sucesión exacta (0, 20.5.12.1)
\[
  \mathfrak J/\mathfrak J^2\longrightarrow\Omega^1_{C/A}\otimes_C B\longrightarrow\Omega^1_{B/A}\longrightarrow 0
\]
y $\Omega^1_{B/A}$ se identifica, por tanto, con el cociente $B/\mathfrak J'$, siendo $\mathfrak J'$ el ideal generado por los elementos $G'(u)$, donde $G$ recorre un sistema de generadores del ideal $\mathfrak J$ (0, 20.5.13); de aquí se sigue inmediatamente el lema.

\begin{proposition}[18.4.2]
\label{IV.18.4.2}
Con las notaciones de (18.4.1), sea $\mathfrak q$ un ideal primo de $B$. Entonces:
\begin{enumerate}
  \item[(i)] Si $\mathfrak q$ no contiene a $u'$, $B_{\mathfrak q}$ es una $A_{\mathfrak p}$-álgebra formalmente no ramificada (siendo $\mathfrak p$ la imagen inversa de $\mathfrak q$ en $A$); en otras palabras, $\Spec(B_{u'})$ es formalmente no ramificado sobre $\Spec(A)$.
  \item[(ii)] Supongamos además que $F$ es mónico y genera $\mathfrak J$. Entonces, para que $\Spec(B)$ sea étale sobre $\Spec(A)$ en el punto $\mathfrak q$, es necesario y suficiente que $u'\notin\mathfrak q$.
\end{enumerate}
\end{proposition}

La hipótesis $u'\notin\mathfrak q$ implica que $\Omega^1_{B_{\mathfrak q}/A_{\mathfrak p}}=0$ (0, 20.5.9), por lo que (i) se sigue de (17.2.1). Además, bajo las hipótesis de (ii), $B$ es un $A$-módulo libre en virtud de la división euclídea; como el anulador $\mathfrak J'$ de $\Omega^1_{B/A}$ es entonces igual a $u'B$ en virtud de (18.4.1) y $\Omega^1_{B/A}$ es un $B$-módulo de presentación finita (16.4.22), el anulador de $\Omega^1_{B_{\mathfrak q}/A_{\mathfrak p}}$ es igual a $u'B_{\mathfrak q}$ (Bourbaki, \emph{Alg. comm.}, cap. II, \textsection~2, n.\textsuperscript{o}~4, fórmula (9)), y (ii) se sigue, por tanto, de (i) y de la implicación $c)\Rightarrow a)$ en (17.6.1).

\begin{corollary}[18.4.3]
\label{IV.18.4.3}
Con las notaciones de (18.4.2), supongamos que $F$ es mónico y genera $\mathfrak J$. Entonces, para que $B$ sea una $A$-álgebra étale, es necesario y suficiente que $u'$ sea invertible en $B$ (o, lo que equivale a lo mismo, que el ideal de $A[T]$ generado por $F$ y $F'$ sea igual a $A[T]$).
\end{corollary}

Habida cuenta de (18.4.2), (ii), decir que $\Spec(B)$ es étale sobre $\Spec(A)$ significa, en efecto, que $u'$ no pertenece a ningún ideal primo de $B$, es decir, que es invertible en $B$.

Se dice que un polinomio mónico $F\in A[T]$ tal que el ideal de $A[T]$ generado por $F$ y $F'$ sea igual al propio $A[T]$ es \emph{separable}; es inmediato que esta definición coincide con la definición usual (Bourbaki, \emph{Alg.}, cap. V, \textsection~7, n.\textsuperscript{o}~6) cuando $A$ es un cuerpo.

\begin{lemma}[18.4.4]
\label{IV.18.4.4}
Sean $A$ un anillo local, $B$ una $A$-álgebra finita monógena y $u$ un generador de la $A$-álgebra $B$. Sean $\mathfrak n_i$ ($1\leq i\leq r$) ideales maximales de $B$ tales que $\Spec(B)$ sea formalmente no ramificado sobre $\Spec(A)$ en los puntos $\mathfrak n_i$. Entonces existe un polinomio mónico $F\in A[T]$ tal que $F(u)=0$ y $F'(u)\notin\mathfrak n_i$ para todo índice $i$ ($1\leq i\leq r$). Además, si $k$ es el cuerpo

\oldpage[IV-4]{119}

residual de $A$, $f\in k[T]$ el polinomio minimal de la imagen de $u$ en $B\otimes_A k$, existe un $F\in A[T]$ cuya imagen canónica es $f$ y tal que $F(u)=0$; tal polinomio $F$ satisface las condiciones $F'(u)\notin\mathfrak n_i$ para $1\leq i\leq r$.
\end{lemma}

El ideal maximal $\mathfrak m$ de $A$ es la imagen inversa de cada uno de los $\mathfrak n_i$ (II, 6.1.10); sea $L=B\otimes_A k$, que es una $k$-álgebra finita. Sea $\xi$ la imagen de $u$ en $L$, y sea $n$ el rango de $L$ sobre $k$; el polinomio minimal $f\in k[T]$ de $\xi$ sobre $k$ es, por tanto, de grado $n$, y $L$ es isomorfa a $k[T]/f.k[T]$. Si $\mathfrak n'_i=\mathfrak n_i/\mathfrak mB$, la hipótesis implica que $\Spec(L)$ es étale sobre $k$ en los puntos $\mathfrak n'_i$ ($1\leq i\leq r$) en virtud de (17.6.1, $c)$), y por tanto $f'(\xi)\notin\mathfrak n'_i$ en virtud de (18.4.2, (ii)). Observemos ahora que $\mathfrak mB$ está contenido en el radical de $B$; como $1,\xi,\ldots,\xi^{n-1}$ forman una base de $L$ sobre $k$, se sigue del lema de Nakayama que $1,u,\ldots,u^{n-1}$ generan el $A$-módulo $B$ y, por consiguiente, existe un polinomio mónico $F\in A[T]$ de grado $n$ tal que $F(u)=0$; además, como $\xi$ es raíz de la imagen canónica de $F$ en $k[T]$, esta imagen es necesariamente igual a $f$. Pero entonces la imagen de $F'(u)$ en $L$ es $f'(\xi)$ y, como $f'(\xi)\notin\mathfrak n'_i$, se tiene $F'(u)\notin\mathfrak n_i$ para todo $i$.

\begin{proposition}[18.4.5]
\label{IV.18.4.5}
Sean $A$ un anillo local, $k$ su cuerpo residual y $B$ una $A$-álgebra finita (resp. finita y de presentación finita). Se supone además, o bien que el cuerpo $k$ es infinito, o bien que $B$ es un anillo local. Sea $n$ el rango de $L=B\otimes_A k$ sobre $k$. Para que $B$ sea una $A$-álgebra formalmente no ramificada (resp. étale), es necesario y suficiente que exista un polinomio mónico separable $F\in A[T]$ (18.4.3) tal que $B$ sea isomorfa a un cociente de $A[T]/F.A[T]$ (resp. isomorfa a $A[T]/F.A[T]$). Además, se puede suponer que $F$ es de grado $n$ (resp. $F$ es necesariamente de grado $n$).
\end{proposition}

Las condiciones son suficientes en virtud de (18.4.2), sin suponer que $k$ sea infinito ni que $B$ sea local. Para ver que las condiciones son necesarias, observemos que, si $B$ es una $A$-álgebra formalmente no ramificada, $L$ es un álgebra finita y separable sobre $k$, y por tanto un producto directo de un número finito de extensiones finitas y separables $k_j$ de $k$ ($1\leq j\leq r$), siendo $k_j$ generada por un elemento $\xi_j$, de polinomio minimal $f_j$ ($1\leq j\leq r$) (Bourbaki, \emph{Alg.}, cap. V, \textsection~11, n.\textsuperscript{o}~4, prop. 4). Demostremos que, en virtud de las hipótesis formuladas sobre $k$ o $B$, existe un elemento $\xi$ de $L$ que genera la $k$-álgebra $L$. Esto es inmediato si $B$ es local, pues entonces $r=1$. En caso contrario, como $k$ se supone infinito, se puede suponer que los polinomios irreducibles $f_j\in k[T]$ son todos \emph{distintos} (sustituyendo, si es preciso, cada $\xi_j$ por $\xi_j+a_j$, para un elemento adecuado $a_j\in k$); si se pone $f=f_1f_2\cdots f_r$, es claro que $L$ es isomorfa a $k[T]/f.k[T]$ en los dos casos considerados y, por tanto, está generada por un elemento $\xi$ de polinomio minimal $f\in k[T]$ de grado $n$. Si $u\in B$ es un elemento de imagen $\xi$ en $L$, el lema de Nakayama muestra que los elementos $1,u,\ldots,u^{n-1}$ generan el $A$-módulo $B$; esto muestra ya que existe un polinomio mónico $F\in A[T]$ de grado $n$ tal que $F(u)=0$, y $u$ genera la $A$-álgebra $B$, que es, por consiguiente, isomorfa a un álgebra cociente de $A[T]/F.A[T]$; además, $B$ es un anillo semilocal y, en cada uno de sus ideales maximales $\mathfrak n_i$, se tiene $F'(u)\notin\mathfrak n_i$ por (18.4.4), lo que demuestra que $F'(u)$ es invertible en $B$ y, por tanto, que $F$ es un polinomio \emph{separable}. Finalmente, si $B$ es una $A$-álgebra étale, como $B$ es un $A$-módulo plano de presentación finita (1.4.7), es un $A$-módulo libre, y $1,u,\ldots,u^{n-1}$ forman una base del $A$-módulo $B$ (Bourbaki, \emph{Alg. comm.}, cap. II, \textsection~3, n.\textsuperscript{o}~2, prop. 5); en otras palabras, la $A$-álgebra $B$ es

\oldpage[IV-4]{120}

isomorfa a $A[T]/F.A[T]$, y para todo otro polinomio mónico $G\in A[T]$ tal que $B$ sea isomorfa a $A[T]/G.A[T]$, $G$ es necesariamente de grado $n$.

\begin{theorem}[18.4.6]
\label{IV.18.4.6}
(Chevalley). ---
\begin{enumerate}
  \item[(i)] Sean $f:X\to Y$ un morfismo localmente de tipo finito, $x$ un punto de $X$, $y=f(x)$, y pongamos $A=\sh{O}_{Y,y}$. Para que $\sh{O}_{X,x}$ sea una $A$-álgebra formalmente étale (resp. formalmente no ramificada), es necesario y suficiente que existan un polinomio mónico $F\in A[T]$ y un ideal maximal $\mathfrak n$ de $B=A[T]/F.A[T]$ (resp. de un álgebra cociente $B$ de $A[T]/F.A[T]$) tales que $\sh{O}_{X,x}$ sea $A$-isomorfa a $B_{\mathfrak n}$ y que, si $u$ es la imagen de $T$ en $B$, se tenga $F'(u)\notin\mathfrak n$.
  \item[(ii)] Supongamos además que $f$ sea localmente de presentación finita. Entonces, para que $f$ sea étale en el punto $x$, es necesario y suficiente, además, que se pueda tomar $B=A[T]/F.A[T]$.
\end{enumerate}
\end{theorem}

Las condiciones son suficientes en virtud de (18.4.2). Para ver que son necesarias, se puede evidentemente reducir al caso en que $X$ e $Y$ son afines y, habida cuenta de la observación (17.4.1.2), al caso en que $f$ es formalmente no ramificado y cuasifinito. Como $f$ es afín, se sigue de (8.12.8) que existen una $A$-álgebra finita $C$ y un ideal maximal $\mathfrak r$ de $C$ (necesariamente sobre el ideal maximal $\mathfrak m$ de $A$) tales que $\sh{O}_{X,x}$ sea $A$-isomorfa a $C_{\mathfrak r}$. Además (17.4.1.2), el cuerpo residual $C/\mathfrak r=C_{\mathfrak r}/\mathfrak rC_{\mathfrak r}$ es una extensión finita separable de $k=A/\mathfrak m$, y por tanto de la forma $k[v]$, donde $v$ es \emph{separable} sobre $k$. Sean $\mathfrak r_i$ ($1\leq i\leq h$) los ideales maximales del anillo semilocal $C$ \emph{distintos} de $\mathfrak r$; existe un elemento $u\in C$ perteneciente a todos los $\mathfrak r_i$ y tal que su imagen en $C/\mathfrak r$ sea igual a $v$ (Bourbaki, \emph{Alg. comm.}, cap. II, \textsection~1, n.\textsuperscript{o}~2, prop. 5). Vamos a demostrar que la sub-$A$-álgebra $B=A[u]$ de $C$ y el ideal (necesariamente maximal puesto que $C$ es finita sobre $B$) $\mathfrak n=\mathfrak r\cap B$ de $B$ satisfacen lo requerido.

Para tratar el caso en que $\sh{O}_{X,x}$ es una $A$-álgebra formalmente no ramificada, bastará demostrar que $B_{\mathfrak n}$ es isomorfa a $C_{\mathfrak r}$; en efecto, $B_{\mathfrak n}$ será entonces formalmente no ramificada sobre $A$ y la existencia del polinomio $F\in A[T]$ que tiene las propiedades del enunciado se seguirá de (18.4.4). Observemos ahora que, puesto que $f$ es formalmente no ramificado, $C/\mathfrak r$ es isomorfa a la $A$-álgebra $C_{\mathfrak r}/\mathfrak mC_{\mathfrak r}$ (17.4.1.2). La cuestión queda así reducida a demostrar el lema siguiente:

\begin{lemma}[18.4.6.1]
\label{IV.18.4.6.1}
Sean $A$ un anillo local, $\mathfrak m$ su ideal maximal, $C$ una $A$-álgebra finita y $\mathfrak r$ un ideal maximal de $C$. Sea $u$ un elemento de $C$ que pertenece a todos los ideales maximales de $C$ distintos de $\mathfrak r$, que no pertenece a $\mathfrak r$, y tal que $C_{\mathfrak r}/\mathfrak mC_{\mathfrak r}$ sea un álgebra monógena sobre $k=A/\mathfrak m$, generada por la imagen de $u$ en $C_{\mathfrak r}/\mathfrak mC_{\mathfrak r}$. Pongamos $B=A[u]$, $\mathfrak n=\mathfrak r\cap B$. Entonces el homomorfismo canónico $B_{\mathfrak n}\to C_{\mathfrak r}$ es un isomorfismo.
\end{lemma}

Pongamos $R=B-\mathfrak n$, $S=C-\mathfrak r$, de modo que $B_{\mathfrak n}=R^{-1}B$ y $C_{\mathfrak r}=S^{-1}C$; el homomorfismo canónico $B_{\mathfrak n}\to C_{\mathfrak r}$ puede escribirse como el compuesto
\[
  R^{-1}B\xrightarrow{g}R^{-1}C\xrightarrow{h}S^{-1}C
\]
y basta demostrar que cada uno de estos dos homomorfismos es biyectivo.

Demostremos primero que $h:R^{-1}C\to S^{-1}C=C_{\mathfrak r}$ es biyectivo; basta ver que las imágenes en $R^{-1}C$ de los elementos de $S$ son invertibles, o también que todo ideal maximal $\mathfrak p$ de $R^{-1}C$ tiene una imagen inversa en $C$ disjunta de $S$ y, por tanto, necesariamente igual a $\mathfrak r$. Ahora bien, como $R^{-1}C$ es una $(R^{-1}B)$-álgebra finita, la imagen inversa de $\mathfrak p$ en

\oldpage[IV-4]{121}

$R^{-1}B=B_{\mathfrak n}$ es el único ideal maximal $\mathfrak nB_{\mathfrak n}$ de este anillo, de modo que la imagen inversa de $\mathfrak p$ en $B$ es igual a $\mathfrak n$. Por otra parte, si $\mathfrak q$ es la imagen inversa de $\mathfrak p$ en $C$, se tiene $\mathfrak q\cap B=\mathfrak n$, y como $\mathfrak n$ es maximal en $B$ y $C$ una $B$-álgebra finita, $\mathfrak q$ es necesariamente uno de los ideales maximales de $C$; además, se tiene $u\notin\mathfrak q$ puesto que $u\notin\mathfrak r$ y $u\in B$, por lo que, por hipótesis, se tiene necesariamente $\mathfrak q=\mathfrak r$.

Por otra parte, como $B\subset C$, el homomorfismo $g:R^{-1}B\to R^{-1}C$ es inyectivo ($0_{\mathrm I}$, 1.3.2); para ver que es sobreyectivo, observemos que $R^{-1}C$ es un $(R^{-1}B)$-módulo de tipo finito y, por otra parte, que $\mathfrak mR^{-1}B$ está contenido en el ideal maximal del anillo local $R^{-1}B$; en virtud del lema de Nakayama, basta demostrar que el homomorfismo
\[
  R^{-1}B/\mathfrak mR^{-1}B\longrightarrow R^{-1}C/\mathfrak mR^{-1}C
\]
es sobreyectivo. Pero, en virtud de la primera parte de la demostración, $R^{-1}C/\mathfrak mR^{-1}C$ se identifica con $C_{\mathfrak r}/\mathfrak mC_{\mathfrak r}$; por hipótesis, esta $k$-álgebra está generada por la imagen de $u$ y, \emph{a fortiori}, es igual a la imagen de $R^{-1}B/\mathfrak mR^{-1}B$.

Consideremos en segundo lugar el caso en que $f$ es étale en el punto $x$. Sustituyendo $X$ por un entorno de $x$, se puede suponer que $X$ es un entorno de $\mathfrak n$ en $\Spec(B)$ (I, 7.2). Pongamos $B'=A[T]/F.A[T]$ y sea $\mathfrak n'$ la imagen inversa de $\mathfrak n$ en $B'$; como la imagen de $F'(T)$ en $B'$ no pertenece a $\mathfrak n'$ por hipótesis, el morfismo $\Spec(B')\to\Spec(A)$ es étale en el punto $\mathfrak n'$ por (18.4.2, (ii)). Como, por hipótesis, $\Spec(B)\to\Spec(A)$ es étale en el punto $\mathfrak n$, se concluye (17.3.4) que $\Spec(B)\to\Spec(B')$ es étale en el punto $\mathfrak n$; pero como este morfismo es una inmersión, solo puede ser étale en un punto si es un isomorfismo local en dicho punto (17.9.1), por lo que $B_{\mathfrak n}$ y $B'_{\mathfrak n'}$ son isomorfos.

Finalmente, supongamos que $\sh{O}_{X,x}$ es una $A$-álgebra formalmente étale; con las notaciones anteriores, se sigue de (17.1.5) que $B_{\mathfrak n}$ es una $B'_{\mathfrak n'}$-álgebra formalmente étale; pero, como el homomorfismo $B'_{\mathfrak n'}\to B_{\mathfrak n}$ es sobreyectivo, esto solo puede ocurrir si dicho homomorfismo es \emph{biyectivo} (0, 19.10.3, (i)). Esto termina la demostración de (18.4.6).

\begin{corollary}[18.4.7]
\label{IV.18.4.7}
Sean $f:X\to Y$ un morfismo localmente de tipo finito y $x$ un punto de $X$. Para que $f$ sea formalmente no ramificado en el punto $x$, es necesario y suficiente que exista un entorno abierto $U$ de $x$ tal que $f|U$ se factorice como $U\xrightarrow{j}X'\xrightarrow{h}Y$, donde $h$ es un morfismo étale y $j$ una inmersión cerrada.
\end{corollary}

Se puede evidentemente reducir al caso en que $Y=\Spec(R)$ es afín y $f$ de tipo finito. Si $A=\sh{O}_{Y,y}$, con $y=f(x)$, la condición para que $f$ sea formalmente no ramificado en el punto $x$ equivale, en virtud de (17.4.1.2), a decir que $\sh{O}_{X,x}$ es una $A$-álgebra formalmente no ramificada. Si así ocurre, se puede aplicar (18.4.6, (i)); sustituyendo, si es preciso, $Y$ por un entorno afín de $y$, se puede suponer (con las notaciones de (18.4.6)) que el polinomio $F$ es la imagen en $A[T]$ de un polinomio mónico $G\in R[T]$. Pongamos entonces $X'=\Spec(R[T]/G.R[T])$; sea $x'$ la imagen del punto $\mathfrak n$ de $\Spec(B)$ por el morfismo correspondiente al homomorfismo compuesto $R[T]/G.R[T]\to A[T]/F.A[T]\to B$. Se sigue de (18.4.2) que el morfismo $h:X'\to Y$ correspondiente al homomorfismo canónico $R\to R[T]/G.R[T]$ es étale en el punto $x'$, de modo que, restringiendo si es preciso $X'$ e $Y$ a entornos abiertos de $x'$ e $y$, respectivamente, se puede suponer que $h$ es étale. Por otra parte, en virtud de (I, 6.5.1, (ii)) y de (I, 7.2), del hecho de que se tiene un homomorfismo local $\varphi:\sh{O}_{X',x'}\to\sh{O}_{X,x}$ se sigue que este homomorfismo corresponde a un morfismo

\oldpage[IV-4]{122}

$j:U\to X'$ de un entorno abierto $U$ de $x$ en $X$; restringiendo, si es preciso, $U$, se puede, aplicando (I, 6.5.1, (i)), suponer que $h\circ j=f|U$, de modo que $j$ es de tipo finito (I, 6.3.4, (v)); finalmente, el homomorfismo $\varphi$ es sobreyectivo por (18.4.6, (i)), por lo que se sigue de (I, 6.5.4, (i)) que se puede, restringiendo aún más $U$ y $X'$, suponer que $j$ es una inmersión cerrada. Esto demuestra la necesidad de la condición enunciada; su suficiencia es inmediata (17.1.3, (i) y (ii)).

\begin{corollary}[18.4.8]
\label{IV.18.4.8}
Sea $f:X\to Y$ un morfismo localmente de presentación finita. Para que $f$ sea no ramificado (resp. étale), es necesario y suficiente que exista una familia de morfismos planos $g_\alpha:Y'_\alpha\to Y$ y, para cada $\alpha$, un abierto $U_\alpha$ en $X'_\alpha=X\times_Y Y'_\alpha$, tales que, si $g'_\alpha:X'_\alpha\to X$ y $f'_\alpha:X'_\alpha\to Y'_\alpha$ son las proyecciones canónicas, los $g'_\alpha(U_\alpha)$ formen un recubrimiento de $X$, y que cada uno de los morfismos compuestos $U_\alpha\to X'_\alpha\xrightarrow{f'_\alpha}Y'_\alpha$ sea una inmersión cerrada (resp. abierta). Se pueden tomar además los $g_\alpha$ étales.
\end{corollary}

La necesidad de la condición para los morfismos étales es trivial, tomando un único $Y'_\alpha$ igual a $X$, siendo el morfismo $g_\alpha$ correspondiente igual a $f$, y el abierto $U\subset X\times_Y X$ la diagonal. Cuando $f$ es no ramificado, la necesidad de la condición se sigue de (18.4.7); se toma un recubrimiento abierto $(V_\alpha)$ de $X$ tal que, para cada $\alpha$, $f|V_\alpha$ se factorice como $V_\alpha\xrightarrow{j_\alpha}Y'_\alpha\xrightarrow{g_\alpha}Y$, donde $j_\alpha$ es una inmersión cerrada y $g_\alpha$ un morfismo étale. Entonces $j_\alpha:V_\alpha\to Y'_\alpha$ se factoriza como $V_\alpha\xrightarrow{s_\alpha}X'_\alpha\xrightarrow{f'_\alpha}Y'_\alpha$, donde $s_\alpha$ es una $V_\alpha$-sección de $X'_\alpha$, y como el morfismo $g'_\alpha:X'_\alpha\to X$ es étale, $s_\alpha$ es una inmersión abierta (17.4.1), y basta tomar $U_\alpha=s_\alpha(V_\alpha)$ para obtener lo requerido.

La suficiencia de las condiciones se sigue de (17.7.1); se concluye que $f$ es no ramificado (resp. étale) en cada punto de $g'_\alpha(U_\alpha)$ y, por tanto, en todo $X$, puesto que los $g'_\alpha(U_\alpha)$ recubren $X$.

\begin{proposition}[18.4.9]
\label{IV.18.4.9}
Sean $S$ un preesquema, $f:X\to S$ un morfismo, $h:Y\to S$ un morfismo localmente de presentación finita, $g:X\to Y$ un $S$-morfismo, $x$ un punto de $X$, y $y=g(x)$. Las condiciones siguientes son equivalentes:
\begin{enumerate}
  \item[a)] $h$ es étale en el punto $y$ y $g$ es plano en el punto $x$.
  \item[b)] $h$ es no ramificado en el punto $y$ y $f$ es plano en el punto $x$.
\end{enumerate}
\end{proposition}

Como $f=h\circ g$, $a)$ implica que $f$ es plano en el punto $x$ (2.1.6), y evidentemente que $h$ es no ramificado en el punto $y$, por lo que $a)$ implica $b)$.

Para demostrar que $b)$ implica $a)$, se puede suponer primero que $h$ es no ramificado (sustituyendo $Y$ por un entorno de $y$); después, sustituyendo $S$ por un entorno abierto de $s=h(y)=f(x)$, se puede suponer que existen un morfismo étale $u:S'\to S$, un punto $y'$ de $Y'=Y_{(S')}$ sobre $y$ y un entorno abierto $V$ de $y'$ en $Y'$ tal que, si $h'=h_{(S')}:Y'\to S'$, la restricción de $h'$ a $V$ sea una inmersión cerrada (18.4.8). Si se demuestra entonces que $h'$ es étale en el punto $y'$, se seguirá que $h$ es étale en el punto $y$ (17.7.1, (ii)); además, $f'=f_{(S')}:X_{(S')}\to S'$ es plano en todo punto $x'$ sobre $x$. Como las proyecciones $v:Y'\to Y$, $w:X'\to X$ son morfismos étales (y por tanto planos), si se demuestra que $g'=g_{(S')}:X_{(S')}\to Y_{(S')}$ es plano en el punto $x'$, se seguirá de (2.2.11, (iv)) que $g$ será plano en el punto $x$. Se puede, por tanto, reducir al caso en que $h$ es una \emph{inmersión cerrada}

\oldpage[IV-4]{123}

de presentación finita, suponiéndose que $f$ es plano en el punto $x$. Sea $\mathscr J$ el ideal cuasicoherente de tipo finito de $\sh{O}_S$ que define el subpreesquema cerrado $Y$ de $S$. La hipótesis de que $f$ es plano en el punto $x$ implica que el homomorfismo $\sh{O}_{S,s}\to\sh{O}_{X,x}$ es inyectivo ($0_{\mathrm I}$, 6.5.1); \emph{a fortiori}, el homomorfismo $\sh{O}_{S,s}\to\sh{O}_{Y,y}$ es inyectivo, lo que significa que $\mathscr J_s=0$, puesto que es el núcleo del homomorfismo anterior. Como $\mathscr J$ es de tipo finito, existe un entorno abierto $U$ de $s$ en $S$ tal que $\mathscr J|U=0$ ($0_{\mathrm I}$, 5.2.2). Se puede, por tanto, suponer que $h$ es una inmersión abierta, y entonces es claro que $g$ es plano en el punto $x$.

\begin{proposition}[18.4.10]
\label{IV.18.4.10}
Designemos por $\mathbf P(f,x)$ una propiedad que satisface las condiciones siguientes:
\begin{enumerate}
  \item[$1^\circ$] Para todo morfismo $f:X\to Y$ y todo isomorfismo local $h:Y\to Z$, $\mathbf P(f,x)$ equivale a $\mathbf P(h\circ f,x)$ para $x\in X$.
  \item[$2^\circ$] Para todo morfismo $f:X\to Y$, todo morfismo étale $g:Y'\to Y$, todo punto $x\in X$, si se pone $X'=X\times_Y Y'$, $f'=f_{(Y')}:X'\to Y'$ y si $x'\in X'$ está sobre $x$, las propiedades $\mathbf P(f,x)$ y $\mathbf P(f',x')$ son equivalentes («invariancia por cambio de base étale»).
\end{enumerate}
Sean entonces $S$ un preesquema, $f:X\to S$ y $h:Y\to S$ dos morfismos, $g:X\to Y$ un $S$-morfismo, $x$ un punto de $X$, $y=g(x)$, y supongamos que $h$ es étale en el punto $y$. Entonces las propiedades $\mathbf P(f,x)$ y $\mathbf P(g,x)$ son equivalentes.
\end{proposition}

Con las mismas notaciones que en la demostración de (18.4.9) y sustituyendo $S$ (resp. $Y$) por un entorno de $s=h(y)$ (resp. de $y$), se puede aquí, en virtud de (18.4.8), encontrar un morfismo étale $u:S'\to S$ tal que $h'$ sea una inmersión abierta; si $x'\in X'$ está sobre $x$, $\mathbf P(f',x')$ (resp. $\mathbf P(g',x')$) es entonces, por hipótesis, equivalente a $\mathbf P(f,x)$ (resp. $\mathbf P(g,x)$). Se puede, por tanto, reducir al caso en que $h$ es una inmersión abierta. Como la hipótesis implica que $\mathbf P(g,x)$ es entonces equivalente a $\mathbf P(h\circ g,x)$, se deduce la conclusión.

\begin{env}[18.4.11]
\label{IV.18.4.11}
\emph{Ejemplos.} --- Se puede tomar como propiedad $\mathbf P(f,x)$ cualquiera de las siguientes, habida cuenta de (17.7.4, (ii)):
\begin{enumerate}
  \item[(i)] $f$ es plano en el punto $x$ (2.2.11, (iv));
  \item[(ii)] $f$ es localmente de presentación finita y de coprofundidad $\leq n$ en el punto $x$ (6.8.1 y 6.7.8);
  \item[(iii)] $f$ es localmente de presentación finita y de Cohen-Macaulay en el punto $x$ (6.8.1 y 6.7.8);
  \item[(iv)] $f$ es localmente de presentación finita y posee la propiedad $(S_n)$ en el punto $x$ (6.8.1 y 6.7.8);
  \item[(v)] $f$ es localmente de presentación finita y posee la propiedad $(R_n)$ en el punto $x$ (6.8.1 y 6.7.8);
  \item[(vi)] $f$ es localmente de presentación finita y normal en el punto $x$ (6.8.1 y 6.7.8);
  \item[(vii)] $f$ es localmente de presentación finita y reducido en el punto $x$ (6.8.1 y 6.7.8);
  \item[(viii)] $f$ es no ramificado en el punto $x$ (17.7.4);
  \item[(ix)] $f$ es liso en el punto $x$ (17.7.4);
  \item[(x)] $f$ es étale en el punto $x$ (17.7.4).
\end{enumerate}
\end{env}

\oldpage[IV-4]{124}

\begin{corollary}[18.4.12]
\label{IV.18.4.12}
\begin{enumerate}
  \item[(i)] Sean $f:X\to Y$ un morfismo localmente de tipo finito y $x$ un punto de $X$. Si $f$ es plano y formalmente no ramificado en el punto $x$, entonces, en toda factorización $f|U:U\xrightarrow{j}X'\xrightarrow{h}Y$, donde $U$ es un entorno abierto de $x$, $j$ una inmersión cerrada y $h$ un morfismo étale (18.4.7), el homomorfismo $\sh{O}_{X',j(x)}\to\sh{O}_{X,x}$ correspondiente a $j$ es biyectivo (en particular, $\sh{O}_{X,x}$ es una $\sh{O}_{Y,f(x)}$-álgebra \emph{esencialmente étale} (18.6.1)).
  \item[(ii)] Para que un morfismo $f:X\to Y$ sea étale, es necesario y suficiente que sea localmente de tipo finito, formalmente no ramificado y plano.
\end{enumerate}
\end{corollary}

(i) Como el homomorfismo $\sh{O}_{X',j(x)}\to\sh{O}_{X,x}$ es sobreyectivo, basta demostrar que es inyectivo y, para ello, que hace de $\sh{O}_{X,x}$ un $\sh{O}_{X',j(x)}$-módulo \emph{fielmente plano}, o solamente \emph{plano} ($0_{\mathrm I}$, 6.5.1 y 6.6.2); en otras palabras, se trata de demostrar que $j$ es un morfismo \emph{plano en el punto} $x$; pero, puesto que $h\circ j=f$ es, por hipótesis, plano en el punto $x$ y que $h$ es étale, esto se sigue de (18.4.10) y (18.4.11, (i)).

(ii) Solo hay que demostrar la suficiencia de las condiciones enunciadas. Para todo $x\in X$, se tiene, por tanto, en un entorno abierto $U$ de $x$ una factorización de $f|U$ con las propiedades consideradas en (i). Como, por hipótesis, $f$ es plano y formalmente no ramificado en \emph{todos} los puntos de $U$, el resultado de (i) se aplica no solo a $x$ sino a todos los puntos de $U$; esto significa que, si $\mathscr J$ es el ideal cuasicoherente de $\sh{O}_{X'}$ correspondiente al subpreesquema cerrado de $X'$ asociado a $j$, se tiene $\mathscr J_{j(z)}=0$ para todo $z\in U$, por lo que $j$ es una inmersión abierta, y $f$ es, por consiguiente, étale en todo punto de $U$, y por tanto en todo punto de $X$.

\begin{corollary}[18.4.13]
\label{IV.18.4.13}
Sean $f:X\to Y$ un morfismo localmente de tipo finito, $x$ un punto de $X$, y $y=f(x)$. Supongamos que $y$ admite un entorno abierto que es un preesquema reducido y solo tiene un número finito de componentes irreducibles. Entonces, para que $f$ sea étale en el punto $x$, es necesario y suficiente que $f$ sea plano y formalmente no ramificado en el punto $x$.
\end{corollary}

Solo hay que demostrar la suficiencia de la condición. Como la cuestión es local sobre $X$ e $Y$, se puede suponer (18.4.7) que $f$ se factoriza como $X\xrightarrow{j}X'\xrightarrow{h}Y$, donde $h$ es étale y $j$ una inmersión cerrada, y además que $Y$ es reducido y solo tiene un número finito de componentes irreducibles. Entonces $X'$ es reducido (17.5.7) y, sustituyendo si es preciso $X'$ por un entorno abierto de $j(x)$, se puede suponer que $X'$ solo tiene un número finito de componentes irreducibles: en efecto, se puede suponer que $h$ es cuasifinito (17.6.1) y, como los puntos maximales de $X'$ están sobre los puntos maximales de $Y$ (2.3.4), su número es finito. Todo se reduce a demostrar, con las notaciones de la demostración de (18.4.12), que se tiene $\mathscr J_{x'}=0$ para todos los puntos $x'$ de un entorno de $j(x)$ en $X'$, sabiendo que $\mathscr J_{j(x)}=0$. Ahora bien, sustituyendo $X'$ por un entorno afín de $j(x)$, se puede suponer que todas las componentes irreducibles de $X'$ contienen a $j(x)$; si $X'=\Spec(A')$, y si $\mathfrak p'$ es el ideal primo de $A'$ correspondiente al punto $j(x)$, el morfismo $\Spec(A'_{\mathfrak p'})\to\Spec(A')$ es dominante, por lo que el homomorfismo correspondiente $A'\to A'_{\mathfrak p'}$ es \emph{inyectivo}, puesto que $A'$ es reducido (I, 1.2.7). Si $\mathscr J=\widetilde{\mathfrak J}$, donde $\mathfrak J$ es un ideal de $A'$, $\mathfrak J$ se identifica, por tanto, con un subconjunto de $\mathfrak J_{\mathfrak p'}$, y la hipótesis $\mathfrak J_{\mathfrak p'}=0$ implica, por tanto, $\mathfrak J=0$.

\begin{corollary}[18.4.14]
\label{IV.18.4.14}
Sean $A$ un anillo local de ideal maximal $\mathfrak m$, de cuerpo residual $k$, y $B$ una $A$-álgebra finita.

\oldpage[IV-4]{125}

\begin{enumerate}
  \item[(i)] Para que $B$ sea una $A$-álgebra formalmente no ramificada, es necesario y suficiente que $B\otimes_A k$ sea una $k$-álgebra étale.
  \item[(ii)] Para que $B$ sea una $A$-álgebra étale, es necesario y suficiente que $B\otimes_A k$ sea una $k$-álgebra étale y que $B$ sea un $A$-módulo plano (lo que equivale a decir ($0_{\mathrm I}$, 6.3.3) que, para todo ideal maximal $\mathfrak n$ de $B$ (necesariamente sobre $\mathfrak m$), $B_{\mathfrak n}$ es un $A$-módulo plano).
\end{enumerate}
\end{corollary}

Evidentemente, solo hay que demostrar la suficiencia de las condiciones enunciadas. Es claro que (ii) se sigue de (i) y de (18.4.12, (ii)), habida cuenta de (17.1.2, (i)). Para demostrar (i), observemos que, si $B\otimes_A k$ es una $k$-álgebra étale, se tiene $\Omega^1_{(B\otimes_A k)/k}=0$ (17.2.1). Ahora bien, se tiene $\Omega^1_{(B\otimes_A k)/k}=\Omega^1_{B/A}\otimes_B(B\otimes_A k)$ (0, 20.5.5) y, puesto que $B$ es una $A$-álgebra de tipo finito, $\Omega^1_{B/A}$ es un $B$-módulo de tipo finito (0, 20.4.7). Pero como $B$ es una $A$-álgebra finita, $\mathfrak mB$ está contenido en el radical de $B$ (Bourbaki, \emph{Alg. comm.}, cap. V, \textsection~2, n.\textsuperscript{o}~1, prop. 1), por lo que el lema de Nakayama demuestra que $\Omega^1_{B/A}=0$ y, por consiguiente, $B$ es una $A$-álgebra formalmente no ramificada (17.2.1).

\subsection{Anillos locales henselianos\texorpdfstring{\protect\footnotemark}{}.}
\footnotetext{La noción de anillo local henseliano se debe a Azumaya, y la de henselización a Nagata, a quien se deben también los principales resultados de esta teoría.}
\label{subsection:IV.18.5}

\begingroup
\newcommand{\shHom}{\mathscr{H}\!om}
\begin{env}[18.5.1]
\label{IV.18.5.1}
Sean $X$ un preesquema y $\sh{E}$ un $\sh{O}_X$-módulo \emph{localmente libre} de rango finito; el dual $\check{\sh{E}}=\shHom_{\sh{O}_X}(\sh{E},\sh{O}_X)$ es entonces un $\sh{O}_X$-módulo localmente libre cuyo rango en todo punto de $X$ es igual al de $\sh{E}$ en dicho punto, y el homomorfismo canónico $\sh{E}\to\shHom_{\sh{O}_X}(\check{\sh{E}},\sh{O}_X)=(\check{\sh{E}})^\vee$ es \emph{biyectivo}. Para todo morfismo $X'\to X$, pongamos $\sh{E}_{(X')}=\sh{E}\otimes_{\sh{O}_X}\sh{O}_{X'}$, que es un $\sh{O}_{X'}$-módulo localmente libre, y consideremos el \emph{conjunto} $\Gamma(X',\sh{E}_{(X')})$ de las secciones de este $\sh{O}_{X'}$-módulo sobre $X'$. Vamos a ver que se define así un \emph{funtor contravariante representable}
\begin{equation}
  {}^t\!\mathbf V_{\sh{E}}:X'\longmapsto\Gamma(X',\sh{E}_{(X')})
  \tag{18.5.1.1}\label{IV.18.5.1.1}
\end{equation}
de la categoría de los $X$-preesquemas en la de los \emph{conjuntos} ($0_{\mathrm{III}}$, 8.1.8).

En primer lugar, se ha definido efectivamente un funtor, pues si $f:X''\to X'$ es un $X$-morfismo de $X$-preesquemas, se tiene $\sh{E}_{(X'')}=f^*(\sh{E}_{(X')})$, de donde ($0_{\mathrm I}$, 4.4.3.2) se obtiene una aplicación $\Gamma(X',\sh{E}_{(X')})\to\Gamma(X'',f_*(\sh{E}_{(X'')}))=\Gamma(X'',\sh{E}_{(X'')})$ que termina de definir el funtor ${}^t\!\mathbf V$. Demostremos a continuación que el $X$-preesquema $\mathbf V(\check{\sh{E}})$ (II, 1.7.8) \emph{representa} el funtor ${}^t\!\mathbf V$. En efecto, es inmediato que se tiene $(\sh{E}_{(X')})^\vee=(\check{\sh{E}})_{(X')}$, y por tanto $\mathbf V(\check{\sh{E}})\times_X X'=\mathbf V((\check{\sh{E}})_{(X')})$; teniendo en cuenta (I, 3.3.14), la cuestión se reduce a definir una biyección $\Gamma(\mathbf V(\check{\sh{E}})/X)\to\Gamma(X,\sh{E})$ y a verificar que la biyección $\Gamma(\mathbf V((\check{\sh{E}})_{(X')})/X')\to\Gamma(X',\sh{E}_{(X')})$ es funtorial en $X'$. Ahora bien, se tiene una biyección canónica de $\Gamma(\mathbf V(\check{\sh{E}})/X)$ sobre $\Hom_{\sh{O}_X}(\check{\sh{E}},\sh{O}_X)$ (II, 1.7.8), y la transposición $u\mapsto{}^t u$ es una biyección canónica de $\Hom_{\sh{O}_X}(\check{\sh{E}},\sh{O}_X)$ sobre $\Hom_{\sh{O}_X}(\sh{O}_X,\sh{E})=\Gamma(X,\sh{E})$, en virtud de la identificación de $(\check{\sh{E}})^\vee$ con $\sh{E}$. La verificación de la propiedad funtorial es inmediata.
\end{env}
\endgroup

\oldpage[IV-4]{126}

Observemos además que, de conformidad con la teoría general ($0_{\mathrm{III}}$, 8.1.6), el automorfismo identidad de $\mathbf V(\check{\sh{E}})$ corresponde canónicamente a una sección $c$ de $\sh{E}_{(\mathbf V(\check{\sh{E}}))}$ sobre $\mathbf V(\check{\sh{E}})$, es decir (II, 1.4.1), a un homomorfismo de $\mathbf S_{\sh{O}_X}(\check{\sh{E}})$-módulos
\[
  u:\mathbf S_{\sh{O}_X}(\check{\sh{E}})\longrightarrow
  \mathbf S_{\sh{O}_X}(\check{\sh{E}})\otimes_{\sh{O}_X}\sh{E};
\]
para todo abierto afín $W$ de $X$, si se pone $\Gamma(W,\sh{O}_X)=A$, si se identifica $\Gamma(W,\mathbf S_{\sh{O}_X}(\check{\sh{E}}))$ con un álgebra de polinomios $C=A[T_1,\ldots,T_n]$, de modo que los $T_i$ formen una base de $\Gamma(W,\check{\sh{E}})$, y si, por último, se denota por $(e_i)$ la base dual de $(T_i)$ en $\Gamma(W,\sh{E})=\Gamma(W,\check{\sh{E}})^\vee$, se ve inmediatamente que $u$ corresponde al homomorfismo de $C$-módulos tal que
\[
  u(1)=\sum_{i=1}^n T_i\otimes e_i.
\]
Si $X'=\Spec(A)$ es afín y tal que $\sh{E}_{(X')}$ sea isomorfo a $\sh{O}_{X'}^n$, $\Gamma(X',\sh{E}_{(X')})$ es un $A$-módulo libre de rango $n$; de manera figurada, puede decirse que el objeto $\mathbf V(\check{\sh{E}})$ representa «el conjunto de puntos del espacio afín torcido sobre $X$ definido por $\sh{E}$».

\begin{env}[18.5.2]
\label{IV.18.5.2}
Recordemos (II, 1.7.8) que, por definición, se tiene $\mathbf V(\check{\sh{E}})=\Spec(\mathbf S_{\sh{O}_X}(\check{\sh{E}}))$. Sea $\mathscr J$ un ideal cuasicoherente de $\mathbf S_{\sh{O}_X}(\check{\sh{E}})$, de modo que $\Spec(\mathbf S_{\sh{O}_X}(\check{\sh{E}})/\mathscr J)$ sea un subpreesquema cerrado de $\mathbf V(\check{\sh{E}})$; vamos a interpretarlo como \emph{representante de un funtor} de la categoría de los $X$-preesquemas en la de los conjuntos. Observemos para ello que una sección $u\in\Gamma(X,\sh{E})$ se identifica canónicamente con un $\sh{O}_X$-homomorfismo $u:\sh{O}_X\to\sh{E}$, al cual corresponde por transposición un $\sh{O}_X$-homomorfismo ${}^t u:\check{\sh{E}}\to\sh{O}_X$ y, por consiguiente, un homomorfismo de $\sh{O}_X$-álgebras $v:\mathbf S_{\sh{O}_X}(\check{\sh{E}})\to\sh{O}_X$. Sea $\operatorname{Al}(X,\sh{E},\mathscr J)$ el conjunto de los $u\in\Gamma(X,\sh{E})$ tales que $\mathscr J$ esté \emph{contenido en el núcleo} de $v$; de estas definiciones resulta inmediatamente que
\[
  X'\longmapsto\operatorname{Al}(X',\sh{E}_{(X')},\mathscr J\otimes_{\sh{O}_X}\sh{O}_{X'})
\]
es un funtor representado por $\Spec(\mathbf S_{\sh{O}_X}(\check{\sh{E}})/\mathscr J)$. Si $X'=\Spec(A)$ es afín y tal que $\sh{E}_{(X')}$ sea isomorfo a $\sh{O}_{X'}^n$, $\Gamma(X',\sh{E}_{(X')})$ puede identificarse con el conjunto $A^n$, y $\mathscr J\otimes_{\sh{O}_X}\sh{O}_{X'}$ con un $\sh{O}_{X'}$-módulo de la forma $\widetilde{\mathfrak J}$, donde $\mathfrak J$ es un ideal del anillo de polinomios $A[T_1,\ldots,T_n]$; el conjunto $\operatorname{Al}(X',\sh{E}_{(X')},\mathscr J\otimes_{\sh{O}_X}\sh{O}_{X'})$ se identifica entonces con la parte de $A^n$ formada por los puntos $(t_1,\ldots,t_n)$ tales que $P(t_1,\ldots,t_n)=0$ para todos los polinomios $P\in\mathfrak J$; de manera figurada, puede decirse por tanto que el objeto $\Spec(\mathbf S_{\sh{O}_X}(\check{\sh{E}})/\mathscr J)$ representa «el subconjunto algebraico del espacio afín torcido $\Gamma(X,\sh{E})$ formado por los puntos que anulan el ideal $\Gamma(X,\mathscr J)$». Este $X$-preesquema se denota también por $\mathbf{Al}(\sh{E},\mathscr J)$. Obsérvese que, si $\mathscr J$ es un ideal de tipo finito de $\mathbf S_{\sh{O}_X}(\check{\sh{E}})$, $\mathbf{Al}(\sh{E},\mathscr J)$ es un $X$-preesquema \emph{de presentación finita}, puesto que $\mathbf S_{\sh{O}_X}(\check{\sh{E}})$ es una $\sh{O}_X$-álgebra de presentación finita.
\end{env}

\begin{lemma}[18.5.3]
\label{IV.18.5.3}
Sean $S$ un preesquema y $f:X\to S$ un morfismo finito y localmente libre (18.2.7). Consideremos el funtor contravariante de la categoría de los $S$-preesquemas en la categoría de los conjuntos
\begin{equation}
  S'\longmapsto\operatorname{Of}(X\times_S S')
  \tag{18.5.3.1}\label{IV.18.5.3.1}
\end{equation}
donde $\operatorname{Of}(X\times_S S')$ es el conjunto de las partes a la vez abiertas y cerradas del espacio subyacente a $X\times_S S'$. Entonces este funtor es representable por un $S$-preesquema $\mathbf{Of}(X)$, que es afín, étale y de presentación finita sobre $S$.
\end{lemma}

\oldpage[IV-4]{127}

Por hipótesis se tiene $X=\Spec(\mathscr B)$, donde $\mathscr B=f_*(\sh{O}_X)$ es una $\sh{O}_S$-álgebra finita y localmente libre. Pongamos $X'=X\times_S S'$, $f'=f_{(S')}$, $\mathscr B'=\mathscr B\otimes_{\sh{O}_S}\sh{O}_{S'}=f'_*(\sh{O}_{X'})$, de modo que $X'=\Spec(\mathscr B')$; existe entonces una biyección canónica, funtorial en $S'$, del conjunto $\operatorname{Of}(X')$ sobre el conjunto $\operatorname{Id}(\mathscr B')$ de los \emph{idempotentes} del anillo $\Gamma(S',\mathscr B')=\Gamma(X',\sh{O}_{X'})$. En efecto, en virtud de la \emph{equivalencia} entre la categoría de los $S'$-esquemas afines sobre $S'$ y la categoría opuesta de la categoría de las $\sh{O}_{S'}$-álgebras cuasicoherentes (II, 1.2.7 y 1.3.1), existe una correspondencia biunívoca canónica entre las descomposiciones de $X'$ en \emph{suma} $X'_1\amalg X'_2$ de dos subpreesquemas inducidos sobre abiertos de $X'$ y las descomposiciones de $\mathscr B'$ en \emph{suma directa} de dos ideales $\mathscr B'_1$, $\mathscr B'_2$; estas últimas, a su vez, forman un conjunto en correspondencia biunívoca canónica (y funtorial en $S'$) con $\operatorname{Id}(\mathscr B')$. Basta, pues, demostrar el lema para el funtor $S'\longmapsto\operatorname{Id}(\mathscr B')$.

Para ello, vamos a demostrar que existe un ideal de tipo finito $\mathscr J$ de $\mathbf S_{\sh{O}_S}(\check{\mathscr B})$ tal que $\operatorname{Id}(\mathscr B')$ sea de la forma $\operatorname{Al}(S',\mathscr B',\mathscr J\otimes_{\sh{O}_S}\sh{O}_{S'})$ (18.5.2). Observemos a este efecto que, puesto que $\mathscr B$ es una $\sh{O}_S$-álgebra, su imagen inversa $\mathscr B_{(\mathbf V(\check{\mathscr B}))}$ es una $\sh{O}_{(\mathbf V(\check{\mathscr B}))}$-álgebra, y puede formarse entonces en esta álgebra el \emph{cuadrado} $c^2$ de la sección canónica $c$ (18.5.1); este corresponde canónicamente a un homomorfismo $u^{(2)}:\mathbf S_{\sh{O}_S}(\check{\mathscr B})\to\mathbf S_{\sh{O}_S}(\check{\mathscr B})\otimes_{\sh{O}_S}\mathscr B$ de $\mathbf S_{\sh{O}_S}(\check{\mathscr B})$-módulos, y se verifica inmediatamente que, para un abierto afín $W$ de $S$, con las notaciones de (18.5.1), $u^{(2)}$ corresponde al homomorfismo de $C$-módulos tal que
\[
  u^{(2)}(1)=\sum_k\Bigl(\sum_{i,j}c_{ijk}T_iT_j\Bigr)\otimes e_k,
\]
donde $(c_{ijk})$ es la tabla de multiplicación del álgebra $\Gamma(W,\mathscr B)$. Demostremos que el ideal $\mathscr J$ de $\mathbf S_{\sh{O}_S}(\check{\mathscr B})$ engendrado por el \emph{núcleo} del homomorfismo $u-u^{(2)}$ responde a la cuestión. Basta, en efecto, observar que el ideal $\Gamma(W,\mathscr J)$ de $C$ está engendrado por los polinomios $P_k(T_1,\ldots,T_n)=T_k-\sum_{i,j}c_{ijk}T_iT_j$ y que los idempotentes de $\Gamma(W,\mathscr B)$ son precisamente los elementos $\sum_i t_i e_i$ de esta álgebra tales que $(t_1,\ldots,t_n)$ anula todos los polinomios $P_k$ ($1\leq i\leq n$). Se deduce entonces de (18.5.2) que el $S$-preesquema afín de presentación finita $\mathbf{Of}(X)=\mathbf{Al}(\mathscr B,\mathscr J)$ representa efectivamente el funtor (18.5.3.1). Queda por demostrar que $\mathbf{Of}(X)$ es étale sobre $S$ o, lo que es equivalente, que es \emph{formalmente étale} sobre $S$. Pero si $S'$ es un $S$-preesquema y $S'_0$ un subpreesquema cerrado de $S'$ definido por un ideal localmente nilpotente de $\sh{O}_{S'}$ (y que tiene, por tanto, el mismo espacio subyacente que $S'$), está claro que $X'=X\times_S S'$ y $X'_0=X\times_S S'_0$ tienen el mismo espacio subyacente; por tanto, la aplicación canónica $\operatorname{Of}(X')\to\operatorname{Of}(X'_0)$ es \emph{biyectiva}, lo que termina la demostración (17.1.1).

\begin{proposition}[18.5.4]
\label{IV.18.5.4}
Sean $S$ un preesquema y $S_0$ un subpreesquema cerrado de $S$; consideremos las propiedades siguientes:
\begin{enumerate}
  \item[a)] Para todo morfismo finito $g:S'\to S$, la aplicación canónica
  \begin{equation}
    \operatorname{Of}(S')\longrightarrow\operatorname{Of}(S'\times_S S_0)
    \tag{18.5.4.1}\label{IV.18.5.4.1}
  \end{equation}
  (cf. (18.5.3)) es biyectiva.
  \item[a')] Para todo morfismo finito y localmente libre $g:S'\to S$, la aplicación (18.5.4.1) es biyectiva.

\oldpage[IV-4]{128}

  \item[b)] Para todo morfismo étale y separado $g:S'\to S$, la aplicación canónica
  \begin{equation}
    \Gamma(S'/S)\longrightarrow\Gamma(S'\times_S S_0/S_0)
    \tag{18.5.4.2}\label{IV.18.5.4.2}
  \end{equation}
  es biyectiva.
\end{enumerate}
La condición $b)$ implica $a')$; si además $S$ es cuasicompacto y cuasiseparado, la condición $a)$ implica $b)$.
\end{proposition}

Demostremos primero que $b)$ implica $a')$. Supongamos, pues, verificada $b)$, y sea $g:S'\to S$ un morfismo finito y localmente libre; pongamos $S'_0=S'\times_S S_0$, de modo que $g_0=g_{(S_0)}:S'_0\to S_0$ es finito y localmente libre. Entonces se deduce de (18.5.3) que $P=\mathbf{Of}(S')$ es un $S$-preesquema \emph{étale y separado}; además, la definición del funtor $\operatorname{Of}$ muestra inmediatamente que, si se pone $P_0=\mathbf{Of}(S'_0)$ (para la categoría de los $S_0$-preesquemas), se tiene $P_0=P\times_S S_0$. Dicho esto, se tiene por definición el diagrama conmutativo
\[
\xymatrix{
  \Gamma(P/S)\ar[r]\ar[d]^{\sim} & \Gamma(P_0/S_0)\ar[d]^{\sim}\\
  \operatorname{Of}(S')\ar[r] & \operatorname{Of}(S'_0)
}
\]
donde las flechas verticales son las biyecciones canónicas. Como la hipótesis $b)$, aplicada al morfismo $P\to S$, implica que la primera fila es una biyección, lo mismo ocurre con la segunda, lo que demuestra nuestra afirmación.

Antes de demostrar que $a)$ implica $b)$ cuando $S$ es cuasicompacto y cuasiseparado, estableceremos el

\begin{lemma}[18.5.4.3]
\label{IV.18.5.4.3}
Si $S$ y $S_0$ verifican la condición $a)$ de (18.5.4), entonces, para todo morfismo finito $g:S'\to S$, $S'$ es la única vecindad de $S'_0=g^{-1}(S_0)=S'\times_S S_0$ en $S'$.
\end{lemma}

En efecto, equivale a decir que, si $T'$ es una parte cerrada de $S'$ tal que $T'\cap S'_0=\varnothing$, entonces $T'=\varnothing$. Ahora bien, si se denota también por $T'$ un subpreesquema cerrado de $S'$ cuyo espacio subyacente es $T'$, el morfismo compuesto $h:T'\to S'\xrightarrow{g}S$ es finito y $h^{-1}(S_0)$ es vacío; la condición $a)$ aplicada al morfismo $h$ implica necesariamente que $T'$ es vacío.

Demostrado este lema, probemos primero que, bajo la hipótesis $a)$, la aplicación (18.5.4.2) es \emph{inyectiva}. En efecto, si $u'$, $u''$ son dos $S$-secciones de $S'$, el hecho de que el morfismo $S'\to S$ sea no ramificado implica que el preesquema de coincidencias de $u'$ y $u''$ está inducido sobre un abierto $U$ de $S$ (17.4.6). Si las restricciones a $S_0$ de $u'$ y $u''$ son iguales, el hecho de que $u'$ y $u''$ sean inmersiones abiertas (17.4.1) implica que $U$ contiene a $S_0$ y, por tanto, es igual a $S$ en virtud del lema (18.5.4.3), aplicado al caso en que $S'=S$.

\oldpage[IV-4]{129}

Queda por demostrar que, bajo la hipótesis $a)$, la aplicación (18.5.4.2) es \emph{sobreyectiva} (siendo $S$ cuasicompacto y cuasiseparado). Sea, pues, $u_0:S_0\to S'_0$ una $S_0$-sección de $S'_0$; como $u_0(S_0)$ es cuasicompacto en $S'$, puede recubrirse por un número finito de abiertos afines $V_i$ tales que la restricción $g|V_i$ sea un morfismo de tipo finito; si $V$ es la unión de los $V_i$, se concluye que $g|V:V\to S$ es un morfismo \emph{de presentación finita}, pues es separado y localmente de presentación finita por hipótesis (1.6.1). Sustituyendo $S'$ por $V$, puede suponerse entonces que $g$ es de presentación finita. Como $S$ es cuasicompacto y cuasiseparado y $g$ es cuasifinito y separado (17.6.1), se deduce del «Main theorem» (8.12.6) que $g$ se factoriza como $S'\xrightarrow{j}S''\xrightarrow{f}S$, donde $j$ es una inmersión abierta y $f$ un morfismo \emph{finito}. Pongamos $S''_0=S''\times_S S_0$, $j_0=j_{(S_0)}:S'_0\to S''_0$, que es una inmersión abierta, y $f_0=f_{(S_0)}:S''_0\to S_0$, que es un morfismo finito. Entonces $u_0$ es también una $S_0$-sección de $S''_0$. Como $g_0:S'_0\to S_0$ es étale, $u_0$ es una inmersión abierta de $S_0$ en $S'_0$ (17.4.1), por lo que $u_0(S_0)$ es abierto en $S'_0$ y \emph{a fortiori} en $S''_0$; pero, por otra parte, como $f_0$ es un morfismo finito, y por tanto separado, $u_0$ es una inmersión cerrada de $S_0$ en $S''_0$ (I, 5.4.6), de modo que $X_0=u_0(S_0)$ es \emph{a la vez abierto y cerrado} en $S''_0$. En virtud de la hipótesis $a)$, existe un subconjunto a la vez abierto y cerrado $X$ de $S''$ tal que $X\cap S''_0=X_0$. Demostremos primero que el morfismo $f:S''\to S$ es \emph{étale en los puntos de $X$}: en efecto, el conjunto $U$ de los puntos de $X$ en los que $f|X$ es étale es abierto y contiene por hipótesis a $X_0\subset S'_0$. Pero el lema (18.5.4.3), aplicado al morfismo finito $f|X$, demuestra que $U=X$. Por otra parte, $S'\cap X$ es abierto en $X$ y contiene a $X_0$ por hipótesis; por tanto, el mismo razonamiento demuestra que $S'\cap X=X$, es decir, $X\subset S'$. Queda por demostrar que, para todo $s\in S$, el número geométrico $n(s)$ de puntos de $X\cap f^{-1}(s)$ es igual a $1$, pues de ello se seguirá que $f|X$ es radicial y sobreyectivo, y, como $f|X$ es étale, se habrá demostrado (17.9.1) que $f|X$ es un isomorfismo del abierto $X\subset S'$ sobre $S$, cuyo isomorfismo inverso $u$ será la $S$-sección buscada que prolonga $u_0$. Pero, como $f|X$ es étale y \emph{finito}, $s\mapsto n(s)$ es continua en $S$ (18.2.8), y, como $X\cap S''_0=X_0$, se tiene $n(s)=1$ en $S_0$; el conjunto de los puntos $s\in S$ tales que $n(s)=1$, al ser abierto en $S$ y contener a $S_0$, es igual a $S$ por (18.5.4.3). c.q.d.

\begin{remark}[18.5.4.4]
\label{IV.18.5.4.4}
Puede demostrarse que el enunciado (18.5.4) sigue siendo válido cuando, en la condición $b)$, se supone únicamente que el morfismo $g$ es étale (pero no necesariamente separado) [43, exp. XII].
\end{remark}

\begin{definition}[18.5.5]
\label{IV.18.5.5}
Se dice que un preesquema $S$ y un subpreesquema cerrado $S_0$ de $S$ forman un \emph{par henseliano} si verifican la condición $a)$ de (18.5.4).
\end{definition}

Teniendo en cuenta (I, 5.1.8), equivale a decir que $(S,S_0)$ es un par henseliano o que $(S_{\mathrm{red}},(S_0)_{\mathrm{red}})$ lo es.

\begin{proposition}[18.5.6]
\label{IV.18.5.6}
\begin{enumerate}
  \item[(i)] Si $(S,S_0)$ es un par henseliano, entonces, para todo morfismo finito $f:S'\to S$, si $S'_0$ es el subpreesquema $f^{-1}(S_0)$ de $S'$, el par $(S',S'_0)$ es henseliano.
  \item[(ii)] Sean $S=\coprod_\alpha S^{(\alpha)}$ una suma de preesquemas y $S_0$ un subpreesquema cerrado de $S$, suma de los subpreesquemas cerrados $S_0^{(\alpha)}$ de los $S^{(\alpha)}$. Para que el par $(S,S_0)$ sea henseliano, es necesario y suficiente que lo sea cada uno de los pares $(S^{(\alpha)},S_0^{(\alpha)})$.
\end{enumerate}
\end{proposition}

La afirmación (i) es consecuencia inmediata de la definición, puesto que, para todo morfismo finito $g:S''\to S'$, $f\circ g:S''\to S$ es un morfismo finito. Del mismo modo, bajo las condiciones de (ii), para que un morfismo $g:S'\to S$ sea finito, es necesario y suficiente que cada una

\oldpage[IV-4]{130}

de sus restricciones $g^{(\alpha)}:S'^{(\alpha)}=g^{-1}(S^{(\alpha)})\to S^{(\alpha)}$ lo sea, y, si se pone $S'_0=g^{-1}(S_0)$, $S'^{(\alpha)}_0=(g^{(\alpha)})^{-1}(S^{(\alpha)}_0)$, existe una correspondencia biunívoca entre las partes abiertas y cerradas $U$ (resp. $U_0$) de $S'$ (resp. $S'_0$) y las familias $(U^{(\alpha)})$ (resp. $U^{(\alpha)}_0$), donde $U^{(\alpha)}$ (resp. $U^{(\alpha)}_0$) es una parte abierta y cerrada de $S'^{(\alpha)}$ (resp. $S'^{(\alpha)}_0$), de donde se sigue la afirmación (ii).

\begin{remark}[18.5.7]
\label{IV.18.5.7}
Sean $S=\Spec(A)$ un esquema afín y $S_0$ un subesquema cerrado de $S$ definido por un ideal $\mathfrak J$ de $A$. Entonces, si el par $(S,S_0)$ es henseliano, el ideal $\mathfrak J$ está necesariamente contenido en el \emph{radical} de $A$. En efecto, si $\mathfrak m$ es un ideal maximal de $A$, $\mathfrak m$ debe pertenecer a $V(\mathfrak J)=S_0$, en virtud de (18.5.4.3); dicho de otro modo, debe tenerse $\mathfrak J\subset\mathfrak m$, de donde se sigue la conclusión. En particular, supongamos que $S_0$ se \emph{reduce a un punto}, es decir, que el ideal $\mathfrak J$ es \emph{maximal}; entonces $\mathfrak J$ debe ser el radical de $A$, dicho de otro modo, $A$ debe ser un anillo \emph{local}, y $S_0$ el único punto cerrado de $\Spec(A)$.
\end{remark}

\begin{definition}[18.5.8]
\label{IV.18.5.8}
Se dice que un anillo $A$ es \emph{henseliano} si es semilocal y si, denotando por $\mathfrak r$ el radical de $A$, el par $(\Spec(A),\Spec(A/\mathfrak r))$ es henseliano. Se llama \emph{esquema local henseliano} a un esquema isomorfo al espectro de un anillo local henseliano.
\end{definition}

\begin{proposition}[18.5.9]
\label{IV.18.5.9}
\begin{enumerate}
  \item[(i)] Para que un anillo semilocal $A$ sea henseliano, es necesario y suficiente que sea un producto directo de anillos locales henselianos.
  \item[(ii)] Para que un anillo local $A$ sea henseliano, es necesario y suficiente que toda $A$-álgebra finita $B$ sea isomorfa a un producto de anillos locales.
\end{enumerate}
\end{proposition}

(i) En efecto, la definición, aplicada a $S=\Spec(A)$ y $S_0=\Spec(A/\mathfrak r)$, muestra, puesto que $S_0$ es un conjunto finito discreto y cerrado en $S$, que $S$ es reunión de un número finito de partes abiertas y cerradas $S_i$ ($1\leq i\leq n$) disjuntas dos a dos, cada una de las cuales contiene exactamente uno de los ideales maximales $\mathfrak m_i$ de $A$; la conclusión se deduce de (18.5.6, (ii)) y de la observación (18.5.7).

(ii) Todo morfismo finito $S'\to\Spec(A)$ es de la forma $\Spec(B)\to\Spec(A)$, donde $B$ es una $A$-álgebra finita. Si $k$ es el cuerpo residual de $A$, $\Spec(B\otimes_A k)$ es el espectro de un anillo artiniano, y por tanto finito y discreto. Decir que el par $(\Spec(A),\Spec(k))$ es henseliano significa entonces que $B$ es un producto directo de anillos $A_i$ tales que $\Spec(A_i\otimes_A k)$ se reduce a un punto, es decir, que $A_i$ (que es una $A$-álgebra \emph{finita}) no debe tener más que un solo ideal maximal (Bourbaki, \emph{Alg. comm.}, cap. V, \textsection~2, n.\textsuperscript{o}~1, prop. 1).

El estudio de los anillos henselianos se reduce así esencialmente al de los anillos henselianos \emph{locales}.

\begin{proposition}[18.5.10]
\label{IV.18.5.10}
Si $A$ es un anillo henseliano, toda $A$-álgebra finita $B$ es un anillo henseliano (y, por tanto, un producto directo de anillos locales henselianos (18.5.9)).
\end{proposition}

En efecto, si $\mathfrak r'$ es el radical de $B$, la imagen inversa de $\mathfrak r'$ en $A$ es el radical $\mathfrak r$ de $A$, y todo ideal de $B$ situado sobre un ideal maximal de $A$ es un ideal maximal de $B$; por tanto, el conjunto $V(\mathfrak r')$ en $\Spec(B)$ es la imagen inversa del conjunto $V(\mathfrak r)$ en $\Spec(A)$. La proposición es entonces consecuencia de (18.5.6, (i)).

\begin{theorem}[18.5.11]
\label{IV.18.5.11}
Sean $A$ un anillo local y $\mathfrak m$ su ideal maximal. Las condiciones siguientes son equivalentes:
\begin{enumerate}
  \item[a)] $A$ es henseliano; dicho de otro modo, toda $A$-álgebra finita $B$ es isomorfa a un producto de anillos locales.

\oldpage[IV-4]{131}

  \item[a')] La condición $a)$ se satisface para todas las $A$-álgebras $B$ de la forma $A[T]/F.A[T]$, donde $F\in A[T]$ es un polinomio mónico.
  \item[b)] Sean $S=\Spec(A)$, $S_0=\Spec(k)$. Para todo morfismo étale $g:S'\to S$, si se pone $S'_0=S'\otimes_A k$, toda $S_0$-sección $u_0$ de $S'_0$ es la restricción de una $S$-sección $u$ de $S'$.
  \item[c)] Para todo morfismo $f:X\to S$, separado y localmente de tipo finito, y todo punto $x\in X$ tal que $f(x)$ sea igual al punto cerrado $s$ de $S$ y que $f$ sea cuasifinito en el punto $x$ (Err$_{\mathrm{III}}$, 20), $X$ es suma de dos preesquemas $X'$, $X''$ tales que $X'=\Spec(\sh{O}_{X,x})$ y que $f|X':X'\to S$ sea un morfismo finito.
  \item[c')] Para todo morfismo localmente de tipo finito $f:X\to S$, y todo punto $x\in X$ tal que $f$ sea cuasifinito en el punto $x$ y que $f(x)$ sea igual al punto cerrado $s$ de $S$, $\sh{O}_{X,x}$ es un álgebra finita sobre $\sh{O}_{S,s}=A$.
  \item[c'')] Para todo morfismo localmente de presentación finita $f:X\to S$, y todo punto $x\in X$ tal que $f$ sea cuasifinito en el punto $x$ y que $f(x)$ sea igual al punto cerrado $s$ de $S$, $\sh{O}_{X,x}$ es un álgebra finita y de presentación finita sobre $\sh{O}_{S,s}=A$.
\end{enumerate}
\end{theorem}

Observemos primero que la condición $c')$ (resp. $c'')$ es equivalente a la misma condición en la que se supone además que $f$ es \emph{separado}, pues la cuestión es local sobre $X$. Del mismo modo, la condición $b)$ es equivalente a la misma condición en la que se supone además que $g$ es \emph{separado}: en efecto, basta aplicar esta última a la restricción de $g$ a una vecindad abierta afín del punto $u_0(S_0)$ en $S'$. Nos limitaremos, pues, de ahora en adelante al caso en que, en $b)$, $c')$ y $c'')$, los morfismos dados son \emph{separados}.

El hecho de que $a)$ implique $b)$ y de que $b)$ implique $a')$ es un caso particular de (18.5.4). Demostremos además que $a')$ implica $a)$. Se trata de probar que, si $e_0$ es un idempotente de $C=B\otimes_A k$, existe un idempotente $e\in B$ cuya imagen canónica es $e_0$. Si $b$ es un elemento de $B$ cuya imagen en $C$ es $e_0$, la sub-$A$-álgebra $A[b]=B'$ de $B$ es finita, y la imagen canónica $C'$ de $B'\otimes_A k$ en $C$ contiene a $e_0$. Ahora bien, $B'\otimes_A k$ es una $k$-álgebra finita, y por tanto un producto directo de $k$-álgebras locales finitas; por consiguiente, $e_0$ es la imagen en $C'$ de un idempotente $e'_0$ de $B'\otimes_A k$. La cuestión se reduce así al caso en que $B$ es monógena y, por tanto, isomorfa a una $A$-álgebra cociente de un álgebra de la forma $A[T]/F.A[T]$, donde $F$ es un polinomio mónico. Ahora bien, en virtud de $a')$, $A[T]/F.A[T]$ es un producto directo de anillos locales; lo mismo sucede entonces con todas sus álgebras cocientes, lo que termina de demostrar la existencia del idempotente $e$.

Es inmediato que $c)$ implica $a)$, como se ve razonando por inducción sobre el número de ideales maximales del anillo semilocal $B$. Para ver que $a)$ implica $c)$, puede uno limitarse, en virtud de (13.1.4), al caso en que $f$ es un morfismo afín y cuasifinito. Entonces, aplicando el «Main theorem» (8.12.8), $f$ puede escribirse como un morfismo compuesto
\[
  X\xrightarrow{j}Y\xrightarrow{g}S,
\]
donde $g$ es un morfismo finito y $j$ una inmersión abierta. Como $Y=\Spec(B)$, donde $B$ es una $A$-álgebra finita, se deduce de $a)$ que $B$ es un producto directo de anillos locales, que son evidentemente $A$-álgebras finitas, y $\sh{O}_{X,x}$ se identifica con uno de estos anillos locales puesto que $f(x)=s$; además, todo abierto de $Y$ que contiene a $x$ contiene necesariamente a $\Spec(\sh{O}_{X,x})$.

Es trivial que $c)$ implica $c')$, en virtud de la observación inicial. Demostremos que $c')$ implica $c'')$. Supongamos, en efecto, verificada $c')$, y probemos que, bajo las condiciones

\oldpage[IV-4]{132}

de $c'')$, el conjunto $Z=\Spec(\sh{O}_{X,x})$ se identifica entonces con una parte \emph{abierta y cerrada} de $X$, lo que establecerá $c'')$ (I, 2.4.2). En primer lugar, el morfismo compuesto $Z\xrightarrow{j}X\xrightarrow{f}S$, donde $j$ es el morfismo canónico (I, 2.4.1), es finito y de presentación finita por hipótesis, y, como $f$ es separado y localmente de presentación finita, $j$ es también un morfismo finito y de presentación finita ((II, 6.1.5) y (I, 4.3)); por consiguiente, $j$ es un morfismo cerrado (II, 6.1.10), lo que demuestra que $Z$ es cerrado en $X$. Se deduce entonces de (I, 2.4.2) y (I, 4.2.2) que $j$ es una inmersión cerrada. Pero si $\mathscr J$ es el ideal de $\sh{O}_X$ que define $Z$, se tiene por hipótesis $\mathscr J_x=0$ y, por tanto, también $\mathscr J_z=0$ en todo punto $z$ de una vecindad abierta $V$ de $x$ en $X$, puesto que, por hipótesis, $\mathscr J$ es un $\sh{O}_X$-módulo cuasicoherente de tipo finito ((1.4.7) y ($0_{\mathrm I}$, 5.2.2)). Ahora bien, una vecindad tal contiene a $Z$, por lo que $Z$ es abierto en $X$, puesto que $\mathscr J|V=0$.

Por último, $c'')$ implica $a')$: en efecto, si $B=A[T]/F.A[T]$, el morfismo $X=\Spec(B)\to\Spec(A)=Y$ es finito y de presentación finita, y la demostración anterior muestra que $B$ es el producto directo de las $A$-álgebras finitas $\sh{O}_{X,x_i}$, donde los $x_i$ son los puntos de la fibra del punto cerrado de $Y$.

\begin{corollary}[18.5.12]
\label{IV.18.5.12}
Sean $A$ un anillo semilocal y $\mathfrak r$ su radical; pongamos $S=\Spec(A)$, $S_0=\Spec(A/\mathfrak r)$. Para que $A$ sea henseliano, es necesario y suficiente que, para todo morfismo finito $f:X\to S$ y todo morfismo étale y separado $g:Y\to S$, si se pone $X_0=X\times_S S_0$ y $Y_0=Y\times_S S_0$, la aplicación canónica
\[
  \Hom_S(X,Y)\longrightarrow\Hom_{S_0}(X_0,Y_0)
\]
sea biyectiva.
\end{corollary}

La suficiencia de la condición se deduce de la equivalencia de $a)$ y $b)$ en (18.5.11), aplicando esta condición al caso en que $f=1_S$. Para ver que la condición es necesaria, observemos que, si $A$ es henseliano, el par $(X,X_0)$ es henseliano por (18.5.6, (i)); además, $\Hom_S(X,Y)=\Gamma(X\times_S Y/X)$ y $\Hom_{S_0}(X_0,Y_0)=\Gamma(X_0\times_{S_0}Y_0/X_0)$; como $X\times_S Y$ es étale y separado sobre $X$, la conclusión se deduce del hecho de que $a)$ implica $b)$ en (18.5.4).

\begin{remark}[18.5.13]
\label{IV.18.5.13}
Las condiciones equivalentes $a)$, $a')$, $b)$ del teorema (18.5.11) son también equivalentes a la siguiente («lema de Hensel»):
\begin{enumerate}
  \item[a'')] Para todo polinomio mónico $F\in A[T]$, de imagen canónica $F_0\in k[T]$, y toda descomposición $F_0=G_0H_0$ de $F_0$ en un producto de dos polinomios mónicos primos entre sí $G_0$, $H_0$ de $k[T]$, existe un único par $(G,H)$ de polinomios mónicos de $A[T]$ que posee las propiedades siguientes: $G_0$ y $H_0$ son las imágenes canónicas respectivas de $G$ y $H$, se tiene $F=GH$, y el ideal de $A[T]$ engendrado por $G$ y $H$ es igual a $A[T]$.
\end{enumerate}
\end{remark}

Estableceremos primero el lema siguiente:

\begin{lemma}[18.5.13.1]
\label{IV.18.5.13.1}
Sean $A$ un anillo local de cuerpo residual $k$, $F\in A[T]$ un polinomio mónico y $B$ la $A$-álgebra $A[T]/F.A[T]$. Existe una correspondencia canónica entre las descomposiciones de $B$ en producto directo de dos $A$-álgebras cocientes $B'$, $B''$ y las descomposiciones $F=GH$ de $F$ en producto de dos polinomios mónicos $G$, $H$ de $A[T]$, tales que el ideal engendrado por $G$ y $H$ sea igual a $A[T]$; las álgebras cocientes $B'$, $B''$ correspondientes a tal par de polinomios $G$, $H$ son respectivamente $A[T]/H.A[T]$ y $A[T]/G.A[T]$.
\end{lemma}

\oldpage[IV-4]{133}

Si $F=GH$ y si $G$ y $H$ engendran el ideal $A[T]$, existen dos polinomios $P$, $Q$ de $A[T]$ tales que $1=PG+QH$. Se deduce que la intersección de los ideales principales $\mathfrak a=G.A[T]$ y $\mathfrak b=H.A[T]$ es igual a $\mathfrak c=F.A[T]$: en efecto, si $R\in G.A[T]\cap H.A[T]$, puede escribirse $R=PRG+QRH$; ahora bien, $RH$ (resp. $RG$) es un múltiplo de $F$, puesto que $R$ es un múltiplo de $G$ (resp. $H$), y por tanto $R\in F.A[T]$. Como $A[T]=\mathfrak a+\mathfrak b$, $A[T]/\mathfrak c$ es suma directa de los ideales $\mathfrak a/(\mathfrak a\cap\mathfrak b)$ y $\mathfrak b/(\mathfrak a\cap\mathfrak b)$, canónicamente isomorfos, respectivamente, a $A[T]/\mathfrak b$ y $A[T]/\mathfrak a$.

Recíprocamente, supongamos dada una descomposición de $B$ en producto directo de dos $A$-álgebras $B'$, $B''$, que se identifican canónicamente con dos ideales $e'B$, $e''B$ de $B$, correspondientes a una descomposición $1=e'+e''$ de $1$ en idempotentes ortogonales $e'$, $e''$ de $B$. Pongamos también $B_0=B\otimes_A k=k[T]/F_0.k[T]$, donde $F_0$ es la imagen canónica de $F$ en $k[T]$, del mismo grado $n$ que $F$; si $e'_0$, $e''_0$ son las imágenes canónicas de $e'$, $e''$ en $B_0$, son dos idempotentes ortogonales tales que $1=e'_0+e''_0$, y $B_0$ es, por tanto, producto directo de $B'_0=e'_0B_0$ y $B''_0=e''_0B_0$. Sean $t$ y $t_0$ las imágenes canónicas de $T$ en $B$ y $B_0$; como $B'_0$ (resp. $B''_0$) es una $k$-álgebra finita engendrada por $t'_0=e'_0t_0$ (resp. $t''_0=e''_0t_0$), admite una base de la forma $\{e'_0,t'_0,t_0'{}^2,\ldots,t_0'{}^{s-1}\}$ (resp. $\{e''_0,t''_0,t_0''{}^2,\ldots,t_0''{}^{r-1}\}$), con $r+s=n$. Por otra parte, al ser $B$ un $A$-módulo libre (de base $\{1,t,\ldots,t^{n-1}\}$), $B'$ y $B''$ son $A$-módulos proyectivos, y por tanto libres puesto que $A$ es un anillo local (Bourbaki, \emph{Alg. comm.}, cap. II, \textsection~5, n.\textsuperscript{o}~3, cor. de la prop. 5); se deduce entonces de lo anterior y de Bourbaki, \emph{Alg. comm.}, cap. II, \textsection~3, n.\textsuperscript{o}~3, prop. 5, que, si se pone $t'=e't$, $t''=e''t$, $\{e',t',\ldots,t'^{s-1}\}$ (resp. $\{e'',t'',\ldots,t''^{r-1}\}$) es una base del $A$-módulo $B'$ (resp. $B''$). Existe entonces un polinomio mónico $H$ (resp. $G$) de grado $s$ (resp. $r$) de $A[T]$ tal que $e'H(t')=0$ y $e''G(t'')=0$; como $t^h=t'^h+t''^h$ para todo entero $h\geq1$, y $t'^h=e't'^h$, $t''^h=e''t''^h$, se tiene también $G(t)=e'G(t')$ y $H(t)=e''H(t'')$, de donde $G(t)H(t)=0$; se concluye que el polinomio $G(T)H(T)$ es divisible por $F(T)$; pero, como los grados de estos dos polinomios mónicos son iguales, se tiene $GH=F$. Además, $B'$ (resp. $B''$) es isomorfa a $A[T]/H.A[T]$ (resp. $A[T]/G.A[T]$). Por último, existen dos polinomios $R$, $S$ de $A[T]$ tales que $e'=R(t)$ y $e''=S(t)$; como $R(t)=e'R(t')+e''R(t'')$, se tiene necesariamente $e''R(t'')=0$, y del mismo modo $e'S(t')=0$, de suerte que, por definición de $G$ y $H$, $R=QH$ y $S=PG$, donde $P$, $Q$ pertenecen a $A[T]$; la relación $1=R(t)+S(t)$ en $B$ da entonces, por definición, $PG+QH=1+LF$ para cierto polinomio $L\in A[T]$, y, como $F=GH$, esto demuestra que el ideal engendrado por $G$ y $H$ es $A[T]$ y termina la demostración del lema.

Establecido este lema, basta aplicarlo, por una parte, al anillo local $A$ y, por otra, al cuerpo $k$, para ver inmediatamente que las condiciones $a')$ y $a'')$ son equivalentes.

\begin{proposition}[18.5.14]
\label{IV.18.5.14}
Todo anillo semilocal $A$, separado y completo para la topología $\mathfrak r$-preádica (donde $\mathfrak r$ es el radical de $A$), es henseliano.
\end{proposition}

En efecto, $A$ es un producto directo de anillos locales separados y completos (Bourbaki, \emph{Alg. comm.}, cap. III, \textsection~2, n.\textsuperscript{o}~13, cor. de la prop. 19), por lo que la cuestión se reduce al caso en que $A$ es un anillo local. Verifiquemos el criterio $a')$ de (18.5.11). Como $B$ es un $A$-módulo libre de tipo finito, es evidentemente separado y completo para la topología $\mathfrak r$-preádica,

\oldpage[IV-4]{134}

que es también la topología $\mathfrak s$-preádica, donde $\mathfrak s$ es el radical del anillo semilocal $B$, pues $B/\mathfrak rB$ es un anillo artiniano de radical $\mathfrak s/\mathfrak rB$. Se sabe entonces (Bourbaki, \emph{Alg. comm.}, cap. III, \textsection~2, n.\textsuperscript{o}~13, cor. de la prop. 19) que $B$ es un producto directo de anillos locales.

\begin{proposition}[18.5.15]
\label{IV.18.5.15}
Sean $A$ un anillo henseliano y $\mathfrak r$ su radical. Entonces el funtor $B\longmapsto B/\mathfrak rB$ es una equivalencia de la categoría de las $A$-álgebras finitas y étales con la categoría de las $(A/\mathfrak r)$-álgebras finitas y étales.
\end{proposition}

El hecho de que el funtor del enunciado sea plenamente fiel es un caso particular de (18.5.12). Para demostrar que este funtor es una equivalencia de categorías, puede uno limitarse al caso en que $A$ es local; basta entonces aplicar (18.1.1) (para los morfismos étales) a $S=\Spec(A)$ y $S_0=\Spec(A/\mathfrak r)$, reducido a un solo punto.

\begin{remarks}[18.5.16]
\label{IV.18.5.16}
\begin{enumerate}
  \item[(i)] Ignoramos si la proposición (18.5.15) se generaliza a un par henseliano $(S,S_0)$, incluso cuando $S$ es afín y noetheriano.
  \item[(ii)] Sean $A$ un anillo noetheriano y $\mathfrak J$ un ideal de $A$ tal que $A$ sea \emph{separado y completo} para la topología $\mathfrak J$-preádica. Entonces el par $(\Spec(A),\Spec(A/\mathfrak J))$ es henseliano: en efecto, para toda $A$-álgebra finita $B$, el anillo $B$ es separado y completo para la topología $\mathfrak J$-preádica ($0_{\mathrm I}$, 7.3.6). Sustituyendo $B$ por $A$ y $\mathfrak J B$ por $\mathfrak J$, todo se reduce entonces a ver que la aplicación que a todo idempotente de $A$ hace corresponder su clase mód. $\mathfrak J$ es \emph{biyectiva}. Ahora bien, se tiene $A=\varprojlim (A/\mathfrak J^n)$. Denotemos por $\operatorname{Idem}(A)$ el conjunto de los idempotentes de $A$ y, para todo homomorfismo de anillos $\varphi:A\to B$, sea $\operatorname{Idem}(\varphi)$ la aplicación de $\operatorname{Idem}(A)$ en $\operatorname{Idem}(B)$ que es la restricción de $\varphi$; de la definición del límite proyectivo se deduce que se tiene
  \[
    \operatorname{Idem}(A)=\varprojlim_n\operatorname{Idem}(A/\mathfrak J^n)
  \]
  para las aplicaciones $\psi_{nm}:\operatorname{Idem}(A/\mathfrak J^m)\to\operatorname{Idem}(A/\mathfrak J^n)$ que son las restricciones de las aplicaciones canónicas $A/\mathfrak J^m\to A/\mathfrak J^n$. Pero, puesto que $\Spec(A/\mathfrak J^n)\to\Spec(A/\mathfrak J^m)$ es un homeomorfismo, las $\psi_{nm}$ son biyecciones (como se ha visto en la demostración de (18.5.3)); esto demuestra nuestra afirmación.
\end{enumerate}
\end{remarks}

\begin{theorem}[18.5.17]
\label{IV.18.5.17}
Sean $A$ un anillo local henseliano, $S=\Spec(A)$ y $s$ el punto cerrado de $S$. Para todo morfismo liso $f:X\to S$, la aplicación canónica
\[
  \Gamma(X/S)\longrightarrow\Gamma(X_s/k(s))
\]
(donde $X_s=f^{-1}(s)$) es sobreyectiva.
\end{theorem}

El dato de una $k(s)$-sección de $X_s$ equivale al de un punto $x\in X$ situado sobre $s$ y racional sobre $k(s)$, y se trata de demostrar que existe una $S$-sección $u:S\to X$ tal que $u(s)=x$. Teniendo en cuenta (17.16.3, (i)), puede suponerse que $f$ es étale. La conclusión se deduce entonces del criterio (18.5.11, $b)$). (El lector observará que, en virtud de este criterio, la validez de (18.5.17) para un anillo local dado $A$ es necesaria y suficiente para que $A$ sea henseliano).

\begin{remark}[18.5.18]
\label{IV.18.5.18}
Las condiciones de (18.5.11) son también equivalentes a la siguiente:
\begin{enumerate}
  \item[d)] Para todo morfismo localmente de tipo finito $f:X\to S$ y todo punto $x\in X$ tal que $f(x)$ sea el punto cerrado $s$ de $S$ y que $\sh{O}_{X,x}/\mathfrak m_s\sh{O}_{X,x}$ sea un cuerpo canónicamente isomorfo a $k(s)=k$, existe una vecindad abierta $U$ de $x$ en $X$ tal que $f|U$ sea una inmersión cerrada.
\end{enumerate}
\end{remark}

\oldpage[IV-4]{135}

Es inmediato que $d)$ implica la condición $b)$ de (18.5.11), pues una inmersión cerrada étale es una inmersión abierta (18.9.1). Recíprocamente, supongamos verificadas las condiciones de (18.5.11) y demostremos $d)$. La hipótesis de $d)$ implica que $f$ es cuasifinito en el punto $x$ (Err$_{\mathrm{III}}$, 20), en virtud de (13.1.4). Por tanto, en virtud de la condición $c)$ de (18.5.11), $\sh{O}_{X,x}$ es una $A$-álgebra finita y $\Spec(\sh{O}_{X,x})$ una vecindad abierta $U$ de $x$ en $X$. Además, como $\sh{O}_{X,x}/\mathfrak m_s\sh{O}_{X,x}$ es isomorfo a $k=A/\mathfrak m_s$, el lema de Nakayama demuestra que el homomorfismo $A\to\sh{O}_{X,x}$ es sobreyectivo y, por tanto, que $f|U$ es una inmersión cerrada.

\begin{proposition}[18.5.19]
\label{IV.18.5.19}
Sean $A$ un anillo local noetheriano y henseliano, $S=\Spec(A)$, $s$ el punto cerrado de $S$ y $f:X\to S$ un morfismo propio; pongamos $X_s=f^{-1}(s)$. Entonces la aplicación $Y\longmapsto Y\cap X_s$ es una biyección del conjunto de las componentes conexas de $X$ sobre el conjunto de las componentes conexas de $X_s$.
\end{proposition}

Considerando un subpreesquema cerrado de $X$ cuyo espacio subyacente es una componente conexa de $X$, la cuestión se reduce a demostrar que, si $X$ es \emph{conexo y no vacío}, entonces $X_s$ es \emph{conexo y no vacío}. El hecho de que $X_s$ no sea vacío se deduce de (II, 7.2.1); para demostrar que $X_s$ es conexo, razonemos por reducción al absurdo, considerando la factorización de Stein $f:X\xrightarrow{f'}S'\xrightarrow{g}S$ del morfismo propio $f$ (III, 4.3.3); por hipótesis, el conjunto discreto finito $g^{-1}(s)$ contendría al menos dos puntos. Puesto que $A$ es henseliano y $g$ separado y de tipo finito, se deduciría entonces de (18.5.11, $c)$ que $S'$ sería suma de dos preesquemas no vacíos $S'_1$, $S'_2$ (pues la intersección de uno de ellos con $g^{-1}(s)$ se reduce a un solo punto). Puesto que $f'$ es sobreyectivo (III, 4.3.1), se concluiría que $X$ es suma de dos preesquemas no vacíos, contrariamente a la hipótesis.

\subsection{Henselización.}
\label{subsection:IV.18.6}

\begin{env}[18.6.1]
\label{IV.18.6.1}
Dado un anillo local $A$, diremos que una $A$-álgebra local $B$ es \emph{esencialmente étale} si existen una $A$-álgebra étale $C$ y un ideal primo $\mathfrak n$ de $C$ tales que $B$ sea $A$-isomorfa a $C_{\mathfrak n}$ y que el homomorfismo compuesto $A\to C\to C_{\mathfrak n}$ sea \emph{local} (dicho de otro modo, que $\mathfrak n$ esté sobre el ideal maximal $\mathfrak m$ de $A$); esto implica (17.6.1) que $\mathfrak n$ es el \emph{único} ideal primo de $B$ sobre $\mathfrak m$, puesto que $B$ es un anillo local; por tanto, si $B$, $B'$ son dos $A$-álgebras locales esencialmente étales, todo $A$-homomorfismo $B\to B'$ es \emph{local}.

Si $A'$ es una $A$-álgebra local esencialmente étale y $A''$ una $A'$-álgebra local esencialmente étale, entonces $A''$ es una $A$-álgebra local esencialmente étale. En efecto, se tiene por hipótesis $A'=B_{\mathfrak n}$, donde $B$ es una $A$-álgebra étale y $\mathfrak n$ un ideal primo de $B$ sobre el ideal maximal de $A$, y $A''=B'_{\mathfrak n'}$, donde $B'$ es una $A'$-álgebra étale y $\mathfrak n'$ un ideal primo de $B'$ sobre el ideal $\mathfrak nB_{\mathfrak n}$. Si se pone $S=B-\mathfrak n$, de modo que $A'=S^{-1}B$, $B'$ es de la forma $S^{-1}C$, donde $C$ es una $B$-álgebra de presentación finita. Por consiguiente, $C$ es una $A$-álgebra de presentación finita y $A''$ es de la forma $C_{\mathfrak r}$, donde $\mathfrak r$ es un ideal primo de $C$ sobre el ideal maximal de $A$. Como $A'$ es una $A$-álgebra formalmente étale y $A''$ una $A'$-álgebra formalmente étale, $A''$ es una $A$-álgebra formalmente

\oldpage[IV-4]{136}

étale (para las topologías discretas (17.1.3)); por tanto, el morfismo $\Spec(C)\to\Spec(A)$ es étale en el punto $\mathfrak r$ (17.6.1) y, por consiguiente, existe un elemento $g\in C$ tal que $C_g$ sea una $A$-álgebra étale, lo que demuestra que $A''$ es una $A$-álgebra esencialmente étale.

Dado un anillo local $A$, existe un \emph{conjunto} $\mathfrak E$ de $A$-álgebras locales esencialmente étales tal que toda $A$-álgebra local esencialmente étale sea $A$-isomorfa a un álgebra perteneciente a $\mathfrak E$. Para verlo, basta evidentemente observar que existe un conjunto $\mathfrak F$ de $A$-álgebras de tipo finito tal que toda $A$-álgebra de tipo finito sea isomorfa a un álgebra perteneciente a $\mathfrak F$; puede tomarse $\mathfrak F$ igual al conjunto de los cocientes de las álgebras de polinomios $A[T_1,\ldots,T_n]$ ($n\in\mathbf N$).

En lo que sigue, si $A$ y $B$ son dos anillos locales, denotaremos por $\Hom.\operatorname{loc}(A,B)$ el conjunto de los homomorfismos \emph{locales} de $A$ en $B$.
\end{env}

\begin{lemma}[18.6.2]
\label{IV.18.6.2}
Sean $A$, $A'$ dos anillos locales, $k$, $k'$ sus cuerpos residuales respectivos, $\varphi:A\to A'$ un homomorfismo (local) que hace de $A'$ una $A$-álgebra esencialmente étale (18.6.1) y tal que el homomorfismo correspondiente $k\to k'$ sea biyectivo. Entonces, para todo anillo local henseliano $B$, la aplicación canónica
\[
  \Hom(\varphi,1_B):\Hom.\operatorname{loc}(A',B)\longrightarrow\Hom.\operatorname{loc}(A,B)
\]
es biyectiva.
\end{lemma}

Pongamos $S=\Spec(A)$, $Y=\Spec(B)$, y sean $s$ e $y$ los puntos cerrados de $S$ e $Y$, respectivamente. Por hipótesis, $A'$ es isomorfa a un anillo local $\sh{O}_{X,x}$ de un $S$-esquema $X$ étale sobre $S$, en un punto $x\in X$ situado sobre $s$. Supongamos dado un homomorfismo local $\psi:A\to B$, que hace de $Y$ un $S$-esquema; se trata de ver que existe \emph{uno y solo un} $S$-morfismo $f:Y\to X$ tal que $f(y)=x$. Pongamos $X'=X\times_S Y$ y observemos que, puesto que $k(x)=k(s)$, existe un solo punto $x'\in X'$ situado sobre $x$ e $y$, y que $k(x')=k(y)$. Hay que demostrar que existe una única $Y$-sección $f'$ de $X'$ tal que $f'(y)=x'$. Ahora bien, el morfismo $g:X'\to Y$ es étale y separado, y la fibra $X'_0=g^{-1}(y)$ tiene por anillo local en el punto $x'$ el cuerpo $k(x')=k(y)$. Si se pone $Y_0=\Spec(k(y))$, existe entonces una única $Y_0$-sección $f'_0$ de $X'_0$ tal que $f'_0(y)=x'$, y la conclusión se deduce de que $B$ se supone henseliano y de (18.5.11, $b)$).

Diremos que una $A$-álgebra local $A'$ que verifica las condiciones de (18.6.2) es \emph{estrictamente esencialmente étale}. Observemos que el criterio (18.5.11, $b)$) significa que, \emph{para que $A$ sea henseliano, es necesario y suficiente que toda $A$-álgebra estrictamente esencialmente étale sea $A$-isomorfa a $A$}.

\begin{lemma}[18.6.3]
\label{IV.18.6.3}
Sean $A$ un anillo local y $A_1$, $A_2$ dos $A$-álgebras locales estrictamente esencialmente étales.
\begin{enumerate}
  \item[(i)] Existe a lo sumo un $A$-homomorfismo (necesariamente local) de $A_1$ en $A_2$.
  \item[(ii)] Existe una $A$-álgebra local estrictamente esencialmente étale $A_3$ y dos $A$-homomorfismos $A_1\to A_3$, $A_2\to A_3$.
\end{enumerate}
\end{lemma}

Pongamos $S=\Spec(A)$; por hipótesis hay dos $S$-esquemas étales sobre $S$, $X_1$, $X_2$, y dos puntos $x_1\in X_1$, $x_2\in X_2$ situados sobre el punto cerrado $s$ de $S$ y tales que $A_1=\sh{O}_{X_1,x_1}$, $A_2=\sh{O}_{X_2,x_2}$. Pongamos $X_3=X_1\times_S X_2$; las hipótesis $k(x_1)=k(x_2)=k(s)$ implican que existe un solo punto $x_3\in X_3$ situado sobre $x_1$ y $x_2$, y que $k(x_3)=k(s)$ (I, 3.4.9). Además, $X_3$ es étale sobre $S$ (17.3.3), por lo que $A_3=\sh{O}_{X_3,x_3}$ verifica las condiciones de (ii). Por otra parte,

\oldpage[IV-4]{137}

se ha visto que un $A$-homomorfismo de $A_1$ en $A_2$ es necesariamente local; corresponde a un $S$-morfismo $f$ de $X'_2=\Spec(A_2)$ en $X_1$ tal que $f(x_2)=x_1$, o también, poniendo $X'_3=X_1\times_S X'_2$, a una $X'_2$-sección $f'$ de $X'_3$ tal que $f'(x_2)=x_3$. Como $k(x_3)=k(x_2)$ y $X'_2$ es conexo, la unicidad de $f'$ se deduce de (17.4.9), de donde se sigue (i).

\begin{env}[18.6.4]
\label{IV.18.6.4}
Denotemos por $\mathfrak S$ el subconjunto del conjunto $\mathfrak E$ definido en (18.6.1) formado por las $A$-álgebras estrictamente esencialmente étales pertenecientes a $\mathfrak E$. Se deduce de (18.6.3) que la relación «existe un $A$-homomorfismo de $A_1$ en $A_2$» es una relación de \emph{preorden} en $\mathfrak S$ y hace de $\mathfrak S$ un conjunto \emph{filtrante creciente}. Indexando $\mathfrak S$ por sí mismo mediante la aplicación identidad, se deduce de (18.6.3, (i)) que, si $\lambda\leq\mu$ en $\mathfrak S$, existe un $A$-homomorfismo \emph{único} $\varphi_{\mu\lambda}:A_\lambda\to A_\mu$ y $(A_\lambda,\varphi_{\mu\lambda})$ es evidentemente un \emph{sistema inductivo} de $A$-álgebras locales, siendo los $\varphi_{\mu\lambda}$ homomorfismos locales. Además, se deduce de (17.3.5) que, para $\lambda\leq\mu$, $A_\mu$ es una $A_\lambda$-álgebra estrictamente esencialmente étale.
\end{env}

\begin{definition}[18.6.5]
\label{IV.18.6.5}
Se llama \emph{henselizado} de un anillo local $A$, y se denota por ${}^h\!A$, a la $A$-álgebra que es el límite inductivo del sistema inductivo $(A_\lambda,\varphi_{\mu\lambda})$ definido en (18.6.4).
\end{definition}

Esta definición solo depende en apariencia de la elección de $\mathfrak E$; si $\mathfrak E'$ es otro conjunto de $A$-álgebras locales esencialmente étales que tiene la misma propiedad que $\mathfrak E$, y $\mathfrak S'$ el subconjunto de $\mathfrak E'$ formado por las $A$-álgebras estrictamente esencialmente étales, se deduce de (18.6.3, (ii)) que, para las relaciones de preorden consideradas, $\mathfrak S$ y $\mathfrak S'$ son \emph{cofinales} en $\mathfrak S''=\mathfrak S\cup\mathfrak S'$ y, por tanto, dan el mismo límite inductivo salvo isomorfismo. Veremos, por lo demás, más adelante (18.6.6, (i) y (ii)) que ${}^h\!A$ y el homomorfismo canónico $A\to{}^h\!A$ forman una solución de un problema universal y, por consiguiente, están determinados salvo isomorfismo único.

Obsérvese que, si $A$ es un cuerpo, se tiene evidentemente ${}^h\!A=A$. En el caso general, ${}^h\!A$ es también el henselizado de todas las $A$-álgebras estrictamente esencialmente étales $A_\lambda$ ($\lambda\in\mathfrak S$), pues, si $B$ es una $A_\lambda$-álgebra estrictamente esencialmente étale, $B$ es también una $A$-álgebra estrictamente esencialmente étale (18.6.1), y por tanto (salvo un $A$-isomorfismo) una de las $A_\mu$ para $\mu\geq\lambda$, y las $A_\mu$ tales que $\mu\geq\lambda$ y que $A_\mu$ sea una $A_\lambda$-álgebra estrictamente esencialmente étale forman un conjunto cofinal en el conjunto preordenado de las $A_\lambda$, en virtud de (18.6.1) y (18.6.3).

\begin{theorem}[18.6.6]
\label{IV.18.6.6}
Sean $A$ un anillo local y ${}^h\!A$ su henselizado.
\begin{enumerate}
  \item[(i)] ${}^h\!A$ es un anillo local henseliano y el homomorfismo estructural $A\to{}^h\!A$ es local.
  \item[(ii)] Para todo anillo local henseliano $B$, la aplicación canónica
  \[
    \Hom.\operatorname{loc}({}^h\!A,B)\longrightarrow\Hom.\operatorname{loc}(A,B)
  \]
  es biyectiva.
  \item[(iii)] ${}^h\!A$ es un $A$-módulo fielmente plano y, si $\mathfrak m$ es el ideal maximal de $A$, $\mathfrak m.{}^h\!A$ es el ideal maximal de ${}^h\!A$, y el homomorfismo $A/\mathfrak mA\to{}^h\!A/\mathfrak m.{}^h\!A$ de los cuerpos residuales es biyectivo.
  \item[(iv)] Si $\widehat A$ y $({}^h\!A)\widehat{\phantom{x}}$ son los completados separados de los anillos locales $A$ y ${}^h\!A$, el homomorfismo $\widehat A\to({}^h\!A)\widehat{\phantom{x}}$ deducido del homomorfismo estructural $A\to{}^h\!A$ por compleción es biyectivo.
  \item[(v)] Para que ${}^h\!A$ sea noetheriano, es necesario y suficiente que $A$ lo sea.
  \item[(vi)] Si $A$ es henseliano, el homomorfismo canónico $A\to{}^h\!A$ es biyectivo.
\end{enumerate}
\end{theorem}

\oldpage[IV-4]{138}

Sea $\mathfrak m_\lambda$ el ideal maximal de $A_\lambda$; del hecho de que, en (17.6.1), $a)$ implica $c')$, se deduce que, para $\lambda\leq\mu$, se tiene $\mathfrak m_\mu=\mathfrak m_\lambda A_\mu$, que el homomorfismo $A_\lambda/\mathfrak m_\lambda\to A_\mu/\mathfrak m_\mu$ es biyectivo y que $A_\mu$ es un $A_\lambda$-módulo plano. El hecho de que ${}^h\!A$ sea local, de que el homomorfismo $A\to{}^h\!A$ sea local, la afirmación (iii) y la suficiencia de (v) se deducen entonces de ($0_{\mathrm{III}}$, 10.3.1.3); la necesidad de (v) se deduce de que ${}^h\!A$ es un $A$-módulo fielmente plano ($0_{\mathrm I}$, 6.5.2).

Para demostrar que ${}^h\!A$ es henseliano, apliquemos el criterio (18.5.11, $b)$). Pongamos, pues, $S=\Spec({}^h\!A)$, $S_0=\Spec({}^h\!A/\mathfrak m.{}^h\!A)$, y sea $g:S'\to S$ un morfismo étale; pongamos $S'_0=g^{-1}(S_0)$, y sea $f_0:S_0\to S'_0$ una $S_0$-sección de $S'_0$. Razonando como en (18.5.4), puede suponerse que $g$ es de presentación finita. Se deduce entonces de (8.8.2) y (17.7.5) que existen un índice $\lambda\in\mathfrak S$, un morfismo étale $g_\lambda:S'^{(\lambda)}\to S^{(\lambda)}=\Spec(A_\lambda)$ y, si se pone $S_0^{(\lambda)}=\Spec(A_\lambda/\mathfrak mA_\lambda)$ y $S_0'^{(\lambda)}=g_\lambda^{-1}(S_0^{(\lambda)})$, una $S_0^{(\lambda)}$-sección $f_0^{(\lambda)}:S_0^{(\lambda)}\to S_0'^{(\lambda)}$ tales que
\[
  S'=S'^{(\lambda)}\times_{S^{(\lambda)}}S,\qquad
  g=g_\lambda\times1,\qquad
  f_0=f_0^{(\lambda)}\times1.
\]
Sean $s$ el punto cerrado de $S$, $x=f_0(s)$ y $x_\lambda$ la proyección de $x$ en $S'^{(\lambda)}$; como $x_\lambda$ está sobre el punto cerrado $s_\lambda$ de $S^{(\lambda)}$, el anillo local $C_\lambda$ de $S'^{(\lambda)}$ en el punto $x_\lambda$ es una $A_\lambda$-álgebra esencialmente étale; además, como $f_0^{(\lambda)}(s_\lambda)=x_\lambda$, se tiene $k(x_\lambda)=k(s_\lambda)$; dicho de otro modo, $C_\lambda$ es una $A_\lambda$-álgebra local estrictamente esencialmente étale y, por consiguiente, (18.6.1) es $A_\lambda$-isomorfa a una $A$-álgebra $A_\mu$ con $\mu\geq\lambda$. Existe entonces un $A_\lambda$-homomorfismo $C_\lambda\to{}^h\!A$, es decir, un $S^{(\lambda)}$-morfismo $h:S\to S'^{(\lambda)}$ tal que $h(s)=x_\lambda$ y, por consiguiente, existe efectivamente una $S$-sección $f$ de $S'$ tal que $f(s)=x$, lo que termina de demostrar (i).

Para demostrar (ii), basta observar que se tiene
\[
  \Hom.\operatorname{loc}({}^h\!A,B)=\varprojlim_\lambda\Hom.\operatorname{loc}(A_\lambda,B)
\]
y que, por (18.6.2), los homomorfismos canónicos $\Hom.\operatorname{loc}(A_\lambda,B)\leftarrow\Hom.\operatorname{loc}(A,B)$ son biyectivos.

Para demostrar (iv), observemos que, para todo $\lambda$ y todo entero $n>0$, $\mathfrak m_\lambda^n=\mathfrak m^nA_\lambda$, y
\[
  (\mathfrak m.{}^h\!A)^n=\mathfrak m^n.{}^h\!A=\varinjlim_\lambda\mathfrak m_\lambda^n.
\]
Se concluye, por la exactitud del funtor $\varinjlim$ en la categoría de los $A$-módulos, que
\[
  {}^h\!A/(\mathfrak m.{}^h\!A)^n=\varinjlim_\lambda(A_\lambda/\mathfrak m_\lambda^n),
\]
y, por consiguiente, basta demostrar que, para todo entero $n$ y todo índice $\lambda$, el homomorfismo $A/\mathfrak m^n\to A_\lambda/\mathfrak m_\lambda^n$ es biyectivo. Ahora bien, esto es cierto por hipótesis para $n=1$; por otra parte, puesto que $A_\lambda$ es un $A$-módulo plano, se tiene
\[
  \mathfrak m_\lambda^n/\mathfrak m_\lambda^{n+1}
  =(\mathfrak m^n/\mathfrak m^{n+1})\otimes_A A_\lambda
  =(\mathfrak m^n/\mathfrak m^{n+1})\otimes_{A/\mathfrak m}(A_\lambda/\mathfrak mA_\lambda)
\]
y, como $A/\mathfrak m\to A_\lambda/\mathfrak mA_\lambda=A_\lambda/\mathfrak m_\lambda$ es biyectivo, el homomorfismo $\mathfrak m^n/\mathfrak m^{n+1}\to\mathfrak m_\lambda^n/\mathfrak m_\lambda^{n+1}$ es también biyectivo; la conclusión se deduce entonces de Bourbaki, \emph{Alg. comm.}, cap. III, \textsection~2, n.\textsuperscript{o}~8, cor. 3 del teor. 1.

Por último, si $A$ es henseliano, se deduce de la observación que precede a (18.6.3) que los homomorfismos $A\to A_\lambda$ son biyectivos, lo que demuestra (vi) por definición de ${}^h\!A$.

\begin{env}[18.6.7]
\label{IV.18.6.7}
Sea ahora $A$ un anillo semilocal y sean $\mathfrak m_i$ ($1\leq i\leq r$) sus ideales maximales. Se llama \emph{henselizado} de $A$, y se denota también por ${}^h\!A$, al anillo producto $\prod_{i=1}^r{}^h(A_{\mathfrak m_i})$ de los henselizados de los anillos locales $A_{\mathfrak m_i}$. Es un $A$-módulo fielmente plano y una $A$-álgebra semilocal, cuyos ideales maximales son los $\mathfrak m_i.{}^h\!A$ en virtud de (18.6.6, (iii)); además, si $\mathfrak r=\bigcap_i\mathfrak m_i$ es el radical de $A$, se deduce de lo anterior que $\mathfrak r.{}^h\!A$

\oldpage[IV-4]{139}

es el radical de ${}^h\!A$ y que la aplicación canónica $A/\mathfrak r\to{}^h\!A/\mathfrak r.{}^h\!A$ es biyectiva. Como el completado separado $\widehat A$ de $A$ para la topología $\mathfrak r$-preádica es el producto de los completados separados $\widehat A_{\mathfrak m_i}$ (Bourbaki, \emph{Alg. comm.}, cap. III, \textsection~2, n.\textsuperscript{o}~13, prop. 18), el homomorfismo canónico $\widehat A\to({}^h\!A)\widehat{\phantom{x}}$ es biyectivo por (18.6.6, (iv)), y está claro por (18.6.6, (v)) que, para que ${}^h\!A$ sea noetheriano, es necesario y suficiente que $A$ lo sea.

Para obtener el análogo de la propiedad universal (18.6.6, (ii)), convengamos, cuando $A$ y $B$ son dos anillos semilocales de radicales respectivos $\mathfrak r$, $\mathfrak s$, en llamar \emph{homomorfismo semilocal} de $A$ en $B$ a todo homomorfismo $\varphi$ tal que $\varphi(\mathfrak r)\subset\mathfrak s$. En el caso particular en que $B$ es un \emph{producto} de anillos locales $B_j$ ($1\leq j\leq n$), el dato de tal homomorfismo equivale al de sus proyecciones $\varphi_j:A\to B_j$, que son homomorfismos sujetos a la única condición de que $\varphi_j(\mathfrak r)\subset\mathfrak n_j$, donde $\mathfrak n_j$ es el ideal maximal de $B_j$. Por otra parte, si $\mathfrak m_i$ ($1\leq i\leq m$) son los ideales maximales de $A$, el ideal primo $\varphi_j^{-1}(\mathfrak n_j)$ debe contener uno de los $\mathfrak m_i$, puesto que contiene su intersección $\mathfrak r$; por tanto, es igual a uno de los $\mathfrak m_i$ y $\varphi_j$ se factoriza como $A\to A_{\mathfrak m_i}\xrightarrow{\psi_{ij}}B_j$, donde $\psi_{ij}$ es un homomorfismo \emph{local}. Si $\Hom.\operatorname{sloc}(A,B)$ denota el conjunto de los homomorfismos semilocales de $A$ en $B$, puede identificarse entonces este conjunto con
\[
  \prod_j\left(\bigcup_i\Hom.\operatorname{loc}(A_{\mathfrak m_i},B_j)\right)
\]
en el caso considerado; como un anillo semilocal henseliano es producto directo de anillos locales henselianos, se ve que la propiedad universal (18.6.6, (ii)) sigue siendo válida para un anillo semilocal $A$ cuando se sustituyen en ella los homomorfismos locales por los homomorfismos semilocales.
\end{env}

\begin{proposition}[18.6.8]
\label{IV.18.6.8}
Sean $A$ un anillo semilocal y $B$ una $A$-álgebra semilocal y entera sobre $A$. Entonces $B\otimes_A({}^h\!A)$ es una $B$-álgebra semilocal isomorfa a ${}^h\!B$.
\end{proposition}

Como $B\otimes_A({}^h\!A)$ es entero sobre ${}^h\!A$, cada uno de sus ideales maximales está sobre uno de los ideales maximales de ${}^h\!A$, y por tanto sobre un ideal maximal de $A$; y, como $B$ es entera sobre $A$, los ideales primos de $B$ situados sobre un ideal maximal de $A$ son ideales maximales de $B$, de modo que, finalmente, un ideal maximal de $B\otimes_A({}^h\!A)$ tiene por proyecciones en $\Spec({}^h\!A)$ y $\Spec(B)$ ideales maximales; teniendo en cuenta (18.6.6, (iii)) y (I, 3.4.9), se concluye que el anillo $B\otimes_A({}^h\!A)$ es semilocal. Además, para todo anillo local henseliano $C$, $\Hom.\operatorname{sloc}(B\otimes_A({}^h\!A),C)$ está en correspondencia biunívoca con el conjunto de los pares $(\varphi,\psi)$ de homomorfismos semilocales $\varphi:B\to C$, $\psi:{}^h\!A\to C$ tales que los compuestos
\[
  A\longrightarrow B\xrightarrow{\varphi}C,\qquad
  A\longrightarrow{}^h\!A\xrightarrow{\psi}C
\]
sean iguales. Pero, en virtud de la correspondencia biunívoca entre $\Hom.\operatorname{sloc}(A,C)$ y $\Hom.\operatorname{sloc}({}^h\!A,C)$, se ve que, para todo $\varphi\in\Hom.\operatorname{sloc}(B,C)$, existe uno y solo un $\psi$ que tiene la propiedad anterior; por tanto, la aplicación
\[
  \Hom.\operatorname{sloc}(B\otimes_A({}^h\!A),C)\longrightarrow\Hom.\operatorname{sloc}(B,C)
\]
es biyectiva, lo que demuestra la proposición en virtud de la unicidad de la solución de un problema universal.

\begin{theorem}[18.6.9]
\label{IV.18.6.9}
Sea $A$ un anillo semilocal.
\begin{enumerate}
  \item[(i)] Para que ${}^h\!A$ sea reducido (resp. normal), es necesario y suficiente que $A$ lo sea.
  \item[(ii)] Supongamos que $A$ es noetheriano. Entonces, para todo ideal primo $\mathfrak p$ de $A$, el anillo $({}^h\!A)_{\mathfrak p}/\mathfrak p.({}^h\!A)_{\mathfrak p}$ es producto directo de un número finito de extensiones algebraicas separables de $k(\mathfrak p)$ (lo que implica que las fibras del morfismo canónico $\Spec({}^h\!A)\to\Spec(A)$ son \emph{geométricamente regulares}).
\end{enumerate}
\end{theorem}

\oldpage[IV-4]{140}

(i) Como ${}^h\!A$ es un $A$-módulo fielmente plano, se deduce de (2.1.13) que, si ${}^h\!A$ es reducido (resp. normal), lo mismo sucede con $A$. Para demostrar el recíproco, puede uno limitarse al caso en que el anillo $A$ es local, de modo que, con las notaciones de (18.6.4), se tiene ${}^h\!A=\varinjlim A_\alpha$. Ahora bien, se deduce de la definición de los $A_\alpha$ y de (17.5.7) que, si $A$ es reducido (resp. íntegro e integralmente cerrado), lo mismo sucede con los $A_\alpha$; además, los homomorfismos $A_\alpha\to A_\beta$ para $\alpha\leq\beta$ son inyectivos en virtud de ($0_{\mathrm I}$, 6.5.1), puesto que $A_\beta$ es un $A_\alpha$-módulo fielmente plano; por tanto, los morfismos $\Spec(A_\beta)\to\Spec(A_\alpha)$ son dominantes, y se concluye de (5.13.2) (resp. (5.13.4)) que ${}^h\!A$ es reducido (resp. íntegro e integralmente cerrado).

(ii) Puede uno limitarse al caso en que $A$ es local. En virtud del hecho de que el funtor $\varinjlim$ conmuta con el producto tensorial, la fibra del morfismo $\Spec({}^h\!A)\to\Spec(A)$ en un punto $\mathfrak p$ es límite inductivo de las fibras de los morfismos $\Spec(A_\alpha)\to\Spec(A)$ en dicho punto. Como ${}^h\!A$ es noetheriano, se ve, teniendo en cuenta (17.6.2), que la cuestión se reduce a demostrar el lema siguiente:

\begin{lemma}[18.6.9.1]
\label{IV.18.6.9.1}
Sean $(B_\alpha)$ un sistema inductivo de anillos artinianos y $B=\varinjlim B_\alpha$ su límite inductivo. Si $B$ es noetheriano, entonces $B$ es artiniano; si además, para $\alpha\leq\beta$, los homomorfismos $\varphi_{\beta\alpha}:B_\alpha\to B_\beta$ son inyectivos, entonces existe $\lambda$ tal que, para $\alpha\geq\lambda$, el número de componentes locales $B_\alpha^{(i)}$ de $B_\alpha$ sea constante, que $\varphi_{\beta\alpha}(B_\alpha^{(i)})\subset B_\beta^{(i)}$ para todo $i$ y que los $B^{(i)}=\varinjlim B_\alpha^{(i)}$ sean las componentes locales de $B$.
\end{lemma}

Sean, en efecto, $Y_\alpha=\Spec(B_\alpha)$, $Y=\Spec(B)$, que, como espacio topológico, es igual a $\varprojlim Y_\alpha$ (8.2.10). Los espacios $Y_\alpha$ son finitos y discretos; además, si los $\varphi_{\beta\alpha}$ son inyectivos, los morfismos canónicos $Y_\beta\to Y_\alpha$ son dominantes y, por tanto, sobreyectivos. El lema se deducirá entonces del lema puramente topológico siguiente:

\begin{lemma}[18.6.9.2]
\label{IV.18.6.9.2}
Sea $(Y_\alpha,\psi_{\alpha\beta})$ un sistema proyectivo de espacios topológicos finitos y discretos. Si $Y=\varprojlim Y_\alpha$ es noetheriano, $Y$ es finito y discreto. Si además los $\psi_{\alpha\beta}$ son sobreyectivos, existe $\lambda$ tal que, para $\alpha\geq\lambda$, el número de elementos de $Y_\alpha$ sea constante e igual al número de elementos de $Y$.
\end{lemma}

En efecto, al ser $Y$ compacto y noetheriano, toda parte de $Y$ es compacta y, por tanto, cerrada, lo que implica que $Y$ es discreto y, por tanto, finito puesto que es compacto. Si además los $\psi_{\alpha\beta}$ son sobreyectivos, lo mismo sucede con $\psi_\alpha:Y\to Y_\alpha$, por lo que $\operatorname{Card}(Y)\geq\operatorname{Card}(Y_\alpha)$; y, como $\operatorname{Card}(Y_\alpha)$ es función creciente de $\alpha$, existe $\lambda$ tal que, para $\alpha\geq\lambda$, $\operatorname{Card}(Y_\alpha)=\operatorname{Card}(Y_\lambda)$; los $\psi_{\alpha\beta}$ son entonces biyectivos para $\alpha\geq\lambda$, y se tiene por tanto $\operatorname{Card}(Y)=\operatorname{Card}(Y_\lambda)$.

\begin{corollary}[18.6.10]
\label{IV.18.6.10}
Sea $A$ un anillo local noetheriano. Para que ${}^h\!A$ posea una de las propiedades siguientes:
\begin{enumerate}
  \item[$a)$] ser un anillo de Cohen-Macaulay (0, 16.5.3);
  \item[$b)$] verificar la propiedad $(S_n)$ (5.7.3);
  \item[$c)$] ser regular;
  \item[$d)$] verificar la propiedad $(R_n)$ (5.8.2);
\end{enumerate}
es necesario y suficiente que $A$ posea esa misma propiedad.
\end{corollary}

\oldpage[IV-4]{141}

Esto se deduce inmediatamente del hecho de que las fibras del morfismo $\Spec({}^h\!A)\to\Spec(A)$ son geométricamente regulares (18.6.9) y de (6.4.1), (6.5.1) y (6.5.3).

\begin{corollary}[18.6.11]
\label{IV.18.6.11}
Sea $A$ un anillo semilocal noetheriano. Para que un ideal primo $\mathfrak p$ de ${}^h\!A$ pertenezca a $\operatorname{Ass}({}^h\!A)$, es necesario y suficiente que $\mathfrak p\cap A$ pertenezca a $\operatorname{Ass}(A)$.
\end{corollary}

Teniendo en cuenta la descripción de las fibras del morfismo $\Spec({}^h\!A)\to\Spec(A)$, esto se deduce inmediatamente de (3.3.1).

\begin{proposition}[18.6.12]
\label{IV.18.6.12}
Sea $A$ un anillo local; las propiedades siguientes son equivalentes:
\begin{enumerate}
  \item[$a)$] $A$ es unirrama (0, 23.2.1) (resp. reducido y unirrama).
  \item[$b)$] Para toda $A$-álgebra local estrictamente esencialmente étale $A'$, $\Spec(A')$ es irreducible (resp. íntegro).
  \item[$c)$] $\Spec({}^h\!A)$ es irreducible (resp. íntegro).
\end{enumerate}
\end{proposition}

Observemos primero que ${}^h(A_{\mathrm{red}})=({}^h\!A)_{\mathrm{red}}$, pues ${}^h(A_{\mathrm{red}})=({}^h\!A)\otimes_A A_{\mathrm{red}}$ (18.6.8) y, como ${}^h(A_{\mathrm{red}})$ es reducido (18.6.9), es igual a $({}^h\!A)_{\mathrm{red}}$. Esto permite considerar en la demostración únicamente el caso en que $A$ es reducido y, por tanto, también $A'$ (17.5.7) y ${}^h\!A$ (18.6.9).

Con las notaciones de (18.6.4), la condición $b)$ significa que todos los $A_\alpha$ son íntegros; se concluye que ${}^h\!A=\varinjlim A_\alpha$ es íntegro (5.13.3), y por tanto $b)$ implica $c)$.

Está claro que, si ${}^h\!A$ es íntegro, lo mismo sucede con los $A_\alpha\subset{}^h\!A$, y por tanto $c)$ implica $b)$.

Demostremos que $c)$ implica $a)$. Observemos primero que $A\subset{}^h\!A$ es íntegro; sean $K$ y $L$ los cuerpos de fracciones de $A$ y ${}^h\!A$, respectivamente. Basta ver que, para toda $A$-álgebra \emph{finita} $B\subset K$, $B$ es un anillo local. Ahora bien, $B$ es un anillo semilocal, cuyo henselizado es entonces $B\otimes_A({}^h\!A)$ (18.6.8). Por planitud, $B\otimes_A({}^h\!A)$ se identifica con un subanillo de $K\otimes_A({}^h\!A)$ y este último con un subanillo de $L$; por tanto, ${}^h\!B$ es íntegro. Pero, como ${}^h\!B$ es un anillo semilocal henseliano, es producto directo de anillos locales (18.5.9), por lo que solo puede ser íntegro si es local, lo que implica efectivamente que $B$ es local.

Demostremos por último que $a)$ implica $c)$. Sea $C$ la clausura integral del anillo íntegro $A$. Como $C$ es un anillo local por hipótesis, $C\otimes_A({}^h\!A)$ es el henselizado de $C$ (18.6.8), y por tanto es íntegro e integralmente cerrado (18.6.9). Como $A$ es un subanillo de $C$, ${}^h\!A$ se identifica por planitud con un subanillo de ${}^h\!C$ y, por tanto, es también íntegro, lo que termina la demostración.

\begin{corollary}[18.6.13]
\label{IV.18.6.13}
Sea $A$ un anillo local henseliano. Entonces, para que $A$ sea unirrama, es necesario y suficiente que $A_{\mathrm{red}}$ sea íntegro.
\end{corollary}

\begin{proposition}[18.6.14]
\label{IV.18.6.14}
\begin{enumerate}
  \item[(i)] Todo límite inductivo filtrante de anillos locales henselianos (donde los homomorfismos de transición son locales) es un anillo local henseliano.
  \item[(ii)] El funtor $A\leadsto{}^h\!A$ en la categoría de los anillos locales (donde los morfismos son los homomorfismos locales) conmuta con los límites inductivos filtrantes.
  \item[(iii)] Todo anillo local henseliano $A$ es límite inductivo de un sistema inductivo filtrante de anillos locales henselianos y noetherianos $A_\lambda$ (siendo locales los homomorfismos de transición $A_\lambda\to A_\mu$).
\end{enumerate}
\end{proposition}

(i) Sea $(A_\lambda)$ un sistema inductivo filtrante de anillos locales henselianos, donde los homomorfismos de transición son locales. Demostremos que $A'=\varinjlim A_\lambda$, que es local ($0_{\mathrm{III}}$, 10.3.1.3), es henseliano. Ahora bien, sean $C$ una $A'$-álgebra étale y $\mathfrak n$ un ideal primo

\oldpage[IV-4]{142}

de $C$ sobre el ideal maximal $\mathfrak m'$ de $A'$, tal que $C_{\mathfrak n}$ sea estrictamente esencialmente étale sobre $A'$ (18.6.1), dicho de otro modo, que $k(\mathfrak m')\to k(\mathfrak n)$ sea un isomorfismo. En virtud de (8.8.2, (ii)) y de (17.7.8), existen un índice $\lambda$ y una $A_\lambda$-álgebra étale $C_\lambda$ tales que $C=C_\lambda\otimes_{A_\lambda}A'$ salvo isomorfismo; además, si $\mathfrak n_\lambda$ es la imagen inversa de $\mathfrak n$ en $C_\lambda$, y $\mathfrak m_\lambda$ el ideal maximal de $A_\lambda$, puede suponerse, en virtud de (8.8.2.4) y de la transitividad de las fibras (I, 3.6.4), que $(C_\lambda)_{\mathfrak n_\lambda}$ es estrictamente esencialmente étale sobre $A_\lambda$; puesto que $A_\lambda$ es henseliano, $(C_\lambda)_{\mathfrak n_\lambda}$ es necesariamente isomorfo a $A_\lambda$ (18.5.11, $b)$); existe por consiguiente una vecindad de $\mathfrak n_\lambda$ en $\Spec(C_\lambda)$ isomorfa a $\Spec(A_\lambda)$ (I, 6.5.4); se concluye que existe una vecindad de $\mathfrak n$ en $\Spec(C)$ isomorfa a $\Spec(A')$ y, por tanto, que $C_{\mathfrak n}$ es isomorfo a $A'$, lo que demuestra que $A'$ es henseliano (18.6.2) y termina la demostración.

(ii) Supongamos que $B=\varinjlim B_\lambda$, donde los $B_\lambda$ son anillos locales y los homomorfismos de transición $B_\lambda\to B_\mu$ ($\lambda\leq\mu$) son locales. Pongamos $A_\lambda={}^h\!B_\lambda$; los $A_\lambda$ son anillos locales henselianos (18.6.6, (vi)) y, en virtud de (18.6.6, (ii)), los homomorfismos de transición $B_\lambda\to B_\mu$ determinan de manera única homomorfismos locales $A_\lambda\to A_\mu$, de modo que $(A_\lambda)$ es un sistema inductivo. Todo se reduce a ver que $A'=\varinjlim A_\lambda$ es canónicamente isomorfo a $A={}^h\!B$. En virtud de (18.6.6, (ii)) y de la definición de los límites inductivos, se tiene, para todo anillo local henseliano $E$,
\[
  \Hom.\operatorname{loc}(A',E)
  =\varprojlim\Hom.\operatorname{loc}(A_\lambda,E)
  =\varprojlim\Hom.\operatorname{loc}(B_\lambda,E)
  =\Hom.\operatorname{loc}(B,E).
\]
Pero, como $A'$ es henseliano por (i), esto demuestra nuestra afirmación.

(iii) El anillo $A$ es límite inductivo filtrante de sus subanillos locales noetherianos $B_\lambda$, siendo locales los homomorfismos de transición (5.13.3, (iii)). Como los ${}^h\!B_\lambda$ son anillos locales henselianos y noetherianos (18.6.6, (v)) y $A={}^h\!A$ (18.6.6, (vi)), basta aplicar (ii) al sistema inductivo $(B_\lambda)$.

\begin{corollary}[18.6.15]
\label{IV.18.6.15}
Sean $A$ un anillo local henseliano y $X$ un $A$-preesquema de presentación finita sobre $A$. Existen entonces un anillo local noetheriano y henseliano $A_0$, un homomorfismo local $A_0\to A$, un preesquema $X_0$ de tipo finito sobre $A_0$ y un $A$-isomorfismo $X_0\otimes_{A_0}A\xrightarrow{\sim}X$.
\end{corollary}

Esto se deduce de (18.6.14) y de (8.8.2, (ii)).

\subsection{Henselización y anillos excelentes.}
\label{subsection:IV.18.7}

\begin{env}[18.7.1]
\label{IV.18.7.1}
En este número denotaremos por $\mathbf P(Z,k)$ una propiedad de la forma considerada en (7.3.1), suponiendo además que la propiedad $\mathbf Q(A,k)$ satisface la condición siguiente:

Para toda extensión algebraica separable $k'$ de $k$, y toda $k'$-álgebra local noetheriana $A$, la propiedad $\mathbf Q(A,k)$ es \emph{equivalente} a $\mathbf Q(A,k')$.

Es inmediato que todas las propiedades $\mathbf P$ consideradas en (7.3.8) verifican la condición anterior, habida cuenta de que, si $K$ es una extensión finita de $k$, $k'\otimes_k K$ es un producto directo de extensiones algebraicas separables de $K$.
\end{env}

\oldpage[IV-4]{143}

\begin{proposition}[18.7.2]
\label{IV.18.7.2}
Sea $A$ un anillo local noetheriano. Para que ${}^h\!A$ sea un $\mathbf P$-anillo (7.3.13), es necesario y suficiente que $A$ lo sea.
\end{proposition}

En efecto, en virtud de (18.6.6, (iv)), el completado $\widehat A$ de $A$ es también el completado de ${}^h\!A$ y se tiene, por tanto, $A\subset{}^h\!A\subset\widehat A$. Según (18.6.9), para todo $x\in\Spec(A)$, la fibra en $x$ del morfismo $\Spec({}^h\!A)\to\Spec(A)$ es discreta y finita, y, en cada uno de los puntos $y_i$ de esta fibra, $k(y_i)$ es una extensión algebraica separable de $k(x)$. La fibra formal de $A$ en el punto $x$ es, por tanto, el preesquema \emph{suma} de las fibras formales de ${}^h\!A$ en los puntos $y_i$. La conclusión resulta entonces de la hipótesis formulada sobre $\mathbf P$ en (18.7.1).

\begin{corollary}[18.7.3]
\label{IV.18.7.3}
Sea $A$ un anillo local noetheriano. Para que ${}^h\!A$ sea universalmente japonés, es necesario y suficiente que $A$ lo sea.
\end{corollary}

En efecto, para un anillo local noetheriano $B$, es equivalente decir que $B$ es universalmente japonés o que las fibras formales de $B$ son geométricamente reducidas (7.6.4 y 7.7.1); la conclusión resulta, pues, de (18.7.2).

\begin{corollary}[18.7.4]
\label{IV.18.7.4}
Sea $A$ un anillo local noetheriano. Para que las fibras formales de ${}^h\!A$ sean geométricamente regulares, es necesario y suficiente que las de $A$ lo sean.
\end{corollary}

Es un caso particular de (18.7.2).

\begin{proposition}[18.7.5]
\label{IV.18.7.5}
Sea $A$ un anillo local noetheriano. Si $A$ es universalmente catenario (5.6.2) (resp. formalmente catenario (7.1.9), resp. estrictamente formalmente catenario (7.2.6)), lo mismo sucede con ${}^h\!A$.
\end{proposition}

Con las notaciones de (18.6.4), si $A$ es universalmente catenario, lo mismo sucede con los $A_\alpha$, que son $A$-álgebras esencialmente de tipo finito (5.6.3). Para demostrar que, en este caso, ${}^h\!A$ es universalmente catenario, bastará, pues, establecer el lema siguiente:

\begin{lemma}[18.7.5.1]
\label{IV.18.7.5.1}
Sea $(B_\alpha,\varphi_{\beta\alpha})$ un sistema inductivo de anillos noetherianos y supongamos que, para $\alpha\leq\beta$, $\varphi_{\beta\alpha}$ haga de $B_\beta$ un $B_\alpha$-módulo fielmente plano y que $B=\varinjlim B_\alpha$ sea noetheriano. Entonces, si los $B_\alpha$ son catenarios (resp. universalmente catenarios), lo mismo sucede con $B$.
\end{lemma}

Decir que $B$ es universalmente catenario equivale a decir que $B[T_1,\ldots,T_n]$ es catenario para todo $n$ (5.6.2), y es inmediato que el sistema inductivo $(B_\alpha[T_1,\ldots,T_n])$ verifica las mismas condiciones que $(B_\alpha)$; basta, pues, demostrar que, si los $B_\alpha$ son catenarios, lo mismo sucede con $B$. Si $\mathfrak p_0\supset\mathfrak p_1\supset\cdots\supset\mathfrak p_r$ es una cadena saturada de ideales primos de $B$, las imágenes inversas $\mathfrak p_{0\alpha}\supset\mathfrak p_{1\alpha}\supset\cdots\supset\mathfrak p_{r\alpha}$ de estos ideales en $B_\alpha$ forman una cadena de ideales primos, y basta demostrar que, para $\alpha$ suficientemente grande, esta cadena es saturada. Ahora bien, cada $\mathfrak p_i$ es un ideal de tipo finito, por lo que existe un $\lambda$ tal que, para $\alpha\geq\lambda$, se tenga $\mathfrak p_i=\mathfrak p_{i\alpha}.B$. Se tiene, pues,
\[
1=\operatorname{codim}(V(\mathfrak p_i),V(\mathfrak p_{i+1}))
=\operatorname{codim}(V(\mathfrak p_{i\alpha}),V(\mathfrak p_{i+1,\alpha}))
\]
puesto que $B$ es un $B_\alpha$-módulo plano ((6.1.4) y ($0_{\mathrm I}$, 6.2.3)), lo que termina de demostrar el lema.

Supongamos ahora que $A$ es formalmente catenario. Sean $\mathfrak p$ un ideal primo de ${}^h\!A$ y $\mathfrak q$ su imagen inversa en $A$, de modo que $\mathfrak q.{}^h\!A\subset\mathfrak p$, y, por consiguiente, ${}^h\!A/\mathfrak p$ es un anillo cociente de ${}^h\!A/\mathfrak q.{}^h\!A$. Basta, pues, demostrar que ${}^h\!A/\mathfrak q.{}^h\!A$ es formalmente catenario para todo ideal primo $\mathfrak q$ de $A$ (7.1.9), y, como ${}^h\!A/\mathfrak q.{}^h\!A={}^h(A/\mathfrak q)$ por (18.6.8), basta (7.1.11) comprobar que, si $A$ es \emph{íntegro y formalmente catenario}, ${}^h\!A$ es \emph{formalmente equidimensional}. Pero, como el completado de ${}^h\!A$ es igual a $\widehat A$, esto resulta de la hipótesis de que $A$ es formalmente catenario.

\oldpage[IV-4]{144}

Supongamos, por último, que $A$ es estrictamente formalmente catenario; acabamos de ver que ${}^h\!A$ es formalmente catenario, y queda por demostrar que las fibras del morfismo $\Spec(\widehat A)\to\Spec({}^h\!A)$ verifican la propiedad $(S_1)$ (7.2.5, $b)$). Pero esto resulta de la hipótesis sobre $A$ y de (18.7.2), aplicado al caso en que $\mathbf P(Z,k)$ es la propiedad $(S_1)$ para $Z$.

\begin{corollary}[18.7.6]
\label{IV.18.7.6}
Si un anillo local noetheriano $A$ es excelente (7.8.2), lo mismo sucede con ${}^h\!A$.
\end{corollary}

\begin{remark}[18.7.7]
\label{IV.18.7.7}
Puede suceder que ${}^h\!A$ sea un anillo local excelente sin que $A$ sea universalmente catenario (ni, por consiguiente, excelente). Para verlo, retomemos el ejemplo del anillo $A$ de (5.6.11), con las mismas notaciones. El anillo $E$ construido en (5.6.11) es excelente cuando $k_0$ es de característica $0$ (7.8.3, (ii) y (iii)); consideremos dos anillos $E_1$, $E_2$ isomorfos a $E$, y «peguemos» $\Spec(E_1)$ y $\Spec(E_2)$ de modo que, si $\mathfrak m_i$ y $\mathfrak m'_i$ son los ideales maximales de $E_i$ ($i=1,2$) correspondientes a $\mathfrak m$ y $\mathfrak m'$, el punto $\mathfrak m_1$ quede «pegado» a $\mathfrak m'_2$ y el punto $\mathfrak m_2$ a $\mathfrak m'_1$; más precisamente, si $\varepsilon_i$ y $\varepsilon'_i$ son los homomorfismos canónicos de $E_i$ sobre $k(\mathfrak m_i)$ y $k(\mathfrak m'_i)$ ($i=1,2$), el esquema obtenido es $\Spec(R)$, donde $R$ es el subanillo de $E_1\times E_2$ formado por los pares $(x_1,x_2)$ tales que $\varepsilon'_2(x_2)=\sigma(\varepsilon_1(x_1))$ y $\varepsilon'_1(x_1)=\sigma(\varepsilon_2(x_2))$. Se comprueba fácilmente (por ejemplo, con ayuda de (17.6.3)) que $R$ es una $C$-álgebra finita y étale, y que hay dos ideales maximales $\mathfrak r_1$, $\mathfrak r_2$ de $R$ sobre el ideal maximal $\mathfrak n$ de $C$, siendo los completados de $R_{\mathfrak r_1}$ y $R_{\mathfrak r_2}$ canónicamente isomorfos al de $A=C_{\mathfrak n}$; $R_{\mathfrak r_1}$ es, pues, una $A$-álgebra estrictamente esencialmente étale y, por consiguiente, su henselizado es igual a ${}^h\!A$. Además, se ha visto (7.8.4, (ii)) que las fibras formales de $A$ son geométricamente regulares, por lo que lo mismo sucede con las de ${}^h\!A$ (18.7.2). Finalmente, el anillo $R$ es universalmente catenario, pues sus cocientes por sus dos ideales primos minimales son isomorfos a $E$, de donde se sigue la conclusión por (5.6.3, (iii)). El anillo $R_{\mathfrak r_1}$ es, pues, excelente y, por consiguiente, lo mismo sucede con ${}^h\!A$ (18.7.6), mientras que $A$ no es universalmente catenario.
\end{remark}

\subsection{Anillos estrictamente locales y henselización estricta.}
\label{subsection:IV.18.8}

\begin{proposition}[18.8.1]
\label{IV.18.8.1}
Sea $A$ un anillo local. Las condiciones siguientes son equivalentes:
\begin{enumerate}
  \item[$a)$] $A$ es un anillo henseliano y su cuerpo residual $k$ es separablemente cerrado.
  \item[$b)$] $A$ es un anillo henseliano y todo recubrimiento étale de $S=\Spec(A)$ es trivial.
  \item[$c)$] Para todo morfismo étale $f:X\to S$ y todo punto $x\in X$ tal que $f(x)$ sea el punto cerrado $s$ de $S$, existe una $S$-sección $u:S\to X$ de $X$ tal que $u(s)=x$.
\end{enumerate}
\end{proposition}

La hipótesis $c)$ implica que, para todo morfismo étale $f:X\to S$, los cuerpos residuales en los puntos de $f^{-1}(s)$ son todos isomorfos a $k$ y, además, que se verifica la condición $b)$ de (18.5.11).
Por tanto, $A$ es henseliano, y el hecho de que todo recubrimiento étale de $S$ sea trivial se deduce de (18.2.10, (ii)); así pues, $c)$ implica $b)$. Como los recubrimientos étale de $\Spec(k)$ son triviales si y solo si $k$ es separablemente cerrado, $b)$ implica $a)$ en virtud de (18.5.11, $b)$). Por último, también se deduce de (18.5.11, $b)$) que $a)$ implica $c)$.

\begin{definition}[18.8.2]
\label{IV.18.8.2}
Se dice que un anillo local es un anillo \emph{estrictamente local} si verifica las condiciones equivalentes de (18.8.1). Se llama \emph{esquema estrictamente local} a un esquema isomorfo al espectro de un anillo estrictamente local.
\end{definition}

\begin{remark}[18.8.3]
\label{IV.18.8.3}
Las condiciones de (18.8.1) son también equivalentes a la siguiente:
\begin{enumerate}
  \item[$d)$] Para todo morfismo localmente de tipo finito $f:X\to S$ y todo punto $x\in X$ tal que $f(x)$ sea el punto cerrado $s$ de $S$ y que $\sh{O}_{X,x}/\mathfrak m_s\sh{O}_{X,x}$ sea un cuerpo, extensión finita separable de $k(s)=k$, existe una vecindad abierta $U$ de $x$ en $X$ tal que $f|U$ sea una inmersión cerrada.
\end{enumerate}
\end{remark}

\oldpage[IV-4]{145}

(Se observará que la hipótesis de $d)$ se verifica si $f$ es \emph{formalmente no ramificado} en el punto $x$ (17.4.1.2).)

Dejamos al lector la demostración, que es sustancialmente la misma que la de (18.5.18).

\begin{lemma}[18.8.4]
\label{IV.18.8.4}
Sean $A$, $A'$ dos anillos locales, $\varphi:A\to A'$ un homomorfismo local que hace de $A'$ una $A$-álgebra esencialmente étale (18.6.1). Entonces, para todo anillo estrictamente local $B$, todo homomorfismo local $\psi:A\to B$ y todo homomorfismo de $k(A)$-álgebras $\alpha:k(A')\to k(B)$, existe un único homomorfismo local $\psi':A'\to B$ tal que $\psi=\psi'\circ\varphi$ y que $\alpha$ sea igual al homomorfismo $\overline{\psi'}$ deducido de $\psi'$ por paso a los cocientes.
\end{lemma}

Con las notaciones de (18.6.2), el homomorfismo $\alpha$ corresponde a un $k(s)$-morfismo $\Spec(k(y))\to\Spec(k(x))$ y, por tanto, a un punto $x'$ bien determinado de $X'$ situado sobre $y$. La existencia de la $Y$-sección $f'$ de $X'$ tal que $f'(y)=x'$ se deduce entonces de (18.8.1, $c)$), puesto que $f'$ es étale y $B$ estrictamente local, y su unicidad se deduce de que $f'$ es separado y de (17.4.9).

\begin{lemma}[18.8.5]
\label{IV.18.8.5}
Sean $A$ un anillo local, $A_1$, $A_2$ dos $A$-álgebras locales esencialmente étale, y $K$ un cuerpo, extensión de $k(A)$.
\begin{enumerate}
  \item[(i)] Para todo $k(A)$-homomorfismo $\gamma:k(A_1)\to k(A_2)$, existe a lo sumo un $A$-homomorfismo (necesariamente local) $\psi:A_1\to A_2$ tal que $\gamma$ sea el homomorfismo $\overline\psi$ deducido de $\psi$ por paso a los cocientes.
  \item[(ii)] Para todo par de $A$-homomorfismos locales $\beta_1:A_1\to K$, $\beta_2:A_2\to K$, existe una $A$-álgebra local esencialmente étale $A_3$, dos $A$-homomorfismos $\varphi_1:A_1\to A_3$, $\varphi_2:A_2\to A_3$ y un $A$-homomorfismo $\beta:A_3\to K$ tales que $\beta_1=\beta\circ\varphi_1$ y $\beta_2=\beta\circ\varphi_2$.
\end{enumerate}
\end{lemma}

Con las notaciones de (18.6.3), los homomorfismos $\beta_1$, $\beta_2$ de (ii) corresponden a dos $S$-morfismos $\Spec(K)\to X_1$, $\Spec(K)\to X_2$, de imágenes respectivas $x_1$, $x_2$; se deduce de ellos un $S$-morfismo bien determinado $\Spec(K)\to X_3$ (I, 3.2.1), de imagen $x_3$ situada sobre $x_1$ y $x_2$; la $A$-álgebra $A_3=\sh{O}_{X_3,x_3}$ y el homomorfismo $A_3\to K$ correspondientes a estos datos satisfacen las condiciones de (ii). Por otra parte, el dato de $\gamma$ en (i) corresponde a un $S$-morfismo $\Spec(k(x_2))\to\Spec(k(x_1))$, o también a un punto $x'_3$ de $X'_3$ situado sobre $x_2$; la unicidad de $\psi$ se deduce de la unicidad de una $X_2$-sección de $X'_3$ que pase por $x'_3$ (17.4.9).

\begin{env}[18.8.6]
\label{IV.18.8.6}
Sean $A$ un anillo local, $\mathfrak m$ su ideal maximal, $k$ su cuerpo residual y $\Omega$ una clausura separable de $k$. Consideremos el conjunto $\mathfrak E$ de $A$-álgebras esencialmente étale definido en (18.6.1), y denotemos por $L$ el conjunto de los $A$-homomorfismos locales $A'\to\Omega$, donde $A'$ recorre $\mathfrak E$; se observará que tal homomorfismo se factoriza en
\[
  \lambda:A'\longrightarrow k(A')\xrightarrow{\omega_\lambda}\Omega,
\]
donde $\omega_\lambda:k(A')\to\Omega$ es un $k$-homomorfismo (puesto que $\mathfrak mA'$ es el ideal maximal de $A'$); recíprocamente, el dato de tal homomorfismo $\omega_\lambda$ determina de manera única un homomorfismo $\lambda\in L$. Para todo $\lambda\in L$, denotaremos por $A_\lambda$ la $A$-álgebra esencialmente étale perteneciente a $\mathfrak E$ sobre la que está definido $\lambda$. El conjunto $L$ está \emph{preordenado} por la relación «existe un $A$-homomorfismo $\varphi_{\mu\lambda}:A_\lambda\to A_\mu$ tal que $\lambda=\mu\circ\varphi_{\mu\lambda}$», y se deduce de (18.8.5, (i)) que este homomorfismo $\varphi_{\mu\lambda}$ es \emph{único}. Además, $L$ es \emph{filtrante creciente} para la relación de preorden anterior (que seguiremos denotando por $\lambda\leq\mu$), en virtud

\oldpage[IV-4]{146}

de (18.8.5, (ii)). Es claro que $(A_\lambda,\varphi_{\mu\lambda})$ es un \emph{sistema inductivo} de $A$-álgebras locales, siendo los $\varphi_{\mu\lambda}$ homomorfismos locales; además, se deduce de (17.3.5) que, para $\lambda\leq\mu$, $A_\mu$ es una $A_\lambda$-álgebra esencialmente étale. Si $\overline\varphi_{\mu\lambda}:k(A_\lambda)\to k(A_\mu)$ es el $k$-homomorfismo deducido de $\varphi_{\mu\lambda}$ por paso a los cocientes, $(k(A_\lambda),\overline\varphi_{\mu\lambda})$ es un \emph{sistema inductivo} de extensiones algebraicas separables de $k$, y los $\omega_\lambda:k(A_\lambda)\to\Omega$ forman un sistema inductivo de $k$-homomorfismos. Además, el límite inductivo
\begin{equation}
  \omega:\varinjlim k(A_\lambda)\longrightarrow\Omega
  \label{IV.18.8.6.1}
  \tag{18.8.6.1}
\end{equation}
es un $k$-\emph{isomorfismo}. Basta, en efecto, demostrar que, para toda extensión finita separable $k'$ de $k$, existe una $A$-álgebra esencialmente étale $A'$ tal que $k'$ sea el cuerpo residual de $A'$ y que el homomorfismo $k\to k'$ se deduzca de $A\to A'$ por paso a los cocientes. Pero esto se deduce de (18.1.1), aplicado a los morfismos étale, teniendo en cuenta que $\Spec(A)$ es la única vecindad de su punto cerrado.

Observemos ahora que, si $\mathfrak m_\lambda$ es el ideal maximal de $A_\lambda$, se deduce de (17.6.1) que, para $\lambda\leq\mu$, se tiene $\mathfrak m_\mu=\mathfrak m_\lambda A_\mu$ y que $A_\mu$ es un $A_\lambda$-módulo plano. Por tanto, se deduce de ($0_{\mathrm{III}}$, 10.3.1.3) que el anillo $\varinjlim A_\lambda$ es \emph{local}, que el homomorfismo canónico $w:A\to\varinjlim A_\lambda$ es local y que se tiene un $k$-isomorfismo canónico
\[
  \varinjlim k(A_\lambda)\xrightarrow{\sim}k(\varinjlim A_\lambda);
\]
identificando estos dos cuerpos, se deduce un $k$-\emph{isomorfismo canónico}
\begin{equation}
  \omega^{-1}:\Omega\xrightarrow{\sim}k(\varinjlim A_\lambda).
  \label{IV.18.8.6.2}
  \tag{18.8.6.2}
\end{equation}
\end{env}

\begin{definition}[18.8.7]
\label{IV.18.8.7}
Sean $A$ un anillo local, $k=k(A)$ su cuerpo residual, e $i:k\to\Omega$ un homomorfismo de $k$ en una clausura separable $\Omega$ de $k$; se llama \emph{henselizado estricto} de $A$ relativo a $i$, y se denota por ${}^{hs}\!A_{(i)}$ (o por ${}^{hs}\!A$ cuando ello no da lugar a confusión), al límite inductivo del sistema inductivo $(A_\lambda,\varphi_{\mu\lambda})$ definido en (18.8.6).
\end{definition}

Como se tiene un $k$-isomorfismo canónico $\omega^{-1}:\Omega\xrightarrow{\sim}k({}^{hs}\!A_{(i)})$, $\Omega$ está dotado canónicamente de una estructura de ${}^{hs}\!A_{(i)}$-álgebra por el homomorfismo local
\[
  {}^{hs}\!A_{(i)}\longrightarrow k({}^{hs}\!A_{(i)})\xrightarrow{\omega}\Omega.
\]

Como en (18.6.5), se ve que la definición dada en (18.8.7) solo depende en apariencia de la elección del conjunto $\mathfrak E$; veremos además más abajo (18.8.8, (i) y (ii)) que ${}^{hs}\!A_{(i)}$ es un objeto que representa un funtor enteramente definido por los datos de $A$ y de $i$ ($0_{\mathrm{III}}$, 8.1.8), y que, por tanto, está definido como tal salvo isomorfismo único.

\begin{proposition}[18.8.8]
\label{IV.18.8.8}
Sean $A$ un anillo local, $k$ su cuerpo residual, $i:k\to\Omega$ un homomorfismo de $k$ en una clausura separable de $k$, y ${}^{hs}\!A_{(i)}$ el henselizado estricto de $A$ correspondiente a $i$.
\begin{enumerate}
  \item[(i)] ${}^{hs}\!A_{(i)}$ es un anillo estrictamente local (18.8.2) y el homomorfismo estructural $w:A\to{}^{hs}\!A_{(i)}$ es local.
  \item[(ii)] Para todo homomorfismo local $u:A\to B$, donde $B$ es un anillo estrictamente local, y todo $k$-homomorfismo $\psi:\Omega\to k(B)$, existe un único $A$-homomorfismo $v:{}^{hs}\!A_{(i)}\to B$ tal que $\psi$ se factoriza en $\Omega\xrightarrow{\omega^{-1}}k({}^{hs}\!A_{(i)})\xrightarrow{\overline v}k(B)$, donde $\overline v$ se deduce de $v$ por paso a los cocientes.

\oldpage[IV-4]{147}

  \item[(iii)] ${}^{hs}\!A_{(i)}$ es un $A$-módulo fielmente plano y, si $\mathfrak m$ es el ideal maximal de $A$, $\mathfrak m.{}^{hs}\!A_{(i)}$ es el ideal maximal de ${}^{hs}\!A_{(i)}$, y $k({}^{hs}\!A_{(i)})$ es una clausura separable de $k$.
  \item[(iv)] Para que ${}^{hs}\!A_{(i)}$ sea noetheriano, es necesario y suficiente que $A$ lo sea.
  \item[(v)] Si $A$ es estrictamente local (y, por tanto, $\Omega=k$), se tiene ${}^{hs}\!A\xrightarrow{\sim}A$.
  \item[(vi)] Si $i':k\to\Omega'$ es un segundo homomorfismo de $k$ en una clausura separable de $k$, para todo $k$-isomorfismo $\sigma:\Omega\xrightarrow{\sim}\Omega'$, el $A$-homomorfismo correspondiente (por (ii)) $v_\sigma:{}^{hs}\!A_{(i)}\to{}^{hs}\!A_{(i')}$ es un isomorfismo.
\end{enumerate}
\end{proposition}

Ya hemos visto más arriba que ${}^{hs}\!A_{(i)}$ es un anillo local y que $w$ es un homomorfismo local; el hecho de que ${}^{hs}\!A_{(i)}$ sea henseliano (y, por tanto, estrictamente local) se demuestra como en (18.6.6). Las afirmaciones de (iii), así como la suficiencia de la condición (iv), se deducen también de ($0_{\mathrm{III}}$, 10.3.1.3); la necesidad de esta condición se deduce de que ${}^{hs}\!A_{(i)}$ es un $A$-módulo fielmente plano ($0_{\mathrm I}$, 6.5.2).

Para demostrar (ii), observemos, con las notaciones de (18.8.6), que, para todo $\lambda\in L$, existe un único $A$-homomorfismo local $v_\lambda:A_\lambda\to B$ tal que el compuesto $k(A_\lambda)\xrightarrow{\omega_\lambda}\Omega\xrightarrow{\psi}k(B)$ se deduce de $v_\lambda$ por paso a los cocientes, en virtud de (18.8.4) y de la hipótesis de que $B$ es estrictamente local. La unicidad de $v_\lambda$ implica además que los $v_\lambda$ forman un sistema inductivo, de donde se deduce la existencia del homomorfismo $v$ que satisface las condiciones de (ii); su unicidad se deduce de (17.4.9). La unicidad de $v$ implica inmediatamente la afirmación (vi), considerando los compuestos $v_\sigma\circ v_{\sigma^{-1}}$ y $v_{\sigma^{-1}}\circ v_\sigma$. Por último, si $A$ es estrictamente local y $A'=B_{\mathfrak n}$ es una $A$-álgebra local esencialmente étale, donde $B$ es una $A$-álgebra étale, se deduce de (18.8.1) que existe una $\Spec(A)$-sección de $\Spec(B)$ que toma el valor $\mathfrak n$ en el punto cerrado de $\Spec(A)$; por tanto, se deduce de (17.4.1) que el homomorfismo $A\to A'$ es biyectivo, lo que demuestra (v).

\begin{env}[18.8.8.1]
\label{IV.18.8.8.1}
Sea $\mathcal C$ la categoría de los anillos estrictamente locales, cuyos morfismos son los homomorfismos locales; para todo $B\in\mathcal C$, denotemos por $F(B)$ el conjunto de los pares $(u,\psi)$ formados por un homomorfismo local $u:A\to B$ y un $k$-homomorfismo $\psi:\Omega\to k(B)$ tales que el compuesto $k(A)=k\xrightarrow{i}\Omega\xrightarrow{\psi}k(B)$ se deduce de $u$ por paso a los cocientes. Puede decirse que el objeto ${}^{hs}\!A_{(i)}$ \emph{representa} el funtor covariante $F:\mathcal C\to\mathbf{Ens}$ ($0_{\mathrm{III}}$, 8.1.11).

Se observará que, en virtud de (18.8.8, (vi)), los henselizados estrictos ${}^{hs}\!A_{(i)}$ son todos $A$-isomorfos (para los diversos $k$-homomorfismos $i$ de $k$ en clausuras separables de $k$), pero, para $i$ e $i'$ dados, hay en general una infinidad de $A$-isomorfismos ${}^{hs}\!A_{(i)}\xrightarrow{\sim}{}^{hs}\!A_{(i')}$. Más precisamente, el grupo de los $A$-automorfismos de ${}^{hs}\!A_{(i)}$ es isomorfo al grupo de Galois de $\Omega$ sobre $k$.
\end{env}

\begin{env}[18.8.9]
\label{IV.18.8.9}
Sea ahora $A$ un anillo \emph{semilocal}, sean $\mathfrak m_j$ ($1\leq j\leq r$) sus ideales maximales y, para todo índice $j$, consideremos un homomorfismo $i_j:k(A_{\mathfrak m_j})\to\Omega_j$ de $k(A_{\mathfrak m_j})$ en una clausura separable de este cuerpo. Se llama \emph{henselizado estricto} de $A$ relativo a los $i_j$, y se denota por ${}^{hs}\!A$ (cuando ello no da lugar a confusión), al producto $\prod_{j=1}^r({}^{hs}(A_{\mathfrak m_j}))$ de los henselizados estrictos de los anillos locales $A_{\mathfrak m_j}$ relativos a los homomorfismos $i_j$.

\oldpage[IV-4]{148}

Es, por tanto, en virtud de (18.8.8), un $A$-módulo fielmente plano y una $A$-álgebra semilocal cuyos ideales maximales son los $\mathfrak m_j.{}^{hs}\!A$, y cuyo radical es $\mathfrak r.{}^{hs}\!A$ (si $\mathfrak r$ denota el radical de $A$). Por último, la propiedad universal (18.8.8, (ii)) sigue siendo válida sustituyendo «homomorfismo local» por «homomorfismo semilocal» (18.6.7), y sustituyendo $\Omega$ por $\prod_j\Omega_j$.
\end{env}

\begin{proposition}[18.8.10]
\label{IV.18.8.10}
Sean $A$ un anillo semilocal y $B$ una $A$-álgebra finita. Sean $\mathfrak n_j$ ($1\leq j\leq r$) los ideales maximales de $B$. Entonces existe, para cada $j$, un ideal maximal $\mathfrak q_j$ de $B\otimes_A({}^{hs}\!A)$ situado sobre $\mathfrak n_j$, tal que ${}^{hs}\!B$ es isomorfo al producto directo de los anillos locales $(B\otimes_A({}^{hs}\!A))_{\mathfrak q_j}$ ($1\leq j\leq r$).
\end{proposition}

Evidentemente, podemos limitarnos al caso en que $A$ es local, sustituyendo $A$ por $A_{\mathfrak m_j}$, donde $\mathfrak m_j$ es el ideal maximal de $A$ imagen inversa de $\mathfrak n_j$. Pongamos, por otra parte, $C=B_{\mathfrak n_j}$; en virtud de la definición de ${}^{hs}\!B$ (18.8.9), todo se reduce a ver que ${}^{hs}\!C$ es $C$-isomorfo a $(C\otimes_A({}^{hs}\!A))_{\mathfrak q}$, donde $\mathfrak q$ es uno de los ideales maximales de $C\otimes_A({}^{hs}\!A)$. Sean $k$ y $K$ los cuerpos residuales de $A$ y $C$, y $\Omega$ el de ${}^{hs}\!A$, que es por definición una clausura separable de $k$; como $K$ es una extensión finita de $k$, podemos suponer que $K$ y $\Omega$ están contenidos en una misma clausura algebraica de $k$; entonces es inmediato que $K\otimes_k\Omega$ es producto directo de un número finito de cuerpos, todos isomorfos a la clausura separable $K\Omega$ de $K$. Por otra parte, como el anillo ${}^{hs}\!A$ es henseliano, $C\otimes_A({}^{hs}\!A)$ es producto directo de los anillos locales $(C\otimes_A({}^{hs}\!A))_{\mathfrak q_i}$, donde $\mathfrak q_i$ ($1\leq i\leq s$) recorre los ideales maximales de $C\otimes_A({}^{hs}\!A)$, y estos anillos locales son henselianos (18.5.10) y tienen $K\Omega$ por cuerpo residual, luego son estrictamente locales (18.8.1). Para abreviar, pongamos $C'=C\otimes_A({}^{hs}\!A)$ y $C'_i=C'_{\mathfrak q_i}$, que es, por tanto, un anillo cociente de $C'$; sean $f:C\to C'$, $g:{}^{hs}\!A\to C'$ y $\varphi_i:C'\to C'_i$ las aplicaciones canónicas.

\[
\xymatrix{
  &&& {}^{hs}\!C \ar@<0.8ex>[dl]_{w'} \\
  && C'_i \ar@<0.8ex>[ur]_{w} \\
  C \ar[uurrr]^{v} \ar[r]_{f} & C' \ar[ur] \\
  A \ar[u] \ar[r] & {}^{hs}\!A \ar[u]^{g} \ar[uuurr]_{u}
}
\]

Sea $v:C\to{}^{hs}\!C$ el homomorfismo canónico; existe un único homomorfismo local $u:{}^{hs}\!A\to{}^{hs}\!C$ que, por paso a los cocientes, da la inyección canónica $\Omega\to K\Omega$ y tal que el compuesto $A\to C\xrightarrow{v}{}^{hs}\!C$ sea igual al compuesto $A\to{}^{hs}\!A\xrightarrow{u}{}^{hs}\!C$ (18.8.8); por tanto, existe un único homomorfismo local $w_0:C'\to{}^{hs}\!C$ tal que $w_0\circ g=u$, $w_0\circ f=v$, y, puesto que ${}^{hs}\!C$ es un anillo local, este homomorfismo se factoriza en $C'\xrightarrow{\varphi_i}C'_i\xrightarrow{w}{}^{hs}\!C$ para un índice $i$ bien determinado, siendo $w$ un homomorfismo local. Por otra parte, como el anillo $C'_i$ es estrictamente local, existe un homomorfismo local $w':{}^{hs}\!C\to C'_i$ que, por paso a los cocientes, da el automorfismo identidad de $K\Omega$ y que satisface $w'\circ v=\varphi_i\circ f$ (18.8.8). Se deduce en primer lugar que $w\circ w'$ es un endomorfismo de ${}^{hs}\!C$ que, por paso a los cocientes, da el automorfismo identidad de $K\Omega$ y, por tanto, es

\oldpage[IV-4]{149}

la identidad (18.8.8). Por otra parte, se tiene $w'\circ w\circ\varphi_i\circ f=w'\circ v=\varphi_i\circ f$, y $w'\circ w\circ\varphi_i\circ g=w'\circ u=\varphi_i\circ g$; de donde $w'\circ w\circ\varphi_i=\varphi_i$, es decir, $w'\circ w$ es el automorfismo identidad de $C'_i$. Q.E.D.

\begin{remark}[18.8.11]
\label{IV.18.8.11}
La demostración de (18.8.10) utiliza el hecho de que $B$ es entera sobre $A$, pero solo utiliza el hecho de que $B$ sea una $A$-álgebra \emph{finita} para establecer que $K\otimes_k\Omega$ es producto directo de un número finito de cuerpos. Ahora bien, esta última propiedad sigue verificándose cuando se supone que $B$ es una $A$-álgebra \emph{entera semilocal} cuyos ideales maximales $\mathfrak n_j$ tienen grados residuales separables $[k(\mathfrak n_j):k]_s$ finitos.
\end{remark}

\begin{proposition}[18.8.12]
\label{IV.18.8.12}
Sea $A$ un anillo semilocal.
\begin{enumerate}
  \item[(i)] Para que ${}^{hs}\!A$ sea reducido (resp. normal), es necesario y suficiente que $A$ lo sea.
  \item[(ii)] Supongamos que $A$ es noetheriano. Entonces, para todo ideal primo $\mathfrak p$ de $A$, el anillo $({}^{hs}\!A)_{\mathfrak p}/\mathfrak p.({}^{hs}\!A)_{\mathfrak p}$ es producto directo de un número finito de extensiones algebraicas separables de $k(\mathfrak p)$ (lo que implica que las fibras del morfismo canónico $\Spec({}^{hs}\!A)\to\Spec(A)$ son \emph{geométricamente regulares} y son espacios \emph{discretos}).
\end{enumerate}
\end{proposition}

\begin{corollary}[18.8.13]
\label{IV.18.8.13}
Sea $A$ un anillo local noetheriano. Para que ${}^{hs}\!A$ posea una de las propiedades siguientes:
\begin{enumerate}
  \item[$a)$] ser un anillo de Cohen--Macaulay (0, 16.5.3);
  \item[$b)$] verificar la condición $(S_n)$ (5.7.3);
  \item[$c)$] ser regular;
  \item[$d)$] verificar la condición $(R_n)$ (5.8.2);
\end{enumerate}
es necesario y suficiente que $A$ posea esta misma propiedad.
\end{corollary}

\begin{corollary}[18.8.14]
\label{IV.18.8.14}
Sea $A$ un anillo semilocal noetheriano. Para que un ideal primo $\mathfrak p$ de ${}^{hs}\!A$ pertenezca a $\operatorname{Ass}({}^{hs}\!A)$, es necesario y suficiente que $\mathfrak p\cap A$ pertenezca a $\operatorname{Ass}(A)$.
\end{corollary}

Las demostraciones son las mismas que las de (18.6.9), (18.6.10) y (18.6.11), respectivamente.

\begin{proposition}[18.8.15]
\label{IV.18.8.15}
Sea $A$ un anillo local; las propiedades siguientes son equivalentes:
\begin{enumerate}
  \item[$a)$] $A$ es geométricamente unirrama (resp. reducido y geométricamente unirrama).
  \item[$b)$] Para toda $A$-álgebra local esencialmente étale $A'$, $\Spec(A')$ es irreducible (resp. íntegro).
  \item[$c)$] $\Spec({}^{hs}\!A)$ es irreducible (resp. íntegro).
\end{enumerate}
\end{proposition}

Como en (18.6.12), se reduce al caso en que $A$ es reducido, utilizando la relación ${}^{hs}(A_{\mathrm{red}})=({}^{hs}\!A)\otimes_A A_{\mathrm{red}}$ (18.8.11), y la equivalencia entre $b)$ y $c)$ se demuestra como en (18.6.12); lo mismo sucede con el hecho de que $a)$ implica $c)$, teniendo en cuenta (18.8.11) y la definición de un anillo local íntegro geométricamente unirrama. Por último, para demostrar que $c)$ implica $a)$, conservemos las mismas notaciones que en (18.6.12); se trata también de ver que ${}^{hs}\!B$ es íntegro. Pero, por (18.8.10), ${}^{hs}\!B$ es un anillo localizado del anillo semilocal $B\otimes_A({}^{hs}\!A)$, que, por planitud, se identifica también con un subanillo del cuerpo $L$ y, por tanto, es íntegro; en consecuencia, lo mismo sucede con ${}^{hs}\!B$, lo que termina la demostración.

\begin{corollary}[18.8.16]
\label{IV.18.8.16}
Sea $A$ un anillo local henseliano (resp. estrictamente local). Para que $A$ sea unirrama (resp. geométricamente unirrama), es necesario y suficiente que $\Spec(A)$ sea irreducible.
\end{corollary}

\oldpage[IV-4]{150}

\begin{proposition}[18.8.17]
\label{IV.18.8.17}
Sea $A$ un anillo local noetheriano. Si $A$ es universalmente catenario, lo mismo sucede con ${}^{hs}\!A$.
\end{proposition}

La demostración es la misma que la de (18.7.5).

\begin{proposition}[18.8.18]
\label{IV.18.8.18}
\begin{enumerate}
  \item[(i)] Todo límite inductivo filtrante de anillos estrictamente locales (donde los homomorfismos de transición son locales) es un anillo estrictamente local.
  \item[(ii)] El funtor $A\leadsto{}^{hs}\!A$ en la categoría de los anillos locales (donde los morfismos son los homomorfismos locales) conmuta con los límites inductivos filtrantes.
  \item[(iii)] Todo anillo estrictamente local $A$ es límite inductivo de un sistema inductivo filtrante de anillos estrictamente locales noetherianos (siendo locales los homomorfismos de transición).
\end{enumerate}
\end{proposition}

Las demostraciones están calcadas de las de (18.6.14).

\subsection{Fibras formales de los anillos noetherianos henselianos.}
\label{subsection:IV.18.9}

\begin{theorem}[18.9.1]
\label{IV.18.9.1}
Sea $A$ un anillo local noetheriano henseliano cuyas fibras formales son geométricamente normales. Entonces las fibras formales de $A$ son geométricamente íntegras (y, por tanto, geométricamente conexas).
\end{theorem}

Las dos afirmaciones del enunciado son, de hecho, equivalentes, pues todo preesquema localmente noetheriano, conexo y normal es íntegro. Como toda $A$-álgebra íntegra finita es un anillo local henseliano (18.5.9 y 18.5.10), se deduce de (7.3.16.2) que la cuestión se reduce a demostrar que, si además se supone que $A$ es \emph{íntegro}, entonces la fibra del morfismo $\Spec(\widehat A)\to\Spec(A)$ en el punto genérico de $\Spec(A)$ es \emph{íntegra}, lo que se deduce del

\begin{corollary}[18.9.2]
\label{IV.18.9.2}
Sea $A$ un anillo local noetheriano henseliano íntegro cuyas fibras formales son geométricamente normales. Entonces $\widehat A$ es íntegro.
\end{corollary}

En efecto, se deduce de (18.8.16) que $A$ es unirrama, por lo que el teorema se deduce de la hipótesis formulada sobre las fibras formales de $A$ y de (7.6.3).

\begin{corollary}[18.9.3]
\label{IV.18.9.3}
Bajo las hipótesis de (18.9.2), el cuerpo de fracciones $L$ de $\widehat A$ es una extensión separable del cuerpo de fracciones $K$ de $A$, y $K$ es algebraicamente cerrado en $L$.
\end{corollary}

Esto se deduce de que la fibra del morfismo $\Spec(\widehat A)\to\Spec(A)$ en el punto genérico de $\Spec(A)$ es geométricamente íntegra, y de (4.3.2) y (4.3.5).

\begin{corollary}[18.9.4]
\label{IV.18.9.4}
Sea $A$ un anillo local noetheriano y henseliano cuyas fibras formales sean geométricamente normales (por ejemplo, un anillo local noetheriano henseliano y excelente), y pongamos $A'=\widehat A$. Sean $X$ un $A$-preesquema, $X'=X\otimes_A A'$, y $g:X'\to X$ la proyección canónica. Entonces la aplicación $U\mapsto g^{-1}(U)$ es una biyección del conjunto de las partes a la vez abiertas y cerradas de $X$ sobre el conjunto de las partes a la vez abiertas y cerradas de $X'$.
\end{corollary}

Como $g$ es un morfismo fielmente plano y cuasicompacto, se sabe por (2.3.12) que la topología de $X$ es \emph{cociente} de la de $X'$ por la relación de equivalencia definida por $g$. Todo se reduce, por tanto, a ver que una parte abierta y cerrada $U'$ de $X'$ está \emph{saturada} para esta relación. Ahora bien, como el morfismo $\Spec(A')\to\Spec(A)$ tiene fibras geométricamente conexas (18.9.1), las fibras de $g$ son conexas, de donde se deduce inmediatamente nuestra afirmación.

En particular, si $X$ es localmente conexo (lo que ocurrirá si $X$ es localmente noetheriano), entonces, para toda componente conexa $U$ de $X$ (que es abierta y cerrada

\oldpage[IV-4]{151}

en $X$), $g^{-1}(U)$ es conexa (Bourbaki, \emph{Top. gén.}, cap.~I\textsuperscript{er}, 3\textsuperscript{e} ed., \textsection~11, n\textsuperscript{o}~3, prop.~7). Si se denota por $\pi_0(X)$ el conjunto de las componentes conexas de $X$, la aplicación $\pi_0(X')\to\pi_0(X)$ deducida canónicamente de $g$ es, por tanto, \emph{biyectiva}.

\begin{corollary}[18.9.5]
\label{IV.18.9.5}
Bajo las hipótesis de (18.9.4), el funtor
\begin{equation}
  Z\longmapsto Z\otimes_A A'
  \tag{18.9.5.1}
  \label{IV.18.9.5.1}
\end{equation}
de la categoría de los preesquemas étales sobre $X$ a la de los preesquemas étales sobre $X'$ es plenamente fiel.
\end{corollary}

Sean $Z_1$, $Z_2$ dos preesquemas étales sobre $X$, y pongamos $Z'_i=Z_i\otimes_A A'$ ($i=1,2$); se trata de demostrar que todo $X'$-morfismo $Z'_1\to Z'_2$ procede, por cambio de base, de un único $X$-morfismo $Z_1\to Z_2$. Supongamos primero que $Z_2$ es \emph{separado} sobre $X$; entonces, como $\Hom_X(Z_1,Z_2)$ se identifica con $\Gamma((Z_1\times_X Z_2)/Z_1)$ y $Z_1\times_X Z_2$ es étale y separado sobre $Z_1$, $\Hom_X(Z_1,Z_2)$ se identifica funtorialmente, en virtud de (17.9.3), con el conjunto de las partes abiertas y cerradas $U$ de $Z_1\times_X Z_2$ tales que la restricción a $U$ de la proyección $p:Z_1\times_X Z_2\to Z_1$ sea un morfismo sobreyectivo y radicial. La afirmación se deduce, por tanto, en el caso considerado de (18.9.4) y de que $Z'_1\times_{X'}Z'_2=(Z_1\times_X Z_2)\otimes_A A'$.

Pasemos al caso general: bastará demostrar que el grafo $\Gamma'$ de un $X'$-morfismo $Z'_1\to Z'_2$, que es abierto (17.9.3) en $Z'_1\times_{X'}Z'_2$, es de la forma $\Gamma\otimes_A A'$, donde $\Gamma$ está inducido sobre una parte \emph{abierta} de $Z=Z_1\times_X Z_2$. En efecto, la restricción a $\Gamma$ de la proyección $p:Z\to Z_1$ es entonces un \emph{isomorfismo}, pues, por el cambio de base $A\to A'$, esta restricción se convierte en el isomorfismo dado por la restricción a $\Gamma'$ de la proyección $p':Z'\to Z'_1$ (con $Z'=Z\otimes_A A'$), y basta aplicar (2.7.1, (viii)); la conclusión se deduce entonces de la caracterización de los grafos de morfismos (I, 5.3).

Si $q:Z'\to Z$ es la proyección canónica, para demostrar que $\Gamma'$ es de la forma $q^{-1}(\Gamma)$, donde $\Gamma$ es abierto en $Z$, basta, puesto que $q$ es un morfismo fielmente plano y cuasicompacto, demostrar que existe un conjunto $U\subset Z$ tal que $\Gamma'=q^{-1}(U)$ (2.3.12). Pongamos $Z''=Z'\times_Z Z'=Z\times_X X''$, donde $X''=X'\times_X X'$, y sean $q_1:Z''\to Z'$, $q_2:Z''\to Z'$ las proyecciones canónicas. Aplicando (4.5.19.1), basta mostrar que $q_1^{-1}(\Gamma')=q_2^{-1}(\Gamma')$. Pero esta propiedad es verdadera si y solo si lo es después de cada cambio de base $\Spec(k(x))\to X$, donde $x$ recorre $X$. En otras palabras, basta demostrar el corolario cuando $X$ es el \emph{espectro de un cuerpo}; pero, como todo $X$-preesquema étale es entonces automáticamente \emph{separado} sobre $X$ (17.6.2, $c'$)), se vuelve al caso considerado al comienzo de la demostración.

\begin{remarks}[18.9.6]
\label{IV.18.9.6}
\begin{enumerate}
  \item[(i)] Es posible que, si el cuerpo residual de $A$ es de característica 0, el funtor (18.9.5.1) induzca incluso una \emph{equivalencia} entre la categoría de los recubrimientos étales de $X$ y la categoría de los recubrimientos étales de $X'$. En todo caso, puede demostrarse cuando $A$ es \emph{excelente}, utilizando la resolución de singularidades de Hironaka (M.~Artin). Es plausible que el enunciado análogo siga siendo verdadero sin restricción sobre la característica del cuerpo residual, a condición de limitarse a los recubrimientos «galoisianos principales» cuyo grupo sea de orden primo con la característica residual (cf. [41]).

\oldpage[IV-4]{152}

  \item[(ii)] Se ignora si las fibras formales de $A$ son geométricamente conexas (o incluso geométricamente irreducibles) cuando $A$ es un anillo local henseliano cualquiera. Bastaría con que, para todo anillo local noetheriano, henseliano e \emph{íntegro} $A$, $\Spec(\widehat A)$ fuera irreducible. Se ignora si esto sucede siempre.

  \item[(iii)] Se ignora si, cuando $A$ es un anillo local excelente, su henselizado estricto ${}^{hs}\!A$ es excelente. Puede verse que, para dilucidar esta cuestión, es posible reducirse al caso en que $A=k[[T_1,\ldots,T_n]]$, siendo $k$ un cuerpo de característica $p>0$. Teniendo en cuenta (18.8.17), la cuestión consiste en saber si las fibras formales de ${}^{hs}\!A$ son geométricamente regulares. La respuesta es afirmativa cuando $n=1$, pero no se conoce para $n=2$. Puede demostrarse que la respuesta es afirmativa siempre que, para todo esquema $Y_1$ finito sobre $Y=\Spec(\widehat A)$, puedan resolverse las singularidades de $Y_1$ (7.9.1).
\end{enumerate}
\end{remarks}

La afirmación de conexidad formulada en (18.9.1) se generaliza de la manera siguiente:

\begin{theorem}[18.9.7]
\label{IV.18.9.7}
Sean $A$, $B$ dos anillos locales noetherianos, $\rho:A\to B$ un homomorfismo local, y $f:\Spec(B)=X\to\Spec(A)=Y$ el morfismo correspondiente. Supongamos que se verifican las condiciones siguientes:
\begin{enumerate}
  \item[(i)] $f$ es plano y todas las fibras $X_y=f^{-1}(y)$ ($y\in Y$) son geométricamente reducidas (4.6.2).
  \item[(ii)] O bien el anillo $A$ es geométricamente unirrama (0, 23.2.1), o bien el anillo $A$ es unirrama (0, 23.2.1) y $k(B)$ es una extensión primaria (4.3.1) de $k(A)$.
\end{enumerate}
Entonces, si $\eta$ es el punto genérico de $Y$, $X_\eta$ es conexo y, para toda parte cerrada $T$ de $X$ tal que $f(T)$ sea rara en $Y$, $X-T$ es conexo.
\end{theorem}

Las dos conclusiones enunciadas son, de hecho, equivalentes; en efecto, la hipótesis de que $f(T)$ sea rara implica que $f(T)$ no contiene a $\eta$; por otra parte, como $\{\eta\}$ es proconstructible (I, 9.6) y $f$ es plano, $X_\eta$ es denso en $X$ (2.3.10), y, por consiguiente, si $X_\eta$ es conexo, lo mismo sucede con $X-T$, que lo contiene. Recíprocamente, cuando $U$ recorre el conjunto de las vecindades abiertas afines de $\eta$ en $Y$, $X_\eta$ es límite proyectivo de los subpreesquemas inducidos sobre los abiertos $f^{-1}(U)$ de $X$ (8.1.2, $a)$); si se supone que $X-T$ es conexo para toda parte cerrada $T$ de $X$ tal que $f(T)$ sea rara en $Y$, los $f^{-1}(U)$ son conexos y, por tanto, lo mismo sucede con $X_\eta$ (8.4.1).

Demostremos primero que podemos limitarnos al caso en que $A$ es íntegro y $Y=\Spec(A)$ es geométricamente unirrama (6.15.1) (lo que implica que $A$ es geométricamente unirrama, pero no equivale a esta condición (6.15.2)). Es claro que podemos suponer primero que $Y$ es reducido, considerando el morfismo $X\times_Y Y_{\mathrm{red}}\to Y_{\mathrm{red}}$, deducido de $f$ por cambio de base, que es plano y tiene las mismas fibras que $f$. Podemos suponer, por tanto, que $A$ es íntegro y unirrama; si $K$ es el cuerpo de fracciones de $A$, existe entonces una sub-$A$-álgebra finita $A''$ de $K$ tal que, si $A'$ es la clausura integral de $A$, el morfismo $\Spec(A')\to\Spec(A'')$ sea radicial (0, 23.2.5); se deduce de ello (6.15.5) que $\Spec(A'')$ es geométricamente unirrama. Como $A$ es unirrama, $A'$ es un anillo local y, por tanto, lo mismo sucede con $A''$; demostremos que el anillo $B''=B\otimes_A A''$ también es local. En efecto, $B''$ es una $B$-álgebra finita y, por tanto, un anillo semilocal noetheriano; si $k''$ es el cuerpo residual de $A''$, los ideales maximales de $B''$ son los puntos de $\Spec(B'')$ situados sobre el punto cerrado de $\Spec(B)$ y, por tanto, los puntos de $\Spec(k(B)\otimes_{k(A)}k'')$ (I, 3.4.9)

\oldpage[IV-4]{153}

puesto que el punto cerrado de $\Spec(A'')$ es el único punto situado sobre el punto cerrado de $\Spec(A)$. Ahora bien, cuando $A$ es geométricamente unirrama, $k''$ es una extensión radicial de $k(A)$, por lo que $\Spec(k(B)\otimes_{k(A)}k'')$ es radicial sobre $\Spec(k(B))$ y, por tanto, solo contiene un punto. Cuando $A$ es unirrama y $k(B)$ es una extensión primaria de $k(A)$, $\Spec(k(B)\otimes_{k(A)}k'')$ es irreducible (4.3.2) y, como es un espacio finito y discreto (I, 6.4.4), sigue reduciéndose a un punto, lo que demuestra que, en los dos casos considerados, $B''$ es local. Si $Y''=\Spec(A'')$, $X''=\Spec(B'')=X\times_Y Y''$, el morfismo $f''=f_{(Y'')}:X''\to Y''$ es plano y tiene fibras geométricamente reducidas; además, si $\eta''$ es el punto genérico de $Y''$, $k(\eta'')=k(\eta)=K$ por definición, por lo que las fibras $X_\eta$ y $X''_{\eta''}$ son isomorfas. Basta, por tanto, demostrar que $X''_{\eta''}$ es conexa.

\begin{remark}[18.9.7.1]
\label{IV.18.9.7.1}
Cuando el anillo $A$ es \emph{japonés}, en lo que precede puede tomarse $A''=A'$ por definición (0, 23.1.1); se ve, por tanto, que en este caso incluso se reduciría la demostración del teorema al caso en que $A$ es \emph{íntegramente cerrado}.
\end{remark}

\begin{env}[18.9.7.2]
\label{IV.18.9.7.2}
Supongamos, pues, en adelante que $Y$ es íntegro y geométricamente unirrama. Observemos que, si $T$ es una parte cerrada de $X$ tal que $f(T)$ sea rara en $Y$, $T$ no contiene ningún punto maximal de $X$, puesto que $f$ es plano (2.3.4), y es, por tanto, rara en $X$; como $X$ es un esquema local y, por consiguiente, conexo, para demostrar que $X-T$ es conexo basta, en virtud de (15.5.6.1), mostrar que, para todo $x\in X$ tal que $f(x)\neq\eta$, $\Spec(\sh{O}_{X,x})-\{x\}$ es conexo. Pongamos $y=f(x)$, $A_1=\sh{O}_{Y,y}$, $B_1=\sh{O}_{X,x}$, $Y_1=\Spec(A_1)$, $X_1=\Spec(B_1)$, y sea $f_1:X_1\to Y_1$ el morfismo correspondiente al homomorfismo local $A_1\to B_1$ deducido de $\rho$; es claro que $f_1$ es plano, y se deduce de (4.6.1) que sus fibras son geométricamente reducidas; además, $A_1$ es íntegro y geométricamente unirrama (6.15.1), y $\dim(A_1)\geq 1$. Se reduce así la cuestión a demostrar el
\end{env}

\begin{lemma}[18.9.7.3]
\label{IV.18.9.7.3}
Sean $A$, $B$ dos anillos locales noetherianos, $Y=\Spec(A)$, $X=\Spec(B)$, $\rho:A\to B$ un homomorfismo local, y $f:X\to Y$ el morfismo correspondiente. Se supone que se verifican las condiciones siguientes:
\begin{enumerate}
  \item[(i)] $f$ es plano y todas las fibras $X_y$ ($y\in Y$) son geométricamente reducidas.
  \item[(ii)] $A$ es íntegro, geométricamente unirrama y $\dim(A)\geq 1$.
\end{enumerate}
Entonces, si $b$ es el punto cerrado de $X$, $X-\{b\}$ es conexo.
\end{lemma}

Observemos primero que, en virtud del teorema de Hartshorne (5.10.7), $X-\{b\}$ es conexo si se tiene $\operatorname{prof}(B)\geq 2$. Ahora bien, se tiene (6.3.1)
\begin{equation}
  \operatorname{prof}(B)=\operatorname{prof}(A)+\operatorname{prof}(B\otimes_A k)
  \tag{18.9.7.4}
  \label{IV.18.9.7.4}
\end{equation}
denotando por $k$ el cuerpo residual de $A$. Por otra parte, puesto que $A$ es íntegro y $\dim(A)\geq 1$, se tiene $\operatorname{prof}(A)\geq 1$ y, por tanto, $\operatorname{prof}(B)\geq 2$ \emph{salvo} cuando $\operatorname{prof}(B\otimes_A k)=0$; además, $B\otimes_A k$ es reducido por (i), de modo que la relación $\operatorname{prof}(B\otimes_A k)=0$ significa que $B\otimes_A k$ es un \emph{cuerpo} (0, 16.4.7). Supondremos, pues, en adelante que se verifica esta última condición. Comenzaremos tratando un caso en el que la demostración es muy sencilla.

A) \emph{Caso en que $A$ es íntegramente cerrado.} Entonces, si $\dim(A)\geq 2$, se tiene $\operatorname{prof}(A)\geq 2$ (0, 16.5.1) y, por tanto, $\operatorname{prof}(B)\geq 2$. Solo queda, pues, considerar el caso en que $\dim(A)=1$ y $B\otimes_A k$ es un cuerpo; $A$ es entonces un anillo de valoración discreta, y por tanto regular, y puesto que $B\otimes_A k$

\oldpage[IV-4]{154}

es un cuerpo y $f$ es plano, $B$ es regular (0, 17.3.3); pero, además, se tiene (6.1.1.1) $\dim(B)=\dim(A)+\dim(B\otimes_A k)=\dim(A)=1$, por lo que $B$ es también un anillo de valoración discreta. Pero entonces $X-\{b\}$ se reduce a un solo punto, de donde se deduce el lema (18.9.7.3) en este caso.

Se observará que esto demuestra el enunciado del teorema (18.9.7) cuando se supone además que, en dicho enunciado, el anillo $A$ es \emph{japonés}, pues, en virtud de la observación (18.9.7.1), se vuelve al caso en que $A$ es \emph{íntegramente cerrado} y entonces, en la reducción precedente (18.9.7.3), el anillo $A_1=\sh{O}_{Y,y}$ es también íntegramente cerrado y puede aplicarse el resultado que acaba de demostrarse.

B) \emph{Caso en que $\dim(A)\geq 2$.} El resultado de (18.9.7.3) se deducirá entonces de los dos lemas siguientes:

\begin{lemma}[18.9.7.5]
\label{IV.18.9.7.5}
Sean $A$ un anillo semilocal noetheriano y reducido, $A'$ su clausura integral en su anillo total de fracciones, $X=\Spec(A)$, $X'=\Spec(A')$; sean $S$ el conjunto de los puntos cerrados de $X$ y $U=X-S$. Supongamos que, para todo punto $x'\in X'$ situado sobre un punto de $S$, se tiene $\dim(\sh{O}_{X',x'})\geq 2$; entonces el anillo $A^{(1)}=\Gamma(U,\sh{O}_X)$ es entero sobre $A$ y es una $A$-álgebra isomorfa a una subálgebra de $A'$.
\end{lemma}

Recordemos que, si los $\mathfrak p_i$ ($1\leq i\leq m$) son los ideales primos minimales de $A$, $A'$ es el producto directo de las clausuras integrales $A'_i$ de los anillos íntegros $A_i=A/\mathfrak p_i$, de modo que $X'$ es la suma de los esquemas $X'_i=\Spec(A'_i)$. Si $X_0$ es la suma de los esquemas $X_i=\Spec(A_i)$ ($1\leq i\leq m$) y $U_0$ la imagen inversa de $U$ en $X_0$, el homomorfismo canónico $\Gamma(U,\sh{O}_X)\to\Gamma(U_0,\sh{O}_{X_0})$ es inyectivo, puesto que $X$ es reducido y el morfismo $X_0\to X$ (y, por tanto, también el morfismo $U_0\to U$) es sobreyectivo. Como $\Gamma(U_0,\sh{O}_{X_0})$ es el producto directo de los $\Gamma(U_i,\sh{O}_{X_i})$ ($1\leq i\leq m$), donde $U_i$ es la imagen inversa de $U$ en $X_i$, todo se reduce a demostrar que, para cada $i$, $\Gamma(U_i,\sh{O}_{X_i})$ es isomorfo a una sub-$A_i$-álgebra de $A'_i$. En otras palabras, se vuelve al caso en que $A$ es un anillo semilocal noetheriano íntegro. Como $X$ es reducido y el morfismo $X'\to X$ es sobreyectivo, el homomorfismo $\Gamma(U,\sh{O}_X)\to\Gamma(U',\sh{O}_{X'})$ (donde $U'$ es la imagen inversa de $U$ en $X'$) es inyectivo, y todo se reduce a ver que el homomorfismo canónico $A'=\Gamma(X',\sh{O}_{X'})\to\Gamma(U',\sh{O}_{X'})$ es \emph{biyectivo}. Ahora bien, por (I, 8.2.1.1), $\Gamma(U',\sh{O}_{X'})$ es la intersección de los anillos locales $A'_{\mathfrak p'}$, donde $\mathfrak p'$ recorre $U'$, y, en virtud de la hipótesis, entre estos anillos locales figuran todos aquellos para los que $\mathfrak p'$ es de altura 1. Pero el razonamiento de (0, 23.2.7) se aplica también a un anillo semilocal íntegro noetheriano, por lo que $A'$ es un \emph{anillo de Krull semilocal} y es, por tanto, la intersección de los anillos locales $A'_{\mathfrak p'}$, donde $\mathfrak p'$ recorre el conjunto de los ideales primos de altura 1 de $A'$ (Bourbaki, \emph{Alg. comm.}, cap.~VII, \textsection~1, n\textsuperscript{o}~6, th.~4); \emph{a fortiori}, $A'$ es la intersección de los $A'_{\mathfrak p'}$ para $\mathfrak p'\in U'$, lo que termina la demostración del lema.

\begin{remark}[18.9.7.6]
\label{IV.18.9.7.6}
Se sabe por (0, 23.2.5) que existe una sub-$A$-álgebra finita $A''$ del cuerpo de fracciones $K$ de $A$ tal que el morfismo $\Spec(A')\to\Spec(A'')$ sea \emph{radicial}; como este morfismo es también entero y dominante, y por tanto sobreyectivo (II, 6.1.10), es un \emph{homeomorfismo} (2.4.5). En virtud de la interpretación geométrica de $\dim(\sh{O}_{X',x'})$ (5.1.2), se tiene, por tanto, para todo punto $x''$ de $X''=\Spec(A'')$ situado sobre un punto $x\in X$, $\dim(\sh{O}_{X'',x''})=\dim(\sh{O}_{X',x'})\geq 2$ si $x'$ es el único punto de $X'$ situado sobre $x$. Dicho

\oldpage[IV-4]{155}

esto, se deduce en primer lugar de (0, 16.1.5) y (0, 16.1.4.1) que la relación $\dim(\sh{O}_{X',x'})\geq 2$ implica $\dim(\sh{O}_{X,x})\geq 2$. Recíprocamente, supongamos que $\dim(\sh{O}_{X,x})\geq 2$ y que $A$ es universalmente catenario (5.6.2); entonces lo mismo sucede con $\sh{O}_{X,x}$ (5.6.3), y por tanto $\dim(\sh{O}_{X'',x''})=\dim(\sh{O}_{X,x})\geq 2$ en virtud de (5.6.10), de modo que se tiene $\dim(\sh{O}_{X',x'})\geq 2$. La misma conclusión es válida si $\sh{O}_{X,x}$ es \emph{geométricamente unirrama}, pues la fibra del morfismo $X'\to X$ en el punto $x$ se reduce entonces a un solo punto, y $\sh{O}_{X',x'}$ es una $\sh{O}_{X,x}$-álgebra \emph{entera}; basta, pues, aplicar (0, 16.1.5).
\end{remark}

\begin{lemma}[18.9.7.7]
\label{IV.18.9.7.7}
Sean $A$ un anillo local noetheriano, de dimensión $\geq 2$, íntegro y geométricamente unirrama, $Y=\Spec(A)$, e $y$ el punto cerrado de $Y$. Sean $X$ un preesquema localmente noetheriano y conexo, y $f:X\to Y$ un morfismo plano. Entonces, para toda parte cerrada $T\subset f^{-1}(y)$, $X-T$ es conexo.
\end{lemma}

El mismo razonamiento que al comienzo de (18.9.7.2) reduce la cuestión a demostrar que, para todo punto $x\in f^{-1}(y)$, $\Spec(\sh{O}_{X,x})-\{x\}$ es conexo. Podemos, por tanto, limitarnos al caso en que $X=\Spec(B)$, donde $B$ es un anillo local noetheriano y el homomorfismo $A\to B$ correspondiente a $f$ es local; pongamos $U=Y-\{y\}$ y observemos que $V=f^{-1}(U)$ es denso en $X$ (2.3.10); basta, por tanto, demostrar que $V$ es \emph{conexo}. Conservemos las notaciones de (18.9.7.5) y pongamos, por tanto, $A^{(1)}=\Gamma(U,\sh{O}_Y)$, $Y_1=\Spec(A^{(1)})$, $X_1=X\times_Y Y_1=\Spec(B\otimes_A A^{(1)})$, de modo que se tiene el diagrama conmutativo
\[
  \xymatrix{
    X \ar[d]_{f} & X_1 \ar[l]^{g_1} \ar[d]^{f_1} \\
    Y & Y_1 \ar[l]_{g}
  }
\]

En virtud de (2.3.1) y de la definición de $A^{(1)}$, se tiene un isomorfismo canónico $\Gamma(V,\sh{O}_X)\xrightarrow{\sim}B\otimes_A A^{(1)}$, y basta, por tanto, demostrar que este último anillo es \emph{local}, pues tal anillo no puede ser producto de dos anillos no reducidos a 0 (III, 7.8.6.1). Pero se ha visto en (18.9.7.5) que $A^{(1)}$ es una sub-$A$-álgebra de la clausura integral $A'$ de $A$, puesto que $\dim(A)\geq 2$. La hipótesis de que $A$ es geométricamente unirrama implica que el morfismo $\Spec(A')\to\Spec(A)$ es radicial en el punto $y$ (6.15.3); \emph{a fortiori}, $g:Y_1\to Y$ es radicial en el punto $y$ y, por consiguiente, $g_1$ es radicial en el punto cerrado $x$ de $X$; como, además, $g_1$ es un morfismo entero, los puntos cerrados de $X_1$ están situados sobre $x$, por lo que solo existe un punto cerrado de $X_1$, lo que termina la demostración de (18.9.7.7).

Es claro que (18.9.7.7) demuestra (18.9.7.3) cuando $\dim(A)\geq 2$ (incluso sin suponer que las fibras $X_y$ sean geométricamente reducidas).

C) \emph{Caso en que $\dim(A)=1$.} Se ha visto al comienzo de la demostración de (18.9.7.3) que puede suponerse $\dim(B\otimes_A k)=0$ y, por planitud (6.1.1.1), se tiene, por tanto,
\[
  \dim(B)=\dim(A)+\dim(B\otimes_A k)=1.
\]
Como $A$ es íntegro, $Y$ se compone de dos puntos, el punto genérico $\eta$ y el punto cerrado $a$. Además, puesto que $f$ es plano, toda componente irreducible de $X$ domina $Y$ (2.3.4), por lo que $X-\{b\}$, que es el conjunto de los puntos maximales $\xi_i$ de $X$, es igual al

\oldpage[IV-4]{156}

conjunto subyacente a $f^{-1}(\eta)$, y la fibra $f^{-1}(\eta)$ es suma de los $\Spec(k(\xi_i))$. La hipótesis sobre los $X_y$ y el hecho de que estas fibras sean de dimensión 0 muestran que son \emph{geométricamente regulares} (6.7.6) y, \emph{a fortiori}, geométricamente normales; dicho de otro modo, $f$ es un \emph{morfismo normal}. Pero, como $A$ es íntegro y geométricamente unirrama, se deduce entonces de (6.15.10) que $B$ es íntegro y geométricamente unirrama, por lo que $f^{-1}(\eta)=X-\{b\}$ se reduce a un solo punto y, \emph{a fortiori}, es conexo; esto termina, por tanto, la demostración de (18.9.7.3) y la de (18.9.7).

\begin{remark}[18.9.7.8]
\label{IV.18.9.7.8}
En el caso C) de la demostración de (18.9.7.3), puede evitarse recurrir al delicado resultado (6.15.10) razonando de la manera siguiente: como $A$ es íntegro y de dimensión 1, se deduce del teorema de Krull--Akizuki (Bourbaki, \emph{Alg. comm.}, cap.~VII, \textsection~2, n\textsuperscript{o}~5, prop.~5) que su clausura integral $A'$ es un anillo \emph{noetheriano}. El mismo razonamiento que al comienzo de la demostración de (18.9.7) muestra entonces que $B'=B\otimes_A A'$ es un anillo local; además, por planitud, $B'$ está contenido en el anillo total de fracciones $R$ del anillo reducido $B$ (3.3.5). Ahora bien, el razonamiento que demuestra el teorema de Krull--Akizuki (Bourbaki, \emph{loc. cit.}) se aplica igualmente a un anillo local noetheriano \emph{reducido} de dimensión 1 y muestra que, para tal anillo $B$, todo anillo comprendido entre $B$ y su anillo total de fracciones es \emph{noetheriano}. Como el anillo $B'$ es un anillo local noetheriano y el morfismo $\Spec(B')\to\Spec(A')$ es normal, $B'$ es un anillo íntegro e íntegramente cerrado (6.5.4), y se concluye como al comienzo de la demostración de (18.9.7).
\end{remark}

\begin{corollary}[18.9.8]
\label{IV.18.9.8}
Con las notaciones de (18.9.7), supongamos que se verifica la condición (i) de (18.9.7) y una de las dos condiciones siguientes:
\begin{enumerate}
  \item[(ii$'$)] $A$ es un anillo estrictamente local (18.8.2).
  \item[(ii$''$)] $A$ es un anillo local henseliano y $k(B)$ es una extensión primaria de $k(A)$.
\end{enumerate}
Entonces todas las fibras $X_y$ ($y\in Y$) son geométricamente conexas.
\end{corollary}

Sea $\mathfrak p$ el ideal primo de $A$ correspondiente al punto $y$, y pongamos $A'=A/\mathfrak p$ y $B'=B\otimes_A A'=B/\mathfrak pB$, que son evidentemente anillos locales noetherianos; es claro que el morfismo $f':\Spec(B')\to\Spec(A')$ verifica la condición (i) de (18.9.7); como $y$ es el punto genérico de $\Spec(A')$ y $k(A')=k(A)$, $k(B')=k(B)$, se ve que basta verificar que, en el caso (ii$'$) (resp. (ii$''$)), $A'$ es geométricamente unirrama (resp. unirrama). Pero $A'$ es íntegro y, en el caso (ii$'$) (resp. (ii$''$)), es estrictamente local (resp. henseliano) en virtud de (18.5.10). Se concluye que $A'$ es geométricamente unirrama (resp. unirrama) en virtud de (18.8.15) (resp. (18.6.12)).

\begin{corollary}[18.9.9]
\label{IV.18.9.9}
Sea $A$ un anillo local noetheriano henseliano, universalmente japonés (\emph{i.e.} (7.6.4) y (7.7.2)), cuyas fibras formales (7.3.13) son geométricamente reducidas; entonces las fibras formales de $A$ son geométricamente conexas.
\end{corollary}

Basta aplicar (18.9.8) al caso en que $B=\widehat A$, puesto que $B$ es un $A$-módulo plano y que $k(B)=k(A)$.

\begin{remark}[18.9.10]
\label{IV.18.9.10}
Las demostraciones de (18.9.4) y (18.9.5) solo utilizan el hecho de que las fibras formales de $A$ son geométricamente conexas. Las conclusiones

\oldpage[IV-4]{157}

de estas dos proposiciones siguen siendo, por tanto, válidas cuando se supone que el anillo local $A$ es noetheriano, henseliano y universalmente japonés.
\end{remark}

\begin{corollary}[18.9.11]
\label{IV.18.9.11}
Sean $X$, $Y$ dos preesquemas localmente noetherianos, suponiéndose que $Y$ es geométricamente unirrama y que $X$ es conexo. Sea $f:X\to Y$ un morfismo reducido (\emph{i.e.}, plano y de fibras geométricamente reducidas (6.8.1)). Entonces, para toda parte cerrada $T$ de $X$ tal que $f(T)$ sea rara en $Y$, $X-T$ es conexo.
\end{corollary}

Como $f$ es plano y $f(T)$, al ser rara, no puede contener ningún punto maximal de $Y$, $T$ no puede contener ningún punto maximal de $X$ (2.3.4) y es, por tanto, rara en $X$. Utilizando (15.5.6.1), basta demostrar que, para todo $x\in T$, $\Spec(\sh{O}_{X,x})-\{x\}$ es conexo. Pongamos $y=f(x)$, que no es un punto maximal de $Y$. Como el anillo $\sh{O}_{Y,y}$ es geométricamente unirrama, puede aplicarse (18.9.7) al morfismo
\[
  f_1:\Spec(\sh{O}_{X,x})\longrightarrow\Spec(\sh{O}_{Y,y}),
\]
que verifica las condiciones (i) y (ii) de este teorema, y a la parte cerrada $\{x\}$ de $\Spec(\sh{O}_{X,x})$,
puesto que $f_1(\{x\})=\{y\}$ es rara en $\Spec(\sh{O}_{Y,y})$.

\subsection{Preesquemas étales sobre un preesquema geométricamente unirrama o normal}
\label{subsection:IV.18.10}

\begin{theorem}[18.10.1]
\label{IV.18.10.1}
Sean $f:X\to Y$ un morfismo localmente de presentación finita, $x$ un punto de $X$, $y=f(x)$. Supongamos que $Y$ sea íntegro y geométricamente unirrama en el punto $y$. Entonces, para que $f$ sea étale en el punto $x$, es necesario y suficiente que se verifiquen las dos condiciones siguientes:
\begin{enumerate}
  \item[(i)] $f$ es no ramificado en el punto $x$;
  \item[(ii)] el homomorfismo $\sh{O}_{Y,y}\to\sh{O}_{X,x}$ es inyectivo.
\end{enumerate}
Además, si así sucede, $X$ es íntegro y geométricamente unirrama en el punto $x$.

Finalmente, si $\sh{O}_{Y,y}$ es noetheriano, en el criterio anterior se puede reemplazar la condición (ii) por la condición

\emph{(ii bis)} Se tiene $\dim\sh{O}_{Y,y}\leq\dim\sh{O}_{X,x}$.
\end{theorem}

El hecho de que, si $f$ es étale, $X$ sea íntegro y geométricamente unirrama en el punto $x$ es un caso particular de (17.5.7). Por otra parte, es claro que las condiciones (i) y (ii) son necesarias para que $f$ sea étale en el punto $x$, pues $\sh{O}_{X,x}$ es entonces un $\sh{O}_{Y,y}$-módulo fielmente plano (17.6.1). Probemos que estas condiciones son suficientes.

Pongamos $A=\sh{O}_{Y,y}$, $A'={}^{hs}\!A$ (18.8.7), $Y'=\Spec(A')$, $X'=X\times_Y Y'$, $f'=f_{(Y')}:X'\to Y'$. Sean $y'$ el punto cerrado de $Y'$ y $x'$ un punto de $X'$ situado sobre $x$ y sobre $y'$; puesto que el morfismo $g:Y'\to Y$ es plano (18.8.8), todo se reduce a ver que $f'$ es étale en el punto $x'$ (17.7.1, (iii)). La hipótesis (ii) implica que el morfismo $\Spec(\sh{O}_{X,x})\to\Spec(\sh{O}_{Y,y})$ es dominante (I, 1.2.7); como $\sh{O}_{Y,y}$ es íntegro, $\Spec(\sh{O}_{Y,y})$ tiene un único punto genérico $y_1$, y por tanto existe una generalización $x_1$ de $x$ en $X$ situada sobre $y_1$. Puesto que el morfismo proyección $X'\to X$ es plano, existe una generalización $x'_1$ de $x'$ situada sobre $x_1$ (2.3.4); pongamos $y'_1=f'(x'_1)$; se tiene entonces $g(y'_1)=y_1$. Ahora bien, la hipótesis sobre $y$ implica que $A'$ es íntegro (18.8.15), de modo que $Y'$ tiene un único punto genérico $\eta$, necesariamente situado sobre $y_1$, puesto que $g$

\oldpage[IV-4]{158}

es plano (2.3.4); por otra parte, $\eta$ es también el punto genérico de $g^{-1}(y_1)$ ($0_{\mathrm I}$, 2.1.8), y se sabe además que las fibras de $g$ son espacios discretos (18.8.12), luego $g^{-1}(y_1)=\{\eta\}$ y, por consiguiente, $y'_1=\eta$. Por tanto, el morfismo $\Spec(\sh{O}_{X',x'})\to\Spec(\sh{O}_{Y',y'})$, restricción de $f'$, es dominante, y como $\sh{O}_{Y',y'}=A'$ es íntegro, el homomorfismo correspondiente $\sh{O}_{Y',y'}\to\sh{O}_{X',x'}$ es inyectivo (I, 1.2.7). Por otra parte, $f'$ es no ramificado en el punto $x'$; por (18.8.3, $d$), existe un entorno $U$ de $x'$ en $X'$ tal que $f'|U$ sea una inmersión cerrada de $U$ en $Y'$; \emph{a fortiori}, el homomorfismo $\sh{O}_{Y',y'}\to\sh{O}_{X',x'}$ es sobreyectivo y, puesto que es inyectivo, es biyectivo. Pero entonces $\sh{O}_{X',x'}$ es un $\sh{O}_{Y',y'}$-módulo plano, y por tanto $f'$ es étale en el punto $x'$ (17.6.1).

Finalmente, cuando $\sh{O}_{Y,y}$ es noetheriano, sabemos que si $f$ es étale en el punto $x$, se tiene $\dim\sh{O}_{Y,y}=\dim\sh{O}_{X,x}$ (17.6.4). Recíprocamente, supongamos verificadas las condiciones (i) y (ii \emph{bis}), y demostremos que implican (ii). Sea $\mathfrak J$ el núcleo del homomorfismo canónico $\sh{O}_{Y,y}\to\sh{O}_{X,x}$; restringiendo $Y$ a un entorno de $y$ si es necesario, podemos suponer que $\mathfrak J=\mathscr{J}_y$, donde $\mathscr{J}$ es un ideal de $\sh{O}_Y$; si $Y_1$ es el subpreesquema cerrado de $Y$ definido por $\mathscr{J}$, restringiendo aún más $X$ e $Y$ a entornos de $x$ y $y$, respectivamente, podemos suponer que $f$ se factoriza como $X\xrightarrow{f_1}Y_1\to Y$ (I, 6.5.1). Es claro que $f_1$ sigue siendo no ramificado en el punto $x$, luego es cuasi-finito en ese punto por (17.4.1); la demostración de (5.4.1, (i)) prueba entonces que $\dim(\sh{O}_{X,x})\leq\dim(\sh{O}_{Y,y}/\mathfrak J)$. Pero si $\mathfrak J\neq0$, se deduciría que $\dim(\sh{O}_{Y,y}/\mathfrak J)<\dim(\sh{O}_{Y,y})$, puesto que $Y$ es íntegro ($0$, 16.1.2.2); se tendría por tanto $\dim\sh{O}_{X,x}<\dim\sh{O}_{Y,y}$, en contradicción con la hipótesis, lo que demuestra (ii).

\begin{remark}[18.10.2]
\label{IV.18.10.2}
\begin{enumerate}
  \item[(i)] Cuando $A=\sh{O}_{Y,y}$ es noetheriano y se sabe que su completado $\widehat A$ es íntegro (por ejemplo, si $A$ es regular), en la demostración anterior se puede reemplazar $A'$ por $\widehat A$.
  \item[(ii)] Cuando $Y$ es localmente noetheriano, se puede dar de (18.10.1) una demostración más rápida, que no utiliza la henselización estricta, pero que hace intervenir los resultados bastante delicados de las \textsection\textsection~14 y 15. Como $f$ es cuasi-finito en el punto $x$ por hipótesis (17.4.1), reemplazando $X$ por un entorno abierto de $x$ podemos suponer que $f^{-1}(y)=\{x\}$. La hipótesis (ii) implica, como hemos visto, que existe una componente irreducible $Z$ de $X$ que contiene a $x$ y domina la única componente irreducible $Y_0$ de $Y$ que contiene a $y$. Puesto que el morfismo $f|Z$ es cuasi-finito en el punto $x$, y por tanto equidimensional en ese punto (13.2.2), se deduce de la hipótesis sobre $Y$ y del criterio de Chevalley (14.4.4) que $g=f|Z$ es universalmente abierto en el punto $x$. Además, como $f$ (y por tanto también $g$) es no ramificado en el punto $x$, $g^{-1}(y)$ es geométricamente reducido sobre $k(y)$ (17.4.1); como $\sh{O}_{Y,y}$ es íntegro, se concluye de (15.2.3) que $g$ es plano en el punto $x$, y por tanto $f$ es étale en ese punto (18.4.9).
  \item[(iii)] Si las notaciones y las hipótesis sobre $Y$ son las de (18.10.1), supongamos solamente que $f$ sea localmente de tipo finito (y no necesariamente localmente de presentación finita). Entonces, para que $\sh{O}_{X,x}$ sea una $\sh{O}_{Y,y}$-álgebra esencialmente étale (18.6.1), es necesario y suficiente que $f$ satisfaga las dos condiciones siguientes:

\oldpage[IV-4]{159}

\begin{enumerate}
  \item[$1^\circ$] $f$ es formalmente no ramificado en el punto $x$;
  \item[$2^\circ$] el homomorfismo $\sh{O}_{Y,y}\to\sh{O}_{X,x}$ es inyectivo.
\end{enumerate}

Solo queda por demostrar la suficiencia de estas condiciones. Teniendo en cuenta (18.4.12), basta demostrar que $f$ es plano en el punto $x$. Ahora bien, retomando la demostración de (18.10.1) y las notaciones utilizadas en ella, basta demostrar que $f'$ es plano en el punto $x'$ (2.5.1); por otra parte, $f'$ es formalmente no ramificado en el punto $x'$ (17.1.3). Por tanto, podemos limitar la demostración al caso en que $Y=\Spec(A)$, donde $A$ es estrictamente local y $y=f(x)$ es el punto cerrado de $Y$. Usando entonces (18.8.3), vemos que el homomorfismo $\sh{O}_{Y,y}\to\sh{O}_{X,x}$ es sobreyectivo, luego biyectivo por la hipótesis, y esto basta para demostrar que $f$ es plano en el punto $x$.

Supongamos además que $Y$ sea localmente íntegro en el punto $y$ (I, 2.1.8). Entonces las condiciones $1^\circ$ y $2^\circ$ anteriores (junto con el hecho de que $Y$ sea geométricamente unirrama en el punto $y$ y que $f$ sea localmente de tipo finito) ya implican que $f$ es étale en el punto $x$. En efecto, esto se deduce de lo anterior y de (18.4.13).
\end{enumerate}
\end{remark}

\begin{corollary}[18.10.3]
\label{IV.18.10.3}
Sean $Y$ un preesquema íntegro y geométricamente unirrama, $\eta$ su punto genérico, $X$ un preesquema conexo, $f:X\to Y$ un morfismo localmente de tipo finito y tal que $f^{-1}(\eta)$ no sea vacío. Entonces, si $f$ es formalmente no ramificado, $f$ es étale y $X$ es íntegro y geométricamente unirrama.
\end{corollary}

De la hipótesis y de (18.4.13) se deduce que $f$ es étale en todos los puntos donde es plano, y en particular en los puntos de $f^{-1}(\eta)$, puesto que $\sh{O}_{Y,\eta}=k(\eta)$ es un cuerpo. Por tanto, el conjunto abierto $U$ de los puntos de $X$ donde $f$ es étale no es vacío. Mostremos, por otra parte, que para todo $x\in\overline U$, el homomorfismo $\sh{O}_{Y,f(x)}\to\sh{O}_{X,x}$ es inyectivo. Sea, en efecto, $s$ una sección de $\sh{O}_Y$ sobre un entorno abierto de $f(x)$, que siempre podemos suponer que es el propio $Y$ (la cuestión es local en $X$ e $Y$), y supongamos que su imagen $t\in\Gamma(Y,f_*(\sh{O}_X))=\Gamma(X,\sh{O}_X)$ tenga germen nulo en el punto $x$, de modo que se anula en un entorno abierto $V$ de $x$; entonces la restricción de $t$ al abierto $U\cap V\neq\varnothing$ es nula. Pero la restricción $f|(U\cap V)$ es étale, luego $W=f(U\cap V)$ es abierto en $Y$ y contiene por consiguiente el punto genérico $\eta$. Como el homomorfismo $\sh{O}_{Y,\eta}\to\sh{O}_{X,z}$ es inyectivo para todo punto $z\in f^{-1}(\eta)$, la hipótesis $t|(U\cap V)=0$ implica que $s_\eta=0$; y como $Y$ es íntegro, esto implica $s=0$, lo que prueba nuestra afirmación. Aplicando ahora (18.10.2, (iii)), concluimos que $f$ es plano en el punto $x$, es decir, $x\in U$, y por tanto $U$ es a la vez abierto y cerrado en $X$; como no es vacío y $X$ es conexo, se tiene $U=X$.

Observemos ahora que, puesto que $f$ es étale, luego localmente cuasi-finito, y $Y$ es irreducible, los puntos máximos de $X$ pertenecen a $f^{-1}(\eta)$ (2.3.4), y para todo $x\in X$ existe un entorno de $x$ que contiene solamente un número finito de esos puntos máximos; en otras palabras, el conjunto de las componentes irreducibles de $X$ es localmente finito. Pero por (18.10.1), $X$ es íntegro y geométricamente unirrama en todo punto; por tanto, cada punto de $X$ pertenece a una sola componente irreducible y, puesto que estas forman una familia localmente finita, son abiertas (y cerradas) en $X$; como $X$ es conexo, es íntegro.

\oldpage[IV-4]{160}

\begin{remark}[18.10.3.1]
\label{IV.18.10.3.1}
Si $f:X\to Y$ es localmente de presentación finita, dominante y cuasicompacto, la fibra $f^{-1}(\eta)$ no es vacía por (I, 1.1.5). En cambio, si no suponemos que $f$ sea cuasicompacto, $f$ puede ser no ramificado y dominante sin que $f$ sea étale (manteniéndose siempre las hipótesis de (18.10.3) para $X$ e $Y$). Esto es lo que muestra el ejemplo siguiente: tomemos para $Y$ el plano afín $\Spec(\mathbf C[T,U])$; consideremos por otra parte una familia $(X_j)_{j\in\mathbf Z}$ de preesquemas isomorfos a la recta afín $\Spec(\mathbf C[T])$, y formemos el preesquema obtenido «pegando» $X_{2j}$ y $X_{2j+1}$ en el punto $-1$, y $X_{2j+1}$ y $X_{2j+2}$ en el punto $+1$. Definamos además $f:X\to Y$ como igual, sobre $X_{2j}$, a la inmersión cerrada cuya imagen es la recta de ecuación $y=2j$, que transforma el punto $-1$ en $(2j-1,2j)$ y el punto $+1$ en $(2j+1,2j)$; sobre $X_{2j+1}$, $f_j$ es la inmersión cerrada cuya imagen es la recta de ecuación $x=2j+1$, que transforma el punto $-1$ en $(2j+1,2j)$ y el punto $+1$ en $(2j+1,2j+2)$.\footnote{La transcripción francesa imprime aquí $-1$ en la segunda aparición; el testigo inglés mantenido y el pegado de las fibras requieren $+1$.}
Es claro que $f$ es una inmersión local (por tanto, no ramificada) y es dominante, pero la fibra del punto genérico de $Y$ es vacía y $f$ no es étale.
\end{remark}

\begin{corollary}[18.10.4]
\label{IV.18.10.4}
Sean $g:Y\to S$, $h:X\to S$ dos morfismos localmente de presentación finita, $f:X\to Y$ un $S$-morfismo, $x$ un punto de $X$, $y=f(x)$, $s=g(y)=h(x)$. Supongamos que $h$ sea plano en el punto $x$ y que $g^{-1}(s)$ sea normal en el punto $y$. Entonces, para que $f$ sea étale en el punto $x$, es necesario y suficiente que $f$ sea no ramificado en el punto $x$ y que $\dim_xh^{-1}(s)=\dim_yg^{-1}(s)$.
\end{corollary}

Las condiciones son necesarias por el hecho de que si $f$ es étale en el punto $x$, lo mismo ocurre con $f_s:h^{-1}(s)\to g^{-1}(s)$, y la conclusión se deduce entonces de que $f_s$ es plano y cuasi-finito en el punto $x$ (6.1.2). Recíprocamente, si se verifican las condiciones del enunciado, para ver que $f$ es étale en el punto $x$ basta, en virtud de (17.8.3), ver que lo es $f_s$. Por tanto, podemos limitarnos al caso en que $S$ es el espectro de un cuerpo. De la hipótesis $\dim_xX=\dim_yY$ y del hecho de que $f$ es no ramificado en el punto $x$ se deduce entonces (por (5.2.3) y (17.4.1)) que $\dim\sh{O}_{X,x}=\dim\sh{O}_{Y,y}$; puesto que $f$ es cuasi-finito en el punto $x$ y $\sh{O}_{Y,y}$ es íntegro, de ($0_{\mathrm I}$, 7.4.4) y ($0$, 16.3.10) se deduce que el homomorfismo $\sh{O}_{Y,y}\to\sh{O}_{X,x}$ es inyectivo; basta entonces aplicar (18.10.1) para concluir.

\begin{corollary}[18.10.5]
\label{IV.18.10.5}
Con las notaciones de (18.10.4), supongamos que $h$ sea plano y que $g^{-1}(s)$ sea normal para todo $s\in S$ (lo que, por ejemplo, sucede si $g$ es liso). Para que $f$ sea una inmersión abierta, es necesario y suficiente que $f$ sea un monomorfismo y que, para todo $x\in X$, se tenga $\dim_xh^{-1}(s)=\dim_yg^{-1}(s)$, donde $y=f(x)$ y $s=g(y)=h(x)$.
\end{corollary}

Solo hay que demostrar la suficiencia de las condiciones; ahora bien, de (17.1.3) se deduce que un monomorfismo localmente de presentación finita es no ramificado, de modo que las hipótesis implican que $f$ es étale (18.10.4); se concluye entonces por (17.9.1).

\begin{lemma}[18.10.6]
\label{IV.18.10.6}
Sea $f:X\to Y$ un morfismo plano y localmente cuasi-finito ($\mathrm{Err}_{\mathrm{III}}$, 20).
\begin{enumerate}
  \item[(i)] El conjunto $\operatorname{Max}(X)$ de los puntos máximos de $X$ está en correspondencia biyectiva canónica con el conjunto $\coprod_{y\in\operatorname{Max}(Y)}f^{-1}(y)$.
  \item[(ii)] Si $Y$ es irreducible y tiene punto genérico $\eta$, el conjunto de las componentes irreducibles de $X$ está en correspondencia biyectiva canónica con $f^{-1}(\eta)$ y, en particular, es finito si y solo si $f^{-1}(\eta)$ es finito.

\oldpage[IV-4]{161}

  \item[(iii)] Si el conjunto de las componentes irreducibles de $Y$ es localmente finito, lo mismo sucede con el conjunto de las componentes irreducibles de $X$ y, en particular, $X$ es localmente conexo.
\end{enumerate}
\end{lemma}

La afirmación (i) se deduce inmediatamente de (2.3.4) y ($0_{\mathrm I}$, 2.1.6), y del hecho de que las fibras $f^{-1}(y)$ son conjuntos discretos. Es claro que (ii) es un caso particular de (i). Finalmente, para demostrar (iii), en virtud de (i) podemos limitarnos al caso en que $f$ es cuasi-finito y $Y$ es irreducible ($0_{\mathrm I}$, 2.1.6); la conclusión se deduce entonces de (ii) y del hecho de que las fibras de $f$ son conjuntos finitos (II, 6.2.2).

\begin{proposition}[18.10.7]
\label{IV.18.10.7}
Sean $Y$ un preesquema geométricamente unirrama e irreducible (resp. íntegro), $f:X\to Y$ un morfismo étale. Entonces $X$ es isomorfo a una suma de preesquemas irreducibles (resp. íntegros) $X_\lambda$, donde $\lambda$ recorre la fibra $f^{-1}(\eta)$ del punto genérico $\eta$ de $Y$. Los $X_\lambda$ son geométricamente unirrama; si $Y$ es normal, los $X_\lambda$ son normales.
\end{proposition}

En efecto, $X$ es geométricamente unirrama (17.5.7), por lo que cada punto de $X$ pertenece a una sola componente irreducible; además, en virtud de (18.10.6) y (17.6.1), el conjunto de las componentes irreducibles de $X$ es localmente finito, luego ($0_{\mathrm I}$, 2.1.6) son las componentes conexas $X_\lambda$ de $X$ ($\lambda\in f^{-1}(\eta)$), y $X$ es la suma de los $X_\lambda$; por (17.5.7), si $Y$ es íntegro (resp. normal), los $X_\lambda$ son íntegros (resp. normales).

\begin{proposition}[18.10.8]
\label{IV.18.10.8}
Sean $Y$ un preesquema normal e íntegro, de punto genérico $\eta$, $K=k(\eta)$ su cuerpo de funciones racionales, y $f:X\to Y$ un morfismo étale. Supongamos que $f^{-1}(\eta)$ sea finita y no vacía, de modo que el $K$-esquema $f^{-1}(\eta)$ sea el espectro de una $K$-álgebra finita separable $L$, producto directo de un número finito de cuerpos $K_i$ ($1\leq i\leq n$), extensiones finitas separables de $K$ (17.6.2). Sea $Y'$ la clausura integral de $Y$ en $L$, isomorfa a la suma de las clausuras integrales $Y_i$ de $Y$ en los $K_i$ (II, 6.3.6). Entonces el morfismo $f$ se factoriza de una única manera como $f:X\xrightarrow{f'}Y'\xrightarrow{g}Y$, de modo que la restricción $f'_\eta:f^{-1}(\eta)\to g^{-1}(\eta)$ sea el isomorfismo canónico. Además, $f'$ es un isomorfismo local; para que $f'$ sea una inmersión abierta, es necesario y suficiente que $f$ sea separado.
\end{proposition}

En virtud de (18.10.7), podemos limitarnos al caso en que $X$ sea íntegro y normal, $L$ su cuerpo de funciones racionales y $f$ dominante. La existencia y unicidad de la factorización considerada de $f$ se deducen entonces de (II, 6.3.9). Como $f'$ es birracional y localmente cuasi-finito, las últimas afirmaciones se deducen de (8.12.10).

\begin{corollary}[18.10.9]
\label{IV.18.10.9}
Bajo las hipótesis de (18.10.8), para que $f$ sea un morfismo finito (es decir, para que $X$ sea un recubrimiento étale de $Y$), es necesario y suficiente que $X$ sea isomorfo a la clausura integral $Y'$ de $Y$ en $L$.
\end{corollary}

Si $f$ es finito, la inyección canónica $j:X\to Y'$ (que es una inmersión abierta) es también un morfismo finito (II, 6.1.5, (v)), luego un morfismo cerrado (II, 6.1.10), y como $j(X)$ es denso en $Y'$ por (18.10.8), se tiene $j(X)=Y'$. Recíprocamente, como $Y$ es normal y $L$ es un producto directo de extensiones finitas separables de $K$, $Y'$ es finito sobre $Y$ (Bourbaki, \emph{Alg. comm.}, cap.~V, \textsection~1, n\textsuperscript{o}~6, cor.~1 de la prop.~18), de donde se sigue el corolario.

\begin{env}[18.10.10]
\label{IV.18.10.10}
Sea $Y$ un preesquema normal e íntegro, $K=R(Y)$ su cuerpo de funciones racionales. Diremos que una $K$-álgebra $L$ de rango finito es no ramificada sobre $Y$ si: $1^\circ$ $L$ es una $K$-álgebra separable, es decir, un producto directo de extensiones finitas separables $K_i$ de $K$ ($1\leq i\leq n$); $2^\circ$ la clausura integral $Y'$ de $Y$ en $L$ (suma de los preesquemas clausuras integrales de $Y$ en los $K_i$) es no ramificada sobre $Y$ (lo que, por (18.10.3),

\oldpage[IV-4]{162}

equivale a decir que $Y'$ es étale sobre $Y$). Cuando $Y=\Spec(A)$, donde $A$ es un anillo íntegro e integralmente cerrado y $K$ su cuerpo de fracciones, diremos también (por abuso de lenguaje, y cuando no haya riesgo de confusión con la terminología de (17.3.2, (ii))) que $L$ es no ramificada sobre $A$ en lugar de «no ramificada sobre $Y$».
\end{env}

\begin{remark}[18.10.11]
\label{IV.18.10.11}
En lugar de decir que $L$ es no ramificada sobre $Y$, algunos autores dicen también que $L$ es «no ramificada sobre $K$»; nos abstendremos de seguir este uso, que da lugar a confusiones.
\end{remark}

Se puede entonces expresar (18.10.9) de la forma siguiente:

\begin{corollary}[18.10.12]
\label{IV.18.10.12}
Sean $Y$ un preesquema normal e íntegro, $K=R(Y)$ su cuerpo de funciones racionales. Entonces el functor $X\rightsquigarrow R(X)$ es una equivalencia entre la categoría de los recubrimientos étales de $Y$ y la categoría de las $K$-álgebras finitas étales no ramificadas sobre $Y$. Se obtiene un functor cuasi-inverso haciendo corresponder a toda $K$-álgebra finita étale $L$, no ramificada sobre $Y$, la clausura integral $Y'$ de $Y$ en $L$.
\end{corollary}

\begin{proposition}[18.10.13]
\label{IV.18.10.13}
Sean $Y$ un preesquema normal e íntegro, $K=R(Y)$ su cuerpo de funciones racionales.
\begin{enumerate}
  \item[(i)] El cuerpo $K$ es una $K$-álgebra no ramificada sobre $Y$.
  \item[(ii)] Sea $L$ una extensión finita de $K$, no ramificada sobre $Y$, y sea $Z$ la clausura integral de $Y$ en $L$. Entonces, si $M$ es una $L$-álgebra no ramificada sobre $Z$, $M$ es también una $K$-álgebra no ramificada sobre $Y$ («transitividad» de la no ramificación).
  \item[(iii)] Sean $Y'$ un preesquema normal e íntegro, $K'$ su cuerpo de funciones racionales, y $g:Y'\to Y$ un morfismo dominante. Si $L$ es una $K$-álgebra no ramificada sobre $Y$, entonces $L\otimes_K K'$ es una $K'$-álgebra no ramificada sobre $Y'$ (propiedad de «traslación»). Además, si $Y=\Spec(A)$ y $Y'=\Spec(A')$ son afines, y $C$ es la clausura integral de $A$ en $L$, entonces $C'=C\otimes_A A'$ es la clausura integral de $A'$ en $L\otimes_K K'$.
\end{enumerate}
\end{proposition}

La afirmación (i) es inmediata, pues $Y$ es por hipótesis su propia normalización en $K$. Para demostrar (ii), sea $Z'$ la clausura integral de $Z$ en $M$, que por hipótesis es un recubrimiento étale de $Z$. Como $Z$ es étale y separado sobre $Y$, lo mismo sucede con $Z'$ (17.3.3), que evidentemente es la clausura integral de $Y$ en $M$. Como $M$ es una $K$-álgebra finita y separable, se deduce de (18.10.10) que $M$ es no ramificada sobre $Y$. Finalmente, para demostrar (iii), observemos que si $Z$ es la clausura integral de $Y$ en $L$, entonces $Z\times_Y Y'$ es étale y separado sobre $Y'$ (17.3.3), finito sobre $Y'$ porque $Z$ es finito sobre $Y$, y tiene como anillo de funciones racionales $L\otimes_K K'$ (I, 3.4.9); como $Y'$ es normal, lo mismo sucede con $Z\times_Y Y'$ (17.5.7), que es por tanto la clausura integral de $Y'$ en $L\otimes_K K'$, lo que termina la demostración.

\begin{corollary}[18.10.14]
\label{IV.18.10.14}
Sean $Y$ y $Y'$ dos preesquemas normales e íntegros, $K$, $K'$ sus respectivos cuerpos de funciones racionales, y $g:Y'\to Y$ un morfismo dominante, de modo que $K'$ es una extensión de $K$. Sean $L$ una extensión finita de $K$ no ramificada sobre $Y$, y $L_1$ una extensión compuesta (Bourbaki, \emph{Alg.}, cap.~VIII, \textsection~8, def.~1) de $L$ y $K'$. Entonces $L_1$ es una extensión de $K'$ no ramificada sobre $Y'$. Además, si $Y=\Spec(A)$, $Y'=\Spec(A')$ son afines, y $C$ es la clausura integral de $A$ en $L$, entonces el subanillo $C_1=A[C,A']$ de $L_1$ es la clausura integral de $A'$ en $L_1$.
\end{corollary}

En efecto, $L\otimes_K K'$ es un producto directo de extensiones finitas separables de $K'$, y $L_1$

\oldpage[IV-4]{163}

es una de estas extensiones; por tanto, la clausura integral de $Y'$ en $L\otimes_K K'$ es una suma de preesquemas, uno de los cuales es la clausura integral de $Y'$ en $L_1$. El corolario se deduce inmediatamente de (18.10.13).

Cuando, por ejemplo, $A$ es el anillo de los enteros de un cuerpo de números algebraicos $K$, y $K'$ y $L$ son extensiones algebraicas de $K$, hay ejemplos clásicos que muestran que la relación $C_1=A[C,A']$ falla cuando $L$ no es no ramificada sobre $A$ [\ ].

\begin{env}[18.10.15]
\label{IV.18.10.15}
Supongamos que $Y=\Spec(A)$ sea afín, que $A$ sea íntegro e integralmente cerrado; sean $K$ su cuerpo de fracciones, $L$ una extensión finita separable de $K$, y supongamos que la clausura integral $C$ de $A$ en $L$ sea un $A$-módulo proyectivo de tipo finito (lo que, por ejemplo, ocurre si $A$ es un anillo de Dedekind (Bourbaki, \emph{Alg. comm.}, cap.~VII, \textsection~4, n\textsuperscript{o}~10, prop.~22)). Decir que $L$ es no ramificada sobre $A$ significa que $C$ es étale (18.10.10) y, en virtud de (18.2.7, (ii)), esto equivale a decir que el discriminante $d_{C/A}$ de $\Spec(C)$ sobre $\Spec(A)$ es invertible en $A$. En el caso particular en que $C$ es un $A$-módulo libre y $(x_i)_{1\leq i\leq n}$ es una base de $C$ sobre $A$, esto significa que $\det(\operatorname{Tr}_{C/A}(x_ix_j))$ es invertible en $A$.
\end{env}

\begin{theorem}[18.10.16]
\label{IV.18.10.16}
Sean $Y$ un preesquema íntegro, $\eta$ su punto genérico, $f:X\to Y$ un morfismo separado y localmente de tipo finito. Supongamos que toda componente irreducible de $X$ domina $Y$ y que la fibra genérica $X_\eta=f^{-1}(\eta)$ es un preesquema finito sobre $K=k(\eta)$, de modo que $(X_\eta)_{\mathrm{red}}$ es igual a $\Spec(L)$, donde $L=\prod_iL_i$ es un producto de extensiones finitas $L_i$ de $K$. Pongamos
\[
  n=[L:K]=\sum_i[L_i:K],\qquad n_s=\sum_i[L_i:K]_s \quad\text{(suma de los grados separables de los $L_i$)}.
\]
Sea $y$ un punto geométricamente unirrama de $Y$, y designemos por $n(y)$ la suma de los grados separables sobre $k(y)$ de los cuerpos residuales de los puntos aislados de $f^{-1}(y)$. Entonces:
\begin{enumerate}
  \item[(i)] Se tiene $n(y)\leq n_s\leq n$.
  \item[(ii)] Supongamos que $X$ sea reducido y que el punto $y$ sea normal en $Y$. Para que $n(y)=n$, es necesario y suficiente que exista un entorno abierto $U$ de $y$ en $Y$ tal que la restricción $f^{-1}(U)\to U$ de $f$ sea un morfismo étale y finito.
\end{enumerate}
\end{theorem}

\emph{A) Reducción al caso en que $f$ es cuasi-finito.} --- Se sabe (13.1.4) que el conjunto $Z$ de los puntos $x\in X$ aislados en $f^{-1}(f(x))$ es abierto en $X$ y, por hipótesis, $Z$ contiene a $f^{-1}(\eta)$. Como $f^{-1}(\eta)$ es finita, $Z$ es la unión filtrante creciente de abiertos cuasicompactos $V_\lambda$ que contienen $f^{-1}(\eta)$ (y por tanto son densos en $X$ por hipótesis). Si $n_\lambda(y)$ es la suma de los grados separables sobre $k(y)$ de los cuerpos residuales de los puntos de $V_\lambda\cap f^{-1}(y)$, se tiene $n(y)=\sup_\lambda n_\lambda(y)$; para demostrar (i), bastará mostrar que $n_\lambda(y)\leq n_s$. Por otra parte, si $n(y)=n$, existe un índice $\lambda$ tal que $n_\lambda(y)=n$. Supongamos demostrado que entonces existe un entorno abierto $U$ de $y$ en $Y$ tal que la restricción $V_\lambda\cap f^{-1}(U)\to U$ de $f$ sea un morfismo étale y finito. Como $f$ es separado, la inyección canónica $V_\lambda\cap f^{-1}(U)\to f^{-1}(U)$ es también un morfismo finito (II, 6.1.5, (v)), de modo que $V_\lambda\cap f^{-1}(U)$ es cerrado en $f^{-1}(U)$ (II, 6.1.10); pero, al ser denso en todo $f^{-1}(U)$, necesariamente coincide con él.

\emph{B) Reducción al caso en que $X$ es reducido, $f$ finito y $Y=\Spec(\sh{O}_{Y,y})$.} --- Observemos que la hipótesis sobre $f$ implica que también se verifica para $f_{\mathrm{red}}$ y que la conclusión de (i) solo concierne a $f_{\mathrm{red}}$. Por tanto, podemos reemplazar en todas partes $f$ por $f_{\mathrm{red}}$ y suponer que $X$ es reducido.

\oldpage[IV-4]{164}

Por otra parte, en virtud de (18.12.13)\footnote{El lector comprobará que (18.10.16) no se utiliza en la demostración de (18.12.13). Si se supone que $f$ es localmente de presentación finita (resp. cuasi-proyectivo), se puede reemplazar (18.12.13) por (8.12.8) (resp. (8.12.6)).}, se puede (después de reemplazar $Y$ por un entorno abierto afín $V$ de $y$ y $f$ por su restricción $f^{-1}(V)\to V$) factorizar $f$ como $X\xrightarrow{j}X'\xrightarrow{g}Y$, donde $g$ es finito y $j$ una inmersión abierta. Como $X$ es reducido, podemos reemplazar $X'$ por el subpreesquema cerrado reducido cuyo espacio subyacente es $\overline{j(X)}$ (I, 5.2.2), y suponer así que $X'$ es reducido y que toda componente irreducible de $X'$ domina $Y$. Además, $j(X)$ es un abierto denso de $X'$ y, por hipótesis, tiene solo un número finito de componentes irreducibles; por tanto, ($0_{\mathrm I}$, 1.2.7), las componentes irreducibles de $X'$ son las clausuras de las de $j(X)$, y la fibra $g^{-1}(\eta)$ se identifica con $f^{-1}(\eta)$. Si suponemos demostrado el teorema cuando además $f$ es un morfismo finito, podremos aplicarlo a $g$, y como $f^{-1}(y)$ se identifica con un subpreesquema de $g^{-1}(y)$, la relación (i) para $g$ implica la misma relación para $f$. Supongamos, por otra parte, que $n(y)=n$; entonces, por lo anterior, podemos aplicar la conclusión de (ii) al morfismo $g$ y, reemplazando $Y$ por un entorno abierto de $y$ si es necesario, suponer que $g$ es étale, y por tanto también $f$; además, se tiene entonces $f^{-1}(y)=g^{-1}(y)$. Como el morfismo $g$ es cerrado y $j(X)$ es abierto en $X'$, existe un entorno $U$ de $y$ en $Y$ tal que $j(X)\cap g^{-1}(U)=j(f^{-1}(U))$, lo que demuestra que (ii) también es cierta para $f$.

Así, nos reducimos al caso en que $f$ es además un morfismo finito. Es inmediato que, para demostrar (i), podemos limitarnos al caso $Y=\Spec(\sh{O}_{Y,y})$. Mostremos que lo mismo vale para demostrar (ii). En efecto, supongamos probado que la condición $n(y)=n$ implica que $X\times_Y\Spec(\sh{O}_{Y,y})$ es étale y finito sobre $\Spec(\sh{O}_{Y,y})$; como esto implica que $f$ es plano y formalmente no ramificado en los puntos de $f^{-1}(y)$, y $Y$ es íntegro, se concluye de (18.4.13) que $f$ es étale en los puntos de $f^{-1}(y)$, y por tanto también en un entorno abierto $V$ de $f^{-1}(y)$ en $X$; pero como $f$ es un morfismo finito, luego cerrado, $V$ contiene un abierto de la forma $f^{-1}(U)$, donde $U$ es un entorno abierto de $y$ en $Y$.

\emph{C) Reducción al caso en que $Y=\Spec(\sh{O}_{Y,y})$ y $\sh{O}_{Y,y}$ es estrictamente local.} --- Pongamos $A=\sh{O}_{Y,y}$, y sea $A'={}^{hs}\!A$ la henselización estricta de $A$ (18.8.7), que es un anillo local íntegro y geométricamente unirrama (18.8.16). Sea $Y'=\Spec(A')$, y sean $y'$ el punto cerrado y $\eta'$ el punto genérico de $Y'$; $y'$ está sobre $y$ y, puesto que el morfismo $g:Y'\to Y$ es plano (18.8.8), $\eta'$ está sobre $\eta$ (2.3.4). Pongamos $X'=X\times_Y Y'$, $f'=f_{(Y')}:X'\to Y'$; $f'$ es finito y, como $g$ es plano, toda componente irreducible de $X'$ domina $Y'$ (2.3.7).

Si $n'(y')$, $n'$ y $n'_s$ son los números definidos para $f'$ de la misma manera que $n(y)$, $n$ y $n_s$ para $f$, se tiene $n'_s=n_s$ y $n'(y')=n(y)$ (I, 6.4.8), y por tanto es equivalente demostrar la afirmación (i) del enunciado para $f$ o para $f'$. Además, si $X$ es reducido y $y$ un punto normal de $Y$ (luego $A$ integralmente cerrado), entonces $A'$ es integralmente cerrado (18.8.12). Por otra parte, se deduce de la definición (18.8.7), de la observación (8.1.2, $a$)) y del teorema del doble límite inductivo que $A'$ es un límite inductivo de anillos $A_\alpha$ tales que los morfismos $\Spec(A_\alpha)\to\Spec(A)=Y$ son étales; por consiguiente, $X'$ es el límite proyectivo de los $X_\alpha=X\otimes_A A_\alpha$, que son étales sobre $X$, y por tanto reducidos puesto que $X$ lo es (17.5.7); pasando al límite,

\oldpage[IV-4]{165}

se deduce que $X'$ es reducido (8.7.1). Ahora bien, si $n(y)=n$, se tiene $n_s=n$ y, por consiguiente, los cuerpos residuales en los puntos de $X_\eta$ son necesariamente separables sobre $k(\eta)$; los cuerpos residuales en los puntos de $X'_{\eta'}$ son por tanto algebraicos y separables sobre $k(\eta')$ (4.6.1), es decir, $n'_s=n'=n$, y por lo anterior $n'(y')=n'$. El morfismo $f'$ tiene así las mismas propiedades que $f$, y si demostramos que $f'$ es étale, se deducirá que $f$ es étale por (17.7.1).

\emph{D) Fin de la demostración.} --- Supongamos por tanto que $A$ sea estrictamente local. Como $A$ es henseliano y el morfismo $f$ es finito, se deduce de (18.5.11) que si $x_j$ ($1\leq j\leq m$) son los puntos distintos de $f^{-1}(y)$, $X$ es la suma de los subpreesquemas abiertos $\Spec(\sh{O}_{X,x_j})=X_j$, que son por tanto las componentes conexas de $X$. Por hipótesis, cada uno de ellos domina $Y$, luego encuentra a $f^{-1}(\eta)$, y por consiguiente el número de puntos de $f^{-1}(\eta)$ es $\geq m$; \emph{a fortiori} se tiene $n_s\geq m$. Pero como $k(y)$ es separablemente cerrado y los $k(x_j)$ son algebraicos sobre $k(y)$, son necesariamente extensiones radiciales de $k(y)$, luego $m=n(y)$, lo que demuestra (i). Supongamos ahora verificadas las hipótesis de (ii); las relaciones $n(y)=n_s=n$ implican, en primer lugar, puesto que toda componente irreducible de $X$ domina $Y$, que cada $X_j$ es irreducible; además, como los $X_j$ son reducidos, son íntegros; finalmente, las relaciones $m=n_s=n$ implican que para el punto genérico $\xi_j$ de $X_j$, $k(\xi_j)$ es una extensión separable de $k(\eta)$ de grado 1, y por tanto isomorfa a $k(\eta)$; en otras palabras, la restricción $f|X_j:X_j\to Y$ de $f$ es un morfismo finito y birracional; además, $X_j$ es íntegro y, por hipótesis, $Y$ es normal; por tanto, (8.12.10.1) $f|X_j$ es un isomorfismo, lo que demuestra que $f$ es étale y termina la demostración de (18.10.16).

\begin{remark}[18.10.17]
\label{IV.18.10.17}
El método de «localización étale» utilizado en la demostración de (18.10.16) permite también mejorar los resultados de (15.5.1), eliminando la hipótesis noetheriana. Tenemos primero el resultado siguiente, que generaliza (15.5.2):
\end{remark}

\begin{env}[18.10.17.1]
\label{IV.18.10.17.1}
Sean $f:X\to Y$ un morfismo separado y cuasi-finito, $y$ un punto de $Y$ tal que $f$ sea universalmente abierto en los puntos de $f^{-1}(y)$. Entonces, para toda generalización $y'$ de $y$, se tiene $n(y)\leq n(y')$.
\end{env}

Se puede reemplazar $Y$ por el subpreesquema cerrado reducido $Z$ cuyo espacio subyacente es $\overline{\{y'\}}$ y $X$ por $f^{-1}(Z)$, y suponer así que $Y$ es íntegro con punto genérico $y'$. Usando (18.12.13) como en (18.10.16, B)), podemos suponer que $X$ es un abierto denso de un preesquema $X'$ y que $f$ es la restricción a $X$ de un morfismo finito $g:X'\to Y$. Mostremos que $g^{-1}(y')=f^{-1}(y')$. En efecto, sea $i:g^{-1}(y')\to X'$ la inyección canónica, que es un morfismo plano porque $y'$ es el punto genérico de $Y$ (I, 3.6.5); se puede escribir $f^{-1}(y')=i^{-1}(X)$. Por otra parte, la inyección canónica $j:X\to X'$ es un morfismo de tipo finito, ya que el compuesto $g\circ j=f$ es de tipo finito y $g$ es separado (1.5.4); por tanto, $X$ es un abierto retrocompacto de $X'$ y, en consecuencia, es proconstructible en $X'$ (1.9.5, (v)). Se deduce entonces de (2.3.10) que $i^{-1}(\overline{X})=\overline{i^{-1}(X)}$, es decir, $f^{-1}(y')$ es denso en $g^{-1}(y')$; pero como $g^{-1}(y')$ es discreto, necesariamente $f^{-1}(y')=g^{-1}(y')$.

Así, nos reducimos a demostrar que la suma de los números $[k(x):k(y)]_s$, donde $x$ recorre el conjunto de los puntos de $g^{-1}(y)$ en los que $g$ es universalmente abierto, es a lo sumo igual

\oldpage[IV-4]{166}

al número $n(y')$ (para el morfismo $g$). Usando después el hecho de que la propiedad de ser universalmente abierto en un punto se conserva por cambio de base, se muestra, como en (18.10.16, C)), que podemos reducirnos al caso $Y=\Spec(\sh{O}_{Y,y})$, con $\sh{O}_{Y,y}$ estrictamente local. Entonces, con las notaciones de (18.10.16, D)), $X$ es la suma de los subpreesquemas abiertos $X_j=\Spec(\sh{O}_{X,x_j})$, donde $x_j$ recorre $g^{-1}(y)$; la hipótesis de que $g$ sea abierto en uno de estos puntos $x_j$ implica que el subpreesquema $X_j$ correspondiente domina $Y$ (I, 10.3); se termina como en (18.10.16, D)).

\emph{Se deduce entonces de (18.10.17.1) que las conclusiones de (15.5.1) siguen siendo válidas cuando ya no se supone que $Y$ sea localmente noetheriano, sino solamente que $f$ sea separado, cuasi-finito y de presentación finita.}

En efecto, para demostrar la afirmación (i) de (15.5.1), se observa que, como $f$ es de presentación finita, el conjunto $E$ de los puntos $y'\in Y$ tales que $n(y')\geq n(y)$ es localmente constructible (9.7.9); para mostrar que $y$ es interior a $E$, basta, en virtud de (I, 10.1), ver que toda generalización $y'$ de $y$ pertenece a $E$, lo que no es otra cosa que (18.10.17.1).

Para demostrar la afirmación (ii) de (15.5.1), como $f$ es de presentación finita, se puede usar (8.10.5, (xii)) aplicado según el método de (8.1.2, a)), y por tanto suponer ya que $Y=\Spec(\sh{O}_{Y,y})$. Aplicando después (2.7.1, (vii)), se ve que mediante el cambio de base fielmente plano $Y'\to Y$, donde $Y'=\Spec({}^{hs}\!\sh{O}_{Y,y})$, podemos suponer que el anillo $\sh{O}_{Y,y}$ es estrictamente local. Aplicando (18.5.11, c)), vemos entonces que si $x_j$ ($1\leq j\leq n$) son los puntos de $f^{-1}(y)$, $X$ es la suma de $n$ subpreesquemas abiertos $X_j=\Spec(\sh{O}_{X,x_j})$ que son finitos sobre $Y$, y de un preesquema abierto $X''$. Si demostramos que $X''=\varnothing$, habremos demostrado que $f$ es finito y, por tanto, propio. Ahora bien, como los cuerpos $k(x_j)$ son extensiones algebraicas de $k(y)$, son radiciales, y por consiguiente $n(y)=n$; como $f$ es abierto en cada uno de los puntos $x_j$, $X_j$ domina $Y$ (I, 10.3), y la restricción de $f$ a $X_j$, al ser un morfismo finito, es sobreyectiva (II, 6.1.10). La hipótesis de que $z\mapsto n(z)$ es constante en $Y$ implica por tanto que $X''\cap f^{-1}(y')=\varnothing$ para todo $y'\in Y$, es decir, $X''=\varnothing$.

Finalmente, la demostración de (iii) reproduce sin cambios la dada en (15.5.1).

\begin{env}[18.10.17.2]
\label{IV.18.10.17.2}
Mediante métodos análogos que utilizan la henselización estricta, en la mayoría de los resultados de las \textsection\textsection~14 y 15 se pueden eliminar las hipótesis noetherianas.
\end{env}

\begin{lemma}[18.10.18]
\label{IV.18.10.18}
Sean $X$, $Y$ dos preesquemas, $f:X\to Y$ un morfismo birracional\footnote{Entendemos por esto la noción definida en (6.15.4), pero sin suponer que $X$ e $Y$ sean reducidos.}, $x$ un punto de $X$. Para que $f$ sea un isomorfismo local en el punto $x$, es necesario y suficiente que $f$ sea étale en el punto $x$. Para que $f$ sea una inmersión abierta, es necesario y suficiente que $f$ sea étale y separado.
\end{lemma}

Las condiciones enunciadas son trivialmente necesarias; todo se reduce a ver que son suficientes. Para la primera afirmación, la cuestión es local en $X$ e $Y$, de modo que podemos suponer que $f$ es étale y que $X$ e $Y$ son afines, con $f$ separado, y nos reducimos a demostrar la segunda afirmación. En virtud de (17.9.1), hay que demostrar que $f$ es radicial. Sea pues $y\in Y$ y demostremos que $f^{-1}(y)$ es radicial sobre $k(y)$. El cambio de base $\Spec(\sh{O}_{Y,y})\to Y$ no cambia el hecho de que $f$ sea étale, separado y birracional, por lo que podemos limitarnos al caso $Y=\Spec(A)$, con $A=\sh{O}_{Y,y}$.

\oldpage[IV-4]{167}

Sea entonces $A'={}^{hs}\!A$, que es un anillo local y un $A$-módulo fielmente plano (18.8.8), siendo local el homomorfismo $A\to A'$. Pongamos $Y'=\Spec(A')$, $X'=X\times_Y Y'$, $f'=f_{(Y')}:X'\to Y'$, de modo que $f'$ es étale y separado; además, como el morfismo $Y'\to Y$ es plano, el razonamiento de (6.15.4.1) muestra que $f'$ es también birracional. Si $y'$ es el punto cerrado de $Y'$, bastará, en virtud de (2.6.1, (v)), demostrar que $f'^{-1}(y')$ es radicial sobre $k(y')$; puesto que $f'$ es étale, bastará demostrar que $f'^{-1}(y')$ contiene a lo sumo un punto (17.6.1, c')).

Supongamos, pues, que $Y=\Spec(A)$, que $A$ es estrictamente local, y que $f$ es separado, étale y birracional, y demostremos que $f^{-1}(y)$ no puede contener más de un punto. En efecto, si en $f^{-1}(y)$ existieran dos puntos distintos $x_1$, $x_2$, existirían dos $Y$-secciones $u_1$, $u_2$ de $f$ tales que $u_1(y)=x_1$ y $u_2(y)=x_2$ (18.8.1); pero, como $f$ es étale y separado, $u_1$ y $u_2$ serían isomorfismos de $Y$ sobre dos componentes conexas (abiertas) de $X$ (17.9.4), y habría por tanto al menos dos puntos máximos de $X$ sobre todo punto máximo de $Y$, lo que contradice la hipótesis de que $f$ es birracional.

\begin{proposition}[18.10.19]
\label{IV.18.10.19}
Sea $Y$ un preesquema reducido cuyo conjunto de componentes irreducibles es localmente finito, $f:X\to Y$ un morfismo separado, localmente de tipo finito y formalmente no ramificado, $g$ una $Y$-sección racional de $f$ (es decir, una $Y$-aplicación racional de $Y$ en $X$), $U$ el dominio de definición de $g$ (I, 7.2.1), y sea $Z$ la clausura de $g(U)$ en $X$. Entonces, para todo punto $y\in Y-U$ tal que $Y$ sea geométricamente unirrama en el punto $y$, se tiene $Z_y=Z\cap f^{-1}(y)=\varnothing$. En particular, si $Y$ es geométricamente unirrama, $g(U)$ es una parte a la vez abierta y cerrada de $X$. Si, para toda parte cerrada irreducible $T$ de $X$, $f(T)$ es cerrada en $Y$, $g$ está definida en todo punto geométricamente unirrama de $Y$.
\end{proposition}

Como $Y$ es reducido y $f$ separado, se deduce de (I, 7.2.2) que existe una $U$-sección $u$ de $f^{-1}(U)$ perteneciente a la clase $g$; además, por (17.4.1.2), puesto que $f$ es formalmente no ramificado, $u$ es un isomorfismo de $U$ sobre el preesquema inducido en el abierto $u(U)$ de $X$. Si seguimos denotando por $Z$ el subpreesquema reducido de $X$ cuyo espacio subyacente es $Z$, $Z|_U=u(U)$, al ser reducido, también está inducido por $Z$; por tanto, la restricción $f'=f|_Z$ de $f$ es un morfismo birracional de $Z$ en $Y$, y además formalmente no ramificado, pues así lo es $f$ (17.1.3). Como la cuestión es local en $Y$, y $Y$ es reducido y geométricamente unirrama (luego íntegro) en el punto $y$, podemos limitarnos (reemplazando $Y$ por un entorno abierto de $y$) al caso en que $Y$ sea íntegro; entonces $u(U)$ es irreducible, y también $Z$, por lo que $Z$ es asimismo íntegro. Por (I, 8.2.7), como $f'$ es dominante, el homomorfismo $\sh{O}_{Y,f(x)}\to\sh{O}_{Z,x}$ es inyectivo para todo $x\in Z_y$, y se concluye de (18.10.2, (iii)) que $f'$ es un morfismo étale; al ser separado y birracional, es un isomorfismo local en el punto $x$ en virtud de (18.10.18); pero esto implicaría que la $U$-sección $u$ se puede prolongar a un abierto estrictamente mayor que $U$, en contradicción con la definición de este último. Por tanto, necesariamente $Z_y=\varnothing$, lo que establece la primera afirmación; la segunda es una consecuencia evidente. Además, bajo las hipótesis de la última afirmación del enunciado, como $Z$ es cerrado e irreducible, $f(Z)$ es cerrado en $Y$ y, por tanto, igual a $Y$, es decir, $Z_y\ne\varnothing$ para todo $y\in Y$, lo que termina la demostración.

\begin{remark}[18.10.20]
\label{IV.18.10.20}
Sea $Y$ un preesquema localmente noetheriano. Se dice que un morfismo $f:X\to Y$ es esencialmente propio si es localmente de tipo finito y si, para todo $Y$-esquema de la forma $Y'=\Spec(A)$, donde $A$ es un anillo de valoración discreta, la aplicación canónica (II, 7.3.2.2) relativa al $Y'$-preesquema $X\times_Y Y'$ es biyectiva.

\oldpage[IV-4]{168}

El razonamiento de (II, 7.3.8) prueba que esto equivale a decir que $f$ (supuesto localmente de tipo finito) es separado y que, para todo cambio de base $Y''\to Y$, donde $Y''$ es localmente noetheriano, la imagen por $f_{(Y'')}:X_{(Y'')}\to Y''$ de toda parte cerrada irreducible de $X''$ es cerrada. Decir que $f$ es propio (para $Y$ localmente noetheriano) significa por tanto que $f$ es esencialmente propio y de tipo finito (II, 7.3.8); pero se encuentran ejemplos importantes de morfismos esencialmente propios que no son de tipo finito (por ejemplo, ciertos «preesquemas de Picard» o ciertos «preesquemas de Néron-Severi» (cap. VI)). Se deduce evidentemente de (18.10.19) que si $Y$ es localmente noetheriano, reducido y geométricamente unirrama, y si el morfismo $f:X\to Y$ es esencialmente propio y no ramificado, entonces toda $Y$-sección racional de $f$ está definida en todas partes.
\end{remark}

\begin{proposition}[18.10.21]
\label{IV.18.10.21}
Sea $f:X\to Y$ un morfismo localmente de presentación finita, $x$ un punto de $X$, $y=f(x)$, y supongamos que el anillo $\sh{O}_{Y,y}$ sea íntegro y geométricamente unirrama. Sea $r$ el número mínimo de generadores del $\sh{O}_{X,x}$-módulo $(\Omega^1_{X/Y})_x$ (igual al rango del espacio vectorial sobre $k(x)$ $(\Omega^1_{X/Y})_x\otimes_{\sh{O}_{X,x}}k(x))$. Entonces, para que $f$ sea liso en el punto $x$, es necesario y suficiente que, si $y'$ es el punto genérico de la única componente irreducible de $Y$ que contiene a $y$, exista una generalización $x'$ de $x$ tal que $f(x')=y'$ y $\dim_{x'}f^{-1}(y')\geq r$.
\end{proposition}

Si $f$ es liso en el punto $x$, también lo es en un entorno abierto $U$ de $x$, y por tanto en toda generalización $x'$ de $x$, y $\dim_{x'}(f^{-1}(f(x')))=r$ en cada uno de esos puntos, en virtud de (17.10.2). Como además $f$ es plano en $U$, existe una generalización de $x$ situada sobre toda generalización de $y$ (2.3.4), lo que demuestra que la condición es necesaria. Para demostrar que es suficiente, observemos primero que por (17.5.1) basta demostrar que el morfismo obtenido de $f$ mediante el cambio de base $\Spec(\sh{O}_{Y,y})\to Y$ es liso en el punto $x$; por tanto, podemos suponer que $Y=\Spec(\sh{O}_{Y,y})$.

La cuestión es local en $X$, por lo que (por (16.4.22) y ($0_{\mathrm I}$, 5.2.2)) podemos suponer que $X$ es afín y que existen $r$ secciones $t_i$ ($1\leq i\leq r$) de $\sh{O}_X$ sobre $X$ tales que las secciones $d_{X/Y}t_i$ generan $\Omega^1_{X/Y}$. Sea
\[
  g:X\to Z=Y[T_1,\ldots,T_r]=\mathbf V^r_Y
\]
el $Y$-morfismo correspondiente al homomorfismo $\sh{O}_Y^r\to f_*(\sh{O}_X)$ de $\sh{O}_Y$-módulos definido por las $t_i$ (consideradas como secciones de $f_*(\sh{O}_X)$ sobre $Y$) (II, 1.2.7). Se tiene la sucesión exacta (16.4.19)
\[
  g^*(\Omega^1_{Z/Y})\longrightarrow\Omega^1_{X/Y}\longrightarrow\Omega^1_{X/Z}\longrightarrow0.
\]
Como los $dT_i$ ($1\leq i\leq r$) generan $\Omega^1_{Z/Y}$, se deduce de la definición de $g$ que el homomorfismo $g^*(\Omega^1_{Z/Y})\to\Omega^1_{X/Y}$ es sobreyectivo, luego $\Omega^1_{X/Z}=0$ y, por tanto, $g$ es no ramificado (17.4.1). Vamos a demostrar que $g$ es étale en el punto $x$, y teniendo en cuenta (17.3.8), esto probará que $f$ es liso en el punto $x$. Como se supone que $\sh{O}_{Y,y}$ es íntegro y geométricamente

\oldpage[IV-4]{169}

unirrama, lo mismo ocurre con $\sh{O}_{Z,z}$, donde $z=g(x)$, pues el morfismo estructural $Z\to Y$ es liso (17.3.8). En virtud de (18.10.1), basta demostrar que el homomorfismo $\sh{O}_{Z,z}\to\sh{O}_{X,x}$ es inyectivo; equivale a ver que el morfismo correspondiente $\Spec(\sh{O}_{X,x})\to\Spec(\sh{O}_{Z,z})$ es dominante, puesto que $\sh{O}_{Z,z}$ es íntegro (I, 1.2.7).

Sea $y'$ el punto genérico de $Y=\Spec(\sh{O}_{Y,y})$ y sea $z'$ el punto genérico de $Z=\Spec(\sh{O}_{Y,y}[T_1,\ldots,T_r])$, que es íntegro; si $h:Z\to Y$ es el morfismo estructural, se tiene $h(z')=y'$. Basta demostrar que existe en $f^{-1}(y')$ una generalización $x'$ de $x$ tal que la imagen de $x'$ por el morfismo $g_{y'}:f^{-1}(y')\to h^{-1}(y')$ sea igual a $z'$. Además, $z'$ es también el punto genérico de $h^{-1}(y')=\Spec(k(y')[T_1,\ldots,T_r])$; como el morfismo $g_{y'}$ es no ramificado, luego cuasi-finito (17.4.2), lo mismo vale para su restricción a toda componente irreducible de $f^{-1}(y')$; como $f^{-1}(y')$ y $h^{-1}(y')$ son localmente noetherianos y $\dim(h^{-1}(y'))=r$, se deduce de (5.4.1) que la restricción de $g_{y'}$ a toda componente irreducible de dimensión $\geq r$ de $f^{-1}(y')$ es un morfismo dominante. Ahora bien, por hipótesis existe una componente de ese tipo que contiene una generalización $x'$ de $x$; \emph{a fortiori}, su punto genérico es una generalización de $x$, de donde se sigue la conclusión.

\subsection{Aplicación a las álgebras locales noetherianas completas sobre un cuerpo.}
\label{subsection:IV.18.11}

El lema siguiente generaliza (0, 21.9.1) y (0, 21.9.2):

\begin{lemma}[18.11.1]
\label{IV.18.11.1}
Sean $k$ un cuerpo; $A$ un anillo local noetheriano completo que es una $k$-álgebra, y $K$ el cuerpo residual de $A$.
\begin{enumerate}
  \item[(i)] Si el espacio vectorial sobre $K$ $\Omega^1_{K/k}$ tiene rango finito, el $A$-módulo $\hat{\Omega}^1_{A/k}$ es de tipo finito.
  \item[(ii)] Si además $A$ es una $k$-álgebra formalmente lisa (para la topología adica), $\hat{\Omega}^1_{A/k}$ es un $A$-módulo libre de rango igual a
\[
    \dim(A)+\operatorname{rg}_K(\Omega^1_{K/k})-\operatorname{rg}_K(\Upsilon_{K/k}).
\]
\end{enumerate}
Además, para todo subcuerpo $k_0$ de $k$ tal que $\Omega^1_{k/k_0}$ sea un espacio vectorial sobre $k$ de rango finito, $\hat{\Omega}^1_{A/k_0}$ es un $A$-módulo libre de rango igual a
\[
  \dim(A)+\operatorname{rg}_K(\Omega^1_{K/k_0})-\operatorname{rg}_K(\Upsilon_{K/k_0})+\operatorname{rg}_k(\Omega^1_{k/k_0}).
\]
\end{lemma}

La afirmación (i) se incluye solo como recuerdo, pues fue demostrada en (0, 20.7.15). Para probar (ii), observemos que esta afirmación ya quedó establecida de hecho por la demostración de (0, 21.9.2), cuyo enunciado es el único que utiliza la finitud del grado de trascendencia de $K$ sobre $k$ (por medio de la igualdad de Cartier (0, 21.7.1)).

\begin{lemma}[18.11.2]
\label{IV.18.11.2}
Sean $k$ un cuerpo de característica $p>0$, $C$ una $k$-álgebra local noetheriana completa, cuyo cuerpo residual es una extensión finita de $k$, y que es íntegra; sean $n$ su dimensión y $L$ su cuerpo de fracciones. Entonces $\Omega^1_{L/k}$ y $\Upsilon_{L/k}$ son espacios vectoriales sobre $L$ de rango finito, y se tiene
\begin{equation}
\label{IV.18.11.2.1}
\tag{18.11.2.1}
  \operatorname{rg}_L(\Omega^1_{L/k})-\operatorname{rg}_L(\Upsilon_{L/k})\geq n.
\end{equation}

\oldpage[IV-4]{170}

Supongamos además que $[k:k^p]<+\infty$. Entonces se tiene la igualdad
\begin{equation}
\label{IV.18.11.2.2}
\tag{18.11.2.2}
  \operatorname{rg}_L(\Omega^1_{L/k})-\operatorname{rg}_L(\Upsilon_{L/k})=n.
\end{equation}
Además, $\Omega^1_{C/k}$ es entonces isomorfo a $\hat{\Omega}^1_{C/k}$, y por tanto, por el lema (18.11.1), es un $C$-módulo de tipo finito.
\end{lemma}

En efecto, se sabe (0, 19.8.9) que existe una sub-$k$-álgebra $C_0$ de $C$ isomorfa a un álgebra de series formales $k[[T_1,\ldots,T_n]]$ y tal que $C$ sea una $C_0$-álgebra finita. Por consiguiente, $L$ es una extensión finita del cuerpo de fracciones $L_0=k((T_1,\ldots,T_n))$ de $C_0$. Ahora bien, se tiene la sucesión exacta de espacios vectoriales sobre $L$ obtenida aplicando (0, 20.6.17.1) al subcuerpo primo de $k$ y a los tres cuerpos $k\subset L_0\subset L$ (y teniendo en cuenta (0, 20.6.21, (i)))
\[
  0\to\Upsilon_{L_0/k}\otimes_{L_0}L\to\Upsilon_{L/k}\to\Upsilon_{L/L_0}
  \to\Omega^1_{L_0/k}\otimes_{L_0}L\to\Omega^1_{L/k}\to\Omega^1_{L/L_0}\to0.
\]
Como $L$ es una extensión finita de $L_0$, $\Omega^1_{L/L_0}$ y $\Upsilon_{L/L_0}$ son espacios vectoriales sobre $L$ de rango finito y del mismo rango, en virtud de la igualdad de Cartier (0, 21.7.1). Como $L_0$ es separable sobre $k$ (0, 21.9.6.4), se tiene $\Upsilon_{L_0/k}=0$ (0, 20.6.3); de la sucesión exacta anterior se deduce ya que $\Upsilon_{L/k}$ tiene rango finito y que, en todos los casos, ($0_{\mathrm{III}}$, 11.10.2)
\[
  \operatorname{rg}_L(\Omega^1_{L/k})-\operatorname{rg}_L(\Upsilon_{L/k})
  =\operatorname{rg}_{L_0}(\Omega^1_{L_0/k})-\operatorname{rg}_{L_0}(\Upsilon_{L_0/k}).
\]

Para demostrar (18.11.2.1) o (18.11.2.2), podemos por tanto limitarnos a demostrar estas relaciones reemplazando $C$ y $L$ por $C_0$ y $L_0$. Como $L_0$ es separable sobre $k$ (0, 21.9.6.4), se tiene $\Upsilon_{L_0/k}=0$ (0, 20.6.3); por otra parte, $\Omega^1_{L_0/k}=\Omega^1_{C_0/k}\otimes_{C_0}L_0$ (0, 20.5.9). Se sabe (0, 21.9.4) que si $C_1=k[[T_1^p,\ldots,T_n^p]]$, $\hat{\Omega}^1_{C_0/k}$ se identifica con $\Omega^1_{C_0/C_1}$; por otra parte, se tiene $\Omega^1_{C_0/k}=\Omega^1_{C_0/k[C_0^p]}$ (0, 21.1.5), y como $C_0^p=k^p[[T_1^p,\ldots,T_n^p]]$, se tiene $k[C_0^p]\subset C_1$, con igualdad cuando $[k:k^p]<+\infty$. Ahora bien, existe una sucesión exacta $\Omega^1_{C_0/k[C_0^p]}\to\Omega^1_{C_0/C_1}\to0$ (0, 20.5.7), de donde se obtiene una sucesión exacta $\Omega^1_{L_0/k}\to\hat{\Omega}^1_{C_0/k}\otimes_{C_0}L_0\to0$. Como $\hat{\Omega}^1_{C_0/k}$ es un $C_0$-módulo libre de rango $n$ (0, 21.9.3), se ve que en todos los casos $\operatorname{rg}_{L_0}(\Omega^1_{L_0/k})\geq n$, con igualdad cuando $[k:k^p]<+\infty$; esto demuestra ya (18.11.2.1) y (18.11.2.2).

Finalmente, para ver que $\Omega^1_{C/k}$ es isomorfo a $\hat{\Omega}^1_{C/k}$ cuando $[k:k^p]<+\infty$, basta demostrar que $\Omega^1_{C/k}$ es un $C$-módulo de tipo finito, puesto que $C$ es un anillo local noetheriano completo ($0_{\mathrm I}$, 7.3.3). Como $\Omega^1_{C/k}$ es isomorfo a $\Omega^1_{C/k[C^p]}$ y $k[C_0^p]\subset k[C^p]$, todo se reduce a demostrar que $C$ es un $k[C_0^p]$-módulo de tipo finito (0, 20.4.7); pero esto se deduce de que $C$ es un $C_0$-módulo de tipo finito y $C_0$ un $k[C_0^p]$-módulo de tipo finito por la hipótesis $[k:k^p]<+\infty$.

\begin{lemma}[18.11.3]
\label{IV.18.11.3}
Sean $k$ un cuerpo, $p$ su exponente característico, y supongamos que $[k:k^p]<+\infty$. Sea $A$ una $k$-álgebra local noetheriana completa, cuyo cuerpo residual es una extensión finita de $k$. Sea $\mathfrak p$ un ideal primo de $A$ tal que $A_{\mathfrak p}$ sea geométricamente regular sobre $k$ (6.7.6), de modo que existe un único ideal primo mínimo $\mathfrak q$ de $A$ contenido en $\mathfrak p$. Entonces $(\hat{\Omega}^1_{A/k})_{\mathfrak p}$ es un $A_{\mathfrak p}$-módulo libre de rango igual a $\dim(A/\mathfrak q)$.
\end{lemma}

Como $A$ es noetheriano y $\Spec(A)$ es regular (y \,\emph{a fortiori} íntegro) en el punto $\mathfrak p$, $\mathfrak p$ pertenece a una sola componente irreducible de $\Spec(A)$, y por tanto no contiene

\oldpage[IV-4]{171}

un solo ideal mínimo $\mathfrak q$ de $A$, y además se tiene $\mathfrak q_{\mathfrak p}=0$. Si ponemos $B=A/\mathfrak q$, la sucesión (0, 20.7.20)
\[
  \mathfrak q/\mathfrak q^2\longrightarrow
  \hat{\Omega}^1_{A/k}\widehat{\otimes}_A B
  \longrightarrow\hat{\Omega}^1_{B/k}\longrightarrow0
\]
es exacta; en efecto, $\hat{\Omega}^1_{A/k}$ es un $A$-módulo de tipo finito (18.11.1), por lo que $\hat{\Omega}^1_{A/k}\widehat{\otimes}_A B=\hat{\Omega}^1_{A/k}\otimes_A B=\hat{\Omega}^1_{A/k}/\mathfrak q\hat{\Omega}^1_{A/k}$, y como este $B$-módulo es de tipo finito, todo sub-$B$-módulo suyo es cerrado ($0_{\mathrm I}$, 7.3.5), teniendo en cuenta (0, 20.4.5). Como la imagen de $\mathfrak q/\mathfrak q^2$ es densa en el núcleo del homomorfismo $\hat{\Omega}^1_{A/k}\widehat{\otimes}_A B\to\hat{\Omega}^1_{B/k}$ (0, 20.7.20), necesariamente es igual a dicho núcleo, de donde se obtiene nuestra afirmación. Al localizar en $\mathfrak p$ la sucesión exacta anterior, se ve además que el homomorfismo canónico $(\hat{\Omega}^1_{A/k})_{\mathfrak p}\to(\hat{\Omega}^1_{B/k})_{\mathfrak p/\mathfrak q}$ es biyectivo. Así, podemos limitarnos al caso $\mathfrak q=0$, es decir, al caso en que $A$ es íntegro. Distingamos entonces dos casos:

\emph{I) $p>1$.} Entonces podemos aplicar (18.11.2) al álgebra cociente $C=A/\mathfrak p$, cuyo cuerpo de fracciones $K$ no es otro que el cuerpo residual de $A_{\mathfrak p}$; se obtiene por tanto
\begin{equation}
\label{IV.18.11.3.1}
\tag{18.11.3.1}
  \operatorname{rg}_K(\Omega^1_{K/k})-\operatorname{rg}_K(\Upsilon_{K/k})=\dim(A/\mathfrak p).
\end{equation}

Observemos ahora (0, 19.6.6) que $A_{\mathfrak p}$ es una $k$-álgebra formalmente lisa para su topología $\mathfrak pA_{\mathfrak p}$-preádica; por tanto, $\Omega^1_{A_{\mathfrak p}/k}$ es formalmente proyectivo (0, 20.4.9) para la topología $\mathfrak p$-preádica (0, 20.4.5); por otra parte, $\Omega^1_{A_{\mathfrak p}/k}=(\Omega^1_{A/k})_{\mathfrak p}$ (0, 20.5.9) es un $A_{\mathfrak p}$-módulo de tipo finito, en virtud de (18.11.2) aplicado a $A$. Para todo entero $j$, $\Omega^1_{A_{\mathfrak p}/k}/\mathfrak p^{j+1}\Omega^1_{A_{\mathfrak p}/k}$ es por tanto un $(A_{\mathfrak p}/\mathfrak p^{j+1}A_{\mathfrak p})$-módulo proyectivo de rango $m=\operatorname{rg}_K(\Omega^1_{A_{\mathfrak p}/k}\otimes_{A_{\mathfrak p}}K)$ (0, 19.2.4); se concluye entonces de ($0_{\mathrm{III}}$, 10.2.1 y 10.1.3) que $\hat{\Omega}^1_{A_{\mathfrak p}/k}$ es un $A_{\mathfrak p}$-módulo libre de rango $m$. Sea $A'=(A_{\mathfrak p})^\wedge$ el álgebra completada de $A_{\mathfrak p}$ para su topología $\mathfrak pA_{\mathfrak p}$-preádica; $A'$ es también una $k$-álgebra formalmente lisa (0, 19.3.6) para su topología adica, y se deduce de (0, 20.7.14) y (0, 20.4.5) que $\hat{\Omega}^1_{A'/k}=\hat{\Omega}^1_{A_{\mathfrak p}/k}$; concluimos que $\hat{\Omega}^1_{A'/k}$ es un $A'$-módulo libre de rango $m$. Pero de (18.11.1) y del hecho de que $\dim(A')=\dim(A_{\mathfrak p})$ (0, 16.2.4) se deduce que
\begin{equation}
\label{IV.18.11.3.2}
\tag{18.11.3.2}
  m=\dim(A_{\mathfrak p})+
  \bigl(\operatorname{rg}_K(\Omega^1_{K/k})-\operatorname{rg}_K(\Upsilon_{K/k})\bigr),
\end{equation}
de donde, en virtud de (18.11.3.1),
\[
  m=\dim(A/\mathfrak p)+\dim(A_{\mathfrak p}).
\]
Pero como $A$ es un anillo local noetheriano completo, es un cociente de un anillo regular (0, 19.8.8), luego por (0, 16.5.12) se tiene $\dim(A)=\dim(A/\mathfrak p)+\dim(A_{\mathfrak p})$, lo que termina la demostración en este caso.

\emph{II) $p=1$.} Como arriba, se tiene $\dim(A)=\dim(A/\mathfrak p)+\dim(A_{\mathfrak p})$; pongamos $n=\dim(A)$, $r=\dim(A/\mathfrak p)$, $s=\dim(A_{\mathfrak p})$; veremos que $(\hat{\Omega}^1_{A/k})_{\mathfrak p}$ es un $A_{\mathfrak p}$-módulo libre de rango $n$. Primero demostraremos el lema siguiente:

\begin{lemma}[18.11.3.3]
\label{IV.18.11.3.3}
Sean $A$ un anillo local noetheriano, $\mathfrak p$ un ideal primo de $A$ tal que $A_{\mathfrak p}$ sea un anillo de Cohen-Macaulay. Para todo sistema de parámetros $(t_i)_{1\leq i\leq s}$ de $A_{\mathfrak p}$, existe una sucesión $(x_i)_{1\leq i\leq s}$ de elementos de $\mathfrak p$ tal que los $x_i$ formen parte de un sistema de parámetros de $A$

\oldpage[IV-4]{172}

y que, para todo $i$, $x_i/1$ sea congruente con $t_i$ módulo el ideal $t_1A_{\mathfrak p}+\cdots+t_{i-1}A_{\mathfrak p}$. En particular, si $A_{\mathfrak p}$ es regular y los $t_i$ forman un sistema regular de parámetros de $A_{\mathfrak p}$, lo mismo ocurre con los $x_i/1$.
\end{lemma}

Este lema será a su vez consecuencia del siguiente:

\begin{lemma}[18.11.3.4]
\label{IV.18.11.3.4}
Sean $A$ un anillo local noetheriano, $\mathfrak p$ un ideal primo de $A$ y $x$ un elemento de $A$; entonces existe un elemento $x'\in A$ tal que $x'/1=x/1$ en $A_{\mathfrak p}$ y $x'$ no pertenezca a ninguno de los ideales primos mínimos de $A$ no contenidos en $\mathfrak p$.
\end{lemma}

Mostremos primero cómo (18.11.3.4) implica (18.11.3.3). Multiplicando, si es necesario, los $t_i$ por elementos invertibles de $A_{\mathfrak p}$, podemos suponer que son de la forma $x_i/1$, con $x_i\in\mathfrak p$ para $1\leq i\leq s$. Razonemos por inducción sobre $s$; como $t_1$ forma parte de un sistema de parámetros de $A_{\mathfrak p}$, no pertenece a ninguno de los ideales primos mínimos de $A_{\mathfrak p}$ (0, 16.5.5), de modo que ningún $x'_1\in A$ tal que $x'_1/1=t_1$ puede pertenecer a un ideal primo mínimo de $A$ contenido en $\mathfrak p$. En virtud de (18.11.3.4), podemos además escoger $x'_1$ tal que $x'_1/1=t_1$ (por tanto $x'_1\in\mathfrak p$) y $x'_1$ no pertenezca a ningún ideal primo mínimo de $A$, y por consiguiente forme parte de un sistema de parámetros de $A$ (0, 16.3.4 y 16.3.7). Se procede después por inducción, considerando el anillo $A'=A/x'_1A$ y el ideal primo $\mathfrak p'=\mathfrak p/x'_1A$ de este anillo; como $A'_{\mathfrak p'}=A_{\mathfrak p}/t_1A_{\mathfrak p}$, $A'_{\mathfrak p'}$ es también un anillo de Cohen-Macaulay (0, 16.5.5), y las imágenes $t'_i$ ($2\leq i\leq s$) de los $t_i$ en $A'_{\mathfrak p'}$ forman un sistema de parámetros de este anillo; basta aplicar a $A'$ y a los $t'_i$ ($2\leq i\leq s$) la hipótesis de inducción. La última afirmación de (18.11.3.3) resulta de que, en $A_{\mathfrak p}$, un sistema de parámetros que genera el ideal maximal es un sistema regular de parámetros (0, 17.1.1).

Probemos, pues, (18.11.3.4). Sea $(\mathfrak p'_k)_{1\leq k\leq r}$ la sucesión de los ideales primos mínimos de $A$ no contenidos en $\mathfrak p$ y tales que $x\in\mathfrak p'_k$, y sea $(\mathfrak p_j)_{1\leq j\leq n}$ la sucesión de los ideales primos mínimos de $A$ distintos de los $\mathfrak p'_k$; podemos suponer que $r\geq1$. Como $\mathfrak p'_k$ no contiene $\bigcap_{1\leq j\leq n}\mathfrak p_j$, existe un elemento $y\in\bigcap_{1\leq j\leq n}\mathfrak p_j$ que no está contenido en ninguno de los $\mathfrak p'_k$ (Bourbaki, \emph{Alg. comm.}, cap. II, \textsection~1, n\textsuperscript{o}~1, prop. 2); además $y/1$ pertenece a todos los ideales primos mínimos de $A_{\mathfrak p}$, por lo que es nilpotente; si $(y/1)^h=0$, el elemento $x'=x+y^h$ satisface las condiciones del enunciado, pues, por una parte, $y^h\notin\mathfrak p'_k$ para $1\leq k\leq r$ y, por definición de los $\mathfrak p'_k$, se tiene $x'\notin\mathfrak p'_k$ para $1\leq k\leq r$; por otra parte, si $\mathfrak p_j$ es un ideal primo mínimo de $A$ no contenido en $\mathfrak p$ pero tal que $x\notin\mathfrak p_j$, también se tiene $x'\notin\mathfrak p_j$ puesto que $y\in\mathfrak p_j$.

Volvamos a la demostración de (18.11.3) cuando $p=1$. En virtud de (18.11.3.3), existe un sistema regular de parámetros $(t_i)_{1\leq i\leq s}$ de $A_{\mathfrak p}$ tal que $t_i=x_i/1$, donde los $x_i$ ($1\leq i\leq s$) pertenecen a $\mathfrak p$ y forman parte de un sistema de parámetros $(x_j)_{1\leq j\leq n}$ de $A$. Pongamos $A_0=k[[T_1,\ldots,T_n]]$; como $A$ es una $k$-álgebra completa y los $x_j$ pertenecen al ideal maximal $\mathfrak m$ de $A$, existe un $k$-homomorfismo local $u:A_0\to A$ tal que $u(T_j)=x_j$ para $1\leq j\leq n$ (Bourbaki, \emph{Alg. comm.}, cap. III, \textsection~4, n\textsuperscript{o}~5, prop. 6); si $\mathfrak n$ es el ideal de $A$ generado por los $x_j$ ($1\leq j\leq n$), por hipótesis es un ideal de definición de $A$; se deduce entonces de ($0_{\mathrm I}$, 7.4.4 y 7.4.3) que $u$ hace de $A$ una $A_0$-álgebra finita.

Pongamos $\mathfrak p_0=\sum_{j=1}^s A_0T_j$, $B_0=(A_0)_{\mathfrak p_0}$ y $B=A\otimes_{A_0}B_0$, de modo que $u_{\mathfrak p_0}:B_0\to B$ hace

\oldpage[IV-4]{173}

de $B$ una $B_0$-álgebra finita; además, si $\mathfrak p'$ es el ideal de $B$ generado por $\mathfrak p$, se tiene $B_{\mathfrak p'}=A_{\mathfrak p}$; como $\mathfrak p'$ contiene $\mathfrak p_0B$ por construcción, está sobre el ideal maximal $\mathfrak p_0B_0$ de $B_0$. Mostremos que el morfismo $\Spec(B)\to\Spec(B_0)$ no es ramificado en el punto $\mathfrak p'$: esto resulta, en efecto, de que $k(\mathfrak p')$ es una extensión finita del cuerpo $k(\mathfrak p_0)$ de característica 0, luego es necesariamente separable, y de que $B_{\mathfrak p'}/\mathfrak p_0B_{\mathfrak p'}=k(\mathfrak p')$ en virtud de la elección de los $x_j$ para $1\leq j\leq s$ (17.4.1).

Anotemos ahora el lema siguiente:

\begin{lemma}[18.11.3.5]
\label{IV.18.11.3.5}
Sean $k$ un cuerpo, $R$, $S$ dos $k$-álgebras locales noetherianas completas, tales que, si $K$ es el cuerpo residual de $S$, $\Omega^1_{K/k}$ tenga rango finito sobre $K$; sea $u:R\to S$ un $k$-homomorfismo que hace de $S$ una $R$-álgebra finita; entonces se tiene $\hat{\Omega}^1_{S/R}=\Omega^1_{S/R}$, y la sucesión
\begin{equation}
\label{IV.18.11.3.6}
\tag{18.11.3.6}
  \hat{\Omega}^1_{R/k}\otimes_R S\xrightarrow{v}\hat{\Omega}^1_{S/k}\xrightarrow{w}\Omega^1_{S/R}\to0
\end{equation}
(cf. (0, 20.7.17.3)) es exacta.
\end{lemma}

La primera afirmación resulta de que $\Omega^1_{S/R}$ es un $S$-módulo de tipo finito (0, 20.4.7) y de ($0_{\mathrm I}$, 7.3.6). Por otra parte, se sabe (0, 20.7.17.3) que la imagen de $v$ es densa en $\operatorname{Ker}(w)$ y que $w$ es sobreyectivo; pero de (18.11.1) resulta que $\hat{\Omega}^1_{S/k}$ es un $S$-módulo de tipo finito; concluimos que todo sub-$S$-módulo de $\hat{\Omega}^1_{S/k}$ es cerrado ($0_{\mathrm I}$, 7.3.5), de donde se sigue el lema.

Apliquemos este lema al caso $R=A_0$, $S=A$, y observemos que $\hat{\Omega}^1_{A_0/k}$ es un $A_0$-módulo libre de rango $n$ (0, 21.9.3); por tanto $\hat{\Omega}^1_{A_0/k}\widehat{\otimes}_{A_0}A=\hat{\Omega}^1_{A_0/k}\otimes_{A_0}A$ y este $A$-módulo es libre de rango $n$. Puesto que $\Spec(B)$ no es ramificado sobre $\Spec(B_0)$ en el punto $\mathfrak p'$, se tiene $(\Omega^1_{A/A_0})_{\mathfrak p}=\Omega^1_{B_{\mathfrak p'}/B_0}=0$ (16.4.15 y 17.4.1). Si localizamos en $\mathfrak p$ la sucesión exacta (18.11.3.6) (aplicada a $A_0$ y $A$), obtenemos por tanto un homomorfismo sobreyectivo
\[
  (\hat{\Omega}^1_{A_0/k}\otimes_{A_0}A)_{\mathfrak p}\longrightarrow(\hat{\Omega}^1_{A/k})_{\mathfrak p},
\]
de donde se concluye que el $A_{\mathfrak p}$-módulo $(\hat{\Omega}^1_{A/k})_{\mathfrak p}$ admite un sistema de $n$ generadores. Pero el $A$-módulo $\hat{\Omega}^1_{A/k}$ tiene rango $n$, en virtud de (0, 21.9.5), aplicable al anillo completo e íntegro $A$ por el teorema de Cohen (0, 19.8.8, (ii)) y por el hecho de que el cuerpo de fracciones de $A$ tiene característica 0. Por tanto, el $A_{\mathfrak p}$-módulo $(\hat{\Omega}^1_{A/k})_{\mathfrak p}$ también tiene rango $n$, y como su cociente por su submódulo de torsión admite un sistema de $n$ generadores, dicho cociente es necesariamente libre; se deduce inmediatamente que los $n$ generadores de $(\hat{\Omega}^1_{A/k})_{\mathfrak p}$ obtenidos arriba forman un sistema libre, de donde se sigue la conclusión.

\begin{lemma}[18.11.4]
\label{IV.18.11.4}
Sean $k$ un cuerpo, $A$ una $k$-álgebra local noetheriana completa, cuyo cuerpo residual es una extensión finita de $k$. Sea $k'$ una extensión cualquiera de $k$, y pongamos $A'=A\widehat{\otimes}_k k'$. Entonces:
\begin{enumerate}
  \item[(i)] $A'$ es un anillo semilocal noetheriano completo, suma directa de anillos locales completos $A'_i$ ($1\leq i\leq r$) que son $A$-módulos fielmente planos y cuyos cuerpos residuales son extensiones finitas de $k'$; si $\mathfrak m$ es el ideal maximal de $A$, $\mathfrak mA'$ es un ideal de definición de $A'$; se tiene $\dim(A'_i)=\dim(A)$ para todo $i$.
  \item[(ii)] Para todo $i$, $\hat{\Omega}^1_{A'_i/k'}$ es canónicamente isomorfo a $\hat{\Omega}^1_{A/k}\otimes_A A'_i$.
\end{enumerate}
\end{lemma}

\oldpage[IV-4]{174}

(i) Las primeras afirmaciones resultan inmediatamente de (7.5.5) y de ($0_{\mathrm I}$, 6.6.2); el hecho de que $\mathfrak mA'$ sea un ideal de definición de $A'$ resulta también de (7.5.5), pues si $K$ es el cuerpo residual de $A$, $K\otimes_k k'$ es una $k'$-álgebra finita. Finalmente, $A'_i/\mathfrak mA'_i$ es uno de los anillos locales artinianos componentes directos de $K\otimes_k k'$ (7.5.5), por tanto tiene dimensión 0; como $A'_i$ es un $A$-módulo plano, la igualdad de las dimensiones de $A$ y de $A'_i$ resulta de (6.1.2).

(ii) Como $\hat{\Omega}^1_{A/k}$ es, en virtud de la hipótesis sobre $K$, un $A$-módulo de tipo finito (18.11.1), $\hat{\Omega}^1_{A/k}\otimes_A A'$ es completo y se identifica con el producto tensorial completado $\hat{\Omega}^1_{A/k}\widehat{\otimes}_A A'$; en virtud de la asociatividad del producto tensorial completado ($0_{\mathrm I}$, 7.7.4), $\hat{\Omega}^1_{A/k}\widehat{\otimes}_A A'=\hat{\Omega}^1_{A/k}\widehat{\otimes}_A(A\widehat{\otimes}_k k')$ se identifica con $\hat{\Omega}^1_{A/k}\widehat{\otimes}_k k'$. Pero $\Omega^1_{A/k}\otimes_k k'$ se identifica con $\Omega^1_{A''/k'}$, donde $A''=A\otimes_k k'$ (0, 20.5.5), y como $A'$ es por definición el completado separado de $A''$, el completado separado de $\Omega^1_{A''/k'}$ se identifica por construcción con el de $\Omega^1_{A'/k'}$ (0, 20.7.4); en otras palabras, se tiene un isomorfismo canónico
\[
  \hat{\Omega}^1_{A/k}\widehat{\otimes}_k k'\xrightarrow{\sim}\hat{\Omega}^1_{A'/k'}.
\]
La conclusión de (ii) resulta ahora de que $\Omega^1_{A'/k'}$ es suma directa de los $\Omega^1_{A'_i/k'}$ (0, 20.4.13), y si $\mathfrak r$ es el radical de $A'$, la topología $\mathfrak r$-preádica sobre $\Omega^1_{A'/k'}$ se identifica con el producto de las topologías $\mathfrak m'_i$-preádicas sobre los $\Omega^1_{A'_i/k'}$ (donde $\mathfrak m'_i$ es el ideal maximal de $A'_i$); finalmente, se deduce que el completado separado $\hat{\Omega}^1_{A'/k'}$ para la topología $\mathfrak r$-preádica se identifica con el producto de los completados separados $\hat{\Omega}^1_{A'_i/k'}$ para las topologías $\mathfrak m'_i$-preádicas, y basta utilizar (0, 20.4.5).

\begin{proposition}[18.11.5]
\label{IV.18.11.5}
Sean $k$ un cuerpo, $A$ una $k$-álgebra local noetheriana y completa cuyo cuerpo residual es una extensión finita de $k$, $\mathfrak p$ un ideal primo de $A$, distinto del ideal maximal $\mathfrak m$, tal que exista un ideal primo mínimo $\mathfrak q\subset\mathfrak p$ para el cual $\dim(A/\mathfrak q)=\dim(A)$ (lo que ocurre en particular cuando $A$ es equidimensional), y $m$ un entero $\geq0$. Las condiciones siguientes son equivalentes:
\begin{enumerate}
  \item[$a)$] El $A_{\mathfrak p}$-módulo $(\hat{\Omega}^1_{A/k})_{\mathfrak p}$ admite un sistema de $m$ generadores.
  \item[$b)$] Existe un $k$-homomorfismo local $u:B\to A$, donde $B=k[[T_1,\ldots,T_m]]$, que hace de $A$ una $B$-álgebra finita, y tal que el morfismo correspondiente $\Spec(A)\to\Spec(B)$ sea no ramificado en el punto $\mathfrak p$.
\end{enumerate}
\end{proposition}

Para probar que $b)$ implica $a)$, observemos que, en virtud del lema (18.11.3.5), se tiene entonces la sucesión exacta
\[
  \hat{\Omega}^1_{B/k}\otimes_B A\longrightarrow\hat{\Omega}^1_{A/k}\longrightarrow\Omega^1_{A/B}\longrightarrow0
\]
puesto que $A$ es una $B$-álgebra finita; localizando en $\mathfrak p$ y observando que por hipótesis $(\Omega^1_{A/B})_{\mathfrak p}=0$ (17.4.1), obtenemos un homomorfismo sobreyectivo $(\hat{\Omega}^1_{B/k}\otimes_B A)_{\mathfrak p}\to(\hat{\Omega}^1_{A/k})_{\mathfrak p}$, y la conclusión resulta de que $\hat{\Omega}^1_{B/k}$ es un $B$-módulo libre de rango $m$ (0, 21.9.3).

Para probar que $a)$ implica $b)$, probemos primero los lemas siguientes.

\begin{lemma}[18.11.5.1]
\label{IV.18.11.5.1}
Sean $k$ un cuerpo, $A$ una $k$-álgebra local noetheriana y completa, cuyo cuerpo residual es una extensión finita de $k$, $\mathfrak p$ un ideal primo de $A$, y $\mathfrak q$ un ideal mínimo de $A$ contenido en $\mathfrak p$. Entonces el número mínimo de generadores del $A_{\mathfrak p}$-módulo $(\hat{\Omega}^1_{A/k})_{\mathfrak p}$ es mayor o igual que $\dim(A/\mathfrak q)$.
\end{lemma}

\oldpage[IV-4]{175}

Pongamos $n=\dim(A/\mathfrak q)$, y sea $m$ el número mínimo de generadores del $A_{\mathfrak p}$-módulo $(\hat{\Omega}^1_{A/k})_{\mathfrak p}$, que es igual a $\operatorname{rg}_{k(\mathfrak p)}((\hat{\Omega}^1_{A/k})_{\mathfrak p}\otimes_{A_{\mathfrak p}}k(\mathfrak p))$ (Bourbaki, \emph{Alg. comm.}, cap. II, \textsection~3, n\textsuperscript{o}~2, cor. 2 de la prop. 4). Observemos que $(\hat{\Omega}^1_{A/k})_{\mathfrak q}=(\hat{\Omega}^1_{A/k})_{\mathfrak p}\otimes_{A_{\mathfrak p}}A_{\mathfrak q}$, por tanto el número mínimo de generadores del $A_{\mathfrak q}$-módulo $(\hat{\Omega}^1_{A/k})_{\mathfrak q}$ es como máximo $m$. Basta, en consecuencia, considerar el caso en que $\mathfrak p=\mathfrak q$ es mínimo. En segundo lugar, mostremos que podemos reducirnos al caso en que $k$ es algebraicamente cerrado. En efecto, sea $k'$ una clausura algebraica de $k$, y pongamos $A'=A\widehat{\otimes}_k k'$, que es suma directa de $k'$-álgebras locales $A'_i$ ($1\leq i\leq r$), cuyos cuerpos residuales son isomorfos a $k'$, y que son $A$-módulos fielmente planos (18.11.4). Aplicando (2.3.4) y (6.1.1), se ve que, en un $A'_i$ dado, existe un ideal mínimo $\mathfrak q'_i$ sobre $\mathfrak q$ tal que $n'=\dim(A'_i/\mathfrak q'_i)\geq\dim(A/\mathfrak q)=n$. Por otra parte, por (18.11.4), se tiene $(\hat{\Omega}^1_{A'_i/k'})_{\mathfrak q'_i}=(\hat{\Omega}^1_{A/k})_{\mathfrak q}\otimes_{A_{\mathfrak q}}(A'_i)_{\mathfrak q'_i}$, por tanto el número mínimo $m'$ de generadores del $(A'_i)_{\mathfrak q'_i}$-módulo $(\hat{\Omega}^1_{A'_i/k'})_{\mathfrak q'_i}$ es como máximo $m$; de aquí se sigue nuestra afirmación.

Supongamos, pues, que $k$ es algebraicamente cerrado, y mostremos además que podemos reducirnos al caso $\mathfrak q=0$. En efecto, se tiene $\Omega^1_{(A/\mathfrak q)/A}=0$ (0, 20.4.12), por lo que la sucesión exacta (18.11.3.6) proporciona un homomorfismo sobreyectivo $\hat{\Omega}^1_{A/k}\otimes_A(A/\mathfrak q)\to\hat{\Omega}^1_{(A/\mathfrak q)/k}$, de donde se deduce inmediatamente que el número mínimo de generadores del $K$-módulo $\hat{\Omega}^1_{(A/\mathfrak q)/k}\otimes_{A/\mathfrak q}K$, donde $K$ es el cuerpo de fracciones de $A/\mathfrak q$, es como máximo $m$.

Pero si $A$ es íntegro y $K$ su cuerpo de fracciones, $K$ es separable sobre el cuerpo algebraicamente cerrado $k$, y por tanto una $k$-álgebra geométricamente regular (6.7.6); podemos aplicar entonces (18.11.3) para $\mathfrak p=0$, observando que $k=k^p$ puesto que $k$ es algebraicamente cerrado, y en este caso se tiene $m=n$.

\begin{lemma}[18.11.5.2]
\label{IV.18.11.5.2}
Sean $k$ un cuerpo, $A$ una $k$-álgebra local noetheriana completa, cuyo cuerpo residual es una extensión finita de $k$, $\mathfrak m$ su ideal maximal, $\mathfrak p\subset\mathfrak m$ un ideal primo de $A$. Sea $m$ el número mínimo de generadores del $A_{\mathfrak p}$-módulo $(\hat{\Omega}^1_{A/k})_{\mathfrak p}$, y sea $(x_i)_{1\leq i\leq m}$ una familia de elementos de $\mathfrak m$ que no pertenecen a $\mathfrak p$. Entonces existen $m$ elementos invertibles $u_i$ ($1\leq i\leq m$) de $A$ tales que, si ponemos $y_i=u_ix_i$, las imágenes canónicas de los $dy_i$ en $(\hat{\Omega}^1_{A/k})_{\mathfrak p}$ generan este $A_{\mathfrak p}$-módulo.
\end{lemma}

Para todo $x\in A$, designemos por $\delta(x)$ la imagen canónica de $dx$ ($=d_{A/k}x$) en $(\hat{\Omega}^1_{A/k})_{\mathfrak p}\otimes_{A_{\mathfrak p}}k(\mathfrak p)=\hat{\Omega}^1_{A/k}\otimes_A k(\mathfrak p)$; como este espacio vectorial sobre $k(\mathfrak p)$ tiene por hipótesis rango $m$, basta (en virtud del lema de Nakayama) probar que se pueden determinar los $u_i$ de modo que los $\delta(u_ix_i)$ formen un sistema libre. Razonemos por inducción y supongamos que, para un entero $r<m$, se han determinado los $u_i$ ($1\leq i\leq r$) de modo que los $\delta(u_ix_i)$ para $1\leq i\leq r$ sean linealmente independientes sobre $k(\mathfrak p)$; si $\delta(x_{r+1})$ no es combinación lineal de los $\delta(u_ix_i)$ para $1\leq i\leq r$, basta tomar $u_{r+1}=1$ para continuar la inducción. En caso contrario, observemos que, para $u\in A$, $\delta(ux_{r+1})$ es la imagen canónica de $(du)x_{r+1}+u(dx_{r+1})$; como la imagen canónica de $u(dx_{r+1})$ es combinación lineal de los $\delta(u_ix_i)$ para $1\leq i\leq r$, y, por otra parte, la imagen canónica de $x_{r+1}$ en $k(\mathfrak p)$ es $\neq0$, basta probar que existe un elemento invertible $u\in A$ tal que $\delta(u)$ no sea combinación lineal de los $\delta(u_ix_i)$ para $1\leq i\leq r$.

\oldpage[IV-4]{176}

Ahora bien, de (0, 20.7.15) y de ($0_{\mathrm I}$, 7.2.9) resulta que los $\delta(x)$ generan el espacio vectorial sobre $k(\mathfrak p)$ $\hat{\Omega}^1_{A/k}\otimes_A k(\mathfrak p)$; como por hipótesis $r<m$, existe por tanto $z\in A$ tal que $\delta(z)$ no sea combinación lineal de los $\delta(u_ix_i)$ para $1\leq i\leq r$. Si $z\notin\mathfrak m$, tomaremos $u=z$; en caso contrario, $1+z$ es invertible y $\delta(1+z)=\delta(z)$ puesto que $d(1+z)=dz$; tomaremos entonces $u=1+z$, lo que termina de demostrar (18.11.5.2).

Volvamos ahora a la demostración de la implicación $a)\Rightarrow b)$ en (18.11.5). Observemos primero que, puesto que $\mathfrak p\neq\mathfrak m$, existe un sistema de parámetros $(x_i)_{1\leq i\leq n}$ de $A$ (con $n=\dim(A)$) tal que $x_i\notin\mathfrak p$ para $1\leq i\leq n$: en efecto, no puede ocurrir que $x_i\in\mathfrak p$ para todo $i$, pues entonces $A/\mathfrak p$ tendría longitud finita y $\mathfrak p$ sería maximal, contra la hipótesis. Pero si, por ejemplo, $x_1\notin\mathfrak p$, basta, para cada índice $i$ tal que $x_i\in\mathfrak p$, sustituir $x_i$ por $x_1+x_i$, para obtener un sistema de parámetros ninguno de cuyos elementos pertenece a $\mathfrak p$. La hipótesis $a)$ y la relación $\dim(A/\mathfrak q)=n$ implican, en virtud de (18.11.5.1), que $m\geq n$; podemos por tanto considerar una familia $(x_i)_{1\leq i\leq m}$ de elementos $x_i\notin\mathfrak p$, cuyos $n$ primeros forman un sistema de parámetros de $A$. Multiplicando además los $x_i$ por elementos invertibles $u_i$ de $A$, podemos, gracias a (18.11.5.2), suponer que las imágenes de los $dx_i$ ($1\leq i\leq m$) en $(\hat{\Omega}^1_{A/k})_{\mathfrak p}$ generan este $A_{\mathfrak p}$-módulo, y la multiplicación por los $u_i$ no ha alterado el hecho de que los $x_i$ para $1\leq i\leq n$ forman un sistema de parámetros.

Consideremos entonces el $k$-homomorfismo local $u:B\to A$ tal que $u(T_i)=x_i$ para $1\leq i\leq m$ (Bourbaki, \emph{Alg. comm.}, cap. III, \textsection~4, n\textsuperscript{o}~5, prop. 6); como los $x_i$ generan un ideal de definición de $A$, se deduce de ($0_{\mathrm I}$, 7.4.4 y 7.4.3) que $u$ hace de $A$ una $B$-álgebra finita. Por tanto, por (18.11.3.5), se tiene la sucesión exacta
\[
  \hat{\Omega}^1_{B/k}\otimes_B A\xrightarrow{v}\hat{\Omega}^1_{A/k}\longrightarrow\Omega^1_{A/B}\longrightarrow0.
\]
Pero los $d_Ax_i$ son las imágenes canónicas por $v$ de los elementos $d_BT_i\otimes1$ (0, 20.5.2.6). Si localizamos en $\mathfrak p$ la sucesión exacta anterior, vemos por tanto que $v_{\mathfrak p}:(\hat{\Omega}^1_{B/k}\otimes_B A)_{\mathfrak p}\to(\hat{\Omega}^1_{A/k})_{\mathfrak p}$ es sobreyectivo, y por consiguiente $0=(\Omega^1_{A/B})_{\mathfrak p}=\Omega^1_{A_{\mathfrak p}/B_{\mathfrak r}}$ (donde $\mathfrak r$ es la imagen recíproca de $\mathfrak p$ en $B$, cf. (16.4.15)).\footnote{La transcripción francesa imprime «dans $A$»; el testigo inglés y la definición $\mathfrak r=u^{-1}(\mathfrak p)$ exigen «en $B$». Se conserva aquí la discrepancia de la fuente.} En virtud de (17.4.1), esto implica la propiedad $b)$ de (18.11.5). C.Q.F.D.

\begin{remark}[18.11.6]
\label{IV.18.11.6}
En el enunciado de (18.11.5) no se puede suprimir la hipótesis $\mathfrak p\neq\mathfrak m$. En efecto, para todo $k$-homomorfismo local $u:B\to A$, donde $B=k[[T_1,\ldots,T_r]]$ ($r$ un entero cualquiera) que hace de $A$ una $B$-álgebra finita, $\mathfrak m$ es el único punto de $\Spec(A)$ sobre el ideal maximal $\mathfrak n$ de $B$; el morfismo $\Spec(A)\to\Spec(B)$ solo puede ser no ramificado en $\mathfrak m$ si $A/\mathfrak nA$ es un cuerpo, extensión separable de $k$ (17.4.1), lo que implica que el cuerpo residual $K$ de $A$ es una extensión separable de $k$. Si no se cumple esta condición, la conclusión de (18.11.5) nunca puede verificarse para $\mathfrak p=\mathfrak m$, cualquiera que sea el entero $m$.
\end{remark}

\begin{corollary}[18.11.7]
\label{IV.18.11.7}
Sean $k$ un cuerpo, $A$ una $k$-álgebra local noetheriana completa equidimensional, cuyo cuerpo residual es una extensión finita de $k$, y $\mathfrak p$ un ideal primo de $A$. Entonces el número mínimo de generadores del $A_{\mathfrak p}$-módulo $(\hat{\Omega}^1_{A/k})_{\mathfrak p}$ es mayor o igual que $\dim(A)$.
\end{corollary}

Este es un caso particular de (18.11.5.1).

Más particularmente:

\oldpage[IV-4]{177}

\begin{corollary}[18.11.8]
\label{IV.18.11.8}
Sean $k$ un cuerpo, $A$ una $k$-álgebra local noetheriana completa íntegra, cuyo cuerpo residual es una extensión finita de $k$. Si $K$ es el cuerpo de fracciones de $A$, se tiene
\[
  \operatorname{rg}_K\bigl((\hat{\Omega}^1_{A/k})\otimes_A K\bigr)\geq\dim(A).
\]
\end{corollary}

Basta tomar $\mathfrak p=0$ en (18.11.7).

\begin{corollary}[18.11.9]
\label{IV.18.11.9}
Sean $k$ un cuerpo, $A$ una $k$-álgebra local noetheriana completa de dimensión $n$, cuyo cuerpo residual es una extensión finita de $k$. Sea $\mathfrak p$ un ideal primo de $A$ distinto del ideal maximal y que contiene un ideal primo mínimo $\mathfrak q$ tal que $\dim(A/\mathfrak q)=n$. Supongamos que el $A_{\mathfrak p}$-módulo $(\hat{\Omega}^1_{A/k})_{\mathfrak p}$ admite $n$ generadores. Entonces:
\begin{enumerate}
  \item[(i)] El $A_{\mathfrak p}$-módulo $(\hat{\Omega}^1_{A/k})_{\mathfrak p}$ es libre de rango $n$.
  \item[(ii)] Existe un $k$-homomorfismo local $u:B\to A$, donde $B=k[[T_1,\ldots,T_n]]$, que hace de $A$ una $B$-álgebra finita, y tal que el morfismo correspondiente $\Spec(A)\to\Spec(B)$ sea étale en el punto $\mathfrak p$.
  \item[(iii)] La $k$-álgebra $A_{\mathfrak p}$ es geométricamente regular.
\end{enumerate}
\end{corollary}

Probemos primero (ii); en virtud de (18.11.5), existe un homomorfismo local
\[
  u:B=k[[T_1,\ldots,T_n]]\longrightarrow A
\]
que hace de $A$ una $B$-álgebra finita, y tal que el morfismo correspondiente $\Spec(A)\to\Spec(B)$ sea no ramificado en el punto $\mathfrak p$; pongamos $\mathfrak r=u^{-1}(\mathfrak p)$. La hipótesis sobre $\mathfrak q\subset\mathfrak p$ y el hecho de que $A$ sea un cociente de un anillo regular (0, 19.8.8) implican que $\dim(A_{\mathfrak p})=n-\dim(A/\mathfrak p)$ (0, 16.5.12). Análogamente, $\dim(B_{\mathfrak r})=n-\dim(B/\mathfrak r)$. Finalmente, puesto que el morfismo $\Spec(A)\to\Spec(B)$ es no ramificado en el punto $\mathfrak p$, la fibra de este morfismo en el punto $\mathfrak r$ tiene dimensión 0, por lo que (0, 16.3.9) se tiene $\dim(A/\mathfrak p)\leq\dim(B/\mathfrak r)$ y, por consiguiente, $\dim(B_{\mathfrak r})\leq\dim(A_{\mathfrak p})$. Concluimos entonces de (18.10.1) que el morfismo $\Spec(A)\to\Spec(B)$ es étale en el punto $\mathfrak p$.

La afirmación (iii) resulta de que $A_{\mathfrak p}$ es formalmente liso sobre $B_{\mathfrak r}$ y que $B_{\mathfrak r}$ es formalmente liso sobre $k$ (para las topologías preádicas) (0, 19.3.4 y 19.3.5); por tanto, $A_{\mathfrak p}$ es formalmente liso sobre $k$ para su topología preádica y, en consecuencia, geométricamente regular sobre $k$ (0, 22.5.8).

Probemos finalmente (i). Sea $k'$ una extensión algebraicamente cerrada de $k$, y consideremos la $k'$-álgebra semilocal $A'=A\widehat{\otimes}_k k'$. Si consideramos $A$ como $B$-módulo de tipo finito mediante el homomorfismo $u$, podemos escribir, salvo isomorfismo canónico, $A'=A\otimes_B(B\widehat{\otimes}_k k')$ ((7.5.7.1), cuyo enunciado recuerda que no es necesario suponer que el cuerpo residual de $B$ sea una extensión finita de $k$). Pongamos $B'=B\widehat{\otimes}_k k'$, que se identifica canónicamente con el álgebra de series formales $k'[[T_1,\ldots,T_n]]$. Puesto que el morfismo $\Spec(A)\to\Spec(B)$ es finito y étale en el punto $\mathfrak p$, el morfismo $\Spec(A')\to\Spec(B')$ es finito y étale en todo punto $\mathfrak p'$ sobre $\mathfrak p$ (17.3.3); además, $A'$ es suma directa de anillos locales $A'_i$ de dimensión $n$ ($1\leq i\leq r$) (18.11.4) y $\mathfrak p'$ se identifica con un ideal primo $\mathfrak p'_i$ de uno de los $A'_i$. El mismo razonamiento que arriba prueba entonces que $(A'_i)_{\mathfrak p'_i}$ es geométricamente regular sobre $k'$, y como $k'$ es perfecto, resulta de (18.11.3) que $(\hat{\Omega}^1_{A'_i/k'})_{\mathfrak p'_i}$ es un $(A'_i)_{\mathfrak p'_i}$-módulo libre de tipo finito. Pero como $(\hat{\Omega}^1_{A'_i/k'})_{\mathfrak p'_i}$ es isomorfo a $(\hat{\Omega}^1_{A/k})_{\mathfrak p}\otimes_{A_{\mathfrak p}}(A'_i)_{\mathfrak p'_i}$ (18.11.4) y $(A'_i)_{\mathfrak p'_i}$ es un $A_{\mathfrak p}$-módulo fielmente plano, se ve que $(\hat{\Omega}^1_{A/k})_{\mathfrak p}$ es un $A_{\mathfrak p}$-módulo libre de tipo finito (2.5.2);

\oldpage[IV-4]{178}

además, en virtud de (18.11.5.1) y de la hipótesis $\dim(A/\mathfrak q)=n$, concluimos que $(\hat{\Omega}^1_{A/k})_{\mathfrak p}$ es un $A_{\mathfrak p}$-módulo libre de rango $n$. C.Q.F.D.

\begin{theorem}[18.11.10]
\label{IV.18.11.10}
Sean $k$ un cuerpo de exponente característico $p$, $A$ una $k$-álgebra local noetheriana completa, cuyo cuerpo residual es una extensión finita de $k$, y $\mathfrak p$ un ideal primo de $A$ distinto del ideal maximal. Las condiciones siguientes son equivalentes:
\begin{enumerate}
  \item[$a)$] Para toda extensión $k'$ de $k$ y todo ideal primo $\mathfrak p'$ de $A'=A\widehat{\otimes}_k k'$ sobre $\mathfrak p$, $A'_{\mathfrak p'}$ es un anillo regular.
  \item[$a')$] Existe una extensión perfecta $k'$ de $k$ y un ideal primo $\mathfrak p'$ de $A'=A\widehat{\otimes}_k k'$ sobre $\mathfrak p$, tales que $A'_{\mathfrak p'}$ sea regular.
  \item[$b)$] Sea $n$ la mayor de las dimensiones de las componentes irreducibles de $\Spec(A)$ a las que pertenece $\mathfrak p$; entonces $(\hat{\Omega}^1_{A/k})_{\mathfrak p}$ es un $A_{\mathfrak p}$-módulo libre de rango $n$.
\end{enumerate}
Supongamos además que $n=\dim(A)$ (lo que ocurre si $A$ es equidimensional); entonces las condiciones anteriores son también equivalentes a:
\begin{enumerate}
  \item[$c)$] Existe un $k$-homomorfismo local $u:B\to A$, donde $B=k[[T_1,\ldots,T_n]]$, que hace de $A$ una $B$-álgebra finita, y tal que el morfismo correspondiente $\Spec(A)\to\Spec(B)$ sea étale en el punto $\mathfrak p$.
\end{enumerate}
Cada una de las condiciones $a)$, $a')$, $b)$ implica la siguiente:
\begin{enumerate}
  \item[$d)$] El anillo $A_{\mathfrak p}$ es geométricamente regular sobre $k$.
\end{enumerate}
Si además $[k:k^p]<+\infty$, la condición $d)$ es equivalente a $a)$, $a')$ y $b)$.
\end{theorem}

El hecho de que $d)$ implique $b)$ cuando $[k:k^p]<+\infty$ no es otro que (18.11.3). Es claro que $a)$ implica trivialmente $a')$; mostremos que $c)$ implica $a)$ cuando $\dim(A)=n$. Con las notaciones de $a)$, se tiene entonces
\[
  A'=A\otimes_B B',\qquad B'=B\widehat{\otimes}_k k'=k'[[T_1,\ldots,T_n]]
\]
(7.5.7.1). El morfismo $\Spec(A')\to\Spec(B')$ es entonces finito y étale en el punto $\mathfrak p'$, y el razonamiento que prueba (18.11.9, (iii)) muestra que $A'_{\mathfrak p'}$ es regular.

El hecho de que $b)$ implique $c)$ cuando $\dim(A)=n$ resulta de (18.11.9). Mostremos que $a')$ implica $b)$ cuando $\dim(A)=n$. Se sabe que $A'$ es un $A$-módulo plano (18.11.4) y $\dim(A')=\dim(A)=n$; de (2.3.4) y (6.1.1) resulta que existe un ideal primo mínimo $\mathfrak q'$ de $A'$ contenido en $\mathfrak p'$, sobre $\mathfrak q$, tal que $\dim(A'/\mathfrak q')=\dim(A/\mathfrak q)=n$. Por otra parte, resulta de (18.11.4) que
\[
  (\hat{\Omega}^1_{A'/k'})_{\mathfrak p'}=(\hat{\Omega}^1_{A/k})_{\mathfrak p}\otimes_{A_{\mathfrak p}}A'_{\mathfrak p'};
\]
como $k'$ es perfecto, la hipótesis de que $A'_{\mathfrak p'}$ es regular implica que es geométricamente regular sobre $k'$ (6.7.7). Puesto que $k'^p=k'$, podemos aplicar a $A'$ y $\mathfrak p'$ el hecho de que $d)$ implica $b)$; por tanto $(\hat{\Omega}^1_{A'/k'})_{\mathfrak p'}$ es un $A'_{\mathfrak p'}$-módulo libre de rango $n$; por fidelidad plana (18.11.4 y 2.5.2), concluimos que $(\hat{\Omega}^1_{A/k})_{\mathfrak p}$ es un $A_{\mathfrak p}$-módulo libre de rango $n$.

Esto demuestra la equivalencia de $a)$, $a')$, $b)$ y $c)$ cuando $\dim(A)=n$. Falta probar que $a)$, $a')$ y $b)$ siguen siendo equivalentes en el caso general. Sea $\mathfrak J$ el ideal de $A$, núcleo del homomorfismo canónico $A\to A_{\mathfrak p}$, y pongamos $A_1=A/\mathfrak J$; necesariamente $\mathfrak J\subset\mathfrak p$, y si ponemos $\mathfrak p_1=\mathfrak p/\mathfrak J$, el homomorfismo canónico $A_{\mathfrak p}\to(A_1)_{\mathfrak p_1}$ es biyectivo; se concluye (I, 6.5.4) que la inclusión canónica $\Spec(A_1)\to\Spec(A)$ es un isomorfismo local en el punto $\mathfrak p_1$, y se tiene $\mathfrak J_{\mathfrak p}=0$. Como en (18.11.3), se obtiene la sucesión exacta
\[
  \mathfrak J/\mathfrak J^2\longrightarrow\hat{\Omega}^1_{A/k}\otimes_A A_1\longrightarrow\hat{\Omega}^1_{A_1/k}\longrightarrow0.
\]

\oldpage[IV-4]{179}

Y al localizar en $\mathfrak p$ se obtiene un isomorfismo
\[
  (\hat{\Omega}^1_{A/k})_{\mathfrak p}\xrightarrow{\sim}(\hat{\Omega}^1_{A_1/k})_{\mathfrak p_1}.
\]
Esto demuestra que la condición $b)$ para el anillo $A$ y el ideal $\mathfrak p$ es equivalente a la condición $b)$ para el anillo $A_1$ y el ideal $\mathfrak p_1$. Por otra parte, con las notaciones de $a)$, se tiene $A'_1=A_1\widehat{\otimes}_k k'=A'/\mathfrak JA'$ salvo isomorfismo ((7.5.7.1), donde la hipótesis sobre el cuerpo residual de $B$ es superflua); si $\mathfrak p'$ es un ideal primo de $A'$ sobre $\mathfrak p$, todo elemento de $\mathfrak JA'$ anula a un elemento de $A'-\mathfrak p'$, luego $\mathfrak JA'\subset\mathfrak p'$, y si ponemos $\mathfrak p'_1=\mathfrak p'/\mathfrak JA'$, $\mathfrak p'_1$ está sobre $\mathfrak p_1$ y $(A'_1)_{\mathfrak p'_1}$ se identifica canónicamente con $A'_{\mathfrak p'}$; esto muestra, pues, que la condición $a)$ (respectivamente $a')$) para el anillo $A$ y el ideal $\mathfrak p$ es equivalente a la condición $a)$ (respectivamente $a')$) para el anillo $A_1$ y el ideal $\mathfrak p_1$. Ahora bien, todos los ideales primos mínimos de $A$ contenidos en $\mathfrak p$ contienen $\mathfrak J$, puesto que $\Spec(A_1)\to\Spec(A)$ es un isomorfismo local en el punto $\mathfrak p_1$; por otra parte, los ideales de $\operatorname{Ass}_A(A/\mathfrak J)$ son los ideales de $\operatorname{Ass}(A)$ que están contenidos en $\mathfrak p$ (Bourbaki, \emph{Alg. comm.}, cap. IV, \textsection~1, n\textsuperscript{o}~2, prop. 6); por tanto, los ideales primos mínimos de $A_1$ están todos contenidos en $\mathfrak p_1$, y se tiene en consecuencia $\dim(A_1)=n$. Basta entonces aplicar a $A_1$ y a $\mathfrak p_1$ lo demostrado arriba.

\begin{remarks}[18.11.11]
\label{IV.18.11.11}
\begin{enumerate}
  \item[(i)] La equivalencia de las condiciones $d)$ y $b)$ en (18.11.10) ya no es válida si no se supone $[k:k^p]<+\infty$. En efecto, en el ejemplo de (0, 22.7.7, (ii)), el anillo $B=A/\mathfrak q$ es íntegro y de dimensión 1; por otra parte, la sucesión
  \[
    (\mathfrak q/\mathfrak q^2)\otimes_A L\xrightarrow{j}\hat{\Omega}^1_{A/k}\otimes_A L\longrightarrow\hat{\Omega}^1_{B/k}\otimes_B L\longrightarrow0
  \]
  es exacta (misma demostración que en (18.11.3)), y se sabe que $j$ no es inyectivo, luego $j=0$ puesto que $(\mathfrak q/\mathfrak q^2)\otimes_A L$ tiene rango 1. Concluimos que $\operatorname{rg}_L(\hat{\Omega}^1_{B/k}\otimes_B L)=2$. Como $B_{\mathfrak p}$ es una $k$-álgebra geométricamente regular, se ve que aquí la condición $d)$ no implica $b)$.
  \item[(ii)] Con las notaciones de (18.11.10), supongamos que $n=\dim(A)$ y sea $(x_i)_{1\leq i\leq n}$ un sistema de parámetros de $A$ que no pertenece a $\mathfrak p$; de ello se deduce un $k$-homomorfismo local $u:B\to A$ tal que $u(T_i)=x_i$ para $1\leq i\leq n$, que hace de $A$ una $B$-álgebra finita. Para que el morfismo correspondiente $\Spec(A)\to\Spec(B)$ sea étale en el punto $\mathfrak p$, es necesario y suficiente que las imágenes de los $d_Ax_i$ en $(\hat{\Omega}^1_{A/k})_{\mathfrak p}$ formen un sistema de generadores de este $A_{\mathfrak p}$-módulo. En efecto, en la demostración de (18.11.5) vimos que, si el morfismo $\Spec(A)\to\Spec(B)$ es no ramificado en el punto $\mathfrak p$, el homomorfismo canónico $(\hat{\Omega}^1_{B/k}\otimes_B A)_{\mathfrak p}\to(\hat{\Omega}^1_{A/k})_{\mathfrak p}$ es sobreyectivo y las imágenes de los elementos $d_BT_i\otimes1$, que generan $(\hat{\Omega}^1_{B/k}\otimes_B A)_{\mathfrak p}$, son los $d_Ax_i$; de aquí la necesidad de la condición. Recíprocamente, el mismo razonamiento muestra que si se cumple esta condición, el morfismo $\Spec(A)\to\Spec(B)$ es no ramificado en $\mathfrak p$, y el razonamiento de (18.11.9) prueba de hecho que este morfismo es étale en $\mathfrak p$.
\end{enumerate}
\end{remarks}

\begin{corollary}[18.11.12]
\label{IV.18.11.12}
Sean $k$ un cuerpo de exponente característico $p$ tal que $[k:k^p]<+\infty$, $A$ una $k$-álgebra local noetheriana completa e íntegra, que no es un cuerpo y cuyo cuerpo residual es una extensión finita de $k$. Existe entonces una extensión finita radicial $k'$ de $k$ tal que, poniendo $A'=A\otimes_k k'$, la $k'$-álgebra $A'_{\mathrm{red}}$ y el ideal primo $0$ de esta álgebra satisfacen las condiciones equivalentes $a)$, $a')$, $b)$, $c)$ y $d)$ de (18.11.10). En particular, si $n=\dim(A)=\dim(A')$,

\oldpage[IV-4]{180}

existe un $k'$-homomorfismo local $B'=k'[[T_1,\ldots,T_n]]\to A'_{\mathrm{red}}$ que hace de $A'_{\mathrm{red}}$ una $B'$-álgebra finita, y tal que el cuerpo de fracciones $K'$ de $A'_{\mathrm{red}}$ sea una extensión finita separable del cuerpo de fracciones $k'((T_1,\ldots,T_n))$ de $B'$.
\end{corollary}

Como el morfismo $\Spec(A')\to\Spec(A)$ es radicial y finito, $A'$ es un anillo local completo y su nilradical es el único ideal primo sobre el ideal $0$ de $A$; por planitud, $A'$ se identifica con un subanillo de $K\otimes_k k'$, que es además el anillo total de fracciones de $A'$; por tanto $K'=(K\otimes_k k')_{\mathrm{red}}$. Hay que probar, teniendo en cuenta la equivalencia $d)\Leftrightarrow c)$ de (18.11.10), que existe una extensión radicial finita $k'$ de $k$ tal que $(K\otimes_k k')_{\mathrm{red}}$ sea una extensión separable de $k'$ (6.7.6). Ahora bien, se sabe (0, 19.8.9) que bajo las hipótesis dadas existe una $k$-subálgebra $C=k[[T_1,\ldots,T_m]]$ de $A$ tal que $A$ sea una $C$-álgebra finita; por tanto $K$ es una extensión finita del cuerpo de fracciones $K_1=k((T_1,\ldots,T_m))$ de $C$, que es separable sobre $k$ (0, 21.9.6.4). Si $p$ es el exponente característico de $k$, podemos escribir entonces
\[
  K\otimes_k k^{p^{-\infty}}=K\otimes_{K_1}(K_1\otimes_k k^{p^{-\infty}})
\]
y $K_1\otimes_k k^{p^{-\infty}}$ es un cuerpo, extensión radicial de $K_1$; se concluye que $K\otimes_k k^{p^{-\infty}}$ es una álgebra finita sobre el cuerpo $K_1\otimes_k k^{p^{-\infty}}$, y por tanto un anillo artiniano. La conclusión resulta, pues, del lema más general siguiente:

\begin{lemma}[18.11.12.1]
\label{IV.18.11.12.1}
Sean $k$ un cuerpo de característica $p>0$ y $K$ una extensión de $k$. Las condiciones siguientes son equivalentes:
\begin{enumerate}
  \item[$a)$] El anillo $K\otimes_k k^{p^{-\infty}}$ es artiniano.
  \item[$b)$] Existe una extensión $k'$ de $k$ de tipo finito tal que $(K\otimes_k k')_{\mathrm{red}}$ sea una $k'$-álgebra separable.
  \item[$c)$] Existe una extensión finita radicial $k'$ de $k$ tal que $(K\otimes_k k')_{\mathrm{red}}$ sea una extensión separable de $k'$.
\end{enumerate}
\end{lemma}

Mostremos primero que $c)$ implica $a)$. El anillo $A=K\otimes_k k'$ es entonces un anillo artiniano local, y si $\mathfrak N$ es su nilradical, el cuerpo residual $L=A/\mathfrak N$ es separable sobre $k'$; por consiguiente $L\otimes_{k'}k^{p^{-\infty}}$ es un cuerpo, y como es igual a
\[
  (A\otimes_{k'}k^{p^{-\infty}})/(\mathfrak N\otimes_{k'}k^{p^{-\infty}}),
\]
se ve que $\mathfrak N\otimes_{k'}k^{p^{-\infty}}$ es el nilradical de $A\otimes_{k'}k^{p^{-\infty}}=K\otimes_k k^{p^{-\infty}}$; como $\mathfrak N$ es un ideal de tipo finito, el nilradical de $K\otimes_k k^{p^{-\infty}}$ es, pues, de tipo finito, lo que implica que el anillo $K\otimes_k k^{p^{-\infty}}$ es artiniano.

Recíprocamente, probemos que $a)$ implica $c)$. Sea $\mathfrak N$ el nilradical del anillo local artiniano $B=K\otimes_k k^{p^{-\infty}}$, que por hipótesis está generado por un número finito de elementos de la forma
\[
  y_i=\sum_j x_{ij}\otimes\xi_{ij},\qquad x_{ij}\in K,\quad \xi_{ij}\in k^{p^{-\infty}}.
\]
Sean $k'$ la extensión radicial finita de $k$ generada por los $\xi_{ij}$, y $\mathfrak N_0$ el ideal de $B_0=K\otimes_k k'$ generado por los $y_i$; es claro que los $y_i$ son nilpotentes en $B_0$; por otra parte, $\mathfrak N_0\otimes_{k'}k^{p^{-\infty}}=\mathfrak N$, y por tanto $\mathfrak N\cap B_0=\mathfrak N_0$; $\mathfrak N_0$ contiene el nilradical de $B_0$, luego es igual a él. Como $B/\mathfrak N=(B_0/\mathfrak N_0)\otimes_{k'}k^{p^{-\infty}}$ es reducido, se concluye que $B_0/\mathfrak N_0=(K\otimes_k k')_{\mathrm{red}}$ es separable sobre $k'$ (4.6.1).\footnote{La transcripción francesa imprime «sobre $k$»; el testigo inglés y la extensión de escala indican «sobre $k'$». Se conserva aquí la discrepancia de la fuente.}

Es claro que $c)$ implica $b)$. Recíprocamente, supongamos que se cumple $b)$, y observemos que existe una extensión separable $k_1$ de $k$ tal que $k'$ sea una extensión radicial finita de $k_1$; pongamos $K_1=K\otimes_k k_1$, que es un cuerpo. Aplicando la equivalencia de $a)$ y $c)$ a la extensión $K_1$ de $k_1$, se ve que $K_1\otimes_{k_1}k_1^{p^{-\infty}}$ es un anillo artiniano; pero este anillo es igual

\oldpage[IV-4]{181}

a
\[
  K\otimes_k k_1^{p^{-\infty}}=(K\otimes_k k^{p^{-\infty}})\otimes_{k^{p^{-\infty}}}k_1^{p^{-\infty}},
\]
por tanto $K\otimes_k k^{p^{-\infty}}$ también es artiniano (Bourbaki, \emph{Alg. comm.}, cap. I, \textsection~3, n\textsuperscript{o}~5, cor. de la prop. 8); así hemos probado que $b)$ implica $a)$, lo que termina la demostración de (18.11.12.1) y de (18.11.12).

\subsection{Aplicaciones de la localización étale a los morfismos cuasi-finitos (generalizaciones de resultados anteriores)}
\label{subsection:IV.18.12}

Los resultados de este número nos fueron comunicados por P.~Deligne.

\begin{theorem}[18.12.1]
\label{IV.18.12.1}
Sean $f:X\to Y$ un morfismo localmente de tipo finito, $x$ un punto de $X$, $y=f(x)$. Supongamos que $x$ sea un punto aislado del espacio $f^{-1}(y)$. Entonces, existe un morfismo étale $Y'\to Y$, un punto $x'$ de $X'=X\times_Y Y'$ por encima de $x$ y un entorno abierto $V'$ de $x'$ en $X'$ tal que, si $f'=f_{(Y')}:X'\to Y'$, $f'|V'$ sea un morfismo finito. Si además $f$ es separado, $V'$ es a la vez abierto y cerrado, y $X'$ es por tanto suma de dos subpreesquemas inducidos sobre abiertos de $X'$, uno de los cuales es finito sobre $Y'$ y contiene $x'$.
\end{theorem}

La última afirmación se deduce de que, si $f$ es separado, lo mismo ocurre con $f'$; por tanto, si $j':V'\to X'$ es la inyección canónica, el hecho de que el morfismo $f'|V'=f'\circ j'$ sea finito implica que lo mismo ocurre con $j'$ (II, 6.1.5), y en consecuencia $V'=j'(V')$ es cerrado en $X'$ (II, 6.1.10).

La cuestión es local en $X$, por lo que podemos suponer $Y=\Spec(A)$ y $X=\Spec(B)$ afines, siendo $B$ una $A$-álgebra de tipo finito, por tanto de la forma $C/\mathfrak J$, donde $C=A[T_1,\ldots,T_r]$ y $\mathfrak J$ es un ideal de $C$; $f$ es entonces separado, y podemos además suponer que $f^{-1}(y)=\{x\}$. Sea $(\mathfrak J_\lambda)$ la familia de los ideales de tipo finito de $C$ contenidos en $\mathfrak J$, de modo que $\mathfrak J$ es la unión filtrante de los $\mathfrak J_\lambda$. Si $X$ y los $X_\lambda=\Spec(C/\mathfrak J_\lambda)$ se consideran como subpreesquemas cerrados de $Z=\Spec(C)$, se tiene por tanto, para los espacios subyacentes, $X=\bigcap_\lambda X_\lambda$; si $f_\lambda:X_\lambda\to Y$ es el morfismo estructural, se deduce que
\[
  f^{-1}(y)=\bigcap_\lambda f_\lambda^{-1}(y),
\]
y como los conjuntos $f_\lambda^{-1}(y)$ son cerrados en el espacio noetheriano $Z_y=\Spec(k(y)[T_1,\ldots,T_r])$, existe un índice $\lambda$ tal que $f_\lambda^{-1}(y)=f^{-1}(y)=\{x\}$. Podemos, pues, suponer que los $X_\lambda$ verifican la misma condición que $X$ en el punto $x$; si se demuestra la proposición para tal $X_\lambda$, se deduce para $X$, pues si $x'$ es un punto por encima de $x$ en $Z'=Z\times_Y Y'$, $X'_\lambda=p^{-1}(X_\lambda)$, donde $p:Z'\to Z$ es la proyección (de modo que también se tiene $X'=p^{-1}(X)$), si existe un entorno abierto $V'_\lambda$ de $x'$ en $X'_\lambda$ que sea finito sobre $Y'$, \emph{a fortiori} la restricción de $f'_\lambda$ al subpreesquema cerrado $X'\cap V'_\lambda$ de $V'_\lambda$ será un morfismo finito. Nos reducimos, pues, al caso en que $f$ es de presentación finita.

Observemos además que el conjunto de puntos de $X$ aislados en su fibra es abierto en $X$ (13.1.4), por lo que $f$ es cuasi-finito en el punto $x$, y podemos restringirnos por tanto al caso en que $f$ es cuasi-finito. Sea $A''$ el anillo local henselizado de $\sh{O}_{Y,y}$, y pongamos $Y''=\Spec(A'')$. Si $X''=X\times_Y Y''$, $f''=f_{(Y'')}:X''\to Y''$ es separado, cuasi-finito y de presentación finita; como $A''$ es henseliano, para todo $x''\in X''$ por encima de $x$ hay un entorno abierto y cerrado $V''$ de $x''$ en $X''$ que es finito sobre $Y''$ (18.5.11, c)); como la inmersión $V''\to X''$ es abierta y cerrada, es cuasi-compacta, por tanto de presentación finita (1.6.2), y en consecuencia $f''|V''$ es de presentación finita.

Dicho esto, por definición, $A''$ es límite inductivo de una familia filtrante $(A_\lambda)$ de $\sh{O}_{Y,y}$-álgebras estrictamente esencialmente étale (18.6.5); cada $A_\lambda$ es a su vez límite inductivo de una familia filtrante $(B_{\lambda\mu})$ de $\sh{O}_{Y,y}$-álgebras étale ((18.6.1) y (8.1.2, a)). Finalmente, como se puede suponer que $B_{\lambda\mu}$ es una $\sh{O}_{Y,y}$-álgebra de presentación finita, se deduce también de (8.1.2, a)) y de (17.7.8) que $B_{\lambda\mu}$ es límite inductivo de una familia filtrante $(C_{\lambda\mu\nu})$ de $A$-álgebras étale. Finalmente, se ve, usando el teorema de la doble límite inductiva, que $Y''$ es límite proyectivo de una familia filtrante $(Y'_\alpha)$ de esquemas afines étale sobre $Y$. Si $x'_\alpha$ es la proyección de $x''$ sobre $X'_\alpha=X\times_Y Y'_\alpha$, podemos suponer que el entorno $V''$, que es cuasi-compacto, es de la forma $V'_\alpha\times_{Y'_\alpha}Y''$ donde $V'_\alpha$ es un entorno abierto de $x'_\alpha$ en $X'_\alpha$ (8.2.11). Finalmente, puesto que $V''$ es de presentación finita sobre $Y''$, se concluye de (8.10.5, (x)) que para un $\alpha$ adecuado, $V'_\alpha$ es finito sobre $Y'_\alpha$. C.Q.F.D.

\begin{remark}[18.12.2]
\label{IV.18.12.2}
En la demostración anterior se ve (teniendo en cuenta (18.6.2)) que hemos construido un $Y'$ que responde a la cuestión, y además que, si $y'=f'(x')$, el homomorfismo $k(y)\to k(y')$ es biyectivo.\footnote{La transcripción francesa imprime «$f'(X')$»; el testigo inglés mantenido y la notación del punto del teorema requieren «$f'(x')$». Se conserva aquí la discrepancia de la fuente.}
\end{remark}

\begin{corollary}[18.12.3]
\label{IV.18.12.3}
Sean $f:X\to Y$ un morfismo localmente de tipo finito y separado, $y$ un punto de $Y$ tal que el subespacio $f^{-1}(y)$ sea finito y discreto. Existe entonces un morfismo étale $Y'\to Y$, un punto $y'\in Y'$ por encima de $y$, tal que $k(y')=k(y)$, y una descomposición de $X'$ en suma de dos subpreesquemas $X'_1$, $X'_2$, inducidos sobre abiertos de $X'$, tales que la restricción de $f'=f_{(Y')}:X'\to Y'$ a $X'_1$ sea un morfismo finito y se tenga $X'_2\cap f'^{-1}(y')=\varnothing$.
\end{corollary}

Si $n$ es el número de puntos de $f^{-1}(y)$, razonamos por inducción sobre $n$, siendo el corolario trivial para $n=0$. Sea $x$ un punto de $f^{-1}(y)$; en virtud de (18.12.1) y (18.12.2) hay un morfismo étale $Y_1\to Y$, un punto $y_1$ de $Y_1$ por encima de $y$ tal que $k(y_1)=k(y)$ y, si ponemos $S_1=X\times_Y Y_1$, $f_1=f_{(Y_1)}:S_1\to Y_1$, existe un punto $x_1$ de $f_1^{-1}(y_1)$ tal que $S_1$ es suma de dos abiertos $V_1$, $X_1$, siendo $V_1$ finito sobre $Y_1$ y un entorno de $x_1$.\footnote{La transcripción francesa imprime «$f_{(S_1)}$»; el testigo inglés mantenido y la notación del cambio de base requieren «$f_{(Y_1)}$». Se conserva aquí la discrepancia de la fuente.}

\oldpage[IV-4]{182}

En virtud de la igualdad $k(y_1)=k(y)$, la fibra $f_1^{-1}(y_1)$ en $S_1$ es isomorfa a $f^{-1}(y)$, por lo que $X_1\cap f_1^{-1}(y_1)$ es finito, discreto y tiene $n-1$ puntos. Como $f_1|X_1$ es localmente de tipo finito y separado, aplicamos la hipótesis de inducción a este morfismo: hay un morfismo étale $Y'\to Y_1$, un punto $y'\in Y'$ por encima de $y_1$ tal que $k(y')=k(y_1)$, y, si $X'=X\times_Y Y'$ y $p:X'\to S_1$ es la proyección canónica, se tiene una descomposición de $p^{-1}(X_1)$ en suma de dos subpreesquemas $U'$ y $X'_2$ inducidos sobre abiertos de $X'$ tales que $X'_2\cap f'^{-1}(y')=\varnothing$ y $U'$ sea finito sobre $Y'$. Por otra parte, $V'=p^{-1}(V_1)$ es finito sobre $Y'$ y $X'$ es suma de $U'$, $V'$ y $X'_2$; responderemos, pues, a la cuestión tomando $X'_1$ como suma de $U'$ y $V'$.

El corolario siguiente mejora (8.11.1):

\begin{corollary}[18.12.4]
\label{IV.18.12.4}
Para que un morfismo $f:X\to Y$ sea finito, es necesario y suficiente que sea separado, universalmente cerrado, localmente de tipo finito y que, para todo $y\in Y$, $f^{-1}(y)$ sea un espacio finito discreto; en particular, un morfismo propio y cuasi-finito es finito.
\end{corollary}

En efecto, se puede aplicar el corolario (18.12.3) a un punto cualquiera $y$ de $Y$. Con las notaciones de este corolario, $f'$ es un morfismo cerrado, por lo que, puesto que $X'_2$ es cerrado en $X'$, $f'(X'_2)$ es cerrado en $Y'$ y no contiene a $y'$; existe por consiguiente un entorno abierto $U'$ de $y'$ en $Y'$ tal que $U'\to Y$ sea de tipo finito y $f'^{-1}(U')=X\times_Y U'$ sea finito sobre $U'$. Sea $U$ la imagen de $U'$ en $Y$, que es un abierto de $Y$ en virtud de (II, 3.1), puesto que el morfismo $Y'\to Y$ es étale, por tanto plano y localmente de presentación finita; se tiene evidentemente todavía $X\times_Y U'=X\times_U U'=f^{-1}(U)\times_U U'$. Puesto que ahora el morfismo $U'\to U$ es fielmente plano y cuasi-compacto, se deduce de (2.7.1, (xv)) que el morfismo $f^{-1}(U)\to U$, restricción de $f$, es finito, lo que demuestra el corolario.

\begin{remark}[18.12.5]
\label{IV.18.12.5}
Obsérvese que en la demostración de (18.12.4) no se utilizó la hipótesis de que $f$ sea universalmente cerrado bajo su forma general, sino solamente que $f_{(Y')}$ sea un morfismo cerrado para todo morfismo étale $Y'\to Y$.
\end{remark}

\begin{corollary}[18.12.6]
\label{IV.18.12.6}
Sea $f:X\to Y$ un morfismo de preesquemas. Las condiciones siguientes son equivalentes:
\begin{enumerate}
  \item[$a)$] $f$ es una inmersión cerrada.
  \item[$b)$] $f$ es un monomorfismo propio.
  \item[$c)$] $f$ es propio y, para todo $y\in Y$, el $k(y)$-preesquema $X_y=f^{-1}(y)$ es radicial y geométricamente reducido sobre $k(y)$ (es decir, vacío o $k(y)$-isomorfo a $\Spec(k(y))$).
\end{enumerate}
\end{corollary}

Esto se deduce de (18.12.4) del mismo modo que (8.11.5) se deduce de (8.11.1).

\begin{proposition}[18.12.7]
\label{IV.18.12.7}
Sean $f:X\to Y$ un morfismo localmente de tipo finito, $y$ un punto de $Y$. Para que exista un entorno abierto $U$ de $y$ tal que la restricción $f^{-1}(U)\to U$ de $f$ sea una inmersión cerrada, es necesario y suficiente que $X_y$ sea un $k(y)$-preesquema radicial y geométricamente reducido sobre $k(y)$ (es decir, vacío o $k(y)$-isomorfo a $\Spec(k(y))$) y que exista un entorno abierto $V$ de $y$ tal que la restricción $f^{-1}(V)\to V$ de $f$ sea un morfismo universalmente cerrado.
\end{proposition}

Las condiciones son evidentemente necesarias; probemos que son suficientes. Podemos restringirnos al caso en que $X_y\neq\varnothing$. En virtud de la segunda condición, ya podemos suponer que $f$ es universalmente cerrado. Mostremos además que podemos suponer que $f$ es afín, por tanto separado; esto se deduce del lema siguiente:

\begin{lemma}[18.12.7.1]
\label{IV.18.12.7.1}
Sean $f:X\to Y$ un morfismo cerrado, $y\in Y$ un punto tal que todo entorno de $X_y$ contenga un entorno afín de $X_y$ (condición siempre satisfecha cuando $X_y$ es vacío o se reduce a un solo punto). Entonces existe un entorno abierto afín $U$ de $y$ tal que la restricción $f^{-1}(U)\to U$ de $f$ sea un morfismo afín.
\end{lemma}

En efecto, sea $U_0$ un entorno abierto afín de $y$ en $Y$ y sea $V$ un entorno abierto afín de $X_y$ contenido en $f^{-1}(U_0)$. Puesto que $f$ es cerrado, existe un entorno abierto afín $U\subset U_0$ de $y$ tal que $f^{-1}(U)\subset V$. Como la restricción $g:V\to U_0$ de $f$ es un morfismo afín, lo mismo ocurre con su restricción $f^{-1}(U)\to U$ (II, 1.2.5).

Supongamos, pues, $f$ afín y universalmente cerrado, y mostremos que existe un entorno abierto $U$ de $y$ tal que la restricción $f^{-1}(U)\to U$ sea un morfismo finito; como podemos suponer $Y$ y $X$ afines, podemos suponer $f$ de tipo finito, por tanto propio, y basta demostrar que existe un entorno $U$ de $y$ tal que $f^{-1}(U)\to U$ sea un morfismo cuasi-finito (18.12.4). Pero existe un entorno $V$ del conjunto $X_y$ (reducido a un solo punto) tal que $f|V$ sea cuasi-finito (13.1.4), y como $f$ es cerrado existe un entorno $U$ de $y$ tal que $f^{-1}(U)\subset V$, lo que concluye la demostración, teniendo en cuenta (18.12.6).

\begin{proposition}[18.12.8]
\label{IV.18.12.8}
Sea $f:X\to Y$ un morfismo de preesquemas. Para que $f$ sea entero, es necesario y suficiente que $f$ sea afín y universalmente cerrado.
\end{proposition}

Las condiciones son necesarias (II, 6.1.10); probemos que son suficientes. Podemos suponer $Y=\Spec(A)$ afín, por tanto también $X=\Spec(B)$, y hay que demostrar que todo elemento $b\in B$ es entero sobre $A$ (II, 6.1.1). Sea $B'$ la sub-$A$-álgebra de $B$ generada por $b$, que es una $A$-álgebra de tipo finito; pongamos $X'=\Spec(B')$, de modo que $f:X\to Y$ se factoriza como
\[
  X\xrightarrow{g}X'\xrightarrow{h}Y,
\]
donde $h$ es de tipo finito, y $g$ es dominante puesto que el homomorfismo $B'\to B$ es inyectivo (I, 1.2.7). Como $h$ es separado y $h\circ g$ universalmente cerrado, $g$ es universalmente cerrado (II, 5.4.3 y

\oldpage[IV-4]{183}

5.4.9), por tanto sobreyectivo puesto que es dominante; se concluye (II, 5.4.3 y 5.4.9) que $h$ también es universalmente cerrado, por tanto propio puesto que es de tipo finito y separado. Pero entonces, para todo $y\in Y$, el morfismo $h^{-1}(y)\to\Spec(k(y))$ inducido por $h$ es propio y afín, por tanto finito (III, 4.4.2), es decir, $h$ es cuasi-finito; se deduce, pues, de (18.12.4) que $h$ es finito, lo que prueba que $B'$ es una $A$-álgebra finita y, en consecuencia, que $b$ es entero sobre $A$. C.Q.F.D.

\begin{remark}[18.12.9]
\label{IV.18.12.9}
Puede ocurrir que un morfismo $f:X\to Y$ sea entero cuando se lo supone solamente separado, universalmente cerrado y tal que para todo $y\in Y$, el morfismo $f^{-1}(y)\to\Spec(k(y))$ inducido por $f$ sea entero. Para ello habría que demostrar que estas condiciones implican que $f$ es afín, o equivalentemente que toda fibra $X_y$ está contenida en un entorno abierto afín.
\end{remark}

\begin{corollary}[18.12.10]
\label{IV.18.12.10}
Un morfismo $f:X\to Y$ que es inyectivo y universalmente cerrado es entero.
\end{corollary}

Basta, por (18.12.8), demostrar que $f$ es afín, lo que se deduce del lema (18.12.7.1) y de la hipótesis.

\begin{corollary}[18.12.11]
\label{IV.18.12.11}
Sea $f:X\to Y$ un morfismo de preesquemas (resp. un morfismo localmente de tipo finito). Para que $f$ sea un homeomorfismo universal (2.4.2), es necesario y suficiente que $f$ sea entero (resp. finito), radicial y sobreyectivo.
\end{corollary}

Se sabe que las condiciones son suficientes (2.4.5); son necesarias en virtud de (18.12.10) y (18.12.4).

La proposición siguiente mejora (8.11.2):

\begin{proposition}[18.12.12]
\label{IV.18.12.12}
Todo morfismo $f:X\to Y$ que sea cuasi-finito y separado es cuasi-afín.
\end{proposition}

Pongamos $\mathscr A=f_*(\sh{O}_X)$, que es una $\sh{O}_Y$-álgebra cuasi-coherente puesto que $f$ es cuasi-compacto y separado (1.7.4); sea $Z=\Spec(\mathscr A)$ (II, 1.3.1), de modo que $f$ se factoriza canónicamente como
\[
  X\xrightarrow{g}Z\xrightarrow{h}Y,
\]
siendo $h$ afín y correspondiendo $g$ al automorfismo identidad de $\mathscr A$ (II, 1.2.7); bastará demostrar que $g$ es una inmersión abierta, o, lo que equivale a lo mismo (17.9.1), que $g$ es étale y radicial. Evidentemente bastará demostrar que, para todo $y\in Y$, el morfismo $g$ es étale y radicial en cada punto de $f^{-1}(y)$. Ahora bien, para todo $y\in Y$, podemos aplicar a $f$ el resultado de (18.12.3), conservando sus notaciones; puesto que el morfismo $Y'\to Y$ es plano, y que $f$ es cuasi-compacto y separado, se tiene, salvo un isomorfismo canónico,
\[
  f'_*(\sh{O}_{X'})=\mathscr A\otimes_{\sh{O}_Y}\sh{O}_{Y'}=\mathscr A'
\]
(2.3.1), por tanto $Z'=Z\times_Y Y'$ se identifica con $\Spec(f'_*(\sh{O}_{X'}))$, y en la factorización canónica
\[
  X'\xrightarrow{g'}Z'\xrightarrow{h'}Y'
\]
de $f'$ (II, 1.2.7), se tiene $h'=h_{(Y')}$ y $g'=g_{(Y')}$. Dicho esto, la descomposición de $X'$ en suma de dos subpreesquemas $X'_1$, $X'_2$ induce la descomposición de $\mathscr A'$ en producto directo de las dos $\sh{O}_{Y'}$-álgebras cuasi-coherentes $\mathscr A'_1$, $\mathscr A'_2$, imágenes directas respectivas de $\sh{O}_{X'_1}$ y $\sh{O}_{X'_2}$, de modo que $Z'$ se identifica con la suma $Z'_1\amalg Z'_2$, donde $Z'_i=\Spec(\mathscr A'_i)$ y $g'(X'_i)\subset Z'_i$ para $i=1,2$. Como $X'_1$ es finito sobre $Y'$, $g'_1=g'|X'_1$ es un isomorfismo de $X'_1$ sobre $Z'_1$, puesto que $f'|X'_1$ es afín; como $X'_2\cap f'^{-1}(y')=\varnothing$, se ve que $g'$ es étale y radicial en cada punto de $f'^{-1}(y')$. Puesto que el morfismo $Y'\to Y$ es plano y localmente de presentación finita, se deduce primero de (17.7.4) que $g$ es étale en todos los puntos de $X$ que son proyecciones de puntos de $f'^{-1}(y')$, es decir, en todos los puntos de $f^{-1}(y)$ (I, 3.5.2). Por otra parte, el morfismo
\[
  g'_{y'}:f'^{-1}(y')\longrightarrow h'^{-1}(y')
\]
inducido por $g'$ es radicial; como $k(y')=k(y)$, el morfismo $g_y:f^{-1}(y)\to h^{-1}(y)$ también es radicial, es decir, $g$ es radicial en cada punto de $f^{-1}(y)$, lo que concluye la demostración de la proposición.

El enunciado siguiente mejora igualmente (8.12.6):

\begin{corollary}[18.12.13]
\label{IV.18.12.13}
("Main Theorem" de Zariski). Sean $Y$ un preesquema cuasi-compacto y cuasi-separado, $f:X\to Y$ un morfismo cuasi-finito y separado. Entonces existe una factorización de $f$
\[
  X\xrightarrow{g}Z\xrightarrow{u}Y
\]
donde $g$ es una inmersión abierta (necesariamente cuasi-compacta) y $u$ un morfismo finito.
\end{corollary}

En efecto, se deduce de (18.12.12) que $f$ es cuasi-afín, por tanto cuasi-proyectivo. Se puede entonces aplicar (8.12.8), donde la hipótesis de existencia de un $\sh{O}_Y$-Módulo amplio puede de hecho reemplazarse por la sola hipótesis de que $Y$ es cuasi-compacto y cuasi-separado: en efecto, se deduce de (8.12.3) que la existencia de la factorización anunciada en (8.12.8) es una propiedad local sobre $Y$, y basta, pues, demostrarla cuando $Y$ es afín.

\begin{remark}[18.12.14]
\label{IV.18.12.14}
Se puede dar de (18.12.13) una demostración análoga a la de (18.12.12) que no utiliza (8.12.8) (pero utiliza (18.12.3), por tanto (18.5.11), que a su vez se sirve del «Main Theorem» en su forma local (8.12.9)). Conservemos, en efecto, las notaciones de la demostración de (18.12.12), y sea $\mathscr B$ la $\sh{O}_Y$-álgebra cuasi-coherente, clausura entera de $\sh{O}_Y$ en $\mathscr A$ (II, 6.3.2). Si ponemos $T=\Spec(\mathscr B)$, basta, en virtud de (8.12.3), demostrar que el $Y$-morfismo $g:X\to T$ correspondiente a la inyección canónica $\mathscr B\to\mathscr A$ es una inmersión, y para ello basta demostrar (17.9.1) que $g$ es étale y radicial. Con las notaciones de (18.12.12), podemos suponer que $Y=\Spec(C)$ y $Y'=\Spec(C')$ son afines, siendo $C'$ una $C$-álgebra étale (17.3.2); por tanto $\mathscr A=\widetilde A$, donde $A$ es una $C$-álgebra, $\mathscr A'=\widetilde{A'}$ con $A'=A\otimes_C C'$, y $\mathscr B=\widetilde B$, donde $B$ es la clausura entera de $C$ en $A$. El álgebra $A'$ es isomorfa al producto $A'_1\times A'_2$, donde $A'_1$ es una $C'$-álgebra finita. Bastará demostrar que $B'=B\otimes_C C'$, que

\oldpage[IV-4]{184}

se identifica por planitud con una sub-$C'$-álgebra de $A'$ y contiene $A'_1$; en efecto, $B'$ se descompondrá entonces en producto $A'_1\times A''_2$, donde $A''_2$ es una sub-$C'$-álgebra de $A'_2$, y si ponemos $T''_2=\Spec(A''_2)$, $T'=T\times_Y Y'$ será suma de $Z'_1=\Spec(A'_1)$ y de $T''_2$, y $g'|X'_1$ será un isomorfismo de $X'_1$ sobre $Z'_1$, lo que permitirá concluir como en (18.12.12).

Para demostrar que $B'$ contiene $A'_1$, basta evidentemente demostrar la proposición siguiente, que extiende parcialmente (6.14.4) cuando ya no se hacen hipótesis noetherianas:
\end{remark}

\begin{proposition}[18.12.15]
\label{IV.18.12.15}
Sean $C$ un anillo, $C'$ una $C$-álgebra étale, $A$ una $C$-álgebra, $B$ la clausura entera de $C$ en $A$; pongamos $A'=A\otimes_C C'$, $B'=B\otimes_C C'$, identificándose $B'$ con una subálgebra de $A'$; entonces $B'$ es la clausura entera de $C'$ en $A'$.
\end{proposition}

Considerando la familia filtrante creciente de las sub-$C$-álgebras de tipo finito de $A$ y razonando como en (6.14.4, II), podemos suponer primero que $A$ es una $C$-álgebra de tipo finito, por tanto de la forma $E/\mathfrak J$, donde $E=C[T_1,\ldots,T_r]$ y $\mathfrak J$ es un ideal de $E$; entonces se tiene $A'=E'/\mathfrak JE'$, donde $E'=C'[T_1,\ldots,T_r]$. Sea $(\mathfrak J_\lambda)$ la familia filtrante creciente de los ideales de tipo finito de $E$ contenidos en $\mathfrak J$, de modo que $A$ es límite inductivo de los $E/\mathfrak J_\lambda$; si $B_\lambda$ es la clausura entera de $C$ en $E/\mathfrak J_\lambda$, $B$ es límite inductivo de los $B_\lambda$, como muestra el razonamiento de (5.13.4). Del mismo modo, la clausura entera de $C'$ en $E'/\mathfrak JE'$ es límite inductivo de las clausuras enteras de $C'$ en $E'/\mathfrak J_\lambda E'$; si se demuestra que esta última es igual a $B_\lambda\otimes_C C'$, se deduce que
\[
  B'=\varinjlim_\lambda(B_\lambda\otimes_C C')
\]
será la clausura entera de $C'$ en $A'$. Nos reducimos así a demostrar la proposición cuando $A$ es una $C$-álgebra de presentación finita.

Mostremos a continuación que podemos reducirnos al caso en que $C$ es noetheriano. En efecto, $C$ es unión filtrante de sus subálgebras de tipo finito sobre $\mathbf Z$, que denotamos $C_\alpha$, por lo que se deduce de (17.7.8) que existe un índice $\alpha$ y una $C_\alpha$-álgebra étale $C'_\alpha$ tales que $C'=C'_\alpha\otimes_{C_\alpha}C$; además, $C'$ es límite inductivo de los $C'_\beta=C'_\alpha\otimes_{C_\alpha}C_\beta$ para $\beta\geq\alpha$. Podemos además suponer, puesto que $A$ es una $C$-álgebra de presentación finita, que $A=A_\alpha\otimes_{C_\alpha}C$, donde $A_\alpha$ es una $C_\alpha$-álgebra de tipo finito, y $A$ es límite inductivo de los $A_\beta=A_\alpha\otimes_{C_\alpha}C_\beta$ para $\beta\geq\alpha$ (8.9.1). El razonamiento de (5.13.4) muestra entonces que $B$ es el límite inductivo de las clausuras enteras $B_\beta$ de $C_\beta$ en $A_\beta$. Del mismo modo, $A'$ es límite inductivo de los
\[
  A'_\beta=A_\beta\otimes_{C_\beta}C'_\beta=A_\alpha\otimes_{C_\alpha}C'_\beta
\]
para $\beta\geq\alpha$, por lo que el mismo razonamiento que antes muestra que bastará demostrar que $B'_\alpha=B_\alpha\otimes_{C_\alpha}C'_\alpha$ es la clausura entera de $C'_\alpha$ en $A'_\alpha$.

Una vez reducido al caso en que $C$ es noetheriano, la proposición es un caso particular de (6.14.4). Sin embargo, hay que observar que, cuando se supone que $C'$ es étale sobre el anillo noetheriano $C$, la demostración de (6.14.4) ya no exige el delicado teorema (6.14.1). En efecto, las reducciones sucesivas de la demostración de (6.14.4) conducen, sin utilizar (6.14.1), al caso en que $C$ es íntegro y $A$ es el cuerpo de fracciones de $C$. Como entonces $B$ es un anillo normal y el morfismo $\Spec(B')\to\Spec(B)$ es étale, se deduce de (17.5.7) que $B'$ es un anillo normal; el razonamiento termina recurriendo solamente al lema elemental (6.14.1.1), pero no a la parte difícil de la demostración de (6.14.1).

\begin{proposition}[18.12.16]
\label{IV.18.12.16}
Sean $g:Y\to S$ un morfismo cuasi-compacto, $f:X\to Y$ un morfismo separado y cuasi-finito. Para todo $\sh{O}_Y$-Módulo invertible $\mathscr L$, amplio relativamente a $S$ (II, 4.6.1), el $\sh{O}_X$-Módulo $f^*(\mathscr L)$ es amplio relativamente a $S$.
\end{proposition}

En efecto, en virtud de (18.12.12), el morfismo $f$ es cuasi-afín, por tanto (II, 5.1.6) el $\sh{O}_X$-Módulo $\sh{O}_X$ es amplio para $f$. Se deduce de (II, 4.6.13, (ii)) que
\[
  \sh{O}_X\otimes_{\sh{O}_X}f^*(\mathscr L^{\otimes n})=f^*(\mathscr L^{\otimes n})
\]
es amplio relativamente a $S$ para un $n$ suficientemente grande. Pero como $f^*(\mathscr L^{\otimes n})=(f^*(\mathscr L))^{\otimes n}$, esto implica que $f^*(\mathscr L)$ es amplio relativamente a $S$ (II, 4.6.9, (i)).

\begin{corollary}[18.12.17]
\label{IV.18.12.17}
Sean $Z$ un esquema afín, $h:X\to Z$ un morfismo de tipo finito, $\mathscr L$ un $\sh{O}_X$-Módulo invertible. Con las notaciones de (II, 4.5.2), las condiciones de (II, 4.5.2) (que equivalen al hecho de que $\mathscr L$ es amplio) son también equivalentes a cada una de las siguientes:
\begin{enumerate}
  \item[$b'')$] $h$ es separado, se tiene $G(\varepsilon)=X$ y el morfismo canónico $u:G(\varepsilon)\to\operatorname{Proj}(S)$ tiene fibras finitas y discretas.
  \item[$b''')$] Se tiene $G(\varepsilon)=X$ y el morfismo canónico $u$ es radicial.
\end{enumerate}
\end{corollary}

Si $Z=\Spec(A)$, el anillo $S$ está provisto canónicamente de una estructura de $A$-álgebra graduada, por lo que se tiene un morfismo estructural $g:\operatorname{Proj}(S)\to Z$ que es separado (II, 2.4.2), y se tiene $h=g\circ u$. Como todo morfismo radicial es separado (I, 1.8.7.1), $b''')$ implica que $h$ es separado, por lo que implica $b'')$. Como la condición $b)$ de (II, 4.5.2) implica evidentemente $b''')$, queda por ver que $b'')$ implica que $\mathscr L$ es amplio. Puesto que $h=g\circ u$ es de tipo finito por hipótesis y que $g$ es separado, $u$ es también de tipo finito (I, 6.3.4), por lo que la hipótesis $b'')$ implica que $u$ es cuasi-finito y separado (I, 5.5.1). Para demostrar que $\mathscr L$ es amplio, aplicaremos (18.12.16). Pongamos $Y=\operatorname{Proj}(S)$. Como $X$ es cuasi-compacto, existe un entero $n>0$ y un número finito de elementos $f_j\in S_n$ tales que las imágenes recíprocas $u^{-1}(D_+(f_j))$ recubren $X$; podemos, pues, considerar $u$ como un morfismo cuasi-finito y separado de $X$ en el abierto
\[
  Y'=\bigcup_j D_+(f_j)
\]
de $Y$. Consideremos entonces el $\sh{O}_{Y'}$-Módulo invertible
\[
  \mathscr L'=\sh{O}_Y(n)|Y'
\]
(II, 2.5.8); es muy amplio relativamente a $Z$ (II, 4.4.3), y $u^*(\mathscr L')$ no es otra cosa, por definición, que $\mathscr L^{\otimes n}$

\oldpage[IV-4]{185}

(II, 3.7.9, (ii)). En virtud de (18.12.16), $\mathscr L^{\otimes n}$ es, pues, amplio relativamente a $Z$ (lo que equivale a decir que es amplio (II, 4.6.6)), y por consiguiente lo mismo ocurre con $\mathscr L$ (II, 4.5.6, (i)).

Dejamos al lector el cuidado de enunciar del mismo modo condiciones equivalentes a las de (II, 4.6.3) para los $\sh{O}_X$-Módulos relativamente amplios.


% Mechanical subsection components; prose is translated in the component files.
\section{Inmersiones regulares y planitud normal}
\label{section:IV.19}

Esta sección está dedicada, por una parte, al estudio de las inmersiones regulares (16.9.2) $Y\to X$ de preesquemas, especialmente cuando $X$ e $Y$ son planos sobre un mismo preesquema $S$; por otra parte, ofrece complementos sobre las sucesiones $M$-regulares y la planitud normal, generalizando en particular resultados de planitud establecidos por H. Hironaka en su teoría de la resolución de singularidades [35].

\subsection{Propiedades de las inmersiones regulares}
\label{subsection:IV.19.1}

\begin{proposition}[19.1.1]
\label{IV.19.1.1}
Sean $X$ un preesquema localmente noetheriano, $j:Y\to X$ una inmersión e $y$ un punto regular de $Y$. Para que $j$ sea una inmersión regular en el punto $y$, es necesario y suficiente que $X$ sea regular en el punto $j(y)$.
\end{proposition}

Teniendo en cuenta (16.9.10), basta demostrar el

\begin{corollary}[19.1.2]
\label{IV.19.1.2}
Sean $A$ un anillo local noetheriano y $\mathfrak J$ un ideal de $A$; supongamos que el anillo cociente $B=A/\mathfrak J$ sea regular. Para que $A$ sea regular, es necesario y suficiente que el ideal $\mathfrak J$ sea regular.
\end{corollary}

En efecto, si $A$ es regular, se sabe ($0$, 17.1.9) que $\mathfrak J$ está engendrado por una parte de un sistema regular de parámetros de $A$, por lo que $\mathfrak J$ es regular. Recíprocamente, si $\mathfrak J$ es regular y $(x_i)_{1\leq i\leq r}$ es una sucesión $A$-regular que engendra $\mathfrak J$, entonces $(x_i)$ forma parte de un sistema de parámetros de $A$ ($0$, 16.4.1), y se deduce de ($0$, 17.1.7) que $A$ es regular.

\begin{definition}[19.1.3]
\label{IV.19.1.3}
Sean $X$ un preesquema, $Y$ un subpreesquema de $X$ y $U$ un abierto de $X$ que contiene a $Y$ y tal que $Y$ esté definido por un ideal $\mathscr J$ de $\sh{O}_U$; supongamos que $\mathscr J/\mathscr J^2$ sea un $(\sh{O}_U/\mathscr J)$-módulo localmente libre (como ocurre si $Y$ está inmerso cuasirregularmente en $X$). Para todo $y\in Y$, se llama \emph{codimensión transversal de $Y$ en $X$ en el punto $y$}, y se denota por $\operatorname{codim}_y^*(Y,X)$, al rango del $(\sh{O}_{U,y}/\mathscr J_y)$-módulo libre $\mathscr J_y/\mathscr J_y^2$. Si $f:Z\to X$ es una inmersión de imagen $Y$, se llama asimismo \emph{codimensión transversal de $f$ en el punto $z\in Z$} a la codimensión transversal en $X$, en el punto $f(z)$, del subpreesquema $Y$.
\end{definition}

\begin{proposition}[19.1.4]
\label{IV.19.1.4}
Sean $X$ un preesquema localmente noetheriano e $Y$ un subpreesquema de $X$ inmerso regularmente; entonces, para todo $y\in Y$, se tiene
\begin{equation}
  \operatorname{codim}_y^*(Y,X)=\operatorname{codim}_y(Y,X).
  \tag{19.1.4.1}\label{IV.19.1.4.1}
\end{equation}
\end{proposition}

En virtud de (5.1.3.2), $\operatorname{codim}_y(Y,X)$ es igual al menor valor de las dimensiones de los anillos locales $A=\sh{O}_{X,z}$, donde $z$ es un punto genérico de una componente

\oldpage[IV-4]{186}

irreducible de $Y$ que contiene a $y$; como tal punto $z$ pertenece a toda vecindad de $y$ en $X$, $\mathscr J_z/\mathscr J_z^2$ es un $(A/\mathscr J_z)$-módulo libre de rango $n=\operatorname{codim}_y^*(Y,X)$; basta, pues, comprobar que $\dim(A)=n$. Ahora bien, puesto que $z$ es un punto maximal del subpreesquema $Y$ definido por $\mathscr J$, se tiene $\dim(\sh{O}_{X,z}/\mathscr J_z)=0$, de modo que $\mathscr J_z$ es un \emph{ideal de definición} del anillo local noetheriano $A$; además, la hipótesis de que $\mathscr J$ es cuasirregular implica que $\mathscr J_z^m/\mathscr J_z^{m+1}$ es un $(A/\mathscr J_z)$-módulo libre de rango $\binom{m+n-1}{n-1}$; si $r$ es la longitud del anillo artiniano $A/\mathscr J_z$, la longitud de $\mathscr J_z^m/\mathscr J_z^{m+1}$ es, por tanto, $r\binom{m+n-1}{n-1}$; el polinomio de Hilbert--Samuel de $A$ para la filtración $\mathscr J_z$-preádica es entonces de grado $n$, lo que demuestra la proposición ($0$, 16.2.3).

En virtud de (19.1.4), en lo sucesivo se dirá «codimensión» en lugar de «codimensión transversal» cuando se trate de un subpreesquema $Y$ inmerso regularmente en $X$ (\emph{incluso cuando $X$ no sea localmente noetheriano}).

\begin{proposition}[19.1.5]
\label{IV.19.1.5}
\begin{enumerate}
  \item[(i)] Para que una inmersión $j:Y\to X$ sea abierta, es necesario y suficiente que sea regular y de codimensión $0$ en todo punto.
  \item[(ii)] Sean $f:Y\to X$ una inmersión regular (resp. cuasirregular), $g:X'\to X$ un morfismo plano y $Y'=Y\times_X X'$; entonces el morfismo $f'=f_{(X')}:Y'\to X'$ es una inmersión regular (resp. cuasirregular), y para todo $y'\in Y'$, si $y$ es la proyección de $y'$ en $Y$, la codimensión de $f'$ en el punto $y'$ es igual a la codimensión de $f$ en el punto $y$. Recíprocamente, si $f$ es una inmersión y $g:X'\to X$ un morfismo fielmente plano y cuasicompacto, entonces, si $f'$ es una inmersión cuasirregular, también lo es $f$.
  \item[(iii)] Si $f:Y\to X$ y $g:Z\to Y$ son inmersiones regulares, entonces $f\circ g:Z\to X$ es una inmersión regular, y para todo $z\in Z$, la codimensión de $f\circ g$ en el punto $z$ es la suma de la codimensión de $g$ en el punto $z$ y de la codimensión de $f$ en el punto $g(z)$. Además, la sucesión de homomorfismos canónicos (16.2.7.1)
\begin{equation}
  0\longrightarrow g^*(\sh{N}_{Y/X})\longrightarrow\sh{N}_{Z/X}\longrightarrow\sh{N}_{Z/Y}\longrightarrow0
  \tag{19.1.5.1}\label{IV.19.1.5.1}
\end{equation}
es exacta, y para todo $z\in Z$ existe una vecindad de $z$ en $Z$ en la cual las restricciones de los homomorfismos de esta sucesión forman una sucesión escindida.
  \item[(iv)] Sean $X$ un preesquema localmente noetheriano, $f:Y\to X$ y $g:Z\to Y$ dos inmersiones, y $z$ un punto de $Z$. Las condiciones siguientes son equivalentes:
  \begin{enumerate}
    \item[a)] $g$ es una inmersión regular en el punto $z$ y $f$ una inmersión regular en el punto $g(z)$.
    \item[b)] $g$ y $f\circ g$ son inmersiones regulares en el punto $z$.
    \item[c)] $f\circ g$ es una inmersión regular en el punto $z$, $f$ es una inmersión regular en el punto $g(z)$, y la sucesión de homomorfismos canónicos (19.1.5.1), restringida a una vecindad abierta suficientemente pequeña de $z$ en $Z$, es exacta y escindida.
  \end{enumerate}
  \item[(v)] Sean $X$ un preesquema localmente noetheriano, $f:Y\to X$ y $g:Z\to Y$ dos inmersiones, y $z$ un punto de $Z$; supongamos que $y=g(z)$ sea un punto regular de $Y$ y que $x=f(g(z))$ sea un punto regular

\oldpage[IV-4]{187}

de $X$ (lo que implica, por (19.1.1), que $f$ es una inmersión regular en el punto $g(z)$); entonces, para que $f\circ g$ sea una inmersión regular en el punto $z$, es necesario y suficiente que $g$ sea una inmersión regular en el punto $z$.
\end{enumerate}
\end{proposition}

La afirmación (i) es otra manera de expresar (16.1.10). Para demostrar la primera afirmación de (ii), podemos limitarnos al caso en que $Y$ es un subpreesquema cerrado de $X$ definido por un ideal regular (resp. cuasirregular) $\mathscr J$; entonces $Y'$ es el subpreesquema cerrado de $X'$ definido por el ideal $g^*(\mathscr J)\sh{O}_{X'}$ (I, 4.4.5), que se identifica aquí con $g^*(\mathscr J)$ puesto que $g$ es plano; reemplazando $X$ por un abierto afín suficientemente pequeño, podemos suponer que $\mathscr J$ admite una sucesión regular (resp. cuasirregular) de generadores $f_i$ (secciones de $\sh{O}_X$ sobre $X$). Si la sucesión $(f_i)$ es regular, también lo es la sucesión $(f_i\otimes1)$ (que engendra $g^*(\mathscr J)\sh{O}_{X'}$), en virtud de ($0$, 15.2.5); el resultado análogo para las sucesiones cuasirregulares se deduce inmediatamente de la planitud de $g$ y del criterio (16.9.4). Del mismo modo, si $g$ es fielmente plano y cuasicompacto, y si $g^*(\mathscr J)$ es cuasirregular, $\mathscr J$ es de tipo finito en virtud de (2.5.2); por planitud, se tiene $g^*(\mathscr J)/g^*(\mathscr J)^2=g^*(\mathscr J/\mathscr J^2)$, por lo que, dado que $g^*(\mathscr J/\mathscr J^2)$ es localmente libre por hipótesis (16.9.4), también lo es $\mathscr J/\mathscr J^2$ (2.5.2). Por último, la condición (iii) de (16.9.4) relativa a $g^*(\mathscr J)$ implica la misma condición para $\mathscr J$ en virtud de la planitud fiel de $g$ (2.2.7). Se concluye, por tanto, de (16.9.4) que $\mathscr J$ es cuasirregular.

Para demostrar (iii), podemos suponer asimismo que $Y$ y $Z$ son subpreesquemas cerrados de $X$, definidos respectivamente por ideales $\mathscr J$, $\mathscr K$ de $\sh{O}_X$ tales que $\mathscr J\subset\mathscr K$; además, podemos suponer que existen secciones $f_i$ ($1\leq i\leq m$) de $\mathscr J$ sobre $X$ que engendran $\mathscr J$ y tales que, para $1\leq i\leq m$, el endomorfismo de $\sh{O}_X/(\sum_{j=1}^{i-1}f_j\sh{O}_X)$ definido por $f_i$ es inyectivo, y secciones $\overline f_i$ ($m+1\leq i\leq n$) de $\mathscr K/\mathscr J$ sobre $X$ que engendran $\mathscr K/\mathscr J$ y tales que, para $m+1\leq i\leq n$, el endomorfismo de
\[
  (\sh{O}_X/\mathscr J)\big/\Bigl(\sum_{j=m+1}^{i-1}\overline f_j(\sh{O}_X/\mathscr J)\Bigr)
\]
definido por $\overline f_i$ es inyectivo. Para todo $z\in Z$, sea $U$ una vecindad abierta de $z$ en $X$ tal que exista una sección $f_i$ de $\mathscr K$ sobre $U$, tal que $\overline f_i|U$ sea su clase en $\Gamma(U,\mathscr K/\mathscr J)$ para $m+1\leq i\leq n$; entonces
\[
  (\sh{O}_X/\mathscr J)\big/\Bigl(\sum_{j=m+1}^{i-1}\overline f_j(\sh{O}_X/\mathscr J)\Bigr),
\]
restringido a $U$ es isomorfo a
\[
  \sh{O}_U\big/\Bigl((\mathscr J|U)+\sum_{j=m+1}^{i-1}f_j\sh{O}_U\Bigr)
  =\sh{O}_U\big/\Bigl(\sum_{j=1}^{i-1}f_j\sh{O}_U\Bigr),
\]
y se ve, por tanto, que la sucesión $(f_i)_{1\leq i\leq n}$ es regular en $\sh{O}_U$; como engendra $\mathscr K|U$, esto demuestra la primera afirmación de (iii). Para demostrar la última afirmación de (iii), podemos limitarnos (con las mismas notaciones) al caso en que las imágenes $t_i$ de los $f_i$ en $\mathscr K/\mathscr K^2$ (para $1\leq i\leq n$) formen una base de este $(\sh{O}_X/\mathscr K)$-módulo (16.9.5); por la misma razón, podemos suponer que los $t_i$ tales que $1\leq i\leq m$ son las imágenes canónicas de los elementos de una base del $(\sh{O}_X/\mathscr K)$-módulo $g^*(\mathscr J/\mathscr J^2)$, y que las imágenes canónicas $\overline t_i$ de los $t_i$ en $(\mathscr K/\mathscr J)/(\mathscr K/\mathscr J)^2$, para $m+1\leq i\leq n$, forman una base de este $(\sh{O}_X/\mathscr K)$-módulo. Esto termina la demostración de (iii).

Pasemos a la demostración de (iv). El hecho de que $a)$ implica $c)$ resulta de (iii). Demostremos que $c)$ implica $a)$. Pongamos de nuevo $A=\sh{O}_{X,f(g(z))}$, de modo que $\sh{O}_{Y,g(z)}=A/\mathfrak J$,

\oldpage[IV-4]{188}

$\sh{O}_{Z,z}=A/\mathfrak R$ con $\mathfrak J\subset\mathfrak R$. Por hipótesis, $\mathfrak J$ y $\mathfrak R$ son ideales regulares de $A$ y se tiene una sucesión exacta escindida
\[
  0\longrightarrow\mathfrak J/\mathfrak J\mathfrak R\longrightarrow\mathfrak R/\mathfrak R^2\longrightarrow(\mathfrak R/\mathfrak J)/(\mathfrak R/\mathfrak J)^2\longrightarrow0
\]
(cf. (16.2.7)). Esto implica que existen elementos $f_j$ ($1\leq j\leq r+s$) de $\mathfrak R$ cuyas imágenes $f'_j$ en $\mathfrak R/\mathfrak R^2$ forman una base de este $(A/\mathfrak R)$-módulo, y tales además que los $f_j$ de índice $j\leq r$ pertenecen a $\mathfrak J$ y sus imágenes en $\mathfrak J/\mathfrak J\mathfrak R$ forman una base de este $(A/\mathfrak R)$-módulo. Los $f_j$ para $1\leq j\leq r+s$ engendran, pues, $\mathfrak R$, y los $f_j$ para $1\leq j\leq r$ engendran $\mathfrak J$ porque $A$ es noetheriano (Bourbaki, \emph{Alg. comm.}, cap. II, \textsection~3, n\textsuperscript{o}~2, prop. 5); como $\mathfrak R$ es un ideal regular, se deduce de (16.9.5) que la sucesión $(f_j)_{1\leq j\leq r+s}$ es regular. Por definición, las imágenes de $f_{r+1},\ldots,f_{r+s}$ en $\mathfrak R/\mathfrak J$ forman, por tanto, una sucesión regular en este $(A/\mathfrak J)$-módulo, lo que termina de demostrar que $c)$ implica $a)$ en (iv).

Por último, teniendo en cuenta (16.9.10), es inmediato que (v), así como la equivalencia de $a)$ y $b)$ en (iv), resultan del

\begin{corollary}[19.1.6]
\label{IV.19.1.6}
Sean $A$ un anillo local noetheriano, $\mathfrak J$, $\mathfrak R$ dos ideales de $A$ contenidos en su ideal maximal $\mathfrak m$ y tales que $\mathfrak J\subset\mathfrak R$, $A'=A/\mathfrak J$, $\mathfrak R'=\mathfrak R/\mathfrak J$. Entonces:
\begin{enumerate}
  \item[(i)] Si $\mathfrak J$ y $\mathfrak R'$ son ideales regulares, también lo es $\mathfrak R$.
  \item[(ii)] Si $\mathfrak R$ y $\mathfrak R'$ son ideales regulares, también lo es $\mathfrak J$.
  \item[(iii)] Supongamos que los anillos $A$ y $A'=A/\mathfrak J$ son regulares. Para que el ideal $\mathfrak R$ sea regular, es necesario y suficiente que $\mathfrak R'$ lo sea.
\end{enumerate}
\end{corollary}

La afirmación (i) es un caso particular de (19.1.5, (iii)). Para demostrar (ii), consideremos el homomorfismo sobreyectivo canónico $u:\mathfrak R/\mathfrak R^2\to\mathfrak R'/\mathfrak R'^2=\mathfrak R/(\mathfrak R^2+\mathfrak J)$. Como, por hipótesis, $\mathfrak R/\mathfrak R^2$ y $\mathfrak R'/\mathfrak R'^2$ son $(A/\mathfrak R)$-módulos libres, cuyos rangos respectivos denotaremos por $p+q$ y $q$, el núcleo de $u$ es un $(A/\mathfrak R)$-módulo \emph{proyectivo} (Bourbaki, \emph{Alg.}, cap. II, 3\textsuperscript{a} ed., \textsection~2, n\textsuperscript{o}~2, prop. 4), y por tanto \emph{libre} de rango $p$, puesto que $A/\mathfrak R$ es un anillo local noetheriano ($0_{\mathrm{III}}$, 10.1.3); en otras palabras, existe una base de $\mathfrak R/\mathfrak R^2$ cuyos $p$ primeros elementos forman una base de $(\mathfrak R^2+\mathfrak J)/\mathfrak R^2$. Esto significa también (en virtud del lema de Nakayama) que existe un sistema de generadores $(f_i)_{1\leq i\leq p+q}$ de $\mathfrak R$, que forma una sucesión regular en $A$, tal que la sucesión $(f_i)_{1\leq i\leq p}$ está formada por elementos de $\mathfrak J$; se trata de demostrar que esta última sucesión engendra $\mathfrak J$. Para ello, denotemos por $\mathfrak J'\subset\mathfrak J$ el ideal que engendra; considerando el anillo $A/\mathfrak J'$, se ve que nos reducimos a demostrar el lema siguiente:

\begin{lemma}[19.1.6.1]
\label{IV.19.1.6.1}
Sean $B$ un anillo local noetheriano de ideal maximal $\mathfrak n$, $(g_j)_{1\leq j\leq q}$ una sucesión regular en $B$ de elementos de $\mathfrak n$, $\mathfrak b$ el ideal que engendra, y $\mathfrak c\subset\mathfrak b$ un ideal tal que las imágenes canónicas de los $g_j$ en $B/\mathfrak c$ sigan formando una sucesión regular en este anillo. Entonces $\mathfrak c=0$.
\end{lemma}

Como el lema es trivial para $q=0$, razonemos por inducción sobre $q$. La hipótesis de inducción, aplicada al anillo $B/g_1B$ y a las imágenes canónicas de los $g_j$ ($2\leq j\leq q$)

\oldpage[IV-4]{189}

y de $\mathfrak c$ en este anillo cociente, muestra que $\mathfrak c\subset g_1B$. Sea $\mathfrak d$ el ideal de $B$ formado por los $x\in B$ tales que $g_1x\in\mathfrak c$; se tiene entonces $\mathfrak c=g_1\mathfrak d$; pero, como por hipótesis $g_1$ es regular en $B/\mathfrak c$, se tiene necesariamente $\mathfrak d=\mathfrak c$, y por tanto $\mathfrak c=g_1\mathfrak c$. Como $\mathfrak c$ es de tipo finito y $g_1$ está contenido en el radical de $B$, el lema de Nakayama muestra que $\mathfrak c=0$.

Demostremos por último la afirmación (iii) de (19.1.6). En virtud de (i) y de (19.1.2), basta ver que, si $\mathfrak R$ es regular, también lo es $\mathfrak R'$. Observemos que $\mathfrak J$ es regular (19.1.2), y razonemos por inducción sobre el rango $q$ del $(A/\mathfrak J)$-módulo libre $\mathfrak J/\mathfrak J^2$. Si $q=0$, se tiene $\mathfrak J=0$ por el lema de Nakayama y la afirmación es trivial. Como $A$ y $A/\mathfrak J$ son regulares, se sabe ($0$, 17.1.9) que $\mathfrak J$ está engendrado por una parte de $q$ elementos de un sistema regular de parámetros de $A$; por tanto, si $f$ es uno de los elementos de este sistema de generadores de $\mathfrak J$, $A_1=A/fA$ es regular y $f\notin\mathfrak m^2$ ($0$, 17.1.8). Sean $\mathfrak J_1$, $\mathfrak R_1$ las imágenes canónicas de $\mathfrak J$, $\mathfrak R$ en $A_1$; se tiene $\mathfrak J_1\subset\mathfrak R_1$, $A_1/\mathfrak J_1=A/\mathfrak J$, luego $A_1/\mathfrak J_1$ es regular y, por consiguiente, $\mathfrak J_1$ es un ideal regular (19.1.2). Demostremos que $\mathfrak R_1$ es un ideal regular en $A_1$; como $\mathfrak R$ es un ideal regular en $A$, basta demostrar que $f$ forma parte de un sistema regular de generadores de $\mathfrak R$, y para ello (16.9.5) basta demostrar que la imagen de $f$ en $\mathfrak R/(\mathfrak R^2+\mathfrak m\mathfrak R)=\mathfrak R/\mathfrak m\mathfrak R$ forma parte de una base de este $(A/\mathfrak m)$-espacio vectorial; en otras palabras, que $f\notin\mathfrak m\mathfrak R$, lo que se deduce, en efecto, de $f\notin\mathfrak m^2$. Entonces, como $\mathfrak J_1/\mathfrak J_1^2$ es un $(A_1/\mathfrak J_1)$-módulo libre de rango $q-1$, la hipótesis de inducción muestra que $\mathfrak R'_1=\mathfrak R_1/\mathfrak J_1$ es un ideal regular de $A_1/\mathfrak J_1$; y como $\mathfrak R'$ se identifica con $\mathfrak R'_1$ en el isomorfismo canónico entre $A/\mathfrak J$ y $A_1/\mathfrak J_1$, $\mathfrak R'$ es regular. Q.E.D.

\begin{remarks}[19.1.7]
\label{IV.19.1.7}
\begin{enumerate}
  \item[(i)] Desconocemos si la compuesta de dos inmersiones cuasirregulares es cuasirregular.
  \item[(ii)] En un espacio anillado en anillos locales $X$, para todo ideal $\mathscr J$ de $\sh{O}_X$, $\sh{O}_X/\mathscr J$ es de tipo finito y tiene, por tanto, soporte cerrado en $X$ ($0_{\mathrm I}$, 5.2.2). Se pueden definir entonces, como en (I, 4.1.3 y 4.2.1), las nociones de subespacio anillado (localmente cerrado) de $X$ definido por un ideal $\mathscr J$ de un haz $\sh{O}_U$, donde $U$ es un abierto de $X$, y la noción de inmersión. Las definiciones de inmersiones cuasirregulares y regulares y de codimensión transversal se aplican entonces sin modificación, y los resultados de (19.1.5, (i) a (iv))\footnote{El texto impreso remite a (19.1.4), que no contiene incisos; la referencia debe ser a la Proposición (19.1.5), inmediatamente anterior.} siguen siendo válidos, suponiendo en la segunda afirmación de (ii) y en (iv) que el haz $\sh{O}_X$ sea coherente y los anillos locales $\sh{O}_{X,x}$ noetherianos; en cuanto a (ii), hay que \emph{definir} $Y'$ como el subespacio anillado definido por el $\sh{O}_{X'}$-ideal $g^*(\mathscr J)$; las demostraciones permanecen sin cambio.
\end{enumerate}
\end{remarks}

\begin{env}[19.1.8]
\label{IV.19.1.8}
Ya hemos señalado (en particular en (16.9.6, (ii))) que, en los anillos no noetherianos, la noción de «sucesión regular» de elementos del anillo no posee las propiedades que la harían utilizable. Investigaciones recientes parecen mostrar que la noción que en este caso debe sustituir a la de sucesión regular es la de una sucesión $\mathbf f=(f_1,\ldots,f_n)$ con la propiedad de que el complejo del álgebra exterior $\mathrm K_0(\mathbf f,M)$ (III, 1.1.3) tenga \emph{homología nula en grados $>0$} (cf. (III, 1.1.4)); véase en particular A. Grothendieck, \emph{Séminaire de géométrie algébrique}, 1966, de próxima aparición.
\end{env}

\subsection{Inmersiones transversalmente regulares}
\label{subsection:IV.19.2}

\oldpage[IV-4]{190}

\begin{definition}[19.2.1]
\label{IV.19.2.1}
Sean $u:X\to S$ un morfismo de preesquemas y $\mathscr F$ un $\sh{O}_X$-módulo cuasicoherente. Se dice que una sucesión $(f_i)_{1\leq i\leq n}$ de secciones de $\sh{O}_X$ sobre $X$ es \emph{transversalmente $\mathscr F$-regular relativamente a $S$} (o relativamente a $u$) si la sucesión $(f_i)$ es $\mathscr F$-regular y los $\sh{O}_X$-módulos
\[
  \mathscr F\big/\Bigl(\sum_{j\leq i}f_j\mathscr F\Bigr)
\]
sean planos sobre $S$ para $0\leq i\leq n$. Se dice que un ideal cuasicoherente $\mathscr J$ de $\sh{O}_X$ es \emph{transversalmente regular relativamente a $S$} (o relativamente a $u$) en un punto $x\in\operatorname{Supp}(\sh{O}_X/\mathscr J)$ si existe una vecindad abierta $U$ de $x$ y una sucesión finita $(f_i)$ de secciones de $\mathscr J$ sobre $U$ que es transversalmente $\sh{O}_U$-regular relativamente a $S$ y engendra $\mathscr J|U$.

En una $A$-álgebra $B$, se dice que un ideal $\mathfrak J$ es \emph{transversalmente regular relativamente a $A$} si $\widetilde{\mathfrak J}$ es transversalmente regular relativamente a $S=\Spec(A)$ en todo punto de $V(\mathfrak J)$ en $\Spec(B)$.
\end{definition}

\begin{definition}[19.2.2]
\label{IV.19.2.2}
Sean $S$ un preesquema, $X$, $Y$ dos $S$-preesquemas, y $u:Y\to X$ un $S$-morfismo que es una inmersión. Se dice que $u$ es una \emph{inmersión transversalmente regular relativamente a $S$ en un punto $y\in Y$} si existe una vecindad $V$ de $u(y)$ en $X$ tal que el subpreesquema inducido sobre $u(Y)\cap V$ por el subpreesquema de $X$ asociado a $u$ esté definido por un ideal de $\sh{O}_V$ transversalmente regular relativamente a $S$ en el punto $u(y)$.

Es claro que el conjunto de puntos de $Y$ en los que una inmersión es transversalmente regular relativamente a $S$ es \emph{abierto} en $Y$; si es igual a $Y$, se dice simplemente que $u$ es una \emph{inmersión transversalmente regular relativamente a $S$}.

Si $S\to T$ es un morfismo \emph{plano}, entonces, si $\mathscr J$ es un ideal de $\sh{O}_X$ transversalmente regular relativamente a $S$ en $x\in X$, también lo es relativamente a $T$ ($0_{\mathrm I}$, 6.2.1). Una $S$-inmersión $Y\to X$ transversalmente regular en un punto $y\in Y$ relativamente a $S$ es, por tanto, también una $T$-inmersión transversalmente regular relativamente a $T$ en ese punto.
\end{definition}

\begin{proposition}[19.2.3]
\label{IV.19.2.3}
Sean $S$ un preesquema, $X$, $Y$ dos $S$-preesquemas, y $u:Y\to X$ un $S$-morfismo que es una inmersión transversalmente regular relativamente a $S$ en un punto $y\in Y$. Entonces, para todo cambio de base $S'\to S$, si se pone $X'=X\times_S S'$, $Y'=Y\times_S S'$, la inmersión $u'=u_{(S')}:Y'\to X'$ es transversalmente regular relativamente a $S'$ en todo punto $y'\in Y'$ situado sobre $y$.
\end{proposition}

Esto resulta inmediatamente de la definición y de ($0$, 15.1.15).

\begin{proposition}[19.2.4]
\label{IV.19.2.4}
Sean $S$ un preesquema, $X$, $Y$ dos $S$-preesquemas, y $u:Y\to X$ un $S$-morfismo que es una inmersión. Se supone, o bien que $S$ y $X$ son localmente noetherianos, o bien que los morfismos estructurales $g:X\to S$, $h:Y\to S$ son localmente de presentación finita. Sean $y$ un punto de $Y$ y $s$ su imagen en $S$. Las condiciones siguientes son equivalentes:
\begin{enumerate}
  \item[a)] $u$ es transversalmente regular relativamente a $S$ en el punto $y$.
  \item[b)] Los morfismos $g$ y $h$ son planos en vecindades de los puntos $u(y)$ e $y$, respectivamente, y, si se pone $X_s=g^{-1}(s)$, $Y_s=h^{-1}(s)$, la inmersión $u_s:Y_s\to X_s$ es regular en el punto $y$.
  \item[b$'$)] Los morfismos $g$ y $h$ son planos en vecindades de los puntos $u(y)$ e $y$, respectivamente, y, para todo cambio de base $S'\to S$, si se pone $X'=X\times_S S'$, $Y'=Y\times_S S'$, la inmersión $u'=u_{(S')}:Y'\to X'$ es regular en todo punto de $Y'$ situado sobre $y$.
  \item[c)] El morfismo $h$ es plano en una vecindad del punto $y$, y la inmersión $u$ es regular en una vecindad del punto $y$.

\oldpage[IV-4]{191}

  \item[c$'$)] El morfismo $h$ es plano en una vecindad del punto $y$, y la inmersión $u$ es cuasirregular en una vecindad del punto $y$.
\end{enumerate}
\end{proposition}

Ya se ha demostrado en (19.2.3) que $a)$ implica $b')$ sin hipótesis de finitud, y es trivial que $b')$ implica $b)$. Demostremos que $b)$ implica $c)$. Para ello, como la cuestión es local sobre $X$ y sobre $S$, podemos suponer que $Y$ es un subpreesquema cerrado de $X$ definido por un ideal cuasicoherente $\mathscr J$ de $\sh{O}_X$, y que $u$ es la inyección canónica. La hipótesis de que $\sh{O}_X/\mathscr J$ es $g$-plano implica que la sucesión
\[
  0\longrightarrow\mathscr J\otimes_{\sh{O}_S}\mathbf k(s)\longrightarrow\sh{O}_X\otimes_{\sh{O}_S}\mathbf k(s)\longrightarrow(\sh{O}_X/\mathscr J)\otimes_{\sh{O}_S}\mathbf k(s)\longrightarrow0
\]
es exacta ($0_{\mathrm I}$, 6.1.2); por tanto, $Y_s$ es el subpreesquema cerrado de $X_s$ definido por el ideal $\mathscr J_s$ de $\sh{O}_{X_s}=\sh{O}_X\otimes_{\sh{O}_S}\mathbf k(s)$, que se identifica con $\mathscr J\otimes_{\sh{O}_S}\mathbf k(s)=\mathscr J/\mathfrak m_s\mathscr J$. Consideremos una sucesión finita $(f_i)$ de secciones de $\mathscr J$ sobre una vecindad de $y$ en $X$, cuyas imágenes en
\[
  \mathscr J_y/(\mathfrak m_s\sh{O}_{X,y}\mathscr J_y+\mathscr J_y^2)
\]
forman una base de este $(\sh{O}_{X,y}/(\mathfrak m_s\sh{O}_{X,y}+\mathscr J_y))$-módulo (que es libre en virtud de la hipótesis $b)$). Puesto que $X_s$ es localmente noetheriano, se deduce de (16.9.11), (16.9.5) y de la hipótesis $b)$ que las imágenes $(f_i)_s=f_i\otimes1$, secciones de $\mathscr J_s=\mathscr J\otimes_{\sh{O}_S}\mathbf k(s)$ sobre una vecindad de $y$ en $X_s$, forman una sucesión regular que engendra la restricción de $\mathscr J_s$ a dicha vecindad. Como $\mathfrak m_s\sh{O}_{X,y}$ está contenido en el radical de $\sh{O}_{X,y}$, y como en los dos casos considerados $\mathscr J_y$ es un ideal de tipo finito de $\sh{O}_{X,y}$, se deduce del lema de Nakayama que los $(f_i)_y$ engendran también $\mathscr J_y$; como $\mathscr J$ es un ideal de tipo finito de $\sh{O}_X$ en ambos casos, existe una vecindad $U$ de $y$ en $X$ tal que los $f_i|U$ engendran $\mathscr J|U$ ($0_{\mathrm I}$, 5.2.2). El hecho de que $b)$ implica $c)$ es entonces consecuencia de ($0$, 15.1.16) cuando $S$ y $X$ son localmente noetherianos, y de (11.3.8) cuando $g$ y $h$ son localmente de presentación finita; en este último caso, (11.3.8) demuestra también la equivalencia de $c)$ y $c')$ (que es trivial cuando $X$ y $S$ son localmente noetherianos ($0$, 15.1.11)). Por último, si se verifica $c)$, para demostrar $a)$ podemos reducirnos de nuevo al caso en que $Y$ es un subpreesquema cerrado de $X$ definido por $\mathscr J$; por hipótesis, existe una sucesión regular $(f_i)$ de secciones de $\mathscr J$ sobre una vecindad de $y$ en $X$ tal que los $f_i$ engendran $\mathscr J|U$ y $h$ es plano en $U$; para ver que $c)$ implica $a)$, basta aplicar de nuevo ($0$, 15.1.16) cuando $S$ y $X$ son localmente noetherianos, y (11.3.8) cuando $g$ y $h$ son localmente de presentación finita.

\begin{remarks}[19.2.5]
\label{IV.19.2.5}
\begin{enumerate}
  \item[(i)] Si en la condición $b)$ se suprime la hipótesis de que $h$ es plano en una vecindad de $y$, la condición así modificada ya no implica $a)$. Basta tomar como $S$ el espectro de un anillo local que no sea un cuerpo, como $s$ el punto cerrado de $S$, y poner $Y=X_s$.
  \item[(ii)] Supongamos que $Y$ es un subpreesquema cerrado de $X$ definido por un ideal cuasicoherente $\mathscr J$ de $\sh{O}_X$. En el curso de la demostración de (19.2.4) se ha probado que, cuando se cumplen las condiciones equivalentes de esta proposición, para todo sistema $(f_i)$ de secciones de $\mathscr J$ sobre una vecindad de $y$ en $X$, tal que las imágenes de los $f_i$ en $\mathscr J_y/(\mathfrak m_s\sh{O}_{X,y}\mathscr J_y+\mathscr J_y^2)$ forman una base de este $(\sh{O}_{X,y}/(\mathfrak m_s\sh{O}_{X,y}+\mathscr J_y))$-módulo, existe una vecindad abierta $U$ de $y$ en $X$ tal que los $f_i|U$ satisfacen las condiciones de la definición (19.2.1).
\end{enumerate}
\end{remarks}

\oldpage[IV-4]{192}

\begin{corollary}[19.2.6]
\label{IV.19.2.6}
Supongamos, o bien que $S$ y $X$ son localmente noetherianos, o bien que $g$ y $h$ son localmente de presentación finita. Supongamos además que las fibras $X_s$ e $Y_s$ son preesquemas regulares en los puntos $u(y)$ e $y$, respectivamente, y que los morfismos $g$ y $h$ son planos en vecindades de $u(y)$ e $y$, respectivamente. Entonces la inmersión $Y\to X$ es transversalmente regular en el punto $y$.\footnote{El texto impreso dice «x» aquí y de nuevo en la demostración; el punto introducido en la Proposición (19.2.4) es y.}
\end{corollary}

En efecto, se deduce de (19.1.1) que la inmersión $Y_s\to X_s$ es regular en el punto $y$, y basta aplicar (19.2.4, $b)$).

\begin{proposition}[19.2.7]
\label{IV.19.2.7}
\begin{enumerate}
  \item[(i)] Toda $S$-inmersión abierta cuyo dominio es plano sobre $S$ es transversalmente regular relativamente a $S$.
  \item[(ii)] Sean $f:Y\to X$ una $S$-inmersión y $g:S'\to S$ un morfismo, y pongamos $X'=X_{(S')}$, $Y'=Y_{(S')}$, $f'=f_{(S')}:Y'\to X'$. Si $f$ es una inmersión transversalmente regular relativamente a $S$, entonces $f'$ es una inmersión transversalmente regular relativamente a $S'$. El recíproco es cierto si $g$ es fielmente plano y, además, $X$ e $Y$ son localmente de presentación finita sobre $S$.
  \item[(iii)] Si $f:Y\to X$ y $g:Z\to Y$ son inmersiones transversalmente regulares relativamente a $S$, también lo es $f\circ g:Z\to X$.
  \item[(iv)] Sean $X$ un $S$-preesquema y $f:Y\to X$, $g:Z\to Y$ dos $S$-inmersiones. Se supone, o bien que $S$ y $X$ son localmente noetherianos, o bien que $X$, $Y$ y $Z$ son localmente de presentación finita sobre $S$. Entonces, si $g$ y $f\circ g$ son transversalmente regulares en el punto $z\in Z$ relativamente a $S$, $f$ es transversalmente regular en el punto $g(z)$ relativamente a $S$.
  \item[(v)] Bajo las condiciones generales de (iv), supongamos que $X$ e $Y$ son planos sobre $S$, que $Y_s$ es regular en el punto $g(z)$ y que $X_s$ es regular en el punto $f(g(z))$; entonces, si $f\circ g$ es transversalmente regular en el punto $z$, también lo es $g$.
\end{enumerate}
\end{proposition}

La afirmación (i) es inmediata, y la primera afirmación de (ii) resulta inmediatamente de (19.2.3). Para demostrar la segunda afirmación de (ii), apliquemos (19.2.4): como $X'$ e $Y'$ son localmente de presentación finita sobre $S'$, este criterio muestra primero que $X'$ e $Y'$ son planos sobre $S'$, y por tanto que $X$ e $Y$ son planos sobre $S$ (2.5.1); por otra parte, para todo $s\in S$, si $s'\in S'$ es un punto situado sobre $s$, $Y'_{s'}\to X'_{s'}$ es por hipótesis una inmersión regular, y se deduce entonces de (19.1.5, (ii)) que lo mismo ocurre con $Y_s\to X_s$, puesto que estos preesquemas son localmente noetherianos y que, en tal caso, decir que la inmersión $Y_s\to X_s$ es regular equivale a decir que es cuasirregular. Para demostrar (iii), nos reducimos al caso en que $Z$ es un subpreesquema cerrado de $Y$ e $Y$ un subpreesquema cerrado de $X$, y formamos, como en (19.1.5, (iii)), un sistema de generadores del ideal de $\sh{O}_X$ que define $Z$. Para demostrar (iv), se observa que la hipótesis, junto con (19.2.4), muestra primero que $X$, $Y$ y $Z$ son planos sobre $S$ en una vecindad de $z$; además, si $s$ es la imagen de $z$ en $S$, basta demostrar que la inmersión $f_s:Y_s\to X_s$ es regular en el punto $g(z)$, sabiendo que ocurre lo mismo con $g_s:Z_s\to Y_s$ y $f_s\circ g_s$; pero esto resulta de (19.1.5, (iv)). Por último, para (v), se razona del mismo modo que para (iv); la hipótesis sobre $f\circ g$ implica ya, por (19.2.4, $b)$), que $Z$ es plano sobre $S$ en una vecindad de $z$, y lo mismo ocurre con $X$ e $Y$ por hipótesis; nos reducimos entonces a demostrar que $g_s$ es regular en el punto $z$, sabiendo que $f_s\circ g_s$ lo es, lo que se deduce de (19.1.5, (v)), dadas las hipótesis formuladas sobre $X_s$ e $Y_s$.

\oldpage[IV-4]{193}

\begin{proposition}[19.2.8]
\label{IV.19.2.8}
Sean $X$ un $S$-preesquema, $Y$ un subpreesquema cerrado de $X$ definido por un ideal cuasicoherente $\mathscr J$ de $\sh{O}_X$, y $u:Y\to X$ la inyección canónica. Si $u$ es transversalmente regular relativamente a $S$, los $\sh{O}_X$-módulos $\mathscr J^n$ y $\mathscr J^n/\mathscr J^m$ ($0\leq n<m$) son planos sobre $S$ en los puntos de $Y$. Para todo cambio de base $S'\to S$, si $X'=X_{(S')}$, $Y'=Y_{(S')}$ y $u'=u_{(S')}$, se tiene
\[
  \sh{G}r_n(u')=\sh{G}r_n(u)\otimes_{\sh{O}_S}\sh{O}_{S'}
\]
salvo isomorfismo para todo $n>0$.
\end{proposition}

En efecto, $\sh{O}_X$ y $\sh{O}_X/\mathscr J$ son, por hipótesis, planos sobre $S$ en los puntos de $Y$, y $\mathscr J^n/\mathscr J^{n+1}$ es un $(\sh{O}_X/\mathscr J)$-módulo localmente libre (16.9.3), luego es plano sobre $S$ en los puntos de $Y$. La primera afirmación resulta entonces de ($0_{\mathrm I}$, 6.1.2), y la segunda de (11.2.9.2).

La proposición siguiente generaliza y precisa (17.16.1):

\begin{proposition}[19.2.9]
\label{IV.19.2.9}
Sean $f:X\to S$ un morfismo, $x$ un punto de $X$ y $s=f(x)$. Supongamos que $f$ es de presentación finita en el punto $x$ y que es un morfismo de Cohen--Macaulay en el punto $x$ (6.8.1). Entonces existe un subpreesquema $Y$ de $X$ que contiene a $x$ y posee las propiedades siguientes:
\begin{enumerate}
  \item[(i)] La inyección canónica $j:Y\to X$ es transversalmente regular relativamente a $S$.
  \item[(ii)] El morfismo $g=f\circ j:Y\to S$ es de Cohen--Macaulay (lo que equivale a decir, teniendo en cuenta (i), que todas las fibras $g^{-1}(g(y))$ de $g$ son preesquemas de Cohen--Macaulay).
  \item[(iii)] El punto $x$ es un punto maximal de $Y_s=g^{-1}(s)$.
\end{enumerate}
\end{proposition}

Pongamos $X_s=f^{-1}(s)$; por hipótesis, el anillo $\sh{O}_{X_s,x}$ es de Cohen--Macaulay; sea $(t_i)_{1\leq i\leq n}$ un sistema de parámetros de este anillo, que es por tanto una sucesión regular ($0$, 16.5.7). Existe una vecindad abierta $U$ de $x$ en $X$ y secciones $h_i$ ($1\leq i\leq n$) de $\sh{O}_X$ sobre $U$ tales que los $t_i$ son las imágenes canónicas de los $(h_i)_x$ en el anillo cociente $\sh{O}_{X_s,x}$ de $\sh{O}_{X,x}$. Sea $\mathscr J$ el ideal de $\sh{O}_U$ engendrado por los $h_i$, y tomemos como $Y$ el subpreesquema cerrado de $U$ definido por $\mathscr J$. Podemos suponer que $f|U$ es localmente de presentación finita, y entonces lo mismo ocurre con $g$ por la definición de $Y$; se deduce de la equivalencia de $c)$ y $a)$ en (11.3.8) que la sucesión $((h_i)_x)$ es regular en $\sh{O}_{X,x}$ y que $\sh{O}_{Y,x}$ es plano sobre $\sh{O}_{S,s}$; en virtud de (11.3.1), $f$ y $g$ son, por tanto, planos en una vecindad de $x$, y se deduce de (19.2.4) que $j$ es una inmersión transversalmente regular relativamente a $S$ (reemplazando, si es necesario, $Y$ por $Y\cap V$, donde $V$ es una vecindad abierta de $x$ en $X$). Así queda probado (i), y se ve que podemos suponer que $g:Y\to S$ es un morfismo plano. Como $\sh{O}_{Y_s,x}$, cociente de $\sh{O}_{X_s,x}$ por el ideal engendrado por los $t_i$, es artiniano por construcción ($0$, 16.3.6), $x$ es un punto maximal de $Y_s$. Por último, como el anillo $\sh{O}_{Y_s,x}$ es artiniano, es un anillo de Cohen--Macaulay ($0$, 16.5.1); por tanto, reemplazando si es necesario $Y$ por su intersección con una vecindad abierta de $x$ en $X$, podemos suponer que $g$ es un morfismo de Cohen--Macaulay, en virtud de (12.1.7, (iii)).

En particular, se recupera la existencia de «cuasisecciones» \emph{planas} $Y$ de un morfismo $f:X\to S$ en un punto $x\in X_s$ que es \emph{cerrado} en $X_s$, tal que $f$ es plano en el punto $x$ y $\sh{O}_{X_s,x}$ es un anillo de Cohen--Macaulay (17.16.1); se ve además que la inmersión $Y\to X$ es \emph{transversalmente regular} relativamente a $S$. Este último complemento vale también para la cuasisección étale definida en (17.16.3, (i)).

\subsection{Intersecciones completas relativas (caso plano).}
\label{subsection:IV.19.3}

\oldpage[IV-4]{194}

\begin{definition}[19.3.1]
\label{IV.19.3.1}
Sea $A$ un anillo local noetheriano. Se dice que $A$ es un \emph{anillo de intersección completa} (o también un \emph{anillo de intersección completa absoluta}, si hay riesgo de confusión) si la completación $\widehat A$ es isomorfa al cociente de un anillo local noetheriano completo regular $B$ por un ideal regular (es decir, (16.9.7), un ideal engendrado por una sucesión regular de elementos de $B$).

Todo anillo local regular es un anillo de intersección completa.

Se dice que un preesquema localmente noetheriano $X$ es una \emph{intersección completa} (absoluta) \emph{en el punto} $x\in X$ si el anillo $\sh{O}_{X,x}$ es un anillo de intersección completa.
\end{definition}

\begin{proposition}[19.3.2]
\label{IV.19.3.2}
Sean $B$ un anillo local noetheriano regular y $\mathfrak J$ un ideal de $B$. Para que $A=B/\mathfrak J$ sea un anillo de intersección completa, es necesario y suficiente que el ideal $\mathfrak J$ sea regular (es decir, (16.9.7), que esté engendrado por una sucesión regular de elementos de $B$).
\end{proposition}

Se puede escribir $\widehat A=\widehat B/\mathfrak J\widehat B$ y, puesto que $\widehat B$ es un $B$-módulo fielmente plano y $B$ es noetheriano, $\mathfrak J\widehat B$ es regular si y solo si $\mathfrak J$ lo es (19.1.5, (ii)). Esto muestra (en virtud de ($0$, 17.1.5)), por una parte, que la condición del enunciado es suficiente y, por otra, que para demostrar que es necesaria se puede restringir al caso en que $A$ y $B$ son completos. Por hipótesis, existe entonces un anillo local noetheriano completo regular $B'$ tal que $A$ es isomorfo a un cociente $B'/\mathfrak J'$, donde $\mathfrak J'$ es un ideal regular de $B'$. Consideremos entonces el \emph{producto fibrado} $B\times_A B'=B''$ para los homomorfismos sobreyectivos $B\to A$, $B'\to A$ ($0$, 18.1.2); demostraremos primero el lema siguiente:

\begin{lemma}[19.3.2.1]
\label{IV.19.3.2.1}
Sean $A$, $E$, $F$ tres anillos, $f:E\to A$, $g:F\to A$ dos homomorfismos, y pongamos $G=E\times_A F$ ($0$, 18.1.2).
\begin{enumerate}
  \item[(i)] Si $f$ es sobreyectivo y $E$ y $F$ son anillos noetherianos, lo mismo ocurre con $G$.
  \item[(ii)] Si $A$, $E$ y $F$ son anillos locales y $f$ y $g$ son homomorfismos locales, $G$ es un anillo local y los homomorfismos canónicos $G\to E$, $G\to F$ son locales.
  \item[(iii)] Si se verifican las condiciones de (i) y (ii), $g$ es sobreyectivo y $E$ y $F$ son anillos locales noetherianos completos, lo mismo ocurre con $G$.
\end{enumerate}
\end{lemma}

\begin{enumerate}
  \item[(i)] Teniendo en cuenta ($0$, 18.1.3 y 18.1.5), $G$ es un $A$-anillo aumentado sobre $F$, cuyo ideal de aumento $\mathfrak J'$ se identifica canónicamente con el núcleo $\mathfrak J$ de $f$, y la estructura de $G$-módulo sobre $\mathfrak J'$ se identifica con la de $E$-módulo sobre $\mathfrak J$ por medio del homomorfismo canónico $G\to E$. Si $\mathfrak a$ es un ideal cualquiera de $G$, $\mathfrak a/(\mathfrak a\cap\mathfrak J')$ es isomorfo a $(\mathfrak a+\mathfrak J')/\mathfrak J'$, y por tanto a un ideal de $F$, de modo que es de tipo finito; por otra parte, $\mathfrak a\cap\mathfrak J'$ es isomorfo a un ideal de $E$ contenido en $\mathfrak J$, y por tanto es también de tipo finito, lo que demuestra que $\mathfrak a$ es de tipo finito y, por consiguiente, que $G$ es noetheriano.
  \item[(ii)] Sean $\mathfrak r$, $\mathfrak s$ los ideales maximales de $E$ y $F$, respectivamente; el conjunto $\mathfrak m=\mathfrak r\times_A\mathfrak s$ es evidentemente un ideal de $G$, imagen inversa de $\mathfrak r$ por la proyección $G\to E$ y de $\mathfrak s$ por la proyección $G\to F$. Se deduce inmediatamente que los elementos de $\mathfrak m$ no son invertibles en $G$. Por otra parte, si $(x,y)\notin\mathfrak m$ y, por ejemplo, $x\notin\mathfrak r$, $f(x)$ no pertenece al ideal maximal de $A$; por tanto, puesto que $g(y)=f(x)$ y $g$ es local, se tiene también $y\notin\mathfrak s$; se

\oldpage[IV-4]{195}

concluye que $x$ e $y$ son invertibles, y por tanto también lo es $(x,y)$, lo que termina de demostrar (ii).
  \item[(iii)] La imagen de $\mathfrak m$ por el homomorfismo sobreyectivo $G\to E$ es el ideal maximal $\mathfrak r$ de $E$, por lo que la imagen de $\mathfrak m^n\cap\mathfrak J'$ es $\mathfrak r^n\cap\mathfrak J$; como $\mathfrak J$ es cerrado (y por tanto completo) para la topología inducida por la topología $\mathfrak r$-ádica de $E$ ($0_{\mathrm I}$, 7.3.5), $\mathfrak J'$ es completo para la topología inducida por la topología $\mathfrak m$-preádica de $G$. Por otra parte, $G/\mathfrak J'$, provisto de la topología $\mathfrak m$-preádica, es isomorfo a $F$ provisto de la topología $\mathfrak s$-ádica, y es, por tanto, completo para la topología cociente de la topología $\mathfrak m$-preádica de $G$. Se deduce inmediatamente que $G$ es completo para la topología $\mathfrak m$-preádica.
\end{enumerate}

Establecido este lema, del teorema de Cohen ($0$, 19.8.8) se sigue que $B''$ es isomorfo a un cociente de un anillo local noetheriano completo \emph{regular} $C$. Se deduce entonces de (19.1.2) que la inmersión $\Spec(B')\to\Spec(C)$ es regular; como, por hipótesis, lo mismo ocurre con la inmersión $\Spec(A)\to\Spec(B')$, se concluye de (19.1.5, (iii)) que la inmersión $\Spec(A)\to\Spec(C)$ es regular; pero esta inmersión se escribe también como la compuesta $\Spec(A)\to\Spec(B)\to\Spec(C)$. Ahora bien, como $B$ es regular por hipótesis, la inmersión $\Spec(B)\to\Spec(C)$ es regular (19.1.2); por tanto, también lo es $\Spec(A)\to\Spec(B)$ por (19.1.5, (iv)).

\begin{corollary}[19.3.3]
\label{IV.19.3.3}
Sea $X$ un preesquema localmente noetheriano localmente inmersible en un preesquema regular. Entonces el conjunto de los $x\in X$ en los que $X$ es una intersección completa es abierto en $X$ y contiene al conjunto de los puntos regulares de $X$.
\end{corollary}

En efecto, se puede restringir al caso en que $X=\Spec(B/\mathfrak J)$, donde $B$ es un anillo regular y $\mathfrak J$ un ideal de $B$. El conjunto de los $\mathfrak p\supset\mathfrak J$ de $\Spec(B)$ tales que $\mathfrak J_{\mathfrak p}$ es regular en $B_{\mathfrak p}$ es entonces abierto (16.9.5), y es el conjunto de los puntos de $\Spec(B)$ en los que $X$ es una intersección completa en virtud de (19.3.2).

\begin{corollary}[19.3.4]
\label{IV.19.3.4}
Sean $k$ un cuerpo, $X$ un preesquema localmente de tipo finito sobre $k$, $k'$ una extensión de $k$ y $X'=X\otimes_k k'$. Sean $x'$ un punto de $X'$ y $x$ su proyección en $X$. Para que $X$ sea una intersección completa en $x$, es necesario y suficiente que $X'$ sea una intersección completa en $x'$.
\end{corollary}

Como la cuestión es local sobre $X$, se puede suponer que $X=\Spec(A)$, donde $A$ es una $k$-álgebra de tipo finito y, por tanto, cociente de un álgebra de polinomios $k[T_1,\ldots,T_n]$, de modo que $X$ es un subpreesquema cerrado del esquema regular $Y=\Spec(k[T_1,\ldots,T_n])$ ($0$, 17.3.7). Si se pone $Y'=Y\otimes_k k'=\Spec(k'[T_1,\ldots,T_n])$, $Y'$ es también un esquema regular. Ahora bien, puesto que $\Spec(k')$ es fielmente plano sobre $\Spec(k)$, se deduce de (19.1.5, (ii)) y del hecho de que se trata de preesquemas noetherianos que la inmersión $X\to Y$ es regular en el punto $x$ si y solo si la inmersión $X'\to Y'$ lo es en el punto $x'$. La conclusión se sigue, pues, del criterio (19.3.2).

\begin{corollary}[19.3.5]
\label{IV.19.3.5}
Sean $A$, $C$ dos anillos locales noetherianos, $\rho:A\to C$ un homomorfismo local, $B=C/\mathfrak J$ un anillo cociente de $C$ y $k$ el cuerpo residual de $A$. Se supone que $\rho$ hace de $C$ un $A$-módulo plano y que el anillo $C\otimes_A k$ es regular. Entonces las condiciones siguientes son equivalentes:
\begin{enumerate}
  \item[a)] El ideal $\mathfrak J$ es transversalmente regular relativamente a $A$.
  \item[b)] $B$ es un $A$-módulo plano y $B\otimes_A k$ es un anillo de intersección completa.
\end{enumerate}
\end{corollary}

\oldpage[IV-4]{196}

En virtud de (19.2.4), la condición $a)$ equivale a decir que $B$ es un $A$-módulo plano y que $\mathfrak J\otimes_A k$ (que se identifica con un ideal de $C\otimes_A k$ en virtud de la planitud de $B$ sobre $A$ ($0_{\mathrm I}$, 6.1.2)) es regular en este anillo. Como el anillo $C\otimes_A k$ es regular, en virtud de (19.3.2) equivale a decir (cuando $B$ es un $A$-módulo plano) que $\mathfrak J\otimes_A k$ es un ideal regular de $C\otimes_A k$, o que $B\otimes_A k=(C\otimes_A k)/(\mathfrak J\otimes_A k)$ es un anillo de intersección completa; de donde se sigue el corolario.

\begin{definition}[19.3.6]
\label{IV.19.3.6}
Sea $f:X\to S$ un morfismo plano localmente de presentación finita. Se dice que $X$ es una \emph{intersección completa relativamente a $S$ en un punto} $x$ si la fibra $f^{-1}(f(x))$ es una intersección completa (absoluta) en el punto $x$. Se dice que $X$ es una \emph{intersección completa relativamente a $S$}, o que el morfismo $f$ es un \emph{morfismo de intersección completa}, si $X$ es una intersección completa relativamente a $S$ en cada uno de sus puntos.
\end{definition}

\begin{proposition}[19.3.7]
\label{IV.19.3.7}
Sean $g:X\to S$, $h:Y\to S$ dos morfismos planos localmente de presentación finita, $f:Y\to X$ una $S$-inmersión, $y$ un punto de $Y$, $x=f(y)$ y $s=g(x)=h(y)$. Se supone que la fibra $X_s=g^{-1}(s)$ es regular en el punto $x$. Entonces, para que $f$ sea una inmersión transversalmente regular en el punto $y$, es necesario y suficiente que $Y$ sea una intersección completa relativamente a $S$ en el punto $y$.
\end{proposition}

Decir que $f$ es transversalmente regular en el punto $y$ equivale, en virtud de (19.2.4), a decir que $f_s:Y_s\to X_s$ es una inmersión regular en el punto $y$. Pero, como $X_s$ es regular en el punto $x=f(y)$, esto equivale, en virtud de (19.3.2), a decir que $Y_s$ es una intersección completa (absoluta) en el punto $y$, de donde se sigue la proposición.

\begin{corollary}[19.3.8]
\label{IV.19.3.8}
Sea $f:X\to S$ un morfismo plano localmente de presentación finita. El conjunto $U$ de los $x\in X$ en los que $X$ es una intersección completa relativamente a $S$ es abierto en $X$. Si además $f$ es propio, el conjunto $E$ de los $s\in S$ tales que $X_s=f^{-1}(s)$ es una intersección completa en cada uno de sus puntos es una parte abierta de $S$.
\end{corollary}

Como $E=S-f(X-U)$,\footnote{El texto impreso dice $Y-f(X-U)$; el conjunto ambiente en el que acaba de definirse $E$ es $S$.} la segunda afirmación se sigue de la primera y del hecho de que $f$ es una aplicación cerrada. La primera afirmación es local sobre $X$, por lo que se puede suponer que $X$ es un subpreesquema cerrado de un preesquema de la forma $Z=S[T_1,\ldots,T_n]$ ($T_i$ indeterminadas). Como este último es liso sobre $S$ (17.3.8), sus fibras $Z_s$ son regulares y, por tanto, en virtud de (19.3.7), equivale a decir que $x\in U$ o que la inmersión $X\to Z$ es transversalmente regular en el punto $x$; se concluye, pues, el corolario por (19.2.2).

\begin{proposition}[19.3.9]
\label{IV.19.3.9}
\begin{enumerate}
  \item[(i)] Toda inmersión abierta es un morfismo de intersección completa.
  \item[(ii)] Sea $f:X\to S$ un morfismo plano y localmente de presentación finita que es un morfismo de intersección completa. Para todo cambio de base $g:S'\to S$, $f'=f_{(S')}:X_{(S')}\to S'$ es un morfismo de intersección completa. El recíproco es cierto si $g$ es fielmente plano y cuasicompacto.
  \item[(iii)] Si $f:X\to Y$ y $g:Y\to Z$ son dos morfismos planos localmente de presentación finita que son morfismos de intersección completa, lo mismo ocurre con $g\circ f:X\to Z$.
\end{enumerate}
\end{proposition}

La afirmación (i) se deduce inmediatamente de la definición (19.3.1). Para demostrar (ii), observemos primero que, si $f$ es plano y localmente de presentación finita, lo mismo ocurre

\oldpage[IV-4]{197}

con $f'$ (2.1.4), y el recíproco es cierto si $g$ es fielmente plano ((2.5.1) y (2.7.1)); además, para todo $s'\in S'$, si $s=g(s')$, equivale a decir que $f^{-1}(s)$ es una intersección completa o que $f'^{-1}(s')$ es una intersección completa (19.3.4); de aquí se sigue (ii), teniendo en cuenta que $g$ es sobreyectivo cuando es fielmente plano.

Por último, bajo las hipótesis de (iii), $g\circ f$ es plano y localmente de presentación finita. Por definición (19.3.6), nos reducimos entonces al caso en que $Z$ es el espectro de un cuerpo, y a demostrar que los anillos locales $\sh{O}_{X,x}$ son anillos de intersección completa, lo cual está contenido en el resultado más general siguiente:

\begin{corollary}[19.3.10]
\label{IV.19.3.10}
Sean $f:X\to S$ un morfismo plano localmente de presentación finita que es un morfismo de intersección completa, $x$ un punto de $X$ y $s=f(x)$. Si $\sh{O}_{S,s}$ es un anillo noetheriano de intersección completa, lo mismo ocurre con $\sh{O}_{X,x}$.
\end{corollary}

Evidentemente, se puede restringir al caso en que $S=\Spec(A)$ con $A=\sh{O}_{S,s}$. Demostremos, además, que se puede restringir al caso en que el anillo local $A$ es \emph{completo}. En efecto, pongamos $A'=\widehat A$, $X'=X\otimes_A A'$, y sea $s'$ el único punto cerrado de $S'=\Spec(A')$, que es el único punto de $S'$ sobre $s$; si $f'=f_{(S')}:X'\to S'$, la fibra $f'^{-1}(s')$ es canónicamente isomorfa a $f^{-1}(s)$, puesto que $A$ y $A'$ tienen el mismo cuerpo residual (I, 3.6.4); hay, pues, un único punto $x'\in X'$ sobre $x$ y $s'$, y por consiguiente se tiene $\sh{O}_{X',x'}=\sh{O}_{X,x}\otimes_A A'$. Se deduce que los anillos locales $\sh{O}_{X,x}$ y $\sh{O}_{X',x'}$ tienen \emph{completaciones isomorfas}, pues, en general, si $E$ es un anillo local que es una $A$-álgebra (siendo local el homomorfismo $A\to E$), la completación separada $(E\otimes_A\widehat A)^{\wedge}$ del anillo $E\otimes_A\widehat A$ provisto de la topología producto tensorial es isomorfa a la completación separada de $\widehat E\widehat\otimes_A\widehat A\cong\widehat E$, y por tanto a la completación separada $\widehat E$ de $E$ ($0_{\mathrm I}$, 7.7.1). En virtud de la definición (19.3.1), equivale, pues, a decir que $\sh{O}_{X,x}$ es un anillo de intersección completa o que $\sh{O}_{X',x'}$ es un anillo de intersección completa.

Supongamos, pues, que $A$ es \emph{completo}; se puede suponer además que $X=\Spec(B)$, siendo $B$ cociente de un álgebra de polinomios $C=A[T_1,\ldots,T_r]$. Como los anillos de polinomios sobre un cuerpo son anillos regulares ($0$, 17.3.7), se deduce de (19.3.7) que se puede suponer que la inmersión $X\to Z=\Spec(C)$ es transversalmente regular relativamente a $S$; por tanto, se puede suponer que $B=C/\mathfrak J$, donde $\mathfrak J$ está engendrado por una sucesión regular $(g_i)_{1\leq i\leq n}$ de elementos de $C$. Por otra parte, puesto que $A$ es completo, se puede, por hipótesis, escribirlo como $A'/\mathfrak R$, donde $A'$ es un anillo local regular y $\mathfrak R$ un ideal de $A'$ ($0$, 19.8.8), y la hipótesis de que $A$ es un anillo de intersección completa implica que $\mathfrak R$ está engendrado por una sucesión regular $(f_j)_{1\leq j\leq m}$ (19.3.2). En el anillo de polinomios $C'=A'[T_1,\ldots,T_r]$, los elementos $f_j$ ($1\leq j\leq m$) siguen formando una sucesión regular que engendra el ideal $\mathfrak R'=\mathfrak R[T_1,\ldots,T_r]$ ($0$, 15.1.4); como $C'/\mathfrak R'=C$, se ve que, si para cada $i$, $g'_i$ es un elemento de $C'$ de imagen $g_i$ en $C$, la sucesión formada por los $f_j$ ($1\leq j\leq m$) y los $g'_i$ ($1\leq i\leq n$) es regular en $C'$; si $\mathfrak J'$ es el ideal de $C'$ que engendra, se tiene $B=C'/\mathfrak J'$, de donde se sigue la conclusión en virtud de (19.3.2), puesto que $C'$ es un anillo regular ($0$, 17.3.7).

\subsection{Aplicación: criterios para que los preesquemas obtenidos por explosión sean regulares o lisos.}
\label{subsection:IV.19.4}

\oldpage[IV-4]{198}

\begin{env}[19.4.1]
\label{IV.19.4.1}
Sean $X$ un preesquema y $Y$ un subpreesquema cerrado de $X$, definido por un ideal cuasicoherente $\mathscr J$ de $\sh{O}_X$. Recordemos (II, 8.1.3) que el $X$-esquema $X'$ obtenido por explosión de $Y$ es el preesquema
\[
  X'=\operatorname{Proj}(\mathscr S),\qquad\text{donde}\qquad \mathscr S=\bigoplus_{n\geq0}\mathscr J^n.
\]
Si $\mathscr J$ es de tipo finito, el morfismo estructural $f:X'\to X$ es proyectivo. Sin ninguna hipótesis sobre $\mathscr J$, el subpreesquema cerrado
\[
  Y'=f^{-1}(Y)
\]
de $X'$ está definido por el ideal cuasicoherente $\mathscr J\sh{O}_{X'}$ de $\sh{O}_{X'}$, canónicamente isomorfo a $\sh{O}_{X'}(1)$; más precisamente, se tiene una sucesión exacta
\begin{equation}
  0\longrightarrow\sh{O}_{X'}(1)\xrightarrow{\widetilde u_0}\sh{O}_{X'}\longrightarrow\sh{O}_{Y'}\longrightarrow0.
  \tag{19.4.1.1}\label{IV.19.4.1.1}
\end{equation}
donde $\mathscr J\sh{O}_{X'}$ es la imagen de $\widetilde u_0$ (II, 8.1.7 y 8.1.8). Como en un entorno de un punto de $X'$ el $\sh{O}_{X'}$-módulo $\sh{O}_{X'}(1)$ es libre de rango $1$, $\mathscr J\sh{O}_{X'}$ está engendrado en tal entorno por una sección que no es divisor de cero\footnote{El texto impreso dice «section inversible»; la inyectividad de la sucesión exacta precedente da un elemento que no es divisor de cero, no una unidad.} de $\sh{O}_{X'}$ y, por consiguiente, $\mathscr J\sh{O}_{X'}/\mathscr J^2\sh{O}_{X'}$ es un $\sh{O}_{X'}/\mathscr J\sh{O}_{X'}$-módulo libre de rango $1$. Esto demuestra la primera afirmación del lema siguiente:
\end{env}

\begin{lemma}[19.4.2]
\label{IV.19.4.2}
La inmersión canónica $Y'\to X'$ es regular y de codimensión $1$, y $Y'$ es canónicamente isomorfo al $Y$-esquema $\operatorname{Proj}(\operatorname{gr}_{\mathscr J}^*(\sh{O}_X))$.
\end{lemma}

La segunda afirmación se deduce de (II, 3.5.3)\footnote{El texto impreso cita (II, 3.3.3); el resultado de cambio de base para Proj utilizado aquí es (II, 3.5.3).} aplicado a la inyección canónica $Y\to X$, teniendo en cuenta que $\mathscr S\otimes_{\sh{O}_X}\sh{O}_Y=\operatorname{gr}_{\mathscr J}^*(\sh{O}_X)$ por definición, puesto que $\sh{O}_Y$ es isomorfo a $\sh{O}_X/\mathscr J$ y $\mathscr J^n\otimes_{\sh{O}_X}(\sh{O}_X/\mathscr J)$ a $\mathscr J^n/\mathscr J^{n+1}$.

\begin{corollary}[19.4.3]
\label{IV.19.4.3}
Si la inmersión canónica $Y\to X$ es cuasirregular, la restricción $g:Y'\to Y$ del morfismo $f$ es lisa.
\end{corollary}

En efecto, $\sh{N}_{Y/X}=\mathscr J/\mathscr J^2$ es entonces un $\sh{O}_Y$-módulo localmente libre de rango finito y $\operatorname{gr}_{\mathscr J}^*(\sh{O}_X)$ es isomorfo a $\mathbf S_{\sh{O}_Y}^*(\sh{N}_{Y/X})$ (16.9.8); por tanto, $Y'$ es $Y$-isomorfo a $\mathbf P(\sh{N}_{Y/X})$ y, en consecuencia, es liso sobre $Y$ (17.3.9).

Veremos más adelante (cap. V) que $g$ sigue siendo liso cuando la inmersión $Y\to X$ «solo presenta singularidades ordinarias».

\begin{proposition}[19.4.4]
\label{IV.19.4.4}
Con las notaciones de (19.4.1), supongamos que $X$ es localmente noetheriano y que el morfismo $g:Y'\to Y$, restricción de $f$, es liso. Sean $x'$ un punto de $Y'$ y $x=g(x')$. Para que $Y$ sea regular en el punto $x$, es necesario y suficiente que $Y'$ sea regular en el punto $x'$; cuando así ocurre, $X'$ es entonces regular en el punto $x'$.
\end{proposition}

La primera afirmación se deduce de (17.5.8), y la segunda de (19.1.1) y (19.4.2).

\begin{corollary}[19.4.5]
\label{IV.19.4.5}
Bajo las hipótesis generales de (19.4.4), supongamos que $Y$ es regular en el punto $x$. Entonces, para toda generización $x_1$ de $x$ en $X$ que no pertenezca a $Y$, $X$ es regular en el punto $x_1$. Si $\operatorname{Reg}(X)$ es abierto ((6.12), por ejemplo, si $X$ es excelente (7.8.6, (iii)))

\oldpage[IV-4]{199}

y si $Y$ es regular, entonces existe un entorno abierto $U$ de $Y$ en $X$ tal que $X$ es regular en los puntos de $U-Y$.
\end{corollary}

Para demostrar la primera afirmación, se puede restringir al caso en que $X$ es un esquema local de punto cerrado $x$; la hipótesis implica entonces que $Y$ es regular ($0$, 17.3.2), y por (19.4.4) lo mismo ocurre con $Y'$. Observemos ahora que el morfismo $f:X'\to X$ es proyectivo, y por tanto propio; en consecuencia (II, 7.2.1), si se considera el único punto $x'_1$ de $X'$ sobre $x_1$ (II, 8.1.3), existe una especialización de $x'_1$ que pertenece a $Y'$; aplicándole (19.4.4) y utilizando ($0$, 17.3.2), se concluye que $X'$ es regular en el punto $x'_1$, y por tanto que $X$ es regular en el punto $x_1$, puesto que $X'-Y'$ es isomorfo a $X-Y$ (II, 8.1.3). Si $Y$ es regular, todo punto de $X-Y$ que sea una generización de un punto de $Y$ pertenece, pues, a $\operatorname{Reg}(X)$; si $\operatorname{Reg}(X)$ es abierto, $U=Y\cup\operatorname{Reg}(X)$ es entonces localmente constructible y estable por generización, y por tanto abierto ($0_{\mathrm{III}}$, 9.2.5), y responde a la cuestión.

\begin{proposition}[19.4.6]
\label{IV.19.4.6}
Sean $g:X\to S$ un morfismo, $Y$ un subpreesquema cerrado de $X$ definido por un ideal cuasicoherente $\mathscr J$, $X'$ el $X$-esquema obtenido por explosión de $Y$ y $Y'$ la imagen inversa de $Y$ en $X'$.
\begin{enumerate}
  \item[(i)] Las condiciones $a)$ y $b)$ siguientes son equivalentes. Implican $c)$; si $X$ es plano sobre $S$, la condición $c)$ también es equivalente a ellas:
  \begin{enumerate}
    \item[a)] $\operatorname{gr}_{\mathscr J}^*(\sh{O}_X)$ es una $\sh{O}_X$-álgebra $g$-plana.
    \item[b)] Para todo $n\geq0$, el $\sh{O}_X$-módulo $\sh{O}_X/\mathscr J^{n+1}$ es $g$-plano (en otras palabras, el $n$-ésimo entorno infinitesimal $Y^{(n)}$ de $Y$ en $X$ (16.1.2) es $S$-plano).
    \item[c)] Para todo cambio de base $S_1\to S$ y todo entero $n\geq1$, si se pone $X_1=X\times_S S_1$, el homomorfismo canónico $\mathscr J^n\otimes_{\sh{O}_S}\sh{O}_{S_1}\to\mathscr J^n\sh{O}_{X_1}$ es biyectivo.
  \end{enumerate}
  Cuando se satisface la condición $a)$, o la condición equivalente $b)$, $Y'$ es $S$-plano y, para todo cambio de base $S_1\to S$, si $Y_1$ es la imagen inversa de $Y$ en $X_1$, el preesquema $X'_1$ obtenido por explosión de $Y_1$ en $X_1$ es canónicamente isomorfo a $X'\times_S S_1$.
  \item[(ii)] Supongamos que se satisfacen las condiciones $a)$ y $b)$ de (i) y, además, que los morfismos $X\to S$ e $Y\to S$ son localmente de presentación finita. Entonces $\mathscr S=\bigoplus_{n\geq0}\mathscr J^n$ es una $\sh{O}_X$-álgebra de presentación finita, el morfismo $X'\to X$ es de presentación finita, la inmersión canónica $Y'\to X'$ es transversalmente regular de codimensión $1$ relativamente a $S$, y el conjunto de los puntos de $X$ (resp. $X'$) en los que $X$ (resp. $X'$) es $S$-plano es un entorno de $Y$ (resp. $Y'$).
\end{enumerate}
\end{proposition}

\begin{enumerate}
  \item[(i)] Podemos trabajar en el caso afín $S=\Spec(A)$, $X=\Spec(B)$, con $\mathscr J=\widetilde{\mathfrak J}$. Pongamos $Q_m=B/\mathfrak J^m$ para $m\geq1$ y $G_r=\mathfrak J^r/\mathfrak J^{r+1}$ para $r\geq0$. El álgebra graduada de $a)$ es plana si y solo si todos los $G_r$ son planos. Como $Q_1=G_0$, las sucesiones exactas $0\to G_r\to Q_{r+1}\to Q_r\to0$ para $r\geq1$ prueban la equivalencia de $a)$ y $b)$: una extensión de módulos planos es plana, y el núcleo de una aplicación sobreyectiva entre módulos planos es plano ($0_{\mathrm I}$, 6.1.2). La sucesión exacta de los Tor asociada a $0\to\mathfrak J^m\to B\to Q_m\to0$ muestra que $b)$ implica $c)$. Recíprocamente, supongamos que $B$ es plano sobre $A$ y que se satisface la condición $c)$. Para toda $A$-álgebra $A'$, la misma sucesión exacta y la inyectividad de $c)$ dan $\operatorname{Tor}_1^A(Q_m,A')=0$. Tomando el álgebra $A'=A\oplus M$ con ideal $M$ de cuadrado nulo, cuya multiplicación es $(a,m)(a',m')=(aa',am'+a'm)$, donde $M$ es un $A$-módulo arbitrario, se obtiene $\operatorname{Tor}_1^A(Q_m,M)=0$. Por tanto, todos los $Q_m$ son planos, lo que prueba la recíproca condicional. Las conclusiones que siguen se establecen bajo la condición $a)$, o bajo la condición equivalente $b)$. Si se pone $\mathscr J_1=\mathscr J\otimes_{\sh{O}_S}\sh{O}_{S_1}$, se tiene entonces $\bigoplus_{n\geq0}\mathscr J_1^n=(\bigoplus_{n\geq0}\mathscr J^n)\otimes_{\sh{O}_S}\sh{O}_{S_1}$, de donde se deduce inmediatamente la afirmación relativa al preesquema obtenido por explosión $X'_1$. Por último, para demostrar que $Y'$ es $S$-plano, se puede restringir también al caso en que $S$ y $X$ son afines; si se pone $C=\operatorname{gr}_{\mathfrak J}^*(B)$, todo punto de $Y'$ admite entonces, en virtud de (19.4.2) y de (II, 2.3.6), un

\oldpage[IV-4]{200}

entorno abierto afín cuyo anillo es de la forma $C_{(t)}$, donde $t$ es un elemento homogéneo de grado $1$ de $C$; como por hipótesis $C$ es un $A$-módulo plano, lo mismo ocurre con su anillo de fracciones $C_t$ y con la componente de grado $0$, $C_{(t)}$, de este $A$-módulo graduado, de donde se sigue nuestra afirmación.
  \item[(ii)] Se puede restringir siempre al caso en que $S$ y $X$ son afines, siendo entonces $B$ una $A$-álgebra de presentación finita y $\mathfrak J$ un ideal de tipo finito de $B$. En virtud de (8.9.1), (8.6.3) y (11.2.9), existen un subanillo noetheriano $A_0$ de $A$, una $A_0$-álgebra de tipo finito $B_0$ y un ideal $\mathfrak J_0$ de $B_0$ tales que $B=B_0\otimes_{A_0}A$, $\mathfrak J=\mathfrak J_0B$, $\operatorname{gr}_{\mathfrak J_0}^*(B_0)$ es un $A_0$-módulo \emph{plano} y $\operatorname{gr}_{\mathfrak J}^*(B)=\operatorname{gr}_{\mathfrak J_0}^*(B_0)\otimes_{A_0}A$. En virtud de (i), se tiene entonces $\bigoplus_{n\geq0}\mathfrak J^n=(\bigoplus_{n\geq0}\mathfrak J_0^n)\otimes_{A_0}A=(\bigoplus_{n\geq0}\mathfrak J_0^n)\otimes_{B_0}B$; como $\bigoplus_{n\geq0}\mathfrak J_0^n$ es una $B_0$-álgebra engendrada por el ideal de tipo finito $\mathfrak J_0$ de $B_0$, y por tanto una $B_0$-álgebra de tipo finito, es de presentación finita puesto que $B_0$ es noetheriano; se deduce que $\bigoplus_{n\geq0}\mathfrak J^n$ es una $B$-álgebra de presentación finita. Del mismo modo, el morfismo $\operatorname{Proj}(\bigoplus_{n\geq0}\mathfrak J_0^n)\to\Spec(B_0)$ es de tipo finito (II, 2.7.1), y por tanto de presentación finita; así, el morfismo $X'\to X$, que se deduce de él por cambio de base, es de presentación finita. Ya se sabe por (19.4.2) que la inmersión canónica $Y'\to X'$ es regular; como $Y'$ es $S$-plano, según se ha visto en (i), se deduce de (19.2.4) que la inmersión $Y'\to X'$ es transversalmente regular relativamente a $S$ y que $X'$ es $S$-plano en los puntos de $Y'$, y por tanto (11.3.1) en un entorno abierto de $Y'$. Por otra parte, el hecho de que $\operatorname{gr}_{\mathscr J}^*(\sh{O}_X)$ sea una $\sh{O}_X$-álgebra $S$-plana implica, por (11.3.4), que $X$ es $S$-plano en los puntos de $Y$, y por tanto (11.3.1) en un entorno de $Y$.
\end{enumerate}

\begin{corollary}[19.4.7]
\label{IV.19.4.7}
Sean $g:X\to S$ un morfismo localmente de presentación finita y $Y$ un subpreesquema cerrado de $X$ tal que el morfismo compuesto $Y\to X\to S$ sea localmente de presentación finita; supongamos además que $Y$ sea $S$-plano y que $\sh{O}_X$ sea normalmente plano a lo largo de $Y$ (11.3.4). Entonces (con las notaciones de (19.4.1)), el morfismo $X'\to X$ es de presentación finita, la inmersión canónica $Y'\to X'$ es transversalmente regular de codimensión $1$ relativamente a $S$, $Y'$ es $Y$-plano (y por tanto $S$-plano), y el conjunto de los puntos de $X$ (resp. $X'$) en los que $X$ (resp. $X'$) es $S$-plano es un entorno de $Y$ (resp. $Y'$).
\end{corollary}

En efecto, por hipótesis $\operatorname{gr}_{\mathscr J}^*(\sh{O}_X)$ es un $(\sh{O}_X/\mathscr J)$-módulo plano, y por tanto es $S$-plano, puesto que se supone que $Y$ es $S$-plano. Se pueden aplicar, por consiguiente, los resultados de (19.4.6, (i) y (ii)). Como $Y'$ es $Y$-isomorfo a $\operatorname{Proj}(\operatorname{gr}_{\mathscr J}^*(\sh{O}_X))$ (19.4.2), el mismo razonamiento que en (19.4.6) muestra que $Y'$ es $Y$-plano.

\begin{proposition}[19.4.8]
\label{IV.19.4.8}
Supongamos que se verifican las hipótesis de (19.4.7) y, además, que el morfismo $Y'\to Y$ es liso. Sean $x'$ un punto de $Y'$ y $x$ su imagen en $Y$. Para que $Y$ sea liso sobre $S$ en el punto $x$, es necesario y suficiente que $Y'$ sea liso sobre $S$ en el punto $x'$; cuando así ocurre, $X'$ es entonces liso sobre $S$ en el punto $x'$.

Las conclusiones anteriores se verifican, en particular, cuando el morfismo $X\to S$ es localmente de presentación finita y la inmersión $Y\to X$ es transversalmente regular relativamente a $S$.
\end{proposition}

Teniendo en cuenta la compatibilidad de la formación de un preesquema obtenido por explosión con los cambios de base en el caso considerado (19.4.6, (i)), así como (17.5.1) y (17.7.1),

\oldpage[IV-4]{201}

basta considerar el caso en que $S$ es el espectro de un cuerpo algebraicamente cerrado. Pero entonces equivale a decir que un $S$-preesquema localmente de tipo finito es liso sobre $S$ en un punto o que es regular en dicho punto (17.15.1); las conclusiones se deducen, pues, de (19.4.4).

Cuando $X\to S$ es localmente de presentación finita y la inmersión canónica $Y\to X$ es cuasirregular, esta inmersión es por definición de presentación finita, puesto que el ideal que define $Y$ es de tipo finito; por tanto, el morfismo compuesto $Y\to X\to S$ es localmente de presentación finita. Ya se ha visto en (19.4.3) que, en estas condiciones, el morfismo $Y'\to Y$ es liso; además, se sabe entonces que $\sh{N}_{Y/X}$ es un $\sh{O}_Y$-módulo localmente libre y que $\operatorname{gr}_{\mathscr J}^*(\sh{O}_X)$ es isomorfo a $\mathbf S_{\sh{O}_Y}^*(\sh{N}_{Y/X})$, y por tanto es un $\sh{O}_Y$-módulo plano. Así pues, se verifican todas las hipótesis de (19.4.7) y se pueden aplicar las conclusiones del enunciado de (19.4.8).

\begin{corollary}[19.4.9]
\label{IV.19.4.9}
Bajo las hipótesis generales de (19.4.8), supongamos además que $Y$ es liso sobre $S$ en el punto $x$. Entonces, para toda generización $x_1$ de $x$ en $X$ que no pertenezca a $Y$, $X$ es liso sobre $S$ en el punto $x_1$. Si $Y$ es liso sobre $S$, existe entonces un entorno abierto $U$ de $Y$ en $X$ tal que $X$ es liso sobre $S$ en los puntos de $U-Y$.
\end{corollary}

Para la primera afirmación, se puede restringir al caso en que $X$ es un esquema local de punto cerrado $x$, y la hipótesis implica entonces que $Y$ es liso sobre $S$, pues todo entorno del punto cerrado de $Y$ en $Y$ es necesariamente el propio $Y$; se concluye, por tanto, mediante el mismo razonamiento que en (19.4.5), utilizando que el conjunto de los puntos en los que $X'$ es liso sobre $S$ es abierto en $X'$. Para demostrar la segunda afirmación, se puede reducir al caso en que $S$ es noetheriano y $X$ de tipo finito sobre $S$, gracias a (8.9.1), (8.6.3), (11.2.9) y (17.7.6); se concluye entonces como en (19.4.5), utilizando que el conjunto de los puntos en los que $X$ es liso sobre $S$ es abierto en $X$.

\begin{remark}[19.4.10]
\label{IV.19.4.10}
En (19.4.6, (ii)), sustituyamos la hipótesis de que $X$ e $Y$ son localmente de presentación finita sobre $S$ por la de que $S$ y $X$ (y por tanto también $Y$) son \emph{localmente noetherianos}. Entonces es claro que $\bigoplus_{n\geq0}\mathscr J^n$ es una $\sh{O}_X$-álgebra de tipo finito, que el morfismo $X'\to X$ es de tipo finito, y se sigue también de (19.2.4), aplicado a los preesquemas localmente noetherianos, que la inmersión $Y'\to X'$ es transversalmente regular relativamente a $S$ y que $X'$ es $S$-plano en los puntos de $Y'$.
\end{remark}

\begin{proposition}[19.4.11]
\label{IV.19.4.11}
Sean $X$ un preesquema, $\mathscr E$ un $\sh{O}_X$-módulo cuasicoherente, $u:\mathscr E\to\sh{O}_X$ un homomorfismo de $\sh{O}_X$-módulos, $\mathscr J=u(\mathscr E)$ el ideal cuasicoherente de $\sh{O}_X$ imagen de $u$, y $Y$ el subpreesquema cerrado de $X$ definido por $\mathscr J$. Sea $X'$ el $X$-esquema obtenido por explosión de $Y$. Por otra parte, sean $P=\mathbf P(\mathscr E)$ el fibrado proyectivo sobre $X$ definido por $\mathscr E$ (II, 4.1.1), $p:P\to X$ el morfismo estructural, $\alpha_1^{\#}:p^*(\mathscr E)\to\sh{O}_P(1)$ el homomorfismo canónico (II, 4.1.5.1) y $\mathscr H$ su núcleo, de modo que se tiene la sucesión exacta
\[
  0\longrightarrow\mathscr H\longrightarrow p^*(\mathscr E)\xrightarrow{\alpha_1^{\#}}\sh{O}_P(1)\longrightarrow0.
\]
Por último, sea $\mathscr K$ el ideal cuasicoherente de $\sh{O}_P$ imagen de la restricción $v:\mathscr H\to\sh{O}_P$ de $p^*(u)$, y sea $Z$ el subpreesquema cerrado de $P$ definido por $\mathscr K$.

\oldpage[IV-4]{202}

\begin{enumerate}
  \item[(i)] La imagen de $X'$ por la inmersión cerrada $j:X'\to\mathbf P(\mathscr E)$ correspondiente al homomorfismo sobreyectivo de $\sh{O}_X$-álgebras graduadas $\mathbf S_{\sh{O}_X}(\mathscr E)\to\mathbf S_{\sh{O}_X}(\mathscr J)\to\bigoplus_{n\geq0}\mathscr J^n$ (II, 3.6.2) es un subpreesquema cerrado de $Z$.
  \item[(ii)] Si $\mathscr E$ es un $\sh{O}_X$-módulo localmente libre de rango $m$, $\mathscr H$ es un $\sh{O}_P$-módulo localmente libre de rango $m-1$.
  \item[(iii)] Se supone que $X$ es localmente noetheriano, que $\mathscr E$ es localmente libre de tipo finito, que el subpreesquema $Y$ es regular y que, en todo punto $x\in Y$, la inmersión canónica $i:Y\to X$ es regular y de codimensión igual a $\operatorname{rg}_{\sh{O}_{X,x}}(\mathscr E_x)$ (lo que expresaremos diciendo que el homomorfismo $u$ es \emph{regular} en el punto $x$ (cf. cap. V)). Entonces la imagen de $X'$ por la inmersión cerrada $j:X'\to P$ es el subpreesquema cerrado $Z$ de $P$; en todos los puntos de $X'-Y'$ ($Y'$ es la imagen inversa de $Y$ en $X'$), la inmersión $j$ es transversalmente regular relativamente a $X$; por último, en todos los puntos $x'$ de $Y'$, $X'$ es regular, la inmersión $j$ es regular y de codimensión $\operatorname{rg}_{\sh{O}_{P,z}}(\mathscr H_z)$, donde $z=j(x')$ (en otras palabras, el homomorfismo $v:\mathscr H\to\sh{O}_P$ es regular en $z$).
\end{enumerate}
\end{proposition}

Como todas las cuestiones son locales sobre $X$, se puede restringir al caso en que $X=\Spec(A)$ es afín y $\mathscr E=\widetilde E$, donde $E$ es un $A$-módulo. Para verificar (i), se puede además examinar únicamente lo que ocurre en un abierto afín de $P$ de la forma $D_+(t)$, donde $t\in E=\mathbf S_A^1(E)$ (II, 2.3.14). Si se pone $S=\mathbf S_A(E)$, se tiene entonces $\Gamma(D_+(t),\sh{O}_P)=S_{(t)}$ (II, 2.4.1), y $p^*(\mathscr E)|D_+(t)=(E\otimes_A S_{(t)})^\sim$. Remitiéndonos a la definición de $\alpha_1^{\#}$ (II, 2.6.2.2), se ve que a todo elemento
\[
  a=\sum_i x^{(i)}\otimes\frac{x_1^{(i)}x_2^{(i)}\cdots x_n^{(i)}}{t^n}
\]
de $E\otimes_A S_{(t)}$ (donde $x^{(i)}$ y los $x_j^{(i)}$ pertenecen a $E$), $\alpha_1^{\#}$ hace corresponder el elemento
\[
  \sum_i\frac{x^{(i)}x_1^{(i)}\cdots x_n^{(i)}}{t^n}
\]
de $(S(1))_{(t)}$. Se trata (para demostrar (i)) de probar que, si este último elemento es nulo, lo mismo ocurre con la imagen de
\[
  a'=v(a)=\sum_i u(x^{(i)})\frac{x_1^{(i)}\cdots x_n^{(i)}}{t^n}
\]
por el homomorfismo canónico $S_{(t)}\to(\bigoplus_{n\geq0}(u(E))^n)_{(u(t))}$ (II, 3.6.2). Ahora bien, esta imagen no es sino
\[
  \sum_i\frac{u(x^{(i)})u(x_1^{(i)})\cdots u(x_n^{(i)})}{(u(t))^n}
\]
en el anillo $A_{u(t)}$, es decir, el producto del elemento $u(t)$ por la imagen canónica de $\sum_i x^{(i)}x_1^{(i)}\cdots x_n^{(i)}/t^{n+1}$ mediante el homomorfismo de álgebras que prolonga canónicamente a $S_{(t)}$ el homomorfismo $u:E\to A$ de $A$-módulos. Esto demuestra (i).

Para establecer (ii), observemos que se puede suponer que $E$ es un $A$-módulo libre de rango $m$; puesto que $\alpha_1^{\#}:E\otimes_A S_{(t)}\to(S(1))_{(t)}$ es sobreyectivo y $(S(1))_{(t)}$ es un $S_{(t)}$-módulo libre de rango $1$, la sucesión exacta
\[
  0\longrightarrow H_{(t)}\longrightarrow E\otimes_A S_{(t)}\xrightarrow{\alpha_1^{\#}}(S(1))_{(t)}\longrightarrow0
\]
es \emph{escindida}, y $H_{(t)}$ es, por tanto, un $S_{(t)}$-módulo proyectivo de rango $m-1$, de donde se sigue (ii).

Para demostrar (iii), examinemos primero lo que ocurre sobre un punto $x\in X-Y$. Entonces $u:\mathscr E\to\sh{O}_X$, restringido a un entorno de $x$, es sobreyectivo, por lo que se puede restringir al caso en que $\mathscr J=\sh{O}_X$, caso en el cual $X'$ se identifica canónicamente con $X$. Para ver entonces que el subpreesquema cerrado de $P$ imagen de $X'$ se identifica con $Z$, se puede restringir, como al principio, al caso en que $X=\Spec(A)$ es afín y, además, $E=A^m$, de modo que $S=A[T_0,\ldots,T_{m-1}]$. Con las notaciones del principio, se puede suponer también que $u(t)\ne0$ (pues los valores de $t\in E$

\oldpage[IV-4]{203}

que verifican esta condición engendran el álgebra simétrica $\mathbf S_A(E)$); cambiando, si es necesario, la base de $E$, se puede suponer, por tanto, que $t=T_0$ y $u(T_i)=0$ para $1\leq i\leq m-1$, de modo que $S_{(t)}$ se identifica canónicamente con $A[T_1,\ldots,T_{m-1}]$. Los elementos de $H$ son entonces los elementos de la forma
\[
  \sum_\alpha L_\alpha(\mathbf T)\otimes\mathbf T^\alpha
\]
de $E\otimes_A A[T_1,\ldots,T_{m-1}]$, donde los $L_\alpha(\mathbf T)=\lambda_{\alpha0}+\lambda_{\alpha1}T_1+\cdots+\lambda_{\alpha,m-1}T_{m-1}$ deben ser tales que $\sum_\alpha L_\alpha(\mathbf T)\mathbf T^\alpha=0$ en $A[T_1,\ldots,T_{m-1}]$. Se comprueba fácilmente que, cuando se satisface la condición $\lambda_{\mathbf0,0}=0$, los valores de los $\lambda_{\alpha0}$ pueden elegirse \emph{arbitrariamente} para $|\alpha|\geq1$, y siempre se pueden determinar entonces los $\lambda_{\alpha i}$ para $i\geq1$ de modo que se verifique la condición $\sum_\alpha L_\alpha(\mathbf T)\mathbf T^\alpha=0$. Los elementos correspondientes del ideal $\mathscr K$ son entonces los polinomios $\sum_\alpha\lambda_{\alpha0}\mathbf T^\alpha$ \emph{sin término constante}, y como $(\mathbf S_A(A))_{(t)}$ se identifica con $A$, estos polinomios son exactamente los elementos del núcleo del homomorfismo $(\mathbf S_A(E))_{(t)}\to(\mathbf S_A(A))_{(t)}$. Se ha demostrado así (19.2.1) que la inmersión $j:X'\to P$ es \emph{transversalmente regular} relativamente a $X$ en todos los puntos de $X'-Y'$ y que $j(X')=Z$ en estos puntos.

Queda por examinar lo que ocurre sobre un punto $x\in Y$. La hipótesis de «regularidad» de $u$ en este punto, formulada en (iii), equivale (en virtud de (19.1.1) y de ($0$, 17.1.7)) a decir que $X$ es \emph{regular} en el punto $x$ y que el homomorfismo
\[
  u_x\otimes1:\mathscr E_x\otimes_{\sh{O}_{X,x}}\mathbf k(x)\longrightarrow\mathfrak m_x/\mathfrak m_x^2
\]
es \emph{inyectivo}. Dicho esto, como el morfismo $p:P\to X$ es liso (17.3.9), para todo $z\in P$ sobre $x$, $P$ es regular en el punto $z$ (17.5.8); además ($0$, 17.3.3), el homomorfismo canónico
\begin{equation}
  (\mathfrak m_x/\mathfrak m_x^2)\otimes_{\mathbf k(x)}\mathbf k(z)\longrightarrow\mathfrak m_z/\mathfrak m_z^2
  \tag{19.4.11.1}\label{IV.19.4.11.1}
\end{equation}
es \emph{inyectivo}. Se deduce que el homomorfismo
\[
  (\mathscr E_x\otimes_{\sh{O}_{X,x}}\sh{O}_{P,z})\otimes_{\sh{O}_{P,z}}\mathbf k(z)\longrightarrow\mathfrak m_z/\mathfrak m_z^2
\]
deducido de $p^*(u)$ es también inyectivo, pues se escribe como el compuesto
\[
  (\mathscr E_x\otimes_{\sh{O}_{X,x}}\mathbf k(x))\otimes_{\mathbf k(x)}\mathbf k(z)\longrightarrow(\mathfrak m_x/\mathfrak m_x^2)\otimes_{\mathbf k(x)}\mathbf k(z)\longrightarrow\mathfrak m_z/\mathfrak m_z^2
\]
donde la primera flecha es inyectiva por planitud y la segunda es el homomorfismo inyectivo (19.4.11.1). Como $\mathscr H$ es (en un entorno de $z$ en $P$) un factor directo de $\mathscr E\otimes_{\sh{O}_X}\sh{O}_P=p^*(\mathscr E)$, el homomorfismo $v_z\otimes1:\mathscr H_z\otimes_{\sh{O}_{P,z}}\mathbf k(z)\to\mathfrak m_z/\mathfrak m_z^2$ es también inyectivo, lo que establece la «regularidad» del homomorfismo $v$ en el punto $z$. Queda por ver que la imagen por $j$ de $X'$ es idéntica al subpreesquema cerrado $Z$. Para ello bastará demostrar, por una parte, que las imágenes por $j$ de los puntos de $q^{-1}(x)$ (donde $q:X'\to X$ es el morfismo estructural) son exactamente los de $Z\cap p^{-1}(x)$ y, por otra, que en cada uno de estos puntos $z=j(x')$, el homomorfismo sobreyectivo $\sh{O}_{Z,z}\to\sh{O}_{X',x'}$ es biyectivo. Pero, teniendo en cuenta la hipótesis y (19.4.3), la restricción $Y'\to Y$ de $q$ es lisa y, puesto que $\sh{O}_{Y,x}$ se supone regular, lo mismo ocurre con $\sh{O}_{Y',x'}$ (17.5.8); además, la inmersión $Y'\to X'$ es regular (19.4.2), por lo que $\sh{O}_{X',x'}$ es también un anillo regular (19.1.1). Para demostrar que el homomorfismo $\sh{O}_{Z,z}\to\sh{O}_{X',x'}$ es biyectivo, basta entonces probar que las

\oldpage[IV-4]{204}

dimensiones de estos dos anillos regulares son iguales ($0$, 17.1.9). Pero, en virtud de ($0$, 17.3.3), esto equivale también a la igualdad de las dimensiones de los dos anillos $\sh{O}_{Z,z}\otimes_{\sh{O}_{X,x}}\mathbf k(x)$ y $\sh{O}_{X',x'}\otimes_{\sh{O}_{X,x}}\mathbf k(x)$, de modo que finalmente se pueden sustituir $X'$ y $Z$ por las \emph{fibras} $q^{-1}(x)$ y $Z\cap p^{-1}(x)$, respectivamente, y reducirse a demostrar que el morfismo $q^{-1}(x)\to Z\cap p^{-1}(x)$ deducido de $j$ es un \emph{isomorfismo}. Ahora bien, aquí se puede restringir al caso en que $X=\Spec(B)$ con $B=\sh{O}_{X,x}$ y $E=B^m$, de donde $S=B[T_1,\ldots,T_m]$; además, la hipótesis significa que existe un sistema regular de parámetros $(t_i)_{1\leq i\leq n}$ de $B$ tal que $u(T_i)=t_i$ para $1\leq i\leq m$ ($0$, 17.1.7); sea $\mathfrak J$ la imagen de $u$; si $\mathfrak m$ es el ideal maximal de $B$, se tiene $\mathfrak J^n\otimes_B B/\mathfrak m=\mathfrak J^n/\mathfrak m\mathfrak J^n=(\mathfrak J^n/\mathfrak J^{n+1})\otimes_B(B/\mathfrak m)$ y, por consiguiente (16.9.4.1), $(\bigoplus_{n\geq0}\mathfrak J^n)\otimes_B\mathbf k(x)$ se identifica con $\mathbf k(x)[T_1,\ldots,T_m]$, de donde se deduce que la inmersión $q^{-1}(x)\to p^{-1}(x)$ es aquí un \emph{isomorfismo}; esto implica \emph{a fortiori} que el subpreesquema cerrado $Z\cap p^{-1}(x)$ es idéntico a la fibra $p^{-1}(x)$, y termina la demostración de (19.4.11).

\subsection{Criterios de M-regularidad.}
\label{subsection:IV.19.5}

Retomamos aquí, completándolos, los criterios para que una sucesión de elementos de un anillo $A$ sea M-regular o M-cuasirregular (siendo $M$ un $A$-módulo), ya estudiados en ($0$, 15.1).

\begin{theorem}[19.5.1]
\label{IV.19.5.1}
Sean $A$ un anillo, $\mathbf f=(f_i)_{1\leq i\leq n}$ una sucesión de elementos de $A$ y $M$ un $A$-módulo. Consideremos las propiedades siguientes:
\begin{enumerate}
  \item[a)] La sucesión $\mathbf f$ es M-regular.
  \item[b)] Se tiene $\operatorname{H}_i(\mathbf f,M)=0$ para todo $i>0$ (III, 1.1.3).
  \item[b$'$)] Se tiene $\operatorname{H}_1(\mathbf f,M)=0$.
  \item[c)] La sucesión $\mathbf f$ es M-cuasirregular.
\end{enumerate}
Entonces se tienen las implicaciones
\[
  a)\Rightarrow b)\Rightarrow b')\qquad\text{et}\qquad a)\Rightarrow c).
\]
Pongamos $\mathfrak J=\sum_{i=1}^n f_iA$. Si los módulos $M_i=M/(\sum_{j=1}^i f_jM)$ son separados para la topología $\mathfrak J$-preádica (resp. si todo módulo cociente de un submódulo de $M$ es separado para la topología $\mathfrak J$-preádica), se tiene también la implicación $c)\Rightarrow a)$ (resp. $b')\Rightarrow a)$.
\end{theorem}

Las implicaciones $a)\Rightarrow c)$ y $c)\Rightarrow a)$ cuando los $M_i$ son separados para la topología $\mathfrak J$-preádica ya se han demostrado ($0$, 15.1.9), y se recuerdan únicamente por completitud. También se ha demostrado (III, 1.1.4 y 1.1.3.3) que $a)$ implica $b)$, y $b)$ implica trivialmente $b')$. Queda, pues, por ver que, cuando todo módulo cociente de un submódulo de $M$ es separado para la topología $\mathfrak J$-preádica, $b')$ implica $a)$. Razonemos por inducción sobre $n$. Para $n=1$, la afirmación se deduce trivialmente de las definiciones, pues entonces $\operatorname{H}_1(\mathbf f,M)$ no es sino el núcleo de la homotecia $z\mapsto f_1z$ en $M$ (III, 1.1.1 y 1.1.2). Supongamos, pues, que $n\geq2$, y pongamos $\mathbf f'=(f_i)_{1\leq i\leq n-1}$; con las notaciones de (III, 1.1.2), se tiene $\operatorname{K}_\bullet(\mathbf f,M)=\operatorname{K}_\bullet(f_n)\otimes_A\operatorname{K}_\bullet(\mathbf f',M)$. Se tiene entonces (III, 1.1.4.1) la sucesión exacta
\begin{equation}
  0\longrightarrow\operatorname{H}_0(f_n,\operatorname{H}_1(\mathbf f',M))\longrightarrow\operatorname{H}_1(\mathbf f,M)\longrightarrow\operatorname{H}_1(f_n,\operatorname{H}_0(\mathbf f',M))\longrightarrow0.
  \tag{19.5.1.1}\label{IV.19.5.1.1}
\end{equation}

\oldpage[IV-4]{205}

La relación $\operatorname{H}_1(\mathbf f,M)=0$ implica, pues, $\operatorname{H}_0(f_n,\operatorname{H}_1(\mathbf f',M))=0$ y $\operatorname{H}_1(f_n,\operatorname{H}_0(\mathbf f',M))=0$. La primera de estas dos relaciones significa que $f_n\operatorname{H}_1(\mathbf f',M)=\operatorname{H}_1(\mathbf f',M)$ (III, 1.1.3.5). Ahora bien, por definición (III, 1.1.1), $\operatorname{H}_1(\mathbf f',M)$ es isomorfo a un cociente de un submódulo $N$ de $M^{n-1}$; tomando sobre $M^{n-1}$ la filtración por los $M^j$ para $j\leq n-1$, sobre $N$ la filtración inducida y sobre $\operatorname{H}_1(\mathbf f',M)$ la filtración cociente de la de $N$, se obtiene sobre $\operatorname{H}_1(\mathbf f',M)$ una filtración finita cuyos cocientes son isomorfos a cocientes de submódulos de $M$, y por tanto son separados para la topología $\mathfrak J$-preádica por hipótesis. Como la relación $f_n\operatorname{H}_1(\mathbf f',M)=\operatorname{H}_1(\mathbf f',M)$ implica \emph{a fortiori} $\mathfrak J\operatorname{H}_1(\mathbf f',M)=\operatorname{H}_1(\mathbf f',M)$, se deduce de la hipótesis y de las observaciones anteriores que $\operatorname{H}_1(\mathbf f',M)=0$. La hipótesis de inducción demuestra entonces que la sucesión $\mathbf f'$ es M-regular. Por otra parte, la relación $\operatorname{H}_1(f_n,\operatorname{H}_0(\mathbf f',M))=0$ significa que $f_n$ es $(M/(\sum_{i=1}^{n-1}f_iM))$-regular; por tanto, la sucesión $\mathbf f$ es M-regular.

\begin{corollary}[19.5.2]
\label{IV.19.5.2}
Supongamos que $A$ es un anillo noetheriano, que los $f_i$ ($1\leq i\leq n$) pertenecen al radical de $A$ y que $M$ es un $A$-módulo de tipo finito. Entonces las cuatro condiciones $a)$, $b)$, $b')$, $c)$ de (19.5.1) son equivalentes.
\end{corollary}

En efecto, se sabe entonces ($0_{\mathrm I}$, 7.3.5) que todo $A$-módulo de tipo finito es separado para la topología $\mathfrak J$-preádica.

\begin{proposition}[19.5.3]
\label{IV.19.5.3}
Sean
\[
  0\longrightarrow M'\longrightarrow M\longrightarrow M''\longrightarrow0
\]
una sucesión exacta de $A$-módulos, $\mathbf f=(f_i)_{1\leq i\leq n}$ una sucesión M-regular y $\mathfrak J=\sum_{i=1}^n f_iA$. Consideremos las propiedades siguientes:
\begin{enumerate}
  \item[a)] La sucesión $\mathbf f$ es $M''$-regular.
  \item[b)] La filtración $\mathfrak J$-preádica de $M'$ está inducida sobre $M'$ por la filtración $\mathfrak J$-preádica de $M$ (en otras palabras, el homomorfismo canónico $\operatorname{gr}_{\mathfrak J}^*(M')\to\operatorname{gr}_{\mathfrak J}^*(M)$ es inyectivo).
  \item[c)] El homomorfismo canónico $M'/(\sum_{i=1}^n f_iM')\to M/(\sum_{i=1}^n f_iM)$ es inyectivo.
\end{enumerate}
Entonces se tienen las implicaciones $a)\Rightarrow b)\Rightarrow c)$, y $a)$ implica además que la sucesión $\mathbf f$ es $M'$-regular. Si todo cociente de un submódulo de $M''$ es separado para la topología $\mathfrak J$-preádica, las condiciones $a)$, $b)$ y $c)$ son equivalentes.
\end{proposition}

Es claro que $c)$ es consecuencia de $b)$. Demostremos que, bajo las hipótesis de la última afirmación, $c)$ implica $a)$: puesto que $\mathbf f$ es M-regular, se tiene, por (19.5.1), $\operatorname{H}_1(\mathbf f,M)=0$, de donde se obtiene una sucesión exacta, porción de la sucesión exacta de homología
\[
  0\longrightarrow\operatorname{H}_1(\mathbf f,M'')\longrightarrow\operatorname{H}_0(\mathbf f,M')\longrightarrow\operatorname{H}_0(\mathbf f,M).
\]
Ahora bien, la condición $c)$ expresa que el homomorfismo $\operatorname{H}_0(\mathbf f,M')\to\operatorname{H}_0(\mathbf f,M)$ es inyectivo (III, 1.1.3.5); equivale, pues, a $\operatorname{H}_1(\mathbf f,M'')=0$, de donde, en virtud de la hipótesis de separación y de (19.5.1), se sigue que $\mathbf f$ es $M''$-regular.

Demostremos a continuación que $a)$ implica $c)$ y que $\mathbf f$ es $M'$-regular, procediendo por inducción sobre $n$. Para $n=1$, como $f_1$ es M-regular, también es $M'$-regular, y la propiedad $c)$ no es sino el lema (3.4.1.4). Para $n>1$, la hipótesis de inducción muestra que

\oldpage[IV-4]{206}

si se pone $\mathbf f'=(f_1,\ldots,f_{n-1})$, la sucesión $\mathbf f'$ es $M'$-regular y se tiene, en virtud de $c)$ aplicada a $\mathbf f'$, una sucesión exacta
\[
  0\longrightarrow M'/(\textstyle\sum_{i=1}^{n-1}f_iM')\longrightarrow M/(\textstyle\sum_{i=1}^{n-1}f_iM)\longrightarrow M''/(\textstyle\sum_{i=1}^{n-1}f_iM'')\longrightarrow0.
\]
Por hipótesis, $f_n$ es $(M/(\sum_{i=1}^{n-1}f_iM))$-regular y $(M''/(\sum_{i=1}^{n-1}f_iM''))$-regular; el mismo razonamiento muestra, por una parte, que $f_n$ es $(M'/(\sum_{i=1}^{n-1}f_iM'))$-regular y, por otra, en virtud de (3.4.1.4), que la sucesión
\[
  0\longrightarrow M'/(\textstyle\sum_{i=1}^{n}f_iM')\longrightarrow M/(\textstyle\sum_{i=1}^{n}f_iM)\longrightarrow M''/(\textstyle\sum_{i=1}^{n}f_iM'')\longrightarrow0
\]
es exacta, de donde se sigue $c)$.

Demostremos, por último, que $a)$ implica $b)$. Como la sucesión $\mathbf f$ es entonces M-regular y $M'$-regular, y por tanto M-cuasirregular y $M'$-cuasirregular, se tienen isomorfismos canónicos
\[
  \operatorname{gr}_{\mathfrak J}^*(M')\xrightarrow{\sim}\operatorname{gr}_{\mathfrak J}^0(M')[T_1,\ldots,T_n],
  \qquad
  \operatorname{gr}_{\mathfrak J}^*(M)\xrightarrow{\sim}\operatorname{gr}_{\mathfrak J}^0(M)[T_1,\ldots,T_n]
\]
($0$, 15.1.7), y como se ha visto anteriormente que $\operatorname{gr}_{\mathfrak J}^0(M')\to\operatorname{gr}_{\mathfrak J}^0(M)$ es inyectivo, lo mismo ocurre con $\operatorname{gr}_{\mathfrak J}^*(M')\to\operatorname{gr}_{\mathfrak J}^*(M)$.

Obsérvese también que la equivalencia de las condiciones $a)$, $b)$, $c)$ de (19.5.3) es válida, en particular, cuando $A$ es noetheriano, $M''$ es un $A$-módulo de tipo finito y los $f_i$ pertenecen al radical de $A$.

\begin{env}[19.5.4]
\label{IV.19.5.4}
En lo que sigue de este número, conservamos las notaciones $A$, $\mathbf f$, $M$, $\mathfrak J$ de (19.5.1), y suponemos además que $M$ está provisto de una filtración decreciente $(M_k)$ formada por sub-$A$-módulos, tal que $M_0=M$. Recordemos ($0$, 15.1.5) que se define entonces sobre $M$ una segunda filtración decreciente formada por los submódulos
\[
  M_k'=M_k+\mathfrak JM_{k-1}+\cdots+\mathfrak J^{k-1}M_1+\mathfrak J^kM_0
\]
y que, si $\operatorname{gr}_\bullet(M)$ y $\operatorname{gr}_\bullet'(M)$ son los $A$-módulos graduados asociados a las filtraciones $(M_k)$ y $(M_k')$, respectivamente, se define un homomorfismo graduado sobreyectivo de grado $0$ ($0$, 15.1.5.2)
\[
  \psi_M:(\operatorname{gr}_\bullet(M)\otimes_A(A/\mathfrak J))[T_1,\ldots,T_n]\longrightarrow\operatorname{gr}_\bullet'(M).
\]
Recordemos también que se tiene el homomorfismo canónico sobreyectivo ($0$, 15.1.1.1)
\[
  \varphi_{\operatorname{gr}_\bullet(M)}:(\operatorname{gr}_\bullet(M)\otimes_A(A/\mathfrak J))[T_1,\ldots,T_n]\longrightarrow\operatorname{gr}_{\mathfrak J}^*(\operatorname{gr}_\bullet(M)).
\]
\end{env}

\begin{theorem}[19.5.5]
\label{IV.19.5.5}
Con las notaciones de (19.5.4), consideremos las propiedades siguientes:
\begin{enumerate}
  \item[a)] La sucesión $\mathbf f$ es $\operatorname{gr}_\bullet(M)$-regular.
  \item[b)] El homomorfismo canónico $\psi_M$ es biyectivo.
  \item[c)] El homomorfismo canónico $\varphi_{\operatorname{gr}_\bullet(M)}$ es biyectivo (en otras palabras ($0$, 15.1.7), la sucesión $\mathbf f$ es $\operatorname{gr}_\bullet(M)$-cuasirregular).
\end{enumerate}

\oldpage[IV-4]{207}

Se tienen entonces las implicaciones\footnote{La implicación $b)\Rightarrow c)$, sin hipótesis de separación sobre los $\operatorname{gr}_k(M)$, se debe a P. Deligne, cuya demostración reproducimos a continuación.}
\[
  a)\Rightarrow b)\Rightarrow c).
\]
Además, si, para todo entero $k\geq0$, los módulos cocientes de los submódulos de $\operatorname{gr}_k(M)$ son separados para la topología $\mathfrak J$-preádica (lo que ocurre cuando $A$ es noetheriano, los $f_i$ pertenecen al radical de $A$ y los $\operatorname{gr}_k(M)$ son $A$-módulos de tipo finito), las condiciones $a)$, $b)$ y $c)$ son equivalentes.
\end{theorem}

La implicación $a)\Rightarrow b)$ se ha demostrado en ($0$, 15.1.8); la implicación $c)\Rightarrow a)$ bajo la hipótesis de separación es un caso particular de ($0$, 15.1.9). Queda por demostrar que $b)$ implica $c)$, lo que se hará en varias etapas.

Denotaremos por $\mathfrak J$ el ideal $\sum_{i=1}^n f_iA$ engendrado por los $f_i$; para toda sucesión $\mathbf q=(q_1,\ldots,q_n)$ de $n$ enteros $\geq0$, pondremos $|\mathbf q|=\sum_{i=1}^nq_i$ y $\mathbf f^{\mathbf q}=f_1^{q_1}f_2^{q_2}\cdots f_n^{q_n}$.

\begin{env}[19.5.5.1]
\label{IV.19.5.5.1}
Para que la aplicación $\psi_M$ sea inyectiva (y por tanto biyectiva), es necesario y suficiente que, para $p\geq0$ y $q\geq0$, se verifique la condición siguiente:

$(\mathrm I_{q,p})$ Para toda familia $(x_{\mathbf q}^p)_{|\mathbf q|=q}$ de elementos de $M_p$, la relación
\[
  \sum_{|\mathbf q|=q}\mathbf f^{\mathbf q}x_{\mathbf q}^p\in\sum_{i=1}^q\mathfrak J^{q-i}M_{p+i}+M_{p+q+1}'
\]
implica $x_{\mathbf q}^p\in M_{p+1}+\mathfrak JM_p$ para todo $\mathbf q$ tal que $|\mathbf q|=q$.

Se procede como en ($0$, 15.1.6), considerando el sub-$A$-módulo $Q_k$ (resp. $Q_k'$) de los términos de grado $k$ del primer (resp. segundo) miembro de ($0$, 15.1.5.2). Se dota a $Q_k$ de la filtración
\[
  (Q_k)_i=\sum_{j\leq k-i}\left(\sum_{|\mathbf l|=j}(\operatorname{gr}_{k-j}(M)\otimes_A(A/\mathfrak J))T^{\mathbf l}\right)
\]
y a $Q_k'$ de la filtración imagen formada por los $\psi_M((Q_k)_i)=(Q_k')_i$; basta de nuevo demostrar que los homomorfismos $(\psi_M)_{ki}:\operatorname{gr}_i(Q_k)\to\operatorname{gr}_i(Q_k')$ son inyectivos. Ahora bien, se tiene
\[
  \operatorname{gr}_i(Q_k)=\sum_{|\mathbf l|=k-i}((M_i/M_{i+1})\otimes_A(A/\mathfrak J))T^{\mathbf l}
  =\sum_{|\mathbf l|=k-i}(M_i/(\mathfrak JM_i+M_{i+1}))T^{\mathbf l};
\]

\oldpage[IV-4]{208}

por otra parte, $(Q_k')_{i+1}$ es la imagen de $\sum_{h=0}^{k-i-1}\mathfrak J^{k-i-h-1}M_{i+h+1}$ en $M_k'/M_{k+1}'$. Decir que $(\psi_M)_{ki}$ es inyectivo equivale, pues, a formular la condición $(\mathrm I_{k-i,i})$; de donde se sigue nuestra afirmación.
\end{env}

\begin{env}[19.5.5.2]
\label{IV.19.5.5.2}
Para que la sucesión $\mathbf f$ sea $\operatorname{gr}_\bullet(M)$-cuasirregular, basta que, para $p\geq0$ y $q\geq0$, se verifique la condición siguiente:

$(\mathrm S_{q,p})$ Para toda familia $(x_{\mathbf q}^p)_{|\mathbf q|=q}$ de elementos de $M_p$, la relación
\[
  \sum_{|\mathbf q|=q}\mathbf f^{\mathbf q}x_{\mathbf q}^p\in M_{p+1}
\]
implica $x_{\mathbf q}^p\in M_{p+1}+\mathfrak JM_p$ para todo $\mathbf q$ tal que $|\mathbf q|=q$.

Por definición, decir que $\mathbf f$ es $\operatorname{gr}_\bullet(M)$-cuasirregular significa que, para todo $p\geq0$ y todo $q\geq0$, una relación
\[
  \sum_{|\mathbf q|=q}\mathbf f^{\mathbf q}x_{\mathbf q}^p\in M_{p+1}+\mathfrak J^{q+1}M_p
\]
para una familia $(x_{\mathbf q}^p)_{|\mathbf q|=q}$ de elementos de $M_p$ implica $x_{\mathbf q}^p\in M_{p+1}+\mathfrak JM_p$ para todo $\mathbf q$ tal que $|\mathbf q|=q$. Ahora bien, la hipótesis sobre la sucesión $(x_{\mathbf q}^p)$ significa que existe una familia $(y_{\mathbf q}^p)_{|\mathbf q|=q}$ de elementos de $\mathfrak JM_p$ tal que $\sum_{|\mathbf q|=q}\mathbf f^{\mathbf q}(x_{\mathbf q}^p-y_{\mathbf q}^p)\in M_{p+1}$. La condición $(\mathrm S_{q,p})$ implica entonces $x_{\mathbf q}^p-y_{\mathbf q}^p\in M_{p+1}+\mathfrak JM_p$ para todo $\mathbf q$ tal que $|\mathbf q|=q$, de donde $x_{\mathbf q}^p\in M_{p+1}+\mathfrak JM_p$; se ve, pues, que las condiciones $(\mathrm S_{q,p})$ implican que la sucesión $\mathbf f$ es $\operatorname{gr}_\bullet(M)$-cuasirregular.
\end{env}

Basta, pues, demostrar la proposición siguiente:

\begin{env}[19.5.5.3]
\label{IV.19.5.5.3}
Para todo $r\geq0$, si las condiciones $(\mathrm I_{q,p})$ se verifican para $p+q<r$, entonces las condiciones $(\mathrm S_{q,p})$ se verifican para $p+q<r$.

Introduzcamos las condiciones

$(\mathrm A_{q,p})$ Para toda familia $(x_{\mathbf q}^p)_{|\mathbf q|=q}$ de elementos de $M_p$, la relación
\[
  \sum_{|\mathbf q|=q}\mathbf f^{\mathbf q}x_{\mathbf q}^p\in M_{p+1}
\]
implica
\[
  \sum_{|\mathbf q|=q}\mathbf f^{\mathbf q}x_{\mathbf q}^p\in\sum_{i=1}^q\mathfrak J^{q-i}M_{p+i}.
\]

Es claro que la conjunción de $(\mathrm A_{q,p})$ y $(\mathrm I_{q,p})$ implica $(\mathrm S_{q,p})$. Por otra parte, las condiciones $(\mathrm A_{1,p})$ y $(\mathrm S_{0,p})$ son triviales. Por consiguiente, (19.5.5.3) se deducirá de la proposición siguiente:
\end{env}

\begin{env}[19.5.5.4]
\label{IV.19.5.5.4}
Para cualesquiera $p\geq0$ y $q\geq1$, la condición $(\mathrm A_{q+1,p})$ es consecuencia de la conjunción de las condiciones $(\mathrm A_{q,p})$, $(\mathrm I_{q-1,p+1})$, $(\mathrm I_{q-2,p+2})$, $\ldots$, $(\mathrm I_{0,p+q})$.

En efecto, una vez demostrada esta proposición, las condiciones $(\mathrm I_{q,p})$ que se suponen verificadas para $p+q<r$ implicarán, para cada $p\leq r$, la condición $(\mathrm A_{q,p})$ para $1\leq q<r-p$, por inducción sobre $q$.

Demostremos, pues, (19.5.5.4). Observemos que la condición $(\mathrm A_{q,p})$ equivale también a
\end{env}

\begin{env}[19.5.5.5]
\label{IV.19.5.5.5}
\[
  \mathfrak J^qM_p\cap M_{p+1}\subset\sum_{i=1}^q\mathfrak J^{q-i}M_{p+i}.
\]

Sea, pues, $m\in\mathfrak J^{q+1}M_p\cap M_{p+1}\subset\mathfrak J^qM_p\cap M_{p+1}$. Aplicando la condición $(\mathrm A_{q,p})$, se ve que existe, para cada $i$ tal que $1\leq i\leq q$, una familia $(y_{\mathbf q}^{p+i})_{|\mathbf q|=q-i}$ de elementos de $M_{p+i}$ tal que
\[
  m=\sum_{i=1}^q\left(\sum_{|\mathbf q|=q-i}\mathbf f^{\mathbf q}y_{\mathbf q}^{p+i}\right).
\]
Supongamos que, para $1\leq i<j$ ($j\geq1$), se ha demostrado que
\[
  y_{\mathbf q}^{p+i}\in\mathfrak JM_{p+i}+M_{p+i+1}\qquad(\text{pour }|\mathbf q|=q-i).
\]
De aquí se deduce, por definición (19.5.4),
\[
  m-\sum_{i<j}\left(\sum_{|\mathbf q|=q-i}\mathbf f^{\mathbf q}y_{\mathbf q}^{p+i}\right)\in M_{p+q+1}'.
\]
Puesto que, por hipótesis, se verifica $(\mathrm I_{q-j,p+j})$, se deduce
\[
  y_{\mathbf q}^{p+j}\in\mathfrak JM_{p+j}+M_{p+j+1}
\]

\oldpage[IV-4]{209}

y, por inducción sobre $j$, se ve que esta condición se verifica para todo $j$ tal que $1\leq j\leq q$. Se tiene, pues,
\[
  m\in\sum_{i=1}^q\mathfrak J^{q-i}(\mathfrak JM_{p+i}+M_{p+i+1})=\sum_{i=1}^{q+1}\mathfrak J^{q-i+1}M_{p+i}
\]
y así se ha demostrado $(\mathrm A_{q+1,p})$ y terminado la demostración de (19.5.5).
\end{env}

\begin{remarks}[19.5.6]
\label{IV.19.5.6}
\begin{enumerate}
  \item[(i)] Los resultados de este número conservan un sentido y son válidos cuando, en lugar de un $A$-módulo $M$ y de elementos $f_i\in A$, se toma un objeto $M$ de una categoría abeliana cualquiera y $n$ endomorfismos $f_i$ de $M$ que conmutan dos a dos; las demostraciones se adaptan sin dificultad a este caso más general. La hipótesis de separación para un $A$-módulo $N$ respecto de la topología $\mathfrak J$-preádica debe sustituirse aquí por la hipótesis de que todo subobjeto de $N$ contenido en todos los $\sum_i f_i^r(N)$ ($r$ entero arbitrario) es necesariamente nulo. La misma observación se aplica a los resultados de 19.7.
  \item[(ii)] Supongamos que $A$ es un anillo \emph{graduado} en grados $\geq0$. Si $M$ (resp. cada $\operatorname{gr}_p(M)$) es un $A$-módulo graduado en grados acotados inferiormente y si los $f_i$ no contienen ninguna componente homogénea de grado $0$, la hipótesis de separación de ($0$, 15.1.9) se verifica \emph{ipso facto} para $M$ (resp. cada $\operatorname{gr}_p(M)$), y se ve, por tanto, que en este caso se tiene la implicación $c)\Rightarrow a)$ en (19.5.1) (resp. (19.5.5)).
  \item[(iii)] Recordemos que, sin hipótesis de separación, la implicación $c)\Rightarrow b)$ deja de ser válida incluso para $n=1$ ($0$, 15.1.12, (iii)).
\end{enumerate}
\end{remarks}

\subsection{Sucesiones regulares relativamente a un módulo filtrado cociente.}
\label{subsection:IV.19.6}

\begin{env}[19.6.1]
\label{IV.19.6.1}
Sean $A$ un anillo, $\mathbf f=(f_1,\ldots,f_n)$ una sucesión finita de elementos de $A$, $\mathfrak J=f_1A+\cdots+f_nA$ el ideal que engendra, y consideremos una sucesión exacta de $A$-módulos
\[
  0\longrightarrow R\xrightarrow{j}N\xrightarrow{p}M\longrightarrow0
\]
donde se supone que $N$ está provisto de una filtración decreciente $(N_k)$ formada por $A$-submódulos, con $N_0=N$. Pongamos $R_k=R\cap N_k$ (filtración inducida por $(N_k)$), $M_k=p(N_k)$ (filtración cociente de $(N_k)$), y denotemos por $\operatorname{gr}_\bullet(N)$, $\operatorname{gr}_\bullet(R)$ y $\operatorname{gr}_\bullet(M)$ los $A$-módulos graduados asociados a estas tres filtraciones. Pongamos además, para todo $k$,
\[
  N_k'=N_k+\mathfrak JN_{k-1}+\cdots+\mathfrak J^{k-1}N_1+\mathfrak J^kN_0
\]
de modo que $M_k'=p(N_k')=M_k+\mathfrak JM_{k-1}+\cdots+\mathfrak J^{k-1}M_1+\mathfrak J^kM_0$; pongamos además $R_k'=R\cap N_k'$ (filtración \emph{inducida} por $(N_k')$), y denotemos por $\operatorname{gr}_\bullet'(N)$, $\operatorname{gr}_\bullet'(M)$ y $\operatorname{gr}_\bullet'(R)$ los $A$-módulos graduados asociados a estas tres filtraciones; se tiene entonces un diagrama conmutativo de sucesiones exactas
\begin{equation}
\xymatrix{
  0 \ar[r] & \operatorname{gr}_\bullet(R) \ar[r]^-{\operatorname{gr}(j)} \ar[d] & \operatorname{gr}_\bullet(N) \ar[r]^-{\operatorname{gr}(p)} \ar[d] & \operatorname{gr}_\bullet(M) \ar[r] \ar[d] & 0 \\
  0 \ar[r] & \operatorname{gr}_\bullet'(R) \ar[r]^-{\operatorname{gr}'(j)} & \operatorname{gr}_\bullet'(N) \ar[r]^-{\operatorname{gr}'(p)} & \operatorname{gr}_\bullet'(M) \ar[r] & 0
}
\tag{19.6.1.1}\label{IV.19.6.1.1}
\end{equation}

\oldpage[IV-4]{210}

donde las flechas verticales son los homomorfismos canónicos de un módulo graduado asociado a una filtración en el módulo graduado asociado a una filtración menos fina.

Obsérvese que, si se pone
\[
  R_k''=R_k+\mathfrak JR_{k-1}+\cdots+\mathfrak J^{k-1}R_1+\mathfrak J^kR_0
\]
se tiene evidentemente $R_k''\subset R_k'$, pero las filtraciones $(R_k'')$ y $(R_k')$ son en general distintas; denotaremos por $\operatorname{gr}_\bullet''(R)$ el $A$-módulo graduado asociado a la filtración $(R_k'')$.
\end{env}

\begin{env}[19.6.2]
\label{IV.19.6.2}
En ($0$, 15.1.5.2) se definieron los homomorfismos graduados sobreyectivos de grado $0$
\begin{align*}
  \psi_N&:(\operatorname{gr}_\bullet(N)\otimes_A(A/\mathfrak J))[T_1,\ldots,T_n]\longrightarrow\operatorname{gr}_\bullet'(N),\\
  \psi_M&:(\operatorname{gr}_\bullet(M)\otimes_A(A/\mathfrak J))[T_1,\ldots,T_n]\longrightarrow\operatorname{gr}_\bullet'(M),\\
  \psi_R&:(\operatorname{gr}_\bullet(R)\otimes_A(A/\mathfrak J))[T_1,\ldots,T_n]\longrightarrow\operatorname{gr}_\bullet''(R),
\end{align*}
de los cuales los dos primeros se deducen de las dos últimas flechas verticales de (19.6.1.1).

Como la filtración $(R_k'')$ es más fina que $(R_k')$ sobre $R$, se tiene un homomorfismo canónico $\gamma:\operatorname{gr}_\bullet''(R)\to\operatorname{gr}_\bullet'(R)$; denotaremos por $\psi_R'$ el homomorfismo compuesto $\gamma\circ\psi_R$. Se deduce entonces inmediatamente de las definiciones que se tiene un diagrama conmutativo
\begin{equation}
\xymatrix{
  & 0 \ar[d] \\
  (\operatorname{gr}_\bullet(R)\otimes_A(A/\mathfrak J))[T_1,\ldots,T_n] \ar[r]^-{\psi_R'} \ar[d]^{j'} & \operatorname{gr}_\bullet'(R) \ar[d]^{\operatorname{gr}'(j)} \\
  (\operatorname{gr}_\bullet(N)\otimes_A(A/\mathfrak J))[T_1,\ldots,T_n] \ar[r]^-{\psi_N} \ar[d]^{p'} & \operatorname{gr}_\bullet'(N) \ar[d]^{\operatorname{gr}'(p)} \\
  (\operatorname{gr}_\bullet(M)\otimes_A(A/\mathfrak J))[T_1,\ldots,T_n] \ar[r]^-{\psi_M} \ar[d] & \operatorname{gr}_\bullet'(M) \ar[d] \\
  0 & 0
}
\tag{19.6.2.1}\label{IV.19.6.2.1}
\end{equation}
donde las columnas son exactas.
\end{env}

\begin{proposition}[19.6.3]
\label{IV.19.6.3}
Con las notaciones anteriores, supongamos que la sucesión $\mathbf f$ es $\operatorname{gr}_\bullet(N)$-regular. Consideremos las propiedades siguientes:
\begin{enumerate}
  \item[a)] $\mathbf f$ es $\operatorname{gr}_\bullet(M)$-regular.
  \item[b)] $\psi_M$ es biyectivo.
  \item[c)] $\psi_R'$ es sobreyectivo (dicho de otro modo, $\operatorname{gr}_\bullet'(R)$ está engendrado, como $\operatorname{gr}_{\mathfrak J}^*(A)$-módulo, por $\operatorname{Im}(\operatorname{gr}_\bullet(R)\to\operatorname{gr}_\bullet'(R))$).
  \item[d)] $\psi_R'$ es inyectivo.
  \item[d$'$)] El homomorfismo $(\psi_R')_0:\operatorname{gr}_\bullet(R)\otimes_A(A/\mathfrak J)\to\operatorname{gr}_\bullet'(R)$ obtenido restringiendo $\psi_R'$ a los polinomios de grado $0$ es inyectivo.
  \item[e)] $\psi_R'$ es biyectivo.
  \item[f)] $j'$ es inyectivo.

\oldpage[IV-4]{211}

  \item[f$'$)] El homomorfismo $\operatorname{gr}_\bullet(R)\otimes_A(A/\mathfrak J)\to\operatorname{gr}_\bullet(N)\otimes_A(A/\mathfrak J)$ obtenido restringiendo $j'$ a los polinomios de grado $0$ es inyectivo.
  \item[g)] La filtración $\mathfrak J$-preádica de $\operatorname{gr}_\bullet(R)$ está inducida por la de $\operatorname{gr}_\bullet(N)$.
  \item[h)] Las filtraciones $(R_k')$ y $(R_k'')$ sobre $R$ son idénticas.
\end{enumerate}
Se tienen entonces las implicaciones siguientes:
\[
\xymatrix@C=2.4em@R=1.8em{
  a) \ar@{=>}[r] \ar@{=>}[dr] & e) \ar@{=>}[r] \ar@{=>}[dr] & b) \ar@{<=>}[r] & c) \ar@{<=>}[r] & h) \\
  & g) \ar@{=>}[r] & d) \ar@{<=>}[r] & d') \ar@{<=>}[r] & f) \ar@{<=>}[r] & f')
}
\]
Además, cuando todo módulo cociente de un submódulo de algún $\operatorname{gr}_k(M)$ es separado para la topología $\mathfrak J$-preádica (lo que ocurre cuando $A$ es noetheriano, los $f_i$ pertenecen al radical de $A$ y los $\operatorname{gr}_k(M)$ son $A$-módulos de tipo finito), las condiciones $a)$ a $h)$ son todas equivalentes.
\end{proposition}

Observemos que, según (19.5.5), de la hipótesis de que $\mathbf f$ es $\operatorname{gr}_\bullet(N)$-regular resulta que $\psi_N$ es \emph{biyectivo}; como $\operatorname{gr}'(j)$ es inyectivo, la equivalencia de $d)$ y $f)$ resulta del diagrama (19.6.2.1), así como la de $d')$ y $f')$; por otra parte, $f)$ y $f')$ son trivialmente equivalentes, lo que demuestra la equivalencia de $d)$, $d')$, $f)$ y $f')$.

Como ya se sabe que $\psi_M$ es sobreyectivo y $\psi_N$ biyectivo, la equivalencia de $b)$ y $c)$ se deduce también del diagrama (19.6.2.1) por un caso particular del lema de los cinco.

La relación $\psi_R'=\gamma\circ\psi_R$ (19.6.2) y el hecho de que $\psi_R$ es sobreyectivo muestran que la condición $c)$ equivale a decir que $\gamma$ es sobreyectivo, lo que equivale también a la condición $h)$; se ha demostrado, pues, la equivalencia de $b)$, $c)$ y $h)$.

Es trivial que $e)$ equivale a la conjunción de $c)$ y $d)$, y por tanto también a la de $b)$ y $d)$. Observemos ahora que $a)$ implica $b)$ y que $b)$ implica $a)$ bajo la hipótesis de separación (19.5.5). Por otra parte, las implicaciones $a)\Rightarrow g)\Rightarrow f')$, y la implicación $f')\Rightarrow a)$ bajo la hipótesis de separación, resultan de (19.5.3). Se ha demostrado, pues, que $a)$ implica $e)$ (equivalente a la conjunción de $b)$ y $f')$), y también que todas las condiciones son equivalentes bajo la hipótesis de separación. Q.E.D.

\begin{corollary}[19.6.4]
\label{IV.19.6.4}
Sean $\mathfrak R$ un ideal de $A$ y $\mathfrak L=\mathfrak J+\mathfrak R$. Supongamos que la filtración $(N_k)$ (y por tanto también $(M_k)$) es la filtración $\mathfrak R$-preádica, de modo que las filtraciones $(M_k')$ y $(N_k')$ son las filtraciones $\mathfrak L$-preádicas. Consideremos $\operatorname{gr}_\bullet(R)$ (resp. $\operatorname{gr}_\bullet'(R)$) como un módulo graduado sobre $\operatorname{gr}_{\mathfrak R}^*(A)$ (resp. $\operatorname{gr}_{\mathfrak L}^*(A)$).
\begin{enumerate}
  \item[(i)] Sea $S$ un subconjunto de $\operatorname{gr}_\bullet(R)$ que engendra $\operatorname{gr}_\bullet(R)$ como $\operatorname{gr}_{\mathfrak R}^*(A)$-módulo. Entonces la condición $c)$ de (19.6.3) equivale a la condición siguiente:
  \item[c$'$)] La imagen de $S$ en $\operatorname{gr}_\bullet'(R)$ por $\psi_R'$ engendra $\operatorname{gr}_\bullet'(R)$ como $\operatorname{gr}_{\mathfrak L}^*(A)$-módulo.
  \item[(ii)] Supongamos satisfecha la condición $e)$ de (19.6.3), y además que los cocientes de los $\operatorname{gr}_k(R)$ sean separados para la topología $\mathfrak J$-preádica (lo que ocurre cuando $A$ es noetheriano, los $f_i$ pertenecen al radical de $A$ y los $\operatorname{gr}_k(N)$ son $A$-módulos de tipo finito). Entonces, para que un subconjunto $S$ de $\operatorname{gr}_\bullet(R)$ formado por elementos homogéneos engendre $\operatorname{gr}_\bullet(R)$ como $\operatorname{gr}_{\mathfrak R}^*(A)$-módulo, es necesario y suficiente que su imagen por $\psi_R'$ engendre $\operatorname{gr}_\bullet'(R)$ como $\operatorname{gr}_{\mathfrak L}^*(A)$-módulo.
\end{enumerate}
\end{corollary}

\oldpage[IV-4]{212}

\begin{enumerate}
  \item[(i)] Consideremos el homomorfismo de grado $0$
  \[
    \psi_A:(\operatorname{gr}_{\mathfrak R}^*(A)\otimes_A(A/\mathfrak J))[T_1,\ldots,T_n]\longrightarrow\operatorname{gr}_\bullet'(A)=\operatorname{gr}_{\mathfrak L}^*(A)
  \]
  ($0$, 15.1.5.2). El hecho de que este homomorfismo sea \emph{sobreyectivo} implica que el $\operatorname{gr}_{\mathfrak L}^*(A)$-submódulo\footnote{El texto impreso dice «sous-algèbre»; como $\operatorname{gr}_\bullet'(R)$ solo se considera aquí como módulo y la conclusión es $\operatorname{Im}(\psi_R')$, debe decir «sous-module».} de $\operatorname{gr}_\bullet'(R)$ engendrado por $S'=\psi_R'(S)$ no es otro que $\operatorname{Im}(\psi_R')$, de donde resulta la equivalencia de $c)$ y $c')$.
  \item[(ii)] Puesto que $e)$ se satisface, también se satisface $c)$, y en virtud de (i) queda por demostrar la suficiencia de la condición del enunciado. Ahora bien, como $\psi_R'$ es biyectivo, decir que $S'$ engendra $\operatorname{gr}_\bullet'(R)$ considerado como $\operatorname{gr}_{\mathfrak L}^*(A)$-módulo equivale a decir que $S$ engendra $(\operatorname{gr}_\bullet(R)\otimes_A(A/\mathfrak J))[T_1,\ldots,T_n]$ como $\operatorname{gr}_{\mathfrak L}^*(A)$-módulo, y, en virtud de la sobreyectividad de $\psi_A$, en este enunciado se puede sustituir el anillo $\operatorname{gr}_{\mathfrak L}^*(A)$ por $(\operatorname{gr}_{\mathfrak R}^*(A)\otimes_A(A/\mathfrak J))[T_1,\ldots,T_n]$. Evidentemente, esto equivale a decir que $S$ engendra $\operatorname{gr}_\bullet(R)\otimes_A(A/\mathfrak J)$ como módulo sobre el anillo $\operatorname{gr}_{\mathfrak R}^*(A)\otimes_A(A/\mathfrak J)$, o también como $\operatorname{gr}_{\mathfrak R}^*(A)$-módulo. Ahora bien, si $T$ es el $\operatorname{gr}_{\mathfrak R}^*(A)$-submódulo graduado de $\operatorname{gr}_\bullet(R)$ engendrado por $S$, y si se pone $T'=\operatorname{gr}_\bullet(R)/T$, se tiene por hipótesis $T'/\mathfrak JT'=0$, o $\mathfrak JT'=T'$; pero $T'$ es un módulo graduado, y cada $T_p'$ es un cociente de $\operatorname{gr}_p(R)$, por tanto separado para la topología $\mathfrak J$-preádica; la hipótesis $\mathfrak JT_p'=T_p'$ implica, pues, $T_p'=0$ para todo entero $p$, y por tanto $T'=0$, lo que termina de demostrar el corolario.
\end{enumerate}

\subsection{Criterio de planitud normal de Hironaka.}
\label{subsection:IV.19.7}
\label{IV.19.7}

\begin{theorem}[19.7.1]
\label{IV.19.7.1}
\textup{(Hironaka).} Sean $A$ un anillo, $\mathfrak R$ un ideal de $A$, $\mathbf f=(f_1,\ldots,f_n)$ una sucesión de elementos de $A$ que es $(A/\mathfrak R)$-regular, $\mathfrak J=f_1A+\cdots+f_nA$ el ideal engendrado por $\mathbf f$, $\mathfrak L=\mathfrak J+\mathfrak R$ y $M$ un $A$-módulo.

En este caso se tiene
\[
  (\operatorname{gr}_{\mathfrak R}^*(M))\otimes_A(A/\mathfrak J)
  =(\operatorname{gr}_{\mathfrak R}^*(M))\otimes_A(A/\mathfrak L)
  =(\operatorname{gr}_{\mathfrak R}^*(M))\otimes_{A/\mathfrak R}(A/\mathfrak L),
\]
pues
\[
  (\mathfrak R^nM/\mathfrak R^{n+1}M)\otimes_A(A/\mathfrak J)
  =\mathfrak R^nM/(\mathfrak R^{n+1}+\mathfrak J\mathfrak R^n)M
  =\mathfrak R^nM/(\mathfrak R^{n+1}+\mathfrak L\mathfrak R^n)M,
\]
y la última igualdad resulta de Bourbaki, \emph{Alg.}, cap. II, 3\textsuperscript{a} ed., § 3, n\textsuperscript{o} 7, cor. 2 de la prop. 6.

Consideremos las condiciones siguientes:
\begin{enumerate}
  \item[a)] $\operatorname{gr}_{\mathfrak R}^*(M)$ es un $(A/\mathfrak R)$-módulo plano.
  \item[b)] $\operatorname{gr}_{\mathfrak L}^*(M)$ es un $(A/\mathfrak L)$-módulo plano y el homomorfismo canónico ($0$, 15.1.5.2)
  \[
    \psi_M:((\operatorname{gr}_{\mathfrak R}^*(M))\otimes_A(A/\mathfrak L))[T_1,\ldots,T_n]\longrightarrow\operatorname{gr}_{\mathfrak L}^*(M)
  \]
  es biyectivo.
  \item[c)] $(\operatorname{gr}_{\mathfrak R}^*(M))\otimes_{A/\mathfrak R}(A/\mathfrak L)$ es un $(A/\mathfrak L)$-módulo plano y la sucesión $\mathbf f$ es $(\operatorname{gr}_{\mathfrak R}^*(M))$-regular.
  \item[d)] $\operatorname{gr}_{\mathfrak L}^*(M)$ es un $(A/\mathfrak L)$-módulo plano y, para todo $\mathfrak p\in\operatorname{Ass}_{A/\mathfrak R}(A/\mathfrak L)$, $(\operatorname{gr}_{\mathfrak R}^*(M))_{\mathfrak p}$ es un $(A/\mathfrak R)_{\mathfrak p}$-módulo plano.
\end{enumerate}
Se tienen entonces las implicaciones
\[
\xymatrix@C=3em@R=1.6em{
  a) \ar@{=>}[r] \ar@{=>}[dr] & c) \ar@{=>}[r] & b) \\
  & d)
}
\]
Cuando todo cociente de un submódulo de algún $\operatorname{gr}_{\mathfrak R}^p(M)$ ($p\geq0$) es idealmente separado respecto de $\mathfrak J$ (Bourbaki, \emph{Alg. comm.}, cap. III, § 5, n\textsuperscript{o} 1), las condiciones $a)$, $b)$ y $c)$ son equivalentes.

\oldpage[IV-4]{213}

Por último, cuando $A$ es noetheriano, los $f_i$ pertenecen al radical de $A$, la sucesión $\mathbf f$ es $\operatorname{gr}_{\mathfrak R}^*(A)$-regular y $M$ es un $A$-módulo de tipo finito, las condiciones $a)$, $b)$, $c)$ y $d)$ son equivalentes.
\end{theorem}

La implicación $a)\Rightarrow c)$ es inmediata ($0$, 15.1.13). Si se cumple $c)$, resulta de (19.5.5) que $\psi_M$ es biyectivo; además, $\operatorname{gr}_{\mathfrak R}^*(M)\otimes_A(A/\mathfrak L)=\operatorname{gr}_{\mathfrak R}^*(M)\otimes_{A/\mathfrak R}(A/\mathfrak L)$ es un $(A/\mathfrak L)$-módulo plano, luego lo mismo ocurre con $(\operatorname{gr}_{\mathfrak R}^*(M)\otimes_A(A/\mathfrak L))[T_1,\ldots,T_n]$ y, por consiguiente, con $\operatorname{gr}_{\mathfrak L}^*(M)$, puesto que $\psi_M$ es un $(A/\mathfrak L)$-isomorfismo; por tanto, $c)$ implica $b)$. El hecho de que $a)$ implique $b)$ muestra inmediatamente que $a)$ implica también $d)$.

Cuando todo cociente de un submódulo\footnote{El texto impreso dice solamente «tout quotient d'un $\operatorname{gr}_{\mathfrak R}^p(M)$»; para aplicar (19.5.5) se requiere, como en el enunciado precedente, todo cociente de un submódulo de $\operatorname{gr}_{\mathfrak R}^p(M)$.} de algún $\operatorname{gr}_{\mathfrak R}^p(M)$ es idealmente separado respecto de $\mathfrak J$, resulta de (19.5.5) que la condición $b)$ implica que $\mathbf f$ es $(\operatorname{gr}_{\mathfrak R}^*(M))$-regular; además, como $(\operatorname{gr}_{\mathfrak R}^*(M)\otimes_{A/\mathfrak R}(A/\mathfrak L))[T_1,\ldots,T_n]$, isomorfo a $\operatorname{gr}_{\mathfrak L}^*(M)$, es entonces un $(A/\mathfrak L)$-módulo plano, lo mismo ocurre con $\operatorname{gr}_{\mathfrak R}^*(M)\otimes_{A/\mathfrak R}(A/\mathfrak L)$, que es un factor directo suyo; esto demuestra, pues, que $b)$ implica entonces $c)$. Por otra parte, bajo la misma hipótesis, $c)$ implica $a)$ en virtud de ($0$, 15.1.21) (donde la hipótesis noetheriana puede sustituirse por la hipótesis de que $M$ es idealmente separado respecto de $\mathfrak J$, en virtud de ($0_{\mathrm{III}}$, 10.2.1)).

Queda, pues, por demostrar que, cuando $A$ es noetheriano, $M$ es de tipo finito y los $f_i$ pertenecen al radical de $A$, $d)$ implica $a)$.

Por hipótesis existe un $A$-módulo \emph{libre de tipo finito} $N$ y una sucesión exacta $0\to R\to N\to M\to0$. Basta demostrar que $d)$ implica $b)$, puesto que todo cociente de un submódulo de algún $\operatorname{gr}_{\mathfrak R}^p(M)$ es entonces idealmente separado respecto de $\mathfrak J$, en virtud de que $\operatorname{gr}_{\mathfrak R}^p(M)$ es un $A$-módulo de tipo finito, $\mathfrak J$ está contenido en el radical de $A$ y $A$ es noetheriano (Bourbaki, \emph{Alg. comm.}, cap. III, § 5, n\textsuperscript{o} 1); en otros términos, se trata de ver que $\psi_M$ es \emph{biyectivo}.

Como, por hipótesis, la sucesión $\mathbf f$ es $\operatorname{gr}_{\mathfrak R}^*(A)$-regular, también es $\operatorname{gr}_{\mathfrak R}^*(N)$-regular; por consiguiente, se puede aplicar (19.6.3), pues todo módulo cociente de un submódulo de algún $\operatorname{gr}_{\mathfrak R}^k(M)$ es entonces separado para la topología $\mathfrak J$-preádica, ya que los $f_i$ pertenecen al radical de $A$ por hipótesis. El diagrama conmutativo (19.6.2.1) se escribe aquí
\begin{equation}
\xymatrix{
  & 0 \ar[d] \\
  (\operatorname{gr}_{\mathfrak R}^*(R,N)\otimes_{A/\mathfrak R}(A/\mathfrak L))[T_1,\ldots,T_n] \ar[r]^-{\psi_R'} \ar[d] & \operatorname{gr}_{\mathfrak L}^*(R,N) \ar[d] \\
  (\operatorname{gr}_{\mathfrak R}^*(N)\otimes_{A/\mathfrak R}(A/\mathfrak L))[T_1,\ldots,T_n] \ar[r]^-{\psi_N} \ar[d] & \operatorname{gr}_{\mathfrak L}^*(N) \ar[d] \\
  (\operatorname{gr}_{\mathfrak R}^*(M)\otimes_{A/\mathfrak R}(A/\mathfrak L))[T_1,\ldots,T_n] \ar[r]^-{\psi_M} \ar[d] & \operatorname{gr}_{\mathfrak L}^*(M) \ar[d] \\
  0 & 0
}
\tag{19.7.1.1}\label{IV.19.7.1.1}
\end{equation}
donde $\operatorname{gr}_{\mathfrak R}^*(R,N)$ (resp. $\operatorname{gr}_{\mathfrak L}^*(R,N)$) es el graduado asociado a $R$ para la filtración por los $R\cap\mathfrak R^mN$ (resp. $R\cap\mathfrak L^mN$); recordemos que en este diagrama las columnas son exactas y que $\psi_N$

\oldpage[IV-4]{214}

es \emph{biyectivo}. Se trata, en virtud de (19.6.3), de demostrar que $\psi_R'$ es \emph{sobreyectivo}, lo que se hará en varias etapas.

\begin{env}[19.7.1.2]
\label{IV.19.7.1.2}
Pongamos $B=\operatorname{gr}_{\mathfrak R}^*A\otimes_{A/\mathfrak R}(A/\mathfrak L)$, $P=\operatorname{gr}_{\mathfrak R}^*(N)\otimes_{A/\mathfrak R}(A/\mathfrak L)$ y $Q=\operatorname{gr}_{\mathfrak L}^*(R,N)$; puesto que $\psi_N$ es biyectivo, $P[T_1,\ldots,T_n]$ se identifica con $\operatorname{gr}_{\mathfrak L}^*(N)$, y $Q$ con un $B[T_1,\ldots,T_n]$-submódulo de $P[T_1,\ldots,T_n]$; veamos primero que se tiene
\end{env}

\begin{env}[19.7.1.3]
\label{IV.19.7.1.3}
\[
  Q=(Q\cap P)[T_1,\ldots,T_n].
\]
\end{env}

Para ello, pongamos $Z=Q/((Q\cap P)[T_1,\ldots,T_n])$; es un $B[T_1,\ldots,T_n]$-submódulo de
\[
  P[T_1,\ldots,T_n]/((Q\cap P)[T_1,\ldots,T_n])=(P/(Q\cap P))[T_1,\ldots,T_n].
\]
Como $(A/\mathfrak L)$-módulo, es por tanto isomorfo a un submódulo de una suma directa de $(A/\mathfrak L)$-módulos isomorfos a $P/(Q\cap P)=(P+Q)/Q$. Pero $(P+Q)/Q$ es un $(A/\mathfrak L)$-módulo isomorfo a un submódulo de $P[T_1,\ldots,T_n]/Q$, que no es otro que $\operatorname{gr}_{\mathfrak L}^*(M)$ en virtud del diagrama (19.7.1.1). Por tanto,
\begin{equation}
  \operatorname{Ass}_{A/\mathfrak L}(Z)\subset\operatorname{Ass}_{A/\mathfrak L}(\operatorname{gr}_{\mathfrak L}^*(M)).
  \tag{19.7.1.4}\label{IV.19.7.1.4}
\end{equation}

Pero cada $\operatorname{gr}_{\mathfrak L}^m(M)$ es, por hipótesis, un $(A/\mathfrak L)$-módulo \emph{plano} de tipo finito, y por tanto proyectivo, puesto que $A/\mathfrak L$ es noetheriano; se tiene, por consiguiente, $\operatorname{Ass}_{A/\mathfrak L}(\operatorname{gr}_{\mathfrak L}^*(M))\subset\operatorname{Ass}_{A/\mathfrak L}(A/\mathfrak L)$. Para ver que $Z=0$, basta, pues, (Bourbaki, \emph{Alg. comm.}, cap. IV, § 1, n\textsuperscript{o} 1, cor. 1 de la prop. 2 y n\textsuperscript{o} 3, prop. 7) ver que, para todo $\mathfrak q\in\operatorname{Ass}_{A/\mathfrak L}(A/\mathfrak L)$, se tiene $Z_{\mathfrak q}=0$. Ahora bien, los ideales de $\operatorname{Ass}_{A/\mathfrak L}(A/\mathfrak L)$ son los ideales de la forma $\mathfrak p/(\mathfrak L/\mathfrak R)$, donde $\mathfrak p\in\operatorname{Ass}_{A/\mathfrak R}(A/\mathfrak L)$, y equivale por tanto a ver que $Z_{\mathfrak p}=0$ para todo $\mathfrak p\in\operatorname{Ass}_{A/\mathfrak R}(A/\mathfrak L)$. Si $\mathfrak p=\mathfrak r/\mathfrak R$, donde $\mathfrak r$ es un ideal primo de $A$, se tiene $(\operatorname{gr}_{\mathfrak R}^*(M))_{\mathfrak p}=\operatorname{gr}_{\mathfrak R_{\mathfrak r}}^*(M_{\mathfrak r})$ y $(A/\mathfrak L)_{\mathfrak p}=A_{\mathfrak r}/\mathfrak L_{\mathfrak r}$, y las imágenes de los $f_i$ en $A_{\mathfrak r}$ forman una sucesión $\operatorname{gr}_{\mathfrak R_{\mathfrak r}}^*(A_{\mathfrak r})$-regular ($0$, 15.1.14); aplicando a $A_{\mathfrak r}$, $\mathfrak R_{\mathfrak r}$ y $M_{\mathfrak r}$ la implicación $a)\Rightarrow b)$ del enunciado, se ve que la hipótesis $d)$ implica que $\psi_{M_{\mathfrak r}}$ es biyectivo, y por tanto también $\psi_{R_{\mathfrak r}}'$ en virtud de (19.6.3); en otros términos, $Q_{\mathfrak r}=Q_{\mathfrak r}'[T_1,\ldots,T_n]$, donde se ha puesto $Q'=\operatorname{gr}_{\mathfrak R}^*(R,N)\otimes_{A/\mathfrak R}(A/\mathfrak L)$, que es aquí un $B$-submódulo de $P$; en particular, $Q_{\mathfrak r}'=Q_{\mathfrak r}\cap P_{\mathfrak r}$, de modo que finalmente se tiene $Z_{\mathfrak r}=Z_{\mathfrak p}=0$, lo que termina de demostrar (19.7.1.3).

\begin{env}[19.7.1.5]
\label{IV.19.7.1.5}
En virtud de (19.7.1.3), para demostrar que $\psi_R'$ es sobreyectivo, basta mostrar que todo elemento homogéneo de grado $m$ de $Q\cap P$ es la imagen de un elemento del mismo grado de $\operatorname{gr}_{\mathfrak R}^*(R,N)\otimes_{A/\mathfrak R}(A/\mathfrak L)$. Ahora bien, los elementos de grado $m$ de $P$ se identifican (por $\psi_N$) con los elementos de $\operatorname{gr}_{\mathfrak L}^m(N)$ que pertenecen a $(\mathfrak R^mN+\mathfrak L^{m+1}N)/\mathfrak L^{m+1}N$; los que son imágenes de elementos de $\operatorname{gr}_{\mathfrak R}^m(R,N)\otimes_{A/\mathfrak R}(A/\mathfrak L)$ se identifican con los elementos de $\operatorname{gr}_{\mathfrak L}^m(N)$ pertenecientes a $((R\cap\mathfrak R^mN)+\mathfrak L^{m+1}N)/\mathfrak L^{m+1}N$. Bastará, pues, demostrar finalmente, para todo entero $m>0$, la inclusión
\[
  ((R\cap\mathfrak L^mN)+\mathfrak L^{m+1}N)\cap(\mathfrak R^mN+\mathfrak L^{m+1}N)\subset(R\cap\mathfrak R^mN)+\mathfrak L^{m+1}N.
\]
Como $\mathfrak R\subset\mathfrak L$, se comprueba inmediatamente que el primer miembro de esta relación es igual a $(R\cap(\mathfrak R^mN+\mathfrak L^{m+1}N))+\mathfrak L^{m+1}N$, de modo que todo se reduce a demostrar
\end{env}

\begin{env}[19.7.1.6]
\label{IV.19.7.1.6}
\[
  R\cap(\mathfrak R^mN+\mathfrak L^{m+1}N)\subset(R\cap\mathfrak R^mN)+\mathfrak L^{m+1}N.
\]
\end{env}

\oldpage[IV-4]{215}

Procederemos por inducción sobre $m$, suponiendo, pues, que (19.7.1.6) se cumple cuando se sustituye $m$ por un entero $m'<m$. Consideraremos, por otra parte, para $m$ fijo y un entero $d\geq0$, la relación
\begin{equation}
  R\cap(\mathfrak R^mN+\mathfrak L^{m+1}N)\subset(R\cap\mathfrak R^dN\cap\mathfrak L^mN)+\mathfrak L^{m+1}N.
  \tag{$*_d$}
\end{equation}

Está claro que basta demostrar $(*_d)$ para $d=m$, y por otra parte $(*_0)$ es trivialmente verdadera. Demostraremos $(*_d)$ por inducción sobre $d$; en otros términos, suponemos, para un $d\geq0$ fijo (con $d<m$), que $(*_d)$ es verdadera, y queremos demostrar
\begin{equation}
  R\cap(\mathfrak R^mN+\mathfrak L^{m+1}N)\subset(R\cap\mathfrak R^{d+1}N\cap\mathfrak L^mN)+\mathfrak L^{m+1}N.
  \tag{$*_{d+1}$}
\end{equation}

Consideremos para ello, para un $h\geq0$, la relación
\begin{equation}
  R\cap(\mathfrak R^mN+\mathfrak L^{m+1}N)\subset(R\cap(\mathfrak R^{d+1}N+\mathfrak R^d\mathfrak L^hN)\cap\mathfrak L^mN)+\mathfrak L^{m+1}N.
  \tag{$**_{d+1,h}$}
\end{equation}

La hipótesis $(*_d)$ implica que $(**_{d+1,0})$ es verdadera. Demostraremos por inducción sobre $h$ que $(**_{d+1,h})$ es verdadera para todo $h\geq0$. Pero se tiene el

\begin{lemma}[19.7.1.7]
\label{IV.19.7.1.7}
Sean $A$ un anillo noetheriano, $N$ un $A$-módulo de tipo finito, $E$ y $F$ dos submódulos de $N$, y $\mathfrak L$ un ideal de $A$. Para todo $k>0$ existe $h>0$ tal que
\[
  E\cap(F+\mathfrak L^hN)\subset(E\cap F)+\mathfrak L^kN.
\]
\end{lemma}

En efecto, sea $\varphi:N\to N/(E\cap F)=N_1$ el homomorfismo canónico, y pongamos $E_1=\varphi(E)$, $F_1=\varphi(F)$, de modo que $E_1\cap F_1=0$ y $\varphi(E\cap(F+\mathfrak L^hN))=E_1\cap(F_1+\mathfrak L^hN_1)$. Se sabe ($0_{\mathrm I}$, 7.3.2) que existe $h$ tal que $(E_1+F_1)\cap\mathfrak L^hN_1\subset\mathfrak L^k(E_1+F_1)$; si se elige así $h$ y $x_1\in E_1$ es tal que existe $y_1\in F_1$ para el cual $x_1-y_1\in\mathfrak L^hN_1$, se deduce de lo anterior que existen $u_1\in\mathfrak L^kE_1$ y $v_1\in\mathfrak L^kF_1$ tales que $x_1=u_1$ e $y_1=v_1$; por tanto, $E_1\cap(F_1+\mathfrak L^hN_1)\subset\mathfrak L^kN_1$, lo que demuestra el lema.

Basta entonces aplicar este lema con $E=R\cap\mathfrak L^mN$, $F=\mathfrak R^{d+1}N$ y tomar $k=m+1$ para deducir que, con el valor correspondiente de $h$, $(**_{d+1,h})$ implica $(*_{d+1})$.

Supondremos, pues, en todo lo que sigue que $(**_{d+1,h})$ se cumple.

\begin{env}[19.7.1.8]
\label{IV.19.7.1.8}
Demostración de $(**_{d+1,h+1})$: Primer caso: $h\leq m-d$.
\end{env}

Partamos, pues, de un elemento
\begin{env}[19.7.1.9]
\label{IV.19.7.1.9}
\[
  g\in R\cap(\mathfrak R^{d+1}N+\mathfrak R^d\mathfrak L^hN)\cap\mathfrak L^mN
\]
\end{env}
congruente módulo $\mathfrak L^{m+1}N$ con un elemento $g_1\in R\cap(\mathfrak R^mN+\mathfrak L^{m+1}N)$; se tiene entonces
\begin{env}[19.7.1.10]
\label{IV.19.7.1.10}
\[
  g\in\mathfrak R^mN+\mathfrak L^{m+1}N.
\]
\end{env}

Bastará demostrar que $g$ es también congruente módulo $\mathfrak L^{m+1}N$ con un elemento $g_2\in R\cap(\mathfrak R^{d+1}N+\mathfrak R^d\mathfrak L^{h+1}N)$, pues esto implicará $g_2\in\mathfrak L^mN$, ya que se tendrá $g\in\mathfrak L^mN$ y $g_2-g\in\mathfrak L^{m+1}N$. En el caso $h\leq m-d$ que consideramos, basta demostrar la relación
\begin{env}[19.7.1.11]
\label{IV.19.7.1.11}
\[
  (\mathfrak R^mN+\mathfrak L^{m+1}N)\cap\mathfrak R^dN\subset\mathfrak L^{m-d+1}\mathfrak R^dN+\mathfrak R^{d+1}N.
\]
\end{env}

\oldpage[IV-4]{216}

pues las relaciones (19.7.1.9) y (19.7.1.10) implican que $g$ pertenece al primer miembro de (19.7.1.11), y como $h\leq m-d$, se tiene $\mathfrak L^{m-d+1}\subset\mathfrak L^{h+1}$, lo que establecerá $(**_{d+1,h+1})$. Por otra parte, se tiene $\mathfrak R^mN\subset\mathfrak R^dN$, y la relación
\begin{env}[19.7.1.12]
\label{IV.19.7.1.12}
\[
  \mathfrak L^{m+1}N\cap\mathfrak R^dN\subset\mathfrak L^{m-d+1}\mathfrak R^dN
\]
\end{env}
implicará (19.7.1.11). Como $N$ es un $A$-módulo libre de tipo finito, se ve que basta demostrar
\begin{env}[19.7.1.13]
\label{IV.19.7.1.13}
\[
  \mathfrak L^{m+1}\cap\mathfrak R^d\subset\mathfrak L^{m-d+1}\mathfrak R^d.
\]
\end{env}

Bastará demostrar que, en general, para $d<m$, se tiene
\begin{env}[19.7.1.14]
\label{IV.19.7.1.14}
\[
  \mathfrak L^{m+1}\cap\mathfrak R^d\subset\mathfrak L^{m-d+1}\mathfrak R^d+\mathfrak R^{d+1}.
\]
\end{env}
En efecto, como $\mathfrak R\subset\mathfrak L$, y por tanto $\mathfrak L^{m-d+1}\mathfrak R^d\subset\mathfrak L^{m+1}$, la relación precedente implicará
\[
  \mathfrak L^{m+1}\cap\mathfrak R^d\subset\mathfrak L^{m-d+1}\mathfrak R^d+\mathfrak L^{m+1}\cap\mathfrak R^{d+1}
\]
y por inducción sobre $d\leq m$ se concluirá
\[
  \mathfrak L^{m+1}\cap\mathfrak R^d\subset\mathfrak L^{m-d+1}\mathfrak R^d+\mathfrak R^{m+1}
\]
de donde resulta (19.7.1.13), puesto que $\mathfrak R^{m+1}\subset\mathfrak L^{m-d+1}\mathfrak R^d$.

Como $\mathfrak L=\mathfrak J+\mathfrak R$, la relación (19.7.1.14) se escribe también
\begin{env}[19.7.1.15]
\label{IV.19.7.1.15}
\[
\begin{aligned}
  &(\mathfrak J^{m+1}+\mathfrak J^m\mathfrak R+\cdots+\mathfrak J\mathfrak R^m+\mathfrak R^{m+1})\cap\mathfrak R^d\\
  &\quad\subset\mathfrak J^{m-d+1}\mathfrak R^d+\mathfrak J^{m-d}\mathfrak R^{d+1}+\cdots+\mathfrak J\mathfrak R^m+\mathfrak R^{m+1}+\mathfrak R^{d+1}\\
  &\quad=\mathfrak J^{m-d+1}\mathfrak R^d+\mathfrak R^{d+1}.
\end{aligned}
\]
\end{env}

Veamos que esta inclusión es a su vez consecuencia de
\begin{env}[19.7.1.16]
\label{IV.19.7.1.16}
\[
  \mathfrak R^a\mathfrak J^b\cap\mathfrak R^{a+1}\subset\mathfrak R^{a+1}\mathfrak J^{b-1}
\]
\end{env}
válida para $a\geq0$, $b\geq1$. En efecto, para demostrar (19.7.1.15) basta mostrar que, para $0<q\leq d$, se tiene
\begin{env}[19.7.1.17]
\label{IV.19.7.1.17}
\[
\begin{aligned}
  &(\mathfrak J^{m+1}+\mathfrak J^m\mathfrak R+\cdots+\mathfrak J\mathfrak R^m+\mathfrak R^{m+1})\cap\mathfrak R^d\\
  &\quad\subset(\mathfrak J^{m+1-q}\mathfrak R^q+\mathfrak J^{m-q}\mathfrak R^{q+1}+\cdots+\mathfrak R^{m+1})\cap\mathfrak R^d
\end{aligned}
\]
\end{env}
pues para $q=d$ esto implicará (19.7.1.15). Ahora bien, para demostrar (19.7.1.17), basta proceder por inducción sobre $q$, suponiendo verdadera la relación para $q<d$; por hipótesis, un elemento del primer miembro se escribe entonces $y+z$, con $y\in\mathfrak J^{m+1-q}\mathfrak R^q$ y $z\in\mathfrak J^{m-q}\mathfrak R^{q+1}+\cdots+\mathfrak R^{m+1}$; como $y+z\in\mathfrak R^d\subset\mathfrak R^{q+1}$, se deduce $y\in\mathfrak R^{q+1}$, puesto que $z\in\mathfrak R^{q+1}$; teniendo en cuenta (19.7.1.16), se tiene $y\in\mathfrak J^{m-q}\mathfrak R^{q+1}$, y por tanto
\[
  y+z\in\mathfrak J^{m-q}\mathfrak R^{q+1}+\cdots+\mathfrak R^{m+1},\qquad\text{et}\qquad y+z\in\mathfrak R^d
\]
por hipótesis, lo que demuestra (19.7.1.17) con $q$ sustituido por $q+1$.

Queda por establecer (19.7.1.16). El caso $b=1$ es trivial; supongamos $b\geq2$. Un elemento del primer miembro se escribe $a_1f_1+\cdots+a_nf_n$, con $a_i\in\mathfrak R^a\mathfrak J^{b-1}$, y su imagen en $Q=\mathfrak R^a/\mathfrak R^{a+1}$ es nula. Como la sucesión $\mathbf f$ es $Q$-regular, las relaciones de Koszul permiten elegir $c_{ij}\in\mathfrak R^a\mathfrak J^{b-2}$, con $c_{ij}=-c_{ji}$ y $c_{ii}=0$, tales que $a_i-\sum_j c_{ij}f_j\in\mathfrak R^{a+1}$.\footnote{Nota de corrección. El control de la filtración necesario para esta elección se demuestra como sigue. Pongamos $I=\mathfrak J$, $s=b-1$ y $K_\bullet=K_\bullet(\mathbf f;Q)$, con diferencial $\delta$. Para $u\in I^sK_1$ tal que $\delta_1u=0$, la regularidad da $u=\delta_2C$. Por ($0$, 15.1.7 y 15.1.9), $\operatorname{gr}_I(Q)\simeq(Q/IQ)[T_1,\ldots,T_n]$; el complejo de Koszul de las indeterminadas sobre este módulo es exacto en grados positivos. Si $C\in I^dK_2$ con $d<s-1$, su clase de grado $d$ es un ciclo. Para $d=0$, esta clase es nula, pues una preimagen por $\delta_3$ tendría grado de coeficiente $-1$. Para $d\geq1$, es el borde de una clase de grado de coeficiente $d-1$; al levantar esta clase a $D\in I^{d-1}K_3$ y sustituir $C$ por $C-\delta_3D$, se obtiene $C\in I^{d+1}K_2$, sin modificar $\delta_2C$. Después de a lo sumo $s-1$ pasos, $C\in I^{s-1}K_2$. Al levantar sus coeficientes a $\mathfrak R^a\mathfrak J^{b-2}$, se obtienen los $c_{ij}$ requeridos. No es necesario pasar al completado. El texto impreso concluye erróneamente que los coeficientes elegidos inicialmente pertenecen necesariamente a $\mathfrak R^{a+1}$.} Sustituyamos los $a_i$ por estos nuevos coeficientes: los términos antisimétricos se cancelan, de modo que la suma $a_1f_1+\cdots+a_nf_n$ no cambia y todos sus coeficientes pertenecen a $\mathfrak R^a\mathfrak J^{b-1}\cap\mathfrak R^{a+1}$; basta ahora razonar por

\oldpage[IV-4]{217}

inducción sobre $b$ para $b\geq1$ (siendo trivial el caso $b=1$) para terminar de demostrar (19.7.1.16) y, por consiguiente, también $(**_{d+1,h+1})$ para $h\leq m-d$.

\begin{env}[19.7.1.18]
\label{IV.19.7.1.18}
Demostración de $(**_{d+1,h+1})$; Segundo caso: $h>m-d$.
\end{env}

Demostremos primero que, para todo $h\geq0$, se tiene
\begin{env}[19.7.1.19]
\label{IV.19.7.1.19}
\[
  R\cap(\mathfrak R^{d+1}N+\mathfrak R^d\mathfrak L^hN)\subset\mathfrak J^h(R\cap\mathfrak R^dN)+\mathfrak L^{h+1}\mathfrak R^dN+\mathfrak R^{d+1}N.
\]
\end{env}

Consideremos para ello un elemento $z$ del primer miembro de (19.7.1.19). Observemos que se tiene $\mathfrak R^{d+1}N+\mathfrak R^d\mathfrak L^hN=\mathfrak R^{d+1}N+\mathfrak R^d\mathfrak J^hN$, puesto que $\mathfrak L=\mathfrak R+\mathfrak J$; consideraremos la clase $\bar z$ de $z$ en $\operatorname{gr}_{\mathfrak J}^h(\operatorname{gr}_{\mathfrak R}^d(N))$, que, en virtud de (19.5.5) y de la hipótesis sobre $\mathbf f$, se identifica con un $(A/\mathfrak L)$-submódulo de $\operatorname{gr}_{\mathfrak L}^{h+d}(N)$. Del hecho de que $z\in R$ mostraremos que se tiene incluso (con la identificación precedente)
\begin{env}[19.7.1.20]
\label{IV.19.7.1.20}
\[
  \bar z\in\operatorname{gr}_{\mathfrak L}^{h+d}(R,N).
\]
\end{env}

En efecto, se observó, en la demostración de (19.7.1.3), que, para todo ideal primo $\mathfrak r$ de $A$ tal que $\mathfrak p=\mathfrak r/\mathfrak R\in\operatorname{Ass}_{A/\mathfrak R}(A/\mathfrak L)$, la aplicación $\psi_{R_{\mathfrak r}}'$ es biyectiva y, por consiguiente, la relación (19.7.1.6), sustituyendo todos los módulos por sus localizados en $\mathfrak r$, es verdadera. Como $z$ pertenece a $R\cap(\mathfrak R^dN+\mathfrak L^{d+1}N)$, su imagen $z/1$ en $R_{\mathfrak r}$ pertenece a $(R_{\mathfrak r}\cap\mathfrak R_{\mathfrak r}^dN_{\mathfrak r})+\mathfrak L_{\mathfrak r}^{d+1}N_{\mathfrak r}$; por tanto, la imagen $\bar z/1$ de $\bar z$ en $\operatorname{gr}_{\mathfrak L_{\mathfrak r}}^{h+d}(N_{\mathfrak r})$ pertenece a
\[
  \operatorname{gr}_{\mathfrak L_{\mathfrak r}}^{h+d}(R_{\mathfrak r},N_{\mathfrak r})
  =(\operatorname{gr}_{\mathfrak L}^{h+d}(R,N))_{\mathfrak r}
  =(\operatorname{gr}_{\mathfrak L}^{h+d}(R,N))_{\mathfrak q},
\]
donde $\mathfrak q=\mathfrak r/\mathfrak L$. En otros términos, la imagen de $\bar z$ por la aplicación canónica
\[
  (\operatorname{gr}_{\mathfrak L}^{h+d}(N))_{\mathfrak q}\longrightarrow(\operatorname{gr}_{\mathfrak L}^{h+d}(M))_{\mathfrak q}
\]
es nula para todo $\mathfrak q\in\operatorname{Ass}_{A/\mathfrak L}(A/\mathfrak L)$. Como se vio en la demostración de (19.7.1.3) que
\[
  \operatorname{Ass}_{A/\mathfrak L}(\operatorname{gr}_{\mathfrak L}^{h+d}(M))\subset\operatorname{Ass}_{A/\mathfrak L}(A/\mathfrak L),
\]
se deduce (Bourbaki, \emph{Alg. comm.}, cap. IV, § 1, n\textsuperscript{o} 1, cor. 1 de la prop. 2 y n\textsuperscript{o} 3, prop. 7) que la imagen de $\bar z$ por la aplicación canónica $\operatorname{gr}_{\mathfrak L}^{h+d}(N)\to\operatorname{gr}_{\mathfrak L}^{h+d}(M)$ es nula, es decir, que se tiene la relación (19.7.1.20).

Sentado esto, por definición se puede escribir
\begin{env}[19.7.1.21]
\label{IV.19.7.1.21}
\[
  z\equiv\sum_{|\mathbf p|=h}c_{\mathbf p}\mathbf f^{\mathbf p}\pmod{\mathfrak R^{d+1}N}
\]
\end{env}
donde se pone, como de costumbre, $\mathbf p=(p_1,\ldots,p_n)$, $\mathbf f^{\mathbf p}=f_1^{p_1}f_2^{p_2}\cdots f_n^{p_n}$, $|\mathbf p|=p_1+\cdots+p_n$, y donde $c_{\mathbf p}\in\mathfrak R^dN$. Además, como $\psi_N$ identifica $\operatorname{gr}_{\mathfrak L}^*(N)$ con $(\operatorname{gr}_{\mathfrak R}^*(N)\otimes_{A/\mathfrak R}(A/\mathfrak L))[T_1,\ldots,T_n]$, los $c_{\mathbf p}$ están determinados módulo $\mathfrak R^d\mathfrak LN$, y si $\bar c_{\mathbf p}$ es la clase de $c_{\mathbf p}$ módulo $\mathfrak R^d\mathfrak LN$, se tiene, por la identificación precedente,
\begin{env}[19.7.1.22]
\label{IV.19.7.1.22}
\[
  \bar z=\sum_{|\mathbf p|=h}\bar c_{\mathbf p}T^{\mathbf p}.
\]
\end{env}
Utilizando (19.7.1.3), se deduce que necesariamente se tiene, para todo $\mathbf p$ tal que $|\mathbf p|=h$ (identificándose esta vez $\bar c_{\mathbf p}$ por $\psi_N$ con un elemento de $\operatorname{gr}_{\mathfrak R}^d(N)$),
\begin{env}[19.7.1.23]
\label{IV.19.7.1.23}
\[
  \bar c_{\mathbf p}\in\operatorname{gr}_{\mathfrak R}^d(R,N).
\]
\end{env}
Pero, puesto que $d<m$, la hipótesis de que $(*_d)$ es verdadera implica que $\bar c_{\mathbf p}$ pertenece a la imagen de $\psi_R'$, y por tanto que se puede suponer que $c_{\mathbf p}\in(R\cap\mathfrak R^dN)+\mathfrak R^d\mathfrak LN$. La relación (19.7.1.21) da entonces
\[
  z\in\mathfrak J^h(R\cap\mathfrak R^dN)+\mathfrak J^h\mathfrak R^d\mathfrak LN+\mathfrak R^{d+1}N
\]
y, como $\mathfrak J\subset\mathfrak L$, esto demuestra (19.7.1.19).

\oldpage[IV-4]{218}

En particular, se puede escribir
\[
  g=t+g_2
\]
con $t\in\mathfrak J^h(R\cap\mathfrak R^dN)$ y $g_2\in\mathfrak R^{d+1}N+\mathfrak R^d\mathfrak L^hN$; además, como $t\in R$ y $g\in R$, se tiene $g_2\in R$. Por otra parte, utilizando la hipótesis $h\geq m-d+1$, se ve que $t\in\mathfrak L^{m+1}N$. El elemento $g_2$ satisface, pues, todas las condiciones enunciadas al principio de (19.7.1.8), lo que termina la demostración de (19.7.1).

\begin{remarks}[19.7.2]
\label{IV.19.7.2}
\begin{enumerate}
  \item[(i)] Cuando se supone que $\operatorname{gr}_{\mathfrak R}^*(A)$ es un $(A/\mathfrak R)$-módulo \emph{plano}, la hipótesis de que la sucesión $\mathbf f$ es $(A/\mathfrak R)$-regular implica que también es $\operatorname{gr}_{\mathfrak R}^*(A)$-regular ($0$, 15.1.14). Obsérvese que este será el caso cuando $A$ y $A/\mathfrak R$ sean anillos regulares ($0$, 17.3.6): en efecto, $\mathfrak R$ es entonces un ideal regular de $A$ (19.1.2), y por tanto cuasirregular; localizando en los ideales maximales de $A$ que contienen $\mathfrak R$ y utilizando ($0$, 15.1.7), se ve que $\operatorname{gr}_{\mathfrak R}^*(A)$ es un $(A/\mathfrak R)$-módulo \emph{proyectivo}.
  \item[(ii)] Sería interesante poder demostrar la última conclusión de (19.7.1) suponiendo solamente que existe una $A$-álgebra $B$ que es un anillo noetheriano, para el cual $\mathfrak JB$ está contenido en el radical de $B$, y que $M$ es un $B$-módulo de tipo finito.
\end{enumerate}
\end{remarks}

\begin{corollary}[19.7.3]
\label{IV.19.7.3}
Sean $A$ un anillo local noetheriano, $\mathfrak m$ su ideal maximal, $\mathfrak R$ un ideal de $A$ tal que el anillo $A/\mathfrak R$ sea regular y de dimensión $n$, y $M$ un $A$-módulo de tipo finito. Las condiciones siguientes son equivalentes:
\begin{enumerate}
  \item[a)] $\operatorname{gr}_{\mathfrak R}^*(M)$ es un $(A/\mathfrak R)$-módulo plano.
  \item[b)] Si $P(T)$ y $Q(T)$ son las series de Poincaré de los $(A/\mathfrak m)$-módulos graduados $\operatorname{gr}_{\mathfrak m}^*(M)$ y $\operatorname{gr}_{\mathfrak R}^*(M)\otimes_{A/\mathfrak R}(A/\mathfrak m)$, se tiene
\end{enumerate}
\begin{equation}
  Q(T)=P(T)(1-T)^n.
  \tag{19.7.3.1}\label{IV.19.7.3.1}
\end{equation}

Sea $\mathbf f=(f_1,\ldots,f_n)$ una sucesión de elementos de $A$ tales que las imágenes de los $f_i$ en $A/\mathfrak R$ formen un sistema regular de parámetros de $A/\mathfrak R$ ($0$, 17.1.6), de modo que $\mathbf f$ es una sucesión $(A/\mathfrak R)$-regular; además, si $\mathfrak J$ es el ideal engendrado por $\mathbf f$, se tiene $\mathfrak J+\mathfrak R=\mathfrak m$, y como $A/\mathfrak m$ es un cuerpo, $\operatorname{gr}_{\mathfrak m}^*(M)$ es un $(A/\mathfrak m)$-módulo plano. Como $A$ es noetheriano y $M$ un $A$-módulo de tipo finito, se puede aplicar la equivalencia de $a)$ y $b)$ en (19.7.1). Además, como $A/\mathfrak m$ es un cuerpo y $\psi_M$ es sobreyectivo, el hecho de que $\psi_M$ sea biyectivo equivale a decir que, para todo $h\geq0$, $\operatorname{gr}_{\mathfrak m}^h(M)$ y el submódulo de los elementos de grado $h$ de $(\operatorname{gr}_{\mathfrak R}^*(M)\otimes_{A/\mathfrak R}(A/\mathfrak m))[T_1,\ldots,T_n]$ tienen el mismo rango sobre $A/\mathfrak m$; pero, para el segundo de estos módulos, el rango en cuestión es evidentemente igual a
\[
  \sum_{r=0}^h\binom{r+n-1}{n-1}\operatorname{rg}(\operatorname{gr}_{\mathfrak R}^{h-r}(M)\otimes_{A/\mathfrak R}(A/\mathfrak m)).
\]
Como
\[
  (1-T)^{-n}=\sum_{r=0}^{\infty}\binom{r+n-1}{n-1}T^r,
\]
esto demuestra que la relación (19.7.3.1) es necesaria y suficiente para que $\psi_M$ sea biyectivo.
\end{corollary}

\begin{corollary}[19.7.4]
\label{IV.19.7.4}
Sean $X$ un preesquema localmente noetheriano, $Y$ un subpreesquema cerrado de $X$, regular y conexo, y $\sh F$ un $\sh O_X$-módulo coherente normalmente plano (6.10.1) a lo largo de $Y$. Entonces existe una serie formal $R\in\mathbb{Q}[[T]]$ (independiente de $x\in Y$) tal que, para todo $x\in Y$, la serie de Poincaré de $\operatorname{gr}_{\mathfrak m_x}^*(\sh F_x)$ viene dada por la fórmula
\begin{equation}
  P_x(\sh F)(T)=R(T)(1-T)^{-n}\qquad\text{donde }n=\dim(\sh O_{Y,x}).
  \tag{19.7.4.1}\label{IV.19.7.4.1}
\end{equation}

\oldpage[IV-4]{219}

Sea $\sh J$ el ideal coherente de $\sh O_X$ que define $Y$; la hipótesis sobre $\sh F$ significa que $\operatorname{gr}_{\sh J}^*(\sh F)$ es un $\sh O_Y$-módulo plano. Para todo $x\in Y$ se puede, pues, aplicar (19.7.3) tomando $A=\sh O_{X,x}$, $\mathfrak R=\sh J_x$ y $M=\sh F_x$; cumpliéndose por hipótesis la condición $a)$, la serie de Poincaré $P_x(\sh F)(T)$ es entonces igual a $R_x(T)(1-T)^{-n}$, donde $R_x(T)$ es la serie de Poincaré de $\operatorname{gr}_{\sh J_x}^*(\sh F_x)\otimes_{\sh O_{Y,x}}k(x)$. Ahora bien, cada uno de los $\sh O_Y$-módulos $\sh J^h\sh F/\sh J^{h+1}\sh F$ es plano y coherente por hipótesis, y por tanto localmente libre (2.1.12); como se supone que $Y$ es conexo, cada uno de los $\sh J^h\sh F/\sh J^{h+1}\sh F$ tiene rango constante en $Y$, lo que demuestra que la serie de Poincaré $R_x(T)$ es independiente de $x\in Y$.
\end{corollary}

\begin{corollary}[19.7.5]
\label{IV.19.7.5}
Sean $X$ un preesquema localmente noetheriano, $Y$ un subpreesquema cerrado de $X$ y $Z$ un subpreesquema cerrado de $Y$; se supone que $Y$ y $Z$ son regulares. Sea $\sh F$ un $\sh O_X$-módulo coherente.
\begin{enumerate}
  \item[(i)] Si $\sh F$ es normalmente plano a lo largo de $Y$ en los puntos de $Z$ (cf. (11.3.4) para la definición de la planitud normal en un punto), entonces $\sh F$ es normalmente plano a lo largo de $Z$.
  \item[(ii)] Supongamos además que $Z$ es conexo y que $\sh O_X$ es normalmente plano a lo largo de $Y$ en los puntos de $Z$ (lo que ocurre, en particular, si $X$ es regular en los puntos de $Z$). Si $\sh F$ es normalmente plano a lo largo de $Z$ y existe un punto $z\in Z$ tal que $\sh F$ es normalmente plano a lo largo de $Y$ en el punto $z$, entonces $\sh F$ es normalmente plano a lo largo de $Y$ en los puntos de $Z$.
\end{enumerate}

Sean $\sh K$ y $\sh L\supset\sh K$ los ideales coherentes de $\sh O_X$ que definen respectivamente $Y$ y $Z$. Para todo punto $z\in Z$, la inmersión canónica $Z\to Y$ es regular en el punto $z$, puesto que $Y$ y $Z$ son regulares en ese punto (19.1.1); en otros términos, existe una sucesión $(f_i)_{1\leq i\leq n}$ de elementos de $\sh O_{X,z}$ que es $\sh K_z$-regular y tal que $\sh L_z=\sh K_z+\sum_i f_i\sh O_{X,z}$.
\begin{enumerate}
  \item[(i)] La hipótesis significa que $\operatorname{gr}_{\sh K_z}^*(\sh F_z)$ es un $\sh O_{Y,z}$-módulo plano; implica, pues, que $\operatorname{gr}_{\sh L_z}^*(\sh F_z)$ es un $\sh O_{Z,z}$-módulo plano, en virtud de que en (19.7.1), $a)$ implica $b)$.
  \item[(ii)] La hipótesis suplementaria sobre $X$ implica que la sucesión $(f_i)$ es $\operatorname{gr}_{\sh K_z}^*(\sh O_{X,z})$-regular\footnote{El texto impreso dice $\operatorname{gr}_{\sh K_z}^*(\sh F_z)$; (19.7.2) y la hipótesis suplementaria sobre $X$ se aplican a $\sh O_{X,z}$.} en todo punto $z\in Z$ (19.7.2). Por otra parte, como $Z$ es regular y conexo, es íntegro; si $\operatorname{gr}_{\sh K_z}^*(\sh F_z)$ es $\sh O_{Y,z}$-plano en un punto $z\in Z$, entonces, como el punto genérico $\zeta$ de $Z$ es una generalización de $z$, se concluye que $\operatorname{gr}_{\sh K_\zeta}^*(\sh F_\zeta)$ es $\sh O_{Y,\zeta}$-plano ($0_{\mathrm I}$, 6.3.1). Se deduce que, si en todo punto $z'$ de $Z$, $\operatorname{gr}_{\sh L_{z'}}^*(\sh F_{z'})$ es un $\sh O_{Z,z'}$-módulo plano, entonces $\operatorname{gr}_{\sh K_{z'}}^*(\sh F_{z'})$ es un $\sh O_{Y,z'}$-módulo plano, pues se puede aplicar la equivalencia de $a)$ y $d)$ en (19.7.1).
\end{enumerate}
\end{corollary}

\subsection{Propiedades del paso al límite proyectivo.}
\label{subsection:IV.19.8}
\label{IV.19.8}

En este número, las notaciones y convenciones sobre los límites proyectivos son las de (8.5.1) y (8.8.1).

\begin{proposition}[19.8.1]
\label{IV.19.8.1}
Supongamos que los morfismos de transición $S_\mu\to S_\lambda$ $(\lambda\leq\mu)$ sean planos y, además, que se cumpla una de las dos hipótesis siguientes:
\begin{enumerate}
  \item[$1^\circ$] Los preesquemas $S_\lambda$ son localmente noetherianos.
  \item[$2^\circ$] Los morfismos de transición $S_\mu\to S_\lambda$ son sobreyectivos (y, por tanto, fielmente planos).
\end{enumerate}

\oldpage[IV-4]{220}

\emph{Bajo estas condiciones:}
\begin{enumerate}
  \item[(i)] Supongamos que $S_\alpha$ sea cuasicompacto. Sea $\sh F_\alpha$ un $\sh O_{S_\alpha}$-módulo cuasicoherente, que suponemos además de tipo finito cuando se cumple la hipótesis $1^\circ$. Sea $(f_{i\alpha})_{1\leq i\leq n}$ una sucesión finita de secciones de $\sh F_\alpha$ sobre $S_\alpha$. Para que la sucesión de secciones $f_i$ del $\sh O_S$-módulo $\sh F$ sobre $S$, correspondiente a las $f_{i\alpha}$ $(1\leq i\leq n)$, sea $\sh F$-regular, es necesario y suficiente que exista $\lambda\geq\alpha$ tal que la sucesión de secciones $f_{i\lambda}$ de $\sh F_\lambda$ sobre $S_\lambda$ sea $\sh F_\lambda$-regular.
  \item[(ii)] Supongamos que $X_\alpha$ sea cuasicompacto, y sea $j_\alpha:X_\alpha\to S_\alpha$ una inmersión, que suponemos localmente de presentación finita cuando se cumple la hipótesis $2^\circ$. Entonces, para que la inmersión correspondiente $j:X\to S$ sea regular, es necesario y suficiente que exista $\lambda\geq\alpha$ tal que $j_\lambda:X_\lambda\to S_\lambda$ sea regular.
\end{enumerate}
\end{proposition}

\begin{enumerate}
  \item[(i)] Si $p_\lambda:S\to S_\lambda$ es la proyección canónica, se tiene
  \[
    \sh F/\Bigl(\sum_{j=1}^{i-1}f_j\sh F\Bigr)
    =p_\lambda^*\Bigl(\sh F_\lambda/\Bigl(\sum_{j=1}^{i-1}f_{j\lambda}\sh F_\lambda\Bigr)\Bigr),
  \]
  por tanto, nos reducimos al caso $n=1$, en cuyo caso se suprime el índice $i$. La suficiencia de la condición resulta de que $p_\lambda$ es plano (8.3.8) y de ($0$, 15.1.5). En el caso $2^\circ$, $p_\lambda$ es fielmente plano (8.3.8), y la necesidad de la condición resulta también de ($0$, 15.2.5), con $\lambda=\alpha$. En el caso $1^\circ$, denotemos por $\sh N_\lambda$ (resp. $\sh N$) el núcleo del endomorfismo $\sh F_\lambda\to\sh F_\lambda$ (resp. $\sh F\to\sh F$) dado por la multiplicación por $f_\lambda$ (resp. por $f$); como $\sh F_\lambda$ es coherente por hipótesis, también lo es $\sh N_\lambda$, de modo que la hipótesis $\sh N=0$ implica $\sh N_\lambda=0$ para algún $\lambda\geq\alpha$ en virtud de (8.5.8, (ii)).
  \item[(ii)] Como $p_\lambda$ es plano, la suficiencia de la condición resulta de (19.1.5, (ii)). Para demostrar su necesidad, observemos que, como $j_\alpha(X_\alpha)$ es cuasicompacto, está contenido en un abierto cuasicompacto de $S_\alpha$, y podemos reducirnos al caso en que $S_\alpha$ también sea cuasicompacto y $j_\alpha$ una inmersión cerrada; así, la imagen de $X_\alpha$ está definida por un ideal cuasicoherente $\sh J_\alpha$ de $\sh O_{S_\alpha}$, que además podemos suponer de tipo finito en los casos $1^\circ$ y $2^\circ$ ($j_\alpha$ es localmente de presentación finita en el caso $2^\circ$); la imagen de $X_\lambda$ (resp. $X$) en $S_\lambda$ (resp. $S$) está entonces definida por $\sh J_\lambda=\sh J_\alpha\sh O_{S_\lambda}$ (resp. $\sh J=\sh J_\alpha\sh O_S$), que sigue siendo de tipo finito. Además, teniendo en cuenta (8.2.11), podemos suponer que $\sh J$ está engendrado por una sucesión regular $(f_i)_{1\leq i\leq n}$ de secciones sobre $S$, que por tanto definen un homomorfismo sobreyectivo $u:\sh O_S^n\to\sh J$. Teniendo en cuenta (8.5.2, (i)) y (8.5.7), existen $\lambda\geq\alpha$ y un homomorfismo sobreyectivo $u_\lambda:\sh O_{S_\lambda}^n\to\sh J_\lambda$ tales que $u=p_\lambda^*(u_\lambda)$; por tanto, los $f_i$ son las imágenes canónicas de secciones $f_{i\lambda}$ de $\sh J_\lambda$ sobre $S_\lambda$ que engendran este ideal. En virtud de (i), existe entonces $\mu\geq\lambda$ para el cual la sucesión $(f_{i\mu})$ es $\sh O_{S_\mu}$-regular, y por tanto la inmersión $j_\mu$ es regular.
\end{enumerate}

\begin{proposition}[19.8.2]
\label{IV.19.8.2}
Sea $X_\alpha$ un preesquema $S_\alpha$ cuasicompacto y localmente de presentación finita.
\begin{enumerate}
  \item[(i)] Sean $\sh F_\alpha$ un $\sh O_{X_\alpha}$-módulo cuasicoherente de presentación finita y $(f_{i\alpha})_{1\leq i\leq n}$ una sucesión de secciones de $\sh O_{X_\alpha}$ sobre $X_\alpha$. Para que la sucesión de secciones $f_i$ de $\sh O_X$ sobre $X$, correspondiente a las $f_{i\alpha}$, sea $\sh F$-transversalmente regular relativamente a $S$ (19.2.1), es necesario y suficiente que exista $\lambda\geq\alpha$ tal que la sucesión de secciones $f_{i\lambda}$ de $\sh O_{X_\lambda}$ sobre $X_\lambda$ sea $\sh F_\lambda$-transversalmente regular relativamente a $S_\lambda$.
  \item[(ii)] Sea $Y_\alpha$ un preesquema $S_\alpha$ cuasicompacto, y sea $j_\alpha:Y_\alpha\to X_\alpha$ una inmersión $S_\alpha$ localmente

\oldpage[IV-4]{221}

  de presentación finita. Para que la inmersión $S$ correspondiente $j:Y\to X$ sea transversalmente regular relativamente a $S$, es necesario y suficiente que exista $\lambda\geq\alpha$ tal que $j_\lambda:Y_\lambda\to X_\lambda$ sea transversalmente regular relativamente a $S_\lambda$.
\end{enumerate}
\end{proposition}

\begin{enumerate}
  \item[(i)] La suficiencia de la condición resulta de ($0$, 15.1.15). Para demostrar que es necesaria, bastará mostrar que todo punto $x\in X$ posee un entorno abierto $V(x)$ para el cual existen un índice $\lambda=\lambda(x)\geq\alpha$ y un entorno abierto $V_\lambda$ de la imagen $x_\lambda$ de $x$ en $X_\lambda$, tales que $V(x)$ sea la imagen inversa de $V_\lambda$ y la sucesión de las restricciones a $V_\lambda$ de los $f_{i\lambda}$ sea $(\sh F_\lambda|V_\lambda)$-transversalmente regular relativamente a $S_\lambda$; en efecto, bastará entonces aplicar (8.3.4) al abierto $U_\lambda$ de $X_\lambda$ definido como el mayor abierto en el que las restricciones a $U_\lambda$ de los $f_{i\lambda}$ forman una sucesión $(\sh F_\lambda|U_\lambda)$-transversalmente regular relativamente a $S_\lambda$.

  Sea $s$ la imagen de $x$ en $S$, y para todo $\lambda\geq\alpha$ sea $s_\lambda$ la imagen de $s$ en $S_\lambda$. Entonces la fibra $X_s$ es el límite proyectivo de las fibras $(X_\lambda)_{s_\lambda}$ y es localmente noetheriana, puesto que el morfismo $X\to S$ es localmente de presentación finita; como además $(X_\alpha)_{s_\alpha}$ es cuasicompacto y los morfismos de transición $(X_\mu)_{s_\mu}\to(X_\lambda)_{s_\lambda}$ son planos (pues también lo es $\Spec(k(s_\mu))\to\Spec(k(s_\lambda))$), podemos aplicar a las secciones $f_i\otimes1$ de $\sh F_s=\sh F\otimes_{\sh O_S}k(s)$ sobre $X_s$ el resultado de (19.8.1, (i)); de aquí se obtiene un $\lambda\geq\alpha$ tal que las secciones $f_{i\lambda}\otimes1$ de $(\sh F_\lambda)_{s_\lambda}=\sh F_\lambda\otimes_{\sh O_{S_\lambda}}k(s_\lambda)$ sobre $(X_\lambda)_{s_\lambda}$ forman una sucesión $(\sh F_\lambda)_{s_\lambda}$-regular. Además, por (11.2.6), podemos suponer que $\sh F_\lambda$ es $S_\lambda$-plano en el punto $x_\lambda$. Se deduce entonces de (11.3.8) que existe un entorno abierto $V_\lambda$ de $x_\lambda$ en $X_\lambda$ que satisface las condiciones requeridas.

  \item[(ii)] La suficiencia de la condición resulta de (19.2.7, (ii)). Para demostrar su necesidad, basta todavía mostrar que todo punto $x\in j(Y)$ admite un entorno abierto $V(x)$ para el cual existen un índice $\lambda=\lambda(x)\geq\alpha$ y un entorno abierto $V_\lambda$ de la imagen $x_\lambda$ de $x$ en $X_\lambda$, tales que $V(x)$ sea la imagen inversa de $V_\lambda$ y la restricción $j_\lambda^{-1}(V_\lambda)\to V_\lambda$ de $j_\lambda$ sea transversalmente regular relativamente a $S_\lambda$. Sea $s$ la imagen de $x$ en $S$, y para todo $\lambda\geq\alpha$ sea $s_\lambda$ la imagen de $s$ en $S_\lambda$. Como $X_s$ (resp. $Y_s$) es el límite proyectivo de las fibras $(X_\lambda)_{s_\lambda}$ (resp. $(Y_\lambda)_{s_\lambda}$), podemos utilizar, como en (i), (19.8.1, (ii)); y como, por hipótesis, la inmersión $Y_s\to X_s$ es regular, existe un índice $\lambda(s)$ tal que, para $\lambda\geq\lambda(s)$, la inmersión $(Y_\lambda)_{s_\lambda}\to(X_\lambda)_{s_\lambda}$ es regular. Además, en virtud de (19.2.4), existe un entorno abierto cuasicompacto $W(x)$ de $x$ en $X$ tal que los morfismos estructurales $W(x)\to S$ y $Y\cap W(x)\to S$ son planos; aplicando (11.2.6) y (8.2.11), podemos suponer que existen $\lambda\geq\lambda(s)$ y un entorno abierto cuasicompacto $W(x_\lambda)$ de $x_\lambda$ en $X_\lambda$, tales que $W(x)$ sea la imagen inversa de $W(x_\lambda)$ y las restricciones a $W(x_\lambda)$ y a $Y_\lambda\cap W(x_\lambda)$, respectivamente, de los morfismos estructurales $X_\lambda\to S_\lambda$ y $Y_\lambda\to S_\lambda$ sean planas. Concluimos, pues, de (19.2.4) que existe un entorno abierto $V_\lambda$ de $x_\lambda$ que responde a la cuestión.
\end{enumerate}

\begin{remark}[19.8.3]
\label{IV.19.8.3}
Las demostraciones precedentes muestran que los enunciados de (19.8.1) y (19.8.2) siguen siendo válidos cuando se sustituyen las condiciones relativas a la regularidad o a la regularidad transversal en todo $X$ (resp. $X_\lambda$) por esas condiciones en un entorno de un punto $x\in X$ (resp. en un entorno de $x_\lambda$, proyección de $x$).
\end{remark}

\subsection{\texorpdfstring{Sucesiones $\sh F$-regulares y profundidad.}{Sucesiones F-regulares y profundidad.}}
\label{subsection:IV.19.9}
\label{IV.19.9}

\oldpage[IV-4]{222}

\begin{env}[19.9.1]
\label{IV.19.9.1}
Recordemos que, si $X$ es un preesquema localmente noetheriano, $\sh F$ un
$\sh O_X$-módulo coherente y $T$ una parte de $X$, se define (5.10.1)
\begin{equation}
  \operatorname{prof}_T(\sh F)=\inf_{x\in T}\operatorname{prof}(\sh F_x).
  \tag{19.9.1.1}\label{IV.19.9.1.1}
\end{equation}

Asimismo, para todo punto $t\in T$ se define
\begin{equation}
  \operatorname{prof}_{T,t}(\sh F)=\inf_{x\in T\cap\Spec(\sh O_{X,t})}\operatorname{prof}(\sh F_x),
  \tag{19.9.1.2}\label{IV.19.9.1.2}
\end{equation}
y se dice que $\operatorname{prof}_{T,t}(\sh F)$ es la
\emph{$T$-profundidad de $\sh F$ en el punto $t$}; es claro que se tiene
\begin{equation}
  \operatorname{prof}_T(\sh F)=\inf_{t\in T}\operatorname{prof}_{T,t}(\sh F).
  \tag{19.9.1.3}\label{IV.19.9.1.3}
\end{equation}
\end{env}

\begin{lemma}[19.9.2]
\label{IV.19.9.2}
Sean $A$ un anillo noetheriano, $M$ un $A$-módulo de tipo finito y
$\mathfrak J$ un ideal de $A$. Para que, para todo ideal primo
$\mathfrak p\supset\mathfrak J$, se tenga
$\operatorname{prof}_{A_{\mathfrak p}}(M_{\mathfrak p})\geq r$, es necesario
y suficiente que exista una sucesión $M$-regular de $r$ elementos
pertenecientes a $\mathfrak J$.
\end{lemma}

La suficiencia de la condición se deduce inmediatamente de ($0$, 16.4.5);
probemos su necesidad. Razonamos por inducción sobre $r$ (no hay nada que
demostrar si $r=0$): puesto que
$\operatorname{prof}_{A_{\mathfrak p}}(M_{\mathfrak p})\geq r-1$ para todo
$\mathfrak p\supset\mathfrak J$, por hipótesis existe una sucesión
$M$-regular $(g_i)_{1\leq i\leq r-1}$ formada por elementos de $\mathfrak J$.
Pongamos $N=M/(g_1M+\cdots+g_{r-1}M)$; para todo
$\mathfrak p\supset\mathfrak J$, las imágenes de los $g_i$ en $A_{\mathfrak p}$
pertenecen por tanto al ideal maximal $\mathfrak pA_{\mathfrak p}$ y forman
una sucesión $M_{\mathfrak p}$-regular ($0$, 15.1.14); se concluye entonces de
($0$, 16.4.6) que
$\operatorname{prof}_{A_{\mathfrak p}}(N_{\mathfrak p})
=\operatorname{prof}_{A_{\mathfrak p}}(M_{\mathfrak p})-(r-1)\geq1$, y se ve
que todo se reduce a demostrar el lema para $r=1$. Pero la hipótesis significa
entonces que, para todo $\mathfrak p\supset\mathfrak J$,
$\mathfrak pA_{\mathfrak p}$ no está asociado a $M_{\mathfrak p}$
($0$, 16.4.6), por lo que (Bourbaki, \emph{Alg. comm.}, cap. IV, § 1,
n\textsuperscript{o} 2, prop. 5), $\mathfrak p$ no está asociado a $M$. En
otras palabras, $\mathfrak J$ no está contenido en ninguno de los ideales
primos de $\operatorname{Ass}(M)$, y por tanto no está contenido en su unión
(Bourbaki, \emph{Alg. comm.}, cap. II, § 1, n\textsuperscript{o} 1, prop. 2).
Pero como esa unión es el conjunto de los elementos que no son $M$-regulares
(Bourbaki, \emph{Alg. comm.}, cap. IV, § 1, n\textsuperscript{o} 1, cor. 2 de
la prop. 2), queda demostrado el lema.

\begin{proposition}[19.9.3]
\label{IV.19.9.3}
Sean $X$ un preesquema localmente noetheriano, $\sh F$ un
$\sh O_X$-módulo coherente, $Y$ una parte cerrada de $X$, $y$ un punto de
$Y$ y $n$ un entero $>0$. Las condiciones siguientes son equivalentes:
\begin{enumerate}
  \item[a)] Existe una vecindad abierta $U$ de $y$ en $X$ tal que
  $\operatorname{prof}_{Y\cap U}(\sh F|U)\geq n$.
  \item[b)] Existe una vecindad abierta $U$ de $y$ en $X$ y una sucesión
  $(\sh F|U)$-regular de $n$ secciones $f_i$ de $\sh O_U$ sobre $U$, tales que
  $f_i(x)=0$ ($0_{\mathrm I}$, 5.5.1) para $1\leq i\leq n$ en todo punto
  $x\in U\cap Y$.
  \item[c)] Se tiene $\operatorname{prof}_{Y,y}(\sh F)\geq n$.
\end{enumerate}
Además, el conjunto de los $y\in Y$ que verifican estas condiciones es abierto
en $Y$.
\end{proposition}

La última afirmación se deduce trivialmente de a). La equivalencia de a) y b)
se sigue de (19.9.2), aplicado a una vecindad abierta afín de $y$ en $X$, y
de ($0$, 15.1.14). Es claro que a) implica c), pues toda vecindad abierta de
$y$ en $X$ contiene $\Spec(\sh O_{X,y})$; demostremos recíprocamente que c)
implica b). De c) y de la definición (19.9.1.1) se deduce que se puede aplicar
(19.9.2) al anillo $\sh O_{X,y}$, al $\sh O_{X,y}$-módulo $\sh F_y$ y al ideal
$\sh J_y$ de $\sh O_{X,y}$, por lo que existe una sucesión $\sh F_y$-regular
$(t_i)_{1\leq i\leq n}$ de $n$ elementos de $\sh J_y$. Existe entonces una
vecindad abierta $U$ de $y$ en $X$ y una sucesión $(f_i)_{1\leq i\leq n}$ de
secciones de $\sh J$ sobre $U$ tales que los $t_i$ son los gérmenes de los
$f_i$ en el punto $y$; usando ($0$, 15.2.4), se deduce que existe una
vecindad abierta $U'\subset U$ de $y$ en $X$ tal que los $f_i|U'$ forman una
sucesión $(\sh F|U')$-regular.

\oldpage[IV-4]{223}

\begin{corollary}[19.9.4]
\label{IV.19.9.4}
Con las notaciones de (19.9.3), la función
$y\mapsto\operatorname{prof}_{Y,y}(\sh F)$ es semicontinua inferiormente en
$Y$.
\end{corollary}

\begin{proposition}[19.9.5]
\label{IV.19.9.5}
Con las notaciones de (19.9.3), sean $X'$ un preesquema localmente noetheriano,
$g:X'\to X$ un morfismo plano, $Y'=g^{-1}(Y)$ y $\sh F'=g^*(\sh F)$. Entonces,
para todo punto $y'\in Y'$ situado sobre $y$, se tiene
$\operatorname{prof}_{Y',y'}(\sh F')=\operatorname{prof}_{Y,y}(\sh F)$.
\end{proposition}

En efecto, para todo $z'\in\Spec(\sh O_{X',y'})$, se deduce de (6.3.1) que
$\operatorname{prof}(\sh F'_{z'})\geq\operatorname{prof}(\sh F_{g(z')})$,
y si $z''$ es un punto maximal de la fibra $g^{-1}(g(z'))$, generalización de
$z'$ (y por tanto perteneciente también a $\Spec(\sh O_{X',y'})$), se tiene
(\emph{loc. cit.})
$\operatorname{prof}(\sh F'_{z''})=\operatorname{prof}(\sh F_{g(z')})$.
La proposición se deduce entonces de
$g(\Spec(\sh O_{X',y'}))=\Spec(\sh O_{X,g(y')})$, puesto que $g$ es plano
(2.3.4).

\begin{proposition}[19.9.6]
\label{IV.19.9.6}
Sean $f:X\to S$ un morfismo localmente de presentación finita, $Y$ una parte
cerrada de $X$, $\sh F$ un $\sh O_X$-módulo de presentación finita y
$f$-plano, y $n$ un entero $>0$. Para todo punto $y\in Y$, las condiciones
siguientes son equivalentes (denotamos por $X_{f(y)}$ la fibra de $f$ en el
punto $f(y)$, por $\sh F_{f(y)}$ la imagen inversa de $\sh F$ en
$X_{f(y)}$, y ponemos $Y_{f(y)}=X_{f(y)}\cap Y$):
\begin{enumerate}
  \item[a)] Se tiene
  $\operatorname{prof}_{Y_{f(y)},y}(\sh F_{f(y)})\geq n$.
  \item[b)] Existe una vecindad abierta $U$ de $y$ en $X$ y una sucesión
  $(g_i)_{1\leq i\leq n}$ de $n$ secciones de $\sh O_X$ sobre $U$ que es
  transversalmente $\sh F$-regular relativamente a $S$ (cf. (19.2.1) para la
  terminología), y tal que $g_i(x)=0$ para $1\leq i\leq n$ y para todo
  $x\in U\cap Y$.
\end{enumerate}
El conjunto $V_n$ de los $y\in Y$ que verifican estas condiciones es abierto en
$Y$; si además $Y$ es localmente constructible en $X$, $V_n$ es
retrocompacto en $X$ ($0_{\mathrm{III}}$, 9.1.1).
\end{proposition}

En virtud de (19.9.3), la condición a) equivale a la existencia de una
vecindad abierta $U$ de $y$ en $X$ y de una sucesión
$(\sh F_{f(y)}|(U\cap X_{f(y)}))$-regular $(h_i)_{1\leq i\leq n}$ de secciones
de $\sh O_{X_{f(y)}}$ sobre $U\cap X_{f(y)}$ tales que $h_i(z)=0$ para todo
$z\in U\cap Y_{f(y)}$. Si se designa todavía por $Y$ un subpreesquema cerrado
de $X$ cuyo espacio subyacente sea $Y$, y por $\sh J$ el ideal cuasi-coherente
de $\sh O_X$ que define $Y$, también se puede decir que
$h_i\in\sh J_{f(y)}$ para $1\leq i\leq n$. Como se puede sustituir $U$ por una
vecindad afín más pequeña de $y$, se puede suponer que los $h_i$ son las
imágenes de una sucesión $(g_i)_{1\leq i\leq n}$ de secciones de $\sh J$ sobre
$U$, de modo que los gérmenes $(g_i)_y$ pertenecen al ideal maximal
$\mathfrak m_y$ de $\sh O_{X,y}$. La equivalencia de a) y b) se deduce entonces
de (11.3.8), que prueba también que el conjunto $V_n$ es abierto en $Y$.

Queda por demostrar la última afirmación; como se trata todavía de propiedades
locales sobre $X$, podemos suponer que $S=\Spec(A)$ es afín y que $X$ es de
presentación finita sobre $S$; existe entonces un subanillo noetheriano $A_0$
de $A$, un preesquema $X_0$ de tipo finito sobre $S_0=\Spec(A_0)$ y un
$\sh O_{X_0}$-módulo coherente $\sh F_0$ tales que
$X=X_0\times_{S_0}S$ y
$\sh F=\sh F_0\otimes_{\sh O_{X_0}}\sh O_X$ (8.9.1); además se puede suponer
que, si $f_0:X_0\to S_0$ es el morfismo estructural, $\sh F_0$ es
$f_0$-plano (11.2.6). Puesto que $Y$ es constructible, se puede suponer además
(8.3.11) que $Y=p^{-1}(Y_0)$, donde $p:X\to X_0$ es la proyección canónica.
Entonces, para todo $y\in Y$, se deduce de (19.9.5) y de la transitividad de
las fibras que
\[
  \operatorname{prof}_{Y_{f(y)},y}(\sh F_{f(y)})
  =\operatorname{prof}_{(Y_0)_{f_0(y_0)},y_0}((\sh F_0)_{f_0(y_0)}),
\]
poniendo $y_0=p(y)$; si $V_n^0$ es el conjunto de los $y_0\in Y_0$ para los
cuales el segundo miembro de esta relación es $\geq n$, se tiene por tanto
$V_n=p^{-1}(V_n^0)$. Como $X_0$ es noetheriano, $V_n^0$ es un abierto
retrocompacto en $X_0$, y por consiguiente (cf. la demostración de (1.8.2))
$V_n$ es retrocompacto en $X$.

\oldpage[IV-4]{224}

\begin{corollary}[19.9.7]
\label{IV.19.9.7}
Bajo las hipótesis generales de (19.9.6), la función
$y\mapsto\operatorname{prof}_{Y_{f(y)},y}(\sh F_{f(y)})$ es semicontinua
inferiormente en $Y$. Si además $Y$ es localmente constructible, esta función
es localmente constructible.
\end{corollary}

\begin{proposition}[19.9.8]
\label{IV.19.9.8}
Sean $f:X\to S$ un morfismo localmente de presentación finita, $Z$ una parte
cerrada de $X$, $\sh F$ un $\sh O_X$-módulo de presentación finita y
$f$-plano, y $\sh G$ un $\sh O_S$-módulo cuasi-coherente. Supongamos que, para
todo $x\in Z$, se tiene
$\operatorname{prof}((\sh F_{f(x)})_x)\geq1$
(resp. $\operatorname{prof}((\sh F_{f(x)})_x)\geq2$). Entonces, si
$j:X-Z\to X$ es la inmersión canónica y se pone
$\sh H=\sh F\otimes_{\sh O_X}f^*(\sh G)$, el homomorfismo canónico
$\rho_{\sh H}:\sh H\to j_*(j^*(\sh H))$ relativo a $j$ es inyectivo (resp.
bijectivo).
\end{proposition}

La cuestión es evidentemente local sobre $X$ y $S$, y basta demostrar la
proposición en una vecindad de un punto $x\in Z$. En otras palabras, podemos
limitarnos al caso en que $X$ y $S$ son afines. Gracias a (19.9.6), podemos
suponer que existe una sección transversalmente $\sh F$-regular $g_1$ de
$\sh O_X$ sobre $X$ relativamente a $S$ (resp. una sucesión transversalmente
$\sh F$-regular $(g_1,g_2)$ de dos secciones de $\sh O_X$ sobre $X$ relativamente
a $S$) tal que, si $Z'$ es el conjunto de los $x'\in X$ para los cuales
$g_1(x')=0$ (resp. $g_1(x')=g_2(x')=0$), se tenga $Z\subset Z'$.

Volviendo a la demostración de (19.9.8), por (19.9.6) se tiene todavía
$\operatorname{prof}((\sh F_{f(x)})_x)\geq1$
(resp. $\operatorname{prof}((\sh F_{f(x)})_x)\geq2$) para todo $x\in Z'$.
Si la proposición está demostrada sustituyendo el par $(X,Z)$ por
$(X,Z')$ y $(X-Z,Z'\cap(X-Z))$, es inmediato que también lo estará para
$(X,Z)$; podemos, pues, limitar la demostración al caso de $X$ y $Z'$. Es
decir, podemos suponer que (19.9.8) se demuestra cuando
$S=\Spec(A)$, $X=\Spec(B)$, donde $B$ es una $A$-álgebra de presentación
finita, $Z=V(\mathfrak J)$, donde $\mathfrak J$ es un ideal de tipo finito del
anillo $B$, $\sh F=\widetilde M$ y $\sh G=\widetilde N$, donde $N$ es un
$A$-módulo y $M$ un $B$-módulo de presentación finita que es un $A$-módulo
plano. Además podemos reducirnos al caso en que $N$ es un $A$-módulo de
presentación finita. En efecto, $N$ es límite inductivo de un sistema inductivo
filtrante de $A$-módulos de presentación finita $N_\alpha$ (por doble límite
inductivo); si se pone $\sh G_\alpha=\widetilde{N_\alpha}$, $\sh H$ es el
límite inductivo de los $\sh H_\alpha=\sh F\otimes_{\sh O_X}f^*(\sh G_\alpha)$;
como $j_*$ y $j^*$ conmutan con los límites inductivos y
$\varinjlim$ es un funtor exacto en la categoría de módulos cuasi-coherentes, si
la proposición se demuestra para cada uno de los $\sh H_\alpha$, también se
demostrará para $\sh H$.

Además basta demostrar que el homomorfismo canónico
\[
  \Gamma(X,\sh H)\longrightarrow\Gamma(X-Z,\sh H)
\]
es inyectivo (resp. biyectivo). Existe entonces un subanillo noetheriano
$A_0$ de $A$, una $A_0$-álgebra de tipo finito $B_0$, un ideal
$\mathfrak J_0$ de $B_0$ y un $B_0$-módulo de tipo finito $M_0$ que es un

\oldpage[IV-4]{225}

$A_0$-módulo plano, y un $A_0$-módulo $N_0$, tales que
$B=B_0\otimes_{A_0}A$, $\mathfrak J=\mathfrak J_0B$,
$M=M_0\otimes_{A_0}A$ y $N=N_0\otimes_{A_0}A$ (8.9.1, 8.5.11 y 11.2.7).
Además, sea $(A_\lambda)$ la familia de las sub-$A_0$-álgebras de tipo finito
de $A$, de modo que $A=\varinjlim A_\lambda$; pongamos
\[
\begin{gathered}
  B_\lambda=B_0\otimes_{A_0}A_\lambda,\qquad
  \mathfrak J_\lambda=\mathfrak J_0B_\lambda,\qquad
  M_\lambda=M_0\otimes_{A_0}A_\lambda=M_0\otimes_{B_0}B_\lambda,\\
  N_\lambda=N_0\otimes_{A_0}A_\lambda,\qquad
  S_\lambda=\Spec(A_\lambda),\quad
  X_\lambda=\Spec(B_\lambda),\quad
  Z_\lambda=V(\mathfrak J_\lambda);
\end{gathered}
\]
si $p_\lambda:X\to X_\lambda$ es la proyección canónica, se tiene
$Z=p_\lambda^{-1}(Z_\lambda)$ y
$X-Z=p_\lambda^{-1}(X_\lambda-Z_\lambda)$. Como
$X_\lambda-Z_\lambda$ es cuasicompacto y cuasi-separado, se deduce (8.5.2)
que, si se pone $\sh F_\lambda=\widetilde{M_\lambda}$ y
$\sh H_\lambda=(M_\lambda\otimes_{A_\lambda}N_\lambda)^\sim$, los
homomorfismos
\[
  \varinjlim\Gamma(X_\lambda,\sh H_\lambda)\longrightarrow\Gamma(X,\sh H)
  \qquad\text{y}\qquad
  \varinjlim\Gamma(X_\lambda-Z_\lambda,\sh H_\lambda)\longrightarrow\Gamma(X-Z,\sh H)
\]
son biyectivos; por la exactitud del funtor $\varinjlim$, basta demostrar que
para todo $\lambda$ suficientemente grande el homomorfismo canónico
$\Gamma(X_\lambda,\sh H_\lambda)\to
\Gamma(X_\lambda-Z_\lambda,\sh H_\lambda)$ es inyectivo (resp. biyectivo).
Pero para todo $x\in Z$, si $x_\lambda$ es la proyección de $x$ en
$Z_\lambda$, se tiene
\[
  \operatorname{prof}((\sh F_{f(x)})_x)
  =\operatorname{prof}(((\sh F_\lambda)_{f_\lambda(x_\lambda)})_{x_\lambda})
\]
en virtud de (4.2.7) y (6.7.1); para reducirnos al caso en que $A$ es
noetheriano, hay que probar que, para $\lambda$ suficientemente grande, se
tiene
$\operatorname{prof}(((\sh F_\lambda)_{f_\lambda(x_\lambda)})_{x_\lambda})
\geq1$ (resp. $\geq2$) para todo punto $x_\lambda\in Z_\lambda$. Ahora bien,
denotemos por $Z_\lambda'$ el conjunto de los puntos $x_\lambda\in Z_\lambda$
para los cuales, para toda generalización $z_\lambda$ de $x_\lambda$ en
$Z_\lambda\cap(X_\lambda)_{f_\lambda(x_\lambda)}$, se tiene
$\operatorname{prof}(((\sh F_\lambda)_{f_\lambda(x_\lambda)})_{z_\lambda})
\geq1$ (resp. $\geq2$); si $p_{\lambda\mu}:Z_\mu\to Z_\lambda$ es la proyección canónica
para $\lambda\leq\mu$, se deduce de la hipótesis y de (19.9.5) que
$Z_\mu'=p_{\lambda\mu}^{-1}(Z_\lambda')$ y
$Z=p_\lambda^{-1}(Z_\lambda')$; por otra parte, se deduce de (19.9.6) que
$Z_\lambda'$ es abierto en $Z_\lambda$, por lo que nuestra afirmación se sigue
de (8.3.4).

Para terminar la demostración cuando $A$ es noetheriano, observemos que, en
virtud de la hipótesis de planitud y de (6.3.1), se tiene entonces
$\operatorname{prof}_Z(\sh H)\geq1$
(resp. $\operatorname{prof}_Z(\sh H)\geq2$); pero la conclusión de (19.9.8)
se deduce en este caso de (5.10.2) (resp. 5.10.4)).

\begin{remark}[19.9.9]
\label{IV.19.9.9}
La misma demostración, junto con los resultados de la parte 3 del cap. III
sobre la profundidad y la cohomología local, permite establecer la siguiente
generalización de (19.9.8):

\emph{Bajo las condiciones generales de (19.9.8), supongamos que, para todo
$x\in Z$, se tiene $\operatorname{prof}((\sh F_{f(x)})_x)\geq k$; entonces el
homomorfismo canónico}
\[
  H^i(X,\sh H)\longrightarrow H^i(X-Z,\sh H|(X-Z))
\]
\emph{es biyectivo para $i\leq k-2$ e inyectivo para $i=k-1$} (lo que,
expresado mediante la noción cohomológica general de profundidad introducida
en el cap. III, se escribe también
$\operatorname{prof}_Z(\sh H)\geq k$).
\end{remark}


\newcommand{\shHom}{\mathcal{H}\!om}
% Mechanical subsection components; prose is translated in the component files.
\section{Funciones meromorfas y pseudomorfismos}
\label{IV.20}

\setcounter{subsection}{-1}

\subsection{Introduction}
\label{IV.20.0}

La mayoría de las nociones y resultados de los \textsection\textsection20 y 21 se vinculan directamente con el cap.~I, y apenas dependen de los caps.~II a IV, salvo por el uso ocasional de las nociones de profundidad y de anillo local regular (en \sref{IV.20.6}, \sref{IV.21.11}, \sref{IV.21.13} y \sref{IV.21.15}), del «teorema principal» de Zariski en \sref{IV.20.4} y \sref{IV.21.12}, y de las propiedades de las inmersiones transversalmente regulares en \sref{IV.20.6} y \sref{IV.21.15}.

\oldpage[IV-4]{226}
En el \textsection20 introducimos varias variantes de la noción de aplicación racional, estudiada ya en \sref[I]{section:I.7} desde un punto de vista aún bastante próximo al clásico y, por ello, poco adaptado al caso de preesquemas no necesariamente reducidos.
Las nociones y resultados del \textsection20 se utilizan para desarrollar en el \textsection21 (n\textsuperscript{os} \sref{IV.21.1} a \sref{IV.21.7}) la noción general de divisor y sus propiedades más elementales.
Esta noción resulta especialmente cómoda cuando los anillos locales de los preesquemas considerados son noetherianos e íntegramente cerrados, y sobre todo cuando son además \emph{factoriales} (\sref{IV.21.6} y \sref{IV.21.7}), debido a que en este último caso se identifica con la noción de \emph{ciclo de codimensión~$1$} (combinación lineal de subpreesquemas irreducibles de codimensión~$1$).
En \sref{IV.21.9} se determinan los divisores sobre un preesquema noetheriano de dimensión~$1$, aunque no necesariamente normal, lo cual es útil para diversas aplicaciones.
Los n\textsuperscript{os} \sref{IV.21.11} y \sref{IV.21.12} proporcionan dos teoremas importantes, debidos respectivamente a Auslander--Buchsbaum y Van der Waerden, relacionados con la noción de anillo factorial (los n\textsuperscript{os} \sref{IV.21.9}, \sref{IV.21.11} y \sref{IV.21.12} son independientes entre sí).
En los n\textsuperscript{os} \sref{IV.21.13} y \sref{IV.21.14}, también independientes de los tres anteriores, estudiamos una variante útil de la noción de anillo local factorial: la de anillo local \emph{parafactorial}, introducida en particular en~\cite{IV-41} al desarrollar teoremas de comparación entre el grupo de Picard de un preesquema proyectivo $X$ sobre un cuerpo $k$ y el de una «sección hiperplana».
Veremos en \sref{IV.21.14.1} (teorema de Ramanujam--Samuel) que los anillos locales parafactoriales son mucho más numerosos de lo que cabría esperar \emph{a priori}.

En \sref{IV.20.5}, \sref{IV.20.6} y \sref{IV.21.15} retomamos las nociones anteriores, pero desde un punto de vista «relativo» a un preesquema base fijo.
Por ahora estas nociones se utilizan con bastante poca frecuencia; en particular, la noción de divisor relativo apenas se emplea salvo cuando se trata de divisores positivos, y en ese caso puede explicitarse ventajosamente sin recurrir a la noción de función meromorfa relativa, mediante la noción de inmersión transversalmente regular de codimensión~$1$.
Por tanto, conviene omitir estas secciones en una primera lectura.

\subsection{Funciones meromorfas}
\label{IV.20.1}

\begin{env}[20.1.1]
\label{IV.20.1.1}
Sea $(X,\sh{O}_X)$ un espacio anillado, y sea $\sh{S}$ un subhaz \emph{de conjuntos} de $\sh{O}_X$.
Para todo abierto $U$ de $X$, consideremos el \emph{anillo de fracciones} $\Gamma(U,\sh{O}_X)[\Gamma(U,\sh{S})^{-1}]$ (Bourbaki, \emph{Alg. comm.}, cap.~II, \textsection2, n\textsuperscript{o}1).
Es inmediato que la aplicación $U\mapsto\Gamma(U,\sh{O}_X)[\Gamma(U,\sh{S})^{-1}]$ es un \emph{prehaz de anillos} (\sref[0]{0.1.5.1} y \sref[0]{0.1.5.7}).
Denotamos por $\sh{O}_X[\sh{S}^{-1}]$ el \emph{haz de anillos} asociado a este prehaz y decimos que es el \emph{haz de anillos de fracciones de $\sh{O}_X$, con denominadores en $\sh{S}$}; es un $\sh{O}_X$-módulo \emph{plano}.
Es inmediato que para todo $x\in X$ se tiene un isomorfismo canónico
\[
\label{IV.20.1.1.1}
  (\sh{O}_X[\sh{S}^{-1}])_x\isoto\sh{O}_{X,x}[\sh{S}_x^{-1}],
  \tag{20.1.1.1}
\]
pues el razonamiento de \sref[0]{0.1.4.5} se generaliza inmediatamente al caso en que se tiene un sistema inductivo $(A_\alpha,\vphi_{\beta\alpha})$ de anillos y, para cada índice $\alpha$, un subconjunto $S_\alpha$ de $A_\alpha$ tal que
\oldpage[IV-4]{227}
$\vphi_{\beta\alpha}(S_\alpha)\subset S_\beta$ para $\alpha\leq\beta$; tomamos entonces $S$ como el límite inductivo en $A=\varinjlim A_\alpha$ del sistema inductivo de subconjuntos $(S_\alpha)$.
\end{env}

\begin{env}[20.1.2]
\label{IV.20.1.2}
Sea ahora $\sh{F}$ un $\sh{O}_X$-módulo.
Ponemos entonces
\[
\label{IV.20.1.2.1}
  \sh{F}[\sh{S}^{-1}]=\sh{F}\otimes_{\sh{O}_X}\sh{O}_X[\sh{S}^{-1}]
  \tag{20.1.2.1}
\]
y decimos que es el \emph{haz de módulos de fracciones de $\sh{F}$ con denominadores en $\sh{S}$}; es inmediato que está asociado al prehaz de módulos $U\mapsto\Gamma(U,\sh{F})[\Gamma(U,\sh{S})^{-1}]$, y que para todo $x\in X$ se tiene un isomorfismo canónico
\[
\label{IV.20.1.2.2}
  (\sh{F}[\sh{S}^{-1}])_x\isoto\sh{F}_x[\sh{S}_x^{-1}].
  \tag{20.1.2.2}
\]
\end{env}

\begin{env}[20.1.3]
\label{IV.20.1.3}
Nos interesaremos aquí por el caso en que $\sh{S}$ es el subhaz $\sh{S}(\sh{O}_X)$ de $\sh{O}_X$ tal que, para todo abierto $U$, $\Gamma(U,\sh{S})$ sea el \emph{conjunto de elementos regulares} del anillo $\Gamma(U,\sh{O}_X)$; es inmediato que se trata de un haz (y no solamente de un prehaz), pues la regularidad de una sección de $\sh{O}_X$ sobre $U$ se verifica «fibra a fibra» (es decir, significa que el germen de la sección en $x$ es regular en $\sh{O}_{X,x}$ para todo $x\in U$); en otras palabras, $\sh{S}(\sh{O}_X)_x$ no es otro que el conjunto de elementos regulares de $\sh{O}_{X,x}$.
El haz de anillos correspondiente
\[
  \sh{M}_X=\sh{O}_X[\sh{S}^{-1}]
\]
se llama el \emph{haz de gérmenes de funciones meromorfas sobre $X$}, y las secciones de $\sh{M}_X$ sobre $X$ se llaman las \emph{funciones meromorfas sobre $X$}; forman un anillo que denotamos por $M(X)$.
Para todo $\sh{O}_X$-módulo $\sh{F}$,
\[
  \sh{F}\otimes_{\sh{O}_X}\sh{M}_X=\sh{F}[\sh{S}^{-1}]
\]
se denota también por $\sh{M}_X(\sh{F})$ y se llama el \emph{haz de gérmenes de secciones meromorfas de $\sh{F}$}; sus secciones sobre $X$ forman un $M(X)$-módulo que se denota por $M(X,\sh{F})$, cuyos elementos se llaman \emph{secciones meromorfas de $\sh{F}$ sobre $X$}.
Estas definiciones implican que para todo abierto $U$ de $X$ se tiene un isomorfismo canónico $\sh{M}_X(\sh{F})|U\isoto\sh{M}_U(\sh{F}|U)$, en particular $\sh{M}_X|U\isoto\sh{M}_U$.
\end{env}

\begin{env}[20.1.3.1]
\label{IV.20.1.3.1}
Si $X$ es un \emph{preesquema reducido}, observemos que si un elemento $s\in\Gamma(U,\sh{O}_X)$ es tal que $s_\xi\neq0$ para todo punto maximal $\xi$ de $U$, entonces $s$ es \emph{regular}.
En efecto, si $st=0$ para un $t\in\Gamma(U,\sh{O}_X)$, tenemos $s_\xi t_\xi=0$, luego $t_\xi=0$ pues $\sh{O}_{X,\xi}$ es un cuerpo, y decir que $t_\xi=0$ para todo punto maximal $\xi$ de $X$ significa que $t=0$: en efecto, nos reducimos inmediatamente al caso en que $U$ es afín, y un elemento de un anillo reducido que pertenece a todos los ideales primos mínimos es nulo por definición.
La recíproca es cierta si el conjunto de componentes irreducibles de $X$ es \emph{localmente finito}.
En efecto, nos reducimos inmediatamente al caso en que $X=\Spec(A)$ es afín; si $\mathfrak p_i$ ($1\leq i\leq n$) son los ideales primos mínimos de $A$ y si $s\in\mathfrak p_i$ para un índice $i$, existe $t\in A$ tal que $t\in\mathfrak p_j$ para $j\neq i$ y $t\notin\mathfrak p_i$ (Bourbaki, \emph{Alg. comm.}, cap.~II, \textsection1, n\textsuperscript{o}1, prop.~1); por tanto $st\in\mathfrak p_i$ para todo $i$, y en consecuencia $st=0$ puesto que $A$ es reducido; por tanto $s$ no es regular.
\end{env}

\begin{env}[20.1.4]
\label{IV.20.1.4}
Para todo abierto $U$ de $X$, el homomorfismo $t\mapsto t/1$ de $\Gamma(U,\sh{O}_X)$ en $\Gamma(U,\sh{O}_X)[\Gamma(U,\sh{S})^{-1}]$ (que no es otro que el \emph{anillo total de fracciones} de
\oldpage[IV-4]{228}
$\Gamma(U,\sh{O}_X)$) es inyectivo; por tanto, estos homomorfismos definen un \emph{homomorfismo canónico inyectivo}
\[
\label{IV.20.1.4.1}
  i:\sh{O}_X\to\sh{M}_X
  \tag{20.1.4.1}
\]
que permite identificar $\sh{O}_X$ con un subhaz de $\sh{M}_X$.
Dada una función meromorfa $\vphi\in M(X)$, decimos que $\vphi$ está \emph{definida} en un abierto $U$ de $X$ si $\vphi|U$ es una \emph{sección de $\sh{O}_U$} sobre $U$; los axiomas de los haces muestran que, para una sección $\vphi$ dada, existe un abierto \emph{mayor} en el que $\vphi$ está definida; lo llamamos \emph{dominio de definición} de $\vphi$ y lo denotamos por $\operatorname{dom}(\vphi)$.
\end{env}

\begin{env}[20.1.5]
\label{IV.20.1.5}
Para todo $\sh{O}_X$-módulo $\sh{F}$, se deduce de \sref{IV.20.1.4.1} un di-homomorfismo formado por $i$ y el homomorfismo de haces de grupos aditivos
\[
\label{IV.20.1.5.1}
  1_{\sh{F}}\otimes i:\sh{F}\to\sh{M}_X(\sh{F})=\sh{F}\otimes_{\sh{O}_X}\sh{M}_X.
  \tag{20.1.5.1}
\]
Obsérvese que este último ya no es inyectivo en general; cuando lo es, se dice que $\sh{F}$ es \emph{estrictamente sin torsión}: esto significa que para todo abierto $U$ de $X$ y toda sección $s\in\Gamma(U,\sh{O}_X)$ que sea un elemento regular de este anillo, la homotecia $z\mapsto sz$ de $\Gamma(U,\sh{F})$ es inyectiva; esta condición se cumple evidentemente si $\sh{F}$ es \emph{localmente libre}.
\end{env}

\begin{proposition}[20.1.6]
\label{IV.20.1.6}
Sean $X$ un preesquema localmente noetheriano y $\sh{F}$ un $\sh{O}_X$-módulo cuasicoherente.
Para que $\sh{F}$ sea estrictamente sin torsión, es necesario y suficiente que $\operatorname{Ass}(\sh{F})\subset\operatorname{Ass}(\sh{O}_X)$.
\end{proposition}

\begin{proof}
En efecto, nos reducimos inmediatamente al caso en que $X=\Spec(A)$ es afín, $\sh{F}=\widetilde M$, y se sabe que los elementos $s$ de $A$ que pertenecen a un ideal de $\operatorname{Ass}(M)$ son exactamente aquellos para los cuales la homotecia $z\mapsto sz$ no es inyectiva (Bourbaki, \emph{Alg. comm.}, cap.~IV, \textsection1, n\textsuperscript{o}1, cor.~2 de la prop.~2).
\end{proof}

\begin{env}[20.1.7]
\label{IV.20.1.7}
Si $u$ es una sección de $\sh{M}_X(\sh{F})$ sobre $X$, decimos que $u$ está \emph{definida} en un punto $x\in X$ si existe un entorno abierto $V$ de $x$ en $X$ tal que $u|V$ es la imagen de una sección de $\sh{F}$ sobre $V$ mediante el di-homomorfismo \sref{IV.20.1.5.1}.
Decimos que $u$ está \emph{definida} en un abierto $U$ de $X$ si está definida en todo punto de $U$; existe todavía un abierto mayor en el que $u$ está definida, llamado \emph{dominio de definición} de $u$ y denotado por $\operatorname{dom}(u)$.
Cuando $\sh{F}$ es estrictamente sin torsión, de modo que $\sh{F}$ se identifica mediante \sref{IV.20.1.5.1} con un subhaz de $\sh{M}_X(\sh{F})$, decir que $u$ está definida en $U$ significa que $u|U$ es una \emph{sección de $\sh{F}$ sobre $U$}.
\end{env}

\begin{env}[20.1.8]
\label{IV.20.1.8}
Conforme a las notaciones generales de \sref[0\textsubscript{I}]{0.5.4.7}, denotamos por $\sh{M}_X^*$ el haz de grupos multiplicativos tal que $\Gamma(U,\sh{M}_X^*)$ sea (para todo abierto $U$ de $X$) el grupo de los \emph{elementos invertibles} de $\Gamma(U,\sh{M}_X)$.
Este haz no es otro que el haz $\sh{S}(\sh{M}_X)$ definido en \sref{IV.20.1.3}: en efecto, si $s\in\Gamma(U,\sh{S}(\sh{M}_X))$, para todo $x\in U$ existe un entorno abierto $V\subset U$ de $x$ tal que $s|V$ es un elemento regular en el \emph{anillo total de fracciones} de $\Gamma(V,\sh{O}_X)$, y sabemos que tal elemento es necesariamente invertible en este anillo de fracciones.
Diremos que las secciones de $\sh{M}_X^*$ sobre $X$ son las \emph{funciones meromorfas regulares} (obsérvese que aquí nos apartamos de la terminología seguida por ciertos autores, que llaman funciones meromorfas «regulares» a las que son \emph{secciones de $\sh{O}_X$}, identificado como un subhaz de $\sh{M}_X$).

Sea $\sh{L}$ un $\sh{O}_X$-módulo invertible \sref[0\textsubscript{I}]{0.5.4.1}; entonces es claro que $\sh{M}_X(\sh{L})=\sh{L}\otimes_{\sh{O}_X}\sh{M}_X$
\oldpage[IV-4]{229}
es un $\sh{M}_X$-módulo invertible.
Sea $U$ un abierto tal que $\sh{L}|U$ sea isomorfo a $\sh{O}_U$; como todo automorfismo de $\sh{M}_U$ es la multiplicación por un elemento invertible de $\Gamma(U,\sh{M}_X)$ \sref[0\textsubscript{I}]{0.5.4.7}, es equivalente decir que una sección $s\in\Gamma(U,\sh{M}_X(\sh{L}))$ tiene una imagen invertible en $\Gamma(U,\sh{M}_X)$ mediante \emph{algún} isomorfismo o mediante \emph{todo} isomorfismo sobre $\Gamma(U,\sh{M}_X)$; en este caso diremos que $s$ es una \emph{sección meromorfa regular de $\sh{L}$ sobre $U$}; una sección $s$ de $\sh{L}$ sobre $X$ se llamará \emph{sección meromorfa regular de $\sh{L}$} si, para todo abierto $U$ tal que $\sh{L}|U$ sea isomorfo a $\sh{O}_U$, $s|U$ es una sección meromorfa regular de $\sh{L}$ sobre $U$.
Denotaremos por $(\sh{M}_X(\sh{L}))^*$ el subhaz de $\sh{M}_X(\sh{L})$ tal que, para todo abierto $U$, $\Gamma(U,(\sh{M}_X(\sh{L}))^*)$ sea el conjunto de secciones meromorfas regulares de $\sh{L}$ sobre $U$.
Sea $s$ una sección meromorfa de $\sh{L}$ sobre $X$ (es decir, una sección de $\sh{M}_X(\sh{L})$); define un homomorfismo $h_s:\sh{M}_X\to\sh{M}_X(\sh{L})$ que a toda sección $t$ de $\sh{M}_X$ sobre un abierto $U$ hace corresponder $(s|U)t$.
Se deduce inmediatamente de lo anterior que, para que $s$ sea \emph{regular}, es necesario y suficiente que $h_s$ sea \emph{inyectivo}, y de hecho entonces $h_s$ es un homomorfismo \emph{biyectivo} de $\sh{M}_X$ sobre $\sh{M}_X(\sh{L})$, y su restricción a $\sh{M}_X^*$ es una biyección sobre $(\sh{M}_X(\sh{L}))^*$.
Concluimos que la homotecia $t\mapsto ts$ es un isomorfismo de $M(X)$ sobre $M(X,\sh{L})$.
\end{env}

\begin{env}[20.1.9]
\label{IV.20.1.9}
Sea $s$ una sección meromorfa regular del $\sh{O}_X$-módulo invertible $\sh{L}$ sobre $X$; entonces, para todo $\sh{O}_X$-módulo $\sh{F}$, $s$ define igualmente un homomorfismo $h_s\otimes1_{\sh{F}}:\sh{M}_X(\sh{F})\to\sh{M}_X(\sh{F}\otimes_{\sh{O}_X}\sh{L})$, que es de nuevo \emph{biyectivo}.
\end{env}

\begin{env}[20.1.10]
\label{IV.20.1.10}
Sea $s$ una sección meromorfa de un $\sh{O}_X$-módulo invertible $\sh{L}$ sobre $X$; para que $s$ sea regular, es necesario y suficiente que exista una sección meromorfa $s'$ de $\sh{L}^{-1}$ sobre $X$ tal que la imagen canónica de $s\otimes s'$ en $\sh{M}_X$ \sref[0\textsubscript{I}]{0.5.4.3} sea la sección unidad, y esta sección $s'$ es entonces única: en efecto, la necesidad de la existencia local de tal sección es evidente, y su unicidad local entraña su existencia (y unicidad) global; además, la existencia de $s'$ es trivialmente suficiente para que $s$ sea regular.
Pondremos $s'=s^{-1}$.

Finalmente, si $\sh{L}'$ es un segundo $\sh{O}_X$-módulo invertible, $s$ (resp. $s'$) una sección meromorfa regular de $\sh{L}$ (resp. $\sh{L}'$) sobre $X$, entonces $s\otimes s'$ es evidentemente una sección meromorfa regular de $\sh{L}\otimes\sh{L}'$ sobre $X$.
\end{env}

\begin{env}[20.1.11]
\label{IV.20.1.11}
Si $f:X'\to X$ es un morfismo de espacios anillados, no existe en general una aplicación natural que asocie a una función meromorfa sobre $X$ una función meromorfa sobre $X'$.
Por ejemplo, si $X$ es el espectro de un anillo local íntegro $A$ y $X'$ el de su cuerpo residual $k$, no existe ningún homomorfismo natural del cuerpo de fracciones $K$ de $A$ en $k$, y solo se puede hacer corresponder a un elemento de $K$ un elemento de $k$ si ya pertenece a $A$.

En general, si $f=(\psi,\theta)$, denotemos, para todo abierto $U$ de $X$, por $\sh{S}_f(U)$ el conjunto de secciones \emph{regulares} $s\in\Gamma(U,\sh{O}_X)$ tales que la imagen de $s$ mediante
\[
  \Gamma(\theta^\#):\Gamma(U,\sh{O}_X)\to\Gamma(f^{-1}(U),\sh{O}_{X'})
\]
sea una sección \emph{regular}.
Es inmediato que $U\mapsto\sh{S}_f(U)$ es un \emph{subhaz} del haz de conjuntos $\sh{S}(\sh{O}_X)$, que denotamos por $\sh{S}_f$.
Ponemos $\sh{M}_f=\sh{O}_X[\sh{S}_f^{-1}]$; es un subhaz
\oldpage[IV-4]{230}
de anillos de $\sh{M}_X$, y se deduce canónicamente de $\theta^\#:\psi^*(\sh{O}_X)\to\sh{O}_{X'}$ un homomorfismo de haces de anillos $\theta^{\prime\#}:\psi^*(\sh{M}_f)\to\sh{M}_{X'}$ que prolonga $\theta^\#$ (Bourbaki, \emph{Alg. comm.}, cap.~II, \textsection2, n\textsuperscript{o}1, prop.~2); de donde, recordando que $f^*(\sh{M}_f)=\psi^*(\sh{M}_f)\otimes_{\psi^*(\sh{O}_X)}\sh{O}_{X'}$, un homomorfismo canónico de $\sh{O}_{X'}$-álgebras
\[
\label{IV.20.1.11.1}
  f^*(\sh{M}_f)\to\sh{M}_{X'}.
  \tag{20.1.11.1}
\]
Para toda función meromorfa $\vphi$ sobre $X$ que \emph{sea una sección de $\sh{M}_f$}, $\Gamma(\theta^{\prime\#})(\vphi)$ es una función meromorfa sobre $X'$, llamada \emph{imagen inversa de $\vphi$ por $f$}, y denotada por $\vphi\circ f$ si no hay riesgo de confusión.

Asimismo, si $\sh{F}$ es un $\sh{O}_X$-módulo, ponemos $\sh{M}_f(\sh{F})=\sh{F}\otimes_{\sh{O}_X}\sh{M}_f$, y deducimos inmediatamente de $\theta^{\prime\#}$ un homomorfismo canónico (que también se escribe $u\mapsto u\circ f$)
\[
  \Gamma(X,\sh{M}_f(\sh{F}))\to\Gamma(X',\sh{M}_{X'}(f^*(\sh{F}))).
\]
Además, si $u\in\Gamma(X,\sh{M}_f(\sh{F}))$ está definida \sref{IV.20.1.7} en un punto $x$, $u$ coincide, en un entorno $U$ de $x$, con una sección de la forma $\sum_i h_i\otimes(t_i/s_i)$, donde los $h_i$ pertenecen a $\Gamma(U,\sh{F})$, los $t_i$ a $\Gamma(U,\sh{O}_X)$ y los $s_i$ a $\Gamma(U,\sh{S}_f)$.
Como, por hipótesis, las imágenes de los $s_i$ en $\Gamma(f^{-1}(U),\sh{O}_{X'})$ son regulares, se ve que $u\circ f$ está definida en todo punto de $f^{-1}(U)$; en otras palabras, se tiene
\[
\label{IV.20.1.11.2}
  f^{-1}(\operatorname{dom}(u))\subset\operatorname{dom}(u\circ f).
  \tag{20.1.11.2}
\]
Veremos más adelante \sref{IV.20.6.5}[(i)] ejemplos (con $\sh{F}=\sh{O}_X$) en los que los dos miembros de \sref{IV.20.1.11.2} pueden ser distintos.

Consideremos en particular el caso en que $\sh{M}_f=\sh{M}_X$; entonces, si $\sh{L}$ es un $\sh{O}_X$-módulo invertible, la imagen en $\sh{M}_{X'}(f^*(\sh{L}))$, mediante $\Gamma(\theta^{\prime\#})$, de una sección meromorfa \emph{regular} de $\sh{L}$ sobre $X$ \sref{IV.20.1.8} es una sección meromorfa \emph{regular} de $f^*(\sh{L})$ sobre $X'$, como se deduce inmediatamente de la definición de estas secciones y del hecho de que un homomorfismo de anillos transforma un elemento invertible en un elemento invertible.

Sea $f':X''\to X'$ un segundo morfismo de espacios anillados, y supongamos que $\sh{M}_f=\sh{M}_X$ y $\sh{M}_{f'}=\sh{M}_{X'}$; entonces, si ponemos $f''=f\circ f'$, también se tiene $\sh{M}_{f''}=\sh{M}_X$, y se ve inmediatamente que para toda sección meromorfa $u$ de $\sh{F}$ sobre $X$ se tiene $u\circ f''=(u\circ f)\circ f'$.
\end{env}

\begin{proposition}[20.1.12]
\label{IV.20.1.12}
Si el morfismo $f:X'\to X$ es plano \sref[0\textsubscript{I}]{0.6.7.1}, se tiene $\sh{M}_f=\sh{M}_X$, y el homomorfismo $\vphi\mapsto\vphi\circ f$ está definido en todo $M(X)$.
Además, si $f$ es un morfismo (plano) de espacios anillados en anillos locales, se tiene $\operatorname{dom}(\vphi\circ f)=f^{-1}(\operatorname{dom}(\vphi))$; si además $f$ es sobreyectivo (por tanto, fielmente plano), el homomorfismo $\vphi\mapsto\vphi\circ f$ es inyectivo.
\end{proposition}

\begin{proof}
La primera afirmación se deduce del hecho de que, si $B$ es una $A$-álgebra que es un $A$-módulo plano, todo elemento de $A$ que no es divisor de $0$ en $A$ tampoco es divisor de $0$ en $B$ \sref[0\textsubscript{I}]{0.6.3.4}.
Para demostrar las otras afirmaciones, observemos que, para todo $x'\in X'$, si $x=f(x')$, $\sh{O}_{X',x'}$ es un $\sh{O}_{X,x}$-módulo plano, y como el homomorfismo $\sh{O}_{X,x}\to\sh{O}_{X',x'}$ es \emph{local} por hipótesis, es inyectivo (\sref[0\textsubscript{I}]{0.6.5.1} y \sref[0\textsubscript{I}]{0.6.6.2}); si ponemos $A=\sh{O}_{X,x}$, $B=\sh{O}_{X',x'}$, de modo que $A$ se identifica con un subanillo de $B$, $(f^*(\sh{M}_X))_{x'}$ es igual a $S^{-1}A\otimes_A B=S^{-1}B$, donde $S$ es el conjunto de elementos regulares de $A$, $(\sh{M}_{X'})_{x'}$ es igual a $T^{-1}B$, donde $T$ es el conjunto
\oldpage[IV-4]{231}
de elementos regulares de $B$, y como se ha visto que $S\subset T$, el homomorfismo $S^{-1}B\to T^{-1}B$ es inyectivo; en otras palabras, esto prueba que el homomorfismo \sref{IV.20.1.11.1} $f^*(\sh{M}_X)\to\sh{M}_{X'}$ es \emph{inyectivo} (de donde se deduce la última afirmación del enunciado).
El cociente $f^*(\sh{M}_X)/\sh{O}_{X'}$ se identifica con un sub-$\sh{O}_{X'}$-módulo de $\sh{M}_{X'}/\sh{O}_{X'}$, y $(f^*(\sh{M}_X)/\sh{O}_{X'})_{x'}$ se identifica con $(\sh{M}_X/\sh{O}_X)_x\otimes_{\sh{O}_{X,x}}\sh{O}_{X',x'}$.
Supongamos entonces que $x\notin\operatorname{dom}(\vphi)$; la imagen de $\vphi_x$ en $(\sh{M}_X/\sh{O}_X)_x$ es por tanto $\neq0$; por fiel planitud, se deduce que lo mismo ocurre con la imagen de $(\vphi\circ f)_{x'}$ en $(\sh{M}_{X'}/\sh{O}_{X'})_{x'}$, luego $x'\notin\operatorname{dom}(\vphi\circ f)$, lo que concluye la demostración.
\end{proof}

\begin{remark}[20.1.13]
\label{IV.20.1.13}
Sea $X$ un espacio analítico complejo \emph{reducido}; entonces la noción de función meromorfa sobre $X$ definida anteriormente coincide con la noción usual.
Consideremos por otra parte un preesquema $Y$, localmente de tipo finito sobre el cuerpo $\bb{C}$; sabemos entonces que se puede asociar a $Y$ un espacio analítico $Y^\mathrm{an}$ que tiene el mismo espacio topológico subyacente, y que el morfismo canónico $f:Y^\mathrm{an}\to Y$ es \emph{plano}~\cite{IV-37}; en virtud de \sref{IV.20.1.12}, el homomorfismo canónico $u\mapsto u\circ f$ de $M(Y)$ en $M(Y^\mathrm{an})$ está por tanto definido en todas partes y es inyectivo; pero no es \emph{sobreyectivo} en general.
Por ejemplo, cuando $Y=\bb{V}_{\bb{C}}^r$ (\textbf{Err}\textsubscript{\textbf{III}},~14) es el espacio afín de dimensión $r$ sobre $\bb{C}$, $M(Y)$ se identifica canónicamente con el cuerpo $R(Y)$ de las funciones racionales sobre $Y$ \sref{IV.20.2.13}[(i)], mientras que $M(Y^\mathrm{an})$ es el cuerpo de las funciones meromorfas usuales sobre $\bb{C}^r$.
En razón de este hecho, a menudo es preferible, en Geometría algebraica, abstenerse de la terminología introducida en esta sección y utilizar la terminología equivalente de «pseudo-función», que se definirá más abajo.
\end{remark}

\subsection{Pseudo-morfismos y pseudo-funciones}
\label{IV.20.2}

Los únicos espacios anillados de los que se tratará en esta sección son los \emph{preesquemas}.

\begin{env}[20.2.1]
\label{IV.20.2.1}
Recordemos \sref{IV.11.10.2} que en un preesquema $X$ se dice que un abierto $U$ es \emph{esquemáticamente denso} si, para todo abierto $V$ de $X$, el homomorfismo canónico $\Gamma(V,\sh{O}_X)\to\Gamma(V\cap U,\sh{O}_X)$ es \emph{inyectivo}.

Consideremos dos preesquemas $X$, $Y$, dos abiertos \emph{esquemáticamente densos} $U$, $U'$ de $X$; decimos que dos morfismos $u:U\to Y$, $u':U'\to Y$ son \emph{equivalentes} si existe un abierto $U''\subset U\cap U'$, \emph{esquemáticamente denso} en $X$, tal que $u|U''=u'|U''$.
Como se desprende inmediatamente de la definición de los abiertos esquemáticamente densos que la intersección de dos de tales abiertos es de nuevo uno de ellos, es inmediato que la relación precedente es efectivamente una relación de equivalencia.
Una clase de equivalencia para esta relación se llama un \emph{pseudo-morfismo} de $X$ en $Y$, o una \emph{aplicación racional estricta} de $X$ en $Y$.

Si $S$ es un preesquema y $X$, $Y$ dos $S$-preesquemas, llamamos \emph{pseudo-$S$-morfismo} a la clase de equivalencia (para la relación precedente) de un $S$-morfismo de un abierto esquemáticamente denso de $X$ en $Y$.
Un pseudo-morfismo no es por tanto otra cosa que un pseudo-$S$-morfismo para $S=\Spec(\bb{Z})$.

Denotamos por $\operatorname{Ps.hom}(X,Y)$ (resp.\ $\operatorname{Ps.hom}_S(X,Y)$) el conjunto de pseudo-morfismos (resp.\ pseudo-$S$-morfismos) de $X$ en $Y$.
\end{env}

\begin{env}[20.2.2]
\label{IV.20.2.2}
\oldpage[IV-4]{232}
Se desprende de la definición recordada arriba que si $U$ es un abierto esquemáticamente denso en $X$, entonces, para todo abierto $V$ de $X$, $U\cap V$ es esquemáticamente denso en $V$.
Si dos morfismos $u:U\to Y$, $u':U'\to Y$ de abiertos esquemáticamente densos de $X$ en $Y$ son equivalentes, se sigue que sus restricciones $u|(V\cap U)$ y $u'|(V\cap U')$ son también morfismos equivalentes (para el preesquema inducido en $V$); su clase se llama la \emph{restricción} a $V$ del pseudo-morfismo $\omega$, clase de $u$, y denotamos este pseudo-morfismo por $\omega|V$.
Si $W\subset V$ es un abierto de $X$, es claro que $(\omega|V)|W=\omega|W$.
Esto muestra que las aplicaciones de restricción definen un \emph{prehaz de conjuntos} $U\mapsto\operatorname{Ps.hom}(U,Y)$; obsérvese que este prehaz \emph{no es} en general un haz (cf.\ \sref{IV.20.2.16}); el haz asociado se denota por $\mathscr{P}\!s.\!\operatorname{hom}(X,Y)$.
Para los pseudo-$S$-morfismos se ve del mismo modo que $U\mapsto\operatorname{Ps.hom}_S(U,Y)$ es un prehaz de conjuntos, cuyo haz asociado denotamos por $\mathscr{P}\!s.\!\operatorname{hom}_S(X,Y)$.

Si $V$ es un abierto \emph{esquemáticamente denso} en $X$, entonces, para todo abierto $U$ esquemáticamente denso en $X$, $U\cap V$ es también esquemáticamente denso en $X$, por lo que la aplicación $\omega\mapsto\omega|V$ es una \emph{biyección} del conjunto $\operatorname{Ps.hom}(X,Y)$ (resp.\ $\operatorname{Ps.hom}_S(X,Y)$) sobre $\operatorname{Ps.hom}(V,Y)$ (resp.\ $\operatorname{Ps.hom}_S(V,Y)$).
\end{env}

\begin{env}[20.2.3]
\label{IV.20.2.3}
Dado un pseudo-$S$-morfismo $\omega$ de $X$ en $Y$, decimos que está \emph{definido} en un punto $x\in X$ si existe un entorno abierto $V$ de $x$ en $X$, un abierto $U\subset V$ \emph{que contiene $x$ y es esquemáticamente denso en $V$}, y finalmente un morfismo $S$ $u:U\to Y$ cuya clase en $\operatorname{Ps.hom}_S(V,Y)$ sea igual a $\omega|V$ \sref{IV.20.2.2}; llamamos \emph{dominio de definición} de $\omega$ y denotamos por $\operatorname{dom}_S(\omega)$ (o simplemente por $\operatorname{dom}(\omega)$ si $S=\Spec(\bb{Z})$) al conjunto de los $x\in X$ donde $\omega$ está definido; es evidentemente un \emph{abierto} de $X$.
Además, para todo abierto $W$ de $X$, se tiene
\[
\label{IV.20.2.3.1}
  \operatorname{dom}_S(\omega|W)=(\operatorname{dom}_S(\omega))\cap W
  \tag{20.2.3.1}
\]
en virtud de la propiedad de los abiertos esquemáticamente densos recordada en \sref{IV.20.2.2}.
\end{env}

\begin{proposition}[20.2.4]
\label{IV.20.2.4}
Supongamos que $X$, $Y$ sean $S$-preesquemas tales que el morfismo estructural $Y\to S$ sea separado; entonces, si $\omega$ es un pseudo-$S$-morfismo de $X$ en $Y$, $\operatorname{dom}_S(\omega)$ es el mayor de los abiertos esquemáticamente densos $U$ de $X$ tales que existe un morfismo $S$ $u:U\to Y$ perteneciente a la clase $\omega$.
\end{proposition}

\begin{proof}
Basta evidentemente demostrar que si $U$, $U'$ son dos abiertos esquemáticamente densos en $X$ tales que dos morfismos $S$ $u:U\to Y$ y $u':U'\to Y$ sean equivalentes, entonces las restricciones de $u$ y $u'$ a $U\cap U'$ son iguales.
Ahora bien, por hipótesis existe un abierto $U''\subset U\cap U'$, esquemáticamente denso en $X$ y en el que $u$ y $u'$ coinciden; como $U''$ es también esquemáticamente denso en $U\cap U'$, nuestra afirmación se deduce de \sref{IV.11.10.1}[(d)].
\end{proof}

\begin{corollary}[20.2.5]
\label{IV.20.2.5}
Sea $S$ un esquema sobre $S_0$, y sean $X$, $Y$ dos $S$-preesquemas tales que el morfismo compuesto $Y\to S\to S_0$ sea separado (lo que implica \sref[I]{I.5.5.1} que $Y\to S$ también lo es).
Para todo pseudo-$S$-morfismo $\omega$ de $X$ en $Y$, se tiene entonces $\operatorname{dom}_S(\omega)=\operatorname{dom}_{S_0}(\omega)$.
En particular, si $Y$ es un esquema, se tiene $\operatorname{dom}_S(\omega)=\operatorname{dom}(\omega)$.
\end{corollary}

\begin{proof}
En efecto, es claro que $\operatorname{dom}_S(\omega)\subset\operatorname{dom}_{S_0}(\omega)$ sin hipótesis de separación; en virtud de \sref{IV.20.2.4}, basta demostrar que si un morfismo $S_0$ $u_0:U_0\to Y$ de un abierto
\oldpage[IV-4]{233}
esquemáticamente denso $U_0$ de $X$ en $Y$ es tal que el compuesto $v_0:U_0\xrightarrow{u_0}Y\to S$ coincide con la restricción del morfismo estructural $w_0:U_0\to S$ en un abierto $U\subset U_0$ esquemáticamente denso en $X$, entonces $v_0=w_0$.
Pero en virtud de la hipótesis de que el morfismo $S\to S_0$ es separado, esto se deduce de nuevo de \sref{IV.11.10.1}[(d)].
\end{proof}

\begin{corollary}[20.2.6]
\label{IV.20.2.6}
Bajo las hipótesis de \sref{IV.20.2.4}, el prehaz $U\mapsto\operatorname{Ps.hom}_S(U,Y)$ es un haz.
\end{corollary}

\begin{proof}
En efecto, sean $U$ un abierto de $X$, $(U_\alpha)$ un recubrimiento de $U$ por abiertos de $U$; supongamos dado para cada $\alpha$ un pseudo-$S$-morfismo $\omega_\alpha$ de $U_\alpha$ en $Y$, de modo que para todo par de índices $\alpha$, $\beta$, las restricciones \sref{IV.20.2.2} de los pseudo-$S$-morfismos $\omega_\alpha$ y $\omega_\beta$ a $U_\alpha\cap U_\beta$ sean iguales; en virtud de \sref{IV.20.2.3.1}, esto implica que $\operatorname{dom}_S(\omega_\alpha)\cap U_\beta=\operatorname{dom}_S(\omega_\beta)\cap U_\alpha$.
Sea entonces $W$ la unión de los abiertos $\operatorname{dom}_S(\omega_\alpha)$, y, para todo $\alpha$, sea $u_\alpha$ el único morfismo $S$ $\operatorname{dom}_S(\omega_\alpha)\to Y$ perteneciente a la clase $\omega_\alpha$ \sref{IV.20.2.4}; en virtud de la hipótesis y de \sref{IV.20.2.4}, las restricciones de $u_\alpha$ y $u_\beta$ a $\operatorname{dom}_S(\omega_\alpha)\cap\operatorname{dom}_S(\omega_\beta)$ son iguales, por lo que existe un morfismo $u:W\to Y$ cuya restricción a cada abierto $\operatorname{dom}_S(\omega_\alpha)$ es igual a $u_\alpha$; es claro que $W$ es esquemáticamente denso en $U$, por lo que define un pseudo-$S$-morfismo $\omega$ de $U$ en $Y$ cuyas restricciones a los $U_\alpha$ son los $\omega_\alpha$.
\end{proof}

\begin{remark}[20.2.7]
\label{IV.20.2.7}
Sabemos \sref{IV.11.10.4} que cuando $X$ es \emph{reducido}, equivale a decir que un abierto de $X$ es denso en $X$ o que es esquemáticamente denso en $X$; la noción de pseudo-morfismo (resp.\ pseudo-$S$-morfismo) de $X$ en $Y$ coincide entonces con la de \emph{aplicación racional} (resp.\ $S$-\emph{aplicación racional}) de $X$ en $Y$ \sref[I]{I.7.1.2}.
En el caso general, la noción de pseudo-morfismo parece más útil que la de aplicación racional.
\end{remark}

\begin{env}[20.2.8]
\label{IV.20.2.8}
Llamamos \emph{pseudo-función} sobre $X$ a un pseudo-morfismo de $X$ en $\Spec(\bb{Z}[T])$ ($T$ indeterminada), o, lo que equivale a lo mismo, a un pseudo-morfismo $X$ de $X$ en $X\otimes_{\bb{Z}}\bb{Z}[T]$; equivale también \sref[I]{I.3.3.15} a decir que una pseudo-función sobre $X$ es una \emph{clase de equivalencia de secciones de $\sh{O}_X$ sobre abiertos esquemáticamente densos de $X$}, siendo equivalentes dos secciones $g\in\Gamma(U,\sh{O}_X)$, $g'\in\Gamma(U',\sh{O}_X)$ sobre tales abiertos si existe un abierto $U''\subset U\cap U'$, esquemáticamente denso en $X$, donde $g$ y $g'$ coinciden.
Podemos aplicar aquí \sref{IV.20.2.4} con $S=X$ e $Y=X\otimes_{\bb{Z}}\bb{Z}[T]$; por tanto, para toda pseudo-función $\vphi$ sobre $X$, existe un abierto \emph{mayor} esquemáticamente denso $\operatorname{dom}(\vphi)$ en $X$ tal que hay una sección de $\sh{O}_X$ sobre $\operatorname{dom}(\vphi)$ perteneciente a la clase $\vphi$.
Es claro que el haz $\mathscr{P}\!s.\!\operatorname{hom}_X(X,X\otimes_{\bb{Z}}\bb{Z}[T])$ es aquí un \emph{haz de anillos}, e incluso una $\sh{O}_X$-\emph{Álgebra}, que denotaremos por $\sh{M}'_X$; se deduce de \sref{IV.20.2.6} que, para todo abierto $U$ de $X$, $\Gamma(U,\sh{M}'_X)$ es igual al conjunto de pseudo-morfismos de $U$ en $\Spec(\bb{Z}[T])$; ponemos $M'(X)=\Gamma(X,\sh{M}'_X)$.
Decir que una sección $\vphi$ de $\sh{M}'_X$ sobre $U$ es \emph{invertible} en el anillo $\Gamma(U,\sh{M}'_X)$ significa que existe un abierto $U'$ esquemáticamente denso en $\operatorname{dom}(\vphi)$, por tanto en $U$, tal que, si $g$ es la única sección de $\sh{O}_X$ sobre $\operatorname{dom}(\vphi)$ perteneciente a la clase $\vphi$, $g|U'$ sea \emph{invertible} en $\Gamma(U',\sh{O}_X)$.
Se deduce de \sref[I]{I.3.3.15} que, en la correspondencia canónica entre $\Gamma(V,\sh{O}_X)$ y $\Hom(V,\bb{Z}[T])$ ($V$ abierto de $X$), los elementos invertibles de $\Gamma(V,\sh{O}_X)$ corresponden a los morfismos que se
\oldpage[IV-4]{234}
\emph{factorizan a través de} $V\to\Spec(\bb{Z}[T,T^{-1}])\to\Spec(\bb{Z}[T])$.
Concluimos que el haz ${\sh{M}'}^*_X$ de gérmenes de secciones invertibles de $\sh{M}'_X$ se identifica canónicamente con el haz $\mathscr{P}\!s.\!\operatorname{hom}_X(X,X\otimes_{\bb{Z}}\bb{Z}[T,T^{-1}])$.
\end{env}

\begin{lemma}[20.2.9]
\label{IV.20.2.9}
Sean $U$ un abierto de $X$ y $s$ un elemento regular del anillo $\Gamma(U,\sh{O}_X)$; entonces el abierto $U_s$ de los $x\in U$ tales que $s(x)\neq0$ es esquemáticamente denso en $U$.
\end{lemma}

\begin{proof}
En efecto, sean $V$ un abierto de $U$ y $t$ una sección de $\sh{O}_X$ sobre $V$ tal que $t|(V\cap U_s)=0$.
Para todo abierto afín $W\subset V$, existe entonces un entero $n$ tal que $s^nt|W=0$ \sref[I]{I.1.4.1}; pero como $s$ es un elemento \emph{regular} de $\Gamma(U,\sh{O}_X)$, esto implica $t|W=0$, de donde $t=0$.
\end{proof}

\begin{env}[20.2.10]
\label{IV.20.2.10}
Consideremos una función meromorfa $f\in M(X)$ \sref{IV.20.1.4}; entonces $\operatorname{dom}(f)$ es \emph{esquemáticamente denso} en $X$: en efecto, todo punto de $X$ admite un entorno abierto $U$ tal que existe una sección $s\in\Gamma(U,\sh{O}_X)$ que es un elemento regular de este anillo y para la cual $s(f|U)\in\Gamma(U,\sh{O}_X)$; como $s|U_s$ es invertible, concluimos que $f|U_s\in\Gamma(U_s,\sh{O}_X)$, luego $U_s\subset\operatorname{dom}(f)$ por definición \sref{IV.20.1.4}, y nuestra afirmación se deduce por tanto del lema \sref{IV.20.2.9}.
Podemos asociar así a $f$ la pseudo-función igual a la clase de la sección de $\sh{O}_X$ sobre $\operatorname{dom}(f)$, restricción de $f$; operando del mismo modo sustituyendo $X$ por un abierto cualquiera de $X$, obtenemos así un \emph{homomorfismo canónico de $\sh{O}_X$-Álgebras}
\[
\label{IV.20.2.10.1}
  \sh{M}_X\to\sh{M}'_X
  \tag{20.2.10.1}
\]
que, por restricción, da evidentemente un homomorfismo de haces de grupos abelianos
\[
\label{IV.20.2.10.2}
  \sh{M}_X^*\to{\sh{M}'}^*_X
  \tag{20.2.10.2}
\]
para los haces de gérmenes de secciones invertibles de $\sh{M}_X$ y $\sh{M}'_X$.
\end{env}

\begin{proposition}[20.2.11]
\label{IV.20.2.11}
\medskip
\begin{enumerate}[label=\emph{(\roman*)}]
  \item El homomorfismo canónico \emph{\sref{IV.20.2.10.1}} (y por consiguiente también \emph{\sref{IV.20.2.10.2}}) es \emph{inyectivo}.
  \item Supongamos, bien que $X$ sea localmente noetheriano, bien que $X$ sea reducido y que el conjunto de sus componentes irreducibles sea localmente finito.
  Entonces el homomorfismo canónico \emph{\sref{IV.20.2.10.1}} (y por consiguiente también \emph{\sref{IV.20.2.10.2}}) es \emph{biyectivo}.
\end{enumerate}
\end{proposition}

\begin{proof}
Las cuestiones consideradas son locales en $X$, por lo que podemos limitarnos al caso en que $X=\Spec(A)$ es afín, y demostrar entonces que el homomorfismo canónico $M(X)\to M'(X)$ es siempre inyectivo, y que es biyectivo en los casos considerados en~(ii).
Teniendo en cuenta la definición del haz $\sh{M}_X$ \sref{IV.20.1.3}, podemos observar además que \sref{IV.20.2.10.1} proviene en realidad de un homomorfismo de \emph{prehaces}
\[
  \Gamma(U,\sh{O}_X)[\Gamma(U,\sh{S})^{-1}]\to\Gamma(U,\sh{M}'_X)
\]
y basta demostrar que, para $U$ afín, este último es inyectivo (resp.\ biyectivo bajo las hipótesis de~(ii)).
Denotando por $S$ el conjunto de elementos regulares de $A$ (de modo que $S^{-1}A$ es el \emph{anillo total de fracciones} de $A$), debemos por tanto considerar el homomorfismo canónico
\[
\label{IV.20.2.11.1}
  S^{-1}A\to M'(X)
  \tag{20.2.11.1}
\]
\oldpage[IV-4]{235}
y demostrar que es inyectivo (resp.\ biyectivo bajo las condiciones de~(ii)).
Ahora bien, se puede escribir $S^{-1}A=\varinjlim A_t$, donde $t$ recorre el conjunto de elementos \emph{regulares} de $A$, ordenado según la relación «$t$ es un divisor de $t'$» \sref[0]{0.1.4.5}; también se puede escribir $A_t=\Gamma(\mathrm{D}(t),\sh{O}_X)$.
Por otra parte, por definición $M'(X)=\varinjlim\Gamma(U,\sh{O}_X)$, donde $U$ recorre el conjunto de abiertos esquemáticamente densos de $X$ (ordenado por la relación $\supset$), y la aplicación \sref{IV.20.2.11.1} no es otra que la aplicación canónica que proviene del hecho de que los $\mathrm{D}(t)$ constituyen una parte del conjunto de abiertos esquemáticamente densos en $X$ \sref{IV.20.2.9}.
Obsérvese que, por definición de abierto esquemáticamente denso, el homomorfismo $\Gamma(U,\sh{O}_X)\to\Gamma(U',\sh{O}_X)$ para dos abiertos de este tipo $U'\subset U$ es siempre \emph{inyectivo}, y por tanto lo mismo ocurre con los homomorfismos $\Gamma(U,\sh{O}_X)\to M'(X)$; concluimos que \sref{IV.20.2.11.1} es inyectivo.
Para demostrar que \sref{IV.20.2.11.1} es biyectivo, basta demostrar que el conjunto de los $\mathrm{D}(t)$ es \emph{cofinal} en el conjunto de abiertos esquemáticamente densos, es decir, que para tal abierto $U$ existe un $t$ regular en $A$ tal que $\mathrm{D}(t)\subset U$.
Ahora bien, cuando $X$ es reducido y las componentes irreducibles de $X$ forman un conjunto localmente finito, este conjunto es finito puesto que se ha supuesto que $X$ es afín, o sea, $A$ solo tiene un número finito de ideales primos mínimos $\mathfrak{p}_i$; como $A$ es reducido, la intersección de los $\mathfrak{p}_i$ se reduce a $0$, y decir que $t$ es regular equivale por tanto a decir que $t$ no pertenece a ninguno de los $\mathfrak{p}_i$; concluimos por el razonamiento de \sref[I]{I.7.1.9.1}.
Cuando $A$ es noetheriano, decir que $U=X-Y$ (donde $Y=V(\mathfrak{J})$ es cerrado en $X$) es esquemáticamente denso significa \sref{IV.5.10.2} que $Y$ \emph{no corta} a $\operatorname{Ass}(\sh{O}_X)$, y en virtud de Bourbaki, \emph{Alg.~comm.}, cap.~IV, \textsection1, n\textsuperscript{o}~4, prop.~8, esto implica la existencia de un $t\in\mathfrak{J}$ tal que $t$ sea $A$-regular, por lo que $U\supset\mathrm{D}(t)$.

Además, en el curso de esta demostración hemos demostrado el
\end{proof}

\begin{corollary}[20.2.12]
\label{IV.20.2.12}
Si $X=\Spec(A)$, donde $A$ es noetheriano, o reducido y con solo un número finito de ideales primos mínimos, los abiertos esquemáticamente densos en $X$ son aquellos que contienen un abierto de la forma $\mathrm{D}(t)$, donde $t$ es regular en $A$, y $M(X)=M'(X)$ es el anillo total de fracciones $S^{-1}A$, donde $S$ es el conjunto de elementos regulares de $A$\footnote{La transcripción francesa dice aquí «elementos regulares de $X$»; como $S\subset A$ y se forma $S^{-1}A$, debe decir «de $A$», como confirma el testigo inglés mantenido.}.
\end{corollary}

\begin{remarks}[20.2.13]
\label{IV.20.2.13}
\begin{enumerate}[label=(\roman*)]
  \item Sea $\vphi$ un elemento de $M(X)$, $\vphi'$ su imagen en $M'(X)$; por definición (\sref{IV.20.1.4} y \sref{IV.20.2.8}) se tiene $\operatorname{dom}(\vphi)\subset\operatorname{dom}(\vphi')$; pero, de hecho, se tiene la igualdad $\operatorname{dom}(\vphi)=\operatorname{dom}(\vphi')$, pues existe una sección de $\sh{O}_X$ sobre $\operatorname{dom}(\vphi')$ perteneciente a la clase $\vphi'$, y la función meromorfa correspondiente es igual a $\vphi$ por \sref{IV.20.2.11}[(i)], por lo que $\operatorname{dom}(\vphi')\subset\operatorname{dom}(\vphi)$.

  \item Ya se ha señalado que cuando $X$ es \emph{reducido}, se tiene $\sh{M}'_X=\sh{R}(X)$ (haz de funciones racionales sobre $X$ \sref[I]{I.7.3.2}); si además las componentes irreducibles de $X$ forman un conjunto localmente finito, se tiene $\sh{M}_X=\sh{M}'_X=\sh{R}(X)$.
  En general, puesto que todo abierto esquemáticamente denso es denso, se tiene un homomorfismo canónico $\sh{M}'_X\to\sh{R}(X)$, pero incluso cuando $X$ es localmente noetheriano, este homomorfismo no es necesariamente inyectivo.
  Por ejemplo, si $X=\Spec(A)$, donde $A$ es un anillo noetheriano tal que $\operatorname{Ass}(A)$ contiene ideales primos inmersos (lo que implica que $A$ no es reducido), se ha visto que $M(X)$ y $M'(X)$ se identifican con el anillo total de fracciones $S^{-1}A$,
  \oldpage[IV-4]{236}
  donde $S$ es el conjunto de elementos regulares de $A$, complemento de la unión de los ideales $\mathfrak{p}\in\operatorname{Ass}(A)$; por el contrario, $R(X)$ se identifica con $Q^{-1}A$, donde $Q$ es el complemento de la unión de los ideales primos \emph{mínimos} de $A$ \sref[I]{I.7.1.9}, y el homomorfismo canónico $A\to Q^{-1}A$ (y \emph{a fortiori} $S^{-1}A\to Q^{-1}A$) no es por tanto inyectivo, pues existen en $A-Q$ elementos $\neq0$ de $A$ anulados por elementos de $Q$ (Bourbaki, \emph{Alg.~comm.}, cap.~IV, \textsection1, n\textsuperscript{o}~1, cor.~2 de la prop.~1).

  \item Obsérvese también que incluso cuando $X$ es localmente noetheriano, el $\sh{O}_X$-Módulo $\sh{M}'_X$ \emph{no es necesariamente cuasicoherente}.
  Considérese por ejemplo un anillo local noetheriano $A$ de dimensión $\geq2$, cuyo ideal maximal $\mathfrak{m}$ satisface $\mathfrak{m}\in\operatorname{Ass}(A)$ y tal que, si se pone $X=\Spec(A)$, el esquema inducido en el abierto $U$ de $X$, complementario de $\{\mathfrak{m}\}$, sea \emph{íntegro}.
  Los únicos elementos regulares de $A$ son entonces los elementos invertibles, por lo que $\Gamma(X,\sh{M}'_X)=M'(X)=A$; si $\sh{M}'_X$ fuese cuasicoherente, sería por tanto igual a $\sh{O}_X$; pero como $U$ es un esquema íntegro, $\sh{M}'_X|U=R(U)$ es un haz \emph{simple} \sref[I]{I.7.3.5}, mientras que $\sh{O}_U$ no es un haz simple puesto que $\dim(A)\geq2$.

  Queda por dar un ejemplo de un anillo $A$ que tenga las propiedades precedentes.
  Basta considerar un anillo local íntegro $B$ de dimensión $\geq2$ y de cuerpo residual $k$, y tomar $A=B\oplus k$ con la ley de multiplicación $(b,x)(b',x')=(bb',bx'+b'x)$.

  \item Si $X$ es localmente noetheriano, se deduce de \sref{IV.5.10.2} que los abiertos esquemáticamente densos en $X$ son aquellos que \emph{contienen} el conjunto $\operatorname{Ass}(\sh{O}_X)$.
\end{enumerate}
\end{remarks}

\begin{env}[20.2.14]
\label{IV.20.2.14}
Sean $X$ un preesquema, $\sh{F}$ un $\sh{O}_X$-Módulo cuasicoherente y \emph{estrictamente sin torsión} \sref{IV.20.1.5}, de modo que $\sh{F}$ se identifica con un sub-$\sh{O}_X$-Módulo de $\sh{M}_X(\sh{F})$.
Para toda sección meromorfa $u$ de $\sh{F}$ sobre $X$, llamamos el \emph{Ideal de los denominadores de $u$} al anulador $\sh{J}$ de la sección $\overline{u}$, imagen de $u$ en $\sh{M}_X(\sh{F})/\sh{F}$.
El Ideal $\sh{J}$ es \emph{cuasicoherente}: en efecto, como la cuestión es local en $X$, podemos limitarnos al caso en que $X$ sea afín, y existe una sección $s\in\Gamma(X,\sh{S}(\sh{O}_X))$ tal que $v=su\in\Gamma(X,\sh{F})$.
Decir que, para un abierto $U\subset X$, una sección $f\in\Gamma(U,\sh{O}_X)$ pertenece a $\Gamma(U,\sh{J})$ significa que $f(u|U)\in\Gamma(U,\sh{F})$, y puesto que $s|U$ es un elemento regular de $\Gamma(U,\sh{O}_X)$ y $\sh{F}$ es estrictamente sin torsión, la relación precedente equivale también a $f((su)|U)\in\Gamma(U,s\sh{F})$; si $\overline{v}$ es la sección de $\sh{F}/s\sh{F}$, imagen canónica de $v$, se ve por tanto que $\sh{J}$ es el núcleo del homomorfismo $\sh{O}_X\to\sh{F}/s\sh{F}$ obtenido por multiplicación por la sección $\overline{v}$.
Como $\sh{F}/s\sh{F}$ es cuasicoherente, lo mismo ocurre con $\sh{J}$.

Se deduce inmediatamente de la definición precedente que $\operatorname{dom}(u)$ es \emph{el abierto complementario del subpreesquema cerrado de $X$ definido por el Ideal de los denominadores de $u$}.
\end{env}

\begin{proposition}[20.2.15]
\label{IV.20.2.15}
Sean $f:X'\to X$ un morfismo, $\sh{F}$ un $\sh{O}_X$-Módulo cuasicoherente, $\vphi$ una sección de $\sh{M}_f(\sh{F})$ sobre $X$ \sref{IV.20.1.11}.
Entonces $f^{-1}(\operatorname{dom}(\vphi))$ es esquemáticamente denso en $X'$.
\end{proposition}

\begin{proof}
Como la cuestión es evidentemente local en $X$ y $X'$, podemos suponer $X=\Spec(A)$, $X'=\Spec(A')$ afines, $\sh{F}=\widetilde{M}$ y $\vphi=h\otimes(1/s)$, donde $h\in M$ y $s$ es un elemento regular de $A$ cuya imagen $s'$ en $A'$ es un elemento regular.
Sabemos entonces \sref{IV.20.2.9} que $\mathrm{D}(s')$ es un abierto esquemáticamente denso en $X'$, y es la imagen inversa por $f$ de $\mathrm{D}(s)$, que está contenido en $\operatorname{dom}(\vphi)$.
\end{proof}

\begin{remark}[20.2.16]
\label{IV.20.2.16}
Cuando $Y$ no es separado, el prehaz $U\mapsto\operatorname{Ps.hom}_S(U,Y)$ sobre $X$ no es necesariamente un haz, incluso cuando $X$ es noetheriano, como lo muestra el ejemplo siguiente.
  Tomaremos para $X$ el espectro de un anillo semilocal noetheriano $A$ de dimensión $\geq2$, que tiene exactamente dos
  \oldpage[IV-4]{237}
  ideales máximos $\mathfrak{m}'$, $\mathfrak{m}''$ (de modo que $X$ tiene exactamente dos puntos cerrados $x'$, $x''$), de suerte que $\mathfrak{m}'$ y $\mathfrak{m}''$ pertenecen a $\operatorname{Ass}(A)$ y que el abierto $U=X-\{x',x''\}$ sea \emph{íntegro}.
Sean $U'=X-\{x'\}$, $U''=X-\{x''\}$, de modo que $U=U'\cap U''$.
Obsérvese que los únicos abiertos esquemáticamente densos en $X$ son los que contienen $x'$ y $x''$ \sref{IV.20.2.13}[(iv)], por lo que necesariamente son iguales a $X$.
Para obtener un contraejemplo, basta definir dos morfismos $S$ $u':U'\to Y$, $u'':U''\to Y$ (con $S=X$) cuyas restricciones a $U$ pertenezcan al mismo pseudo-morfismo de $U$ en $Y$, pero tales que, para todo entorno $V'$ de $x''$ en $U'$ y todo entorno $V''$ de $x'$ en $U''$, las restricciones de $u'$ y $u''$ a $V'\cap V''$ sean distintas.
Para ello, consideremos una parte cerrada irreducible $Z$ de $X$ que contiene $x'$ y $x''$, distinta de $X$; sea $Y$ el $X$-preesquema obtenido pegando dos preesquemas $Y'$, $Y''$ isomorfos a $X$ a lo largo del abierto por doquier denso $X-Z$ \sref[I]{I.2.3.1}.
Tomaremos entonces para $u'$ y $u''$ las restricciones respectivas a $U'$ y $U''$ de los isomorfismos inversos de los isomorfismos estructurales $Y'\to X$, $Y''\to X$.
Como $V'$ y $V''$ contienen el punto genérico de $Z$, las restricciones de $u'$ y $u''$ a $V'\cap V''$ son distintas, pero las restricciones de $u'$ y $u''$ a $U-(U\cap Z)$ son iguales, y $U-(U\cap Z)$ es un abierto denso en $U$, por tanto esquemáticamente denso puesto que $U$ es reducido.

Solo queda definir $A$ y $Z$ de modo que tengan las propiedades precedentes.
Sean $X_0=\Spec(B)$ un esquema afín íntegro (por ejemplo, el plano afín sobre un cuerpo $k$), $Y$ una parte cerrada irreducible de $X_0$ (por ejemplo, una recta afín) que contiene dos puntos cerrados distintos $x'$ y $x''$, correspondientes a los ideales máximos $\mathfrak{n}'$, $\mathfrak{n}''$ de $B$.
Consideremos el anillo $C=B\oplus(B/\mathfrak{n}'\oplus B/\mathfrak{n}'')$ con la multiplicación $(b,z)(b',z')=(bb',bz'+b'z)$.
Si $X_1=\Spec(C)$, se tiene $X_0=(X_1)_\mathrm{red}$ y $X_1$ es reducido salvo en los puntos $x'$, $x''$.
Basta entonces tomar $A=R^{-1}C$, donde $R$ es el complemento de la unión de los ideales máximos de $C$ en los puntos $x'$, $x''$, y para $Z$ la traza de $Y$ sobre $X=\Spec(A)$.
\end{remark}

\subsection{Composición de los pseudomorfismos}
\label{IV.20.3}

\begin{env}[20.3.1]
\label{IV.20.3.1}
Sean $X$, $Y$, $Z$ tres preesquemas, $\omega$ un pseudomorfismo de $X$ en $Y$, y $f:Y\to Z$ un morfismo.
Es claro que si $U'$, $U''$ son dos abiertos esquemáticamente densos en $X$, y $u':U'\to Y$, $u'':U''\to Y$ son dos morfismos de la clase $\omega$, los morfismos $f\circ u'$ y $f\circ u''$ son equivalentes (para la relación definida en \sref{IV.20.2.1}); por tanto, para todos los morfismos $u$ de la clase $\omega$, los morfismos $f\circ u$ pertenecen a una misma clase de equivalencia, que denotamos $f\circ\omega$ y llamamos el pseudomorfismo de $X$ en $Z$ \emph{compuesto de $f$ y de $\omega$}.
Se tiene $\operatorname{dom}(f\circ\omega)=\operatorname{dom}(\omega)$.
Si $g:Z\to T$ es un morfismo, es claro que $g\circ(f\circ\omega)=(g\circ f)\circ\omega$.
\end{env}

\begin{env}[20.3.2]
\label{IV.20.3.2}
Supongamos ahora dado un pseudomorfismo-$S$ $\omega$ de $X$ en $Y$, donde $Y$ es \emph{separado sobre $S$}, de modo que existe un $S$-morfismo $u:\operatorname{dom}_S(\omega)\to Y$ de la clase $\omega$ \sref{IV.20.2.4}.
Sea $f:X'\to X$ un $S$-morfismo tal que el abierto $V'=f^{-1}(\operatorname{dom}_S(\omega))$ sea esquemáticamente denso en $X'$; entonces se dice que la clase (para la relación de equivalencia de \sref{IV.20.2.1}) del $S$-morfismo $u\circ(f|V')$\footnote{La transcripción francesa imprime $f|U'$, pero el abierto introducido en esta misma frase es $V'=f^{-1}(\operatorname{dom}_S(\omega))$; el testigo inglés mantenido restituye $f|V'$.} (clase que está definida en virtud de la hipótesis anterior) es el pseudomorfismo-$S$ \emph{compuesto de $\omega$ y de $f$}, y se denota $\omega\circ f$; es claro que se tiene
\[
\label{IV.20.3.2.1}
  f^{-1}(\operatorname{dom}_S(\omega))\subset\operatorname{dom}_S(\omega\circ f)
  \tag{20.3.2.1}
\]

Para $f:X'\to X$ dado, denotamos por $\operatorname{Ps.hom}_S(X,Y)^f$ el conjunto de los pseudomorfismos-$S$ $\omega$ de $X$ en $Y$ que satisfacen la condición anterior.
Si $\omega$ es tal pseudomorfismo-$S$, es claro que para todo abierto $V$ de $X$,
\[
  f^{-1}(\operatorname{dom}_S(\omega|V))=f^{-1}(V\cap\operatorname{dom}_S(\omega))=f^{-1}(V)\cap f^{-1}(\operatorname{dom}_S(\omega))
\]
es esquemáticamente denso en $f^{-1}(V)$; por tanto, si $f^V:f^{-1}(V)\to V$ es la restricción de $f$, el compuesto $(\omega|V)\circ f^V$ está definido y es igual a $(\omega\circ f)|f^{-1}(V)$.
Se define así una
\oldpage[IV-4]{238}
aplicación canónica de restricción $\operatorname{Ps.hom}_S(X,Y)^f\to\operatorname{Ps.hom}_S(V,Y)^{f^V}$, lo que permite definir un \emph{prehaz} de conjuntos $V\mapsto\operatorname{Ps.hom}_S(V,Y)^{f^V}$ sobre $X$, subprehaz de $V\mapsto\operatorname{Ps.hom}_S(V,Y)$, y por tanto un haz de conjuntos asociado que denotamos $\mathscr{P}\!s.\!\operatorname{hom}_S(X,Y)^f$.
Además, para todo abierto $V$ de $X$, se tiene una aplicación $\omega\mapsto\omega\circ f^V$ de $\operatorname{Ps.hom}_S(V,Y)^{f^V}$ en $\operatorname{Ps.hom}_S(f^{-1}(V),Y)$, que define por tanto un $f$-morfismo de haces de conjuntos
\[
  \mathscr{P}\!s.\!\operatorname{hom}_S(X,Y)^f\to\mathscr{P}\!s.\!\operatorname{hom}_S(X',Y).
\]
\end{env}

\begin{env}[20.3.3]
\label{IV.20.3.3}
Sea ahora $f':X''\to X'$ un $S$-morfismo tal que el abierto $f'^{-1}(f^{-1}(\operatorname{dom}_S(\omega)))$ sea esquemáticamente denso en $X''$; entonces $\omega\circ(f\circ f')$ está definido y $u\circ f\circ f'$ pertenece a este pseudomorfismo; además, en virtud de \sref{IV.20.3.2.1}, $f'^{-1}(\operatorname{dom}_S(\omega\circ f))$ es \emph{a fortiori} esquemáticamente denso en $X''$, de modo que $(\omega\circ f)\circ f'$ también está definido y $u\circ f\circ f'$ pertenece a este pseudomorfismo; por tanto $(\omega\circ f)\circ f'=\omega\circ(f\circ f')$.
Por otra parte, para todo $S$-morfismo $g:Y\to Z$, se tiene $\operatorname{dom}_S(g\circ\omega)=\operatorname{dom}_S(\omega)$ \sref{IV.20.3.1}; por tanto $(g\circ\omega)\circ f$ está definido y $g\circ u\circ f$ pertenece a este pseudomorfismo, lo que muestra que $(g\circ\omega)\circ f=g\circ(\omega\circ f)$.
\end{env}

\begin{env}[20.3.4]
\label{IV.20.3.4}
El caso más importante en la definición \sref{IV.20.3.2} es aquel en que $\mathscr{P}\!s.\!\operatorname{hom}_S(X,Y)^f=\mathscr{P}\!s.\!\operatorname{hom}_S(X,Y)$: para ello basta que, para todo abierto $U$ de $X$ y todo abierto $V$ esquemáticamente denso en $U$, $f^{-1}(V)$ sea esquemáticamente denso en $f^{-1}(U)$; cuando esto ocurre, para \emph{todo} abierto $U$ de $X$ y \emph{todo} pseudomorfismo-$S$ $\omega:U\to Y$ se puede definir el compuesto $\omega\circ f^U$, \emph{incluso si $Y$ no es separado sobre $S$}.
En efecto, si $u':U'\to Y$ y $u'':U''\to Y$ son dos morfismos de la clase $\omega$, coinciden en un abierto $U_0$ esquemáticamente denso en $U$; por tanto, los morfismos compuestos $f^{-1}(U')\to U'\xrightarrow{u'}Y$ y $f^{-1}(U'')\to U''\xrightarrow{u''}Y$ coinciden en $f^{-1}(U_0)$, que por hipótesis es esquemáticamente denso en $f^{-1}(U)$; así, se puede definir $\omega\circ f^U$ como la clase de cualquiera de los morfismos $f^{-1}(U')\to U'\xrightarrow{u'}Y$, donde $u'$ pertenece a $\omega$.
\end{env}

\begin{proposition}[20.3.5]
\label{IV.20.3.5}
Sean $X$, $X'$ dos preesquemas, y $f:X'\to X$ un morfismo.
Entonces, en cada uno de los tres casos siguientes, para todo abierto $U$ de $X$ y todo abierto $V$ esquemáticamente denso en $U$, $f^{-1}(V)$ es esquemáticamente denso en $f^{-1}(U)$:
\begin{enumerate}[label=\emph{(\roman*)}]
  \item $f$ es plano y localmente de presentación finita.
  \item $f$ es plano y el espacio subyacente a $X$ es localmente noetheriano.
  \item $X'$ es reducido, el conjunto de las componentes irreducibles de $X$ es localmente finito, y toda componente irreducible de $X'$ domina una componente irreducible de $X$.
\end{enumerate}
\end{proposition}

\begin{proof}
La afirmación~(i) resulta de \sref{IV.11.10.5}[(ii), (b)]; la afirmación~(ii) resulta de \sref{IV.11.10.5}[(ii), (a)], pues entonces todo abierto $V$ en $U$ es retrocompacto, es decir, la inclusión canónica $j:V\to U$ es un morfismo cuasicompacto.
Finalmente, para probar~(iii), observemos que, puesto que $X'$ es reducido, basta mostrar que para todo abierto $U$ de $X$ y todo abierto $V$ denso en $U$, $f^{-1}(V)$ es \emph{denso} en $f^{-1}(U)$.
Ahora bien, para que $f^{-1}(V)$ sea denso en $f^{-1}(U)$, basta que $f^{-1}(V)$ contenga todos los puntos máximos de $X'$ que pertenecen a $f^{-1}(U)$; la conclusión resulta, pues, de la hipótesis de que la imagen por $f$ de todo punto
\oldpage[IV-4]{239}
máximo de $X'$ perteneciente a $f^{-1}(U)$ es un punto máximo de $X$ perteneciente a $U$, y por tanto a $V$, puesto que el conjunto de las componentes irreducibles de $X$ es localmente finito.
\end{proof}

\begin{env}[20.3.6]
\label{IV.20.3.6}
Sean $X$, $Y$ dos $S$-preesquemas; supongamos que $X$ satisface una de las dos hipótesis siguientes:
\begin{enumerate}[label=\emph{\alph*)}]
  \item $X$ es localmente noetheriano;
  \item $X$ es reducido y el conjunto de sus componentes irreducibles es localmente finito.
\end{enumerate}
Entonces, para todo $x\in X$, el $S$-morfismo canónico $j:\Spec(\sh{O}_{X,x})\to X$ es plano y satisface la condición~(ii) de \sref{IV.20.3.5} en el caso~a), y la condición~(iii) de \sref{IV.20.3.5} en el caso~b).
Para todo pseudomorfismo-$S$ $\omega$ de $X$ en $Y$, el compuesto $\omega\circ j$ está, pues, definido y es un pseudomorfismo-$S$ de $\Spec(\sh{O}_{X,x})$ en $Y$, llamado la \emph{restricción} de $\omega$ a $\Spec(\sh{O}_{X,x})$.
Observemos ahora que si $X$ satisface la condición~a) (respectivamente b)) de \sref{IV.20.3.6}, lo mismo ocurre con todo preesquema inducido sobre un abierto $U$ de $X$ que contenga $x$.
Por paso al límite inductivo, se deduce de las aplicaciones canónicas $\operatorname{Ps.hom}_S(U,Y)\to\operatorname{Ps.hom}_S(\Spec(\sh{O}_{X,x}),Y)$ así obtenidas una aplicación canónica
\[
\label{IV.20.3.6.1}
  (\mathscr{P}\!s.\!\operatorname{hom}_S(X,Y))_x\to\operatorname{Ps.hom}_S(\Spec(\sh{O}_{X,x}),Y)
  \tag{20.3.6.1}
\]
donde el primer miembro es la fibra en el punto $x$ del haz $\mathscr{P}\!s.\!\operatorname{hom}_S(X,Y)$, conjunto de los gérmenes en $x$ de pseudomorfismos-$S$ de entornos abiertos de $x$ en $Y$.
\end{env}

\begin{proposition}[20.3.7]
\label{IV.20.3.7}
Bajo las hipótesis de \sref{IV.20.3.6}, la aplicación canónica \sref{IV.20.3.6.1} es inyectiva.
Si además $Y$ es localmente de presentación finita sobre $S$, la aplicación \sref{IV.20.3.6.1} es biyectiva.
\end{proposition}

\begin{proof}
Por aplicación del método de \sref{IV.8.1.2}[(a)], esta proposición será un caso particular de la siguiente:
\end{proof}

\begin{proposition}[20.3.8]
\label{IV.20.3.8}
Con las notaciones de \sref{IV.8.8.1}, supongamos que $X_\alpha$ sea cuasicompacto (respectivamente, cuasicompacto y cuasiseparado) y que $Y_\alpha$ sea localmente de tipo finito (respectivamente, localmente de presentación finita) sobre $S_\alpha$.
Supongamos además que se verifique una de las condiciones siguientes:
\begin{enumerate}[label=\emph{(\roman*)}]
  \item Los morfismos de transición $S_\mu\to S_\lambda$ (para $\lambda\leq\mu$) son planos, y los $X_\lambda$ y $X$ son noetherianos.
  \item Los $X_\lambda$ son reducidos, el conjunto de las componentes irreducibles de $X$ y de cada uno de los $X_\lambda$ es finito, y, para $\lambda\leq\mu$, toda componente irreducible de $X_\mu$ domina una componente irreducible de $X_\lambda$.
\end{enumerate}
Entonces, la aplicación canónica
\[
\label{IV.20.3.8.1}
  \varinjlim\operatorname{Ps.hom}_{S_\lambda}(X_\lambda,Y_\lambda)\to\operatorname{Ps.hom}_S(X,Y)
  \tag{20.3.8.1}
\]
es inyectiva (respectivamente, biyectiva).
\end{proposition}

\begin{proof}
Observemos primero que, en el caso~(i), los morfismos $X_\mu\to X_\lambda$ (para $\lambda\leq\mu$) y $X\to X_\lambda$ son planos; por tanto, de \sref{IV.20.3.4} y \sref{IV.20.3.5} resulta que las aplicaciones canónicas
\[
  \operatorname{Ps.hom}_{S_\lambda}(X_\lambda,Y_\lambda)\to\operatorname{Ps.hom}_{S_\mu}(X_\mu,Y_\mu)
\]
\oldpage[IV-4]{240}
para $\lambda\leq\mu$, y $\operatorname{Ps.hom}_{S_\lambda}(X_\lambda,Y_\lambda)\to\operatorname{Ps.hom}_S(X,Y)$ están definidas (sin hipótesis de separación sobre los $Y_\lambda$ ni sobre $Y$); por consiguiente, lo mismo ocurre con la aplicación \sref{IV.20.3.8.1}.
La proposición se deducirá del lema siguiente:

\begin{lemma}[20.3.8.2]
\label{IV.20.3.8.2}
Con las notaciones de \sref{IV.8.8.1}, supongamos que $X_\alpha$ sea cuasicompacto, y sea $U_\alpha$ un abierto de $X_\alpha$; sean $U_\lambda$ y $U$ sus imágenes recíprocas en $X_\lambda$ y $X$ para $\lambda\geq\alpha$.
Supongamos que se verifica una de las condiciones~(i),~(ii) de \sref{IV.20.3.8}.
Entonces, para que $U$ sea esquemáticamente denso en $X$, es necesario y suficiente que exista $\lambda\geq\alpha$ tal que $U_\lambda$ sea esquemáticamente denso en $X_\lambda$; en ese caso, $U_\mu$ es esquemáticamente denso en $X_\mu$ para $\mu\geq\lambda$.
\end{lemma}

Supongamos primero verificada la condición~(i); designemos por $j_\lambda:U_\lambda\to X_\lambda$ y $j:U\to X$ las inclusiones canónicas, y por $\sh{J}_\lambda$ el núcleo del homomorfismo canónico $\sh{O}_{X_\lambda}\to(j_\lambda)_*(\sh{O}_{U_\lambda})$.
Como la inmersión $j_\lambda$ es cuasicompacta y cuasiseparada, $(j_\lambda)_*(\sh{O}_{U_\lambda})$ es un $\sh{O}_{X_\lambda}$-módulo cuasicoherente; por tanto $\sh{J}_\lambda$ es un ideal cuasicoherente de $\sh{O}_{X_\lambda}$, y puesto que $X_\lambda$ es noetheriano, $\sh{J}_\lambda$ es coherente y, por consiguiente, de tipo finito.
Por otra parte, como el morfismo de transición $p_{\lambda\mu}:X_\mu\to X_\lambda$ (respectivamente $p_\lambda:X\to X_\lambda$) es plano, resulta de \sref{IV.2.3.1} que se puede identificar $(j_\mu)_*(\sh{O}_{U_\mu})$ con $p_{\lambda\mu}^*((j_\lambda)_*(\sh{O}_{U_\lambda}))$ (respectivamente, $j_*(\sh{O}_U)$ con $p_\lambda^*((j_\lambda)_*(\sh{O}_{U_\lambda}))$).
La afirmación resulta entonces de la definición de abierto esquemáticamente denso y de \sref{IV.8.5.8}[(ii)].

Para establecer \sref{IV.20.3.8.2} cuando se verifica la condición~(ii), probaremos primero dos lemas.

\begin{lemma}[20.3.8.3]
\label{IV.20.3.8.3}
Bajo las hipótesis de \sref{IV.8.2.2}, sea $M_\lambda$ (respectivamente $M$) el conjunto de los puntos máximos de $S_\lambda$ (respectivamente $S$).
Supongamos que para todo $\lambda$ el conjunto $M_\lambda$ sea finito, y que los $M_\lambda$ formen un sistema proyectivo de conjuntos.
Entonces $M$ es el límite proyectivo del sistema de los $M_\lambda$.
\end{lemma}

Mostremos primero que un punto $s\in\varprojlim M_\lambda$ es máximo en $S$: en efecto, si $s'$ es una generalización de $s$, la imagen $s'_\lambda$ de $s'$ en $S_\lambda$ es una generalización de la imagen $s_\lambda$ de $s$, y por tanto es igual a $s_\lambda$ para todo $\lambda$; esto implica $s'=s$, puesto que el conjunto subyacente a $S$ es el límite proyectivo de los conjuntos subyacentes a los $S_\lambda$ \sref{IV.8.2.9}.
Recíprocamente, sea $s$ un punto máximo de $S$ y probemos que pertenece a $\varprojlim M_\lambda$.
Sean $s_\lambda$ la imagen de $s$ en $S_\lambda$, y $M'_\lambda$ el conjunto de los puntos máximos de $S_\lambda$ que son generalizaciones de $s_\lambda$; los $M'_\lambda$ son conjuntos finitos no vacíos que forman un sistema proyectivo, por tanto $M'=\varprojlim M'_\lambda$ no es vacío y la aplicación $M'\to M'_\lambda$ es sobreyectiva (Bourbaki, \emph{Ens.}, cap.~III, 2\textsuperscript{a} ed., \textsection~7, n\textsuperscript{o}~4, Ejemplo~I).
Por otra parte, se tiene $\Spec(\sh{O}_{S,s})=\varprojlim\Spec(\sh{O}_{S_\lambda,s_\lambda})$ en virtud de \sref{IV.8.2.12} y \sref{IV.8.2.9}; por tanto, los puntos de $\varprojlim M'_\lambda$ son también puntos máximos de $\Spec(\sh{O}_{S,s})$, por la primera parte del razonamiento.
Así, $M'=\varprojlim M'_\lambda$ se reduce necesariamente al punto $s$; se concluye que $M'_\lambda$ se reduce al punto $s_\lambda$ y, por consiguiente, los $s_\lambda$ son máximos en los $S_\lambda$, lo que termina la demostración del lema.

\begin{lemma}[20.3.8.4]
\label{IV.20.3.8.4}
Siendo las hipótesis las de \sref{IV.20.3.8.3}, supongamos además que $S_\alpha$ sea cuasicompacto; sea $U_\alpha$ un conjunto abierto de $S_\alpha$, y sean $U_\lambda$ y $U$ sus imágenes recíprocas en $S_\lambda$ para $\lambda\geq\alpha$ y en $S$\footnote{La transcripción francesa introduce aquí $V_\alpha,V_\lambda,V$, pero el enunciado y toda la demostración emplean inmediatamente y de manera coherente $U_\alpha,U_\lambda,U$; el testigo inglés mantenido corrige la terna.}.
Si $U_\alpha$ es denso en $S_\alpha$, $U_\lambda$ es denso en $S_\lambda$ para $\lambda\geq\alpha$ y $U$ es denso en $S$.
Recíprocamente, si $U$ es denso en $S$ y además el conjunto $M$ de los puntos máximos de $S$ es finito, existe $\lambda\geq\alpha$ tal que $U_\lambda$ sea denso en $S_\lambda$.
\end{lemma}

En efecto, puesto que $M_\alpha$ es finito, la hipótesis de que $U_\alpha$ es denso en $S_\alpha$ implica que
\oldpage[IV-4]{241}
 $M_\alpha\subset U_\alpha$, luego $M_\lambda\subset U_\lambda$ para $\lambda\geq\alpha$ y $M\subset U$ en virtud de \sref{IV.20.3.8.3}, lo que prueba la primera afirmación.
Recíprocamente, supongamos $M$ finito y $U$ denso en $S$, luego $M\subset U$; observemos que los $U_\lambda$ son abiertos, por tanto ind-constructibles, y los $M_\lambda$ son finitos, por tanto pro-constructibles \sref{IV.1.9.6}.
La segunda afirmación resulta entonces de \sref{IV.8.3.2}.

\phantomsection\label{IV.20.3.8.5}
\emph{20.3.8.5. Fin de la demostración de \sref{IV.20.3.8.2}.}
Supongamos verificada la condición~(ii), y observemos que entonces $X$ es \emph{reducido} en virtud de \sref{IV.8.7.1}; por tanto, es equivalente decir que $U_\lambda$ (respectivamente $U$) es esquemáticamente denso en $X_\lambda$ (respectivamente $X$) o que es denso en $X_\lambda$ (respectivamente $X$), y la conclusión resulta de \sref{IV.20.3.8.4} aplicado a los $X_\lambda$ y a $X$.

\phantomsection\label{IV.20.3.8.6}
\emph{20.3.8.6. Fin de la demostración de \sref{IV.20.3.8}.}
Para mostrar que la aplicación \sref{IV.20.3.8.1} es inyectiva, consideremos dos morfismos $u'_\lambda$, $u''_\lambda$ de un abierto esquemáticamente denso $U_\lambda\subset X_\lambda$ en $Y_\lambda$, y supongamos que sus imágenes $u'$, $u''$, morfismos de $U$ en $Y$, coinciden en un abierto esquemáticamente denso $V$ de $U$.
Además, bajo cualquiera de las hipótesis~(i),~(ii), podemos suponer que $V$ es \emph{cuasicompacto}; esto es evidente si $X$ es noetheriano; en caso contrario, como $X$ solo tiene un número finito de puntos máximos y es reducido, basta sustituir $V$ por la unión de un número finito de entornos abiertos afines de esos puntos máximos (contenidos en $V$ por hipótesis).
Entonces existe un $\lambda\geq\alpha$ tal que $V$ sea la imagen recíproca de un abierto cuasicompacto $V_\lambda$ de $X_\lambda$ \sref{IV.8.2.11}, y resulta de \sref{IV.20.3.8.2} que, tomando $\lambda$ suficientemente grande, podemos suponer que $V_\lambda$ es esquemáticamente denso en $X_\lambda$.
Resulta entonces de \sref{IV.8.8.2}[(i)] que, tomando $\lambda$ suficientemente grande, $u'_\lambda$ y $u''_\lambda$ coinciden en $V_\lambda$ y, por tanto, pertenecen al mismo pseudomorfismo.\footnote{La transcripción francesa dice aquí por única vez «pseudo-homomorfismo»; el objeto considerado es el pseudomorfismo definido en esta sección, como confirma el testigo inglés mantenido.}

Probemos finalmente que la aplicación \sref{IV.20.3.8.1} es sobreyectiva.
Consideremos ahora un morfismo $u$ de un abierto esquemáticamente denso $U$ de $X$ en $Y$; como arriba, podemos suponer que $U$ es cuasicompacto y cuasiseparado ($U$ puede sustituirse, en el caso~(ii), por una unión de un número finito de abiertos afines dos a dos disjuntos).
Podemos, pues, suponer además que existe un $\lambda\geq\alpha$ tal que $U$ sea la imagen recíproca de un abierto cuasicompacto $U_\lambda$ de $X_\lambda$ \sref{IV.8.2.11}, automáticamente cuasiseparado, y por \sref{IV.20.3.8.2} podemos suponer además que $U_\lambda$ sea esquemáticamente denso en $X_\lambda$.
Puesto que se supone que los $Y_\lambda$ son localmente de presentación finita, resulta de \sref{IV.8.8.2}[(i)] que existe un $\lambda$ tal que $u$ sea la imagen de un morfismo $u_\lambda:U_\lambda\to Y_\lambda$; de donde se sigue la conclusión.

\begin{remarks}[20.3.8.7]
\label{IV.20.3.8.7}
\begin{enumerate}[label=(\roman*)]
  \item Para probar que la aplicación \sref{IV.20.3.8.1} es \emph{inyectiva}, no es necesario, bajo la hipótesis~(i) de \sref{IV.20.3.8}, suponer que $X$ sea noetheriano.
  En efecto, el lema \sref{IV.20.3.8.2} no utiliza esta hipótesis.
  Con las notaciones de \sref{IV.20.3.8.6}, sea $Z_\lambda$ el subpreesquema de coincidencias de $u'_\lambda$ y $u''_\lambda$, y sea $Z$ el subpreesquema de coincidencias de $u'$ y $u''$; resulta de la definición \sref{IV.17.4.5} y de \sref[I]{I.3.3.10.1} que $Z$ es el límite proyectivo de los $Z_\mu$ para $\mu\geq\lambda$.
  Ahora bien, por hipótesis, $Z$ contiene un abierto esquemáticamente denso en $X$; se deduce que $Z$ mismo es inducido sobre un abierto de $X$ en virtud del lema siguiente:

\begin{lemma}[20.3.8.8]
\label{IV.20.3.8.8}
Sea $T$ un preesquema.
Entonces todo subpreesquema $W$ de $T$ que contiene un abierto esquemáticamente denso $V$ de $T$ es inducido sobre un abierto (esquemáticamente denso) de $T$.
\end{lemma}

  En efecto, el subespacio de $T$ subyacente a $W$ es localmente cerrado, y por tanto abierto
  \oldpage[IV-4]{242}
  en su clausura, lo que prueba ya que el espacio subyacente a $W$ es abierto en $T$; la conclusión resulta entonces de \sref{IV.11.10.1}[(c)].

  Establecido este lema, se concluye que para $\mu\geq\lambda$ suficientemente grande, $Z_\mu$ es inducido sobre un abierto de $X_\mu$ en virtud de \sref{IV.8.6.3}, pues $Z_\lambda$, como subpreesquema de un preesquema noetheriano, es de presentación finita sobre $X_\lambda$ \sref{IV.1.6.2}, y por tanto lo mismo ocurre con los $Z_\mu$ para $\mu\geq\lambda$ sobre $X_\mu$ y con $Z$ sobre $X$.
  Podemos aplicar ahora \sref{IV.20.3.8.2}, que muestra que para $\mu\geq\lambda$ suficientemente grande, $Z_\mu$ es esquemáticamente denso en $X_\mu$, de donde se sigue la conclusión.

  \item Si, en la hipótesis~(i) de \sref{IV.20.3.8}, se suprime la condición de que $X$ sea noetheriano, el razonamiento de \sref{IV.20.3.8.6} muestra todavía que la imagen de \sref{IV.20.3.8.1} está formada por los pseudomorfismos-$S$ que tienen un representante que es un $S$-morfismo $U\to Y$, donde $U$ es esquemáticamente denso en $X$ y \emph{cuasicompacto y cuasiseparado}.
\end{enumerate}
\end{remarks}
\end{proof}

\begin{corollary}[20.3.9]
\label{IV.20.3.9}
Supongamos verificada en $X$ una u otra de las hipótesis~a),~b) de \sref{IV.20.3.6}, y que $Y$ sea localmente de presentación finita sobre $S$.
Entonces, para que un pseudomorfismo-$S$ $\omega$ de $X$ en $Y$ esté definido en el punto $x$ \sref{IV.20.2.3}, es necesario y suficiente que su restricción a $\Spec(\sh{O}_{X,x})$ esté definida en todas partes (es decir, sea un $S$-morfismo de $\Spec(\sh{O}_{X,x})$ en $Y$).
\end{corollary}

El resultado siguiente, que utilizaremos en la demostración de \sref{IV.20.3.11}, utiliza la teoría del descenso fielmente plano del cap.~VI.
El lector podrá comprobar que los resultados de \sref{IV.20.3} no se utilizarán en esta teoría.

\begin{lemma}[20.3.10]
\label{IV.20.3.10}
Sean $f:X'\to X$ un $S$-morfismo fielmente plano y cuasicompacto, $X''=X'\times_X X'$, $p_1$ y $p_2$ las proyecciones canónicas de $X''$ en $X'$, y $Y$ un preesquema separado sobre $S$.
Sean $U$ un abierto de $X$, $U'=f^{-1}(U)$, $U''=p_1^{-1}(U')=p_2^{-1}(U')$, y supongamos que $U''$ sea esquemáticamente denso en $X''$.
Sea $g:U\to Y$ un $S$-morfismo; entonces, si $g\circ(f|U')$ se prolonga a un $S$-morfismo $X'\to Y$, $g$ se prolonga a un $S$-morfismo $X\to Y$.
\end{lemma}

Obsérvese que las hipótesis implican que $U$ (respectivamente $U'$) es esquemáticamente denso en $X$ (respectivamente $X'$) \sref{IV.11.10.5}[(i)]; por tanto, también puede decirse que si $\omega$ denota el pseudomorfismo-$S$ clase de $g$, el enunciado de \sref{IV.20.3.10} significa que \emph{si $\omega\circ f$ está definido en todas partes, lo mismo ocurre con $\omega$}.

\begin{proof}
Para probar \sref{IV.20.3.10}, designemos por $g'$ un $S$-morfismo $X'\to Y$ que prolonga $g\circ(f|U')$, y pongamos $g''_i=g'\circ p_i:X''\to Y$ ($i=1,2$).
Si ponemos $f''=f\circ p_1=f\circ p_2:X''\to X$, es claro que $g''_1$ y $g''_2$ coinciden en $U''$ con $g\circ(f''|U'')$.
Pero como $Y$ es separado sobre $S$ y $U''$ es esquemáticamente denso en $X''$, se tiene $g''_1=g''_2$ \sref{IV.11.10.1}[(d)].
Como $f$ es fielmente plano y cuasicompacto, resulta de la teoría del descenso (cap.~VI) que existe un único $S$-morfismo $h:X\to Y$ tal que $h\circ f=g'$; como la restricción $U'\to U$ de $f$ es un morfismo fielmente plano y cuasicompacto y $U''=U'\times_U U'$, el resultado de unicidad anterior, aplicado a $U$ en lugar de $X$, muestra que $h|U=g$, lo que prueba el lema.
\end{proof}

\begin{proposition}[20.3.11]
\label{IV.20.3.11}
Sean $Y$ un $S$-preesquema separado sobre $S$, $\omega$ un pseudomorfismo-$S$ de $X$ en $Y$, y $f:X'\to X$ un $S$-morfismo.
Supongamos que $f$ sea plano y que se cumpla una de las condiciones siguientes:
\begin{enumerate}[label=\emph{(\roman*)}]
\oldpage[IV-4]{243}
  \item $f$ es un morfismo abierto, o sobreyectivo y cuasicompacto, y $\operatorname{dom}_S(\omega)$ contiene un abierto $U$ esquemáticamente denso en $X$ y retrocompacto en $X$.
  \item $f$ es localmente de presentación finita.
  \item $Y$ es localmente de presentación finita sobre $S$, y $X$ satisface una de las condiciones~a),~b) de \sref{IV.20.3.6}.
\end{enumerate}
Entonces, $f^{-1}(\operatorname{dom}_S(\omega))$ es esquemáticamente denso en $X'$, de modo que $\omega\circ f$ está definido \sref{IV.20.3.2} y se tiene
\[
\label{IV.20.3.11.1}
  \operatorname{dom}_S(\omega\circ f)=f^{-1}(\operatorname{dom}_S(\omega)).
  \tag{20.3.11.1}
\]
\end{proposition}

\begin{proof}
Probemos primero que $f^{-1}(\operatorname{dom}_S(\omega))$ es esquemáticamente denso en $X'$.
Como la cuestión es local en $X$ y $X'$, podemos suponer que $X$ y $X'$ son afines, y puesto que $f$ es plano, basta ver, en virtud de \sref{IV.11.10.5}[(ii), (a)], que $\operatorname{dom}_S(\omega)$ contiene un conjunto abierto $U$ retrocompacto y esquemáticamente denso en $X$.
Esto resulta de la hipótesis en el caso~(i), y de \sref{IV.20.2.12} en el caso~(iii), teniendo en cuenta que en un esquema afín todo abierto de la forma $\mathrm{D}(t)$ es retrocompacto; finalmente, en el caso~(ii), se ve directamente que $f^{-1}(\operatorname{dom}_S(\omega))$ es esquemáticamente denso en $X'$ aplicando \sref{IV.20.3.5}[(i)].

Probemos ahora \sref{IV.20.3.11.1}, es decir, que para todo punto $x'\in\operatorname{dom}_S(\omega\circ f)$ se tiene $x=f(x')\in\operatorname{dom}_S(\omega)$.
Observemos primero que podemos reducirnos al caso en que $f$ sea \emph{fielmente plano y cuasicompacto}.
Esta es, en efecto, la hipótesis en el segundo caso de~(i); en los demás casos, la cuestión es local en $X'$, y podemos suponer ya que $X$ y $X'$ son afines, de modo que $f$ es cuasicompacto.
En el primer caso de~(i) y en el caso~(ii), $f$ es un morfismo abierto \sref{IV.11.3.1}; por tanto, sustituyendo $X$ por el abierto $f(X')$, podemos suponer que $f$ es sobreyectivo y, por consiguiente, fielmente plano.
En el caso~(iii), utilizando \sref{IV.20.3.9}, podemos limitarnos a probar que la restricción de $\omega$ a $\Spec(\sh{O}_{X,x})$ está definida en todas partes, y podemos sustituir $X$ por $\Spec(\sh{O}_{X,x})$, $X'$ por $X'\times_X\Spec(\sh{O}_{X,x})$ y $f$ por su restricción a este último preesquema, que es un morfismo sobreyectivo \sref{IV.2.3.4}, y por tanto fielmente plano.

Supongamos, pues, que $f$ sea fielmente plano y cuasicompacto; con las notaciones del lema \sref{IV.20.3.10}, basta ver entonces que $U''$ es esquemáticamente denso en $X''$, tomando para $U$ un abierto esquemáticamente denso en $X$, contenido en $\operatorname{dom}_S(\omega)$; esto será así, en virtud de \sref{IV.11.10.5}[(ii), (a)], si $U$ se toma retrocompacto en $X$ (puesto que el morfismo $f'':X''\to X$ es plano).
Ahora bien, la existencia de tal abierto $U$ forma parte de la hipótesis en el caso~(i); en el caso~(iii) resulta de \sref{IV.20.2.12} y del hecho de que en un esquema afín $\Spec(A)$ todo abierto de la forma $\mathrm{D}(t)$ es retrocompacto.
Finalmente, en el caso~(ii), tomemos $U=\operatorname{dom}_S(\omega)$ y mostremos directamente que $U''$ es esquemáticamente denso en $X''$; para ello basta observar que $f'':X''\to X$ es plano y localmente de presentación finita, y aplicar \sref{IV.11.10.5}[(ii), (b)].
\end{proof}

\begin{corollary}[20.3.12]
\label{IV.20.3.12}
Sea $\vphi$ una pseudofunción sobre un preesquema $X$.
Entonces, para todo morfismo plano y localmente de presentación finita $f:X'\to X$, la pseudofunción $\vphi\circ f$ está definida y se tiene $\operatorname{dom}(\vphi\circ f)=f^{-1}(\operatorname{dom}(\vphi))$.
\end{corollary}

\begin{remark}[20.3.13]
\label{IV.20.3.13}
\oldpage[IV-4]{244}
Cuando $X$ satisface una de las condiciones~a),~b) de \sref{IV.20.3.6}, vimos en \sref{IV.20.2.11} que las pseudofunciones sobre $X$ se identifican con las funciones meromorfas sobre $X$.
En virtud de \sref{IV.20.1.12} y de \sref{IV.20.2.13}[(i)], si se supone solamente que el morfismo $f:X'\to X$ es \emph{plano}, entonces, para toda pseudofunción $\vphi$ sobre $X$, $\vphi\circ f$ está definida y se tiene $\operatorname{dom}(\vphi\circ f)=f^{-1}(\operatorname{dom}(\vphi))$.
\end{remark}

\subsection{Propiedades de los dominios de definición de las funciones racionales}
\label{IV.20.4}

\begin{env}[20.4.1]
\label{IV.20.4.1}
Sean $X$, $Y$ dos $S$-preesquemas, y sea $\omega$ un pseudo-$S$-morfismo de $X$ en $Y$.
Sea $u$ un $S$-morfismo $U\to Y$ perteneciente a $\omega$, donde $U$ es esquemáticamente denso en $X$, y consideremos el \emph{grafo} $\Gamma_u$, subpreesquema de $U\times_S Y$ \sref[I]{I.5.3.11}, y por tanto subpreesquema de $X\times_S Y$.
Supongamos que este subpreesquema admite una \emph{clausura} \sref[I]{I.9.5.11} en $X\times_S Y$; si $j:\Gamma_u\to X\times_S Y$ es la inyección canónica, así ocurre cuando el $\sh{O}_{X\times_S Y}$-módulo $j_*(\sh{O}_{\Gamma_u})$ es cuasicoherente; de la definición de la clase de equivalencia $\omega$ \sref{IV.20.2.1} se deduce que $j_*(\sh{O}_{\Gamma_u})$ no depende del representante $u$ considerado, por lo que la clausura preesquemática de $\Gamma_u$ queda entonces bien determinada por $\omega$; se dice que este subpreesquema cerrado de $X\times_S Y$ es el \emph{grafo} del pseudo-$S$-morfismo $\omega$, y se denota por $\Gamma_\omega$.
Obsérvese que $\Gamma_\omega$ está definido si existe un morfismo $u:U\to Y$ de la clase $\omega$ tal que $U$ sea \emph{retrocompacto} en $X$ y si, además, $Y$ es cuasiseparado sobre $S$, pues entonces la inyección $j$ es un morfismo cuasicompacto y cuasiseparado (\sref[I]{I.2.2} y \sref[I]{I.7.4}); la primera hipótesis se cumple siempre cuando $X$ es \emph{localmente noetheriano}.
Obsérvese también que, cuando $X$ es \emph{reducido}, lo mismo ocurre con $\Gamma_u$, que es isomorfo a $U$ \sref[I]{I.5.3.11}; entonces $\Gamma_\omega$ existe y no es sino el subpreesquema \emph{reducido} de $X\times_S Y$ cuyo espacio subyacente es la clausura, en $X\times_S Y$, del espacio subyacente a $\Gamma_u$ (\sref[I]{I.5.2.1} y \sref[I]{I.5.2.2}).
Obsérvese finalmente que, si $Y$ es \emph{separado} sobre $S$, $\Gamma_u$ es \emph{cerrado} en $U\times_S Y$ \sref[I]{I.5.4.3}, y por tanto está \emph{inducido por} $\Gamma_\omega$ (cuando este último existe) sobre el abierto $\Gamma_\omega\cap(U\times_S Y)$ \sref[I]{I.9.5.10}; por el contrario, este preesquema inducido es en general distinto de $\Gamma_u$ cuando $Y$ no es separado.
En particular, si $v:X\to Y$ es un $S$-morfismo, el grafo $\Gamma_\omega$ de la clase $\omega$ de $v$ puede ser distinto del grafo $\Gamma_v$.
Por ello, en lo que sigue, cuando se trate del grafo de un pseudo-$S$-morfismo, nos limitaremos al caso en que \emph{$Y$ es separado sobre $S$}.
\end{env}

\begin{env}[20.4.2]
\label{IV.20.4.2}
Supongamos, pues, que $\Gamma_\omega$ está definido y que $Y$ es separado sobre $S$; denotemos por $p$ y $q$ las restricciones a $\Gamma_\omega$ de las proyecciones canónicas
\[
  \xymatrix{
  & \Gamma_\omega \ar[dl]_p \ar[dr]^q & \\
  X & & Y
  }
\]
Entonces, si $U\subset\operatorname{dom}_S(\omega)$, la restricción $p^{-1}(U)\to U$ de $p$ es un isomorfismo \sref[I]{I.5.3.11}; recíprocamente, si $U$ es un abierto de $X$ que tiene esta propiedad, y si $u$ es el isomorfismo inverso de la restricción $p^{-1}(U)\to U$ de $p$, $q\circ u$ es un $S$-morfismo de $U$ en $Y$ que coincide con un morfismo de la clase $\omega$ en $U\cap\operatorname{dom}_S(\omega)$.
Se concluye que $\operatorname{dom}_S(\omega)$ es el \emph{mayor abierto} $U$ de $X$ tal que la restricción $p^{-1}(U)\to U$ de $p$ sea un
\oldpage[IV-4]{245}
isomorfismo.
Sea $v:\operatorname{dom}_S(\omega)\to\Gamma_\omega$ el $S$-morfismo inverso de la restricción $p^{-1}(\operatorname{dom}_S(\omega))\to\operatorname{dom}_S(\omega)$ de $p$; se denota a veces por $p^{-1}$ el pseudo-$S$-morfismo de $X$ en $\Gamma_\omega$, clase de $v$ (cuya aplicación racional asociada \sref{IV.20.2.13}[(ii)] es entonces \emph{birracional}); como $p^{-1}(\operatorname{dom}_S(\omega))$ es el grafo del $S$-morfismo $u=q\circ v:\operatorname{dom}_S(\omega)\to Y$, es esquemáticamente denso en $\Gamma_\omega$ \sref{IV.11.10.3}[(iv)], por lo que $\omega$ puede considerarse como el \emph{compuesto} $q\circ p^{-1}$ \sref{IV.20.3.2}.
Para toda parte $M$ del espacio subyacente a $X$, se pone a veces $\omega(M)=q(p^{-1}(M))$, lo que equivale a considerar $\omega$ como una aplicación de $X$ en $\mathfrak{P}(Y)$ (o, como dicen algunos autores, una «función multiforme»).
Obsérvese que, cuando $x\in\operatorname{dom}_S(\omega)$, $\omega(\{x\})$ es el conjunto $\{u(x)\}$; en general, para un $x\in X$ cualquiera, si se denota por $s$ la imagen de $x$ en $S$ y por $Y_s$ la fibra en el punto $s$ del morfismo estructural $Y\to S$, el conjunto $\omega(\{x\})$ es una parte del preesquema $Y_s$.
\end{env}

\begin{env}[20.4.3]
\label{IV.20.4.3}
En todo el resto de este número nos limitaremos al caso en que $X$ sea \emph{reducido}, de modo que hay entonces identidad entre los pseudo-$S$-morfismos y las $S$-\emph{aplicaciones racionales} \sref{IV.20.2.7}.
Además, el \emph{grafo} de toda $S$-aplicación racional de $X$ en $Y$ está entonces definido \sref{IV.20.4.1}.
\end{env}

\begin{proposition}[20.4.4]
\label{IV.20.4.4}
Sean $X$ un $S$-preesquema localmente íntegro, $Y$ un $S$-preesquema localmente de tipo finito y separado sobre $S$, $\omega$ una $S$-aplicación racional de $X$ en $Y$, y $p:\Gamma_\omega\to X$ la proyección canónica.
Entonces, si $x\in X$ es un punto normal tal que el conjunto $p^{-1}(x)$ contiene un punto aislado, $\omega$ está definida en el punto $x$.
\end{proposition}

\begin{proof}
En efecto, la primera proyección $p_1:X\times_S Y\to X$ es entonces un morfismo separado y localmente de tipo finito, por lo que lo mismo sucede con su restricción $p:\Gamma_\omega\to X$, que es además \emph{birracional}; y, como $\Gamma_\omega$ es reducido y $X$ íntegro, $\Gamma_\omega$ es \emph{íntegro}. Se deduce, pues, de (\textbf{Err}\textsubscript{\textbf{IV}}, 30) que la hipótesis sobre $x$ implica la existencia de un entorno abierto $U$ de $x$ tal que la restricción $p^{-1}(U)\to U$ de $p$ sea un isomorfismo; de ahí la conclusión por \sref{IV.20.4.2}.

Obsérvese que el enunciado \sref{IV.20.4.4} es la formulación original dada por Zariski de su «teorema principal» (salvo que originalmente estaba restringido a los esquemas algebraicos sobre un cuerpo base).
\end{proof}

\begin{proposition}[20.4.5]
\label{IV.20.4.5}
Bajo las hipótesis y notaciones de \sref{IV.20.4.4}, supongamos además que $X$ es normal, y sean $X'$ un preesquema reducido y $f:X'\to X$ un morfismo localmente de tipo finito y universalmente abierto.
Entonces $\omega\circ f$ está definida, y se tiene $\operatorname{dom}_S(\omega\circ f)=f^{-1}(\operatorname{dom}_S(\omega))$ (en otras palabras, si $x'\in X'$ y $x=f(x')$, para que $\omega$ esté definida en el punto $x$ es necesario y suficiente que $\omega\circ f$ lo esté en el punto $x'$).
\end{proposition}

\begin{proof}
Como $X'$ es reducido y $f^{-1}(\operatorname{dom}_S(\omega))$ es denso en todo $X'$ en virtud de \sref{IV.2.4.11}, el compuesto $\omega\circ f$ está definido; para demostrar que, cuando $\omega\circ f$ está definida en el punto $x'$, $\omega$ lo está en el punto $x$, podemos sustituir evidentemente $X'$ por un entorno abierto de $x'$ y suponer, por tanto, que $\omega\circ f$ está \emph{definida en todas partes}; además, como $f(X')$ es abierto en $X$, podemos suponer que $f$ es \emph{sobreyectivo}.
En virtud de la hipótesis sobre $f$, el morfismo $f_{(Y)}:X'\times_S Y\to X\times_S Y$ es entonces abierto, por lo que $\overline{f_{(Y)}^{-1}(M)}=f_{(Y)}^{-1}(\overline{M})$ para toda parte $M$ de $X\times_S Y$; tomando como $M$ el conjunto $\Gamma_u$, donde $u:\operatorname{dom}_S(\omega)\to Y$ es la restricción de $\omega$ a $\operatorname{dom}_S(\omega)$, se deduce de la relación anterior y de \sref[I]{I.3.3.12} que el conjunto subyacente a $\Gamma_{\omega\circ f}$ es igual a $f_{(Y)}^{-1}(\Gamma_\omega)$; como se sabe
\oldpage[IV-4]{246}
que $\Gamma_{\omega\circ f}$ es un preesquema reducido \sref{IV.20.4.1}, se ve que el $X'$-preesquema $\Gamma_{\omega\circ f}$ es igual a $(\Gamma_\omega\times_X X')_\mathrm{red}$.
Pero, como por hipótesis el morfismo compuesto $\Gamma_{\omega\circ f}\to\Gamma_\omega\times_X X'\xrightarrow{p^{(X')}}X'$ es un isomorfismo, $p^{(X')}$ es necesariamente \emph{radicial}; como $f$ es sobreyectivo, lo mismo sucede con $p:\Gamma_\omega\to X$ \sref[I]{I.3.4.8}, por lo que $p^{-1}(x)$ es un conjunto reducido a un punto \sref[I]{I.3.6.4}; se deduce entonces de \sref{IV.20.4.4} que $\omega$ está definida en el punto $x$.
\end{proof}

La proposición siguiente da un \emph{criterio valuativo} para que una aplicación racional esté definida en un punto:

\begin{proposition}[20.4.6]
\label{IV.20.4.6}
Sean $S$ un preesquema, y $X$, $Y$ dos $S$-preesquemas; se supone que $X$ es localmente noetheriano y que $Y$ es localmente de tipo finito sobre $S$.
Sean $U$ un abierto denso de $X$, $h:U\to Y$ un $S$-morfismo, $x\in X-U$ un punto normal de $X$, y $h'_x:\Spec(k(x))\to Y$ un $S$-morfismo.
Para que $h$ pueda prolongarse a un $S$-morfismo $h':U'\to Y$, donde $U'$ es un abierto de $X$ que contiene $U$ y $x$, y donde el morfismo compuesto $\Spec(k(x))\to U'\xrightarrow{h'}Y$ es el $S$-morfismo dado $h'_x$, es necesario y suficiente que se cumpla la condición siguiente:
\begin{quote}
\emph{(P) Para todo espectro $S_1$ de un anillo de valoración discreta, con punto cerrado $a$ y punto genérico $b$, y todo morfismo $u:S_1\to X$ tal que $u(a)=x$ y $u(b)\in U$, el morfismo compuesto $h_1=h\circ(u|\{b\}):\{b\}=u^{-1}(U)\to Y$ se prolonga a un morfismo $h'_1:S_1\to Y$ tal que el diagrama}
\[
  \xymatrix{
  \Spec(k(a)) \ar[r] & S_1 \ar[d]^{h'_1} \\
  \Spec(k(x)) \ar[r]_{h'_x} & Y
  }
\]
\emph{sea conmutativo.}
\end{quote}
Además, si se cumple esta condición y si $h'':U''\to Y$ es un morfismo que satisface las mismas condiciones que $h'$, existe un conjunto abierto que contiene $U\cup\{x\}$ y en el cual $h'$ y $h''$ coinciden.
\end{proposition}

\begin{proof}
Demostremos primero la última afirmación; podemos suponer que $h'$ y $h''$ están definidos en el mismo abierto $U_0\supset U\cup\{x\}$.
El subpreesquema $Z$ de coincidencia de $h'$ y $h''$ \sref{IV.17.4.5} contiene $U$ y $x$, por lo que existe un entorno abierto $V$ de $x$ en $U_0$ tal que el subpreesquema inducido por $Z$ sobre el abierto $Z\cap V$ sea un subpreesquema \emph{cerrado} del preesquema inducido por $X$ sobre $V$; como este preesquema $Z\cap V$ contiene el subpreesquema inducido por $V$ sobre el abierto $U\cap V$, y este último es esquemáticamente denso en $V$, $Z\cap V$ es necesariamente igual a $V$ \sref{IV.20.3.8.8}, lo que demuestra que $h'$ y $h''$ coinciden en $U\cup V$.

Siendo evidente la necesidad de la condición del enunciado, demostremos que es suficiente.
En virtud de la correspondencia biyectiva entre los $S$-morfismos de $X$ en $Y$ y las $X$-secciones de $X\times_S Y$ \sref[I]{I.3.3.15}, podemos limitarnos al caso en que $S=X$ y en que $h$ es, por tanto, una $U$-sección de $Y$; obsérvese que entonces $Y$ es localmente noetheriano, y podemos evidentemente (puesto que $X$ es localmente íntegro y localmente noetheriano) suponer que $X=S$ es irreducible.
Además, teniendo en cuenta \sref{IV.20.3.7}, podemos suponer que $X=\Spec(A)$, donde $A$ es un anillo \emph{local} noetheriano, íntegro e integralmente cerrado.

\oldpage[IV-4]{247}
Pongamos $y=h'_x(x)$. Para todo $x'\in U$, $y$ es una especialización de $h(x')$.\footnote{La transcripción francesa invierte aquí la relación de especialización; el testigo inglés mantenido y las igualdades $h(x')=h'_1(b)$, $y=h'_1(a)$ muestran que $y$ es la especialización de $h(x')$.}
En efecto, existe un espectro $S_1$ de un anillo de valoración discreta y un morfismo $u:S_1\to X$ tal que $u(a)=x$ y $u(b)=x'$ \sref[II]{II.7.1.9}.
Aplicando la condición del enunciado, se obtiene inmediatamente nuestra afirmación, pues $h(x')=h_1(b)=h'_1(b)$ e $y=h'_1(a)$.
Si $Y'$ es un entorno abierto afín de $y$, se tiene, por tanto, $h(X\cap U)\subset Y'$; podemos sustituir entonces $Y$ por $Y'$, es decir, suponer que $Y$ es \emph{afín}, y por tanto separado sobre $X$.
Sea $\omega$ la $X$-sección racional de $Y$ a la que pertenece $h$, de modo que su grafo tiene aquí por conjunto subyacente la clausura de $h(U)$ en $Y$.
Como $Y$ es separado sobre $X$, podemos aplicar \sref{IV.20.4.4} a $\omega$: bastará demostrar que, si $p:\Gamma_\omega\to X$ es la proyección canónica, $p^{-1}(x)$ se reduce al único punto $y$ y que $h'_x(x)=y$.
En efecto, por \sref{IV.20.4.4}, $h$ se prolongará a una sección $h'$ sobre un abierto $U'$ que contiene $U\cup\{x\}$, tal que $h'(x)=y$; y, como no existe más que un $X$-morfismo $\Spec(k(x))\to Y$ que transforma $x$ en $y$, esto demostrará la identidad de $h'_x$ con el compuesto de $h'$ y $\Spec(k(x))\to U'$.

Como, para todo $x'\in X\cap U$, $y$ es una especialización de $h(x')$, se tiene $y\in p^{-1}(x)$.
Supongamos ahora que $z$ es un punto cualquiera de $p^{-1}(x)$.
Como $\Gamma_\omega$ es la clausura de $h(U)$ y $h(U)$ está formado por puntos adherentes a $h(\xi)$, donde $\xi$ es el punto genérico de $X$, $\Gamma_\omega$ es la clausura de $h(\xi)$ en $Y$.
Se toma entonces un espectro $S_1$ de un anillo de valoración discreta y un morfismo $v:S_1\to Y$ tales que $v(a)=z$ y $v(b)=h(\xi)$ \sref[II]{II.7.1.9}, y se pone $u=p\circ v$, de modo que $u(a)=x$ y $u(b)=\xi$.
Aplicando a $u$ la condición del enunciado, se obtiene un morfismo $h'_1:S_1\to Y$ tal que $h'_1(a)=y$ y $h'_1(b)=h(\xi)$; pero esto implica $z=y$ en virtud de \sref[II]{II.7.2.3}, puesto que $Y$ es separado sobre $X$ y, por tanto, $v$ y $h'_1$ deben coincidir, ya que son iguales en el punto $b$.
C.Q.F.D.
\end{proof}

\begin{corollary}[20.4.7]
\label{IV.20.4.7}
Bajo las hipótesis sobre $S$, $X$, $Y$, $U$ y $h$ de \sref{IV.20.4.6}, sea $E$ una parte de $X-U$ tal que $X$ sea normal en todo punto de $E$, y, para cada $x\in E$, sea $h'_x:\Spec(k(x))\to Y$ un $S$-morfismo tal que se cumpla la condición~(P) de \sref{IV.20.4.6}.
Supongamos además que la unión $F$ de los grafos de los $h'_x$ (identificados con partes de $X\times_S Y$), para $x\in E$, está contenida en la unión de un número finito de abiertos $V_i$ de $X\times_S Y$ que son \emph{separados} sobre $X$ (condición que se cumple automáticamente si $Y$ es \emph{separado} sobre $S$, o si $X$ es \emph{cuasicompacto} e $Y$ de \emph{tipo finito} sobre $S$).
Entonces existe un $S$-morfismo $h':U'\to Y$, donde $U'$ es un abierto de $X$ que contiene $U\cup E$, tal que, para todo $x\in E$, el compuesto
\[
  \Spec(k(x))\to U'\xrightarrow{h'}Y
\]
sea igual a $h'_x$.
\end{corollary}

\begin{proof}
Obsérvese primero que, en \sref{IV.20.4.6}, si se supone que $Y$ es \emph{separado}, existe un mayor abierto $U_0\supset U$, igual al dominio de la $S$-aplicación racional correspondiente a $h$, en el que $h$ puede prolongarse, y este prolongamiento es único \sref[I]{I.7.2.2}; de ahí la conclusión en este caso, en virtud de \sref{IV.20.4.6}.
En el caso general, sea $E_i$ el conjunto de los $x\in E$ tales que $(x,h'_x(x))\in V_i$.
En virtud de \sref{IV.20.4.6}, para todo $x\in E$ existe un prolongamiento $h^{(x)}$ de $h$ a $U\cup W^{(x)}$, donde $W^{(x)}$ es un entorno abierto de $x$ en $X$ tal que el grafo de $h^{(x)}|W^{(x)}$ está contenido en \emph{todos} los $V_i$ tales que $x\in E_i$.
Como $V_i$ es separado sobre $X$ y $X$ es reducido, para dos puntos $x'$, $x''$ de $E_i$, $h^{(x')}$ y $h^{(x'')}$ coinciden en $W^{(x')}\cap W^{(x'')}$, puesto que
\oldpage[IV-4]{248}
coinciden en el abierto denso en todas partes $U\cap W^{(x')}\cap W^{(x'')}$ de este conjunto.
Existe, pues, un morfismo $h_i:U\cup U_i\to Y$ que prolonga $h$ a un abierto $U\cup U_i$ que contiene $E_i$; además, para todo par de índices $i$, $j$, los grafos de las restricciones $h_i|(U_i\cap U_j)$ y $h_j|(U_i\cap U_j)$ están contenidos en $V_i\cap V_j$; como $V_i\cap V_j$ es separado sobre $X$ y los morfismos anteriores coinciden en un abierto denso en todas partes $U\cap U_i\cap U_j$ de $U_i\cap U_j$, son iguales.
El morfismo $h'$, igual a $h$ en $U$ y a $h_i$ en cada uno de los $U_i$, responde a la cuestión.
\end{proof}

\begin{remark}[20.4.8]
\label{IV.20.4.8}
Cuando $E=X-U$, es claro que, si para todo abierto afín $T$ de $X$ existe un $S$-morfismo $h'_T:T\to Y$ que prolonga $h|(U\cap T)$ y tal que el compuesto $\Spec(k(x))\to T\xrightarrow{h'_T}Y$ es igual a $h'_x$ para todo $x\in T-(U\cap T)$, entonces los $h'_T$ son las restricciones de un $S$-morfismo $h':X\to Y$ definido en todas partes y que prolonga $h$, en virtud de la afirmación de unicidad de \sref{IV.20.4.6}.
Para demostrar la existencia de $h'$, nos reducimos, pues, al caso en que $X$ es \emph{afín}, y basta entonces que el conjunto $F$, unión de los grafos de los $h'_x$, sea \emph{cuasicompacto} en $X\times_S Y$ para poder aplicar \sref{IV.20.4.7}.
Así sucede cuando los $h'_x$ son de la forma $\Spec(k(x))\to Z\xrightarrow{h_0}Y$, donde $Z$ es un subpreesquema cerrado de $X$ cuyo espacio subyacente es $X-U$, y $h_0$ es un $S$-morfismo.
\end{remark}

\begin{corollary}[20.4.9]
\label{IV.20.4.9}
Sean $S$ un preesquema localmente noetheriano, $X$ un preesquema localmente noetheriano, $f:X\to S$ un morfismo plano y $g:Y\to S$ un morfismo localmente de tipo finito.
Sean $U$ un abierto denso de $S$, $h:f^{-1}(U)\to Y$ un $S$-morfismo, $T=S-U$, $Z$ el subpreesquema cerrado reducido de $X$ cuyo espacio subyacente es $f^{-1}(T)$, y $h_0:Z\to Y$ un $S$-morfismo.
Supongamos que $X$ es normal en todos los puntos de $Z$.
Para que exista un $S$-morfismo (necesariamente único) $h':X\to Y$ que prolongue $h$ y $h_0$, es necesario y suficiente que se cumpla la condición siguiente:

\emph{Para todo espectro $S_1$ de un anillo de valoración discreta, con punto cerrado $a$ y punto genérico $b$, y todo morfismo $u:S_1\to S$ tal que $u(a)\in T$ y $u(b)\in U$, existe un $S_1$-morfismo $h'_1:X_{(S_1)}\to Y_{(S_1)}$ que prolonga $h_{(S_1)}:f_{(S_1)}^{-1}(b)\to Y_{(S_1)}$ y es tal que el diagrama}
\[
  \xymatrix{
  \Spec(k(z_1)) \ar[r] \ar[d] & Z_{(S_1)} \ar[d]^{(h_0)_{(S_1)}} \\
  X_{(S_1)} \ar[r]_{h'_1} & Y_{(S_1)}
  }
\]
\emph{sea conmutativo para todo $z_1\in Z_{(S_1)}$.}

La hipótesis de que $f$ es plano implica, en efecto, que $f^{-1}(U)$ es denso en $X$ \sref{IV.2.3.10}, y basta entonces aplicar \sref{IV.20.4.8}.
\end{corollary}

\begin{corollary}[20.4.10]
\label{IV.20.4.10}
Bajo las hipótesis de \sref{IV.20.4.6}, supongamos además que $Y$ es separado y localmente cuasifinito sobre $S$.
Sean $U$ un abierto denso de $X$, $h:U\to Y$ un $S$-morfismo, $\omega$ la $S$-aplicación racional correspondiente, y $x\in X-U$ un punto normal de $X$.
Para que $\omega$ esté definida en el punto $x$, es necesario y suficiente que se cumpla la condición siguiente: existe un espectro $S_1$ de un anillo de valoración discreta, con punto cerrado $a$ y punto genérico $b$, y un morfismo $u:S_1\to X$
\oldpage[IV-4]{249}
tales que $u(a)=x$, $u(b)\in U$ y la $S$-aplicación racional $\omega\circ u$ está definida en todas partes.

En efecto, por hipótesis todas las fibras del morfismo de proyección $X\times_S Y\to X$ están formadas por puntos aislados, y para aplicar \sref{IV.20.4.4} basta saber que la fibra $p^{-1}(x)$ no es vacía en $\Gamma_\omega$.
Ahora bien, si $h_1$ es el único morfismo de la clase $\omega\circ u$, el único punto de $X\times_S Y$ situado sobre $x$ y sobre $h_1(a)$ pertenece a $\Gamma_\omega$, de donde se obtiene la conclusión.
\end{corollary}

\begin{proposition}[20.4.11]
\label{IV.20.4.11}
Sean $X$ un $S$-preesquema localmente noetheriano,\footnote{La transcripción francesa omite aquí la estructura de $S$-preesquema sobre $X$, necesaria para que los morfismos del enunciado sean $S$-morfismos y usada explícitamente en la demostración.} $Y$ un $S$-preesquema afín sobre $S$, $U$ un abierto de $X$, y $Z=X-U$.
Supongamos que $\operatorname{prof}_Z(\sh{O}_X)\geq 2$ \sref{IV.5.10.1}; entonces todo $S$-morfismo $f:U\to Y$ se prolonga de manera única a un $S$-morfismo de $X$ en $Y$.
\end{proposition}

\begin{proof}
Podemos limitarnos al caso en que $S$ y $X$ sean afines y (en virtud de \sref[I]{I.3.3.14}) al caso en que $S=X$; se tiene, pues, $X=\Spec(A)$, $Y=\Spec(B)$, siendo $B$ una $A$-álgebra.\footnote{La transcripción francesa y el testigo inglés añaden aquí «de tipo finito», hipótesis que no figura en el enunciado ni se necesita: toda $A$-álgebra es cociente de un álgebra de polinomios en una familia posiblemente infinita de indeterminadas.}
Como $B$ es cociente de un álgebra de polinomios $A[(T_\lambda)]_{\lambda\in L}$, $Y$ es un subpreesquema cerrado de $Y'=X[(T_\lambda)]_{\lambda\in L}$.
Por otra parte, la hipótesis sobre $Z$ implica que $U$ es esquemáticamente denso en $X$ en virtud de \sref{IV.20.2.13}[(iv)] y \sref{IV.5.10.2}.
Si se demuestra que todo $X$-morfismo $u:U\to Y$ se prolonga de manera única a un $X$-morfismo $v':X\to Y'$, se deducirá que $v'$ se factoriza como $X\xrightarrow{v}Y\to Y'$; en efecto, el subpreesquema $v'^{-1}(Y)$ es cerrado y contiene $U$ \sref[I]{I.4.4.1}, por lo que es idéntico a $X$ \sref{IV.20.3.8.8}.
En estas condiciones, $v$ será el único $S$-morfismo de $X$ en $Y$ que prolonga $u$.
Podemos, pues, limitarnos al caso en que $Y=Y'$.
Pero entonces hay una correspondencia biyectiva entre los $X$-morfismos de un abierto $U\subset X$ en $Y'$ y las familias $(s_\lambda)_{\lambda\in L}$ de secciones de $\sh{O}_X$ sobre $U$ \sref[II]{II.1.7.9}; la conclusión se deduce, pues, de \sref{IV.5.10.5}.
\end{proof}

\begin{corollary}[20.4.12]
\label{IV.20.4.12}
Sean $X$ un $S$-preesquema localmente noetheriano, reducido y que satisface la condición $(S_2)$ \sref{IV.5.7.2} (por ejemplo, un $S$-preesquema localmente noetheriano y normal \sref{IV.5.8.6}), $Y$ un $S$-preesquema afín sobre $S$, y $f$ una $S$-aplicación racional de $X$ en $Y$; entonces toda componente irreducible de $X-\operatorname{dom}(f)$ es de codimensión $1$ en $X$.
\end{corollary}

\begin{proof}
Esto equivale a decir que, si $Z^{(2)}$ es el conjunto de los $x\in X$ tales que $\dim(\sh{O}_{X,x})\geq 2$, entonces, para toda parte cerrada $Z\subset Z^{(2)}$ de $X$, todo $S$-morfismo de $X-Z$ en $Y$ se prolonga a un $S$-morfismo de $X$ en $Y$; ahora bien, esto se deduce de la hipótesis sobre $X$ \sref{IV.5.7.2} y de \sref{IV.20.4.11}.
\end{proof}

\subsection{Pseudomorfismos relativos}
\label{IV.20.5}

\begin{env}[20.5.1]
\label{IV.20.5.1}
Sean $X$, $Y$ dos preesquemas sobre $S$.
Se deduce de las definiciones \sref{IV.11.10.8} que la intersección de dos abiertos $U$, $U'$ de $X$, \emph{universalmente esquemáticamente densos con respecto a $S$}, posee también esta propiedad.
Definimos, pues, una relación de equivalencia entre los $S$-morfismos $u:U\to Y$ reemplazando en \sref{IV.20.2.1} «esquemáticamente denso» por «universalmente esquemáticamente denso con respecto a $S$».
Una clase de equivalencia para esta relación se llama un \emph{pseudo-morfismo de $X$ en $Y$ relativamente a $S$}, y el conjunto de estas clases se denota $\operatorname{Ps.hom}_{X/S}(X,Y)$.
\end{env}

\begin{env}[20.5.2]
\label{IV.20.5.2}
Como todo abierto de $X$ universalmente esquemáticamente denso con respecto a $S$ es en particular esquemáticamente denso, los elementos de un pseudo-morfismo
\oldpage[IV-4]{250}
de $X$ en $Y$ relativamente a $S$ son equivalentes en el sentido de \sref{IV.20.2.1}, de donde una aplicación canónica
\[
\label{IV.20.5.2.1}
  \operatorname{Ps.hom}_{X/S}(X,Y)\to\operatorname{Ps.hom}_S(X,Y).
  \tag{20.5.2.1}
\]
\end{env}

\begin{proposition}[20.5.3]
\label{IV.20.5.3}
Supongamos que $Y$ es separado sobre $S$.
Entonces:
\begin{enumerate}[label=\emph{(\roman*)}]
  \item La aplicación \sref{IV.20.5.2.1} es inyectiva e identifica $\operatorname{Ps.hom}_{X/S}(X,Y)$ con el subconjunto de $\operatorname{Ps.hom}_S(X,Y)$ formado por los pseudo-$S$-morfismos $\omega$ tales que $\operatorname{dom}_S(\omega)$ sea universalmente esquemáticamente denso relativamente a $S$.
  \item El prehaz $U\mapsto\operatorname{Ps.hom}_{U/S}(U,Y)$ sobre $X$ es un haz.
\end{enumerate}
\end{proposition}

\begin{proof}
La afirmación~(i) es inmediata, pues decir que dos $S$-morfismos $u:U\to Y$, $u':U'\to Y$ son equivalentes en el sentido de \sref{IV.20.2.1} significa que son las restricciones de un mismo morfismo $\operatorname{dom}_S(\omega)\to Y$ \sref{IV.20.2.4}, y si $U$ y $U'$ son universalmente esquemáticamente densos con respecto a $S$, lo mismo ocurre \emph{a fortiori} con $\operatorname{dom}_S(\omega)$.
Para demostrar~(ii), observemos que $U\mapsto\operatorname{Ps.hom}_S(U,Y)$ es entonces un haz \sref{IV.20.2.6}; por otra parte, dado un recubrimiento abierto $(U_\alpha)$ de $U$ y un pseudo-$S$-morfismo $\omega$ de $U$ en $Y$, se deduce inmediatamente de las definiciones (cf.\ \sref{IV.20.2.1}) que $\operatorname{dom}_S(\omega)$ es universalmente esquemáticamente denso en $U$ relativamente a $S$ si y solo si cada uno de los conjuntos $\operatorname{dom}_S(\omega)\cap U_\alpha=\operatorname{dom}_S(\omega|U_\alpha)$ es universalmente esquemáticamente denso en $U_\alpha$ relativamente a $S$; de donde~(ii).
\end{proof}

Cuando $Y$ es separado, denotaremos por $\mathscr{P}\!s.\!\operatorname{hom}_{X/S}(X,Y)$ el subhaz
\[
  U\mapsto\operatorname{Ps.hom}_{U/S}(U,Y)
\]
de $\mathscr{P}\!s.\!\operatorname{hom}_S(X,Y)$.

Cuando $X$ es \emph{plano y localmente de presentación finita} sobre $S$ y $Y$ separado sobre $S$, los pseudo-morfismos de $X$ en $Y$ relativamente a $S$ se identifican todavía con los pseudo-$S$-morfismos $\omega$ tales que, para toda fibra $X_s$ del morfismo $X\to S$, $\operatorname{dom}_S(\omega)\cap X_s$ sea \emph{esquemáticamente denso} en $X_s$ \sref{IV.11.10.10}.

\begin{env}[20.5.4]
\label{IV.20.5.4}
Un caso particularmente importante en que $Y$ es separado sobre $S$ es el caso en que $Y=S[T]=S\otimes_{\bb{Z}}\Spec(\bb{Z}[T])$ ($T$ indeterminada).
Hay entonces una correspondencia biyectiva entre los pseudo-$S$-morfismos de $X$ en $Y$ y las \emph{pseudo-funciones} sobre $X$ \sref{IV.20.2.8}, en virtud de la definición de un producto de preesquemas sobre $\bb{Z}$.
Los pseudo-morfismos de $X$ en $Y$ relativamente a $S$ se identifican entonces, en virtud de \sref{IV.20.5.3}, con las \emph{pseudo-funciones $\varphi$ sobre $X$ tales que $\operatorname{dom}(\varphi)$ sea universalmente esquemáticamente denso relativamente a $S$}.
El haz $\mathscr{P}\!s.\!\operatorname{hom}_{X/S}(X,Y)$ es un \emph{subhaz de anillos} de $\sh{M}'_X$, que denotamos $\sh{M}'_{X/S}$.

Ponemos entonces $\operatorname{Ps.hom}_{X/S}(X,Y)=M'(X/S)$ y decimos que sus elementos son las \emph{pseudo-funciones sobre $X$ relativas a $S$}.
\end{env}

\begin{env}[20.5.5]
\label{IV.20.5.5}
Sean $X$, $Y$, $Z$ tres preesquemas sobre $S$, $f:Y\to Z$ un $S$-morfismo.
En el razonamiento de \sref{IV.20.3.1}, podemos reemplazar en todas partes «esquemáticamente denso» por «universalmente esquemáticamente denso relativamente a $S$»; para todo pseudo-morfismo $\omega\in\operatorname{Ps.hom}_{X/S}(X,Y)$, los morfismos $f\circ u$, donde $u$ recorre $\omega$, pertenecen, pues, a una misma clase de equivalencia (para la relación definida en \sref{IV.20.5.1}), que llamamos el pseudo-morfismo de $X$ en $Z$, relativo a $S$, \emph{compuesto} de $f$ y de $\omega$, y denotamos $f\circ\omega$.
\oldpage[IV-4]{251}
Si $g:Z\to T$ es un morfismo, es inmediato que $g\circ(f\circ\omega)=(g\circ f)\circ\omega$.
\end{env}

\begin{env}[20.5.6]
\label{IV.20.5.6}
Supongamos que $Y$ sea \emph{separado} sobre $S$, y sea $\omega$ un pseudo-morfismo de $X$ en $Y$ relativamente a $S$.
Si $f:X'\to X$ es un $S$-morfismo tal que $f^{-1}(\operatorname{dom}_S(\omega))$ sea \emph{universalmente esquemáticamente denso en $X'$ relativamente a $S$}, se deduce de \sref{IV.20.3.2} que el pseudo-$S$-morfismo $\omega\circ f$ tiene un dominio universalmente esquemáticamente denso relativamente a $S$, por lo que, en virtud de \sref{IV.20.5.3}\footnote{La transcripción francesa omite el sujeto $\omega\circ f$ y hace gramaticalmente de esta referencia el sujeto de «puede considerarse»; el testigo inglés mantenido restituye la construcción correcta.}, puede considerarse como un pseudo-morfismo \emph{relativo a $S$}.
Cuando $X'$ es plano y localmente de presentación finita sobre $S$, la condición de que $f^{-1}(\operatorname{dom}_S(\omega))$ sea universalmente esquemáticamente denso relativamente a $S$ equivale también a decir que para todo $s\in S$, $(f^{-1}(\operatorname{dom}_S(\omega)))_s$ (notación de \sref{IV.11.10.10}) sea esquemáticamente denso en $X'_s$, o, denotando $f_s:X'_s\to X_s$ el morfismo deducido de $f$ por cambio de base, que la imagen recíproca por $f_s$ de $(\operatorname{dom}_S(\omega))_s$ sea esquemáticamente densa en $X'_s$.
Esta última condición se cumplirá en particular si, para todo $s\in S$, $X_s$, $X'_s$ y $f_s$ verifican una de las tres condiciones~(i),~(ii),~(iii) de \sref{IV.20.3.5}.
\end{env}

\begin{env}[20.5.7]
\label{IV.20.5.7}
Supongamos ahora que $X$ y $X'$ sean ambos preesquemas sobre $S$ \emph{planos y localmente de presentación finita sobre $S$}, y que $f:X'\to X$ sea un $S$-morfismo \emph{plano} (o, lo que equivale a lo mismo \sref{IV.11.3.10}, que para todo $s\in S$, $f_s:X'_s\to X_s$ sea un morfismo plano).
Entonces, para todo abierto $U$ de $X$ y todo abierto $V\subset U$ universalmente esquemáticamente denso en $U$ relativamente a $S$, se deduce de \sref{IV.11.10.5} y \sref{IV.11.10.9} que $f^{-1}(V)$ es universalmente esquemáticamente denso en $f^{-1}(U)$ relativamente a $S$.
Para todo pseudo-morfismo $\omega$ de $X$ en $Y$ relativamente a $S$, se deduce de \sref{IV.20.3.4} que el pseudo-$S$-morfismo $\omega\circ f$ está definido y es un pseudo-morfismo de $X'$ en $Y$ relativamente a $S$, \emph{incluso cuando no se supone que $Y$ sea separado sobre $S$}.
Se deduce de esto que en este caso, para todo $S$-morfismo $g:Y\to Z$, $(g\circ\omega)\circ f$ sigue estando definido y es igual a $g\circ(\omega\circ f)$ (con las definiciones de \sref{IV.20.5.1}), y es por tanto un pseudo-morfismo \emph{relativo} a $S$.
\end{env}

\begin{env}[20.5.8]
\label{IV.20.5.8}
Sea $X$ un preesquema sobre $S$, $S'\to S$ un morfismo cualquiera, $X'=X_{(S')}$, $g:X'\to X$ la proyección canónica.
Entonces, para todo abierto $U$ de $X$ y todo abierto $V\subset U$ universalmente esquemáticamente denso en $U$ relativamente a $S$, $V'=g^{-1}(V)$ es universalmente esquemáticamente denso en $U'=g^{-1}(U)$ relativamente a $S'$ por definición \sref{IV.11.10.8}.
Sea entonces $\omega$ un pseudo-morfismo de $X$ en un preesquema $Y$ sobre $S$ relativamente a $S$; si $u_1$, $u_2$ son $S$-morfismos de la clase $\omega$, definidos respectivamente en abiertos $U_1$, $U_2$ de $X$, universalmente esquemáticamente densos en $X$ relativamente a $S$, se deduce de lo anterior que los morfismos $u'_1=u_1\circ(g|g^{-1}(U_1))$ y $u'_2=u_2\circ(g|g^{-1}(U_2))$ coinciden en un abierto $U'_3$ universalmente esquemáticamente denso relativamente a $S'$.
Ahora bien, si $Y'=Y_{(S')}$ y si $h:Y'\to Y$ es la proyección canónica, $u'_1$ y $u'_2$ se factorizan canónicamente como $u'_1=h\circ v_1$, $u'_2=h\circ v_2$, y $v_1$ y $v_2$ son dos $S'$-morfismos en $Y'$ que coinciden en $U'_3$.
Se ve, pues, que cuando $u_1$ recorre la clase $\omega$, los $S'$-morfismos $v_1$ correspondientes pertenecen a un mismo pseudo-morfismo relativamente a $S'$, llamado \emph{imagen recíproca} de $\omega$ mediante el morfismo cambio de base $S'\to S$ y denotado $\omega_{(S')}$.
Es claro que si $S''\to S'$ es un segundo morfismo, se tiene $(\omega_{(S')})_{(S'')}=\omega_{(S'')}$ (para el morfismo cambio de base compuesto $S''\to S'\to S$).
\end{env}

\oldpage[IV-4]{252}

\subsection{Funciones meromorfas relativas}
\label{IV.20.6}

\begin{env}[20.6.1]
\label{IV.20.6.1}
Sean $S$ un preesquema, $X$ un $S$-preesquema que es \emph{plano y localmente de presentación finita} sobre $S$; para todo $s\in S$, denotaremos por $X_s$ la fibra en el punto $s$ del morfismo estructural $X\to S$.
En general, si $\varphi$ es una función meromorfa sobre $X$, no es posible asociarle de manera «natural», para cada $s\in S$, una función meromorfa «inducida» sobre $X_s$ \sref{IV.20.1.11}.
Para todo abierto $U$ de $X$, denotemos por $T_{X/S}(U)$ el conjunto de las secciones $t\in\Gamma(U,\sh{O}_X)$ tales que, para todo $s\in S$, la imagen $t_s$ de $t$ por el homomorfismo canónico $\Gamma(U,\sh{O}_X)\to\Gamma(U\cap X_s,\sh{O}_{X_s})$ sea una sección \emph{regular}; esto implica además, por la equivalencia de a) y b) en \sref{IV.11.3.7}, que $t$ es ella misma una sección \emph{regular}.
Está claro que $U\mapsto T_{X/S}(U)$ es un \emph{subhaz} del haz de conjuntos $\sh{S}=\sh{S}(\sh{O}_X)$ (notación de \sref{IV.20.1.3}), que se denota por $\sh{T}_{X/S}$; se pone
\[
\label{IV.20.6.1.1}
  \sh{M}_{X/S}=\sh{O}_X[\sh{T}_{X/S}^{-1}]
  \tag{20.6.1.1}
\]
(notación de \sref{IV.20.1.1}) y se dice que este haz de anillos es el \emph{haz de gérmenes de funciones meromorfas sobre $X$ relativas a $S$}; sus secciones sobre $X$ se llaman \emph{funciones meromorfas sobre $X$ relativas a $S$} y su conjunto se denota por $M(X/S)$.
Está claro que $\sh{M}_{X/S}$ es un \emph{subhaz} de $\sh{M}_X=\sh{O}_X[\sh{S}^{-1}]$; para toda función meromorfa $\varphi\in M(X/S)$ y todo $s\in S$, queda entonces definida la imagen inversa de $\varphi$ por la inyección canónica $j_s:X_s\to X$ \sref{IV.20.1.11}, y se denota por $\varphi_s$.
\end{env}

\begin{env}[20.6.2]
\label{IV.20.6.2}
Sea ahora $\sh{F}$ un $\sh{O}_X$-módulo cuasicoherente; se pone
\[
\label{IV.20.6.2.1}
  \sh{M}_{X/S}(\sh{F})=\sh{F}[\sh{T}_{X/S}^{-1}]=\sh{F}\otimes_{\sh{O}_X}\sh{M}_{X/S};
  \tag{20.6.2.1}
\]
las secciones de $\sh{M}_{X/S}(\sh{F})$ se llaman \emph{secciones meromorfas de $\sh{F}$ sobre $X$, relativas a $S$} y su conjunto se denota por $M(X/S,F)$.
El homomorfismo canónico $\sh{F}\to\sh{M}_{X/S}(\sh{F})$ no es necesariamente inyectivo; cuando lo es, se dice que $\sh{F}$ es \emph{sin torsión relativamente a $S$}: esto significa que para todo abierto $U$ de $X$ y toda sección $t\in T_{X/S}(U)$, $t$ es $(\sh{F}|U)$-regular; esta condición se cumple \emph{a fortiori} cuando $\sh{F}$ es \emph{estrictamente sin torsión} \sref{IV.20.1.5}.
En este último caso, de las definiciones \sref{IV.20.1.2} se deduce inmediatamente que el homomorfismo canónico de $\sh{O}_X$-módulos
\[
\label{IV.20.6.2.2}
  \sh{M}_{X/S}(\sh{F})\to\sh{M}_X(\sh{F})
  \tag{20.6.2.2}
\]
es \emph{inyectivo}, de modo que las secciones meromorfas de $\sh{F}$ relativas a $S$ son secciones meromorfas de $\sh{F}$ en el sentido de \sref{IV.20.1.3}.
\end{env}

\begin{proposition}[20.6.3]
\label{IV.20.6.3}
La imagen por el homomorfismo inyectivo \sref{IV.20.2.10.1} del subhaz $\sh{M}_{X/S}$ de $\sh{M}_X$ es el subhaz $\sh{M}'_{X/S}$ de las pseudo-funciones sobre $X$ relativas a $S$ (es decir, de las pseudo-funciones cuyo dominio de definición es universalmente esquemáticamente denso relativamente a $S$ \sref{IV.20.5.4}).
\end{proposition}

\begin{proof}
Evidentemente, podemos limitarnos a probar que la imagen de $M(X/S)$ por el homomorfismo canónico $M(X)\to M'(X)$ es igual a $M'(X/S)$; la proposición es entonces consecuencia de la proposición más general siguiente:
\end{proof}

\oldpage[IV-4]{253}
\begin{proposition}[20.6.4]
\label{IV.20.6.4}
Sea $\sh{F}$ un $\sh{O}_X$-módulo cuasicoherente de presentación finita y estrictamente sin torsión.
Entonces, para que una sección meromorfa $\varphi$ de $\sh{F}$ sobre $X$ sea una sección meromorfa relativa a $S$, es necesario y suficiente que $\operatorname{dom}(\varphi)$ sea universalmente esquemáticamente denso relativamente a $S$.
\end{proposition}

\begin{proof}
La necesidad de la condición se deduce de \sref{IV.20.2.15} aplicado a cada inyección canónica $X_s\to X$ ($s\in S$), teniendo en cuenta \sref{IV.11.10.9}.
Para ver que la condición es suficiente, hay que probar que para todo $x\in X$ existe un entorno abierto $U$ de $x$ en $X$ y una sección de $\sh{M}_{X/S}(\sh{F})$ sobre $U$ cuya restricción a $U\cap\operatorname{dom}(\varphi)$ coincida con $\varphi$ en un abierto esquemáticamente denso de $U\cap\operatorname{dom}(\varphi)$.
Consideremos el ideal de los denominadores $\sh{J}$ de $\varphi$ \sref{IV.20.2.14}, que es cuasicoherente y define un subpreesquema cerrado de $X$ cuyo espacio subyacente es $X-\operatorname{dom}(\varphi)$.
Por hipótesis, si $s$ es la imagen de $x$ en $S$, $\operatorname{dom}(\varphi)\cap X_s$ es esquemáticamente denso en el preesquema localmente noetheriano $X_s$, y por tanto, en virtud de \sref{IV.20.2.13}[(iv)], contiene a $\operatorname{Ass}(\sh{O}_{X_s})$; esto implica que la imagen del ideal $\sh{J}_x$ de $\sh{O}_{X,x}$ en $\sh{O}_{X_s,x}=\sh{O}_{X,x}\otimes_{\sh{O}_{S,s}}k(s)$ no está contenida en ninguno de los ideales primos $\mathfrak{p}_i\in\operatorname{Ass}(\sh{O}_{X_s,x})$ (que son finitos); por tanto (Bourbaki, \emph{Alg.~comm.}, cap.~II, \textsection1, n\textsuperscript{o}~1, prop.~2), existe un elemento $t_x\in\sh{J}_x$ cuya imagen en $\sh{O}_{X_s,x}$ no pertenece a ninguno de los $\mathfrak{p}_i$ y es, por consiguiente, regular en este anillo noetheriano.
Sea $t$ una sección de $\sh{J}$ sobre un entorno abierto afín $U$ de $x$, cuyo germen en el punto $x$ sea $t_x$; puesto que $X$ es plano y localmente de presentación finita sobre $S$, podemos suponer \sref{IV.11.3.8} que $t$ es una sección \emph{regular} de $\sh{O}_X$ sobre $U$ y que, para todo $s'\in S$, la imagen de $t$ en $\Gamma(U\cap X_{s'},\sh{O}_{X_{s'}})$ es también \emph{regular}; en otros términos, se tiene $t\in T_{X/S}(U)$.
Pero entonces, por definición de $\sh{J}$, como $\sh{F}$ es estrictamente sin torsión, $t(\varphi|(U\cap\operatorname{dom}(\varphi)))$ es una sección $u$ de $\sh{F}$ sobre $U\cap\operatorname{dom}(\varphi)$; por otra parte, $U\cap\operatorname{dom}(\varphi)$ contiene el conjunto abierto $U_t$ de los puntos $x'\in U$ en los que $t(x')\neq0$, y este último contiene a $x$ y es esquemáticamente denso en $U$ \sref{IV.20.2.9}.
Se ve, pues, que en $U_t$, $\varphi$ coincide con la restricción a $U_t$ de la sección $u/t$ de $\sh{M}_{X/S}(\sh{F})$ sobre $U$.
C.Q.F.D.
\end{proof}

\begin{remarks}[20.6.5]
\label{IV.20.6.5}
\begin{enumerate}[label=(\roman*)]
  \item Sea $\varphi$ una función meromorfa sobre $X$ relativa a $S$, de modo que para todo $s\in S$, $\varphi_s$ es una función meromorfa sobre $X_s$ \sref{IV.20.6.1}; en virtud de \sref{IV.20.1.11.1}, se tiene
  \[
  \label{IV.20.6.5.1}
    \operatorname{dom}(\varphi)\cap X_s\subset\operatorname{dom}(\varphi_s).
    \tag{20.6.5.1}
  \]
  Pero conviene observar que, aun cuando $S$ sea el espectro de un anillo de valoración discreta $A$ y $X=S[T]$ ($T$ indeterminada), los dos miembros de \sref{IV.20.6.5.1} no son necesariamente iguales: por ejemplo, si $\pi$ es un uniformizante de $A$, es inmediato que $\varphi=\pi/T$ es una función meromorfa sobre $X$ relativa a $S$\footnote{La transcripción francesa dice aquí «sobre $S$», aunque $T$ es la indeterminada de $X=S[T]$; el testigo inglés mantenido restituye «sobre $X$».}, pues si $a$ y $b$ son el punto cerrado y el punto genérico de $S$, $T$ es regular en $\Gamma(X_a,\sh{O}_{X_a})=k[T]$ y en $\Gamma(X_b,\sh{O}_{X_b})=K[T]$, siendo $k$ y $K$ el cuerpo residual y el cuerpo de fracciones de $A$.
  Se tiene $\operatorname{dom}(\varphi)=\mathrm{D}(T)$ en $X$, pero $\operatorname{dom}(\varphi_a)=X_a$ puesto que $\varphi_a=0$.

  \item Para que una función meromorfa $\varphi$ relativa a $S$ sea \emph{invertible en el anillo} $M(X/S)$, es necesario y suficiente que, para todo $s\in S$, $\varphi_s$ sea invertible en $M(X_s)$ (en otros términos, que $\varphi_s$ sea una función meromorfa regular sobre $X_s$ \sref{IV.20.1.8}).
  La condición es
  \oldpage[IV-4]{254}
  en efecto trivialmente necesaria.
  Recíprocamente, si se cumple, y si $x$ es un punto cualquiera de $X$, $s$ su imagen en $S$, existen por hipótesis un entorno abierto $U$ de $x$ en $X$ y dos secciones $u$, $t$ de $\sh{O}_X$ sobre $U$ tales que $t\in T_{X/S}(U)$ y $\varphi|U=u/t$; la hipótesis implica que, si $u_s$ es la imagen de $u$ en $\Gamma(U\cap X_s,\sh{O}_{X_s})$, $u_s$ es regular en el punto $x$.
  Restringiendo $U$, podemos suponer, por tanto, en virtud de \sref{IV.11.3.8}, que $u\in T_{X/S}(U)$, de donde la conclusión.

  Cuando $\varphi$ es invertible en $M(X/S)$, se dice también que $\varphi$ es una \emph{función meromorfa regular relativa a $S$}.
  Obsérvese que $\varphi\in M(X/S)$ puede ser invertible en $M(X)$ (en otros términos, ser \emph{regular} en el sentido de \sref{IV.20.1.8}) sin serlo en $M(X/S)$, como muestra inmediatamente el ejemplo de~(i).

  \item Sea $\sh{L}$ un $\sh{O}_X$-módulo \emph{invertible}, y sea $\varphi$ una sección meromorfa \emph{regular} de $\sh{L}$ sobre $X$ \sref{IV.20.1.8}; se dice que $\varphi$ es \emph{regular relativamente a $S$} si, para todo abierto $U$ de $X$ tal que $\sh{L}|U$ sea isomorfo a $\sh{O}_U$, $\varphi|U$ corresponde a un elemento de $\Gamma(U,\sh{M}_X)$ que es \emph{regular relativamente a $S$}; en virtud de~(ii), es inmediato que para ello es necesario y suficiente que, para todo $s\in S$, $\varphi_s$ sea una sección meromorfa \emph{regular} \sref{IV.20.1.8} del $\sh{O}_{X_s}$-módulo invertible $\sh{L}_s=\sh{L}\otimes_{\sh{O}_X}k(s)$.
  Si $\varphi'$ es la inversa de $\varphi$ en $\sh{L}^{-1}$ \sref{IV.20.1.10}, $\varphi'$ es también regular relativamente a $S$.
  Si $\sh{L}_1$ es un segundo $\sh{O}_X$-módulo invertible, $\varphi_1$ una sección meromorfa de $\sh{L}_1$ sobre $X$, regular relativamente a $S$, entonces $\varphi\otimes\varphi_1$ es una sección meromorfa de $\sh{L}\otimes\sh{L}_1$ sobre $X$, regular relativamente a $S$.
\end{enumerate}
\end{remarks}

\begin{proposition}[20.6.6]
\label{IV.20.6.6}
Sean $X$ un $S$-preesquema plano y localmente de presentación finita sobre $S$, y $\sh{F}$ un $\sh{O}_X$-módulo localmente libre de tipo finito; para todo $s\in S$, se denota por $X_s$ la fibra en el punto $s$ del morfismo estructural $f:X\to S$.
Sea $\varphi$ una sección meromorfa de $\sh{F}$ sobre $X$, relativa a $S$, y supongamos que $\varphi$ está definida en todos los puntos $x\in X$ tales que $\operatorname{prof}(\sh{O}_{X_{f(x)},x})=1$.
Entonces $\varphi$ está definida en todas partes.
\end{proposition}

\begin{proof}
Por hipótesis, $\operatorname{dom}(\varphi)\cap X_s$ es esquemáticamente denso en $X_s$ para todo $s\in S$, y por tanto contiene los puntos $x$ de $X_s$ tales que $\operatorname{prof}(\sh{O}_{X_s,x})=0$ \sref{IV.5.10.2}; la hipótesis significa, pues, que si se pone $Z=X-\operatorname{dom}(\varphi)$, se tiene $\operatorname{prof}(\sh{O}_{X_{f(x)},x})\geq2$ en todos los puntos de $Z$.
Basta entonces aplicar \sref{IV.19.9.8}.
\end{proof}

\begin{env}[20.6.7]
\label{IV.20.6.7}
Sean $X$, $X'$ dos $S$-preesquemas planos y localmente de presentación finita sobre $S$, y $f=(\psi,\theta):X'\to X$ un $S$-morfismo.
Para todo abierto $U$ de $X$, denotemos por $T_f(U)$ el conjunto de las secciones $t\in T_{X/S}(U)$ cuya imagen en $\Gamma(f^{-1}(U),\sh{O}_{X'})$ pertenezca a $T_{X'/S}(f^{-1}(U))$; es inmediato que $U\mapsto T_f(U)$ es un subhaz del haz de conjuntos $\sh{T}_{X/S}$, que se denota por $\sh{T}_f$.
Se pone $\sh{M}_{X/S,f}=\sh{O}_X[\sh{T}_f^{-1}]$; es un subhaz de anillos de $\sh{M}_{X/S}$, y de $\theta^\#: \psi^*(\sh{O}_X)\to\sh{O}_{X'}$ se deduce canónicamente un homomorfismo de haces de anillos $\theta'^\#: \psi^*(\sh{M}_{X/S,f})\to\sh{M}_{X'/S}$ que prolonga $\theta^\#$.
Si una función meromorfa $\varphi$ sobre $X$, relativa a $S$, es \emph{una sección de $\sh{M}_{X/S,f}$}, $\Gamma(\theta'^\#)(\varphi)$ es una función meromorfa sobre $X'$, llamada \emph{imagen inversa de $\varphi$ por $f$}, y se denota por $\varphi\circ f$ si ello no da lugar a confusión.
Las definiciones de \sref{IV.20.1.11} relativas a los $\sh{O}_X$-módulos se extienden inmediatamente del mismo modo.
\end{env}

\begin{proposition}[20.6.8]
\label{IV.20.6.8}
Con las notaciones de \sref{IV.20.6.7}, si el $S$-morfismo $f:X'\to X$ es \emph{plano}, se tiene $\sh{M}_{X/S,f}=\sh{M}_{X/S}$, y el homomorfismo $\varphi\mapsto\varphi\circ f$ está definido en todo $M(X/S)$.
\end{proposition}

\begin{proof}
\oldpage[IV-4]{255}
En efecto, la hipótesis implica, en virtud de \sref{IV.11.3.10}, que para todo $s\in S$, $f_s:X'_s\to X_s$ es plano; por tanto, para toda sección $t\in T_{X/S}(U)$, si $t'$ es su imagen inversa en $\Gamma(f^{-1}(U),\sh{O}_{X'})$, $t'_{s}$, que es la imagen inversa de $t_s$, es una sección \emph{regular} de $\sh{O}_{X'_s}$ sobre $f^{-1}(U)\cap X'_s$, en virtud de \sref{IV.20.1.12}; se concluye que, por definición, $t'\in T_{X'/S}(f^{-1}(U))$, de donde la proposición.
\end{proof}

De aquí se deduce, como en \sref{IV.20.1.12}, un homomorfismo canónico de $\sh{O}_{X'}$-álgebras
\[
\label{IV.20.6.8.1}
  f^*(\sh{M}_{X/S})\to\sh{M}_{X'/S}.
  \tag{20.6.8.1}
\]

\begin{env}[20.6.9]
\label{IV.20.6.9}
Consideremos por último un morfismo \emph{arbitrario} $S'\to S$, y pongamos $X'=X\times_S S'$, que es plano y localmente de presentación finita sobre $S'$; sea $p:X'\to X$ la proyección canónica.
Sean $U$ un abierto de $X$, $t$ una sección perteneciente a $T_{X/S}(U)$, y $t'$ su imagen en $\Gamma(p^{-1}(U),\sh{O}_{X'})$; para todo $s'\in S'$, si $s\in S$ es la imagen de $s'$, se tiene $X'_{s'}=X_s\otimes_{k(s)}k(s')$, por lo que el morfismo $X'_{s'}\to X_s$ es plano, y, por consiguiente, en virtud de \sref{IV.20.1.12}, la imagen inversa $t'_{s'}$ de $t_s$ en $\Gamma(p^{-1}(U)\cap X'_{s'},\sh{O}_{X'_{s'}})$ es regular; esto prueba que $t'\in T_{X'/S'}(p^{-1}(U))$.
Esto permite definir canónicamente, como en \sref{IV.20.6.8}, un homomorfismo canónico de $\sh{O}_{X'}$-álgebras
\[
\label{IV.20.6.9.1}
  p^*(\sh{M}_{X/S})\to\sh{M}_{X'/S'}.
  \tag{20.6.9.1}
\]

Mediante la identificación permitida por \sref{IV.20.6.3}, esta noción de cambio de base para las funciones meromorfas relativas es un caso particular de la noción análoga para los pseudo-morfismos relativos \sref{IV.20.5.8}.
\end{env}


% Mechanical subsection components; prose is translated in the component files.
\section{Divisores}
\label{section:IV.21}

Sobre el contenido de esta sección, véanse los comentarios de la introducción del \textsection~20.
Para las propiedades \emph{globales} de los divisores, remitimos al lector a la sección que se les dedica en el cap.~V.

\subsection{Divisores sobre un espacio anillado}
\label{subsection:IV.21.1}
\label{IV.21.1}

\begin{env}[21.1.1]
\label{IV.21.1.1}
Sean $(X,\sh{O}_X)$ un espacio anillado, $\sh{M}_X$ el haz de gérmenes de funciones meromorfas sobre $X$ \sref{IV.20.1.3}, y $\sh{M}_X^*$ el haz (de grupos multiplicativos) de gérmenes de funciones meromorfas regulares sobre $X$ \sref{IV.20.1.8}.
Es claro que el haz (de grupos multiplicativos) $\sh{O}_X^*$ de gérmenes de secciones invertibles de $\sh{O}_X$ se identifica canónicamente con un subhaz (de grupos multiplicativos) de $\sh{M}_X^*$.
\end{env}

\begin{definition}[21.1.2]
\label{IV.21.1.2}
Se llama \emph{haz de divisores sobre $X$} y se denota $\sh{D}iv_X$ el haz cociente (de grupos conmutativos) $\sh{M}_X^*/\sh{O}_X^*$;
las secciones de este haz sobre $X$ se llaman los \emph{divisores sobre $X$}; forman un grupo conmutativo denotado $\operatorname{Div}(X)$.
Para toda sección $f$ de $\sh{M}_X^*$ sobre $X$ (es decir, toda función meromorfa regular sobre $X$ \sref{IV.20.1.8}), se llama \emph{divisor de $f$} y se denota $\operatorname{div}(f)$ (o $\operatorname{div}_X(f)$) el divisor sobre $X$ imagen de $f$ mediante el homomorfismo canónico
\[
  \Gamma(X,\sh{M}_X^*)\longrightarrow\Gamma(X,\sh{D}iv_X).
\]
\end{definition}

\oldpage[IV-4]{256}

El \emph{soporte} de un divisor $D$ es el conjunto cerrado de los $x\in X$ tales que $D_x\neq0$.
Se denota $\operatorname{Supp}(D)$.

Para todo abierto $U$ de $X$, se tiene evidentemente $\sh{M}_X^*|U=\sh{M}_U^*$, $\sh{O}_X^*|U=\sh{O}_U^*$, por tanto $\sh{D}iv_X|U=\sh{D}iv_U$, y en consecuencia el haz $\sh{D}iv_X$ es igual al prehaz $U\mapsto\operatorname{Div}(U)$.

Cuando $X=\operatorname{Spec}(A)$ es afín, se escribe $\operatorname{Div}(A)$ en lugar de $\operatorname{Div}(\operatorname{Spec}(A))$.

\begin{env}[21.1.3]
\label{IV.21.1.3}
Denotaremos siempre \emph{aditivamente} el grupo $\operatorname{Div}(X)$ de divisores sobre $X$.
Para dos funciones meromorfas regulares $f$, $g$ sobre $X$, se tiene, pues,
\[
\label{IV.21.1.3.1}
  \operatorname{div}(fg)=\operatorname{div}(f)+\operatorname{div}(g),
  \tag{21.1.3.1}
\]
\[
\label{IV.21.1.3.2}
  \operatorname{div}(f^{-1})=-\operatorname{div}(f).
  \tag{21.1.3.2}
\]
Por definición, para toda función meromorfa regular $f$ sobre $X$, se tiene la equivalencia
\[
\label{IV.21.1.3.3}
  \operatorname{div}(f)=0 \iff f\in\Gamma(X,\sh{O}_X^*),
  \tag{21.1.3.3}
\]
de donde, para dos funciones meromorfas regulares $f$, $g$ sobre $X$,
\[
\label{IV.21.1.3.4}
  \operatorname{div}(f)=\operatorname{div}(g) \iff fg^{-1}\in\Gamma(X,\sh{O}_X^*).
  \tag{21.1.3.4}
\]
\end{env}

\begin{env}[21.1.4]
\label{IV.21.1.4}
Sea ahora $\sh{L}$ un $\sh{O}_X$-Módulo invertible, y sea $s$ una \emph{sección meromorfa regular} \sref{IV.20.1.8} de $\sh{L}$ sobre $X$.
Todo $x\in X$ posee un entorno abierto $U$ tal que $\sh{L}|U$ sea isomorfo a $\sh{O}_U$, y por tanto $\sh{M}_X(\sh{L})|U$ sea isomorfo a $\sh{M}_X|U$;
mediante uno de estos isomorfismos, $s|U$ corresponde a una sección $f\in\Gamma(U,\sh{M}_X^*)$, y como dos isomorfismos de $\sh{L}|U$ sobre $\sh{O}_U$ solo difieren por la multiplicación por un elemento de $\Gamma(U,\sh{O}_X^*)$ (\EGAReadLink{0.5.4.3}{$\mathbf{0}_{\mathrm I}$, 5.4.3}), el elemento $\operatorname{div}_U(f)$ de $\Gamma(U,\sh{D}iv_X)$ es \emph{independiente} del isomorfismo elegido;
es claro que estos elementos (para $U$ variable) son las restricciones de una sección de $\sh{D}iv_X$ sobre $X$, que se llama el \emph{divisor de $s$} y se denota $\operatorname{div}(s)$ (tal divisor no es necesariamente de la forma $\operatorname{div}(g)$ para una función meromorfa regular $g$ sobre $X$; véase \sref{IV.21.2.9}).
Para $\sh{L}=\sh{O}_X$, la definición de $\operatorname{div}(s)$ coincide con la de \sref{IV.21.1.2}.
Si $\sh{L}$, $\sh{L}'$ son dos $\sh{O}_X$-Módulos invertibles, y $s$ (resp. $s'$) es una sección meromorfa regular de $\sh{L}$ (resp. $\sh{L}'$) sobre $X$, se tiene inmediatamente
\[
\label{IV.21.1.4.1}
  \operatorname{div}(s\otimes s')=\operatorname{div}(s)+\operatorname{div}(s')
  \tag{21.1.4.1}
\]
\[
\label{IV.21.1.4.2}
  \operatorname{div}(s^{\otimes n})=n.\operatorname{div}(s) \quad\text{para todo}\quad n\in\mathbf{Z}
  \tag{21.1.4.2}
\]
(siendo $s^{-1}$ la sección meromorfa regular de $\sh{L}^{-1}$ sobre $X$ definida en \sref{IV.20.1.10})
y, para dos secciones meromorfas regulares $s$, $s'$ de $\sh{L}$ sobre $X$, se tiene la relación
\[
\label{IV.21.1.4.3}
  \operatorname{div}(s)=\operatorname{div}(s') \iff s'=ts \quad\text{con}\quad t\in\Gamma(X,\sh{O}_X^*).
  \tag{21.1.4.3}
\]
\end{env}

\begin{env}[21.1.5]
\label{IV.21.1.5}
El haz $\sh{S}(\sh{O}_X)$ \sref{IV.20.1.3}, cuyas secciones sobre un abierto $U$ de $X$ son los elementos regulares de $\Gamma(U,\sh{O}_X)$, es un \emph{subhaz de monoides} de $\sh{M}_X^*$; se puede escribir
\[
\label{IV.21.1.5.1}
  \sh{S}(\sh{O}_X)=\sh{O}_X\cap\sh{M}_X^*.
  \tag{21.1.5.1}
\]
\end{env}

\oldpage[IV-4]{257}

Si se denota por $\sh{S}(\sh{O}_X)^{-1}$ el haz cuyas secciones sobre $U$ son los inversos en $\Gamma(U,\sh{M}_X^*)$ de los elementos de $\Gamma(U,\sh{S}(\sh{O}_X))$, es claro que se tiene $\Gamma(U,\sh{S}(\sh{O}_X))\cap\Gamma(U,\sh{S}(\sh{O}_X)^{-1})=\Gamma(U,\sh{O}_X^*)$, y por tanto
\[
\label{IV.21.1.5.2}
  \sh{S}(\sh{O}_X)\cap\sh{S}(\sh{O}_X)^{-1}=\sh{O}_X^*.
  \tag{21.1.5.2}
\]

\begin{definition}[21.1.6]
\label{IV.21.1.6}
El subhaz de conjuntos de $\sh{D}iv_X$, imagen canónica del subhaz $\sh{S}(\sh{O}_X)$ de $\sh{M}_X^*$, se denota $\sh{D}iv_X^+$; sus secciones sobre $X$ se llaman \emph{divisores positivos sobre $X$}, y su conjunto se denota $\operatorname{Div}^+(X)$.
\end{definition}

Puesto que $\sh{S}(\sh{O}_X)$ es un haz de monoides (multiplicativos), se tiene
\[
\label{IV.21.1.6.1}
  \operatorname{Div}^+(X)+\operatorname{Div}^+(X)\subset\operatorname{Div}^+(X)
  \tag{21.1.6.1}
\]
y, por otra parte, en virtud de \sref{IV.21.1.5.2} y \sref{IV.21.1.3.3},
\[
\label{IV.21.1.6.2}
  \operatorname{Div}^+(X)\cap(-\operatorname{Div}^+(X))=\{0\}.
  \tag{21.1.6.2}
\]
Estas dos relaciones muestran que $\operatorname{Div}^+(X)$ es el conjunto de los elementos positivos para una \emph{estructura de orden} sobre el grupo $\operatorname{Div}(X)$, compatible con esta estructura de grupo; esta relación de orden se denota $D\leq D'$; es decir, se tiene
\[
\label{IV.21.1.6.3}
  D\geq0 \iff D\in\operatorname{Div}^+(X).
  \tag{21.1.6.3}
\]

En lo que sigue supondremos siempre que $\operatorname{Div}(X)$ está provisto de esta estructura de orden;
es claro que $\sh{D}iv_X^+|U=\sh{D}iv_U^+$, por tanto $\Gamma(U,\sh{D}iv_X^+)=\operatorname{Div}^+(U)$, y se puede decir que $\sh{D}iv_X^+$ define sobre $\sh{D}iv_X$ una estructura de haz de grupos ordenados.
La fibra $(\sh{D}iv_X^+)_x$ en un punto $x$ del haz $\sh{D}iv_X^+$ es un submonoide del grupo $(\sh{D}iv_X)_x$, constituido por los elementos $\geq0$ para una estructura de orden compatible con la estructura de grupo;
para un divisor $D$ sobre $X$, decir que $D\geq0$ equivale a decir que $D_x\geq0$ para todo $x\in X$.

Por definición, para toda función meromorfa regular $f$ sobre $X$, se tiene la relación
\[
\label{IV.21.1.6.4}
  \operatorname{div}(f)\geq0 \iff f\in\Gamma(X,\sh{S}(\sh{O}_X))
  \tag{21.1.6.4}
\]
es decir, $\operatorname{div}(f)\geq0$ significa que $f$ es una \emph{sección regular} de $\sh{O}_X$, o también una sección de $\sh{O}_X$ \emph{invertible en} $M(X)$.

Más generalmente, dado un divisor $D$ sobre $X$, la relación $\operatorname{div}(f)\geq D$ equivale a la siguiente: para todo abierto $U$ de $X$ tal que $D|U=\operatorname{div}(g)$, donde $g\in\Gamma(U,\sh{M}_X^*)$, existe un elemento regular $t$ de $\Gamma(U,\sh{O}_X)$ tal que $f|U=tg$.

\begin{env}[21.1.7]
\label{IV.21.1.7}
Sean $\sh{L}$ un $\sh{O}_X$-Módulo invertible y $s$ una sección meromorfa regular de $\sh{L}$ sobre $X$; se tiene la relación
\[
\label{IV.21.1.7.1}
  \operatorname{div}(s)\geq0 \iff s\in\Gamma(X,\sh{L})\cap\Gamma(X,(\sh{M}_X(\sh{L}))^*)
  \tag{21.1.7.1}
\]
como se deduce inmediatamente de las definiciones \sref{IV.21.1.4} y \sref{IV.21.1.6}.
\end{env}

\begin{proposition}[21.1.8]
\label{IV.21.1.8}
Sean $X$ un preesquema localmente noetheriano y $D$ un divisor sobre $X$.
Supongamos que para todo $x\in X$ tal que $\operatorname{prof}(\sh{O}_{X,x})=1$ se tenga $D_x\geq0$ (resp. $D_x=0$).
Entonces se tiene $D\geq0$ (resp. $D=0$).
\end{proposition}

\begin{proof}
Como la cuestión es local sobre $X$, se puede suponer que $D=\operatorname{div}(f)$, siendo $f$ una
\oldpage[IV-4]{258}
función meromorfa regular sobre $X$;
la relación $D_x\geq0$ equivale a $x\in\operatorname{dom}(f)$, por lo que la hipótesis significa que, si $T=X-\operatorname{dom}(f)$, se tiene $\operatorname{prof}(\sh{O}_{X,x})\geq2$ para todo $x\in T$ (pues $\operatorname{dom}(f)$ contiene los puntos maximales de $X$).
Por consiguiente, de \sref{IV.5.10.5} se deduce que el homomorfismo de restricción $\Gamma(X,\sh{O}_X)\to\Gamma(X-T,\sh{O}_X)$ es biyectivo, lo que muestra que existe una sección $s$ de $\sh{O}_X$ sobre $X$ tal que $f=s|(X-T)$.
Pero, por definición de $T$, esto implica $T=\varnothing$, y por tanto $f=s$ y $D\geq0$.
La afirmación relativa a la relación $D_x=0$ se deduce inmediatamente aplicando lo anterior a $-D$, en virtud de \sref{IV.21.1.6.2}.
\end{proof}


\begin{corollary}[21.1.9]
\label{IV.21.1.9}
Sean $X$ un preesquema localmente noetheriano y $D$ un divisor sobre $X$.
Sea $S$ el soporte de $D$.
Entonces, para todo punto maximal $x$ de $S$, se tiene $\operatorname{prof}(\sh{O}_{X,x})=1$.
\end{corollary}

\begin{proof}
Pongamos, en efecto, $X_1=\operatorname{Spec}(\sh{O}_{X,x})$;
teniendo en cuenta \sref{IV.20.2.11} y \sref{IV.20.3.6}, el haz $\sh{M}_{X_1}^*$ es inducido sobre $X_1$ por $\sh{M}_X^*$, por lo que podemos limitarnos al caso en que $X=X_1$; en este caso, como $x\in S$ y $x$ es un punto maximal de $S$, se tiene necesariamente $S=\{x\}$.
Si se tuviera entonces $\operatorname{prof}(\sh{O}_{X,x})\neq1$, se concluiría, en virtud de \sref{IV.21.1.8}, que $D=0$, lo que contradice la definición de $S$.
\end{proof}

\begin{proposition}[21.1.10]
\label{IV.21.1.10}
Sea $A$ un anillo local noetheriano;
para que $\operatorname{Div}(A)=0$, es necesario y suficiente que $\operatorname{prof}(A)=0$ (es decir, que el ideal maximal $\mathfrak{m}$ de $A$ esté asociado a $A$ (\EGAReadLink{0.16.4.6}{$\mathbf{0}$, 16.4.6})).
\end{proposition}

\begin{proof}
En efecto, decir que $\operatorname{Div}(A)=0$ significa que en $A$ todo elemento regular es invertible, o también que todos los elementos de $\mathfrak{m}$ son divisores de cero, lo que significa que $\mathfrak{m}\in\operatorname{Ass}(A)$ (Bourbaki, \emph{Alg. comm.}, cap.~IV, \textsection~1, n\textsuperscript{o}~1, cor.~3 de la prop.~2).
\end{proof}

\subsection{Divisores e ideales fraccionarios invertibles}
\label{subsection:IV.21.2}

\begin{env}[21.2.1]
\label{IV.21.2.1}
Sea $(X,\sh{O}_X)$ un espacio anillado.
Se llama \emph{ideal fraccionario sobre $X$} a un sub-$\sh{O}_X$-módulo del $\sh{O}_X$-módulo $\sh{M}_X$ de los gérmenes de funciones meromorfas sobre $X$.
Un ideal fraccionario $\sh{J}$ sobre $X$ que sea un $\sh{O}_X$-módulo invertible se llama \emph{ideal fraccionario invertible}.
\end{env}

\begin{proposition}[21.2.2]
\label{IV.21.2.2}
Para que un ideal fraccionario $\sh{J}$ sobre $X$ sea invertible, es necesario y suficiente que para todo $x\in X$ exista un entorno abierto $U$ de $x$ y una sección $f\in\Gamma(U,\sh{M}_X^*)$ tales que $\sh{J}|U=\sh{O}_U.f$.
\end{proposition}

\begin{proof}
La condición es evidentemente suficiente, pues la aplicación $s\mapsto s(f|V)$ de $\Gamma(V,\sh{O}_X)$ en $\Gamma(V,\sh{J})$ es evidentemente biyectiva para todo abierto $V\subset U$.
Para ver que es necesaria, observemos que por hipótesis existe un entorno abierto $U$ de $x$ y un isomorfismo de $\sh{O}_X$-módulos $\sh{O}_U\xrightarrow{\sim}\sh{J}|U$.
Si $f$ es la imagen de la sección $1\in\Gamma(U,\sh{O}_X)$ por este isomorfismo, se puede suponer, restringiendo $U$, que $f=u/s$, donde $u\in\Gamma(U,\sh{O}_X)$ y $s\in\Gamma(U,\sh{S}(\sh{O}_X))$, y el isomorfismo considerado hace corresponder entonces, a toda sección $v\in\Gamma(V,\sh{O}_X)$ (donde $V$ es un abierto contenido en $U$), la sección $v(u|V)/(s|V)$;
decir que la aplicación así definida es biyectiva significa que $u|V$ es un elemento regular de $\Gamma(V,\sh{O}_X)$, y por tanto $f\in\Gamma(U,\sh{M}_X^*)$.
\end{proof}

Obsérvese que para todo abierto $U$ de $X$ tal que $\sh{J}|U=\sh{O}_U.f$, con $f\in\Gamma(U,\sh{M}_X^*)$, la sección $f$ está determinada \emph{salvo multiplicación} por un elemento de $\Gamma(U,\sh{O}_X^*)$, puesto que la
\oldpage[IV-4]{259}
multiplicación por estos elementos proporciona todos los automorfismos del $\sh{O}_U$-módulo $\sh{O}_U$ (\EGAReadLink{0.5.4.3}{$\mathbf{0}_{\mathrm I}$, 5.4.3}).

\begin{corollary}[21.2.3]
\label{IV.21.2.3}
\begin{enumerate}[label=(\roman*)]
  \item Sea $\sh{J}$ un ideal fraccionario invertible;
  entonces el $\sh{O}_X$-módulo invertible $\sh{J}^{-1}$ se identifica canónicamente con el ideal fraccionario $\sh{J}'$ (\emph{transportador} de $\sh{J}$ en $\sh{O}_X$) definido del modo siguiente: para todo abierto $U$ de $X$ tal que $\sh{J}|U=\sh{O}_U.f$, donde $f\in\Gamma(U,\sh{M}_X^*)$, se tiene $\sh{J}'|U=\sh{O}_U.f^{-1}$.
  \item Si $\sh{J}_1$ y $\sh{J}_2$ son dos ideales fraccionarios invertibles, la aplicación canónica $\sh{J}_1\otimes\sh{J}_2\to\sh{J}_1\sh{J}_2$ es un isomorfismo de $\sh{O}_X$-módulos.
\end{enumerate}
\end{corollary}

\begin{proof}
La afirmación (ii) resulta inmediatamente de \sref{IV.21.2.2}.
Por otra parte, la observación hecha al final de \sref{IV.21.2.2} prueba que existe efectivamente un único ideal fraccionario $\sh{J}'$ definido por la condición del enunciado;
el isomorfismo canónico de $\sh{J}'$ sobre $\sh{J}^{-1}=\shHom_{\sh{O}_X}(\sh{J},\sh{O}_X)$ se obtiene haciendo corresponder a toda sección $s(f^{-1}|V)$ de $\Gamma(V,\sh{J}')$ (donde $V$ es un abierto contenido en $U$ y $s\in\Gamma(V,\sh{O}_X)$) el homomorfismo $t(f|V)\mapsto st$ de $\Gamma(V,\sh{J})$ en $\Gamma(V,\sh{O}_X)$.
\end{proof}

En virtud de \sref{IV.21.2.3},~(i), se identificarán generalmente los $\sh{O}_X$-módulos invertibles $\sh{J}'$ y $\sh{J}^{-1}$, considerando por tanto $\sh{J}^{-1}$ como un sub-$\sh{O}_X$-módulo de $\sh{M}_X$.

\begin{env}[21.2.4]
\label{IV.21.2.4}
De \sref{IV.21.2.3} resulta que el conjunto $\operatorname{Id.inv}(X)$ de los ideales fraccionarios invertibles sobre $X$ está dotado de una estructura de \emph{grupo conmutativo} para la ley de composición $(\sh{J}_1,\sh{J}_2)\mapsto\sh{J}_1\sh{J}_2$, siendo $\sh{O}_X$ el elemento neutro de este grupo.
Es claro que para todo abierto $U$ de $X$ se tiene $(\sh{J}_1\sh{J}_2)|U=(\sh{J}_1|U)(\sh{J}_2|U)$, por lo que la aplicación de restricción $\sh{J}\mapsto\sh{J}|U$ es un homomorfismo de grupos $\operatorname{Id.inv}(X)\to\operatorname{Id.inv}(U)$;
se define así un \emph{prehaz de grupos conmutativos} $U\mapsto\operatorname{Id.inv}(U)$;
es inmediato que, de hecho, este prehaz es un \emph{haz de grupos conmutativos}, que se denota $\sh{J}d.inv_X$.
\end{env}

\begin{env}[21.2.5]
\label{IV.21.2.5}
Para toda función meromorfa regular $f\in\Gamma(X,\sh{M}_X^*)$, de \sref{IV.21.2.2} resulta que $\sh{J}(f)=\sh{O}_X.f$ es un ideal fraccionario invertible, y se tiene evidentemente $\sh{J}(f_1f_2)=\sh{J}(f_1)\sh{J}(f_2)$; en otras palabras, la aplicación $f\mapsto\sh{J}(f)$ es un homomorfismo del grupo conmutativo $\Gamma(X,\sh{M}_X^*)$ en el grupo conmutativo $\operatorname{Id.inv}(X)$.
Sustituyendo $X$ por un abierto cualquiera $U$ de $X$ y observando que los homomorfismos obtenidos son compatibles con las operaciones de restricción, se obtiene un homomorfismo canónico de haces de grupos conmutativos:
\[
\label{IV.21.2.5.1}
  I_0:\sh{M}_X^*\to\sh{J}d.inv_X.
  \tag{21.2.5.1}
\]
Obsérvese que si $f\in\Gamma(X,\sh{O}_X^*)$, se tiene $\sh{J}(f)=\sh{O}_X$;
de ello se deduce inmediatamente que el homomorfismo $I_0$ se factoriza como
\[
\label{IV.21.2.5.2}
  \sh{M}_X^*\to\sh{M}_X^*/\sh{O}_X^*=\sh{D}iv_X\xrightarrow{I}\sh{J}d.inv_X
  \tag{21.2.5.2}
\]
donde $I$ es un homomorfismo del haz de grupos aditivos $\sh{D}iv_X$ en el haz de grupos multiplicativos $\sh{J}d.inv_X$;
se tiene, por tanto, para todo abierto $U$ de $X$ un homomorfismo $\sh{J}_U:\operatorname{Div}(U)\to\operatorname{Id.inv}(U)$ de grupos conmutativos tal que, para toda sección $f\in\Gamma(U,\sh{M}_X^*)$, se tiene
\[
\label{IV.21.2.5.3}
  \sh{J}_U(\operatorname{div}_U(f))=\sh{O}_U.f.
  \tag{21.2.5.3}
\]

\oldpage[IV-4]{260}

Se concluye que $\sh{J}_X(D)$, para todo divisor $D\in\operatorname{Div}(X)$, es el ideal fraccionario invertible definido del modo siguiente: para todo abierto $U$ de $X$ tal que $D|U=\operatorname{div}_U(f)$, donde $f\in\Gamma(U,\sh{M}_X^*)$, $\sh{J}_X(D)|U$ es el ideal fraccionario invertible $\sh{O}_U.f$.
Se tiene, por tanto, en virtud de \sref{IV.21.1.6}, para toda función meromorfa regular $f\in\Gamma(X,\sh{M}_X^*)$, la relación
\[
\label{IV.21.2.5.4}
  f\in\Gamma(X,\sh{J}_X(D)) \iff \operatorname{div}(f)\geq D.
  \tag{21.2.5.4}
\]
\end{env}

\begin{proposition}[21.2.6]
\label{IV.21.2.6}
El homomorfismo $I:\sh{D}iv_X\to\sh{J}d.inv_X$ es biyectivo.
\end{proposition}

\begin{proof}
Se define, en efecto, un homomorfismo $I_X'$ de $\operatorname{Id.inv}(X)$ en $\operatorname{Div}(X)$ haciendo corresponder a todo ideal fraccionario invertible $\sh{J}$ sobre $X$ el divisor $I_X'(\sh{J})$ definido del modo siguiente: para todo abierto $U$ de $X$ tal que $\sh{J}|U=\sh{O}_U.f$, donde $f\in\Gamma(U,\sh{M}_X^*)$ \sref{IV.21.2.2}, se toma $I_X'(\sh{J})|U=\operatorname{div}_U(f)$;
en virtud de la observación que sigue a \sref{IV.21.2.2}, esta definición es independiente del generador $f$ elegido en $\sh{J}|U$ y determina efectivamente un divisor sobre $X$.
Además, esta definición muestra inmediatamente que los homomorfismos $\sh{J}_X$ e $I_X'$ son inversos entre sí.
Sustituyendo $X$ por un abierto cualquiera $U$, se deduce la definición del isomorfismo $I':\sh{J}d.inv_X\to\sh{D}iv_X$ inverso de $I$, de donde resulta la proposición.
Se pondrá $I_X'(\sh{J})=\operatorname{div}(\sh{J})$, de modo que, para toda función meromorfa regular $f$ sobre $X$, se tiene
\[
\label{IV.21.2.6.1}
  \operatorname{div}(\sh{O}_X.f)=\operatorname{div}(f).
  \tag{21.2.6.1}
\]
\end{proof}

\begin{env}[21.2.7]
\label{IV.21.2.7}
Se identificarán a menudo los haces $\sh{D}iv_X$ y $\sh{J}d.inv_X$ (respectivamente, los grupos $\operatorname{Div}(X)$ y $\operatorname{Id.inv}(X)$) mediante los isomorfismos anteriores $I$ e $I'$ (respectivamente, $\sh{J}_X$ e $I_X'$).
Obsérvese que se tiene la relación
\[
\label{IV.21.2.7.1}
  D\geq0 \iff \sh{J}_X(D)\subset\sh{O}_X \quad\text{para}\quad D\in\operatorname{Div}(X)
  \tag{21.2.7.1}
\]
como resulta inmediatamente de las definiciones \sref{IV.21.1.6} y \sref{IV.21.1.5.1};
en otras palabras, la imagen $\sh{J}_X(\operatorname{Div}^+(X))$ es el conjunto de los \emph{ideales} de $\sh{O}_X$ (llamados también a veces \emph{ideales enteros}) que son $\sh{O}_X$-módulos invertibles: tal ideal $\sh{J}$ se caracteriza además por el hecho de que para todo $x\in X$ existe un entorno abierto $U$ de $x$ en $X$ tal que $\sh{J}|U=\sh{O}_U.f$, donde $f$ es un elemento regular de $\Gamma(U,\sh{O}_X)$.
El conjunto $\sh{J}_X(\operatorname{Div}^+(X))$ de estos ideales es, por tanto, un submonoide de $\operatorname{Id.inv}(X)$, igual al conjunto de los elementos positivos para una relación de orden compatible con la estructura de grupo de $\operatorname{Id.inv}(X)$, y es inmediato que esta relación no es otra que la relación \emph{opuesta} a la inclusión;
en otras palabras, se tiene
\[
\label{IV.21.2.7.2}
  \sh{J}_1\subset\sh{J}_2 \iff \operatorname{div}(\sh{J}_1)\geq\operatorname{div}(\sh{J}_2)
  \tag{21.2.7.2}
\]
para $\sh{J}_1$, $\sh{J}_2$ en $\operatorname{Id.inv}(X)$.
\end{env}

\begin{env}[21.2.8]
\label{IV.21.2.8}
Para todo divisor $D$ sobre $X$, se pone
\[
\label{IV.21.2.8.1}
  \sh{O}_X(D)=(\sh{J}_X(D))^{-1};
  \tag{21.2.8.1}
\]
$\sh{O}_X(D)$ es, por tanto, un ideal fraccionario invertible definido del modo siguiente: para todo abierto $U$ de $X$ tal que $D|U=\operatorname{div}_U(f)$, donde $f\in\Gamma(U,\sh{M}_X^*)$, $\sh{O}_X(D)|U$ es el ideal
\oldpage[IV-4]{261}
fraccionario invertible $\sh{O}_U.f^{-1}$;
en virtud de \sref{IV.21.1.6}, para toda función meromorfa regular $f$, se tiene la relación
\[
\label{IV.21.2.8.2}
  f\in\Gamma(X,\sh{O}_X(D)) \iff \operatorname{div}(f)\geq-D.
  \tag{21.2.8.2}
\]
Además, es claro que se tienen los isomorfismos canónicos \sref{IV.21.2.3}
\[
\label{IV.21.2.8.3}
  \left\{\begin{aligned}
    &\sh{O}_X(0)=\sh{O}_X,\quad \sh{O}_X(D+D')=\sh{O}_X(D)\sh{O}_X(D')\xrightarrow{\sim}\sh{O}_X(D)\otimes\sh{O}_X(D')\\
    &\sh{O}_X(nD)=(\sh{O}_X(D))^n\xrightarrow{\sim}(\sh{O}_X(D))^{\otimes n}
  \end{aligned}\right.
  \tag{21.2.8.3}
\]
para todo entero $n\in\mathbf{Z}$ y para dos divisores cualesquiera $D$, $D'$ sobre $X$.
\end{env}

\begin{env}[21.2.9]
\label{IV.21.2.9}
Sea $\sh{J}$ un ideal fraccionario invertible sobre $X$.
La inyección canónica $\sh{J}\hookrightarrow\sh{M}_X$ define por tensorización un homomorfismo de $\sh{O}_X$-módulos
\[
\label{IV.21.2.9.1}
  \sh{M}_X(\sh{J})=\sh{J}\otimes_{\sh{O}_X}\sh{M}_X\to\sh{M}_X\otimes_{\sh{O}_X}\sh{M}_X=\sh{M}_X
  \tag{21.2.9.1}
\]
que es un \emph{isomorfismo}: en efecto, si $U$ es un abierto de $X$ tal que $\sh{J}|U=\sh{O}_U.f$, donde $f\in\Gamma(U,\sh{M}_X^*)$, el homomorfismo \sref{IV.21.2.9.1} restringido a $U$ no es otro que el isomorfismo que a toda sección $t$ de $\sh{M}_X(\sh{J})|U=\sh{M}_X|U$ sobre $V\subset U$ hace corresponder la sección $t(f|V)^{-1}$ del mismo haz.
En el isomorfismo \sref{IV.21.2.9.1}, a las secciones meromorfas regulares de $\sh{J}$ sobre $X$ corresponden las funciones meromorfas regulares sobre $X$.

Consideremos en particular el caso en que $\sh{J}=\sh{O}_X(D)$, donde $D$ es un divisor sobre $X$;
se tiene entonces un isomorfismo canónico
\[
\label{IV.21.2.9.2}
  \sh{M}_X(\sh{O}_X(D))\xrightarrow{\sim}\sh{M}_X
  \tag{21.2.9.2}
\]
y se denota por $s_D$ la sección meromorfa regular de $\sh{O}_X(D)$ sobre $X$ que corresponde mediante este isomorfismo a la sección $1$ de $\sh{M}_X$.
Si $U$ es un abierto de $X$ tal que $\sh{O}_X(D)|U=\sh{O}_U.f^{-1}$, donde $f\in\Gamma(U,\sh{M}_X^*)$, se tiene $s_D|U=1$ en $\Gamma(U,\sh{M}_X)$;
como entonces se tiene $D|U=\operatorname{div}_U(f)$, se deduce de \sref{IV.21.1.4} que
\[
\label{IV.21.2.9.3}
  \operatorname{div}(s_D)=D.
  \tag{21.2.9.3}
\]
Por otra parte, de los isomorfismos canónicos \sref{IV.21.2.8.3} se deducen inmediatamente las fórmulas
\[
\label{IV.21.2.9.4}
  s_0=1,\quad s_{D+D'}=s_D\otimes s_{D'},\quad s_{nD}=s_D^{\otimes n}\quad(n\in\mathbf{Z}).
  \tag{21.2.9.4}
\]
\end{env}

\begin{env}[21.2.10]
\label{IV.21.2.10}
Consideremos, entre dos pares $(\sh{L},s)$, $(\sh{L}',s')$, donde $\sh{L}$ y $\sh{L}'$ son dos $\sh{O}_X$-módulos invertibles, y $s$ (respectivamente, $s'$) es una sección meromorfa regular de $\sh{L}$ (respectivamente, $\sh{L}'$) sobre $X$, la relación: «existe un isomorfismo $u:\sh{L}\xrightarrow{\sim}\sh{L}'$ tal que $\overline{u}(s)=s'$», donde $\overline{u}:\Gamma(X,\sh{M}_X(\sh{L}))\xrightarrow{\sim}\Gamma(X,\sh{M}_X(\sh{L}'))$ es el isomorfismo deducido de $u$ (obsérvese que el isomorfismo $u$ que satisface esta condición queda entonces determinado de manera única).
Es claro que se trata de una relación de equivalencia, y puesto que existe un conjunto de $\sh{O}_X$-módulos invertibles tal que todo $\sh{O}_X$-módulo invertible es isomorfo a un elemento de este conjunto (\EGAReadLink{0.5.4.7}{$\mathbf{0}_{\mathrm I}$, 5.4.7}), se puede hablar del conjunto $D(X)$ de las clases de equivalencia de los pares $(\sh{L},s)$ para la relación anterior.
Para todo $\sh{O}_X$-módulo invertible $\sh{L}$ y toda sección
\oldpage[IV-4]{262}
meromorfa regular $s$ de $\sh{L}$ sobre $X$, se denotará por $\operatorname{cl}(\sh{L},s)$ el elemento de $D(X)$ correspondiente al par $(\sh{L},s)$.
De (\EGAReadLink{0.5.4.3}{$\mathbf{0}_{\mathrm I}$, 5.4.3}) resulta que si $s$, $s'$ son dos secciones meromorfas regulares de $\sh{L}$ sobre $X$, la relación $\operatorname{cl}(\sh{L},s)=\operatorname{cl}(\sh{L},s')$ equivale a la existencia de una sección $t\in\Gamma(X,\sh{O}_X^*)$ tal que $s'=ts$.

Es inmediato que si $(\sh{L},s)$ es equivalente a $(\sh{L}_1,s_1)$ y $(\sh{L}',s')$ a $(\sh{L}_1',s_1')$, los pares $(\sh{L}\otimes\sh{L}',s\otimes s')$ y $(\sh{L}_1\otimes\sh{L}_1',s_1\otimes s_1')$ son equivalentes;
se define, por tanto, en $D(X)$ una ley de composición poniendo
\[
\label{IV.21.2.10.1}
  \operatorname{cl}(\sh{L},s)\operatorname{cl}(\sh{L}',s')=\operatorname{cl}(\sh{L}\otimes\sh{L}',s\otimes s');
  \tag{21.2.10.1}
\]
es inmediato que se trata de una ley de grupo conmutativo, cuyo elemento neutro es $\operatorname{cl}(\sh{O}_X,1)$ y en el que el inverso de $\operatorname{cl}(\sh{L},s)$ es $\operatorname{cl}(\sh{L}^{-1},s^{\otimes(-1)})$.
\end{env}

\begin{proposition}[21.2.11]
\label{IV.21.2.11}
Las aplicaciones
\[
\label{IV.21.2.11.1}
  D\mapsto\operatorname{cl}(\sh{O}_X(D),s_D),\qquad \operatorname{cl}(\sh{L},s)\mapsto\operatorname{div}(s)
  \tag{21.2.11.1}
\]
son isomorfismos inversos, respectivamente, de $\operatorname{Div}(X)$ sobre $D(X)$ y de $D(X)$ sobre $\operatorname{Div}(X)$.
\end{proposition}

\begin{proof}
Teniendo en cuenta \sref{IV.21.2.8.3}, \sref{IV.21.2.9.4} y \sref{IV.21.2.10.1}, basta ver que los compuestos de estas dos aplicaciones son las identidades en $\operatorname{Div}(X)$ y $D(X)$, respectivamente.
La primera afirmación no es otra que \sref{IV.21.2.9.3}.
Por otra parte, sea $D=\operatorname{div}(s)$, donde $s$ es una sección meromorfa regular de $\sh{L}$ sobre $X$, y sea $U$ un abierto de $X$ tal que exista un isomorfismo de $\sh{L}|U$ sobre $\sh{O}_U$ que transforme $s|U$ en $f\in\Gamma(U,\sh{M}_X^*)$, de modo que $D|U=\operatorname{div}_U(f)$, $\sh{O}_X(D)|U=\sh{O}_U.f^{-1}$ y $s_D|U$ es el elemento unidad de $\Gamma(U,\sh{M}_X)$.
Existe, por tanto, un isomorfismo $v_U:\sh{L}|U\to\sh{O}_U.f^{-1}=\sh{O}_X(D)|U$ tal que $\overline{v}_U$ (notación de \sref{IV.21.2.10}) transforma $s|U$ en $f.f^{-1}=1$;
se ve inmediatamente que estos isomorfismos son compatibles con las operaciones de restricción, por lo que definen un isomorfismo $v:\sh{L}\to\sh{O}_X(D)$ tal que $\overline{v}(s)=s_D$.
\end{proof}

Se puede transportar mediante el primero de los isomorfismos \sref{IV.21.2.11.1} la estructura de grupo ordenado de $\operatorname{Div}(X)$ a $D(X)$;
los elementos $\geq0$ de $D(X)$ son, por tanto, las clases $\operatorname{cl}(\sh{L},s)$ tales que $\operatorname{div}(s)\geq0$, es decir, \sref{IV.21.1.7.1}, tales que
\[
  s\in\Gamma(X,\sh{L})\cap\Gamma(X,(\sh{M}_X(\sh{L}))^*).
\]

\begin{env}[21.2.12]
\label{IV.21.2.12}
Sea $D$ un divisor \emph{positivo} sobre un \emph{preesquema} $X$;
el ideal fraccionario $\sh{J}_X(D)$ es entonces un ideal de $\sh{O}_X$ que es un $\sh{O}_X$-módulo invertible;
sea $Y(D)$ el subpreesquema cerrado de $X$ que este define.
Para todo $x\in Y(D)$ existe, por hipótesis, un entorno abierto $U$ de $x$ en $X$ y una sección regular $t\in\Gamma(U,\sh{O}_X)$ tal que $\sh{J}_X(D)|U=\sh{O}_U.t$ \sref{IV.21.2.7};
en otras palabras, la inmersión canónica $Y(D)\to X$ es \emph{regular y de codimensión $1$} en todo punto de $Y(D)$ \sref{IV.19.1.4}.
Recíprocamente, si $Y$ es un subpreesquema cerrado de $X$, inmerso regularmente en $X$ y de codimensión $1$ en todo punto de $Y$, existe un único divisor positivo $D$ tal que $Y(D)=Y$, pues para todo $x\in Y$ existe un entorno $U$ de $x$ en $X$ tal que $Y\cap U$ está definido por un ideal de $\sh{O}_U$ de la forma $\sh{O}_U.t$, donde $t$ es regular en $\Gamma(U,\sh{O}_X)$.
Obsérvese que entonces se tiene $\operatorname{Supp}(D)=Y(D)$, pues decir que $D_x\neq0$ significa que $t_x$ (con las notaciones anteriores) no es invertible, es decir, que $x\in Y(D)$.
\end{env}

\oldpage[IV-4]{263}


\subsection{Equivalencia lineal de los divisores}
\label{subsection:IV.21.3}

\begin{env}[21.3.1]
\label{IV.21.3.1}
Se dice que un divisor $D$ sobre $X$ es \emph{principal} si es de la forma $\operatorname{div}(f)$, donde $f$ es una función meromorfa regular sobre $X$;
las funciones meromorfas regulares $f'$ tales que $\operatorname{div}(f')=D$ son entonces todas las de la forma $tf$, donde $t\in\Gamma(X,\sh{O}_X^*)$ \sref{IV.21.1.3.4}.
El conjunto de los divisores principales es un subgrupo de $\operatorname{Div}(X)$, denotado por $\operatorname{Div.princ}(X)$, isomorfo a $\Gamma(X,\sh{M}_X^*)/\Gamma(X,\sh{O}_X^*)$.
Dos divisores $D$, $D'$ se llaman \emph{linealmente equivalentes} si $D-D'$ es un divisor principal;
los divisores principales son, por tanto, los divisores linealmente equivalentes a $0$.
\end{env}

\begin{env}[21.3.2]
\label{IV.21.3.2}
Recordemos (\EGAReadLink{0.5.4.7}{$\mathbf{0}_{\mathrm I}$, 5.4.7}) que se puede hablar del conjunto de las clases de equivalencia de los $\sh{O}_X$-módulos invertibles para la relación de isomorfismo;
este conjunto se denota por $\operatorname{Pic}(X)$, y para todo $\sh{O}_X$-módulo invertible $\sh{L}$, se denota por $\operatorname{cl}(\sh{L})$ la clase de equivalencia de los $\sh{O}_X$-módulos isomorfos a $\sh{L}$;
además, $\operatorname{Pic}(X)$ es un grupo conmutativo para la multiplicación definida por $\operatorname{cl}(\sh{L})\operatorname{cl}(\sh{L}')=\operatorname{cl}(\sh{L}\otimes\sh{L}')$.
Está claro que la aplicación
\[
\label{IV.21.3.2.1}
  r:\operatorname{cl}(\sh{L},s)\mapsto\operatorname{cl}(\sh{L})
  \tag{21.3.2.1}
\]
es un homomorfismo del grupo $D(X)$ \sref{IV.21.2.10} en el grupo $\operatorname{Pic}(X)$.
Por composición, se deduce, pues, un homomorfismo
\[
\label{IV.21.3.2.2}
  l:\operatorname{Div}(X)\xrightarrow{\sim}D(X)\xrightarrow{r}\operatorname{Pic}(X)
  \tag{21.3.2.2}
\]
(denotado también por $l_X$) tal que, para todo divisor $D$, se tiene
\[
\label{IV.21.3.2.3}
  l(D)=\operatorname{cl}(\sh{O}_X(D)).
  \tag{21.3.2.3}
\]

Observemos por último que, si $u:X'\to X$ es un morfismo de espacios anillados, y $\sh{L}_1$, $\sh{L}_2$ son dos $\sh{O}_X$-módulos invertibles isomorfos, los $\sh{O}_{X'}$-módulos invertibles $u^*(\sh{L}_1)$ y $u^*(\sh{L}_2)$ (\EGAReadLink{0.5.4.5}{$\mathbf{0}_{\mathrm I}$, 5.4.5}) son isomorfos;
como, por otra parte, para dos $\sh{O}_X$-módulos invertibles cualesquiera $\sh{L}_1$, $\sh{L}_2$, se tiene $u^*(\sh{L}_1\otimes\sh{L}_2)=u^*(\sh{L}_1)\otimes u^*(\sh{L}_2)$ salvo isomorfismo canónico (\EGAReadLink{0.4.3.3}{$\mathbf{0}_{\mathrm I}$, 4.3.3}), se ve que el morfismo $u$ define canónicamente un homomorfismo de grupos conmutativos
\[
\label{IV.21.3.2.4}
  \operatorname{Pic}(u):\operatorname{Pic}(X)\to\operatorname{Pic}(X').
  \tag{21.3.2.4}
\]
\end{env}

\begin{proposition}[21.3.3]
\label{IV.21.3.3}
\begin{enumerate}[label=(\roman*)]
  \item El núcleo del homomorfismo canónico $l:\operatorname{Div}(X)\to\operatorname{Pic}(X)$ es el subgrupo $\operatorname{Div.princ}(X)$; en otros términos, para que $\sh{O}_X(D)$ y $\sh{O}_X(D')$ sean isomorfos, es necesario y suficiente que $D$ y $D'$ sean linealmente equivalentes.
  Se tiene, por tanto, un homomorfismo inyectivo canónico
  \[
  \label{IV.21.3.3.1}
    \operatorname{Div}(X)/\operatorname{Div.princ}(X)\to\operatorname{Pic}(X)
    \tag{21.3.3.1}
  \]
  deducido de $l$.
  \item Para que un $\sh{O}_X$-módulo invertible $\sh{L}$ sea tal que $\operatorname{cl}(\sh{L})$ sea de la forma $l(D)$, o también para que $\sh{L}$ sea isomorfo a un $\sh{O}_X$-módulo de la forma $\sh{O}_X(D)$, es necesario y suficiente que exista una sección meromorfa regular de $\sh{L}$.
\end{enumerate}
\end{proposition}

\begin{proof}
\oldpage[IV-4]{264}
La proposición se deduce inmediatamente de las definiciones y de \sref{IV.21.2.10}.
\end{proof}

\begin{proposition}[21.3.4]
\label{IV.21.3.4}
Sea $X$ un preesquema que satisface una de las dos hipótesis siguientes:
\begin{enumerate}[label=\emph{\alph*)}]
  \item $X$ es localmente noetheriano y $\operatorname{Ass}(\sh{O}_X)$ está contenido en un abierto afín de $X$.
  \item $X$ es reducido y el conjunto de sus componentes irreducibles es localmente finito.
\end{enumerate}
Entonces el homomorfismo canónico $l:\operatorname{Div}(X)\to\operatorname{Pic}(X)$ es sobreyectivo y, al pasar al cociente, da un isomorfismo
\[
  \operatorname{Div}(X)/\operatorname{Div.princ}(X)\xrightarrow{\sim}\operatorname{Pic}(X).
\]
\end{proposition}

\begin{proof}
Basta mostrar que todo $\sh{O}_X$-módulo invertible $\sh{L}$ admite una sección meromorfa regular sobre $X$ \sref{IV.21.3.3}.
En ambos casos bastará, gracias a \sref{IV.20.2.11},~(ii), definir una sección $s$ de $(\sh{M}_X(\sh{L}))^*$ sobre un abierto esquemáticamente denso $U$ de $X$, o bien, en el caso a), sobre un abierto $U$ que contenga $\operatorname{Ass}(\sh{O}_X)$ \sref{IV.20.2.13},~(iv).
En efecto, sea entonces $V$ un abierto cualquiera de $X$ tal que $\sh{L}|V$ sea isomorfo a $\sh{O}_V$, de modo que existe un isomorfismo de $\sh{M}_X(\sh{L})|V$ sobre $\sh{M}_X|V$ que transforma $s|(U\cap V)$ en una sección $f_V$ de $\sh{M}_X^*$ sobre $U\cap V$.
Como $U\cap V$ es esquemáticamente denso en $V$, se deduce de \sref{IV.20.2.11} que existe una única función meromorfa regular $g_V$ en $V$ tal que $g_V|(U\cap V)=f_V$, y esta sección corresponde, por el isomorfismo considerado, a una sección meromorfa regular $u_V$ de $\sh{L}$ sobre $V$ tal que $u_V|(U\cap V)=s|(U\cap V)$.
Además, si $V'$ es un segundo abierto de $X$ tal que $\sh{L}|V'$ sea isomorfo a $\sh{O}_{V'}$, las restricciones de $u_V$ y $u_{V'}$ a $V\cap V'$ son iguales, pues corresponden por isomorfismo a dos funciones meromorfas que coinciden en un abierto esquemáticamente denso $U\cap V\cap V'$, y se concluye de nuevo por \sref{IV.20.2.11};
los $u_V$ son, por tanto, efectivamente las restricciones de una sección de $(\sh{M}_X(\sh{L}))^*$ sobre $X$.

Dicho esto, en el caso b), para cada uno de los puntos maximales $x_\lambda$ de $X$ se toma un abierto $U_\lambda$ que contiene a $x_\lambda$, que no interseca ninguna componente irreducible de $X$ distinta de $\overline{\{x_\lambda\}}$ y tal que $\sh{L}|U_\lambda$ sea isomorfo a $\sh{O}_{U_\lambda}$;
se tomará como $s$ la sección tal que $s|U_\lambda$ sea la sección de $\sh{L}|U_\lambda$ que corresponde, por el isomorfismo anterior, a la sección unidad de $\sh{O}_{U_\lambda}$.

En el caso a), por hipótesis podemos tomar como $U$ un abierto afín (y por tanto noetheriano); en otros términos, podemos suponer que $X=\operatorname{Spec}(A)$, donde $A$ es noetheriano, y $\sh{L}=\widetilde{P}$, donde $P$ es un $A$-módulo proyectivo de rango $1$.
Si $S$ es el conjunto de los elementos regulares de $A$, se tiene $\Gamma(X,\sh{M}_X)=S^{-1}A$ \sref{IV.20.2.12} y $\Gamma(X,\sh{M}_X(\sh{L}))=S^{-1}P$;
ahora bien, $S$ es el conjunto de los elementos que no pertenecen a ninguno de los ideales asociados a $A$, por lo que $S^{-1}A$ es un anillo semilocal cuyos ideales maximales provienen de los elementos maximales de $\operatorname{Ass}(A)$, y $S^{-1}P$ es un $S^{-1}A$-módulo proyectivo de rango $1$, y por tanto es libre de rango $1$ (Bourbaki, \emph{Alg. comm.}, cap.~II, \textsection~5, n\textsuperscript{o}~3, prop.~5);
un elemento que forme una base de este $S^{-1}A$-módulo es, por consiguiente, en virtud de \sref{IV.20.1.8}, una sección meromorfa regular de $\sh{L}$ sobre $X$.
\end{proof}

\begin{corollary}[21.3.5]
\label{IV.21.3.5}
Si $X$ es un preesquema noetheriano tal que existe un $\sh{O}_X$-módulo invertible amplio (\textbf{II}, 4.5.3) (por ejemplo, un preesquema cuasiproyectivo sobre un espectro de anillo noetheriano (\textbf{II}, 5.3.1 y 4.6.6)), el homomorfismo canónico $l:\operatorname{Div}(X)\to\operatorname{Pic}(X)$ es sobreyectivo.
\end{corollary}

\begin{proof}
\oldpage[IV-4]{265}
En efecto (\textbf{II}, 4.5.4), existe entonces un entorno abierto afín del conjunto finito $\operatorname{Ass}(\sh{O}_X)$.
\end{proof}

\begin{remark}[21.3.6]
\label{IV.21.3.6}
Recordemos (\EGAReadLink{0.5.4.7}{$\mathbf{0}_{\mathrm I}$, 5.4.7}) que se tiene un isomorfismo canónico $\pi:\mathrm{H}^1(X,\sh{O}_X^*)\xrightarrow{\sim}\operatorname{Pic}(X)$ definido de la manera siguiente.
Se parte de un $1$-cociclo $(c_{\alpha\beta})$ con valores en $\sh{O}_X^*$ correspondiente a un recubrimiento abierto $(U_\alpha)$ de $X$, siendo $c_{\alpha\beta}$ un elemento de $\Gamma(U_\alpha\cap U_\beta,\sh{O}_X^*)$, y se le asocia la clase del $\sh{O}_X$-módulo invertible obtenido pegando los $\sh{O}_{U_\alpha}$ según los isomorfismos $\sh{O}_{U_\alpha}|(U_\alpha\cap U_\beta)\xrightarrow{\sim}\sh{O}_{U_\beta}|(U_\alpha\cap U_\beta)$ definidos por la multiplicación por $c_{\alpha\beta}$.
Por otra parte, de la sucesión exacta de haces de grupos conmutativos
\[
  1\to\sh{O}_X^*\to\sh{M}_X^*\to\sh{D}iv_X\to0
\]
se deduce el homomorfismo de borde de la sucesión exacta de cohomología
\[
\label{IV.21.3.6.1}
  \partial:\operatorname{Div}(X)\to\mathrm{H}^1(X,\sh{O}_X^*).
  \tag{21.3.6.1}
\]
Mostremos que el homomorfismo compuesto
\[
  \operatorname{Div}(X)\xrightarrow{\partial}\mathrm{H}^1(X,\sh{O}_X^*)\xrightarrow{\pi}\operatorname{Pic}(X)
\]
no es otro que el homomorfismo $l$ definido en \sref{IV.21.3.2.2}.
En efecto, hay que partir de un divisor $D$ y de un recubrimiento abierto $(U_\alpha)$ de $X$ tal que $D|U_\alpha=\operatorname{div}_{U_\alpha}(g_\alpha)$, donde $g_\alpha$ es una función meromorfa regular sobre $U_\alpha$;
$\partial(D)$ es la clase de cohomología del cociclo $(c_{\alpha\beta})$, donde $c_{\alpha\beta}=g_{\alpha|\beta}g_{\beta|\alpha}^{-1}$, designando $g_{\alpha|\beta}$ la restricción de $g_\alpha$ a $U_\alpha\cap U_\beta$.
Está claro que la imagen por $\pi$ de esta clase de cohomología es la clase del ideal fraccionario invertible $\sh{L}$ tal que para todo $\alpha$, $\sh{L}|U_\alpha=\sh{O}_{U_\alpha}.g_\alpha^{-1}$, que por definición no es otro que $\sh{O}_X(D)$ \sref{IV.21.2.8}.
\end{remark}

\subsection{Imágenes recíprocas de divisores}
\label{subsection:IV.21.4}
\label{IV.21.4}

\begin{env}[21.4.1]
\label{IV.21.4.1}
Sea $f:X'\to X$ un morfismo de espacios anillados;
nos proponemos dar condiciones que permitan asociar a un divisor $D$ sobre $X$ un divisor $D'$ sobre $X'$, \emph{imagen recíproca} de $D$ por $f$.
Observemos primero que, para toda sección $t\in\Gamma(X,\sh{O}_X^*)$, la imagen de $t$ mediante el homomorfismo canónico $\Gamma(X,\sh{O}_X)\to\Gamma(X',\sh{O}_{X'})$ sigue siendo invertible; es decir, pertenece a $\Gamma(X',\sh{O}_{X'}^*)$.
Consideremos entonces $D$ dado por la clase de equivalencia de un par $(\sh{L},s)$, donde $\sh{L}$ es un $\sh{O}_X$-Módulo invertible y $s$ una sección meromorfa regular de $\sh{L}$ sobre $X$ \sref{IV.21.2.11}.
Formemos el $\sh{O}_{X'}$-Módulo invertible $f^*(\sh{L})=\sh{L}'$;
decir que la imagen recíproca $s\circ f$ de $s$ por $f$ existe \sref{IV.20.1.11} y es una sección meromorfa regular de $\sh{L}'$ sobre $X'$ equivale a decir que existen las imágenes recíprocas por $f$ de $s$ y de $s^{\otimes(-1)}$; es decir, que $s\in\Gamma(X,\sh{M}_f(\sh{L}))$ y $s^{\otimes(-1)}\in\Gamma(X,\sh{M}_f(\sh{L}^{-1}))$;
la observación anterior muestra entonces que, si $(\sh{L}_1,s_1)$ es un par equivalente a $(\sh{L},s)$ en el sentido de \sref{IV.21.2.10}, la imagen recíproca $s_1\circ f$ existe y es una sección meromorfa regular de $\sh{L}_1'=f^*(\sh{L}_1)$ sobre $X'$, y los pares $(\sh{L}',s\circ f)$ y $(\sh{L}_1',s_1\circ f)$ son equivalentes.
Se puede, pues, dar la siguiente definición:
\end{env}

\begin{definition}[21.4.2]
\label{IV.21.4.2}
Dado un morfismo $f:X'\to X$ de espacios anillados, se dice que la \emph{imagen recíproca por $f$} de un divisor $D$ sobre $X$ existe si se tiene $s_D\in\Gamma(X,\sh{M}_f(\sh{O}_X(D)))$ y $s_{-D}\in\Gamma(X,\sh{M}_f(\sh{O}_X(-D)))$ (cf. \sref{IV.20.1.11}).
Se llama entonces \emph{imagen recíproca} de $D$ por $f$ y se denota por $f^*(D)$ al divisor sobre $X'$ que corresponde canónicamente \sref{IV.21.2.11} a la clase del par $(f^*(\sh{O}_X(D)),s_D\circ f)$.
\end{definition}

De esta definición se deduce inmediatamente que, si $D$ y $D'$ tienen imágenes recíprocas por $f$, lo mismo ocurre con $-D$ y con $D+D'$ (teniendo en cuenta \sref{IV.21.2.9.4}), y que se tiene
\oldpage[IV-4]{266}
$f^*(D+D')=f^*(D)+f^*(D')$.
En otras palabras, el conjunto $\operatorname{Div}^f(X)$ de los divisores sobre $X$ cuya imagen recíproca por $f$ existe es un subgrupo de $\operatorname{Div}(X)$, y la aplicación $D\mapsto f^*(D)$ es un homomorfismo creciente del subgrupo ordenado $\operatorname{Div}^f(X)$ al grupo ordenado $\operatorname{Div}(X')$, que hace conmutativo el diagrama
\[
\label{IV.21.4.2.1}
  \xymatrix{
    \operatorname{Div}^f(X) \ar[r]^{l_X} \ar[d]_{f^*} & \operatorname{Pic}(X) \ar[d]^{\operatorname{Pic}(f)} \\
    \operatorname{Div}(X') \ar[r]^{l_{X'}} & \operatorname{Pic}(X')
  }
  \tag{21.4.2.1}
\]

\begin{env}[21.4.3]
\label{IV.21.4.3}
La definición \sref{IV.21.4.2} muestra inmediatamente que, para que $f^*(D)$ exista, es necesario y suficiente que, para todo abierto $U$ de $X$, exista la imagen recíproca por $f^U:f^{-1}(U)\to U$ (restricción de $f$) de $D|U$.
Ahora bien, si $D=\operatorname{div}(g)$, donde $g$ es una función meromorfa regular sobre $X$, decir que la imagen recíproca de $s_D$ por $f$ existe y es una sección meromorfa regular de $f^*(\sh{O}_X(D))$ significa \sref{IV.21.2.9} que la imagen recíproca de $g$ por $f$ existe y es una función meromorfa regular sobre $X'$.
De ello se deduce inmediatamente una segunda descripción de $\operatorname{Div}^f(X)$ y de $f^*(D)$: se considera el subhaz de grupos de $\sh{M}_X^*$, denotado $\sh{M}_f^*$, formado por los gérmenes de funciones meromorfas regulares definidas sobre un abierto de $X$ cuyas imágenes recíprocas por $f$ existen y son regulares sobre el abierto obtenido por imagen recíproca \sref{IV.20.1.11}.
Entonces, si $f=(\psi,\theta)$, el homomorfismo canónico \sref{IV.20.1.11} $\psi^*(\sh{M}_f)\to\sh{M}_{X'}$ induce por restricción un homomorfismo de haces de grupos $\psi^*(\sh{M}_f^*)\to\sh{M}_{X'}^*$.
Poniendo $\sh{D}iv_X'=\sh{M}_f^*/\sh{O}_X^*$, se tiene $\operatorname{Div}^f(X)=\Gamma(X,\sh{D}iv_X')$, y la aplicación $D\mapsto f^*(D)$ corresponde al homomorfismo de haces de grupos $\psi^*(\sh{M}_f^*)/\psi^*(\sh{O}_X^*)\to\sh{M}_{X'}^*/\sh{O}_{X'}^*=\sh{D}iv_{X'}$ deducido del anterior por paso a los cocientes.
\end{env}

\begin{env}[21.4.4]
\label{IV.21.4.4}
De las definiciones anteriores se deduce inmediatamente que, si $f':X''\to X'$ es un segundo morfismo de espacios anillados y $D$ es un divisor sobre $X$ tal que existen las imágenes recíprocas $f^*(D)$ y $f'^*(f^*(D))$, entonces $(f\circ f')^*(D)$ existe y es igual a $f'^*(f^*(D))$.
\end{env}

\begin{proposition}[21.4.5]
\label{IV.21.4.5}
Sea $f:X'\to X$ un morfismo de espacios anillados.
En cada uno de los tres casos siguientes, está definida la imagen recíproca por $f$ de todo divisor sobre $X$:
\begin{enumerate}[label=(\roman*)]
  \item $f$ es plano.
  \item $X$ y $X'$ son preesquemas localmente noetherianos y se tiene $f(\operatorname{Ass}(\sh{O}_{X'}))\subset\operatorname{Ass}(\sh{O}_X)$.
  \item $X$ y $X'$ son preesquemas, el conjunto de las componentes irreducibles de $X$ es localmente finito, $X'$ es reducido y toda componente irreducible de $X'$ domina una componente irreducible de $X$.
\end{enumerate}
\end{proposition}

\begin{proof}
Basta, en efecto, mostrar que en los tres casos se tiene $\sh{M}_f=\sh{M}_X$;
en el caso (i), esto se deduce de \sref{IV.20.1.12}.
En el caso (iii), podemos limitarnos al caso en que $X=\operatorname{Spec}(A)$ y $X'=\operatorname{Spec}(A')$ son afines;
si $s\in A$ es regular, no pertenece a ningún ideal primo minimal de $A$ \sref{IV.20.1.3.1}, por lo que la hipótesis implica que su imagen en $A'$ no pertenece a ningún ideal primo minimal de $A'$, y es, por consiguiente, un elemento regular de $A'$ \sref{IV.20.1.3.1}.
En el caso (ii), las funciones meromorfas sobre $X'$ se identifican con las pseudo-funciones sobre $X'$ \sref{IV.20.2.11}, y la hipótesis, junto con \sref{IV.20.2.13},~(iv), asegura que la imagen recíproca
\oldpage[IV-4]{267}
por $f$ de todo abierto esquemáticamente denso en $X$ es un abierto esquemáticamente denso en $X'$;
se concluye, pues, por \sref{IV.20.3.12}.
\end{proof}

\begin{corollary}[21.4.6]
\label{IV.21.4.6}
Sea $X$ un preesquema que posea una de las propiedades siguientes:
\begin{enumerate}[label=(\roman*)]
  \item $X$ es localmente noetheriano.
  \item $X$ es reducido y el conjunto de sus componentes irreducibles es localmente finito.
\end{enumerate}
Entonces, para todo $x\in X$, se tiene un isomorfismo canónico
\[
\label{IV.21.4.6.1}
  (\sh{D}iv_X)_x\xrightarrow{\sim}\operatorname{Div}(\sh{O}_{X,x}).
  \tag{21.4.6.1}
\]
\end{corollary}

\begin{proof}
Esto se deduce, en efecto, de \sref{IV.20.2.11}, \sref{IV.20.3.7} y del hecho de que $(\sh{O}_X^*)_x$ se identifica con el grupo de los elementos invertibles del anillo $\sh{O}_{X,x}$.
\end{proof}

\begin{env}[21.4.7]
\label{IV.21.4.7}
Sean $X$, $X'$ dos preesquemas y $f:X'\to X$ un morfismo.
Si $D$ es un divisor \emph{positivo} sobre $X$ tal que la imagen recíproca $f^*(D)$ esté definida \sref{IV.21.4.2}, entonces el subpreesquema cerrado $Y(f^*(D))$ de $X'$ no es sino la imagen recíproca $f^{-1}(Y(D))$;
esto se deduce inmediatamente de las definiciones \sref{IV.21.4.2} y \sref{IV.21.2.12}.
\end{env}

\begin{proposition}[21.4.8]
\label{IV.21.4.8}
Sean $X$, $Y$ dos preesquemas;
sea $f:X\to Y$ un morfismo fielmente plano.
Entonces, si un divisor $D$ sobre $Y$ es tal que $f^*(D)\geq0$ (la existencia de $f^*(D)$ se deduce de \sref{IV.21.4.5}), se tiene $D\geq0$.
En particular, la aplicación $D\mapsto f^*(D)$ de $\operatorname{Div}(Y)$ en $\operatorname{Div}(X)$ es inyectiva.
\end{proposition}

\begin{proof}
Como la cuestión es local sobre $Y$, podemos limitarnos al caso en que $D=\operatorname{div}(w)$, con $w=uv^{-1}$, siendo $u$ y $v$ dos secciones regulares de $\sh{O}_Y$ sobre $Y$.
Por hipótesis se tiene $u\sh{O}_X\subset v\sh{O}_X$, por lo que, para todo $x\in X$, si se pone $y=f(x)$, se tiene $u_y\sh{O}_{X,x}\subset v_y\sh{O}_{X,x}$;
se concluye que $u_y\sh{O}_{Y,y}\subset v_y\sh{O}_{Y,y}$ en virtud de la hipótesis de que $\sh{O}_{X,x}$ es un $\sh{O}_{Y,y}$-módulo fielmente plano y de Bourbaki, \emph{Alg. comm.}, cap.~I, \textsection~3, n\textsuperscript{o}~5, prop.~10,~(ii);
de donde $u\sh{O}_Y\subset v\sh{O}_Y$, puesto que $f$ es sobreyectivo, y por consiguiente $D\geq0$.
\end{proof}

% EGA IV, Errata y Addenda, lista 3, n.º 53.
% Este componente se inserta después de 21.4.8 en el texto de lectura corregido
% y se repite literalmente bajo Err_IV, 53 en el apéndice de erratas impreso.

La proposición siguiente y sus corolarios nos fueron señalados por
M.~Raynaud.

\begin{proposition}[21.4.9]
\ifdefined\EGAIVSuppressErratumLabels\else\label{IV.21.4.9}\fi
Sean $X$ e $Y$ dos preesquemas localmente noetherianos, y sea
$f:X\to Y$ un morfismo fielmente plano. Supongamos además que $f$ es
cuasicompacto o que $f$ es abierto. Un divisor $D'$ sobre $X$ es de la
forma $f^*(D)$, con $D$ un divisor sobre $Y$, si y solo si se cumple la
condición siguiente: para todo $y\in Y$ tal que
$\operatorname{prof}(\sh{O}_{Y,y})\leq 1$, pongamos
\[
 Y_1=\operatorname{Spec}(\sh{O}_{Y,y}),\qquad
 X_1=X\times_Y Y_1,
 \qquad f_1=f_{(Y_1)}:X_1\to Y_1;
\]
entonces la imagen recíproca $D_1'$ de $D'$ sobre $X_1$ es de la forma
$f_1^*(D_1)$ para algún $D_1\in\operatorname{Div}(Y_1)$.
\end{proposition}

\oldpage[IV]{359}
\begin{proof}
La necesidad de la condición resulta inmediatamente de la transitividad de las
imágenes recíprocas de divisores, pues todos los morfismos considerados son planos
{\EGAReadLink{IV.21.4.4}{(21.4.4)}}. Para todo abierto $U$ de $Y$, denotemos por $f_U$ la
restricción $f^{-1}(U)\to U$ de $f$. Por {\EGAReadLink{IV.21.4.8}{(21.4.8)}}, la aplicación
\[
 D_U\longmapsto f_U^*(D_U)
\]
de $\operatorname{Div}(U)$ en
$\operatorname{Div}(f^{-1}(U))$ es inyectiva. Así, si $U_1,U_2$ son
dos abiertos de $Y$ y
\[
 D'|f^{-1}(U_1)=f_{U_1}^*(D_1),\qquad
 D'|f^{-1}(U_2)=f_{U_2}^*(D_2),
\]
entonces necesariamente
$D_1|U_1\cap U_2=D_2|U_1\cap U_2$. Se sigue inmediatamente que existe un
mayor abierto $U\subset Y$ para el cual
\[
 D'|f^{-1}(U)=f_U^*(D_U)
 \quad\text{con}\quad D_U\in\operatorname{Div}(U).
\]
Debemos demostrar que $U=Y$. Supongamos lo contrario, y sea $y$ un punto
maximal de $Y\setminus U$. Encontraremos un entorno abierto $V$ de
$y$ en $Y$ tal que
  $D'|f^{-1}(V)=f_V^*(D_V)$ para algún
$D_V\in\operatorname{Div}(V)$, contradiciendo la maximalidad de $U$.

\textup{I) Caso en que $f$ es cuasicompacto.}
Utilizando la cuasicompacidad de $f$ y {\EGAReadLink{IV.20.3.8}{(20.3.8)}}, por el método de
({\EGAReadLink{IV.8.1.2}{8.1.2}}, a), nos reducimos a demostrar ---con las notaciones del
enunciado, pero sin imponer inicialmente una cota a
$\operatorname{prof}(\sh{O}_{Y,y})$--- que $D_1'$ es de la forma
$f_1^*(D_1)$. Si $\operatorname{prof}(\sh{O}_{Y,y})\leq 1$, esto es
precisamente la hipótesis. Supongamos, pues, que
$\operatorname{prof}(\sh{O}_{Y,y})\geq 2$. Como $y$ es un punto
maximal de $Y\setminus U$, el abierto $W=Y_1\cap U$ de $Y_1$ es
$Y_1\setminus\{y\}$. Podemos, por tanto, limitarnos al caso
$Y=Y_1$ y $U=Y\setminus\{y\}$.

Por hipótesis existe un divisor $D_U\in\operatorname{Div}(U)$ tal
que $f_U^*(D_U)=D'|f^{-1}(U)$. El divisor $D_U$ es la clase de
equivalencia de un par $(\sh{L}_U,s)$, donde $\sh{L}_U$ es un
$\sh{O}_U$-módulo invertible y $s$ una sección meromorfa regular de
$\sh{L}_U$ sobre $U$ {\EGAReadLink{IV.21.2.11}{(21.2.11)}}. Por tanto, en virtud de
{\EGAReadLink{IV.21.4.1}{(21.4.1)}}, $D'|f^{-1}(U)$ es la clase de equivalencia de
$(\sh{L}'_U,s')$, donde
\[
 \sh{L}'_U=f_U^*(\sh{L}_U),\qquad s'=s\circ f_U.
\]
Sean $i:U\to Y$ y $j:f^{-1}(U)\to X$ las inmersiones canónicas,
y pongamos
\[
 \sh{L}=i_*(\sh{L}_U),\qquad
 \sh{L}'=j_*(\sh{L}'_U).
\]
Como $f$ es plano, $f^*(\sh{L})$ es canónicamente isomorfo a
$\sh{L}'$ {\EGAReadLink{IV.2.3.1}{(2.3.1)}}. Además, la planitud y la desigualdad
$\operatorname{prof}(\sh{O}_{Y,y})\geq 2$ implican
$\operatorname{prof}(\sh{O}_{X,x})\geq 2$ para todo
$x\in f^{-1}(y)$ {\EGAReadLink{IV.6.3.1}{(6.3.1)}}.

El divisor $D'$ es la clase de equivalencia de un par
$(\sh{L}'',s'')$, donde $\sh{L}''$ es un $\sh{O}_X$-módulo invertible
y $s''$ una sección meromorfa regular de $\sh{L}''$ sobre $X$.
Por consiguiente, $\sh{L}'_U$ es isomorfo a
$\sh{L}''|f^{-1}(U)$, y las observaciones precedentes, junto con
{\EGAReadLink{IV.5.9.8}{(5.9.8)}} y {\EGAReadLink{IV.5.10.5}{(5.10.5)}}, muestran que $\sh{L}'$ es un
$\sh{O}_X$-módulo invertible isomorfo a $\sh{L}''$. Como $f$ es
fielmente plano y cuasicompacto, se sigue que $\sh{L}$ es un
$\sh{O}_Y$-módulo invertible {\EGAReadLink{IV.2.5.2}{(2.5.2)}} y, por tanto, isomorfo a $\sh{O}_Y$, pues
$Y$ es un esquema local.

Además, como $\operatorname{prof}(\sh{O}_{Y,y})\geq 1$, el abierto
$U=Y\setminus\{y\}$ es esquemáticamente denso en $Y$
({\EGAReadLink{IV.20.2.13}{20.2.13}}, (iv)). La sección meromorfa regular $s$ de $\sh{L}_U$
sobre $U$ se prolonga, por tanto, a una sección meromorfa regular $s_0$ de
$\sh{L}$ sobre $Y$ {\EGAReadLink{IV.20.2.11}{(20.2.11)}}. La clase de equivalencia de
$(\sh{L},s_0)$ es así un divisor $D$ sobre $Y$. Evidentemente, $f^*(D)$ y
$D'$ inducen el mismo divisor sobre
$f^{-1}(U)=X\setminus f^{-1}(y)$; y, como
$\operatorname{prof}(\sh{O}_{X,x})\geq 2$ para todo
$x\in f^{-1}(y)$, {\EGAReadLink{IV.21.1.8}{(21.1.8)}} da $D'=f^*(D)$. Esto prueba el caso I.

\textup{II) Caso en que $f$ es abierto.}
Se trata de encontrar un entorno abierto $V$ de $y$ en $Y$ tal
que $D'|f^{-1}(V)=f_V^*(D_V)$ para algún
$D_V\in\operatorname{Div}(V)$, por lo que podemos suponer primero que $Y$ es afín.
Como $f$ es abierto, las imágenes $f(W)$, cuando $W$ recorre los
abiertos afines de $X$, forman un recubrimiento abierto de $Y$. Como $Y$
es cuasicompacto, existe un abierto cuasicompacto $T\subset X$
tal que $f(T)=Y$, y $f|T$ es cuasicompacto (I, 1.1.1). Con mayor razón,
para todo abierto cuasicompacto $T_\alpha\supset T$ de $X$,
\oldpage[IV]{360}
se tiene
$f(T_\alpha)=Y$ y $f|T_\alpha$ es cuasicompacto.

Aplicando el caso I a $f|T_\alpha$, obtenemos un divisor $D_\alpha$ sobre
$Y$ tal que
\[
 (f|T_\alpha)^*(D_\alpha)=D'|T_\alpha.
\]
Por transitividad de las imágenes recíprocas de divisores por morfismos planos
{\EGAReadLink{IV.21.4.4}{(21.4.4)}}, también se tiene
$(f|T)^*(D_\alpha)=D'|T$. Por tanto, para cualesquiera $\alpha,\beta$,
\[
 (f|T)^*(D_\alpha)=(f|T)^*(D_\beta),
\]
y {\EGAReadLink{IV.21.4.8}{(21.4.8)}} da $D_\alpha=D_\beta$. Si $D$ denota su valor
común, entonces $f^*(D)|T_\alpha=D'|T_\alpha$ para todo $\alpha$, y
por consiguiente $f^*(D)=D'$.
\end{proof}

\begin{corollary}[21.4.10]
\ifdefined\EGAIVSuppressErratumLabels\else\label{IV.21.4.10}\fi
Supongamos las hipótesis generales de {\EGAReadLink{IV.21.4.9}{(21.4.9)}} y, además, que
$Y$ es normal y que, para todo $y\in Y$ con
$\dim(\sh{O}_{Y,y})=1$, la fibra $X_y=f^{-1}(y)$ no tiene ningún
ciclo primo asociado inmerso. Un divisor $D'$ sobre $X$ es de la forma $f^*(D)$,
con $D\in\operatorname{Div}(Y)$, si y solo si se cumplen las dos
condiciones siguientes:
\begin{enumerate}[label=\textup{\arabic*\textdegree}]
\item para todo punto maximal $y$ de $Y$, la imagen recíproca de $D'$ en
$\operatorname{Div}(X_y)$ es cero;
\item para todo $y\in Y$ tal que $\dim(\sh{O}_{Y,y})=1$, existe un
entero $n_y$ tal que, para todo punto maximal $x$ de $X_y$, la
imagen recíproca de $D'$ en
$\operatorname{Div}(\operatorname{Spec}(\sh{O}_{X,x}))$ es igual a
$\operatorname{div}(t_y^{n_y})$, donde $t_y$ es la imagen por
$\sh{O}_{Y,y}\to\sh{O}_{X,x}$ de un uniformizante del anillo de
valoración discreta $\sh{O}_{Y,y}$.
\end{enumerate}
\end{corollary}

\begin{proof}
Aplicamos el criterio {\EGAReadLink{IV.21.4.9}{(21.4.9)}}. En un punto maximal $y$ de $Y$,
$\operatorname{Div}(Y_1)=0$ porque $\sh{O}_{Y,y}$ es artiniano;
la condición 1\textdegree{} es, por tanto, equivalente al criterio de
{\EGAReadLink{IV.21.4.9}{(21.4.9)}} en tal punto. Queda por demostrar que las condiciones
1\textdegree{} y 2\textdegree{} implican asimismo el criterio en un
punto $y\in Y$ con $\operatorname{prof}(\sh{O}_{Y,y})=1$. Como $Y$
es normal, estos son precisamente los puntos para los cuales
$\dim(\sh{O}_{Y,y})=1$, equivalentemente, aquellos para los cuales
$\sh{O}_{Y,y}$ es un anillo de valoración discreta {\EGAReadLink{IV.5.8.6}{(5.8.6)}}.

Podemos, pues, limitarnos a $Y=\operatorname{Spec}(\sh{O}_{Y,y})$. Si $t$
es un uniformizante de $\sh{O}_{Y,y}$, la aplicación
$n\mapsto\operatorname{div}(t^n)$ es un isomorfismo de $\mathbb Z$
sobre $\operatorname{Div}(Y)$. La condición de {\EGAReadLink{IV.21.4.9}{(21.4.9)}} dice que
$D'$ es de la forma $\operatorname{div}(t^n\circ f)$ para algún
$n\in\mathbb Z$, e implica evidentemente la condición 2\textdegree{}.
Recíprocamente, supongamos que se cumple 2\textdegree{}, así como 1\textdegree{}, y
escribamos $n=n_y$. Queda por demostrar
$D'=\operatorname{div}(t^n\circ f)$. Si $\eta$ es el punto genérico de
$Y$, los gérmenes de estos dos divisores en todo punto de $X_\eta$ son
cero por la condición 1\textdegree{}. Sus gérmenes en todo punto maximal
de $X_y$ son iguales por la condición 2\textdegree{}. En un punto no maximal
$x\in X_y$, se tiene
$\operatorname{prof}(\sh{O}_{X_y,x})\geq 1$, porque $X_y$ no tiene
ciclo primo asociado inmerso y, por tanto,
$\operatorname{prof}(\sh{O}_{X,x})\geq 2$ por {\EGAReadLink{IV.6.3.1}{(6.3.1)}}. La
igualdad buscada resulta de {\EGAReadLink{IV.21.1.8}{(21.1.8)}}.
\end{proof}

\begin{corollary}[21.4.11]
\ifdefined\EGAIVSuppressErratumLabels\else\label{IV.21.4.11}\fi
Supongamos las hipótesis generales de {\EGAReadLink{IV.21.4.9}{(21.4.9)}} y, además, que
$Y$ es normal y que, para todo $y\in Y$ con
$\dim(\sh{O}_{Y,y})=1$, la fibra $X_y$ es íntegra. Un divisor $D'$ es
de la forma $f^*(D)$, con $D\in\operatorname{Div}(Y)$, si y solo si,
para todo punto maximal $y$ de $Y$, la imagen recíproca de $D'$ en
$\operatorname{Div}(X_y)$ es cero.
\end{corollary}

\begin{proof}
Es suficiente comprobar que la condición 2\textdegree{} de {\EGAReadLink{IV.21.4.10}{(21.4.10)}} es entonces
automática. Si $\dim(\sh{O}_{Y,y})=1$ y $X_y$ es íntegra, entonces en
el único punto maximal $x$ de $X_y$ se tiene
$\dim(\sh{O}_{X,x})=1$ por {\EGAReadLink{IV.6.1.3}{(6.1.3)}}. Además,
$\sh{O}_{Y,y}$ es regular y $\sh{O}_{X_y,x}$ es un cuerpo, por lo que
$\sh{O}_{X,x}$ es regular {\EGAReadLink{IV.6.5.1}{(6.5.1)}} y, por tanto, un anillo de valoración discreta.
Si $t$ es un uniformizante de $\sh{O}_{Y,y}$, su imagen $t_y$ en
$\sh{O}_{X,x}$ es un uniformizante de $\sh{O}_{X,x}$
($0$, 17.3.3). La imagen recíproca de $D'$ en
$\operatorname{Div}(\operatorname{Spec}(\sh{O}_{X,x}))$ es, por tanto,
\oldpage[IV]{361}
de la forma $\operatorname{div}(t_y^n)$, lo que es exactamente la condición
2\textdegree{} de {\EGAReadLink{IV.21.4.10}{(21.4.10)}}, puesto que $X_y$ solo tiene un punto maximal.
\end{proof}

\begin{remark}[21.4.12]
\ifdefined\EGAIVSuppressErratumLabels\else\label{IV.21.4.12}\fi
Si en {\EGAReadLink{IV.21.4.11}{(21.4.11)}} se supone solamente que los $X_y$ son reducidos, hay que
añadir a la condición del enunciado que las imágenes recíprocas de $D'$
en cada uno de los grupos
$\operatorname{Div}(\operatorname{Spec}(\sh{O}_{X,x}))$, cuando $x$
recorre los puntos maximales de $X_y$, sean todas de la forma
$\operatorname{div}((t_x')^n)$, donde $t_x'$ es un uniformizante de
$\sh{O}_{X,x}$ y el entero $n$ es el mismo para todos estos puntos.
\end{remark}

\begin{corollary}[21.4.13]
\ifdefined\EGAIVSuppressErratumLabels\else\label{IV.21.4.13}\fi
Supongamos las hipótesis de {\EGAReadLink{IV.21.4.11}{(21.4.11)}}. Para todo
$\sh{O}_X$-módulo invertible $\sh{L}$ tal que, para todo punto maximal $\eta$ de
$Y$, la imagen recíproca
\[
 \sh{L}_\eta
 =\sh{L}\otimes_{\sh{O}_X}\sh{O}_{X_\eta}
\]
sea isomorfa a $\sh{O}_{X_\eta}$, existen un
$\sh{O}_Y$-módulo invertible $\sh{M}$ y un isomorfismo
$f^*(\sh{M})\simeq\sh{L}$.
\end{corollary}

\begin{proof}
Basta encontrar una sección meromorfa regular $s$ de $\sh{L}$ sobre
$X$ tal que el divisor $D'$ correspondiente a $(\sh{L},s)$ tenga
$D'_x=0$ en todo punto $x$ de una fibra $X_\eta$ sobre un punto maximal
$\eta$ de $Y$. Aplicando entonces {\EGAReadLink{IV.21.4.11}{(21.4.11)}} a $D'$, se obtiene el resultado por
{\EGAReadLink{IV.21.4.2}{(21.4.2)}}.

Los puntos de $\operatorname{Ass}(\sh{O}_X)$ están sobre los puntos
maximales de $Y$ {\EGAReadLink{IV.3.3.1}{(3.3.1)}}. Como $Y$ es reducido y localmente noetheriano,
existe un abierto denso $V$ de $Y$ que es la unión de una
familia disjunta de abiertos afines, cada uno de los cuales contiene exactamente un
punto maximal de $Y$; y $f^{-1}(V)$ es esquemáticamente denso en $X$
{\EGAReadLink{IV.5.10.2}{(5.10.2)}}. Podemos, por tanto, limitarnos al caso en que $Y$ es
afín e íntegro.

Para el único punto maximal $\eta$ de $Y$, la fibra $X_\eta$ es entonces
íntegra, de modo que $\operatorname{Ass}(\sh{O}_X)$ consta de un solo punto
$z$, el punto genérico de $X$. Sea $s_\eta$ la sección de
$\sh{L}_\eta$ sobre $X_\eta$ que corresponde, mediante un isomorfismo,
a la sección unidad de $\sh{O}_{X_\eta}$. Sea $U$ un entorno abierto
afín en $X$ de un punto $x\in X_\eta$, y sea $s$ una sección
de $\sh{L}$ sobre $U$ con el mismo germen en $x$ que $s_\eta$. Como $U$
es esquemáticamente denso en $X$ ({\EGAReadLink{IV.20.2.13}{20.2.13}}, (iv)), $s$ se identifica con una
sección meromorfa regular unívocamente determinada de $\sh{L}$ sobre $X$.
Resulta claramente de {\EGAReadLink{IV.20.2.11}{(20.2.11)}} y {\EGAReadLink{IV.20.3.5}{(20.3.5)}} que esta sección induce sobre
$X_\eta$ una sección meromorfa regular de $\sh{L}_\eta$ igual a
$s_\eta$ y, por tanto, tiene la propiedad requerida.
\end{proof}


\subsection{Imágenes directas de divisores}
\label{subsection:IV.21.5}
\label{IV.21.5}

\begin{env}[21.5.1]
\label{IV.21.5.1}
Sean $X$, $X'$ dos preesquemas, y $f:X'\to X$ un morfismo.
En este número daremos condiciones suficientes para que se pueda asociar a todo divisor $D'$ sobre $X'$ un divisor $D$ sobre $X$, \emph{imagen directa} de $D'$ por $f$.
Nos limitaremos al caso en que $f$ sea un morfismo \emph{finito} (para condiciones más generales, véase el capítulo de este Tratado dedicado a la teoría de intersecciones).
\end{env}

\begin{lemma}[21.5.2]
\label{IV.21.5.2}
Sean $A$ un anillo y $E$ un $A$-módulo libre de rango finito.
Para que un endomorfismo $u$ de $E$ sea inyectivo, es necesario y suficiente que $\det(u)$ sea un elemento regular de $A$.
\end{lemma}

\begin{proof}
Esto se prueba en Bourbaki, \emph{Alg.}, cap.~III, 3\textsuperscript{a} ed., \textsection~8, n\textsuperscript{o}~2, prop.~3.
\end{proof}

\begin{env}[21.5.3]
\label{IV.21.5.3}
Supongamos ahora que el morfismo $f:X'\to X$ sea finito y que, además, $f$ satisfaga una de las dos propiedades siguientes:
\begin{enumerate}[label=\Roman*)]
  \item $f$ es un morfismo finito \emph{localmente libre}; en otras palabras, \sref{IV.18.2.7}, el $\sh{O}_X$-módulo cuasicoherente de tipo finito $f_*(\sh{O}_{X'})$ es localmente libre.
  \item $X$ es un preesquema localmente noetheriano reducido, el $\sh{R}(X)$-módulo $f_*(\sh{O}_{X'})\otimes_{\sh{O}_X}\sh{R}(X)$ es localmente libre y, para toda sección $s'\in\Gamma(f^{-1}(U),\sh{O}_{X'})$ ($U$ abierto
  \oldpage[IV-4]{268}
  en $X$), $\mathrm{N}_{f_*(\sh{O}_{X'})|\sh{O}_X}(s')$ es una sección de $\sh{O}_X$ sobre $U$ (cf. (\textbf{II}, 6.5.1)).
  (Se recuerda que la condición (II) se satisface para todo morfismo finito $f$ cuando $X$ es un preesquema localmente noetheriano normal (\emph{loc. cit.}).)
\end{enumerate}

Se sabe entonces (\textbf{II}, 6.5.5) que para todo $\sh{O}_{X'}$-módulo invertible $\sh{L}'$ se define la norma $\sh{L}=\mathrm{N}_{X'/X}(\sh{L}')$ (que escribiremos también $\mathrm{N}(\sh{L}')$), que es un $\sh{O}_X$-módulo invertible.
Además, para todo abierto $U$ de $X$ y toda sección regular $s'\in\Gamma(f^{-1}(U),\sh{O}_{X'})$, la norma $\mathrm{N}_{X'/X}(s')=\mathrm{N}_{f_*(\sh{O}_{X'})|\sh{O}_X}(s')$ (que escribiremos también $\mathrm{N}(s')$) es un elemento regular de $\Gamma(U,\sh{O}_X)$;
en efecto, se reduce inmediatamente al caso en que $U=X$ es afín, y entonces la conclusión resulta de \sref{IV.21.5.2} bajo la hipótesis (I);
por otra parte, bajo la hipótesis (II), el hecho de que $\sh{R}(X)$ sea un $\sh{O}_X$-módulo plano implica que la sección $s'\otimes1$ de $\Gamma(U,f_*(\sh{O}_{X'})\otimes_{\sh{O}_X}\sh{R}(X))$ es también regular (\EGAReadLink{0.6.1.2}{$\mathbf{0}_{\mathrm I}$, 6.1.2}), y la conclusión resulta de nuevo de aplicar \sref{IV.21.5.2} al anillo $\Gamma(U,\sh{R}(X))$, teniendo en cuenta la definición de la norma de una sección (\textbf{II}, 6.5.3).
Dicho esto, sea $u'$ una sección meromorfa de $\sh{L}'$ sobre $X'$;
como el morfismo $f$ es afín, todo punto $x\in X$ admite un entorno abierto $U$ tal que $u'|f^{-1}(U)$ sea de la forma $t'/s'$, donde $t'\in\Gamma(f^{-1}(U),\sh{L}')$ y $s'$ es una sección regular en $\Gamma(f^{-1}(U),\sh{O}_{X'})$;
el elemento $\mathrm{N}(t')/\mathrm{N}(s')$ (donde $\mathrm{N}(t')$ es la sección de $\sh{L}$ definida en (\textbf{II}, 6.5.3)) es entonces una sección meromorfa de $\sh{L}$ sobre $U$ en virtud de lo anterior, y de las propiedades multiplicativas de la norma (\textbf{II}, 6.5.3.1) resulta que esta sección solo depende de $u'|f^{-1}(U)$ y no de su escritura en la forma $t'/s'$;
por la misma razón, al variar $U$, las secciones meromorfas de $\sh{L}$ sobre $U$ así definidas son las restricciones de una misma sección meromorfa de $\sh{L}$ sobre $X$, que se llama la \emph{norma} de $u'$ y se denota $\mathrm{N}_{X'/X}(u')$ (o simplemente $\mathrm{N}(u')$).
La aplicación así definida
\[
\label{IV.21.5.3.1}
  \mathrm{N}_{X'/X}:\Gamma(X',\sh{M}_{X'}(\sh{L}'))\to\Gamma(X,\sh{M}_X(\mathrm{N}_{X'/X}(\sh{L}')))
  \tag{21.5.3.1}
\]
prolonga la norma definida en (\textbf{II}, 6.5.5.3);
si $u'$ es una sección meromorfa regular de $\sh{L}'$ sobre $X'$, de lo anterior resulta inmediatamente que $\mathrm{N}(u')$ es una sección meromorfa regular de $\sh{L}$ sobre $X$, pues (con las mismas notaciones) $\mathrm{N}(t')$ es regular si $t'$ lo es.
Finalmente, si $\sh{L}_1'$, $\sh{L}_2'$ son dos $\sh{O}_{X'}$-módulos invertibles, y $s_1'$ (respectivamente, $s_2'$) es una sección meromorfa regular de $\sh{L}_1'$ (respectivamente, $\sh{L}_2'$) sobre $X'$, se tiene, en virtud de lo anterior y de (\textbf{II}, 6.5.3.1),
\[
\label{IV.21.5.3.2}
  \mathrm{N}(s_1'\otimes s_2')=\mathrm{N}(s_1')\otimes\mathrm{N}(s_2').
  \tag{21.5.3.2}
\]
\end{env}

\begin{env}[21.5.4]
\label{IV.21.5.4}
Supongamos siempre que $f$ satisface una de las hipótesis I), II) de \sref{IV.21.5.3}.
Si $(\sh{L}_1',s_1')$, $(\sh{L}_2',s_2')$ son dos pares formados cada uno por un $\sh{O}_{X'}$-módulo invertible y una sección meromorfa regular de este módulo sobre $X'$, que además son equivalentes en el sentido de \sref{IV.21.2.10}, entonces los pares $(\mathrm{N}(\sh{L}_1'),\mathrm{N}(s_1'))$ y $(\mathrm{N}(\sh{L}_2'),\mathrm{N}(s_2'))$ son también equivalentes, pues la norma de un isomorfismo de $\sh{O}_{X'}$-módulos invertibles es un isomorfismo de sus normas (\textbf{II}, 6.5.3), y se ha visto anteriormente que $\mathrm{N}_{X'/X}$ transforma las secciones de $\Gamma(X',\sh{O}_{X'}^*)$ en secciones de $\Gamma(X,\sh{O}_X^*)$.
Se puede, por tanto, dar la definición siguiente:
\end{env}

\oldpage[IV-4]{269}

\begin{definition}[21.5.5]
\label{IV.21.5.5}
Dado un morfismo finito $f:X'\to X$ de preesquemas que satisfaga una de las condiciones I), II) de \sref{IV.21.5.3}, se llama \emph{imagen directa} (o \emph{norma}) de un divisor $D'$ sobre $X'$ por $f$, y se denota $f_*(D')$ (o $\mathrm{N}_{X'/X}(D')$), el divisor sobre $X$ que corresponde canónicamente \sref{IV.21.2.11} a la clase del par $(\mathrm{N}_{X'/X}(\sh{O}_{X'}(D')),\mathrm{N}_{X'/X}(s_{D'}))$.
\end{definition}

De esta definición resulta inmediatamente, teniendo en cuenta \sref{IV.21.2.9.4} y \sref{IV.21.5.3.2}, que si $D_1'$, $D_2'$, $D'$ son divisores sobre $X'$, se tiene
\[
\label{IV.21.5.5.1}
  f_*(D_1'+D_2')=f_*(D_1')+f_*(D_2')
  \tag{21.5.5.1}
\]
y $D'\geq0$ implica $f_*(D')\geq0$; en otras palabras, $D'\mapsto f_*(D')$ es un homomorfismo creciente del grupo ordenado $\operatorname{Div}(X')$ en el grupo ordenado $\operatorname{Div}(X)$.
La definición \sref{IV.21.5.5} muestra además inmediatamente que para todo abierto $U$ de $X$ se tiene $f_*^U(D'|f^{-1}(U))=f_*(D')|U$ (siendo $f^U$ la restricción $f^{-1}(U)\to U$ de $f$), y los homomorfismos $f_*^U$, para $U$ variable, definen por tanto un homomorfismo de haces de grupos ordenados
\[
\label{IV.21.5.5.2}
  \mathrm{N}_{X'/X}:f_*(\sh{D}iv_{X'})\to\sh{D}iv_X.
  \tag{21.5.5.2}
\]
Además, para todo abierto $U$ de $X$, todo $\sh{O}_{X'}$-módulo invertible $\sh{L}'$ y toda sección meromorfa regular $s'$ de $\sh{L}'$ sobre $f^{-1}(U)$, se tiene, según las definiciones anteriores y \sref{IV.21.1.4},
\[
\label{IV.21.5.5.3}
  \operatorname{div}_U(\mathrm{N}(s'))=f_*^U(\operatorname{div}_{f^{-1}(U)}(s')).
  \tag{21.5.5.3}
\]

\begin{proposition}[21.5.6]
\label{IV.21.5.6}
Sea $f:X'\to X$ un morfismo finito localmente libre y supongamos que $f_*(\sh{O}_{X'})$ sea de rango constante $n$.
Entonces, para todo divisor $D$ sobre $X$, $f^*(D)$ está definido\footnote{El texto francés impreso da $f_*(D)$; la fórmula siguiente y la demostración requieren $f^*(D)$.} y se tiene
\[
\label{IV.21.5.6.1}
  f_*(f^*(D))=nD.
  \tag{21.5.6.1}
\]
\end{proposition}

\begin{proof}
La primera afirmación resulta del hecho de que $f$ es plano \sref{IV.21.4.5}, y la segunda es consecuencia inmediata de las definiciones y de (\textbf{II}, 6.5.3.2).
\end{proof}

\begin{proposition}[21.5.7]
\label{IV.21.5.7}
Sean $f:X'\to X$ un morfismo finito que satisface una de las hipótesis I), II) de \sref{IV.21.5.3}, y $f':X''\to X'$ un morfismo finito localmente libre de rango constante $n$.
Entonces $f''=f\circ f':X''\to X$ satisface la misma hipótesis que $f$, y para todo divisor $D''$ sobre $X''$ se tiene
\[
\label{IV.21.5.7.1}
  f_*''(D'')=f_*(f_*'(D'')).
  \tag{21.5.7.1}
\]
\end{proposition}

Teniendo en cuenta la definición \sref{IV.21.5.5}, basta probar el resultado siguiente:

\begin{lemma}[21.5.7.2]
\label{IV.21.5.7.2}
Bajo las hipótesis de \sref{IV.21.5.7}, se tiene un isomorfismo functorial
\[
\label{IV.21.5.7.3}
  \mathrm{N}_{X''/X}(\sh{L}'')\xrightarrow{\sim}\mathrm{N}_{X'/X}(\mathrm{N}_{X''/X'}(\sh{L}''))
  \tag{21.5.7.3}
\]
en la categoría de los $\sh{O}_X$-módulos invertibles.
\end{lemma}

En efecto, teniendo en cuenta la definición de la norma de una sección de $\sh{L}''$ (\textbf{II}, 6.5.3) y la definición \sref{IV.21.5.5}, se obtiene inmediatamente \sref{IV.21.5.7.1}.
Para probar \sref{IV.21.5.7.2}, basta, teniendo en cuenta las definiciones de (\textbf{II}, 6.5.2 y 6.5.3), probar que para toda sección $s$ de $f_*''(\sh{O}_{X''})=f_*(f_*'(\sh{O}_{X''}))$ sobre un abierto $U$ de $X$, se tiene
\[
\label{IV.21.5.7.4}
  \mathrm{N}_{f_*''(\sh{O}_{X''})|\sh{O}_X}(s)=\mathrm{N}_{f_*(\sh{O}_{X'})|\sh{O}_X}(\mathrm{N}_{f_*'(\sh{O}_{X''})|\sh{O}_{X'}}(s)).
  \tag{21.5.7.4}
\]

\oldpage[IV-4]{270}

La cuestión es evidentemente local sobre $X$, por lo que podemos limitarnos al caso en que $X=\operatorname{Spec}(A)$ sea afín;
se tiene entonces $X'=\operatorname{Spec}(A')$ y $X''=\operatorname{Spec}(A'')$, y puede suponerse que $A''$ es un $A'$-módulo proyectivo de rango $n$.
Cuando $f$ es localmente libre, puede suponerse que $A'$ es un $A$-módulo libre de rango $m$, y entonces $A''$ es un $A$-módulo proyectivo de rango $mn$; restringiendo $X$ a un abierto conveniente, puede suponerse que $A''$ es un $A$-módulo libre de rango $mn$;
la fórmula \sref{IV.21.5.7.4} resulta entonces de la transitividad de la norma (Bourbaki, \emph{Alg.}, cap.~VIII, \textsection~12, n\textsuperscript{o}~2, prop.~7).
Cuando $f$ satisface la hipótesis II), $A$ es noetheriano reducido y, si $R$ es su anillo total de fracciones, $A'\otimes_A R$ es un $R$-módulo libre de rango $m$; por tanto, $A''\otimes_A R=A''\otimes_{A'}(A'\otimes_A R)$ es un $R$-módulo proyectivo de rango $mn$, y como $R$ es entonces un anillo semilocal, este $R$-módulo es libre;
la proposición resulta entonces de nuevo de la transitividad de las normas.

\begin{proposition}[21.5.8]
\label{IV.21.5.8}
Sean $f:X'\to X$ un morfismo finito y $g:Y\to X$ un morfismo;
se pone $Y'=X'\times_XY$, $f'=f_{(Y)}:Y'\to Y$, $g'=g_{(X')}:Y'\to X'$.
Supongamos que se satisface una de las condiciones siguientes:
\begin{enumerate}[label=(\roman*)]
  \item $f$ es localmente libre y $g$ es plano.
  \item $f$ es localmente libre, $X$ e $Y$ son localmente noetherianos, $g(\operatorname{Ass}(\sh{O}_Y))\subset\operatorname{Ass}(\sh{O}_X)$ y $g'(\operatorname{Ass}(\sh{O}_{Y'}))\subset\operatorname{Ass}(\sh{O}_{X'})$.
  \item $f$ satisface la hipótesis II) de \sref{IV.21.5.3}, $Y$ es localmente noetheriano, $Y$ e $Y'$ son reducidos y toda componente irreducible de $Y$ (respectivamente, $Y'$) domina una componente irreducible de $X$ (respectivamente, $X'$).
\end{enumerate}
Entonces, para todo divisor $D'$ sobre $X'$, $g'^*(D')$ está definido, $g^*(f_*(D'))$ está definido y se tiene
\[
\label{IV.21.5.8.1}
  g^*(f_*(D'))=f_*'(g'^*(D')).
  \tag{21.5.8.1}
\]
\end{proposition}

\begin{proof}
En efecto, en todos los casos resulta de (\textbf{II}, 6.5.8) que se tiene un isomorfismo functorial
\[
  g^*(\mathrm{N}_{X'/X}(\sh{L}'))\xrightarrow{\sim}\mathrm{N}_{Y'/Y}(g'^*(\sh{L}'))
\]
en la categoría de los $\sh{O}_Y$-módulos invertibles;
además (\textbf{II}, 6.5.4), si $s'$ es una sección de $\sh{O}_{X'}$ sobre $f^{-1}(U)$ y $s''$ la sección correspondiente de $\sh{O}_{Y'}$ sobre $g'^{-1}(f^{-1}(U))$ ($U$ abierto de $X$), $\mathrm{N}_{Y'/Y}(s'')$ es la sección de $\sh{O}_Y$ sobre $g^{-1}(U)$ que corresponde a la sección $\mathrm{N}_{X'/X}(s')$ de $\sh{O}_X$ sobre $U$.
La fórmula \sref{IV.21.5.8.1} resultará, por tanto, de las definiciones si se prueba que $g'^*(D')$ y $g^*(D)$ están definidos, cualesquiera que sean los divisores $D'$ sobre $X'$ y $D$ sobre $X$.
En lo que respecta a $D$, esto resulta de las hipótesis formuladas y de \sref{IV.21.4.5}.
En lo que respecta a $D'$, en el caso~(i) $g'$ es plano, mientras que en los casos~(ii) y~(iii) se satisfacen las hipótesis correspondientes de \sref{IV.21.4.5}\footnote{El texto francés impreso omite esta cláusula y pasa directamente del caso~(i) a la conclusión para todos los casos.}; por tanto, en todos los casos $g'^*(D')$ está definido.
\end{proof}


% Mechanical subsection components; prose is translated in the component files.
\subsection{Ciclo 1-codimensional asociado a un divisor}
\label{subsection:IV.21.6}
\label{IV.21.6}

\begin{env}[21.6.1]
\label{IV.21.6.1}
Sea $X$ un preesquema localmente noetheriano, y sea $\mathfrak{I}(X)$ el conjunto de los \emph{conjuntos cerrados irreducibles} de $X$ (que está en correspondencia biunívoca con $X$ por la aplicación $x\mapsto\overline{\{x\}}$).
En el grupo producto $\bb{Z}^X$, consideremos el subgrupo $\mathfrak{Z}(X)$ de los elementos $(n_x)_{x\in X}$ tales que el conjunto de los $\overline{\{x\}}\in\mathfrak{I}(X)$ para los cuales $n_x\neq0$ (o, lo que es lo mismo, el conjunto de los $x\in X$ tales que $n_x\neq0$) sea localmente finito.
Está claro que $\mathfrak{Z}(X)$ es un subgrupo de $\bb{Z}^X$ que contiene el grupo suma directa $\bb{Z}^{(X)}$ (grupo libre que tiene por base $\mathfrak{I}(X)$), y es igual a él cuando $X$ es noetheriano.
Los elementos de $\mathfrak{Z}(X)$ se llaman \emph{ciclos} sobre $X$, y los de $\mathfrak{I}(X)$, \emph{ciclos primos} (en general no forman una base de $\mathfrak{Z}(X)$ cuando $X$ no es noetheriano).
Siempre se considera $\mathfrak{Z}(X)$ ordenado por el orden inducido sobre este subgrupo por el orden producto de $\bb{Z}^X$, y se denotará por $\mathfrak{Z}^+(X)$ el conjunto de los ciclos $\geq0$.

\oldpage[IV-4]{271}

Para todo ciclo $Z\in\mathfrak{Z}(X)$, igual a $(n_x)_{x\in X}$, se escribirá, por abuso de lenguaje,
\[
  Z=\sum_{x\in X}n_x\cdot\overline{\{x\}}.
\]
Se llama \emph{multiplicidad de $Z$ en el punto $x$} y se denota por $\operatorname{mult}_x(Z)$ el entero $n_x$ (positivo o negativo).
Decir que $Z\geq0$ significa que $\operatorname{mult}_x(Z)\geq0$ para todo $x\in X$.
Se llama \emph{soporte de $Z$} y se denota por $\Supp(Z)$ la unión de los $\overline{\{x\}}$ tales que $\operatorname{mult}_x(Z)\neq0$;
por definición de $\mathfrak{Z}(X)$, es un conjunto cerrado de $X$, como unión de una familia localmente finita de conjuntos cerrados.
Se llama \emph{dimensión} (resp.\ \emph{codimensión}) de $Z$ y se denota por $\dim(Z)$ (resp.\ $\codim(Z)$) la dimensión (resp.\ codimensión en $X$) de $\Supp(Z)$.
\end{env}

\begin{env}[21.6.2]
\label{IV.21.6.2}
Se dice que un conjunto cerrado $Y$ de $X$ es \emph{puramente de codimensión $p$} (en $X$) si \emph{cada una} de sus componentes irreducibles es de codimensión $p$ en $X$.
Se dice que un ciclo es \emph{puramente de codimensión $p$}, o (por abuso de lenguaje) es un \emph{ciclo $p$-codimensional}, si su soporte es puramente de codimensión $p$.
Se denota por $X^{(p)}$ el conjunto de los $x\in X$ tales que $\codim(\overline{\{x\}},X)=p$, o, lo que es lo mismo \sref{IV.5.1.2}, $\dim(\sh{O}_{X,x})=p$:
los ciclos puramente de codimensión $p$ forman un subgrupo $\mathfrak{Z}^p(X)$ de $\mathfrak{Z}(X)$, isomorfo al subgrupo de $\bb{Z}^{X^{(p)}}$ formado por los $(n_x)_{x\in X^{(p)}}$ tales que el conjunto de los $\overline{\{x\}}$ (o el conjunto de los $x$) para los cuales $n_x\neq0$ sea localmente finito;
este subgrupo contiene el grupo libre $\bb{Z}^{(X^{(p)})}$ y es idéntico a él cuando $X$ es noetheriano.
Se denota por $\mathfrak{Z}^{p+}(X)$ el conjunto de los elementos $\geq0$ de $\mathfrak{Z}^p(X)$.
Está claro que el grupo ordenado $\mathfrak{Z}(X)$ contiene la suma directa de los subgrupos ordenados $\mathfrak{Z}^p(X)$, y es idéntico a esta suma directa cuando $X$ es noetheriano;
en este último caso, se considera $\mathfrak{Z}(X)$ \emph{graduado} por los $\mathfrak{Z}^p(X)$ para $p\in\bb{N}$.
\end{env}

\begin{env}[21.6.3]
\label{IV.21.6.3}
Sean $Z=\sum_{x\in X}n_x\cdot\overline{\{x\}}$ un ciclo sobre $X$ y $U$ un abierto de $X$;
se llama \emph{restricción} de $Z$ a $U$ y se denota por $Z|U$ el ciclo $\sum_{x\in U}n_x\cdot(U\cap\overline{\{x\}})$ sobre $U$;
se tiene $\Supp(Z|U)=\Supp(Z)\cap U$.
Está claro que $Z\mapsto Z|U$ es un homomorfismo de grupos ordenados de $\mathfrak{Z}(X)$ en $\mathfrak{Z}(U)$ (resp.\ de $\mathfrak{Z}^p(X)$ en $\mathfrak{Z}^p(U)$), y que si $V$ es un abierto contenido en $U$, se tiene $Z|V=(Z|U)|V$;
la aplicación $U\mapsto\mathfrak{Z}(U)$ (resp.\ $U\mapsto\mathfrak{Z}^p(U)$) es, por tanto, un \emph{prehaz} sobre $X$ de grupos ordenados conmutativos.
De hecho, este prehaz es un \emph{haz}, suma directa de los haces $(i_x)_*(\bb{Z}_{\overline{\{x\}}})$, donde $x$ recorre $X$ (resp.\ $X^{(p)}$), y para cada $x\in X$, $i_x$ es la inyección canónica $\overline{\{x\}}\to X$ y $\bb{Z}_{\overline{\{x\}}}$ es el haz simple asociado al haz constante $\bb{Z}$ sobre el espacio $\overline{\{x\}}$:
esto se deduce inmediatamente de la descripción de las secciones de una suma directa de haces de grupos conmutativos (G, II, 2.7).
Este haz se denota por $\mathscr{Z}_X$ (resp.\ $\mathscr{Z}_X^p$).
Se denota por $\mathscr{Z}_X^+$ (resp.\ $\mathscr{Z}_X^{p+}$) el subhaz
\oldpage[IV-4]{272}
de monoides $U\mapsto\mathfrak{Z}^+(U)$ (resp.\ $U\mapsto\mathfrak{Z}^{p+}(U)$) de $\mathscr{Z}_X$ (resp.\ $\mathscr{Z}_X^p$).
El haz $\mathscr{Z}_X$ es evidentemente \emph{suma directa} de los haces $\mathscr{Z}_X^p$.

Señalemos finalmente que los haces $\mathscr{Z}_X^p$ (y, por tanto, también $\mathscr{Z}_X$) son \emph{flácidos}.
Supongamos, en efecto, dada una sección $Z$ de $\mathscr{Z}_X^p$ sobre un abierto $U$, de modo que el conjunto $S$ de los $x\in U$ tales que $\operatorname{mult}_x(Z)\neq0$ sea localmente finito \emph{en} $U$;
este conjunto también es localmente finito \emph{en} $X$, pues para todo $z\in X$ y todo entorno abierto noetheriano $V$ de $z$, $U\cap V$ es noetheriano y, por tanto, solo contiene un número finito de puntos de $S$.
Se define entonces una sección $Z'$ de $\mathscr{Z}_X^p$ sobre $X$ poniendo $Z'=\sum_{x\in U}\operatorname{mult}_x(Z)\cdot\overline{\{x\}}$ y, puesto que la adherencia de $x$ en $U$ es $\overline{\{x\}}\cap U$, se tiene efectivamente $Z'|U=Z$, de donde resulta nuestra afirmación.
\end{env}

\begin{env}[21.6.4]
\label{IV.21.6.4}
Nos proponemos definir un homomorfismo canónico
\[
\label{IV.21.6.4.1}
  c:\sh{D}iv_X\to\mathscr{Z}_X^1
  \tag{21.6.4.1}
\]
de haces de grupos conmutativos.
Evidentemente, basta definir un homomorfismo de $\sh{M}_X^*$ en $\mathscr{Z}_X^1$ que se anule en $\sh{O}_X^*$ \sref{IV.21.1.2};
ahora bien, por definición, $\sh{M}_X^*$ es el \emph{simetrizado} del haz de monoides $\sh{O}_X\cap\sh{M}_X^*$, y un homomorfismo $\sh{M}_X^*\to\mathscr{Z}_X^1$ queda determinado de manera única por su restricción $\sh{S}(\sh{O}_X)\to\mathscr{Z}_X^1$, que es un homomorfismo de haces de \emph{monoides} sometido a la única condición de ser nulo en $\sh{O}_X^*$;
equivalentemente, para definir $c$, basta definir un homomorfismo de \emph{haces de monoides}
\[
\label{IV.21.6.4.2}
  c:\sh{D}iv_X^+\to\mathscr{Z}_X^{1+}.
  \tag{21.6.4.2}
\]
\end{env}

\begin{env}[21.6.5]
\label{IV.21.6.5}
Ahora bien, se ha visto que a todo divisor positivo $D$ sobre $X$ se hace corresponder el subpreesquema cerrado $Y(D)$ de $X$, definido por el ideal $\sh{J}_X(D)\subset\sh{O}_X$, y que está regularmente inmerso y es de codimensión $1$.
En cada uno de los puntos maximales $x$ de $Y(D)$ se tiene, por tanto, (\sref{IV.19.1.4} y \sref{IV.5.1.2}) $\codim(\overline{\{x\}},X)=\dim(\sh{O}_{X,x})=1$, es decir, $x\in X^{(1)}$;
además, el conjunto de estos puntos maximales es localmente finito en $X$ \sref{IV.3.1.6}, y $\sh{O}_{Y(D),x}$ es un anillo \emph{artiniano}.
En todo punto $x\in X^{(1)}$ que no sea punto maximal de $Y(D)$, se tiene necesariamente $x\notin Y(D)$ y, por tanto, $\sh{O}_{Y(D),x}=0$.
Pongamos
\[
\label{IV.21.6.5.1}
  \operatorname{cyc}(D)=\sum_{x\in X^{(1)}}\operatorname{long}(\sh{O}_{Y(D),x})\cdot\overline{\{x\}}
  \tag{21.6.5.1}
\]
que es, por tanto, un elemento de $\mathfrak{Z}^1(X)$.
\end{env}

\begin{proposition}[21.6.6]
\label{IV.21.6.6}
La aplicación $D\mapsto\operatorname{cyc}(D)$ es un homomorfismo del monoide $\operatorname{Div}^+(X)$ en $\mathfrak{Z}^{1+}(X)$.
\end{proposition}

\begin{proof}
Todo se reduce a ver que para dos divisores positivos $D$, $D'$, se tiene
\[
  \operatorname{cyc}(D+D')=\operatorname{cyc}(D)+\operatorname{cyc}(D').
\]
Ahora bien, se tiene \sref{IV.21.2.5} $\sh{J}_X(D+D')=\sh{J}_X(D)\sh{J}_X(D')$;
todo se reduce a ver que si $x\in X^{(1)}$, si se pone $A=\sh{O}_{X,x}$, y si $t$ y $t'$ son dos elementos regulares de $A$, se tiene $\operatorname{long}(A/tt'A)=\operatorname{long}(A/tA)+\operatorname{long}(A/t'A)$;
pero, puesto que $t$ es regular, $tA/tt'A$ es isomorfo a $A/t'A$, de donde resulta inmediatamente la proposición.
\end{proof}

Se deduce inmediatamente de las definiciones que, para todo abierto $U\subset X$, se tiene
\[
  Y(D|U)=Y(D)\cap U,
\]
\oldpage[IV-4]{273}
por tanto, $\operatorname{cyc}(D|U)=\operatorname{cyc}(D)|U$, y los homomorfismos $\operatorname{cyc}:\operatorname{Div}^+(U)\to\mathfrak{Z}^1(U)$ definen así un homomorfismo de haces de monoides \sref{IV.21.6.4.2}, de donde se obtiene el homomorfismo buscado \sref{IV.21.6.4.1} de haces de grupos.

Se tiene $\Supp(\operatorname{cyc}(D))\subset\Supp(D)$ para todo divisor $D$ y
\begin{env}[21.6.6.2]
\label{IV.21.6.6.2}
\[
  \Supp(\operatorname{cyc}(D))=\Supp(D)\quad\text{cuando}\quad D\geq0;
\]
\end{env}
en efecto, ya se ha visto la segunda relación \sref{IV.21.2.12};
cuando $D$ es arbitrario, la relación $D_x=0$ implica la existencia de un entorno abierto $U$ de $x$ tal que $D|U=0$, de donde $\operatorname{cyc}(D)|U=0$, lo que demuestra nuestra afirmación.

\begin{env}[21.6.7]
\label{IV.21.6.7}
El homomorfismo \sref{IV.21.6.4.1} da, al pasar a los grupos de secciones sobre $X$, un \emph{homomorfismo creciente de grupos ordenados} $\operatorname{Div}(X)\to\mathfrak{Z}^1(X)$, que se seguirá denotando por $D\mapsto\operatorname{cyc}(D)$;
el número $\operatorname{mult}_x(\operatorname{cyc}(D))$ para $x\in X^{(1)}$ se denota también por $\operatorname{mult}_x(D)$ o $\operatorname{mult}_{\overline{\{x\}}}(D)$ y se llama \emph{multiplicidad del divisor $D$ en el punto $x$}, o \emph{multiplicidad del ciclo primo $\overline{\{x\}}$ en $D$};
es un entero positivo o negativo, igual, como se ha visto, a $\operatorname{long}(\sh{O}_{Y(D),x})$ cuando $D$ es un divisor \emph{positivo};
se tiene, por tanto, por definición
\[
\label{IV.21.6.7.1}
  \operatorname{cyc}(D)=\sum_{x\in X^{(1)}}\operatorname{mult}_x(D)\cdot\overline{\{x\}},
  \tag{21.6.7.1}
\]
y, en virtud de que $D\mapsto\operatorname{cyc}(D)$ es un homomorfismo, se tiene
\[
\label{IV.21.6.7.2}
  \operatorname{mult}_x(-D)=-\operatorname{mult}_x(D),\qquad
  \operatorname{mult}_x(D+D')=\operatorname{mult}_x(D)+\operatorname{mult}_x(D')
  \tag{21.6.7.2}
\]
para dos divisores cualesquiera $D$, $D'$.

Para toda función meromorfa regular $f$ sobre $X$, se pone en particular, para $x\in X^{(1)}$,
\[
\label{IV.21.6.7.3}
  o_x(f)=\operatorname{mult}_x(\operatorname{div}(f))
  \tag{21.6.7.3}
\]
y se dice que $o_x(f)$ (entero positivo o negativo) es \emph{el orden de $f$ en el punto} $x\in X^{(1)}$.
Si $f_x\in\sh{O}_{X,x}$ (elemento regular de este anillo local por hipótesis), se tiene, por tanto,
\[
\label{IV.21.6.7.4}
  o_x(f)=\operatorname{long}(\sh{O}_{X,x}/f_x\sh{O}_{X,x});
  \tag{21.6.7.4}
\]
para dos funciones meromorfas regulares $f$, $f'$ sobre $X$, se tiene
\[
\label{IV.21.6.7.5}
  o_x(ff')=o_x(f)+o_x(f'),\qquad o_x(f^{-1})=-o_x(f)
  \tag{21.6.7.5}
\]
para todo $x\in X^{(1)}$.
Los ciclos $1$-codimensionales
\[
  Z^+(f)=\sum_{x\in X^{(1)}}(o_x(f))^+\cdot\overline{\{x\}},\qquad
  Z^-(f)=\sum_{x\in X^{(1)}}(o_x(f))^-\cdot\overline{\{x\}}
\]
se llaman respectivamente el \emph{ciclo de los ceros} y el \emph{ciclo de los polos} (o \emph{ciclo polar}) de $f$, y se tiene evidentemente $\operatorname{cyc}(\operatorname{div}(f))=Z^+(f)-Z^-(f)$.
Los ciclos $1$-codimensionales de la forma $\operatorname{cyc}(\operatorname{div}(f))$ se llaman \emph{principales} (o \emph{linealmente equivalentes a cero}) y forman un subgrupo $\mathfrak{Z}^1_\mathrm{princ}(X)$ de $\mathfrak{Z}^1(X)$.
Las secciones sobre $X$ de la imagen $c(\sh{D}iv_X)\subset\mathscr{Z}_X^1$ se
\oldpage[IV-4]{274}
llaman \emph{ciclos localmente principales};
tal ciclo $Z$ se caracteriza, por tanto, por el hecho de que, para todo $y\in X$, existe un entorno abierto $U$ de $y$ en $X$ y una función meromorfa regular $f$ sobre $U$ tal que $Z|U=\sum_{x\in U\cap X^{(1)}}o_x(f)(\overline{\{x\}}\cap U)$.
Si $Z=\sum_{x\in X^{(1)}}n_x\cdot\overline{\{x\}}$, es equivalente decir que, para todo $y\in X$, si se pone $T_y=\Spec(\sh{O}_{X,y})$ y $Z_y=\sum_{x\in X^{(1)}\cap T_y}n_x(\overline{\{x\}}\cap T_y)$, $Z_y$ es \emph{principal}, es decir, existe una función meromorfa regular $g$ sobre $T_y$ tal que $Z_y=\operatorname{cyc}(\operatorname{div}(g))$.
Esto se deduce inmediatamente, en efecto, de la definición anterior y de \sref{IV.20.3.7}, que establece una correspondencia biunívoca entre las funciones meromorfas regulares sobre $T_y$ y los gérmenes de funciones meromorfas regulares sobre los entornos abiertos de $y$ en $X$ cuando $X$ es localmente noetheriano.
El hecho de que $Z_y$ sea principal se expresa también diciendo que $Z$ es \emph{principal en el punto $y$}.
El conjunto de los puntos en los que $Z$ es principal es evidentemente \emph{abierto}, en virtud de lo anterior.

Del homomorfismo canónico $\operatorname{cyc}:\operatorname{Div}(X)\to\mathfrak{Z}^1(X)$ se deduce un homomorfismo $\operatorname{Div}(X)/\operatorname{Div.princ}(X)\to\mathfrak{Z}^1(X)/\mathfrak{Z}^1_\mathrm{princ}(X)$, en virtud de la definición de $\mathfrak{Z}^1_\mathrm{princ}(X)$.
Se dice que el grupo $\mathfrak{Z}^1(X)/\mathfrak{Z}^1_\mathrm{princ}(X)$ es el \emph{grupo de las clases de ciclos $1$-codimensionales} de $X$ y se denota por $\mathfrak{Cl}(X)$.
Dos elementos de $\mathfrak{Z}^1(X)$ que tienen la misma imagen en $\mathfrak{Cl}(X)$ se llaman \emph{linealmente equivalentes}.
\end{env}

\begin{env}[21.6.8]
\label{IV.21.6.8}
Consideremos en particular el caso en que $X=\Spec(A)$, donde $A$ es un anillo \emph{noetheriano, íntegro e íntegramente cerrado}.
Entonces $X^{(1)}$ es el conjunto de los \emph{ideales primos de altura $1$} de $A$, y $\mathfrak{Z}^1(X)$ se identifica, por tanto, con el \emph{grupo de los divisores} (en el sentido de N.~Bourbaki) del anillo de Krull $A$ (Bourbaki, \emph{Álg.\ conmut.}, cap.~VII, \textsection~1, n\textsuperscript{o}~3, cor.\ del teorema~2 y n\textsuperscript{o}~6, teorema~3).

Como, por otra parte, las funciones meromorfas regulares sobre $X$ se identifican entonces con los elementos no nulos del cuerpo de fracciones $K$ de $A$, la aplicación $f\mapsto\operatorname{cyc}(\operatorname{div}(f))$ de $M(X)$ en $\mathfrak{Z}^1(X)$ se identifica con la aplicación denotada por $f\mapsto\operatorname{div}(f)$ en Bourbaki (\emph{loc.\ cit.}, \textsection~1, n\textsuperscript{o}~1);
$\mathfrak{Z}^1_\mathrm{princ}(X)$ se identifica, por tanto, con el \emph{grupo de los divisores principales} de $A$ en el sentido de Bourbaki, y $\mathfrak{Cl}(X)$ con el \emph{grupo de las clases de divisores} de $A$ en el sentido de Bourbaki (\emph{loc.\ cit.}, \textsection~1, n\textsuperscript{o}~2 y n\textsuperscript{o}~10).
\end{env}

\begin{theorem}[21.6.9]
\label{IV.21.6.9}
Sea $X$ un preesquema localmente noetheriano normal.
\begin{enumerate}[label=(\roman*)]
  \item El homomorfismo canónico $\operatorname{cyc}:\operatorname{Div}(X)\to\mathfrak{Z}^1(X)$ es inyectivo y su imagen está formada por los ciclos localmente principales.
  \item Las condiciones siguientes son equivalentes:
  \begin{enumerate}[label=\emph{\alph*)}]
    \item El homomorfismo $\operatorname{cyc}:\operatorname{Div}(X)\to\mathfrak{Z}^1(X)$ es biyectivo.
    \item Todo ciclo $1$-codimensional es localmente principal.
    \item Para todo $x\in X$, el anillo local $\sh{O}_{X,x}$ es factorial.
  \end{enumerate}
  (Se dice entonces que $X$ es un preesquema \emph{localmente factorial}.)
\end{enumerate}
\end{theorem}

\begin{proof}
\begin{enumerate}[label=(\roman*)]
  \item Basta demostrar que
\[
\label{IV.21.6.9.1}
  \operatorname{cyc}^{-1}(\mathfrak{Z}^{1+}(X))=\operatorname{Div}^+(X),
  \tag{21.6.9.1}
\]
puesto que se tiene $\operatorname{Div}^+(X)\cap(-\operatorname{Div}^+(X))=0$ y $\mathfrak{Z}^{1+}(X)\cap(-\mathfrak{Z}^{1+}(X))=0$.
Se reduce, por tanto, a demostrar que si $D$ es un divisor sobre $X$ tal que $\operatorname{mult}_x(D)\geq0$ para todo $x\in X^{(1)}$, entonces $D\geq0$.
Ahora bien, para todo $x\notin X^{(1)}$, el anillo local $\sh{O}_{X,x}$ es íntegro e íntegramente cerrado y
\oldpage[IV-4]{275}
de dimensión $0$ o $\geq2$, por lo que se tiene $\operatorname{prof}(\sh{O}_{X,x})=0$ o $\operatorname{prof}(\sh{O}_{X,x})\geq2$ \sref[0]{0.16.5.1}.
Por otra parte, en los puntos $x\in X^{(1)}$, el anillo $\sh{O}_{X,x}$ es un anillo de valoración discreta (\sref[II]{II.7.1.6}) y, por tanto, $\operatorname{prof}(\sh{O}_{X,x})=1$ \sref[0]{0.16.5.1};
los únicos puntos de $X$ tales que $\operatorname{prof}(\sh{O}_{X,x})=1$ son, por consiguiente, los puntos de $X^{(1)}$, y la conclusión se deduce entonces de \sref{IV.21.1.8}.

  \item La equivalencia de \emph{a)} y \emph{b)} es clara en virtud de (i).
Según la caracterización de los ciclos localmente principales dada en \sref{IV.21.6.7}, y la relación entre ciclos $1$-codimensionales sobre $\Spec(A)$ y divisores (en el sentido de Bourbaki) del anillo $A$ cuando $A$ es un anillo noetheriano íntegramente cerrado \sref{IV.21.6.8}, la condición \emph{b)} equivale a decir que, para todo $x\in X$, todo divisor del anillo $\sh{O}_{X,x}$ es principal, es decir, que el anillo $\sh{O}_{X,x}$ es factorial (Bourbaki, \emph{Álg.\ conmut.}, cap.~VII, \textsection~3, n\textsuperscript{o}~1), de donde resulta la equivalencia de \emph{b)} y \emph{c)}.
\end{enumerate}
\end{proof}

\begin{env}[21.6.9.2]
\label{IV.21.6.9.2}
Cuando $A$ es un anillo noetheriano \emph{factorial}, está claro que $X=\Spec(A)$ es localmente factorial (Bourbaki, \emph{Álg.\ conmut.}, cap.~VII, \textsection~3, n\textsuperscript{o}~4, prop.~3).
Si, en este caso, se escribe un elemento $f\neq0$ del cuerpo de fracciones $K$ de $A$ en la forma $r/s$, donde $r$ y $s$ son dos elementos \emph{primos entre sí} de $A$, cuyos divisores están bien determinados (\emph{loc.\ cit.}, \textsection~3, n\textsuperscript{o}~3), estos divisores se identifican respectivamente con el \emph{divisor de los ceros} y el \emph{divisor de los polos} de $f$ \sref{IV.21.6.7}.
\end{env}

\begin{corollary}[21.6.10]
\label{IV.21.6.10}
Sea $X$ un preesquema localmente noetheriano normal.
\begin{enumerate}[label=(\roman*)]
  \item Existe un homomorfismo canónico inyectivo
\[
\label{IV.21.6.10.1}
  \operatorname{Pic}(X)\to\mathfrak{Cl}(X).
  \tag{21.6.10.1}
\]

  \item Si $X$ es localmente factorial, el homomorfismo \sref{IV.21.6.10.1} es biyectivo, y recíprocamente.
\end{enumerate}
\end{corollary}

\begin{proof}
Se ha visto \sref{IV.21.6.7} que la imagen de $\operatorname{Div.princ}(X)$ por el homomorfismo $\operatorname{cyc}:\operatorname{Div}(X)\to\mathfrak{Z}^1(X)$ es $\mathfrak{Z}^1_\mathrm{princ}(X)$;
se deduce, por tanto, de \sref{IV.21.6.9} que el homomorfismo $\operatorname{Div}(X)/\operatorname{Div.princ}(X)\to\mathfrak{Z}^1(X)/\mathfrak{Z}^1_\mathrm{princ}(X)=\mathfrak{Cl}(X)$ deducido de $\operatorname{cyc}$ es inyectivo, y que es biyectivo si y solo si $X$ es localmente factorial.
La conclusión se deduce entonces del hecho de que, cuando $X$ es localmente noetheriano y reducido, el homomorfismo canónico $\operatorname{Div}(X)/\operatorname{Div.princ}(X)\to\operatorname{Pic}(X)$ \sref{IV.21.3.3.1} es biyectivo \sref{IV.21.3.4}, (ii).
\end{proof}

\begin{corollary}[21.6.11]
\label{IV.21.6.11}
Sea $X$ un preesquema localmente noetheriano y localmente factorial.
Entonces el haz $\sh{D}iv_X$ es flácido, y para todo abierto $U$ de $X$, el homomorfismo canónico $\operatorname{Pic}(X)\to\operatorname{Pic}(U)$ es sobreyectivo.
\end{corollary}

\begin{proof}
Teniendo en cuenta \sref{IV.21.6.9}, (ii), la primera afirmación equivale a decir que el haz $\mathscr{Z}_X^1$ es flácido, lo que se ha visto más arriba \sref{IV.21.6.3}.
Para todo abierto $U$ de $X$, el homomorfismo canónico $\mathfrak{Z}^1(X)\to\mathfrak{Z}^1(U)$ es, por tanto, sobreyectivo;
como, en virtud de \sref{IV.21.6.10}, el homomorfismo $\operatorname{Pic}(X)\to\operatorname{Pic}(U)$ se identifica canónicamente con $\mathfrak{Z}^1(X)/\mathfrak{Z}^1_\mathrm{princ}(X)\to\mathfrak{Z}^1(U)/\mathfrak{Z}^1_\mathrm{princ}(U)$, también es sobreyectivo.
\end{proof}

\begin{proposition}[21.6.12]
\label{IV.21.6.12}
Sea $X$ un preesquema noetheriano reducido.
Sea $(U_\lambda)_{\lambda\in L}$ una familia filtrante decreciente de abiertos de $X$ que satisface las condiciones siguientes:
\begin{enumerate}
  \item[\rm 1\textsuperscript{o}] Si $Y_\lambda=X-U_\lambda$, se tiene $\codim(Y_\lambda,X)\geq2$ para todo $\lambda\in L$.
  \item[\rm 2\textsuperscript{o}] Para todo $x\in\bigcap_{\lambda\in L}U_\lambda$, el anillo $\sh{O}_{X,x}$ es factorial.
\end{enumerate}
Entonces existen isomorfismos canónicos
\[
\label{IV.21.6.12.1}
  \varinjlim\operatorname{Div}(U_\lambda)\simto\mathfrak{Z}^1(X),\qquad
  \varinjlim\operatorname{Pic}(U_\lambda)\simto\mathfrak{Cl}(X)
  \tag{21.6.12.1}
\]
tales que, para todo abierto $V$ de $X$, los diagramas
\[
\label{IV.21.6.12.2}
  \xymatrix{
    \varinjlim\operatorname{Div}(U_\lambda) \ar[r]^-\sim \ar[d] & \mathfrak{Z}^1(X) \ar[d] & & \varinjlim\operatorname{Pic}(U_\lambda) \ar[r]^-\sim \ar[d] & \mathfrak{Cl}(X) \ar[d]\\
    \varinjlim\operatorname{Div}(U_\lambda\cap V) \ar[r]^-\sim & \mathfrak{Z}^1(V) & & \varinjlim\operatorname{Pic}(U_\lambda\cap V) \ar[r]^-\sim & \mathfrak{Cl}(V)
  }
  \tag{21.6.12.2}
\]
sean conmutativos.
\end{proposition}

\begin{proof}
La hipótesis 1\textsuperscript{o} sobre $U_\lambda$ implica que, para todo $\lambda\in L$, se tiene $X^{(1)}\subset U_\lambda$ \sref{IV.5.1.3.1}; por tanto, el homomorfismo de restricción $\mathfrak{Z}^1(X)\to\mathfrak{Z}^1(U_\lambda)$ es biyectivo y, en consecuencia, se tiene un isomorfismo canónico $\varinjlim\mathfrak{Z}^1(U_\lambda)\simto\mathfrak{Z}^1(X)$.
Los homomorfismos canónicos $\operatorname{cyc}:\operatorname{Div}(U_\lambda)\to\mathfrak{Z}^1(U_\lambda)$ definen, por tanto, pasando al límite inductivo, el primero de los homomorfismos canónicos \sref{IV.21.6.12.1}, y el segundo se deduce de él pasando a los cocientes.
Además, de la condición 1\textsuperscript{o} se deduce que los $U_\lambda$ son densos en $X$ y, por tanto, esquemáticamente densos, puesto que $X$ es reducido \sref{IV.11.10.4}; en consecuencia, toda función meromorfa sobre $X$ queda enteramente determinada por su restricción a cada $U_\lambda$;
de esto se deduce inmediatamente que, en el isomorfismo $\mathfrak{Z}^1(X)\simto\mathfrak{Z}^1(U_\lambda)$, la imagen de $\mathfrak{Z}^1_\mathrm{princ}(X)$ es $\mathfrak{Z}^1_\mathrm{princ}(U_\lambda)$, por lo que se tiene también un isomorfismo canónico $\varinjlim\mathfrak{Cl}(U_\lambda)\simto\mathfrak{Cl}(X)$.
El segundo de los homomorfismos canónicos \sref{IV.21.6.12.1} será, por tanto, un isomorfismo si lo es el primero, y queda entonces por demostrar este último punto, siendo trivial la conmutatividad de los diagramas \sref{IV.21.6.12.2}.

\oldpage[IV-4]{276}
Demostremos primero que el homomorfismo $\varinjlim\operatorname{Div}(U_\lambda)\to\mathfrak{Z}^1(X)$ es \emph{inyectivo}.
Pongamos $T=\bigcap_\lambda U_\lambda$, y señalemos que los $U_\lambda$ forman un \emph{sistema fundamental de entornos} de $T$.
En efecto, para todo abierto $V\supset T$, $X-V$ es un subespacio cerrado de $X$ y, por tanto, un espacio noetheriano en el que todo conjunto cerrado irreducible admite un punto genérico;
como los conjuntos $(X-V)\cap(X-U_\lambda)$ son cerrados en $X-V$, forman una familia filtrante creciente y tienen por unión $X-V$, nuestra afirmación se deduce de \sref[0]{0.9.2.4}.
Dicho esto, se trata de ver que si $D\in\operatorname{Div}(U_\lambda)$ es tal que $\operatorname{cyc}(D)=0$ en $\mathfrak{Z}^1(U_\lambda)$, entonces existe $\mu\geq\lambda$ tal que $D|U_\mu=0$.
En virtud de lo anterior, bastará ver que para todo $x\in T$, se tiene $D_x=0$;
en efecto, habrá entonces un entorno $W(x)$ de $x$ en $X$ tal que $D|W(x)=0$, y la unión de los $W(x)$ para $x\in T$ contiene algún $U_\mu$.
Teniendo en cuenta \sref{IV.21.4.6}, se reduce, por tanto, al caso en que $X=\Spec(\sh{O}_{X,x})$ con $x\in T$;
pero, por hipótesis, $\sh{O}_{X,x}$ es factorial, luego íntegro e íntegramente cerrado, y $\Spec(\sh{O}_{X,x})$ es entonces normal;
por consiguiente, la conclusión se deduce de \sref{IV.21.6.9}, (i).

\oldpage[IV-4]{277}
Para demostrar que el homomorfismo $\varinjlim\operatorname{Div}(U_\lambda)\to\mathfrak{Z}^1(X)$ es biyectivo, bastará demostrar del mismo modo que, para todo $Z\in\mathfrak{Z}^1(X)$ y todo $x\in T$, existe un entorno $W(x)$ de $x$ en $X$ y un divisor $D^{(x)}$ sobre $W(x)$ tal que $\operatorname{cyc}(D^{(x)})=Z|W(x)$.
En efecto, se podrá entonces recubrir el conjunto cuasicompacto $T$ por un número finito de $W(x_i)$ y, en virtud de la primera parte de la demostración, puesto que el número de pares $(i,j)$ es finito, existirá un índice $\lambda$ tal que, en cada uno de los $W(x_i)\cap W(x_j)\cap U_\lambda$, las restricciones de $D^{(x_i)}$ y $D^{(x_j)}$ coincidan;
como se puede suponer además que $U_\lambda$ está contenido en la unión de los $W(x_i)$, se ve que las restricciones $D^{(x_i)}|(W(x_i)\cap U_\lambda)$ son las restricciones de un mismo divisor $D\in\operatorname{Div}(U_\lambda)$ que será tal que $\operatorname{cyc}(D)=Z|U_\lambda$.
Se reduce de nuevo al caso en que $X=\Spec(\sh{O}_{X,x})$ con $x\in T$, pero como $\sh{O}_{X,x}$ es factorial, también lo son sus localizaciones (Bourbaki, \emph{Álg.\ conmut.}, cap.~VII, \textsection~3, n\textsuperscript{o}~4, prop.~3), y basta aplicar \sref{IV.21.6.9}, (ii).
\end{proof}

\begin{corollary}[21.6.13]
\label{IV.21.6.13}
Sea $A$ un anillo local noetheriano, íntegro e íntegramente cerrado, de dimensión $\geq2$.
Pongamos $X=\Spec(A)$, y sean $a$ el punto cerrado de $X$ y $U=X-\{a\}$.
Para que $A$ sea factorial, es necesario y suficiente que $U$ sea localmente factorial y que $\operatorname{Pic}(U)=0$.
\end{corollary}

\begin{proof}
En efecto, decir que $A$ es factorial equivale a decir que $\mathfrak{Cl}(X)=0$ (Bourbaki, \emph{Álg.\ conmut.}, cap.~VII, \textsection~1, n\textsuperscript{o}~4, cor.\ del teorema~2 y \textsection~3, n\textsuperscript{o}~2, teorema~1);
basta entonces utilizar la existencia del segundo isomorfismo \sref{IV.21.6.12.1}, tomando la familia $(U_\lambda)$ restringida al único abierto $U$.
\end{proof}

\begin{corollary}[21.6.14]
\label{IV.21.6.14}
Sean $A$ un anillo local noetheriano de dimensión $\geq2$, $X=\Spec(A)$, $a$ el punto cerrado de $X$ y $U=X-\{a\}$.
Las condiciones siguientes son equivalentes:
\begin{enumerate}[label=\emph{\alph*)}]
  \item $A$ es factorial.
  \item $\operatorname{Pic}(U)=0$ y $\operatorname{prof}(A)\geq2$ (condiciones que expresaremos más adelante \sref{IV.21.13} diciendo que el anillo $A$ es \emph{parafactorial}), y $U$ es localmente factorial.
\end{enumerate}
\end{corollary}

\begin{proof}
Está claro que si $A$ es factorial, es normal y, por tanto, $\operatorname{prof}(A)\geq2$, puesto que $\dim(A)\geq2$ \sref[0]{0.16.5.1}, y \sref{IV.21.6.13} muestra que \emph{a)} implica \emph{b)}.
Recíprocamente, si se satisface \emph{b)}, basta ver que $A$ es íntegramente cerrado para concluir, mediante \sref{IV.21.6.13}, que \emph{b)} implica \emph{a)}.
En virtud del criterio de Serre \sref{IV.5.8.6}, basta verificar para $X$ las condiciones $(\mathrm{R}_1)$ y $(\mathrm{S}_2)$.
Ahora bien, al ser $U$ localmente factorial, satisface estas condiciones, y la hipótesis $\operatorname{prof}(A)\geq2$ implica que $X$ también las satisface.
\end{proof}

\subsection{Interpretación de los ciclos positivos de codimensión 1 en términos de subpreesquemas}
\label{subsection:IV.21.7}
\label{IV.21.7}

\begin{env}[21.7.1]
\label{IV.21.7.1}
Sean $X$ un preesquema localmente noetheriano y $C=\sum_{x\in X^{(1)}}n_x\cdot\overline{\{x\}}$ un ciclo positivo de codimensión $1$ (de modo que se tiene $n_x\geq0$ para todo $x\in X^{(1)}$, y $n_x=0$ salvo en un conjunto localmente finito de puntos).
Denotemos por $Y(C)$ el subpreesquema cerrado de $X$, \emph{imagen cerrada} (\sref[I]{I.9.5.3} y \sref{IV.9.5.1}) del preesquema \emph{suma} $Y'(C)=\coprod_{x\in X^{(1)}}\Spec(\sh{O}_{X,x}/\mathfrak{m}_x^{n_x})$ por el morfismo canónico, y por $\sh{J}_X(C)$ (o $\sh{J}(C)$) el ideal de $\sh{O}_X$ que define $Y(C)$.
\end{env}

\begin{proposition}[21.7.2]
\label{IV.21.7.2}
Sea $X$ un preesquema localmente noetheriano que satisface la condición $(\mathrm{R}_1)$ \sref{IV.5.8.2}.
Para que un subpreesquema cerrado $Y$ de $X$ sea de la forma $Y(C)$, donde $C$ es un ciclo positivo de codimensión $1$, es necesario y suficiente que satisfaga las dos condiciones siguientes:
\begin{enumerate}[label=(\roman*)]
  \item $Y$ es puramente de codimensión $1$.
\oldpage[IV-4]{278}
  \item $Y$ satisface la condición $(\mathrm{S}_1)$; en otros términos \sref{IV.5.7.5}, no tiene ningún ciclo primo asociado inmerso.
\end{enumerate}
Se tiene entonces
\[
\label{IV.21.7.2.1}
  \operatorname{mult}_x(C)=\operatorname{long}\sh{O}_{Y,x}.
  \tag{21.7.2.1}
\]
La aplicación $C\mapsto Y(C)$ es una biyección de $\mathfrak{Z}^{1+}(X)$ sobre el conjunto de los subpreesquemas cerrados de $X$ que satisfacen las condiciones (i) y (ii).
\end{proposition}

\begin{proof}
Las condiciones son necesarias (sin suponer que $X$ satisface $(\mathrm{R}_1)$).
En efecto, como la cuestión es local sobre $X$, se puede suponer que $X=\Spec(A)$, donde $A$ es un anillo noetheriano;
se tiene entonces $Y(C)=\Spec(A/\mathfrak{J})$, donde $\mathfrak{J}$ es, por definición \sref[I]{I.9.5.1}, el núcleo del homomorfismo canónico $A\to\bigoplus_i A_{\mathfrak{p}_i}/\mathfrak{p}_i^{n_i}A_{\mathfrak{p}_i}$, donde los $\mathfrak{p}_i$ son los ideales primos correspondientes a los puntos $x\in X^{(1)}$ para los que $n_x>0$, y donde se ha puesto $n_i=n_x$.
La imagen inversa $\mathfrak{q}_i$ de $\mathfrak{p}_i^{n_i}A_{\mathfrak{p}_i}$ en $A$ es un ideal $\mathfrak{p}_i$-primario y se tiene $\mathfrak{J}=\bigcap_i\mathfrak{q}_i$.
Además, como los $x\in X^{(1)}$ son tales que $\overline{\{x\}}$ es de codimensión $1$, ningún punto de $X^{(1)}$ puede ser adherente a otro punto de $X^{(1)}$.
Por tanto, los $\mathfrak{p}_i$ son los ideales primos minimales de $A/\mathfrak{J}$ y el conjunto de los $\mathfrak{p}_i$ es igual a $\operatorname{Ass}(A/\mathfrak{J})$, lo que prueba las condiciones (i) y (ii).

Las condiciones son suficientes.
Denotando por $\sh{J}$ el ideal de $\sh{O}_X$ que define $Y$, la hipótesis (ii) implica que $\operatorname{Ass}(\sh{O}_X/\sh{J})$ es idéntico al conjunto $F$ de los puntos maximales de $Y$, y la hipótesis (i) implica que $F$ está contenido en $X^{(1)}$;
por tanto \sref{IV.3.2.6}, $\sh{O}_X/\sh{J}$ se identifica con un sub-$\sh{O}_X$-módulo de $\bigoplus_{x\in F}\mathscr{G}(x)$, donde, para cada $x\in F$, $\mathscr{G}(x)$ es un $\sh{O}_X$-módulo irredundante tal que $\operatorname{Ass}(\mathscr{G}(x))=\{x\}$.
Ahora bien, como $X$ satisface $(\mathrm{R}_1)$, para cada $x\in F$, $\sh{O}_{X,x}$ es un anillo de valoración discreta y, por consiguiente, los ideales primarios no nulos de este anillo son las potencias $\mathfrak{m}_x^{n_x}$ del ideal maximal;
suponiendo de nuevo que $X=\Spec(A)$, se concluye que $\mathfrak{J}=\widetilde{\mathfrak{J}}$\footnote{La transcripción francesa imprime aquí el haz $\sh{J}$, pero tras suponer $X=\Spec(A)$ la igualdad es entre los ideales $\mathfrak{J}$ y $\widetilde{\mathfrak{J}}$ de $A$, como confirma el testigo inglés mantenido.}, donde $\widetilde{\mathfrak{J}}=\bigcap_i\mathfrak{q}_i$, siendo los $\mathfrak{q}_i$ las imágenes inversas en $A$ de los ideales $\mathfrak{p}_i^{n_i}A_{\mathfrak{p}_i}$, donde los $\mathfrak{p}_i$ corresponden a los puntos de $F$.
Se ve, pues, que $Y$ es efectivamente de la forma $Y(C)$.
\end{proof}

\begin{corollary}[21.7.3]
\label{IV.21.7.3}
Sea $X$ un preesquema localmente noetheriano.
Para todo divisor positivo $D$ sobre $X$, tal que $X$ satisfaga $(\mathrm{R}_1)$ en los puntos maximales del subpreesquema cerrado $Y(D)$ \sref{IV.21.2.12}, $Y(D)$ mayora al subpreesquema cerrado $Y(\operatorname{cyc}(D))$ \sref{IV.21.7.1}, y tiene el mismo espacio subyacente;
para que estos dos subpreesquemas sean iguales, es necesario y suficiente que $Y(D)$ satisfaga la condición $(\mathrm{S}_1)$, o también que, para todo $z\in Y(D)$ distinto de un punto maximal de $Y(D)$, se tenga $\operatorname{prof}(\sh{O}_{X,z})\geq2$ (condición que se satisface siempre cuando $X$ es \emph{normal}).
\end{corollary}

\begin{proof}
En efecto, como la cuestión es local, siempre se puede suponer que $\sh{J}_X(D)=t\sh{O}_X$, siendo $t$ una sección regular de $\sh{O}_X$ sobre $X$;
en todo punto maximal $x$ de $Y(D)$, que pertenece necesariamente a $X^{(1)}$, $\sh{O}_{X,x}$ es por hipótesis un anillo de valoración discreta, por lo que $t_x\sh{O}_{X,x}=\mathfrak{m}_x^{n_x}$, donde $n_x=\operatorname{mult}_x(D)$.
Se puede suponer que $X=\Spec(A)$, y entonces, si $\sh{J}_X(\operatorname{cyc}(D))=\widetilde{\mathfrak{J}}$, se deduce de lo anterior y de \sref{IV.21.7.1} que $\mathfrak{J}=tA$ y $\widetilde{\mathfrak{J}}$ tienen los mismos \emph{ideales primarios no inmersos}, y por tanto $\widetilde{\mathfrak{J}}\subset\mathfrak{J}$, puesto que $\widetilde{\mathfrak{J}}$ es la intersección
\oldpage[IV-4]{279}
de estos ideales primarios \sref{IV.21.7.2}.
Esto prueba que $Y(D)$ mayora a $Y(\operatorname{cyc}(D))$ y que estos dos subpreesquemas son iguales si y solo si $Y(D)$ no tiene ningún ciclo primo asociado inmerso (en otros términos, si satisface $(\mathrm{S}_1)$).
Como para todo $x\in Y(D)$ existe, por hipótesis, un entorno abierto $U$ de $x$ en $X$ y un elemento regular $t\in\Gamma(U,\sh{O}_X)$ tal que $\sh{O}_{Y(D)}|(U\cap Y(D))$ sea la restricción a $Y(D)\cap U$ de $\sh{O}_U/t\sh{O}_U$, decir que $Y(D)$ satisface $(\mathrm{S}_1)$ significa también que en todo punto $z\in Y(D)$ no maximal se tiene $\operatorname{prof}(\sh{O}_{X,z}/t_z\sh{O}_{X,z})\geq1$, es decir \sref[0]{0.16.4.6}, $\operatorname{prof}(\sh{O}_{X,z})\geq2$.
La afirmación relativa al caso en que $X$ es normal es entonces inmediata, puesto que $X$ satisface $(\mathrm{S}_2)$ y que en los puntos no maximales $z$ de $Y(D)$ se tiene $\dim(\sh{O}_{X,z})\geq2$ \sref[0]{0.16.3.4}.
\end{proof}

\begin{env}[21.7.3.1]
\label{IV.21.7.3.1}
Se ve, pues, que cuando $X$ es \emph{normal}, $\operatorname{Div}^+(X)$ se identifica canónicamente con el conjunto de los subpreesquemas cerrados $Y$ de $X$ que satisfacen las condiciones (i) y (ii) de \sref{IV.21.7.2} y están \emph{regularmente inmersos} en $X$.
\end{env}

\begin{proposition}[21.7.4]
\label{IV.21.7.4}
Sea $X$ un preesquema localmente noetheriano, reducido en cada uno de sus puntos aislados.
Las condiciones siguientes son equivalentes:
\begin{enumerate}
  \item[\emph{a)}] El homomorfismo canónico $c:\sh{D}iv_X\to\mathscr{Z}_X^1$ \sref{IV.21.6.4} es un isomorfismo de haces de grupos ordenados.
  \item[\emph{a$'$)}] Todo ciclo primo de codimensión $1$ de $X$ es la imagen por $\operatorname{cyc}$ de un divisor positivo, y el homomorfismo canónico $c:\sh{D}iv_X\to\mathscr{Z}_X^1$ es inyectivo.
  \item[\emph{a$''$)}] Para todo subpreesquema cerrado íntegro $Y$ de $X$, de codimensión $1$, la inmersión canónica $Y\to X$ es regular.
  \item[\emph{b)}] $X$ es normal y el homomorfismo $\operatorname{cyc}:\operatorname{Div}(X)\to\mathfrak{Z}^1(X)$ es biyectivo.
  \item[\emph{c)}] $X$ es localmente factorial.
\end{enumerate}
\end{proposition}

\begin{proof}
La equivalencia de \emph{b)} y \emph{c)} se ha demostrado en \sref{IV.21.6.9}, así como el hecho de que \emph{c)} implica \emph{a)}.
Además, \emph{b)} implica \emph{a$''$)} en virtud de \sref{IV.21.7.3.1}, y \emph{a)} implica trivialmente \emph{a$'$)}.
Queda por ver que \emph{a$'$)} o \emph{a$''$)} implica \emph{c)}.

Supongamos, pues, que se satisface la condición \emph{a$'$)}, y mostremos primero que $X$ es normal.
Observemos en primer lugar que, si $X$ satisface \emph{a$'$)}, ocurre lo mismo con todo esquema local $\Spec(\sh{O}_{X,x})$.
Consideremos entonces $x\in X^{(1)}$, de modo que $A=\sh{O}_{X,x}$ es un anillo local noetheriano de dimensión $1$.
Aplicando la condición \emph{a$'$)} a $\Spec(A)$ y al ciclo primo de codimensión $1$ formado por el punto cerrado de $\Spec(A)$, se ve que en $A$ el ideal maximal está engendrado por un solo elemento regular de $A$;
por tanto \sref[0]{0.17.1.1}, $A$ es un anillo regular.
Como los anillos localizados $A_{\mathfrak p}$ son también regulares \sref[0]{0.17.3.2}, se ve que $X$ es regular en todos sus puntos maximales no aislados;
como también es reducido (y por tanto regular) en sus puntos aislados por hipótesis, se concluye que $X$ satisface la condición $(\mathrm{R}_1)$.
Mostremos además que $X$ satisface también $(\mathrm{S}_1)$; en otros términos, que para todo $x\in X$ tal que $\dim(\sh{O}_{X,x})\geq1$, se tiene $\mathfrak{m}_x\notin\operatorname{Ass}(\sh{O}_{X,x})$ \sref[0]{0.16.4.6}.
En efecto, la hipótesis \emph{a$'$)} aplicada a $X_1=\Spec(\sh{O}_{X,x})$ muestra que existe sobre este preesquema al menos un divisor no nulo; en otros términos, se tiene $\sh{M}_{X_1}^*\neq\sh{O}_{X_1}^*$, y basta aplicar \sref{IV.21.1.10}.
De estos resultados se deduce ya que $X$ es reducido \sref{IV.5.8.5}.
Probemos a continuación que satisface la condición $(\mathrm{S}_2)$ (lo que establecerá que $X$ es normal, en virtud del criterio de Serre \sref{IV.5.8.6}).
Razonemos por reducción al absurdo, suponiendo que el conjunto $E$ de los puntos $x\in X$ en los que $X$ no satisface $(\mathrm{S}_2)$ no es vacío, y sea $x\in E$ un punto para el que $\dim(\sh{O}_{X,x})$ sea la menor posible;
como $X$ satisface $(\mathrm{S}_1)$, se tiene $\dim(\sh{O}_{X,x})\geq2$.
En $X_1=\Spec(\sh{O}_{X,x})$, el abierto $U=X_1-\{x\}$ satisface $(\mathrm{S}_2)$;
mostremos que $X_1$ satisface $(\mathrm{S}_2)$, con lo cual se llegará a una contradicción.
Basta, en virtud de \sref{IV.5.10.5}, mostrar que toda sección $f$ de $\sh{O}_{X_1}$ sobre $U$ se prolonga a una sección de $\sh{O}_{X_1}$ sobre $X_1$.
Como $X_1$ es reducido y $U$ es denso en $X_1$, $U$ es esquemáticamente denso en $X_1$ y, por consiguiente, $f$ es la restricción a $U$ de una sección meromorfa regular $g\in M(X_1)$.
Además, como $\dim(\sh{O}_{X,x})\geq2$, se tiene $\operatorname{cyc}(\operatorname{div}(g))\geq0$ puesto que $g|U=f$;
como el homomorfismo $\operatorname{cyc}$ es inyectivo, se tiene necesariamente $\operatorname{div}(g)\geq0$, luego \sref{IV.21.6.4}, $g$ es una sección de $\sh{O}_{X_1}$ sobre $X_1$, lo que establece nuestra afirmación.

Para mostrar que \emph{a$'$)} implica \emph{c)}, observemos entonces que la condición \emph{a$'$)} implica que el homomorfismo canónico $\operatorname{cyc}:\operatorname{Div}(X)\to\mathfrak{Z}^1(X)$ es sobreyectivo;
como $X$ es normal, basta aplicar \sref{IV.21.6.9}.

Probemos ahora que \emph{a$''$)} implica \emph{c)}.
Se deduce de \sref{IV.21.6.5} que \emph{a$''$)} implica que todo ciclo primo de codimensión $1$ sobre $X$ es la imagen por $\operatorname{cyc}$ de un divisor positivo;
la primera parte del razonamiento anterior prueba, por tanto, de nuevo que $X$ satisface $(\mathrm{R}_1)$ y $(\mathrm{S}_1)$.
Queda aún por ver que $X$ es normal (siendo entonces idéntico el final del razonamiento); es decir, que si $x\in X$ es tal que $\dim(\sh{O}_{X,x})\geq2$, se tiene $\operatorname{prof}(\sh{O}_{X,x})\geq2$.
Ahora bien, si $y$ es una generalización de $x$ tal que $\overline{\{y\}}$ sea de codimensión $1$ en $X$, el subpreesquema reducido $Y$ de $X$ que tiene a $\overline{\{y\}}$ por espacio
\oldpage[IV-4]{280}
subyacente es íntegro, y por tanto está regularmente inmerso en $X$ por hipótesis.
Esto implica que existe un elemento regular $t$ de $\sh{O}_{X,x}$ tal que el ideal $t\sh{O}_{X,x}$ sea el ideal primo que define el preesquema $Y_1$, imagen inversa de $Y$ en $X_1=\Spec(\sh{O}_{X,x})$.
Como $\sh{O}_{X,x}/t\sh{O}_{X,x}$ es íntegro, esto prueba que $\operatorname{prof}(\sh{O}_{X,x})\geq2$ y termina la demostración de \sref{IV.21.7.4}.
\end{proof}

\begin{remark}[21.7.5]
\label{IV.21.7.5}
En \sref{IV.21.7.4} no se puede sustituir la condición \emph{a$'$)} por la condición más débil de que todo ciclo primo de codimensión $1$ de $X$ sea la imagen por $\operatorname{cyc}$ de un divisor positivo.
Así lo muestra el ejemplo en que $X$ es el esquema afín definido en \sref{IV.14.1.5} («unión de un plano y de una recta que lo corta») \emph{que no es normal};
con las notaciones introducidas en \sref{IV.14.1.5}, los ciclos primos de codimensión $1$ de $X$ son los del plano $X_1$ y los puntos cerrados de la recta $X_2$ distintos del punto común a $X_1$ y $X_2$;
si $t_1$, $t_2$, $t_3$ son las imágenes de $T_1$, $T_2$, $T_3$ en $A/\mathfrak a$, se ve, pues, que los ciclos primos de codimensión $1$ de $X$ están definidos por los ideales primos principales generados por $P(t_1,t_2)$ ($P$ irreducible en $K[T_1,T_2]$) o por $t_3-a$, con $a\neq0$ en $K$;
como estos elementos son regulares en $A/\mathfrak a$, todo ciclo de codimensión $1$ es efectivamente la imagen canónica de un divisor positivo.
\end{remark}

\begin{corollary}[21.7.6]
\label{IV.21.7.6}
Sea $X$ un preesquema localmente noetheriano, reducido en cada uno de sus puntos aislados.
Las condiciones siguientes son equivalentes:
\begin{enumerate}
  \item[\emph{a)}] El homomorfismo canónico $c:\sh{D}iv_X\to\mathscr{Z}_X^1$ es un isomorfismo de haces de grupos ordenados, y se tiene $\operatorname{Div}(X)=\operatorname{Div.princ}(X)$.
  \item[\emph{a$'$)}] El homomorfismo canónico $c:\sh{D}iv_X\to\mathscr{Z}_X^1$ es inyectivo, y todo ciclo primo de codimensión $1$ es la imagen por $\operatorname{cyc}$ de un divisor positivo principal; en otros términos, es de la forma $\operatorname{cyc}(\operatorname{div}(f))$, donde $f$ es una sección regular de $\sh{O}_X$ sobre $X$.
  \item[\emph{a$''$)}] Para todo subpreesquema cerrado íntegro $Y$ de $X$, de codimensión $1$, existe una sección regular $f$ de $\sh{O}_X$ sobre $X$ tal que $Y$ esté definido por el ideal $f\sh{O}_X$.
  \item[\emph{b)}] $X$ es normal y todo ciclo primo de codimensión $1$ sobre $X$ es principal.
  \item[\emph{c)}] $X$ es localmente factorial y $\operatorname{Pic}(X)=0$.
  \item[\emph{d)}] $X$ es normal, y para todo abierto $U$ de $X$, se tiene $\operatorname{Pic}(U)=0$.
\end{enumerate}
\end{corollary}

\begin{proof}
La equivalencia de \emph{a)}, \emph{a$'$)}, \emph{a$''$)}, \emph{b)} y \emph{c)} se deduce inmediatamente de \sref{IV.21.7.4} y de \sref{IV.21.6.11}.
Además, se deduce inmediatamente de \emph{a$''$)} que todo abierto no vacío $U$ de $X$ satisface también las mismas condiciones; en otros términos, estas condiciones implican \emph{d)}.
Queda por ver que \emph{d)} implica \emph{b)}.
Ahora bien, denotemos por $U_\lambda$ los abiertos complementarios de las uniones finitas de conjuntos de la forma $\overline{\{x\}}$, donde $\dim(\sh{O}_{X,x})\geq2$;
está claro que los $U_\lambda$ forman una familia filtrante decreciente y que, para todo $x\in\bigcap_\lambda U_\lambda$, se tiene $\dim(\sh{O}_{X,x})\leq1$, por lo que $\sh{O}_{X,x}$ es un anillo principal en virtud de la hipótesis.
Se puede, pues, aplicar a la familia $(U_\lambda)$ la proposición \sref{IV.21.6.12}, que muestra que $\mathfrak{Cl}(X)$ es isomorfo a $\varinjlim\operatorname{Pic}(U_\lambda)$, luego $\mathfrak{Cl}(X)=0$ en virtud de la hipótesis, lo que prueba \emph{b)} por definición.
\end{proof}

\begin{remark}[21.7.7]
\label{IV.21.7.7}
Cuando $X=\Spec(A)$, donde $A$ es un anillo noetheriano íntegro, las condiciones equivalentes de \sref{IV.21.7.6} equivalen a decir que $A$ es un anillo \emph{factorial}.
\end{remark}

\subsection{Divisores y normalización}
\label{subsection:IV.21.8}

\begin{lemma}[21.8.1]
\label{IV.21.8.1}
Sea $f:X'\to X$ un morfismo entero de preesquemas.
Para todo $\sh{O}_{X'}$-Módulo localmente libre $E'$ de rango constante $n$ y todo $x\in X$, existe un entorno abierto $U$ de $x$ en $X$ tal que, si se pone $U'=f^{-1}(U)$, $E'|U'$ sea isomorfo a $\sh{O}_{U'}^n$.
\end{lemma}

\begin{proof}
Como la cuestión es local sobre $X$, podemos limitarnos al caso en que $X=\Spec(A)$ es afín, en cuyo caso se tiene $X'=\Spec(A')$, donde $A'$ es una $A$-álgebra entera.
Entonces $A'$ es límite inductivo de sus sub-$A$-álgebras \emph{finitas} $A'_\lambda$.
Poniendo $X'_\lambda=\Spec(A'_\lambda)$ y denotando por $p_\lambda:X'\to X'_\lambda$ el morfismo estructural, se deduce de \sref{IV.8.5.2} y \sref{IV.8.5.5} que hay un índice $\lambda$ y un $\sh{O}_{X'_\lambda}$-Módulo localmente libre $E'_\lambda$ de rango $n$ tal que $E'=p_\lambda^*(E'_\lambda)$, y bastará evidentemente demostrar el lema para $X'_\lambda$ y $E'_\lambda$;
en otras palabras, podemos limitarnos al caso en que el morfismo $f$ es \emph{finito}.
Pongamos $B=\sh{O}_{X,x}$ y sea $B'$ el anillo del esquema afín $X'_0=X'\times_X\Spec(\sh{O}_{X,x})$;
como $B$ es un anillo local y $B'$ es una $B$-álgebra finita, $B'$ es un anillo semilocal (Bourbaki, \emph{Alg.\ comm.}, cap.~V, \textsection~2, n\textsuperscript{o}~1, prop.~3);
se concluye que el $\sh{O}_{X'_0}$-Módulo localmente libre $E'_0=E'\otimes_{\sh{O}_{X'}}\sh{O}_{X'_0}$ es isomorfo a $\sh{O}_{X'_0}^n$ (Bourbaki, \emph{Alg.\ comm.}, cap.~II, \textsection~5, n\textsuperscript{o}~3, prop.~5).
Considerando $X'_0$ como límite
\oldpage[IV-4]{281}
proyectivo de los preesquemas inducidos por $X'$ sobre los abiertos $f^{-1}(U)$, donde $U$ recorre el conjunto filtrante de los entornos abiertos afines de $x$ en $X$, siguiendo el método de \sref{IV.8.1.2}, \emph{a)}), y aplicando \sref{IV.8.5.2.5}, se obtiene la conclusión buscada.
\end{proof}

\begin{corollary}[21.8.2]
\label{IV.21.8.2}
Sea $f:X'\to X$ un morfismo entero de preesquemas.
Entonces:
\begin{enumerate}[label=(\roman*)]
  \item Se tiene $\RR^1 f_*(\sh{O}^*_{X'})=0$ \sref[0]{0.12.2.1}.
  \item El grupo $\operatorname{Pic}(X')$ es canónicamente isomorfo a $\HH^1(X,f_*(\sh{O}^*_{X'}))$.
\end{enumerate}
\end{corollary}

\begin{proof}
\begin{enumerate}[label=(\roman*)]
  \item $\RR^1 f_*(\sh{O}^*_{X'})$ es el haz asociado al prehaz de grupos conmutativos $U\mapsto\HH^1(f^{-1}(U),\sh{O}^*_{X'})$ sobre $X$ (\emph{loc.\ cit.});
la fibra de este haz en un punto $x$ puede, pues, identificarse \sref[0]{0.5.4.7} con el grupo conmutativo $\varinjlim\operatorname{Pic}(f^{-1}(U))$, donde $U$ recorre el conjunto de los entornos abiertos de $x$, y los homomorfismos de transición $\operatorname{Pic}(f^{-1}(U'))\to\operatorname{Pic}(f^{-1}(U))$ para dos abiertos $U\subset U'$ están definidos por \sref{IV.21.3.2.4}.
Ahora bien, para todo $x\in X$, todo entorno abierto $U'$ de $x$ en $X$ y todo elemento $\zeta\in\operatorname{Pic}(f^{-1}(U'))$, se deduce de \sref{IV.21.8.1} que existe un entorno abierto $U\subset U'$ de $x$ tal que la imagen de $\zeta$ en $\operatorname{Pic}(f^{-1}(U))$ sea nula.
Por tanto, la fibra de $\RR^1f_*(\sh{O}^*_{X'})$ en el punto $x$ es nula.

  \item La sucesión espectral de Leray para el functor compuesto $\sh{F}'\mapsto\Gamma(X,f_*(\sh{F}'))$ de haces de grupos conmutativos sobre $X'$ (T, 2.4) da la sucesión exacta de términos de grado bajo
\[
  0\to\HH^1(X,f_*(\sh{O}^*_{X'}))\to\HH^1(X',\sh{O}^*_{X'})\to\HH^0(X,\RR^1f_*(\sh{O}^*_{X'}))
\]
y la conclusión se deduce de (i) y del isomorfismo entre $\operatorname{Pic}(X')$ y $\HH^1(X',\sh{O}^*_{X'})$ \sref[0]{0.5.4.7}.
\end{enumerate}
\end{proof}

\begin{proposition}[21.8.3]
\label{IV.21.8.3}
Sea $f=(\psi,\theta):X'\to X$ un morfismo entero de preesquemas;
supongamos que, para todo abierto $U$ de $X$, el homomorfismo $\Gamma(\theta):\Gamma(U,\sh{O}_X)\to\Gamma(U,f_*(\sh{O}_{X'}))$ transforma los elementos regulares en elementos regulares, lo que permite prolongar canónicamente el homomorfismo $\theta:\sh{O}_X\to f_*(\sh{O}_{X'})$ a un homomorfismo de haces de anillos $\theta':\sh{M}_X\to f_*(\sh{M}_{X'})$, de donde se obtienen homomorfismos de haces de grupos multiplicativos $\theta^*:\sh{O}^*_X\to f_*(\sh{O}^*_{X'})$ y $\theta'^*:\sh{M}^*_X\to f_*(\sh{M}^*_{X'})$, de los que se obtiene por paso a los cocientes un homomorfismo $\theta''^*:\sh{D}iv_X\to f_*(\sh{D}iv_{X'})$.
Se tiene entonces un diagrama conmutativo
\[
\label{IV.21.8.3.1}
  \xymatrix{
    1 \ar[r] & \sh{O}^*_X \ar[r] \ar[d]^{\theta^*} & \sh{M}^*_X \ar[r] \ar[d]^{\theta'^*} & \sh{D}iv_X \ar[r] \ar[d]^{\theta''^*} & 0\\
    1 \ar[r] & f_*(\sh{O}^*_{X'}) \ar[r] & f_*(\sh{M}^*_{X'}) \ar[r] & f_*(\sh{D}iv_{X'}) \ar[r] & 0
  }
  \tag{21.8.3.1}
\]
en el que las dos filas son exactas.
\end{proposition}

\begin{proof}
El único punto que hay que verificar es la exactitud de la segunda fila del diagrama, que se deduce de aplicar la sucesión exacta de cohomología para el functor $f_*$ a la sucesión exacta de haces de grupos conmutativos sobre $X'$
\[
  1\to\sh{O}^*_{X'}\to\sh{M}^*_{X'}\to\sh{D}iv_{X'}\to0
\]
y de \sref{IV.21.8.2}.
\end{proof}

\begin{corollary}[21.8.4]
\label{IV.21.8.4}
Si, además de las hipótesis de \sref{IV.21.8.3}, el homomorfismo $\theta'$ es un isomorfismo de haces de anillos, entonces $\theta^*:\sh{O}^*_X\to f_*(\sh{O}^*_{X'})$ es inyectivo, $\theta''^*:\sh{D}iv_X\to f_*(\sh{D}iv_{X'})$ es sobreyectivo y $\operatorname{Ker}(\theta''^*)$ es isomorfo a $\operatorname{Coker}(\theta^*)$.
\end{corollary}

\begin{proof}
Es una consecuencia inmediata del diagrama de la serpiente (Bourbaki, \emph{Alg.\ comm.}, cap.~I, \textsection~2, n\textsuperscript{o}~4, prop.~2) aplicado al diagrama \sref{IV.21.8.3.1}.
\end{proof}

\oldpage[IV-4]{282}

\begin{proposition}[21.8.5]
\label{IV.21.8.5}
Sean $X$, $X'$ dos preesquemas localmente noetherianos y $f:X'\to X$ un morfismo entero;
se supone que existe un abierto esquemáticamente denso $U$ en $X$ tal que $f^{-1}(U)$ sea esquemáticamente denso en $X'$ (cf.\ \sref{IV.20.3.5}), y que el morfismo $f^{-1}(U)\to U$, restricción de $f$, sea un isomorfismo.
Entonces:
\begin{enumerate}[label=(\roman*)]
  \item El homomorfismo $\theta':\sh{M}_X\to f_*(\sh{M}_{X'})$ está definido y es biyectivo, el homomorfismo
\[
  \theta^*:\sh{O}^*_X\to f_*(\sh{O}^*_{X'})
\]
es inyectivo, el homomorfismo $\theta''^*:\sh{D}iv_X\to f_*(\sh{D}iv_{X'})$ es sobreyectivo, $\operatorname{Ker}(\theta''^*)$ es isomorfo a $\sh{N}=\operatorname{Coker}(\theta^*)$ y el soporte del haz de grupos multiplicativos $\sh{N}$ está contenido en $S=X-U$.
  \item Supongamos además que el conjunto $S$ es discreto y pongamos, para abreviar, $\sh{O}'_X=f_*(\sh{O}_{X'})$.
Entonces se tiene un diagrama conmutativo
\[
\label{IV.21.8.5.1}
  \xymatrix{
    1 \ar[r] & \prod_{s\in S}\sh{O}'^{*}_{X,s}/\sh{O}^*_{X,s} \ar[r]^-{j} \ar[d] & \operatorname{Div}(X) \ar[r] \ar[d]^{k_X} & \operatorname{Div}(X') \ar[r] \ar[d]^{k_{X'}} & 0\\
    1 \ar[r] & \bigl(\prod_{s\in S}\sh{O}'^{*}_{X,s}/\sh{O}^*_{X,s}\bigr)/\operatorname{Im}\Gamma(X',\sh{O}^*_{X'}) \ar[r]^-{j'} & \operatorname{Pic}(X) \ar[r] & \operatorname{Pic}(X') \ar[r] & 0
  }
  \tag{21.8.5.1}
\]
donde las filas son exactas y la flecha vertical de la izquierda es el homomorfismo canónico.
\end{enumerate}
\end{proposition}

\begin{proof}
\begin{enumerate}[label=(\roman*)]
  \item La hipótesis implica que se tiene un isomorfismo canónico $\sh{M}''_X\simto f_*(\sh{M}''_{X'})$ para los haces de gérmenes de pseudo-funciones \sref{IV.20.2.2}, de donde se deduce la afirmación relativa a $\theta'$, dada la existencia de los isomorfismos canónicos $\sh{M}_X\simto\sh{M}''_X$ y $\sh{M}_{X'}\simto\sh{M}''_{X'}$ \sref{IV.20.2.11}.
El resto de la afirmación (i) se deduce de \sref{IV.21.8.4}, salvo lo relativo al soporte de $\sh{N}$, que resulta directamente de la hipótesis sobre $U$.

  \item Aplicando la sucesión exacta de cohomología a las dos sucesiones exactas de haces de grupos conmutativos
\[
\label{IV.21.8.5.2}
\begin{gathered}
  1\to\sh{N}\to\sh{D}iv_X\to f_*(\sh{D}iv_{X'})\to0,\\
  1\to\sh{O}^*_X\to f_*(\sh{O}^*_{X'})\to\sh{N}\to1,
\end{gathered}
  \tag{21.8.5.2}
\]
se obtienen respectivamente las sucesiones exactas
\[
\label{IV.21.8.5.3}
\begin{gathered}
  1\to\Gamma(X,\sh{N})\xrightarrow{j}\operatorname{Div}(X)\to\operatorname{Div}(X')\to\HH^1(X,\sh{N}),\\
  \Gamma(X',\sh{O}^*_{X'})\to\Gamma(X,\sh{N})\xrightarrow{\partial}\operatorname{Pic}(X)\to\operatorname{Pic}(X')\to\HH^1(X,\sh{N}),
\end{gathered}
  \tag{21.8.5.3}
\]
\oldpage[IV-4]{283}
teniendo en cuenta el isomorfismo canónico $\operatorname{Pic}(X')\simto\HH^1(X,f_*(\sh{O}^*_{X'}))$ \sref{IV.21.8.2}.
Ahora bien, como el soporte de $\sh{N}$ está contenido en el conjunto discreto $S$, cerrado en $X$, se tiene
\[
  \HH^1(X,\sh{N})=\HH^1(S,\sh{N}|S)\quad\text{(G, II, 4.9.2)},\qquad
  \HH^1(S,\sh{N}|S)=\prod_{s\in S}\HH^1(\{s\},\sh{N}_s)=0
\]
por definición de la cohomología (G, II, 4.4).
Asimismo, se tiene $\Gamma(X,\sh{N})=\prod_{s\in S}\sh{N}_s$ y $\sh{N}_s=\sh{O}'^{*}_{X,s}/\sh{O}^*_{X,s}$ en virtud de la segunda sucesión exacta \sref{IV.21.8.5.2}.
Queda precisar las inyecciones $j$ y $j'$.
Ahora bien, se puede describir una sección $t$ de $\sh{N}$ sobre $X$ tomando un recubrimiento de $X$ formado por $U$ y por abiertos $U_i$ tales que $U_i\cap S$ se reduzca a un solo punto $s_i$, y que $U_i\cap U_j\subset U$ para $i\neq j$, y considerando después, para cada $i$, una sección $t_i$ de $\sh{N}$ sobre $U_i$, necesariamente tal que $(t_i)_{s_i}\in\sh{O}'^{*}_{X,s_i}/\sh{O}^*_{X,s_i}$ y $(t_i)_x=0$ en los puntos $x\in U_i-\{s_i\}$.
El germen $(t_i)_{s_i}$ proviene de una sección $u_i$ de $\sh{O}^*_{X'}$ sobre $f^{-1}(U_i)$, que también puede considerarse como una sección de $\sh{M}^*_{X'}$ sobre $f^{-1}(U_i)$ y, por tanto, como una sección de $\sh{M}^*_X$ sobre $U_i$ (en virtud de la hipótesis);
a esta sección le corresponde canónicamente una sección $d_i\in\Gamma(U_i,\sh{D}iv_X)$ y, como en $f^{-1}(U\cap U_i)$ la restricción de $u_i$ se identifica con una sección de $\sh{O}^*_X$ sobre $U\cap U_i$ en virtud de la hipótesis, las restricciones de los $d_i$ a $U\cap U_i$ son todas nulas; por tanto, los $d_i$ son las restricciones de un mismo divisor perteneciente a $\operatorname{Div}(X)$, que es la imagen de la sección $t$ por $j$.
Asimismo, la imagen de $t$ por $\partial:\Gamma(X,\sh{N})\to\HH^1(X,\sh{O}^*_X)$ proviene del cociclo igual a la restricción de $u_i$ en $U\cap U_i$ para cada $i$, y a $(u_i|(U_i\cap U_j))(u_j|(U_i\cap U_j))^{-1}$ en $U_i\cap U_j$, de donde se obtienen la expresión de $j'(t)$ y la conmutatividad del diagrama \sref{IV.21.8.5.1}.
\end{enumerate}
\end{proof}

\begin{remark}[21.8.6]
\label{IV.21.8.6}
\begin{enumerate}[label=(\roman*)]
  \item Las condiciones de aplicación de \sref{IV.21.8.5}, (i) se satisfacen en particular cuando $X$ es un preesquema localmente noetheriano reducido que solo tiene un número finito de componentes irreducibles, $X'$ es su normalización y $X'$ es también localmente noetheriano (\sref[II]{II.6.3.8});
el $\sh{O}_X$-Módulo $\sh{D}iv_X$ es entonces una \emph{extensión} de $f_*(\sh{D}iv_{X'})$ por el conúcleo $\sh{N}$ de $\sh{O}^{*}_X\to f_*(\sh{O}^*_{X'})$, que puede considerarse conocido.
Si además $X$ es una curva algebraica (reducida) sobre un cuerpo $k$, se cumplen las condiciones de aplicación de \sref{IV.21.8.5}, (ii).

  \item Cuando se cumplen las condiciones de aplicación de \sref{IV.21.8.5}, (ii), y además el homomorfismo canónico global $\Gamma(X,\sh{O}_X)\to\Gamma(X',\sh{O}_{X'})$ es \emph{biyectivo}, se ve que los núcleos de los homomorfismos sobreyectivos $\operatorname{Div}(X)\to\operatorname{Div}(X')$ y $\operatorname{Pic}(X)\to\operatorname{Pic}(X')$ son \emph{isomorfos}.

  \item Cuando se cumplen las condiciones de aplicación de \sref{IV.21.8.5}, (ii), se ve que el homomorfismo $\operatorname{Div}(X)\to\operatorname{Div}(X')$ solo puede ser inyectivo (en cuyo caso es biyectivo) si $\sh{O}'^{*}_{X,s}=\sh{O}^*_{X,s}$ para todo $s\in S$.
Para un $s\in S$ tal que el anillo $\sh{O}'_{X,s}$ sea \emph{local} (lo que, teniendo en cuenta el hecho de que $X'=\Spec(\sh{O}'_X)$ \sref[II]{II.1.3}, significa que solo existe un punto $s'\in X'$ sobre $s$), esto implica necesariamente que los cuerpos residuales $\boldsymbol{k}(s)$ y $\boldsymbol{k}(s')$ son iguales y que $1+\mathfrak{m}_s=1+\mathfrak{m}_{s'}$, y por tanto equivale finalmente a la relación $\sh{O}'_{X,s}=\sh{O}_{X,s}$, o también (teniendo en cuenta la hipótesis) al hecho de que existe un entorno abierto $V$ de $s$ en $X$ tal que el morfismo $f^{-1}(V)\to V$ sea un isomorfismo.
En el caso general, el anillo $\sh{O}'_{X,s}$ es semilocal (esto es evidente cuando el morfismo $f$ es finito y, en el caso general, se deduce de \sref[0]{0.23.2.6}); el homomorfismo canónico $\sh{O}^{*}_{X,s}\to\sh{O}'^{*}_{X,s}$ define, por paso a los cocientes, un homomorfismo del grupo multiplicativo $\boldsymbol{k}(s)^*$ en el producto de los grupos multiplicativos $\boldsymbol{k}(s'_j)^*$, siendo los $s'_j$ los puntos (en
\oldpage[IV-4]{284}
número $>1$) de $X'$ situados sobre $s$.
Es inmediato que este homomorfismo solo puede ser biyectivo si $\boldsymbol{k}(s)$ y todos los $\boldsymbol{k}(s'_j)$ son iguales al cuerpo $\bb{F}_2$ de $2$ elementos;
además, si se cumple esta condición, es necesario también que el grupo multiplicativo $1+\mathfrak{m}_s$ tenga por imagen $1+\mathfrak{r}$, donde $\mathfrak{r}$ es el radical de $\sh{O}'_{X,s}$, o, equivalentemente, que $\mathfrak{m}_s=\mathfrak{r}$, lo que implica que $\sh{O}'_{X,s}\otimes_{\sh{O}_{X,s}}\boldsymbol{k}(s)$ es un producto directo de cuerpos isomorfos a $\boldsymbol{k}(s)$ (lo que, cuando el morfismo $f$ es finito, implica que no está ramificado en los puntos $s'_j$ \sref{IV.17.4.1}).
Si, por ejemplo, \emph{ningún} cuerpo residual de $X$ es isomorfo a $\bb{F}_2$, el homomorfismo canónico $\operatorname{Div}(X)\to\operatorname{Div}(X')$ solo puede ser biyectivo si $f$ es un \emph{isomorfismo}.
En el caso en que el homomorfismo canónico $\Gamma(X,\sh{O}_X)\to\Gamma(X',\sh{O}_{X'})$ es biyectivo, se concluye de lo anterior y de (ii) que, en las consideraciones precedentes, se puede reemplazar el homomorfismo $\operatorname{Div}(X)\to\operatorname{Div}(X')$ por el homomorfismo $\operatorname{Pic}(X)\to\operatorname{Pic}(X')$.
\end{enumerate}
\end{remark}

\subsection{Divisores sobre los preesquemas de dimensión 1}
\label{subsection:IV.21.9}
\label{IV.21.9}

\begin{env}[21.9.1]
\label{IV.21.9.1}
Sean $X$ un espacio topológico, $x$ un punto de $X$, $i_x:\{x\}\to X$ la inyección canónica.
Si $A(x)$ es un grupo conmutativo, se lo puede considerar como un haz de grupos conmutativos sobre el espacio $\{x\}$ reducido a un solo punto, y tomar su imagen directa $(i_x)_*(A(x))$, que es un haz de grupos conmutativos sobre $X$ \sref[0]{0.3.4.1};
de las definiciones resulta inmediatamente que, para todo abierto $U$ de $X$, $\Gamma(U,(i_x)_*(A(x)))=A(x)$ si $x\in U$, y se reduce a $0$ en caso contrario;
por tanto, para $y\in\overline{\{x\}}$ se tiene $((i_x)_*(A(x)))_y=A(x)$, y para $y\notin\overline{\{x\}}$, $((i_x)_*(A(x)))_y=0$.

Si ahora $\sh{F}$ es un haz de grupos conmutativos sobre $X$ e $Y$ un subconjunto de $X$, para todo abierto $U$ de $X$ se tiene un homomorfismo canónico
\[
\label{IV.21.9.1.1}
  \Gamma(U,\sh{F})\to\prod_{x\in U\cap Y}\sh{F}_x
  =\prod_{x\in Y}\Gamma(U,(i_x)_*(\sh{F}_x))
  \tag{21.9.1.1}
\]
y, como estos homomorfismos conmutan con las aplicaciones de restricción, definen un homomorfismo canónico de \emph{haces}
\[
\label{IV.21.9.1.2}
  j_Y:\sh{F}\to\prod_{x\in Y}(i_x)_*(\sh{F}_x).
  \tag{21.9.1.2}
\]
\end{env}

\begin{lemma}[21.9.2]
\label{IV.21.9.2}
Sean $X$ un espacio topológico localmente noetheriano, $X_0$ el conjunto de sus puntos cerrados y $\sh{F}$ un haz de grupos conmutativos sobre $X$.
Las condiciones siguientes son equivalentes:
\begin{enumerate}[label=\emph{\alph*)}]
  \item El homomorfismo canónico $j_{X_0}:\sh{F}\to\prod_{x\in X_0}(i_x)_*(\sh{F}_x)$ es inyectivo y tiene por imagen $\bigoplus_{x\in X_0}(i_x)_*(\sh{F}_x)$.
  \item[\rm a$'$)] Existe una familia de grupos conmutativos $(A(x))_{x\in X_0}$ tal que $\sh{F}$ sea isomorfo a $\bigoplus_{x\in X_0}(i_x)_*(A(x))$.
  \item Toda sección de $\sh{F}$ sobre un abierto $U$ de $X$ tiene un soporte discreto contenido en $X_0$.
\end{enumerate}
Cuando esto ocurre, para todo $x\in X_0$ el grupo $A(x)$ está determinado salvo isomorfismo único y es isomorfo a $\sh{F}_x$.
Además, el haz $\sh{F}$ es entonces flácido.
\end{lemma}

\begin{proof}
\oldpage[IV-4]{285}
Como los puntos de $X_0$ son cerrados en $X$, se tiene $((i_x)_*(A(x)))_y=0$ para todo $y\neq x\in X_0$;
esta observación prueba la afirmación de unicidad relativa a los grupos $A(x)$, y por otra parte es claro que \emph{a)} implica \emph{a$'$)}.
Para ver que \emph{a$'$)} implica \emph{b)}, podemos limitarnos al caso en que $U$ sea noetheriano;
entonces (G, II, 3.10) se tiene
\[
  \Gamma(U,\bigoplus_{x\in X_0}(i_x)_*(A(x)))
  =\bigoplus_{x\in X_0}\Gamma(U,(i_x)_*(A(x))),
\]
por lo que toda sección $s$ de $\sh{F}$ sobre $U$ es suma directa de un número \emph{finito} de secciones $s_k\in\Gamma(U,(i_{x_k})_*(A(x_k)))$ ($1\leq k\leq n$), con $x_k\in X_0$.
Pero, como $x_k$ es cerrado en $X$, el soporte de $s_k$ se reduce a $\{x_k\}$; por tanto, el soporte de $s$ está contenido en el conjunto finito de los $x_k$, que es evidentemente discreto puesto que los puntos $x_k$ son cerrados en $X$.
Demostremos finalmente que \emph{b)} implica \emph{a)};
para todo abierto noetheriano $U$, todo soporte de una sección de $\sh{F}$ sobre $U$, por ser discreto y cuasicompacto, es \emph{finito}.
Se deduce, pues, de la hipótesis \emph{b)} que, para todo abierto noetheriano $U$, la imagen del homomorfismo \sref{IV.21.9.1.1} (para $Y=X_0$) está contenida en $\bigoplus_{x\in X_0}\Gamma(U,(i_x)_*(\sh{F}_x))$ y que este homomorfismo es inyectivo, lo que establece \emph{a)}.

Por último, para demostrar que $\sh{F}$ es flácido, consideremos una sección $s$ de $\sh{F}$ sobre un abierto $U$ de $X$;
como el soporte de $s$ es un subespacio discreto y cerrado de $X$, se prolonga $s$ a una sección $s'$ de $\sh{F}$ sobre $X$ poniendo $s'_x=0$ en $X-U$.
\end{proof}

\begin{remark}[21.9.3]
\label{IV.21.9.3}
\begin{enumerate}[label=(\roman*)]
  \item En la condición \emph{b)} de \sref{IV.21.9.2}, no basta suponer que el soporte de toda sección de $\sh{F}$ sobre un abierto cualquiera de $X$ sea \emph{discreto};
así lo muestra el ejemplo en que se toma como $X$ el espectro de un anillo de valoración discreta, con un punto cerrado $b$ y un punto genérico $a$, y como $\sh{F}$ el haz de grupos conmutativos tal que $\sh{F}_b=0$ y $\sh{F}_a=\bb{Z}$.

  \item Supuestas satisfechas las condiciones de \sref{IV.21.9.2}, sea $E$ un subconjunto discreto de $X_0$ y, para todo $x\in E$, sea $a(x)$ un elemento de $A(x)$.
Existe entonces una sección única $t$ de $\bigoplus_{x\in X_0}(i_x)_*(A(x))$ sobre $X$ tal que $t_x=a(x)$ para todo $x\in E$ y que el soporte de $t$ esté contenido en $E$.
En efecto, para todo abierto noetheriano $U$, $E\cap U$ es finito, y basta tomar como $t$ la sección de $\bigoplus_{x\in E}(i_x)_*(A(x))$ cuya restricción a cada abierto noetheriano $U$ sea la suma de los $a(x)$ para $x\in E\cap U$.
\end{enumerate}
\end{remark}

\begin{proposition}[21.9.4]
\label{IV.21.9.4}
Sean $X$ un preesquema localmente noetheriano de dimensión $\leq1$ y $X^{(1)}$ el conjunto de los puntos $x\in X$ tales que $\dim(\sh{O}_{X,x})=1$.
Se tiene entonces un isomorfismo canónico
\[
\label{IV.21.9.4.1}
  \sh{D}iv_X\simto\bigoplus_{x\in X^{(1)}}(i_x)_*(\operatorname{Div}(\sh{O}_{X,x}))
  \tag{21.9.4.1}
\]
y $\sh{D}iv_X$ es flácido.
\end{proposition}

\begin{proof}
Teniendo en cuenta el isomorfismo \sref{IV.21.4.6.1}, el homomorfismo \sref{IV.21.9.4.1} está definido por \sref{IV.21.9.1.2}:
demostremos que es una biyección.
Como $\dim(X)\leq1$, los puntos de $X^{(1)}$ son los \emph{puntos cerrados no aislados} de $X$.
Nos reducimos a demostrar que:
1\textsuperscript{o} $\sh{D}iv_X$ satisface la condición \emph{b)} de \sref{IV.21.9.2};
2\textsuperscript{o} para todo punto aislado $x\in X$, se tiene $(\sh{D}iv_X)_x=0$.
El segundo punto resulta de que $\sh{O}_{X,x}$ es entonces un anillo artiniano y de que todo elemento regular de $\sh{O}_{X,x}$ es, por tanto, invertible.
\oldpage[IV-4]{286}
Para 1\textsuperscript{o}, basta observar que, para todo abierto $U$ de $X$ y todo divisor $D\in\operatorname{Div}(U)$, los puntos maximales del soporte $S$ de $D$ son tales que $\operatorname{prof}(\sh{O}_{X,x})=1$ \sref{IV.21.1.9}, y por tanto, \emph{a fortiori}, pertenecen a $X^{(1)}\cap U$;
puesto que $\dim(X)\leq1$, el conjunto de estos puntos es igual a $S$, y $S$ es, por consiguiente, discreto, ya que el conjunto de las componentes irreducibles de $S$ es localmente finito.
\end{proof}

\begin{corollary}[21.9.5]
\label{IV.21.9.5}
Sea $X$ un preesquema localmente noetheriano de dimensión $\leq1$.
Para todo subconjunto discreto $E\subset X^{(1)}$ y toda familia $(D(x))_{x\in E}$ con $D(x)\in\operatorname{Div}(\sh{O}_{X,x})$, existe un divisor $D$ sobre $X$, y solo uno, tal que el soporte de $D$ esté contenido en $E$ y que $D_x=D(x)$ para todo $x\in E$.
\end{corollary}

\begin{proof}
Esto resulta de \sref{IV.21.9.4} y de \sref{IV.21.9.3}, (ii).
\end{proof}

\begin{corollary}[21.9.6]
\label{IV.21.9.6}
Sea $X$ un preesquema localmente noetheriano de dimensión $\leq1$;
todo divisor $D$ sobre $X$ es de la forma $D'-D''$, donde $D'$ y $D''$ son dos divisores $\geq0$ cuyos soportes están contenidos en el de $D$.
\end{corollary}

\begin{proof}
En virtud de \sref{IV.21.9.5}, nos reducimos inmediatamente al caso en que $X=\Spec(A)$, donde $A$ es un anillo local noetheriano;
basta entonces observar que $M(X)$ es el anillo total de fracciones de $A$, y que una sección de $\sh{M}^*_X$ sobre $X$ se escribe, por definición, como cociente $b/a$ de dos elementos regulares de $A$.
\end{proof}

\begin{corollary}[21.9.7]
\label{IV.21.9.7}
Sean $X$ un preesquema localmente noetheriano de dimensión $\leq1$, sin puntos aislados, y $U$ un abierto denso de $X$.
Existe entonces un divisor $D\geq0$ sobre $X$, cuyo soporte está contenido en $U$ y que corta cada una de las componentes irreducibles de $X$.
\end{corollary}

\begin{proof}
Podemos suponer que $U$ es unión de abiertos disjuntos no vacíos $U_\alpha$, cada uno de los cuales está contenido en una sola componente irreducible de $X$;
basta entonces tomar en cada $U_\alpha$ un punto $x_\alpha$ cerrado en $X$ (existe uno, pues el único punto genérico de $U_\alpha$ no puede ser aislado por hipótesis), y aplicar \sref{IV.21.9.5} al conjunto discreto de los $x_\alpha$.
\end{proof}

El interés de este corolario es que permitirá demostrar que una curva algebraica separada $X$ sobre un cuerpo $k$ es cuasiproyectiva, pues el divisor $D\geq0$ definido en \sref{IV.21.9.7} es entonces necesariamente \emph{amplio} en virtud del teorema de Riemann--Roch para las curvas (cap.~V).

\begin{env}[21.9.8]
\label{IV.21.9.8}
Para los preesquemas localmente noetherianos $X$ de dimensión $\leq1$, la proposición \sref{IV.21.9.4} reduce la determinación de $\sh{D}iv_X$ al caso en que $X=\Spec(A)$, siendo $A$ un anillo local noetheriano de dimensión $1$.

Cuando $A$ es un anillo local regular de dimensión $1$ (es decir, un anillo de valoración discreta), el grupo $\operatorname{Div}(A)$ se identifica canónicamente con $\bb{Z}$ \sref{IV.21.6.8};
por tanto, en virtud de \sref{IV.21.9.2}:
\end{env}

\begin{proposition}[21.9.9]
\label{IV.21.9.9}
Sea $X$ un preesquema localmente noetheriano regular de dimensión $\leq1$.
Entonces el haz $\sh{D}iv_X$ es canónicamente isomorfo a $\bigoplus_{x\in X^{(1)}}(i_x)_*(\bb{Z}(x))$, donde $\bb{Z}(x)$ es el grupo aditivo $\bb{Z}$ considerado como haz de grupos sobre el espacio $\{x\}$.
\end{proposition}

\begin{env}[21.9.10]
\label{IV.21.9.10}
Supongamos solamente que el anillo local noetheriano $A$ de dimensión $1$ sea \emph{reducido};
entonces, si $A'$ es el \emph{normalizado} de $A$ (la clausura integral de $A$ en su anillo total de fracciones), $A'$ es noetheriano en virtud del teorema de Krull--Akizuki (Bourbaki, \emph{Alg.\ comm.}, cap.~VII, \textsection~2, n\textsuperscript{o}~5, prop.~5), y se ha visto en \sref{IV.21.8.6} cómo
\oldpage[IV-4]{287}
$\operatorname{Div}(A)$ puede obtenerse como \emph{extensión} de $\operatorname{Div}(A')$, siendo este último grupo de la forma $\bb{Z}^r$.
\end{env}

\begin{proposition}[21.9.11]
\label{IV.21.9.11}
Sean $X$ un preesquema noetheriano y $X_0$ un subpreesquema cerrado de $X$ que tenga las propiedades siguientes:
\begin{enumerate}
  \item[\rm 1\textsuperscript{o}] $\dim(X_0)\leq1$.
  \item[\rm 2\textsuperscript{o}] Para todo subconjunto localmente cerrado $Y$ de $X$ tal que $Y\cap X_0$ sea discreto, existe un subconjunto $Y'$ de $Y$, cerrado en $X$ y abierto en $Y$, que contiene a $Y\cap X_0$.
\end{enumerate}
En estas condiciones:
\begin{enumerate}[label=(\roman*)]
  \item Sea $Z_0$ la unión de los conjuntos $\overline{\{x\}}\cap X_0$ cuando $x$ recorre el subconjunto de $\operatorname{Ass}(\sh{O}_X)$ formado por los puntos tales que $\overline{\{x\}}\cap X_0$ sea finito.
Entonces, para todo divisor $D_0$ sobre $X_0$ cuyo soporte no corte a $Z_0$, existe un divisor $D$ sobre $X$ cuya imagen inversa por la inyección canónica $X_0\to X$ existe \sref{IV.21.4} y es igual a $D_0$;
si además $D_0\geq0$, se puede suponer $D\geq0$.

  \item Supongamos además que existe en $X_0$ un abierto afín $U_0$ que contiene a $\operatorname{Ass}(\sh{O}_{X_0})\cup Z_0$ (condición que se verifica automáticamente cuando existe un $\sh{O}_{X_0}$-módulo amplio \sref[II]{II.4.5.4}).
Entonces el homomorfismo canónico \sref{IV.21.3.2.4} $\operatorname{Pic}(X)\to\operatorname{Pic}(X_0)$ es sobreyectivo.
\end{enumerate}
\end{proposition}

\begin{proof}
\begin{enumerate}[label=(\roman*)]
  \item Teniendo en cuenta \sref{IV.21.9.6}, podemos limitarnos a demostrar la afirmación relativa a los divisores $D_0\geq0$;
en virtud de \sref{IV.21.9.4}, el soporte $T$ de $D_0$ es un conjunto discreto y cerrado en $X_0$.
Podemos suponer $D_0\neq0$, es decir, $T\neq\emp$, pues de lo contrario no hay nada que demostrar.
Para todo $x\in T$, existe un elemento $s_x\in\sh{O}_{X_0,x}$ que no es divisor de cero en este anillo, pertenece a su ideal maximal y cuya imagen en $(\sh{D}iv_{X_0})_x$ es $(D_0)_x$.
Existe un entorno abierto afín $U^{(x)}$ de $x$ en $X$ y una sección $g^{(x)}$ de $\sh{O}_X$ sobre $U^{(x)}$ tales que $s_x$ sea la imagen en $\sh{O}_{X_0,x}$ del germen $(g^{(x)})_x\in\sh{O}_{X,x}$;
veremos que, tomando $U^{(x)}$ suficientemente pequeño, se puede conseguir que $g^{(x)}$ no sea divisor de cero en $\Gamma(U^{(x)},\sh{O}_X)$, es decir, \sref{IV.3.1.9}, que, si $V^{(x)}$ denota el subconjunto cerrado de $U^{(x)}$ formado por los $y$ tales que $g^{(x)}(y)=0$, se tenga $V^{(x)}\cap\operatorname{Ass}(\sh{O}_X)=\emp$.

Ahora bien, si $z\in\operatorname{Ass}(\sh{O}_X)$, la hipótesis $x\notin Z_0$ implica que, o bien $x\notin\overline{\{z\}}$, o bien $x$ no está aislado en $\overline{\{z\}}\cap X_0$.
Reemplazando $U^{(x)}$ por un abierto más pequeño, se puede suponer que el segundo caso es el que se produce para todo $z\in V^{(x)}\cap\operatorname{Ass}(\sh{O}_X)$.
$V^{(x)}$ contendría entonces una componente irreducible de dimensión $1$ de $X_0\cap U^{(x)}$ que contiene a $x$;
pero esto significaría que la imagen $\overline{g}^{(x)}$ de $g^{(x)}$ en $\Gamma(U^{(x)}\cap X_0,\sh{O}_{X_0})$ estaría en el nilradical de este anillo y, por consiguiente, $s_x$ pertenecería al nilradical de $\sh{O}_{X_0,x}$, lo que es absurdo.
Se tiene, pues, $x\notin\overline{\{z\}}$ para todo punto $z\in V^{(x)}\cap\operatorname{Ass}(\sh{O}_X)$ y, como este conjunto es finito, se puede suponer, reemplazando $U^{(x)}$ por un entorno abierto más pequeño de $x$, que $V^{(x)}\cap\operatorname{Ass}(\sh{O}_X)=\emp$.

En virtud de la elección de $U^{(x)}$, se puede definir un divisor $D'^{(x)}$ sobre $U^{(x)}$ mediante $D'^{(x)}=\operatorname{div}(g^{(x)})$;
por otra parte, se ha visto más arriba que $x$ está necesariamente aislado en $V^{(x)}\cap X_0$, por lo que, reemplazando de nuevo $U^{(x)}$ por un entorno abierto más pequeño de $x$, se puede suponer que $V^{(x)}\cap X_0$ se reduce al punto $x$.
Pero existe entonces, en virtud de la condición 2\textsuperscript{o}, un subconjunto $W^{(x)}$ de $V^{(x)}$, cerrado en $X$ y abierto en $V^{(x)}$, tal que $W^{(x)}\cap X_0=\{x\}$.
Reemplazando de nuevo $U^{(x)}$ por un entorno abierto más pequeño de $x$, se puede suponer
\oldpage[IV-4]{288}
que $W^{(x)}=V^{(x)}$, es decir, que $V^{(x)}$ es cerrado en $X$.
Se puede definir entonces un divisor $D^{(x)}$ sobre $X$ mediante la condición de que $D^{(x)}|U^{(x)}=D'^{(x)}$ y $D^{(x)}|(X-V^{(x)})=0$, lo que tiene sentido porque, en $U^{(x)}\cap(X-V^{(x)})$, la restricción de $g^{(x)}$ es una sección invertible, por lo que las restricciones a este abierto de los dos divisores considerados sobre $U^{(x)}$ y $X-V^{(x)}$ son iguales.
Basta entonces, para responder a la cuestión, tomar $D=\sum_{x\in T}D^{(x)}$, lo que tiene sentido puesto que $T$ es finito.

  \item Teniendo en cuenta el diagrama conmutativo \sref{IV.21.4.2.1}, la conclusión resultará de (i) si se demuestra que todo $\sh{O}_{X_0}$-módulo invertible $\sh{L}_0$ es isomorfo a un $\sh{O}_{X_0}$-módulo de la forma $\sh{O}_{X_0}(D_0)$ \sref{IV.21.2.8}, donde $D_0$ es un divisor sobre $X_0$ cuyo soporte no corta a $Z_0$.
Ahora bien, como $U_0$ es esquemáticamente denso en $X_0$ \sref{IV.20.2.13}, (iv), basta para ello demostrar que existe una sección $t\in\Gamma(U_0,\sh{L}_0)$ tal que $t(z_0)\neq0$ en los puntos de $\operatorname{Ass}(\sh{O}_{X_0})\cup Z_0$ \sref{IV.3.1.9}.
Podemos, pues, limitarnos al caso en que $X_0=\Spec(A_0)$ sea afín;
pero entonces $\sh{L}_0$ es amplio \sref[II]{II.5.1.4} y, como el conjunto $\operatorname{Ass}(\sh{O}_{X_0})\cup Z_0$ es finito, la conclusión resulta de \sref[II]{II.4.5.4}.
\end{enumerate}
\end{proof}

\begin{corollary}[21.9.12]
\label{IV.21.9.12}
Sean $A$ un anillo local henseliano \sref{IV.18.5.8}, $S=\Spec(A)$, $s_0$ el punto cerrado de $S$, $f:X\to S$ un morfismo separado de presentación finita, y supongamos que el conjunto $X_0=f^{-1}(s_0)$ sea de dimensión $\leq1$.
Entonces, para todo subesquema cerrado $X'_0$ de $X$, que tenga a $X_0$ como conjunto subyacente y sea de presentación finita sobre $S$, el homomorfismo canónico \sref{IV.21.3.2.4} $\operatorname{Pic}(X)\to\operatorname{Pic}(X'_0)$ es sobreyectivo.
\end{corollary}

\begin{proof}
Demostremos primero que basta probar el corolario cuando el anillo local henseliano $A$ sea además noetheriano.
En efecto, se sabe \sref{IV.18.6.15} que se puede escribir $A=\varinjlim A_\lambda$, donde los $A_\lambda$ son anillos locales noetherianos y henselianos, siendo locales los homomorfismos $A_\lambda\to A$;
existe además un índice $\alpha$ y un morfismo separado de presentación finita $f_\alpha:X_\alpha\to S_\alpha=\Spec(A_\alpha)$ tales que $X$ y $f$ se deduzcan de $X_\alpha$ y $f_\alpha$ por el cambio de base $S\to S_\alpha$ \sref{IV.8.10.5}, (v);
con las notaciones habituales \sref{IV.8.8.1}, los morfismos $f_\lambda:X_\lambda\to S_\lambda$ serán separados para $\lambda\geq\alpha$ y se tendrá $X=X_\lambda\times_{S_\lambda}S$.
Además, se puede suponer que, si $s_{0\alpha}$ es el único punto cerrado de $S_\alpha$, existe un subesquema cerrado $X'_{0\alpha}$ de $X_\alpha$, que tiene el mismo conjunto subyacente que $f_\alpha^{-1}(s_{0\alpha})$, tal que $X'_0=X'_{0\alpha}\times_{S_\alpha}S$ \sref{IV.8.6.3};
para $\lambda\geq\alpha$, pongamos $X'_{0\lambda}=X'_{0\alpha}\times_{S_\alpha}S_\lambda$;
se tiene $\dim(X'_{0\lambda})=\dim(X'_0)\leq1$ por transitividad de las fibras y \sref{IV.4.1.4}.
Se trata de ver que, si se ha demostrado que el homomorfismo $\operatorname{Pic}(X_\lambda)\to\operatorname{Pic}(X'_{0\lambda})$ es sobreyectivo para todo $\lambda\geq\alpha$, lo mismo ocurre con $\operatorname{Pic}(X)\to\operatorname{Pic}(X'_0)$.
Se tiene evidentemente el diagrama conmutativo de homomorfismos canónicos
\[
  \xymatrix{
    \operatorname{Pic}(X_\lambda) \ar[r] \ar[d] & \operatorname{Pic}(X) \ar[d]\\
    \operatorname{Pic}(X'_{0\lambda}) \ar[r] & \operatorname{Pic}(X'_0)
  }
\]
Para todo $\sh{O}_{X'_0}$-módulo invertible $\sh{L}'_0$, existe un $\lambda\geq\alpha$ y un $\sh{O}_{X'_{0\lambda}}$-módulo invertible $\sh{L}'_{0\lambda}$ tal que $\sh{L}'_0$ se deduzca de él por el cambio de base $S\to S_\lambda$ \sref{IV.8.5.2} y \sref{IV.8.5.5};
basta considerar un $\sh{O}_{X_\lambda}$-módulo invertible $\sh{L}_\lambda$ tal que $\operatorname{cl}(\sh{L}'_{0\lambda})$ sea la imagen de $\operatorname{cl}(\sh{L}_\lambda)$ en $\operatorname{Pic}(X'_{0\lambda})$, y tomar luego el $\sh{O}_X$-módulo invertible $\sh{L}$ deducido de $\sh{L}_\lambda$ por cambio de base;
en virtud de la conmutatividad del diagrama anterior, $\operatorname{cl}(\sh{L}'_0)$ será la imagen de $\operatorname{cl}(\sh{L})$.

Supongamos, pues, que $A$ es noetheriano, y por tanto también $X$, y comprobemos que $X$ y $X'_0$ satisfacen las condiciones de \sref{IV.21.9.11}, (ii).
Por hipótesis, se tiene $\dim(X'_0)\leq1$;
por otra parte, para verificar la condición 2\textsuperscript{o} de \sref{IV.21.9.11}, consideremos el subpreesquema reducido $Y$ de $X$ que tenga el subconjunto localmente cerrado $Y$ como espacio subyacente;
como el morfismo $g:Y\to S$, restricción de $f$, es cuasifinito en cada uno de los puntos $x_i$ de $Y\cap X_0$ ($1\leq i\leq n$), se puede aplicar \sref{IV.18.5.11}, c), y se ve que $Y$ es suma de los subpreesquemas abiertos $Y_i=\Spec(\sh{O}_{Y,x_i})$ ($1\leq i\leq n$) y de un preesquema $Y''$ tal que $Y''\cap X_0=\emp$, y además que las inyecciones canónicas $j_i:Y_i\to X$ son tales que $f\circ j_i$ sea un morfismo finito.
Como $f$ es separado, $j_i$ es también un morfismo finito, por lo que $Y_i$ es cerrado en $X$;
se responde, pues, a la cuestión tomando $Y'=\bigcup_{1\leq i\leq n}Y_i$.

Queda por verificar la hipótesis suplementaria de \sref{IV.21.9.11}, (ii).
Ahora bien, como $X'_0$ es una curva separada sobre el cuerpo $\boldsymbol{k}(s_0)$, es un $\boldsymbol{k}(s_0)$-esquema cuasiproyectivo (cap.~V)\footnote{El lector comprobará que \sref{IV.21.9.12} no se utiliza para demostrar esta propiedad en el cap.~V.}, por lo que existe un $\sh{O}_{X'_0}$-módulo amplio \sref[II]{II.5.3.1}, y esto termina la demostración.
\end{proof}

\begin{remark}[21.9.13]
\label{IV.21.9.13}
Bajo las condiciones de \sref{IV.21.9.12}, suponiendo además que $f$ sea propio, el morfismo $f$ es incluso \emph{proyectivo}:
en efecto, si $\sh{L}_0$ es un $\sh{O}_{X'_0}$-módulo amplio, existe, en virtud de \sref{IV.21.9.12}, un $\sh{O}_X$-módulo invertible $\sh{L}$ cuya imagen inversa en $X'_0$ es isomorfa a $\sh{L}_0$, y por tanto es amplia.
Como todo entorno de $s_0$ en $S$ coincide necesariamente con $S$, se deduce entonces de \sref{IV.9.6.4} que $\sh{L}$ es un $\sh{O}_X$-módulo amplio, de donde resulta la conclusión (\sref[II]{II.5.3.1} y \sref[II]{II.5.5.3}).
\end{remark}

\subsection{Imágenes inversas e imágenes directas de ciclos 1-codimensionales}
\label{subsection:IV.21.10}

En un capítulo posterior, dedicado a la teoría de intersecciones, se desarrollarán sistemáticamente las nociones de imagen inversa e imagen directa de ciclos.
En el presente número nos limitaremos a definir estas nociones en ciertos casos particulares útiles y, para los ciclos $1$-codimensionales, estas definiciones se elegirán de modo que sean compatibles con las definiciones correspondientes para los divisores \sref{IV.21.4} y \sref{IV.21.5}, habida cuenta del homomorfismo $\operatorname{cyc}$ definido en \sref{IV.21.6}.

\begin{env}[21.10.1]
\label{IV.21.10.1}
\oldpage[IV-4]{289}
Sean $X$, $X'$ dos preesquemas localmente noetherianos, $f:X'\to X$ un morfismo y $T$ una parte de $X$;
supongamos que la imagen por $f$ de todo punto maximal de $X'$ sea un punto maximal de $X$, y que, para todo $x'\in X'^{(1)}$ (es decir, tal que $\dim(\sh{O}_{X',x'})=1$), se cumpla una de las tres condiciones siguientes para el punto $x=f(x')$:
\begin{enumerate}[label=(\roman*)]
  \item $x\notin T$;
  \item $x\in X^{(1)}$ y $\sh{O}_{X',x'}$ es un $\sh{O}_{X,x}$-módulo plano;
  \item el anillo $\sh{O}_{X,x}$ es factorial y $\mathfrak{m}_{x'}\notin\operatorname{Ass}(\sh{O}_{X',x'})$.
\end{enumerate}
Bajo estas condiciones nos proponemos, para todo ciclo $1$-codimensional $Z$ sobre $X$ cuyo soporte esté contenido en $T$, definir un ciclo $1$-codimensional $Z'=f^*(Z)$ de modo que $Z\mapsto f^*(Z)$ sea un homomorfismo de grupos ordenados desde el subgrupo de $\mathfrak{Z}^1(X)$ formado por los ciclos con soporte contenido en $T$ hacia el grupo ordenado $\mathfrak{Z}^1(X')$.
Para ello, sea
\[
\label{IV.21.10.1.1}
  Z=\sum_{x\in T\cap X^{(1)}}n_x\cdot\overline{\{x\}}
  \tag{21.10.1.1}
\]
donde la familia de los $x\in T\cap X^{(1)}$ tales que $n_x\neq0$ es localmente finita.
Para todo $x'\in X'^{(1)}$, definamos un entero $n_{x'}$ del modo siguiente, poniendo $x=f(x')$:
\begin{enumerate}[label=\arabic*\textsuperscript{o}]
  \item si $x\notin T$, se toma $n_{x'}=0$;
  \item si $x\in X^{(1)}$ y si $\sh{O}_{X',x'}$ es un $\sh{O}_{X,x}$-módulo plano, se sabe por \sref{IV.6.1.1} que $\dim(\sh{O}_{X',x'}/\mathfrak{m}_x\sh{O}_{X',x'})=0$; dicho de otro modo, $\sh{O}_{X',x'}/\mathfrak{m}_x\sh{O}_{X',x'}$ es un $\sh{O}_{X',x'}$-módulo de longitud finita $\lambda_{x'}$, y se toma $n_{x'}=\lambda_{x'}n_x$;
  \item si $\sh{O}_{X,x}$ es factorial y $\mathfrak{m}_{x'}\notin\operatorname{Ass}(\sh{O}_{X',x'})$, se sabe por \sref{IV.21.6.9} que el homomorfismo canónico $\operatorname{cyc}:\operatorname{Div}(\sh{O}_{X,x})\to\mathfrak{Z}^1(\Spec(\sh{O}_{X,x}))$ es biyectivo; por otra parte, como $\dim(\sh{O}_{X',x'})=1$ y $\mathfrak{m}_{x'}\notin\operatorname{Ass}(\sh{O}_{X',x'})$, $\operatorname{Ass}(\sh{O}_{X',x'})$ consta únicamente de los puntos maximales de $\Spec(\sh{O}_{X',x'})$, de modo que la hipótesis sobre $f$ implica $f(\operatorname{Ass}(\sh{O}_{X',x'}))\subset\operatorname{Ass}(\sh{O}_{X,x})$, y de \sref{IV.21.4.5}, (ii), resulta que está definido el homomorfismo $f^*:\operatorname{Div}(\sh{O}_{X,x})\to\operatorname{Div}(\sh{O}_{X',x'})$;
por último, como $x'$ es el único punto cerrado de $\Spec(\sh{O}_{X',x'})$, $\mathfrak{Z}^1(\Spec(\sh{O}_{X',x'}))$ es canónicamente isomorfo a $\bb{Z}$.
\end{enumerate}
Se tiene, pues, un homomorfismo canónico compuesto
\[
\label{IV.21.10.1.2}
  \mathfrak{Z}^1(\Spec(\sh{O}_{X,x}))
  \xrightarrow{\operatorname{cyc}^{-1}}\operatorname{Div}(\sh{O}_{X,x})
  \xrightarrow{f^*}\operatorname{Div}(\sh{O}_{X',x'})
  \xrightarrow{\operatorname{cyc}}\mathfrak{Z}^1(\Spec(\sh{O}_{X',x'}))
  \xrightarrow{\sim}\bb{Z}
  \tag{21.10.1.2}
\]
Si $Z_x$ es el ciclo $\sum_{y\in\Spec(\sh{O}_{X,x})\cap X^{(1)}}n_y(\overline{\{y\}}\cap\Spec(\sh{O}_{X,x}))$ sobre $\Spec(\sh{O}_{X,x})$, se toma entonces como $n_{x'}$ la \emph{imagen} de $Z_x$ por el homomorfismo \sref{IV.21.10.1.2}.
\end{env}

\oldpage[IV-4]{290}
\begin{env}[21.10.2]
\label{IV.21.10.2}
Nos proponemos demostrar que:
\begin{enumerate}[label=\emph{\Alph*)}]
  \item Cuando dos de las condiciones 1\textsuperscript{o}, 2\textsuperscript{o}, 3\textsuperscript{o} de \sref{IV.21.10.1} se satisfacen \emph{simultáneamente}, los valores correspondientes de $n_{x'}$ \emph{coinciden}.
  \item El conjunto de los $x'\in X'^{(1)}$ tales que $n_{x'}\neq0$ es \emph{localmente finito} en $X'$.
\end{enumerate}

Para demostrar A), supongamos primero que $x\notin T$ y que $x$ satisface una de las condiciones 2\textsuperscript{o} o 3\textsuperscript{o} de \sref{IV.21.10.1};
entonces $n_x=0$ y, si estamos en el caso 2\textsuperscript{o}, se tiene $n_{x'}=0$;
si estamos en el caso 3\textsuperscript{o}, se tiene $Z_x=0$ puesto que $\Supp(Z)\subset T$, y por tanto también $n_{x'}=0$.
Queda por considerar el caso en que se está a la vez en los casos 2\textsuperscript{o} y 3\textsuperscript{o};
entonces, puesto que $\dim(\sh{O}_{X,x})=1$, $\sh{O}_{X,x}$ es un anillo de valoración discreta;
si $t$ es un uniformizador de este anillo, el divisor correspondiente a $Z_x$ es $\operatorname{div}(t^{n_x})$ en $\operatorname{Div}(\sh{O}_{X,x})$, y su imagen en $\operatorname{Div}(\sh{O}_{X',x'})$ es el divisor de $t'^{n_x}$, donde $t'$ es el elemento (regular) de $\sh{O}_{X',x'}$ imagen de $t$.
Evidentemente, podemos limitarnos al caso $n_x=1$; entonces la definición \sref{IV.21.6.5.1} muestra que la imagen de $Z_x$ por \sref{IV.21.10.1.2} es el número $\lambda_{x'}$, lo que termina la demostración de A).

Demostremos ahora B).
Pongamos $T_0=\Supp(Z)\subset T$, $T'_0=f^{-1}(T_0)$;
basta demostrar que la relación $n_{x'}\neq0$ implica que $x'$ pertenece al conjunto de los puntos maximales de $T'_0$, siendo este último localmente finito en $X'$.
Es inmediato que necesariamente se tiene $x'\in T'_0$;
si $x'$ no fuera maximal en el conjunto cerrado $T'_0$, existiría
\oldpage[IV-4]{291}
una generalización $y'$ de $x'$ en $T'_0$, distinta de $x'$, y como $x'\in X'^{(1)}$, $y'$ sería necesariamente un punto maximal de $X'$;
por consiguiente, $y=f(y')$ sería un punto maximal de $X$ por hipótesis;
pero esto es absurdo, puesto que $y\in T_0$ y $T_0$ es puramente de codimensión $1$ en $X$.
\end{env}

\begin{env}[21.10.3]
\label{IV.21.10.3}
Ahora podemos poner
\[
  f^*(Z)=\sum_{x'\in X'^{(1)}}n_{x'}\overline{\{x'\}},
\]
teniendo sentido la suma del segundo miembro en virtud de lo demostrado en \sref{IV.21.10.2};
se dice que el ciclo $1$-codimensional $f^*(Z)$ es la \emph{imagen inversa} de $Z$ por $f$.
Es inmediato que la aplicación $f^*$ así definida es un homomorfismo de grupos ordenados.
Además, si $U$ es un abierto de $X$, $V$ un abierto de $X'$ tal que $f(V)\subset U$ y $f':V\to U$ la restricción de $f$, de las definiciones resulta inmediatamente que
\[
\label{IV.21.10.3.1}
  f'^*(Z|U)=f^*(Z)|V.
  \tag{21.10.3.1}
\]

Denotemos por $\Gamma_T(\mathscr{Z}^1_X)$ el mayor subhaz de grupos conmutativos de $\mathscr{Z}^1_X$ cuyo soporte está contenido en $T$;
de la relación \sref{IV.21.10.3.1} resulta que las aplicaciones $\Gamma(U,\Gamma_T(\mathscr{Z}^1_X))\to\Gamma(V,\mathscr{Z}^1_{X'})$ que acabamos de definir determinan un homomorfismo de haces de grupos conmutativos ordenados sobre $X'$
\[
  \psi^*(\Gamma_T(\mathscr{Z}^1_X))\to\mathscr{Z}^1_{X'}
\]
donde $\psi$ es la aplicación continua subyacente al morfismo $f$.
\end{env}

\begin{proposition}[21.10.4]
\label{IV.21.10.4}
Supongamos satisfechas las condiciones de \sref{IV.21.10.1}.
Entonces, para todo divisor $D$ sobre $X$ tal que $\Supp(D)\subset T$ y esté definido $f^*(D)$ \sref{IV.21.4.2}, se tiene
\[
\label{IV.21.10.4.1}
  \operatorname{cyc}(f^*(D))=f^*(\operatorname{cyc}(D)).
  \tag{21.10.4.1}
\]
\end{proposition}

\begin{proof}
Como la cuestión es local sobre $X$, podemos limitarnos al caso en que $X=\Spec(A)$ es afín, $D=\operatorname{div}(t)$, $t$ es un elemento regular no invertible de $A$ y el subpreesquema $Y(D)$ \sref{IV.21.2.12} tiene un único punto maximal $y$, de modo que $\operatorname{cyc}(D)=n_y\cdot\overline{\{y\}}$, donde $n_y$ es la longitud de $\sh{O}_{X,y}/t_y\sh{O}_{X,y}$ \sref{IV.21.6.5.1}.
En la demostración de \sref{IV.21.10.2}, B), hemos visto que los puntos $x'\in X'$ tales que $\operatorname{mult}_{x'}(f^*(\operatorname{cyc}(D)))\neq0$ son puntos maximales de $f^{-1}(\overline{\{y\}})$.
Si el punto $x'$ está en el caso 3\textsuperscript{o} de \sref{IV.21.10.1}, la igualdad de las multiplicidades en $x'$ de los dos miembros de \sref{IV.21.10.4.1} resulta de la definición de $n_{x'}$ mediante el homomorfismo \sref{IV.21.10.1.2}.
Supongamos, por el contrario, que $x'$ está en el caso 2\textsuperscript{o} de \sref{IV.21.10.1}, y sea $x=f(x')$;
puesto que $x\in X^{(1)}\cap\overline{\{y\}}$, se tiene necesariamente $x=y$.
Observemos ahora que $f^*(D)=\operatorname{div}(t')$, donde $t'$ es la imagen de $t$ en $\Gamma(X',\sh{O}_{X'})$, y $Y(f^*(D))=f^{-1}(Y(D))$;
la multiplicidad de $f^*(D)$ en el punto $x'$ es, pues, la longitud $n'_{x'}$ del $\sh{O}_{X',x'}$-módulo
\[
  \sh{O}_{X',x'}/t'_{x'}\sh{O}_{X',x'}
  =(\sh{O}_{X,y}/t_y\sh{O}_{X,y})\otimes_{\sh{O}_{X,y}}\sh{O}_{X',x'};
\]
como de \sref{IV.4.7.1} resulta que $n'_{x'}=\lambda_{x'}n_y$, las multiplicidades en el punto $x'$ de los dos miembros de \sref{IV.21.10.4.1} son también iguales, lo que termina la demostración.
\end{proof}

\begin{env}[21.10.5]
\label{IV.21.10.5}
\oldpage[IV-4]{292}
Supongamos ahora que $f:X'\to X$ es un morfismo de preesquemas localmente noetherianos que transforma todo punto maximal de $X'$ en un punto maximal de $X$, y supongamos además que, para toda parte cerrada rara $T$ de $X$, $f$ satisface las condiciones de \sref{IV.21.10.1};
esto significa también que, para todo $x'\in X'^{(1)}$, o bien $x=f(x')$ es un punto maximal de $X$, o bien $x'$ satisface una de las condiciones (ii), (iii) de \sref{IV.21.10.1}.
Teniendo en cuenta que todo ciclo $1$-codimensional tiene soporte raro en $X$, se ve que $f^*(Z)$ está definido para todo ciclo $1$-codimensional $Z$ sobre $X$;
en virtud de \sref{IV.21.10.3.1}, se ha definido así un homomorfismo de haces de grupos conmutativos ordenados
\[
  \psi^*(\mathscr{Z}^1_X)\to\mathscr{Z}^1_{X'}.
\]
Si además, para todo divisor $D$ sobre $X$, está definido $f^*(D)$ \sref{IV.21.4.5}, el hecho de que el soporte de $D$ sea raro en $X$ \sref{IV.21.6.6} implica que \sref{IV.21.10.4} es aplicable y, por tanto, se tiene la fórmula \sref{IV.21.10.4.1} para todo divisor $D$ sobre $X$.
En particular:
\end{env}

\begin{proposition}[21.10.6]
\label{IV.21.10.6}
Sean $X$, $X'$ dos preesquemas localmente noetherianos y $f:X'\to X$ un morfismo plano.
Entonces $f^*(Z)$ está definido para todo ciclo $1$-codimensional $Z$ sobre $X$, $f^*(D)$ está definido para todo divisor $D$ sobre $X$ y se tiene la relación \sref{IV.21.10.4.1}.
\end{proposition}

\begin{proof}
En efecto, si $x'\in X'^{(1)}$ es tal que $x=f(x')$ no es maximal, de \sref{IV.6.1.1} resulta necesariamente que $x\in X^{(1)}$, de modo que estamos en el caso (ii) de \sref{IV.21.10.1}.
Podemos, pues, aplicar \sref{IV.21.10.5}, teniendo en cuenta \sref{IV.21.4.5} y \sref{IV.2.3.4}.
\end{proof}

\begin{remark}[21.10.7]
\label{IV.21.10.7}
La existencia de $f^*(Z)$ para todo ciclo $1$-codimensional $Z$ sobre $X$ resulta ya de la hipótesis de que $f$ es \emph{plano en los puntos $x'$ de $X'$ de codimensión $\leq1$ en $X'$} (es decir, tales que $\dim\sh{O}_{X',x'}\leq1$);
en efecto, para todo punto maximal $z'\in X'$, de \sref{IV.6.1.1} resulta que $z=f(z')$ es un punto maximal de $X$, puesto que $\dim\sh{O}_{X,z}\leq\dim\sh{O}_{X',z'}=0$.
Asimismo, si $x'\in X'^{(1)}$, $x=f(x')$ es maximal o pertenece a $X^{(1)}$ por \sref{IV.6.1.1};
podemos, pues, aplicar de nuevo \sref{IV.21.10.5}.
\end{remark}

\begin{proposition}[21.10.8]
\label{IV.21.10.8}
Sean $X$, $X'$, $X''$ tres preesquemas localmente noetherianos y $f:X'\to X$, $g:X''\to X'$ dos morfismos;
se supone que $f$ (resp.\ $g$) es plano en todo punto de codimensión $\leq1$ en $X'$ (resp.\ $X''$).
Entonces $f\circ g$ es plano en todo punto de codimensión $\leq1$ en $X''$ y, para todo ciclo $1$-codimensional $Z$ sobre $X$, se tiene $(f\circ g)^*(Z)=g^*(f^*(Z))$.
\end{proposition}

\begin{proof}
La primera afirmación resulta de \sref{IV.6.1.1} y \sref{IV.2.1.6}.
La segunda resulta de que $f$ (resp.\ $g$, resp.\ $f\circ g$) transforma los puntos maximales de $X'$ (resp.\ $X''$, resp.\ $X''$) en puntos maximales de $X$ (resp.\ $X'$, resp.\ $X$), y de que, si $x''\in X''^{(1)}$ es tal que $x'=g(x'')\in X'^{(1)}$ y $x=f(x')\in X^{(1)}$, se tiene
\[
  \operatorname{long}(\sh{O}_{X'',x''}/\mathfrak{m}_x\sh{O}_{X'',x''})
  =\operatorname{long}(\sh{O}_{X'',x''}/\mathfrak{m}_{x'}\sh{O}_{X'',x''})
   \cdot\operatorname{long}(\sh{O}_{X',x'}/\mathfrak{m}_x\sh{O}_{X',x'}).
\]
En efecto, si se pone $A=\sh{O}_{X,x}$, $A'=\sh{O}_{X',x'}$, $A''=\sh{O}_{X'',x''}$, $\mathfrak{m}=\mathfrak{m}_x$, $\mathfrak{m}'=\mathfrak{m}_{x'}$, se tiene $A''/\mathfrak{m}A''=(A'/\mathfrak{m}A')\otimes_{A'}A''$, y la fórmula
\[
  \operatorname{long}_{A''}(A''/\mathfrak{m}A'')
  =\operatorname{long}_{A'}(A'/\mathfrak{m}A')
   \cdot\operatorname{long}_{A''}(A''/\mathfrak{m}'A'')
\]
resulta de \sref{IV.4.7.1} y de la hipótesis de planitud de $A''$ sobre $A'$.
\end{proof}

\begin{env}[21.10.9]
\label{IV.21.10.9}
Sean $X$ un preesquema localmente noetheriano y $\Lambda$ un grupo conmutativo, escrito aditivamente y considerado, por tanto, como $\bb{Z}$-módulo;
denotaremos también por $\bb{Z}$ y $\Lambda$ los haces simples sobre $X$ asociados a los prehaces constantes iguales respectivamente a $\bb{Z}$ y $\Lambda$ \sref[0]{0.3.6}.
Se llama \emph{haz de gérmenes de ciclos $1$-codimensionales con coeficientes en $\Lambda$} al haz de grupos conmutativos $\mathscr{Z}^1_X\otimes_{\bb{Z}}\Lambda$.
Si $\Lambda=\bb{Q}$, se dice que las secciones
\oldpage[IV-4]{293}
de $\mathscr{Z}^1_X\otimes_{\bb{Z}}\bb{Q}$ sobre $X$ son los \emph{ciclos $1$-codimensionales con coeficientes racionales}.
Como las fibras de $\mathscr{Z}^1_X$ son $\bb{Z}$-módulos sin torsión \sref{IV.21.6.3}, el homomorfismo canónico $\mathscr{Z}^1_X\to\mathscr{Z}^1_X\otimes_{\bb{Z}}\bb{Q}$ es \emph{inyectivo}, de modo que los ciclos $1$-codimensionales se identifican con ciclos $1$-codimensionales con coeficientes racionales.
\end{env}

\begin{env}[21.10.10]
\label{IV.21.10.10}
Vamos a ver que, bajo ciertas condiciones, se puede ampliar la definición de $f^*(Z)$ dada en \sref{IV.21.10.3} para un ciclo $1$-codimensional $Z$ sobre $X$, pero a condición de tomar como $f^*(Z)$ un ciclo $1$-codimensional con coeficientes racionales sobre $X'$.
El caso más general en que nos situamos es aquel en el que $f$ transforma todo punto maximal de $X'$ en un punto maximal de $X$ y, en todo punto $x'\in X'^{(1)}$, se cumple una de las condiciones (i), (ii), (iii) de \sref{IV.21.10.1} \emph{o una cuarta condición} (poniendo $x=f(x')$):
\begin{enumerate}
  \item[(iv)] $x\in X^{(1)}$, $\mathfrak{m}_x\notin\operatorname{Ass}(\sh{O}_{X,x})$ y, \emph{además, si se pone} $A=\widehat{\sh{O}}_{X,x}$, $A'=\widehat{\sh{O}}_{X',x'}$, \emph{y si $K$ denota el anillo total de fracciones de $A$, entonces $A'$ es una $A$-álgebra finita y $K'=A'\otimes_AK$ es un $K$-módulo} libre.
\end{enumerate}

Denotemos entonces por $r_{x'}$ el rango del $K$-módulo libre $K'$, por $q_{x'}$ el grado de $\boldsymbol{k}(x')=\boldsymbol{k}(A')$ sobre $\boldsymbol{k}(x)=\boldsymbol{k}(A)$, y pongamos $\mu_{x'}=r_{x'}/q_{x'}$.
Para un ciclo $1$-codimensional con soporte contenido en $T$, dado por \sref{IV.21.10.1.1}, y para todo $x'\in X'^{(1)}$, se define $c_{x'}\in\bb{Q}$ igual al número $n_{x'}\in\bb{Z}$ cuando se está en uno de los casos $1\textsuperscript{o}$, $2\textsuperscript{o}$, $3\textsuperscript{o}$ de \sref{IV.21.10.1}; pero queda aquí una cuarta posibilidad:
\begin{enumerate}
  \item[$4\textsuperscript{o}$] si $x'$ satisface la condición (iv) anterior, se toma $c_{x'}=\mu_{x'}n_x\in\bb{Q}$.
\end{enumerate}
\end{env}

\begin{env}[21.10.11]
\label{IV.21.10.11}
Queda por demostrar que, cuando la condición $4\textsuperscript{o}$ de \sref{IV.21.10.10} se satisface \emph{al mismo tiempo} que una de las condiciones $1\textsuperscript{o}$, $2\textsuperscript{o}$, $3\textsuperscript{o}$ de \sref{IV.21.10.1}, se tiene $n_{x'}=c_{x'}$.
Esto es evidente si $x\notin T$, pues entonces $n_x=0$.
Para estudiar los otros dos casos, observemos que $\mathfrak{m}_x\sh{O}_{X',x'}$ es cerrado para la topología $\mathfrak{m}_{x'}$-preádica, de modo que el completado de $\sh{O}_{X',x'}/\mathfrak{m}_x\sh{O}_{X',x'}$ para esta última topología es $A'/\mathfrak{m}_xA'$.
Si se está a la vez en los casos $2\textsuperscript{o}$ y $4\textsuperscript{o}$, $\sh{O}_{X',x'}/\mathfrak{m}_x\sh{O}_{X',x'}$ es discreto para la topología $\mathfrak{m}_{x'}$-preádica y, por tanto, isomorfo a $A'/\mathfrak{m}_xA'$.
Como $A'$ es una $A$-álgebra finita y \emph{plana} $(\mathbf{0}_{\mathrm{III}},10.2.3)$, es un $A$-módulo \emph{libre} $(\mathbf{0}_{\mathrm{III}},10.1.3)$, y el rango de $A'/\mathfrak{m}_xA'$ sobre $A/\mathfrak{m}_xA=\boldsymbol{k}(x)$ es igual al rango de $A'$ sobre $A$ y, por tanto, también al de $K'$ sobre $K$.
Por otra parte, este rango es también igual al producto de la longitud $\lambda_{x'}$ de $\sh{O}_{X',x'}/\mathfrak{m}_x\sh{O}_{X',x'}$ (como $\sh{O}_{X',x'}$-módulo, o como $\boldsymbol{k}(x')$-módulo) por el rango $[\boldsymbol{k}(x'): \boldsymbol{k}(x)]=q_{x'}$, lo que demuestra la relación $\mu_{x'}=r_{x'}/q_{x'}=\lambda_{x'}$ en este caso.

Supongamos, por último, que se está a la vez en los casos $3\textsuperscript{o}$ y $4\textsuperscript{o}$.
Entonces, puesto que $\dim(\sh{O}_{X,x})=1$, $\sh{O}_{X,x}$ es un anillo de valoración discreta y, por tanto, regular.
Por otra parte, $\sh{O}_{X',x'}$ tiene dimensión $1$ y, puesto que $\mathfrak{m}_{x'}\notin\operatorname{Ass}(\sh{O}_{X',x'})$, $\sh{O}_{X',x'}$ es un anillo de Cohen--Macaulay $(\mathbf{0},16.4.6)$;
por último, puesto que $A'$ es un $A$-módulo de tipo finito, $A'/\mathfrak{m}_xA'$ es un $\boldsymbol{k}(x)$-espacio vectorial de rango finito;
como $\sh{O}_{X',x'}/\mathfrak{m}_x\sh{O}_{X',x'}$ está contenido en $A'/\mathfrak{m}_xA'$, es también un $\boldsymbol{k}(x)$-espacio vectorial de rango finito y, por tanto, un anillo artiniano.
La aplicación de \sref{IV.6.1.5} muestra entonces que $\sh{O}_{X',x'}$ es un $\sh{O}_{X,x}$-módulo \emph{plano}; por tanto, se está también en el caso $2\textsuperscript{o}$ y se concluye por lo visto anteriormente.

Así pues, en el caso que hemos considerado, se pondrá
\[
  f^*(Z)=\sum_{x'\in X'^{(1)}}c_{x'}\cdot\overline{\{x'\}}.
\]
\oldpage[IV-4]{294}
que es, por tanto, un ciclo $1$-codimensional sobre $X'$ con coeficientes \emph{racionales}.
Se ha definido así también un homomorfismo $f^*$ de grupos ordenados que satisface \sref{IV.21.10.3.1} y, por consiguiente, un homomorfismo de haces de grupos conmutativos
\[
  \psi^*(\Gamma_T(\mathscr{Z}^1_X))\to\mathscr{Z}^1_{X'}\otimes_{\bb{Z}}\bb{Q}.
\]

Cuando $f$ satisface las condiciones anteriores para \emph{toda parte cerrada $T$ rara} en $X$, es decir, cuando para todo $x'\in X'^{(1)}$ o bien $x=f(x')$ es un punto maximal de $X$, o bien $x'$ satisface una de las condiciones (ii), (iii) o (iv), entonces $f^*(Z)$ está definido para todo ciclo $1$-codimensional $Z$ sobre $X$, y se ha definido un homomorfismo de haces de grupos conmutativos ordenados
\[
  \psi^*(\mathscr{Z}^1_X)\to\mathscr{Z}^1_{X'}\otimes_{\bb{Z}}\bb{Q},
\]
de donde, por tensorización, se obtiene un homomorfismo de haces de $\bb{Q}$-espacios vectoriales ordenados
\[
\label{IV.21.10.11.1}
  \psi^*(\mathscr{Z}^1_X\otimes_{\bb{Z}}\bb{Q})\to\mathscr{Z}^1_{X'}\otimes_{\bb{Z}}\bb{Q}.
  \tag{21.10.11.1}
\]
\end{env}

\begin{remark}[21.10.12]
\label{IV.21.10.12}
Cuando se está en la situación de \sref{IV.21.10.10}, puede ocurrir efectivamente que, para ciclos $1$-codimensionales $Z$ sobre $X$, $f^*(Z)$ sea un ciclo $1$-codimensional con coeficientes \emph{no enteros}; dicho de otro modo, los números $\mu_{x'}$ pueden no ser enteros.
Se obtiene un ejemplo tomando el anillo íntegro completo $A$ de $(\mathbf{0}_{\mathrm{IV}},6.15.11)$, (ii), y su clausura integral $A'$: el punto cerrado $x'$ de $\Spec(A')$ satisface la condición (iv) de \sref{IV.21.10.10} y se tiene $\mu_{x'}=\frac12$.
\end{remark}

\begin{lemma}[21.10.13]
\label{IV.21.10.13}
Sean $A$ un anillo local noetheriano de dimensión $1$ y $t$ un elemento regular de $A$ perteneciente al ideal maximal $\mathfrak{m}$ (lo que implica que $\mathfrak{m}\notin\operatorname{Ass}(A)$).
\begin{enumerate}[label=(\roman*)]
  \item Para todo $A$-módulo $M$ de tipo finito, el núcleo $N_t(M)$ y el conúcleo $P_t(M)$ de la homotecia $t_M:M\to M$ de razón $t$ tienen longitud finita. Se pone $d_t(M)=\operatorname{long}P_t(M)-\operatorname{long}N_t(M)$.
  \item Si $0\to M'\to M\to M''\to0$ es una sucesión exacta de $A$-módulos de tipo finito, se tiene $d_t(M)=d_t(M')+d_t(M'')$.
  \item Se tiene $d_t(M)\geq0$ para todo $A$-módulo $M$ de tipo finito; para que $d_t(M)=0$, es necesario y suficiente que $M$ sea de longitud finita.
  \item Si $K$ es el anillo total de fracciones de $A$ y si $M\otimes_AK$ es un $K$-módulo libre de rango $n$, se tiene $d_t(M)=n\cdot d_t(A)=n\cdot\operatorname{long}(A/tA)$.
  \item Si $M$ satisface las hipótesis de (iv) y si, además, $t$ es $M$-regular, se tiene $\operatorname{long}(M/tM)=n\cdot\operatorname{long}(A/tA)$.
\end{enumerate}
\end{lemma}

\begin{proof}
(i) $\Spec(A)$ está formado por el punto $\mathfrak{m}$ y por los ideales primos minimales $\mathfrak{p}_i$;
como, por hipótesis, $t\notin\mathfrak{p}_i$ para todo $i$ (Bourbaki, \emph{Alg. comm.}, cap.~IV, \textsection~1, n\textsuperscript{o}~1, cor.~3 de la prop.~2), la imagen de $t$ en cada $A_{\mathfrak{p}_i}$ es invertible y, por tanto, los soportes de los $A$-módulos de tipo finito $N_t(M)$ y $P_t(M)$ son vacíos o están reducidos a $\mathfrak{m}$;
de $(\mathbf{0},16.1.10)$ se concluye que estos módulos tienen longitud finita.

(ii) Puesto que $t$ es regular, se tiene una sucesión exacta
\[
  0\to A\xrightarrow{t_A}A\to A/tA\to0
\]
\oldpage[IV-4]{295}
y puesto que $\operatorname{Tor}_i^A(M,A)=0$ para $i\geq1$, la sucesión exacta de los Tor da, por una parte, la sucesión exacta
\[
  0\to\operatorname{Tor}_1^A(M,A/tA)\to M\xrightarrow{t_M}M\to M/tM\to0
\]
y, por otra parte, para $i\geq2$,
\[
  0=\operatorname{Tor}_i^A(M,A)\to\operatorname{Tor}_i^A(M,A/tA)\to\operatorname{Tor}_{i-1}^A(M,A)=0,
\]
por tanto, se tiene $N_t(M)=\operatorname{Tor}_1^A(M,A/tA)$ y $\operatorname{Tor}_i^A(M,A/tA)=0$ para $i\geq2$;
la sucesión exacta de los Tor da una sucesión exacta
\[
  0\to\operatorname{Tor}_1^A(M',A/tA)\to\operatorname{Tor}_1^A(M,A/tA)\to\operatorname{Tor}_1^A(M'',A/tA)\to
\]
\[
  M'\otimes_A(A/tA)\to M\otimes_A(A/tA)\to M''\otimes_A(A/tA)\to0
\]
y esta sucesión se escribe, según lo anterior,
\[
  0\to N_t(M')\to N_t(M)\to N_t(M'')\to P_t(M')\to P_t(M)\to P_t(M'')\to0,
\]
lo que demuestra (ii).

Para demostrar (iii), observemos que existe una serie de composición $(M_h)$ de $M$ cuyos cocientes $M_h/M_{h+1}$ son isomorfos a $A/\mathfrak{m}$ o a uno de los $A/\mathfrak{p}_i$ (Bourbaki, \emph{Alg. comm.}, cap.~IV, \textsection~1, n\textsuperscript{o}~4, teorema~1), y para que $M$ sea de longitud finita es necesario y suficiente que todos estos cocientes sean isomorfos a $A/\mathfrak{m}$.
Todo se reduce, pues, en virtud de (ii), a demostrar que $d_t(A/\mathfrak{m})=0$ y $d_t(A/\mathfrak{p}_i)>0$.
Ahora bien, como la imagen de $t$ en $A/\mathfrak{m}$ es $0$, se tiene $N_t(A/\mathfrak{m})=A/\mathfrak{m}$ y $P_t(A/\mathfrak{m})=A/\mathfrak{m}$, de donde resulta la primera afirmación;
por otra parte, la imagen de $t$ en $A/\mathfrak{p}_i$ es regular, de modo que $N_t(A/\mathfrak{p}_i)=0$ y $P_t(A/\mathfrak{p}_i)=A/(tA+\mathfrak{p}_i)$, que no está reducido a $0$, de donde resulta la segunda afirmación.

(iv) Existe una base de $M\otimes_AK$ de la forma $(x_j/s)_{1\leq j\leq n}$, donde $s$ es un elemento regular de $A$ y $x_j\in M$.
Consideremos el homomorfismo $u:A^n\to M$ que transforma el elemento $e_j$ $(1\leq j\leq n)$ de la base canónica de $A^n$ en $x_j$, y demostremos que $\operatorname{Ker}(u)$ y $\operatorname{Coker}(u)$ tienen \emph{longitud finita};
en efecto, para todo $i$, la imagen de $s$ en $A_{\mathfrak{p}_i}$ es invertible y, como $K_{\mathfrak{p}_i}=A_{\mathfrak{p}_i}$, las imágenes de los $x_j/s$ en $M_{\mathfrak{p}_i}$ forman una base de este $A_{\mathfrak{p}_i}$-módulo;
por tanto, $u_{\mathfrak{p}_i}:A_{\mathfrak{p}_i}^n\to M_{\mathfrak{p}_i}$ es biyectivo.
Como en (i), se concluye que los soportes de $\operatorname{Ker}(u)$ y $\operatorname{Coker}(u)$ están contenidos en $\{\mathfrak{m}\}$ y, por tanto, que estos módulos tienen longitud finita.
Así pues, de (ii) y (iii) resulta que $d_t(M)=d_t(A^n)=n\cdot d_t(A)=n\cdot\operatorname{long}(A/tA)$, puesto que $t$ es regular.

Por último, está claro que (v) se deduce inmediatamente de (iv), pues entonces $N_t(M)=0$.
\end{proof}

Este lema permite generalizar \sref{IV.21.10.4}:

\begin{proposition}[21.10.13]
\label{IV.21.10.13bis}
Supongamos que $f$ satisface las condiciones de \sref{IV.21.10.10}.
Entonces, para todo divisor $D$ sobre $X$ tal que $\Supp(D)\subset T$ y esté definido $f^*(D)$, se tiene
\[
\label{IV.21.10.13.1}
  \operatorname{cyc}(f^*(D))=f^*(\operatorname{cyc}(D)).
  \tag{21.10.13.1}
\]
\end{proposition}

\begin{proof}
Razonando como en \sref{IV.21.10.4}, todo se reduce a ver, con las mismas notaciones, que si $x'$ está en el caso $4\textsuperscript{o}$ de \sref{IV.21.10.10} y si $n_y$ es la longitud de $\sh{O}_{X,y}/t_y\sh{O}_{X,y}$, entonces la
\oldpage[IV-4]{296}
longitud $n_{x'}$ de $\sh{O}_{X',x'}/t'_{x'}\sh{O}_{X',x'}$ es igual a $\mu_{x'}n_y$, siendo $\mu_{x'}$ el número racional definido en \sref{IV.21.10.10}.
Como $\sh{O}_{X',x'}/t'_{x'}\sh{O}_{X',x'}$ tiene longitud finita, tiene la misma longitud que su completado $\mathfrak{m}_{x'}$-preádico $A'/t'_{x'}A'$, que también puede escribirse $A'/t_yA'$;
además, como $t'_{x'}$ es regular por hipótesis en $\sh{O}_{X',x'}$, también lo es en $A'$ por planitud $(\mathbf{0}_{\mathrm{I}},6.3.4)$ y, cuando $A'$ se considera como $A$-módulo, puede decirse también que $t_y$ es $A'$-regular.
Como $A$ tiene dimensión $1$ y $A'$ es un $A$-módulo de tipo finito tal que $A'\otimes_AK$ es un $K$-módulo libre de rango $r_{x'}$, se puede aplicar \sref{IV.21.10.13}, (v), a $M=A'$ y a $t_y$, y la longitud de $A'/t_yA'$ como $A$-módulo es, por tanto, $r_{x'}n_y$.
Como $\boldsymbol{k}(x')$ es un $\boldsymbol{k}(x)$-espacio vectorial de rango $q_{x'}$, la longitud de $A'/t_yA'$ como $A'$-módulo es, por tanto, $r_{x'}n_y/q_{x'}=\mu_{x'}n_y$, lo que termina la demostración.
\end{proof}

\begin{env}[21.10.14]
\label{IV.21.10.14}
Sean ahora $X$, $X'$ dos preesquemas localmente noetherianos y $f:X'\to X$ un morfismo que tiene las \emph{dos propiedades} siguientes:
\begin{enumerate}[label=\emph{\alph*)}]
  \item $f$ es finito;
  \item la imagen por $f$ de todo punto maximal de $X'$ es un punto maximal de $X$.
\end{enumerate}
Para todo $x\in X^{(1)}$, todos los puntos $x'\in f^{-1}(x)$ pertenecen entonces a $X'^{(1)}$, como resulta de la hipótesis \emph{b)} y de la desigualdad $(\mathbf{0},16.3.9.1)$, por ser discreta la fibra $f^{-1}(x)$.
Sea entonces
\[
  Z'=\sum_{x'\in X'^{(1)}}n_{x'}\cdot\overline{\{x'\}}
\]
un ciclo $1$-codimensional sobre $X'$.
Para todo $x\in X^{(1)}$, definamos un entero $n_x$ mediante la fórmula
\[
  n_x=\sum_{x'\in f^{-1}(x)}n_{x'}\cdot[\boldsymbol{k}(x'): \boldsymbol{k}(x)],
\]
lo que tiene sentido, pues los puntos de $f^{-1}(x)$ son finitos y $\boldsymbol{k}(x')$ es un cuerpo de grado finito sobre $\boldsymbol{k}(x)$ \sref[I]{I.6.4.4}.
Además, el conjunto de los $x\in X^{(1)}$ tales que $n_x\neq0$ es \emph{localmente finito} en $X$, pues está contenido en la imagen por $f$ del conjunto de los $x'\in X'^{(1)}$ tales que $n_{x'}\neq0$, y la conclusión resulta de que el morfismo $f$ es cuasicompacto.
Se puede, pues, definir un ciclo $1$-codimensional sobre $X$ poniendo
\[
\label{IV.21.10.14.1}
  f_*(Z')=\sum_{x\in X^{(1)}}n_x\cdot\overline{\{x\}}
  \tag{21.10.14.1}
\]
y se dice que $f_*(Z')$ es la \emph{imagen directa} de $Z'$ por $f$.
Está claro que la aplicación $f_*:\mathfrak{Z}^1(X')\to\mathfrak{Z}^1(X)$ así definida es un homomorfismo de grupos ordenados.
Además, si $U$ es un abierto de $X$ y $f_U:f^{-1}(U)\to U$ es la restricción de $f$, de las definiciones resulta inmediatamente que
\[
\label{IV.21.10.14.2}
  (f_U)_*(Z'|f^{-1}(U))=f_*(Z')|U,
  \tag{21.10.14.2}
\]
por tanto, denotando por $\psi$ la aplicación continua subyacente al morfismo $f$, las aplicaciones $\Gamma(f^{-1}(U),\mathscr{Z}^1_{X'})\to\Gamma(U,\mathscr{Z}^1_X)$ que acabamos de definir determinan un homomorfismo de haces de grupos conmutativos ordenados sobre $X$
\[
  \psi_*(\mathscr{Z}^1_{X'})\to\mathscr{Z}^1_X.
\]
\end{env}

\begin{proposition}[21.10.15]
\label{IV.21.10.15}
\oldpage[IV-4]{297}
Sean $X$, $X'$, $X''$ tres preesquemas localmente noetherianos y $f:X'\to X$, $g:X''\to X'$ dos morfismos que satisfacen las condiciones \emph{a)} y \emph{b)} de \sref{IV.21.10.14}.
Entonces $f\circ g$ satisface las mismas condiciones y, para todo ciclo $1$-codimensional $Z''$ sobre $X''$, se tiene
\[
  (f\circ g)_*(Z'')=f_*(g_*(Z'')).
\]
\end{proposition}

\begin{proof}
Esto resulta inmediatamente de las definiciones.
\end{proof}

\begin{proposition}[21.10.16]
\label{IV.21.10.16}
Sean $X$, $X'$, $X_1$ tres preesquemas localmente noetherianos, $f:X'\to X$ un morfismo que satisface las condiciones \emph{a)} y \emph{b)} de \sref{IV.21.10.14}, y $g:X_1\to X$ un morfismo plano.
Se pone $X'_1=X'\times_XX_1$ (de modo que $X'_1$ es localmente noetheriano) y se denotan por $f_1:X'_1\to X_1$ y $g_1:X'_1\to X'$ las proyecciones canónicas.
Entonces $f_1$ satisface las condiciones \emph{a)} y \emph{b)} de \sref{IV.21.10.14} y, para todo ciclo $1$-codimensional $Z'$ sobre $X'$, se tiene
\[
\label{IV.21.10.16.1}
  g^*(f_*(Z'))=(f_1)_*(g_1^*(Z')).
  \tag{21.10.16.1}
\]
\end{proposition}

\begin{proof}
Está claro que $f_1$ es finito y satisface la condición \emph{b)} de \sref{IV.21.10.14} en virtud de \sref{IV.2.3.7}.
Para demostrar \sref{IV.21.10.16.1}, nos reducimos inmediatamente al caso en que $X$, $X'$ y $X_1$ son los espectros de anillos locales noetherianos de dimensión $1$, $A$, $A'$ y $A_1$, siendo $A'$ un $A$-módulo finito y $A_1$ un $A$-módulo plano.
Denotando por $x$, $x'$ y $x_1$ los puntos cerrados de $X$, $X'$ y $X_1$, respectivamente, se trata de ver que
\[
\label{IV.21.10.16.2}
  \sum_{x'_1}\lambda_{x'_1}[\boldsymbol{k}(x'_1):\boldsymbol{k}(x_1)]
  =\lambda_{x_1}[\boldsymbol{k}(x'):\boldsymbol{k}(x)]
  \tag{21.10.16.2}
\]
donde $x'_1$ recorre el conjunto de los puntos cerrados de $X'_1$ (es decir, el conjunto de los puntos situados a la vez sobre $x'$ y sobre $x_1$), y donde $\lambda_{x'_1}=\operatorname{long}(\sh{O}_{X'_1,x'_1}/\mathfrak{m}_{x'}\sh{O}_{X'_1,x'_1})$ y $\lambda_{x_1}=\operatorname{long}(A_1/\mathfrak{m}_xA_1)$.
El primer miembro de \sref{IV.21.10.16.2} se escribe también
\[
  \operatorname{long}_{A_1}(A'_1/\mathfrak{m}_{x'}A'_1),
\]
donde se ha puesto $A'_1=A'\otimes_AA_1$.
Se tiene, pues, $A'_1/\mathfrak{m}_{x'}A'_1=(A'/\mathfrak{m}_{x'}A')\otimes_AA_1$ y, como $A_1$ es un $A$-módulo plano, por \sref{IV.4.7.1} se tiene
\[
  \operatorname{long}_{A_1}(A'_1/\mathfrak{m}_{x'}A'_1)
  =\operatorname{long}_{A}(A'/\mathfrak{m}_{x'}A')\operatorname{long}_{A_1}(A_1/\mathfrak{m}_xA_1)
  =[\boldsymbol{k}(x'): \boldsymbol{k}(x)]\cdot\lambda_{x_1},
\]
lo que termina la demostración.
\end{proof}

\begin{proposition}[21.10.17]
\label{IV.21.10.17}
Sean $X$, $X'$ dos preesquemas localmente noetherianos y $f:X'\to X$ un morfismo finito localmente libre.
Entonces $f$ satisface la condición \emph{b)} de \sref{IV.21.10.14} y, para todo divisor $D'$ sobre $X'$, se tiene
\[
\label{IV.21.10.17.1}
  f_*(\operatorname{cyc}(D'))=\operatorname{cyc}(f_*(D')).
  \tag{21.10.17.1}
\]
\end{proposition}

\begin{proof}
Como $f$ es plano y finito, la condición \emph{b)} de \sref{IV.21.10.14} y la relación $f(X'^{(1)})\subset X^{(1)}$ resultan de \sref{IV.6.1.1}.
La definición \sref{IV.21.10.14.1} muestra que podemos reducirnos al caso en que $X=\Spec(A)=\Spec(\sh{O}_{X,x})$ para un $x\in X^{(1)}$;
entonces $X'=\Spec(A')$, donde $A'$ es una
\oldpage[IV-4]{298}
$A$-álgebra que es un $A$-módulo libre de rango finito;
además, podemos suponer que $D'=\operatorname{div}(t')$, donde $t'$ es un elemento regular de $A'$;
entonces se tiene $f_*(D')=\operatorname{div}(t)$, donde $t=\mathrm{N}_{A'/A}(t')$ es un elemento regular de $A$ \sref{IV.21.5.2}.
Podemos limitarnos al caso en que $t'$ no es invertible en $A'$, lo que equivale a que $t$ no lo sea en $A$;
el anillo $A/tA$ tiene entonces longitud finita y es $\neq0$, y $A'/tA'$ es un producto directo de anillos locales artinianos $A'_i$ $(1\leq i\leq r)$, siendo el cuerpo residual de $A'_i$ igual a $\boldsymbol{k}(x'_i)$, donde los $x'_i$ $(1\leq i\leq r)$ son los puntos de $X'$ situados sobre $x$.
Si $t'_i$ $(1\leq i\leq r)$ es la imagen de $t'$ en $A'_i$, $A'/t'A'$ es el producto directo de los $A'_i/t'_iA'_i$;
como el producto $\operatorname{long}_{A'_i}(A'_i/t'_iA'_i)\cdot[\boldsymbol{k}(x'_i):\boldsymbol{k}(x)]$ es igual a $\operatorname{long}_A(A'_i/t'_iA'_i)$, se ve que la multiplicidad en el punto $x$ del primer miembro de \sref{IV.21.10.17.1} es $\operatorname{long}_A(A'/t'A')$, de modo que la fórmula que hay que demostrar se reduce a
\[
\label{IV.21.10.17.2}
  \operatorname{long}_A(A'/t'A')=\operatorname{long}_A(A/tA).
  \tag{21.10.17.2}
\]

Esta relación resulta del lema más general siguiente.
\end{proof}

\begin{lemma}[21.10.17.3]
\label{IV.21.10.17.3}
Sean $A$ un anillo local noetheriano de dimensión $1$, $M$ un $A$-módulo libre de rango finito y $u$ un endomorfismo inyectivo de $M$.
Entonces se tiene
\[
\label{IV.21.10.17.4}
  \operatorname{long}_A(\operatorname{Coker}u)=\operatorname{long}_A(A/(\det u)A).
  \tag{21.10.17.4}
\]
\end{lemma}

\begin{proof}
Distinguamos varios casos.

I) $A$ es un anillo de valoración discreta.
Sea, en efecto, $\pi$ un uniformizador de $A$, y observemos que $\operatorname{long}_A(A/\pi^kA)=k$;
si $n$ es el rango de $M$ y si $\pi^{m_i}$ $(1\leq i\leq n)$ son los factores invariantes de $u$, $\operatorname{Coker}(u)$ es la suma directa de los $A$-módulos $A/\pi^{m_i}A$, por lo que tiene longitud $m=\sum_{i=1}^nm_i$, y $\det(u)$ es el producto de un elemento invertible por $\pi^m$, de donde resulta la conclusión en este caso (Bourbaki, \emph{Alg.}, cap.~VII, \textsection~4, n\textsuperscript{o}~5, cor.~1 de la prop.~4).

II) $A$ es un anillo íntegro completo (de dimensión $1$).
Se sabe entonces por $(\mathbf{0},19.8.8)$, (ii), que existe un subanillo $B$ de $A$, que es un anillo de valoración discreta, tal que $B\to A$ es un homomorfismo local que hace de $A$ un $B$-módulo de tipo finito;
como este $B$-módulo es evidentemente sin torsión, es libre (Bourbaki, \emph{Alg. comm.}, cap.~VI, \textsection~3, n\textsuperscript{o}~6, lema~1).
Denotemos por $M'$ el conjunto $M$ provisto de su estructura de $B$-módulo (libre), y por $u'$ el endomorfismo $u$ considerado como $B$-endomorfismo.
De I) resulta que
\[
\label{IV.21.10.17.5}
  \operatorname{long}_B(\operatorname{Coker}u')=\operatorname{long}_B(B/(\det u')B).
  \tag{21.10.17.5}
\]
Pero se tiene (Bourbaki, \emph{Alg.}, cap.~VIII, \textsection~12, n\textsuperscript{o}~2, prop.~7)
\[
  \det u'=\mathrm{N}_{A/B}(\det u),
\]
por tanto, aplicando \sref{IV.21.10.17.4} a la homotecia $x\mapsto(\det u)x$ del $B$-módulo libre $A$, se obtiene
\[
  \operatorname{long}_B(B/(\det u')B)=\operatorname{long}_B(A/(\det u)A),
\]
de donde, sustituyendo en \sref{IV.21.10.17.5} y dividiendo por $[\boldsymbol{k}(A):\boldsymbol{k}(B)]$, longitud del $B$-módulo $\boldsymbol{k}(A)$, se obtiene \sref{IV.21.10.17.4} en el caso considerado.

\oldpage[IV-4]{299}
III) $A$ es un anillo completo.
Observemos que podemos suponer, además, que $\mathfrak{m}\notin\operatorname{Ass}(A)$;
en efecto, puesto que $\det(u)$ es un elemento regular de $A$, si se tuviera $\mathfrak{m}\in\operatorname{Ass}(A)$, se deduciría que $\det(u)\notin\mathfrak{m}$, de modo que $\det(u)$ sería invertible, $u$ sería un automorfismo de $M$, y la fórmula \sref{IV.21.10.17.4} sería entonces trivial, pues sus dos miembros serían nulos.

En lo que sigue, para un endomorfismo $v$ de un módulo $N$ sobre un anillo $R$, tal que $\operatorname{Ker}v$ y $\operatorname{Coker}v$ tengan longitud finita, se pondrá
\[
  \chi(N,v)=\operatorname{long}_R(\operatorname{Ker}(v))-\operatorname{long}_R(\operatorname{Coker}(v)).
\]
Obsérvese que la hipótesis sobre $v$ equivale a decir que el complejo
\[
  K^0:0\to N\xrightarrow{v}N\to0
\]
tiene sus módulos de cohomología de longitud finita y que $\chi(N,v)=\chi(H^0(K^0))$ $(\mathbf{0}_{\mathrm{III}},11.10)$.
De ello se deduce que, si $N'$, $N$, $N''$ son tres $R$-módulos y se tiene un diagrama conmutativo
\[
  \xymatrix{
    0 \ar[r] & N' \ar[r]^r \ar[d]^{v'} & N \ar[r]^s \ar[d]^v & N'' \ar[r] \ar[d]^{v''} & 0\\
    0 \ar[r] & N' \ar[r]_r & N \ar[r]_s & N'' \ar[r] & 0
  }
\]
cuyas filas son exactas, y si dos de los tres números $\chi(N,v)$, $\chi(N',v')$ y $\chi(N'',v'')$ están definidos, entonces también lo está el tercero y se tiene
\[
\label{IV.21.10.17.6}
  \chi(N,v)=\chi(N',v')+\chi(N'',v'').
  \tag{21.10.17.6}
\]
En efecto, esto resulta inmediatamente de la sucesión exacta de cohomología.

Por último, si $N$ es un $R$-módulo de longitud finita, se tiene $\chi(N,v)=0$.

Con estas notaciones se tiene el lema siguiente.
\end{proof}

\begin{lemma}[21.10.17.7]
\label{IV.21.10.17.7}
Sea $A$ un anillo local noetheriano de dimensión $1$ cuyo ideal maximal $\mathfrak{m}$ es tal que $\mathfrak{m}\notin\operatorname{Ass}(A)$;
sean $\mathfrak{p}_i$ $(1\leq i\leq n)$ los ideales primos minimales de $A$.
Sea $M$ un $A$-módulo libre de tipo finito y $u$ un endomorfismo de $M$ tal que $\chi(M,u)$ está definido;
para cada $i$, pongamos $M_i=M\otimes_A(A/\mathfrak{p}_i)$ y sea $u_i$ el endomorfismo $u\otimes1_{A/\mathfrak{p}_i}$ de $M_i$;
entonces, si $\chi(M_i,u_i)$ está definido para cada $i$, se tiene
\[
\label{IV.21.10.17.8}
  \chi(M,u)=\sum_{i=1}^n\operatorname{long}(A_{\mathfrak{p}_i})\cdot\chi(M_i,u_i).
  \tag{21.10.17.8}
\]
\end{lemma}

\begin{proof}
Puesto que $\mathfrak{m}\notin\operatorname{Ass}(A)$, se tiene una descomposición primaria reducida única $(0)=\bigcap_{1\leq i\leq n}\mathfrak{q}_i$, donde el ideal $\mathfrak{q}_i$ es $\mathfrak{p}_i$-primario para $1\leq i\leq n$.
Si se pone $M'_i=M\otimes_A(A/\mathfrak{q}_i)$, se tiene, pues, una sucesión exacta
\[
  0\to M\to\bigoplus_iM'_i\to M''\to0
\]
de $A$-módulos, donde $M''$ tiene longitud finita:
en efecto, localizando la sucesión exacta anterior en cada uno de los $\mathfrak{p}_i$, se obtiene $M''_{\mathfrak{p}_i}=0$, pues $(\mathfrak{q}_i)_{\mathfrak{p}_i}=0$ y $(\mathfrak{q}_j)_{\mathfrak{p}_i}=A_{\mathfrak{p}_i}$ para $j\neq i$ (Bourbaki, \emph{Alg. comm.}, cap.~IV, \textsection~2, n\textsuperscript{o}~4, prop.~6), de modo que $(M'_i)_{\mathfrak{p}_i}=M_{\mathfrak{p}_i}$ y $(M'_j)_{\mathfrak{p}_i}=0$;
\oldpage[IV-4]{300}
como el soporte de $M''$ está, pues, reducido a $\mathfrak{m}$, $M''$ tiene longitud finita $(\mathbf{0},16.1.10)$.
Si se pone $u'_i=u\otimes1_{A/\mathfrak{q}_i}$, y si $u''$ es el endomorfismo de $M''$ deducido de $\bigoplus u'_i$ pasando a los cocientes, de \sref{IV.21.10.17.6} se deduce que $\chi(M,u)+\chi(M'',u'')=\sum_i\chi(M'_i,u'_i)$ y, puesto que $M''$ tiene longitud finita, $\chi(M'',u'')=0$.
Para demostrar \sref{IV.21.10.17.8}, podemos, pues, reducirnos al caso en que $n=1$.
Denotaremos entonces por $\mathfrak{p}$ el único ideal primo minimal, que es el nilradical de $A$;
si $A_0=A/\mathfrak{p}$, $M_0=M\otimes_AA_0$, $u_0=u\otimes1_{A_0}$, se trata de demostrar que, si $\chi(M_0,u_0)$ está definido, se tiene
\[
\label{IV.21.10.17.9}
  \chi(M,u)=\operatorname{long}A_{\mathfrak{p}}\cdot\chi(M_0,u_0).
  \tag{21.10.17.9}
\]

Sea $\mathfrak{n}_j$ $(0\leq j\leq r)$ la «potencia simbólica $j$-ésima» de $\mathfrak{p}$, imagen inversa en $A$ del ideal $(\mathfrak{p}A_{\mathfrak{p}})^j$ de $A_{\mathfrak{p}}$ $(1\leq j\leq r)$, con $\mathfrak{n}_0=A$ y $\mathfrak{n}_r=0$;
pongamos
\[
  M_j=\mathfrak{n}_jM/\mathfrak{n}_{j+1}M\qquad(0\leq j\leq r-1),
\]
y denotemos por $v_j$ el endomorfismo de $M_j$ deducido de $u$ por restricción a $\mathfrak{n}_jM$ y paso a los cocientes.
Vamos a demostrar primero que cada uno de los números $\chi(M_j,v_j)$ está definido y que se tiene
\[
\label{IV.21.10.17.10}
  \chi(M,u)=\sum_j\chi(M_j,v_j).
  \tag{21.10.17.10}
\]
La primera afirmación implicará la segunda, aplicando \sref{IV.21.10.17.6} a cada una de las sucesiones exactas
\[
  0\to\mathfrak{n}_jM/\mathfrak{n}_{j+1}M\to M/\mathfrak{n}_{j+1}M\to M/\mathfrak{n}_jM\to0.
\]

Para demostrar la primera afirmación, observemos que, si $m$ es el rango del $A$-módulo libre $M$, $M_j$ es isomorfo a $(\mathfrak{n}_j/\mathfrak{n}_{j+1})^m$; dicho de otro modo, puesto que $\mathfrak{p}$ anula cada uno de los cocientes $\mathfrak{n}_j/\mathfrak{n}_{j+1}$, $M_j$ es un $A_0$-módulo isomorfo a $M_0\otimes_{A_0}(\mathfrak{n}_j/\mathfrak{n}_{j+1})$.
Denotemos por $l_j$ el rango del $A_0$-módulo $\mathfrak{n}_j/\mathfrak{n}_{j+1}$;
como el cuerpo de fracciones $K_0$ de $A_0$ es el cuerpo residual de $A_{\mathfrak{p}}$, $l_j$ es también la longitud del $A_{\mathfrak{p}}$-módulo $(\mathfrak{p}A_{\mathfrak{p}})^j/(\mathfrak{p}A_{\mathfrak{p}})^{j+1}$.
Existe un sistema de generadores del $A_0$-módulo $M_j$ que contiene una base de $M_j\otimes_{A_0}K_0$;
por tanto, existe un $A_0$-homomorfismo
\[
  w_j:M_0^{l_j}\to M_j
\]
cuya localización en el ideal $(0)$ de $A_0$ es un isomorfismo, de modo que el soporte de $\operatorname{Ker}(w_j)$ y de $\operatorname{Coker}(w_j)$ está reducido al ideal maximal $\mathfrak{m}/\mathfrak{p}$ de $A_0$;
$\operatorname{Ker}(w_j)$ y $\operatorname{Coker}(w_j)$ son, por tanto, $A_0$-módulos de longitud finita $(\mathbf{0},16.1.10)$.
Como por hipótesis $\chi(M_0,u_0)$ está definido, también lo está $\chi(M_0^{l_j},u_0^{l_j})=l_j\chi(M_0,u_0)$ y, en virtud de \sref{IV.21.10.17.6} y del hecho de que $\operatorname{Ker}(w_j)$ y $\operatorname{Coker}(w_j)$ tienen longitud finita, se ve que $\chi(M_j,v_j)$ está definido y es igual a $l_j\chi(M_0,u_0)$;
la relación \sref{IV.21.10.17.10} da entonces
\[
  \chi(M,u)=\left(\sum_jl_j\right)\chi(M_0,u_0),
\]
y, en virtud de una observación anterior, $\sum_jl_j$ no es sino la longitud de $A_{\mathfrak{p}}$, lo que termina la demostración del lema \sref{IV.21.10.17.7}.
\end{proof}

\oldpage[IV-4]{301}
Para aplicar este lema cuando $A$ es un anillo completo y $\mathfrak{m}\notin\operatorname{Ass}(A)$, se observa que, si $u$ es inyectivo, también lo son los $u_i$ (con las notaciones del lema):
en efecto, $\det u$ es entonces un elemento regular de $A$, por lo que no pertenece a ninguno de los $\mathfrak{p}_i$, y su imagen $\det u_i$ en $A/\mathfrak{p}_i$ es, por consiguiente, un elemento $\neq0$ de este anillo íntegro, lo que demuestra que $u_i$ es inyectivo.
Como $\operatorname{Coker}(u_i)$ es imagen de $\operatorname{Coker}(u)$, también tiene longitud finita y, por tanto, $\chi(M_i,u_i)$ está definido para todo $i$;
se tiene, por consiguiente, la fórmula \sref{IV.21.10.17.8}.
Por otra parte, puesto que $\det(u)$ es un elemento regular de $A$, no está contenido en ninguno de los $\mathfrak{p}_i$;
el ideal $(\det u)A$ es, pues, $\mathfrak{m}$-primario y el cociente $A/(\det u)A$ tiene longitud finita.
Aplicando el mismo razonamiento anterior a la homotecia inyectiva $t:\xi\mapsto(\det u)\xi$ de $A$ y a sus imágenes $t_i=\det u_i$ en los $A/\mathfrak{p}_i$, se obtiene
\[
  \chi(A,t)=\sum_i\operatorname{long}(A_{\mathfrak{p}_i})\cdot\chi(A/\mathfrak{p}_i,t_i).
\]
Pero los anillos $A/\mathfrak{p}_i$ son íntegros y completos y, aplicando el resultado de II) a cada uno de ellos, se obtiene de nuevo la relación \sref{IV.21.10.17.4} para $M$ y $u$.

IV) Caso general.
Pongamos $A'=\widehat{A}$, $M'=M\otimes_AA'$, $u'=u\otimes1_{A'}$;
se tiene $\det(u')=\det(u)$ y, por planitud, $\operatorname{Coker}(u')=(\operatorname{Coker}(u))_{(A')}$ y $A'/(\det u)A'=(A/(\det u)A)_{(A')}$;
como la fórmula \sref{IV.21.10.17.4} es verdadera para $A'$ y $u'$ en virtud de III), también lo es para $A$ y $u$ en virtud de \sref{IV.4.7.1}.

Esto termina la demostración de \sref{IV.21.10.17}.

\begin{proposition}[21.10.18]
\label{IV.21.10.18}
Sean $X$, $X'$ dos preesquemas localmente noetherianos y $f:X'\to X$ un morfismo finito que transforma todo punto maximal de $X'$ en un punto maximal de $X$ y satisface, para todo $x'\in X'^{(1)}$\footnote{Corrección de fuente: el francés impreso y el inglés mantenido dicen $x'\in X'$. Las condiciones (ii)--(iv) de \sref{IV.21.10.10} se formulan para $x'\in X'^{(1)}$; con la lectura impresa, ningún punto maximal de un $X'$ no vacío puede satisfacerlas.}, una de las condiciones (ii), (iii), (iv) de \sref{IV.21.10.10}.
Se supone además que existe un entero $n$ tal que, para todo punto maximal $x$ de $X$, $(f_*(\sh{O}_{X'}))_x$ es un $\sh{O}_{X,x}$-módulo libre de rango $n$.
Entonces, para todo ciclo $1$-codimensional $Z$ sobre $X$, se tiene
\[
\label{IV.21.10.18.1}
  f_*(f^*(Z))=n\cdot Z
  \tag{21.10.18.1}
\]
(«fórmula de proyección»).
\end{proposition}

\begin{proof}
En virtud de las definiciones, nos reducimos inmediatamente al caso en que $X$ es el espectro de un anillo local noetheriano $A$ de dimensión $1$, con punto cerrado $x$, y $Z=\overline{\{x\}}$;
pongamos $X'=\Spec(A')$, donde $A'$ es una $A$-álgebra finita y, para todo ideal primo minimal $\mathfrak{p}_i$ de $A$, $A'_{\mathfrak{p}_i}$ es un $A_{\mathfrak{p}_i}$-módulo libre de rango $n$.
Demostremos que podemos limitarnos, además, al caso en que $A$ es completo.
Hagamos, en efecto, el cambio de base $h:Y\to X$, donde $Y=\Spec(B)$, con $B=\widehat{A}$, y pongamos
\[
  Y'=X'\times_XY=\Spec(B\otimes_AA'),\qquad g=f_{\{Y\}}:Y'\to Y;
\]
el morfismo $g$ es entonces finito y, como $h$ es plano, los puntos maximales de $Y$ están situados sobre los de $X$;
sobre cada $\mathfrak{p}_i$ solo hay un número finito de ideales minimales $\mathfrak{p}_{ij}$ de $B$, y $(B\otimes_AA')_{\mathfrak{p}_{ij}}$ es un $B_{\mathfrak{p}_{ij}}$-módulo libre de rango $n$.
Por último, si $f$ satisface una de las condiciones (ii), (iii) o (iv) de \sref{IV.21.10.10} en cada punto $x'$ de $f^{-1}(x)$, $g$ satisface la condición correspondiente en el único punto $y'$ de $g^{-1}(y)$ situado sobre $x'$ (denotando por $y$ el punto cerrado de $Y$);
esto es inmediato para las condiciones (ii) y (iv);
para la condición (iii), esta implica que $A$ es un anillo de valoración discreta, por lo que lo mismo ocurre con $B$, y la condición
\oldpage[IV-4]{302}
$\mathfrak{m}_{y'}\notin\operatorname{Ass}(\sh{O}_{Y',y'})$ resulta, por planitud, de la condición $\mathfrak{m}_{x'}\notin\operatorname{Ass}(\sh{O}_{X',x'})$ $(\mathbf{0},3.3.1)$.
El morfismo $g$ satisface, pues, las mismas condiciones que $f$;
si se demuestra la fórmula \sref{IV.21.10.18.1} para $g$ y $\overline{\{y\}}$, la misma fórmula será válida para $f$ y $\overline{\{x\}}$ en virtud de \sref{IV.21.10.16.1}.

Podemos, pues, suponer que $A$ es \emph{completo};
entonces lo mismo ocurre con $A'$, que es, por tanto, un producto directo de anillos locales completos;
podemos, por consiguiente, limitarnos al caso en que $A'$ es un anillo \emph{local}, y queda verificar \sref{IV.21.10.18.1} en cada uno de los casos (ii), (iii), (iv) para el punto cerrado $x'$ de $X'$.
En el caso (ii), como $A'$ es un $A$-módulo plano de tipo finito, es un $A$-módulo \emph{libre} $(\mathbf{0}_{\mathrm{III}},10.1.3)$, de rango $n$ en virtud de la hipótesis.
Ahora bien, por definición \sref{IV.21.10.1} y \sref{IV.21.10.3}, se tiene
\[
  f^*(Z)=\lambda_{x'}\cdot\overline{\{x'\}},
\]
donde $\lambda_{x'}$ es la longitud del $A'$-módulo $A'/\mathfrak{m}_xA'$; entonces,
\[
  f_*(f^*(Z))=(\lambda_{x'}[\boldsymbol{k}(x'): \boldsymbol{k}(x)])\cdot\overline{\{x\}};
\]
pero $\lambda_{x'}[\boldsymbol{k}(x'): \boldsymbol{k}(x)]$ es la longitud de $A'/\mathfrak{m}_xA'=A'\otimes_A\boldsymbol{k}(x)$ como $A$-módulo, o también su rango como $\boldsymbol{k}(x)$-espacio vectorial, y por tanto es igual a $n$.

En el caso (iii), $A$ es un anillo de valoración discreta y, por tanto, regular, y la hipótesis $\mathfrak{m}_{x'}\notin\operatorname{Ass}(\sh{O}_{X',x'})$ implica que $A'$ es un anillo de Cohen--Macaulay $(\mathbf{0},16.4.6)$;
como $\dim(A')=\dim(A)=1$ y $A'/\mathfrak{m}_xA'$ es un anillo artiniano, de \sref{IV.6.1.5} resulta que $A'$ es un $A$-módulo plano y nos reducimos al caso (ii).

En el caso (iv), si $K$ es el anillo total de fracciones de $A$, $K'=A'\otimes_AK$ es por hipótesis un $K$-módulo libre de rango $n$ y, por definición \sref{IV.21.10.10}, se tiene
\[
  f^*(Z)=\frac{n}{[\boldsymbol{k}(x'): \boldsymbol{k}(x)]}\cdot\overline{\{x'\}},
\]
de donde $f_*(f^*(Z))=n\cdot\overline{\{x\}}$.
\end{proof}

Obsérvese que la fórmula \sref{IV.21.10.18.1} es aplicable, en particular, cuando el morfismo $f$ es \emph{finito} y \emph{plano} y tal que, para todo punto maximal $x$ de $X$, $(f_*(\sh{O}_{X'}))_x$ es un $\sh{O}_{X,x}$-módulo libre de rango $n$.

\begin{corollary}[21.10.19]
\label{IV.21.10.19}
Bajo las hipótesis de \sref{IV.21.10.18}, sea $D$ un divisor sobre $X$ tal que esté definido $f^*(D)$ \sref{IV.21.4.5}; entonces se tiene
\[
\label{IV.21.10.19.1}
  f_*(\operatorname{cyc}(f^*(D)))=n\cdot\operatorname{cyc}(D).
  \tag{21.10.19.1}
\]
\end{corollary}

\begin{proof}
Esto resulta de \sref{IV.21.10.18} y \sref{IV.21.10.13bis}.
\end{proof}


% Mechanical subsection components; prose is translated in the component files.
\subsection{Factorialidad de los anillos locales regulares}
\label{subsection:IV.21.11}
\label{IV.21.11}

\begin{theorem}[21.11.1]
\label{IV.21.11.1}
\textup{(Auslander-Buchsbaum)} ---
Un anillo local noetheriano regular es factorial.
\end{theorem}

\begin{proof}
La demostración que sigue se debe a I.~Kaplansky.

Sea $A$ un anillo local noetheriano regular de dimensión $n$; razonaremos
por inducción sobre $n$. Para $n=0$, $A$ es un cuerpo y, para $n=1$, un anillo de
valoración discreta, por tanto principal (y \emph{a fortiori} factorial). Supongamos, pues,
$n\geq2$ y demostrado el teorema para los anillos regulares de dimensión $<n$.
Pongamos $X=\operatorname{Spec}(A)$, denotemos por $a$ el punto cerrado de $X$ y pongamos
$U=X-\{a\}$. En todo punto $y\in U$, se tiene
$\dim(\mathcal O_{X,y})\leq n-1$ y, puesto que $A$ es regular, los anillos
$\mathcal O_{X,y}$ también lo son $(\mathbf 0, 17.3.2)$; por tanto, la hipótesis de
inducción implica que son factoriales. Además, se tiene
$\operatorname{prof}(A)=\dim(A)\geq2$, ya que $A$ es regular y, por tanto, de
Cohen--Macaulay $(\mathbf 0, 17.1.3)$.
\oldpage[IV-4]{303}
Utilizando (21.6.14), nos reducimos a demostrar que $\operatorname{Pic}(U)=0$.
Consideremos, pues, un $\mathcal O_U$-módulo invertible $\mathscr L$, y demostremos que es
isomorfo a $\mathcal O_U$. Resulta de $(\mathbf I, 9.4.5)$ que existe un
$\mathcal O_X$-módulo \emph{coherente} $\mathscr F$ tal que
$\mathscr F|U=\mathscr L$. Como $A$ es regular, y por tanto de dimensión cohomológica
finita $(\mathbf 0, 17.3.1)$, existe una resolución libre por la izquierda \emph{finita} de
$\mathscr F$:
\[
  0\leftarrow\mathscr F\leftarrow\mathcal O_X^{n_1}
  \leftarrow\mathcal O_X^{n_2}\leftarrow\cdots
  \leftarrow\mathcal O_X^{n_h}\leftarrow0
\]
$(\mathbf 0, 17.2.8$ y $\mathbf 0, 17.2.2$, (iii)). Por restricción a $U$, se tiene
entonces una resolución finita
\[
\label{IV.21.11.1.1}
  0\leftarrow\mathscr L\leftarrow\mathcal O_U^{n_1}
  \leftarrow\mathcal O_U^{n_2}\leftarrow\cdots
  \leftarrow\mathcal O_U^{n_h}\leftarrow0.
  \tag{21.11.1.1}
\]
\end{proof}

El teorema resultará entonces de las consideraciones generales siguientes. Sea,
sobre un espacio anillado $X$, $\mathscr E$ un $\mathcal O_X$-módulo localmente libre de rango finito;
denotaremos por $\Lambda^{\max}\mathscr E$ el $\mathcal O_X$-módulo invertible que,
en un entorno de cada punto de $X$, es igual a $\Lambda^p\mathscr E$, donde
$p$ denota el rango de $\mathscr E$ en este entorno (que puede variar con la componente
conexa de $X$). Con esta notación, se tiene el

\begin{lemma}[21.11.1.2]
\label{IV.21.11.1.2}
Sean $X$ un espacio anillado en anillos locales, y
\[
  0\leftarrow\mathscr E_0\xleftarrow{u_0}\mathscr E_1
  \leftarrow\cdots\leftarrow\mathscr E_h\leftarrow0
\]
una sucesión exacta de $\mathcal O_X$-módulos localmente libres de rango finito; entonces el
$\mathcal O_X$-módulo invertible
\[
  \bigotimes_{0\leq i\leq h}
  (\Lambda^{\max}\mathscr E_i)^{\otimes(-1)^i}
\]
es isomorfo a $\mathcal O_X$.
\end{lemma}

\begin{proof}
Demostremos cómo este lema permitirá concluir la demostración de (21.11.1). Para ello basta
observar que, para todo entero $n$,
$\Lambda^{\max}(\mathcal O_U^n)=\Lambda^n\mathcal O_U^n$ es isomorfo a
$\mathcal O_U$, al igual que $\mathcal O_U^{\otimes(-1)}$. Como, por otra parte,
$\Lambda^{\max}\mathscr L=\mathscr L$ para todo $\mathcal O_U$-módulo invertible
$\mathscr L$, el lema (21.11.1.2), aplicado a la sucesión exacta (21.11.1.1), muestra
que $\mathscr L$ es isomorfo a $\mathcal O_U$.

Queda por demostrar (21.11.1.2); se procede por inducción sobre $h$, siendo trivial el lema
para $h=1$. Para $h>1$, $\mathscr N=\operatorname{Ker}(u_0)$ es un
$\mathcal O_X$-módulo localmente libre de rango finito $(\mathbf 0_{\mathrm I}, 5.5.5)$,
y se tienen las dos sucesiones exactas
\[
  0\leftarrow\mathscr E_0\leftarrow\mathscr E_1
  \leftarrow\mathscr N\leftarrow0
\]
\[
  0\leftarrow\mathscr N\leftarrow\mathscr E_2
  \leftarrow\cdots\leftarrow\mathscr E_h\leftarrow0.
\]
En virtud de la hipótesis de inducción,
$(\Lambda^{\max}\mathscr N)\otimes
(\bigotimes_{2\leq i\leq h}(\Lambda^{\max}\mathscr E_i)^{\otimes(-1)^{i-1}})$
es isomorfo a $\mathcal O_X$; por tanto, bastará definir un isomorfismo canónico
\[
  (\Lambda^{\max}\mathscr N)\otimes(\Lambda^{\max}\mathscr E_0)
  \xrightarrow{\sim}\Lambda^{\max}\mathscr E_1
\]
para terminar la demostración. Ahora bien, existe un recubrimiento abierto $(U_\alpha)$ de
$X$ tal que, en todo $U_\alpha$, $\mathscr E_1|U_\alpha$ sea suma directa de
$\mathscr N|U_\alpha$ y de un $\mathcal O_{U_\alpha}$-módulo localmente libre
$\mathscr M_\alpha$ $(\mathbf 0_{\mathrm I}, 5.5.5)$, de donde resulta un isomorfismo canónico
$v_\alpha:\mathscr M_\alpha\xrightarrow{\sim}\mathscr E_0|U_\alpha$.
Como se tiene un isomorfismo canónico
\[
  r_\alpha:(\Lambda^{\max}\mathscr N|U_\alpha)
  \otimes(\Lambda^{\max}\mathscr M_\alpha)
  \xrightarrow{\sim}(\Lambda^{\max}\mathscr E_1)|U_\alpha,
\]
se deduce de él mediante $v_\alpha$ un isomorfismo
\[
  u_\alpha:(\Lambda^{\max}\mathscr N|U_\alpha)
  \otimes(\Lambda^{\max}\mathscr E_0|U_\alpha)
  \xrightarrow{\sim}\Lambda^{\max}\mathscr E_1|U_\alpha.
\]
\oldpage[IV-4]{304}
Se trata de demostrar que $u_\alpha$ y $u_\beta$ coinciden en
$U_\alpha\cap U_\beta$ para dos índices cualesquiera $\alpha$, $\beta$. Ahora bien, si
$v'_\alpha$ y $v'_\beta$ son las restricciones a $U_\alpha\cap U_\beta$ de
$v_\alpha$ y $v_\beta$ respectivamente, se tiene
$v'_\alpha=v'_\beta\circ w_{\beta\alpha}$, donde
\[
  w_{\beta\alpha}:\mathscr M_\alpha|(U_\alpha\cap U_\beta)
  \xrightarrow{\sim}\mathscr M_\beta|(U_\alpha\cap U_\beta)
\]
es la «proyección paralela a $\mathscr N$» tal que, para toda sección
$s\in\Gamma(U_\alpha\cap U_\beta,\mathscr M_\alpha)$,
$w_{\beta\alpha}(s)=s+t$, con
$t\in\Gamma(U_\alpha\cap U_\beta,\mathscr N)$; la identidad de $u_\alpha$ y de
$u_\beta$ resulta inmediatamente de este hecho y de la definición del isomorfismo canónico
$r_\alpha$ (Bourbaki, \emph{Alg.}, cap.~III, 3\textsuperscript{a} ed.).
\end{proof}

\subsection{El teorema de pureza de Van der Waerden para el conjunto de
ramificación de un morfismo birracional}
\label{subsection:IV.21.12}
\label{IV.21.12}

\begin{env}[21.12.1]
\label{IV.21.12.1}
Sean $X$ y $U$ dos preesquemas, y sea $f:U\to X$ un morfismo cuasicompacto y
cuasiseparado, de modo que $f_*(\mathcal O_U)$ es una $\mathcal O_X$-álgebra
cuasicoherente (1.7.4). Se llama \emph{envolvente afín} del $X$-preesquema $U$ al
$X$-preesquema afín sobre $X$
\[
  U^0=\operatorname{Aff}(U/X)=\operatorname{Spec}(f_*(\mathcal O_U))
  =\operatorname{Spec}(\mathscr A(U))\qquad(\mathbf{II}, 1.1.1).
\]
Si $f^0:U^0\to X$ es el morfismo estructural, se tiene, por definición,
\[
  \mathscr A(U^0)=f^0_*(\mathcal O_{U^0})=f_*(\mathcal O_U)=\mathscr A(U),
\]
y al isomorfismo identidad de $f_*(\mathcal O_U)$ le corresponde, por
$(\mathbf{II}, 1.2.7)$, un $X$-morfismo canónico
\[
\label{IV.21.12.1.1}
  i_U:U\to U^0.
  \tag{21.12.1.1}
\]

Para que $i_U$ sea un \emph{isomorfismo}, es evidentemente necesario y suficiente que el
morfismo $f:U\to X$ sea \emph{afín}. Para todo $X$-preesquema $V$ \emph{afín sobre
$X$}, la aplicación $u\mapsto u\circ i_U$ es una \emph{biyección}
\[
  \operatorname{Hom}_X(U^0,V)\xrightarrow{\sim}\operatorname{Hom}_X(U,V)
\]
functorial en $V$: esto resulta de la existencia de las biyecciones canónicas
\[
  \operatorname{Hom}_X(U,V)\xrightarrow{\sim}
  \operatorname{Hom}_{\mathcal O_X}(\mathscr A(V),\mathscr A(U))
\]
y
$\operatorname{Hom}_X(U^0,V)\xrightarrow{\sim}
\operatorname{Hom}_{\mathcal O_X}(\mathscr A(V),\mathscr A(U^0))$
$(\mathbf{II}, 1.2.7)$.

Se puede decir, pues, que $U^0$ \emph{representa} el funtor covariante
$V\mapsto\operatorname{Hom}_X(U,V)$ en la categoría de los $X$-preesquemas afines sobre
$X$ $(\mathbf 0_{\mathrm{III}}, 8.1.11)$. Se deduce de ello
$(\mathbf 0_{\mathrm{III}}, 8.1.7)$ que, para $X$ fijo,
$U\mapsto\operatorname{Aff}(U/X)$ es un funtor covariante de la categoría de los
$X$-preesquemas cuasicompactos y cuasiseparados sobre $X$ en la categoría de los
$X$-preesquemas afines sobre $X$; más precisamente, si $U_1$, $U_2$ son dos
$X$-preesquemas cuasicompactos y cuasiseparados sobre $X$, a todo $X$-morfismo
$g:U_1\to U_2$ le corresponde el único $X$-morfismo $g^0:U_1^0\to U_2^0$ que hace
conmutativo el diagrama
\[
  \xymatrix{
    U_1 \ar[r]^g \ar[d]_{i_{U_1}} & U_2 \ar[d]^{i_{U_2}} \\
    U_1^0 \ar[r]_{g^0} & U_2^0.
  }
\]
\end{env}

\oldpage[IV-4]{305}
Más generalmente, consideremos un diagrama conmutativo
\[
  \xymatrix{
    U \ar[r]^v \ar[d]_f & U' \ar[d]^{f'} \\
    X \ar[r]_u & X'
  }
\]
donde los morfismos $f$, $f'$ son cuasicompactos y cuasiseparados y el morfismo $u$ es
\emph{afín}. Si se pone $h=u\circ f$, se tiene
$h_*(\mathcal O_U)=u_*(f_*(\mathcal O_U))=u_*(f^0_*(\mathcal O_{U^0}))$ y
$u\circ f^0$ es un morfismo afín; se tiene, por consiguiente,
$U^0=\operatorname{Aff}(U/X')$ (relativamente al morfismo $h$), de donde resulta un único
$X'$-morfismo $v^0:U^0\to U'^0$ que hace conmutativo el diagrama
\[
  \xymatrix{
    U \ar[r]^v \ar[d]_{i_U} & U' \ar[d]^{i_{U'}} \\
    U^0 \ar[r]_{v^0} & U'^0.
  }
\]

\begin{proposition}[21.12.2]
\label{IV.21.12.2}
Sean $f:U\to X$ un morfismo cuasicompacto y cuasiseparado y
$h:X'\to X$ un morfismo plano; se pone $U'=U\times_X X'$,
$f'=f_{\{X'\}}:U'\to X'$. Entonces se tiene un $X'$-isomorfismo canónico
\[
\label{IV.21.12.2.1}
  \operatorname{Aff}(U'/X')\xrightarrow{\sim}
  \operatorname{Aff}(U/X)\times_X X'.
  \tag{21.12.2.1}
\]
\end{proposition}

\begin{proof}
En efecto, se tiene
$\operatorname{Aff}(U'/X')=\operatorname{Spec}(f'_*(\mathcal O_{U'}))$, y
$\operatorname{Aff}(U/X)\times_X X'=\operatorname{Spec}(h^*(f_*(\mathcal O_U)))$
$(\mathbf{II}, 1.5.1)$; el isomorfismo del enunciado procede del isomorfismo canónico
$h^*(f_*(\mathcal O_U))\xrightarrow{\sim}f'_*(\mathcal O_{U'})$ (2.3.1).
\end{proof}

\begin{corollary}[21.12.3]
\label{IV.21.12.3}
Para todo morfismo cuasicompacto y cuasiseparado $f:U\to X$ y todo $x\in X$, se tiene,
salvo isomorfismo canónico,
\[
  U^0\times_X\operatorname{Spec}(\mathcal O_{X,x})
  =(U\times_X\operatorname{Spec}(\mathcal O_{X,x}))^0.
\]
\end{corollary}

Resulta asimismo de (21.12.2) que se tiene, salvo isomorfismo canónico, para todo
abierto $V$ de $X$,
\[
\label{IV.21.12.4}
  (f^{-1}(V))^0=(f^0)^{-1}(V).
  \tag{21.12.4}
\]

\begin{env}[21.12.5]
\label{IV.21.12.5}
Consideraremos, en particular, el caso en que $f:U\to X$ es una \emph{inmersión
abierta}, identificándose así $U$ con un preesquema inducido sobre un abierto de $X$.
Como el
morfismo $f^{-1}(U)\to U$ es la identidad, resulta de (21.12.4) que el morfismo
$(f^0)^{-1}(U)\to U$, restricción de $f^0$, es un isomorfismo, siendo entonces
$i_U:U\to U^0$ una \emph{inmersión abierta} que permite identificar $U$ con un
preesquema inducido sobre un abierto de $U^0$.

Más generalmente, para un abierto $V$ de $X$, la restricción $(f^0)^{-1}(V)\to V$ de
$f^0$ es un isomorfismo de $(f^0)^{-1}(V)$ sobre el preesquema inducido sobre el abierto
$V\cap U$ si y solo si la inmersión abierta $U\cap V\to V$ es un morfismo
\emph{afín}. Es claro $(\mathbf{II}, 1.2.1)$
\oldpage[IV-4]{306}
que la unión de estos abiertos $V$ es el \emph{mayor} de ellos, $U_1$, que
contiene a $U$; en virtud de lo precedente, $U_1$ es también el mayor abierto que no
interseca el conjunto
\[
  \operatorname{Daf}(U/X)=f^0(U^0)-U
\]
(«defecto de afinidad» del abierto $U$ relativamente a $X$, que es vacío si y solo
si $U$ es afín sobre $X$); en otras palabras, el conjunto cerrado $Z=X-U_1$ es la
\emph{adherencia} del conjunto $\operatorname{Daf}(U/X)$.

Obsérvese que, para todo morfismo \emph{plano} $h:X'\to X$, si se pone
$U'=h^{-1}(U)$, se tiene
\[
\label{IV.21.12.5.1}
  \operatorname{Daf}(U'/X')=h^{-1}(\operatorname{Daf}(U/X))
  \tag{21.12.5.1}
\]
como resulta inmediatamente de (21.12.2.1) y de $(\mathbf I, 3.4.8)$. En particular,
para todo abierto $V$ de $X$, se tiene
\[
\label{IV.21.12.5.2}
  \operatorname{Daf}((U\cap V)/V)=\operatorname{Daf}(U/X)\cap V
  \tag{21.12.5.2}
\]
y, para todo $x\in X$,
\[
\label{IV.21.12.5.3}
  \operatorname{Daf}((U\cap\operatorname{Spec}(\mathcal O_{X,x}))
  /\operatorname{Spec}(\mathcal O_{X,x}))
  =\operatorname{Daf}(U/X)\cap\operatorname{Spec}(\mathcal O_{X,x}).
  \tag{21.12.5.3}
\]
\end{env}

Cuando $X$ sea localmente noetheriano, precisaremos la naturaleza del conjunto
$\operatorname{Daf}(U/X)$, lo cual mostrará, por ejemplo, que cuando $U$ es denso en
todo $X$, $\operatorname{Daf}(U/X)$ no es un conjunto cerrado raro
arbitrario:

\begin{theorem}[21.12.6]
\label{IV.21.12.6}
Sean $X$ un preesquema localmente noetheriano, $U$ un abierto no vacío de $X$ y
$f:U\to X$ la inyección canónica. Entonces:
\begin{enumerate}[label=(\roman*)]
  \item La adherencia $Z=X-U_1$ de $\operatorname{Daf}(U/X)$ es de codimensión
  $\geq2$ en $X$.
  \item Si $T=X-U\supset Z$ es de codimensión $\geq2$, el morfismo
  $f^0:U^0\to X$ es sobreyectivo.
\end{enumerate}
\end{theorem}

\begin{proof}
(i) Demostremos primero que, para todo punto $x\in\operatorname{Daf}(U/X)$, se tiene
necesariamente $\dim(\mathcal O_{X,x})>1$. En efecto, $x$ evidentemente no puede ser un punto
maximal de $X$, pues está contenido en $\overline U-U$; se trata, pues, de ver que no se
puede tener $\dim(\mathcal O_{X,x})=1$. Ahora bien, esta relación implicaría, por (21.12.5.3),
$x\in\operatorname{Daf}(U^{(x)}/\operatorname{Spec}(\mathcal O_{X,x}))$, donde
$U^{(x)}=U\cap\operatorname{Spec}(\mathcal O_{X,x})$. Pero los únicos abiertos de
$\operatorname{Spec}(\mathcal O_{X,x})$ son el propio $\operatorname{Spec}(\mathcal O_{X,x})$
y los subconjuntos del conjunto (finito) de los puntos maximales de
$\operatorname{Spec}(\mathcal O_{X,x})$. Ahora bien, el conjunto abierto reducido a un punto maximal
de $\operatorname{Spec}(\mathcal O_{X,x})$ es afín, por definición de los preesquemas; se
concluye de ello que \emph{todos} los conjuntos abiertos de
$\operatorname{Spec}(\mathcal O_{X,x})$ son afines y, por tanto,
$\operatorname{Daf}(U^{(x)}/\operatorname{Spec}(\mathcal O_{X,x}))=\varnothing$
contrariamente a la hipótesis formulada.

Para demostrar (i) hay que probar algo más, a saber, que si $x\in X$ es tal que
$\dim(\mathcal O_{X,x})=1$, existe un entorno abierto $W$ de $x$ en $X$ tal que
$W$ no interseca $\operatorname{Daf}(U/X)$, es decir, tal que la inyección
canónica $f_W:U\cap W\to W$ sea afín. Pero, con las mismas notaciones anteriores,
acabamos de ver que la inyección canónica
$f^{(x)}:U^{(x)}\to\operatorname{Spec}(\mathcal O_{X,x})$ es afín. Evidentemente, se
puede limitar al caso en que $X$ es noetheriano; como el morfismo $f$ es de
presentación finita, la conclusión resulta de (8.10.5, (viii)), aplicado según el
método de (8.1.2, a)).

\oldpage[IV-4]{307}
(ii) Se trata de demostrar que, para todo punto $x\in T$, se tiene $x\in f^0(U^0)$;
demostraremos primero que se puede reducir al caso en que $X=\operatorname{Spec}(A)$, donde
$A$ es un anillo local noetheriano \emph{completo}, y $x$ es el punto cerrado de $X$. Para
ello basta efectuar el cambio de base
$h:X'=\operatorname{Spec}(\widehat{\mathcal O}_{X,x})\to X$; si se pone
$U'=h^{-1}(U)$; $f'=f_{\{X'\}}$ es la inyección canónica $U'\to X'$ y, puesto que el
morfismo $h$ es plano, resulta de (21.12.2) que, si se demuestra que $x$ pertenece a
$f^{\prime 0}(U^{\prime 0})$, se deducirá que $x\in f^0(U^0)$. En virtud de (6.1.1), se ha
efectuado así la reducción buscada.

Sea entonces $X_1$ un subpreesquema cerrado reducido de $X$ cuyo espacio subyacente
sea una componente irreducible de $X$, de dimensión \emph{máxima} entre las que
contienen una componente irreducible de $T$, y pongamos
$U_1=U\cap X_1$, $T_1=T\cap X_1=X_1-U_1$; se tiene entonces
$\operatorname{codim}(T_1,X_1)\geq2$ $(\mathbf 0, 14.2.1)$, y el par $(U_1,X_1)$
satisface las mismas hipótesis que el par $(U,X)$ (pues $X_1$ es el espectro de un
anillo cociente de $A$ y, por tanto, local noetheriano y completo). Se tiene, además, un
diagrama conmutativo
\[
  \xymatrix{
    U_1 \ar[r]^j \ar[d]_{f_1} & U \ar[d]^f \\
    X_1 \ar[r]_i & X
  }
\]
donde $i$ y $j$ son las inyecciones canónicas, siendo $i$, por tanto, un morfismo afín; se
deduce de (21.12.1) la existencia de un morfismo $j^0:U_1^0\to U^0$ que hace conmutativo el
diagrama
\[
  \xymatrix{
    U_1^0 \ar[r]^{j^0} \ar[d]_{f_1^0} & U^0 \ar[d]^{f^0} \\
    X_1 \ar[r]_i & X
  }
\]
y, por consiguiente, para demostrar que $x\in f^0(U^0)$ basta demostrar que
$x\in f_1^0(U_1^0)$. Se puede, pues, sustituir $X$ por $X_1$ y suponer, por consiguiente,
que el anillo $A$ es además \emph{íntegro}. Pero $A$ es cociente de un
anillo noetheriano regular en virtud del teorema de Cohen $(\mathbf 0, 19.8.8$, (i)) y,
como es íntegro, se puede aplicar (5.11.1) con la familia $(x_\alpha)$ reducida al
único punto maximal de $X$; como $\operatorname{codim}(T,X)\geq2$ por hipótesis, se
ve que $f_*(\mathcal O_U)$ es un $\mathcal O_X$-módulo \emph{coherente}, de modo que el
morfismo $f^0:U^0\to X$ es finito; como este morfismo es dominante (pues $U$ es denso en
todas partes), es sobreyectivo $(\mathbf{II}, 6.1.10)$ y, por tanto, $x\in f^0(U^0)$. C.Q.D.
\end{proof}

\begin{corollary}[21.12.7]
\label{IV.21.12.7}
Sean $X$ un preesquema localmente noetheriano, $U$ un subpreesquema inducido sobre un
abierto de $X$ y $j:U\to X$ la inmersión canónica; supongamos que $j$ sea un morfismo
afín. Entonces toda componente irreducible de $T=X-U$ es de codimensión $\leq1$
(y, por tanto, de codimensión $1$ si $U$ es denso en todas partes).
\end{corollary}

\begin{proof}
Supongamos, en efecto, que una de las componentes irreducibles $T_1$ de $T$ sea de
codimensión $\geq2$ en $X$. Sustituyendo, si es preciso, $X$ por un entorno abierto del
punto genérico de $T_1$, se puede suponer que $T$ es irreducible y de codimensión $\geq2$.
Pero entonces la hipótesis
\oldpage[IV-4]{308}
de que $j:U\to X$ es afín implica que $U^0$ se identifica con $U$ y $j^0$ con $j$,
es decir, que $j^0(U^0)=U$; ahora bien, esto contradice la conclusión de (21.12.6), que,
bajo la hipótesis de dimensión formulada, implica que $j^0$ debe ser sobreyectivo.
\end{proof}

\begin{env}[21.12.8]
\label{IV.21.12.8}
Sean $A$ un anillo local noetheriano, $Y=\operatorname{Spec}(A)$, $y$ el único punto
cerrado de $Y$ y $Y'=Y-\{y\}$. Consideremos la condición siguiente:
\begin{itemize}
  \item[\textup{(W)}] \emph{Para todo abierto $U$ de $Y$ contenido en $Y'$, que no
  contenga ninguna componente irreducible de $Y'$ y tal que la inmersión canónica
  $U\to Y'$ sea afín, $U$ es a su vez un abierto afín.}
\end{itemize}
Esta condición se deduce de la siguiente:
\begin{itemize}
  \item[\textup{(W')}] \emph{Para toda parte cerrada $T$ de $Y$, cada una de cuyas componentes
  irreducibles sea de codimensión $1$ en $Y$, y tal que para toda componente
  irreducible $Y_i$ de $Y$ se tenga $Y_i\cap T\ne\{y\}$, el abierto $U=Y-T$ es afín.}
\end{itemize}
En efecto, si se satisface (W') y si el abierto $U$ de $Y$ satisface la hipótesis de (W),
resulta de (21.12.7) que ninguna componente irreducible de $T=Y-U$ puede ser de
codimensión $\geq2$; como, además, $U$ no contiene ninguna de las componentes
irreducibles $Y_i-\{y\}$ de $Y'$, la condición (W') muestra que $U$ es afín.

Obsérvese que la condición (W') se simplifica cuando $Y$ es \emph{irreducible} y es
entonces equivalente a la siguiente:
\begin{itemize}
  \item[\textup{(W'')}] \emph{Para toda parte cerrada irreducible $T$ de $Y$, de
  codimensión $1$ en $Y$, $Y-T$ es un abierto afín.}
\end{itemize}
En efecto, es claro que (W') implica (W'') cuando $Y$ es irreducible, y la
recíproca es cierta considerando las componentes irreducibles de $T$ y observando
que la intersección de un número finito de abiertos afines es un abierto afín
$(\mathbf I, 5.5.6)$.
\end{env}

\begin{env}[21.12.9]
\label{IV.21.12.9}
\emph{Ejemplos}. ---
Si $A$ es un anillo local noetheriano \emph{factorial}, satisface (W') (y
\emph{a fortiori} (W)), puesto que todo ideal primo de altura $1$ es principal
$(\mathbf I, 1.3.6)$. Pero existen anillos locales noetherianos no factoriales que
satisfacen (W''), por ejemplo los de dimensión $\leq1$: en efecto, se observó en la
demostración de (21.12.6, (i)) que \emph{todos} los abiertos de $Y$ son entonces afines.
Por otra parte, utilizando la teoría de la dualidad local (cap.~III,
3\textsuperscript{a} parte), se puede demostrar que todo anillo local noetheriano de dimensión $2$
satisface (W').
\end{env}

El interés de la condición (W) reside en el resultado siguiente:

\begin{proposition}[21.12.10]
\label{IV.21.12.10}
Sean $X$, $Y$ dos preesquemas localmente noetherianos, $g:X\to Y$ un morfismo,
$y$ un punto cerrado de $Y$, $Y'=Y-\{y\}$ y $X'=g^{-1}(Y')$; supongamos que el morfismo
$g':X'\to Y'$, restricción de $g$, sea una inmersión abierta, y que el anillo local
$\mathcal O_{Y,y}$ satisfaga la condición (W) (21.12.8). Entonces, para toda componente
irreducible $Z$ de $X_y=g^{-1}(y)$, o bien $Z$ es de codimensión $\leq1$ en $X$,
o bien su punto genérico está aislado en $Z$. Si $Z$ es localmente de tipo finito sobre
$k(y)$, la segunda alternativa implica que $Z$ está reducido a un solo punto.
\end{proposition}

\begin{proof}
La última afirmación del enunciado resulta de que $Z$ es entonces un preesquema de
Jacobson (10.4.7) y, como el conjunto de los puntos cerrados de $Z$ es entonces denso en
$Z$, el punto genérico de $Z$ solo puede estar aislado en $Z$ si $Z$ está reducido a un
solo punto.

\oldpage[IV-4]{309}
Supongamos que existe en $X_y$ una componente irreducible $Z$ cuyo punto
genérico $z$ no está aislado en $Z$, y tal que $\operatorname{codim}(Z,X)\geq2$.
Como la cuestión es local sobre $X$ y $Y$, puesto que $z$ no está aislado en $Z$, se puede,
sustituyendo $X$ por un entorno abierto de $z$ en $X$, suponer que $X$ e $Y$ son afines,
que $X_y=Z$ es irreducible, no está reducido a un punto y satisface
$\operatorname{codim}(X_y,X)\geq2$. La imagen $g(X-X_y)=U$ es un abierto de $Y'$
isomorfo a $X-X_y$ por hipótesis; demostremos que, sustituyendo si es preciso $X$ por un
entorno abierto de $z$ en $X$, se puede suponer que $U$ no contiene ninguna de las
componentes irreducibles de $Y'$ cuya adherencia contiene a $y$.

En efecto, se puede suponer primero que todos los puntos maximales de $X$ son
generalizaciones de $z$, pues el conjunto de estos puntos es finito; por tanto, el conjunto de los puntos
maximales de $U$ es el conjunto de las imágenes $y_i$ por $g$ de los puntos maximales $x_i$ de
$X$ (ya que, por hipótesis, ningún punto maximal de $X$ puede estar contenido en $X_y$,
pues $\dim(\mathcal O_{X,z})\geq2$). Pongamos $X_i=\overline{\{x_i\}}$. Por
hipótesis, se tiene $z\in X_i$ y, como $z$ no está aislado en $X_y$, existe en
$X_y$ un punto $x\ne z$. Como $X_i\ne Z$, existe en $X_i$ un punto $t_i$ tal
que, si se pone $T_i=\overline{\{t_i\}}$, se tenga $\dim(\mathcal O_{T_i,x})=1$
(10.5.9); esto implica, por hipótesis, $t_i\notin Z$, de modo que $t_i$ no es una
generalización de $z$. Sustituyendo $X$ por $X'=X-\bigcup_iT_i$, se ve que la imagen por
$g$ de $X'-X_y$ no contiene las imágenes $g(t_i)$ y, por tanto, no contiene ninguna de las
componentes irreducibles de $Y'$ cuya adherencia contiene a $y$.

Hecho esto, como $X$ e $Y$ son afines, el morfismo $g:X\to Y$ es afín, y lo
mismo sucede con la restricción $g':X'\to Y'$; puesto que $g'$ es una inmersión abierta,
la inmersión canónica $U\to Y'$ es afín. Pongamos
$Y_1=\operatorname{Spec}(\mathcal O_{Y,y})$, $Y'_1=Y_1-\{y\}=Y'\cap Y_1$,
$U_1=U\cap Y_1$. Lo precedente demuestra que la inmersión canónica $U_1\to Y'_1$ es
afín; por tanto, $U_1$ es un abierto afín de $Y_1$ en virtud de (W), es decir, la
inmersión canónica $U_1\to Y_1$ es afín. Pero, como esta inmersión es de
presentación finita (1.6.2), resulta de (8.10.5, (viii)), aplicado según el método de
(8.1.2, a)), que, restringiendo si es preciso $Y$ a un entorno abierto afín de $y$, se puede
suponer que la inmersión $U\to Y$ es afín. Se concluiría que el abierto $X-X_y$ de
$X$, isomorfo a $U$, sería afín, lo que contradiría (21.12.7); queda así demostrada la
proposición.
\end{proof}

\begin{corollary}[21.12.11]
\label{IV.21.12.11}
Sean $X$, $Y$ dos preesquemas localmente noetherianos, suponiéndose que $X$ es
irreducible; sea $g:X\to Y$ un morfismo localmente de tipo finito, y sea $V$ el mayor
abierto de $X$ tal que la restricción $g|V:V\to Y$ sea un isomorfismo local.
Supongamos que, para todo punto $y\in Y$, el anillo local $\mathcal O_{Y,y}$ satisface la
condición (W) (21.12.8). Entonces toda componente irreducible de $T=X-V$, o bien
es de codimensión $\leq1$, o bien tiene su punto genérico $z$ aislado en
$g^{-1}(g(z))$.
\end{corollary}

\begin{proof}
Pongamos $g(z)=y$. Como la cuestión solo depende de la fibra $g^{-1}(y)$ y del anillo
$\mathcal O_{Y,y}$, mediante un cambio de base se puede limitar al caso en que
$Y=\operatorname{Spec}(\mathcal O_{Y,y})$ y $X$ es afín
$((\mathbf I, 3.6.5)$ y $(\mathbf I, 4.5.5))$; sustituyendo eventualmente $X$ por un
entorno abierto de $z$ en $X$, se puede suponer que $T$ es irreducible; además, se
puede limitar al caso en que $V$ no es vacío. Se puede, pues, suponer que el morfismo $g$ es
separado y que $y$ es cerrado en $Y$; puesto que la restricción $g|V:V\to Y$ es un isomorfismo
local y el abierto no vacío $V$ de $X$ es irreducible, resulta de
$(\mathbf I, 8.2.8)$ que $g|V:V\to Y$ es una inmersión abierta. La restricción
$g|(V\cap X_y):V\cap X_y\to\{y\}=\operatorname{Spec}(k(y))$ es, por tanto, también una
inmersión abierta, lo que demuestra que $V\cap X_y$ es, o bien vacío, o bien está reducido a un punto
$x$ racional
\oldpage[IV-4]{310}
sobre $k(y)$ y, por tanto, cerrado en $X_y$ $(\mathbf I, 6.4.2)$. Sustituyendo nuevamente,
si es preciso, $X$ por un entorno abierto más pequeño de $z$ en $X$, se puede, pues,
limitar al caso en que $V\cap X_y=\varnothing$, es decir, $T\supset X_y$; como, además,
$z\in X_y$ es el punto genérico de $T$, se tiene
$g(T)\subset\overline{\{y\}}=\{y\}$, es decir, $T=X_y$. Se está entonces exactamente
en las condiciones de aplicación de (21.12.10), de donde resulta la conclusión.
\end{proof}

\begin{theorem}[21.12.12]
\label{IV.21.12.12}
\textup{(van der Waerden)}. ---
Sean $X$, $Y$ dos preesquemas localmente noetherianos íntegros, y sea $g:X\to Y$ un
morfismo birracional y localmente de tipo finito. Se supone además que $Y$ es
normal y que, para todo $y\in Y$, el anillo $\mathcal O_{Y,y}$ satisface la condición (W)
de (21.12.8) (condiciones que se cumplen, en particular, cuando $Y$ es localmente factorial
(21.12.9)). Si $V$ es el mayor abierto de $X$ tal que la restricción
$g|V:V\to Y$ sea un isomorfismo local, todas las componentes irreducibles de
$T=X-V$ son de codimensión $1$ en $X$.
\end{theorem}

\begin{proof}
Obsérvese que, puesto que $g$ es birracional, el conjunto abierto $V$ no es vacío; basta,
pues, demostrar que un punto maximal $z$ de $T$ no puede estar aislado en $X_y$,
donde $y=g(z)$. Pero, restringiéndose a un entorno abierto de $z$, se puede
limitar al caso en que $T=X_y$; entonces todas las fibras $g^{-1}(y')$ ($y'\in Y$) serían
vacías o estarían reducidas a un punto, y resultaría del «Main theorem» (8.12.10) que $g$
sería un isomorfismo local, contrariamente a la hipótesis $z\in T$.
\end{proof}

\begin{corollary}[21.12.13]
\label{IV.21.12.13}
Supongamos que se satisfacen las hipótesis de (21.12.12) y, además, que $g$ es
cuasifinito en cada uno de los puntos de $X^{(1)}$ (se recuerda que estos son los puntos en los que
$\dim(\mathcal O_{X,x})=1$); entonces $g$ es un isomorfismo local y, si además $g$ es
separado, $g$ es una inmersión abierta.
\end{corollary}

\begin{proof}
Basta demostrar la primera afirmación, pues la segunda es consecuencia de la
primera y de $(\mathbf I, 8.2.8)$. Todo se reduce a demostrar, con las notaciones de
(21.12.12), que $T=\varnothing$; en caso contrario, un punto genérico $z$ de una
componente irreducible de $T$ pertenecería a $X^{(1)}$ en virtud de (21.12.12) y, por
hipótesis, estaría aislado en $X_y$, donde $y=g(z)$; pero se ha visto en la demostración
de (21.12.12) que esto no es posible.
\end{proof}

\begin{remarks}[21.12.14]
\label{IV.21.12.14}
(i) La conclusión de (21.12.12) se aplica cuando $Y$ es regular e íntegro y $X$ es
íntegro, pues, en virtud del teorema de Auslander--Buchsbaum (21.11.1), $Y$ es entonces
localmente factorial. En cambio, se conocen ejemplos en los que $X$ e $Y$ son esquemas
algebraicos normales de dimensión $3$, sobre un cuerpo algebraicamente cerrado de
característica arbitraria, y en los que la conclusión de (21.12.12) ya no es válida.

(ii) El conjunto $T$ de (21.12.12) es también el conjunto de los puntos en los que $g$ está
\emph{ramificado}. En efecto, si $g$ no está ramificado en un punto $x\in X$, tampoco lo está
en un entorno abierto afín $U$ de $x$ (17.3.7), de modo que $g|U$ es un morfismo
separado, cuasifinito y birracional (17.4.3). Puesto que $Y$ es normal, se concluye del «Main
theorem» $(\mathbf{III}, 4.4.9)$ que $g$ es un isomorfismo local en el punto $x$ y, por tanto,
$x\in X-T$. Recíprocamente, $g$ evidentemente no está ramificado en ningún punto en el que es un
isomorfismo local. Esto justifica el título dado a esta sección.

\oldpage[IV-4]{311}
(iii) Sin suponer que $Y$ sea localmente factorial, pero suponiendo, en cambio, que
$X$ es intersección completa relativamente a $Y$ (cap.~V), veremos en el cap.~V
que se tiene un resultado más preciso que (21.12.12), expresando $T$ como ciclo
1-codimensional de una sección de un $\mathcal O_X$-módulo invertible, canónicamente
asociada a $g$. Esto se aplicará, en particular, cuando $X$ e $Y$ sean ambos regulares,
o cuando $g$ sea un $S$-morfismo, siendo $X$ e $Y$ preesquemas lisos sobre $S$.

(iv) Cabe preguntarse si (21.12.10) admite una recíproca; en otras palabras, dado un
anillo local noetheriano $A$, si se pone $Y=\operatorname{Spec}(A)$, designando $y$
el punto cerrado de $Y$, y si, para todo preesquema localmente noetheriano $X$ y
todo morfismo $g:X\to Y$ tal que $g':X'\to Y'$ sea una inmersión abierta, toda
componente irreducible de $X_y$ es de codimensión $\leq1$ en $X$ o está reducida a
un punto, ¿es entonces cierto que $A$ satisface la condición (W) de (21.12.8)?

(v) Sean $Y$ un preesquema íntegro \emph{regular}, $X$ un preesquema localmente
noetheriano \emph{normal} y $g:X\to Y$ un morfismo localmente de tipo finito; supongamos
además que, si $\eta$ es el punto genérico de $Y$, la fibra $X_\eta$ sea \emph{étale}
sobre $k(\eta)$, y sea $T'$ el complementario del mayor abierto $V'$ de $X$ tal que
$g|V':V'\to Y$ sea étale; ¿es cierto entonces que todas las componentes irreducibles de
$T'$ sean de codimensión $1$? Así ocurre cuando se supone además que $g$
es \emph{localmente cuasifinito} («teorema de pureza» de Zariski--Nagata). Se puede
demostrar que la conjetura precedente sería consecuencia de la siguiente (y sería
incluso equivalente a ella si se satisficiera la conjetura de (iv)): para un anillo local regular
$A$ contenido en un anillo local íntegro e íntegramente cerrado $B$, tal que $B$ sea un
$A$-módulo finito, el abierto $U$ de los puntos de $\operatorname{Spec}(B)$ en los que
$\operatorname{Spec}(B)$ no está ramificado sobre $\operatorname{Spec}(A)$ (o, lo que
equivale por (18.10.1), es étale sobre $\operatorname{Spec}(A)$) es afín.
\end{remarks}

\begin{proposition}[21.12.15]
\label{IV.21.12.15}
Sean $S$ un preesquema, $g:X\to S$ un morfismo plano y localmente de presentación
finita, cuyas fibras $X_s=g^{-1}(s)$ son geométricamente irreducibles (4.5.2), y
$h:Y\to S$ un morfismo liso; se denotan por $Y_s=h^{-1}(s)$ las fibras de este morfismo.
Sea $f:X\to Y$ un $S$-morfismo propio tal que, para todo punto maximal $\eta$ de
$S$, el morfismo $f_\eta:X_\eta\to Y_\eta$ sea un isomorfismo. Entonces $f$ es un
isomorfismo.
\end{proposition}

\begin{proof}
Puesto que $h$ es plano, todo punto maximal de $Y$ está sobre un punto maximal de $S$
(2.3.4), de modo que la unión de los $Y_\eta$, cuando $\eta$ recorre el conjunto de los puntos
maximales de $S$, es densa en $Y$; como los $f_\eta$ son isomorfismos, $f$
es dominante y, por tanto, \emph{sobreyectivo}, puesto que es propio. Basta, pues, demostrar
que $f$ es una \emph{inmersión abierta}. Teniendo en cuenta (17.9.5), basta probar
que, para todo $s\in S$, $f_s:X_s\to Y_s$ es una inmersión abierta.

Como todo $s\in S$ es una especialización de un punto maximal $\eta$\footnote{Corrección de fuente: el francés impreso y el inglés mantenido dicen que $s$ es una generalización de $\eta$. La construcción siguiente toma $s\in\overline{\{\eta\}}$, por lo que $s$ es una especialización de $\eta$, o equivalentemente $\eta$ una generalización de $s$.}, mediante el
cambio de base $\operatorname{Spec}(\mathcal O_{S',s})\to S$, donde $S'$ es el
subpreesquema reducido de $S$ cuyo espacio subyacente es $\overline{\{\eta\}}$,
se puede reducir desde ahora al caso en que $S$ es un esquema local e íntegro de punto cerrado $s$:
en efecto, las fibras en los puntos $\eta$ y $s$ no cambian y las propiedades de los
morfismos $g$, $h$ y $f$ se conservan por cambio de base. Además, la cuestión
es local sobre $Y$; sustituyendo $Y$ por un entorno abierto afín $U$ en $Y$ de un punto
de $Y_s$, y $X$ por $f^{-1}(U)$, que es cuasicompacto puesto que $f$ es propio, se ve
que se puede suponer, además, que $X$ e $Y$ son de presentación finita sobre $S$.

Existen entonces un subanillo local noetheriano $A_0$ de
$A=\mathcal O_{S,s}$ tal que $A_0\to A$ sea un homomorfismo local, dos morfismos
de tipo finito $g_0:X_0\to S_0=\operatorname{Spec}(A_0)$,
$h_0:Y_0\to S_0$ y un $S_0$-morfismo
\oldpage[IV-4]{312}
$f_0:X_0\to Y_0$ tales que $g$, $h$ y $f$ se deducen de $g_0$, $h_0$, $f_0$ por el
cambio de base $S\to S_0$ ((8.9.1) y (5.13.3)); se puede, además, suponer que $g_0$ es
plano (11.2.7), $h_0$ liso (17.7.9) y $f_0$ propio (8.10.5, (xii)). Por otra parte, la
proyección sobre $S_0$ del punto genérico $\eta$ de $S$ es el punto genérico $\eta_0$ de
$S_0$, y resulta de (2.7.1, (viii)) que el morfismo
$(f_0)_{\eta_0}:(X_0)_{\eta_0}\to(Y_0)_{\eta_0}$ es un isomorfismo; por último, el punto
cerrado $s$ de $S$ está sobre el punto cerrado $s_0$ de $S_0$, de modo que la fibra
$(X_0)_{s_0}$ de $g_0$ es geométricamente irreducible (4.5.6).

Se ve, pues, que se puede limitar al caso en que $S$ es el espectro de un anillo local íntegro
noetheriano, de punto cerrado $s$, debilitar la hipótesis sobre las fibras suponiendo solamente
que $X_s$ y $X_\eta$ son geométricamente irreducibles, y demostrar que, bajo estas
hipótesis, $f_s$ es una inmersión abierta. Existe entonces un anillo de valoración discreta
$A'$ y un morfismo $S'=\operatorname{Spec}(A')\to S$ que transforma el punto cerrado
$a$ de $S'$ en $s$ y el punto genérico $b$ de $S'$ en $\eta$ $(\mathbf{II}, 7.1.9)$;
sean $g':X'\to S'$, $h':Y'\to S'$, $f':X'\to Y'$ los morfismos deducidos de $g$,
$h$, $f$ por el cambio de base $S'\to S$, que satisfacen de nuevo las hipótesis
satisfechas por $g$, $h$ y $f$ en el enunciado (21.12.15); si se demuestra que
$f'_a:X'_a\to Y'_a$ es una inmersión abierta, resultará de (2.7.1, (x)) que también
lo es $f_s:X_s\to Y_s$. Se puede, pues, limitar finalmente al caso en que $S$ es
el \emph{espectro de un anillo de valoración discreta}.

Entonces, puesto que $h:Y\to S$ es liso, $Y$ es regular (17.5.8), y como la cuestión
es local sobre $Y$, se puede limitar al caso en que $Y$ es íntegro. Como $g:X\to S$ es
plano, los puntos maximales de $X$ están contenidos en $X_\eta$ (2.3.4) y, puesto que
$f_\eta$ es un isomorfismo, $X$ es irreducible y el anillo local en su punto
genérico es un cuerpo; además, por planitud (3.3.2), se ve que $X$ no tiene ciclos
primos inmersos, de modo que $X$ es reducido (3.2.1), por tanto íntegro, y $f$ es un
morfismo birracional y separado. Para demostrar que $f$ es una inmersión abierta, se puede
aplicar, pues, el criterio (21.12.13); para ver que $f$ es cuasifinito en los puntos de
$X^{(1)}$, se observa que la afirmación es evidente en aquellos de estos puntos que
pertenecen a $X_\eta$; el único punto de $X^{(1)}$ que pertenece a $X_s$ es el punto
genérico $x$ de $X_s$, en virtud de (6.1.1). Ahora bien, resulta de (2.4.6) y de (14.2.2)
que los morfismos $g$ y $h$ son equidimensionales; como $X_\eta$ e $Y_\eta$ son
isomorfos, las componentes irreducibles de $X_s$ y de $Y_s$ tienen todas la misma
dimensión. Pero, por hipótesis, $X_s$ es irreducible y se ha visto que $f$ es sobreyectivo,
por lo que también lo es $f_s:X_s\to Y_s$, lo cual implica que $Y_s$ es también
irreducible; si $y$ es el punto genérico de $Y_s$, se tiene, pues, $f(x)=y$, y la relación
$\dim(X_s)=\dim(Y_s)$ implica entonces, en virtud de (5.6.6), que
$\dim(f_s^{-1}(y))=0$, es decir, $f_s$ (y, por tanto, también $f$) es efectivamente cuasifinito
en el punto $x$, lo que termina la demostración.
\end{proof}

\begin{corollary}[21.12.16]
\label{IV.21.12.16}
Sean $g:X\to S$ un morfismo propio, plano y de presentación finita, $h:Y\to S$ un
morfismo liso, propio y de presentación finita, y $f:X\to Y$ un $S$-morfismo. Se
supone que las fibras $X_s=g^{-1}(s)$ de $g$ son geométricamente irreducibles. Sea
$U$ el conjunto de los $s\in S$ tales que $f_s:X_s\to Y_s$ sea un isomorfismo. Entonces
$U$ es abierto y cerrado en $S$, y el morfismo $g^{-1}(U)\to h^{-1}(U)$,
restricción de $f$, es un isomorfismo.
\end{corollary}

\begin{proof}
La última afirmación resultará de la primera y de (17.9.5). Ya se sabe (9.6.1,
(xi)) que $U$ es localmente constructible en $S$. En virtud de (1.10.1), basta
demostrar que $U$ es estable por \emph{especialización} y por \emph{generalización}. Para
demostrar estas afirmaciones se
\oldpage[IV-4]{313}
puede, mediante un cambio de base $\operatorname{Spec}(\mathcal O_{S,s})\to S$,
reducir al caso en que $S$ es un esquema local de punto cerrado $s$ y punto genérico
$\eta$, y hay que demostrar que, para que $f_s$ sea un isomorfismo, es necesario y suficiente
que $f_\eta$ lo sea. Ahora bien, la suficiencia de esta condición resulta de (21.12.15). Para
demostrar que es necesaria, se razona como en la demostración de (21.12.15)
(observando, con las mismas notaciones, que la proyección del punto cerrado de $S$ es el
punto cerrado de $S_0$) para reducir al caso en que, además, $S$ es \emph{noetheriano};
pero entonces la conclusión resulta de $(\mathbf{III}, 4.6.7$, (ii)).
\end{proof}

\begin{remarks}[21.12.17]
\label{IV.21.12.17}
(i) La conclusión de la proposición (21.12.15) ya no es válida si se elimina
la hipótesis de que las fibras $X_s$ son irreducibles. Tomemos, en efecto,
$S=\operatorname{Spec}(A)$, donde $A$ es un anillo de valoración discreta, y
$Y=\mathbb P^1_S$, que es propio y liso sobre $S$ (17.3.9). Denotemos nuevamente por $s$ y
$\eta$ el punto cerrado y el punto genérico de $S$; sea $z$ un punto cerrado de $Y_s$,
por ejemplo un punto racional sobre $k=k(s)$, y sea $X$ el $Y$-preesquema obtenido por
explosión del punto $z$. Como el anillo de polinomios $A[T]$ es regular
$(\mathbf 0, 17.3.7)$ y de dimensión $2$, se ve, como en la demostración de
(15.1.1.6), que, si $f:X\to Y$ es el morfismo estructural, $f^{-1}(z)$ es isomorfo
a $\mathbb P^1_k$ y, por otra parte, el complementario de $f^{-1}(z)$ en $X_s$ es
isomorfo a $Y_s-\{z\}$; por tanto, $Z_1=f^{-1}(z)$ y $Z_2$, adherencia en $X_s$ del
complementario de $Z_1$, son las dos componentes irreducibles de $X_s$. Por
otra parte, $f$ es evidentemente propio y $g=h\circ f$ es plano, pues $X$ es íntegro
$(\mathbf{II}, 8.1.4)$ y, para todo abierto afín $U$ de $X$, el homomorfismo
$A\to\Gamma(U,\mathcal O_X)$ es inyectivo $(\mathbf I, 1.2.7)$, de modo que
$\Gamma(U,\mathcal O_X)$ es un $A$-módulo sin torsión y, por tanto, plano, puesto que $A$
es un anillo de valoración discreta $(\mathbf 0_{\mathrm I}, 6.3.4)$. Es claro que
$f_\eta:X_\eta\to Y_\eta$ es un isomorfismo, aunque
$f_s:X_s\to Y_s$ no lo sea.

(ii) La conclusión de (21.12.15) también falla si, entre las hipótesis, se
suprime la hipótesis de que $f$ es \emph{propio}. Para verlo, basta definir $S$ e
$Y$ como en (i), pero sustituir $X$ por el complementario $X'$ de $Z_2$ en el
preesquema $X$ definido en (i), y $f$ por la restricción $f':X'\to Y$ de $f$; el morfismo
$g':X'\to S$, restricción de $g$, sigue siendo plano, y esta vez $X'_s$ es
geométricamente irreducible; además, $X'_s$ es isomorfo al complementario de un
punto cerrado en $\mathbb P^1_k$, por lo que es un esquema afín isomorfo a $\mathbb A^1_k$;
como su imagen en $Y_s$ está reducida al punto cerrado $z$, $f'$ no es propio
$(\mathbf{III}, 4.4.2)$ y $f'_s$ no es un isomorfismo, aunque $f'_\eta$ sí lo sea.

(iii) Es posible que el enunciado de la proposición (21.12.15) siga siendo válido cuando se
sustituya en él la palabra «isomorfismo» por «morfismo étale» (cf. (21.12.14, (v))). En
tal caso, lo mismo seguirá siendo cierto para (21.12.16).
\end{remarks}

\subsection{Pares parafactoriales. Anillos locales parafactoriales}
\label{subsection:IV.21.13}
\label{IV.21.13}

\begin{definition}[21.13.1]
\label{IV.21.13.1}
Sean $X$ un espacio anillado, $Y$ una parte cerrada de $X$ y $U=X-Y$.
Se dice que el par $(X,Y)$ es \emph{parafactorial} si, para todo abierto $V$ de $X$,
el functor de restricción
$\mathscr L\mapsto\mathscr L|(U\cap V)$, de la categoría de los
$\mathcal O_V$-módulos invertibles a la de los $\mathcal O_{U\cap V}$-módulos
invertibles, es una \emph{equivalencia de categorías}.
\end{definition}

Decir que el par $(X,Y)$ es parafactorial significa, pues, que para todo abierto $V$
de $X$ se satisfacen las condiciones siguientes:
\begin{enumerate}[label=\arabic*\textsuperscript{o}]
  \item el functor $\mathscr L\mapsto\mathscr L|(U\cap V)$ es \emph{plenamente
  fiel}; dicho de otro modo, para dos $\mathcal O_V$-módulos invertibles $\mathscr L$,
  $\mathscr L'$, la aplicación de restricción
  \[
    \operatorname{Hom}_{\mathcal O_V}(\mathscr L,\mathscr L')
    \longrightarrow
    \operatorname{Hom}_{\mathcal O_{U\cap V}}
    (\mathscr L|(U\cap V),\mathscr L'|(U\cap V))
  \]
  es \emph{biyectiva};
  \item el functor $\mathscr L\mapsto\mathscr L|(U\cap V)$ es
  \emph{esencialmente sobreyectivo}; dicho de otro modo, para todo
  $\mathcal O_{U\cap V}$-módulo invertible $\mathscr L_0$, existe un
  $\mathcal O_V$-módulo invertible $\mathscr L$ tal que $\mathscr L_0$ sea
  \oldpage[IV-4]{314}
  isomorfo a $\mathscr L|(U\cap V)$; esto se expresa también diciendo que
  el homomorfismo canónico (21.3.2.4)
  \[
    \operatorname{Pic}(V)\longrightarrow\operatorname{Pic}(U\cap V)
  \]
  es \emph{sobreyectivo}.
\end{enumerate}

\begin{lemma}[21.13.2]
\label{IV.21.13.2}
Sea $f:X'\to X$ un morfismo de espacios anillados; para todo abierto $V$ de $X$,
se denota por $f_V:f^{-1}(V)\to V$ la restricción de $f$. Las condiciones siguientes son
equivalentes:
\begin{enumerate}[label=\emph{\alph*)}]
  \item Para todo abierto $V$ de $X$, el functor
  $\mathscr E\mapsto f_V^*(\mathscr E)$ de la categoría de los $\mathcal O_V$-módulos
  localmente libres de rango finito a la categoría de los
  $\mathcal O_{f^{-1}(V)}$-módulos localmente libres de rango finito es fiel
  (resp.\ plenamente fiel).
  \item[\emph{a')}] Para todo abierto $V$ de $X$, el functor
  $\mathscr L\mapsto f_V^*(\mathscr L)$ de la categoría de los $\mathcal O_V$-módulos
  localmente libres de rango $1$ a la categoría de los
  $\mathcal O_{f^{-1}(V)}$-módulos localmente libres de rango $1$ es fiel
  (resp.\ plenamente fiel).
  \item Para todo abierto $V$ de $X$, el homomorfismo canónico
  $(\mathbf 0_{\mathrm I}, 4.4.3.2)$
  $\mathscr E\to(f_V)_*(f_V^*(\mathscr E))$ es un monomorfismo
  (resp.\ un isomorfismo) para todo $\mathcal O_V$-módulo localmente libre
  $\mathscr E$ de rango finito.
  \item[\emph{b')}] El homomorfismo canónico
  $\mathcal O_X\to f_*(\mathcal O_{X'})$ es un monomorfismo
  (resp.\ un isomorfismo).
\end{enumerate}

Supongamos que el homomorfismo canónico $\mathcal O_X\to f_*(\mathcal O_{X'})$ es
biyectivo; entonces, para que un $\mathcal O_{X'}$-módulo localmente libre de rango finito
$\mathscr E'$ sea isomorfo a un $\mathcal O_{X'}$-módulo de la forma
$f^*(\mathscr E)$, donde $\mathscr E$ es un $\mathcal O_X$-módulo localmente libre de
rango finito, es necesario y suficiente que se satisfagan las dos condiciones siguientes:
\begin{enumerate}[label=(\roman*)]
  \item $f_*(\mathscr E')$ es un $\mathcal O_X$-módulo localmente libre;
  \item el homomorfismo canónico $(\mathbf 0_{\mathrm I}, 4.4.3.3)$
  $f^*(f_*(\mathscr E'))\to\mathscr E'$ es un isomorfismo.
\end{enumerate}
Cuando se satisfacen estas dos condiciones, $\mathscr E$ es isomorfo a
$f_*(\mathscr E')$.
\end{lemma}

\begin{proof}
Decir que el functor $\mathscr E\mapsto f_V^*(\mathscr E)$ es fiel
(resp.\ plenamente fiel) significa que, para dos $\mathcal O_V$-módulos localmente
libres de rango finito $\mathscr E_1$, $\mathscr E_2$, la aplicación
\[
  \operatorname{Hom}_{\mathcal O_V}(\mathscr E_1,\mathscr E_2)
  \longrightarrow
  \operatorname{Hom}_{\mathcal O_{f^{-1}(V)}}
  (f_V^*(\mathscr E_1),f_V^*(\mathscr E_2))
\]
es inyectiva (resp.\ biyectiva); como esto debe cumplirse sustituyendo $V$ por un
abierto $W\subset V$ y $\mathscr E_1$, $\mathscr E_2$ por $\mathscr E_1|W$,
$\mathscr E_2|W$, y como
\[
  \operatorname{Hom}_{\mathcal O_W}(\mathscr E_1|W,\mathscr E_2|W)
  =\Gamma(W,\mathscr Hom_{\mathcal O_V}(\mathscr E_1,\mathscr E_2)),
\]
la condición significa también que el homomorfismo canónico de haces
\[
\label{IV.21.13.2.1}
  \mathscr Hom_{\mathcal O_V}(\mathscr E_1,\mathscr E_2)
  \longrightarrow
  (f_V)_*(f_V^*(\mathscr Hom_{\mathcal O_V}(\mathscr E_1,\mathscr E_2)))
  \tag{21.13.2.1}
\]
es inyectivo (resp.\ biyectivo). Pero, puesto que $\mathscr E_1$ y $\mathscr E_2$ son
localmente libres de rango finito,
$\mathscr Hom_{\mathcal O_V}(\mathscr E_1,\mathscr E_2)$, isomorfo a
$\mathscr E_1^\vee\otimes_{\mathcal O_V}\mathscr E_2$
$(\mathbf 0_{\mathrm I}, 5.4.2)$, es también localmente libre de rango finito; esto muestra
ya que $b)$ implica $a)$ y, recíprocamente, $b)$ es un caso particular de $a)$,
teniendo en cuenta el isomorfismo
$\mathscr E\xrightarrow{\sim}\mathscr Hom_{\mathcal O_V}(\mathcal O_V,\mathscr E)$.
Está claro que $a')$ es un caso particular de $a)$ y $b')$ un caso particular de
$a')$. Recíprocamente, como $b)$ es una propiedad local, para verificarla podemos
limitarnos al caso en que $\mathscr E=\mathcal O_V^n$, y esto muestra que $b')$ implica
$b)$.

Si el homomorfismo canónico $\mathcal O_X\to f_*(\mathcal O_{X'})$ es biyectivo,
y si se satisfacen las condiciones (i) y (ii), está claro que
$\mathscr E'=f^*(\mathscr E)$ con $\mathscr E=f_*(\mathscr E')$, salvo isomorfismo.
\oldpage[IV-4]{315}
Recíprocamente, supongamos que $\mathscr E'=f^*(\mathscr E)$, con $\mathscr E$ localmente
libre; como la cuestión es local sobre $X$, podemos suponer que
$\mathscr E=\mathcal O_X^n$, de donde $\mathscr E'=\mathcal O_{X'}^n$, y las condiciones
(i) y (ii) resultan de la hipótesis de que
$\mathcal O_X\to f_*(\mathcal O_{X'})$ es un isomorfismo.
\end{proof}

En este número aplicaremos el lema anterior a una inyección canónica
$j:U\to X$, donde $U$ es el espacio anillado inducido por $X$ sobre un \emph{abierto} de $X$.
Con estas notaciones, (21.13.2) se traduce en:

\begin{corollary}[21.13.3]
\label{IV.21.13.3}
Sean $X$ un espacio anillado, $Y$ una parte cerrada de $X$, $U=X-Y$ y
$j:U\to X$ la inyección canónica. Las condiciones siguientes son equivalentes:
\begin{enumerate}[label=\emph{\alph*)}]
  \item Para todo abierto $V$ de $X$, el functor de restricción
  $\mathscr E\mapsto\mathscr E|(U\cap V)$ de la categoría de los
  $\mathcal O_V$-módulos localmente libres de rango finito a la categoría de los
  $\mathcal O_{U\cap V}$-módulos localmente libres de rango finito es fiel
  (resp.\ plenamente fiel).
  \item[\emph{a')}] Para todo abierto $V$ de $X$, el functor de restricción
  $\mathscr L\mapsto\mathscr L|(U\cap V)$ de la categoría de los
  $\mathcal O_V$-módulos localmente libres de rango $1$ a la categoría de los
  $\mathcal O_{U\cap V}$-módulos localmente libres de rango $1$ es fiel
  (resp.\ plenamente fiel).
  \item Para todo abierto $V$ de $X$, el homomorfismo canónico
  $\mathscr E\to(j_V)_*(\mathscr E|(U\cap V))$ es inyectivo (resp.\ biyectivo) para
  todo $\mathcal O_V$-módulo localmente libre de rango finito $\mathscr E$.
  \item[\emph{b')}] El homomorfismo canónico
  $\mathcal O_X\to j_*(\mathcal O_U)$ es inyectivo (resp.\ biyectivo).
\end{enumerate}

Supongamos que el homomorfismo canónico $\mathcal O_X\to j_*(\mathcal O_U)$ es
biyectivo. Entonces, para que un $\mathcal O_U$-módulo localmente libre de rango finito
$\mathscr F$ sea de la forma $\mathscr E|U$, donde $\mathscr E$ es un
$\mathcal O_X$-módulo localmente libre de rango finito, es necesario y suficiente que
$j_*(\mathscr F)$ sea un $\mathcal O_X$-módulo localmente libre, y en este caso se
puede tomar $\mathscr E=j_*(\mathscr F)$.
\end{corollary}

\begin{lemma}[21.13.4]
\label{IV.21.13.4}
Sean $X$ un preesquema localmente noetheriano, $Y$ una parte cerrada de $X$,
$U=X-Y$ y $j:U\to X$ la inyección canónica. Para que el homomorfismo canónico
$\mathcal O_X\to j_*(\mathcal O_U)$ sea inyectivo (resp.\ biyectivo), es necesario y
suficiente que $\operatorname{prof}_Y(\mathcal O_X)\geq1$
(resp.\ $\operatorname{prof}_Y(\mathcal O_X)\geq2$) (dicho de otro modo, que
$\operatorname{prof}(\mathcal O_{X,x})\geq1$
(resp.\ $\operatorname{prof}(\mathcal O_{X,x})\geq2$) para todo $x\in Y$).
\end{lemma}

\begin{proof}
Estas afirmaciones son, en efecto, casos particulares de (5.10.2) y (5.10.5).
\end{proof}

\begin{proposition}[21.13.5]
\label{IV.21.13.5}
Sean $X$ un espacio anillado, $Y$ una parte cerrada de $X$, $U=X-Y$ y
$j:U\to X$ la inyección canónica. Para que el par $(X,Y)$ sea parafactorial, es
necesario y suficiente que el homomorfismo canónico $\mathcal O_X\to j_*(\mathcal O_U)$
sea biyectivo, y que para todo abierto $V$ de $X$ y todo
$\mathcal O_{U\cap V}$-módulo invertible $\mathscr L$, $(j_V)_*(\mathscr L)$ sea un
$\mathcal O_V$-módulo invertible (notación de (21.13.3)).
\end{proposition}

\begin{proof}
Es una consecuencia inmediata de la definición (21.13.1) y de (21.13.3).
\end{proof}

\begin{corollary}[21.13.6]
\label{IV.21.13.6}
Sean $X$ un espacio anillado e $Y$ una parte cerrada de $X$.
\begin{enumerate}[label=(\roman*)]
  \item Si el par $(X,Y)$ es parafactorial, también lo es
  $(W,Y\cap W)$ para todo abierto $W$ de $X$. Recíprocamente, si $(W_\alpha)$ es un
  recubrimiento abierto de $X$ tal que cada uno de los pares
  $(W_\alpha,Y\cap W_\alpha)$ es parafactorial, entonces el par $(X,Y)$ es
  parafactorial.
  \item Si el par $(X,Y)$ es parafactorial, también lo es $(X,Y')$ para
  toda parte cerrada $Y'$ de $Y$.
  \item Supongamos que $X$ es un preesquema, y sean $X'$ un preesquema y
  $f:X'\to X$ un morfismo fielmente plano y cuasicompacto; pongamos
  $Y'=f^{-1}(Y)$. Supongamos que el par $(X',Y')$ es parafactorial y que
  el abierto $U=X-Y$ es retrocompacto en $X$ $(\mathbf 0_{\mathrm{III}}, 9.1.1)$.
  Entonces el par $(X,Y)$ es parafactorial.
\end{enumerate}
\end{corollary}

\begin{proof}
\oldpage[IV-4]{316}
(i) Como el hecho de que el homomorfismo canónico
$\mathcal O_X\to j_*(\mathcal O_U)$ sea biyectivo es una propiedad local sobre $X$,
todo se reduce a ver (en virtud de (21.13.5)), designando por
$j_\alpha:U\cap W_\alpha\to W_\alpha$ la inyección canónica, que se tiene la propiedad
siguiente: si, para todo $\alpha$ y para todo $\mathcal O_U$-módulo invertible
$\mathscr L$, $(j_\alpha)_*(\mathscr L|(U\cap W_\alpha))$ es un
$\mathcal O_{W_\alpha}$-módulo invertible, entonces $j_*(\mathscr L)$ es un
$\mathcal O_X$-módulo invertible. Pero la propiedad de ser un $\mathcal O_X$-módulo
invertible es local sobre $X$, y $(W_\alpha)$ es un recubrimiento abierto de $X$;
como se tiene
\[
  j_*(\mathscr L)|W_\alpha
  =(j_\alpha)_*(\mathscr L|(U\cap W_\alpha)),
\]
esto demuestra nuestra afirmación.

(ii) Pongamos $U'=X-Y'$, y sea $j':U'\to X$ la inyección canónica. Decir que
el homomorfismo $\mathcal O_X\to j'_*(\mathcal O_{U'})$ es biyectivo equivale a decir
que para todo abierto $V$ de $X$, el homomorfismo
$\Gamma(V,\mathcal O_X)\to\Gamma(V\cap U',\mathcal O_X)$ es biyectivo; pero
el homomorfismo compuesto
\[
  \Gamma(V,\mathcal O_X)\longrightarrow\Gamma(V\cap U',\mathcal O_X)
  \longrightarrow\Gamma(V\cap U,\mathcal O_X)
\]
es biyectivo por hipótesis (21.13.5), y al sustituir $V$ por $V\cap U'$, se ve que
$\Gamma(V\cap U',\mathcal O_X)\to\Gamma(V\cap U,\mathcal O_X)$ es biyectivo, por lo que
$\Gamma(V,\mathcal O_X)\to\Gamma(V\cap U',\mathcal O_X)$ es biyectivo. Observemos ahora
que si $j'':U\to U'$ es la inyección canónica y si $\mathscr L'$ es un
$\mathcal O_{U'}$-módulo invertible, $\mathscr L'|U$ es un $\mathcal O_U$-módulo
invertible, luego
$j_*(\mathscr L'|U)=j'_*(j''_*(\mathscr L'|U))$ es por hipótesis un
$\mathcal O_X$-módulo invertible. Como el par $(U',U'\cap Y)$ es parafactorial
en virtud de (i), se tiene $j''_*(\mathscr L'|U)\cong\mathscr L'$, luego
$j'_*(\mathscr L')$ es un $\mathcal O_X$-módulo invertible. Para concluir basta
sustituir en el razonamiento anterior $X$, $U$ y $U'$ por $V$,
$V\cap U$ y $V\cap U'$, donde $V$ es un abierto de $X$.

(iii) Pongamos $U'=X'-Y'=f^{-1}(U)$ y observemos que, puesto que se trata de preesquemas, se
puede escribir $U'=U\times_X X'$; sea $f_U:f^{-1}(U)\to U$ la restricción de $f$ y
sea $j':U'\to X'$ la inyección canónica, que también se escribe $j_{(X')}$. Demostremos
primero que el homomorfismo canónico
$\rho:\mathcal O_X\to j_*(\mathcal O_U)$ es biyectivo; para ello, observemos que por
hipótesis el homomorfismo canónico
$\mathcal O_{X'}\to j'_*(\mathcal O_{U'})$ es biyectivo; pero, puesto que el morfismo
$j$ es cuasicompacto y separado por hipótesis, y el morfismo $f$ es plano,
$j'_*(\mathcal O_{U'})=j'_*(f_U^*(\mathcal O_U))$ se identifica canónicamente con
$f^*(j_*(\mathcal O_U))$ en virtud de (2.3.1); como
$\mathcal O_{X'}=f^*(\mathcal O_X)$, se ve que el homomorfismo
$f^*(\rho):f^*(\mathcal O_X)\to f^*(j_*(\mathcal O_U))$ es biyectivo; se concluye
que $\rho$ es biyectivo puesto que $f$ es fielmente plano (2.2.7). Consideremos ahora
un $\mathcal O_U$-módulo invertible $\mathscr L$, y sea
$\mathscr L'=f_U^*(\mathscr L)$, que es un $\mathcal O_{U'}$-módulo invertible. Por
hipótesis $j'_*(f_U^*(\mathscr L))$ es un $\mathcal O_{X'}$-módulo invertible, que
es isomorfo a $f^*(j_*(\mathscr L))$ por (2.3.1), como antes. Pero, puesto que $f$
es fielmente plano y cuasicompacto, esto implica que $j_*(\mathscr L)$ es un
$\mathcal O_X$-módulo localmente libre (2.5.2). Para terminar la demostración basta
sustituir $X$ por un abierto $V$ y $X'$ por $f^{-1}(V)$ en lo que
precede.
\end{proof}

\begin{definition}[21.13.7]
\label{IV.21.13.7}
Se dice que un anillo local $A$ es \emph{parafactorial} si, poniendo
$X=\operatorname{Spec}(A)$ y designando por $a$ el punto cerrado de $X$, el par
$(X,\{a\})$ es parafactorial.
\end{definition}

\begin{proposition}[21.13.8]
\label{IV.21.13.8}
Con las notaciones de (21.13.7), se pone $U=X-\{a\}$. Para que $A$ sea
parafactorial, es necesario y suficiente que verifique las condiciones siguientes:
\begin{enumerate}[label=(\roman*)]
  \item El homomorfismo canónico
  $A=\Gamma(X,\mathcal O_X)\to\Gamma(U,\mathcal O_X)$ es biyectivo.
  \item $\operatorname{Pic}(U)=0$.
\end{enumerate}
\oldpage[IV-4]{317}
Si además $A$ es noetheriano, la condición (i) equivale a
\begin{enumerate}
  \item[\textup{(i bis)}] $\operatorname{prof}(A)\geq2$.
\end{enumerate}
\end{proposition}

\begin{proof}
En efecto, el único abierto de $X$ que contiene $a$ es el propio $X$, y por tanto todo
$\mathcal O_X$-módulo invertible es isomorfo a $\mathcal O_X$, dicho de otro modo,
$\operatorname{Pic}(X)=0$; la primera afirmación resulta entonces de la definición
(21.13.1) y de (21.13.3). La segunda afirmación no es sino un caso particular de
(21.13.4).
\end{proof}

\begin{env}[21.13.9]
\label{IV.21.13.9}
\emph{Ejemplos}. ---

(i) Un anillo local noetheriano parafactorial es necesariamente de dimensión
$\geq2$ en virtud de (21.13.8); dicho de otro modo, un anillo local noetheriano de
dimensión $\leq1$ no es parafactorial.

(ii) Un anillo local noetheriano \emph{factorial} $A$ de dimensión $\geq2$ es
parafactorial, como resulta de (21.13.8) y de (21.6.14).

(iii) Si $A$ es un anillo local noetheriano de dimensión $\geq3$ y parafactorial, no
es necesariamente normal, ni siquiera reducido. Consideremos un anillo local regular
$B$ de dimensión $\geq3$, y sea $A=D_B(B)$ $(\mathbf 0, 18.2.3)$, isomorfo a
$B[T]/(T^2)$; se ve inmediatamente que se tiene
$\operatorname{prof}(A)=\operatorname{prof}(B)=\dim(B)\geq3$. Para demostrar que $A$
es parafactorial basta, con las notaciones de (21.13.8), probar que
$\operatorname{Pic}(U)=0$. Sea $\mathfrak J$ el núcleo del aumento $A\to B$,
tal que $\mathfrak J^2=0$ y que, como $B$-módulo, es isomorfo a $B$.
Como $B=A/\mathfrak J$, $X=\operatorname{Spec}(A)$ y
$X_0=\operatorname{Spec}(B)$ tienen el mismo espacio subyacente; si
$\mathscr J=\widetilde{\mathfrak J}$, $X_0$ es el subesquema definido por
$\mathscr J$; denotaremos por $U_0$ el subpreesquema inducido por $X_0$ sobre el abierto
$U$. Para todo $z\in\mathfrak J$, pongamos $\varphi(z)=1+z$; puesto que
$(1+z)(1-z)=1$, $1+z$ es invertible en $A$ y $\varphi(\mathfrak J)$ es el núcleo
del homomorfismo sobreyectivo canónico de grupos multiplicativos $A^*\to B^*$; dicho
de otro modo, se tiene una sucesión exacta de grupos conmutativos
\[
  0\longrightarrow\mathfrak J\xrightarrow{\varphi}A^*\longrightarrow B^*
  \longrightarrow1
\]
(siendo multiplicativos los tres últimos grupos y aditivos los dos primeros). Por la
misma razón, para todo $t\in A$, se tiene la sucesión exacta
\[
  0\longrightarrow\mathfrak J_t\xrightarrow{\varphi_t}(A_t)^*
  \longrightarrow(B_{t_0})^*\longrightarrow1
\]
designando por $t_0$ la imagen de $t$ en $B$, puesto que
$B_{t_0}=A_t/\mathfrak J_t$; dicho de otro modo, se tiene una sucesión exacta de haces de
grupos conmutativos sobre el espacio topológico $X$
\[
  0\longrightarrow\mathscr J\xrightarrow{u}\mathcal O_X^*
  \longrightarrow\mathcal O_{X_0}^*\longrightarrow1
\]
de donde, por restricción al abierto $U$, una sucesión exacta
\[
  0\longrightarrow\mathscr J|U\longrightarrow\mathcal O_U^*
  \longrightarrow\mathcal O_{U_0}^*\longrightarrow1.
\]
De la sucesión exacta de cohomología se deduce una sucesión exacta
\[
\label{IV.21.13.9.1}
  \mathrm H^1(U,\mathscr J)\longrightarrow\mathrm H^1(U,\mathcal O_U^*)
  \xrightarrow{u_1}\mathrm H^1(U,\mathcal O_{U_0}^*)
  \tag{21.13.9.1}
\]
Pero, puesto que $\mathfrak J$ es un $B$-módulo isomorfo a $B$, se tiene
$\mathrm H^1(U,\mathscr J)=\mathrm H^1(U_0,\mathcal O_{U_0})$, y resulta del
cap.~III, 3\textsuperscript{a} parte (véase también [41, III, Ejemplo III-1]) que se tiene
\oldpage[IV-4]{318}
$\mathrm H^1(U_0,\mathcal O_{U_0})=0$ en razón de la relación
$\operatorname{prof}(B)\geq3$. Por otra parte, como $B$ es factorial, se tiene
$\mathrm H^1(U,\mathcal O_{U_0}^*)=\operatorname{Pic}(U_0)=0$; de la
sucesión exacta anterior se deduce, pues, que
$\operatorname{Pic}(U)=\mathrm H^1(U,\mathcal O_U^*)=0$.

(iv) Existen también anillos locales noetherianos de dimensión $3$ que son
íntegros, íntegramente cerrados y parafactoriales, pero no factoriales. Sea, en efecto, $B$
un anillo local noetheriano, completo, íntegro e íntegramente cerrado, de dimensión
$\geq2$ y no factorial (por ejemplo, el completado del anillo localizado del álgebra
íntegra $k[U,V,W]/(W^2-UV)$ en el ideal imagen de $(U)+(V)+(W)$). Resulta entonces de
lo que se verá más adelante (21.14.2) que el anillo local $A=B[[T]]$ de las series formales
en $T$ sobre $B$ es parafactorial, pero no es factorial, pues de lo contrario $B$ lo sería
en virtud de (21.13.12) más abajo.

(v) Se puede demostrar que un anillo local de intersección completa absoluta (19.3.1) de
dimensión $\geq4$ es parafactorial (cf. [41, XI, 3.13 (i)]).

(vi) Se ha visto (Observación (ii)) que todo anillo local noetheriano factorial $A$ de
dimensión $2$ es parafactorial. Pero hay anillos locales noetherianos de
dimensión $2$ que son parafactoriales y no factoriales. Se puede demostrar que, para
que un anillo local noetheriano $A$ de dimensión $2$ sea parafactorial y no
factorial, es necesario y suficiente que satisfaga las tres condiciones siguientes:
\begin{enumerate}[label=\arabic*\textsuperscript{o}]
  \item $A$ es un anillo de Cohen--Macaulay (dicho de otro modo,
  $\operatorname{prof}(A)=2$).
  \item $A$ es íntegro y, si $A'$ es su clausura integral, $A'$ es factorial y es
  un $A$-álgebra finita.
  \item Sea $\mathfrak J$ el conductor de $A$ en $A'$ (anulador del $A$-módulo
  $A'/A$, o también el mayor ideal de $A'$ contenido en $A$); pongamos
  $B=A/\mathfrak J$, $B'=A'/\mathfrak J$; entonces $\dim(B)=1$ (lo que implica
  $A'\neq A$, dicho de otro modo, $A$ no es íntegramente cerrado), y la aplicación
  canónica $\operatorname{Div}(B)\to\operatorname{Div}(B')$ (21.4.5) es
  \emph{sobreyectiva}.
\end{enumerate}
Se puede demostrar además que estas condiciones implican la propiedad siguiente:

$4^{\mathrm o}$ El anillo $B'$ (y \emph{a fortiori} $B\subset B'$) es reducido, y el
morfismo $\operatorname{Spec}(B')\to\operatorname{Spec}(B)$ es biyectivo.

Si se pone $X=\operatorname{Spec}(A)$, $X'=\operatorname{Spec}(A')$,
$Y=\operatorname{Spec}(B)$, $Y'=\operatorname{Spec}(B')$, de modo que $Y$
(resp.\ $Y'$) está definido por el ideal $\mathfrak J$ de $A$ (resp.\ $A'$), el morfismo
estructural $f:X'\to X$ es un isomorfismo de $X'-Y'$ sobre $X-Y$ (Bourbaki,
\emph{Alg. comm.}, cap.~V, \S1, n\textsuperscript{o}~5, cor.~5 de la prop.~16). Se
ve, pues, en virtud de $4^{\mathrm o}$ que $f$ es un morfismo \emph{biyectivo}; dicho
de otro modo, $X$ es un preesquema \emph{unirrama} (6.15.1) (y en particular
$A'$ es un anillo local noetheriano); en general $X$ no es geométricamente
unirrama. El espacio $Y$, de dimensión $1$, está constituido por el punto cerrado $x$ de
$X$ y los puntos maximales $y_i$ ($1\leq i\leq r$) de $Y$, y el único punto $y_i'$
de $Y'$ sobre $y_i$ es también un punto maximal de $Y'$; como $B$ y $B'$ son
reducidos, se deduce que se tiene
\[
  \mathfrak m_{y_i}\mathcal O_{X',y_i'}=\mathfrak m_{y_i'}.
\]
Para $z=y_i$ y $z'=y_i'$, esta relación también se verifica evidentemente para
$z\notin Y$ y $z'$ el único punto de $f^{-1}(z)$, por tanto para todo
$z\in U=X-\{x\}$ y $z'$ el único punto de $f^{-1}(z)$. En particular, si el anillo
$A$ es de característica $0$ (0, 21.1.1), se ve que $U'$ es \emph{no ramificado} sobre
$U$ (pero no étale en general).

Nos limitaremos aquí a demostrar que las condiciones $1^{\mathrm o}$,
$2^{\mathrm o}$ y $3^{\mathrm o}$ anteriores son \emph{suficientes} para que $A$ sea
parafactorial. Ahora bien, en virtud de la condición $1^{\mathrm o}$ y de (21.13.8), basta
demostrar que $\operatorname{Pic}(U)=0$. Puesto que $A'$ es factorial en virtud de
$2^{\mathrm o}$, se tiene $\operatorname{Pic}(U')=0$, y de (21.8.5, (ii)) se deduce una
sucesión exacta
\[
  1\longrightarrow
  \left(\prod_{i=1}^r
  \mathcal O_{U',y_i'}^*/\mathcal O_{U,y_i}^*\right)
  \big/\operatorname{Im}\Gamma(U',\mathcal O_{U'}^*)
  \longrightarrow\operatorname{Pic}(U)\longrightarrow\operatorname{Pic}(U')
  \longrightarrow0
\]
y se trata, pues, de demostrar que el segundo término de esta sucesión se reduce al elemento
neutro, tomando como $S$ el conjunto de los $y_i$ ($1\leq i\leq r$) y observando
que aquí $f_*(\mathcal O_{X'})$ es la $\mathcal O_X$-álgebra $\widetilde{A'}$. Sean
$\mathfrak p_i$, $\mathfrak p_i'$ los ideales primos de $A$ y $A'$ correspondientes
a los puntos $y_i$, $y_i'$, y observemos que $\mathfrak p_i'\cap A=\mathfrak p_i$;
para cada $i$, démonos un elemento invertible $a_i'/s_i$ de
$A'_{\mathfrak p_i'}$, con $a_i'\in A'$, $s_i\notin\mathfrak p_i'$; como solo
tenemos que examinar los grupos cociente
$(A'_{\mathfrak p_i'})^*/(A_{\mathfrak p_i})^*$, se puede suponer $s_i=1$ para todo
$i$; sea $b_i'\in B'=A'/\mathfrak J$ la imagen canónica de $a_i'$, de modo que
$b_i'/1$ es un elemento no nulo de $k(y_i')=\mathcal O_{Y',y_i'}$. El conjunto
$U\cap Y$ formado por los $y_i$ es un abierto afín de la forma $D(t)$ en $Y$, con
$t\in\mathfrak m$, ideal maximal de $A$; existe, pues, un elemento invertible
$b'\in B'_t$ cuyas imágenes canónicas son los $b_i'/1$, y como $t$ es regular
en el anillo íntegro $A'$, se puede suponer que $b'$ es de la forma $b''/t^m$, donde
$b''\in B'$ es invertible. Sea $a''$ un elemento de $A'$ de clase $b''$, que
es, pues, necesariamente invertible; $a'=a''/t^m$ es invertible en $A'_t$ y para
todo $i$ se tiene
\[
  u_i=a''-a_i't^m\in\mathfrak J,
  \qquad
  t^m a_i'=a''-u_i=a''(1-a''^{-1}u_i);
\]
pero $a''^{-1}u_i\in\mathfrak J\subset\mathfrak m$, luego $1-a''^{-1}u_i$ es un
elemento invertible de $A$, y las clases de $a'$ y $a_i'$ en
$(A'_{\mathfrak p_i'})^*/(A_{\mathfrak p_i})^*$ son las mismas, lo que termina la
demostración.

\oldpage[IV-4]{319}
Para obtener un ejemplo explícito de anillo parafactorial de dimensión $2$ obtenido de
esta manera y \emph{no factorial}, consideremos el anillo
\[
  E=\mathbb R[[U,V]]/(U^2+V^2),
\]
cuya clausura integral $E'$ se identifica con $\mathbb C[[U]]$ (6.15.11). Si se pone
$A=E[[T]]$, la clausura integral de $A$ es el anillo $A'=E'[[T]]$ (Bourbaki,
\emph{Alg. comm.}, cap.~V, \S1, n\textsuperscript{o}~4, prop.~14); se comprueba
inmediatamente que el conductor de $E$ en $E'$ es el ideal maximal $\mathfrak n$ de
$E$, luego el conductor $\mathfrak J$ de $A$ en $A'$ es $\mathfrak n[[T]]$, y se
tiene $B=A/\mathfrak J\cong\mathbb R[[T]]$, $B'=A'/\mathfrak J\cong\mathbb C[[T]]$.
Es entonces inmediato que las condiciones $1^{\mathrm o}$, $2^{\mathrm o}$ y
$3^{\mathrm o}$ enunciadas anteriormente se verifican, pero $A$ ni siquiera es
íntegramente cerrado.

Se puede variar este ejemplo, y el lector verá fácilmente que, si $k$ es un cuerpo
algebraicamente cerrado, el anillo $A$, localización del anillo
$k[U,V,W]/(U^2-WV^2)$ en el ideal maximal engendrado por las imágenes de $U$, $V$ y $W$,
verifica también las condiciones $1^{\mathrm o}$, $2^{\mathrm o}$ y $3^{\mathrm o}$
anteriores.

Podría suceder incluso que estas tres condiciones implicaran que cada uno de los anillos
$B'/\mathfrak p_i'$ fuese un anillo de valoración discreta, lo que implicaría que
$A'$ fuese incluso un anillo \emph{regular}.
\end{env}

\begin{proposition}[21.13.10]
\label{IV.21.13.10}
Sean $X$ un preesquema localmente noetheriano e $Y$ una parte cerrada de $X$. Para que
el par $(X,Y)$ sea parafactorial, es necesario y suficiente que, para todo $y\in Y$,
el anillo local $\mathcal O_{X,y}$ sea parafactorial.
\end{proposition}

\begin{proof}
Pongamos $U=X-Y$ y sea $j:U\to X$ la inyección canónica.

En virtud de (21.13.4), decir que el homomorfismo canónico
$\mathcal O_X\to j_*(\mathcal O_U)$ es biyectivo equivale a decir que cada uno de los
anillos locales $\mathcal O_{X,y}$ para $y\in Y$ satisface la condición (i bis) de
(21.13.8).

Demostremos primero que, si el par $(X,Y)$ es parafactorial, cada uno de los anillos
$\mathcal O_{X,y}$ ($y\in Y$) es parafactorial. Teniendo en cuenta la observación
anterior, todo se reduce a ver que, si se pone
$T_y=\operatorname{Spec}(\mathcal O_{X,y})$ y $U_y=T_y-\{y\}$, todo
$\mathcal O_{U_y}$-módulo invertible $\mathscr L_0$ es isomorfo a
$\mathcal O_{U_y}$. Ahora bien, cuando $V$ recorre el conjunto de los entornos abiertos
afines de $y$ en $X$, resulta de (8.2.13) y $(\mathbf I, 2.4.2)$ que el
preesquema $U_y$ es límite proyectivo de los preesquemas inducidos por $X$ sobre los abiertos
$V-(V\cap\overline{\{y\}})$, que son cuasicompactos puesto que $X$ es localmente
noetheriano. Como los $V-(V\cap\overline{\{y\}})$ son separados, resulta entonces de
(8.5.2, (ii)) y (8.5.5) que hay un entorno abierto afín $V$ de $y$ en $X$ y
un $\mathcal O_{U_V}$-módulo invertible $\mathscr L'_V$ tal que $\mathscr L_0$ es el
haz inducido por $\mathscr L'_V$ sobre $U_y$. Pero la hipótesis implica, en virtud de
(21.13.6, (i) y (ii))\footnote{El francés impreso cita (i) y (iii), y el inglés mantenido solo (i); aquí (ii) es la propiedad
hereditaria para una parte cerrada menor.}, que el par $(V,V\cap\overline{\{y\}})$ es parafactorial;
se concluye por la definición (21.13.1) que existe un $\mathcal O_V$-módulo invertible
$\mathscr L_V$ que induce $\mathscr L'_V$ sobre $U_V$. Sustituyendo si es preciso $V$ por un
entorno abierto más pequeño de $y$, se puede suponer que $\mathscr L_V$ es isomorfo
a $\mathcal O_V$, de donde se concluye que $\mathscr L_0$ es isomorfo a
$\mathcal O_{U_y}$.

Recíprocamente, demostremos que, si todos los $\mathcal O_{X,y}$ ($y\in Y$) son
parafactoriales, el par $(X,Y)$ es parafactorial. Teniendo en cuenta la observación
inicial, todo se reduce evidentemente a demostrar que para todo $\mathcal O_U$-módulo invertible
$\mathscr L$, $j_*(\mathscr L)$ es un $\mathcal O_X$-módulo invertible (21.13.5).
Como la cuestión es local sobre $X$, se puede suponer que $X$ es noetheriano. El conjunto $V$ de los
$x\in X$ tales que la restricción de $j_*(\mathscr L)$ a un entorno abierto de $x$
en $X$ es invertible es evidentemente abierto en $X$; se trata de demostrar que
$V=X$. Para ello, supongamos lo contrario y pongamos $Z=X-V\neq\varnothing$. Sea $y$
un punto \emph{maximal} de $Z$; puesto que $Z\subset Y$ por definición,
$\mathcal O_{X,y}$ es por hipótesis un anillo parafactorial. Sustituyendo si es preciso
$X$ por un entorno abierto de $y$ en $X$, se puede suponer que la restricción de
$j_*(\mathscr L)$
\oldpage[IV-4]{320}
a $X-\overline{\{y\}}$\footnote{El francés impreso y el inglés mantenido escriben aquí el locus con $V$; como $y\notin V$, tras la reducción
al entorno se restablece el abierto punteado con $X$.} es invertible; por tanto, con las notaciones introducidas
anteriormente, el $\mathcal O_{U_y}$-módulo $\mathscr L_0$ inducido por $j_*(\mathscr L)$
sobre $U_y$ es invertible. Puesto que $\mathcal O_{X,y}$ es parafactorial,
$\mathscr L_0$ está inducido por un $\mathcal O_{T_y}$-módulo invertible $\mathscr L'$;
pero como $T_y$ es límite proyectivo de los entornos abiertos $W$ de $y$ en $X$,
se pueden aplicar de nuevo (8.5.2, (ii)) y (8.5.5), lo que prueba la existencia de tal
entorno $W$ y de un $\mathcal O_W$-módulo invertible $\mathscr L''$ que induce
$\mathscr L'$ sobre $T_y$, y por tanto $\mathscr L_0$ sobre $U_y$; aplicando esta
vez (8.5.2.5) y (8.5.2, (i)), se deduce que, sustituyendo si es preciso $W$ por un
entorno abierto más pequeño de $y$, se puede suponer que las restricciones de
$j_*(\mathscr L)$ y de $\mathscr L''$ a
$W-(W\cap\overline{\{y\}})$ son iguales. Si se pone entonces $V_1=V\cup W$ y se
designa por $\mathscr L_1$ el $\mathcal O_{V_1}$-módulo invertible igual a
$\mathscr L''$ en $W$ y a $j_*(\mathscr L)|V$ en $V$, se tiene $U\subset V_1$ y
$\mathscr L_1|U=\mathscr L$. Se concluye entonces del hecho de que el homomorfismo
$\mathcal O_X\to j_*(\mathcal O_U)$ es biyectivo y de (21.13.2, $b$)) que
$j_*(\mathscr L)|V_1$ es isomorfo a $\mathscr L_1$, luego invertible. Pero esto
contradice la definición de $V$, y concluye así la demostración de (21.13.10).
\end{proof}

\begin{corollary}[21.13.11]
\label{IV.21.13.11}
Sean $X$ un preesquema localmente noetheriano, $Y$ una parte cerrada de $X$ y
$U=X-Y$. Supongamos que el par $(X,Y)$ es parafactorial y que $U$ es
localmente factorial; dicho de otro modo (21.13.10), para todo $x\in X$, el anillo
$\mathcal O_{X,x}$ es: $1^{\mathrm o}$ parafactorial si $x\in Y$;
$2^{\mathrm o}$ factorial si $x\notin Y$. Entonces $X$ es localmente factorial
(dicho de otro modo, $\mathcal O_{X,x}$ es de hecho factorial para todo $x\in X$).
\end{corollary}

\begin{proof}
Supongamos, en efecto, lo contrario, y sea $y$ un punto de $Y$ tal que
$\mathcal O_{X,y}$ no es factorial y tiene dimensión \emph{mínima} entre
todos los puntos de $Y$ que tienen esta propiedad. Entonces, si se pone
$T_y=\operatorname{Spec}(\mathcal O_{X,y})$, $U_y=T_y-\{y\}$, la hipótesis de que
$\mathcal O_{X,y}$ es parafactorial implica $\operatorname{Pic}(U_y)=0$ y
$\operatorname{prof}(\mathcal O_{X,y})\geq2$ (21.13.8); además, la elección de $y$
implica que $U_y$ es localmente factorial. Pero entonces (21.6.14) prueba que
$\mathcal O_{X,y}$ es factorial, contrariamente a la hipótesis.
\end{proof}

\begin{proposition}[21.13.12]
\label{IV.21.13.12}
Sean $A$, $B$ dos anillos locales noetherianos y
$\rho:A\to B$ un homomorfismo local que hace de $B$ un $A$-módulo plano. Entonces, si
$B$ es factorial, también lo es $A$.
\end{proposition}

\begin{proof}
Ya se sabe que $A$ es íntegro e íntegramente cerrado (6.5.4), por lo que se puede limitar
al caso en que $\dim(A)\geq2$, y razonar por inducción sobre $n=\dim(A)$. Sean
$X=\operatorname{Spec}(A)$, $a$ el punto cerrado de $X$,
$X'=\operatorname{Spec}(B)$, $f:X'\to X$ el morfismo estructural,
$U=X-\{a\}$ y $U'=f^{-1}(U)$.

Los anillos locales $\mathcal O_{X',x'}$ en los puntos $x'\in f^{-1}(a)$ son factoriales
y de dimensión $\geq2$ (6.1.2), luego parafactoriales (21.13.9, (ii)); se concluye de
(21.13.10) que el par $(X',f^{-1}(a))$ es parafactorial. En virtud de (21.13.6,
(iii)), esto demuestra que el anillo $A$ es parafactorial, luego
$\operatorname{prof}(A)\geq2$ y $\operatorname{Pic}(U)=0$ (21.13.8). Finalmente, como los
anillos locales de $U$ son de dimensión $<n$, la hipótesis de inducción prueba que
$U$ es \emph{localmente factorial}; la conclusión resulta entonces de (21.6.14).
\end{proof}

\begin{env}[21.13.12.1]
\label{IV.21.13.12.1}
Obsérvese que el enunciado (21.13.12) en el que se sustituye «factorial» por «parafactorial»
ya no es exacto, como muestra el ejemplo en que $A=k$ es un cuerpo y $B$ el anillo
local factorial $k[[T_1,T_2]]$. Sin embargo, resulta de (21.13.6, (iii)) que, bajo
las condiciones de (21.13.12), si $\mathfrak m$ designa el ideal maximal de $A$ y si
$\mathfrak mB$ es un \emph{ideal}
\oldpage[IV-4]{321}
\emph{de definición} de $B$ (dicho de otro modo, si $\dim(B/\mathfrak mB)=0$), entonces,
cuando $B$ es parafactorial, también lo es $A$. En particular, si el
completado $\widehat A$ de un anillo local noetheriano es parafactorial, también lo es
$A$.
\end{env}

\begin{remark}[21.13.13]
\label{IV.21.13.13}
Una parte de las definiciones y resultados anteriores está vinculada a consideraciones
más generales sobre la cohomología de los haces de grupos conmutativos sobre un espacio
topológico, desarrolladas en el cap.~III, 3\textsuperscript{a} parte, y solo hemos
dado aquí una exposición independiente de ellas para comodidad del lector. En efecto, si $X$
es un espacio anillado, se tiene una \emph{equivalencia} canónica entre la categoría de los
$\mathcal O_X$-módulos invertibles y la de los \emph{haces homogéneos principales
bajo el haz de grupos} $\mathcal O_X^*$ (16.5.15). Esta equivalencia se define
haciendo corresponder a todo $\mathcal O_X$-módulo invertible $\mathscr L$ el
haz $\mathscr L^*=\mathscr Isom_{\mathcal O_X}(\mathcal O_X,\mathscr L)$ de los
gérmenes de isomorfismos de $\mathcal O_X$ sobre $\mathscr L$; es inmediato que
$\mathscr L^*$ está provisto canónicamente de una estructura de haz homogéneo principal
bajo $\mathcal O_X^*\cong\mathscr Isom_{\mathcal O_X}(\mathcal O_X,\mathcal O_X)$.
La verificación de que el functor $\mathscr L\mapsto\mathscr L^*$ es una
equivalencia de categorías es inmediata.

Dicho esto, consideremos en general un espacio topológico $X$ y un haz de
grupos $\mathscr G$ sobre $X$; sea $Y$ una parte cerrada de $X$, pongamos $U=X-Y$, y
sea $i$ un entero tal que $0\leq i\leq2$; se dirá que el par $(X,Y)$ es
\emph{$i$-puro para $\mathscr G$} si, para todo abierto $V$ de $X$, el functor de
restricción $\mathscr P\mapsto\mathscr P|(U\cap V)$, de la categoría de haces
homogéneos principales bajo el haz de grupos $\mathscr G|V$ a la categoría
análoga bajo el haz de grupos $\mathscr G|(U\cap V)$, es \emph{fiel}
($i=0$), resp.\ \emph{plenamente fiel} ($i=1$), resp.\ una \emph{equivalencia de
categorías} ($i=2$). En el caso en que $X$ es un espacio anillado y
$\mathscr G=\mathcal O_X^*$, decir que $(X,Y)$ es $i$-puro significa, pues, para $i=0$,
que el homomorfismo $\mathcal O_X\to j_*(\mathcal O_U)$ es \emph{inyectivo}; para
$i=1$, que este homomorfismo es \emph{biyectivo}; y finalmente, para $i=2$, que el par
$(X,Y)$ es \emph{parafactorial} (21.13.3).

Volviendo al caso general, recordemos que los morfismos de la categoría de los haces
principales bajo $\mathscr G$ son por definición los isomorfismos. Decir que el par
$(X,Y)$ es $0$-puro significa, pues, que para todo abierto $V$ de $X$, el homomorfismo
canónico $\mathrm H^0(V,\mathscr G)\to\mathrm H^0(U\cap V,\mathscr G)$ es inyectivo;
decir que el par $(X,Y)$ es $1$-puro significa que el homomorfismo canónico
$\mathrm H^0(V,\mathscr G)\to\mathrm H^0(U\cap V,\mathscr G)$ es biyectivo (y entonces
el homomorfismo canónico
$\mathrm H^1(V,\mathscr G)\to\mathrm H^1(U\cap V,\mathscr G)$ es inyectivo); finalmente,
se puede demostrar que para que el par $(X,Y)$ sea $2$-puro es necesario y suficiente que
los homomorfismos
$\mathrm H^i(V,\mathscr G)\to\mathrm H^i(U\cap V,\mathscr G)$ sean biyectivos para
$i=0$ e $i=1$. Cuando $\mathscr G$ es un haz de grupos conmutativos, al
introducir los haces de cohomología $\mathscr H_Y^p(\mathscr G)$ definidos en el
cap.~III, 3\textsuperscript{a} parte, decir que $(X,Y)$ es $i$-puro para $i\leq2$
significa que $\mathscr H_Y^p(\mathscr G)=0$ para $p\leq i$; en esta forma, la
noción se generaliza inmediatamente a todo entero $i$.
\end{remark}

\begin{proposition}[21.13.14]
\label{IV.21.13.14}
Sean $X$ un preesquema localmente noetheriano \emph{normal} y
$(U_\lambda)_{\lambda\in L}$ una familia filtrante decreciente de abiertos de $X$. Las
condiciones siguientes son equivalentes:
\begin{enumerate}[label=\emph{\alph*)}]
  \item Todo ciclo $1$-codimensional $Z$ sobre $X$ tal que exista un índice
  $\lambda$ para el cual $Z|U_\lambda$ sea localmente principal (21.6.7), es
  localmente principal.
  \item Para todo $x\in X$ tal que $\dim(\mathcal O_{X,x})\geq2$ y que $x$
  no pertenezca a $\bigcap_{\lambda\in L}U_\lambda$, el anillo
  $\mathcal O_{X,x}$ es parafactorial.
\end{enumerate}
\end{proposition}

\begin{proof}
\oldpage[IV-4]{322}
La condición $b)$ puede expresarse también del modo siguiente:

\emph{$b')$ Para toda parte cerrada $Y$ de $X$ tal que
$\operatorname{codim}(Y,X)\geq2$ y tal que exista $\lambda$ para el cual
$Y\subset X-U_\lambda$, el par $(X,Y)$ es parafactorial.}

La propiedad $a)$ puede expresarse también del modo siguiente: si el conjunto cerrado
$N$ de los puntos $x\in X$ en los que $Z$ no es principal está contenido en uno de los
$X-U_\lambda$, entonces $Z$ es localmente principal. Como $X$ es normal, la condición
$\dim(\mathcal O_{X,x})\leq1$ implica que $\mathcal O_{X,x}$ es un cuerpo o un
anillo de valoración discreta, luego $Z_x$ es principal; dicho de otro modo, se tiene
necesariamente $\operatorname{codim}(N,X)\geq2$ (5.1.3). La propiedad $a)$ es, pues,
equivalente a la siguiente:

\emph{$a')$ Si existe un conjunto cerrado $Y$ contenido en uno de los $X-U_\lambda$, tal
que $\operatorname{codim}(Y,X)\geq2$ y que $Z|(X-Y)$ sea localmente principal,
entonces $Z$ es localmente principal.}

Observemos, por otra parte, que $b)$ implica $b')$ en virtud de (21.13.10); recíprocamente, si
$b')$ se verifica y si $x\notin U_\lambda$, entonces $Y=\overline{\{x\}}$\footnote{El francés impreso y el inglés mantenido escriben $Y=\{x\}$,
que no es necesariamente cerrado; se restablece aquí su clausura.} está contenida en
$X-U_\lambda$ y se tiene $\operatorname{codim}(Y,X)\geq2$, luego resulta nuevamente de
(21.13.10) que $\mathcal O_{X,x}$ es parafactorial, lo que prueba la equivalencia de
$b)$ y $b')$. Se queda así reducido a probar la equivalencia de $a')$ y $b')$.
Observemos que, puesto que $X$ es normal, se tiene, para toda parte cerrada $Y$ de $X$ tal que
$\operatorname{codim}(Y,X)\geq2$, $\operatorname{prof}_Y(\mathcal O_X)\geq2$
(5.8.6); si se pone $U=X-Y$, el homomorfismo
$\mathcal O_X\to j_*(\mathcal O_U)$ es, pues, biyectivo (21.13.4); como las condiciones
$a')$ y $b')$ son locales, se ve que basta demostrar la equivalencia de las
condiciones siguientes, cuando $X$ es noetheriano y normal e $Y$ es cerrado en $X$ y
tal que $\operatorname{codim}(Y,X)\geq2$:

\emph{$a'')$ Para todo ciclo $1$-codimensional $Z$ sobre $X$, la hipótesis de que $Z|U$
(donde $U=X-Y$) es localmente principal implica que el propio $Z$ es localmente
principal.}

\emph{$b'')$ El homomorfismo canónico
$\operatorname{Pic}(X)\to\operatorname{Pic}(U)$ es sobreyectivo.}

Demostremos primero que $a'')$ implica $b'')$; un elemento
$\zeta_0\in\operatorname{Pic}(U)$ tiene por imagen en $\mathfrak{Cl}(U)$ (21.6.11) la
clase de un ciclo $1$-codimensional $Z_0$ sobre $U$, que es localmente principal.
La hipótesis $\operatorname{codim}(Y,X)\geq2$ implica que el homomorfismo de
restricción $\mathscr Z^1(X)\to\mathscr Z^1(U)$ es biyectivo, luego $Z_0=Z|U$, donde $Z$
es un ciclo $1$-codimensional sobre $X$; puesto que $Z_0$ es localmente principal, también lo
es $Z$ en virtud de $a'')$; resulta de (21.6.11) que la imagen de $Z$ en
$\mathfrak{Cl}(X)$ es la imagen de un único elemento $\zeta\in\operatorname{Pic}(X)$,
y entonces está claro que $\zeta_0$ es la imagen de $\zeta$.

Recíprocamente, demostremos que $b'')$ implica $a'')$. Sea, pues, $Z$ un ciclo
$1$-codimensional sobre $X$ tal que $Z|U$ sea localmente principal; la imagen de $Z|U$
en $\mathfrak{Cl}(U)$ es, pues, la imagen de un único elemento
$\zeta_0\in\operatorname{Pic}(U)$ (21.6.11). Por hipótesis existe
$\zeta\in\operatorname{Pic}(X)$ cuya imagen en $\operatorname{Pic}(U)$ es
$\zeta_0$; la imagen de $\zeta$ en $\mathfrak{Cl}(X)$ es entonces la clase de un ciclo
$1$-codimensional $Z'$ sobre $X$ tal que $Z|U$ y $Z'|U$ son linealmente
equivalentes. Pero, puesto que $U$ es esquemáticamente denso en $X$, la imagen de
$\mathscr Z^1_{\mathrm{princ}}(X)$ por el isomorfismo
$\mathscr Z^1(X)\to\mathscr Z^1(U)$ es $\mathscr Z^1_{\mathrm{princ}}(U)$, luego $Z$
y $Z'$ son linealmente equivalentes, y puesto que $Z'$ es localmente principal, también lo
es $Z$.
\end{proof}

\begin{corollary}[21.13.15]
\label{IV.21.13.15}
Sean $X$ un preesquema localmente noetheriano y normal y $S$ una parte de $X$. Las
condiciones siguientes son equivalentes:
\begin{enumerate}[label=\emph{\alph*)}]
  \item Todo ciclo $1$-codimensional sobre $X$ que es principal en los puntos de $S$ es
  localmente principal.
  \oldpage[IV-4]{323}
  \item Para todo $x\in X$ que no es generalización de un punto de $S$ (dicho de otro
  modo, tal que $\overline{\{x\}}\cap S=\varnothing$) y que es tal que
  $\dim(\mathcal O_{X,x})\geq2$, el anillo $\mathcal O_{X,x}$ es parafactorial.
\end{enumerate}
\end{corollary}

\begin{proof}
Está claro que el conjunto $S'$ de los puntos de $X$ que son generalizaciones de puntos de $S$ es
la intersección de los entornos abiertos de $S$. Se puede limitar al caso en que $X$ es
noetheriano (pues las propiedades $a)$ y $b)$ son locales sobre $X$); como todo ciclo
$1$-codimensional sobre $X$ que es principal en los puntos de $S$ lo es también en los puntos
de un entorno abierto de $S$, se ve que la condición $a)$ de (21.13.15) equivale a
la condición $a)$ de (21.13.14) aplicada a la familia $(U_\lambda)$ de los entornos
abiertos de $S$. El corolario resulta entonces de la equivalencia de $a)$ y $b)$ en
(21.13.14).
\end{proof}

\begin{remark}[21.13.16]
\label{IV.21.13.16}
Bajo las condiciones generales de (21.13.14), supongamos además que se tiene
$\varinjlim\operatorname{Pic}(U_\lambda)=0$. Entonces la condición $a)$ de (21.13.14)
es también equivalente a la siguiente:
\begin{enumerate}
  \item[c)] \emph{Todo ciclo $1$-codimensional sobre $X$ cuyo soporte está contenido
  en uno de los conjuntos $X-U_\lambda$ es localmente principal.}
\end{enumerate}

Se puede limitar al caso en que $X$ es irreducible y todos los $U_\lambda$ son no
vacíos. Está claro que $a)$ implica $c)$, pues si el ciclo $1$-codimensional $Z$ sobre
$X$ tiene su soporte en $X-U_\lambda$, $Z|U_\lambda=0$, luego $Z|U_\lambda$ es
localmente principal. Recíprocamente, $c)$ implica $a)$: sea, en efecto, $Z$ un ciclo
$1$-codimensional sobre $X$ tal que $Z|U_\lambda$ sea localmente principal; como $X$
es normal, $Z|U_\lambda=\operatorname{cyc}(D_\lambda)$, donde $D_\lambda$ es un
divisor sobre $U_\lambda$ (21.6.10, (ii)), y la hipótesis implica (en virtud de
(21.3.4)) que hay un conjunto $U_\mu\subset U_\lambda$ tal que
$D_\mu=D_\lambda|U_\mu$ es equivalente a $0$. Si
$D_\mu=\operatorname{div}(f_\mu)$, donde $f_\mu$ es una función racional regular
sobre $U_\mu$, $f_\mu$ puede considerarse como una función racional regular
sobre $X$. Si $Z'=\operatorname{cyc}(\operatorname{div}(f_\mu))$, se ve entonces que
$Z''=Z-Z'$ tiene su soporte en $X-U_\mu$; en virtud de la hipótesis, $Z''$ es
localmente principal, de donde la conclusión.

La condición $\varinjlim\operatorname{Pic}(U_\lambda)=0$ se cumplirá en particular
si $(U_\lambda)$ es la familia de los entornos abiertos de un \emph{punto} $z\in X$,
pues para todo $\mathcal O_{U_\lambda}$-módulo invertible $\mathscr L_\lambda$, existe
por definición un $U_\mu\subset U_\lambda$ tal que
$\mathscr L_\lambda|U_\mu$ es isomorfo a $\mathcal O_{U_\mu}$. Se ve, pues, que,
en el enunciado de (21.13.15), si $S=\{z\}$, las condiciones $a)$ y $b)$ equivalen
también a la siguiente:
\begin{enumerate}
  \item[c)] \emph{Todo ciclo $1$-codimensional sobre $X$ cuyo soporte no contiene
  a $z$ es localmente principal.}
\end{enumerate}
Si además $\operatorname{Pic}(X)=0$, se concluirá que esta condición implica que
todo ciclo $1$-codimensional cuyo soporte no contiene a $z$ es principal.
\end{remark}

\subsection{El teorema de Ramanujam--Samuel}
\label{subsection:IV.21.14}
\label{IV.21.14}

\begin{theorem}[21.14.1]
\label{IV.21.14.1}
\textup{(Ramanujam--Samuel)}. ---
Sea $A$ un anillo local noetheriano con ideal maximal $\mathfrak m$, tal que su completado
$\widehat A$ sea íntegro e integralmente cerrado. Sean $B$ un anillo local noetheriano
tal que $\dim(B)\geq\dim(A)$ y $\rho:A\to B$ un homomorfismo local que hace de $B$
una $A$-álgebra formalmente lisa (para las topologías preádicas
$(\mathbf 0, 19.3.1)$), y tal que el cuerpo residual de $B$ sea finito sobre el de $A$.
Entonces todo ciclo $1$-codimensional sobre $\operatorname{Spec}(B)$ que sea principal en el
punto $\mathfrak p=\mathfrak mB$ es un ciclo $1$-codimensional principal.
\end{theorem}

\begin{proof}
Si $k=A/\mathfrak m$ es el cuerpo residual de $A$, $B/\mathfrak mB$ es una
$k$-álgebra formalmente lisa (para su topología preádica) $(\mathbf 0, 19.3.5)$,
\oldpage[IV-4]{324}
por tanto regular y, en particular, íntegra; en otras palabras,
$\mathfrak p=\mathfrak mB$ es efectivamente un ideal primo de $B$, lo que justifica
el enunciado. Basta evidentemente probar que todo ideal primo $\mathfrak q$ de altura $1$\footnote{Corrección de fuente: el francés impreso y el inglés mantenido omiten «de altura $1$». La conclusión concierne a los ciclos de codimensión $1$, y un ideal primo principal no nulo tiene altura $1$.} de
$B$ no contenido en $\mathfrak p$ es \emph{principal}.

Sean $\widehat A$ y $\widehat B$ los completados de $A$ y $B$, respectivamente, de
modo que el ideal maximal de $\widehat A$ es $\mathfrak m\widehat A$; se sabe
$(\mathbf 0, 19.3.6)$ que $\widehat B$ es una $\widehat A$-álgebra formalmente
lisa para las topologías ádicas. Sea $k'$ el cuerpo residual de $\widehat B$,
extensión finita de $k$; existe un homomorfismo local $\widehat A\to C$, donde $C$
es un anillo local noetheriano que es un $\widehat A$-módulo \emph{finito} y
\emph{plano} (por tanto libre), y tal que $C/\mathfrak mC$ es isomorfo a $k'$
$(\mathbf 0_{\mathrm{III}}, 10.3.1)$; de ello se deduce que $C$ es \emph{completo}, y
resulta entonces de (7.5.1), (7.5.3) y (6.5.4, (ii)) que $C$ es \emph{íntegro e
integralmente cerrado}. Además,
\[
  D=\widehat B\widehat\otimes_{\widehat A}C
\]
es un anillo semilocal completo, producto directo de anillos locales completos, uno de los cuales,
$D_0$, tiene por cuerpo residual $k'$ (puesto que $k'\otimes_k k'$ es producto directo
de anillos locales, uno de los cuales es isomorfo a $k'$). Como $D$ es una $C$-álgebra formalmente lisa,
lo mismo ocurre con $D_0$; por consiguiente, $D_0/\mathfrak mD_0$ es una
$k'$-álgebra formalmente lisa, con cuerpo residual $k'$, lo que implica que es
$k'$-isomorfa a un álgebra de series formales
$k'[[T_1,\ldots,T_n]]$ $(\mathbf 0, 19.6.4)$; se concluye, por
$(\mathbf 0, 19.7.1.5)$, que $D_0$ es $C$-isomorfo a
$C[[T_1,\ldots,T_n]]$ y, por tanto, \emph{íntegro e integralmente cerrado} (Bourbaki,
\emph{Alg. comm.}, cap.~V, \S1, n\textsuperscript{o}~4, prop.~14). Como los
morfismos $\operatorname{Spec}(D_0)\to\operatorname{Spec}(\widehat B)$ y
$\operatorname{Spec}(D_0)\to\operatorname{Spec}(B)$ son fielmente planos, se
deduce que $B$ y $\widehat B$ son también \emph{íntegros e integralmente cerrados}
(2.1.13). Esto prueba que el ideal $\mathfrak qD_0$ es divisorial (Bourbaki,
\emph{Alg. comm.}, cap.~VII, \S1, n\textsuperscript{o}~10, prop.~15) y no está contenido
en $\mathfrak mD_0$, pues de lo contrario se tendría
\[
  \mathfrak q=(\mathfrak qD_0)\cap B
  \subset(\mathfrak mD_0)\cap B=\mathfrak mB=\mathfrak p
\]
por planitud fiel $(\mathbf 0_{\mathrm I}, 6.5.1)$. Se concluye que ninguno de los
ideales primos $\mathfrak r_h$ de altura $1$ de $D_0$ que contienen
$\mathfrak qD_0$ puede estar contenido en $\mathfrak mD_0$; si se prueba que estos
ideales son principales, resultará que $\mathfrak qD_0$ es principal, pues los
divisores (en el sentido de Bourbaki) de $\mathfrak qD_0$ y de un producto de potencias de los
$\mathfrak r_h$ son iguales (Bourbaki, \emph{Alg. comm.}, cap.~VII, \S1,
n\textsuperscript{o}~4, prop.~5). Como $\mathfrak qD_0\cap B=\mathfrak q$, se
deduce por planitud fiel que $\mathfrak q$ es principal (2.5.2), utilizando el
hecho de que $B$ es un anillo local.

Quedamos así reducidos a probar el teorema en el caso particular en que $A$ es
\emph{completo} y $B=A[[T_1,\ldots,T_n]]$. Probemos primero que se puede reducir al
caso en que $n=1$. En efecto, procedamos por inducción sobre $n$ y sea
$f\in\mathfrak q$ un elemento que no pertenece a $\mathfrak p=\mathfrak mB$; bastará
probar que existe un $A$-automorfismo $\sigma$ de $B$ tal que
$\sigma(f)$ no pertenece al ideal
\[
  \mathfrak p'=\mathfrak p+BT_1+\cdots+BT_{n-1};
\]
en efecto, si $B'=A[[T_1,\ldots,T_{n-1}]]$, $B'$ es completo y tiene por ideal maximal
$\mathfrak mB'+B'T_1+\cdots+B'T_{n-1}=\mathfrak m'$. Como
$B=B'[[T_n]]$ y $\mathfrak p'=\mathfrak m'B$, resultará de la hipótesis de
inducción (teniendo en cuenta que $B'$ es íntegro e integralmente cerrado) que el ideal
$\sigma(\mathfrak q)$ es principal en $B$, y por tanto también $\mathfrak q$. Ahora bien, si
$\overline f\in k[[T_1,\ldots,T_n]]$ es la imagen de $f$ («serie reducida» de $f$),
se sabe (Bourbaki, \emph{Alg. comm.}, cap.~VII, \S3, n\textsuperscript{o}~7,
lema~3) que existe un $A$-automorfismo $\sigma$ de $B$ tal que
$(\overline{\sigma(f)})(0,\ldots,0,T_n)\neq0$, lo que implica evidentemente
$\sigma(f)\notin\mathfrak p'$.

Supongamos, pues, que $n=1$ y escribamos $T$ en lugar de $T_1$, de modo que $B=A[[T]]$.
Basta probar que, en el anillo de polinomios $A[T]$, el ideal
$\mathfrak q\cap A[T]$ es principal;
\oldpage[IV-4]{325}
en efecto, sea $f_0$ un elemento de $\mathfrak q$ que no pertenece a $\mathfrak mB$;
es una serie formal cuya serie reducida $\overline f_0$ no es nula; por tanto,
en virtud del teorema de preparación (Bourbaki, \emph{Alg. comm.}, cap.~VII, \S3,
n\textsuperscript{o}~8, prop.~5), para todo $f\in\mathfrak q$ existen $g\in B$ y
un polinomio $r\in A[T]$ tales que $f=gf_0+r$, y se tiene entonces
$r\in\mathfrak q\cap A[T]$; por otra parte (\emph{loc. cit.}, prop.~6), existen un
polinomio distinguido no constante $F_0\in A[T]$ y un elemento invertible $u\in B$ tales
que $f_0=uF_0$, por lo que también se tiene $F_0\in\mathfrak q\cap A[T]$, lo que prueba que
$\mathfrak q$ está generado por $\mathfrak q\cap A[T]=\mathfrak q_1$. Puesto que $B$
es plano sobre $A[T]$ $(\mathbf 0_{\mathrm I}, 7.3.3)$, resulta de Bourbaki,
\emph{Alg. comm.}, cap.~VII, \S1, n\textsuperscript{o}~10, prop.~15, que
$\mathfrak q_1$ es un ideal primo de altura $1$ en $A[T]$.

Además, se tiene necesariamente $\mathfrak q_1\cap A=0$; de lo contrario,
$\mathfrak q_1\cap A$ tendría necesariamente altura $\geq1$, y resultaría de
(5.5.3) que $\mathfrak q_1=(\mathfrak q_1\cap A)A[T]$. Pero entonces,
puesto que $\mathfrak q_1\cap A\subset\mathfrak m$, se tendría
$\mathfrak q_1\subset\mathfrak mA[T]$, en contradicción con la hipótesis sobre
$\mathfrak q$. Si $K$ es el cuerpo de fracciones de $A$, $\mathfrak q_1K[T]$ es
entonces un ideal primo distinto de $0$ y de $K[T]$ en $K[T]$, y por tanto de la forma
$hK[T]$, donde
\[
  h(T)=T^m+a_1T^{m-1}+\cdots+a_m,
  \qquad m\geq1,
\]
con $a_i\in K$, siendo $h$ irreducible en $K[T]$. Pero se ha visto más arriba que
$\mathfrak q_1$ contiene un polinomio distinguido no constante $F_0\in A[T]$. Si $t$
es la clase de $T$ en $K[T]/\mathfrak q_1K[T]$, $t$ es, pues, raíz del polinomio
$F_0$ en una extensión de $K$ y, por consiguiente, $h$ divide a $F_0$ en $K[T]$; pero,
puesto que $h$ y $F_0$ son mónicos, ello implica que los coeficientes $a_i$ de $h$
son enteros sobre $A$ (Bourbaki, \emph{Alg. comm.}, cap.~V, \S1,
n\textsuperscript{o}~3, prop.~11), por lo que pertenecen a $A$, ya que $A$ es
integralmente cerrado. En otras palabras, se tiene
\[
  h\in A[T]\cap\mathfrak q_1K[T]=\mathfrak q_1
\]
(Bourbaki, \emph{Alg. comm.}, cap.~II, \S1, n\textsuperscript{o}~5, prop.~11);
como todo polinomio $g\in\mathfrak q_1$ es divisible por $h$ en $K[T]$ y $h$
es mónico, los coeficientes de $g/h$ pertenecen a $A$, por lo que
$\mathfrak q_1=hA[T]$. C.Q.D.
\end{proof}

El enunciado (21.14.1) es \emph{equivalente} al siguiente:

\begin{corollary}[21.14.2]
\label{IV.21.14.2}
Bajo las hipótesis de (21.14.1) sobre $A$, $B$ y $\rho$, para todo ideal primo
$\mathfrak q$ de $B$ no contenido en $\mathfrak p=\mathfrak mB$ y tal que
$\dim(B_\mathfrak q)\geq2$, el anillo $B_\mathfrak q$ es parafactorial. En
particular, si $\dim(B\otimes_A k)>0$ (\emph{i.e.} $\dim B>\dim A$
$((\mathbf 0, 19.7.1)$ y $(\mathbf 0, 6.1.1)))$, el anillo $B$ es parafactorial.
\end{corollary}

\begin{proof}
En efecto, se ha visto en la demostración de (21.14.1) que las condiciones del enunciado
implican que $B$ es íntegro e integralmente cerrado; la equivalencia de los enunciados
(21.14.1) y (21.14.2) resulta entonces de (21.13.15), aplicado a
$X=\operatorname{Spec}(B)$ y $S=\{\mathfrak p\}$.
\end{proof}

\begin{proposition}[21.14.3]
\label{IV.21.14.3}
Sean $S$ un preesquema normal y $f:X\to S$ un morfismo liso.
\begin{enumerate}[label=(\roman*)]
  \item Si $S$ es localmente noetheriano (y por tanto también $X$), entonces, para todo $x\in X$
  tal que $\dim(\mathcal O_{X,x})\geq2$ y que $x$ no sea punto maximal de su
  fibra $f^{-1}(f(x))$, el anillo $\mathcal O_{X,x}$ es parafactorial. Todo ciclo
  $1$-codimensional $Z$ sobre $X$ tal que $\operatorname{Supp}(Z)$ no contenga ninguna
  componente irreducible de una fibra $f^{-1}(s)=X_s$ es localmente principal.
  \item Sea $Y$ un subconjunto cerrado de $X$ que no contenga ninguna componente irreducible
  de una fibra $X_s$, y tal que, para todo punto maximal $\eta$ de $S$,
  $Y_\eta=Y\cap X_\eta$ sea de codimensión $\geq2$ en $X_\eta$. Entonces el par
  $(X,Y)$ es parafactorial.
\end{enumerate}
\end{proposition}

\begin{proof}
Obsérvese que las condiciones de (ii) se cumplen si $Y$ es un subconjunto cerrado de $X$
tal que, para todo $s\in S$, $Y_s=Y\cap X_s$ sea de codimensión $\geq2$ en
$X_s$.

Para probar (21.14.3), observemos primero que, bajo las hipótesis de (i), si se pone
$Y=\overline{\{x\}}$, existe un entorno abierto $V$ de $x$ en $X$ tal que
$Y\cap V$ y $V$ satisfacen las condiciones de (ii): en efecto, resulta de la hipótesis
y de (9.5.3) que se puede tomar $V$ de modo que $V\cap Y$ no contenga ninguna
componente irreducible de un $X_s$; por otra parte, para todo punto maximal $\eta$ de
$S$ tal que $Y\cap X_\eta\neq\varnothing$, los puntos de $Y\cap X_\eta$ son
especializaciones de $x$, por lo que en un punto tal $z$ se tiene
$\dim(\mathcal O_{X,z})\geq2$ y, como
$\mathcal O_{X,z}=\mathcal O_{X_\eta,z}$ ($S$ es reducido en el punto $\eta$ en virtud
de (2.1.13)), $Y\cap X_\eta$ es efectivamente de codimensión $\geq2$ en $X_\eta$. Como el
par $(V,Y\cap V)$ es entonces parafactorial en virtud de (ii), se concluye de (21.13.10)
que el anillo $\mathcal O_{X,x}$ es parafactorial, es decir, la primera afirmación
de (i). Para probar la segunda, se puede reducir al caso en que
$Z=\overline{\{z\}}$, donde $z\in X^{(1)}$; como $\mathcal O_{X,z}$ es un anillo local
integralmente cerrado y de dimensión $1$, es un anillo de valoración discreta y $Z$
es, por tanto, principal en el punto $z$; para todo otro punto $x$ de
$\operatorname{Supp}(Z)$, se tiene entonces $\dim(\mathcal O_{X,x})\geq2$ y, por hipótesis, $x$ no es
punto maximal de su fibra, por lo que $\mathcal O_{X,x}$ es parafactorial.
Apliquemos entonces (21.13.15), tomando como $S$ el conjunto formado por $z$ y los puntos
maximales de las fibras de $f$; es claro que $Z$ es principal en los puntos de $S$
puesto que $z$ es el único punto de $S$ que pertenece a $\operatorname{Supp}(Z)$; como,
por otra parte, la condición $b)$ de (21.13.15) se cumple evidentemente, se concluye que
$Z$ es localmente principal.

Quedamos, pues, reducidos a probar (ii). Pongamos $U=X-Y$ y sea $j:U\to X$ la inyección
canónica.
\oldpage[IV-4]{326}
Como las hipótesis y la conclusión son locales sobre $S$ y sobre $X$ en virtud de
(21.13.6, (i)), se puede reducir al caso en que $S$ y $X$ son afines, siendo entonces $f$ de
presentación finita. Puesto que las fibras de $f$ son regulares, la hipótesis sobre $Y$
implica que, para todo punto $x\in Y$, se tiene
\[
  \operatorname{prof}(\mathcal O_{X_{f(x)},x})
  =\dim(\mathcal O_{X_{f(x)},x})\geq1;
\]
resulta de (19.9.8) que el homomorfismo canónico
$\mathcal O_X\to j_*(\mathcal O_U)$ es \emph{inyectivo}. Además, sustituyendo $Y$ por
un subconjunto cerrado mayor según el método de la demostración de (19.9.8), y
utilizando (21.13.6, (ii)), se puede suponer que $Y$ está definido por un ideal \emph{de
tipo finito} del anillo de $X$; por tanto, el abierto $U=X-Y$ es \emph{cuasicompacto y
cuasiseparado} y el conjunto cerrado $Y$ es \emph{constructible}.

Para probar (ii), basta, dados un $\mathcal O_U$-módulo invertible $\mathscr L$
y un punto $x\in Y$, establecer la existencia de un entorno abierto $V$ de $x$
en $X$ tal que: $1^{\mathrm o}$ el homomorfismo canónico
$\mathcal O_V\to(j_V)_*(\mathcal O_{U\cap V})$ sea sobreyectivo;
$2^{\mathrm o}$ exista un $\mathcal O_V$-módulo invertible $\mathscr L'_V$ tal que
$\mathscr L'_V|(U\cap V)=\mathscr L|(U\cap V)$ ((21.13.5) y (21.13.3)).

Pongamos $s=f(x)$; veremos que se puede reducir al caso en que
$S=S_0=\operatorname{Spec}(\mathcal O_{S,s})$. Sea, en efecto, $(S_\nu)$ un sistema
fundamental de entornos abiertos afines de $s$ en $S$, y pongamos
$X_\nu=f^{-1}(S_\nu)$, $Y_\nu=Y\cap X_\nu$, $U_\nu=X_\nu-Y_\nu$,
siendo $j_\nu:U_\nu\to X_\nu$ la inyección canónica; $X_0=X\times_S S_0$ es entonces
límite proyectivo de los $X_\nu$ (8.1.2, $a$)); supongamos verdadera la proposición para el
morfismo $f_0:X_0\to S_0$ y para $Y_0=Y\cap X_0$, y pongamos $U_0=X_0-Y_0$,
siendo $j_0:U_0\to X_0$ la inyección canónica. Como la proyección
$p_\nu:X_0\to X_\nu$ es un morfismo afín, se tiene
\[
  (j_0)_*(\mathcal O_{U_0})
  =p_\nu^*((j_\nu)_*(\mathcal O_{U_\nu}))
  \qquad(\mathbf{II}, 1.5.2)
\]
y el homomorfismo canónico
$\rho_0:\mathcal O_{X_0}\to(j_0)_*(\mathcal O_{U_0})$ no es otro que
$p_\nu^*(\rho_\nu)$, donde
$\rho_\nu:\mathcal O_{X_\nu}\to(j_\nu)_*(\mathcal O_{U_\nu})$ es el homomorfismo
canónico; como, por hipótesis, $\rho_0$ es sobreyectivo, lo mismo ocurre con $\rho_\nu$
para $\nu$ suficientemente grande (8.5.7). Sea, por otra parte, $\mathscr L_0$ el
$\mathcal O_{U_0}$-módulo invertible obtenido por restricción de $\mathscr L$, y observemos que los
$X_\nu$ son afines, y por tanto cuasicompactos y cuasiseparados; puesto que, por hipótesis,
existe un $\mathcal O_{X_0}$-módulo invertible $\mathscr L'_0$ tal que
$\mathscr L'_0|U_0=\mathscr L_0$, existen un $\nu$ suficientemente grande y un
$\mathcal O_{X_\nu}$-módulo cuasicoherente $\mathscr L'_\nu$ de presentación finita
tales que $\mathscr L'_0$ sea isomorfo a $p_\nu^*(\mathscr L'_\nu)$ (8.5.2, (ii));
además, por (8.5.5), se puede suponer que $\mathscr L'_\nu$ es invertible. Finalmente, como
los $U_\nu$ son cuasicompactos y cuasiseparados y
$\mathscr L'_0|U_0=\mathscr L_0$, existe un $\nu$ suficientemente grande tal que
$\mathscr L'_\nu|U_\nu$ y $\mathscr L|U_\nu$ sean isomorfos (8.5.2.5).

\oldpage[IV-4]{327}
Quedamos así reducidos al caso en que $S=\operatorname{Spec}(A)$,
$X=\operatorname{Spec}(B)$, donde $A$ es un anillo \emph{local}; puesto que $X$ es
normal y $f$ liso, $S$ es normal (17.5.7), y por tanto $A$ es \emph{íntegro e
integralmente cerrado}. Consideremos entonces $A$ como límite inductivo de sus
\emph{sub-$\mathbb Z$-álgebras de tipo finito}; como la clausura integral de una
sub-$\mathbb Z$-álgebra de este tipo es un subanillo de $A$ y es también una
$\mathbb Z$-álgebra de tipo finito (7.8.3, (ii), (iii) y (vi)), $A$ es límite inductivo
filtrante de las sub-$\mathbb Z$-álgebras de tipo finito $A_\lambda$ de $A$ que son
anillos \emph{integralmente cerrados}. En virtud de (1.8.4.2), existen un índice
$\lambda$ y una $A_\lambda$-álgebra de tipo finito $B_\lambda$ tales que
\[
  B=B_\lambda\otimes_{A_\lambda}A
\]
salvo $A$-isomorfismo. Pongamos $S_\lambda=\operatorname{Spec}(A_\lambda)$,
$X_\lambda=\operatorname{Spec}(B_\lambda)$; si $p_\lambda:X\to X_\lambda$ es la
proyección canónica, se puede además suponer (puesto que $Y$ es
\emph{constructible}) que $Y=p_\lambda^{-1}(Y_\lambda)$, donde $Y_\lambda$ es un
subconjunto cerrado de $X_\lambda$ (8.3.11). Sea $f_\lambda$ el morfismo
$X_\lambda\to S_\lambda$; en virtud de la transitividad de las fibras y de (4.2.6),
$Y_\lambda$ no contiene ninguna componente irreducible de las fibras
$f_\lambda^{-1}(s_\lambda)$ para ningún $s_\lambda\in S_\lambda$. Puesto que $f$ es
liso, se puede también suponer que $f_\lambda$ es liso (17.7.8); finalmente, la imagen del
punto genérico $\eta$ de $S$ en $S_\lambda$ es el punto genérico $\eta_\lambda$
de $S_\lambda$, y se tiene
\[
  \operatorname{codim}((Y_\lambda)_{\eta_\lambda},
  (X_\lambda)_{\eta_\lambda})
  =\operatorname{codim}(Y_\eta,X_\eta)\geq2
\]
en virtud de (6.1.4). Se ve, pues, que $S_\lambda$, $X_\lambda$ y $Y_\lambda$
satisfacen todas las hipótesis de (ii), y, poniendo
$X_\mu=X_\lambda\times_{S_\lambda}S_\mu$ para $\lambda\leq\mu$, lo mismo ocurre con
$S_\mu$, $X_\mu$ y $Y_\mu$. Probemos que, si se demuestra que el par
$(X_\mu,Y_\mu)$ es parafactorial para todo $\mu\gg\lambda$, lo mismo ocurre con
$(X,Y)$. En efecto, sean $U=X-Y$, $U_\mu=X_\mu-Y_\mu$, y sean $j:U\to X$ y
$j_\mu:U_\mu\to X_\mu$ las inyecciones canónicas; como la proyección
$p_\mu:X\to X_\mu$ es un morfismo afín, se tiene
\[
  j_*(\mathcal O_U)=p_\mu^*((j_\mu)_*(\mathcal O_{U_\mu}))
  \qquad(\mathbf{II}, 1.5.2),
\]
y por consiguiente el homomorfismo canónico
$\rho:\mathcal O_X\to j_*(\mathcal O_U)$ no es otro que $p_\mu^*(\rho_\mu)$, donde
$\rho_\mu:\mathcal O_{X_\mu}\to(j_\mu)_*(\mathcal O_{U_\mu})$ es el homomorfismo
canónico; como este último es biyectivo por hipótesis, lo mismo ocurre con $\rho$.
Por otra parte, para todo $\mathcal O_U$-módulo invertible $\mathscr L$, existen un
$\mu\gg\lambda$ y un $\mathcal O_{U_\mu}$-módulo invertible $\mathscr L_\mu$ tales
que $\mathscr L=p_\mu'^*(\mathscr L_\mu)$ (denotando por
$p_\mu':U\to U_\mu$ la restricción de $p_\mu$) (8.5.2 y 8.5.5), pues $U_\lambda$ es
cuasicompacto ya que $X_\lambda$ es noetheriano. Por hipótesis, existe un
$\mathcal O_{X_\mu}$-módulo invertible $\mathscr L'_\mu$ tal que
$\mathscr L'_\mu|U_\mu$ sea isomorfo a $\mathscr L_\mu$;
$\mathscr L'=p_\mu^*(\mathscr L'_\mu)$ es entonces un $\mathcal O_X$-módulo
invertible tal que $\mathscr L'|U$ sea isomorfo a $\mathscr L$, lo que prueba nuestra
afirmación.

Quedamos así reducidos a probar (ii) cuando el anillo $A$ es una
$\mathbb Z$-álgebra de tipo finito integralmente cerrada; como los anillos locales
$\mathcal O_{S,s}$ de $S$ son entonces anillos excelentes integralmente cerrados, sus
\emph{completados} son también integralmente cerrados (7.8.3, (ii), (iii) y (vii)). Para
probar que el par $(X,Y)$ es parafactorial, basta mostrar que, para todo
$x\in Y$, el anillo $\mathcal O_{X,x}$ es parafactorial (21.13.10). Sea $y$ un punto
cerrado de $Y_{f(x)}$ que sea especialización de $x$ (5.1.11); si se pone
$s=f(x)=f(y)$, $\mathcal O_{S,s}$ tiene un
\oldpage[IV-4]{328}
completado integralmente cerrado, y $\mathcal O_{X,y}$ es una
$\mathcal O_{S,s}$-álgebra formalmente lisa para las topologías preádicas
(17.5.3); puesto que $Y$ no contiene ninguna componente irreducible de $X_s$, se tiene
$\dim(\mathcal O_{X,y})>\dim(\mathcal O_{S,s})$ (17.5.8). Si $s\neq\eta$, se tiene
$\dim(\mathcal O_{S,s})\geq1$; si, por el contrario, $s=\eta$, se tiene por hipótesis
$\dim(\mathcal O_{X,x})\geq2$; por tanto, en todos los casos,
$\dim(\mathcal O_{X,y})\geq2$. Además, el ideal primo de $\mathcal O_{X,y}$
correspondiente al punto $x$ no está contenido en
$\mathfrak m_s\mathcal O_{X,y}$, puesto que $Y$ no contiene ninguna componente
irreducible de $X_s$. Finalmente, puesto que $y$ es cerrado en $X_s$, $k(y)$ es una extensión
finita de $k(s)$ $(\mathbf I, 6.4.2)$. En todos los casos, se puede aplicar a
$A=\mathcal O_{S,s}$, $B=\mathcal O_{X,y}$ y
$B_\mathfrak q=\mathcal O_{X,x}$ el resultado de (21.14.2), lo que termina la
demostración.
\end{proof}

\begin{env}[21.14.4]
\label{IV.21.14.4}
\emph{Observaciones}. ---

(i) Es posible que los enunciados (21.14.1) y (21.14.2) sigan siendo válidos \emph{sin
hipótesis} sobre el cuerpo residual de $B$. Formulado con esta generalidad, el resultado
sería equivalente, en virtud de (21.13.15) y (21.13.12.1), al siguiente: sean $A$ un
anillo local noetheriano completo, íntegro e integralmente cerrado, y $B$ un anillo local
noetheriano que sea una $A$-álgebra formalmente lisa, tal que
$\dim(B)>\dim(A)$; entonces $B$ es parafactorial.

(ii) Veremos en el cap.~VI, empleando una técnica de «descenso finito»
aplicada al morfismo $Y'\to Y=\operatorname{Spec}(A)$, donde $Y'$ es la normalización de
$Y$, que la conclusión de (21.14.1) (o de (21.14.2)) sigue siendo válida si se sustituye
la hipótesis de que $\widehat A$ es integralmente cerrado por la hipótesis de que
$\widehat A$ es \emph{reducido}, siempre que $\dim B\geq\dim A+2$. Si se sustituye esta
última condición por $\dim B\geq\dim A+3$, se puede incluso suprimir la hipótesis de que
$\widehat A$ es reducido. Del mismo modo, la conclusión de (21.14.3) sigue siendo válida si se
sustituye la hipótesis de que $X$ es normal por la hipótesis de que $X$ es reducido, siempre
que $\dim(\mathcal O_{X,x})\geq\dim(\mathcal O_{S,f(x)})+2$.

(iii) En el cap.~III, 3\textsuperscript{a} parte, se prueba que, si $f:X\to S$ es un
morfismo liso y $Y$ es un subconjunto cerrado localmente constructible de $X$ tal que,
para todo $s\in S$, se tenga (con las notaciones de (21.14.3))
$\operatorname{codim}(Y_s,X_s)\geq3$, entonces el par $(X,Y)$ es parafactorial
([41], XII, 4.8). Esta conclusión deja de ser válida cuando solo se supone que
$\operatorname{codim}(Y_s,X_s)\geq2$ para todo $s\in S$ y ya no se supone que
$X$ sea reducido. Por ejemplo, sean $k$ un cuerpo y $A=k[T]/(T^2)$, el álgebra de los números
duales sobre $k$, $S=\operatorname{Spec}(A)$,
$X=\operatorname{Spec}(A[T_1,T_2])$ ($T_1$, $T_2$ indeterminadas), de modo que
$X=\mathbb V_S^2$, siendo $Y$ la «sección nula» de este fibrado; si $s$ es el único
punto cerrado de $S$, se tiene $X_s=\mathbb V_k^2$ y $Y_s$ es la «sección nula» de
$X_s$. Para ver que el par $(X,Y)$ no es parafactorial, basta probar
que el anillo $B=A[T_1,T_2]_\mathfrak m$, donde $\mathfrak m$ es el ideal
$(T_1)+(T_2)$ que define $Y$, no es parafactorial (21.13.10). Sea
$B_0=B_{\mathrm{red}}$, y denotemos por $U$ y $U_0$ los complementos del punto cerrado
en $\operatorname{Spec}(B)$ y $\operatorname{Spec}(B_0)$; razonando como
en (21.13.9), se tiene la sucesión exacta, extensión de (21.13.9.1),
\[
  \Gamma(U,\mathcal O_U)^*\longrightarrow\Gamma(U_0,\mathcal O_{U_0})^*
  \longrightarrow\mathrm H^1(U_0,\mathcal O_{U_0})
  \longrightarrow\operatorname{Pic}(U)\longrightarrow\operatorname{Pic}(U_0).
\]
Ahora bien, se tiene $\Gamma(U,\mathcal O_U)=B$, $\Gamma(U_0,\mathcal O_{U_0})=B_0$ puesto que
$\operatorname{prof}(B)=\operatorname{prof}(B_0)=2$ (5.10.5). Por otra parte,
$\operatorname{Pic}(U_0)=0$ puesto que $B_0$ es factorial, y
$\mathrm H^1(U_0,\mathcal O_{U_0})\neq0$ [41, 3, Ejemplo III-1]; como
el homomorfismo $B^*\to B_0^*$ es sobreyectivo, se concluye que
$\operatorname{Pic}(U)\neq0$, y por tanto que $B$ no es parafactorial.
\oldpage[IV-4]{329}
(iv) El resultado (21.14.3) fue demostrado por primera vez por C.~Seshadri [44] en el caso
particular en que $S$ es un esquema algebraico normal sobre un cuerpo algebraicamente cerrado
$k$ y $X=S\times_k T$, donde $T$ es un $k$-preesquema algebraico liso sobre $k$. La
demostración de Seshadri [44, p.~188--189] es de naturaleza global y utiliza la
teoría de los esquemas de Picard. Proporciona además (\emph{loc. cit.}) resultados
como el siguiente (para el que actualmente no se dispone de una demostración
por métodos locales). Sean $S$ y $T$ dos preesquemas localmente de tipo finito sobre un
cuerpo $k$, $X=S\times_k T$, y $Z$ un ciclo $1$-codimensional sobre $X$ (considerado como
$S$-preesquema); supongamos que se cumplen las condiciones siguientes:
\begin{enumerate}[label=\arabic*\textsuperscript{o}]
  \item $S$ y $T$ son geométricamente normales sobre $k$ (6.7.6);
  \item Para todo punto maximal $\eta$ de $S$, el ciclo $1$-codimensional $Z_\eta$
  sobre la fibra $X_\eta$, que tiene la misma multiplicidad que $Z$ en todo punto de
  $X^{(1)}\cap X_\eta$, es localmente principal (en otras palabras, es la imagen de un
  divisor de $X_\eta$, puesto que $X_\eta$ es normal);
  \item Para todo $s\in S$, $Z$ es principal en los puntos maximales de la fibra $X_s$.
\end{enumerate}
Entonces $Z$ es localmente principal. En otras palabras, como $X$ es normal (6.14.1),
para todo $x\in X$ que no sea maximal en su fibra $X_s$ y que no pertenezca a
ninguna de las «fibras genéricas» $X_\eta$ (lo que implica
$\dim(\mathcal O_{X,x})\geq2$ en virtud de (6.1.1)), el anillo local
$\mathcal O_{X,x}$ es parafactorial, en virtud de (21.13.15)\footnote{Corrección de referencia: el francés impreso y el inglés mantenido dan (21.12.15), que trata de un criterio de isomorfía para un morfismo propio; el criterio de parafactorialidad pertinente es (21.13.15).}.
\footnote{En realidad, en el artículo citado, Seshadri supone que $k$ es algebraicamente
cerrado, que $T$ es separado y «semicompleto» (i.e.\ tal que
$\Gamma(T,\mathcal O_T)$ sea $k$-isomorfo a $k$), y sustituye la hipótesis
$3^{\mathrm o}$ por la hipótesis más fuerte de que $\operatorname{Supp}(Z)$ no contiene
ninguna de las fibras $X_s$ para $s\in S$. Pero, como el enunciado es local sobre $S$, se
concluye inmediatamente que basta formular la hipótesis $3^{\mathrm o}$, y esto prueba que
la conclusión (interpretada como antes en términos de la propiedad de parafactorialidad
de los anillos $\mathcal O_{X,x}$) es local sobre $S$ y sobre $T$, lo que permite
eliminar por completo la hipótesis de que $T$ es «semicompleto» y de que $k$ es
algebraicamente cerrado, puesto que (después de pasar, si es necesario, a la clausura algebraica de $k$)
se puede suponer inicialmente que $T$ es afín, lo que permite sumergirlo como abierto de un
esquema proyectivo y normal sobre $k$, al que se aplica el resultado de Seshadri. Obsérvese
también que, gracias a esta reducción, basta hacer la demostración de Seshadri en
el caso en que $T$ es \emph{proyectivo} (y no solo «semicompleto»), caso en el que la
teoría del esquema de Picard utilizada por Seshadri está contenida en la que se
desarrollará en el cap.~VI de nuestro Tratado.}
\end{env}

\subsection{Divisores relativos}
\label{subsection:IV.21.15}
\label{IV.21.15}

\begin{env}[21.15.1]
\label{IV.21.15.1}
Sean $S$ un preesquema y $f:X\to S$ un morfismo \emph{plano y localmente de
presentación finita}. Se ha definido en (20.6.1) el haz de anillos
$\mathscr M_{X/S}$ de los \emph{gérmenes de funciones meromorfas sobre $X$ relativas a
$S$}, subhaz de $\mathscr M_X$; es claro que la inyección canónica
$\mathcal O_X\to\mathscr M_X$ (20.1.4.1) aplica $\mathcal O_X$ en un
subhaz de $\mathscr M_{X/S}$, con el cual se lo identifica. Sea
$\mathscr M_{X/S}^*$ el haz (en grupos multiplicativos) de los gérmenes de secciones
invertibles de $\mathscr M_{X/S}$, que es, por tanto, un subhaz de $\mathscr M_X^*$
y contiene a $\mathcal O_X^*$ como subhaz.
\end{env}

\begin{definition}[21.15.2]
\label{IV.21.15.2}
Se llama \emph{haz de los divisores sobre $X$ relativos a $S$}, o
\emph{haz de los divisores sobre $X$ transversales a $f$}, y se denota por
$\mathscr{D}iv_{X/S}$ el haz cociente (de grupos conmutativos)
$\mathscr M_{X/S}^*/\mathcal O_X^*$; las secciones de este haz sobre $X$ se
denominan \emph{divisores sobre $X$ relativos a $S$}, o \emph{divisores sobre $X$
transversales a $f$}; forman un grupo conmutativo denotado por $\operatorname{Div}(X/S)$.
\end{definition}

Es claro que $\mathscr{D}iv_{X/S}$ se identifica canónicamente con un subhaz de
$\mathscr{D}iv_X$, y por tanto $\operatorname{Div}(X/S)$ con un subgrupo de
$\operatorname{Div}(X)$, que se denota también \emph{aditivamente}. Para
\oldpage[IV-4]{330}
todo abierto $U$ de $X$, se tiene
$\mathscr M_{X/S}^*|U=\mathscr M_{U/S}^*$, y por tanto
$\mathscr{D}iv_{X/S}|U=\mathscr{D}iv_{U/S}$, de modo que el haz
$\mathscr{D}iv_{X/S}$ es igual al prehaz
$U\mapsto\operatorname{Div}(U/S)$.

Puesto que $\mathscr M_{X/S}^*$ es un subhaz de $\mathscr M_X^*$, las
definiciones, notaciones y fórmulas relativas a los divisores de las secciones de
$\mathscr M_X^*$ sobre $X$ (21.1.3) se aplican sin cambio a las secciones
de $\mathscr M_{X/S}^*$ sobre $X$.

\begin{env}[21.15.3]
\label{IV.21.15.3}
La estructura de haz de grupos ordenados sobre $\mathscr{D}iv_X$ (21.1.6) induce
sobre el subhaz $\mathscr{D}iv_{X/S}$ una estructura de haz de grupos
ordenados, para la cual el haz de monoides de los gérmenes de secciones positivas es
$\mathscr{D}iv_X^+\cap\mathscr{D}iv_{X/S}$, que se denota por
$\mathscr{D}iv_{X/S}^+$. Se tiene
$\Gamma(X,\mathscr{D}iv_{X/S}^+)=\operatorname{Div}^+(X)\cap
\operatorname{Div}(X/S)$; se denota este submonoide de $\operatorname{Div}(X/S)$ por
$\operatorname{Div}^+(X/S)$, y está formado por los elementos $\geq0$ para una estructura
de grupo ordenado sobre $\operatorname{Div}(X/S)$. Resulta de (21.1.5.1) que
$\mathscr{D}iv_{X/S}^+$ es la imagen en $\mathscr{D}iv_{X/S}$ del subhaz de
monoides
\[
\label{IV.21.15.3.1}
  \mathcal O_{X/S}^{\mathrm{reg}}=\mathcal O_X\cap\mathscr M_{X/S}^*.
  \tag{21.15.3.1}
\]
\end{env}

Para todo abierto $U$ de $X$, $\Gamma(U,\mathcal O_{X/S}^{\mathrm{reg}})$ es
el conjunto de las secciones $t$ de $\mathcal O_X$ sobre $U$ tales que $t$ sea
regular y que $1/t$ pertenezca a $\Gamma(U,\mathscr M_{X/S})$, lo que significa,
con las notaciones de (20.6.1), que $t\in\mathscr T_{X/S}(U)$, de modo que el haz
$\mathcal O_{X/S}^{\mathrm{reg}}$ no es sino el haz denotado por $\mathscr T_{X/S}$
en (20.6.1). Se puede escribir, por tanto,
\[
\label{IV.21.15.3.2}
  \mathscr{D}iv_{X/S}^+=\mathcal O_{X/S}^{\mathrm{reg}}/\mathcal O_X^*
  \tag{21.15.3.2}
\]
salvo isomorfismo canónico.

\begin{env}[21.15.3.3]
\label{IV.21.15.3.3}
Sea $D\in\operatorname{Div}^+(X/S)$ y consideremos el subpreesquema cerrado $Y(D)$
de $X$ definido por el Ideal $\mathscr J_X(D)$ de $\mathcal O_X$ (21.6.5); en virtud de
lo anterior, para todo $x\in Y(D)$, existe un entorno abierto $U$ de $x$ en $X$
y una sección $t\in\mathscr T_{X/S}(U)$ tal que
$\mathscr J_X(D)|U=\mathcal O_U\cdot t$; puesto que $x\in Y(D)$, la imagen $t_{f(x)}$ de
$t$ en $\Gamma(U\cap X_{f(x)},\mathcal O_{X_{f(x)}})$ pertenece al ideal maximal
de $\mathcal O_{X_{f(x)},x}$ y además, por definición, para todo $s\in S$, la imagen
$t_s$ de $t$ en $\Gamma(U\cap X_s,\mathcal O_{X_s})$ es una sección
\emph{regular}. Se deduce, pues, de (11.3.8) y (19.2.4) que la inmersión canónica
$Y(D)\to X$ es \emph{transversalmente regular con respecto a $S$ y de
codimensión $1$} en el punto $x$. Siendo inmediata la recíproca, se ve que se pueden
\emph{identificar canónicamente} los divisores positivos sobre $X$ relativos a $S$ con los
subpreesquemas cerrados $Y$ de $X$ tales que la inyección canónica $Y\to X$ sea una
inmersión transversalmente regular con respecto a $S$ y de codimensión $1$. Haremos
habitualmente esta identificación.
\end{env}

\begin{proposition}[21.15.4]
\label{IV.21.15.4}
Sean $D$ un divisor sobre $X$ y $\mathscr J_X(D)$ el Ideal fraccionario invertible
correspondiente (21.2.5). Para que $D\in\operatorname{Div}(X/S)$, es necesario y suficiente
que, para todo $x\in X$, se tenga
$(\mathscr J_X(D)\cap\mathscr M_{X/S}^*)_x\ne0$ (o, equivalentemente, que
se tenga $\mathscr J_X(D)\subset\mathscr M_{X/S}$ y
$\mathscr J_X(D)^{-1}\subset\mathscr M_{X/S}$).
\end{proposition}

\begin{proof}
En efecto, decir que $D\in\operatorname{Div}(X/S)$ significa que, para todo $x\in X$,
existe un entorno abierto $U$ de $x$ y una sección
$f\in\Gamma(U,\mathscr M_{X/S}^*)$ tales que $D|U=\operatorname{div}_U(f)$; como
$\mathscr J_X(D)|U$ es el Ideal fraccionario invertible $\mathcal O_Uf$, la
proposición resulta inmediatamente.
\end{proof}

\oldpage[IV-4]{331}
\begin{proposition}[21.15.5]
\label{IV.21.15.5}
Sean $D$ un divisor sobre $X$, $\mathcal O_X(D)$ el Ideal fraccionario invertible y
$s_D$ la sección meromorfa regular de $\mathcal O_X(D)$ definidos canónicamente por
$D$ (21.2.8 y 21.2.9). Para que $D\in\operatorname{Div}(X/S)$, es necesario y suficiente
que $s_D\in\Gamma(X,\mathscr M_{X/S}(\mathcal O_X(D)))$.
\end{proposition}

\begin{proof}
En efecto, si $U$ es un abierto de $X$ tal que $D|U=\operatorname{div}_U(f)$, donde
$f\in\Gamma(U,\mathscr M_X^*)$, decir que
$s_D\in\Gamma(U,\mathscr M_{X/S}(\mathcal O_X(D)))$ significa, en virtud de las
definiciones (20.6.2), que $f^{-1}\in\Gamma(U,\mathscr M_{X/S})$, de donde resulta la
proposición.
\end{proof}

La interpretación de los divisores sobre $X$ mediante las clases
$\operatorname{cl}(\mathscr L,s)$ (21.2.11) permite, por tanto, interpretar los elementos de
$\operatorname{Div}(X/S)$ como los pares (salvo isomorfismo) $(\mathscr L,s)$
tales que $\mathscr L$ sea un $\mathcal O_X$-módulo invertible y que $s$ sea una
sección meromorfa de $\mathscr L$ sobre $X$, \emph{regular con respecto a
$S$} (20.6.5, (iii)).

\begin{proposition}[21.15.6]
\label{IV.21.15.6}
Sea $D$ un divisor sobre $X$ relativo a $S$, y supongamos que, para todo $x\in X$ tal
que $\operatorname{prof}(\mathcal O_{X_{f(x)},x})=1$, se tenga $D_x\geq0$
(resp. $D_x=0$). Entonces se tiene $D\geq0$ (resp. $D=0$).
\end{proposition}

\begin{proof}
Como en (21.1.8), podemos limitarnos al caso en que $D=\operatorname{div}(\varphi)$, donde
$\varphi$ es una función meromorfa regular relativa a $S$, y todo se reduce a ver
que, si $D_x\geq0$ en todo punto $x\in X$ tal que
$\operatorname{prof}(\mathcal O_{X_{f(x)},x})=1$, $\varphi$ está definida en todo
$X$. Pero esta hipótesis significa que, si $T=X-\operatorname{dom}(\varphi)$, se tiene
$\operatorname{prof}(\mathcal O_{X_{f(x)},x})\geq2$ para todo $x\in T$, y para
concluir basta aplicar (20.6.6).
\end{proof}

\begin{env}[21.15.7]
\label{IV.21.15.7}
Sea $X'$ un segundo $S$-preesquema plano y localmente de presentación finita sobre $S$, y
sea $g:X'\to X$ un $S$-morfismo. Si el $S$-morfismo $g$ es \emph{plano}, se sabe
(21.4.5) que la imagen inversa por $g$ de todo divisor sobre $X$ está definida; si
además $D\in\operatorname{Div}(X/S)$, resulta de la definición (21.15.2) y de
(20.6.8) que se tiene entonces $g^*(D)\in\operatorname{Div}(X'/S)$.
\end{env}

\begin{env}[21.15.8]
\label{IV.21.15.8}
Sea $X'$ un segundo $S$-preesquema plano y localmente de presentación finita sobre $S$, y
sea $g:X'\to X$ un $S$-morfismo \emph{finito y plano}. Obsérvese que $g$ es entonces
necesariamente de presentación finita (1.4.3, 1.4.6 y 1.6.3), por lo que
$g_*(\mathcal O_{X'})$ es un $\mathcal O_X$-módulo plano y de presentación finita, y
por tanto \emph{localmente libre} (2.1.12), es decir, $g$ es un morfismo
localmente libre (18.2.7); para todo $s\in S$, el morfismo correspondiente
$g_s:X'_s\to X_s$ es entonces también finito y localmente libre. Se deduce de (21.5.2)
y (20.6.1) que, para todo abierto $U$ de $X$ y toda sección
$t'\in\Gamma(g^{-1}(U),\mathscr T_{X'/S})$, la norma $N_{X'/X}(t')$ pertenece a
$\Gamma(U,\mathscr T_{X/S})$; el razonamiento de (21.5.3) demuestra luego que, para todo
$\mathcal O_{X'}$-módulo invertible $\mathscr L'$ y toda sección meromorfa $u'$ de
$\mathscr L'$ sobre $X'$, \emph{regular con respecto a $S$}, la norma
$N_{X'/X}(u')$ es una sección meromorfa de
$\mathscr L=N_{X'/X}(\mathscr L')$ sobre $X$, regular con respecto a $S$.
La interpretación de los divisores relativos a $S$ dada en (21.15.5) y la definición de
la imagen directa de un divisor (21.5.5) demuestran entonces que \emph{para todo divisor}
$D'\in\operatorname{Div}(X'/S)$, se tiene $g_*(D')\in\operatorname{Div}(X/S)$.
\end{env}

\begin{env}[21.15.9]
\label{IV.21.15.9}
Consideremos, por último, un morfismo cualquiera $S'\to S$, y (bajo las hipótesis de
(21.15.1)) pongamos $X'=X\times_S S'$, que es plano y localmente de presentación finita
sobre $S'$; si $p:X'\to X$ es la proyección canónica, se ha visto (20.6.9) que se tiene un
homomorfismo canónico
$p^*(\mathscr M_{X/S})\to\mathscr M_{X'/S'}$, que transforma evidentemente toda
sección de $\mathscr M_{X/S}$ sobre un abierto $U$, \emph{regular con respecto
a $S$}, en una sección de $\mathscr M_{X'/S'}$ sobre $p^{-1}(U)$, regular
con respecto a $S'$\footnote{La transcripción francesa dice aquí «relativa a $S$»; tras el cambio de base, la sección pertenece a $\mathscr M_{X'/S'}$ y es regular con respecto a $S'$, como confirma el testigo inglés mantenido.} (20.6.5, (iii)); se concluye entonces de la
\oldpage[IV-4]{332}
definición (21.15.2) y de la exactitud por la derecha del funtor $p^*$ que
el homomorfismo anterior define, pasando a los cocientes, un homomorfismo canónico
\[
\label{IV.21.15.9.1}
  \operatorname{Div}(X/S)\longrightarrow\operatorname{Div}(X'/S')
  \tag{21.15.9.1}
\]
que transforma evidentemente los elementos de $\operatorname{Div}^+(X/S)$ en elementos de
$\operatorname{Div}^+(X'/S')$. Se pone también
$\operatorname{Div}(X'/S')=\operatorname{Div}_{X/S}(S')$ (resp.
$\operatorname{Div}^+(X'/S')=\operatorname{Div}_{X/S}^+(S')$), y se ve inmediatamente
que se han definido así dos funtores contravariantes
\[
  \operatorname{Div}_{X/S}:\mathbf{Sch}_{/S}\longrightarrow\mathbf{Ab},\qquad
  \operatorname{Div}_{X/S}^+:\mathbf{Sch}_{/S}\longrightarrow\mathbf{Ens}
\]
de la categoría de los $S$-preesquemas en la de los grupos conmutativos (resp. de los
conjuntos). Se verán más adelante (cap.~VI) casos importantes en los que el funtor
$\operatorname{Div}_{X/S}^+$ es \emph{representable} $(\mathbf 0_{\mathrm{III}},
8.1.8)$.

Para todo divisor $D\in\operatorname{Div}(X/S)$, la imagen de $D$ por
el homomorfismo (21.15.9.1) no es otra que la \emph{imagen inversa}
$p^*(D)$ (en el sentido de (21.4.2)): la existencia de esta imagen inversa y la afirmación
anterior son, en efecto, consecuencias inmediatas de (20.6.5, (iii)) y (20.6.9).

Los elementos de $\operatorname{Div}(X'/S')$ se llaman a menudo, por abuso de lenguaje,
«\emph{familias de divisores sobre $X$ relativos a $S$, parametrizadas por el
$S$-preesquema $S'$}»; esta terminología se utiliza sobre todo cuando se trata de
divisores \emph{positivos}.
\end{env}


% Bibliografía de EGA IV, cuarta parte, p. 333; los títulos citados se conservan en su lengua bibliográfica.
\clearpage
\oldpage[IV]{333}
\phantomsection
\addcontentsline{toc}{section}{Bibliografía}
\section*{BIBLIOGRAFÍA}

\begingroup
\makeatletter
% The book-class bibliography environment normally creates its own chapter
% heading.  This leaf already carries the printed section heading above.
\expandafter\def\csname chapter\endcsname{\@ifstar{\@gobble}{\@gobble}}
\makeatother
\begin{thebibliography}{44}

\bibitem[41]{IV-41}
A.~Grothendieck, \emph{Séminaire de Géométrie algébrique, 1962,
Cohomologie locale des faisceaux cohérents, et théorèmes de Lefschetz locaux
et globaux} (Institut des Hautes Études scientifiques, Bures-sur-Yvette).

\bibitem[42]{IV-42}
M.~Demazure et A.~Grothendieck, \emph{Séminaire de Géométrie algébrique,
1963--1964, Schémas en groupes} (Institut des Hautes Études scientifiques,
Bures-sur-Yvette).

\bibitem[43]{IV-43}
M.~Artin et A.~Grothendieck, \emph{Séminaire de Géométrie algébrique
1963--1964, Cohomologie étale des schémas} (Institut des Hautes Études
scientifiques, Bures-sur-Yvette).

\bibitem[44]{IV-44}
C.~S.~Seshadri, \emph{Quotient space by an abelian variety},
\emph{Mathematische Annalen}, t.~152 (1962), p.~185--194.

\end{thebibliography}
\endgroup

% Índice de notaciones de EGA IV, cuarta parte, pp. 334--335.
\clearpage
\oldpage[IV]{334}
\phantomsection
\addcontentsline{toc}{section}{Índice de notaciones}
\section*{ÍNDICE DE NOTACIONES}
\markboth{ÍNDICE DE NOTACIONES}{ÍNDICE DE NOTACIONES}

\begingroup
\small
\setlength{\parindent}{0pt}
\setlength{\parskip}{1.5pt}

$Y_f^{(n)}$, $Y^{(n)}$, $\mathcal{G}r_{\bullet}(f)$,
$\mathcal{G}r_n(f)$, $\mathcal{N}_{Y/X}$
($f:Y\to X$ morfismo de espacios anillados) : 16.1.2.

$\mathcal{O}_{Y^{(\infty)}}$ : 16.1.11.

$\mathcal{P}^{n}_{f}$, $\mathcal{P}^{n}_{X/S}$,
$\mathcal{P}^{\infty}_{f}$, $\mathcal{P}^{\infty}_{X/S}$,
$\Omega^{1}_{f}$, $\Omega^{1}_{X/S}$
($f:X\to S$ morfismo de preesquemas) : 16.3.1.

$d^{n}_{f}$, $d^{n}_{X/S}$, $d^{\infty}_{f}$,
$d^{\infty}_{X/S}$, $d^{n}$, $d^{\infty}$, $dt$, $d_{X/S}(t)$
($f:X\to S$ morfismo de preesquemas) : 16.3.6.

$P^{n}(u)$, $\Gr_{\bullet}(u)$, $\Gr_{1}(u)$
($u:X'\to X$ morfismo de preesquemas) : 16.4.3.

$\operatorname{Der}_{S}(\mathcal{O}_{X},\mathcal{F})$ : 16.5.1.

$\mathcal{D}er_{S}(\mathcal{O}_{X},\mathcal{F})$ : 16.5.4.

$\mathcal{T}_{X/S}$ : 16.5.7.

$D_i$, $\partial/\partial s_i$ : 16.5.8.

$T_{X/S}$ : 16.5.12.

$T_{X/S}(x)$, $T_Y(s)$ : 16.5.13.

$\Omega^p_{X/S}$, $\Omega^\bullet_{X/S}$, $\Omega^p_{B/A}$ : 16.6.1.

$d$, $d_{X/S}$ : 16.6.3.

$\mathcal{P}^n_{X/S}(\mathcal{F})$ : 16.7.2.

$d^n_{X/S,\mathcal{F}}$, $d^n_{X/S}$ : 16.7.5.

$\mathcal{P}^{\infty}_{X/S}(\mathcal{F})$,
$d^{\infty}_{X/S,\mathcal{F}}$ : 16.7.7.

$\operatorname{Diff}^{n}_{X/S}(\mathcal{F},\mathcal{G})$,
$\operatorname{Diff}^{n}_{X/S}$,
$\mathcal{D}iff^{n}_{X/S}(\mathcal{F},\mathcal{G})$ : 16.8.3.

$\mathcal{D}iff^{n}_{X/S}$, $\langle t,D\rangle$ : 16.8.4.

$aD$ ($D$ operador diferencial, $a\in\Gamma(U,\mathcal{O}_X)$) : 16.8.5.

$\operatorname{Diff}_{X/S}(\mathcal{F},\mathcal{G})$,
$\mathcal{D}iff_{X/S}(\mathcal{F},\mathcal{G})$ : 16.8.7.

$\mathcal{D}iff_{X/S}$ : 16.8.10.

$|\mathbf p|$, $\mathbf p!$, $\binom{\mathbf p}{\mathbf q}$,
$\mathbf z^{\mathbf p}$, $\boldsymbol\zeta^{\mathbf p}$
($\mathbf p,\mathbf q$ multiíndices) : 16.11.1.

$\dim_x(f)$ ($f:X\to Y$ morfismo localmente de tipo finito, $x\in X$) : 17.10.1.

$\operatorname{Tr}_{B/A}$, $\operatorname{astr}_{B/A}$
($B$ $A$-álgebra finita y libre) : 18.2.1.

$\operatorname{Tr}_{\mathcal B/\mathcal O_X,U}$,
$\operatorname{astr}_{\mathcal B/\mathcal O_X,U}$,
$\operatorname{Tr}_{\mathcal B/\mathcal O_X}$,
$\operatorname{astr}_{\mathcal B/\mathcal O_X}$
($X$ espacio anillado, $\mathcal B$ $\mathcal O_X$-álgebra localmente libre de rango finito) : 18.2.2.

$\operatorname{Tr}_f$ ($f$ morfismo finito y étale) : 18.2.6.

$d_{X/Y}$, $\mathcal D_{X/Y}$ : 18.2.7.

${}^{t}V_{\mathcal E}$ : 18.5.1.

$Al(X,\mathcal E,\mathcal J)$, $Al(\mathcal E,\mathcal J)$ : 18.5.2.

$Of(S)$, $Of(X)$, $Id(\mathcal B')$ : 18.5.3.

$\operatorname{Hom.loc}(A,B)$ : 18.6.1.

${}^{h}A$ : 18.6.5 y 18.6.7.

$\operatorname{Hom.sloc}(A,B)$ : 18.6.7.

${}^{hs}A_{(i)}$, ${}^{hs}A$ : 18.8.7 y 18.8.9.

$\operatorname{codim}^{t}_{y}(Y,X)$ : 19.1.3.

$\operatorname{prof}_{T,t}(\mathcal F)$ : 19.9.1.

$\mathcal O_X[\mathcal S^{-1}]$, $\mathcal F[\mathcal S^{-1}]$ : 20.1.1.

$\mathcal S(\mathcal O_X)$, $\mathcal M_X$, $M(X)$,
$\mathcal M_X(\mathcal F)$, $M(X,\mathcal F)$ : 20.1.3.

$\operatorname{dom}(\varphi)$ (función meromorfa) : 20.1.4.

$\operatorname{dom}(u)$ ($u$ sección meromorfa de $\mathcal F$) : 20.1.7.

$(\mathcal M_X(\mathcal L))^{*}$
($\mathcal L$ $\mathcal O_X$-módulo invertible) : 20.1.8.

$s^{-1}$ ($s$ sección meromorfa de un $\mathcal O_X$-módulo invertible) : 20.1.10.

$\mathcal S_f(U)$, $\mathcal S_f$, $\mathcal M_f$
($f$ morfismo de espacios anillados), $\varphi\circ f$ ($\varphi$ función meromorfa),
$u\circ f$ ($u$ sección meromorfa de $\mathcal F$), $\mathcal M_f(\mathcal F)$ : 20.1.11.

\oldpage[IV]{335}
$\operatorname{Ps.hom}(X,Y)$, $\operatorname{Ps.hom}_S(X,Y)$ : 20.2.1.

$\omega|U$ ($\omega$ pseudomorfismo),
$\mathcal P\!s.\operatorname{hom}(X,Y)$,
$\mathcal P\!s.\operatorname{hom}_S(X,Y)$ : 20.2.2.

$\operatorname{dom}_S(\omega)$, $\operatorname{dom}(\omega)$
($\omega$ pseudomorfismo) : 20.2.3.

$\mathcal M'_X$, $M'(X)$ : 20.2.8.

$f\circ\omega$ ($f$ morfismo, $\omega$ pseudomorfismo) : 20.3.1.

$\omega\circ f$ ($f$ morfismo, $\omega$ pseudomorfismo) : 20.3.2.

$\operatorname{Ps.hom}_S(X,Y)^f$,
$\mathcal P\!s.\operatorname{hom}_S(X,Y)^f$ : 20.3.2.

$\Gamma_\omega$ ($\omega$ pseudomorfismo) : 20.4.1.

$p^{-1}$ ($p:\Gamma_\omega\to X$ restricción de la primera proyección) : 20.4.2.

$\omega(M)$ ($\omega$ pseudomorfismo, $M$ parte de $X$) : 20.4.2.

$\operatorname{Ps.hom}_{X/S}(X,Y)$ : 20.5.1.

$\mathcal P\!s.\operatorname{hom}_{X/S}(X,Y)$ : 20.5.3.

$\mathcal M'_{X/S}$, $M'(X/S)$ : 20.5.4.

$\mathcal S_{X/S}(U)$, $\mathcal S_{X/S}$, $\mathcal M_{X/S}$,
$M(X/S)$ : 20.6.1.

$\mathcal M_{X/S}(\mathcal F)$, $M(X/S,\mathcal F)$ : 20.6.2.

$\mathcal S_f(U)$, $\mathcal S_f$, $\mathcal M_{X/S,f}$ : 20.6.7.

$\mathcal{D}iv_X$, $\operatorname{Div}(X)$, $\operatorname{div}(f)$,
$\operatorname{div}_X(f)$ ($X$ espacio anillado) : 21.1.2.

$\operatorname{Supp}(D)$ ($D$ divisor) : 21.1.2.

$\operatorname{Div}(A)$ ($A$ anillo) : 21.2.1.

$\operatorname{div}(s)$ ($s$ sección de un módulo invertible) : 21.1.4.

$\mathcal{D}iv_X^+$, $\operatorname{Div}^+(X)$, $D\leq D'$
($D,D'$ divisores) : 21.1.6.

$\operatorname{Id.inv}(X)$, $\sh{J}d.inv_X$ : 21.2.4.

$\mathcal J_U$, $\mathcal J_X(D)$ : 21.2.5.

$\operatorname{div}(\mathcal I)$ ($\mathcal I$ ideal fraccionario) : 21.2.6.

$\mathcal O_X(D)$ ($D$ divisor) : 21.2.8.

$s_D$ ($D$ divisor) : 21.2.9.

$D(X)$, $\operatorname{cl}(\mathcal L,s)$ : 21.2.10.

$Y(D)$ ($D$ divisor positivo) : 21.2.12.

$\operatorname{Div.princ}(X)$ : 21.3.1.

$\operatorname{Pic}(X)$, $\operatorname{cl}(\mathcal L)$, $l$, $l_X$,
$\operatorname{Pic}(u)$ : 21.3.2.

$f^*(D)$ ($f$ morfismo de espacios anillados, $D$ divisor),
$\operatorname{Div}^f(X)$ : 21.4.2.

$\mathcal M_f^{**}$, $\mathcal{D}iv_X^f$ : 21.4.3.

$N_{X'/X}(u')$, $N(u')$
($u'$ sección meromorfa de un $\mathcal O_{X'}$-módulo invertible) : 21.5.3.

$f_*(D')$, $N_{X'/X}(D')$
($f$ morfismo finito, $D'$ divisor sobre $X'$) : 21.5.5.

$\mathfrak I(X)$, $\mathfrak Z(X)$, $\mathfrak Z^+(X)$
($X$ preesquema localmente noetheriano) : 21.6.1.

$\operatorname{mult}_x(Z)$, $\operatorname{Supp}(Z)$,
$\dim(Z)$, $\operatorname{codim}(Z)$ ($Z$ ciclo) : 21.6.1.

$X^{(p)}$, $\mathfrak Z^p(X)$, $\mathfrak Z^{p+}(X)$ : 21.6.2.

$Z|U$ ($Z$ ciclo), $\mathscr Z_X$, $\mathscr Z_X^p$,
$\mathscr Z_X^+$, $\mathscr Z_X^{p+}$ : 21.6.3.

$c$ : 21.6.4.

$\operatorname{cyc}(D)$ ($D$ divisor) : 21.6.5 y 21.6.7.

$\operatorname{mult}_x(D)$, $\operatorname{mult}_{\{x\}}(D)$
($D$ divisor), $o_x(f)$ ($f$ función meromorfa regular) : 21.6.7.

$\mathrm Z^+(f)$, $\mathrm Z^-(f)$, $\mathfrak Z^1_{\mathrm{princ}}(X)$,
$\mathfrak{Cl}(X)$ : 21.6.7.

$Y(C)$, $Y'(C)$, $\sh{J}_X(C)$, $\sh{J}(C)$
($C$ ciclo positivo de codimensión 1) : 21.7.1.

$f^*(Z)$ ($Z$ ciclo de codimensión 1, $f$ morfismo) : 21.10.3 y 21.10.11.

$\Gamma_T(\mathscr Z_X^1)$ : 21.10.3.

$f_*(Z')$ ($Z'$ ciclo de codimensión 1, $f$ morfismo finito) : 21.10.14.

$\Lambda^{\max}\mathcal E$ ($\mathcal E$ módulo localmente libre) : 21.11.1.

$U'$, $\operatorname{Aff}(U/X)$ : 21.12.1.

$\operatorname{Daf}(U/X)$ : 21.12.5.

$\mathcal{D}iv_{X/S}$, $\operatorname{Div}(X/S)$ : 21.15.2.

$\mathcal{D}iv_{X/S}^+$, $\operatorname{Div}^+(X/S)$ : 21.15.3.

$\mathcal O_{X/S}^{\mathrm{reg}}$ : 21.15.3.

$\operatorname{Div}_{X/S}(S')$, $\operatorname{Div}_{X/S}^+(S')$
($S'$ $S$-preesquema) : 21.15.9.

\endgroup

% Proyección española de la transcripción diplomática francesa de EGA IV, cuarta parte, pp. 336--339.
% Las entradas y sus remisiones reproducen el índice impreso.
\clearpage
\oldpage[IV]{336}
\phantomsection
\addcontentsline{toc}{section}{Índice terminológico}
\section*{ÍNDICE TERMINOLÓGICO}
\markboth{ÍNDICE TERMINOLÓGICO}{ÍNDICE TERMINOLÓGICO}

\begingroup
\small
\setlength{\parindent}{0pt}
\setlength{\parskip}{2pt}
\newcommand{\EGAIVterm}[2]{\textbf{#1}: #2.\par}
\newcommand{\EGAIVidxref}[1]{\hyperref[IV.#1]{#1}}

\EGAIVterm{Álgebra étale}{{\EGAReadLink{IV.17.3.2}{17.3.2}}}
\EGAIVterm{Álgebra esencialmente étale}{{\EGAReadLink{IV.18.6.1}{18.6.1}}}
\EGAIVterm{$\sh{O}_X$-álgebra finita y étale}{{\EGAReadLink{IV.18.2.3}{18.2.3}}}
\EGAIVterm{Álgebra lisa}{{\EGAReadLink{IV.17.3.2}{17.3.2}}}
\EGAIVterm{Álgebra no ramificada}{{\EGAReadLink{IV.17.3.2}{17.3.2}}}
\EGAIVterm{$K$-álgebra no ramificada sobre $Y$, sobre $A$}{{\EGAReadLink{IV.18.10.11}{18.10.11}}}
\EGAIVterm{Álgebra estrictamente esencialmente étale}{{\EGAReadLink{IV.18.6.2}{18.6.2}}}
\EGAIVterm{Anillo de intersección completa, anillo de intersección completa absoluta}{{\EGAReadLink{IV.19.3.1}{19.3.1}}}
\EGAIVterm{Anillo henseliano}{{\EGAReadLink{IV.18.5.8}{18.5.8}}}
\EGAIVterm{Anillo local parafactorial}{{\EGAReadLink{IV.21.13.7}{21.13.7}}}
\EGAIVterm{Anillo estrictamente local}{{\EGAReadLink{IV.18.8.2}{18.8.2}}}
\EGAIVterm{Aplicación lineal tangente}{{\EGAReadLink{IV.16.5.13}{16.5.13}}}
\EGAIVterm{Aplicación racional estricta}{{\EGAReadLink{IV.20.2.1}{20.2.1}}}
\EGAIVterm{Codimensión de un ciclo}{{\EGAReadLink{IV.21.6.1}{21.6.1}}}
\EGAIVterm{Codimensión transversal en un punto de un subpreesquema}{{\EGAReadLink{IV.19.1.3}{19.1.3}}}
\EGAIVterm{Codimensión transversal de un morfismo de inmersión en un punto}{{\EGAReadLink{IV.19.1.3}{19.1.3}}}
\EGAIVterm{Par de morfismos transversales en un punto}{{\EGAReadLink{IV.17.13.6}{17.13.6}}}
\EGAIVterm{Par henseliano}{{\EGAReadLink{IV.18.5.5}{18.5.5}}}
\EGAIVterm{Par parafactorial}{{\EGAReadLink{IV.21.13.1}{21.13.1}}}
\EGAIVterm{Ciclo}{{\EGAReadLink{IV.21.6.1}{21.6.1}}}
\EGAIVterm{Ciclo 1-codimensional con coeficientes racionales}{{\EGAReadLink{IV.21.10.9}{21.10.9}}}
\EGAIVterm{Ciclo $p$-codimensional}{{\EGAReadLink{IV.21.6.2}{21.6.2}}}
\EGAIVterm{Ciclo de los polos, ciclo de los ceros de una función meromorfa}{{\EGAReadLink{IV.21.6.7}{21.6.7}}}
\EGAIVterm{Ciclo linealmente equivalente a cero}{{\EGAReadLink{IV.21.6.7}{21.6.7}}}
\EGAIVterm{Ciclo localmente principal}{{\EGAReadLink{IV.21.6.7}{21.6.7}}}
\EGAIVterm{Ciclo polar de una función meromorfa}{{\EGAReadLink{IV.21.6.7}{21.6.7}}}
\EGAIVterm{Ciclo primo}{{\EGAReadLink{IV.21.6.1}{21.6.1}}}
\EGAIVterm{Ciclo principal}{{\EGAReadLink{IV.21.6.7}{21.6.7}}}
\EGAIVterm{Ciclo principal en un punto}{{\EGAReadLink{IV.21.6.7}{21.6.7}}}
\EGAIVterm{Ciclo puramente de codimensión $p$}{{\EGAReadLink{IV.21.6.2}{21.6.2}}}
\EGAIVterm{Defecto de afinidad de un abierto en un preesquema}{{\EGAReadLink{IV.21.12.5}{21.12.5}}}
\EGAIVterm{$S$-derivación, $(X/S)$-derivación, $f$-derivación}{{\EGAReadLink{IV.16.5.1}{16.5.1}}}
\EGAIVterm{Diferencial de una sección de $\sh{O}_X$}{{\EGAReadLink{IV.16.3.6}{16.3.6}}}
\EGAIVterm{Diferencial exterior}{{\EGAReadLink{IV.16.6.3}{16.6.3}}}
\EGAIVterm{Dimensión de un ciclo}{{\EGAReadLink{IV.21.6.1}{21.6.1}}}
\EGAIVterm{Dimensión relativa de un morfismo en un punto, de un $Y$-preesquema en un punto}{{\EGAReadLink{IV.17.10.1}{17.10.1}}}
\EGAIVterm{Discriminante de un $Y$-preesquema finito y localmente libre}{{\EGAReadLink{IV.18.2.7}{18.2.7}}}
\EGAIVterm{Divisor sobre un espacio anillado}{{\EGAReadLink{IV.21.1.2}{21.1.2}}}
\EGAIVterm{Divisor de una función meromorfa}{{\EGAReadLink{IV.21.1.2}{21.1.2}}}
\EGAIVterm{Divisor de una sección meromorfa de un módulo invertible}{{\EGAReadLink{IV.21.1.4}{21.1.4}}}
\EGAIVterm{Divisor efectivo}{{\EGAReadLink{IV.21.1.6}{21.1.6}}}
\EGAIVterm{Divisor principal}{{\EGAReadLink{IV.21.3.1}{21.3.1}}}
\EGAIVterm{Divisor relativo a $S$, divisor transversal a $f$}{{\EGAReadLink{IV.21.15.2}{21.15.2}}}
\EGAIVterm{Divisores linealmente equivalentes}{{\EGAReadLink{IV.21.3.1}{21.3.1}}}

\oldpage[IV]{337}
\EGAIVterm{Dominio de definición de una función meromorfa}{{\EGAReadLink{IV.20.1.4}{20.1.4}}}
\EGAIVterm{Dominio de definición de un pseudomorfismo}{{\EGAReadLink{IV.20.2.3}{20.2.3}}}
\EGAIVterm{Dominio de definición de una sección meromorfa de un módulo}{{\EGAReadLink{IV.20.1.7}{20.1.7}}}
\EGAIVterm{Envolvente afín de un $X$-preesquema}{{\EGAReadLink{IV.21.12.1}{21.12.1}}}
\EGAIVterm{Espacio tangente a un preesquema en un punto}{{\EGAReadLink{IV.16.5.13}{16.5.13}}}
\EGAIVterm{Haz conormal}{{\EGAReadLink{IV.16.1.2}{16.1.2}}}
\EGAIVterm{Haz de anillos de fracciones con denominadores en $\sh{S}$}{{\EGAReadLink{IV.20.1.1}{20.1.1}}}
\EGAIVterm{Haz de anillos graduados asociado a un morfismo}{{\EGAReadLink{IV.16.1.2}{16.1.2}}}
\EGAIVterm{Haz de módulos de fracciones con denominadores en $\sh{S}$}{{\EGAReadLink{IV.20.1.2}{20.1.2}}}
\EGAIVterm{Haz de $S$-derivaciones de $\sh{O}_X$}{{\EGAReadLink{IV.16.5.7}{16.5.7}}}
\EGAIVterm{Haz de $S$-derivaciones de $\sh{O}_X$ en $\sh{F}$}{{\EGAReadLink{IV.16.5.4}{16.5.4}}}
\EGAIVterm{Haz de $p$-diferenciales de $X$ relativo a $S$}{{\EGAReadLink{IV.16.6.1}{16.6.1}}}
\EGAIVterm{Haz de los divisores sobre un espacio anillado}{{\EGAReadLink{IV.21.1.2}{21.1.2}}}
\EGAIVterm{Haz de los divisores relativos a $S$, o transversales a $f$}{{\EGAReadLink{IV.21.15.2}{21.15.2}}}
\EGAIVterm{Haz de los gérmenes de ciclos 1-codimensionales con coeficientes en $\Lambda$}{{\EGAReadLink{IV.21.10.9}{21.10.9}}}
\EGAIVterm{Haz de los gérmenes de funciones meromorfas sobre $X$}{{\EGAReadLink{IV.20.1.4}{20.1.4}}}
\EGAIVterm{Haz de los gérmenes de funciones meromorfas relativas a $S$}{{\EGAReadLink{IV.20.6.1}{20.6.1}}}
\EGAIVterm{Haz de los gérmenes de secciones meromorfas de un módulo}{{\EGAReadLink{IV.20.1.3}{20.1.3}}}
\EGAIVterm{Haz de partes principales de orden $n$ del $S$-preesquema $X$}{{\EGAReadLink{IV.16.3.1}{16.3.1}}}
\EGAIVterm{Haz homogéneo formalmente principal (o pseudotorsor) bajo un haz de grupos}{{\EGAReadLink{IV.16.5.15}{16.5.15}}}
\EGAIVterm{Haz homogéneo principal (o torsor) bajo un haz de grupos}{{\EGAReadLink{IV.16.5.15}{16.5.15}}}
\EGAIVterm{Haz homogéneo principal (o torsor) trivial bajo un haz de grupos}{{\EGAReadLink{IV.16.5.15}{16.5.15}}}
\EGAIVterm{Haz tangente de $X$ relativo a $S$}{{\EGAReadLink{IV.16.5.7}{16.5.7}}}
\EGAIVterm{Familia de divisores sobre $X$ relativos a $S$, parametrizados por $S'$}{{\EGAReadLink{IV.21.15.9}{21.15.9}}}
\EGAIVterm{Familia de morfismos transversales en un punto}{{\EGAReadLink{IV.17.13.10}{17.13.10}}}
\EGAIVterm{Familia de subpreesquemas que se cortan transversalmente en un punto}{{\EGAReadLink{IV.17.13.11}{17.13.11}}}
\EGAIVterm{Fibrado tangente de $X$ relativo a $S$}{{\EGAReadLink{IV.16.5.12}{16.5.12}}}
\EGAIVterm{Función meromorfa sobre un espacio anillado}{{\EGAReadLink{IV.20.1.3}{20.1.3}}}
\EGAIVterm{Función meromorfa definida en un punto}{{\EGAReadLink{IV.20.1.4}{20.1.4}}}
\EGAIVterm{Función meromorfa regular}{{\EGAReadLink{IV.20.1.8}{20.1.8}}}
\EGAIVterm{Función meromorfa regular relativa a $S$}{{\EGAReadLink{IV.20.6.5}{20.6.5}}}
\EGAIVterm{Función meromorfa relativa a $S$}{{\EGAReadLink{IV.20.6.1}{20.6.1}}}
\EGAIVterm{Forma traza}{{\EGAReadLink{IV.18.2.1}{18.2.1}}}
\EGAIVterm{Grafo de un pseudomorfismo sobre $S$}{{\EGAReadLink{IV.20.4.1}{20.4.1}}}
\EGAIVterm{Grupo de las clases de ciclos 1-codimensionales}{{\EGAReadLink{IV.21.6.7}{21.6.7}}}
\EGAIVterm{Henselizado de un anillo local}{{\EGAReadLink{IV.18.6.5}{18.6.5}}}
\EGAIVterm{Henselizado de un anillo semilocal}{{\EGAReadLink{IV.18.6.7}{18.6.7}}}
\EGAIVterm{Henselizado estricto de un anillo local}{{\EGAReadLink{IV.18.8.7}{18.8.7}}}
\EGAIVterm{Henselizado estricto de un anillo semilocal}{{\EGAReadLink{IV.18.8.9}{18.8.9}}}
\EGAIVterm{Homomorfismo asociado a la traza}{{\EGAReadLink{IV.18.2.2}{18.2.2}}}
\EGAIVterm{Homomorfismo regular de $\sh{O}_X$-módulos}{{\EGAReadLink{IV.19.4.11}{19.4.11}}}
\EGAIVterm{Homomorfismo semilocal}{{\EGAReadLink{IV.18.6.7}{18.6.7}}}
\EGAIVterm{Homomorfismo traza}{{\EGAReadLink{IV.18.2.2}{18.2.2}}}
\EGAIVterm{Ideal de los denominadores de una sección meromorfa}{{\EGAReadLink{IV.20.2.14}{20.2.14}}}
\EGAIVterm{Ideal discriminante de un $Y$-preesquema finito y localmente libre}{{\EGAReadLink{IV.18.2.7}{18.2.7}}}
\EGAIVterm{Ideal entero}{{\EGAReadLink{IV.21.2.7}{21.2.7}}}
\EGAIVterm{Ideal fraccionario}{{\EGAReadLink{IV.21.2.1}{21.2.1}}}
\EGAIVterm{Ideal fraccionario invertible}{{\EGAReadLink{IV.21.2.1}{21.2.1}}}
\EGAIVterm{Ideal cuasirregular de un haz de anillos}{{\EGAReadLink{IV.16.9.1}{16.9.1}}}
\EGAIVterm{Ideal cuasirregular de un anillo}{{\EGAReadLink{IV.16.9.7}{16.9.7}}}
\EGAIVterm{Ideal regular de un haz de anillos}{{\EGAReadLink{IV.16.9.1}{16.9.1}}}
\EGAIVterm{Ideal regular de un anillo}{{\EGAReadLink{IV.16.9.7}{16.9.7}}}
\EGAIVterm{Ideal transversalmente regular en un punto de un preesquema}{{\EGAReadLink{IV.19.2.1}{19.2.1}}}

\oldpage[IV]{338}
\EGAIVterm{Ideal transversalmente regular de un álgebra}{{\EGAReadLink{IV.19.2.1}{19.2.1}}}
\EGAIVterm{Imagen directa de un ciclo 1-codimensional}{{\EGAReadLink{IV.21.10.14}{21.10.14}}}
\EGAIVterm{Imagen directa de un divisor}{{\EGAReadLink{IV.21.5.5}{21.5.5}}}
\EGAIVterm{Imagen inversa de un ciclo 1-codimensional}{{\EGAReadLink{IV.21.10.3}{21.10.3}} y {\EGAReadLink{IV.21.10.11}{21.10.11}}}
\EGAIVterm{Imagen inversa de un divisor}{{\EGAReadLink{IV.21.4.2}{21.4.2}}}
\EGAIVterm{Imagen inversa de una función meromorfa}{{\EGAReadLink{IV.20.1.11}{20.1.11}}}
\EGAIVterm{Imagen inversa de una función meromorfa relativa a $S$}{{\EGAReadLink{IV.20.6.7}{20.6.7}}}
\EGAIVterm{Inmersión cuasirregular}{{\EGAReadLink{IV.16.9.2}{16.9.2}}}
\EGAIVterm{Inmersión regular}{{\EGAReadLink{IV.16.9.2}{16.9.2}}}
\EGAIVterm{Inmersión regular en un punto}{{\EGAReadLink{IV.16.9.10}{16.9.10}}}
\EGAIVterm{Inmersión transversalmente regular en un punto relativamente a $S$, inmersión transversalmente regular relativamente a $S$}{{\EGAReadLink{IV.19.2.2}{19.2.2}}}
\EGAIVterm{Intersección completa en un punto relativamente a $S$, intersección completa relativamente a $S$ (para un preesquema plano y localmente de presentación finita)}{{\EGAReadLink{IV.19.3.6}{19.3.6}}}
\EGAIVterm{$n$-ésimo invariante normal de un morfismo}{{\EGAReadLink{IV.16.1.2}{16.1.2}}}
\EGAIVterm{Invariante normal de orden infinito de un morfismo}{{\EGAReadLink{IV.16.1.11}{16.1.11}}}
\EGAIVterm{$\sh{O}_X$-módulo de diferenciales de primer orden de $f$ (o de $X$ respecto de $S$, o del $S$-preesquema $X$)}{{\EGAReadLink{IV.16.3.1}{16.3.1}}}
\EGAIVterm{$\sh{O}_X$-módulo sin torsión relativamente a $S$}{{\EGAReadLink{IV.20.6.2}{20.6.2}}}
\EGAIVterm{$\sh{O}_X$-módulo sin torsión}{{\EGAReadLink{IV.20.1.5}{20.1.5}}}
\EGAIVterm{Morfismo diferencialmente liso}{{\EGAReadLink{IV.16.10.1}{16.10.1}}}
\EGAIVterm{Morfismo diferencialmente liso hasta el orden $r$}{{\EGAReadLink{IV.16.11.3}{16.11.3}}}
\EGAIVterm{Morfismo (plano) de intersección completa}{{\EGAReadLink{IV.19.3.6}{19.3.6}}}
\EGAIVterm{Morfismo esencialmente propio}{{\EGAReadLink{IV.18.10.20}{18.10.20}}}
\EGAIVterm{Morfismo étale}{{\EGAReadLink{IV.17.3.1}{17.3.1}}}
\EGAIVterm{Morfismo étale en un punto}{{\EGAReadLink{IV.17.3.7}{17.3.7}}}
\EGAIVterm{Morfismo finito localmente libre de rango $n$, morfismo finito localmente libre}{{\EGAReadLink{IV.18.2.7}{18.2.7}}}
\EGAIVterm{Morfismo formalmente étale, formalmente liso, formalmente neto, formalmente no ramificado}{{\EGAReadLink{IV.17.1.1}{17.1.1}}}
\EGAIVterm{Morfismo liso}{{\EGAReadLink{IV.17.3.1}{17.3.1}}}
\EGAIVterm{Morfismo liso en un punto}{{\EGAReadLink{IV.17.3.7}{17.3.7}}}
\EGAIVterm{Morfismo no ramificado, morfismo neto}{{\EGAReadLink{IV.17.3.1}{17.3.1}}}
\EGAIVterm{Morfismo no ramificado en un punto}{{\EGAReadLink{IV.17.3.7}{17.3.7}}}
\EGAIVterm{Morfismo transversal a $Y$ en el punto $x$ relativamente a $S$}{{\EGAReadLink{IV.17.13.3}{17.13.3}}}
\EGAIVterm{Morfismos equivalentes}{{\EGAReadLink{IV.20.2.1}{20.2.1}}}
\EGAIVterm{Multiplicidad de un ciclo en un punto}{{\EGAReadLink{IV.21.6.1}{21.6.1}}}
\EGAIVterm{Norma de un divisor}{{\EGAReadLink{IV.21.5.5}{21.5.5}}}
\EGAIVterm{Norma de una sección meromorfa}{{\EGAReadLink{IV.21.5.3}{21.5.3}}}
\EGAIVterm{Operador diferencial relativamente a $S$}{{\EGAReadLink{IV.16.8.1}{16.8.1}}}
\EGAIVterm{Orden de un operador diferencial}{{\EGAReadLink{IV.16.8.1}{16.8.1}}}
\EGAIVterm{Orden de una función meromorfa en un punto $x\in X^{(1)}$}{{\EGAReadLink{IV.21.6.7}{21.6.7}}}
\EGAIVterm{Parte principal de orden $n$, de orden infinito, de una sección de $\sh{O}_X$}{{\EGAReadLink{IV.16.3.6}{16.3.6}}}
\EGAIVterm{Parte cerrada puramente de codimensión $p$}{{\EGAReadLink{IV.21.6.2}{21.6.2}}}
\EGAIVterm{Polinomio separable sobre un anillo}{{\EGAReadLink{IV.18.4.3}{18.4.3}}}
\EGAIVterm{Preesquema de característica $p$}{{\EGAReadLink{IV.16.12.1}{16.12.1}}}
\EGAIVterm{Preesquema de coincidencias de dos morfismos}{{\EGAReadLink{IV.17.4.5}{17.4.5}}}
\EGAIVterm{Preesquema diferencialmente liso sobre $S$}{{\EGAReadLink{IV.16.10.1}{16.10.1}}}
\EGAIVterm{Preesquema étale sobre $S$}{{\EGAReadLink{IV.17.3.1}{17.3.1}}}
\EGAIVterm{Preesquema étale sobre $S$ en un punto}{{\EGAReadLink{IV.17.3.7}{17.3.7}}}
\EGAIVterm{Preesquema formalmente étale, formalmente liso, formalmente no ramificado sobre $S$}{{\EGAReadLink{IV.17.1.1}{17.1.1}}}
\EGAIVterm{Preesquema que es una intersección completa absoluta en un punto}{{\EGAReadLink{IV.19.3.1}{19.3.1}}}
\EGAIVterm{Preesquema liso sobre $S$}{{\EGAReadLink{IV.17.3.1}{17.3.1}}}
\EGAIVterm{Preesquema liso sobre $S$ en un punto}{{\EGAReadLink{IV.17.3.7}{17.3.7}}}
\EGAIVterm{Preesquema localmente factorial}{{\EGAReadLink{IV.21.6.9}{21.6.9}}}
\EGAIVterm{Preesquema neto sobre $S$, no ramificado sobre $S$}{{\EGAReadLink{IV.17.3.1}{17.3.1}}}

\oldpage[IV]{339}
\EGAIVterm{Preesquema neto sobre $S$ en un punto, no ramificado sobre $S$ en un punto}{{\EGAReadLink{IV.17.3.7}{17.3.7}}}
\EGAIVterm{$T$-profundidad de un $\sh{O}_X$-módulo coherente en un punto}{{\EGAReadLink{IV.19.9.1}{19.9.1}}}
\EGAIVterm{Pseudofunción}{{\EGAReadLink{IV.20.2.8}{20.2.8}}}
\EGAIVterm{Pseudofunción relativa a $S$}{{\EGAReadLink{IV.20.5.4}{20.5.4}}}
\EGAIVterm{Pseudomorfismo, pseudomorfismo sobre $S$}{{\EGAReadLink{IV.20.2.1}{20.2.1}}}
\EGAIVterm{Pseudomorfismo compuesto de un morfismo y de un pseudomorfismo}{{\EGAReadLink{IV.20.3.1}{20.3.1}}}
\EGAIVterm{Pseudomorfismo compuesto de un pseudomorfismo y de un morfismo}{{\EGAReadLink{IV.20.3.2}{20.3.2}}}
\EGAIVterm{Pseudomorfismo definido en un punto}{{\EGAReadLink{IV.20.2.3}{20.2.3}}}
\EGAIVterm{Pseudomorfismo sobre $S$}{{\EGAReadLink{IV.20.5.1}{20.5.1}}}
\EGAIVterm{Pseudotorsor}{{\EGAReadLink{IV.16.5.15}{16.5.15}}}
\EGAIVterm{Restricción de un ciclo a un abierto}{{\EGAReadLink{IV.21.6.3}{21.6.3}}}
\EGAIVterm{Restricción de un pseudomorfismo a un abierto}{{\EGAReadLink{IV.20.2.2}{20.2.2}}}
\EGAIVterm{Restricción de un pseudomorfismo a un esquema local}{{\EGAReadLink{IV.20.3.6}{20.3.6}}}
\EGAIVterm{Recubrimiento étale, recubrimiento étale localmente trivial, recubrimiento étale trivial}{{\EGAReadLink{IV.18.2.7}{18.2.7}}}
\EGAIVterm{Esquema local henseliano}{{\EGAReadLink{IV.18.5.8}{18.5.8}}}
\EGAIVterm{Esquema estrictamente local}{{\EGAReadLink{IV.18.8.2}{18.8.2}}}
\EGAIVterm{Sección meromorfa de un $\sh{O}_X$-módulo}{{\EGAReadLink{IV.20.1.3}{20.1.3}}}
\EGAIVterm{Sección meromorfa de un $\sh{O}_X$-módulo definida en un punto}{{\EGAReadLink{IV.20.1.7}{20.1.7}}}
\EGAIVterm{Sección meromorfa regular de un $\sh{O}_X$-módulo invertible}{{\EGAReadLink{IV.20.1.8}{20.1.8}}}
\EGAIVterm{Sección meromorfa de un $\sh{O}_X$-módulo relativa a $S$}{{\EGAReadLink{IV.20.6.2}{20.6.2}}}
\EGAIVterm{Sección meromorfa de un $\sh{O}_X$-módulo invertible, regular relativamente a $S$}{{\EGAReadLink{IV.20.6.5}{20.6.5}}}
\EGAIVterm{Subpreesquema inmerso cuasirregularmente, inmerso regularmente}{{\EGAReadLink{IV.16.9.2}{16.9.2}}}
\EGAIVterm{Subpreesquemas que se cortan transversalmente en un punto relativamente a $S$}{{\EGAReadLink{IV.17.13.7}{17.13.7}}}
\EGAIVterm{Soporte de un ciclo}{{\EGAReadLink{IV.21.6.1}{21.6.1}}}
\EGAIVterm{Soporte de un divisor}{{\EGAReadLink{IV.21.1.2}{21.1.2}}}
\EGAIVterm{Sucesión de secciones transversalmente regular relativamente a $S$}{{\EGAReadLink{IV.19.1.2}{19.1.2}}}
\EGAIVterm{Simetría canónica de $X\times_S X$}{{\EGAReadLink{IV.16.3.3}{16.3.3}}}
\EGAIVterm{Simetría canónica de $\sh{P}^{n}_{X/S}$, de $\sh{P}^{\infty}_{X/S}$}{{\EGAReadLink{IV.16.3.4}{16.3.4}}}
\EGAIVterm{Sistema cuasirregular (regular) de generadores de un ideal}{{\EGAReadLink{IV.16.9.1}{16.9.1}}}
\EGAIVterm{Torsor}{{\EGAReadLink{IV.16.5.15}{16.5.15}}}
\EGAIVterm{Vecindad formal de un subpreesquema}{{\EGAReadLink{IV.16.1.11}{16.1.11}}}
\EGAIVterm{$n$-ésima vecindad infinitesimal}{{\EGAReadLink{IV.16.1.2}{16.1.2}}}

\endgroup

% Traducción española de la tabla de materias de EGA IV, cuarta parte, pp. 341--343.
\clearpage
\oldpage[IV]{341}
\phantomsection
\addcontentsline{toc}{section}{Índice}
\section*{ÍNDICE}
\markboth{ÍNDICE}{ÍNDICE}

\begingroup
\small
\makeatletter
\newcommand{\EGAIVcontentsline}[3]{%
  \@dottedtocline{1}{0em}{3.5em}{\textbf{#1}\quad #2}{#3}}
\makeatother

\begin{center}\textbf{CAPÍTULO IV. --- Estudio local de los esquemas y de los morfismos de esquemas (fin)}\end{center}

\EGAIVcontentsline{\S16}{Invariantes diferenciales. Morfismos diferencialmente lisos}{5}
\EGAIVcontentsline{{\texorpdfstring{\protect\EGAReadLink{IV.16.1}{16.1}}{16.1}}}{Invariantes normales de una inmersión}{5}
\EGAIVcontentsline{{\texorpdfstring{\protect\EGAReadLink{IV.16.2}{16.2}}{16.2}}}{Propiedades functoriales de los invariantes normales de una inmersión}{9}
\EGAIVcontentsline{{\texorpdfstring{\protect\EGAReadLink{IV.16.3}{16.3}}{16.3}}}{Invariantes diferenciales fundamentales de un morfismo de preesquemas}{14}
\EGAIVcontentsline{{\texorpdfstring{\protect\EGAReadLink{IV.16.4}{16.4}}{16.4}}}{Propiedades functoriales de los invariantes diferenciales}{16}
\EGAIVcontentsline{{\texorpdfstring{\protect\EGAReadLink{IV.16.5}{16.5}}{16.5}}}{Haces y fibrados tangentes relativos; derivaciones}{27}
\EGAIVcontentsline{{\texorpdfstring{\protect\EGAReadLink{IV.16.6}{16.6}}{16.6}}}{Haces de $p$-diferenciales y diferencial exterior}{34}
\EGAIVcontentsline{{\texorpdfstring{\protect\EGAReadLink{IV.16.7}{16.7}}{16.7}}}{Los $\sh{P}^{n}_{X/S}(\sh{F})$}{36}
\EGAIVcontentsline{{\texorpdfstring{\protect\EGAReadLink{IV.16.8}{16.8}}{16.8}}}{Operadores diferenciales}{39}
\EGAIVcontentsline{{\texorpdfstring{\protect\EGAReadLink{IV.16.9}{16.9}}{16.9}}}{Inmersiones regulares y cuasirregulares}{46}
\EGAIVcontentsline{{\texorpdfstring{\protect\EGAReadLink{IV.16.10}{16.10}}{16.10}}}{Morfismos diferencialmente lisos}{51}
\EGAIVcontentsline{{\texorpdfstring{\protect\EGAReadLink{IV.16.11}{16.11}}{16.11}}}{Operadores diferenciales sobre un $S$-preesquema diferencialmente liso}{53}
\EGAIVcontentsline{{\texorpdfstring{\protect\EGAReadLink{IV.16.12}{16.12}}{16.12}}}{Caso de característica cero: criterio jacobiano para los morfismos diferencialmente lisos}{55}

\EGAIVcontentsline{\S17}{Morfismos lisos, morfismos no ramificados, morfismos étale}{56}
\EGAIVcontentsline{{\texorpdfstring{\protect\EGAReadLink{IV.17.1}{17.1}}{17.1}}}{Morfismos formalmente lisos, morfismos formalmente no ramificados, morfismos formalmente étale}{56}
\EGAIVcontentsline{{\texorpdfstring{\protect\EGAReadLink{IV.17.2}{17.2}}{17.2}}}{Propiedades generales de los diferenciales}{59}
\EGAIVcontentsline{{\texorpdfstring{\protect\EGAReadLink{IV.17.3}{17.3}}{17.3}}}{Morfismos lisos, morfismos no ramificados, morfismos étale}{61}
\EGAIVcontentsline{{\texorpdfstring{\protect\EGAReadLink{IV.17.4}{17.4}}{17.4}}}{Caracterizaciones de los morfismos no ramificados}{65}
\EGAIVcontentsline{{\texorpdfstring{\protect\EGAReadLink{IV.17.5}{17.5}}{17.5}}}{Caracterizaciones de los morfismos lisos}{67}
\EGAIVcontentsline{{\texorpdfstring{\protect\EGAReadLink{IV.17.6}{17.6}}{17.6}}}{Caracterizaciones de los morfismos étale}{70}
\EGAIVcontentsline{{\texorpdfstring{\protect\EGAReadLink{IV.17.7}{17.7}}{17.7}}}{Propiedades de descenso y de paso al límite}{72}
\EGAIVcontentsline{{\texorpdfstring{\protect\EGAReadLink{IV.17.8}{17.8}}{17.8}}}{Criterios de lisitud y de no ramificación por fibras}{79}
\EGAIVcontentsline{{\texorpdfstring{\protect\EGAReadLink{IV.17.9}{17.9}}{17.9}}}{Morfismos étale e inmersiones abiertas}{79}
\EGAIVcontentsline{{\texorpdfstring{\protect\EGAReadLink{IV.17.10}{17.10}}{17.10}}}{Dimensión relativa de un preesquema liso sobre otro}{81}
\EGAIVcontentsline{{\texorpdfstring{\protect\EGAReadLink{IV.17.11}{17.11}}{17.11}}}{Morfismos lisos de preesquemas lisos}{82}
\EGAIVcontentsline{{\texorpdfstring{\protect\EGAReadLink{IV.17.12}{17.12}}{17.12}}}{Subpreesquemas lisos de un preesquema liso. Morfismos lisos y morfismos diferencialmente lisos}{85}

\oldpage[IV]{342}
\EGAIVcontentsline{{\texorpdfstring{\protect\EGAReadLink{IV.17.13}{17.13}}{17.13}}}{Morfismos transversales}{89}
\EGAIVcontentsline{{\texorpdfstring{\protect\EGAReadLink{IV.17.14}{17.14}}{17.14}}}{Caracterizaciones locales e infinitesimales de los morfismos lisos, de los morfismos no ramificados y de los morfismos étale}{98}
\EGAIVcontentsline{{\texorpdfstring{\protect\EGAReadLink{IV.17.15}{17.15}}{17.15}}}{Caso de los preesquemas sobre un cuerpo de base}{99}
\EGAIVcontentsline{{\texorpdfstring{\protect\EGAReadLink{IV.17.16}{17.16}}{17.16}}}{Cuasi-secciones de morfismos planos o lisos}{105}

\EGAIVcontentsline{\S18}{Complementos sobre los morfismos étale. Anillos locales henselianos y anillos estrictamente locales}{109}
\EGAIVcontentsline{{\texorpdfstring{\protect\EGAReadLink{IV.18.1}{18.1}}{18.1}}}{Una equivalencia notable de categorías}{109}
\EGAIVcontentsline{{\texorpdfstring{\protect\EGAReadLink{IV.18.2}{18.2}}{18.2}}}{Recubrimientos étale}{111}
\EGAIVcontentsline{{\texorpdfstring{\protect\EGAReadLink{IV.18.3}{18.3}}{18.3}}}{Álgebras finitas y étale}{114}
\EGAIVcontentsline{{\texorpdfstring{\protect\EGAReadLink{IV.18.4}{18.4}}{18.4}}}{Estructura local de los morfismos no ramificados y de los morfismos étale}{118}
\EGAIVcontentsline{{\texorpdfstring{\protect\EGAReadLink{IV.18.5}{18.5}}{18.5}}}{Anillos locales henselianos}{125}
\EGAIVcontentsline{{\texorpdfstring{\protect\EGAReadLink{IV.18.6}{18.6}}{18.6}}}{Henselización}{135}
\EGAIVcontentsline{{\texorpdfstring{\protect\EGAReadLink{IV.18.7}{18.7}}{18.7}}}{Henselización y anillos excelentes}{142}
\EGAIVcontentsline{{\texorpdfstring{\protect\EGAReadLink{IV.18.8}{18.8}}{18.8}}}{Anillos estrictamente locales y henselización estricta}{144}
\EGAIVcontentsline{{\texorpdfstring{\protect\EGAReadLink{IV.18.9}{18.9}}{18.9}}}{Fibras formales de los anillos noetherianos henselianos}{150}
\EGAIVcontentsline{{\texorpdfstring{\protect\EGAReadLink{IV.18.10}{18.10}}{18.10}}}{Preesquemas étale sobre un preesquema geométricamente unirrama o normal}{157}
\EGAIVcontentsline{{\texorpdfstring{\protect\EGAReadLink{IV.18.11}{18.11}}{18.11}}}{Aplicación a las álgebras locales noetherianas completas sobre un cuerpo}{169}
\EGAIVcontentsline{{\texorpdfstring{\protect\EGAReadLink{IV.18.12}{18.12}}{18.12}}}{Aplicaciones de la localización étale a los morfismos cuasi-finitos (generalizaciones de resultados anteriores)}{181}

\EGAIVcontentsline{\S19}{Inmersiones regulares y planitud normal}{185}
\EGAIVcontentsline{{\texorpdfstring{\protect\EGAReadLink{IV.19.1}{19.1}}{19.1}}}{Propiedades de las inmersiones regulares}{185}
\EGAIVcontentsline{{\texorpdfstring{\protect\EGAReadLink{IV.19.2}{19.2}}{19.2}}}{Inmersiones transversalmente regulares}{190}
\EGAIVcontentsline{{\texorpdfstring{\protect\EGAReadLink{IV.19.3}{19.3}}{19.3}}}{Intersecciones completas relativas (caso plano)}{194}
\EGAIVcontentsline{{\texorpdfstring{\protect\EGAReadLink{IV.19.4}{19.4}}{19.4}}}{Aplicación: criterios de regularidad y de lisitud para los preesquemas obtenidos por explosión}{198}
\EGAIVcontentsline{{\texorpdfstring{\protect\EGAReadLink{IV.19.5}{19.5}}{19.5}}}{Criterios de $M$-regularidad}{204}
\EGAIVcontentsline{{\texorpdfstring{\protect\EGAReadLink{IV.19.6}{19.6}}{19.6}}}{Sucesiones regulares relativamente a un módulo filtrado cociente}{209}
\EGAIVcontentsline{{\texorpdfstring{\protect\EGAReadLink{IV.19.7}{19.7}}{19.7}}}{Criterio de planitud normal de Hironaka}{212}
\EGAIVcontentsline{{\texorpdfstring{\protect\EGAReadLink{IV.19.8}{19.8}}{19.8}}}{Propiedades del paso al límite proyectivo}{219}
\EGAIVcontentsline{{\texorpdfstring{\protect\EGAReadLink{IV.19.9}{19.9}}{19.9}}}{Sucesiones $\sh{F}$-regulares y profundidad}{222}

\EGAIVcontentsline{\S20}{Funciones meromorfas y pseudomorfismos}{223}
\EGAIVcontentsline{{\texorpdfstring{\protect\EGAReadLink{IV.20.0}{20.0}}{20.0}}}{Introducción}{223}
\EGAIVcontentsline{{\texorpdfstring{\protect\EGAReadLink{IV.20.1}{20.1}}{20.1}}}{Funciones meromorfas}{224}
\EGAIVcontentsline{{\texorpdfstring{\protect\EGAReadLink{IV.20.2}{20.2}}{20.2}}}{Pseudomorfismos y pseudofunciones}{231}

\oldpage[IV]{343}
\EGAIVcontentsline{{\texorpdfstring{\protect\EGAReadLink{IV.20.3}{20.3}}{20.3}}}{Composición de los pseudomorfismos}{237}
\EGAIVcontentsline{{\texorpdfstring{\protect\EGAReadLink{IV.20.4}{20.4}}{20.4}}}{Propiedades de los dominios de definición de las funciones racionales}{244}
\EGAIVcontentsline{{\texorpdfstring{\protect\EGAReadLink{IV.20.5}{20.5}}{20.5}}}{Pseudomorfismos relativos}{249}
\EGAIVcontentsline{{\texorpdfstring{\protect\EGAReadLink{IV.20.6}{20.6}}{20.6}}}{Funciones meromorfas relativas}{252}

\EGAIVcontentsline{\S21}{Divisores}{255}
\EGAIVcontentsline{{\texorpdfstring{\protect\EGAReadLink{IV.21.1}{21.1}}{21.1}}}{Divisores sobre un espacio anillado}{255}
\EGAIVcontentsline{{\texorpdfstring{\protect\EGAReadLink{IV.21.2}{21.2}}{21.2}}}{Divisores e ideales fraccionarios invertibles}{258}
\EGAIVcontentsline{{\texorpdfstring{\protect\EGAReadLink{IV.21.3}{21.3}}{21.3}}}{Equivalencia lineal de los divisores}{263}
\EGAIVcontentsline{{\texorpdfstring{\protect\EGAReadLink{IV.21.4}{21.4}}{21.4}}}{Imágenes recíprocas de divisores}{265}
\EGAIVcontentsline{{\texorpdfstring{\protect\EGAReadLink{IV.21.5}{21.5}}{21.5}}}{Imágenes directas de divisores}{267}
\EGAIVcontentsline{{\texorpdfstring{\protect\EGAReadLink{IV.21.6}{21.6}}{21.6}}}{Ciclo 1-codimensional asociado a un divisor}{270}
\EGAIVcontentsline{{\texorpdfstring{\protect\EGAReadLink{IV.21.7}{21.7}}{21.7}}}{Interpretación de los ciclos positivos de codimensión 1 en términos de subpreesquemas}{277}
\EGAIVcontentsline{{\texorpdfstring{\protect\EGAReadLink{IV.21.8}{21.8}}{21.8}}}{Divisores y normalización}{280}
\EGAIVcontentsline{{\texorpdfstring{\protect\EGAReadLink{IV.21.9}{21.9}}{21.9}}}{Divisores sobre los preesquemas de dimensión 1}{284}
\EGAIVcontentsline{{\texorpdfstring{\protect\EGAReadLink{IV.21.10}{21.10}}{21.10}}}{Imágenes recíprocas e imágenes directas de ciclos 1-codimensionales}{289}
\EGAIVcontentsline{{\texorpdfstring{\protect\EGAReadLink{IV.21.11}{21.11}}{21.11}}}{Factorialidad de los anillos regulares}{302}
\EGAIVcontentsline{{\texorpdfstring{\protect\EGAReadLink{IV.21.12}{21.12}}{21.12}}}{El teorema de pureza de Van der Waerden para el conjunto de ramificación de un morfismo birracional}{304}
\EGAIVcontentsline{{\texorpdfstring{\protect\EGAReadLink{IV.21.13}{21.13}}{21.13}}}{Pares parafactoriales. Anillos locales parafactoriales}{313}
\EGAIVcontentsline{{\texorpdfstring{\protect\EGAReadLink{IV.21.14}{21.14}}{21.14}}}{El teorema de Ramanujam--Samuel}{323}
\EGAIVcontentsline{{\texorpdfstring{\protect\EGAReadLink{IV.21.15}{21.15}}{21.15}}}{Divisores relativos}{329}
\EGAIVcontentsline{}{BIBLIOGRAFÍA}{333}
\EGAIVcontentsline{}{ÍNDICE DE NOTACIONES}{334}
\EGAIVcontentsline{}{ÍNDICE TERMINOLÓGICO}{336}

\endgroup

% Traducción española de la transcripción diplomática francesa de EGA IV, cuarta parte, pp. 345--361.
\clearpage
\oldpage[IV]{345}
\phantomsection
\addcontentsline{toc}{section}{Erratas y adendas (Lista 3)}
\section*{ERRATAS Y ADENDAS}
\markboth{ERRATAS Y ADENDAS, LISTA 3}{ERRATAS Y ADENDAS, LISTA 3}
\begin{center}\large (Lista 3)\end{center}

\subsection*{A) Errores tipográficos}

\paragraph{(II, 3.3.2).}
Línea 1 de la p.~58, sustituir $\sh{S}$ por $\sh{M}$. Línea 6 de la
p.~58, sustituir $S(n)$ por $M(n)$.

\paragraph{(II, 4.1.1).}
Línea 8 de la p.~71, sustituir $\sh{O}$-$P$-Module por
$\sh{O}_P$-Module.

\paragraph{(II, 5.4.5).}
Línea 6 de la p.~102, sustituir $u$ por $f\circ u$.

\paragraph{(II, 5.4.8).}
Antes de la línea 4 desde abajo de la p.~102, sustituir $f\times1_{S'}$ por
$f_{(Y')}$ en el diagrama. Línea 4 desde abajo de la p.~102, sustituir
$f\times1_{S'}$ por $f_{(Y')}$.

\paragraph{(II, 6.4.4).}
Línea 1 desde arriba de la p.~122, sustituir normal por intégralement clos.

\paragraph{(II, 6.4.7).}
Línea 14 desde abajo de la p.~122, sustituir dimension por rang.

\paragraph{(II, 6.5.2).}
Línea 5 de la p.~127, sustituir $(\sh{B}|U_\lambda)$ por
$(B|U_\lambda)$. Línea 15 de la p.~127, suprimir el primer
$(N_{\sh{B}/\sh{A}}\omega_{\lambda\mu})$.

\paragraph{($0_{\mathrm{III}}$, 8.1.11).}
Línea 8 desde abajo de la p.~8, sustituir $h_X(u)$ por $h'_X(u)$.

\paragraph{($0_{\mathrm{III}}$, 9.1.2).}
Línea 10 desde abajo de la p.~12, sustituir $\sh{F}$ por $X$.

\paragraph{($0_{\mathrm{III}}$, 10.2.1).}
Línea 6 desde abajo de la p.~18, sustituir « en outre si » por « si en
outre ».

\paragraph{(III, 4.4.12).}
Línea 10 de la p.~138, sustituir chap.~V por chap.~IV (8.11.2).

\paragraph{(III, 5.2.5).}
Línea 12 de la p.~153, sustituir (I, 5.5.4) por (II, 5.5.4).

\paragraph{(III, 7.4.7).}
Línea 15 de la p.~57, sustituir tout $A$-module por tout
$A$-module $M$.

\paragraph{(III, 7.8.5).}
Línea 18 desde abajo de la p.~74, sustituir $F$ por $\sh{F}$.

\paragraph{(III, 7.8.9).}
Línea 9 desde abajo de la p.~75, sustituir $M$ por $\sh{M}$.

\paragraph{(III, 7.9.11).}
Línea 17 desde abajo de la p.~79, sustituir chap.~V por chap.~VI.

\paragraph{($0_{\mathrm{IV}}$, 16.2.2.2).}
Línea 13 de la p.~26, sustituir $M'\cap\mathfrak q^nM'$ por
$M'\cap\mathfrak q^nM$.

\paragraph{($0_{\mathrm{IV}}$, 17.3.7).}
Línea 16 desde abajo de la p.~50, sustituir (17.1.10) por (17.3.3).

\paragraph{($0_{\mathrm{IV}}$, 18.1.3).}
En el diagrama (18.1.3.1), $p_2$ debe estar junto a la flecha
$G\to F$.

\paragraph{($0_{\mathrm{IV}}$, 19.8.7).}
Líneas 7, 8 y 9 desde abajo de la p.~111, sustituir (21.5.5) por (21.5.3).

\paragraph{($0_{\mathrm{IV}}$, 19.8.8).}
Línea 2 desde abajo de la p.~112, sustituir $A$-module por $B$-module.

\paragraph{($0_{\mathrm{IV}}$, 20.5.4).}
Línea 5 desde abajo de la p.~130, sustituir $\Omega_{B'/A'}$ por
$\Omega^1_{B'/A'}$.

\paragraph{($0_{\mathrm{IV}}$, 20.7.8).}
Línea 9 desde abajo de la p.~149, sustituir (18.4.3) por (18.5.4). Línea 4
desde abajo, sustituir (18.4.1) por (18.5.1).

\paragraph{($0_{\mathrm{IV}}$, 21.5.3).}
Línea 1 desde abajo de la p.~166, sustituir ramifiée por ramifié.

\paragraph{($0_{\mathrm{IV}}$, 21.7.1).}
Línea 4 de la p.~170, sustituir (20.6.15.1) por (20.6.17.1).

\paragraph{($0_{\mathrm{IV}}$, 22.1.1).}
Línea 12 de la p.~183, sustituir (20.6.20) por (20.6.10).

\oldpage[IV]{346}
\paragraph{($0_{\mathrm{IV}}$, 22.3.3).}
Línea 13 desde abajo de la p.~192, sustituir $\widehat B_{\mathfrak p}$ por
$\widehat B_{\mathfrak q}$.

\paragraph{($0_{\mathrm{IV}}$, 22.7.3).}
Línea 21 desde abajo de la p.~212, sustituir $\operatorname{Ker}(i)$ por
$\operatorname{Ker}(i_0)$. Línea 12 desde abajo de la p.~212, sustituir
$(\mathfrak p/\mathfrak p^2)\otimes_L L$ por
$(\mathfrak p/\mathfrak p^2)\otimes_A L$.

\paragraph{($0_{\mathrm{IV}}$, 23.1.3).}
Línea 4 de la p.~215, sustituir $x'^q\in K'$ por $x'\in K'$. Línea 9
de la p.~215, sustituir $x'^q\in A'$ por $x'^q\in A$.

\paragraph{({\texorpdfstring{\protect\EGAReadLink{IV.2.1.1}{IV, 2.1.1}}{IV, 2.1.1}}).}
Línea 4 desde abajo de la p.~5, sustituir (III, 1.4.15.1) por
(III, 1.4.15.5).

\paragraph{({\texorpdfstring{\protect\EGAReadLink{IV.2.4.3}{IV, 2.4.3}}{IV, 2.4.3}}).}
Línea 12 de la p.~20, sustituir
$(X_{\mathrm{red}}\times_{Y_{\mathrm{red}}}Y_{\mathrm{red}})_{\mathrm{red}}$
por
$(X_{\mathrm{red}}\times_{Y_{\mathrm{red}}}Y'_{\mathrm{red}})_{\mathrm{red}}$.

\paragraph{({\texorpdfstring{\protect\EGAReadLink{IV.2.6.1}{IV, 2.6.1}}{IV, 2.6.1}}).}
Línea 10 de la p.~27, sustituir $X'\to X_{(S')}$ por
$X'=X_{(S')}$. 

\paragraph{({\texorpdfstring{\protect\EGAReadLink{IV.4.2.1}{IV, 4.2.1}}{IV, 4.2.1}}).}
Línea 20 desde abajo de la p.~55, sustituir
$\operatorname{deg.tr.}_kE$ por $\operatorname{deg.tr.}_LE$ y
$\operatorname{deg.tr.}_kL(t)$ por $\operatorname{deg.tr.}_LL(t)$;
línea 19 desde abajo, intercambiar « première » y « seconde ».

\paragraph{({\texorpdfstring{\protect\EGAReadLink{IV.4.4.2}{IV, 4.4.2}}{IV, 4.4.2}}).}
Línea 1 desde abajo de la p.~59, sustituir $X$ por $X'$.

\paragraph{({\texorpdfstring{\protect\EGAReadLink{IV.4.7.3}{IV, 4.7.3}}{IV, 4.7.3}}).}
Línea 10 desde abajo de la p.~75, sustituir préschéma por $k$-préschéma.

\paragraph{({\texorpdfstring{\protect\EGAReadLink{IV.5.6.1}{IV, 5.6.1}}{IV, 5.6.1}}).}
Línea 16 de la p.~98, sustituir (en el subíndice) $k(\mathfrak p($ por
$k(\mathfrak p)$.

\paragraph{({\texorpdfstring{\protect\EGAReadLink{IV.5.6.5}{IV, 5.6.5}}{IV, 5.6.5}}).}
Línea 20 de la p.~99, suprimir el segundo paréntesis en $k(\xi))$.

\paragraph{({\texorpdfstring{\protect\EGAReadLink{IV.5.6.11.1}{IV, 5.6.11.1}}{IV, 5.6.11.1}}).}
Línea 16 de la p.~102, sustituir $K$ por $\mathfrak K$.

\paragraph{({\texorpdfstring{\protect\EGAReadLink{IV.5.11.2}{IV, 5.11.2}}{IV, 5.11.2}}).}
Nota al pie de la p.~123, sustituir 5.11.2) por (5.11.2).

\paragraph{({\texorpdfstring{\protect\EGAReadLink{IV.6.1.5}{IV, 6.1.5}}{IV, 6.1.5}}).}
Línea 2 desde abajo de la p.~136, sustituir (5.5.1.2) por
($0$, 16.3.9.2).

\paragraph{({\texorpdfstring{\protect\EGAReadLink{IV.6.3.1}{IV, 6.3.1}}{IV, 6.3.1}}).}
Línea 9 de la p.~139, sustituir $y\in\mathfrak m$ por
$y\in\mathfrak n$.

\paragraph{({\texorpdfstring{\protect\EGAReadLink{IV.6.7.2}{IV, 6.7.2}}{IV, 6.7.2}}).}
Línea 18 de la p.~146, sustituir $\sh{O}'_{x'}$ por $\sh{O}_{x'}$.

\paragraph{({\texorpdfstring{\protect\EGAReadLink{IV.6.10.1}{IV, 6.10.1}}{IV, 6.10.1}}).}
Línea 13 de la p.~155, sustituir
$\sh{F}\otimes_{\sh{O}_Y}\sh{O}_Y$ por
$\sh{F}\otimes_{\sh{O}_X}\sh{O}_Y$.

\paragraph{({\texorpdfstring{\protect\EGAReadLink{IV.6.14.1}{IV, 6.14.1}}{IV, 6.14.1}}).}
Línea 21 de la p.~170, sustituir n\textsuperscript{o}~3, prop.~12
por n\textsuperscript{o}~2, prop.~8.

\paragraph{({\texorpdfstring{\protect\EGAReadLink{IV.6.15.2}{IV, 6.15.2}}{IV, 6.15.2}}).}
Línea 13 de la p.~177, sustituir $t(t-v)$ por $t^2(t-v)^2$.

\paragraph{({\texorpdfstring{\protect\EGAReadLink{IV.6.15.7}{IV, 6.15.7}}{IV, 6.15.7}}).}
Línea 6 desde abajo de la p.~179, sustituir $X\otimes_Ak'$ por
$X\otimes_kk'$.

\paragraph{({\texorpdfstring{\protect\EGAReadLink{IV.7.5.4}{IV, 7.5.4}}{IV, 7.5.4}}).}
Línea 14 desde abajo de la p.~206, sustituir
$\operatorname{Spec}(\widehat B\otimes k)$ por
$\operatorname{Spec}(\widehat B\otimes_{\widehat A}k)$.

\paragraph{({\texorpdfstring{\protect\EGAReadLink{IV.7.8.4}{IV, 7.8.4}}{IV, 7.8.4}}).}
Línea 10 desde abajo de la p.~216, sustituir « vérifiant »
({\EGAReadLink{IV.7.8.2}{7.8.2}}, (ii) y (iii)) por
« vérifiant ({\EGAReadLink{IV.7.8.2}{7.8.2}}, (ii) et (iii)) ».

\paragraph{(IV, 2\textsuperscript{e} Partie, Bibliographie).}
Línea 6 de la p.~224, sustituir \emph{Amer. J. of Math.} por
\emph{Annals of Math.}

\paragraph{({\texorpdfstring{\protect\EGAReadLink{IV.8.5.8}{IV, 8.5.8}}{IV, 8.5.8}}).}
Línea 6 de la p.~24, sustituir $\operatorname{Ker}g_\mu$ al final de
la línea por $\operatorname{Ker}g_\lambda$.

\paragraph{({\texorpdfstring{\protect\EGAReadLink{IV.8.10.1}{IV, 8.10.1}}{IV, 8.10.1}}).}
Línea 13 desde abajo de la p.~36, sustituir (2.3.11) por
{\EGAReadLink{IV.8.2.11}{(8.2.11)}}.

\paragraph{({\texorpdfstring{\protect\EGAReadLink{IV.10.2.1}{IV, 10.2.1}}{IV, 10.2.1}}).}
Línea 14 desde abajo de la p.~98, sustituir $y$ por $Y$.

\paragraph{({\texorpdfstring{\protect\EGAReadLink{IV.10.4.11}{IV, 10.4.11}}{IV, 10.4.11}}).}
Línea 10 desde abajo de la p.~104, sustituir (17.9.7) por
{\EGAReadLink{IV.17.9.6}{(17.9.6)}}.

\paragraph{({\texorpdfstring{\protect\EGAReadLink{IV.11.2.9}{IV, 11.2.9}}{IV, 11.2.9}}).}
Línea 3 desde abajo de la p.~126, sustituir $A_\lambda$ por $A_\mu$.

\paragraph{({\texorpdfstring{\protect\EGAReadLink{IV.11.3.4}{IV, 11.3.4}}{IV, 11.3.4}}).}
Línea 16 de la p.~134, sustituir $F$ por $\sh{F}$.

\paragraph{({\texorpdfstring{\protect\EGAReadLink{IV.11.3.10}{IV, 11.3.10}}{IV, 11.3.10}}).}
Línea 9 desde abajo de la p.~139, sustituir (11.3.7) por
{\EGAReadLink{IV.11.3.10}{(11.3.10)}}.

\paragraph{({\texorpdfstring{\protect\EGAReadLink{IV.11.7.5}{IV, 11.7.5}}{IV, 11.7.5}}).}
Línea 11 de la p.~159, sustituir (6.5.11) por (6.15.11).

\paragraph{({\texorpdfstring{\protect\EGAReadLink{IV.12.1.1}{IV, 12.1.1}}{IV, 12.1.1}}).}
Línea 13 de la p.~174, añadir al final: pour $k\geq1$.

\paragraph{({\texorpdfstring{\protect\EGAReadLink{IV.14.2.1}{IV, 14.2.1}}{IV, 14.2.1}}).}
Línea 11 de la p.~202, sustituir (6.1.2) por (6.1.1).

\oldpage[IV]{347}
\paragraph{({\texorpdfstring{\protect\EGAReadLink{IV.14.3.8}{IV, 14.3.8}}{IV, 14.3.8}}).}
Línea 10 desde abajo de la p.~206, sustituir $(f^{-1}y)$ por $f^{-1}(y)$.

\paragraph{({\texorpdfstring{\protect\EGAReadLink{IV.14.5.7}{IV, 14.5.7}}{IV, 14.5.7}}).}
Línea 16 de la p.~218, sustituir (14.5.5) por (14.5.6).

\paragraph{({\texorpdfstring{\protect\EGAReadLink{IV.14.5.11}{IV, 14.5.11}}{IV, 14.5.11}}).}
En el diagrama de la parte superior de la página 223, todas las flechas
horizontales deben ir de derecha a izquierda.

\paragraph{({\texorpdfstring{\protect\EGAReadLink{IV.15.4.2}{IV, 15.4.2}}{IV, 15.4.2}}).}
Línea 23 de la p.~230, sustituir (6.1.4) por (6.1.5).

\paragraph{({\texorpdfstring{\protect\EGAReadLink{IV.15.5.2}{IV, 15.5.2}}{IV, 15.5.2}}).}
Línea 16 de la p.~233, sustituir $f^{-1}(y)$ por $f^{-1}(y')$.

\paragraph{({\texorpdfstring{\protect\EGAReadLink{IV.15.6.10}{IV, 15.6.10}}{IV, 15.6.10}}).}
Línea 6 de la p.~242, sustituir ${}^{-1}(y')$ por $f^{-1}(y')$.
Línea 11 de la p.~242, sustituir $\alpha$) et $\beta$) por
$\alpha$) ou $\beta$).

\subsection*{B) Modificaciones del texto}

\paragraph{(Err\textsubscript{IV}, 1).}
En ($0_{\mathrm{III}}$, 11.1.6), el enunciado (que reproduce un enunciado de
(G, I, 4.4.1)) es incorrecto cuando $q(n)$ es un entero arbitrario;
solo es correcto si $q(n)=0$ o $q(n)=n$. En efecto, en el
primer caso, el núcleo de $d_r^{p,0}$ es entonces $E_r^{p,0}$, pues
la imagen de $d_r^{p,0}$ es $E_r^{p+r,-r+1}=0$ por hipótesis; asimismo,
la imagen de $d_r^{p-r,r-1}$ es nula puesto que
$E_r^{p-r,r-1}=0$; se tiene, por tanto, $E_{r+1}^{p,0}=E_r^{p,0}$. Se
razona del mismo modo cuando $q(n)=n$ para todo $n$.

\paragraph{(Err\textsubscript{IV}, 2).}
En (III, 3.4), suprimir las líneas 2 y 3 de la p.~119.

\oldpage[IV]{348}
\paragraph{(Err\textsubscript{IV}, 3).}
En (III, 4.4.9), línea 15 desde abajo, p.~137, sustituir finies por
finies en tant qu'ensembles.

\paragraph{(Err\textsubscript{IV}, 4).}
En (III, 7.8.10), línea 9 de la p.~76, sustituir simple por lisse.

\paragraph{(Err\textsubscript{IV}, 5).}
En ($0_{\mathrm{IV}}$, 16.1.4), línea 8 desde abajo de la p.~23, sustituir
Dire que $A$ est caténaire por Dire qu'un anneau intègre $A$ est
caténaire.

\paragraph{(Err\textsubscript{IV}, 6).}
En ($0_{\mathrm{IV}}$, 16.1.9), puede suprimirse la hipótesis de que
$A$ sea noetheriano. En efecto, para todo elemento $u\in B_M$, $A[u]$
es una sub-$A$-álgebra de $C$, luego $u$ es entero sobre $A$ (Bourbaki,
\emph{Alg. comm.}, cap.~V, \S1, n\textsuperscript{o}~1, cor. de la
prop.~1), y por consiguiente también sobre $A_M$; se concluye igualmente mediante
{\EGAReadLink{IV.16.1.5}{(16.1.5)}}.

\paragraph{(Err\textsubscript{IV}, 7).}
En ($0_{\mathrm{IV}}$, 16.5.12), obsérvese que la relación
(16.5.11.1) equivale a la misma relación en la que se sustituye $A$ por los
anillos cociente $A/\mathfrak r$, al recorrer $\mathfrak r$ el conjunto
de los ideales primos minimales de $A$. Por tanto, en la demostración de
{\EGAReadLink{IV.16.5.12}{(16.5.12)}} podemos limitarnos al caso
en que $A$ sea íntegro y aplicar (16.1.4.2).

\paragraph{(Err\textsubscript{IV}, 8).}
Suprimir ($0_{\mathrm{IV}}$, 19.3.12), cuya demostración es errónea
(los conjuntos $W_\lambda$, $U_\lambda$ y $V_\lambda$ no son
variedades lineales).

\paragraph{(Err\textsubscript{IV}, 9).}
En ($0_{\mathrm{IV}}$, 19.8.3), líneas 6 y 8 desde abajo de la p.~110,
sustituir homomorphisme por homomorphisme local; lo mismo en
($0_{\mathrm{IV}}$, 19.8.4), línea 1 de la p.~111.

\paragraph{(Err\textsubscript{IV}, 10).}
En ($0_{\mathrm{IV}}$, 20.1.2), línea 2 desde abajo de la p.~117,
sustituir $a\in A$ por $b\in B$ y $aD$ por $bD$; línea 1 desde abajo de la
p.~117, sustituir $A$-module por $B$-module.

\paragraph{(Err\textsubscript{IV}, 11).}
En ($0_{\mathrm{IV}}$, 20.2.1), línea 8 de la p.~119, sustituir
$A$-modules por $B$-modules. Líneas 21 y 22 de la p.~119, sustituir
« di-homomorphismes de modules (relatif à $u$) » por homomorphisme
de $B$-modules.

\paragraph{(Err\textsubscript{IV}, 12).}
En ($0_{\mathrm{IV}}$, 20.4.6.1), sustituir $p_2-p_1$ por $p_1-p_2$;
línea 10 desde abajo de la p.~125, sustituir $x\otimes1-1\otimes x$ por
$1\otimes x-x\otimes1$.

\paragraph{(Err\textsubscript{IV}, 13).}
En ($0_{\mathrm{IV}}$, 20.4.8), línea 14 de la p.~126, sustituir
$p_2=p_1-d_{B/A}$ por $p_1=p_2-d_{B/A}$.

\paragraph{(Err\textsubscript{IV}, 14).}
En ($0_{\mathrm{IV}}$, 20.4.13), sustituir las líneas 6 a 11 desde abajo
de la p.~127 por el texto siguiente:

\medskip
\noindent
(iv) Sean $A$, $B$ dos anillos íntegros tales que $A\subset B$, siendo $B$
entero sobre $A$, y que el cuerpo de fracciones de $B$ sea una
extensión separable del de $A$. Entonces el $B$-módulo
$\Omega^1_{B/A}$ es un módulo de torsión. En efecto, para todo $x\in B$,
sea $f$ el polinomio minimal de $x$ respecto del cuerpo de fracciones
$K$ de $A$; existe un $a\in A$ no nulo tal que $af[T]\in A[T]$; como
$x$ es separable sobre $K$, se tiene $f'(x)\neq0$, y, por otra parte, de la
relación $af(x)=0$ se deduce $af'(x)d_{B/A}x=0$, de donde se sigue nuestra afirmación
en virtud de {\EGAReadLink{IV.20.4.7}{(20.4.7)}}.

\paragraph{(Err\textsubscript{IV}, 15).}
En ($0_{\mathrm{IV}}$, 22.2.7), línea 5 de la p.~189, en lugar de
« formellement lisses », debe decir « formellement lisses pour les
topologies préadiques ».

\paragraph{(Err\textsubscript{IV}, 16).}
En el Sumario del cap.~IV, línea 4 desde abajo de la p.~222, añadir:
et anneaux strictement locaux. Sustituir las líneas 1 a 3 desde abajo por:
\begin{quote}
\S19. Inmersiones regulares y planitud normal.\\
\S20. Funciones meromorfas y seudomorfismos.\\
\S21. Divisores.
\end{quote}
Línea 16 de la p.~223, sustituir « de présentation finie » por
« localement de présentation finie »; línea 19 de la p.~223, sustituir
« comme » por « comme étant ».

\paragraph{(Err\textsubscript{IV}, 17).}
En (IV, 3.1.10), añadir después de la línea 16 de la p.~39:

\begin{remark}[3.1.10.2]
\label{IV.3.1.10.2}
Sea $f:X\to Y$ un morfismo cuasicompacto y cuasiseparado, y sea
$\sh{F}$ un $\sh{O}_X$-módulo cuasicoherente, de modo que
$f_*(\sh{F})$ es un $\sh{O}_Y$-módulo cuasicoherente
{\EGAReadLink{IV.1.7.4}{(1.7.4)}}. Entonces se tiene
\begin{align}
 \operatorname{Ass}(f_*(\sh{F}))&\subset f(X),
 \label{IV.3.1.10.3}\tag{3.1.10.3}\\
 \operatorname{Ass}(f_*(\sh{F}))&\subset f(\operatorname{Ass}(\sh{F}))
 \quad\text{si $X$ es localmente noetheriano}.
 \label{IV.3.1.10.4}\tag{3.1.10.4}
\end{align}

Sea, en efecto, $y\in\operatorname{Ass}(f_*(\sh{F}))$; en virtud de
{\EGAReadLink{IV.2.3.1}{(2.3.1)}} y del hecho de que el morfismo
$\operatorname{Spec}(\sh{O}_{Y,y})\to Y$ es plano,
$(f_*(\sh{F}))_y$ es igual a $(f'_*(\sh{F}'))_y$, donde
\[
 f':X\times_Y\operatorname{Spec}(\sh{O}_{Y,y})
 \longrightarrow\operatorname{Spec}(\sh{O}_{Y,y})
\]
se deduce de $f$ por cambio de base y
$\sh{F}'=\sh{F}\otimes_{\sh{O}_Y}\sh{O}_{Y,y}$. Quedamos, por tanto, reducidos a
demostrar (3.1.10.3) y (3.1.10.4), sustituyendo en ellas $Y$ por
$\operatorname{Spec}(\sh{O}_{Y,y})$, $f$ por $f'$ y $\sh{F}$ por
$\sh{F}'$, teniendo en cuenta (I, 3.6.5) y la definición de
$\operatorname{Ass}(\sh{F}')$.

Por definición, existe un elemento
$t_y\in(f_*(\sh{F}))_y$ no nulo cuyo anulador es $\mathfrak m_y$, y como
$(f_*(\sh{F}))_y$ se identifica con
$\Gamma(Y,f_*(\sh{F}))$ puesto que $Y$ es local, $t_y$ se identifica asimismo
con una sección $s\in\Gamma(X,\sh{F})$. El ideal de $\sh{O}_X$ anulador
de $s$ contiene, por tanto, $\mathfrak m_y\sh{O}_X$, y se deduce de
($0_{\mathrm I}$, 1.7.4) que el soporte de $s$ está contenido en
$f^{-1}(y)$; si este último fuese vacío, se tendría $s=0$, lo cual es
absurdo. Si además $X$ es localmente noetheriano, el
sub-$\sh{O}_X$-módulo coherente $\sh{O}_Xs$ de $\sh{F}$ no es nulo,
luego $\operatorname{Ass}(\sh{O}_Xs)\neq\varnothing$
{\EGAReadLink{IV.3.1.5}{(3.1.5)}}; pero
\[
 \operatorname{Ass}(\sh{O}_Xs)\subset
 \operatorname{Supp}(\sh{O}_Xs)\subset f^{-1}(y),\qquad
 \operatorname{Ass}(\sh{O}_Xs)\subset\operatorname{Ass}(\sh{F})
\]
{\EGAReadLink{IV.3.1.7}{(3.1.7)}}, por tanto
$y\in f(\operatorname{Ass}(\sh{F}))$.
\end{remark}

\oldpage[IV]{349}
\begin{remark}[3.1.10.5]
\label{IV.3.1.10.5}
Obsérvese que, aun cuando $f$ es un morfismo propio, no se tiene
necesariamente
\[
 f(\operatorname{Ass}(\sh{F}))\subset
 \operatorname{Ass}(f_*(\sh{F})).
\]
Esto se ve tomando como $Y$ el espectro de un cuerpo y como $X$ un
espacio proyectivo sobre $k$, por ejemplo $\mathbf P^1_k$; entonces
$\sh{F}=\sh{O}_X(-1)$ es un $\sh{O}_X$-módulo coherente $\neq0$ tal
que $\Gamma(X,\sh{F})=0$ (III, 2.1.13), luego
$\operatorname{Ass}(f_*(\sh{F}))=\varnothing$ y
$\operatorname{Ass}(\sh{F})\neq\varnothing$.
\end{remark}

\paragraph{(Err\textsubscript{IV}, 18).}
Después de (IV, 3.2.8), insertar después de la línea 17 de la p.~43:

\begin{remark}[3.2.9]
\label{IV.3.2.9}
Para todo punto $s$ de un preesquema $X$, denotemos por
$i_s:\operatorname{Spec}(k(s))\to X$ el morfismo canónico, que es
un monomorfismo cuasicompacto. Para todo $k(s)$-módulo $G(s)$,
considerado como módulo sobre $\operatorname{Spec}(k(s))$,
$(i_s)_*(G(s))$ es, por tanto, un $\sh{O}_X$-módulo cuasicoherente
{\EGAReadLink{IV.1.7.4}{(1.7.4)}}. Se comprueba inmediatamente
que la fibra de $(i_s)_*(G(s))$ es nula en todo punto
$x\notin\overline{\{s\}}$ e isomorfa a $G(s)$ en caso contrario.
Si $X$ es localmente noetheriano y $G(s)$ es de longitud finita, está
claro que $\operatorname{Ass}((i_s)_*(G(s)))$ se reduce al único punto
$s$; en otras palabras, $(i_s)_*(G(s))$ es irredundante.

Supongamos ahora que $X$ es localmente noetheriano, y sea $\sh{F}$ un
$\sh{O}_X$-módulo coherente. Sea, por otra parte, $S$ una parte discreta
de $X$, y, para todo $s\in S$, pongamos
$\sh{F}(s)=i_s^*(\sh{F})$, que es un $k(s)$-módulo; consideremos para
cada $s\in S$ un módulo cociente $G(s)$ de $\sh{F}(s)$ de longitud
finita. Se tiene entonces un homomorfismo canónico
\[
 u:\sh{F}\longrightarrow
 \bigoplus_{s\in S}(i_s)_*(\sh{F}(s))\longrightarrow
 \bigoplus_{s\in S}(i_s)_*(G(s))
\]
de $\sh{O}_X$-módulos coherentes (pues $S$, al ser discreto, es
localmente finito en $X$). Si $\sh{H}=\operatorname{Ker}u$,
$\sh{G}=\sh{F}/\sh{H}$ es un $\sh{O}_X$-módulo coherente y se deduce
de $u$ un homomorfismo inyectivo
\[
 v:\sh{G}\longrightarrow\bigoplus_{s\in S}(i_s)_*(G(s)).
\]
Si $\sh{G}^{(s)}$ es la imagen de $\sh{G}$ en $(i_s)_*(G(s))$,
$(\sh{G}^{(s)})_{s\in S}$ es una descomposición irredundante, que es
reducida cuando todos los $G(s)$ son $\neq0$, y tal que
$\operatorname{Ass}(\sh{G})=S$ (por tanto, $\sh{G}$ no tiene ciclo primo
asociado sumergido). Se obtiene así una biyección
$(G(s))_{s\in S}\leftrightarrow\sh{G}$, donde $S$ recorre el conjunto de las
partes discretas de $X$ para las cuales $\sh{F}(s)\neq0$ para todo
$s\in S$ y, para cada $S$, $(G(s))_{s\in S}$ recorre el conjunto
de los cocientes $\neq0$ y de longitud finita de los $\sh{F}(s)$; por otra
parte, $\sh{G}$ recorre el conjunto de los $\sh{O}_X$-módulos cocientes
coherentes de $\sh{F}$, sin ciclo primo sumergido. La biyección
recíproca asocia a cada $\sh{G}$ su descomposición irredundante
reducida.
\end{remark}

\paragraph{(Err\textsubscript{IV}, 19).}
En (IV, 4.2.1), antes de la línea 17 desde abajo de la p.~55, insertar:

\begin{remark}[4.2.1.4]
\label{IV.4.2.1.4}
Bajo las hipótesis de {\EGAReadLink{IV.4.2.1}{(4.2.1)}}, se tiene
además
\begin{equation}
 \dim(K\otimes_kL)=
 \inf(\operatorname{deg.tr.}_kK,\operatorname{deg.tr.}_kL).
 \label{IV.4.2.1.5}\tag{4.2.1.5}
\end{equation}
donde debe entenderse que, si los dos cuerpos $K$, $L$ son de grado de
trascendencia infinito (que pueden ser cardinales arbitrarios), el
primer miembro debe tomarse igual a $+\infty$ (la dimensión considerada
aquí es la definida en ($0$, 16.1.1)).

Supongamos, en efecto, por ejemplo, que
$\operatorname{deg.tr.}_kK\leq\operatorname{deg.tr.}_kL$ y, en primer
lugar, que $\operatorname{deg.tr.}_kK=n<+\infty$. Se ha visto en la demostración
de {\EGAReadLink{IV.4.2.1}{(4.2.1)}} que $K\otimes_kL$ es una
extensión algebraica de $K'\otimes_kL$ y tiene, por tanto, la misma dimensión
($0$, 16.1.5). Podemos, pues, limitarnos al caso en que
$K=k(T_1,\ldots,T_n)$\footnote{La lista de erratas impresa y el inglés
mantenido usan corchetes en esta definición; la lectura corregida usa
paréntesis porque $K$ es el cuerpo de funciones racionales.}; entonces
$K\otimes_kL$ es un anillo de fracciones
de
\[
 k[T_1,\ldots,T_n]\otimes_kL=L[T_1,\ldots,T_n],
\]
que es de dimensión $n$ {\EGAReadLink{IV.5.5.4}{(5.5.4)}}, y
por consiguiente $\dim(K\otimes_kL)\leq n$ ($0$, 16.1.3.1). Para demostrar que
$\dim(K\otimes_kL)=n$, basta evidentemente mostrar que existe en
$\operatorname{Spec}(L[T_1,\ldots,T_n])$ un punto racional sobre $L$
que pertenece a $\operatorname{Spec}(K\otimes_kL)$
{\EGAReadLink{IV.5.2.3}{(5.2.3)}}, o, equivalentemente (I, 1.7.4), que
existe un $L$-homomorfismo $u:K\otimes_kL\to L$. Pero, por definición
del producto tensorial, esto equivale a demostrar la existencia de un
$k$-homomorfismo $v:K\to L$, lo cual es evidente en virtud de
la hipótesis
$\operatorname{deg.tr.}_kK\leq\operatorname{deg.tr.}_kL$.

Si $\operatorname{deg.tr.}_kK\leq\operatorname{deg.tr.}_kL$ y estos
dos cardinales son infinitos, se ve del mismo modo que existe un punto
racional en $\operatorname{Spec}(L[(t_\alpha)_{\alpha\in I}])$ que
pertenece a $\operatorname{Spec}(K\otimes_kL)$; pero el núcleo de todo
$L$-homomorfismo de $L[(t_\alpha)_{\alpha\in I}]$ en $L$ es un
ideal maximal engendrado por una familia de la forma
$(t_\alpha-a_\alpha)_{\alpha\in I}$ con $a_\alpha\in L$. Está claro
que tal ideal maximal contiene una sucesión infinita decreciente
de ideales primos distintos, luego el anillo local correspondiente de
$L[(t_\alpha)_{\alpha\in I}]$ no es noetheriano y tiene
dimensión infinita; como es un anillo local de $K\otimes_kL$, se
ve que este último es no noetheriano y de dimensión infinita.
\end{remark}

\oldpage[IV]{350}
\paragraph{(Err\textsubscript{IV}, 20).}
Después de (IV, 4.5.19), insertar, antes de la línea 12 desde abajo de la p.~67:

\begin{corollary}[4.5.19.3]
\label{IV.4.5.19.3}
Sean $k$ un cuerpo y $X$ un $k$-preesquema de tipo finito sobre $k$, que sea
irreducible pero no geométricamente irreducible. Entonces existe una
parte cerrada $Y$ de $X$, rara en $X$, que contiene todos los puntos
$x\in X$ tales que $k(x)$ sea radicial sobre $k$ (y, en particular, todos
los puntos de $X$ racionales sobre $k$).
\end{corollary}

Sean $k'$ una clausura algebraica de $k$, $X'=X\otimes_kk'$,
$p:X'\to X$ la proyección canónica. Si $x\in X$ es tal que $k(x)$
sea radicial sobre $k$, solo existe un punto $x'\in X'$ sobre
$x$ (I, 3.4.9); este punto pertenece necesariamente a dos componentes
irreducibles distintas de $X'$, pues de otro modo $X$ sería geométricamente
irreducible en virtud de
{\EGAReadLink{IV.4.5.19}{(4.5.19)}}, en contra de
la hipótesis. Sean entonces $X'_i$ ($1\leq i\leq n$) las componentes
irreducibles del preesquema noetheriano $X'$ y, para
$1\leq i<j\leq n$, pongamos $X'_{ij}=X'_i\cap X'_j$; el conjunto
\[
 Y=p\left(\bigcup_{i<j}X'_{ij}\right)
\]
responde a la cuestión. En efecto, como $p$ es un morfismo entero y
los $X'_{ij}$ son cerrados en $X'$, $Y$ es cerrado en $X$.
Por otra parte, si $\xi$ es el punto genérico de $X$, $p^{-1}(\xi)$
es discreto y está formado por los puntos genéricos de los $X'_i$
{\EGAReadLink{IV.4.5.11}{(4.5.11)}}, por lo que no corta ninguno
de los $X'_{ij}$; por consiguiente, el conjunto cerrado $Y$ es raro.

\paragraph{(Err\textsubscript{IV}, 21).}
En (IV, 5.6.11.1), línea 7 desde abajo de la p.~102, después de « intègre »,
añadir « de dimension 2 ».

\paragraph{(Err\textsubscript{IV}, 22).}
Después de (IV, 6.2.3), después de la línea 18 de la p.~138, insertar:

\begin{proposition}[6.2.4]
\label{IV.6.2.4}
Bajo las hipótesis de {\EGAReadLink{IV.6.2.3}{(6.2.3)}}, sea
además $M$ un $A$-módulo de tipo finito y supongamos $M\neq0$ y
$N\neq0$. Entonces se tiene
\begin{equation}
 \operatorname{dim.proj}_B(M\otimes_AN)=
 \operatorname{dim.proj}_A(M)+
 \operatorname{dim.proj}_{B\otimes_Ak}(N\otimes_Ak).
 \label{IV.6.2.4.1}\tag{6.2.4.1}
\end{equation}
\end{proposition}

Consideremos la resolución por la izquierda $L_\bullet$ de $N$ introducida en
{\EGAReadLink{IV.6.2.3}{(6.2.3)}} y una resolución por la izquierda
$K_\bullet$ de $M$ por $A$-módulos libres de tipo finito:
\[
 \cdots\longrightarrow K_i\longrightarrow K_{i-1}\longrightarrow
 \cdots\longrightarrow K_0\longrightarrow M\longrightarrow0.
\]
Se ha visto que los $L_i$, $N$, los $Z_i(L_\bullet)$ son $A$-módulos
planos; lo mismo ocurre con los $B_i(L_\bullet)$, que son iguales a los
$Z_i(L_\bullet)$ salvo para $i=0$, y en este último caso
$B_0(L_\bullet)$ es

\oldpage[IV]{351}
plano por ($0_{\mathrm I}$, 6.1.2); por último, los $H_i(L_\bullet)$ son nulos
por definición salvo $H_0(L_\bullet)=N$, que es un $A$-módulo plano. Se
sabe que, en estas condiciones,
\[
 H_n(K_\bullet\otimes_AL_\bullet)\simeq
 \bigoplus_{p+q=n}
 H_p(K_\bullet)\otimes_AH_q(L_\bullet)
\]
(G, I, 5.5), luego $K_\bullet\otimes_AL_\bullet$ (graduado como de
costumbre) es una resolución de $M\otimes_AN$ por $B$-módulos
libres. Pongamos
\[
 r=\operatorname{dim.proj}_A(M),\qquad
 s=\operatorname{dim.proj}_B(N)=
 \operatorname{dim.proj}_{B\otimes_Ak}(N\otimes_Ak).
\]
Existe, por tanto, una resolución libre $K_\bullet$ de longitud $r$ y una
resolución libre $L_\bullet$ de longitud $s$, y
$K_\bullet\otimes_AL_\bullet$ es entonces una resolución libre de longitud
$r+s$, lo cual demuestra ya que
$\operatorname{dim.proj}_B(M\otimes_AN)\leq r+s$. Queda por demostrar,
según ($0$, 17.2.6), que se tiene
\[
 \operatorname{Tor}^{B}_{r+s}(M\otimes_AN,k(B))\neq0.
\]
Utilizando la resolución precedente $K_\bullet\otimes_AL_\bullet$ para
calcular este funtor, se obtiene
\[
 \operatorname{Tor}^{B}_{r+s}(M\otimes_AN,k(B))=
 H_{r+s}((K_\bullet\otimes_AL_\bullet)\otimes_Bk(B)).
\]
Pero puede escribirse
\[
 (K_\bullet\otimes_AL_\bullet)\otimes_Bk(B)=
 (K_\bullet\otimes_Ak)\otimes_k
 (L_\bullet\otimes_Bk(B)),
\]
que es, por tanto, un producto tensorial de dos complejos de espacios
vectoriales sobre $k$; como todos los $k$-módulos son planos, la fórmula de
Künneth da
\[
 H_{r+s}((K_\bullet\otimes_Ak)\otimes_k
 (L_\bullet\otimes_Bk(B)))\simeq
 H_r(K_\bullet\otimes_Ak)\otimes_k
 H_s(L_\bullet\otimes_Bk(B)),
\]
siendo nula la homología del primer (resp. segundo) complejo en grados
$>r$ (resp. $>s$). Se tiene, por tanto,
\[
 \operatorname{Tor}^{B}_{r+s}(M\otimes_AN,k(B))=
 \operatorname{Tor}^{A}_{r}(M,k)\otimes_k
 \operatorname{Tor}^{B}_{s}(N,k(B))
\]
salvo isomorfismo, y como cada uno de los factores es $\neq0$ por
hipótesis, lo mismo ocurre con su producto tensorial sobre $k$.

\paragraph{(Err\textsubscript{IV}, 23).}
En (IV, 6.10.1), línea 8 de la p.~155, después de canonique, añadir:
$\sh{F}$ un $\sh{O}_X$-Module cohérent.

\paragraph{(Err\textsubscript{IV}, 24).}
En ({\EGAReadLink{IV.6.14.4}{IV, 6.14.4}}), la referencia
dada en la línea 12 desde abajo de la p.~175 es insuficiente; hay que razonar
del modo siguiente. En primer lugar, si $M$ es una base de trascendencia
de $K'$ sobre $A$, $A(M)$ es algebraicamente cerrado en $B(A(M))$
(Bourbaki, \emph{Alg.}, cap.~V, \S6, ejerc.~8 c)). Podemos, por tanto,
limitarnos al caso en que $K'$ sea una extensión algebraica separable de $A$.
Sea $p$ el exponente característico de $A$; se deduce de Bourbaki,
\emph{loc. cit.}, \S9, ejerc.~2, que la clausura algebraica de $K'$
en $B(K')$ es radicial sobre $K'$; demostremos que todo elemento
$x\in B(K')$ tal que $x^{p^r}\in K'$ pertenece ya a $K'$. Para ello,
consideremos una base $(e_\lambda)$ de $K'$ sobre $A$; puesto que $K'$ es una
extensión algebraica separable de $A$, los $e_\lambda^{p^r}$ forman
también una base de $K'$ sobre $A$ (Bourbaki, \emph{loc. cit.}, \S8,
n\textsuperscript{o}~3, cor. de la prop.~4). Se tiene
$x=\sum_\lambda e_\lambda b_\lambda$, con $b_\lambda\in B$ y, por
hipótesis, existen elementos $a_\lambda\in A$ tales que
\[
 \sum_\lambda e_\lambda^{p^r}b_\lambda^{p^r}=
 \sum_\lambda e_\lambda^{p^r}a_\lambda;
\]
pero esto implica $b_\lambda^{p^r}=a_\lambda$ para todo $\lambda$;
como por hipótesis $A$ es algebraicamente cerrado en $B$, se tiene, por tanto,
$b_\lambda\in A$ para todo $\lambda$ y, por consiguiente, $x\in K'$.

\paragraph{(Err\textsubscript{IV}, 25).}
En (IV, 7.4.8, B)), p.~203, la cuestión ha sido resuelta afirmativamente
por R.~Kiehl, \emph{Ausgezeichnete Ringe der nichtarchimedischen
analytischen Geometrie}, de próxima aparición en \emph{Inventiones
Mathematicae}.

\paragraph{(Err\textsubscript{IV}, 26).}
En (IV, 7.8.4), línea 9 desde abajo de la p.~217, sustituir:
\begin{quote}
Veremos más adelante (\S18) que si $k$ es un cuerpo valuado completo no
discreto\ldots
\end{quote}
por:
\begin{quote}
R.~Kiehl ha demostrado (en el artículo citado en
(Err\textsubscript{IV}, 25)) que si $k$ es un cuerpo valuado completo no
discreto\ldots
\end{quote}

\oldpage[IV]{352}
\paragraph{(Err\textsubscript{IV}, 27).}
En (IV, 7.8.4), añadir, antes de la línea 6 desde abajo de la p.~217:

\smallskip
\noindent
\textup{(vi)} Sea $A$ un anillo local noetheriano íntegro de dimensión
1; es entonces universalmente catenario
{\EGAReadLink{IV.7.2.9}{(7.2.9)}} y, para que sea excelente,
es necesario y suficiente, en virtud de
{\EGAReadLink{IV.7.3.19}{(7.3.19)}}, que sus fibras
formales sean geométricamente reducidas, o también que su completado
$\widehat A$ sea reducido y que los cuerpos componentes de su anillo total
de fracciones sean extensiones separables del cuerpo de fracciones
$K$ de $A$. En virtud de {\EGAReadLink{IV.7.6.4}{(7.6.4)}} y
{\EGAReadLink{IV.7.7.2}{(7.7.2)}}, esto equivale asimismo a
decir que $A$ es universalmente japonés. Cuando $A$ es un anillo de
valoración discreta, lo mismo ocurre con $\widehat A$; decir que $A$ es
excelente significa entonces que el cuerpo de fracciones de $\widehat A$
es una extensión separable de $K$, condición siempre satisfecha si
$K$ es de característica 0.

\paragraph{(Err\textsubscript{IV}, 28).}
Después de (IV, 8.10.3), después de la línea 14 de la p.~37, insertar:

\begin{remark}[8.10.3.1]
\label{IV.8.10.3.1}
En {\EGAReadLink{IV.8.10.3}{(8.10.3)}} no puede suprimirse
la hipótesis de que los $\overline{f_\alpha(X_\alpha)}$ sean
constructibles. Tomemos, por ejemplo, $S=\operatorname{Spec}(A)$, donde $A$
es un anillo tal que $S$ sea compacto, totalmente disconexo y sin
punto aislado, siendo maximal todo ideal primo $\mathfrak p$ de $A$ y
$A_\mathfrak p$ un cuerpo (Bourbaki, \emph{Alg. comm.}, cap.~II,
\S4, ejerc.~16 y 17). Sea $s\in S$, y tomemos como
$Y_\alpha=S_\alpha$ los entornos abiertos afines de $s$ en $S$; se
toma $X_\alpha=\operatorname{Spec}(k(s))$ para todo $\alpha$,
siendo $f_\alpha$ y $f$ los morfismos canónicos; está claro que
$f$, que es la identidad, es dominante, pero ninguno de los $f_\alpha$ lo
es.
\end{remark}

\paragraph{(Err\textsubscript{IV}, 29).}
En (IV, 8.12.8), línea 7 de la p.~47, sustituir « tel qu'il existe un
$\sh{O}_Y$-Module ample » por « et quasi-séparé ». Línea 12 de la
p.~47, sustituir « L'hypothèse » por « Il résulte de
{\EGAReadLink{IV.8.12.3}{(8.12.3)}} que l'existence de la
factorisation de l'énoncé est une propriété locale sur $Y$; on peut
donc supposer $Y$ affine, ce qui ».

\paragraph{(Err\textsubscript{IV}, 30).}
Después de (IV, 8.12.11), añadir al final de la p.~48:

\begin{corollary}[8.12.12]
\label{IV.8.12.12}
Sean $X$, $Y$ dos preesquemas íntegros,
$f:X\to Y$ un morfismo separado, birracional y localmente de tipo finito.
Sea $y\in Y$ un punto normal de $Y$ y supongamos que un punto
$x\in f^{-1}(y)$ esté aislado en $f^{-1}(y)$. Entonces existe un
entorno abierto $U$ de $y$ en $Y$ tal que la restricción
$f^{-1}(U)\to U$ de $f$ sea un isomorfismo.
\end{corollary}

\begin{proof}
Basta demostrar que existe un entorno abierto $U$ de $y$ en
$Y$ y una $U$-sección $g:U\to f^{-1}(U)$ de $f^{-1}(U)$ tal que
$g(y)=x$. En efecto, como $U$ contiene el punto genérico $\eta$ de
$Y$ y $f^{-1}(U)$ el punto genérico $\xi$ de $X$, se tiene
necesariamente $g(\eta)=\xi$ puesto que $f$ es birracional
{\EGAReadLink{IV.6.15.4}{(6.15.4)}}; como $g$ es una
inmersión, $g(U)$ es localmente cerrado en $f^{-1}(U)$, luego abierto
en su adherencia en $f^{-1}(U)$, y como esta última es igual
a $f^{-1}(U)$ por lo anterior, $g(U)$ es abierto en
$f^{-1}(U)$; como el subpreesquema de $f^{-1}(U)$ asociado a $g$ es
isomorfo a $U$, luego reducido, y $f^{-1}(U)$ es reducido, se deduce
de (I, 5.2.1) que $g$ es una inmersión abierta. Por otra parte, puesto que
$f$ es separado, $g$ es una inmersión cerrada (I, 5.4.6), luego $g(U)$
es un conjunto a la vez abierto y cerrado en $f^{-1}(U)$. Pero, como
$f^{-1}(U)$ es íntegro, esto implica $g(U)=f^{-1}(U)$, de donde resulta la
conclusión.

Consideremos entonces dos casos:

\textup{I)} Supongamos primero que $Y$ es normal, y sea $V$ el conjunto de los
puntos de $X$ en los cuales $f$ es cuasifinito; es un conjunto abierto
en $X$ {\EGAReadLink{IV.13.1.4}{(13.1.4)}}.\footnote{El
lector comprobará que ni {\EGAReadLink{IV.8.12.12}{(8.12.12)}}
ni ninguna de sus consecuencias se utiliza en la demostración de
{\EGAReadLink{IV.13.1.4}{(13.1.4)}}.} Se deduce entonces de
{\EGAReadLink{IV.8.12.10}{(8.12.10)}}

\oldpage[IV]{353}
que la restricción $f|V:V\to Y$ es una inmersión abierta; existe, por tanto,
un entorno abierto $U$ de $y\in Y$ tal que $f|V$ sea un isomorfismo
de $V$ sobre el subpreesquema inducido por $Y$ sobre $U$; el isomorfismo
recíproco $g$ es entonces una $U$-sección de $f^{-1}(U)$ tal que
$g(y)=x$.

\textup{II)} Como la obtención de una $U$-sección $g$ tal que $g(y)=x$ es
un problema local sobre $X$ y sobre $Y$, podemos suponer que $X$ e $Y$
son afines y que $f$ es de tipo finito; por tanto, $X$ es un
subpreesquema cerrado de un $Y$-preesquema $Z$ de presentación finita sobre
$Y$ (por ejemplo, de la forma $Y[T_1,\ldots,T_n]$). Pongamos
\[
 Y'=\operatorname{Spec}(\sh{O}_{Y,y}),\qquad
 X'=X_{(Y')},\qquad Z'=Z_{(Y')},
\]
de modo que $X'$ es un subpreesquema cerrado de $Z'$. Además, si
$f'=f_{(Y')}$, se tiene $f'^{-1}(y)=f^{-1}(y)$ (I, 3.6.5) y $x$ está, por tanto,
aislado en $f'^{-1}(y)$. Puesto que $Y'$ es normal, puede aplicarse el
resultado de I) a $f'$ y, puesto que $Y'$ es el único entorno de $y$ en
$Y'$, existe una $Y'$-sección $g'$ de $X'$ tal que $g'(y)=x$, que
evidentemente puede considerarse también como una $Y'$-sección de $Z'$.
Puesto que el morfismo estructural $h:Z\to Y$ es de presentación finita,
puede aplicarse ({\EGAReadLink{IV.8.8.2}{8.8.2}}, (i) y
(ii)) siguiendo el método de
({\EGAReadLink{IV.8.1.2}{8.1.2}}, a)), de donde se concluye
la existencia de un entorno abierto $U$ de $y$ en $Y$ y de una
$U$-sección $g$ de $h^{-1}(U)$ tal que $g(y)=x$. Queda por ver
que, de hecho, $g$ es una $U$-sección de $f^{-1}(U)$. Pero, puesto que
$f^{-1}(U)$ es un subpreesquema cerrado de $h^{-1}(U)$ y $U$ es
reducido, basta para ello (I, 5.2.2) comprobar que
$g(U)\subset f^{-1}(U)$. Ahora bien, puesto que $U$ es íntegro y $f^{-1}(U)$
es cerrado en $h^{-1}(U)$, basta demostrar que, si $\eta$ es el punto
genérico de $Y$ (y de $U$), se tiene $g(\eta)\in f^{-1}(U)$. Esto
resulta de que la restricción de $g$ a $Y'$ es una $Y'$-sección de
$X'\subset f^{-1}(U)$ y de que $\eta\in Y'$.
\end{proof}

\paragraph{(Err\textsubscript{IV}, 31).}
En ({\EGAReadLink{IV.10.4}{IV, 10.4}}), añadir, después de la
proposición {\EGAReadLink{IV.10.4.11}{(10.4.11)}}:

\begin{corollary}[10.4.12]
\label{IV.10.4.12}
Sean $S$ un preesquema, $X$ un $S$-preesquema de presentación finita y
$f$ un $S$-endomorfismo radicial de $X$.
\begin{enumerate}[label=\textup{(\roman*)}]
\item El endomorfismo $f$ es biyectivo y finito (por tanto, un homeomorfismo
universal {\EGAReadLink{IV.2.4.5}{(2.4.5)}}).
\item Supongamos además que $X$ sea reducido, que los anillos locales de
$S$ sean noetherianos y universalmente japoneses, y que $f$ sea
birracional (condición automáticamente satisfecha si, para todo punto
maximal $x$ de $X$, $k(x)$ es un cuerpo perfecto). Entonces $f$ es un
$S$-automorfismo.
\end{enumerate}
\end{corollary}

\begin{proof}
\textup{(i)} Como la cuestión es local sobre $S$, podemos limitarnos al caso
en que $S=\operatorname{Spec}(A)$ es afín. Existe entonces una
sub-$\mathbb Z$-álgebra de tipo finito $A_0$ de $A$ tal que, poniendo
$S_0=\operatorname{Spec}(A_0)$, se tenga
$X=X_0\times_{S_0}S$, donde $X_0$ es un preesquema de tipo finito sobre $S_0$,
y que $f$ se deduzca de un $S_0$-endomorfismo $f_0$ de $X_0$ mediante el
cambio de base $S\to S_0$
{\EGAReadLink{IV.8.9.1}{(8.9.1)}}; además, puede suponerse
$f_0$ radicial ({\EGAReadLink{IV.8.10.5}{8.10.5}}, (vii)).
Bastará demostrar que $f_0$ es biyectivo y finito, y podemos, por tanto,
suponer en lo sucesivo que $A$ es una $\mathbb Z$-álgebra de tipo finito, lo
cual implica que $S$ y $X$ son preesquemas excelentes
({\EGAReadLink{IV.7.8.3}{(7.8.3)}} y
{\EGAReadLink{IV.7.8.6}{(7.8.6)}}) y, en particular, que $X$
es un preesquema universalmente japonés
({\EGAReadLink{IV.5.11.4}{(5.11.4)}} y
({\EGAReadLink{IV.7.8.3}{7.8.3}}, (vi))). Además, puede
suponerse $X$ reducido; en efecto,
$f_{\mathrm{red}}:X_{\mathrm{red}}\to X_{\mathrm{red}}$ es un
$S_{\mathrm{red}}$-endomorfismo radicial (I, 5.1.6); si se demuestra que
$f_{\mathrm{red}}$ es finito, lo mismo ocurrirá con $f$: en efecto,
puesto que $f$ es radicial, es separado
{\EGAReadLink{IV.1.8.7.1}{(1.8.7.1)}}; como
$f_{\mathrm{red}}$ es propio, lo mismo ocurre con $f$ (II, 5.4.6);
pero, puesto que $f$ es evidentemente cuasifinito, es finito (III, 4.4.2).

\oldpage[IV]{354}
Sea entonces $X'$ la normalización de $X$ (II, 6.3.8); si $X_i$
($1\leq i\leq n$) son los subpreesquemas reducidos de $X$ cuyos
espacios subyacentes son las componentes irreducibles de $X$, $X'$ es
la suma de las normalizaciones $X'_i$ de los preesquemas íntegros $X_i$;
denotaremos por $x_i$ (resp. $x'_i$) el punto genérico de $X_i$
(resp. $X'_i$). El morfismo canónico $p:X'\to X$ es finito puesto que
$X$ es universalmente japonés y se tiene
$p^{-1}(x_i)=\{x'_i\}$, siendo biyectivo el homomorfismo
$k(x_i)\to k(x'_i)$.

Sea $f=(\psi,\theta)$; ya se sabe
{\EGAReadLink{IV.10.4.11}{(10.4.11)}} que $\psi$ es una
biyección continua del espacio $X$ sobre sí mismo; si se pone
$y_i=\psi(x_i)$, $y_i$ es también un punto maximal de $X$, pues si $y'_i$
fuese una generalización de $y_i$ distinta de $y_i$, $x_i$ sería una
generalización $\neq x_i$ de $\psi^{-1}(y'_i)$. Puesto que el conjunto $M$
de los puntos maximales de $X$ es finito, la restricción $\psi|M$ es una
biyección de $M$ sobre sí mismo, y puede escribirse
$\psi(x_i)=x_{\sigma(i)}$, donde $\sigma$ es una permutación de
$[1,n]$. Entonces, para todo $i$, la restricción $f|X_i$ de $f$ se
factoriza como
\[
 X_i\xrightarrow{f_i}X_{\sigma(i)}\longrightarrow X
\]
(I, 5.2.2), donde la segunda flecha es la inyección canónica. Como los
$X_i$ son íntegros, existe (II, 6.3.9) un único morfismo
$f'_i:X'_i\to X'_{\sigma(i)}$ que hace conmutativo el diagrama
\[
\begin{array}{ccc}
X'_i&\xrightarrow{f'_i}&X'_{\sigma(i)}\\
\downarrow p_i&&\downarrow p_{\sigma(i)}\\
X_i&\xrightarrow{f_i}&X_{\sigma(i)}
\end{array}
\]
y si $f'$ coincide con $f'_i$ en cada uno de los $X'_i$, se tiene, por tanto, un
diagrama conmutativo
\phantomsection\label{IV.10.4.12.1}
\begin{equation}
\label{IV.10.4.12.1.errata}
\begin{array}{ccc}
X'&\xrightarrow{f'}&X'\\
\downarrow p&&\downarrow p\\
X&\xrightarrow{f}&X.
\end{array}
\tag{10.4.12.1}
\end{equation}

Por otra parte, como $f$ es separado, lo mismo ocurre con
$f\circ p=p\circ f'$, luego $f'$ es separado (I, 5.5.1); por último, puesto que
$p$ es finito y $f$ biyectivo, se deduce de
{\EGAReadLink{IV.10.4.12.1}{(10.4.12.1)}} que $f'$ es
cuasifinito. Puesto que los $X'_i$ son íntegros y normales, se deduce de
{\EGAReadLink{IV.8.12.11}{(8.12.11)}} que $f'_i$ se factoriza como
\[
 X'_i\xrightarrow{f''_i}X''_{\sigma(i)}
 \xrightarrow{u_{\sigma(i)}}X'_{\sigma(i)},
\]
donde $X''_{\sigma(i)}$ es la clausura integral de
$X'_{\sigma(i)}$ relativa a la extensión $k(x'_i)$ de
$k(x'_{\sigma(i)})$, y $f''_i$ una inmersión abierta. Pero, como
$k(x'_i)$ es una extensión radicial de $k(x'_{\sigma(i)})$,
$u_{\sigma(i)}$ es un morfismo radicial (Bourbaki,
\emph{Alg. comm.}, cap.~V, \S2, n\textsuperscript{o}~3, lema 4);
luego $f'$ es radicial, y se deduce entonces de
{\EGAReadLink{IV.10.4.11}{(10.4.11)}} que $f'$ es
sobreyectivo; se concluye (puesto que los $u_{\sigma(i)}$ son inyectivos)
que $f''_i$ es biyectivo, luego un isomorfismo; por último, se deduce que
el morfismo $f'$ es finito (puesto que lo es $u_{\sigma(i)}$,
siendo $X'$ universalmente japonés
{\EGAReadLink{IV.7.8.3}{(7.8.3)}}). Ahora bien, como
$p\circ f'=f\circ p$ es entonces finito, $p$ es sobreyectivo y $f$
separado, $f$ es propio (II, 5.4.3) y cuasifinito, luego finito
{\EGAReadLink{IV.8.11.1}{(8.11.1)}}.

\oldpage[IV]{355}
\textup{(ii)} Teniendo en cuenta
({\EGAReadLink{IV.8.10.5}{8.10.5}}, (i)), aplicado siguiendo el
método de ({\EGAReadLink{IV.8.1.2}{8.1.2}}, a)), podemos
limitarnos al caso en que $S$ sea un esquema local. Las reducciones previas
realizadas en la demostración de (i) son, por tanto, innecesarias aquí, y con las mismas
notaciones se ve que se tiene un diagrama conmutativo
{\EGAReadLink{IV.10.4.12.1}{(10.4.12.1)}} en el cual $X$ es
un $S$-preesquema de tipo finito, $p$ es finito, $f'$ es un isomorfismo y
el homomorfismo canónico $\sh{O}_X\to p_*(\sh{O}_{X'})$ es inyectivo.
Pero se demostrará en el cap.~VI, en la teoría de los cocientes de
preesquemas por relaciones de equivalencia, que, bajo estas condiciones, $f$
es él mismo un isomorfismo (el lector comprobará que no se utilizará
{\EGAReadLink{IV.10.4.12}{(10.4.12)}} en la teoría de las
relaciones de equivalencia del cap.~VI).
\end{proof}

\begin{corollary}[10.4.13]
\label{IV.10.4.13}
Sean $k$ un cuerpo y $X$ un preesquema algebraico sobre $k$. Entonces todo
$k$-endomorfismo radicial $f$ de $X$ es finito y biyectivo, luego un
homeomorfismo universal. Si además $X$ es reducido y $f$ birracional
(por ejemplo, si $k$ es de característica 0), $f$ es un
$k$-automorfismo.
\end{corollary}

\begin{remark}[10.4.14]
\label{IV.10.4.14}
\begin{enumerate}[label=\textup{(\roman*)}]
\item En {\EGAReadLink{IV.10.4.12}{(10.4.12)}} y
{\EGAReadLink{IV.10.4.13}{(10.4.13)}} no puede suprimirse
la hipótesis de que $X$ es reducido para la segunda conclusión; así lo
muestra el ejemplo del $B$-endomorfismo correspondiente al aumento
$D_B(L)\to B$ del anillo $D_B(L)=A$ definido en ($0$, 18.2.3), donde
$B$ se identifica con un subanillo de $A$.

\item Sea $X$ un preesquema noetheriano. Se ignora si todo
endomorfismo radicial $f$ de $X$ es necesariamente sobreyectivo y finito.
Si se supone que $f$ es sobreyectivo y, además, $X$ es un preesquema
japonés, el razonamiento de
{\EGAReadLink{IV.10.4.12}{(10.4.12)}} muestra que $f$ es finito;
si además $X$ es reducido y $f$ es birracional, por ejemplo si,
para todo punto maximal $x\in X$, $k(x)$ es perfecto, el mismo
razonamiento muestra que $f$ es un automorfismo.

\item Pueden darse ejemplos de endomorfismos radiciales no
sobreyectivos de esquemas afines íntegros no noetherianos: si $k$ es un
anillo, $A=k[(T_\lambda)_{\lambda\in L}]$ un anillo de polinomios donde $L$
es infinito numerable, y $\mathfrak J$ el ideal engendrado en $A$ por
una subfamilia $(T_\lambda)_{\lambda\in H}$, donde $L-H$ es infinito y
$H\neq\varnothing$, se observa que $A/\mathfrak J$ es isomorfo a
$k[(T_\lambda)_{\lambda\in L-H}]$, luego a $A$; el endomorfismo de
$\operatorname{Spec}(A)$ correspondiente al endomorfismo de anillos
compuesto
\[
 A\xrightarrow{p}A/\mathfrak J\xrightarrow{u}A
\]
donde $p$ es la proyección canónica y $u$ un isomorfismo, responde a la
cuestión, siendo su imagen el subpreesquema cerrado de
$\operatorname{Spec}(A)$ definido por $\mathfrak J$.
\end{enumerate}
\end{remark}

\paragraph{(Err\textsubscript{IV}, 32).}
En (IV, 11.3.14), línea 7 de la p.~141, antes de « Soient », insertar
(i). Después de la línea 9 de la p.~141, añadir:

\smallskip
\noindent
\textup{(ii)} Recíprocamente, sea $f:X\to Y$ un morfismo plano. Si $X$
es unirrama (resp. geométricamente unirrama) en un punto $x$, $Y$
es unirrama (resp. geométricamente unirrama) en el punto $y=f(x)$.

Al comienzo de la línea 10 de la p.~141, insertar (i). Después de la línea 20
de la p.~141, añadir:

\smallskip
\noindent
\textup{(ii)} Observemos que, si $Y_1=Y_{\mathrm{red}}$,
$X_1=X\times_Y Y_1$ tiene como espacio subyacente el de $X$ y el
morfismo $f_1:X_1\to Y_1$ deducido de $f$ por cambio de base es
plano. Podemos, por tanto, limitarnos al caso en que $Y=\operatorname{Spec}(A)$,
siendo $A$ un anillo local reducido y $y$ el punto cerrado de $Y$. Pongamos
$B=\sh{O}_{X,x}$; por hipótesis $B_{\mathrm{red}}$ es unirrama, luego
$\operatorname{Spec}(B)$ tiene un único punto maximal, y por planitud
{\EGAReadLink{IV.2.3.4}{(2.3.4)}} se deduce que lo mismo ocurre

\oldpage[IV]{356}
con $\operatorname{Spec}(A)$; puesto que $A$ es reducido, es, por tanto,
íntegro. Sean $K$ el cuerpo de fracciones de $A$ y $A'$ la clausura
integral de $A$; por planitud, $B\otimes_AA'$ se identifica con un
subanillo de $B\otimes_AK$ que contiene $B$, y como los elementos no
nulos de $A$ son $B$-regulares por planitud, $B\otimes_AK$ se identifica
con un subanillo del anillo total de fracciones $R$ de $B$. Para
demostrar que $A$ es unirrama, basta ver que $A'$ no puede tener
dos ideales maximales distintos (necesariamente sobre $y$
(Bourbaki, \emph{Alg. comm.}, cap.~V, \S2,
n\textsuperscript{o}~1, prop.~1)); si así fuese,
$B'=B\otimes_AA'$ tendría dos ideales primos distintos sobre
$x$ (I, 3.4.7), y basta ver que esto no es posible.

Sea $\mathfrak N$ el nilradical de $B$; si $S$ es el conjunto de los
elementos regulares de $B$, se tiene $R=S^{-1}B$, y
$S^{-1}\mathfrak N=R\mathfrak N$ es el nilradical de $R$, luego
\[
 R/R\mathfrak N=S^{-1}B/S^{-1}\mathfrak N=S^{-1}(B/\mathfrak N)
\]
está contenido en el cuerpo de fracciones $L$ del anillo local íntegro
$B/\mathfrak N=B_{\mathrm{red}}$; lo mismo ocurre, por tanto, con
\[
 B'_{\mathrm{red}}=B'/(B'\cap R\mathfrak N)
\]
que es entero sobre $B_{\mathrm{red}}$; por hipótesis,
$B'_{\mathrm{red}}$, al estar contenido en la clausura integral de
$B_{\mathrm{red}}$, es un anillo local, luego lo mismo ocurre con
$B'$, lo cual establece nuestra afirmación.

Si se supone además que $B$ es geométricamente unirrama, el cuerpo
residual de $B'_{\mathrm{red}}$, al estar contenido en el cuerpo residual de
la clausura integral de $B_{\mathrm{red}}$, es una extensión radicial
del cuerpo residual $k(x)$ de $B$. Observemos ahora que el cuerpo residual
de $B'_{\mathrm{red}}$ es también el de $B'=B\otimes_AA'$, luego es
igual a $k(x)\otimes_{k(y)}k'$, donde $k'$ es el cuerpo residual de $A'$;
se deduce entonces de {\EGAReadLink{IV.2.6.1}{(2.6.1)}} que $k'$
es una extensión radicial de $k(y)$, luego $A$ es también
geométricamente unirrama.

\paragraph{(Err\textsubscript{IV}, 33).}
En ({\EGAReadLink{IV.12.1.2}{IV, 12.1.2}}), el ejemplo dado
en (i) no es correcto, pues el morfismo considerado es abierto. Se
construye del modo siguiente un ejemplo de un morfismo plano $f:X\to Y$ cuya restricción
a una componente irreducible de $X$ no es un morfismo abierto.
Sea $Z=\operatorname{Spec}(\mathbb C[U,V])$ el
plano afín sobre el cuerpo de los números complejos, y consideremos el
esquema $Y$ obtenido pegando dos puntos cerrados distintos $a$, $b$ de
$Z$; se tiene $Y=\operatorname{Spec}(B)$, donde $B$ es el subanillo de
$\mathbb C[U,V]$ formado por los polinomios que toman el mismo valor en los puntos
$a$ y $b$ (identificando el conjunto de los puntos cerrados de $Z$ con
$\mathbb C^2$); es inmediato que $B$ es un sub-$\mathbb C$-espacio
vectorial de $\mathbb C[U,V]$ de codimensión 1, luego $\mathbb C[U,V]$
es un $B$-módulo de tipo finito y, por consiguiente, el morfismo $Z\to Y$ es
finito, siendo la imagen inversa del punto $c$ de $Y$ obtenido pegando
$a$ y $b$ el conjunto $\{a,b\}$.

Consideremos entonces el esquema $X$ obtenido pegando dos ejemplares de
$Z$, sean $Z_1$, $Z_2$, de modo que, si $a_i$, $b_i$ ($i=1,2$) son
los puntos de $Z_i$ correspondientes a $a$ y $b$, $a_1$ se pegue con
$b_2$ y $b_1$ con $a_2$ (cf. ({\EGAReadLink{IV.6.5.5}{6.5.5}},
(ii))). Se ve fácilmente (\emph{loc. cit.}) que el morfismo $X\to Y$
así obtenido es plano y finito (e incluso étale) y que $Z_1$ y $Z_2$
se identifican con las dos componentes irreducibles de $X$. Sin embargo, si
$D$ es una recta en $Z_1$ que pasa por $a_1$ pero no por $b_1$,
la imagen $V$ en $Y$ del abierto $Z_1-D$ de $Z_1$ contiene la imagen de
$b_1$, es decir, el punto $c$, pero no contiene la imagen del punto
genérico de $D$, luego $V$ no es un entorno de $c$.

\paragraph{(Err\textsubscript{IV}, 34).}
En el título de (IV, 13.2), línea 5 desde abajo de la p.~190, sustituir
« équidimensionnels » por « localement équidimensionnels ». Del mismo modo,
en (IV, 13.2.2), línea 4 de la p.~191, sustituir dos veces
« équidimensionnel » por « localement équidimensionnel ».

\oldpage[IV]{357}
Línea 5 de la p.~191, añadir: \emph{Si además las fibras
$f^{-1}(f(x))$ de $f$ son equidimensionales ($0_{\mathrm{IV}}$,
14.1.3), se dice que $f$ es equidimensional.} Línea 6 desde abajo de la
p.~194, sustituir « équidimensionnelles » por « localement
équidimensionnelles ».

\paragraph{(Err\textsubscript{IV}, 35).}
En el título de (IV, 13.3), línea 4 desde abajo de la
p.~194\footnote{La lista de erratas impresa y el inglés mantenido dan
p.~94; el título de (IV, 13.3) se encuentra en la p.~194.}, sustituir
« équidimensionnels » por « localement équidimensionnels ». En
(IV, 13.3.2), línea 6 desde abajo de la p.~196, sustituir dos veces
« équidimensionnel » por « localement équidimensionnel ». Línea 5 desde abajo
de la p.~196, añadir: \emph{Si además las fibras $f^{-1}(f(x))$ de
$f$ son equidimensionales ($0_{\mathrm{IV}}$, 14.1.3), se dice que
$f$ es equidimensional.} En (IV, 13.3.7), línea 14 de la p.~198,
sustituir « équidimensionnel » por « localement équidimensionnel ».
En ({\EGAReadLink{IV.13.3.9}{IV, 13.3.9}}), línea 3 desde
abajo de la p.~198, sustituir « équidimensionalité » por
« équidimensionalité locale ». Líneas 10, 12, 20 y 24 de la p.~199,
sustituir « équidimensionnel » por « localement équidimensionnel ».

\paragraph{(Err\textsubscript{IV}, 36).}
En la introducción del \S14, línea 6 desde abajo de la p.~199, sustituir
« équidimensionnel » por « localement équidimensionnel ».

\paragraph{(Err\textsubscript{IV}, 37).}
En (IV, 14.4.1), línea 5 desde abajo de la p.~209, sustituir « localement
de type fini » por « localement de présentation finie ». El enunciado
actual de ({\EGAReadLink{IV.14.4.1}{IV, 14.4.1}}) es inexacto,
incluso cuando $f:X\to Y$ es una inmersión cerrada, $X$ está reducido a
un solo punto $y$ (cerrado) de $Y$ y $\sh{O}_{Y,y}$ es un cuerpo. En
efecto, existen anillos $A$ tales que $Y=\operatorname{Spec}(A)$ es
compacto, no tiene puntos aislados y es totalmente disconexo, todo ideal primo
$\mathfrak p$ de $A$ es maximal y $A_\mathfrak p$ es un cuerpo
(Bourbaki, \emph{Alg. comm.}, cap.~II, \S4, ejerc.~16 y 17); es
claro que la inmersión cerrada $\{y\}\to Y$ no es un morfismo
abierto, pero cumple las condiciones b) y b') de
({\EGAReadLink{IV.14.4.1}{IV, 14.4.1}}).

\paragraph{(Err\textsubscript{IV}, 38).}
En (IV, 14.4.1.3), línea 14 de la p.~211, sustituir « localement
quasi-fini et dominant » por « localement quasi-fini, localement de
présentation finie et dominant ».

\paragraph{(Err\textsubscript{IV}, 39).}
En (IV, 14.4.2), línea 10 desde abajo de la p.~211, sustituir « localement
de type fini » por « localement de présentation finie ».

\paragraph{(Err\textsubscript{IV}, 40).}
En (IV, 14.4.3), línea 18 de la p.~212, sustituir « localement de
type fini » por « localement de présentation finie ».

\paragraph{(Err\textsubscript{IV}, 41).}
En ({\EGAReadLink{IV.14.4.4}{IV, 14.4.4}}), líneas 15 y 16 desde
abajo de la p.~212, sustituir « localement de type fini » por
« localement de présentation finie ». Líneas 10 y 11 desde abajo, sustituir
« équidimensionnel » por « localement équidimensionnel ».

\paragraph{(Err\textsubscript{IV}, 42).}
En (IV, 14.4.8), después de la línea 1 de la p.~214, añadir:
\begin{quote}
\emph{Si además el conjunto de las componentes irreducibles de $Y$ es
localmente finito, estas dos condiciones son equivalentes a las siguientes:}
\end{quote}

\paragraph{(Err\textsubscript{IV}, 43).}
En (IV, 14.4.10), línea 15 de la p.~215, sustituir « quelconque » por
« dont l'ensemble des composantes irréductibles est localement fini ».
Línea 16 de la p.~215, sustituir « de type fini » por « de présentation
finie ».

\paragraph{(Err\textsubscript{IV}, 44).}
En (IV, 15.5.6.1), línea 15 desde abajo de la p.~234, después de la primera
frase, añadir:

\begin{quote}
Podemos limitarnos al caso en que toda parte cerrada irreducible de $X$
admite un único punto genérico. En efecto, en caso contrario consideramos
el espacio $Y=X/R$,

\oldpage[IV]{358}
cociente de $X$ por la relación de equivalencia
$\overline{\{x\}}=\overline{\{x'\}}$. De esta definición se sigue
inmediatamente que todo abierto de $X$ es saturado para $R$; por tanto, la
topología de $X$ es la imagen recíproca de la de $Y$ por la aplicación
canónica $\pi:X\to X/R$. Un abierto denso de $X$ es de la forma
$\pi^{-1}(U)$, donde $U$ es un abierto denso de $Y$; por consiguiente, una
parte cerrada rara $Z$ de $X$ es de la forma $\pi^{-1}(S)$, donde
$S=\pi(Z)$ es cerrada y rara en $Y$, y decir que $X-Z$ es conexo equivale
a decir que $Y-S$ es conexo. Del mismo modo, $\pi(T_x)$ es el conjunto de
las generalizaciones de $\pi(x)$, y decir que $T_x-\{x\}$ es conexo
equivale a decir que $\pi(T_x)-\{\pi(x)\}$ lo es. Finalmente, por
construcción, $Y$ es un espacio de Kolmogorov localmente noetheriano, por lo
que toda parte cerrada irreducible admite un único punto genérico.
Se ve entonces que demostrar el lema para $X$ equivale a demostrarlo para
$Y$, y podemos suponer que toda parte cerrada irreducible de $X$ tiene un
único punto genérico.
\end{quote}

\paragraph{(Err\textsubscript{IV}, 45).}
En (IV, 15.6.1), línea 18 desde abajo de la p.~236, suprimir « localement
noethérien » (hipótesis no utilizada en la demostración).

\paragraph{(Err\textsubscript{IV}, 46).}
En (IV, 15.6.3), línea 11 de la p.~237, en lugar de « supposons que
pour tout $y\in Y$ », leer: « supposons que $f$ soit localement de
présentation finie et que pour tout $y\in Y$ ». Sustituir la
demostración (líneas 14 a 24 de la p.~237) por la siguiente:
\begin{quote}
En virtud de {\EGAReadLink{IV.1.10.3}{(1.10.3)}}, basta demostrar
que para todo $x\in X_y^0$ y toda generalización $y'$ de $y$, existe una
generalización $x'$ de $x$ tal que $f(x')=y'$. Ahora bien, basta tomar
como $x'$ el punto genérico de $X_{y'}^0$, en virtud de
{\EGAReadLink{IV.15.6.2}{(15.6.2)}}.
\end{quote}

\paragraph{(Err\textsubscript{IV}, 47).}
En (IV, 15.6.7), línea 2 de la p.~241, después de « noethérien », añadir:
« ou si la condition $\beta$) de
({\EGAReadLink{IV.15.6.6}{15.6.6}}, (iii)) est vérifiée ».
Después de la línea 9 de la p.~241, añadir:
\begin{quote}
Cuando se cumple la condición $\beta$) de
({\EGAReadLink{IV.15.6.6}{15.6.6}}, (iii)), el hecho de que
a) implica b') resulta de
({\EGAReadLink{IV.15.6.6}{15.6.6}}, (ii)), y el hecho de que
b') implica a) resulta de
({\EGAReadLink{IV.15.6.6}{15.6.6}}, (iii)).
\end{quote}

\paragraph{(Err\textsubscript{IV}, 48).}
En (IV, 20.1.5), línea 13 de la p.~228, (IV, 20.1.6), línea 17 de la
p.~228, y (IV, 20.1.7), línea 13 desde abajo de la p.~228, suprimir
« strictement ».

\paragraph{(Err\textsubscript{IV}, 49).}
En (IV, 20.2.13), añadir después de la línea 6 de la p.~236: En
particular, si $X$ es reducido y solo tiene un número finito de
componentes irreducibles, la noción de $\sh{O}_X$-módulo sin torsión
definida en ({\EGAReadLink{IV.20.1.5}{IV, 20.1.5}})
coincide con la de (I, 7.4.1).

\paragraph{(Err\textsubscript{IV}, 50).}
En (IV, 20.2.14), líneas 18 y 25 de la p.~236, suprimir
« strictement ».

\paragraph{(Err\textsubscript{IV}, 51).}
En (IV, 20.6.2), línea 13 desde abajo de la p.~252, suprimir
« strictement ».

\paragraph{(Err\textsubscript{IV}, 52).}
En (IV, 20.6.4), líneas 23 y 22 de la p.~253, suprimir
« strictement ».

\paragraph{(Err\textsubscript{IV}, 53).}
En (IV, 21.4), después de la línea 23 de la p.~267, añadir:

\begingroup
\def\EGAIVSuppressErratumLabels{1}
% EGA IV, Errata y Addenda, lista 3, n.º 53.
% Este componente se inserta después de 21.4.8 en el texto de lectura corregido
% y se repite literalmente bajo Err_IV, 53 en el apéndice de erratas impreso.

La proposición siguiente y sus corolarios nos fueron señalados por
M.~Raynaud.

\begin{proposition}[21.4.9]
\ifdefined\EGAIVSuppressErratumLabels\else\label{IV.21.4.9}\fi
Sean $X$ e $Y$ dos preesquemas localmente noetherianos, y sea
$f:X\to Y$ un morfismo fielmente plano. Supongamos además que $f$ es
cuasicompacto o que $f$ es abierto. Un divisor $D'$ sobre $X$ es de la
forma $f^*(D)$, con $D$ un divisor sobre $Y$, si y solo si se cumple la
condición siguiente: para todo $y\in Y$ tal que
$\operatorname{prof}(\sh{O}_{Y,y})\leq 1$, pongamos
\[
 Y_1=\operatorname{Spec}(\sh{O}_{Y,y}),\qquad
 X_1=X\times_Y Y_1,
 \qquad f_1=f_{(Y_1)}:X_1\to Y_1;
\]
entonces la imagen recíproca $D_1'$ de $D'$ sobre $X_1$ es de la forma
$f_1^*(D_1)$ para algún $D_1\in\operatorname{Div}(Y_1)$.
\end{proposition}

\oldpage[IV]{359}
\begin{proof}
La necesidad de la condición resulta inmediatamente de la transitividad de las
imágenes recíprocas de divisores, pues todos los morfismos considerados son planos
{\EGAReadLink{IV.21.4.4}{(21.4.4)}}. Para todo abierto $U$ de $Y$, denotemos por $f_U$ la
restricción $f^{-1}(U)\to U$ de $f$. Por {\EGAReadLink{IV.21.4.8}{(21.4.8)}}, la aplicación
\[
 D_U\longmapsto f_U^*(D_U)
\]
de $\operatorname{Div}(U)$ en
$\operatorname{Div}(f^{-1}(U))$ es inyectiva. Así, si $U_1,U_2$ son
dos abiertos de $Y$ y
\[
 D'|f^{-1}(U_1)=f_{U_1}^*(D_1),\qquad
 D'|f^{-1}(U_2)=f_{U_2}^*(D_2),
\]
entonces necesariamente
$D_1|U_1\cap U_2=D_2|U_1\cap U_2$. Se sigue inmediatamente que existe un
mayor abierto $U\subset Y$ para el cual
\[
 D'|f^{-1}(U)=f_U^*(D_U)
 \quad\text{con}\quad D_U\in\operatorname{Div}(U).
\]
Debemos demostrar que $U=Y$. Supongamos lo contrario, y sea $y$ un punto
maximal de $Y\setminus U$. Encontraremos un entorno abierto $V$ de
$y$ en $Y$ tal que
  $D'|f^{-1}(V)=f_V^*(D_V)$ para algún
$D_V\in\operatorname{Div}(V)$, contradiciendo la maximalidad de $U$.

\textup{I) Caso en que $f$ es cuasicompacto.}
Utilizando la cuasicompacidad de $f$ y {\EGAReadLink{IV.20.3.8}{(20.3.8)}}, por el método de
({\EGAReadLink{IV.8.1.2}{8.1.2}}, a), nos reducimos a demostrar ---con las notaciones del
enunciado, pero sin imponer inicialmente una cota a
$\operatorname{prof}(\sh{O}_{Y,y})$--- que $D_1'$ es de la forma
$f_1^*(D_1)$. Si $\operatorname{prof}(\sh{O}_{Y,y})\leq 1$, esto es
precisamente la hipótesis. Supongamos, pues, que
$\operatorname{prof}(\sh{O}_{Y,y})\geq 2$. Como $y$ es un punto
maximal de $Y\setminus U$, el abierto $W=Y_1\cap U$ de $Y_1$ es
$Y_1\setminus\{y\}$. Podemos, por tanto, limitarnos al caso
$Y=Y_1$ y $U=Y\setminus\{y\}$.

Por hipótesis existe un divisor $D_U\in\operatorname{Div}(U)$ tal
que $f_U^*(D_U)=D'|f^{-1}(U)$. El divisor $D_U$ es la clase de
equivalencia de un par $(\sh{L}_U,s)$, donde $\sh{L}_U$ es un
$\sh{O}_U$-módulo invertible y $s$ una sección meromorfa regular de
$\sh{L}_U$ sobre $U$ {\EGAReadLink{IV.21.2.11}{(21.2.11)}}. Por tanto, en virtud de
{\EGAReadLink{IV.21.4.1}{(21.4.1)}}, $D'|f^{-1}(U)$ es la clase de equivalencia de
$(\sh{L}'_U,s')$, donde
\[
 \sh{L}'_U=f_U^*(\sh{L}_U),\qquad s'=s\circ f_U.
\]
Sean $i:U\to Y$ y $j:f^{-1}(U)\to X$ las inmersiones canónicas,
y pongamos
\[
 \sh{L}=i_*(\sh{L}_U),\qquad
 \sh{L}'=j_*(\sh{L}'_U).
\]
Como $f$ es plano, $f^*(\sh{L})$ es canónicamente isomorfo a
$\sh{L}'$ {\EGAReadLink{IV.2.3.1}{(2.3.1)}}. Además, la planitud y la desigualdad
$\operatorname{prof}(\sh{O}_{Y,y})\geq 2$ implican
$\operatorname{prof}(\sh{O}_{X,x})\geq 2$ para todo
$x\in f^{-1}(y)$ {\EGAReadLink{IV.6.3.1}{(6.3.1)}}.

El divisor $D'$ es la clase de equivalencia de un par
$(\sh{L}'',s'')$, donde $\sh{L}''$ es un $\sh{O}_X$-módulo invertible
y $s''$ una sección meromorfa regular de $\sh{L}''$ sobre $X$.
Por consiguiente, $\sh{L}'_U$ es isomorfo a
$\sh{L}''|f^{-1}(U)$, y las observaciones precedentes, junto con
{\EGAReadLink{IV.5.9.8}{(5.9.8)}} y {\EGAReadLink{IV.5.10.5}{(5.10.5)}}, muestran que $\sh{L}'$ es un
$\sh{O}_X$-módulo invertible isomorfo a $\sh{L}''$. Como $f$ es
fielmente plano y cuasicompacto, se sigue que $\sh{L}$ es un
$\sh{O}_Y$-módulo invertible {\EGAReadLink{IV.2.5.2}{(2.5.2)}} y, por tanto, isomorfo a $\sh{O}_Y$, pues
$Y$ es un esquema local.

Además, como $\operatorname{prof}(\sh{O}_{Y,y})\geq 1$, el abierto
$U=Y\setminus\{y\}$ es esquemáticamente denso en $Y$
({\EGAReadLink{IV.20.2.13}{20.2.13}}, (iv)). La sección meromorfa regular $s$ de $\sh{L}_U$
sobre $U$ se prolonga, por tanto, a una sección meromorfa regular $s_0$ de
$\sh{L}$ sobre $Y$ {\EGAReadLink{IV.20.2.11}{(20.2.11)}}. La clase de equivalencia de
$(\sh{L},s_0)$ es así un divisor $D$ sobre $Y$. Evidentemente, $f^*(D)$ y
$D'$ inducen el mismo divisor sobre
$f^{-1}(U)=X\setminus f^{-1}(y)$; y, como
$\operatorname{prof}(\sh{O}_{X,x})\geq 2$ para todo
$x\in f^{-1}(y)$, {\EGAReadLink{IV.21.1.8}{(21.1.8)}} da $D'=f^*(D)$. Esto prueba el caso I.

\textup{II) Caso en que $f$ es abierto.}
Se trata de encontrar un entorno abierto $V$ de $y$ en $Y$ tal
que $D'|f^{-1}(V)=f_V^*(D_V)$ para algún
$D_V\in\operatorname{Div}(V)$, por lo que podemos suponer primero que $Y$ es afín.
Como $f$ es abierto, las imágenes $f(W)$, cuando $W$ recorre los
abiertos afines de $X$, forman un recubrimiento abierto de $Y$. Como $Y$
es cuasicompacto, existe un abierto cuasicompacto $T\subset X$
tal que $f(T)=Y$, y $f|T$ es cuasicompacto (I, 1.1.1). Con mayor razón,
para todo abierto cuasicompacto $T_\alpha\supset T$ de $X$,
\oldpage[IV]{360}
se tiene
$f(T_\alpha)=Y$ y $f|T_\alpha$ es cuasicompacto.

Aplicando el caso I a $f|T_\alpha$, obtenemos un divisor $D_\alpha$ sobre
$Y$ tal que
\[
 (f|T_\alpha)^*(D_\alpha)=D'|T_\alpha.
\]
Por transitividad de las imágenes recíprocas de divisores por morfismos planos
{\EGAReadLink{IV.21.4.4}{(21.4.4)}}, también se tiene
$(f|T)^*(D_\alpha)=D'|T$. Por tanto, para cualesquiera $\alpha,\beta$,
\[
 (f|T)^*(D_\alpha)=(f|T)^*(D_\beta),
\]
y {\EGAReadLink{IV.21.4.8}{(21.4.8)}} da $D_\alpha=D_\beta$. Si $D$ denota su valor
común, entonces $f^*(D)|T_\alpha=D'|T_\alpha$ para todo $\alpha$, y
por consiguiente $f^*(D)=D'$.
\end{proof}

\begin{corollary}[21.4.10]
\ifdefined\EGAIVSuppressErratumLabels\else\label{IV.21.4.10}\fi
Supongamos las hipótesis generales de {\EGAReadLink{IV.21.4.9}{(21.4.9)}} y, además, que
$Y$ es normal y que, para todo $y\in Y$ con
$\dim(\sh{O}_{Y,y})=1$, la fibra $X_y=f^{-1}(y)$ no tiene ningún
ciclo primo asociado inmerso. Un divisor $D'$ sobre $X$ es de la forma $f^*(D)$,
con $D\in\operatorname{Div}(Y)$, si y solo si se cumplen las dos
condiciones siguientes:
\begin{enumerate}[label=\textup{\arabic*\textdegree}]
\item para todo punto maximal $y$ de $Y$, la imagen recíproca de $D'$ en
$\operatorname{Div}(X_y)$ es cero;
\item para todo $y\in Y$ tal que $\dim(\sh{O}_{Y,y})=1$, existe un
entero $n_y$ tal que, para todo punto maximal $x$ de $X_y$, la
imagen recíproca de $D'$ en
$\operatorname{Div}(\operatorname{Spec}(\sh{O}_{X,x}))$ es igual a
$\operatorname{div}(t_y^{n_y})$, donde $t_y$ es la imagen por
$\sh{O}_{Y,y}\to\sh{O}_{X,x}$ de un uniformizante del anillo de
valoración discreta $\sh{O}_{Y,y}$.
\end{enumerate}
\end{corollary}

\begin{proof}
Aplicamos el criterio {\EGAReadLink{IV.21.4.9}{(21.4.9)}}. En un punto maximal $y$ de $Y$,
$\operatorname{Div}(Y_1)=0$ porque $\sh{O}_{Y,y}$ es artiniano;
la condición 1\textdegree{} es, por tanto, equivalente al criterio de
{\EGAReadLink{IV.21.4.9}{(21.4.9)}} en tal punto. Queda por demostrar que las condiciones
1\textdegree{} y 2\textdegree{} implican asimismo el criterio en un
punto $y\in Y$ con $\operatorname{prof}(\sh{O}_{Y,y})=1$. Como $Y$
es normal, estos son precisamente los puntos para los cuales
$\dim(\sh{O}_{Y,y})=1$, equivalentemente, aquellos para los cuales
$\sh{O}_{Y,y}$ es un anillo de valoración discreta {\EGAReadLink{IV.5.8.6}{(5.8.6)}}.

Podemos, pues, limitarnos a $Y=\operatorname{Spec}(\sh{O}_{Y,y})$. Si $t$
es un uniformizante de $\sh{O}_{Y,y}$, la aplicación
$n\mapsto\operatorname{div}(t^n)$ es un isomorfismo de $\mathbb Z$
sobre $\operatorname{Div}(Y)$. La condición de {\EGAReadLink{IV.21.4.9}{(21.4.9)}} dice que
$D'$ es de la forma $\operatorname{div}(t^n\circ f)$ para algún
$n\in\mathbb Z$, e implica evidentemente la condición 2\textdegree{}.
Recíprocamente, supongamos que se cumple 2\textdegree{}, así como 1\textdegree{}, y
escribamos $n=n_y$. Queda por demostrar
$D'=\operatorname{div}(t^n\circ f)$. Si $\eta$ es el punto genérico de
$Y$, los gérmenes de estos dos divisores en todo punto de $X_\eta$ son
cero por la condición 1\textdegree{}. Sus gérmenes en todo punto maximal
de $X_y$ son iguales por la condición 2\textdegree{}. En un punto no maximal
$x\in X_y$, se tiene
$\operatorname{prof}(\sh{O}_{X_y,x})\geq 1$, porque $X_y$ no tiene
ciclo primo asociado inmerso y, por tanto,
$\operatorname{prof}(\sh{O}_{X,x})\geq 2$ por {\EGAReadLink{IV.6.3.1}{(6.3.1)}}. La
igualdad buscada resulta de {\EGAReadLink{IV.21.1.8}{(21.1.8)}}.
\end{proof}

\begin{corollary}[21.4.11]
\ifdefined\EGAIVSuppressErratumLabels\else\label{IV.21.4.11}\fi
Supongamos las hipótesis generales de {\EGAReadLink{IV.21.4.9}{(21.4.9)}} y, además, que
$Y$ es normal y que, para todo $y\in Y$ con
$\dim(\sh{O}_{Y,y})=1$, la fibra $X_y$ es íntegra. Un divisor $D'$ es
de la forma $f^*(D)$, con $D\in\operatorname{Div}(Y)$, si y solo si,
para todo punto maximal $y$ de $Y$, la imagen recíproca de $D'$ en
$\operatorname{Div}(X_y)$ es cero.
\end{corollary}

\begin{proof}
Es suficiente comprobar que la condición 2\textdegree{} de {\EGAReadLink{IV.21.4.10}{(21.4.10)}} es entonces
automática. Si $\dim(\sh{O}_{Y,y})=1$ y $X_y$ es íntegra, entonces en
el único punto maximal $x$ de $X_y$ se tiene
$\dim(\sh{O}_{X,x})=1$ por {\EGAReadLink{IV.6.1.3}{(6.1.3)}}. Además,
$\sh{O}_{Y,y}$ es regular y $\sh{O}_{X_y,x}$ es un cuerpo, por lo que
$\sh{O}_{X,x}$ es regular {\EGAReadLink{IV.6.5.1}{(6.5.1)}} y, por tanto, un anillo de valoración discreta.
Si $t$ es un uniformizante de $\sh{O}_{Y,y}$, su imagen $t_y$ en
$\sh{O}_{X,x}$ es un uniformizante de $\sh{O}_{X,x}$
($0$, 17.3.3). La imagen recíproca de $D'$ en
$\operatorname{Div}(\operatorname{Spec}(\sh{O}_{X,x}))$ es, por tanto,
\oldpage[IV]{361}
de la forma $\operatorname{div}(t_y^n)$, lo que es exactamente la condición
2\textdegree{} de {\EGAReadLink{IV.21.4.10}{(21.4.10)}}, puesto que $X_y$ solo tiene un punto maximal.
\end{proof}

\begin{remark}[21.4.12]
\ifdefined\EGAIVSuppressErratumLabels\else\label{IV.21.4.12}\fi
Si en {\EGAReadLink{IV.21.4.11}{(21.4.11)}} se supone solamente que los $X_y$ son reducidos, hay que
añadir a la condición del enunciado que las imágenes recíprocas de $D'$
en cada uno de los grupos
$\operatorname{Div}(\operatorname{Spec}(\sh{O}_{X,x}))$, cuando $x$
recorre los puntos maximales de $X_y$, sean todas de la forma
$\operatorname{div}((t_x')^n)$, donde $t_x'$ es un uniformizante de
$\sh{O}_{X,x}$ y el entero $n$ es el mismo para todos estos puntos.
\end{remark}

\begin{corollary}[21.4.13]
\ifdefined\EGAIVSuppressErratumLabels\else\label{IV.21.4.13}\fi
Supongamos las hipótesis de {\EGAReadLink{IV.21.4.11}{(21.4.11)}}. Para todo
$\sh{O}_X$-módulo invertible $\sh{L}$ tal que, para todo punto maximal $\eta$ de
$Y$, la imagen recíproca
\[
 \sh{L}_\eta
 =\sh{L}\otimes_{\sh{O}_X}\sh{O}_{X_\eta}
\]
sea isomorfa a $\sh{O}_{X_\eta}$, existen un
$\sh{O}_Y$-módulo invertible $\sh{M}$ y un isomorfismo
$f^*(\sh{M})\simeq\sh{L}$.
\end{corollary}

\begin{proof}
Basta encontrar una sección meromorfa regular $s$ de $\sh{L}$ sobre
$X$ tal que el divisor $D'$ correspondiente a $(\sh{L},s)$ tenga
$D'_x=0$ en todo punto $x$ de una fibra $X_\eta$ sobre un punto maximal
$\eta$ de $Y$. Aplicando entonces {\EGAReadLink{IV.21.4.11}{(21.4.11)}} a $D'$, se obtiene el resultado por
{\EGAReadLink{IV.21.4.2}{(21.4.2)}}.

Los puntos de $\operatorname{Ass}(\sh{O}_X)$ están sobre los puntos
maximales de $Y$ {\EGAReadLink{IV.3.3.1}{(3.3.1)}}. Como $Y$ es reducido y localmente noetheriano,
existe un abierto denso $V$ de $Y$ que es la unión de una
familia disjunta de abiertos afines, cada uno de los cuales contiene exactamente un
punto maximal de $Y$; y $f^{-1}(V)$ es esquemáticamente denso en $X$
{\EGAReadLink{IV.5.10.2}{(5.10.2)}}. Podemos, por tanto, limitarnos al caso en que $Y$ es
afín e íntegro.

Para el único punto maximal $\eta$ de $Y$, la fibra $X_\eta$ es entonces
íntegra, de modo que $\operatorname{Ass}(\sh{O}_X)$ consta de un solo punto
$z$, el punto genérico de $X$. Sea $s_\eta$ la sección de
$\sh{L}_\eta$ sobre $X_\eta$ que corresponde, mediante un isomorfismo,
a la sección unidad de $\sh{O}_{X_\eta}$. Sea $U$ un entorno abierto
afín en $X$ de un punto $x\in X_\eta$, y sea $s$ una sección
de $\sh{L}$ sobre $U$ con el mismo germen en $x$ que $s_\eta$. Como $U$
es esquemáticamente denso en $X$ ({\EGAReadLink{IV.20.2.13}{20.2.13}}, (iv)), $s$ se identifica con una
sección meromorfa regular unívocamente determinada de $\sh{L}$ sobre $X$.
Resulta claramente de {\EGAReadLink{IV.20.2.11}{(20.2.11)}} y {\EGAReadLink{IV.20.3.5}{(20.3.5)}} que esta sección induce sobre
$X_\eta$ una sección meromorfa regular de $\sh{L}_\eta$ igual a
$s_\eta$ y, por tanto, tiene la propiedad requerida.
\end{proof}

\endgroup

\bigskip
\begin{center}
\small
1967. --- Imprimerie des Presses Universitaires de France. --- Vendôme
(France)\\
ÉDIT. N\textsuperscript{o} 29 470\qquad
IMPRIMÉ EN FRANCE\qquad
IMP. N\textsuperscript{o} 20 298
\end{center}

\end{document}
